\documentclass[11pt]{article}
\usepackage[utf8]{inputenc}
\usepackage[T1]{fontenc}
\usepackage[a4paper,margin=1in]{geometry}
\usepackage{lmodern}
\usepackage{microtype}
\usepackage{amsmath,amssymb,mathtools}
\usepackage{mathrsfs}
\usepackage{tikz-cd}
\usepackage{enumitem}
\usepackage{longtable,booktabs,array}
\usepackage[hidelinks]{hyperref}
\usepackage{hyperxmp}
\newcommand{\sgaStableTarget}[1]{\hypertarget{#1}{}}
\newcommand{\sgaStableTargetInline}[1]{\pdfdest name {#1} xyz\relax}
\setlength{\parindent}{0pt}
\setlength{\parskip}{0.55em}
\emergencystretch=3em
\hbadness=10000
\hfuzz=20pt
\newcommand{\K}{K}
\newcommand{\Gr}{\operatorname{Gr}}
\newcommand{\Fil}{\operatorname{Fil}}
\newcommand{\ch}{\operatorname{ch}}
\newcommand{\Todd}{\operatorname{Todd}}
\newcommand{\cl}{\operatorname{cl}}
\newcommand{\Parf}{\operatorname{Parf}}
\newcommand{\Tor}{\operatorname{Tor}}
\newcommand{\Spec}{\operatorname{Spec}}
\newcommand{\RRR}{\mathrm{RRR}}
\providecommand{\codim}{\operatorname{codim}}
\providecommand{\Inv}{\operatorname{Inv}}
\providecommand{\Pic}{\operatorname{Pic}}
\providecommand{\Z}{\mathbb Z}
\providecommand{\Q}{\mathbb Q}
\providecommand{\R}{\mathbb R}
\providecommand{\C}{\mathbb C}
\providecommand{\Ahat}{\widehat A}
\providecommand{\Ahp}{\widehat A^{+}}
\providecommand{\That}{\widehat T}
\providecommand{\rad}{\operatorname{rad}}
\providecommand{\Gl}{\operatorname{Gl}}
\providecommand{\GL}{\operatorname{GL}}
\providecommand{\ch}{\operatorname{ch}}
\providecommand{\Todd}{\operatorname{Todd}}
\providecommand{\Gr}{\operatorname{Gr}}

\providecommand{\Atil}{\widetilde A}
\providecommand{\Ctil}{\widetilde C}
\providecommand{\K}{\mathrm K}
\providecommand{\cC}{\mathcal C}
\providecommand{\cF}{\mathscr F}
\providecommand{\eps}{\varepsilon}
\providecommand{\ellc}{\ell}
\providecommand{\Tor}{\operatorname{Tor}}
\providecommand{\Coker}{\operatorname{Coker}}
\providecommand{\Ker}{\operatorname{Ker}}
\providecommand{\Supp}{\operatorname{Supp}}
\providecommand{\codim}{\operatorname{codim}}
\providecommand{\Gr}{\operatorname{G}}
\long\def\SGAArchiveOnly#1{}
\providecommand{\rank}{\operatorname{rank}}
\providecommand{\rang}{\operatorname{rang}}
\providecommand{\cO}{\mathcal O}

\title{SGA 6: Intersection Theory and the Riemann--Roch Theorem}
\author{P. Berthelot, A. Grothendieck, L. Illusie}
\date{22 July 2026}

\hypersetup{
  hypertexnames=false,
  pdftitle={SGA 6: Intersection Theory and the Riemann--Roch Theorem -- English Translation},
  pdfauthor={P. Berthelot, A. Grothendieck, L. Illusie},
  pdfsubject={English translation of SGA 6},
  pdfkeywords={SGA 6, intersection theory, Riemann--Roch, English translation}
}
\providecommand{\N}{\mathbb N}
\providecommand{\Zell}{\mathbb Z_\ell}
\providecommand{\Qell}{\mathbb Q_\ell}
\providecommand{\C}{\mathbb C}
\providecommand{\Ch}{\operatorname{Ch}}
\providecommand{\Sym}{\operatorname{Sym}}
\providecommand{\Td}{\operatorname{Td}}

\usepackage[normalem]{ulem}

% Single annotated master; 0 diplomatic, 1 corrected with apparatus, 2 reading.
\providecommand{\SGAReaderForm}{1}
\newcommand{\SGACaptureAppState}[1]{\csname SGACapture@#1\endcsname}
\newcommand{\SGARestoreAppState}[1]{\csname SGARestore@#1\endcsname}
\newcommand{\SGAAppMark}[1]{\ifmmode{}^{[\text{\csname SGAAppNumber@#1\endcsname}]}\else\textsuperscript{[\csname SGAAppNumber@#1\endcsname]}\fi}
\long\def\var#1#2#3{\SGACaptureAppState{#3}\ifnum\SGAReaderForm=0 #1\else #2\ifnum\SGAReaderForm=1\SGAAppMark{#3}\fi\fi}
\newcommand{\SGAPrintPage}[1]{\ifnum\SGAReaderForm=2\marginpar{\scriptsize #1}\else\leavevmode\textsuperscript{\scriptsize[print #1]}\fi}
\expandafter\def\csname SGAAuthorLabel@raisedstar\endcsname{(\textsuperscript{*})}
\expandafter\def\csname SGAAuthorLabel@unmarked\endcsname{}
\newcommand{\SGAAuthorLabel}[1]{\ifcsname SGAAuthorLabel@#1\endcsname\csname SGAAuthorLabel@#1\endcsname\else(#1)\fi}
\makeatletter
\newcommand{\SGAAuthorFootnote}[2]{\begingroup\renewcommand{\thefootnote}{#1}\def\@makefnmark{\hbox{\normalfont\SGAAuthorLabel{#1}}}\def\@makefntext##1{\parindent1em\noindent\hb@xt@1.8em{\@makefnmark\hss}##1}\footnote{#2}\endgroup}
\newcommand{\SGAAuthorFootMark}[1]{\begingroup\renewcommand{\thefootnote}{#1}\def\@makefnmark{\hbox{\normalfont\SGAAuthorLabel{#1}}}\footnotemark[\value{footnote}]\endgroup}
\newcommand{\SGAAuthorFootText}[2]{\begingroup\renewcommand{\thefootnote}{#1}\def\@makefnmark{\hbox{\normalfont\SGAAuthorLabel{#1}}}\def\@makefntext##1{\parindent1em\noindent\hb@xt@1.8em{\@makefnmark\hss}##1}\footnotetext[\value{footnote}]{#2}\endgroup}
\makeatother
\expandafter\def\csname SGACapture@EN-VAR-001\endcsname{%
}
\expandafter\def\csname SGARestore@EN-VAR-001\endcsname{%
}
\expandafter\def\csname SGACapture@EN-VAR-002\endcsname{%
\global\expandafter\let\csname SGAState@EN-VAR-002@Todd\endcsname\Todd%
\global\expandafter\let\csname SGAState@EN-VAR-002@ch\endcsname\ch%
\global\expandafter\let\csname SGAState@EN-VAR-002@cl\endcsname\cl%
}
\expandafter\def\csname SGARestore@EN-VAR-002\endcsname{%
\expandafter\let\expandafter\Todd\csname SGAState@EN-VAR-002@Todd\endcsname%
\expandafter\let\expandafter\ch\csname SGAState@EN-VAR-002@ch\endcsname%
\expandafter\let\expandafter\cl\csname SGAState@EN-VAR-002@cl\endcsname%
}
\expandafter\def\csname SGACapture@EN-VAR-003\endcsname{%
}
\expandafter\def\csname SGARestore@EN-VAR-003\endcsname{%
}
\expandafter\def\csname SGACapture@EN-VAR-004\endcsname{%
\global\expandafter\let\csname SGAState@EN-VAR-004@cl\endcsname\cl%
}
\expandafter\def\csname SGARestore@EN-VAR-004\endcsname{%
\expandafter\let\expandafter\cl\csname SGAState@EN-VAR-004@cl\endcsname%
}
\expandafter\def\csname SGACapture@EN-VAR-005\endcsname{%
\global\expandafter\let\csname SGAState@EN-VAR-005@Parf\endcsname\Parf%
}
\expandafter\def\csname SGARestore@EN-VAR-005\endcsname{%
\expandafter\let\expandafter\Parf\csname SGAState@EN-VAR-005@Parf\endcsname%
}
\expandafter\def\csname SGACapture@EN-VAR-006\endcsname{%
}
\expandafter\def\csname SGARestore@EN-VAR-006\endcsname{%
}
\expandafter\def\csname SGACapture@EN-VAR-007\endcsname{%
}
\expandafter\def\csname SGARestore@EN-VAR-007\endcsname{%
}
\expandafter\def\csname SGACapture@EN-VAR-008\endcsname{%
}
\expandafter\def\csname SGARestore@EN-VAR-008\endcsname{%
}
\expandafter\def\csname SGACapture@EN-VAR-009\endcsname{%
}
\expandafter\def\csname SGARestore@EN-VAR-009\endcsname{%
}
\expandafter\def\csname SGACapture@EN-VAR-010\endcsname{%
}
\expandafter\def\csname SGARestore@EN-VAR-010\endcsname{%
}
\expandafter\def\csname SGACapture@EN-VAR-011\endcsname{%
}
\expandafter\def\csname SGARestore@EN-VAR-011\endcsname{%
}
\expandafter\def\csname SGACapture@EN-VAR-012\endcsname{%
}
\expandafter\def\csname SGARestore@EN-VAR-012\endcsname{%
}
\expandafter\def\csname SGACapture@EN-VAR-013\endcsname{%
}
\expandafter\def\csname SGARestore@EN-VAR-013\endcsname{%
}
\expandafter\def\csname SGACapture@EN-VAR-014\endcsname{%
}
\expandafter\def\csname SGARestore@EN-VAR-014\endcsname{%
}
\expandafter\def\csname SGACapture@EN-VAR-015\endcsname{%
}
\expandafter\def\csname SGARestore@EN-VAR-015\endcsname{%
}
\expandafter\def\csname SGACapture@EN-VAR-016\endcsname{%
\global\expandafter\let\csname SGAState@EN-VAR-016@Ahp\endcsname\Ahp%
}
\expandafter\def\csname SGARestore@EN-VAR-016\endcsname{%
\expandafter\let\expandafter\Ahp\csname SGAState@EN-VAR-016@Ahp\endcsname%
}
\expandafter\def\csname SGACapture@EN-VAR-017\endcsname{%
}
\expandafter\def\csname SGARestore@EN-VAR-017\endcsname{%
}
\expandafter\def\csname SGACapture@EN-VAR-018\endcsname{%
}
\expandafter\def\csname SGARestore@EN-VAR-018\endcsname{%
}
\expandafter\def\csname SGACapture@EN-VAR-019\endcsname{%
}
\expandafter\def\csname SGARestore@EN-VAR-019\endcsname{%
}
\expandafter\def\csname SGACapture@EN-VAR-020\endcsname{%
}
\expandafter\def\csname SGARestore@EN-VAR-020\endcsname{%
}
\expandafter\def\csname SGACapture@EN-VAR-021\endcsname{%
}
\expandafter\def\csname SGARestore@EN-VAR-021\endcsname{%
}
\expandafter\def\csname SGACapture@EN-VAR-022\endcsname{%
}
\expandafter\def\csname SGARestore@EN-VAR-022\endcsname{%
}
\expandafter\def\csname SGACapture@EN-VAR-023\endcsname{%
}
\expandafter\def\csname SGARestore@EN-VAR-023\endcsname{%
}
\expandafter\def\csname SGACapture@EN-VAR-024\endcsname{%
}
\expandafter\def\csname SGARestore@EN-VAR-024\endcsname{%
}
\expandafter\def\csname SGACapture@EN-VAR-025\endcsname{%
}
\expandafter\def\csname SGARestore@EN-VAR-025\endcsname{%
}
\expandafter\def\csname SGACapture@EN-VAR-026\endcsname{%
}
\expandafter\def\csname SGARestore@EN-VAR-026\endcsname{%
}
\expandafter\def\csname SGACapture@EN-VAR-027\endcsname{%
\global\expandafter\let\csname SGAState@EN-VAR-027@Atil\endcsname\Atil%
}
\expandafter\def\csname SGARestore@EN-VAR-027\endcsname{%
\expandafter\let\expandafter\Atil\csname SGAState@EN-VAR-027@Atil\endcsname%
}
\expandafter\def\csname SGACapture@EN-VAR-028\endcsname{%
\global\expandafter\let\csname SGAState@EN-VAR-028@Atil\endcsname\Atil%
}
\expandafter\def\csname SGARestore@EN-VAR-028\endcsname{%
\expandafter\let\expandafter\Atil\csname SGAState@EN-VAR-028@Atil\endcsname%
}
\expandafter\def\csname SGACapture@EN-VAR-029\endcsname{%
}
\expandafter\def\csname SGARestore@EN-VAR-029\endcsname{%
}
\expandafter\def\csname SGACapture@EN-VAR-030\endcsname{%
}
\expandafter\def\csname SGARestore@EN-VAR-030\endcsname{%
}
\expandafter\def\csname SGACapture@EN-VAR-031\endcsname{%
}
\expandafter\def\csname SGARestore@EN-VAR-031\endcsname{%
}
\expandafter\def\csname SGACapture@EN-VAR-032\endcsname{%
}
\expandafter\def\csname SGARestore@EN-VAR-032\endcsname{%
}
\expandafter\def\csname SGACapture@EN-VAR-033\endcsname{%
\global\expandafter\let\csname SGAState@EN-VAR-033@Tor\endcsname\Tor%
}
\expandafter\def\csname SGARestore@EN-VAR-033\endcsname{%
\expandafter\let\expandafter\Tor\csname SGAState@EN-VAR-033@Tor\endcsname%
}
\expandafter\def\csname SGACapture@EN-VAR-034\endcsname{%
}
\expandafter\def\csname SGARestore@EN-VAR-034\endcsname{%
}
\expandafter\def\csname SGACapture@EN-VAR-035\endcsname{%
}
\expandafter\def\csname SGARestore@EN-VAR-035\endcsname{%
}
\expandafter\def\csname SGACapture@EN-VAR-036\endcsname{%
}
\expandafter\def\csname SGARestore@EN-VAR-036\endcsname{%
}
\expandafter\def\csname SGACapture@EN-VAR-037\endcsname{%
}
\expandafter\def\csname SGARestore@EN-VAR-037\endcsname{%
}
\expandafter\def\csname SGACapture@EN-VAR-038\endcsname{%
}
\expandafter\def\csname SGARestore@EN-VAR-038\endcsname{%
}
\expandafter\def\csname SGACapture@EN-VAR-039\endcsname{%
}
\expandafter\def\csname SGARestore@EN-VAR-039\endcsname{%
}
\expandafter\def\csname SGACapture@EN-VAR-040\endcsname{%
}
\expandafter\def\csname SGARestore@EN-VAR-040\endcsname{%
}
\expandafter\def\csname SGACapture@EN-VAR-041\endcsname{%
}
\expandafter\def\csname SGARestore@EN-VAR-041\endcsname{%
}
\expandafter\def\csname SGACapture@EN-VAR-042\endcsname{%
}
\expandafter\def\csname SGARestore@EN-VAR-042\endcsname{%
}
\expandafter\def\csname SGACapture@EN-VAR-043\endcsname{%
}
\expandafter\def\csname SGARestore@EN-VAR-043\endcsname{%
}
\expandafter\def\csname SGACapture@EN-VAR-044\endcsname{%
}
\expandafter\def\csname SGARestore@EN-VAR-044\endcsname{%
}
\expandafter\def\csname SGACapture@EN-VAR-045\endcsname{%
}
\expandafter\def\csname SGARestore@EN-VAR-045\endcsname{%
}
\expandafter\def\csname SGACapture@EN-VAR-046\endcsname{%
}
\expandafter\def\csname SGARestore@EN-VAR-046\endcsname{%
}
\expandafter\def\csname SGACapture@EN-VAR-047\endcsname{%
}
\expandafter\def\csname SGARestore@EN-VAR-047\endcsname{%
}
\expandafter\def\csname SGACapture@EN-VAR-048\endcsname{%
}
\expandafter\def\csname SGARestore@EN-VAR-048\endcsname{%
}
\expandafter\def\csname SGACapture@EN-VAR-049\endcsname{%
}
\expandafter\def\csname SGARestore@EN-VAR-049\endcsname{%
}
\expandafter\def\csname SGACapture@EN-VAR-050\endcsname{%
}
\expandafter\def\csname SGARestore@EN-VAR-050\endcsname{%
}
\expandafter\def\csname SGACapture@EN-VAR-051\endcsname{%
}
\expandafter\def\csname SGARestore@EN-VAR-051\endcsname{%
}
\expandafter\def\csname SGACapture@EN-VAR-052\endcsname{%
}
\expandafter\def\csname SGARestore@EN-VAR-052\endcsname{%
}
\expandafter\def\csname SGACapture@EN-VAR-053\endcsname{%
\global\expandafter\let\csname SGAState@EN-VAR-053@Atil\endcsname\Atil%
\global\expandafter\let\csname SGAState@EN-VAR-053@Ctil\endcsname\Ctil%
}
\expandafter\def\csname SGARestore@EN-VAR-053\endcsname{%
\expandafter\let\expandafter\Atil\csname SGAState@EN-VAR-053@Atil\endcsname%
\expandafter\let\expandafter\Ctil\csname SGAState@EN-VAR-053@Ctil\endcsname%
}
\expandafter\def\csname SGACapture@EN-VAR-054\endcsname{%
}
\expandafter\def\csname SGARestore@EN-VAR-054\endcsname{%
}
\expandafter\def\csname SGACapture@EN-VAR-055\endcsname{%
}
\expandafter\def\csname SGARestore@EN-VAR-055\endcsname{%
}
\expandafter\def\csname SGACapture@EN-VAR-056\endcsname{%
}
\expandafter\def\csname SGARestore@EN-VAR-056\endcsname{%
}
\expandafter\def\csname SGACapture@EN-VAR-057\endcsname{%
}
\expandafter\def\csname SGARestore@EN-VAR-057\endcsname{%
}
\expandafter\def\csname SGACapture@EN-VAR-058\endcsname{%
\global\expandafter\let\csname SGAState@EN-VAR-058@ch\endcsname\ch%
}
\expandafter\def\csname SGARestore@EN-VAR-058\endcsname{%
\expandafter\let\expandafter\ch\csname SGAState@EN-VAR-058@ch\endcsname%
}
\expandafter\def\csname SGACapture@EN-VAR-059\endcsname{%
}
\expandafter\def\csname SGARestore@EN-VAR-059\endcsname{%
}
\expandafter\def\csname SGACapture@EN-VAR-060\endcsname{%
}
\expandafter\def\csname SGARestore@EN-VAR-060\endcsname{%
}
\expandafter\def\csname SGACapture@EN-VAR-061\endcsname{%
}
\expandafter\def\csname SGARestore@EN-VAR-061\endcsname{%
}
\expandafter\def\csname SGACapture@EN-VAR-062\endcsname{%
}
\expandafter\def\csname SGARestore@EN-VAR-062\endcsname{%
}
\expandafter\def\csname SGACapture@EN-VAR-063\endcsname{%
}
\expandafter\def\csname SGARestore@EN-VAR-063\endcsname{%
}
\expandafter\def\csname SGACapture@EN-VAR-064\endcsname{%
}
\expandafter\def\csname SGARestore@EN-VAR-064\endcsname{%
}
\expandafter\def\csname SGACapture@EN-VAR-065\endcsname{%
}
\expandafter\def\csname SGARestore@EN-VAR-065\endcsname{%
}
\expandafter\def\csname SGACapture@EN-VAR-066\endcsname{%
}
\expandafter\def\csname SGARestore@EN-VAR-066\endcsname{%
}
\expandafter\def\csname SGACapture@EN-VAR-067\endcsname{%
}
\expandafter\def\csname SGARestore@EN-VAR-067\endcsname{%
}
\expandafter\def\csname SGACapture@EN-VAR-068\endcsname{%
}
\expandafter\def\csname SGARestore@EN-VAR-068\endcsname{%
}
\expandafter\def\csname SGACapture@EN-VAR-069\endcsname{%
}
\expandafter\def\csname SGARestore@EN-VAR-069\endcsname{%
}
\expandafter\def\csname SGACapture@EN-VAR-070\endcsname{%
}
\expandafter\def\csname SGARestore@EN-VAR-070\endcsname{%
}
\expandafter\def\csname SGACapture@EN-VAR-071\endcsname{%
}
\expandafter\def\csname SGARestore@EN-VAR-071\endcsname{%
}
\expandafter\def\csname SGACapture@EN-VAR-072\endcsname{%
}
\expandafter\def\csname SGARestore@EN-VAR-072\endcsname{%
}
\expandafter\def\csname SGACapture@EN-VAR-073\endcsname{%
}
\expandafter\def\csname SGARestore@EN-VAR-073\endcsname{%
}
\expandafter\def\csname SGACapture@EN-VAR-074\endcsname{%
}
\expandafter\def\csname SGARestore@EN-VAR-074\endcsname{%
}
\expandafter\def\csname SGACapture@EN-VAR-075\endcsname{%
}
\expandafter\def\csname SGARestore@EN-VAR-075\endcsname{%
}
\expandafter\def\csname SGACapture@EN-VAR-076\endcsname{%
}
\expandafter\def\csname SGARestore@EN-VAR-076\endcsname{%
}
\expandafter\def\csname SGACapture@EN-VAR-077\endcsname{%
}
\expandafter\def\csname SGARestore@EN-VAR-077\endcsname{%
}
\expandafter\def\csname SGACapture@EN-VAR-078\endcsname{%
}
\expandafter\def\csname SGARestore@EN-VAR-078\endcsname{%
}
\expandafter\def\csname SGACapture@EN-VAR-079\endcsname{%
}
\expandafter\def\csname SGARestore@EN-VAR-079\endcsname{%
}
\expandafter\def\csname SGACapture@EN-VAR-080\endcsname{%
}
\expandafter\def\csname SGARestore@EN-VAR-080\endcsname{%
}
\expandafter\def\csname SGACapture@EN-VAR-081\endcsname{%
}
\expandafter\def\csname SGARestore@EN-VAR-081\endcsname{%
}
\expandafter\def\csname SGACapture@EN-VAR-082\endcsname{%
}
\expandafter\def\csname SGARestore@EN-VAR-082\endcsname{%
}
\expandafter\def\csname SGACapture@EN-VAR-083\endcsname{%
\global\expandafter\let\csname SGAState@EN-VAR-083@Ob\endcsname\Ob%
}
\expandafter\def\csname SGARestore@EN-VAR-083\endcsname{%
\expandafter\let\expandafter\Ob\csname SGAState@EN-VAR-083@Ob\endcsname%
}
\expandafter\def\csname SGACapture@EN-VAR-084\endcsname{%
}
\expandafter\def\csname SGARestore@EN-VAR-084\endcsname{%
}
\expandafter\def\csname SGACapture@EN-VAR-085\endcsname{%
}
\expandafter\def\csname SGARestore@EN-VAR-085\endcsname{%
}
\expandafter\def\csname SGACapture@EN-VAR-086\endcsname{%
}
\expandafter\def\csname SGARestore@EN-VAR-086\endcsname{%
}
\expandafter\def\csname SGACapture@EN-VAR-087\endcsname{%
}
\expandafter\def\csname SGARestore@EN-VAR-087\endcsname{%
}
\expandafter\def\csname SGACapture@EN-VAR-088\endcsname{%
}
\expandafter\def\csname SGARestore@EN-VAR-088\endcsname{%
}
\expandafter\def\csname SGACapture@EN-VAR-089\endcsname{%
}
\expandafter\def\csname SGARestore@EN-VAR-089\endcsname{%
}
\expandafter\def\csname SGACapture@EN-VAR-090\endcsname{%
}
\expandafter\def\csname SGARestore@EN-VAR-090\endcsname{%
}
\expandafter\def\csname SGACapture@EN-VAR-091\endcsname{%
}
\expandafter\def\csname SGARestore@EN-VAR-091\endcsname{%
}
\expandafter\def\csname SGACapture@EN-VAR-092\endcsname{%
}
\expandafter\def\csname SGARestore@EN-VAR-092\endcsname{%
}
\expandafter\def\csname SGACapture@EN-VAR-093\endcsname{%
}
\expandafter\def\csname SGARestore@EN-VAR-093\endcsname{%
}
\expandafter\def\csname SGACapture@EN-VAR-094\endcsname{%
}
\expandafter\def\csname SGARestore@EN-VAR-094\endcsname{%
}
\expandafter\def\csname SGACapture@EN-VAR-095\endcsname{%
}
\expandafter\def\csname SGARestore@EN-VAR-095\endcsname{%
}
\expandafter\def\csname SGACapture@EN-VAR-096\endcsname{%
}
\expandafter\def\csname SGARestore@EN-VAR-096\endcsname{%
}
\expandafter\def\csname SGACapture@EN-VAR-097\endcsname{%
}
\expandafter\def\csname SGARestore@EN-VAR-097\endcsname{%
}
\expandafter\def\csname SGACapture@EN-VAR-098\endcsname{%
\global\expandafter\let\csname SGAState@EN-VAR-098@Ob\endcsname\Ob%
}
\expandafter\def\csname SGARestore@EN-VAR-098\endcsname{%
\expandafter\let\expandafter\Ob\csname SGAState@EN-VAR-098@Ob\endcsname%
}
\expandafter\def\csname SGACapture@EN-VAR-099\endcsname{%
\global\expandafter\let\csname SGAState@EN-VAR-099@Dpx\endcsname\Dpx%
\global\expandafter\let\csname SGAState@EN-VAR-099@Tr\endcsname\Tr%
}
\expandafter\def\csname SGARestore@EN-VAR-099\endcsname{%
\expandafter\let\expandafter\Dpx\csname SGAState@EN-VAR-099@Dpx\endcsname%
\expandafter\let\expandafter\Tr\csname SGAState@EN-VAR-099@Tr\endcsname%
}
\expandafter\def\csname SGACapture@EN-VAR-100\endcsname{%
\global\expandafter\let\csname SGAState@EN-VAR-100@Kpx\endcsname\Kpx%
\global\expandafter\let\csname SGAState@EN-VAR-100@Tr\endcsname\Tr%
}
\expandafter\def\csname SGARestore@EN-VAR-100\endcsname{%
\expandafter\let\expandafter\Kpx\csname SGAState@EN-VAR-100@Kpx\endcsname%
\expandafter\let\expandafter\Tr\csname SGAState@EN-VAR-100@Tr\endcsname%
}
\expandafter\def\csname SGACapture@EN-VAR-101\endcsname{%
\global\expandafter\let\csname SGAState@EN-VAR-101@Dpx\endcsname\Dpx%
\global\expandafter\let\csname SGAState@EN-VAR-101@Kpx\endcsname\Kpx%
}
\expandafter\def\csname SGARestore@EN-VAR-101\endcsname{%
\expandafter\let\expandafter\Dpx\csname SGAState@EN-VAR-101@Dpx\endcsname%
\expandafter\let\expandafter\Kpx\csname SGAState@EN-VAR-101@Kpx\endcsname%
}
\expandafter\def\csname SGACapture@EN-VAR-102\endcsname{%
\global\expandafter\let\csname SGAState@EN-VAR-102@Kpx\endcsname\Kpx%
}
\expandafter\def\csname SGARestore@EN-VAR-102\endcsname{%
\expandafter\let\expandafter\Kpx\csname SGAState@EN-VAR-102@Kpx\endcsname%
}
\expandafter\def\csname SGACapture@EN-VAR-103\endcsname{%
\global\expandafter\let\csname SGAState@EN-VAR-103@perfamp\endcsname\perfamp%
\global\expandafter\let\csname SGAState@EN-VAR-103@perfdim\endcsname\perfdim%
}
\expandafter\def\csname SGARestore@EN-VAR-103\endcsname{%
\expandafter\let\expandafter\perfamp\csname SGAState@EN-VAR-103@perfamp\endcsname%
\expandafter\let\expandafter\perfdim\csname SGAState@EN-VAR-103@perfdim\endcsname%
}
\expandafter\def\csname SGACapture@EN-VAR-104\endcsname{%
\global\expandafter\let\csname SGAState@EN-VAR-104@Dpx\endcsname\Dpx%
}
\expandafter\def\csname SGARestore@EN-VAR-104\endcsname{%
\expandafter\let\expandafter\Dpx\csname SGAState@EN-VAR-104@Dpx\endcsname%
}
\expandafter\def\csname SGACapture@EN-VAR-105\endcsname{%
}
\expandafter\def\csname SGARestore@EN-VAR-105\endcsname{%
}
\expandafter\def\csname SGACapture@EN-VAR-106\endcsname{%
}
\expandafter\def\csname SGARestore@EN-VAR-106\endcsname{%
}
\expandafter\def\csname SGACapture@EN-VAR-107\endcsname{%
\global\expandafter\let\csname SGAState@EN-VAR-107@Ob\endcsname\Ob%
}
\expandafter\def\csname SGARestore@EN-VAR-107\endcsname{%
\expandafter\let\expandafter\Ob\csname SGAState@EN-VAR-107@Ob\endcsname%
}
\expandafter\def\csname SGACapture@EN-VAR-108\endcsname{%
\global\expandafter\let\csname SGAState@EN-VAR-108@Ob\endcsname\Ob%
}
\expandafter\def\csname SGARestore@EN-VAR-108\endcsname{%
\expandafter\let\expandafter\Ob\csname SGAState@EN-VAR-108@Ob\endcsname%
}
\expandafter\def\csname SGACapture@EN-VAR-109\endcsname{%
\global\expandafter\let\csname SGAState@EN-VAR-109@Cart\endcsname\Cart%
\global\expandafter\let\csname SGAState@EN-VAR-109@is\endcsname\is%
}
\expandafter\def\csname SGARestore@EN-VAR-109\endcsname{%
\expandafter\let\expandafter\Cart\csname SGAState@EN-VAR-109@Cart\endcsname%
\expandafter\let\expandafter\is\csname SGAState@EN-VAR-109@is\endcsname%
}
\expandafter\def\csname SGACapture@EN-VAR-110\endcsname{%
\global\expandafter\let\csname SGAState@EN-VAR-110@Id\endcsname\Id%
\global\expandafter\let\csname SGAState@EN-VAR-110@otL\endcsname\otL%
}
\expandafter\def\csname SGARestore@EN-VAR-110\endcsname{%
\expandafter\let\expandafter\Id\csname SGAState@EN-VAR-110@Id\endcsname%
\expandafter\let\expandafter\otL\csname SGAState@EN-VAR-110@otL\endcsname%
}
\expandafter\def\csname SGACapture@EN-VAR-111\endcsname{%
}
\expandafter\def\csname SGARestore@EN-VAR-111\endcsname{%
}
\expandafter\def\csname SGACapture@EN-VAR-112\endcsname{%
}
\expandafter\def\csname SGARestore@EN-VAR-112\endcsname{%
}
\expandafter\def\csname SGACapture@EN-VAR-113\endcsname{%
}
\expandafter\def\csname SGARestore@EN-VAR-113\endcsname{%
}
\expandafter\def\csname SGACapture@EN-VAR-114\endcsname{%
}
\expandafter\def\csname SGARestore@EN-VAR-114\endcsname{%
}
\expandafter\def\csname SGACapture@EN-VAR-115\endcsname{%
}
\expandafter\def\csname SGARestore@EN-VAR-115\endcsname{%
}
\expandafter\def\csname SGACapture@EN-VAR-116\endcsname{%
\global\expandafter\let\csname SGAState@EN-VAR-116@Qcoh\endcsname\Qcoh%
\global\expandafter\let\csname SGAState@EN-VAR-116@coh\endcsname\coh%
}
\expandafter\def\csname SGARestore@EN-VAR-116\endcsname{%
\expandafter\let\expandafter\Qcoh\csname SGAState@EN-VAR-116@Qcoh\endcsname%
\expandafter\let\expandafter\coh\csname SGAState@EN-VAR-116@coh\endcsname%
}
\expandafter\def\csname SGACapture@EN-VAR-117\endcsname{%
}
\expandafter\def\csname SGARestore@EN-VAR-117\endcsname{%
}
\expandafter\def\csname SGACapture@EN-VAR-118\endcsname{%
}
\expandafter\def\csname SGARestore@EN-VAR-118\endcsname{%
}
\expandafter\def\csname SGACapture@EN-VAR-119\endcsname{%
}
\expandafter\def\csname SGARestore@EN-VAR-119\endcsname{%
}
\expandafter\def\csname SGACapture@EN-VAR-120\endcsname{%
}
\expandafter\def\csname SGARestore@EN-VAR-120\endcsname{%
}
\expandafter\def\csname SGACapture@EN-VAR-121\endcsname{%
}
\expandafter\def\csname SGARestore@EN-VAR-121\endcsname{%
}
\expandafter\def\csname SGACapture@EN-VAR-122\endcsname{%
\global\expandafter\let\csname SGAState@EN-VAR-122@Ob\endcsname\Ob%
}
\expandafter\def\csname SGARestore@EN-VAR-122\endcsname{%
\expandafter\let\expandafter\Ob\csname SGAState@EN-VAR-122@Ob\endcsname%
}
\expandafter\def\csname SGACapture@EN-VAR-123\endcsname{%
}
\expandafter\def\csname SGARestore@EN-VAR-123\endcsname{%
}
\expandafter\def\csname SGACapture@EN-VAR-124\endcsname{%
}
\expandafter\def\csname SGARestore@EN-VAR-124\endcsname{%
}
\expandafter\def\csname SGACapture@EN-VAR-125\endcsname{%
}
\expandafter\def\csname SGARestore@EN-VAR-125\endcsname{%
}
\expandafter\def\csname SGACapture@EN-VAR-126\endcsname{%
}
\expandafter\def\csname SGARestore@EN-VAR-126\endcsname{%
}
\expandafter\def\csname SGACapture@EN-VAR-127\endcsname{%
\global\expandafter\let\csname SGAState@EN-VAR-127@coh\endcsname\coh%
}
\expandafter\def\csname SGARestore@EN-VAR-127\endcsname{%
\expandafter\let\expandafter\coh\csname SGAState@EN-VAR-127@coh\endcsname%
}
\expandafter\def\csname SGACapture@EN-VAR-128\endcsname{%
}
\expandafter\def\csname SGARestore@EN-VAR-128\endcsname{%
}
\expandafter\def\csname SGACapture@EN-VAR-129\endcsname{%
\global\expandafter\let\csname SGAState@EN-VAR-129@Loclib\endcsname\Loclib%
\global\expandafter\let\csname SGAState@EN-VAR-129@Ob\endcsname\Ob%
}
\expandafter\def\csname SGARestore@EN-VAR-129\endcsname{%
\expandafter\let\expandafter\Loclib\csname SGAState@EN-VAR-129@Loclib\endcsname%
\expandafter\let\expandafter\Ob\csname SGAState@EN-VAR-129@Ob\endcsname%
}
\expandafter\def\csname SGACapture@EN-VAR-130\endcsname{%
\global\expandafter\let\csname SGAState@EN-VAR-130@Loclib\endcsname\Loclib%
}
\expandafter\def\csname SGARestore@EN-VAR-130\endcsname{%
\expandafter\let\expandafter\Loclib\csname SGAState@EN-VAR-130@Loclib\endcsname%
}
\expandafter\def\csname SGACapture@EN-VAR-131\endcsname{%
\global\expandafter\let\csname SGAState@EN-VAR-131@parf\endcsname\parf%
}
\expandafter\def\csname SGARestore@EN-VAR-131\endcsname{%
\expandafter\let\expandafter\parf\csname SGAState@EN-VAR-131@parf\endcsname%
}
\expandafter\def\csname SGACapture@EN-VAR-132\endcsname{%
}
\expandafter\def\csname SGARestore@EN-VAR-132\endcsname{%
}
\expandafter\def\csname SGACapture@EN-VAR-133\endcsname{%
}
\expandafter\def\csname SGARestore@EN-VAR-133\endcsname{%
}
\expandafter\def\csname SGACapture@EN-VAR-134\endcsname{%
}
\expandafter\def\csname SGARestore@EN-VAR-134\endcsname{%
}
\expandafter\def\csname SGACapture@EN-VAR-135\endcsname{%
}
\expandafter\def\csname SGARestore@EN-VAR-135\endcsname{%
}
\expandafter\def\csname SGACapture@EN-VAR-136\endcsname{%
}
\expandafter\def\csname SGARestore@EN-VAR-136\endcsname{%
}
\expandafter\def\csname SGACapture@EN-VAR-137\endcsname{%
}
\expandafter\def\csname SGARestore@EN-VAR-137\endcsname{%
}
\expandafter\def\csname SGACapture@EN-VAR-138\endcsname{%
}
\expandafter\def\csname SGARestore@EN-VAR-138\endcsname{%
}
\expandafter\def\csname SGACapture@EN-VAR-139\endcsname{%
}
\expandafter\def\csname SGARestore@EN-VAR-139\endcsname{%
}
\expandafter\def\csname SGACapture@EN-VAR-140\endcsname{%
}
\expandafter\def\csname SGARestore@EN-VAR-140\endcsname{%
}
\expandafter\def\csname SGACapture@EN-VAR-141\endcsname{%
}
\expandafter\def\csname SGARestore@EN-VAR-141\endcsname{%
}
\expandafter\def\csname SGACapture@EN-VAR-142\endcsname{%
}
\expandafter\def\csname SGARestore@EN-VAR-142\endcsname{%
}
\expandafter\def\csname SGACapture@EN-VAR-143\endcsname{%
}
\expandafter\def\csname SGARestore@EN-VAR-143\endcsname{%
}
\expandafter\def\csname SGACapture@EN-VAR-144\endcsname{%
}
\expandafter\def\csname SGARestore@EN-VAR-144\endcsname{%
}
\expandafter\def\csname SGACapture@EN-VAR-145\endcsname{%
}
\expandafter\def\csname SGARestore@EN-VAR-145\endcsname{%
}
\expandafter\def\csname SGACapture@EN-VAR-146\endcsname{%
}
\expandafter\def\csname SGARestore@EN-VAR-146\endcsname{%
}
\expandafter\def\csname SGACapture@EN-VAR-147\endcsname{%
}
\expandafter\def\csname SGARestore@EN-VAR-147\endcsname{%
}
\expandafter\def\csname SGACapture@EN-VAR-148\endcsname{%
}
\expandafter\def\csname SGARestore@EN-VAR-148\endcsname{%
}
\expandafter\def\csname SGACapture@EN-VAR-149\endcsname{%
}
\expandafter\def\csname SGARestore@EN-VAR-149\endcsname{%
}
\expandafter\def\csname SGACapture@EN-VAR-150\endcsname{%
}
\expandafter\def\csname SGARestore@EN-VAR-150\endcsname{%
}
\expandafter\def\csname SGACapture@EN-VAR-151\endcsname{%
}
\expandafter\def\csname SGARestore@EN-VAR-151\endcsname{%
}
\expandafter\def\csname SGACapture@EN-VAR-152\endcsname{%
}
\expandafter\def\csname SGARestore@EN-VAR-152\endcsname{%
}
\expandafter\def\csname SGACapture@EN-VAR-153\endcsname{%
}
\expandafter\def\csname SGARestore@EN-VAR-153\endcsname{%
}
\expandafter\def\csname SGACapture@EN-VAR-154\endcsname{%
}
\expandafter\def\csname SGARestore@EN-VAR-154\endcsname{%
}
\expandafter\def\csname SGACapture@EN-VAR-155\endcsname{%
}
\expandafter\def\csname SGARestore@EN-VAR-155\endcsname{%
}
\expandafter\def\csname SGACapture@EN-VAR-156\endcsname{%
}
\expandafter\def\csname SGARestore@EN-VAR-156\endcsname{%
}
\expandafter\def\csname SGACapture@EN-VAR-157\endcsname{%
\global\expandafter\let\csname SGAState@EN-VAR-157@Tor\endcsname\Tor%
}
\expandafter\def\csname SGARestore@EN-VAR-157\endcsname{%
\expandafter\let\expandafter\Tor\csname SGAState@EN-VAR-157@Tor\endcsname%
}
\expandafter\def\csname SGACapture@EN-VAR-158\endcsname{%
}
\expandafter\def\csname SGARestore@EN-VAR-158\endcsname{%
}
\expandafter\def\csname SGACapture@EN-VAR-159\endcsname{%
}
\expandafter\def\csname SGARestore@EN-VAR-159\endcsname{%
}
\expandafter\def\csname SGACapture@EN-VAR-160\endcsname{%
}
\expandafter\def\csname SGARestore@EN-VAR-160\endcsname{%
}
\expandafter\def\csname SGACapture@EN-VAR-161\endcsname{%
}
\expandafter\def\csname SGARestore@EN-VAR-161\endcsname{%
}
\expandafter\def\csname SGACapture@EN-VAR-162\endcsname{%
}
\expandafter\def\csname SGARestore@EN-VAR-162\endcsname{%
}
\expandafter\def\csname SGACapture@EN-VAR-163\endcsname{%
}
\expandafter\def\csname SGARestore@EN-VAR-163\endcsname{%
}
\expandafter\def\csname SGACapture@EN-VAR-164\endcsname{%
}
\expandafter\def\csname SGARestore@EN-VAR-164\endcsname{%
}
\expandafter\def\csname SGACapture@EN-VAR-165\endcsname{%
}
\expandafter\def\csname SGARestore@EN-VAR-165\endcsname{%
}
\expandafter\def\csname SGACapture@EN-VAR-166\endcsname{%
}
\expandafter\def\csname SGARestore@EN-VAR-166\endcsname{%
}
\expandafter\def\csname SGACapture@EN-VAR-167\endcsname{%
\global\expandafter\let\csname SGAState@EN-VAR-167@parfamp\endcsname\parfamp%
}
\expandafter\def\csname SGARestore@EN-VAR-167\endcsname{%
\expandafter\let\expandafter\parfamp\csname SGAState@EN-VAR-167@parfamp\endcsname%
}
\expandafter\def\csname SGACapture@EN-VAR-168\endcsname{%
}
\expandafter\def\csname SGARestore@EN-VAR-168\endcsname{%
}
\expandafter\def\csname SGACapture@EN-VAR-169\endcsname{%
}
\expandafter\def\csname SGARestore@EN-VAR-169\endcsname{%
}
\expandafter\def\csname SGACapture@EN-VAR-170\endcsname{%
}
\expandafter\def\csname SGARestore@EN-VAR-170\endcsname{%
}
\expandafter\def\csname SGACapture@EN-VAR-171\endcsname{%
}
\expandafter\def\csname SGARestore@EN-VAR-171\endcsname{%
}
\expandafter\def\csname SGACapture@EN-VAR-172\endcsname{%
}
\expandafter\def\csname SGARestore@EN-VAR-172\endcsname{%
}
\expandafter\def\csname SGACapture@EN-VAR-173\endcsname{%
}
\expandafter\def\csname SGARestore@EN-VAR-173\endcsname{%
}
\expandafter\def\csname SGACapture@EN-VAR-174\endcsname{%
}
\expandafter\def\csname SGARestore@EN-VAR-174\endcsname{%
}
\expandafter\def\csname SGACapture@EN-VAR-175\endcsname{%
}
\expandafter\def\csname SGARestore@EN-VAR-175\endcsname{%
}
\expandafter\def\csname SGACapture@EN-VAR-176\endcsname{%
}
\expandafter\def\csname SGARestore@EN-VAR-176\endcsname{%
}
\expandafter\def\csname SGACapture@EN-VAR-177\endcsname{%
}
\expandafter\def\csname SGARestore@EN-VAR-177\endcsname{%
}
\expandafter\def\csname SGACapture@EN-VAR-178\endcsname{%
}
\expandafter\def\csname SGARestore@EN-VAR-178\endcsname{%
}
\expandafter\def\csname SGACapture@EN-VAR-179\endcsname{%
}
\expandafter\def\csname SGARestore@EN-VAR-179\endcsname{%
}
\expandafter\def\csname SGACapture@EN-VAR-180\endcsname{%
}
\expandafter\def\csname SGARestore@EN-VAR-180\endcsname{%
}
\expandafter\def\csname SGACapture@EN-VAR-181\endcsname{%
}
\expandafter\def\csname SGARestore@EN-VAR-181\endcsname{%
}
\expandafter\def\csname SGACapture@EN-VAR-182\endcsname{%
}
\expandafter\def\csname SGARestore@EN-VAR-182\endcsname{%
}
\expandafter\def\csname SGACapture@EN-VAR-183\endcsname{%
}
\expandafter\def\csname SGARestore@EN-VAR-183\endcsname{%
}
\expandafter\def\csname SGACapture@EN-VAR-184\endcsname{%
}
\expandafter\def\csname SGARestore@EN-VAR-184\endcsname{%
}
\expandafter\def\csname SGACapture@EN-VAR-185\endcsname{%
}
\expandafter\def\csname SGARestore@EN-VAR-185\endcsname{%
}
\expandafter\def\csname SGACapture@EN-VAR-186\endcsname{%
}
\expandafter\def\csname SGARestore@EN-VAR-186\endcsname{%
}
\expandafter\def\csname SGACapture@EN-VAR-187\endcsname{%
\global\expandafter\let\csname SGAState@EN-VAR-187@parf\endcsname\parf%
}
\expandafter\def\csname SGARestore@EN-VAR-187\endcsname{%
\expandafter\let\expandafter\parf\csname SGAState@EN-VAR-187@parf\endcsname%
}
\expandafter\def\csname SGACapture@EN-VAR-188\endcsname{%
}
\expandafter\def\csname SGARestore@EN-VAR-188\endcsname{%
}
\expandafter\def\csname SGACapture@EN-VAR-189\endcsname{%
\global\expandafter\let\csname SGAState@EN-VAR-189@naif\endcsname\naif%
}
\expandafter\def\csname SGARestore@EN-VAR-189\endcsname{%
\expandafter\let\expandafter\naif\csname SGAState@EN-VAR-189@naif\endcsname%
}
\expandafter\def\csname SGACapture@EN-VAR-190\endcsname{%
}
\expandafter\def\csname SGARestore@EN-VAR-190\endcsname{%
}
\expandafter\def\csname SGACapture@EN-VAR-191\endcsname{%
}
\expandafter\def\csname SGARestore@EN-VAR-191\endcsname{%
}
\expandafter\def\csname SGACapture@EN-VAR-192\endcsname{%
}
\expandafter\def\csname SGARestore@EN-VAR-192\endcsname{%
}
\expandafter\def\csname SGACapture@EN-VAR-193\endcsname{%
}
\expandafter\def\csname SGARestore@EN-VAR-193\endcsname{%
}
\expandafter\def\csname SGACapture@EN-VAR-194\endcsname{%
}
\expandafter\def\csname SGARestore@EN-VAR-194\endcsname{%
}
\expandafter\def\csname SGACapture@EN-VAR-195\endcsname{%
\global\expandafter\let\csname SGAState@EN-VAR-195@Loclib\endcsname\Loclib%
\global\expandafter\let\csname SGAState@EN-VAR-195@Ob\endcsname\Ob%
\global\expandafter\let\csname SGAState@EN-VAR-195@naif\endcsname\naif%
\global\expandafter\let\csname SGAState@EN-VAR-195@varinjlim\endcsname\varinjlim%
}
\expandafter\def\csname SGARestore@EN-VAR-195\endcsname{%
\expandafter\let\expandafter\Loclib\csname SGAState@EN-VAR-195@Loclib\endcsname%
\expandafter\let\expandafter\Ob\csname SGAState@EN-VAR-195@Ob\endcsname%
\expandafter\let\expandafter\naif\csname SGAState@EN-VAR-195@naif\endcsname%
\expandafter\let\expandafter\varinjlim\csname SGAState@EN-VAR-195@varinjlim\endcsname%
}
\expandafter\def\csname SGACapture@EN-VAR-196\endcsname{%
\global\expandafter\let\csname SGAState@EN-VAR-196@varinjlim\endcsname\varinjlim%
}
\expandafter\def\csname SGARestore@EN-VAR-196\endcsname{%
\expandafter\let\expandafter\varinjlim\csname SGAState@EN-VAR-196@varinjlim\endcsname%
}
\expandafter\def\csname SGACapture@EN-VAR-197\endcsname{%
}
\expandafter\def\csname SGARestore@EN-VAR-197\endcsname{%
}
\expandafter\def\csname SGACapture@EN-VAR-198\endcsname{%
}
\expandafter\def\csname SGARestore@EN-VAR-198\endcsname{%
}
\expandafter\def\csname SGACapture@EN-VAR-199\endcsname{%
}
\expandafter\def\csname SGARestore@EN-VAR-199\endcsname{%
}
\expandafter\def\csname SGACapture@EN-VAR-200\endcsname{%
}
\expandafter\def\csname SGARestore@EN-VAR-200\endcsname{%
}
\expandafter\def\csname SGACapture@EN-VAR-201\endcsname{%
}
\expandafter\def\csname SGARestore@EN-VAR-201\endcsname{%
}
\expandafter\def\csname SGAAuthorLabel@1\endcsname{(\textsuperscript{1})}
\expandafter\def\csname SGAAuthorLabel@2\endcsname{(\textsuperscript{2})}
\expandafter\def\csname SGAAuthorLabel@3\endcsname{(\textsuperscript{3})}
\expandafter\def\csname SGAAuthorLabel@4\endcsname{(\textsuperscript{4})}
\expandafter\def\csname SGAAppNumber@EN-VAR-001\endcsname{1}
\expandafter\def\csname SGAAppNumber@EN-VAR-002\endcsname{2}
\expandafter\def\csname SGAAppNumber@EN-VAR-003\endcsname{3}
\expandafter\def\csname SGAAppNumber@EN-VAR-004\endcsname{4}
\expandafter\def\csname SGAAppNumber@EN-VAR-005\endcsname{5}
\expandafter\def\csname SGAAppNumber@EN-VAR-006\endcsname{6}
\expandafter\def\csname SGAAppNumber@EN-VAR-007\endcsname{7}
\expandafter\def\csname SGAAppNumber@EN-VAR-008\endcsname{8}
\expandafter\def\csname SGAAppNumber@EN-VAR-009\endcsname{9}
\expandafter\def\csname SGAAppNumber@EN-VAR-010\endcsname{10}
\expandafter\def\csname SGAAppNumber@EN-VAR-011\endcsname{11}
\expandafter\def\csname SGAAppNumber@EN-VAR-012\endcsname{12}
\expandafter\def\csname SGAAppNumber@EN-VAR-013\endcsname{13}
\expandafter\def\csname SGAAppNumber@EN-VAR-014\endcsname{14}
\expandafter\def\csname SGAAppNumber@EN-VAR-015\endcsname{15}
\expandafter\def\csname SGAAppNumber@EN-VAR-016\endcsname{16}
\expandafter\def\csname SGAAppNumber@EN-VAR-017\endcsname{17}
\expandafter\def\csname SGAAppNumber@EN-VAR-018\endcsname{18}
\expandafter\def\csname SGAAppNumber@EN-VAR-019\endcsname{19}
\expandafter\def\csname SGAAppNumber@EN-VAR-020\endcsname{20}
\expandafter\def\csname SGAAppNumber@EN-VAR-021\endcsname{21}
\expandafter\def\csname SGAAppNumber@EN-VAR-022\endcsname{22}
\expandafter\def\csname SGAAppNumber@EN-VAR-023\endcsname{23}
\expandafter\def\csname SGAAppNumber@EN-VAR-024\endcsname{24}
\expandafter\def\csname SGAAppNumber@EN-VAR-025\endcsname{25}
\expandafter\def\csname SGAAppNumber@EN-VAR-026\endcsname{26}
\expandafter\def\csname SGAAppNumber@EN-VAR-027\endcsname{27}
\expandafter\def\csname SGAAppNumber@EN-VAR-028\endcsname{28}
\expandafter\def\csname SGAAppNumber@EN-VAR-029\endcsname{29}
\expandafter\def\csname SGAAppNumber@EN-VAR-030\endcsname{30}
\expandafter\def\csname SGAAppNumber@EN-VAR-031\endcsname{31}
\expandafter\def\csname SGAAppNumber@EN-VAR-032\endcsname{32}
\expandafter\def\csname SGAAppNumber@EN-VAR-033\endcsname{33}
\expandafter\def\csname SGAAppNumber@EN-VAR-034\endcsname{34}
\expandafter\def\csname SGAAppNumber@EN-VAR-035\endcsname{35}
\expandafter\def\csname SGAAppNumber@EN-VAR-036\endcsname{36}
\expandafter\def\csname SGAAppNumber@EN-VAR-037\endcsname{37}
\expandafter\def\csname SGAAppNumber@EN-VAR-038\endcsname{38}
\expandafter\def\csname SGAAppNumber@EN-VAR-039\endcsname{39}
\expandafter\def\csname SGAAppNumber@EN-VAR-040\endcsname{40}
\expandafter\def\csname SGAAppNumber@EN-VAR-041\endcsname{41}
\expandafter\def\csname SGAAppNumber@EN-VAR-042\endcsname{42}
\expandafter\def\csname SGAAppNumber@EN-VAR-043\endcsname{43}
\expandafter\def\csname SGAAppNumber@EN-VAR-044\endcsname{44}
\expandafter\def\csname SGAAppNumber@EN-VAR-045\endcsname{45}
\expandafter\def\csname SGAAppNumber@EN-VAR-046\endcsname{46}
\expandafter\def\csname SGAAppNumber@EN-VAR-047\endcsname{47}
\expandafter\def\csname SGAAppNumber@EN-VAR-048\endcsname{48}
\expandafter\def\csname SGAAppNumber@EN-VAR-049\endcsname{49}
\expandafter\def\csname SGAAppNumber@EN-VAR-050\endcsname{50}
\expandafter\def\csname SGAAppNumber@EN-VAR-051\endcsname{51}
\expandafter\def\csname SGAAppNumber@EN-VAR-052\endcsname{52}
\expandafter\def\csname SGAAppNumber@EN-VAR-053\endcsname{53}
\expandafter\def\csname SGAAppNumber@EN-VAR-054\endcsname{54}
\expandafter\def\csname SGAAppNumber@EN-VAR-055\endcsname{55}
\expandafter\def\csname SGAAppNumber@EN-VAR-056\endcsname{56}
\expandafter\def\csname SGAAppNumber@EN-VAR-057\endcsname{57}
\expandafter\def\csname SGAAppNumber@EN-VAR-058\endcsname{58}
\expandafter\def\csname SGAAppNumber@EN-VAR-059\endcsname{59}
\expandafter\def\csname SGAAppNumber@EN-VAR-060\endcsname{60}
\expandafter\def\csname SGAAppNumber@EN-VAR-061\endcsname{61}
\expandafter\def\csname SGAAppNumber@EN-VAR-062\endcsname{62}
\expandafter\def\csname SGAAppNumber@EN-VAR-063\endcsname{63}
\expandafter\def\csname SGAAppNumber@EN-VAR-064\endcsname{64}
\expandafter\def\csname SGAAppNumber@EN-VAR-065\endcsname{65}
\expandafter\def\csname SGAAppNumber@EN-VAR-066\endcsname{66}
\expandafter\def\csname SGAAppNumber@EN-VAR-067\endcsname{67}
\expandafter\def\csname SGAAppNumber@EN-VAR-068\endcsname{68}
\expandafter\def\csname SGAAppNumber@EN-VAR-069\endcsname{69}
\expandafter\def\csname SGAAppNumber@EN-VAR-070\endcsname{70}
\expandafter\def\csname SGAAppNumber@EN-VAR-071\endcsname{71}
\expandafter\def\csname SGAAppNumber@EN-VAR-072\endcsname{72}
\expandafter\def\csname SGAAppNumber@EN-VAR-073\endcsname{73}
\expandafter\def\csname SGAAppNumber@EN-VAR-074\endcsname{74}
\expandafter\def\csname SGAAppNumber@EN-VAR-075\endcsname{75}
\expandafter\def\csname SGAAppNumber@EN-VAR-076\endcsname{76}
\expandafter\def\csname SGAAppNumber@EN-VAR-077\endcsname{77}
\expandafter\def\csname SGAAppNumber@EN-VAR-078\endcsname{78}
\expandafter\def\csname SGAAppNumber@EN-VAR-079\endcsname{79}
\expandafter\def\csname SGAAppNumber@EN-VAR-080\endcsname{80}
\expandafter\def\csname SGAAppNumber@EN-VAR-081\endcsname{81}
\expandafter\def\csname SGAAppNumber@EN-VAR-082\endcsname{82}
\expandafter\def\csname SGAAppNumber@EN-VAR-083\endcsname{83}
\expandafter\def\csname SGAAppNumber@EN-VAR-084\endcsname{84}
\expandafter\def\csname SGAAppNumber@EN-VAR-085\endcsname{85}
\expandafter\def\csname SGAAppNumber@EN-VAR-086\endcsname{86}
\expandafter\def\csname SGAAppNumber@EN-VAR-087\endcsname{87}
\expandafter\def\csname SGAAppNumber@EN-VAR-088\endcsname{88}
\expandafter\def\csname SGAAppNumber@EN-VAR-089\endcsname{89}
\expandafter\def\csname SGAAppNumber@EN-VAR-090\endcsname{90}
\expandafter\def\csname SGAAppNumber@EN-VAR-091\endcsname{91}
\expandafter\def\csname SGAAppNumber@EN-VAR-092\endcsname{92}
\expandafter\def\csname SGAAppNumber@EN-VAR-093\endcsname{93}
\expandafter\def\csname SGAAppNumber@EN-VAR-094\endcsname{94}
\expandafter\def\csname SGAAppNumber@EN-VAR-095\endcsname{95}
\expandafter\def\csname SGAAppNumber@EN-VAR-096\endcsname{96}
\expandafter\def\csname SGAAppNumber@EN-VAR-097\endcsname{97}
\expandafter\def\csname SGAAppNumber@EN-VAR-098\endcsname{98}
\expandafter\def\csname SGAAppNumber@EN-VAR-099\endcsname{99}
\expandafter\def\csname SGAAppNumber@EN-VAR-100\endcsname{100}
\expandafter\def\csname SGAAppNumber@EN-VAR-101\endcsname{101}
\expandafter\def\csname SGAAppNumber@EN-VAR-102\endcsname{102}
\expandafter\def\csname SGAAppNumber@EN-VAR-103\endcsname{103}
\expandafter\def\csname SGAAppNumber@EN-VAR-104\endcsname{104}
\expandafter\def\csname SGAAppNumber@EN-VAR-105\endcsname{105}
\expandafter\def\csname SGAAppNumber@EN-VAR-106\endcsname{106}
\expandafter\def\csname SGAAppNumber@EN-VAR-107\endcsname{107}
\expandafter\def\csname SGAAppNumber@EN-VAR-108\endcsname{108}
\expandafter\def\csname SGAAppNumber@EN-VAR-109\endcsname{109}
\expandafter\def\csname SGAAppNumber@EN-VAR-110\endcsname{110}
\expandafter\def\csname SGAAppNumber@EN-VAR-111\endcsname{111}
\expandafter\def\csname SGAAppNumber@EN-VAR-112\endcsname{112}
\expandafter\def\csname SGAAppNumber@EN-VAR-113\endcsname{113}
\expandafter\def\csname SGAAppNumber@EN-VAR-114\endcsname{114}
\expandafter\def\csname SGAAppNumber@EN-VAR-115\endcsname{115}
\expandafter\def\csname SGAAppNumber@EN-VAR-116\endcsname{116}
\expandafter\def\csname SGAAppNumber@EN-VAR-117\endcsname{117}
\expandafter\def\csname SGAAppNumber@EN-VAR-118\endcsname{118}
\expandafter\def\csname SGAAppNumber@EN-VAR-119\endcsname{119}
\expandafter\def\csname SGAAppNumber@EN-VAR-120\endcsname{120}
\expandafter\def\csname SGAAppNumber@EN-VAR-121\endcsname{121}
\expandafter\def\csname SGAAppNumber@EN-VAR-122\endcsname{122}
\expandafter\def\csname SGAAppNumber@EN-VAR-123\endcsname{123}
\expandafter\def\csname SGAAppNumber@EN-VAR-124\endcsname{124}
\expandafter\def\csname SGAAppNumber@EN-VAR-125\endcsname{125}
\expandafter\def\csname SGAAppNumber@EN-VAR-126\endcsname{126}
\expandafter\def\csname SGAAppNumber@EN-VAR-127\endcsname{127}
\expandafter\def\csname SGAAppNumber@EN-VAR-128\endcsname{128}
\expandafter\def\csname SGAAppNumber@EN-VAR-129\endcsname{129}
\expandafter\def\csname SGAAppNumber@EN-VAR-130\endcsname{130}
\expandafter\def\csname SGAAppNumber@EN-VAR-131\endcsname{131}
\expandafter\def\csname SGAAppNumber@EN-VAR-132\endcsname{132}
\expandafter\def\csname SGAAppNumber@EN-VAR-133\endcsname{133}
\expandafter\def\csname SGAAppNumber@EN-VAR-134\endcsname{134}
\expandafter\def\csname SGAAppNumber@EN-VAR-135\endcsname{135}
\expandafter\def\csname SGAAppNumber@EN-VAR-136\endcsname{136}
\expandafter\def\csname SGAAppNumber@EN-VAR-137\endcsname{137}
\expandafter\def\csname SGAAppNumber@EN-VAR-138\endcsname{138}
\expandafter\def\csname SGAAppNumber@EN-VAR-139\endcsname{139}
\expandafter\def\csname SGAAppNumber@EN-VAR-140\endcsname{140}
\expandafter\def\csname SGAAppNumber@EN-VAR-141\endcsname{141}
\expandafter\def\csname SGAAppNumber@EN-VAR-142\endcsname{142}
\expandafter\def\csname SGAAppNumber@EN-VAR-143\endcsname{143}
\expandafter\def\csname SGAAppNumber@EN-VAR-144\endcsname{144}
\expandafter\def\csname SGAAppNumber@EN-VAR-145\endcsname{145}
\expandafter\def\csname SGAAppNumber@EN-VAR-146\endcsname{146}
\expandafter\def\csname SGAAppNumber@EN-VAR-147\endcsname{147}
\expandafter\def\csname SGAAppNumber@EN-VAR-148\endcsname{148}
\expandafter\def\csname SGAAppNumber@EN-VAR-149\endcsname{149}
\expandafter\def\csname SGAAppNumber@EN-VAR-150\endcsname{150}
\expandafter\def\csname SGAAppNumber@EN-VAR-151\endcsname{151}
\expandafter\def\csname SGAAppNumber@EN-VAR-152\endcsname{152}
\expandafter\def\csname SGAAppNumber@EN-VAR-153\endcsname{153}
\expandafter\def\csname SGAAppNumber@EN-VAR-154\endcsname{154}
\expandafter\def\csname SGAAppNumber@EN-VAR-155\endcsname{155}
\expandafter\def\csname SGAAppNumber@EN-VAR-156\endcsname{156}
\expandafter\def\csname SGAAppNumber@EN-VAR-157\endcsname{157}
\expandafter\def\csname SGAAppNumber@EN-VAR-158\endcsname{158}
\expandafter\def\csname SGAAppNumber@EN-VAR-159\endcsname{159}
\expandafter\def\csname SGAAppNumber@EN-VAR-160\endcsname{160}
\expandafter\def\csname SGAAppNumber@EN-VAR-161\endcsname{161}
\expandafter\def\csname SGAAppNumber@EN-VAR-162\endcsname{162}
\expandafter\def\csname SGAAppNumber@EN-VAR-163\endcsname{163}
\expandafter\def\csname SGAAppNumber@EN-VAR-164\endcsname{164}
\expandafter\def\csname SGAAppNumber@EN-VAR-165\endcsname{165}
\expandafter\def\csname SGAAppNumber@EN-VAR-166\endcsname{166}
\expandafter\def\csname SGAAppNumber@EN-VAR-167\endcsname{167}
\expandafter\def\csname SGAAppNumber@EN-VAR-168\endcsname{168}
\expandafter\def\csname SGAAppNumber@EN-VAR-169\endcsname{169}
\expandafter\def\csname SGAAppNumber@EN-VAR-170\endcsname{170}
\expandafter\def\csname SGAAppNumber@EN-VAR-171\endcsname{171}
\expandafter\def\csname SGAAppNumber@EN-VAR-172\endcsname{172}
\expandafter\def\csname SGAAppNumber@EN-VAR-173\endcsname{173}
\expandafter\def\csname SGAAppNumber@EN-VAR-174\endcsname{174}
\expandafter\def\csname SGAAppNumber@EN-VAR-175\endcsname{175}
\expandafter\def\csname SGAAppNumber@EN-VAR-176\endcsname{176}
\expandafter\def\csname SGAAppNumber@EN-VAR-177\endcsname{177}
\expandafter\def\csname SGAAppNumber@EN-VAR-178\endcsname{178}
\expandafter\def\csname SGAAppNumber@EN-VAR-179\endcsname{179}
\expandafter\def\csname SGAAppNumber@EN-VAR-180\endcsname{180}
\expandafter\def\csname SGAAppNumber@EN-VAR-181\endcsname{181}
\expandafter\def\csname SGAAppNumber@EN-VAR-182\endcsname{182}
\expandafter\def\csname SGAAppNumber@EN-VAR-183\endcsname{183}
\expandafter\def\csname SGAAppNumber@EN-VAR-184\endcsname{184}
\expandafter\def\csname SGAAppNumber@EN-VAR-185\endcsname{185}
\expandafter\def\csname SGAAppNumber@EN-VAR-186\endcsname{186}
\expandafter\def\csname SGAAppNumber@EN-VAR-187\endcsname{187}
\expandafter\def\csname SGAAppNumber@EN-VAR-188\endcsname{188}
\expandafter\def\csname SGAAppNumber@EN-VAR-189\endcsname{189}
\expandafter\def\csname SGAAppNumber@EN-VAR-190\endcsname{190}
\expandafter\def\csname SGAAppNumber@EN-VAR-191\endcsname{191}
\expandafter\def\csname SGAAppNumber@EN-VAR-192\endcsname{192}
\expandafter\def\csname SGAAppNumber@EN-VAR-193\endcsname{193}
\expandafter\def\csname SGAAppNumber@EN-VAR-194\endcsname{194}
\expandafter\def\csname SGAAppNumber@EN-VAR-195\endcsname{195}
\expandafter\def\csname SGAAppNumber@EN-VAR-196\endcsname{196}
\expandafter\def\csname SGAAppNumber@EN-VAR-197\endcsname{197}
\expandafter\def\csname SGAAppNumber@EN-VAR-198\endcsname{198}
\expandafter\def\csname SGAAppNumber@EN-VAR-199\endcsname{199}
\expandafter\def\csname SGAAppNumber@EN-VAR-200\endcsname{200}
\expandafter\def\csname SGAAppNumber@EN-VAR-201\endcsname{201}

\expandafter\def\csname SGACapture@EN-CONTRIB-VII-001\endcsname{%
}
\expandafter\def\csname SGARestore@EN-CONTRIB-VII-001\endcsname{%
}
\expandafter\def\csname SGACapture@EN-CONTRIB-VII-002\endcsname{%
}
\expandafter\def\csname SGARestore@EN-CONTRIB-VII-002\endcsname{%
}
\expandafter\def\csname SGACapture@EN-CONTRIB-VII-003\endcsname{%
\global\expandafter\let\csname SGAState@EN-CONTRIB-VII-003@Tor\endcsname\Tor%
}
\expandafter\def\csname SGARestore@EN-CONTRIB-VII-003\endcsname{%
\expandafter\let\expandafter\Tor\csname SGAState@EN-CONTRIB-VII-003@Tor\endcsname%
}
\expandafter\def\csname SGACapture@EN-CONTRIB-VII-004\endcsname{%
}
\expandafter\def\csname SGARestore@EN-CONTRIB-VII-004\endcsname{%
}
\expandafter\def\csname SGACapture@EN-CONTRIB-VII-005\endcsname{%
}
\expandafter\def\csname SGARestore@EN-CONTRIB-VII-005\endcsname{%
}
\expandafter\def\csname SGACapture@EN-CONTRIB-VII-006\endcsname{%
}
\expandafter\def\csname SGARestore@EN-CONTRIB-VII-006\endcsname{%
}
\expandafter\def\csname SGACapture@EN-CONTRIB-VII-007\endcsname{%
}
\expandafter\def\csname SGARestore@EN-CONTRIB-VII-007\endcsname{%
}
\expandafter\def\csname SGACapture@EN-CONTRIB-VII-008\endcsname{%
}
\expandafter\def\csname SGARestore@EN-CONTRIB-VII-008\endcsname{%
}
\expandafter\def\csname SGACapture@EN-CONTRIB-VII-009\endcsname{%
}
\expandafter\def\csname SGARestore@EN-CONTRIB-VII-009\endcsname{%
}
\expandafter\def\csname SGACapture@EN-CONTRIB-VII-010\endcsname{%
}
\expandafter\def\csname SGARestore@EN-CONTRIB-VII-010\endcsname{%
}
\expandafter\def\csname SGACapture@EN-CONTRIB-VII-011\endcsname{%
}
\expandafter\def\csname SGARestore@EN-CONTRIB-VII-011\endcsname{%
}
\expandafter\def\csname SGACapture@EN-CONTRIB-VII-012\endcsname{%
}
\expandafter\def\csname SGARestore@EN-CONTRIB-VII-012\endcsname{%
}
\expandafter\def\csname SGACapture@EN-CONTRIB-VII-013\endcsname{%
}
\expandafter\def\csname SGARestore@EN-CONTRIB-VII-013\endcsname{%
}
\expandafter\def\csname SGACapture@EN-CONTRIB-VII-014\endcsname{%
}
\expandafter\def\csname SGARestore@EN-CONTRIB-VII-014\endcsname{%
}
\expandafter\def\csname SGACapture@EN-CONTRIB-VII-015\endcsname{%
}
\expandafter\def\csname SGARestore@EN-CONTRIB-VII-015\endcsname{%
}
\expandafter\def\csname SGACapture@EN-CONTRIB-VII-016\endcsname{%
}
\expandafter\def\csname SGARestore@EN-CONTRIB-VII-016\endcsname{%
}
\expandafter\def\csname SGACapture@EN-CONTRIB-VII-017\endcsname{%
\global\expandafter\let\csname SGAState@EN-CONTRIB-VII-017@bfP\endcsname\bfP%
}
\expandafter\def\csname SGARestore@EN-CONTRIB-VII-017\endcsname{%
\expandafter\let\expandafter\bfP\csname SGAState@EN-CONTRIB-VII-017@bfP\endcsname%
}
\expandafter\def\csname SGACapture@EN-CONTRIB-VII-018\endcsname{%
\global\expandafter\let\csname SGAState@EN-CONTRIB-VII-018@bfP\endcsname\bfP%
}
\expandafter\def\csname SGARestore@EN-CONTRIB-VII-018\endcsname{%
\expandafter\let\expandafter\bfP\csname SGAState@EN-CONTRIB-VII-018@bfP\endcsname%
}
\expandafter\def\csname SGACapture@EN-CONTRIB-VII-019\endcsname{%
}
\expandafter\def\csname SGARestore@EN-CONTRIB-VII-019\endcsname{%
}
\expandafter\def\csname SGACapture@EN-CONTRIB-VII-020\endcsname{%
}
\expandafter\def\csname SGARestore@EN-CONTRIB-VII-020\endcsname{%
}
\expandafter\def\csname SGACapture@EN-CONTRIB-VII-021\endcsname{%
\global\expandafter\let\csname SGAState@EN-CONTRIB-VII-021@Tor\endcsname\Tor%
}
\expandafter\def\csname SGARestore@EN-CONTRIB-VII-021\endcsname{%
\expandafter\let\expandafter\Tor\csname SGAState@EN-CONTRIB-VII-021@Tor\endcsname%
}
\expandafter\def\csname SGACapture@EN-CONTRIB-VII-022\endcsname{%
\global\expandafter\let\csname SGAState@EN-CONTRIB-VII-022@Tor\endcsname\Tor%
}
\expandafter\def\csname SGARestore@EN-CONTRIB-VII-022\endcsname{%
\expandafter\let\expandafter\Tor\csname SGAState@EN-CONTRIB-VII-022@Tor\endcsname%
}
\expandafter\def\csname SGACapture@EN-CONTRIB-VII-023\endcsname{%
\global\expandafter\let\csname SGAState@EN-CONTRIB-VII-023@Tor\endcsname\Tor%
}
\expandafter\def\csname SGARestore@EN-CONTRIB-VII-023\endcsname{%
\expandafter\let\expandafter\Tor\csname SGAState@EN-CONTRIB-VII-023@Tor\endcsname%
}
\expandafter\def\csname SGACapture@EN-CONTRIB-VII-024\endcsname{%
\global\expandafter\let\csname SGAState@EN-CONTRIB-VII-024@Tor\endcsname\Tor%
}
\expandafter\def\csname SGARestore@EN-CONTRIB-VII-024\endcsname{%
\expandafter\let\expandafter\Tor\csname SGAState@EN-CONTRIB-VII-024@Tor\endcsname%
}
\expandafter\def\csname SGACapture@EN-CONTRIB-VII-025\endcsname{%
}
\expandafter\def\csname SGARestore@EN-CONTRIB-VII-025\endcsname{%
}
\expandafter\def\csname SGACapture@EN-CONTRIB-VII-026\endcsname{%
}
\expandafter\def\csname SGARestore@EN-CONTRIB-VII-026\endcsname{%
}
\expandafter\def\csname SGACapture@EN-CONTRIB-VII-027\endcsname{%
}
\expandafter\def\csname SGARestore@EN-CONTRIB-VII-027\endcsname{%
}
\expandafter\def\csname SGACapture@EN-CONTRIB-VII-028\endcsname{%
}
\expandafter\def\csname SGARestore@EN-CONTRIB-VII-028\endcsname{%
}
\expandafter\def\csname SGACapture@EN-CONTRIB-VII-029\endcsname{%
}
\expandafter\def\csname SGARestore@EN-CONTRIB-VII-029\endcsname{%
}
\expandafter\def\csname SGACapture@EN-CONTRIB-VII-030\endcsname{%
}
\expandafter\def\csname SGARestore@EN-CONTRIB-VII-030\endcsname{%
}
\expandafter\def\csname SGACapture@EN-CONTRIB-VII-031\endcsname{%
\global\expandafter\let\csname SGAState@EN-CONTRIB-VII-031@bfP\endcsname\bfP%
}
\expandafter\def\csname SGARestore@EN-CONTRIB-VII-031\endcsname{%
\expandafter\let\expandafter\bfP\csname SGAState@EN-CONTRIB-VII-031@bfP\endcsname%
}
\expandafter\def\csname SGACapture@EN-CONTRIB-VII-032\endcsname{%
}
\expandafter\def\csname SGARestore@EN-CONTRIB-VII-032\endcsname{%
}
\expandafter\def\csname SGACapture@EN-CONTRIB-VII-033\endcsname{%
\global\expandafter\let\csname SGAState@EN-CONTRIB-VII-033@bfP\endcsname\bfP%
}
\expandafter\def\csname SGARestore@EN-CONTRIB-VII-033\endcsname{%
\expandafter\let\expandafter\bfP\csname SGAState@EN-CONTRIB-VII-033@bfP\endcsname%
}
\expandafter\def\csname SGACapture@EN-CONTRIB-VII-034\endcsname{%
}
\expandafter\def\csname SGARestore@EN-CONTRIB-VII-034\endcsname{%
}
\expandafter\def\csname SGACapture@EN-CONTRIB-VII-035\endcsname{%
}
\expandafter\def\csname SGARestore@EN-CONTRIB-VII-035\endcsname{%
}
\expandafter\def\csname SGACapture@EN-CONTRIB-VII-036\endcsname{%
}
\expandafter\def\csname SGARestore@EN-CONTRIB-VII-036\endcsname{%
}
\expandafter\def\csname SGACapture@EN-CONTRIB-VII-037\endcsname{%
}
\expandafter\def\csname SGARestore@EN-CONTRIB-VII-037\endcsname{%
}
\expandafter\def\csname SGACapture@EN-CONTRIB-VII-038\endcsname{%
}
\expandafter\def\csname SGARestore@EN-CONTRIB-VII-038\endcsname{%
}
\expandafter\def\csname SGACapture@EN-CONTRIB-VII-039\endcsname{%
}
\expandafter\def\csname SGARestore@EN-CONTRIB-VII-039\endcsname{%
}
\expandafter\def\csname SGACapture@EN-CONTRIB-VII-040\endcsname{%
\global\expandafter\let\csname SGAState@EN-CONTRIB-VII-040@bfP\endcsname\bfP%
}
\expandafter\def\csname SGARestore@EN-CONTRIB-VII-040\endcsname{%
\expandafter\let\expandafter\bfP\csname SGAState@EN-CONTRIB-VII-040@bfP\endcsname%
}
\expandafter\def\csname SGACapture@EN-CONTRIB-VII-041\endcsname{%
\global\expandafter\let\csname SGAState@EN-CONTRIB-VII-041@bfP\endcsname\bfP%
}
\expandafter\def\csname SGARestore@EN-CONTRIB-VII-041\endcsname{%
\expandafter\let\expandafter\bfP\csname SGAState@EN-CONTRIB-VII-041@bfP\endcsname%
}
\expandafter\def\csname SGACapture@EN-CONTRIB-VII-042\endcsname{%
\global\expandafter\let\csname SGAState@EN-CONTRIB-VII-042@bfP\endcsname\bfP%
}
\expandafter\def\csname SGARestore@EN-CONTRIB-VII-042\endcsname{%
\expandafter\let\expandafter\bfP\csname SGAState@EN-CONTRIB-VII-042@bfP\endcsname%
}
\expandafter\def\csname SGACapture@EN-CONTRIB-VII-043\endcsname{%
}
\expandafter\def\csname SGARestore@EN-CONTRIB-VII-043\endcsname{%
}
\expandafter\def\csname SGACapture@EN-CONTRIB-VII-044\endcsname{%
\global\expandafter\let\csname SGAState@EN-CONTRIB-VII-044@bfP\endcsname\bfP%
}
\expandafter\def\csname SGARestore@EN-CONTRIB-VII-044\endcsname{%
\expandafter\let\expandafter\bfP\csname SGAState@EN-CONTRIB-VII-044@bfP\endcsname%
}
\expandafter\def\csname SGACapture@EN-CONTRIB-VII-045\endcsname{%
\global\expandafter\let\csname SGAState@EN-CONTRIB-VII-045@bfP\endcsname\bfP%
}
\expandafter\def\csname SGARestore@EN-CONTRIB-VII-045\endcsname{%
\expandafter\let\expandafter\bfP\csname SGAState@EN-CONTRIB-VII-045@bfP\endcsname%
}
\expandafter\def\csname SGACapture@EN-CONTRIB-VII-046\endcsname{%
}
\expandafter\def\csname SGARestore@EN-CONTRIB-VII-046\endcsname{%
}
\expandafter\def\csname SGACapture@EN-CONTRIB-VII-047\endcsname{%
}
\expandafter\def\csname SGARestore@EN-CONTRIB-VII-047\endcsname{%
}
\expandafter\def\csname SGACapture@EN-CONTRIB-VII-048\endcsname{%
}
\expandafter\def\csname SGARestore@EN-CONTRIB-VII-048\endcsname{%
}
\expandafter\def\csname SGACapture@EN-CONTRIB-VII-049\endcsname{%
}
\expandafter\def\csname SGARestore@EN-CONTRIB-VII-049\endcsname{%
}
\expandafter\def\csname SGACapture@EN-CONTRIB-VII-050\endcsname{%
}
\expandafter\def\csname SGARestore@EN-CONTRIB-VII-050\endcsname{%
}
\expandafter\def\csname SGACapture@EN-CONTRIB-VII-051\endcsname{%
}
\expandafter\def\csname SGARestore@EN-CONTRIB-VII-051\endcsname{%
}
\expandafter\def\csname SGACapture@EN-CONTRIB-VII-052\endcsname{%
}
\expandafter\def\csname SGARestore@EN-CONTRIB-VII-052\endcsname{%
}
\expandafter\def\csname SGACapture@EN-CONTRIB-VII-053\endcsname{%
}
\expandafter\def\csname SGARestore@EN-CONTRIB-VII-053\endcsname{%
}
\expandafter\def\csname SGACapture@EN-CONTRIB-VII-054\endcsname{%
}
\expandafter\def\csname SGARestore@EN-CONTRIB-VII-054\endcsname{%
}
\expandafter\def\csname SGACapture@EN-CONTRIB-VII-055\endcsname{%
}
\expandafter\def\csname SGARestore@EN-CONTRIB-VII-055\endcsname{%
}
\expandafter\def\csname SGACapture@EN-CONTRIB-VII-056\endcsname{%
}
\expandafter\def\csname SGARestore@EN-CONTRIB-VII-056\endcsname{%
}
\expandafter\def\csname SGACapture@EN-CONTRIB-VII-057\endcsname{%
}
\expandafter\def\csname SGARestore@EN-CONTRIB-VII-057\endcsname{%
}
\expandafter\def\csname SGACapture@EN-CONTRIB-VII-058\endcsname{%
}
\expandafter\def\csname SGARestore@EN-CONTRIB-VII-058\endcsname{%
}
\expandafter\def\csname SGACapture@EN-CONTRIB-VII-059\endcsname{%
}
\expandafter\def\csname SGARestore@EN-CONTRIB-VII-059\endcsname{%
}
\expandafter\def\csname SGACapture@EN-CONTRIB-VII-060\endcsname{%
}
\expandafter\def\csname SGARestore@EN-CONTRIB-VII-060\endcsname{%
}
\expandafter\def\csname SGACapture@EN-CONTRIB-VII-061\endcsname{%
}
\expandafter\def\csname SGARestore@EN-CONTRIB-VII-061\endcsname{%
}
\expandafter\def\csname SGACapture@EN-CONTRIB-VII-062\endcsname{%
}
\expandafter\def\csname SGARestore@EN-CONTRIB-VII-062\endcsname{%
}
\expandafter\def\csname SGACapture@EN-CONTRIB-VII-063\endcsname{%
}
\expandafter\def\csname SGARestore@EN-CONTRIB-VII-063\endcsname{%
}
\expandafter\def\csname SGACapture@EN-CONTRIB-VII-064\endcsname{%
}
\expandafter\def\csname SGARestore@EN-CONTRIB-VII-064\endcsname{%
}
\expandafter\def\csname SGACapture@EN-CONTRIB-VII-065\endcsname{%
}
\expandafter\def\csname SGARestore@EN-CONTRIB-VII-065\endcsname{%
}
\expandafter\def\csname SGACapture@EN-CONTRIB-VII-066\endcsname{%
}
\expandafter\def\csname SGARestore@EN-CONTRIB-VII-066\endcsname{%
}
\expandafter\def\csname SGACapture@EN-CONTRIB-VII-067\endcsname{%
\global\expandafter\let\csname SGAState@EN-CONTRIB-VII-067@Fil\endcsname\Fil%
}
\expandafter\def\csname SGARestore@EN-CONTRIB-VII-067\endcsname{%
\expandafter\let\expandafter\Fil\csname SGAState@EN-CONTRIB-VII-067@Fil\endcsname%
}
\expandafter\def\csname SGACapture@EN-CONTRIB-VII-068\endcsname{%
\global\expandafter\let\csname SGAState@EN-CONTRIB-VII-068@Fil\endcsname\Fil%
}
\expandafter\def\csname SGARestore@EN-CONTRIB-VII-068\endcsname{%
\expandafter\let\expandafter\Fil\csname SGAState@EN-CONTRIB-VII-068@Fil\endcsname%
}
\expandafter\def\csname SGACapture@EN-CONTRIB-VII-069\endcsname{%
}
\expandafter\def\csname SGARestore@EN-CONTRIB-VII-069\endcsname{%
}
\expandafter\def\csname SGACapture@EN-CONTRIB-VII-070\endcsname{%
}
\expandafter\def\csname SGARestore@EN-CONTRIB-VII-070\endcsname{%
}
\expandafter\def\csname SGACapture@EN-CONTRIB-VII-071\endcsname{%
}
\expandafter\def\csname SGARestore@EN-CONTRIB-VII-071\endcsname{%
}
\expandafter\def\csname SGACapture@EN-CONTRIB-VII-072\endcsname{%
}
\expandafter\def\csname SGARestore@EN-CONTRIB-VII-072\endcsname{%
}
\expandafter\def\csname SGACapture@EN-CONTRIB-VII-073\endcsname{%
}
\expandafter\def\csname SGARestore@EN-CONTRIB-VII-073\endcsname{%
}
\expandafter\def\csname SGACapture@EN-CONTRIB-VII-074\endcsname{%
}
\expandafter\def\csname SGARestore@EN-CONTRIB-VII-074\endcsname{%
}
\expandafter\def\csname SGACapture@EN-CONTRIB-VII-075\endcsname{%
}
\expandafter\def\csname SGARestore@EN-CONTRIB-VII-075\endcsname{%
}
\expandafter\def\csname SGACapture@EN-CONTRIB-VII-076\endcsname{%
\global\expandafter\let\csname SGAState@EN-CONTRIB-VII-076@Fil\endcsname\Fil%
}
\expandafter\def\csname SGARestore@EN-CONTRIB-VII-076\endcsname{%
\expandafter\let\expandafter\Fil\csname SGAState@EN-CONTRIB-VII-076@Fil\endcsname%
}
\expandafter\def\csname SGACapture@EN-CONTRIB-VII-077\endcsname{%
\global\expandafter\let\csname SGAState@EN-CONTRIB-VII-077@Fil\endcsname\Fil%
}
\expandafter\def\csname SGARestore@EN-CONTRIB-VII-077\endcsname{%
\expandafter\let\expandafter\Fil\csname SGAState@EN-CONTRIB-VII-077@Fil\endcsname%
}
\expandafter\def\csname SGACapture@EN-CONTRIB-VII-078\endcsname{%
}
\expandafter\def\csname SGARestore@EN-CONTRIB-VII-078\endcsname{%
}
\expandafter\def\csname SGACapture@EN-CONTRIB-VII-079\endcsname{%
}
\expandafter\def\csname SGARestore@EN-CONTRIB-VII-079\endcsname{%
}
\expandafter\def\csname SGACapture@EN-CONTRIB-VII-080\endcsname{%
}
\expandafter\def\csname SGARestore@EN-CONTRIB-VII-080\endcsname{%
}
\expandafter\def\csname SGACapture@EN-CONTRIB-VII-081\endcsname{%
\global\expandafter\let\csname SGAState@EN-CONTRIB-VII-081@Fil\endcsname\Fil%
}
\expandafter\def\csname SGARestore@EN-CONTRIB-VII-081\endcsname{%
\expandafter\let\expandafter\Fil\csname SGAState@EN-CONTRIB-VII-081@Fil\endcsname%
}
\expandafter\def\csname SGAAppNumber@EN-CONTRIB-VII-001\endcsname{VII 1}
\expandafter\def\csname SGAAppNumber@EN-CONTRIB-VII-002\endcsname{VII 2}
\expandafter\def\csname SGAAppNumber@EN-CONTRIB-VII-003\endcsname{VII 3}
\expandafter\def\csname SGAAppNumber@EN-CONTRIB-VII-004\endcsname{VII 4}
\expandafter\def\csname SGAAppNumber@EN-CONTRIB-VII-005\endcsname{VII 5}
\expandafter\def\csname SGAAppNumber@EN-CONTRIB-VII-006\endcsname{VII 6}
\expandafter\def\csname SGAAppNumber@EN-CONTRIB-VII-007\endcsname{VII 7}
\expandafter\def\csname SGAAppNumber@EN-CONTRIB-VII-008\endcsname{VII 8}
\expandafter\def\csname SGAAppNumber@EN-CONTRIB-VII-009\endcsname{VII 9}
\expandafter\def\csname SGAAppNumber@EN-CONTRIB-VII-010\endcsname{VII 10}
\expandafter\def\csname SGAAppNumber@EN-CONTRIB-VII-011\endcsname{VII 11}
\expandafter\def\csname SGAAppNumber@EN-CONTRIB-VII-012\endcsname{VII 12}
\expandafter\def\csname SGAAppNumber@EN-CONTRIB-VII-013\endcsname{VII 13}
\expandafter\def\csname SGAAppNumber@EN-CONTRIB-VII-014\endcsname{VII 14}
\expandafter\def\csname SGAAppNumber@EN-CONTRIB-VII-015\endcsname{VII 15}
\expandafter\def\csname SGAAppNumber@EN-CONTRIB-VII-016\endcsname{VII 16}
\expandafter\def\csname SGAAppNumber@EN-CONTRIB-VII-017\endcsname{VII 17}
\expandafter\def\csname SGAAppNumber@EN-CONTRIB-VII-018\endcsname{VII 18}
\expandafter\def\csname SGAAppNumber@EN-CONTRIB-VII-019\endcsname{VII 19}
\expandafter\def\csname SGAAppNumber@EN-CONTRIB-VII-020\endcsname{VII 20}
\expandafter\def\csname SGAAppNumber@EN-CONTRIB-VII-021\endcsname{VII 21}
\expandafter\def\csname SGAAppNumber@EN-CONTRIB-VII-022\endcsname{VII 22}
\expandafter\def\csname SGAAppNumber@EN-CONTRIB-VII-023\endcsname{VII 23}
\expandafter\def\csname SGAAppNumber@EN-CONTRIB-VII-024\endcsname{VII 24}
\expandafter\def\csname SGAAppNumber@EN-CONTRIB-VII-025\endcsname{VII 25}
\expandafter\def\csname SGAAppNumber@EN-CONTRIB-VII-026\endcsname{VII 26}
\expandafter\def\csname SGAAppNumber@EN-CONTRIB-VII-027\endcsname{VII 27}
\expandafter\def\csname SGAAppNumber@EN-CONTRIB-VII-028\endcsname{VII 28}
\expandafter\def\csname SGAAppNumber@EN-CONTRIB-VII-029\endcsname{VII 29}
\expandafter\def\csname SGAAppNumber@EN-CONTRIB-VII-030\endcsname{VII 30}
\expandafter\def\csname SGAAppNumber@EN-CONTRIB-VII-031\endcsname{VII 31}
\expandafter\def\csname SGAAppNumber@EN-CONTRIB-VII-032\endcsname{VII 32}
\expandafter\def\csname SGAAppNumber@EN-CONTRIB-VII-033\endcsname{VII 33}
\expandafter\def\csname SGAAppNumber@EN-CONTRIB-VII-034\endcsname{VII 34}
\expandafter\def\csname SGAAppNumber@EN-CONTRIB-VII-035\endcsname{VII 35}
\expandafter\def\csname SGAAppNumber@EN-CONTRIB-VII-036\endcsname{VII 36}
\expandafter\def\csname SGAAppNumber@EN-CONTRIB-VII-037\endcsname{VII 37}
\expandafter\def\csname SGAAppNumber@EN-CONTRIB-VII-038\endcsname{VII 38}
\expandafter\def\csname SGAAppNumber@EN-CONTRIB-VII-039\endcsname{VII 39}
\expandafter\def\csname SGAAppNumber@EN-CONTRIB-VII-040\endcsname{VII 40}
\expandafter\def\csname SGAAppNumber@EN-CONTRIB-VII-041\endcsname{VII 41}
\expandafter\def\csname SGAAppNumber@EN-CONTRIB-VII-042\endcsname{VII 42}
\expandafter\def\csname SGAAppNumber@EN-CONTRIB-VII-043\endcsname{VII 43}
\expandafter\def\csname SGAAppNumber@EN-CONTRIB-VII-044\endcsname{VII 44}
\expandafter\def\csname SGAAppNumber@EN-CONTRIB-VII-045\endcsname{VII 45}
\expandafter\def\csname SGAAppNumber@EN-CONTRIB-VII-046\endcsname{VII 46}
\expandafter\def\csname SGAAppNumber@EN-CONTRIB-VII-047\endcsname{VII 47}
\expandafter\def\csname SGAAppNumber@EN-CONTRIB-VII-048\endcsname{VII 48}
\expandafter\def\csname SGAAppNumber@EN-CONTRIB-VII-049\endcsname{VII 49}
\expandafter\def\csname SGAAppNumber@EN-CONTRIB-VII-050\endcsname{VII 50}
\expandafter\def\csname SGAAppNumber@EN-CONTRIB-VII-051\endcsname{VII 51}
\expandafter\def\csname SGAAppNumber@EN-CONTRIB-VII-052\endcsname{VII 52}
\expandafter\def\csname SGAAppNumber@EN-CONTRIB-VII-053\endcsname{VII 53}
\expandafter\def\csname SGAAppNumber@EN-CONTRIB-VII-054\endcsname{VII 54}
\expandafter\def\csname SGAAppNumber@EN-CONTRIB-VII-055\endcsname{VII 55}
\expandafter\def\csname SGAAppNumber@EN-CONTRIB-VII-056\endcsname{VII 56}
\expandafter\def\csname SGAAppNumber@EN-CONTRIB-VII-057\endcsname{VII 57}
\expandafter\def\csname SGAAppNumber@EN-CONTRIB-VII-058\endcsname{VII 58}
\expandafter\def\csname SGAAppNumber@EN-CONTRIB-VII-059\endcsname{VII 59}
\expandafter\def\csname SGAAppNumber@EN-CONTRIB-VII-060\endcsname{VII 60}
\expandafter\def\csname SGAAppNumber@EN-CONTRIB-VII-061\endcsname{VII 61}
\expandafter\def\csname SGAAppNumber@EN-CONTRIB-VII-062\endcsname{VII 62}
\expandafter\def\csname SGAAppNumber@EN-CONTRIB-VII-063\endcsname{VII 63}
\expandafter\def\csname SGAAppNumber@EN-CONTRIB-VII-064\endcsname{VII 64}
\expandafter\def\csname SGAAppNumber@EN-CONTRIB-VII-065\endcsname{VII 65}
\expandafter\def\csname SGAAppNumber@EN-CONTRIB-VII-066\endcsname{VII 66}
\expandafter\def\csname SGAAppNumber@EN-CONTRIB-VII-067\endcsname{VII 67}
\expandafter\def\csname SGAAppNumber@EN-CONTRIB-VII-068\endcsname{VII 68}
\expandafter\def\csname SGAAppNumber@EN-CONTRIB-VII-069\endcsname{VII 69}
\expandafter\def\csname SGAAppNumber@EN-CONTRIB-VII-070\endcsname{VII 70}
\expandafter\def\csname SGAAppNumber@EN-CONTRIB-VII-071\endcsname{VII 71}
\expandafter\def\csname SGAAppNumber@EN-CONTRIB-VII-072\endcsname{VII 72}
\expandafter\def\csname SGAAppNumber@EN-CONTRIB-VII-073\endcsname{VII 73}
\expandafter\def\csname SGAAppNumber@EN-CONTRIB-VII-074\endcsname{VII 74}
\expandafter\def\csname SGAAppNumber@EN-CONTRIB-VII-075\endcsname{VII 75}
\expandafter\def\csname SGAAppNumber@EN-CONTRIB-VII-076\endcsname{VII 76}
\expandafter\def\csname SGAAppNumber@EN-CONTRIB-VII-077\endcsname{VII 77}
\expandafter\def\csname SGAAppNumber@EN-CONTRIB-VII-078\endcsname{VII 78}
\expandafter\def\csname SGAAppNumber@EN-CONTRIB-VII-079\endcsname{VII 79}
\expandafter\def\csname SGAAppNumber@EN-CONTRIB-VII-080\endcsname{VII 80}
\expandafter\def\csname SGAAppNumber@EN-CONTRIB-VII-081\endcsname{VII 81}

\begin{document}


% Title page.
\thispagestyle{empty}
\vspace*{3cm}
\begin{center}
{\Large Seminar on Algebraic Geometry of Bois Marie}\par
{\large 1966--67}\par
\vspace{1.5em}
{\Large Intersection Theory and the Riemann--Roch Theorem}\par
{\large (SGA 6)}\par
\vspace{1.5em}
a seminar directed by\par
P. Berthelot, A. Grothendieck, L. Illusie\par
with the collaboration of\par
D. Ferrand, J.-P. Jouanolou, O. Jussila, S. Kleiman, M. Raynaud, J.-P. Serre
\end{center}
\clearpage

\begin{center}\Large\bfseries In-progress collation: unverified candidate\end{center}
This candidate retains the supplied English corpus with the specifically
recorded accepted changes. The complete second scan collation, French-to-English
correction transfer, printed-page alignment and English diplomatic variants
remain unfinished. The form switch is implemented, but these files do not yet
meet the requested completion criteria.
\ifnum\SGAReaderForm=2
\par This reading form keeps the Introduction and the authors' and editors'
text. The completed reading form will omit only the new variant apparatus;
this candidate's verification notice remains visible.
\fi

% New editorial matter for the requested edition. Do not count as original front matter.
\section*{About this edition}
\textit{Théorie des intersections et théorème de Riemann--Roch},
Séminaire de Géométrie Algébrique du Bois-Marie 1966--67 (SGA~6),
directed by P.~Berthelot, A.~Grothendieck and L.~Illusie, with the collaboration
of D.~Ferrand, J.-P.~Jouanolou, O.~Jussila, S.~Kleiman, M.~Raynaud and
J.-P.~Serre. Lecture Notes in Mathematics~225, Springer, 1971.

The French layers are AI-transcribed. The English is AI-translated.
Codex's layers and corrections of August 2026, Claude Opus~5.5's check at
effort max of every numbered page on 10--13 November 2026, and the corrections
and assembly in this build are AI-produced corrections. This build is credited
to GPT-6 Astra (OpenAI, ChatGPT Pro mode). The accompanying decision records
state the actual verification status of this build.

\subsection*{Sources and authorship retained from the supplied records}
The original SGA volumes retain their authors, seminar contributors and
editors, named in the corresponding texts. The collection includes work of
Alexander Grothendieck, Michèle Raynaud, Michel Raynaud, Pierre Deligne,
Nicholas Katz, Michel Demazure, Michael Artin, Jean-Louis Verdier, Luc Illusie
and their collaborators; this list does not replace the detailed credits in
each volume.

SGA~1 retains the French retypesetting associated with Bas Edixhoven's
volunteer project, the Société mathématique de France and Michel Raynaud's
notes, from arXiv math/0206203v2.

SGA~2 is AI-translated from the corrected French TeX of the re-edition edited
by Yves Laszlo, \textit{Cohomologie locale des faisceaux cohérents et théorèmes
de Lefschetz locaux et globaux (SGA~2)}, Documents Mathématiques~4, Société
mathématique de France, 2005 (arXiv math/0511279v1), and follows its numbering
and labels. Most of the initial French typesetting was done by volunteers of
Bas Edixhoven's project (L.~Bayle, N.~Borne, O.~Brinon, J.~Buresi, M.~Chardin,
F.~Ducrot, P.~Graftiaux, F.~Han, P.~Karwasz, L.~Koelblen, D.~Madore, S.~Morel,
D.~Naie, B.~Osserman, J.~Riou and V.~Sécherre); C.~Sabbah adapted it to the
SMF format.

The SGA~2 and SGA~3 English sources acknowledge Jacob C.~Reinhold's English
translation, snapshot \texttt{e7a259f3f8608ad3edf9bf6eead3fd504dd2d23e}, used
for comparison and drafting. His stated CC BY~4.0 notice applies to that
contribution and is preserved; it does not license the underlying French
works or every contribution here. Inherited English wording remains
attributed to its human source. In SGA~3, lemmas, proofs and corrections added
in the Philippe Gille--Patrick Polo re-edition are given in Translator's
notes and Additions that cite it.

SGA~4 is AI-translated from the French LaTeX edition of SGA~4, Tomes~1--3,
distributed by Fabrice Orgogozo (source snapshot at git commit
\texttt{71766d9}), and uses its TeX configuration files. Its footnotes headed
``N.D.E.'' are that edition's editorial notes, translated.

SGA~4\textonehalf{} preserves Daniel Miller's LaTeX typesetting and its
original typesetter's note. It is not attributed to AI. His repository
(\url{https://github.com/dkmiller/sga4.5}) is no longer online; the French
text is in sga-fr, and a copy survives at NomiL/sga4.5.

Additional English translation is AI-translated; subsequent corrections,
editing, reconstruction and typesetting are AI-produced. The publisher did
not translate the text. The source-specific human credits and genuine
third-party notices remain applicable; attribution does not establish
permission or endorsement.

\subsection*{Rights}
The rights in the original work belong to its author or the author's
representatives or heirs. Where we hold rights, in the AI-produced
transcription, translation, corrections and typesetting, they are dedicated
to the public domain under CC0~1.0. This edition is published for educational,
non-profit purposes, and we claim no rights we do not hold. If you hold
rights in this work and object to its publication here, please open a GitHub
issue and we will take it down.

\clearpage
\sgaStableTarget{sga6.front.preface}\section*{Preface}
The present edition of \emph{Séminaire de Géométrie Algébrique 6} reproduces, up to minor changes, the first edition, distributed by the I.H.E.S. in the form of mimeographed exposés. The corrections made concern essentially printing errors, except in \hyperlink{sga6.e01.expose}{Exposés I} and III, where slight modifications were made by L. Illusie to a few statements which were incorrect in the primitive version. Footnotes have also been added, especially in \hyperlink{sga6.e14.expose}{Exposé XIV} by A. Grothendieck, in order to point out recent progress on questions that were still open at the time of the seminar.

Finally, note that Exposé XI, by D. Ferrand, on determinant theory for perfect complexes, was not written up and therefore will not be published; to avoid errors in references, we have left unchanged the numbering of the following exposés.

\begin{flushright}June 1971\end{flushright}

\sgaStableTarget{sga6.front.introduction}\section*{Introduction}
The purpose of the present seminar is to develop a global theory of intersections, and a Riemann--Roch formula, for arbitrary schemes. The reader will find a description of the programme of the seminar in \hyperlink{sga6.e00.expose}{Exposé 0}. The possibility in principle of proving a Riemann--Roch formula for general regular schemes, along the lines of the well-known Borel--Serre report of 1958, was already clear at that time, at least for a projective morphism. The extension to the general case is less evident; the programme in which it belongs, and which is partly carried out in this seminar, goes back to 1960. As was also the case for the duality theory of coherent sheaves (cf. R. Hartshorne, \emph{Residues and Duality}, Lecture Notes in Mathematics no.~20, Springer), an essential tool for formulating a satisfactory theory is Verdier's theory of derived categories; familiarity with it is indispensable for understanding the seminar.

Apart from this theory, the study of the seminar requires little more than a general knowledge of the foundations of algebraic geometry as set out in EGA, chapters I, II, III. We shall only occasionally refer to EGA IV, for certain facts concerning flat or smooth morphisms, which for the most part also appear in the first exposés of SGA 1. Finally, to develop some results with the full generality desirable for applications, we sometimes use the notion of a ringed site and of a ringed topos, for which we refer to SGA 4, Exposés I to IV. A reader unfamiliar with the language of sites and topoi may everywhere replace these objects by ordinary topological spaces, the objects of the topos then being replaced by open subsets of these spaces; nevertheless, we recommend that the reader assimilate the language of topoi, which provides an extremely convenient unifying principle.

The fundamental notion for the theory presented here is that of a perfect complex of modules, together with its various variants, developed in \hyperlink{sga6.e01.expose}{Exposés I}, II, III. It seems clear that these notions, also important in other contexts in algebraic geometry, notably for formulas of Lefschetz--Verdier type in cohomology with discrete coefficients (SGA 5) or coherent coefficients, will also have a role to play outside algebraic geometry, especially in the formulation of a complex analytic variant of the Riemann--Roch theorem presented here, or of suitable variants of the Atiyah--Singer index theorem. Some indications in this direction will be given in \hyperlink{sga6.e02.expose}{Exposé II}. Unfortunately, at present there is still no statement, even heuristic, which would encompass these two theorems, the first of which remains conjectural for the moment. A new idea seems to be missing, just as one is also missing in algebraic geometry for proving Riemann--Roch beyond projective hypotheses (cf. \hyperlink{sga6.e14.sec.2}{Exposé XIV, no. 2)}.

We shall make no use in this seminar of the local theory of intersections, set out in J.-P. Serre, \emph{Algèbre locale. Multiplicités} (Lecture Notes in Mathematics no. 11, Springer), and merely point out here that Serre's course had an evident influence on the genesis of the theory in 1957. Nor shall we use the theory of the Chow ring developed in the Chevalley Seminar of 1958, Exposés 2 and 3 (École Normale Supérieure). That theory is essentially tied to nonsingularity hypotheses, whereas the aim of our seminar is, on the contrary, to develop a theory of intersections on general schemes, and even on general locally ringed topoi. On this matter see the comments in \hyperlink{sga6.e14.expose}{Exposé XIV}, nos. 4 and 8, which give the relations between our theory and that of the Chow ring and raise the question of a generalization of the latter. Thus one may say that the seminar presents a coherent and self-contained theory of intersections, but one which is by no means exhaustive of the different points of view useful, indeed indispensable, in intersection theory; consequently it should be complemented by the cited exposés of Serre and Chevalley.

We have appended to the seminar, as an appendix to \hyperlink{sga6.e00.expose}{Exposé 0}, Grothendieck's 1957 report, ``Classes of sheaves and the Riemann--Roch theorem'', which served as the basis for the Borel--Serre seminar at Princeton in the same year and for their already cited article. That report sketches two proofs of the Riemann--Roch theorem for nonsingular quasi-projective varieties. The first, valid for the time being only in characteristic zero but giving, in return, a slightly more precise result in the case of an immersion, has not yet appeared in the literature (including the work of Borel--Serre and the present seminar). The report presupposes, of course, none of the other exposés of the seminar, and can even serve, like \hyperlink{sga6.e00.expose}{Exposé 0}, as an introduction to the study of the latter, especially for a reader not yet familiar with Borel--Serre. The proof just alluded to applies essentially unchanged to the case of a projective morphism of complex analytic spaces, and will doubtless also apply in arbitrary characteristic once the problem of the operations ``exterior powers'' in the derived category is solved (cf. \hyperlink{sga6.e14.expose}{Exposé XIV}, nos. 1, 2, 3). The absence of a study of this question is no doubt the most serious gap in this seminar, from the point of view we have adopted.

One will notice the absence, in the table of contents of the present seminar, of two of the participants mentioned on the cover page. One of them, J.-P. Jouanolou, took an active part in the technical elaboration of the first part of the seminar, but was prevented from taking part in the oral exposés; in particular, he is responsible for the linear-independence assertion contained in the important statement \hyperlink{sga6.e06.thm.1-1}{VI 1.1}, giving the structure of the ring $K^\circ$ of classes of locally free modules on a projective bundle, which had previously been proved only under the assumption that the base scheme has an ample invertible module. The other, J.-P. Serre, gave two oral exposés which were logically independent of the rest of the seminar, and therefore preferred to publish them in the form of a separate article (cf. J.-P. Serre, ``Grothendieck groups of split reductive group schemes'', to appear in \emph{Publications Mathématiques}, no. 34).

As in the other parts of the \emph{Séminaire de Géométrie Algébrique du Bois Marie}, the sigla SGA 1 to SGA 6 refer to the different parts of that seminar; the Roman numeral following the siglum indicates the number of the exposé, and the Arabic numerals following it correspond to the internal numbering of the exposé in question. For internal references within the present seminar the siglum SGA 6 is omitted, and references of the form I 4.9 are used, meaning: Exposé I, statement 4.9. The siglum EGA refers to the \emph{Éléments de Géométrie Algébrique} of J. Dieudonné and A. Grothendieck.

\sgaStableTarget{sga6.front.table-of-contents}\section*{Table of contents}
\begin{longtable}{@{}p{0.16\linewidth}p{0.64\linewidth}r@{}}
\toprule
Exposé & Title and author & Page \\
\midrule
0 & Outline of a programme for a theory of intersections on general schemes, by A. Grothendieck & 1 \\
Appendix & Classes of sheaves and the Riemann--Roch theorem, by A. Grothendieck & 20 \\
I & Generalities on finiteness conditions in derived categories, by L. Illusie & 78 \\
II & Existence of global resolutions, by L. Illusie & 160 \\
III & Relative finiteness conditions, by L. Illusie & 222 \\
IV & Grothendieck groups of ringed topoi, by L. Illusie & 274 \\
V & Generalities on $\lambda$-rings, by P. Berthelot & 297 \\
VI & The $K^\bullet$ of a projective bundle: computations and consequences, by P. Berthelot & 365 \\
VII & Regular immersions and the computation of $K^\bullet$ of a blow-up, by P. Berthelot & 416 \\
VIII & The Riemann--Roch theorem, by P. Berthelot & 466 \\
IX & Some computations of groups $K$. By P. Berthelot & 498 \\
X & Formalism of intersections on proper algebraic schemes, by O. Jussila, with an appendix by A. Grothendieck, ``Specialization in intersection theory'' & 519 \\
XI & Not written up & -- \\
XII & A relative representability theorem for the Picard functor, by M. Raynaud, written by S. Kleiman & 595 \\
XIII & Finiteness theorems for the Picard functor, by S. Kleiman & 616 \\
XIV & Open problems in intersection theory, by A. Grothendieck & 667 \\
\bottomrule
\end{longtable}

%%%%%%%%%%  scan idx 14  |  folio 1  %%%%%%%%%%
\SGAPrintPage{1}
\sgaStableTarget{sga6.e00.expose}\section*{\underline{Exposé 0}}
\begin{center}
\underline{\underline{OUTLINE OF A PROGRAMME FOR A THEORY OF INTERSECTIONS}}\\[3pt]
\underline{\underline{ON GENERAL SCHEMES}}\\[8pt]
\underline{by A. Grothendieck}
\end{center}

The present exposé is introductory in nature, and reading it is not logically indispensable for the study of the Seminar. It is addressed more particularly to readers familiar with the provisional version of the Riemann--Roch theorem, as it is presented in the report of Borel--Serre [2] or in the report of Grothendieck already mentioned in the Introduction (subsequently cited as [RRR]), which is reproduced as an appendix at the end of the present exposé.

\sgaStableTarget{sga6.e00.sec.1}\subsection*{1.}
Recall the Riemann--Roch formula for a proper morphism
\[
\var{f:X.\longrightarrow Y}{f:X\longrightarrow Y}{EN-VAR-001}
\]
of quasi-projective smooth schemes over a field $k$, and a coherent sheaf $F$ on $X$, whose class in the group $K(X)$ of classes of coherent sheaves on $X$ is denoted by $\cl(F)$:
\sgaStableTarget{sga6.e00.eq.1-1}\begin{equation}\tag{1.1}
\Todd(T_Y)\,\ch_Y\bigl(f_*(\cl(F))\bigr)=f_*\bigl(\Todd(T_X)\,\ch_X(F)\bigr).
\end{equation}
This formula holds in $A(Y)\otimes_{\mathbb Z}\mathbb Q$, where $A(Y)$ is the Chow ring of $Y$, the $f_*$ on the right-hand side being obtained by tensoring with $\mathbb Q$ the homomorphism ``direct image of cycles''
\[
f_*:A(X)\longrightarrow A(Y),
\]
and that on the left-hand side being the Euler--Poincaré characteristic of $F$ relative to $f$:
\[
f_*(\cl(F))=\sum_i(-1)^i\cl\bigl(R^if_*(F)\bigr)~;
\]
$\ch_X$, $\ch_Y$ denote the Chern characters on $X$, respectively on $Y$, and $T_X$, $T_Y$ the tangent bundles of $X$, respectively of $Y$. As is known, $\Todd(-)$ and $\ch(-)$ are universal polynomials with coefficients in $\mathbb Q$ in the Chern classes of the argument. Since the constant term of $\Todd(-)$ is $1$, it is an invertible element for every value of the argument, so that, after multiplication by $\Todd(T_Y)^{-1}$, formula (1.1) takes the equivalent form, more convenient for our purpose,
\sgaStableTarget{sga6.e00.eq.1-2}\begin{equation}\tag{1.2}
\var{\ch_Y\bigl(f_*(\cl(F))=f_*\bigl(\Todd(T_f)\cdot\ch_X(F)\bigr),}{\ch_Y\bigl(f_*(\cl(F))\bigr)=f_*\bigl(\Todd(T_f)\cdot\ch_X(F)\bigr),}{EN-VAR-002}
\end{equation}
where we have put
\sgaStableTarget{sga6.e00.eq.1-3}\begin{equation}\tag{1.3}
T_f=T_X-f^*(T_Y)\in K(X),
\end{equation}

%%%%%%%%%%  scan idx 15  |  folio 2  %%%%%%%%%%
\SGAPrintPage{2}
so that $T_f$ plays the role of a \underline{virtual relative tangent bundle} of $X$ over $Y$. When the morphism $f$ is smooth (i.e. its tangent map is everywhere surjective), we simply have
\[
T_f=T_{X/Y}\qquad\text{(tangent bundle along the fibres)},
\]
whereas when $f:X\to Y$ is an immersion, we find
\[
T_f=-\check N_{X/Y},
\]
where $\check N_{X/Y}$ denotes the normal sheaf of $X$ in $Y$.

One of the main aims of the Seminar is to generalize (1.2) simultaneously in two directions:
\begin{enumerate}[label=\alph*)]
\item To remove the hypothesis of the existence of a base field $k$.
\item To replace the regularity hypotheses on $Y,X$ by a hypothesis of ``local regularity'' of $f$.
\end{enumerate}
Finally, along the way, we shall also consider the problem:
\begin{enumerate}[label=\alph*),start=3]
\item To eliminate the quasi-projectivity hypotheses which, in the absence of a base field, are expressed by the existence of ample invertible Modules on $X$ and on $Y$.
\end{enumerate}

\sgaStableTarget{sga6.e00.sec.2}\subsection*{2.}
Let us first examine generalization a), while retaining the hypotheses of regularity and of the existence of ample invertible Modules on $X$ and on $Y$.

The definition of $K(X)$, $K(Y)$ and of the homomorphism $f_*:K(X)\to K(Y)$ then presents no new problem, because $X$ and $Y$ are regular. The most natural way of giving a meaning to (1.2) therefore seems to consist in defining \underline{Chow rings} $A(X),A(Y)$ and a group homomorphism
\[
f_*:A(X)\longrightarrow A(Y),
\]
in establishing a theory of Chern classes, providing maps
\[
c_i:K(X)\longrightarrow A(X),
\]
and similarly for $Y$, and finally in describing a virtual relative tangent bundle element
\[
T_f\in K(X).
\]

%%%%%%%%%%  scan idx 16  |  folio 3  %%%%%%%%%%
\SGAPrintPage{3}
\sgaStableTarget{sga6.e00.sec.2.1}\subsubsection*{2.1}
As for the latter, a definition presents itself fairly evidently. We use the fact that, by the hypothesis of the existence of an ample invertible Module on $X$ and on $Y$, the morphism $f$ can be factored as
\sgaStableTarget{sga6.e00.eq.2-1}\begin{equation}\tag{2.1}
X\xrightarrow{i}X'\xrightarrow{f'}Y,
\end{equation}
where $i$ is a closed immersion and $f'$ is a smooth morphism; for example, we may take $X'=\mathbb P^r_Y$ for suitable $r$. We then put
\sgaStableTarget{sga6.e00.opening.eq.2-2}\begin{equation}\tag{2.2}
\check T_f=\cl({\underline\Omega}^1_{X'/Y})-\cl({\underline N}_{X/X'}),
\end{equation}
where ${\underline\Omega}^1_{X'/Y}$ is the locally free Module of relative $1$-differentials of $X'$ over $Y$ (or relative cotangent Module of $X'$ over $Y$), and ${\underline N}_{X/X'}$ is the conormal Module of $X$ in $X'$, equal by definition to ${\underline J}/{\underline J}^2$, where ${\underline J}$ is the Ideal on $X'$ defining the closed subscheme $X$; since $X$ and $X'$ are regular, we also conclude that ${\underline N}_{X/X'}$ is locally free. Finally, $\check T_f$ denotes the dual of $T_f$. We readily verify (Exp. VIII) that the element $T_f$ defined by (2.1) does not depend on the chosen factorization of $f$.

\sgaStableTarget{sga6.e00.sec.2.2}\subsubsection*{2.2.}
As for a definition, under the conditions considered, of a Chow ring and a corresponding theory of Chern classes, although it is plausible that such a theory should be developable on the model of the theory already known for smooth quasi-projective algebraic schemes, it seems that this work has not been done at present, and a fortiori it has not been undertaken in the Seminar. Instead, we introduce another ring which will play a role analogous to that of the Chow ring, and whose tensor product with $Q$ is indeed isomorphic to $A(X)\otimes_{\mathbb Z}\mathbb Q$ when $X$ is quasi-projective and smooth over a field $k$. This ring is the associated graded ring of $K(X)$ for a suitable filtration of that ring. A first natural filtration of $K(X)$, the one considered in [RRR], is the ``\underline{topological filtration}'' by the codimension of the support of coherent sheaves, $\operatorname{Filt}^{i}_{\mathrm{top}}K(X)$ being generated by the $\cl(F)$ with $F$ coherent and $\operatorname{codim}(\operatorname{supp}F,X)\ge i$. A first question is the compatibility of this filtration with the ring structure, a question resolved in [RRR], when $X$ is smooth over a field $k$, by means of Chow's ``moving lemma''. Granting this point, we therefore obtain an associated graded ring, which we shall denote here by $\Gr^{\bullet}_{\mathrm{top}}(X)$. In [RRR] a surjective ring homomorphism is defined:
\[
\phi:A(X)\longrightarrow \Gr^{\bullet}_{\mathrm{top}}(X),
\]

%%%%%%%%%%  scan idx 17  |  folio 4  %%%%%%%%%%
\SGAPrintPage{4}
(sending the class of a prime cycle $Z$ to $\cl(\underline O_Z)$,) whose kernel subsequently turns out (Exp. XIV 4.2) to be a torsion subgroup. Moreover, [RRR] shows how the images in $\Gr^{\bullet}_{\mathrm{top}}(X)$ of the Chern classes $c_i(F)$ can be calculated directly in terms of $\cl(F)\in K(X)$ and the exterior-power operations $\lambda^i$ in $K(X)$.

This shows that, for $X$ a noetherian regular scheme equipped with an ample invertible Module, the ring $\Gr^{\bullet}_{\mathrm{top}}(X)$ is a convenient substitute for the hypothetical Chow ring $A(X)$; it also has the advantage over the latter that certain calculations in [RRR] take place directly in this ring and lead to important formulas in it, formulas which, even for $X$ projective and smooth over a field, have not at present been established in $A(X)$ itself (Exp. XIV, 4.3 and 4.4). To justify this point of view completely, however, one would have to resolve the problem mentioned above of the compatibility of the topological filtration of $K(X)$ with its ring structure, a problem which seems essentially connected with Chow's ``moving lemma'', which is itself the essential technical ingredient of the theory of the Chow ring that we have claimed to be able to avoid.

\sgaStableTarget{sga6.e00.sec.2.3}\subsubsection*{2.3.}
To avoid this problem, we give $K(X)$ a second filtration, finer than the topological filtration, which can be described in purely algebraic terms from the augmented $\lambda$-ring structure of $K(X)$. The definition of this filtration is suggested by the expression given in [RRR] for the Chern classes of an $x\in K(X)$ in $\Gr^{\bullet}_{\mathrm{top}}(X)$,
\sgaStableTarget{sga6.e00.eq.2-3}\begin{equation}\tag{2.3}
c_i(x)=\gamma^i(x-\epsilon(x))\mod \operatorname{Filt}^{i+1}_{\mathrm{top}}K(X),
\end{equation}
where $\gamma^i$ denotes a certain linear combination of the $\lambda^j$ \var{($j=i$)}{($j\leq i$)}{EN-VAR-003} made explicit in loc. cit., $\epsilon$ being the augmentation. We show that $\gamma^i(x-\epsilon(x))$ does have topological filtration $\ge i$, which gives a meaning to formula (2.3). We then define the $\lambda$-\underline{filtration} of $K(X)$ by the ideals $\operatorname{Filt}^i K(X)$, where $\operatorname{Filt}^i K(X)$ consists of sums of expressions of the form
\[
\gamma^{i_1}(x_1-\epsilon(x_1))\cdots \gamma^{i_k}(x_k-\epsilon(x_k)),\qquad i_1+\cdots+i_k\ge i.
\]
The associated graded ring of this filtration will be denoted by $\Gr^\bullet(X)$, and it is this ring that will replace the Chow ring of $X$ in the Seminar. We can moreover show that, in the case considered here, the canonical homomorphism
\sgaStableTarget{sga6.e00.eq.2-4}\begin{equation}\tag{2.4}
\Gr^\bullet(X)\longrightarrow \Gr^{\bullet}_{\mathrm{top}}(X)
\end{equation}

%%%%%%%%%%  scan idx 18  |  folio 5  %%%%%%%%%%
\SGAPrintPage{5}
is an isomorphism modulo torsion, which justifies the point of view that this ring tensored with $\mathbb Q$ can play in every respect the same role as the Chow ring tensored with $\mathbb Q$. On the other hand, we have done exactly what is needed for formula (2.3) (with $\operatorname{Filt}^i_{\mathrm{top}}$ replaced by $\operatorname{Filt}^i$) to define a theory of Chern classes with values in $\Gr^\bullet(X)$, from which the Chern classes in $\Gr^{\bullet}_{\mathrm{top}}(X)$ are deduced using the homomorphism (2.4). We show (Exp. V), using the fact (Exp. VI) that $K(X)$ is a \underline{special} $\lambda$-ring in the sense of [RRR] (i.e. a $\lambda$-ring in the sense of Exp. V), that this notion of Chern class satisfies all the usual formal properties reviewed in [RRR].

\sgaStableTarget{sga6.e00.sec.2.4}\subsubsection*{2.4.}
We note a very appreciable advantage of the definition of $\Gr^\bullet(X)$ over that of $A(X)$ and $\Gr^{\bullet}_{\mathrm{top}}(X)$: it remains meaningful for an absolutely arbitrary scheme $X$, and more generally for an arbitrary locally ringed topos $X$, taking for $K(X)$ the group of classes constructed from locally free Modules on $X$, which is still an augmented $\lambda$-ring, hence equipped with a corresponding filtration allowing an associated graded $\Gr^\bullet(X)$ to be defined. As the price of this generality, however, we must observe that, when we are unwilling to disregard torsion phenomena, i.e. to tensor with $\mathbb Q$, the properties of $\Gr^\bullet(X)$ are less satisfactory in various respects than those of its two competitors (Exp. XIV no.~4). Thus, supposing again that $f:X\to Y$ is a proper morphism of noetherian regular schemes equipped with ample invertible Modules, it is not true in general that the homomorphism $f_*:K(X)\longrightarrow K(Y)$ maps $\operatorname{Filt}^i K(X)$ into $\operatorname{Filt}^{i-d} K(Y)$, where $d$ is the relative dimension of $X$ over $Y$, and hence defines, by passage to the associated graded rings, a homomorphism of degree $-d$ from $\Gr^\bullet(X)$ to $\Gr^\bullet(Y)$. However, we can prove (Exp. VII) that this holds after tensoring with $\mathbb Q$, so that we obtain a homomorphism of degree $-d$
\sgaStableTarget{sga6.e00.eq.2-5}\begin{equation}\tag{2.5}
f_*:\Gr^\bullet(X)_{\mathbb Q}\longrightarrow \Gr^\bullet(Y)_{\mathbb Q},
\end{equation}
playing the role of the homomorphism ``direct image of cycles'' $A(X)\to A(Y)$.

With this understood, we see that we have all the structural elements needed to give a meaning to formula (1.2) without postulating the existence of a base field, $f:X\to Y$ being a proper morphism of noetherian regular schemes admitting ample invertible Modules.

%%%%%%%%%%  scan idx 19  |  folio 6  %%%%%%%%%%
\SGAPrintPage{6}
\sgaStableTarget{sga6.e00.sec.3}\subsection*{3.}
Let us now show how we can still give a meaning to formula (1.2) when we abandon the regularity hypothesis on $X,Y$, replacing it by a hypothesis of a local nature on the morphism $f$, which we shall specify.

\sgaStableTarget{sga6.e00.sec.3.1}\subsubsection*{3.1}
A first difficulty arises because the observation that served as the starting point of the theory [RRR], namely that the group $K(X)$ of classes of sheaves on $X$ is the same whether it is defined in terms of coherent sheaves or in terms of locally free sheaves, is essentially tied to the regularity hypothesis on $X$. When $X$ is a general scheme (say noetherian, for simplicity), we must distinguish these two groups, which we shall denote respectively by $K_\bullet(X)$ and $K^\bullet(X)$, the first being covariant in $X$ for proper morphisms, and the second contravariant in $X$ for arbitrary morphisms. We note that $K^\bullet(X)$ has a natural ring structure (and even an augmented $\lambda$-ring structure), whereas $K_\bullet(X)$ in general no longer has a ring structure, but only a module structure over $K^\bullet(X)$. The homomorphism $x\to x\,\cl_\bullet(\underline O_X)$ (where $\cl_\bullet(\underline F)$ denotes the class of $\underline F$ in $K_\bullet(X)$) from $K^\bullet(X)$ to $K_\bullet(X)$ is in general no longer an isomorphism. \uline{For our formulation of the generalized Riemann--Roch theorem, we shall work exclusively with $K^\bullet(X)$, $K^\bullet(Y)$, to the exclusion of $K_\bullet(X)$, $K_\bullet(Y)$}. As we noted in 2.3., these augmented $\lambda$-rings, equipped with their natural $\lambda$-filtration, give rise to graded rings $\Gr^\bullet(X)$, $\Gr^\bullet(Y)$, and ``Chern class'' homomorphisms
\[
c_i:K^\bullet(X)\longrightarrow \Gr^\bullet(X),
\]
and similarly for $Y$. The difficulty, however, arises in defining a direct-image homomorphism
\sgaStableTarget{sga6.e00.eq.3-1}\begin{equation}\tag{3.1}
f_*:K^\bullet(X)\longrightarrow K^\bullet(Y),
\end{equation}
which can no longer be defined by the naive formula
\[
\var{f_*(\cl^\bullet(\underline E))=\sum_i(-1)^i\cl^\bullet\bigl(R^if_*(\underline F)\bigr),}{f_*(\cl^\bullet(\underline F))=\sum_i(-1)^i\cl^\bullet\bigl(R^if_*(\underline F)\bigr),}{EN-VAR-004}
\]
because, even when $\underline F$ is locally free, the Modules $R^if_*(\underline F)$ are in general no longer locally free, nor even of finite tor-dimension, and the preceding formula loses its meaning. A solution to this difficulty is provided by the point of view of \underline{derived categories}. Indeed, although the $R^if_*(\underline F)$ are in general ``bad'' sheaves, say of infinite tor-dimension,

%%%%%%%%%%  scan idx 20  |  folio 7  %%%%%%%%%%
\SGAPrintPage{7}
to which we cannot associate a class in $K^\bullet(Y)$, the complex $Rf_*(\underline F)$ on $Y$, the direct image of $\underline F$ in the sense of derived categories, whose cohomology sheaves are the $R^if_*(\underline F)$, nevertheless tends to be excellent. More precisely, suppose that $\underline F$ has finite tor-dimension over $\underline O_Y$, which is the case when $\underline O_X$ has finite tor-dimension over $\underline O_Y$ and, in addition, $F$ is locally free. We then show (Exp. III) that the complex $Rf_*(F)$ is ``\underline{perfect}'', by which we mean here: it has coherent cohomology sheaves (which we already knew in any case) \underline{and} finite global tor-dimension; or, equivalently: on every affine open $U$ of $Y$, this complex is isomorphic, in the derived category $D(U)$, to a complex bounded in degree and locally free of finite type in each degree. From the point of view of Homological Algebra, perfect complexes are ``as good'' as locally free Modules, of which they are the natural generalization. They moreover form a ``triangulated subcategory'' of the derived category $D(Y)$, i.e. a subcategory stable under cones (or ``mapping cylinders'') of morphisms. On the other hand, to every triangulated category $C$ we associate a group $K(C)$ of classes, a natural variant of the construction of the group of classes of an abelian category, by taking the free group generated by the objects of $C$, modulo the relations
\[
\cl(K)-\cl(K')-\cl(K'')
\]
for every distinguished triangle
\[
\begin{tikzcd}[column sep=large]
 & K'' \arrow[dl] & \\
K' \arrow[rr] & & K \arrow[ul] .
\end{tikzcd}
\]
What replaces the fundamental observation of [RRR] in the present case is the fact that, when the scheme $Y$ admits an ample invertible Module, every perfect complex $L_\bullet$ on $Y$ is globally (and not only locally) isomorphic in $D(Y)$ to a complex $L'_\bullet$ bounded in degree and locally free of finite type in each degree (Exp. II); it follows readily (Exp. IV) that the canonical homomorphism
\[
\var{K^\bullet(X)\longrightarrow K(\Parf(Y)),}{K^\bullet(Y)\longrightarrow K(\Parf(Y)),}{EN-VAR-005}
\]
where $\Parf(Y)$ is the triangulated category of perfect complexes on $Y$, is an isomorphism. Thus to every perfect complex $L_\bullet$ on $Y$ we can associate a class
\[
\cl^\bullet(L_\bullet)\in K^\bullet(Y),
\]

%%%%%%%%%%  scan idx 21  |  folio 8  %%%%%%%%%%
\SGAPrintPage{8}
which is moreover simply
\[
\cl^\bullet(L_\bullet)=\sum_i(-1)^i\cl^\bullet(L'_i).
\]
where $L'_\bullet$ is as above, so that the $L'_i$ are indeed locally free Modules on $Y$ whose classes we can take in $K^\bullet(Y)$.

With this understood, we define a homomorphism (3.1) whenever $f:X\to Y$ is a \underline{proper morphism of finite tor-dimension} of noetherian schemes such that $Y$ admits an ample invertible Module, simply by putting
\sgaStableTarget{sga6.e00.eq.3-2}\begin{equation}\tag{3.2}
f_*(\cl^\bullet(\underline F))=\cl^\bullet(Rf_*(\underline F))
\end{equation}
for every locally free Module $\underline F$ on $X$ (so that $Rf_*(\underline F)$ is indeed a perfect complex on $Y$, and the right-hand side is defined). When $X$ also admits an ample invertible Module, $K^\bullet(X)$ can likewise be interpreted as $K(\Parf(X))$, and the homomorphism (3.1) is precisely the one associated, by functoriality, with the exact functor of triangulated categories induced by
\[
Rf_*:\Parf(X)\longrightarrow \Parf(Y).
\]

\sgaStableTarget{sga6.e00.sec.3.2}\subsubsection*{3.2.}
To give a meaning to (1.2), we must still impose on $f$ a hypothesis allowing a meaning to be given to the homomorphism (2.5)
\sgaStableTarget{sga6.e00.eq.3-3}\begin{equation}\tag{3.3}
f_*:\Gr^\bullet(X)_{\mathbb Q}\longrightarrow \Gr^\bullet(Y)_{\mathbb Q},
\end{equation}
i.e. ensuring that the homomorphism (3.1) just defined is (up to a shift and modulo torsion) compatible with the filtrations of these rings; and finally we must still define a ``virtual relative tangent bundle'' element
\[
T_f\in K^\bullet(X).
\]
For the latter, returning to the definition proposed in 2.1, we see that we should make a hypothesis ensuring that (with the notation of loc. cit.) the conormal Module $\underline N_{X/X'}$ is locally free. The most natural hypothesis for this purpose is that $i$ is a ``regular immersion'', i.e. identifies $X$ with a subscheme of $X'$ defined by an Ideal which locally can be defined by an $\underline O_{X'}$-regular sequence of sections of $\underline O_{X'}$. This hypothesis, of a purely local nature, does not in fact depend on the choice of the factorization of $f$ into an immersion followed by a smooth morphism (Exp. VIII), and, as above, the same holds for \var{the class $T_f$ defined by formula (2.1)}{the class $T_f$ defined by formula (2.2)}{EN-VAR-067}. When the hypothesis just considered holds, we say that the morphism $f$ is a ``\underline{locally complete intersection}'' \underline{morphism}\SGAAuthorFootnote{raisedstar}{We shall also allow ourselves to omit the word ``locally'' in this terminology, which entails no risk of confusion.}. Such a morphism is moreover locally of finite tor-dimension.

%%%%%%%%%%  scan idx 22  |  folio 9  %%%%%%%%%%
\SGAPrintPage{9}
For a locally complete intersection morphism $f$, we introduce a natural notion of ``virtual relative dimension'' $d(f)$ by the formula
\[
d(f)(x)=\dim_x\bigl(f'^{-1}f(x)\bigr)-\codim_x(X,X').
\]
This is an integer-valued function, locally constant in the point $x\in X$, hence constant when $X$ is connected; it likewise does not depend on the choice of the factorization (2.1). When this virtual relative dimension $d$ is constant, $f$ is proper, and $X$ and $Y$ admit ample invertible Modules, we can show (Exp. VIII), although this is far from a trivial result, that
\sgaStableTarget{sga6.e00.eq.3-4}\begin{equation}\tag{3.4}
f_*\bigl(\operatorname{Filt}^iK^\bullet(X)_{\mathbb Q}\bigr)\subset \operatorname{Filt}^{i-d}K^\bullet(Y)_{\mathbb Q},
\end{equation}
hence a homomorphism (3.3) of degree $-d$.

\sgaStableTarget{sga6.e00.sec.3.3}\subsubsection*{3.3.}
We have thus assembled all the structural elements needed to give a meaning to formula (1.2) \uline{for a proper locally complete intersection morphism of noetherian schemes $X,Y$ both admitting ample invertible Modules}. With a little additional care, this formulation moreover remains meaningful independently of any noetherian hypothesis.

In this general form, this Riemann--Roch formula will be proved in Exp. VIII. The reader will see that the proof given in this Seminar combines elements of the two proofs in [RRR] which were discussed in the Introduction. The need to use the methods of the first of these proofs, involving more $\lambda$-formalism than the second (followed in Borel--Serre), is due to the need for a proof of inclusion (3.4), a question that did not arise in the more restricted setting of [RRR] and Borel--Serre [2], where we could work with the Chow ring or the ring $\Gr^\bullet_{\mathrm{top}}$. From the second proof, we retain the introduction of the blow-up of $Y$ along $X$ when $f$ is an immersion. This technical device will doubtless disappear, and the proof for an immersion will extend to the case where we no longer postulate the existence of an ample invertible Module on $Y$, once the question raised in Exp. XIV no.~1 concerning the introduction of exterior-power operations in the derived category $D^-(Y)$ has been resolved.

\sgaStableTarget{sga6.e00.sec.4}\subsection*{4.}
Finally, we must give some indications of how a meaning can still be given to (1.2) without a hypothesis of the existence

%%%%%%%%%%  scan idx 23  |  folio 10  %%%%%%%%%%
\SGAPrintPage{10}
of ample invertible Modules on $X$ and $Y$.

\sgaStableTarget{sga6.e00.sec.4.1}\subsubsection*{4.1.}
First, we note that \uline{for an arbitrary scheme $X$ (or, more generally, locally ringed topos), it seems beyond doubt that the ``right'' invariant to replace the ring $K(X)$ of [RRR]} is not \uline{the ring $K^\bullet(X)$ considered above, defined in terms of locally free Modules on $X$, but the group $K(\Parf(X))$, considered in 3.1., defined in terms of perfect complexes on $X$, a group that we shall denote here by} $K^\bullet(X)$, and which we shall distinguish from the ``naive'' group of 3.1, denoted by $K^\bullet_{\mathrm{naive}}(X)$. This point of view is forced by the need to define a direct-image homomorphism (3.1) for a proper morphism $f:X\to Y$ of finite tor-dimension between, say, noetherian schemes (the non-noetherian case requires somewhat stronger local finiteness properties of $f$, cf. Exp. III). With the definition of $K^\bullet(X)$ now adopted, this homomorphism $f_*$ is indeed defined trivially by formula (3.2), where $\underline F$ is now an arbitrary perfect complex on $X$, which implies that $Rf_*(\underline F)$ is also perfect.

\sgaStableTarget{sga6.e00.sec.4.2}\subsubsection*{4.2}
This new definition of $K^\bullet(X)$ nevertheless raises a new problem of homological algebra: that of defining a $\lambda$-structure on $K^\bullet(X)$ compatible with its natural ring structure arising from the tensor product in the derived category. (NB When $L_\bullet$ and $L'_\bullet$ are two perfect complexes, their total tensor product $L_\bullet\otimes^{\mathbf L}L'_\bullet$ is also perfect). More precisely, we should define ``exterior-power'' operations in the derived category $D^-(X)$, taking perfect complexes to perfect complexes. This question has not yet been clarified at present, and, as was stated in the Introduction, it constitutes the most serious gap in the Seminar. Once this question has been resolved, and proceeding as in 2.3, we obtain a natural filtration on $K^\bullet(X)$, allowing us to form the associated graded ring $\Gr^\bullet(X)$, which will play the role of the Chow ring, and to define a theory of Chern classes
\[
c_i:K^\bullet(X)\longrightarrow \Gr^\bullet(X).
\]

\sgaStableTarget{sga6.e00.sec.4.3}\subsubsection*{4.3}
To give a meaning to formula (1.2) for a proper locally complete intersection morphism $f:X\to Y$ of noetherian schemes, we must still include in the statement of this formula the nontrivial inclusions (3.4) (where $d$ is the virtual relative dimension) to allow (3.3) to be defined, and finally define the virtual relative tangent class

%%%%%%%%%%  scan idx 24  |  folio 11  %%%%%%%%%%
\SGAPrintPage{11}
\[
T_f\in K^\bullet(X).
\]
When $f$ admits a factorization (2.1) through an immersion into a smooth relative scheme, the definition of $T_f$ presents no new difficulty. We can moreover reformulate it in this case, in a manner more in keeping with the ``derived categories'' point of view, by considering the familiar canonical homomorphism
\sgaStableTarget{sga6.e00.eq.4-1}\begin{equation}\tag{4.1}
\underline N_{X/X'}\xrightarrow{d}\Omega^1_{X'/Y}\otimes_{\underline O_{X'}}\underline O_X
\end{equation}
induced by the exterior differential
\[
\var{d_{X'/X}:\underline J\subset\underline O_{X'}\longrightarrow\Omega^1_{X'/Y},}{d_{X'/Y}:\underline J\subset\underline O_{X'}\longrightarrow\Omega^1_{X'/Y},}{EN-VAR-006}
\]
and the chain complex
\sgaStableTarget{sga6.e00.opening.eq.4-2}\begin{equation}\tag{4.2}
L^{X/Y}_\bullet\in\operatorname{Ob}D(X)
\end{equation}
defined from (4.1) ``by extending by zeros'', a complex equal in degree $1$, respectively $0$, to the source, respectively the target, of the homomorphism (4.1). Formula (2.2) then takes the more pleasant form
\sgaStableTarget{sga6.e00.opening.eq.4-3}\begin{equation}\tag{4.3}
\check T_f=\cl^\bullet\bigl(L^{X/Y}_\bullet\bigr),
\end{equation}
a formula which makes sense in $K^\bullet(X)$ when $f$ is a locally complete intersection morphism, since in this case the complex $L^{X/Y}_\bullet$ is manifestly perfect. The assertion of invariance of $T_f$ relative to the chosen factorization of $f$ can then readily be made more precise as an assertion of invariance of the complex $L^{X/Y}_\bullet$: up to canonical isomorphism in the derived category $D(X)$, this complex is likewise independent of the chosen factorization (Exp. VIII).

\sgaStableTarget{sga6.e00.sec.4.4}\subsubsection*{4.4.}
When $f:X\to Y$ is an arbitrary morphism of schemes, we can still construct a canonical object $L^{X/Y}_\bullet$ of the derived category $D(X)$, defined up to unique isomorphism, which on each affine open of $X$ over an affine open of $Y$ is isomorphic to the complex constructed according to the principle explained in the preceding paragraph (where, however, in the general case ``smooth'' should be replaced by ``formally smooth'', so that every algebra $B$ over a ring $A$ occurs as a quotient of a formally smooth algebra, for example a polynomial algebra). For such a definition, we cannot proceed by simple ``gluing'' from the affine constructions, because the objects of the derived category $D(X)$ are by their nature

%%%%%%%%%%  scan idx 25  |  folio 12  %%%%%%%%%%
\SGAPrintPage{12}
essentially not glueable. Here is a general construction of $L^{X/Y}_\bullet$, valid more generally for every morphism $f:X\to Y$ of (commutatively) ringed topoi: let $B=\underline O_X$, $A=f^{-1}(\underline O_Y)$, so that $A$ is a Ring on the space $X$, and $B$ is an $A$-Algebra. Let $T\to B$ be a homomorphism of sheaves of \underline{sets} generating $B$ as an $A$-Algebra, i.e. such that the corresponding homomorphism of $A$-Algebras
\[
\phi:C=A[T]\longrightarrow B
\]
is an epimorphism of sheaves of sets, where $A[T]$ denotes the free commutative $A$-Algebra generated by $T$ (i.e. the sheaf associated with the presheaf $U\longmapsto A(U)[T(U)]$ = free commutative $A(U)$-algebra generated by $T(U)$).
Let $J$ be the kernel of $\phi$, and consider, as in 4.3, the chain complex of length $1$ defined by the natural homomorphism
\[
J/J^2\xrightarrow{d}\Omega^1_{C/A}\otimes_C B.
\]
It is not difficult to see that, up to canonical isomorphism in the category $D(X)$, the complex so constructed is independent of the choice of $T$ and of $T\to B$, so that it deserves the name \uline{relative cotangent complex $L^{X/Y}_\bullet$ of $X$ over $Y$}. We also verify that, on every affine open of $X$ over an affine open of $Y$, it is given (up to canonical isomorphism) by the procedure of 4.3.

When $f$ is locally complete intersection, we moreover see that we even obtain a \underline{perfect} complex, which then allows us again to define an element $T_f$ in $K^\bullet(X)$ by formula (4.3). We point out that here again, in general, it would not have been possible to define $T_f$ as an element of $K^\bullet_{\mathrm{naive}}(X)$, and nothing permits us to suppose that $T_f$ is \underline{globally} isomorphic in $D(X)$ to a bounded complex with locally free components of finite type. This therefore provides further justification for the point of view advanced in 4.1 in favour of the ``sophisticated'' $K^\bullet(X)$, in preference to $K^\bullet_{\mathrm{naive}}(X)$.

\sgaStableTarget{sga6.e00.sec.4.5}\subsubsection*{4.5.}
Subject to the question raised in 4.2, we have therefore again assembled all the structural elements needed to give a meaning to (1.2) for a proper locally complete intersection morphism of arbitrary schemes (the noetherian hypothesis implicitly made in the preceding sections can in fact be removed without essential difficulty). Unfortunately, this does not mean that, even subject to this reservation, the formula in question has been established, even when $Y$ is the spectrum of a field

%%%%%%%%%%  scan idx 26  |  folio 13  %%%%%%%%%%
\SGAPrintPage{13}
that is algebraically closed of characteristic $p>0$, and $X$ is, say, a proper smooth algebraic scheme of dimension $3$. See the comments in Exp. XIV no.~2.

\sgaStableTarget{sga6.e00.sec.5}\subsection*{5.}
In the form presented in no.~4, the Riemann--Roch formula (1.2) also has a meaning for a proper locally complete intersection morphism of complex analytic spaces (or of rigid analytic spaces in the sense of Tate). But of course the proof given in [RRR] applies to this case only when we also assume that the morphism $f$ is \underline{projective}, and, as in the algebraic case, it seems that an essentially new idea is needed to treat the general case. The formula we envisage here has values in an ``analytic Chow ring'' $\Gr^\bullet(X)$, or rather $\Gr^\bullet(X)_{\mathbb Q}$, constructed according to the principle explained in 4.2 (the difficulty encountered there moreover disappears because we are in characteristic zero, as noted in Exp. XIV no.~1).

When $X$ and $Y$ are nonsingular, another version of the Riemann--Roch formula can be given as follows. First let $X$ be an arbitrary complex analytic space, and $X^{\mathrm{top}}$ the same space equipped with the sheaf of continuous complex-valued functions. We thus have a canonical morphism of ringed spaces
\[
X^{\mathrm{top}}\longrightarrow X,
\]
which, by the evident contravariance of the functor $K^\bullet(X)$ in the variable ringed space $X$, defines a ring homomorphism
\sgaStableTarget{sga6.e00.eq.5-1}\begin{equation}\tag{5.1}
K^\bullet(X)\longrightarrow K^\bullet(X^{\mathrm{top}}).
\end{equation}
For simplicity, suppose that $X^{\mathrm{top}}$ is compact; then, as in the case (already mentioned in 3.1) of a scheme admitting an ample invertible module, we see (Exp. II) that every perfect complex on $X^{\mathrm{top}}$ is globally isomorphic to a bounded complex locally free in each degree, i.e. to a complex of complex vector bundles on $X^{\mathrm{top}}$, and consequently $K^\bullet(X^{\mathrm{top}})$ is canonically isomorphic to the topologists' familiar group $K(X^{\mathrm{top}})$ [1], defined in terms of complex vector bundles on the compact space $X^{\mathrm{top}}$. Thus we can interpret the homomorphism (5.1) as a homomorphism from $K^\bullet(X)$, defined in terms of the analytic structure, to the topologists' familiar ring $K(X^{\mathrm{top}})$. Now, for a morphism $f:X\longrightarrow Y$ of compact nonsingular analytic spaces, we know how to associate, by Atiyah--Hirzebruch

%%%%%%%%%%  scan idx 27  |  folio 14  %%%%%%%%%%
\SGAPrintPage{14}
[1], a group homomorphism
\sgaStableTarget{sga6.e00.eq.5-2}\begin{equation}\tag{5.2}
f^{\mathrm{top}}_*:K^\bullet(X^{\mathrm{top}})\longrightarrow K^\bullet(Y^{\mathrm{top}}).
\end{equation}
Using likewise the homomorphism $f_*$ of (3.1) defined by formula (3.2) (which here implicitly uses \var{Grauert's finiteness theorem [4]}{Grauert's finiteness theorem [3]}{EN-VAR-007} for a proper morphism of analytic spaces), we obtain a square of natural homomorphisms involving only groups of type $K^\bullet$ (to the exclusion of rings of Chow-ring or cohomology type):
\sgaStableTarget{sga6.e00.eq.5-3}\begin{equation}\tag{5.3}
\begin{array}{ccc}
K^\bullet(X) & \longrightarrow & K^\bullet(X^{\mathrm{top}}) \\
f_*\downarrow && \downarrow f^{\mathrm{top}}_* \\
K^\bullet(Y) & \longrightarrow & K^\bullet(Y^{\mathrm{top}}).
\end{array}
\end{equation}
The Riemann--Roch type formula to which we alluded consists in asserting that this square commutes. We note the difference in nature between this assertion (which does not disregard torsion phenomena, and lies in a group $K^\bullet(Y^{\mathrm{top}})$ of an essentially discrete nature) and the Riemann--Roch formula of the type considered in no.~4, which lies in a group of ``Chow'' type that is no longer discrete in nature but must instead be tensored with $\mathbb Q$. We note that formula (5.3) is proved, in any case, when $X$ and $Y$ are nonsingular projective varieties [7], as is the formula of the type considered in no.~4, which is proved in this case in [RRR]. This is a plausible variant of the Atiyah--Singer index theorem (and of its generalizations due to \var{Shih [7]}{Shih [6]}{EN-VAR-008} and Atiyah), which plainly calls for a common generalization. Once the commutativity of (5.3) has been proved, the ``differentiable Riemann--Roch theorem'' of Atiyah--Hirzebruch [1] would immediately give a cohomological variant of the complex analytic Riemann--Roch formula, namely the commutativity of the square
\sgaStableTarget{sga6.e00.eq.5-4}\begin{equation}\tag{5.4}
\begin{array}{ccc}
K^\bullet(X) & \xrightarrow{\;x\longrightarrow \ch(x)\Todd(T_X)\;} & H^{2*}(X,\mathbb Q) \\
\downarrow && \downarrow \\
K^\bullet(Y) & \xrightarrow{\;y\longrightarrow \ch(y)\Todd(T_Y)\;} & H^{2*}(Y,\mathbb Q),
\end{array}
\end{equation}
where the horizontal arrows are obtained from the cohomological Chern classes, obtained from the homomorphism (5.1) and the ordinary Chern class homomorphism
\[
K^\bullet(X^{\mathrm{top}})\longrightarrow H^{2*}(X,\mathbb Z),
\]

%%%%%%%%%%  scan idx 28  |  folio 15  %%%%%%%%%%
\SGAPrintPage{15}
the first vertical arrow is the same as in (5.3), and the second is the Gysin homomorphism in rational cohomology.

\sgaStableTarget{sga6.e00.sec.6}\subsection*{6.}
Throughout this discussion, we have considered almost exclusively the contravariant invariant $K^\bullet(X)$, to the exclusion of the covariant invariant $K_\bullet(X)$ whose existence was mentioned in 3.1. Regarding $K^\bullet(X)$ as an analogue of the integral \underline{cohomology} ring of a topological space, we must regard $K_\bullet(X)$ as an analogue of integral \underline{homology}, the Module structure of $K_\bullet(X)$ over $K^\bullet(X)$ then corresponding to the cap-product operation\SGAAuthorFootnote{raisedstar}{This analogy is moreover not purely verbal: after passing to the associated graded objects $\Gr^\bullet(X)$, $\Gr_\bullet(X)$, it can be made precise by an explicit compatibility relation between the module structure on $\Gr_\bullet(X)$ and the ``cap'' operation.}. The fact that, when $X$ is regular, the natural homomorphism $K^\bullet(X)\to K_\bullet(X)$ is an isomorphism must be regarded as the analogue of the Poincaré duality theorem on oriented topological manifolds, asserting that cap product with the fundamental homology class defines an isomorphism from cohomology to homology. These analogies also suggest the desirability of defining, say for a finite polyhedron $X$, an invariant $K_\bullet(X)$ of a covariant nature, whose relations to integral homology (spectral sequence, Chern homomorphism\dots) would be analogous to those between $K^\bullet(X)$ and integral cohomology; nor does it seem excluded that there might exist, in terms of these groups $K_\bullet(X)$, variants of the differentiable Riemann--Roch theorem valid for spaces more general than compact differentiable manifolds.

Returning to an arbitrary noetherian scheme $X$, we should equip $K_\bullet(X)$ with an increasing filtration, $\operatorname{Filt}_iK_\bullet(X)$ being generated by classes $\cl_\bullet(F)$ with $F$ a coherent sheaf on $X$ such that $\dim \operatorname{supp}.F\le i$. We shall prove (Exp. X) the compatibility of this filtration with the module structure and the filtration of $K^\bullet(X)$, i.e. the formula
\sgaStableTarget{sga6.e00.eq.6-1}\begin{equation}\tag{6.1}
\var{\operatorname{Filt}^iK^\bullet(X)\cdot \operatorname{Filt}_j(K_\bullet(X)\subset \operatorname{Filt}_{j-i}K_\bullet(X),}{\operatorname{Filt}^iK^\bullet(X)\cdot \operatorname{Filt}_j K_\bullet(X)\subset \operatorname{Filt}_{j-i}K_\bullet(X),}{EN-VAR-009}
\end{equation}
which equips the associated graded object of the filtered module $K_\bullet(X)$, denoted by
\[
\Gr_\bullet(X),
\]
with a module structure over $\Gr^\bullet(X)$. Its elements can moreover be interpreted as \uline{classes of algebraic cycles} on $X$, (for an equivalence relation less fine than rational equivalence, when $X$ is of finite type over a base field $K$). If

%%%%%%%%%%  scan idx 29  |  folio 16  %%%%%%%%%%
\SGAPrintPage{16}
\[
f:X\longrightarrow Y
\]
is a proper morphism of noetherian schemes, then
\sgaStableTarget{sga6.e00.eq.6-2}\begin{equation}\tag{6.2}
f_*:K_\bullet(X)\longrightarrow K_\bullet(Y)
\end{equation}
is compatible with the filtrations and defines a homomorphism of graded groups
\sgaStableTarget{sga6.e00.eq.6-3}\begin{equation}\tag{6.3}
f_*:\Gr_\bullet(X)\longrightarrow \Gr_\bullet(Y),
\end{equation}
giving a ``projection formula'', i.e. it is $\Gr^\bullet(Y)$-linear, a formula which \var{is reduced by passage to the associated graded objects from}{is deduced by passage to the associated graded objects from}{EN-VAR-010} the analogous projection formula for the homomorphism (6.2) (which is $K^\bullet(Y)$-linear).

\uline{From the point of view proposed in the Seminar}, \uline{a workable theory of intersections}, \uline{on noetherian schemes} $X$ \uline{that are not necessarily regular}, \uline{must consist in the combined study of the contravariant invariants} $K^\bullet(X),\Gr^\bullet(X)$ \uline{and the covariant invariants} $K_\bullet(X),\Gr_\bullet(X)$, \uline{which play roles of an essentially different nature}, \uline{and whose properties complement one another}. Thus the covariant groups $K_\bullet(X),\Gr_\bullet(X)$ can be calculated more readily for certain nonproper fibrations, by means of the ``homotopy exact sequence'' that can be established for them ([RRR] and Exp. IX); for example, we prove (Exp. IX) canonical isomorphisms such as
\[
K_\bullet(X[t])\simeq K_\bullet(X),\qquad \Gr_\bullet(X[t])\simeq \Gr_\bullet(X),
\]
which fail in general for $K^\bullet$ when $X$ is a nonregular scheme. On the other hand, $K^\bullet(X)$ and $\Gr^\bullet(X)$ lend themselves well to calculation because of their ring structures.

For schemes $X$ proper over a field, the projection
\[
f^X:X\longrightarrow(\operatorname{point})
\]
defines an ``\uline{Euler--Poincaré characteristic}'' homomorphism
\sgaStableTarget{sga6.e00.eq.6-4}\begin{equation}\tag{6.4}
\lambda_X=f^X_*:K_\bullet(X)\longrightarrow K_\bullet(\operatorname{point})=\mathbb Z,
\end{equation}
inducing on the subgroup $\Gr_0(X)=\operatorname{Filt}_0K_\bullet(X)$ the ``degree of $0$-cycles'' homomorphism
\sgaStableTarget{sga6.e00.eq.6-5}\begin{equation}\tag{6.5}
\deg_X:\Gr_0(X)\longrightarrow \mathbb Z.
\end{equation}
The multiplication homomorphism
\[
\Gr^i(X)\times\Gr_i(X)\longrightarrow \Gr_0(X),
\]
composed with the degree homomorphism (6.5), then gives a pairing
\sgaStableTarget{sga6.e00.eq.6-6}\begin{equation}\tag{6.6}
\Gr^i(X)\times \Gr_i(X)\longrightarrow \mathbb Z,
\end{equation}
which explains why, in certain respects, $\Gr^\bullet(X)$ and $\Gr_\bullet(X)$ play dual roles,

%%%%%%%%%%  scan idx 30  |  folio 17  %%%%%%%%%%
\SGAPrintPage{17}
just as cohomology and homology do. Thus, for a morphism $f:X\to Y$ of proper algebraic schemes, the two homomorphisms
\[
f^*:\Gr^\bullet(Y)\longrightarrow \Gr^\bullet(X)\qquad\text{and}\qquad f_*:\Gr_\bullet(X)\longrightarrow \Gr_\bullet(Y)
\]
are formally transposes of one another, with respect to the pairings (6.6).

The homomorphisms (6.4) and (6.5), and the pairing (6.6) obtained from them, can be regarded as the source of \uline{numerical relations} in intersection theory on proper algebraic schemes over a given field $k$. This point of view will be developed in some detail in Exp. X, which is largely independent of Exp. VI, VII, VIII (centred on the proof of the Riemann--Roch theorem) and of Exp. IX. Some notions concerning the \uline{behaviour of} $K^\bullet(X),\Gr^\bullet(X)$ \uline{under specialization} will also be given there.

\sgaStableTarget{sga6.e00.sec.7}\subsection*{7.}
As in every theory of intersections, it is important to specify the place of the \uline{Picard group} in the theory (or the group of \uline{divisor classes} in the classical terminology). Subject to a slight restriction, automatically satisfied for example when $X$ is quasi-compact, Exp. X establishes a canonical isomorphism
\sgaStableTarget{sga6.e00.eq.7-1}\begin{equation}\tag{7.1}
c_1:\Pic(X)\xrightarrow{\sim}\Gr^1(X),
\end{equation}
which shows the exact role of the Picard group in the theory we propose. This role justifies a more detailed study of this Picard group in Exposés XII, XIII, which are logically independent of the text of the Seminar. A study of \uline{numerical equivalence} and certain \uline{finiteness theorems for the Picard group} (\var{announced without proof in [6]}{announced without proof in [5]}{EN-VAR-011}) is made by Kleiman in Exp. XIII, using techniques of a purely numerical nature (which do not use the theory of the Picard functor and its representability), due to himself and to Mumford, markedly simpler than Grothendieck's initial proofs. The only result refractory to Kleiman's methods (and which he is obliged to use) \var{is [6, C-07 (i)]}{is [5, C-07 (i)]}{EN-VAR-012}, whose proof and certain refinements (including certain representability theorems for the Picard functor) occupy Exposé XII. The latter moreover has a spirit and a technique entirely foreign to the rest of the Seminar, being connected with [5]; this ex%
%%%%%%%%%%  scan idx 31  |  folio 18  %%%%%%%%%%
\SGAPrintPage{18}%
posé is logically independent of all the others.

Finally, in Exp. XI, of a distinctly more general nature than XII and XIII and closer to the general spirit of the Seminar, we study, on an arbitrary locally ringed topos, a remarkable functor called the ``\uline{determinant functor}'':
\sgaStableTarget{sga6.e00.eq.7-2}\begin{equation}\tag{7.2}
\det^*:\underline{\operatorname{Parfis}}(X)\longrightarrow \underline{\Inv}(X)
\end{equation}
from the category of perfect complexes on $X$, whose morphisms consist only of the isomorphisms in the sense of the derived category $D(X)$, to the category of invertible Modules on $X$. This functor is described almost completely by the requirement that it be of a ``local nature'' and that, for a bounded complex $L_\bullet$ with locally free components, it be given by the formula
\sgaStableTarget{sga6.e00.eq.7-3}\begin{equation}\tag{7.3}
\det^*(L_\bullet)\simeq\bigotimes_i \det(L_i)^{(-1)^i},
\end{equation}
where $\det(L_i)$ denotes the exterior power of maximal degree of the locally free Module $L_i$. Although this exposé too is logically independent of the rest of the Seminar except for Exposé I, it nevertheless forms an important part of the general ``yoga'' proposed in the Seminar and developed mainly in Exposés I to XI.

%%%%%%%%%%  scan idx 32  |  folio 19  %%%%%%%%%%
\SGAPrintPage{19}
\sgaStableTarget{sga6.e00.bibliography}\section*{\underline{BIBLIOGRAPHY}}
\begin{enumerate}[label={[\arabic*]},leftmargin=2.5em]
\item M. Atiyah and F. Hirzebruch, Vector bundles and homogeneous spaces, Symposia in Pure Mathematics, vol. 3, \var{Différential Geometry}{Differential Geometry}{EN-VAR-013} p. 7--38, 1961.
\item A. Borel and J.P. Serre, Le théorème de Riemann--Roch, \var{Bull. Soc. math France}{Bull. Soc. math. France}{EN-VAR-014}, t. 86, p. 97--136 (1958).
\item H. Grauert, Ein Theorem der analytischen Garbentheorie, Pub. Math. no.~5, p. 5--64 (1960).
\item A. Grothendieck, Classes de faisceaux et théorème de Riemann--Roch, mimeographed notes, Princeton 1957, (cited as [RRR] and reproduced as an Appendix to Exp. 0 of the Seminar).
\item A. Grothendieck, Technique de descente et théorèmes d'existence en Géométrie Algébrique, in Fondements de la Géométrie Algébrique, Extracts from the Bourbaki Seminar 1957--1962, Secrétariat Mathématique, Paris.
\item W. Shih, Séminaire de Topologie à l'IHES, 1966/67.
\item M. Atiyah and F. Hirzebruch, The Riemann--Roch theorem for analytic embeddings, Topology, vol. 1, p. 151--166 (1962).
\end{enumerate}

\clearpage

%%%%%%%%%%  scan idx 33  |  folio 20  %%%%%%%%%%
\SGAPrintPage{20}
\sgaStableTarget{sga6.e00.appendix-to-expose-0-rrr}
\begin{center}
{\Large \underline{\underline{CLASSES OF SHEAVES AND THE RIEMANN--ROCH THEOREM}}}\par
\vspace{0.5em}
\underline{by A. GROTHENDIECK}\SGAAuthorFootnote{*}{This is the verbatim reproduction of the report cited in the Introduction to the present Seminar. The footnotes indicated by the signs (*), (**) were added in October 1967.}
\end{center}

\sgaStableTarget{sga6.e00.chapter-i-rings-formal-preliminaries}\subsection*{CHAP. I -- $\lambda$-\uline{Rings} (\uline{formal preliminaries})}
\begin{longtable}{@{}p{0.74\linewidth}r@{}}
1. Definitions & 1\\
2. Examples & 4\\
3. The $\lambda$-ring defined by a graded ring & 8\\
4. The operations $\lambda^p(N,x)$ & 13\\
\end{longtable}

\sgaStableTarget{sga6.e00.chapter-ii-classes-of-coherent-algebraic-sheaves}\subsection*{CHAP. II -- \uline{Classes of coherent algebraic sheaves and Chern classes}\hfill 18}
\begin{longtable}{@{}p{0.74\linewidth}r@{}}
1. Chow's theory & 18\\
2. Definition of the Chern classes of coherent algebraic sheaves & 22\\
3. Functorial generalities on $K(X)$ & 27\\
4. Some technical results & 34\\
5. Sheaf-theoretic definition of Chern classes. Application to the study of injection morphisms & 43\\
6. The Riemann--Roch theorem & 49
\end{longtable}

\begin{center}
Proof of the Riemann--Roch theorem (1 November 1957).
\end{center}

\clearpage

%%%%%%%%%%  scan idx 34  |  folio 21  %%%%%%%%%%
\SGAPrintPage{21}
\sgaStableTarget{sga6.e00.chapter-i-rings-formal-preliminaries.occ2}\section*{CHAPTER I}
\begin{center}
$\lambda$-Rings (Formal preliminaries)
\end{center}

\sgaStableTarget{sga6.e00.sec.1.occ2}\subsection*{\S 1. Definitions}

A $\lambda$\uline{-ring}\SGAAuthorFootnote{*}{In this seminar we instead say \uline{pre-}$\lambda$\uline{-ring} (V 2.1), reserving the name $\lambda$-ring for those satisfying the additional condition on page 3 (cf V 2.4).} is a commutative ring $K$ with unit, equipped with a family of maps
\sgaStableTarget{sga6.e00.eq.1-1.occ2}\begin{equation}\tag{1.1}
\lambda^i:K\longrightarrow K
\end{equation}
($i$ an integer $\geq 0$) satisfying the following conditions:
\sgaStableTarget{sga6.e00.eq.1-2.occ2}\begin{equation}\tag{1.2}
\lambda^0x=1,\qquad \lambda^1x=x,\qquad
\lambda^n(x+y)=\sum_{i=0}^n\lambda^ix\cdot\lambda^{n-i}y.
\end{equation}
For every $x\in K$, put
\sgaStableTarget{sga6.e00.eq.1-3.occ2}\begin{equation}\tag{1.3}
\lambda_t(x)=\sum_{n\geq 0}\lambda^n(x)t^n\in K[[t]] ;
\end{equation}
the preceding relations then express that $x\longmapsto\lambda_t(x)$ is a homomorphism from the additive group $K$ to the multiplicative group $1+K[[t]]^+$ of formal series over $K$ with augmentation $1$, lifting the canonical homomorphism $1+\sum_{i>0}x_it^i\longrightarrow x_1$.

Thus, on the ring $\Z$, there exists a unique $\lambda$-structure compatible with its ring structure and such that

%%%%%%%%%%  scan idx 35  |  folio 22  %%%%%%%%%%
\SGAPrintPage{22}
\sgaStableTarget{sga6.e00.eq.1-4}\begin{equation}\tag{1.4}
\lambda_t(1)=1+t.
\end{equation}
We then have
\sgaStableTarget{sga6.e00.eq.1-5}\begin{equation}\tag{1.5}
\lambda_t(n)=(1+t)^n,
\end{equation}
hence $\lambda^i(n)=\binom{n}{i}$.

The notion of a $\lambda$-homomorphism of $\lambda$-rings is clear; an \uline{augmented} $\lambda$\uline{-ring} is a $\lambda$-ring equipped with a $\lambda$-homomorphism to $\Z$.

To define the notion of a \uline{special} $\lambda$\uline{-ring}, the following considerations are useful. Let $k$ be a commutative ring with unit, and let $K=1+k[[t]]^+$ be the group of formal series with coefficients in $k$ and augmentation $1$; we shall introduce on it a $\lambda$-ring law whose additive structure is multiplication of formal series. The multiplication of this ring will be denoted by $f\circ g$. It is completely determined by the condition that the coefficients $c_n=c_n(f\circ g)$ of $f\circ g$ are given by \uline{universal} polynomials with integer coefficients in the coefficients $a_i=c_i(f)$ and $b_i=c_i(g)$ of $f$ and $g$, respectively:
\sgaStableTarget{sga6.e00.eq.1-6}\begin{equation}\tag{1.6}
c_n=P_n(a_1,\ldots,a_n;b_1,\ldots,b_n),
\end{equation}
that it is distributive with respect to the ``addition'' $fg$, and is given for linear factors by:
\sgaStableTarget{sga6.e00.eq.1-7}\begin{equation}\tag{1.7}
(1+at)\circ(1+bt)=1+abt.
\end{equation}
The $P_n$ are calculated by calculations of symmetric functions; $P_n$ is isobaric of weight $n$ in the $a_i$ and isobaric of weight $n$ in the $b_i$ ($a_i$ and $b_i$ having weight $i$), and is moreover symmetric in $a=(a_i)$ and $b=(b_i)$. Likewise, the $\lambda^i(f)$ are determined unambiguously by the condition that

%%%%%%%%%%  scan idx 36  |  folio 23  %%%%%%%%%%
\SGAPrintPage{23}
the coefficients $c_n(\lambda^i(f))=c_n\text{ of }\lambda^i(f)$ are calculated using \uline{universal} polynomials with integer coefficients in the coefficients $c_n(f)=a_n$ of $f$:
\sgaStableTarget{sga6.e00.eq.1-8}\begin{equation}\tag{1.8}
c_n=P_{i,n}(a_1,\ldots,a_{in}),
\end{equation}
that relations (1.2) hold, and that, for $i>1$, we have:
\sgaStableTarget{sga6.e00.eq.1-9}\begin{equation}\tag{1.9}
\lambda^i(1+at)=1
\end{equation}
($1$ is the zero element of the ring $K=1+k[[t]]^+$). The $P_{i,n}$ are again calculated by calculations of symmetric functions, and we find that $P_{i,n}$ is an isobaric polynomial of weight $in$ in the $a_i$. Equipped with the preceding operations, $K$ becomes a $\lambda$-ring; its zero element is $1$, and its unit element is $1+t$. We also note the relation:
\sgaStableTarget{sga6.e00.eq.1-7bis}\begin{equation}\tag{1.7 bis}
(1+at)\circ f=f(at)
\end{equation}

Now let $K$ be an arbitrary $\lambda$-ring. We say that $K$ is \uline{special} when the additive homomorphism $\lambda_t$ from $K$ to $1+K[[t]]^+$ is a $\lambda$-homomorphism. This therefore means that we have the formulas:
\sgaStableTarget{sga6.e00.eq.1-10}\begin{equation}\tag{1.10}
\lambda_t(1)=1+t,\quad
\lambda_t(xy)=\lambda_t(x)\circ\lambda_t(y),\quad
\lambda_t(\lambda^i(x))=\lambda^i(\lambda_t(x)),
\end{equation}
or, explicitly:
\sgaStableTarget{sga6.e00.eq.1-11}\begin{equation}\tag{1.11}
\begin{cases}
\lambda^i(1)=0 & \text{for }i>1\\
\lambda^n(xy)=P_n(\lambda^1x,\ldots,\lambda^n x;\lambda^1y,\ldots,\lambda^n y)\\
\var{\lambda^n(\lambda^i(x))=P_{i,n}(\lambda^1x,\ldots,\lambda^n x),}{\lambda^n(\lambda^i(x))=P_{i,n}(\lambda^1x,\ldots,\lambda^{in} x),}{EN-VAR-015}
\end{cases}
\end{equation}
where the $P_n$ and $P_{i,n}$ are the universal polynomials considered above. In particular, formulas (1.5) follow. (We can show that the $\lambda$-ring $1+k[[t]]^+$ constructed from a commutative ring $k$ is spe%
%%%%%%%%%%  scan idx 37  |  folio 24  %%%%%%%%%%
\SGAPrintPage{24}%
cial; but we shall not need this fact).

\sgaStableTarget{sga6.e00.sec.2.occ2}\subsection*{\S 2. Examples}

\sgaStableTarget{sga6.e00.ex.2-1}1) Let $G$ be a group and $k$ a commutative ring; we consider the group $\underline{K}(G)$, the quotient of the free group generated by the classes of representations of $G$ in $k$-modules of finite type by the subgroup generated by elements of the form $(\rho\oplus\rho')-(\rho)-(\rho')$ (we write $\oplus$ for direct sum). Tensor product and exterior powers define a $\lambda$-ring structure on $\underline{K}(G)$. We can also pass to the quotient by elements of the form $(\rho)-(\rho')-(\rho'')$, where $\rho$ is a representation that is an extension of $\rho'$ by $\rho''$, trivial as an extension of $k$-modules. We obtain the group $\underline{K}_r(G)$ of reduced classes of representations of $G$, and we find that the $\lambda$-ring law passes to the quotient.

When $k$ is an \uline{algebraically closed field of characteristic} $0$\SGAAuthorFootnote{*}{It is unnecessary for $k$ to be algebraically closed, cf. (VI 3.3) for a still more general result (including all those that will be used in the present report).}, we can show that $\underline{K}(G)$, and a fortiori $\underline{K}_r(G)$, is a \uline{special} $\lambda$-ring; this no longer holds when the characteristic is $\neq 0$, but nevertheless, when $k$ is algebraically closed\SGAAuthorFootMark{*}, $\underline{K}_r(G)$ is still a special $\lambda$-ring. To prove this, we reduce to the case where $G$ is a product $\Gl(m,k)\times\Gl(n,k)$ and only rational representations of $G$ are considered; our assertion follows from the \uline{explicit determination} of $\underline{K}_r(G)$ \uline{to be made below} (example 3); finally, when $k$ has characteristic $0$, complete reducibility of the rational representations of $G$ shows that $\underline{K}(G)=\underline{K}_r(G)$.

%%%%%%%%%%  scan idx 38  |  folio 25  %%%%%%%%%%
\SGAPrintPage{25}
Variants of the preceding example: $G$ is an algebraic group defined over $k$, and only the rational representations of $G$ defined over $k$ are considered; we can consider a topological group, or a Lie group, and continuous representations, etc\dots\SGAAuthorFootnote{1}{Other examples of $\lambda$-rings can be constructed, for example the set of equivalence classes of quadratic forms, the set of classes of representations of a Lie algebra, etc\dots We promise that we shall not need all these examples to treat the Riemann--Roch theorem!}

\sgaStableTarget{sga6.e00.ex.2-2}2) Let $X$ be an algebraic variety, irreducible or not, defined over the field $k$. We can consider the quotient of the free group generated by the classes of vector bundles on $X$ (defined over $k$) by the subgroup generated by elements of the form $(E\oplus E')-(E)-(E')$, on which the $\lambda$-ring operations will be defined by tensor product and exterior powers. When $X$ is complete, this is a free group generated by the indecomposable vector bundles on $X$ (by Krull--Remak--Schmidt); knowing this ring and its basis is equivalent to the complete classification of vector bundles on $X$ and knowing the elementary tensor operations on those bundles. When $k$ is algebraically closed of characteristic $0$, all the other tensor operations on classes of bundles are then known, using what follows in 3). We can make a more drastic passage to the quotient by the group generated by the elements $(E)-(E')-(E'')$, where $E$ is a vector bundle that is an extension of $E'$ by $E''$; we then obtain a $\lambda$-ring, denoted by $\underline{K}(\xi(X))$, where $\xi(X)$ denotes the additive category of vector bundles on $X$ defined over $k$. Using the results of example 1), we show that, when $k$ is algebraically closed\SGAAuthorFootnote{*}{Unnecessary condition, cf. the preceding footnote.}, the $\lambda$-ring $\underline{K}(\xi(X))$ is special, and the same holds for the

%%%%%%%%%%  scan idx 39  |  folio 26  %%%%%%%%%%
\SGAPrintPage{26}%
first $\lambda$-ring defined here, if the characteristic of $k$ is $0$, but not in the general case.

\sgaStableTarget{sga6.e00.ex.2-3}3) \uline{Determination of} $\underline{K}(G)$ \uline{and of} $\underline{K}_r(G)$ \uline{in certain important cases}: We suppose that the field $k$ is algebraically closed, that the group $G$ is algebraic, \uline{affine connected} and ``globally \uline{reductive}'' (the radical of $G$ is isomorphic to $k^{*s}$), and we consider only rational linear representations. In characteristic $0$, the complete reducibility of the rational linear representations of $G$ shows that $\underline{K}(G)=\underline{K}_r(G)$ is the free group generated by the classes of irreducible rational representations. Now, if $T$ is a maximal torus of $G$, we have $\underline{K}(T)=\underline{K}_r(T)=\Z(\That)$, the algebra of the group $\That$ of rational characters of $T$; we have $\That\simeq\Z^r$ if $r$ is the rank of $G$ (this is independent of the characteristic). The Weyl group $W$ operates on $T$, hence on $\underline{K}(T)$, in a way that depends only on the operations of $W$ on the free group $\That$. We have an evident homomorphism of $\lambda$-rings, deduced from the injection of $T$ into $G$:
\sgaStableTarget{sga6.e00.eq.1-13}\begin{equation}\tag{1.13}
\underline{K}_r(G)\longrightarrow \underline{K}_r(T)^W=\underline{K}(T)^W,
\end{equation}
where $\underline{K}(T)^W$ denotes the set of fixed points of $W$ in $\underline{K}(T)$. The theory of linear representations of $G$\SGAAuthorFootnote{*}{Cf.\ Séminaire Chevalley 1956/58 \var{Groupes de Lie albébriques}{Groupes de Lie algébriques}{EN-VAR-017} (École Normale Supérieure).} then proves the following theorem:

\sgaStableTarget{sga6.e00.thm.1-1}\textbf{Theorem 1.1.} \uline{The homomorphism} (1.13) \uline{is an isomorphism.}

We thus see that $\underline{K}_r(G)$, once the ``type'' of $G$ is known (characterized by $W$ operating on a group $\Z^r$), does not depend on the characteristic, and that it is a special $\lambda$-ring. We have a canonical basis of $\underline{K}_r(G)$ corresponding to the orbits of $W$ in $\That$; another basis is obtained as follows:
%%%%%%%%%%  scan idx 40  |  folio 27  %%%%%%%%%%
\SGAPrintPage{27}%
let $(\rho_i)_{1\leq i\leq r'}$ be the representation classes of $G$ corresponding to the \uline{fundamental} representations of $G/\rad G$, and let $(\sigma_i)_{1\leq i\leq r''}$ be a basis of the group of rational homomorphisms from $G$ to $k^*$; we thus have $r'+r''=r$ (the rank of $G$).

\sgaStableTarget{sga6.e00.cor.unnamed}\textbf{Corollary:} \uline{The ring} $\underline{K}_r(G)$ \uline{is identified with the subring}
\[
\Z[\rho_1,\ldots,\rho_{r'};\sigma_1,\ldots,\sigma_{r''};\sigma_1^{-1},\ldots,\sigma_{r''}^{-1}]
\]
\uline{of the field of rational functions in the} $\rho_i$ \uline{and the} $\sigma_j$ \uline{with coefficients in} $\Q$\SGAAuthorFootnote{*}{The derived group $G'$ of $G$ must be assumed simply connected for this statement to be correct.}.

In particular, $\underline{K}_r(G)$ has no divisors of $0$; if $G=\prod_i\Gl(n_i,k)$, the ring $\underline{K}_r(G)$ is made explicit as follows: let $\rho_i$ be the representation of $G$ deduced from the identity representation of $\Gl(n_i,k)$; if we then consider the field of rational functions over $\Q$ in indeterminates $\lambda^j(\rho_i)$ ($1\leq j\leq n_i$), then $\underline{K}_r(G)$ is identified with the subring generated by these indeterminates and the inverses of the $\lambda^{n_i}(\rho_i)$. This justifies our assertion that, at least in characteristic $0$, the operations $\lambda^i$ allow all tensor operations on vector bundles to be reconstructed; this is still true over an algebraically closed field of arbitrary characteristic, provided reduced classes are taken.

\sgaStableTarget{sga6.e00.rem.unnamed}\textbf{Remarks:} a) In example 1), the ring $\underline{K}(G)$ is augmented, the augmentation being defined by the degree of the representations, and, by passage to the quotient, this augmentation defines an augmentation on $\underline{K}_r(G)$. Likewise, the $\lambda$-rings considered in example 2) are augmented when $X$ is $k$-irreducible; in $\underline{K}(\xi(X))$, the augmentation is defined by the \uline{rank of a vector bundle}, that is, the dimension of its fibres.

%%%%%%%%%%  scan idx 41  |  folio 28  %%%%%%%%%%
\SGAPrintPage{28}%
b)
When the field $k$ is not algebraically closed, it is plausible that $\underline{K}(G)$ is a special $\lambda$-ring in characteristic $0$, and likewise for $\underline{K}_r(G)$ in all characteristics.

c) In nonzero characteristic, theorem 1.1 was kindly supplied to me by Chevalley.

\sgaStableTarget{sga6.e00.sec.3.occ2}\subsection*{\S 3. The $\lambda$-ring defined by a graded ring}

Let $A$ be a commutative graded ring with unit; for simplicity we shall assume that the degrees of $A$ are positive and that $A^0=\Z$. We denote by $\Ahat$ the product of the $A^i$ (it is a commutative ring with unit), and by $\Ahp$ the kernel of the evident augmentation $\Ahat\to A^0=\Z$. Then $1+\Ahp$ is a subgroup of the multiplicative group of invertible elements of $\Ahat$. Consider the abelian group
\sgaStableTarget{sga6.e00.eq.1-14}\begin{equation}\tag{1.14}
\widetilde A=\Z\times(1+\Ahp),
\end{equation}
whose elements will be denoted $[n,x]$ with $n\in\Z$ and $\var{x=1+\sum_{i\geq 1}x^i_i\in1+\Ahp}{x=1+\sum_{i\geq 1}x^i\in1+\Ahp}{EN-VAR-016}$ ($x^i\in A^i$). We shall write this group additively, thus:
\sgaStableTarget{sga6.e00.eq.1-15}\begin{equation}\tag{1.15}
[n,x]+[n',x']=[n+n',xx'].
\end{equation}
The zero element of this group is $[0,1]$. We are going to introduce on it an augmented $\lambda$-ring structure, compatible with its additive structure, the augmentation being $\varepsilon:[n,x]\longrightarrow n$. The product $[m,x][n,y]$ is completely characterized by being written $[mn,x^n y^m(x*y)]$, where $x*y$ is an element of $1+\Ahp$ whose components are expressed as universal polynomials with integer coefficients in the $x^i$ and the $y^i$:
\sgaStableTarget{sga6.e00.eq.1-16}\begin{equation}\tag{1.16}
(x*y)^i=Q_i(x^1,\ldots,x^i;y^1,\ldots,y^i)\qquad (i\geq1)
\end{equation}
%%%%%%%%%%  scan idx 42  |  folio 29  %%%%%%%%%%
\SGAPrintPage{29}%
the
polynomial $Q_i$ being isobaric of weight $i$. Moreover, $x*y$ must be distributive with respect to the ``addition'' $xy$ and, for ``linear'' factors, must reduce to the following:
\sgaStableTarget{sga6.e00.eq.1-17}\begin{equation}\tag{1.17}
[1,1+x^1][1,1+y^1]=[1,1+x^1+y^1].
\end{equation}
The polynomials $Q_i$ are evaluated by calculations with elementary symmetric functions (this is the calculation of the Chern classes of a tensor product of vector bundles by means of the Chern classes $x^i$ and $y^i$ of the factors and their respective dimensions $m$ and $n$). From (1.17) we deduce:
\sgaStableTarget{sga6.e00.eq.1-17bis}\begin{equation}\tag{1.17 bis}
(1+x^1)*(1+y^1)=(1+x^1+y^1)/(1+x^1)(1+y^1),
\end{equation}
and an arbitrary product $x*y$ is calculated term by term by decomposing ``formally'' $x$ into a product $\prod_{i=1}^n(1+\alpha_i)$ and $y$ into a product $\prod_{i=1}^n(1+\beta_i)$, where the $\alpha_i$ and the $\beta_i$ have degree $1$.

In this way, $\widetilde A$ becomes an augmented commutative ring, with a unit $[1,1]$, and it can be interpreted as the ring obtained by adjoining a unit to the ring $1+\Ahp$. Moreover, on $\widetilde A$, we define a \uline{filtration} by putting $\widetilde A^0=\widetilde A$ and denoting by $\widetilde A^n$ the set of the $[0,x]$ with $x^i=0$ for $1\leq i<n$; this filtration is compatible with the product, for we have
\sgaStableTarget{sga6.e00.eq.1-18}\begin{equation}\tag{1.18}
\left(1+\sum_{i\geq m}x^i\right)*\left(1+\sum_{j\geq n}y^j\right)
=1-\frac{(m+n-1)!}{(m-1)!(n-1)!}x^m y^n+\cdots,
\end{equation}
the unwritten terms having degree $>m+n$. This formula shows that the graded ring $G(\widetilde A)$ associated with $\widetilde A$ is identified with $A$ from the additive point of view, but not from the multiplicative point of view; we can only say that the map $\eta'$ from $A$ to $G(\widetilde A)$ defined by
%%%%%%%%%%  scan idx 43  |  folio 30  %%%%%%%%%%
\SGAPrintPage{30}%
\sgaStableTarget{sga6.e00.eq.1-19}\begin{equation}\tag{1.19}
\eta'(x^i)=(-1)^{i-1}(i-1)!\,x^i\qquad (x^i\in A^i)
\end{equation}
is a \uline{ring homomorphism}.

We introduce on $\widetilde A$ operations $\lambda^i$ satisfying the axioms of $\lambda$-rings (\S 1, formulas (1.2)), well determined by the condition that the components of $\lambda^i[0,x]$ be deduced from the components $x^i$ of $x$ by \uline{universal} polynomials with integer coefficients
\sgaStableTarget{sga6.e00.eq.1-20}\begin{equation}\tag{1.20}
\lambda^i[0,x]=[0,\lambda^i x]\qquad (\lambda^i x)^{(n)}=Q_{i,n}(x^1,\ldots,x^n)
\end{equation}
($Q_{i,n}$ is isobaric of weight $n$), and that we have
\sgaStableTarget{sga6.e00.eq.1-21}\begin{equation}\tag{1.21}
\lambda^i[1,1+x^1]=0\qquad\text{for }i>1.
\end{equation}
The polynomials $Q_{i,n}$ are evaluated by a calculation with elementary symmetric functions, which is the calculation of the Chern classes of the exterior powers of a vector bundle, by means of the Chern classes and the rank of the latter. We then find that $\widetilde A$ is a \uline{special} augmented $\lambda$-ring and that the $\widetilde A^i$ are stable under the operations $\lambda^j$.

The augmented $\lambda$-ring $\widetilde A$ may be considered as a functor in $A$; consequently, the principal automorphism
\[
x=\sum_{i\geq0}x^i\longrightarrow \check x=\sum_{i\geq0}(-1)^i x^i
\]
of $A$ defines an automorphism of $\widetilde A$, denoted by the same symbol $\check{\ }$.

\sgaStableTarget{sga6.e00.prop.1-2}\textbf{Proposition 1.2.} \uline{For} $x\in\widetilde A$ \uline{of the form} $[n,1+x^1+\cdots+x^n]$ (we therefore have $x^i=0$ for $i>n$), \uline{we have} $\lambda^i(x)=0$ \uline{for} $i>n$\uline{, we have} $\lambda^n(x)=[1,1+x^1]$\uline{, and finally}
\sgaStableTarget{sga6.e00.eq.1-22}\begin{equation}\tag{1.22}
\sum_{i=0}^n(-1)^i\lambda^i(x)=[0,1-(n-1)!x^n+\cdots]
\end{equation}
From this proposition and from the fact that the $\widetilde A^i$ are stable under the operations $\lambda^j$, we deduce that formula (1.22) is valid whenever $\varepsilon(x)=n$, whether or not we have $x^i=0$ for $i>n$.

%%%%%%%%%%  scan idx 44  |  folio 31  %%%%%%%%%%
\SGAPrintPage{31}%
Formula
(1.22) admits the following generalizations: let $K$ be any $\lambda$-ring, otherwise arbitrary; for $x\in K$, put
\sgaStableTarget{sga6.e00.eq.1-23}\begin{equation}\tag{1.23}
\gamma^n(x)=(-1)^n\sum_{i=0}^n(-1)^i\lambda^i(x+n)=\lambda^n(x+n-1)
\end{equation}
(the identity of the last two members follows from an immediate induction). In particular, we have $\gamma^0(x)=1$ and $\gamma^1(x)=x$; next put:
\sgaStableTarget{sga6.e00.eq.1-24}\begin{equation}\tag{1.24}
\gamma_t(x)=\sum_{n\geq0}\gamma^n(x)t^n\in K[[t]]
\end{equation}
whence we deduce the formula:
\sgaStableTarget{sga6.e00.eq.1-25}\begin{equation}\tag{1.25}
\gamma_t(x)=\lambda_{t/(1-t)}(x),
\end{equation}
which is equivalent to
\sgaStableTarget{sga6.e00.eq.1-25bis}\begin{equation}\tag{1.25 bis}
\lambda_s(x)=\gamma_{s/(1+s)}(x).
\end{equation}
This proves that the $\lambda^i$ are expressed by means of the $\gamma^i$ and that $\gamma_t(x+y)=\gamma_t(x)\gamma_t(y)$, and consequently that the maps $\gamma^i$ define a new $\lambda$-ring structure on $K$. Formula (1.22) then shows that, for $K=\widetilde A$ and every $x\in\widetilde A$, we have
\sgaStableTarget{sga6.e00.eq.1-22bis}\begin{equation}\tag{1.22 bis}
\gamma^n(x-\varepsilon(x))=[0,1+(-1)^{n-1}(n-1)!x^n+\cdots]
\end{equation}
and consequently, the \uline{component} $x^n$ \uline{is obtained, up to the factor} $(-1)^{n-1}(n-1)!$, \uline{by reducing modulo} $\widetilde A^{n+1}$ \uline{the element} $\gamma^n(x-\varepsilon(x))$ \uline{of} $\widetilde A^n$.

\sgaStableTarget{sga6.e00.the-chern-homomorphism}\paragraph{The Chern homomorphism:}
For every graded ring $A$ satisfying the preceding conditions, we are going to define a homomorphism of augmented rings:
\sgaStableTarget{sga6.e00.eq.1-26}\begin{equation}\tag{1.26}
\ch:\var{\widetilde A\relbar\joinrel\relbar\joinrel\relbar}{\widetilde A\longrightarrow}{EN-VAR-018} \widehat{A\otimes\Q}.
\end{equation}
This homomorphism is defined unambiguously by the conditions that it be an additive homomorphism, take $1$ to $1$, be such that the com%
%%%%%%%%%%  scan idx 45  |  folio 32  %%%%%%%%%%
\SGAPrintPage{32}%
ponents of $\ch(x)$ in degree $>0$ be given by universal isobaric polynomials with rational coefficients in the components of $x$, and that
\sgaStableTarget{sga6.e00.eq.1-27}\begin{equation}\tag{1.27}
\ch([1,1+x^1])=\exp x^1=\sum_{n\geq0}(x^1)^n/n!.
\end{equation}
We then immediately obtain the explicit expression for $\ch$: let $\eta$ be the additive endomorphism of $\widehat{A\otimes\Q}$ defined, for $x^i$ of degree $i$, by
\sgaStableTarget{sga6.e00.eq.1-28}\begin{equation}\tag{1.28}
\eta(x^i)=1/(-1)^{i-1}(i-1)!\,x^i
\end{equation}
we then have
\sgaStableTarget{sga6.e00.eq.1-29}\begin{equation}\tag{1.29}
\ch\left([n,1+\sum_{i\geq1}x^i]\right)=n+\eta\left(\log\left(1+\sum_{i\geq1}x^i\right)\right).
\end{equation}
As we can ``invert'' this formula, at least when $A$ has only finitely many homogeneous components, we see that, in this case, we have an \uline{isomorphism} $\ch$ \uline{from the ring} $\widetilde A\otimes\Q$ \uline{onto the ring} $\widehat{A\otimes\Q}$ (which completely elucidates the structure of $\widetilde A\otimes\Q$).

Recall that, according to Hirzebruch [1], a formal series $f\in\Q[[t]]$ defines, by universal polynomial formulas, an additive homomorphism $1+\Ahp\longrightarrow 1+(\widehat{A\otimes\Q})^+$, denoted $\mathcal T_f$; when $f(t)=t/(1-\exp(-t))$, we simply put $\mathcal T_f=\mathcal T$ ($\mathcal T$ is the initial of Todd). We extend the definition of $\mathcal T(x)$ to $x\in\widetilde A$. This said, we have the following result:

\sgaStableTarget{sga6.e00.prop.1-3}\textbf{Proposition 1.3.} \uline{Let} $N\in\widetilde A$ \uline{be of the form} $[q,1+N^1+\cdots+N^q]$ (we have $N^i=0$ for $i>q$). \uline{If we put:}
\[
\lambda_{-1}(N)=\sum_{i\geq0}(-1)^i\lambda^i(N)=\var{\sum_{i=0}^{n}(-1)^i\lambda^i(N)}{\sum_{i=0}^{q}(-1)^i\lambda^i(N)}{EN-VAR-019},
\]
\uline{we shall have:}
\sgaStableTarget{sga6.e00.eq.1-30}\begin{equation}\tag{1.30}
\ch(\lambda_{-1}(N))=(-1)^q N^q\mathcal T(\check N)^{-1},
\end{equation}
%%%%%%%%%%  scan idx 46  |  folio 33  %%%%%%%%%%
\SGAPrintPage{33}%
\uline{which is equivalent to:}
\sgaStableTarget{sga6.e00.eq.1-30bis}\begin{equation}\tag{1.30 bis}
\ch(\lambda_{-1}(\check N))=N^q\mathcal T(N)^{-1}.
\end{equation}

\sgaStableTarget{sga6.e00.sec.4.occ2}\subsection*{\S 4. The operations $\lambda^p(N,x)$}

Let $A$ be the ring of formal series, with coefficients in $\Z$, in two infinite sequences of indeterminates, denoted $\lambda^iN$ and $\lambda^i x$ ($i\geq1$). Put $\lambda^1N=N$ and $\lambda^1x=x$, and introduce on $A$ the unique special $\lambda$-ring structure for which we have $\lambda^i(N)=\lambda^iN$ and $\lambda^i(x)=\lambda^ix$ and for which the operations $\lambda^i$ are ``continuous'' (if we restricted ourselves to \uline{polynomials} in indeterminates $\lambda^iN$ and $\lambda^ix$, we would have the ``free $\lambda$-ring generated by $N$ and $x$''; our ring $A$ is a completion of it). We can then form
\sgaStableTarget{sga6.e00.eq.1-31}\begin{equation}\tag{1.31}
\lambda_{-1}(N)=\sum_{i\geq0}(-1)^i\lambda^iN,
\end{equation}
which is a formal series with constant term $1$, hence invertible; we can therefore define an element $\lambda^p(N,x)$ of $A$ unambiguously by putting
\sgaStableTarget{sga6.e00.eq.1-32}\begin{equation}\tag{1.32}
\lambda^p(N,x)\lambda_{-1}(N)=\lambda^p(x\lambda_{-1}(N))
\end{equation}

\sgaStableTarget{sga6.e00.thm.1-4}\textbf{Theorem 1.4.} \uline{Let} $J_q$ \uline{be the closed ideal of} $A$ \uline{generated by the} $\lambda^iN$ \uline{with} $i>q$\uline{; identify} $A/J_q$ \uline{with the ring of formal series in the} $\lambda^i x$ \uline{and the} $\lambda^jN$ \uline{(}$1\leq j\leq q$\uline{). Then} $\lambda^p(N,x)$ \uline{reduced modulo} $J_q$ \uline{is a} polynomial, \uline{isobaric of weight} $p$ \uline{with respect to the variables} $\lambda^i x$\uline{.}

\uline{Let} $N$ \uline{and} $F$ \uline{be two vector spaces of finite dimensions} $q,q'$ \uline{over an algebraically closed field} $k$ \uline{of characteristic} $0$\uline{; let} $N^{(p)}$ \uline{be the kernel of the map} $(a_1,\ldots,a_p)\longrightarrow a_1+\cdots+a_p$ from $N^p$ \uline{to} $N$.
We
%%%%%%%%%%  scan idx 47  |  folio 34  %%%%%%%%%%
\SGAPrintPage{34}%
\uline{let the symmetric group} $\mathfrak S_p$ \uline{operate on} $N^p$\uline{, hence on} $N^{(p)}$ \uline{hence on} $\Lambda(N^{(p)})$\uline{, then on} $\bigotimes^p F$ \uline{and on} $\Lambda(N^{(p)})\otimes\bigotimes^p F$\uline{. Consider this vector space as a graded module over} $\mathfrak S_p\times\Gl(N)\times\Gl(F)$ \uline{and take its ``alternating part'', which is a graded module over} $\Gl(N)\times\Gl(F)\eqqcolon G$\uline{. We then have}
\sgaStableTarget{sga6.e00.eq.1-33}\begin{equation}\tag{1.33}
\lambda^p(N,F)=\underline{E}_G\left(\left(\Lambda(N^{(p)})\otimes\bigotimes^p F\right)^{\mathrm{alt}}\right).
\end{equation}

N.B. If $M$ is a complex of $G$-modules, we denote by $\underline{E}_G(M)$ the alternating sum in $\underline{K}(G)$ of the classes of the components of $M$; this is the Euler--Poincaré characteristic of $M$; in formula (1.33), the first member $\lambda^p(N,F)$ denotes the element of the ring $\underline{K}(G)$ obtained by substituting for the variables $\lambda^iN$ and $\lambda^iF$ the classes in $\underline{K}(G)$ of the $G$-modules $\Lambda^i(N)$ and $\Lambda^i(F)$ respectively.

\textbf{Proof:} By expanding the second member of (1.32) by the second formula (1.11), we see that the first assertion is proved if we prove the analogous assertion obtained by putting $x=1$, i.e. if we prove that, in the $\lambda$-ring $\Z[\lambda^1N,\ldots,\lambda^qN]$ where $\lambda^iN=0$ for $i>q$, the element $\lambda^p(\lambda_{-1}(N))$ is divisible by $\lambda_{-1}(N)$. A priori, the quotient $\lambda^p(N,1)=\lambda^p(\lambda_{-1}(N))/\lambda_{-1}(N)$ is a formal series contained in the field of fractions $K$ of the ring of formal series in the $\lambda^iN$ for $1\leq i\leq q$. On the other hand, also denote by $N$ a vector space of dimension $q$ over an algebraically closed field $k$ of characteristic $0$, and consider the element $\ell^p(N)=\var{\underline{E}_G((\Lambda(N^{(p)})^{\mathrm{alt}})}{\underline{E}_G((\Lambda(N^{(p)}))^{\mathrm{alt}})}{EN-VAR-020}\in \underline{K}(G)$, where we now have $G=\Gl(N)$. We know \var{(\S 2, theorem 1.3)}{(\S 2, theorem 1.1)}{EN-VAR-021} that $\underline{K}(G)$ is identified
%%%%%%%%%%  scan idx 48  |  folio 35  %%%%%%%%%%
\SGAPrintPage{35}%
with the
subring of $K$ generated by the $\lambda^iN$ for $1\leq i\leq q$ and by the inverse of $\lambda^qN$. On the other hand, we readily verify in this ring that we have $\ell^p(N)\lambda_{-1}(N)=\lambda^p(\lambda_{-1}(N))$. As we are in a large field $K$ and $\lambda_{-1}(N)\neq0$, we conclude that $\ell^p(N)=\lambda^p(N,1)$; both members are in the intersection of $\underline{K}(G)$ and the ring of formal series in the $\lambda^iN$, and are therefore polynomials in the $\lambda^iN$.

To prove (1.33) in the general case, we verify that the second member satisfies in $\underline{K}(G)$ the formula analogous to (1.32), that is, its product by $\lambda_{-1}(N)$ is $\lambda^p(F\lambda_{-1}(N))$. Likewise, the product of $\lambda^p(N,F)\in \underline{K}(G)$ by $\lambda_{-1}(N)$ is $\lambda^p(F\lambda_{-1}(N))$, as follows from the fact that $\underline{K}(G)$ is \uline{special}. As $\underline{K}(G)$ is an integral domain (cf.\ \S 2) and $\lambda_{-1}(N)\neq0$ (same reference), we have indeed proved formula (1.33). q.e.d.

\sgaStableTarget{sga6.e00.rem.unnamed.occ2}\textbf{Remarks:} a) When the dimension $q'$ of $F$ is $\geq p$, then (1.33) characterizes the polynomial $\lambda^p(N,x)$ in the $\lambda^iN$ ($1\leq i\leq q$) and the $\lambda^jx$ ($1\leq j\leq p$).

b) It is plausible that (1.33) still holds when the field $k$ is not algebraically closed\SGAAuthorFootnote{*}{This is indeed proved by the same method.}. This question would have to be resolved to prove the Riemann--Roch theorem over a field of characteristic $0$ that is not algebraically closed by the method we are going to explain. Moreover, the passage from an $\mathfrak S_p$-module to its alternating part has a reasonable meaning only when the base field has characteristic $0$.

Now let $K$ be any \uline{special} $\lambda$-ring and let $N$ be an
ele%
%%%%%%%%%%  scan idx 49  |  folio 36  %%%%%%%%%%
\SGAPrintPage{36}%
ment of $K$ such that $\lambda^iN=0$ for sufficiently large $i$. We can then, for every $x\in K$ and every integer $p\geq1$, consider the elements $\lambda^p(N,x)$, which are expressed by universal polynomials with integer coefficients in the $\lambda^iN$ and the $\lambda^jx$. We have $\lambda^1(N,x)=x$ and we have formula (1.32) \uline{in} $K$. To express the properties of the operations $x\mapsto\lambda^p(N,x)$, it is convenient to introduce a new $\lambda$-ring denoted $K_N$. As a commutative ring, $K_N$ is obtained by adjoining a unit to the ring whose underlying additive group is $K$ and whose multiplication is given by
\sgaStableTarget{sga6.e00.eq.1-34}\begin{equation}\tag{1.34}
x_Ny=xy\mu\qquad (\mu=\lambda_{-1}(N)).
\end{equation}
The map $\lambda_t$ from $K_N$ to $K_N[[t]]$ is defined by:
\sgaStableTarget{sga6.e00.eq.1-35}\begin{equation}\tag{1.35}
\lambda_t(n,x)=(1+t)^n\sum_{p\geq0}\lambda^p(N,x)t^p,
\end{equation}
where we have put $\lambda^0(N,x)=(1,0)$ (the unit element of $K_N$). I say that, \uline{equipped with these operations,} $K_N$ \uline{is a special} $\lambda$\uline{-ring}. We therefore have to verify the formulas $\lambda_t(X+Y)=\lambda_t(X)\lambda_t(Y)$, $\lambda_t(XY)=\lambda_t(X).\lambda_t(Y)$, and $\lambda_t(\lambda^i(X))=\lambda^i(\lambda_t(X))$ for $X,Y\in K_N$. We can restrict ourselves to verifying these formulas when $X$ and $Y$ are in $K$, and we find that it suffices to verify them when $K$ is the free special $\lambda$-ring generated by $X,Y$ and $N$ subject to the relations $\lambda^iN=0$ for $i>q$. But the group homomorphism
\sgaStableTarget{sga6.e00.eq.1-36}\begin{equation}\tag{1.36}
K_N\longrightarrow K,
\end{equation}
which takes unit to unit and on $K$ reduces to the homomorphism $x\longrightarrow x\mu$, is a homomorphism compatible with the ring structures and the maps $\lambda^i$, as we verify immediately. In the present case,
%%%%%%%%%%  scan idx 50  |  folio 37  %%%%%%%%%%
\SGAPrintPage{37}%
its
restriction to $K$ is \uline{injective}; as $K$ is a special $\lambda$-ring, it follows readily that the above relations do hold for $X,Y\in K$, and this completes the proof of our assertion.

According to formula (1.23), it is natural, for $x\in K$ and $n\geq1$, to put
\sgaStableTarget{sga6.e00.eq.1-37}\begin{equation}\tag{1.37}
\gamma^n(N,x)=\sum_{i=0}^{n-1}\var{\binom{n-1}{i}\lambda^i(N,x)}{\binom{n-1}{i}\lambda^{n-i}(N,x)}{EN-VAR-022},
\end{equation}
in other words, we put $\gamma^n(N,x)=\gamma^n(x_N)$, where $x_N$ denotes $x$ considered as an element of the $\lambda$-ring $K_N$.

\sgaStableTarget{sga6.e00.prop.1-5}\textbf{Proposition 1.5.} \uline{Let} $A$ \uline{be a commutative graded ring with unit,} $N$ \uline{an element of} $\widetilde A$ \uline{of the form} $N=[q,1+N^1+\cdots+N^q]$\uline{, and} $x$ \uline{an arbitrary element of} $\widetilde A$\uline{:} $x=[\nu,1+\sum_{i\geq1}x^i]$\uline{. Then, for every integer} $\var{j\geq0}{j\geq1}{EN-VAR-023}$\uline{, we have} $\gamma^{q+j}(N,x)\in\widetilde A^j$\uline{, and the component of degree} $j$ \uline{of} $\gamma^{q+j}(N,x)$ \uline{is of the form}
\sgaStableTarget{sga6.e00.eq.1-38}\begin{equation}\tag{1.38}
\left(\gamma^{q+j}(N,x)\right)^{(j)}=(-1)^{j-1}(j-1)!G_{q,j}(\nu,\var{x^1,..,x^j;N^1,..,N^j}{x^1,\ldots,x^j;N^1,\ldots,N^j}{EN-VAR-024}),
\end{equation}
\uline{where} $G_{q,j}$ \uline{is a universal polynomial with integer coefficients characterized by the condition}
\sgaStableTarget{sga6.e00.eq.1-39}\begin{equation}\tag{1.39}
G_{q,j}(\nu,\var{x^1,..,N^j}{x^1,\ldots,N^j}{EN-VAR-025})N^q=(-1)^q(x*\lambda_{-1}(N))^{(q+j)}.
\end{equation}

\textbf{Proof.} We can assume that $A$ is the ring of polynomials with integer coefficients in indeterminates $x^i$ and $N^j$ (with $1\leq j\leq q$). As we evidently have, combining (1.32) and (1.37),
\sgaStableTarget{sga6.e00.eq.1-40}\begin{equation}\tag{1.40}
\gamma^n(N,x)\lambda_{-1}(N)=\gamma^n(x\lambda_{-1}(N)),
\end{equation}
and, for $n=q+j$, the second member has filtration $\geq q+j$ (use formula (1.22 bis), taking into account that $x\lambda_{-1}(N)$ has
%%%%%%%%%%  scan idx 51  |  folio 38  %%%%%%%%%%
\SGAPrintPage{38}%
zero augmentation), whereas $\lambda_{-1}(N)$ has filtration $q$ (formula (1.22)), it follows that $\gamma^n(N,x)$ has filtration $j$, since, by (1.18), the associated graded ring of $A$ is here an \uline{integral domain}. Using the preceding formulas, we immediately obtain the indicated form for $(\gamma^{q+j}(N,x))^{(j)}$.
\begin{flushright}q.e.d.\end{flushright}

\sgaStableTarget{sga6.e00.chapter-ii-classes-of-coherent-algebraic-sheaves.occ2}\section*{CHAPTER II}
\begin{center}
Classes of coherent algebraic sheaves\par
and Chern classes
\end{center}

\sgaStableTarget{sga6.e00.sec.1.occ3}\subsection*{\S 1. Chow's theory}

Throughout this chapter, we assume a \uline{base field} $k$ fixed, which we shall assume algebraically closed, for simplicity. A large part of what follows, and perhaps all of it, is nevertheless valid without this restriction. The algebraic spaces and regular maps considered will be defined over $k$.

An algebraic space is called \uline{quasi-projective} if it is isomorphic to a locally closed subset of a projective space.

Let $X$ be a connected nonsingular quasi-projective variety of dimension $n$. We denote by $A(X)$, and call the \uline{Chow ring} of $X$, the graded ring of classes of cycles on $X$ modulo rational equivalence. ($A^i(X)$ consists of the classes of cycles of dimension $n-i$). The fact that $A(X)$
%%%%%%%%%%  scan idx 52  |  folio 39  %%%%%%%%%%
\SGAPrintPage{39}%
is indeed a graded ring follows from the following theorem of Chow:

\sgaStableTarget{sga6.e00.thm.2-1}\textbf{Theorem 2.1.}\SGAAuthorFootnote{*}{As a result of a misunderstanding on the author's part, the theorem states more than is currently known. In fact, only Chow's ``moving lemma'' is used, saying that two cycles can be moved within their classes so that their intersection has no excess; cf.\ Séminaire Chevalley 1958, Anneau de Chow et applications.} \uline{Every cycle on} $X$ \uline{is rationally equivalent to a nonsingular cycle. Given elements of} $A^i(X)$ \uline{and} $A^j(X)$\uline{, we can find nonsingular representatives of these elements, say} $Y^{n-i}$ \uline{and} $Z^{n-j}$\uline{, which meet transversely everywhere.}

(A cycle is called nonsingular if its components are nonsingular subvarieties. We say that two cycles meet transversely everywhere if every component of one meets every component of the other transversely everywhere, i.e. if the two intersecting varieties are nonsingular at their points of intersection and their tangent varieties there are in general position \dots)

Let $f$ be a morphism (i.e. a regular map) from $X$ to $Y$ ($X$ and $Y$ being quasi-projective and nonsingular); then $f$ defines (using the preceding theorem) a homomorphism of graded rings:
\sgaStableTarget{sga6.e00.eq.2-1.occ2}\begin{equation}\tag{2.1}
f^*:A(Y)\longrightarrow A(X)
\end{equation}
(\uline{inverse image of cycle classes}). Thus $A(X)$ becomes a contravariant functor in $X$.

Let $f$ be a \uline{proper} morphism from $X^n$ to $Y^m$ (i.e. in Weil's terminology, the graph of $f$ is complete over $Y$, cf.\ [3]). Then $f$ also defines an additive homomorphism
\sgaStableTarget{sga6.e00.eq.2-2}\begin{equation}\tag{2.2}
f_*:A(X^n)\longrightarrow A(Y^m)
\end{equation}
preserving the dimension of cycles, and therefore increasing the degree by $m-n$.

% --- Source pages 47--75 ---
%%%%%%%%%%  scan idx 53  |  folio 40  %%%%%%%%%%
\SGAPrintPage{40}%
Thus, with respect to proper morphisms, and disregarding the multiplicative structures and gradings, $A(X)$ is a covariant functor in $X$. Let us note the classical formula
\sgaStableTarget{sga6.e00.eq.2-3.occ2}\begin{equation}\tag{2.3}
 f_*(x.f^*(y))=f_*(x).y .
\end{equation}

We shall also use Chow's theory of Chern classes. It consists in attaching to every vector bundle $E$ on $X$ (still assumed quasi-projective and nonsingular) \uline{Chern classes} $C^i(E)\in A^i(X)$, ($i\geq 1$), in such a way as to satisfy the following conditions; put $C^0(E)=1$ and
\sgaStableTarget{sga6.e00.eq.2-4.occ2}\begin{equation}\tag{2.4}
 \Ctil(E)=\left[\rang E,\sum_{i\geq 0}C^i(E)\right]\in \Atil(X)
\end{equation}
(see chapter I, \S 3 for the definition of $\Atil(X)$). The conditions to be satisfied by a theory of Chern classes are then expressed as follows:
\sgaStableTarget{sga6.e00.eq.2-5.occ2}\begin{equation}\tag{2.5}
 \Ctil(E)=\Ctil(E')+\Ctil(E'')
\end{equation}
if the vector bundle $E$ is an extension of $E'$ by $E''$ (putting $C(E)=\sum_{i\geq0}C^i(E)$, this condition is expressed by $C(E)=C(E').C(E'')$). Formula (2.5) allows $\Ctil$ to be extended to an additive homomorphism from $\underline{K}(\xi(X))$ (cf.\ Chap.\ I, \S 2, example 2) to $\Atil(X)$:
\sgaStableTarget{sga6.e00.eq.2-6}\begin{equation}\tag{2.6}
 \Ctil:\underline{K}(\xi(X))\longrightarrow \Atil(X),
\end{equation}
and we want $\Ctil$ to be a $\lambda$\uline{-homomorphism of} $\lambda$\uline{-rings}; this last condition is also written
\sgaStableTarget{sga6.e00.eq.2-7}\begin{equation}\tag{2.7}
 \Ctil(E\otimes F)=\Ctil(E)\Ctil(F)\qquad
 \Ctil(\Lambda^iE)=\lambda^i\Ctil(E),
\end{equation}
for two vector bundles $E$ and $F$ (of course, $\Ctil$ is automatically compatible with the augmentations). Let $D$ be a divisor on $X$ and $L(D)$ the vector bundle it defines; we want
%%%%%%%%%%  scan idx 54  |  folio 41  %%%%%%%%%%
\SGAPrintPage{41}%
\sgaStableTarget{sga6.e00.eq.2-8}\begin{equation}\tag{2.8}
 C^1(L(D))=\ellc(D)\qquad C^i(L(D))=0\quad (i>1)
\end{equation}
(where,
for an arbitrary cycle $Z$, we denote by $\ellc(Z)$ its class in $A(X)$). Finally, we require the $C^i$ to be \uline{functorial}, i.e. that, when $f$ is a morphism from $X$ to $Y$ and $E$ is a vector bundle on $Y$, we have
\sgaStableTarget{sga6.e00.eq.2-9}\begin{equation}\tag{2.9}
 C^i(f^!(E))=f^*(C^i(E)),
\end{equation}
where $f^!(E)$ is the inverse image of $E$ on $X$. We can also write (2.9) in the following form: the notion of inverse image of a bundle defines a $\lambda$-homomorphism of augmented $\lambda$-rings
\sgaStableTarget{sga6.e00.eq.2-10}\begin{equation}\tag{2.10}
 f^!:\underline{K}(Y)\longrightarrow \underline{K}(X)\SGAAuthorFootnote{*}{The notation $f^!$ used here conflicts with that used in duality theory (cf.\ for example R. Hartshorne, Residues and Duality, Lecture Notes in Math. no.\ 20, Springer, 1966). This is why in this Seminar we replace it by the notation $f^*$.}
\end{equation}
(here we write $\underline{K}(X)$ instead of $\var{\underline{K}(X)}{\underline{K}(\xi(X))}{EN-VAR-026}$), and (2.9) means that the diagram commutes:
\[
\begin{tikzcd}[column sep=large,row sep=large]
K(Y) \arrow[r,"f^!"] \arrow[d,"\Ctil_Y"'] & K(X) \arrow[d,"\Ctil_X"]\\
\var{A(Y)}{\Atil(Y)}{EN-VAR-027} \arrow[r,"f^*"] & \var{A(X)}{\Atil(X)}{EN-VAR-028}.
\end{tikzcd}
\]
(The space on which the symbol $\Ctil$ is considered has been put as a subscript.)

\sgaStableTarget{sga6.e00.thm.2-2}\textbf{Theorem 2.2.} \uline{There exists a theory of Chern classes satisfying conditions (2.5), (2.7), (2.8) and (2.9).}

We shall also admit the following proposition:

\sgaStableTarget{sga6.e00.prop.2-3}\textbf{Proposition 2.3.} \uline{Let} $Y$ \uline{be a closed subset of} $X$ \uline{and} $U$ \uline{its complement; if} $x$ \uline{is an element of} $A(X)$ \uline{inducing} $0$ \uline{on} $U$\uline{, there exists a representative of} $x$ \uline{that is a cycle supported on} $Y$\uline{.}

%%%%%%%%%%  scan idx 55  |  folio 42  %%%%%%%%%%
\SGAPrintPage{42}%
\sgaStableTarget{sga6.e00.cor.unnamed.occ2}\textbf{Corollary.} \uline{Let} $p=\dim X-\dim Y$\uline{. Then, if} $i<p$\uline{, the homomorphism} $A^i(X)\longrightarrow A^i(U)$ \uline{is bijective. Every element of the kernel of} $A^p(X)\longrightarrow A^p(U)$ \uline{is a linear combination of components of dimension} $n-p$ \uline{of} $Y$ \uline{(and is therefore proportional to} $\ellc(Y)$ \uline{if} $Y$ \uline{is irreducible)}

\sgaStableTarget{sga6.e00.sec.2.occ3}\subsection*{\S 2. Definition of the Chern classes of coherent algebraic sheaves}

Let $\mathbb{C}$ be an abelian category, and let $\mathbb{C}_0$ be a subclass of $\mathbb{C}$. We denote by $K(\mathbb{C}_0)$ the quotient group of the free group generated by the classes (up to isomorphism) of elements of $\mathbb{C}_0$, by the subgroup generated by elements of the form $(E)-(E')-(E'')$, where $E$ is an \uline{extension} of $E'$ by $E''$ ($E,E',E''\in\mathbb{C}_0$). We must assume that the classes of elements of $\mathbb{C}_0$ for the isomorphism relation form a set. We denote by $\gamma_{\mathbb{C}_0}(E)$ the class of an $E\in\mathbb{C}_0$ in $K(\mathbb{C}_0)$, and omit the subscript $\mathbb{C}_0$ when there is no risk of confusion. The group $\underline{K}(\mathbb{C}_0)$ is the solution of a \uline{universal problem} concerning functions $E\longrightarrow f(E)$ defined on $\mathbb{C}_0$, with values in an arbitrary abelian group $A$, which are \uline{additive} in the sense that $f(E)=f(E')+f(E'')$ whenever $E$ is an extension of $E'$ by $E''$.

Let $F$ be an additive functor from $\mathbb{C}$ to $\mathbb{C}'$, and $\mathbb{C}'_0$ a subclass of $\mathbb{C}'$ satisfying the same condition as $\mathbb{C}_0$, such that, for every $A\in\mathbb{C}_0$, the $R^qF(A)$ are in $\mathbb{C}'_0$ and are zero for sufficiently large $q$ (we assume that injective resolutions exist in $\mathbb{C}$, so that the derived functors $R^qF$ of $F$ exist); then $F$ defines a homomorphism from $\underline{K}(\mathbb{C}_0)$ to $\underline{K}(\mathbb{C}'_0)$, still denoted $F$, by the formula
\sgaStableTarget{sga6.e00.eq.2-11}\begin{equation}\tag{2.11}
 F(\gamma_{\mathbb{C}_0}(E))=\sum_{i\geq 0}(-1)^i\gamma_{\mathbb{C}'_0}(R^iF(E)).
\end{equation}
More generally, when we have a cohomological functor $(F^i)$ from $\mathbb{C}$ to $\mathbb{C}'$
%%%%%%%%%%  scan idx 56  |  folio 43  %%%%%%%%%%
\SGAPrintPage{43}%
such that,
for every $E\in\mathbb{C}_0$, the $F^i(E)$ are in $\mathbb{C}'_0$ and are zero except for finitely many, we deduce a homomorphism from $\underline{K}(\mathbb{C}_0)$ to $\underline{K}(\mathbb{C}'_0)$. These remarks extend to multifunctors; moreover, the ``classical'' spectral sequences give transitivity properties that we shall not make explicit.

In particular, when $F$ is an exact functor, formula (2.11) simplifies:
\sgaStableTarget{sga6.e00.eq.2-11bis}\begin{equation}\tag{2.11 bis}
 F(\gamma_{\mathbb{C}_0}(E))=\gamma_{\mathbb{C}'_0}(F(E)).
\end{equation}
This applies in particular when $\mathbb{C}=\mathbb{C}'$, when $F$ is the identity functor, and when $\mathbb{C}_0\subset\mathbb{C}'_0$: we have a canonical homomorphism from $\underline{K}(\mathbb{C}_0)$ to $\underline{K}(\mathbb{C}'_0)$.

\sgaStableTarget{sga6.e00.thm.2-3}\textbf{Theorem 2.3.} \uline{Let} $\mathbb{C}$ \uline{be an abelian category and} $\mathbb{C}_0\subset\mathbb{C}_1$ \uline{two subclasses such that the classes} (up to isomorphism) \uline{of objects of} $\mathbb{C}_1$ \uline{form a set. We assume that the following conditions hold}
\begin{enumerate}[label=(\roman*)]
\item \uline{For every exact sequence} $0\to A'\to A\to A''\to0$ \uline{in} $\mathbb{C}$ \uline{such that} $A$ \uline{and} $A''$ \uline{are in} $\mathbb{C}_0$ \uline{(resp.\ }$\mathbb{C}_1$\uline{), we have} $A'\in\mathbb{C}_0$ \uline{(resp.\ }$A'\in\mathbb{C}_1$\uline{).}
\item \uline{Let} $u:A\to B$ \uline{and} $v:C\to B$ \uline{with} $A,B,C\in\mathbb{C}_1$ \uline{and} $u$ \uline{surjective; then there exist} $L\in\mathbb{C}_0$ \uline{and homomorphisms} $u':L\to C$ \uline{and} $v':L\to A$ \uline{such that} $u'$ \uline{is surjective and} $vu'=uv'$\uline{.}
\item \uline{For every} $A\in\mathbb{C}_1$\uline{, there exists an integer} $d=d(A)$ \uline{such that, for every left resolution} $\underline{L}$ \uline{of} $A$ \uline{by objects of} $\mathbb{C}_0$\uline{, we have} $Z_d(\underline{L})\in\mathbb{C}_0$ ($Z_d$ denotes the cycles of dimension $d$).
\end{enumerate}
\uline{Under these conditions, the canonical homomorphism from} $K(\mathbb{C}_0)$ \uline{to} $K(\mathbb{C}_1)$ \uline{is an isomorphism.}

\textbf{Principle of the proof:} let $A\in\mathbb{C}_1$; by (ii) and (iii) there exists a finite resolution $\underline{L}$ of $A$ by objects of $\mathbb{C}_0$; put
%%%%%%%%%%  scan idx 57  |  folio 44  %%%%%%%%%%
\SGAPrintPage{44}%
then $f(A)=\sum_{i\geq 0}(-1)^i\gamma_{\mathbb{C}_0}(L^i)$.
We show that the element of $\underline{K}(\mathbb{C}_0)$ thus constructed does not depend on the chosen resolution $\underline{L}$: given a second resolution $\underline{L}'$, we ``cover'' $\underline{L}$ and $\underline{L}'$ by a third resolution $\underline{L}''$, with \uline{surjective} homomorphisms from $\underline{L}''$ to $\underline{L}$ and $\underline{L}'$ \var{(which is possible by (i) and (ii);}{(which is possible by (i) and (ii));}{EN-VAR-029} we are then reduced to proving $\underline{E}_{\mathbb{C}_0}(\underline{L})=\underline{E}_{\mathbb{C}_0}(\underline{L}'')$. Now the difference of the two is $\underline{E}_{\mathbb{C}_0}(\underline{N})$, where $\underline{N}$ is the kernel of the surjective homomorphism from $\underline{L}''$ onto $\underline{L}$, a kernel that is in $\mathbb{C}_0$ by (i). Now $\underline{N}$ is an acyclic complex in all dimensions, whence $\underline{E}_{\mathbb{C}_0}(\underline{N})=0$, whence the conclusion (N.B. $\underline{E}_{\mathbb{C}_0}(\underline{L})$, for a finite complex with values in $\mathbb{C}_0$, denotes the Euler--Poincaré characteristic $\sum_{i\geq0}(-1)^i\gamma_{\mathbb{C}_0}(L^i)$ of $\underline{L}$).

We prove analogously that the function $f:\mathbb{C}_1\longrightarrow\underline{K}(\mathbb{C}_0)$ thus constructed is \uline{additive}, and therefore defines a homomorphism from $\underline{K}(\mathbb{C}_1)$ to $\underline{K}(\mathbb{C}_0)$, and it is immediate that this homomorphism is inverse to the canonical homomorphism from $\underline{K}(\mathbb{C}_0)$ to $\underline{K}(\mathbb{C}_1)$. q.e.d.

\sgaStableTarget{sga6.e00.rem.unnamed.occ3}\textbf{Remark:} To verify (ii), it suffices to verify that $\mathbb{C}_1$ is contained in a subclass $\mathbb{C}_2$ of $\mathbb{C}$ satisfying the two conditions
\begin{enumerate}[label=(ii \alph*)]
\item Every $A\in\mathbb{C}_2$ is isomorphic to a quotient of an $L\in\mathbb{C}_0$.
\item If $u:A\to B$, with $A,B\in\mathbb{C}_2$, then $\Ker u\in\mathbb{C}_2$.
\end{enumerate}

\sgaStableTarget{sga6.e00.cor.unnamed.occ3}\textbf{Corollary:}
\uline{Let} $X$ \uline{be a quasi-projective algebraic space. Let} $\cF(X)$ \uline{be the category of coherent algebraic sheaves on} $X$\uline{,} $\xi(X)$ \uline{the category of locally free coherent algebraic sheaves} (equivalent to the category of vector bundles on $X$), \uline{and} $\cF_0(X)$ \uline{the category of coherent algebraic sheaves} $F$ \uline{such that, for every} $x\in X$\uline{,} $F_x$ \uline{is an} $\underline{O}_x$\uline{-module}
%%%%%%%%%%  scan idx 58  |  folio 45  %%%%%%%%%%
\SGAPrintPage{45}%
\uline{of finite cohomological dimension (we therefore have} $\xi(X)\subset\cF_0(X)\subset\cF(X)$\uline{). Then the canonical homomorphism}
\sgaStableTarget{sga6.e00.eq.2-12}\begin{equation}\tag{2.12}
 K(\xi(X))\longrightarrow K(\cF_0(X))
\end{equation}
\uline{is an isomorphism.}

We are going to verify the conditions of theorem 2.3. Conditions (i) and (iii) follow from the definition of cohomological dimension and from the fact that a projective module of finite type over the noetherian local ring $\underline{O}_x$ is free. To verify (ii), we shall show that $\mathbb{C}=\cF(X)$ satisfies conditions (ii a) and (ii b) of the remark. The second is trivial; it therefore remains to prove the following result:

\sgaStableTarget{sga6.e00.prop.2-4}\textbf{Proposition 2.4. (Serre).} \uline{Let} $X$ \uline{be a quasi-projective algebraic space. Then every coherent algebraic sheaf on} $X$ \uline{is isomorphic to the quotient of a locally free algebraic sheaf.}

This proposition holds when $X$ is a closed subset of projective space, by [2], since, for sufficiently large $n$, $F(n)$ is generated by its sections, which means that $F$ is a quotient of $\underline{O}(-n)^k$. In the general case, $X$ is an open subset of a closed subset of a projective space, and it suffices to use the following result (due to Cartier and Serre in particular cases):

\sgaStableTarget{sga6.e00.prop.2-5}\textbf{Proposition 2.5.} \uline{Let} $F$ \uline{be a coherent algebraic sheaf on an open subset} $U$ \uline{of an algebraic space} $X$\uline{. Then} $F$ \uline{is isomorphic to the restriction of a coherent algebraic sheaf on} $X$\uline{.}

By applying Zorn's lemma to the open subsets containing $U$ to which $F$ can be extended, we are reduced to the case where $X$ is affine. Let $\overline F$ be the
alge%
%%%%%%%%%%  scan idx 59  |  folio 46  %%%%%%%%%%
\SGAPrintPage{46}%
braic sheaf on $X$ whose sections on an open subset $V$ are the sections of $F$ on $U\cap V$. We readily find that $\overline F$ is \uline{quasi-coherent} (i.e. defined on every affine open subset $V$ -- hence here on all of $X$ -- by a module, of finite type or otherwise, over the coordinate ring of the open subset in question). This is immediate when $U$ itself is affine, for $\overline F$ is then defined on the open subset $V$ by $\Gamma(U\cap V,F)$ considered as a module over $\Gamma(V,\underline{O}_X)$. In the general case, $U$ is a finite union of affine open subsets $U_i$, so $\overline F$ is the kernel of an evident homomorphism of quasi-coherent sheaves:
\[
 \sum_i \overline{F|U_i}\longrightarrow \sum_{i,j}\overline{F|U_i\cap U_j}
\]
and is therefore itself quasi-coherent. As $X$ is assumed affine, $\overline F$ is therefore generated by its sections. As $\overline F|U=F$ is coherent, there exist finitely many sections of $\overline F$ generating $\overline F|U$. These sections generate a coherent subsheaf of $\overline F$, whose restriction to $U$ is isomorphic to $F$.
\begin{flushright}q.e.d.\end{flushright}

Now assume $X$ quasi-projective and \uline{nonsingular}. Then, by the syzygy theorem, we have, with the notation of the corollary to theorem 2.3, $\cF_0(X)=\cF(X)$. In this case we therefore have:
\sgaStableTarget{sga6.e00.eq.2-13}\begin{equation}\tag{2.13}
 \underline{K}(\xi(X))=\underline{K}(\cF(X)).
\end{equation}
We can express this isomorphism as follows:

\sgaStableTarget{sga6.e00.cor.2}\textbf{Corollary 2 to theorem 2.3.} \uline{Let} $X$ \uline{be a nonsingular quasi-projective variety. For every abelian group} $A$\uline{, there is a one-to-one correspondence between additive functions} $f$ \uline{of vector bundles on} $X$\uline{, with values in} $A$\uline{, and additive functions} $g$ \uline{of coherent algebraic sheaves on} $X$\uline{, with values in} $A$\uline{. To every} $f$ \uline{corresponds the unique} $g$ \uline{such that we have:}
%%%%%%%%%%  scan idx 60  |  folio 47  %%%%%%%%%%
\SGAPrintPage{47}%
\sgaStableTarget{sga6.e00.eq.2-14}\begin{equation}\tag{2.14}
 f(E)=g(\underline{O}_X(E))
\end{equation}
\uline{for every vector bundle} $E$ \uline{on} $X$\uline{, denoting by} $\underline{O}_X(E)$ \uline{the sheaf of germs of regular sections of} $E$ \uline{on} $X$\uline{.}

To calculate $g(F)$ for an arbitrary coherent algebraic sheaf, we consider a finite resolution of $F$ by sheaves of the form $\underline{O}_X(E_i)$, and we then have:
\sgaStableTarget{sga6.e00.eq.2-14bis}\begin{equation}\tag{2.14 bis}
 g(F)=\sum_{i\geq0}(-1)^if(E_i).
\end{equation}

In particular, when $X$ is connected, we considered in \S 1 an additive function $\Ctil(E)$ of the vector bundle $E$, with values in the group $\Atil=\Atil(X)$. We deduce from it an additive function of coherent sheaves, say $\Ctil(F)$, and, by passing to components, functions $F\longrightarrow C^i(F)\in A^i(X)$. The $C^i(F)$ are called the \uline{Chern classes of the coherent sheaf} $F$; they are theoretically calculated from the Chern classes of vector bundles by (2.14 bis) (where, however, we must pass to multiplicative notation in the second member). As for the component of $\Ctil(F)$ along $\Z$, it is the \uline{rank} of the sheaf $F$, which corresponds to the rank of a module over an integral domain: when $X$ is affine (a case to which we can always reduce), $F$ is defined by a module of finite type over the coordinate ring of $X$, and we take the rank of this module\dots.

\sgaStableTarget{sga6.e00.sec.3.occ3}\subsection*{\S 3. Functorial generalities on $\underline{K}(X)$}

When $X$ is an algebraic space, we put, for brevity, $\underline{K}(X)=\underline{K}(\cF(X))$, and, when $X$ is nonsingular and quasi-projective, we identify this group with $\underline{K}(\xi(X))$ by corollary 2 to theorem 2.3. As $\underline{K}(\xi(X))$ is
%%%%%%%%%%  scan idx 61  |  folio 48  %%%%%%%%%%
\SGAPrintPage{48}%
not
only an abelian group, but a $\lambda$-ring, it follows by transport of structure that $\underline{K}(X)$ is also a $\lambda$-ring, and is therefore equipped with a multiplicative structure and maps $\lambda^i$. It is important to define these structures directly, without passing to vector bundles; this will be done below.

If $E$ is a vector bundle on $X$, we denote by $\gamma(E)$ its class in $\underline{K}(X)$; if $F$ is a coherent algebraic sheaf on $X$, \var{$\gamma(\Gamma)$ likewise denotes}{$\gamma(F)$ likewise denotes}{EN-VAR-030} its class in $\underline{K}(X)$. If $Y$ is an irreducible subvariety of $X$, we put $\gamma(Y)=\gamma(\underline{O}_Y)$, where $\underline{O}_Y$ denotes the sheaf of local rings of $Y$, considered as a coherent algebraic sheaf on $X$. We define the symbol $\gamma(Y)$ by linearity when $Y$ is an arbitrary cycle on $X$. When there is a risk of confusion, we shall write $\gamma_X$ for $\gamma$.

Let $X$ be an arbitrary algebraic space. We introduce a \uline{filtration} on $\underline{K}(X)$, denoting by $\underline{K}(X)^i$ the subgroup of $\underline{K}(X)$ generated by the $\gamma(F)$ with dim.supp.\ $F\leq n-i$ (where $n$ is the dimension of $X$). From Serre's ``dévissage lemma'', in the form given in [3], follows the result:

\sgaStableTarget{sga6.e00.prop.2-6}\textbf{Proposition 2.6.} \uline{The group} $\underline{K}(X)^i$ \uline{is generated by the} $\gamma(Y)$\uline{, where} $Y$ \uline{runs through the set of subvarieties of} $X$ \uline{of dimension} $\leq n-i$\uline{. If} $X$ \uline{is irreducible, we have} $G^0(\underline{K}(X))=\Z$\uline{.}

The latter isomorphism is given by the homomorphism from $\underline{K}(X)$ to $\Z$ defined by means of the notion of \uline{rank} of a coherent algebraic sheaf.

Now assume $X$ \uline{nonsingular}. Following the developments
%%%%%%%%%%  scan idx 62  |  folio 49  %%%%%%%%%%
\SGAPrintPage{49}%
at the beginning of \S 2, we
define a multiplication on $\underline{K}(X)$ by putting
\sgaStableTarget{sga6.e00.eq.2-15}\begin{equation}\tag{2.15}
 \gamma(F)\gamma(G)=\sum_{i\geq0}(-1)^i\gamma(\underline{\Tor}_i(F,G))
\end{equation}
(the $\underline{\Tor}_i$ being taken, of course, with respect to the local rings of $X$). The exact sequence of the $\underline{\Tor}_i$ shows that formula (2.15) does define a biadditive map from $\underline{K}(X)\times\underline{K}(X)$ to $\underline{K}(X)$; associativity follows from the spectral sequence of multiple $\Tor$. We also verify in this way that the product of several factors can be calculated by the formula:
\sgaStableTarget{sga6.e00.eq.2-16}\begin{equation}\tag{2.16}
 \gamma(F_1)....\gamma(F_p)=\sum_{i\geq0}(-1)^i\gamma(\underline{\Tor}_i(F_1,\ldots,F_p))
\end{equation}
where, in the second member, the $\underline{\Tor}_i$ are ``simultaneous $\underline{\Tor}$'', calculated by simultaneous projective resolution of the arguments $F_i$.

If, in (2.15), we assume $G$ \uline{locally free}, we find:
\sgaStableTarget{sga6.e00.eq.2-15bis}\begin{equation}\tag{2.15 bis}
 \gamma(F)\gamma(G)=\gamma(F\otimes G),
\end{equation}
which shows in particular that $\gamma(\underline{O}_X)$ is the unit element of $\underline{K}(X)$, and that the canonical homomorphism from $\var{\underline{K}(\xi(X)}{\underline{K}(\xi(X))}{EN-VAR-031}$ to $\underline{K}(X)$ is a ring homomorphism.

Formula (2.15) is evidently suggested by Serre's ``$\Tor$ formula'' in intersection theory. This formula proves the second part of the following proposition, the first following immediately from formula (2.15) and a calculation of local $\underline{\Tor}_i$:

\sgaStableTarget{sga6.e00.prop.2-7}\textbf{Proposition 2.7.} \uline{Let} $X$ \uline{be a nonsingular variety, and} $Y$ \uline{and} $Z$ \uline{two cycles in} $X$\uline{. If} $Y$ \uline{and} $Z$ \uline{meet transversely everywhere, we have:}
\sgaStableTarget{sga6.e00.eq.2-17}\begin{equation}\tag{2.17}
 \gamma(Y)\gamma(Z)=\gamma(Y.Z).
\end{equation}
\uline{If} $Y$ \uline{and} $Z$ \uline{are homogeneous and have respective dimensions} $n-i$ \uline{and} $n-j$\uline{, and if they meet without an excess component, we then have:}
%%%%%%%%%%  scan idx 63  |  folio 50  %%%%%%%%%%
\SGAPrintPage{50}%
\sgaStableTarget{sga6.e00.eq.2-18}\begin{equation}\tag{2.18}
 \gamma(Y)\gamma(Z)=\gamma(Y.Z)\qquad\mathrm{mod.}\ \underline K(X)^{i+j+1}.
\end{equation}

We
must take care that equality does not hold in general in (2.18). We shall see in \S 4 that, when $X$ is \uline{nonsingular and quasi-projective, the ring structure of} $\underline K(X)$ \uline{is compatible with its filtration and that the associated graded ring} $G(\underline K(X))$ \uline{is identified with a quotient of} $A(X)$ \uline{by a torsion subgroup (perhaps always reduced to} $0$\SGAAuthorFootnote{*}{No; for a counterexample cf.\ this Seminar XIV 4.7.}\uline{)}.

I know how to define the $\lambda^i(F)$ directly in terms of sheaves only when the field $k$ has characteristic $0$ (this is the reason that has prevented us from proving the Riemann--Roch theorem when $k$ \var{is not of characteristic $\neq0$}{is not of characteristic $0$}{EN-VAR-032}). Let $F$ be a coherent algebraic sheaf on $X$ on which the group $\mathfrak S_p$ of permutations of $p$ letters operates; we then denote by $(F)^{\mathrm{alt}}$ the set of the $x\in F$ such that $s.x=\varepsilon(s)x$ for every $s\in\mathfrak S_p$ ($\varepsilon(s)$ is the \uline{signature} of $s$). This is evidently a coherent algebraic subsheaf of $F$. This said, we have:

\sgaStableTarget{sga6.e00.thm.2-8}\textbf{Theorem 2.8.} \uline{Let} $X$ \uline{be a nonsingular quasi-projective variety over the field} $k$ of characteristic $0$. \uline{Let} $F$ \uline{be a coherent algebraic sheaf on} $X$\uline{; we then have:}
\sgaStableTarget{sga6.e00.eq.2-19}\begin{equation}\tag{2.19}
 \lambda^p(\gamma(F))=\sum_{i\geq0}(-1)^i\gamma\left((\underline{\Tor}_i(F,F,\ldots,F))^{\mathrm{alt}}\right)
\end{equation}

N.B. It is clear that the group $\mathfrak S_p$ acts on the sheaf $\var{\underline{\Tor}_i(F,..,F)}{\underline{\Tor}_i(F,\ldots,F)}{EN-VAR-033}$.

\uline{Principle of the proof}: we consider a finite resolution $\underline L$ of $F$ by locally free sheaves $L_i$, and use the following general lemma:

%%%%%%%%%%  scan idx 64  |  folio 51  %%%%%%%%%%
\SGAPrintPage{51}%
\sgaStableTarget{sga6.e00.lem.2-9}\textbf{Lemma 2.9.} \uline{Let} $\underline L=(L_i)$ \uline{be a graded locally free coherent algebraic sheaf, with finitely many nonzero components. In} $\underline K(X)$ \uline{we have the formula:}
\sgaStableTarget{sga6.e00.eq.2-20}\begin{equation}\tag{2.20}
 \underline E\left((\overset{p}{\otimes}\underline L)^{\mathrm{alt}}\right)=\wedge^p(\underline E(\underline L)).
\end{equation}
In this formula, we consider $\overset{p}{\otimes}\underline L$ as a sheaf on which the group $\mathfrak S_p$ operates, taking the gradings into account (which therefore introduce signs in the well-known way, cf.\ Cartan--Eilenberg); moreover, if $\underline M$ is a graded coherent algebraic sheaf on $X$, we put $\underline E(\underline M)=\sum_{i\geq0}(-1)^i\gamma(M_i)$.

To prove formula (2.20), we are reduced to proving the analogous formula in $\underline K(G)$, where $G$ is an algebraic group (a product of groups $\operatorname{Gl}(n,k)$, if desired) and $\underline L$ a $G$-module of finite dimension over $k$ (cf.\ Chap.\ I, \S 2, \var{example 2, for the definition}{example 1, for the definition}{EN-VAR-034} of $\underline K(G)$). In this case, this formula has already been used in theorem 1.4 to show that the second member of (1.33) satisfies the same functional equation as the first member; its verification is elementary.
\begin{flushright}q.e.d.\end{flushright}

Let $f$ be a morphism from $X$ to $Y$, where $Y$ is an algebraic space. The notion of inverse image of a vector bundle defines a homomorphism $f^!:\underline K(\xi(Y))\longrightarrow\underline K(\xi(X))$ of $\lambda$-rings. Identifying vector bundles with locally free algebraic sheaves, the notion of inverse image of a vector bundle corresponds to the notion of inverse image of an algebraic sheaf (already used in [3, no.\ 8]). When $Y$ is nonsingular, we define more generally a group homomorphism:
\sgaStableTarget{sga6.e00.eq.2-21}\begin{equation}\tag{2.21}
 f^!:\underline K(Y)\longrightarrow\underline K(X)
\end{equation}
%%%%%%%%%%  scan idx 65  |  folio 52  %%%%%%%%%%
\SGAPrintPage{52}%
by a
formula involving an alternating sum of $\Tor$, inspired by the developments at the beginning of \S 2, which the reader will make explicit. The following diagram then commutes:
\[
\begin{tikzcd}[column sep=large,row sep=large]
\underline K(\xi(Y)) \arrow[r,"f^!"] \arrow[d] & \underline K(\xi(X)) \arrow[d]\\
\underline K(Y) \arrow[r,"f^!"] & \underline K(X){,}
\end{tikzcd}
\]
which shows in particular, when $X$ and $Y$ are both nonsingular and quasi-projective, that the two definitions of $f^!$ coincide under the identification of $\underline K(X)$ with $\underline K(\xi(X))$ and of $\underline K(Y)$ with $\underline K(\xi(Y))$.

Another case in which a homomorphism (2.21) can be defined, without $Y$ necessarily being nonsingular, is that in which the functor $F\to f^!(F)$ from $\cF(Y)$ to $\cF(X)$ is exact, i.e.\ when the local rings of $X$ are \uline{flat} modules over the local rings of $Y$. It then suffices to put $f^!(\gamma_Y(F))=\gamma_X(f^!(F))$.
This case occurs in particular when $X$ is a locally trivial fibre space over $Y$ (to see this, we are reduced to the case of a product $X=T\times Y$). Note that we also verify that, in this case, we have $f^!(\underline{O}_Z)=\underline{O}_{f^{-1}(Z)}$ if $Z$ is an irreducible subvariety of $Y$, whence:
\sgaStableTarget{sga6.e00.eq.2-22}\begin{equation}\tag{2.22}
 f^!(\gamma_Y(Z))=\gamma_X(f^{-1}(Z))
\end{equation}
\uline{if} $X$ \uline{is ``flat'' over} $Y$. Without this last condition, formula (2.22) is in general incorrect, but we have the following analogue of proposition 2.7:

\sgaStableTarget{sga6.e00.prop.2-8}\textbf{Proposition 2.8.}
\uline{Let} $f:X\longrightarrow Y$ \uline{be a morphism of nonsingular varieties, and} $Z$ \uline{a cycle on} $Y$\uline{. If} $Z$ \uline{is everywhere transverse to} $f$\uline{, formula}
%%%%%%%%%%  scan idx 66  |  folio 53  %%%%%%%%%%
\SGAPrintPage{53}%
\uline{(2.22) is correct; if} $Z$ \uline{is homogeneous of codimension} $i$ \uline{and} $f^{-1}(Z)$ \uline{has no excess component, we have:}
\sgaStableTarget{sga6.e00.eq.2-23}\begin{equation}\tag{2.23}
 f^!(\gamma_Y(Z))\equiv\gamma_X(f^{-1}(Z)).\qquad\mathrm{mod.}\ \underline K(X)^{i+1}
\end{equation}
Proof analogous to that of proposition 2.7.

Finally, let us mention that we shall see in \S 4 that, \uline{when} $X$ \uline{and} $Y$ \uline{are quasi-projective and nonsingular,} $f^!$ \uline{is compatible with the filtrations of} $\underline K(X)$ \uline{and of} $\underline K(Y)$.

Again let $X$ and $Y$ be two arbitrary algebraic spaces, and $f$ a \uline{proper} morphism from $X$ to $Y$. For every coherent algebraic sheaf $F$ on $X$, the direct image sheaf $f_!(F)$ is coherent algebraic, and the same holds more generally for the sheaves $R^qf_!(F)$ (cf.\ [3]). Recall that, by definition, $R^qf_!(F)$ is the sheaf on $Y$ associated with the presheaf assigning to an open subset $U$ the group $H^q(f^{-1}(U),F)$; in particular, we have:
\sgaStableTarget{sga6.e00.eq.2-24}\begin{equation}\tag{2.24}
 R^0f_!(F)=f_!(F),\qquad R^qf_!(F)=0\quad(q>\dim X),
\end{equation}
moreover:
\sgaStableTarget{sga6.e00.eq.2-25}\begin{equation}\tag{2.25}
 \Gamma(U,R^qf_!(F))=H^q(f^{-1}(U),F)
\end{equation}
\uline{if the open subset} $U$ \uline{is affine}. By \S 2, we define a homomorphism from $\underline K(X)$ to $\underline K(Y)$, still denoted $f_!$, by the formula
\sgaStableTarget{sga6.e00.eq.2-26}\begin{equation}\tag{2.26}
 f_!(\gamma_X(F))=\sum_{q\geq0}(-1)^q\gamma_Y(R^qf_!(F)).
\end{equation}
Thus $\underline K(X)$ becomes a \uline{covariant functor in} $X$, with respect to \uline{proper} morphisms. The relation $\var{(fg)=f_!g_!}{(fg)_!=f_!g_!}{EN-VAR-035}$ (just like the relation $(fg)^!=g^!f^!$ that we forgot to mention in its place) is a parti%
%%%%%%%%%%  scan idx 67  |  folio 54  %%%%%%%%%%
\SGAPrintPage{54}%
cular case of the transitivity relations to which we alluded in \S 2, which follow from the spectral sequences of composite functors.

Let us note the following formula, valid for $Y$ nonsingular (and particularly easy to verify when $Y$ is also quasi-projective, by taking for $y$ the class of a vector bundle on $Y$):
\sgaStableTarget{sga6.e00.eq.2-27}\begin{equation}\tag{2.27}
 f_!(x.f^!(y))=f_!(x).y\qquad(\var{x.\in\underline K(X)}{x\in\underline K(X)}{EN-VAR-036},\ y\in\underline K(Y)).
\end{equation}
In the particular case where the inverse images under $f$ of the points of $Y$ are finite and contained in an affine open subset\SGAAuthorFootnote{*}{This last condition is evidently a consequence of the first (something the author still did not know in 1957!).}, we find that the functor $F\longrightarrow f_!(F)$ is exact and that $R^qf_!(F)=0$ for $q>0$; in this case, we therefore have $f_!(\gamma_X(F))=\gamma_Y(f_!(F))$. In particular, if $X$ is an algebraic subspace of $Y$ and $i$ is the \var{canonical injection from $Y$ to $X$}{canonical injection from $X$ to $Y$}{EN-VAR-037}, we have:
\sgaStableTarget{sga6.e00.eq.2-28}\begin{equation}\tag{2.28}
 i_!(\gamma_X(F))=\gamma_Y(F),
\end{equation}
where, in the second member, $F$ is considered as a sheaf on $Y$ that is zero outside $X$.

\sgaStableTarget{sga6.e00.sec.4.occ3}\subsection*{\S 4. Some technical results}

Let $X$ be an algebraic space; we saw in \S 3 that the projection $p$ from $X\times k$ onto $X$ defines a homomorphism
\sgaStableTarget{sga6.e00.eq.2-29}\begin{equation}\tag{2.29}
 p^!:\underline K(X)\longrightarrow\underline K(X\times k).
\end{equation}

\sgaStableTarget{sga6.e00.prop.2-9}\textbf{Proposition 2.9.} \uline{The homomorphism (2.29) is surjective.}

We argue by induction on $n=\dim X$; the assertion is trivial if $n<0$; assume $n\geq0$; by proposition 2.6, it suffices to
%%%%%%%%%%  scan idx 68  |  folio 55  %%%%%%%%%%
\SGAPrintPage{55}%
verify that, for every irreducible closed subset $Y$ of $X\times k$, $\gamma(Y)$ is in the image of $p^!$. It suffices to continue the reasoning assuming $X$ irreducible; if $p(Y)$ is not dense in $X$, our assertion follows from the induction hypothesis; otherwise, either $Y=X\times k$ and the assertion is trivial, or $Y$ is a hypersurface in $X\times k$. Now let $U$ be an affine open subset of $X$ consisting of nonsingular points; if $D$ is the divisor on $U\times k$ induced by $Y$, we know (commutative algebra) that it is linearly equivalent to a divisor of the form $D'\times k$, where $D'$ is a divisor on $U$. Then $\gamma_{U\times k}(D)-\gamma_{U\times k}(D'\times k)$ belongs to $\underline K(U\times k)^2$ (corollary of proposition 2.10), and, by what we have already proved, this element is of the form $p^!(x')$ with $x'\in\underline K(U)$. Thus $\gamma_{U\times k}(D)=p^!(y')$ with $y'=x'+\gamma_U(D')$; but $y'$ is the restriction of an element $y$ of $\underline K(X)$ (prop.\ 2.10), and consequently the restriction of $\var{\gamma_{X\times k}(D)-p^!(y)}{\gamma_{X\times k}(Y)-p^!(y)}{EN-VAR-038}$ to $U\times k$ is zero; by proposition 2.10, $\var{\gamma_{X\times k}(D)-p^!(y))}{\gamma_{X\times k}(Y)-p^!(y)}{EN-VAR-039}$ comes from an element of $\underline K(X\times k-U\times k)=\underline K(X'\times k)$ with $X'=X-U$. The proof is completed by applying the induction hypothesis again.
\begin{flushright}q.e.d.\end{flushright}

\sgaStableTarget{sga6.e00.cor.unnamed.occ4}\textbf{Corollary:} \uline{Let} $X$ \uline{be an algebraic space and} $p$ \uline{the projection from} $X\times k^n$ \uline{onto} $X$\uline{; then the homomorphism} $p^!:\underline K(X)\longrightarrow\underline K(X\times k^n)$ \uline{is surjective; it is even bijective when} $X$ \uline{is nonsingular.}

Induction on $n$ for the first assertion. For the second, we put $f(x)=(x,0)$, whence $pf=\underline{1}$ and $f^!p^!=1$, which proves that $p^!$ is injective.

\sgaStableTarget{sga6.e00.prop.2-10}\textbf{Proposition 2.10.} \uline{Let} $X$ \uline{be an algebraic space,} $U$ \uline{an open subset of} $X$\uline{, and} $Y$ \uline{its complement in} $X$\uline{. Then the sequence of natural homomorphisms:}
%%%%%%%%%%  scan idx 69  |  folio 56  %%%%%%%%%%
\SGAPrintPage{56}%
\sgaStableTarget{sga6.e00.eq.2-30}\begin{equation}\tag{2.30}
 \underline K(Y)\longrightarrow\underline K(X)\longrightarrow\underline K(U)\longrightarrow 0
\end{equation}
\uline{is exact.}

N.B. Compare with proposition 2.3.

The fact that $\underline K(X)\longrightarrow\underline K(U)$ is surjective follows immediately from proposition 2.5 (or also from proposition 2.6). It remains to prove that, if $x\in\underline K(X)$ has zero restriction to $\underline K(U)$, it is in the image of $\underline K(Y)$. For this, we return to the explicit definition of $\underline K(X)$ and $\underline K(U)$: we have two coherent algebraic sheaves $F$ and $G$ on $X$ with $x=\gamma_X(F)-\gamma_X(G)$ and $\gamma_U(F)-\gamma_U(G)=0$; this last relation means that we can find coherent sheaves $P'$ and $Q'$ on $U$, equipped with filtrations whose factors are pairwise isomorphic on $U$, and such that we have an isomorphism on $U$:
\[
 P=F+Q'\longrightarrow Q=G+P'.
\]
By proposition 2.5, the sheaves $P'$ and $Q'$ are restrictions to $U$ of coherent sheaves on $X$ (denoted by the same letters); likewise, their filtrations come from filtrations on the corresponding sheaves on $X$, by lemma 2.11 below. We then have:
\[
 \gamma_X(F)-\gamma_X(G)=(\gamma_X(P)-\gamma_X(Q))+(\gamma_X(P')-\gamma_X(Q')).
\]
Introducing the factors $P'_i$ and $Q'_j$ of the filtrations of $P'$ and $Q'$ on $X$, the last term is a sum of terms $\gamma_X(P'_i)-\gamma_X(Q'_i)$, where $P'_i|U\simeq Q'_i|U$. We are thus reduced to proving the following particular case: if $R$ and $S$ are coherent sheaves on $X$ such that $R|U\simeq S|U$, then $\gamma_X(R)-\gamma_X(S)$ is in the image of $\underline K(Y)$. Now let $T$ be the graph of the given isomorphism from $R|U$ onto $S|U$; it is a coherent subsheaf of $(R\times S)|U$, and therefore, by lemma 2.11, it is the restriction to $U$ of a coherent subsheaf,
%%%%%%%%%%  scan idx 70  |  folio 57  %%%%%%%%%%
\SGAPrintPage{57}%
again
denoted $T$, of $R\times S$. We then have $\gamma_X(R)-\gamma_X(S)=(\gamma_X(R)-\gamma_X(T))-(\gamma_X(S)-\gamma_X(T))$, and it suffices to prove, for example, that the first term of the second member is in the image of $\underline K(Y)$. Now the projection $u:T\longrightarrow R$ is an isomorphism over $U$, so $\gamma_X(R)-\gamma_X(T)=\gamma_X(\operatorname{Coker}u)-\gamma_X(\operatorname{Ker}u)$, and $\operatorname{Ker}u$ and $\operatorname{Coker}u$ have their support in $Y$.

It therefore remains to prove the following lemma:

\sgaStableTarget{sga6.e00.lem.2-11}\textbf{Lemma 2.11.} \uline{Let} $X$ \uline{be an algebraic space,} $F$ \uline{a coherent algebraic sheaf on} $X$\uline{,} $U$ \uline{an open subset of} $X$\uline{, and} $G$ \uline{a coherent subsheaf of} $F|U$\uline{. Then} $G$ \uline{is the restriction to} $U$ \uline{of a coherent subsheaf of} $F$\uline{.}

Let $\overline G$ be the subsheaf of $F$ whose sections on the open subset $V$ of $X$ are the sections of $F$ on $V$ whose restriction to $U\cap V$ is \var{a section of $V$}{a section of $G$}{EN-VAR-040}. Everything reduces to proving that this sheaf $\overline G$ is coherent. For this, we can assume $X$ and $U$ affine (for $U$ is a union of affine open subsets $U_i$, and $\overline G$ is then the intersection of the analogous sheaves $\overline{G|U_i}$). Then the coordinate ring of $U$ is a ring of fractions $A_S$ of the coordinate ring $A$ of $X$, and we verify immediately that, if $F$ is defined by the $A$-module $M$, then $F|U$ is defined by the $A_S$-module $M_S=A_S\otimes_A M$. The subsheaf $G$ of $F|U$ is then defined by a submodule $N'$ of $M_S$, and we see that $\overline G$ is precisely defined by the submodule $N$ of $M$ that is the inverse image of $N'$ under the canonical homomorphism from $M$ to $M_S$.
\begin{flushright}q.e.d.\end{flushright}

\sgaStableTarget{sga6.e00.cor.1}\textbf{Corollary 1 to proposition 2.10:} \uline{Let} $x\in\underline K(X)$\uline{. For} $x$ \uline{to have} \var{\uline{filtration} $\leq i$}{\uline{filtration} $\geq i$}{EN-VAR-041}\uline{, it is necessary and sufficient that there exist a closed subset} $Y$ \uline{of} $X$\uline{, of dimension} $\leq n-i$\uline{, such that the restriction of} $x$ \uline{to} $X-Y$ \uline{is zero.}

%%%%%%%%%%  scan idx 71  |  folio 58  %%%%%%%%%%
\SGAPrintPage{58}%
\sgaStableTarget{sga6.e00.cor.2.occ2}\textbf{Corollary 2.} \uline{Assume} $X$ \uline{connected of dimension} $n$\uline{, without singularities in dimension} $n-1$\uline{. If} $D$ \uline{and} $D'$ \uline{are two linearly equivalent divisors on} $X$\uline{, we have} $\gamma(D)\equiv\gamma(D')\quad\mathrm{mod.}\ \underline K(X)^2$\uline{.}

By corollary 1, we can assume $X$ nonsingular; the corollary then follows from the general formula:
\sgaStableTarget{sga6.e00.eq.2-31}\begin{equation}\tag{2.31}
 \gamma(D)=1-\gamma(L(-D))\qquad\mathrm{mod}\ \underline K(X)^2,
\end{equation}
where $L(-D)$ denotes the vector bundle of rank $1$ defined by the divisor $-D$. To verify (2.31), we write it in the form:
\sgaStableTarget{sga6.e00.eq.2-31bis}\begin{equation}\tag{2.31 bis}
 \gamma(L(-D))=1-\gamma(D)\qquad\mathrm{mod}\ \underline K(X)^2,
\end{equation}
and note that both members are homomorphisms from the group of divisors to the multiplicative group of invertible elements of the ring $\underline K(X)/\underline K(X)^2$ (note that $\underline K(X)^1/\underline K(X)^2$ has square zero, which is very easily verified from the explicit formula (2.15), using corollary 1). We are thus reduced to the case where $D$ is an irreducible hypersurface. More generally, if $D$ is a hypersurface whose components are disjoint, we have the equality:
\sgaStableTarget{sga6.e00.eq.2-32}\begin{equation}\tag{2.32}
 \gamma(D)=1-\gamma(L(-D))
\end{equation}
as follows from the exact sequence of sheaves
\[
0\longrightarrow\underline{O}_X(L(-D))\longrightarrow\underline{O}_X\longrightarrow\underline{O}_D\longrightarrow0.
\]

\sgaStableTarget{sga6.e00.thm.2-12}\textbf{Theorem 2.12.} \uline{Let} $X$ \uline{be a nonsingular quasi-projective variety, and} $Z$ \uline{and} $Z'$ \uline{two homogeneous cycles of codimension} $p$ \uline{on} $X$ \uline{that are linearly equivalent. We then have:}
\sgaStableTarget{sga6.e00.eq.2-33}\begin{equation}\tag{2.33}
 \gamma_X(Z)\equiv\gamma_X(Z')\qquad\mathrm{mod}\ \underline K(X)^{p+1}.
\end{equation}
By hypothesis, we can find a cycle $Y$ of dimension $n-p+1$ on
%%%%%%%%%%  scan idx 72  |  folio 59  %%%%%%%%%%
\SGAPrintPage{59}%
$X\times k$ ($n$ is the
dimension of $X$), and two points $a,a'\in k$ such that we have:
\[
 Z=f_a^{-1}(Y),\qquad Z'=f_{a'}^{-1}(Y),
\]
where we put $f_u(x)=(x,u)$. By formula (2.23), we then have:
\[
 \gamma_X(Z)\equiv f_a^!(\gamma_{X\times k}(Y))\qquad\mathrm{mod}\ \underline K(X)^{p+1}
\]
and an analogous formula for $\gamma_X(Z')$. It then suffices to prove that $\var{f_a^!=f_a^!}{f_a^!=f_{a'}^!}{EN-VAR-042}$, which follows immediately from proposition 2.9.
\begin{flushright}q.e.d.\end{flushright}

\sgaStableTarget{sga6.e00.cor.1.occ2}\textbf{Corollary 1:} \uline{The filtration of} $\underline K(X)$ \uline{is compatible with its ring structure, i.e.\ we have} $\underline K(X)^i\underline K(X)^j\subset\underline K(X)^{i+j}$ \uline{for all integers} $i,j\geq0$\uline{. Moreover, we have a natural surjective homomorphism of graded rings:}
\sgaStableTarget{sga6.e00.eq.2-34}\begin{equation}\tag{2.34}
 \var{\varphi:A(X)\longrightarrow G(K(X))}{\varphi:A(X)\longrightarrow G(\underline K(X))}{EN-VAR-043}
\end{equation}
(where the second member denotes the graded ring associated with $\underline K(X)$ equipped with the filtration $(\underline K(X)^i)_{i\geq0}$) \uline{defined by the condition}:
\[
 \varphi(Z^{n-p})=\gamma(Z^{n-p})\qquad\mathrm{mod}\ \underline K(X)^{p+1}.
\]
This corollary follows immediately from theorem 2.12, propositions 2.6 and 2.7, and Chow's theorem 2.1, by induction on the integer $i+j$.

\sgaStableTarget{sga6.e00.cor.2.occ3}\textbf{Corollary 2:} \uline{Let} $X$ \uline{and} $Y$ \uline{be two nonsingular quasi-projective varieties and} $f$ \uline{a morphism from} $X$ \uline{to} $Y$\uline{. Then the homomorphism} $f^!$ \var{\uline{from} $\underline K(Y)$ \uline{to} $K(X)$}{\uline{from} $\underline K(Y)$ \uline{to} $\underline K(X)$}{EN-VAR-044} \uline{is compatible with the filtrations, and we have a commutative diagram:}
\[
\begin{tikzcd}[column sep=large,row sep=large]
G(\underline K(Y)) \arrow[r,"G(f^!)"] & G(\underline K(X))\\
A(Y) \arrow[r,"f^*"] \arrow[u,"\varphi_Y"'] & A(X) \arrow[u,"\varphi_X"'] .
\end{tikzcd}
\]
%%%%%%%%%%  scan idx 73  |  folio 60  %%%%%%%%%%
\SGAPrintPage{60}%
Using the graph of $f$, we decompose $f$ into the product of an \uline{immersion} of $X$ into $X\times Y$ and a \uline{projection} from $X\times Y$ onto $Y$. For the latter the corollary is trivial, and for an immersion we proceed as for corollary 1, using proposition 2.8 in place of proposition 2.7.

\sgaStableTarget{sga6.e00.prop.2-13}\textbf{Proposition 2.13 (``cellular decomposition''):} \uline{Let} $P$ \uline{be an algebraic space (irreducible or not), equipped with an increasing family} $P^0\subset P^1\subset\cdots\subset P^n=P$ \uline{of closed subsets such that, for every} $i$ \uline{with} $0\leq i\leq n$\uline{, the connected components of} $P^i-P^{i-1}$ \uline{are algebraic spaces} $V_{i,\alpha}$ \uline{isomorphic to} $k^i$ \uline{(put} $P^{-1}=\varnothing$\uline{). Then:}
\begin{enumerate}[label=\alph*)]
\item \uline{For every algebraic space} $X$\uline{, the natural map from} $\underline K(X)\otimes\underline K(P)$ \uline{to} $\underline K(X\times P)$ \uline{deduced from the tensor product of sheaves on} $X$ \uline{and} $P$ \uline{is surjective.}
\item \uline{The} $\gamma_P(V_{i,\alpha})$ \uline{form a system of generators of the group} $\underline K(P)$\uline{.}
\end{enumerate}
The proof proceeds by induction on $n$, using the corollary of proposition 2.9 and proposition 2.10. (N.B. Unfortunately, I do not know \var{whether the $\gamma_P(V_{i;\alpha})$ form a basis}{whether the $\gamma_P(V_{i,\alpha})$ form a basis}{EN-VAR-045} of $\underline K(P)$ and whether the homomorphism considered in a) is bijective. We shall not need to know this in what follows).

\sgaStableTarget{sga6.e00.thm.2-14}\textbf{Theorem 2.14.}
\uline{Let} $x_i$ \uline{(}$1\leq i\leq q$\uline{) and} $E_j$ \uline{(}$1\leq j\leq r$\uline{) be letters, and} $n_j$ \uline{(}$1\leq j\leq r$\uline{) integers} $>0$\uline{; consider the symbols} $\lambda^k(x_i)$ \uline{with} $k\geq1$ \uline{and} $\lambda^k(E_j)$ \uline{with} $1\leq k\leq n_j$ \uline{as indeterminates, and let} $R$ \uline{be a polynomial with integer coefficients in these indeterminates. Assume that, for every graded ring} $A$\uline{, when in} $R$ \uline{we substitute for the} $x_i$ \uline{elements} $\xi_i$ \uline{and for the} $E_j$ \uline{elements} $\eta_j$ \uline{of}
%%%%%%%%%%  scan idx 74  |  folio 61  %%%%%%%%%%
\SGAPrintPage{61}%
$\Atil$ \uline{subject to the restrictions} $\varepsilon(\xi_i)=0$ \uline{and} $\eta_j=[n_j,1+\eta_j^{(1)}+\cdots+\var{\eta_j^{(nj)}}{\eta_j^{(n_j)}}{EN-VAR-046}]$ (thus $\eta_j^{(k)}=0$ for $k>n_j$), \uline{we have}:
\[
 R((\lambda^k(\xi_i)),(\lambda^k(\eta_j)))\in\Atil^p
\]
\uline{Now let} $X$ \uline{be an irreducible nonsingular quasi-projective variety. For arbitrary elements} $x'_i$ \uline{of} $\underline K(X)^1$ \uline{and vector bundles} $E'_j$ \uline{of dimension} $n_j$ \uline{on} $X$ \uline{(}$1\leq i\leq q$\uline{,} $1\leq j\leq r$\uline{), we have}
\sgaStableTarget{sga6.e00.eq.unnamed}\[
 \var{R((\lambda^k(x'_i)),(\lambda^k(\gamma(E'_j)))}{R((\lambda^k(x'_i)),(\lambda^k(\gamma(E'_j))))}{EN-VAR-047}\in\underline K(X)^p\qquad\text{\SGAAuthorFootMark{*}}\ .
\]
\SGAAuthorFootText{*}{We readily prove that the hypothesis made implies that, if $R$ is expressed in terms of the $\gamma^k(\xi_i)$, $\var{\gamma^k(E_i-\varepsilon(E_i))}{\gamma^k(E_j-\varepsilon(E_j))}{EN-VAR-048}$, then all terms of $R$ have weight $\geq p$ (where $\gamma^k$ has weight $k$). Then 2.14 is a particular case of X 1.3.2.}

\uline{Proof}: By \S 2, the $x'_i$ themselves are differences of elements of the form $\gamma(E)$, where $E$ is a vector bundle. On the other hand, assuming $X$ locally closed in a projective space $P$, it follows from \S 2 and proposition 2.5 that, for every $F\in\cF(X)$, the sheaf $F(n)$ is generated by its sections for sufficiently large $n$. If $F$ is defined by a vector bundle $E$, this proves that $E(n)$ is the inverse image of the canonical vector bundle on a Grassmannian variety $G$ under a morphism from $X$ to $G$. Considering the corresponding morphism from $X$ to $P\times G$, we see that $E$ is the inverse image of a vector bundle on $P\times G$. Thus, applying this to all the vector bundles involved, and using corollary 2 of theorem 2.12, we see that \uline{we are reduced to the case where} $X$ \uline{is a product of Grassmannians}. In this case, the conclusion will follow from the following lemma:

\sgaStableTarget{sga6.e00.lem.2-15}\textbf{Lemma 2.15.}
\uline{Assume that} $X$ \uline{is a product of Grassmannians. Then we can find a} $\lambda$\uline{-homomorphism from} $\underline K(X)$ \uline{to a} $\lambda$\uline{-ring without torsion of the form}
%%%%%%%%%%  scan idx 75  |  folio 62  %%%%%%%%%%
\SGAPrintPage{62}%
$\Atil$ \uline{(where} $A$ \uline{is a graded ring), compatible with the filtrations, and such that the corresponding homomorphism} $\var{G(K(X))\longrightarrow G(A)}{G(\underline K(X))\longrightarrow G(A)}{EN-VAR-049}$ \uline{is injective.}

To prove the lemma, we shall use Chow's theory of Chern classes; a direct proof of the lemma, or of theorem 2.14, would yield a purely ``sheaf-theoretic'' theory of Chern classes (cf.\ \S 5).

We shall consider the homomorphism $\Ctil$ from $\underline K(X)$ to $\Atil(X)$ and put $A=A(X)$; we must prove that the homomorphism $\psi=G(\Ctil)$ from $G(\underline K(X))$ to $G(A(X))\simeq A(X)$ is injective. The fact that $\Ctil$ is compatible with the filtrations holds for every nonsingular quasi-projective variety and follows readily from the corollary to proposition 2.3. On the other hand, as $X$ is a product of Grassmannians $C_j$, we have on $X$ vector bundles $E_j$, inverse images of the canonical bundles on the factors. We shall use the following (known) facts:
\begin{enumerate}[label=\alph*)]
\item On $X$, linear equivalence of cycles is equivalent to numerical equivalence, and $A(X)$ is a group of finite rank;
\item The $C^i(E_j)$ generate the ring $A(X)$.
\end{enumerate}

From a) it follows that $\varphi:A(X)\longrightarrow G(\underline K(X))$ is injective (hence bijective), by the following lemma:

\sgaStableTarget{sga6.e00.lem.2-15.occ2}\textbf{Lemma 2.15.} \uline{Let} $X^n$ \uline{be a nonsingular projective variety, and} $Z^{n-p}$ \uline{a cycle on} $X$ \uline{such that} $\varphi(Z)\in G^p(\underline K(X))$ \uline{is zero; then} $Z^{n-p}$ \uline{is numerically equivalent to zero.}

To prove this lemma, it suffices to use the additive function $\chi(F)=\sum_{i\geq0}(-1)^i\dim H^i(X,F)$ and the fact that this function equals
%%%%%%%%%%  scan idx 76  |  folio 63  %%%%%%%%%%
\SGAPrintPage{63}%
$1$
for $F=\underline{O}_{(x)}$ for every $x\in X$.

Thus, in the case considered, $A(X)$ and $G(\underline K(X))$ have the same finite rank in every degree, and $G(\underline K(X))$ has no torsion, since this holds for $A(X)$. From assertion b) above it follows, in a ``formal'' way, that the map $\Ctil$ from $\underline K(X)$ to $\widetilde{A(X)}$ is surjective, modulo finite groups. It readily follows, since $G(\underline K(X))$ and $A(X)$ have the same finite rank in every degree, that the homomorphism $G(\Ctil)$ from $G(\underline K(X))$ to $G(\widetilde{A(X)})$ is bijective modulo finite groups in every degree (use induction on the degree). As $G(\underline K(X))$ has no torsion, this homomorphism $G(\Ctil)$ is therefore injective.
\begin{flushright}q.e.d.\end{flushright}

N.B. The homomorphism $\Ctil$ is not surjective in general, even when $X$ is a projective space of dimension $2$.

\sgaStableTarget{sga6.e00.sec.5.occ2}\subsection*{\S 5. Sheaf-theoretic definition of Chern classes. Application to the study of injection morphisms.}

\sgaStableTarget{sga6.e00.thm.2-16}\textbf{Theorem 2.16.} \uline{Let} $X$ \uline{be a nonsingular quasi-projective variety. Then, for every} $\var{x\in K(X)}{x\in\underline K(X)}{EN-VAR-050}$ \uline{and every integer} $i\geq0$\uline{, the element} $\gamma^i(x-\varepsilon(x))$ \uline{has filtration} $\geq i$\uline{. Let} $c^i(x)\in G^i(\underline K(X))$ \uline{be the class of this element modulo} $\underline K(X)^{i+1}$\uline{; if we put:}
\[
 \widetilde c(x)=\left[\varepsilon(x),1+\var{\sum_{i\geq\varphi}c^i(x)}{\sum_{i\geq1}c^i(x)}{EN-VAR-051}\right]\in\widetilde{G(\underline K(X))},
\]
\uline{then} $\widetilde c$ \uline{satisfies the conditions analogous to those of theorem 2.2.}

The first assertion follows from theorem 2.14, as does the fact that $x\longrightarrow\widetilde c(x)$ is a homomorphism of $\lambda$-rings (formula (1.25) implies $\gamma_t(x+y)=\gamma_t(x)\gamma_t(y)$, which already shows that $\widetilde c$ is an additive homomorphism). The ``functoriality'' of $\widetilde c$ is verified immediately, and
%%%%%%%%%%  scan idx 77  |  folio 64  %%%%%%%%%%
\SGAPrintPage{64}%
(2.8) follows from (2.31).
\begin{flushright}q.e.d.\end{flushright}

\sgaStableTarget{sga6.e00.cor.1.occ3}\textbf{Corollary 1:} \uline{We have} $c^i(x)=\varphi(C^i(x))$\uline{.}

This follows from a uniqueness theorem for Chern classes, proved as in Hirzebruch by passage to flag varieties. We need only show that, if $Y$ is a flag variety associated with a vector bundle on $X$, then $G(\underline K(X))\longrightarrow G(\underline K(Y))$ is injective, which is verified by explicitly determining $\underline K(Y)$ with all its structures from $\underline K(X)$ (I shall not give the details here).

\sgaStableTarget{sga6.e00.cor.2.occ4}\textbf{Corollary 2:} \uline{Assume that, for every element} $Z\in A(X)$\uline{, there exists a morphism} $f$ \uline{from} $X$ \uline{to a product} $X'$ \uline{of Grassmannians such that} $Z$ \uline{is in the image of} $f^*$\uline{. Then the canonical homomorphisms}
\[
 \varphi:A(X)\longrightarrow G(\underline K(X)),\qquad \psi:G(\underline K(X))\longrightarrow A(X),
\]
\uline{satisfy in dimension} $i$ \uline{the following relations:}
\sgaStableTarget{sga6.e00.eq.2-35}\begin{equation}\tag{2.35}
 \psi\varphi=(-1)^{i-1}(i-1)!\,\underline 1,\qquad
 \varphi\psi=(-1)^{i-1}(i-1)!\,\underline 1.
\end{equation}

As $\varphi$ is surjective, it evidently suffices to prove the first formula, and for this we can assume $X$ a product of Grassmannians. As $\varphi$ is then injective, it suffices to prove the second formula (2.35), which, by corollary 1, is written
\[
 \var{\gamma^i(x-\varepsilon(X))}{\gamma^i(x-\varepsilon(x))}{EN-VAR-052}\equiv(-1)^{i-1}(i-1)!\,x\qquad\mathrm{mod}\ \underline K(X)^{i+1}
\]
for every $x\in\underline K(X)^i$. As $\psi$ is \uline{injective}, it suffices to prove the analogous relation for $\var{\Ctil(x)\in A(X)^i}{\Ctil(x)\in\Atil(X)^i}{EN-VAR-053}$, and this is then a consequence of (1.22 bis).

\sgaStableTarget{sga6.e00.rem.unnamed.occ4}\textbf{Remark:} Corollary 2 makes the validity of (2.35) in all cases plausible; there should be a direct proof of it. In
%%%%%%%%%%  scan idx 78  |  folio 65  %%%%%%%%%%
\SGAPrintPage{65}%
any case, we readily prove the second formula (2.35) modulo the torsion subgroup of $G(\underline K(X))$, and, in characteristic $0$, this same formula is a consequence of the Riemann--Roch theorem. The interest of (2.35) (when it holds) is that it shows that $\varphi$ is injective modulo torsion groups (perhaps always zero!\SGAAuthorFootnote{*}{No, cf.\ the note on page 30.}), so that we ``do not lose much'' by defining the Chern classes as elements of $G(\underline K(X))$; after taking the tensor product with $\Q$, as we shall do later for the Riemann--Roch theorem, the two definitions of Chern classes would then be identical.

Let $X$ still be nonsingular and quasi-projective; let also $Y$ be a nonsingular subvariety of $X$, and $i$ the injection of $Y$ into $X$. For $i^!$ and $i_!$, we then have the following formulas:
\begin{enumerate}[label=(\roman*)]
\item $i^!$ is a homomorphism of $\lambda$-rings
\item $i_!(x.i^!(y))=i_!(x).y$, whence in particular $i_!i^!(y)=\xi y$ with $\xi=i_!i^!(1)=\gamma_X(Y)$,
\item $i^!(i_!(x))=x.\lambda_{-1}(N)$, whence in particular $i^!(\xi)=i^!i_!(1)=\lambda_{-1}(N)$,
\item $i_!(x)i_!(x')=\var{i_!(xx'\lambda_{-1}(N)}{i_!(xx'\lambda_{-1}(N))}{EN-VAR-054}$,
\end{enumerate}
where, by definition, $N$ is the class in $\underline K(Y)$ of the conormal bundle of $Y$ in $X$. Formulas (i) and (ii) hold for arbitrary morphisms, not necessarily injections, and (iii) and (iv) are verified from definition (2.15) and the analogous formula for inverse images (taking for $x$ and $x'$ classes of vector bundles, and using the fact that
\sgaStableTarget{sga6.e00.eq.2-36}\begin{equation}\tag{2.36}
 \Tor^{\underline{O}_X}_i(\underline{O}_Y,\underline{O}_Y)=\underline{O}_Y(\wedge^i N).
\end{equation}
)

%%%%%%%%%%  scan idx 79  |  folio 66  %%%%%%%%%%
\SGAPrintPage{66}%
Formula (iv) expresses that $i_!$ is a ring homomorphism when we introduce on $\underline K(Y)$ a new multiplicative structure, by the definition $x_N x'=xx'\lambda_{-1}(N)$ (cf.\ Chapter I, \S 4). It is completed by the formula:
\begin{enumerate}[label=(\roman*)]
\setcounter{enumi}{4}
\item $i_!(\lambda^p(N,x))=\lambda^p(i_!(x))$ for $p\geq1$,
\end{enumerate}
which means that $i_!$ defines a $\lambda$-homomorphism from $\underline K(Y)_N$ to $\underline K(X)$. This formula (v), certainly true in all cases, has for the moment been proved only in characteristic $0$, since it follows from the conjunction of theorems 1.4 and 2.8.

We have analogous formulas for $i_*$ and $i^*$:
\begin{enumerate}[label=(\roman* bis)]
\item $i^*$ is a homomorphism of graded rings.
\item $i_*(x.i^*(y))=i_*(x).y$, whence in particular $i_*(i^*(y))=\xi' y$, with $\xi'=i_*(1)$,
\item $i^*(i_*(x))=x.(-1)^q N^q$, whence in particular $i^*(\xi')=(-1)^q N^q$,
\item $i_*(x)i_*(x')=i_*(xx'(-1)^q N^q)$,
\item $i_*$ increases degrees by $q$ units.
\end{enumerate}

In these formulas, $q$ is the codimension of $Y$ in $X$, and we put $N^i=C^i(N)$. Only formulas (iii bis) and (iv bis) require a proof; I do not know whether they are well known (or even whether they are true!!!), and merely note that they follow \var{by reduction of formulas (ii) and (iv)}{by reduction of formulas (iii) and (iv)}{EN-VAR-055}, \uline{provided that we work in} $G(\underline K(Y))$ \uline{and} $G(\underline K(X))$ \uline{and not in} $A(Y)$ \uline{and} $A(X)$. We then use the fact that $i_*$ and $i^*$ arise by passage to the quotient from $i_!$ and $i^!$, noting that $i_!$ shifts the filtration by $q$ units.

%%%%%%%%%%  scan idx 80  |  folio 67  %%%%%%%%%%
\SGAPrintPage{67}%
N.B. All the ``universal'' formulas relating $i_!$ and $i^!$ (resp.\ $i_*$ and $i^*$) are formal consequences of the preceding ones, as one can convince oneself without much difficulty.

From formula (v) follows the formula
\[
 i_!(\gamma^p(N,x))=\gamma^p(i_!(x)).
\]
The conjunction of theorem 2.14 and proposition 1.5 shows that $\gamma^p(N,x)$ has, for $p=q+j$, filtration $\geq j$; denoting by $G_p(N,x)$ its class in $G^j(\underline K(Y))$, theorem 2.16 shows that we have
\[
 i_*(G_p(N,x))=c^p(i_!(x)).
\]
Finally, proposition 1.5 together with theorem 2.14 allows us to express $G_p(N,x)$ as a polynomial with integer coefficients in the $c^i(N)$, the $c^i(x)$ and $\varepsilon(x)$; we deduce the theorem:

\sgaStableTarget{sga6.e00.thm.2-17}\textbf{Theorem 2.17.} \uline{Assume that the field} $k$ \uline{has characteristic} $0$ \uline{and let} $i:Y\longrightarrow X$ \uline{be an injection morphism of nonsingular quasi-projective varieties. We then have formulas (i) to (v), and, moreover, for every integer} $j\geq0$\uline{, the formula:}
\sgaStableTarget{sga6.e00.eq.2-37}\begin{equation}\tag{2.37}
 c^{q+j}(i_!(x))=i_*\left(\frac{\widetilde c(x)*\lambda_{-1}(\widetilde c(N))}{(-1)^qN^q}\right)^{(j)}.
\end{equation}

The right-hand side of formula (2.37) is an abbreviated way of denoting the polynomial in the components of $\widetilde c(x)$ and $\widetilde c(N)$ described in proposition 1.5.

\sgaStableTarget{sga6.e00.cor.unnamed.occ5}\textbf{Corollary:} \uline{Under the preceding conditions, we have:}
\sgaStableTarget{sga6.e00.eq.2-38}\begin{equation}\tag{2.38}
 \ch(\widetilde c(i_!(x)))=i_*\left(\ch(\widetilde c(x))\,\mathcal T(\check N)^{-1}\right)
\end{equation}
\uline{where} $\ch$ \uline{is the Chern homomorphism defined by formula} (1.26).

%%%%%%%%%%  scan idx 81  |  folio 68  %%%%%%%%%%
\SGAPrintPage{68}%
It suffices to use (2.37) written in the form:
\sgaStableTarget{sga6.e00.eq.2-37bis}\begin{equation}\tag{2.37 bis}
 \widetilde c(i_!(x))=1+i_*\left(\frac{\widetilde c(x)*\var{\lambda_{-1}(c(N))}{\lambda_{-1}(\widetilde c(N))}{EN-VAR-056}-1}{(-1)^qN^q}\right),
\end{equation}
the explicit formula (1.29), formula (iv bis), and finally (1.30).

We can write (2.38) in the form
\sgaStableTarget{sga6.e00.eq.2-38bis}\begin{equation}\tag{2.38 bis}
 \ch_X(F)=i_*\left(\ch_Y(F)\,\mathcal T(T_{X/Y})^{-1}\right),
\end{equation}
where $F$ is a coherent algebraic sheaf on $Y$ (considered on the left-hand side as a sheaf on $X$ zero outside $Y$), $\ch_X$ and $\ch_Y$ are the Chern homomorphisms taken respectively on $X$ and $Y$, and $T_{X/Y}$ is the normal bundle of $Y$ in $X$.

\sgaStableTarget{sga6.e00.rem.unnamed.occ5}\textbf{Remarks:} Unlike (2.38), formula (2.37) is a ``formula without denominators''. It has been proved (even in characteristic $0$) only by working in $G(\underline K(X))$ and not in $A(X)$, which may not be the same thing. It is certainly true in every characteristic and in $A(X)$. Let us note that we immediately verify that the images of the two sides under $i^*$ are the same, whence it follows, on applying $i_*$, that \var{the product of their difference by $N^q$}{the product of their difference by $\xi'$}{EN-VAR-057} is zero. This suggests a completely different method of proof, by trying to obtain $Y$ as the transverse inverse image of a subvariety $Y'$ in a variety $X'$ such that, in dimension $\leq n=\dim X$, the class of $Y'$ in $A(X')$ is not a divisor of $0$\dots. Unfortunately, we cannot hope that $X'$ is a Grassmannian or something close to one.

We shall note that, if we disregard torsion elements, formula (2.38) is equivalent to (2.37).




% Additional macro definitions
\providecommand{\cS}{\mathcal S}
\providecommand{\cD}{\mathcal D}
\providecommand{\cE}{\mathscr E}
\providecommand{\cG}{\mathscr G}
\providecommand{\cL}{\mathscr L}
\providecommand{\cM}{\mathscr M}
\providecommand{\cP}{\mathscr P}
\providecommand{\cK}{\mathcal K}
\providecommand{\Db}{\mathrm D}
\providecommand{\Kb}{\mathrm K}
\providecommand{\Trc}{\operatorname{Tr}}
\providecommand{\ob}{\operatorname{ob}}
\providecommand{\Coh}{\operatorname{Coh}}
\providecommand{\Cone}{\operatorname{Cone}}
\providecommand{\im}{\operatorname{Im}}
\providecommand{\Id}{\operatorname{Id}}
\providecommand{\RR}{\mathrm{R.R.}}
\providecommand{\RRR}{\mathrm{R.R.R.}}
\providecommand{\GrR}{\mathrm{G}}
\providecommand{\pr}{\operatorname{pr}}


%%%%%%%%%%  scan idx 82  |  folio 69  %%%%%%%%%%
\SGAPrintPage{69}%
\sgaStableTarget{sga6.e00.sec.6.occ2}\section*{\S 6. The Riemann--Roch theorem}

Its statement is as follows: \uline{let} \(f\) \uline{be a} proper \uline{morphism from a nonsingular quasi-projective variety} \(Y\) \uline{to a variety} \(X\) \uline{of the same kind. Then we have:}
\sgaStableTarget{sga6.e00.eq.2-39}\begin{equation}\tag{2.39}
 f_*\bigl(\ch_Y(x)\,\mathcal T(Y)\bigr)
 =
 \var{\ch_X\bigl(f_*(x)\bigr)}{\ch_X\bigl(f_!(x)\bigr)}{EN-VAR-058}\,\mathcal T(X).
\end{equation}
(Both sides are in \(A(X)\otimes\mathbb Q\) and \(\mathcal T(X)=\mathcal T(C(T_X))\).)

It allows us to calculate, up to torsion elements, the Chern classes of \(f_!(x)\) from the Chern classes of \(x\), those of \(T_X\) and \(T_Y\), and the homomorphism \(f_*\). Or, again, it allows us to determine the homomorphism
\(f_!\) from \(\underline K(Y)\otimes\mathbb Q\) to \(\underline K(X)\otimes\mathbb Q\), when we identify these groups by means of the homomorphisms \(\ch_Y\) and \(\ch_X\) respectively with \(A(Y)\otimes\mathbb Q\) and \(A(X)\otimes\mathbb Q\) (if formula (2.35) is true!), using \(f_*\) and the Chern classes of \(T_X\) and \(T_Y\).

I know how to prove formula (2.39) only in characteristic \(0\)\SGAAuthorFootnote{*}{See below for a proof valid in every characteristic, based on a different principle (p.\,51).}, and as an equality in \(G(\underline K(X))\otimes\mathbb Q\). For this, since \(Y\) is locally closed in a projective space \(P\), we decompose \(f\) into the product of an injection morphism from \(Y\) to \(P\times X\), and a projection morphism from \(P\times X\) onto \(X\), and prove (2.39) for \var{each of these two morphisms..}{each of these two morphisms.}{EN-VAR-059}

For an injection morphism, we see immediately that formulas (2.38) and (2.39) are equivalent (the characteristic \(0\) hypothesis entered only through theorem 2.8, which allows us to express the \(\lambda^i(\gamma(F))\), for a coherent algebraic sheaf \(F\), by \uline{local} cohomological invariants of \(F\)).

For a projection morphism \(f:P\times X\longrightarrow X\), we can, by proposition 2.13, assume that \(x\) is of the form \(y\otimes z\), with \(y\in\underline K(P)\) and
%%%%%%%%%%  scan idx 83  |  folio 70  %%%%%%%%%%
\SGAPrintPage{70}%
\(z\in\underline K(X)\).
Then the \var{K\"Unneth}{K\"unneth}{EN-VAR-062} formula gives \(f_!(x)=\chi(y)z\), where, for a coherent algebraic sheaf \(F\) on \(P\), we put:
\[
  \chi(F)=\sum_{i\geq 0}(-1)^i\dim H^i(P,F).
\]
We are then immediately reduced to proving the formula
\sgaStableTarget{sga6.e00.eq.2-40}\begin{equation}\tag{2.40}
  \chi(F)=\varkappa_r\bigl(\ch_P(F)\,\mathcal T(P)\bigr)
\end{equation}
where \(\varkappa_r\) denotes the component of degree \(r=\dim P\). By proposition 2.13, we can assume that \(F=\underline{O}_Q\), where \(Q\) is a linear subvariety of \(P\). Then \(\chi(F)=1\), since it is the arithmetic genus of \(Q\); it remains to see that the right-hand side of formula (2.40) equals \(1\).

Indeed, if \(V\) is a hyperplane of \(P\) and we put \(\xi=\gamma_P(\underline{O}_V)\) and \(L=\ell(V)\), we shall have:
\[
  1-\xi=\gamma_P(E(-L))
\]
(where \(E(-L)\) is the vector bundle defined by \(-L\)), whence:
\[
  C_P(1-\xi)=[1,1-L],
\]
and consequently
\[
  \ch_P(1-\xi)=e^{-L},\qquad\text{whence}\qquad \ch_P(\xi)=1-e^{-L},
\]
and therefore
\[
  \ch_P(\xi^p)=(1-e^{-L})^p.
\]
Now, if \(p\) is \var{the dimension of \(Q\) in \(P\)}{the codimension of \(Q\) in \(P\)}{EN-VAR-060}, we have \(\xi^p=\gamma_P(Q)\); we deduce
\[
  \ch_P(F)\,\mathcal T(P)
  =
  \bigl(L/(1-e^{-L})\bigr)^{r+1}(1-e^{-L})^p
\]
and we must verify that the coefficient of \(L^r\) in this formal series in \(L\) equals \(1\), which follows from the characteristic property of the formal series \(L/(1-e^{-L})\), namely that the coefficient of \(L^q\) in \(\bigl(L/(1-e^{-L})\bigr)^{q+1}\) equals \(1\) (put \(\var{q=r+1-p}{q=r-p}{EN-VAR-061}\)).


%%%%%%%%%%  scan idx 84  |  folio 71  %%%%%%%%%%
\SGAPrintPage{71}%
\sgaStableTarget{sga6.e00.bibliographical-references}\subsection*{Bibliographical references:}
\begin{enumerate}[label={[\arabic*]}]
\item \uline{Hirzebruch}: \var{Neue topologischen Methoden}{Neue topologische Methoden}{EN-VAR-063} in der algebraischen Geometrie.
\item \uline{Serre}: Faisceaux algébriques cohérents, Ann. of Maths.
\item \uline{Grothendieck}: Sur les faisceaux algébriques cohérents, in Séminaire Cartan 1956/57.
\end{enumerate}

\sgaStableTarget{sga6.e00.proof-of-the-riemann-roch-theorem}\section*{Proof of the Riemann--Roch theorem}
\begin{center}
(Grothendieck, 1 November 1957)
\end{center}

\sgaStableTarget{sga6.e00.reminder-of-notation}\subsection*{* Reminder of notation}
\SGAAuthorFootText{unmarked}{The passages between asterisks \textsuperscript{*}\dots\textsubscript{*} are interpolations by the copyist.}
[These are those of the ``report on Riemann--Roch'', cited as R.R.R. in what follows.]

All varieties considered are nonsingular and quasi-projective (i.e.\ isomorphic to a locally closed subvariety of a projective space). The base field is algebraically closed, of arbitrary characteristic.

For such a variety \(X\), we denote by \(K(X)\) the abelian group generated by the classes of coherent sheaves on \(X\) (modulo the identification of an extension with a sum); the alternating sum of Tor makes \(K(X)\) a ring. If \(\underline F\) is a sheaf (resp.\ a vector bundle \(E\), resp.\ a sub%
%%%%%%%%%%  scan idx 85  |  folio 72  %%%%%%%%%%
\SGAPrintPage{72}%
variety \(Y\) of \(X\)), we denote by \(\gamma_X(\underline F)\) (resp.\ \(\gamma_X(E)\), resp.\ \(\gamma_X(Y)\)) the element of \(K(X)\) corresponding to \(\underline F\) (resp.\ to \(\underline E\), resp.\ to \(\underline{O}_Y\)). (In fact, we often write \(\underline F\), \(E\) instead of \(\gamma_X(\underline F)\)\dots).

We denote by \(A(X)\) the graded ring of classes of cycles on \(X\) (for rational equivalence---cf.\ Chow); \(A^i(X)\) consists of cycles of codimension \(i\).

For a vector bundle \(E\), we denote by \(C(E)=\sum C^i(E)\), \(C^i(E)\in A^i(X)\), its Chern class (defined as the inverse image of Schubert cycles); by linearity we deduce the \(C^i(y)\), for every \(y\in K(X)\). If \(C(y)=\prod(1+a_i)\), we put, following Hirzebruch:
\[
  \ch(y)=\sum e^{a_i},\qquad
  T(y)=\prod\frac{a_i}{1-e^{-a_i}}.
\]
These are elements of \(A(X)\otimes\mathbb Q\).

When \(y\) is the tangent bundle of \(X\), we write \(T(X)\) instead of \(T(y)\).

For a vector bundle \(G\), we denote by \(\lambda^pG\) its \(p\)-th exterior power, and put \(\lambda_{-1}(G)=\sum(-1)^p\lambda^p(G)\); this is an element of \(K(X)\).

For a proper regular map \(f\) from a variety \(Y\) to a variety \(X\), we define (additive) homomorphisms:
\[
  f^*:A(X)\longrightarrow A(Y)\qquad\text{and}\qquad
  f_*:A(Y)\longrightarrow A(X).
\]

The first is homogeneous of degree \(0\) and multiplicative. The second is homogeneous of degree \(\dim.X-\dim.Y\).

We also define homomorphisms:
\[
  f^!:K(X)\longrightarrow K(Y)\qquad\text{and}\qquad
  f_!:K(Y)\longrightarrow K(X).
\]

For a coherent sheaf \(\underline F\) on \(X\), \(f^!(\underline F)\) equals, by definition, the alternating sum of the \(\Tor_p^{\underline{O}_X}(\underline{O}_Y,\underline F)\). For vector bundles, this is simply the inverse-image operation.

%%%%%%%%%%  scan idx 86  |  folio 73  %%%%%%%%%%
\SGAPrintPage{73}%
For a coherent sheaf \(G\) on \(Y\), \(f_*(\underline G)\) is coherent on \(X\). Denoting by \(R^pf_*\) the \(p\)-th derived functor of \var{the functor \(f\) thus defined}{the functor \(f_*\) thus defined}{EN-VAR-064}, we have, by definition, \(f_!(\underline G)=\sum(-1)^p R^pf_*(\underline G)\).\quad\textsubscript{*}

\sgaStableTarget{sga6.e00.riemann-roch-theorem}\textbf{RIEMANN--ROCH THEOREM.} \uline{Let} \(f:Y\longrightarrow X\) \uline{be a proper morphism} (\textsuperscript{*} regular map \textsubscript{*}) \uline{from a nonsingular quasi-projective variety} \(Y\) \uline{to another} \(X\)\uline{, and let} \(y\in K(Y)\)\uline{. We have:}
\sgaStableTarget{sga6.e00.eq.1}\begin{equation}\tag{1}
  f_*\bigl(\ch(y)T(Y)\bigr)=\ch(f_!(y))T(X)
  \quad\text{\uline{in} } A(X)\otimes\mathbb Q.
\end{equation}

\uline{Proof} (\textsuperscript{*} abridged \textsubscript{*}).

\sgaStableTarget{sga6.e00.lem.1}\textbf{LEMMA 1.} \uline{Consider proper morphisms of varieties} \(Z\xrightarrow{g}Y\xrightarrow{f}X\)\uline{. If the theorem} \(\RR\) \uline{is true for} \(f\) \uline{and} \(g\)\uline{, it is true for} \(f\circ g\)\uline{. If it is true for} \(fg\) \uline{and} \(g\)\uline{, it is true for} \(f\)\uline{. If} \(f_*\) \uline{is injective, and the theorem} \(\RR\) \uline{is true for} \(fg\) \uline{and} \(f\)\uline{, it is true for} \(g\)\uline{.}

Trivial. (\textsuperscript{*} In fact, when we assume \(\RR\) true for \(fg\) and \(g\), we must also assume that \(g_!\) is surjective to be able to assert that \(\RR\) is true for \(f\). \textsubscript{*})

As a proper morphism is the composite of an injection and a projection \(X\times P\longrightarrow X\) (where \(P\) is a projective space), we are reduced to treating these two cases separately. The second, as we have seen, reduces to a standard calculation, by proposition 2.13 (``cellular decomposition'') of R.R.R. We shall therefore assume \(Y\subset X\), with the injection morphism \(i\). Formula (1) then becomes:
\sgaStableTarget{sga6.e00.eq.2}\begin{equation}\tag{2}
  \ch(i_!(y))=i_*\bigl(\ch(y)T(E)^{-1}\bigr),
\end{equation}
where \(E=T_{X/Y}\) is the normal bundle of \(Y\) in \(X\).

%%%%%%%%%%  scan idx 87  |  folio 74  %%%%%%%%%%
\SGAPrintPage{74}%
\sgaStableTarget{sga6.e00.lem.2}\textbf{LEMMA 2.} \uline{Assume that there exists on} \(X\) \uline{a vector bundle} \(G\) \uline{of rank} \(p=\codim_XY\)\uline{, such that:}
\sgaStableTarget{sga6.e00.eq.3}\begin{equation}\tag{3}
  \gamma_X(Y)=i_!(1)=\lambda_{-1}(G)\,;\qquad
  i^!(G)=\check E\,;\qquad
  (-1)^pC^p(G)=Y.
\end{equation}
\uline{If we then have:}
\sgaStableTarget{sga6.e00.eq.4}\begin{equation}\tag{4}
  y=i^!(x),\qquad \text{\uline{with} }x\in K(X),
\end{equation}
\uline{formula (2) is valid.}

The proof is immediate from the fact that \(\ch\) is a ring homomorphism, and that we have:
\sgaStableTarget{sga6.e00.eq.5}\begin{equation}\tag{5}
  \ch(\lambda_{-1}(G))=(-1)^pC^p(G)T(\check G)^{-1}.
\end{equation}

\sgaStableTarget{sga6.e00.cor.unnamed.occ6}\textbf{COROLLARY.} \(\RR\) \uline{is true if} \(Y\) \uline{is a hypersurface and} \(y\) \uline{is of the form} \(i^!(x)\)\uline{, with} \(x\in K(X)\)\uline{.}

\sgaStableTarget{sga6.e00.lem.3}\textbf{LEMMA 3.} \uline{Assume that} \(X=Y\times P\)\uline{, with} \(P\) \uline{a projective space, and that} \(i\) \uline{is the injection} \(u\longrightarrow (u,v_0)\) \(v_0\in P\)\uline{. Then} \(i_*\) \uline{is injective and} \(\RR\) \uline{is true.}

This is pure standard calculation.

Assume again that \(X,Y,y\) are arbitrary (still with \(Y\subset X\)). Let \(X'\) be the variety obtained by blowing up \(X\) along \(Y\), and \(Y'\) the inverse-image hypersurface of \(Y\). We have a diagram of canonical maps
\[
\begin{tikzcd}
Y \arrow[r,"i"] & X\\
Y' \arrow[u,"g"] \arrow[r,"j"] & X' \arrow[u,"f"']
\end{tikzcd}
\]
(\(j\) being the canonical map). The inverse image of \(E\) on \(Y'\) has a subbundle of rank \(1\), whose dual is denoted by \(L\). Moreover, \(L^{-1}\) is also the restriction to \(Y'\) of the bundle \(L(Y')\) defined on \(X'\) by the divisor \(Y'\). We put \(F=E/L^{-1}\). Elementary calculations give:

%%%%%%%%%%  scan idx 88  |  folio 75  %%%%%%%%%%
\SGAPrintPage{75}%
\sgaStableTarget{sga6.e00.lem.4}\textbf{LEMMA 4.} \uline{We have the formulas (where} \(p=\codim_XY\)\uline{):}
\sgaStableTarget{sga6.e00.eq.6}\sgaStableTarget{sga6.e00.eq.7}\sgaStableTarget{sga6.e00.eq.8}\sgaStableTarget{sga6.e00.eq.9}\begin{align}
  f^!i_!(y)&=j_!\bigl(\lambda_{-1}(\check F)y\bigr). \tag{6}\\
  \lambda_{-1}(\check F)&\equiv \gamma^{p-1}(E)\lambda^p(\check E)
    \quad\text{mod.}(1-L)\text{\SGAAuthorFootMark{*}}. \tag{7}\\
  f_*(1)&=1,\quad\text{\uline{whence} }f_*f^*=1. \tag{8}\\
  f^*i_*(\alpha)&\equiv j_*\bigl(\alpha C^{p-1}(F)\bigr)
    \quad\text{mod.}\Ker f_*\quad\text{\uline{for every} }\alpha\in A(Y). \tag{9}
\end{align}
\SGAAuthorFootText{*}{This unnecessarily complicated formula can be replaced, for example, by VII 4.1.}

[To simplify notation, we have made \(K(Y)\) operate on \(K(Y')\) and \(A(Y)\) on \(A(Y')\) by means of the homomorphisms \(g^!\) and \(g^*\). {*}In formula (7), it seems that \(\gamma^{p-1}(E)\) denotes \(\lambda^{p-1}(E-2)\), that is, \(\sum_{i+j=p-1}\binom{-p}{j}\lambda^i(E)\); the congruence holds in \(K(Y')\).\textsubscript{*}]

Formula (6) follows from the equality:
\sgaStableTarget{sga6.e00.eq.6bis}\begin{equation}\tag{6 bis}
  \Tor_i^{\underline{O}_X(v)}\bigl(\underline{O}_Y(v),\underline{O}_{Y'}(v')\bigr)
  =
  \wedge^i\check F(v')
  \quad (v'\in Y',\ v=g(v')).
\end{equation}

Formula (7) follows from the formulas
\sgaStableTarget{sga6.e00.eq.10}\begin{equation}\tag{10}
  g_!(L^{-i})=0\quad\text{for }0<i\leq p-1,\qquad g_!(1)=1,
\end{equation}
which themselves follow from \var{K\"Unneth}{K\"unneth}{EN-VAR-065} and the cohomology of projective spaces.

Formula (8) is trivial; finally, (9) follows from \(g_*(C^{p-1}(F))=1\), which is deduced, for example, from the immediate formulas:
\sgaStableTarget{sga6.e00.eq.10bis}\begin{equation}\tag{10 bis}
  g_*(\xi^i)=0\quad\text{if }0\leq i<p-1,\qquad
  g_*(\xi^{p-1})=1,
\end{equation}
with \(\xi=C^1(L)\in A^1(Y')\).

[In fact, (9) is very probably an identity, and not merely a congruence.]


%%%%%%%%%%  scan idx 89  |  folio 76  %%%%%%%%%%
\SGAPrintPage{76}%
Formula (7) immediately gives \(y\lambda_{-1}(\check F)\in j^!(K(X'))\), provided that we have:
\sgaStableTarget{sga6.e00.eq.11}\begin{equation}\tag{11}
  y\,\gamma^{p-1}(E)\lambda^p(\check E)\in i^!(K(X)).
\end{equation}
Assume that this condition (11) is satisfied. By the corollary to lemma 2, \(\RR\) is true for \((Y',X',y\lambda_{-1}(\check F))\). On the other hand, to prove (2), it suffices, by formula (8), to prove that
\sgaStableTarget{sga6.e00.eq.12}\begin{equation}\tag{12}
  f^*\ch(i_!(y))
  \equiv
  f^*\bigl(i_*(\ch(y)T(E)^{-1})\bigr)
  \quad\text{mod.}\Ker f_*.
\end{equation}
Now, this last formula is verified immediately, using (6), (9), \(\RR\) for \(\var{(Y',X',y\lambda_{-1}(F))}{(Y',X',y\lambda_{-1}(\check F))}{EN-VAR-066}\) and the formula:
\sgaStableTarget{sga6.e00.eq.13}\begin{equation}\tag{13}
  \ch(\lambda_{-1}(\check F))=C^{p-1}(F)T(F)^{-1}\quad\text{(cf.\ formula (5)).}
\end{equation}

Thus, \(\RR\) is proved whenever \(y\,\gamma^{p-1}(E)=0\), in particular (since \(\gamma^{p-1}(E)\) has filtration \(\geq p-1\) \textsuperscript{*}in \(K(Y)\) filtered by the codimension of supports\textsubscript{*}) whenever \(p-1>\dim.Y\), i.e.\ whenever \(\dim.Y<\frac{\dim.X-1}{2}\). Now, we reduce to this case by combining lemma 1 and lemma 3 (taking for \(P\) a projective space of sufficiently large dimension \(r\), for example \(r>\dim.X+1\)). Q.E.D.

\sgaStableTarget{sga6.e00.rem.unnamed.occ6}N.B. We have used the fact that \(\gamma^{p-1}(E)\) has filtration \(\geq p-1\). We can prove this fact directly, without passing through Grassmannians as in \(\mathrm{R.R.R.}\) 2.16, by studying the bundle \(Y'\) of projective spaces associated with \(E\). This method gives, more generally, th.\ 2.14 of \(\mathrm{R.R.R.}\)

\sgaStableTarget{sga6.e00.extract-from-a-letter-of-grothendieck}\subsection*{Extract from a letter of Grothendieck}
\dots{} ``morally'', the new method rests on the determination of \(K(X)\) and \(A(X)\) when \(X\) is either a projective bundle associated with a vector bundle, or a variety obtained by blowing up a non%
%%%%%%%%%%  scan idx 90  |  folio 77  %%%%%%%%%%
\SGAPrintPage{77}%
singular subvariety. In this second case, I have not made the complete determination; a little something is missing. With the notation of the short paper, we must prove formula (9) (but as an \uline{equality}):
\sgaStableTarget{sga6.e00.eq.9.occ2}\begin{equation}\tag{9}
  f^*(i_*(\alpha))=j_*(\alpha C^{p-1}(F)),\qquad \alpha\in A(Y).
\end{equation}
We observe that the two sides have the same image under \(f_*\), and the same image under \(j^*\), so that the question is equivalent to:
\sgaStableTarget{sga6.e00.eq.9bis}\begin{equation}\tag{9 bis}
  \Ker f_*\cap\Ker j^*=0.
\end{equation}
In any case, (9) is true upon reduction to \(G(K(X'))\) (which destroys a little torsion, at most), since it then follows from (6); there is therefore hardly any doubt that it is true.

\dots I also point out to you that the \uline{equality} (9) follows very easily from (6) without passing through the troublesome (6 bis)\dots

\begin{flushleft}
p.c.c.J-P.S.
\end{flushleft}

\newpage

%%%%%%%%%%  scan idx 91  |  folio 78  %%%%%%%%%%
\SGAPrintPage{78}%
\sgaStableTarget{sga6.e01.expose}\section*{\underline{\underline{GENERALITIES ON FINITENESS CONDITIONS IN}}\\[3pt] \underline{\underline{DERIVED CATEGORIES}}}
\begin{center}
\underline{by L. ILLUSIE}
\end{center}

\sgaStableTarget{sga6.e01.sec.0}\subsection*{0. Introduction}
The purpose of Exposés I to IV is to develop, with the generality appropriate to this Seminar, the formalism of finiteness conditions in the derived categories of ringed topoi.

As was said in Exp.\ 0, it was the need to define, on arbitrary schemes, ``Grothendieck groups'' with good variance properties that required us to generalize and relax the notions of finiteness used until now.

A classical notion such as that of a coherent sheaf on a ringed space \((X,\underline{O}_X)\) becomes uninteresting as soon as \(\underline{O}_X\) ceases to be coherent. One might think of replacing it by the notion of finite presentation, but this has the drawback that the kernel of an epimorphism of Modules of finite presentation is in general no longer of finite presentation. We arrive at a satisfactory notion by observing that, when \(\underline{O}_X\) is coherent, a coherent sheaf \(F\) is not only of finite presentation but of \(n\)-finite presentation for every \(n\in\mathbb N\), ``\(n\)-finite presentation'' meaning that there locally exists an exact sequence

%%%%%%%%%%  scan idx 92  |  folio 79  %%%%%%%%%%
\SGAPrintPage{79}%
\[
  L_n\longrightarrow L_{n-1}\longrightarrow \ldots \longrightarrow L_0\longrightarrow F\longrightarrow 0,
\]
where the \(L_i\) are free of finite type. With the sheaf \(\underline{O}_X\) no longer assumed coherent, let us say that an \(\underline{O}_X\)-Module \(F\) is \uline{pseudo-coherent} when it satisfies the preceding condition, i.e.\ is of \(n\)-finite presentation for every \(n\). Then the notion of pseudo-coherence has, with respect to short exact sequences, the same stability property as coherence: when two of the terms of a short exact sequence of \(\underline{O}_X\)-Modules are pseudo-coherent, the third is too. We readily generalize the preceding notion to complexes of sheaves. A \uline{complex} \(F\) of \(\underline{O}_X\)-Modules is called \uline{pseudo-coherent} when it is \(n\)-pseudo-coherent for every \(n\in\mathbb Z\), ``\(n\)-pseudo-coherent'' meaning that there locally exists a quasi-isomorphism \(L\to F\), where \(L\) is a complex with degrees bounded above and components free of finite type in degree \(\geq n\). For a complex \(F\) concentrated in degree \(0\), it amounts to the same thing to say that \(F\) is pseudo-coherent as a complex or as a Module. The notion of a pseudo-coherent complex has excellent stability properties, described in Exp.\ I \S\S~1 and 2. First, it is stable under isomorphism in the derived category \(\Db(X)\) of the category of \(\underline{O}_X\)-Modules; moreover, when two of the objects of a distinguished triangle of \(\Db(X)\) are pseudo-coherent, the third is too: in other words, the full subcategory \(\Db(X)_{\mathrm{coh}}\) of \(\Db(X)\) formed by pseudo-coherent complexes is a triangulated subcategory. Furthermore, pseudo-coherence is preserved by derived inverse image and tensor product. When \(\underline{O}_X\) is coherent, saying that a complex \(F\) is pseudo-coherent simply means that \(F\) has locally bounded-above cohomology and that all the \(\underline H^i(F)\) are coherent.

When \(\underline{O}_X\) is a sheaf of regular local rings, every \(\underline{O}_X\)-Module
%%%%%%%%%%  scan idx 93  |  folio 80  %%%%%%%%%%
\SGAPrintPage{80}%
\(F\) that is coherent locally admits a \uline{finite} resolution to the left by free Modules of finite type, i.e.\ there locally exists an exact sequence
\[
  0\longrightarrow L_n\longrightarrow \ldots \longrightarrow L_0\longrightarrow F\longrightarrow 0,
\]
where the \(L_i\) are free of finite type. It is therefore natural, when \(\underline{O}_X\) is an arbitrary sheaf of local rings, to consider Modules with the preceding property. Such Modules are called \uline{perfect}. More generally, a \uline{complex} \(F\) of \(\underline{O}_X\)-Modules is called \uline{perfect} when there locally exists a quasi-isomorphism \(L\to F\), where \(L\) is a complex with bounded degrees and components free of finite type. It is shown in Exp.\ I \S~3 that the full subcategory \(\Db(X)_{\mathrm{parf}}\) of \(\Db(X)\) formed by perfect complexes is, like \(\Db(X)_{\mathrm{coh}}\), a triangulated subcategory, ``stable'' under derived inverse image and tensor product. We evidently have \(\Db(X)_{\mathrm{parf}}\subset\Db(X)_{\mathrm{coh}}\), but the inclusion is strict in general. Pseudo-coherence can be redefined ``by passage to the limit'' from perfection: a complex \(F\) is pseudo-coherent precisely when, for every \(n\in\mathbb Z\), \(F\) can be ``locally approximated to order \(n\) by a perfect complex'', by which we mean that there locally exists a distinguished triangle
\[
\begin{tikzcd}
& R \arrow[dl] &\\
L \arrow[rr] && F \arrow[ul] ,
\end{tikzcd}
\]
where \(L\) is perfect and \(R\) acyclic in degree \(\geq n\). Conversely, perfection is recovered from pseudo-coherence through an additional regularity condition: a complex is perfect precisely when it is pseudo-coherent and locally of finite tor-dimension (Exp.\ I \S~5). We thus recover the fact that, when the local rings are regular, every coherent sheaf is perfect, and, more generally, every pseudo-coherent complex with locally bounded cohomology is perfect.


%%%%%%%%%%  scan idx 94  |  folio 81  %%%%%%%%%%
\SGAPrintPage{81}%
Pseudo-coherence and perfection are the two fundamental notions of finiteness with which we shall work in this Seminar. Sections~1 to 5 of Exposé I are devoted to their definition and to establishing their elementary stability properties. For this, we use only two or three local properties of the category of free Modules of finite type with respect to the category of all Modules. It therefore proves convenient---and useful, particularly in Exp.\ II---to axiomatize the situation by introducing a notion of pseudo-coherence (resp.\ perfection) in a category fibered over a site relative to a fibered subcategory satisfying suitable conditions.

Sections~6 to 8 of Exp.\ I generalize to perfect complexes a number of familiar notions for locally free sheaves of finite type: rank, duality, trace of an endomorphism.

In Exp.\ II, we examine the problem of ``the existence of global resolutions'': under what conditions, on a ringed topos \((X,\underline{O}_X)\), is a pseudo-coherent (resp.\ perfect) complex \uline{globally} isomorphic in \(\Db(X)\) to a complex of locally free Modules of finite type? The importance of this question lies essentially in the fact that, at present, we do not know how to generalize to perfect complexes certain usual constructions on locally free Modules: exterior and symmetric powers, Chern classes (for more details on this subject, see Exp.\ XIV). This is why it is convenient to have workable sufficient criteria for the existence of global resolutions, which allow us to reduce certain questions about pseudo-coherent (resp.\ perfect) complexes to analogous questions about locally free sheaves of finite type. Such criteria
%%%%%%%%%%  scan idx 95  |  folio 82  %%%%%%%%%%
\SGAPrintPage{82}%
are given in Exp.\ II. In particular, it is shown that global resolutions exist in the following cases: \(X\) is the Zariski topos of an affine scheme, or more generally of a quasi-compact scheme with an ample invertible sheaf, or even an ample family of invertible sheaves in Kleiman's sense (for example a regular scheme); or again \(X\) is the topos of sheaves of sets on a compact topological space, \(\underline{O}_X\) the sheaf of continuous complex-valued functions. As an illustration---and to bring out the flexibility of the notions introduced---the appendix gives a ``purely sheaf-theoretic'' definition of the index of a family of elliptic operators.

Exposé III studies the stability of finiteness conditions under derived direct image. To obtain reasonable statements, we must ``relativize'' pseudo-coherence and perfection. We work here in the setting of ordinary schemes, which suffices for the Seminar, but there is no doubt that sooner or later an analogous theory will have to be developed for relative analytic schemes or spaces. Let \(p:X\to S\) be an \(S\)-scheme locally of finite type and \(F\) a complex of \(\underline{O}_X\)-Modules; we say that \(F\) is \uline{pseudo-coherent} (\uline{resp.\ perfect}) \uline{relative to} \(p\) when we can locally embed (by a closed immersion) \(X\) in a smooth \(S\)-scheme \(X'\) so that \(F\), extended by zero to \(X'\), is pseudo-coherent (resp.\ perfect). Pseudo-coherence relative to \(p\) is especially interesting when \(S\) is not locally noetherian, since in the opposite case it coincides with ordinary pseudo-coherence. Moreover, as one might expect, saying that \(F\) is perfect rel.\ to \(p\)
%%%%%%%%%%  scan idx 96  |  folio 83  %%%%%%%%%%
\SGAPrintPage{83}%
amounts to saying that \(F\) is pseudo-coherent rel.\ to \(p\) and locally of finite tor-dimension rel.\ to \(p\) (i.e.\ rel.\ to the sheaf of rings \(p^{-1}(\underline{O}_S)\)). The central theorem of Exp.\ III is the \uline{finiteness theorem}, which asserts (in a slightly more precise form) that, for a \uline{proper} morphism \(f:X\to Y\) of \(S\)-schemes locally of finite type, the functor \(\mathbb{R}f_*\) sends complexes pseudo-coherent rel.\ to \(S\) to complexes pseudo-coherent rel.\ to \(S\); this theorem is in fact a conjecture, but is nevertheless proved in the following two special cases: a)~\(S\) is locally noetherian b)~\(f\) is projective; unfortunately, it seems that extending it to the general case is of the same order of difficulty as the analogous theorem in Analytic Geometry (Grauert's ``theorem''). Combining the finiteness theorem with an essentially trivial formula called the \uline{projection formula}, we obtain workable criteria for the stability of relative perfection under direct image. As a corollary, we recover ``Grauert's continuity and semicontinuity theorems'' (EGA III 7.6).

Exposé IV translates the results of Exposés I to III into the language of ``Grothendieck groups''. On a ringed topos \((X,\underline{O}_X)\), we denote by \(K^\bullet(X)\) (resp.\ \(K_\bullet(X)\)) the ``Grothendieck group'' of the category of complexes that are perfect and of finite tor-dimension (resp.\ pseudo-coherent and with bounded cohomology). The group \(K^\bullet(X)\) is a ring, and \(K_\bullet(X)\) a module over \(K^\bullet(X)\). \(K^\bullet\) is a contravariant functor (as the raised dot suggests), while \(K_\bullet\) is a covariant functor for proper pseudo-coherent morphisms (of schemes or analytic spaces, and subject to having the finiteness theorem\dots). There are also unusual variances,
%%%%%%%%%%  scan idx 97  |  folio 84  %%%%%%%%%%
\SGAPrintPage{84}%
covariance of \(K^\bullet\), contravariance of \(K_\bullet\), for morphisms satisfying suitable regularity hypotheses. Finally, the globalization criteria of Exp.\ II allow us to relate these to the ``Grothendieck groups'' defined naively from coherent sheaves or locally free sheaves of finite type.

\sgaStableTarget{sga6.e01.sec.1}\section*{1. Preliminary definitions}

\sgaStableTarget{sga6.e01.sec.1.1}\subsection*{1.1. Fibered categories with additive (resp. abelian, resp. triangulated) fibres}

\sgaStableTarget{sga6.e01.para.1-1-1}1.1.1. Let \(\underline{S}\) be a category, and \(\underline{C}\) a fibered \(\underline{S}\)-category (SGA 1 VI 6.1) whose fibres are additive (resp. triangulated) categories. We shall say that \(\underline{C}\) is an \uline{additive} \(\underline{S}\)\uline{-category} (\uline{resp.\ triangulated}) (or that \(\underline{C}\) is additive (resp.\ triangulated) over \(\underline{S}\)) when, for every arrow \(f:X\to Y\) of \(\underline{S}\), the inverse-image functor (determined up to unique isomorphism) \(f^*:\underline{C}_Y\longrightarrow \underline{C}_X\) is additive (resp.\ exact). We shall say that \(\underline{C}\) is an \uline{abelian} \(\underline{S}\)-category when \(\underline{C}\) is an additive \(\underline{S}\)-category with abelian fibres, and that \(\underline{C}\) is a \uline{flat} abelian \(\underline{S}\)-category when, moreover, the inverse-image functors are exact.

\sgaStableTarget{sga6.e01.para.1-1-2}1.1.2. Let \(\underline{C}\) be a triangulated \(\underline{S}\)-category. We define, for example using a cleavage of \(\underline{C}\) \var{(SGA 1 VI 7,1)}{(SGA 1 VI 7.1)}{EN-VAR-068}, a fibered \(\underline{S}\)-category \(\Trc(\underline{C})\) whose fibre at each \(X\in \ob(\underline{S})\) is the category \(\Trc(\underline{C}_X)\) of distinguished triangles of \(\underline{C}_X\).

\sgaStableTarget{sga6.e01.para.1-1-3}1.1.3. Let \(\underline{C}\) and \(\underline{C}'\) be additive (resp. triangulated) \(\underline{S}\)-categories and \(F:\underline{C}\to\underline{C}'\) a functor. We call \(F\) an \uline{additive} (resp.\ \uline{exact}) \(\underline{S}\)\uline{-functor} when \(F\) is a cartesian \(\underline{S}\)-functor that is additive (resp.\ exact) on each fibre.
% Additional macro definitions
\providecommand{\cC}{\mathcal C}
\providecommand{\cD}{\mathcal D}
\providecommand{\cK}{\mathcal K}
\providecommand{\cO}{\mathcal O}
\providecommand{\cS}{\mathcal S}
\providecommand{\Cpx}{\mathrm C}
\providecommand{\Kpx}{\mathrm K}
\providecommand{\Dpx}{\mathrm D}
\providecommand{\D}{\mathrm D}
\providecommand{\Ho}{\mathrm H}
\providecommand{\RHom}{\mathrm R\!\operatorname{Hom}}
\providecommand{\Ob}{\operatorname{ob}}
\providecommand{\Coker}{\operatorname{Coker}}
\providecommand{\Ker}{\operatorname{Ker}}
\renewcommand{\Im}{\operatorname{Im}}
\providecommand{\Hom}{\operatorname{Hom}}
\providecommand{\Cone}{\operatorname{Cone}}
\providecommand{\id}{\operatorname{Id}}
\providecommand{\qis}{\xrightarrow{\sim}}
%%%%%%%%%%  scan idx 98  |  folio 85  %%%%%%%%%%
\SGAPrintPage{85}%
\sgaStableTarget{sga6.e01.para.1-1-3.occ2}When \(\underline{C}\) and \(\underline{C}'\) are triangulated \(\underline{S}\)-categories and \(F\) is an exact \(\underline{S}\)-functor, \(F\) defines a cartesian \(\underline{S}\)-functor \(\Trc(\underline{C})\longrightarrow \Trc(\underline{C}')\).

\sgaStableTarget{sga6.e01.para.1-1-4}1.1.4. Let \(\underline{C}\) be a triangulated \(\underline{S}\)-category, and \(\underline{C}'\) a full fibered sub-\(\underline{S}\)-category of \(\underline{C}\). We call \(\underline{C}'\) a \uline{triangulated sub-}\(\underline{S}\)\uline{-category} of \(\underline{C}\) when, for every \(X\in\Ob \underline{S}\), \(\underline{C}'_X\) is a triangulated subcategory of \(\underline{C}_X\) ([V] chap.\ I \S~1 2.3).

\sgaStableTarget{sga6.e01.para.1-1-5}1.1.5. Let \(\underline{C}\) be an additive \(\underline{S}\)-category. We define (SGA 4 XVII 2.2, 2.4) an additive (resp.\ triangulated) \(\underline{S}\)-category \(\Cpx(\underline{C})\), (resp.\ \(\Kpx(\underline{C})\)), whose fibre at \(X\in\Ob \underline{S}\) is the category \(\Cpx(\underline{C}_X)\) (resp.\ \(\Kpx(\underline{C}_X)\)) of complexes\SGAAuthorFootnote{1}{with differential of degree \(+1\).} (resp.\ complexes ``up to homotopy'') of \(\underline{C}_X\). When \(\underline{C}\) is a flat abelian \(\underline{S}\)-category, we define (loc.\ cit.) a triangulated \(\underline{S}\)-category \(\Dpx(\underline{C})\), whose fibre at \(X\in\Ob \underline{S}\) is the derived category \(\Dpx(\underline{C}_X)\); for an arrow \(f:X\to Y\) of \(\underline{S}\), we shall denote by \(f^*:\Dpx(\underline{C}_Y)\longrightarrow \Dpx(\underline{C}_X)\) the inverse-image functor, obtained from \(f^*:\Kpx(\underline{C}_Y)\longrightarrow \Kpx(\underline{C}_X)\) by simple passage to the quotient\SGAAuthorFootnote{2}{for \(a\in\Ob \Dpx(\underline{C}_Y)\) (resp.\ \(\mathrm{Fl}\,\Dpx(\underline{C}_Y)\)), we shall often write, in abbreviated form, \(a_X\) instead of \(f^*(a)\).}. We define full triangulated sub-\(\underline{S}\)-categories \(\Kpx^*(\underline{C})\) (resp.\ \(\Dpx^*(\underline{C})\)) \((*=+,-,b)\) by the usual degree conditions.

\sgaStableTarget{sga6.e01.sec.1.2}1.2. \uline{Terminology.} Let \(\underline{S}\) be a site, \(\underline{C}\) a flat abelian \(\underline{S}\)-category (1.1), and \(\underline{C}_0\) a strictly full fibered sub-\(\underline{S}\)-category.

Let \(F\in\Ob\underline{C}_X\), for \(X\in\Ob \underline{S}\). The set of \(f:U\to X\) in \(\underline{S}_{/X}\) such that there exists an epimorphism \(E\longrightarrow f^*F\) with \(E\in\Ob\underline{C}_{0U}\) is a sieve on \(X\). When it is covering, we shall say that \(F\) is of \(\underline{C}_0\)-\uline{finite type} (or simply of finite type when there is no risk of confusion).


%%%%%%%%%%  scan idx 99  |  folio 86  %%%%%%%%%%
\SGAPrintPage{86}%
We shall say that \(\underline{C}_0\) is \uline{locally stable under kernels of epimorphisms} when the following condition is satisfied: for every \(X\in\Ob \underline{S}\) and every epimorphism \(u:E\longrightarrow F\) of \(\underline{C}_X\) with \(E\) and \(F\in\Ob\underline{C}_{0X}\), the kernel of \(u\) is locally in \(\underline{C}_0\) (i.e.\ the set of \(f:U\to X\) in \(\underline{S}_{/X}\) such that \(f^*\Ker(u)\in\Ob\underline{C}_{0U}\) is a refinement of \(X\)).

By the standard procedure of decomposing a long exact sequence into short exact sequences, we see that the preceding condition is equivalent to the following (apparently stronger) condition: for every \(X\in\Ob \underline{S}\) and every exact sequence of \(\underline{C}_X\)
\sgaStableTarget{sga6.e01.eq.unnamed}\begin{equation}\tag{*}
  0\longrightarrow E^0\longrightarrow E^1\longrightarrow\ldots\longrightarrow E^n\longrightarrow 0
\end{equation}
with \(E^i\in\Ob\underline{C}_{0X}\) for \(1\leq i\leq n\), \(E^0\) is locally in \(\underline{C}_0\).

We shall say that \(\underline{C}_0\) is \uline{locally quasi-stable under kernels of epimorphisms} when, for every exact sequence such as \((*)\), \(E^0\) is of \(\underline{C}_0\)-finite type.

On the other hand, we shall say that \(\underline{C}_0\) is \uline{locally liftable in} \(\underline{C}\) when, for every \(X\in\Ob \underline{S}\), every epimorphism \(u:E\longrightarrow F\) of \(\underline{C}_X\) with \(F\in\Ob\underline{C}_{0X}\) locally admits a section. For \(\underline{C}_0\) to be locally liftable in \(\underline{C}\), it is necessary and sufficient that, for every diagram of \(\underline{C}_X\) (\(X\in\Ob \underline{S}\))
\[
\begin{tikzcd}
& H \arrow[d]\\
E \arrow[r,"u"] & F
\end{tikzcd}
\]
where \(u\) is an epimorphism and \(H\in\Ob\underline{C}_0\), the set of \(U\) over \(X\) such that there exists a commutative diagram
% SGA display variant EN-VAR-069
\SGACaptureAppState{EN-VAR-069}%
\ifnum\SGAReaderForm=1\relax\SGAAppMark{EN-VAR-069}\fi
\ifnum\SGAReaderForm=0\relax
\[
\begin{tikzcd}
& H_U \arrow[d] \arrow[dl,"u"']\\
E_U \arrow[r,"U"] & F_U
\end{tikzcd}
\]
\else
\[
\begin{tikzcd}
& H_U \arrow[d] \arrow[dl]\\
E_U \arrow[r,"u_U"] & F_U
\end{tikzcd}
\]
\fi
is a refinement of \(X\).

%%%%%%%%%%  scan idx 100  |  folio 87  %%%%%%%%%%
\SGAPrintPage{87}%
We shall say that \(\underline{C}_0\) is \uline{locally quasi-liftable in} \(\underline{C}\) when the following weaker condition is satisfied: for every \(X\in\Ob \underline{S}\), and every epimorphism \(u:E\longrightarrow F\) of \(\underline{C}_X\) with \(F\in\Ob\underline{C}_{0X}\), the set of \(f:U\to X\) in \(\underline{S}_{/X}\) such that there exists \(v:F'\longrightarrow f^*E\) with \(F'\in\Ob\underline{C}_{0U}\) such that \(f^*(u)v\) is an epimorphism is a refinement of \(X\).

This condition is equivalent to the following: for every diagram of \(\underline{C}_X\)
\[
\begin{tikzcd}
& H \arrow[d]\\
E \arrow[r,"u"] & F ,
\end{tikzcd}
\]
where \(u\) is an epimorphism, and \(H\in\Ob\underline{C}_0\), the set of \(U\) over \(X\) such that there exists a commutative diagram
\[
\begin{tikzcd}
G \arrow[r,"v"] \arrow[d] & H_U \arrow[d]\\
E_U \arrow[r,"u_U"] & F_U ,
\end{tikzcd}
\]
with \(G\in\Ob\underline{C}_0\) and \(v\) an epimorphism is a refinement of \(X\).

When \(\underline{S}\) is a chaotic site, for example the punctual site (one object, one morphism), we shall omit the adverb ``locally'' in the preceding definitions.

\sgaStableTarget{sga6.e01.sec.1.3}\subsection*{1.3. Examples}

a) Let \(\underline{S}\) be a punctual site, \(A\) a ring, and \(\underline{C}\) the category of left \(A\)-modules. The full subcategory \(\underline{C}_0\) of \(\underline{C}\) formed by the projective (left) \(A\)-modules of finite type is stable under kernels of epimorphisms and liftable in \(\underline{C}\); an \(A\)-module is of \(\underline{C}_0\)-finite type if and only if it is of finite type in the
%%%%%%%%%%  scan idx 101  |  folio 88  %%%%%%%%%%
\SGAPrintPage{88}%
ordinary sense. The full subcategory \(\underline{C}_1\) of \(\underline{C}\) formed by the \(A\)-modules of finite type is quasi-liftable in \(\underline{C}\) (but not liftable in general); when \(A\) is noetherian, \(\underline{C}_1\) is stable under kernels of epimorphisms. The full subcategory \(\underline{C}'_0\) of \(\underline{C}\) formed by the free \(A\)-modules of finite type is liftable in \(\underline{C}\) and quasi-stable (not stable in general) under kernels of epimorphisms.

b) Let \((S,\underline{O}_S)\) be a space ringed by local rings, \(\underline{S}\) the site of open subsets of \(S\), and \(\underline{C}\) the fibered \(\underline{S}\)-category (even split (SGA 1 VI 9)) defined by the functor \(U\longmapsto\) (category of left \(\underline{O}_U\)-Modules); \(\underline{C}\) is a flat abelian \(\underline{S}\)-category (1.1). Let \(\underline{C}_0\) be the strictly full additive sub-\(\underline{S}\)-category of \(\underline{C}\) defined by \(U\longmapsto\) (category of free \(\underline{O}_U\)-Modules of finite type). An \(\underline{O}_S\)-Module \(M\) is locally of \(\underline{C}_0\)-finite type if and only if \(M\) is of finite type in the ordinary sense (EGA \(0_{\mathrm I}\) 5.2). It follows from (EGA \(0_{\mathrm I}\) 5.2.7 and 5.4.9) that \(\underline{C}_0\) is locally stable under kernels of epimorphisms and locally liftable in \(\underline{C}\).

c) Let \(G\) be a group, \(\underline{S}\) the topos of \(G\)-sets \var{(SGA 4 V 3.11) \(k\)}{(SGA 4 V 3.11), \(k\)}{EN-VAR-070} a field (regarded as a constant Ring of \(\underline{S}\)), and \(\underline{C}\) the split \(\underline{S}\)-category defined by the functor \(U\longmapsto\) (category of left \(k_U\)-Modules); \(\underline{C}\) is a flat abelian \(\underline{S}\)-category. Let \(\underline{C}_0\) be the strictly full additive sub-\(\underline{S}\)-category of \(\underline{C}\) defined by \(U\longmapsto\) (category of free \(k_U\)-Modules of finite type).
Since localization in \(\underline{S}\) is equivalent to forgetting the operations of \(G\), it is clear that a \(G\)-\(k\)-module \(M\) is locally of \(\underline{C}_0\)-finite type if and only if \(M\) is finite-dimensional as a vector space over \(k\), and that \(\underline{C}_0\) is locally stable under kernels of epimorphisms and locally liftable in \(\underline{C}\).

The reader is left to examine the ``Galois'' variant of c) in which \(G\) is a profinite group and \(\underline{S}\) the topos of sets on which \(G\) operates continuously.

%%%%%%%%%%  scan idx 102  |  folio 89  %%%%%%%%%%
\SGAPrintPage{89}%
In fact, examples b) and c) are special cases of the following example:

d) Let \((\underline{S},\underline{O}_{\underline{S}})\) be a ringed topos (cf.\ SGA 4 IV 2.10)\SGAAuthorFootnote{*}{Here \(\underline{O}_{\underline{S}}\) is a ring of \(\underline{S}\), with identity element, but not necessarily commutative.}, and \(\underline{C}\) the split \(\underline{S}\)-category defined by \(U\longmapsto\) (category of left \(\underline{O}_U\)-Modules); \(\underline{C}\) is a flat abelian \(\underline{S}\)-category. Let \(\underline{C}_0\) be the strictly full additive sub-\(\underline{S}\)-category of \(\underline{C}\) defined by \(U\longmapsto\) (category of free \(\underline{O}_U\)-Modules of finite type). An \(\underline{O}_{\underline{S}}\)-Module locally of \(\underline{C}_0\)-finite type will simply be said to be of \uline{finite type}. The category \(\underline{C}_0\) is locally quasi-stable under kernels of epimorphisms and locally liftable in \(\underline{C}\), as follows from

\sgaStableTarget{sga6.e01.lem.1-3-1}\var{\textbf{Lemma 1.3.1,}}{\textbf{Lemma 1.3.1.}}{EN-VAR-071} \uline{Let} \(P\) \uline{be an} \(\underline{O}_{\underline{S}}\)\var{\uline{-Module, Equivalent conditions}}{\uline{-Module. Equivalent conditions}}{EN-VAR-072}:
\begin{enumerate}[label=(\roman*)]
\item \(P\) \uline{is of finite type and the functor} \(\underline{\Hom}(P,.)\) \uline{is exact};
\item \(P\) \uline{is of finite type and every epimorphism} \(M\to P\) \uline{locally admits a section};
\item \(P\) \uline{is locally a direct factor of a free Module of finite type}.
\end{enumerate}
The proof is immediate and is left as an exercise to the reader.

\sgaStableTarget{sga6.e01.sec.1.4}\subsection*{1.4. An ``extension'' lemma (cf. EGA \(0_{\mathrm{III}}\) 11.9.1)}

\sgaStableTarget{sga6.e01.prop.1-4-1}\textbf{Proposition 1.4.1.} \uline{Let} \var{\(S\) \uline{be a site,}}{\(\underline{S}\) \uline{be a site,}}{EN-VAR-073} \var{\(C\) \uline{a flat abelian}}{\(\underline{C}\) \uline{a flat abelian}}{EN-VAR-074} \var{\(S\)\uline{-category (1.1),}}{\(\underline{S}\)\uline{-category (1.1),}}{EN-VAR-075} \(\underline{C}_0\) \var{\uline{a strictly full additive sub-}\(S\)}{\uline{a strictly full additive sub-}\(\underline{S}\)}{EN-VAR-076}\uline{-category (1.1) of} \var{\(C\)\uline{; assume that}}{\(\underline{C}\)\uline{; assume that}}{EN-VAR-077} \(\underline{C}_0\) \uline{is locally quasi-liftable in} \var{\(C\) \uline{(1.2). Let}}{\(\underline{C}\) \uline{(1.2). Let}}{EN-VAR-078} \(S\in\Ob \underline{S}\)\uline{,} \(u:E\longrightarrow F\) \uline{be a morphism (of degree} \(0\)\uline{) of complexes of} \(\underline{C}_S\) \uline{(differentials of degree} \(+1\)\uline{),} \(G\) \uline{the cone of} \(u\) \uline{(cf.\ [V] chap.\ I \S~1 2.1),} \(n\in Z\)\uline{; assume that} \(\Ho^n(G)\) \uline{is}
%%%%%%%%%%  scan idx 103  |  folio 90  %%%%%%%%%%
\SGAPrintPage{90}%
\uline{of} \(\underline{C}_0\)-%
\uline{finite type (1.2). Then the set of} \(X\in\Ob \underline{S}_{/S}\) \uline{satisfying property} \((P)\) \uline{below is a refinement of} \(S\).

\((P)\): there exist an object \(M\) of \(\underline{C}_{0X}\) and arrows of \(\underline{C}_X\)
\[
  r:M\longrightarrow Z^{n+1}(E_X)=
  \Ker\!\left(E_X^{n+1}\xrightarrow{d^{n+1}}E_X^{n+2}\right)
\]
\[
  \text{and}\quad s:M\longrightarrow F_X^n
\]
such that the following diagram of \(\underline{C}_X\)
\[
\begin{tikzcd}[column sep=3em,row sep=4em]
\ldots \arrow[r] & E_X^{n-1}
 \arrow[r,"{\binom{d_{E_X}^{n-1}}{0}}"]
 \arrow[d,"u_X^{n-1}"] &
E_X^n\oplus M
 \arrow[r,"{(d_{E_X}^n,r)}"]
 \arrow[d,"{(u_X^n,s)}"] &
E_X^{n+1} \arrow[r] \arrow[d,"u_X^{n+1}"] &
\ldots\\
\ldots \arrow[r] & F_X^{n-1} \arrow[r] &
F_X^n \arrow[r] &
F_X^{n+1} \arrow[r] &
\ldots
\end{tikzcd}
\]
is a morphism of complexes \(u':E'\to F_X\) (i.e.\ \(d_{F_X}^n s=u_X^{n+1}r\)), and that
\begin{enumerate}[label=(\roman*)]
\item \(\Ho^{n+1}(u')\) is a monomorphism, with \(\Im\Ho^{n+1}(u')=\Im\Ho^{n+1}(u_X)\),
\item \(\Ho^n(u')\) is an epimorphism.
\end{enumerate}

In pictorial terms, we may say that, by modifying \(E\) and \(u\) locally, \uline{in degree \(n\) only}, by adding to \(E^n\) an object of \(\underline{C}_0\), we can locally make \(\Ho^{n+1}(u)\) injective, without changing its image, and \(\Ho^n(u)\) surjective.

\uline{Proof}. The procedure is as follows: first we show that we can choose \(r\) and \(s\) locally so as to make \(\Ho^{n+1}(u)\) injective; we verify 
%%%%%%%%%%  scan idx 104  |  folio 91  %%%%%%%%%%
\SGAPrintPage{91}%
that the finiteness hypothesis on \(\Ho^n(G)\) is preserved; a new local correction in degree \(n\) then allows us to make \(\Ho^n(u)\) surjective.

The exact cohomology sequence
\[
  \longrightarrow \Ho^n(G)\longrightarrow
  \Ho^{n+1}(E)\longrightarrow
  \Ho^{n+1}(F)\longrightarrow
\]
shows that the kernel of \(\Ho^{n+1}(u)\) is of \(\underline{C}_0\)-finite type. In other words, the set of \(X\) over \(S\) such that there exists an exact sequence
\[
  M\longrightarrow \Ho^{n+1}(E_X)
  \xrightarrow{\Ho^{n+1}(u)}
  \Ho^{n+1}(F_X),
\]
with \(M\in\Ob\underline{C}_{0X}\), is a refinement of \(S\). Since \(\underline{C}_0\) is locally quasi-liftable in \(\underline{C}\), we can moreover require the arrow \(M\longrightarrow \Ho^{n+1}(E_X)\) to factor through \(Z^{n+1}(E_X)\) as
\[
  M\xrightarrow{r}Z^{n+1}(E_X)\longrightarrow \Ho^{n+1}(E_X).
\]
The composite arrow
\[
  M\xrightarrow{u_X^{n+1}r}Z^{n+1}(F_X)\longrightarrow \Ho^{n+1}(F_X)
\]
is zero by construction, so \(u_X^{n+1}r\) factors through \(B^{n+1}(F_X)=\operatorname{Im}\!\left(F_X^n\xrightarrow{d^n}F_X^{n+1}\right)\). Applying quasi-liftability again, we can require \(u_X^{n+1}r\) to factor through \(F_X^n\) as
\[
  M\xrightarrow{s}F_X^n\longrightarrow B^{n+1}(F_X)\longrightarrow Z^{n+1}(F_X).
\]
We then check trivially that \(u'\), defined as in (1.4.1), is a morphism of complexes satisfying (i).

Let \(G'\) be the cone of \(u':E'\to F_X\). Let us show that \(\Ho^n(G')\) is of \(\underline{C}_0\)-finite type. By construction, there is an exact sequence
\[
  0\longrightarrow E_X\xrightarrow{v}E'\longrightarrow M[-n]\longrightarrow 0,
\]
%%%%%%%%%%  scan idx 105  |  folio 92  %%%%%%%%%%
\SGAPrintPage{92}%
with \(u'v=u_X\). To compare \(G_X\) and \(G'\), we can, for example, form the ``octahedron'' of \(\Kpx(\underline{C}_X)\)
% SGA display variant EN-VAR-079
\SGACaptureAppState{EN-VAR-079}%
\ifnum\SGAReaderForm=1\relax\SGAAppMark{EN-VAR-079}\fi
\ifnum\SGAReaderForm=0\relax
\[
\begin{tikzpicture}[>=stealth,baseline=(GX.base),node distance=1.0cm]
  \node (EX) at (0,0) {$E_X$};
  \node (Ep) at (2.6,0) {$E'$};
  \node (FX) at (5.2,0) {$F_X$};
  \node (M)  at (1.35,1.25) {$M[-n]$};
  \node (GX) at (2.65,2.55) {$G_X$};
  \node (Gp) at (3.95,1.25) {$G'$};
  \node at (3.22,1.72) {$(*)$};
  \draw[->] (EX) -- node[below] {$v$} (Ep);
  \draw (Ep) -- node[below] {$u'$} (FX);
  \draw[->] (Ep) -- (M);
  \draw[->] (M) -- (GX);
  \draw[->] (GX) -- (Gp);
  \draw[->] (FX) -- (Gp);
  \draw[->] (M) to[bend left=38] node[left] {$+1$} (EX);
  \draw[->] (Gp) to[bend left=28] node[above] {$+1$} (M);
  \draw[->] (Gp) to[bend left=28] node[right] {$+1$} (Ep);
  \draw[->] (GX) to[bend left=55] node[left] {$+1$} (EX);
  \draw[->] (FX) to[bend right=42] (GX);
  \node[anchor=base] at (6.3,0 |- FX.base) {.};
\end{tikzpicture}
\]
\else
\[
\begin{tikzpicture}[>=stealth,baseline=(GX.base),node distance=1.0cm]
  \node (EX) at (0,0) {$E_X$};
  \node (Ep) at (2.6,0) {$E'$};
  \node (FX) at (5.2,0) {$F_X$};
  \node (M)  at (1.35,1.25) {$M[-n]$};
  \node (GX) at (2.65,2.55) {$G_X$};
  \node (Gp) at (3.95,1.25) {$G'$};
  \node at (3.22,1.72) {$(*)$};
  \draw[->] (EX) -- node[below] {$v$} (Ep);
  \draw[->] (Ep) -- node[below] {$u'$} (FX);
  \draw[->] (Ep) -- (M);
  \draw[->] (M) -- (GX);
  \draw[->] (GX) -- (Gp);
  \draw[->] (FX) -- (Gp);
  \draw[->] (M) to[bend left=38] node[left] {$+1$} (EX);
  \draw[->] (Gp) to[bend left=28] node[above] {$+1$} (M);
  \draw[->] (Gp) to[bend left=28] node[right] {$+1$} (Ep);
  \draw[->] (GX) to[bend left=55] node[left] {$+1$} (EX);
  \draw[->] (FX) to[bend right=42] (GX);
  \node[anchor=base] at (6.3,0 |- FX.base) {.};
\end{tikzpicture}
\]
\fi
The exact cohomology sequence of the distinguished triangle \((*)\) shows that \(\Ho^n(G_X)\longrightarrow \Ho^n(G')\) is an epimorphism, hence that \(\Ho^n(G')\) is indeed of \(\underline{C}_0\)-finite type.

The preceding correction therefore reduces us to the case where \(\Ho^{n+1}(u)\) is a monomorphism, which we shall now assume. The exact cohomology sequence of \(E\to F\to G\) then gives an exact sequence
\begin{equation}\tag{.}
  \Ho^n(E)\xrightarrow{\Ho^n(u)}
  \Ho^n(F)\longrightarrow \Ho^n(G)\longrightarrow 0.
\end{equation}
Since \(\Ho^n(G)\) is of \(\underline{C}_0\)-finite type, the set of \(X\) over \(S\) such that there exists an epimorphism \(M\longrightarrow \Ho^n(G_X)\) with \(M\in\Ob\underline{C}_{0X}\) is, by definition, a refinement of \(S\). Using the quasi-liftability of \(\underline{C}_0\) in \(\underline{C}\), we can require this epimorphism to factor through \(Z^n(F_X)\) as
\[
  M\xrightarrow{s}Z^n(F_X)\longrightarrow \Ho^n(F_X)\longrightarrow \Ho^n(G_X).
\]
It is clear that \(u'\), defined as in (1.4.1) with \((r=0)\), is a morphism of complexes and that \(\Ho^{n+1}(u')=\Ho^{n+1}(u_X)\). On the other hand, it follows immediately from the exact sequence (.) that
%%%%%%%%%%  scan idx 106  |  folio 93  %%%%%%%%%%
\SGAPrintPage{93}%
\[
  \Ho^n(u'):\Ho^n(E_X)\oplus M\longrightarrow \Ho^n(F_X)
\]
is an epimorphism. This completes the proof of (1.4.1).

\sgaStableTarget{sga6.e01.lem.1-4-2}\var{\textbf{Lemma 1.4.2.} \uline{Let} \(C\)}{\textbf{Lemma 1.4.2.} \uline{Let} \(\underline{C}\)}{EN-VAR-080} \uline{be an abelian category,} \(u:E\to F\) \var{\uline{a morphism of complexes of} \(C\)}{\uline{a morphism of complexes of} \(\underline{C}\)}{EN-VAR-081} \uline{(resp.\ an arrow of} \(\Dpx(\underline{C})\)\uline{),} \(G\) \uline{the cone of} \(u\)\uline{. Equivalent conditions for a given} \(n\in\mathbb Z\):
\begin{enumerate}[label=(\roman*)]
\item \(\Ho^i(G)=0\) \uline{for} \(i\geq n\);
\item \(\Ho^i(u)\) \uline{is an isomorphism for} \(i>n\) \uline{and an epimorphism for} \(i=n\).
\end{enumerate}
The proof follows immediately from the exact cohomology sequence.

\sgaStableTarget{sga6.e01.def.1-4-3}\textbf{Definition 1.4.3.} \uline{We shall say that} \(u\) \uline{is an} \(n\)-quasi-isomorphism (\uline{resp.\ an} \(n\)-isomorphism) \uline{when} \(u\) \uline{satisfies the equivalent conditions of (1.4.2).}

That being so, we deduce at once from (1.4.1) the

\sgaStableTarget{sga6.e01.cor.1-4-4}\textbf{Corollary 1.4.4.} \uline{Let} \(\underline{S},S,\underline{C},\underline{C}_0\) \uline{be as in (1.4.1). Let} \(u:E\to F\) \uline{be an} \((n+1)\)\uline{-quasi-isomorphism of complexes of} \(\underline{C}_S\) \uline{(}\(n\in\mathbb Z\) \uline{given) such that} \(\Ho^n(G)\) \uline{is of} \(\underline{C}_0\)\uline{-finite type. Then the set of} \(X\) \uline{over} \(S\) \uline{such that there exist} \((M,r,s)\) \uline{as in (1.4.1) such that} \(u'\) \uline{is an} \(n\)\uline{-quasi-isomorphism is a refinement of} \(S\)\uline{.}

In pictorial terms, using the finiteness condition on \(\Ho^n(G)\), we can locally ``gain one notch'', i.e.\ modify \(u\) locally in degree \(n\) by adding to \(E^n\) an object of \(\underline{C}_0\) so as to make \(u\) an \(n\)-quasi-isomorphism.

\sgaStableTarget{sga6.e01.cor.1-4-5}\textbf{Corollary 1.4.5.} \uline{Let} \(\underline{C}\) \uline{be an abelian category,} \(n\in\mathbb Z\)\uline{, and} \(u:E\to F\) \uline{an} \(n\)\uline{-quasi-isomorphism of complexes of} \(\underline{C}\)\uline{. There exists a commutative diagram of} \(\Cpx(\underline{C})\)
\[
\begin{tikzcd}
E \arrow[r,"v"] \arrow[d,"u"'] & E' \arrow[dl,"u'"]\\
F & ,
\end{tikzcd}
\]
%%%%%%%%%%  scan idx 107  |  folio 94  %%%%%%%%%%
\SGAPrintPage{94}%
\uline{where} \(u'\) \uline{is a quasi-isomorphism and} \(v\) \uline{a monomorphism inducing the identity in degree} \(\geq n\) (i.e.\ such that \(E'^i=E^i\) and \(v^i=\id(E^i)\) for \(i\geq n\)).

\uline{Proof.} Applying (1.4.4) in the case where \(\underline{S}\) is the punctual site and \(\underline{C}_0=\underline{C}\), we obtain a commutative triangle of \(\Cpx(\underline{C})\)
\[
\begin{tikzcd}
E \arrow[r,"v_{n-1}"] \arrow[d,"u"'] & E'_{\,n-1} \arrow[dl,"u'_{n-1}"]\\
F & ,
\end{tikzcd}
\]
where \(u'_{n-1}\) is an \((n-1)\)-quasi-isomorphism and \(v_{n-1}\) a monomorphism inducing the identity in degree \(\geq n\). Iterating the procedure, we find, for \(i\leq n-1\), a sequence of commutative triangles of \(\Cpx(\underline{C})\)
\[
\begin{tikzcd}
E'_i \arrow[r,"v_{i-1}"] \arrow[d,"u'_i"'] & E'_{\,i-1} \arrow[dl,"u'_{i-1}"]\\
F & ,
\end{tikzcd}
\]
where \(u'_i\) is an \(i\)-quasi-isomorphism and \(v_{i-1}\) a monomorphism inducing the identity in degree \(\geq i\). The corollary follows by passage to the limit. More precisely, denoting by \(d'_i\) the differential of \(E'_i\), we define \(E'\) by
\[
  E'^i=E^i,\quad d'^i=d^i\quad\text{for }i\geq n,
\]
and
\[
  E'^i=(E'_i)^i,\quad d'^i=(d'_i)^i\quad\text{for }i<n,
\]
then \(v\) and \(u'\) by
\[
  \var{v^j=(v_i)^j}{v^j=(v_i v_{i+1}\ldots v_{n-1})^j}{EN-VAR-082},\qquad u'^j=(u'_i)^j\qquad\text{for }i\leq j,
\]
and we check at once that \(E'\), \(v\), \(u'\) so defined answer the question.

%%%%%%%%%%  scan idx 108  |  folio 95  %%%%%%%%%%
\SGAPrintPage{95}%
In other words, briefly, we can extend an \(n\)-quasi-isomorphism to a quasi-isomorphism ``without changing anything given in degree \(\geq n\)''.

\sgaStableTarget{sga6.e01.sec.2}\section*{2. Pseudo-coherent complexes}

\sgaStableTarget{sga6.e01.para.2-0}2.0. In this section, \(\underline{S}\) denotes a site, \(S\) an object of \(\underline{S}\), \(\underline{C}\) a flat abelian \(\underline{S}\)-category (1.1), and \(\underline{C}_0\) a strictly full additive sub-\(\underline{S}\)-category of \(\underline{C}\). We assume that \(\underline{C}_0\) is locally quasi-stable under kernels of epimorphisms (1.2) and locally quasi-liftable in \(\underline{C}\) (1.2).

\sgaStableTarget{sga6.e01.def.2-1}\textbf{Definition 2.1.} \uline{Let} \(X\in\Ob \underline{S}\)\uline{,} \(E\) \uline{be a complex of} \(\underline{C}_X\)\uline{, and} \(n\in\mathbb Z\)\uline{. We shall say that} \(E\) \uline{is} strictly \(\underline{C}_0\)-\(n\)-pseudo-coherent (resp.\ \(\underline{C}_0\)-pseudo-coherent, resp.\ \(\underline{C}_0\)-perfect) \uline{when} \(E^i=0\) \uline{for} \(i\) \uline{sufficiently large and} \(\var{E^i\in\Ob\underline{C}_{0S}}{E^i\in\Ob\underline{C}_{0X}}{EN-VAR-083}\) \uline{for} \(i\geq n\) \uline{(resp.\ for every} \(i\)\uline{, resp.\ for every} \(i\) \uline{and} \(E^i=0\) \uline{for almost every} \(i\)\uline{).}

When there is no risk of confusion, we shall sometimes say ``strictly \(n\)-pseudo-coherent (resp\dots)'' instead of ``strictly \(\underline{C}_0\)-\(n\)-pseudo-coherent (resp\dots)''.

\sgaStableTarget{sga6.e01.lem.2-1-1}\textbf{Lemma 2.1.1.} \uline{Let} \(E\in\Ob\Cpx(\underline{C}_S)\) \uline{and} \(n\in\mathbb Z\)\uline{. When} \(E\) \uline{is strictly} \(n\)\uline{-pseudo-coherent and} \(\Ho^i(E)=0\) \uline{for} \(i>n\)\uline{, then} \(\Ho^n(E)\) \uline{is of finite type.}

\uline{Proof.} The sequence
\[
  0\longrightarrow Z^n(E)\longrightarrow E^n\longrightarrow E^{n+1}\longrightarrow\ldots\longrightarrow 0
\]
is exact and \(E^i\in\Ob\underline{C}_0\) for \(i\geq n\). Since \(\underline{C}_0\) is locally quasi-stable under kernels of epimorphisms, it follows that \(Z^n(E)\) is of finite type, hence so is \(\Ho^n(E)\), which is a quotient of it.

%%%%%%%%%%  scan idx 109  |  folio 96  %%%%%%%%%%
\SGAPrintPage{96}%
\sgaStableTarget{sga6.e01.prop.2-2}\textbf{Proposition 2.2.} \uline{Let} \(F\in\Ob\Dpx(\underline{C}_S)\) \uline{and} \(n\in\mathbb Z\)\uline{. Then}:

a) \uline{When} \(F\) \uline{is strictly} \(n\)\uline{-pseudo-coherent (relative to} \(\underline{C}_0\)\uline{) and} \(u:E\to F\) \uline{is a quasi-isomorphism, the set of} \(X\) \uline{over} \(S\) \uline{such that there exists a quasi-isomorphism} \(v:F'\longrightarrow E_X\) \uline{with} \(F'\) \uline{strictly} \(n\)\uline{-pseudo-coherent is a refinement of} \(S\)\uline{.}

b) \uline{The following conditions are equivalent}:
\begin{enumerate}[label=(\roman*)]
\item \uline{the set of} \(X\) \uline{over} \(S\) \uline{such that there exists a quasi-isomorphism} \(u:E\longrightarrow F_X\) \uline{of} \(\Cpx(\underline{C}_X)\)\uline{, with} \(E\) \uline{strictly} \(n\)\uline{-pseudo-coherent, is a refinement of} \(S\);
\item[\((\mathrm{i}')\)] \uline{the same condition as (i), but with} \(u\) \uline{an isomorphism of} \(\Dpx(\underline{C}_X)\);
\item \uline{the set of} \(X\) \uline{over} \(S\) \uline{such that there exists an} \(n\)\uline{-quasi-isomorphism (1.4.3)} \(u:E\longrightarrow F_X\) \uline{of} \(\Cpx(\underline{C}_X)\)\uline{, with} \(E\) \uline{strictly perfect, is a refinement of} \(S\);
\item[\((\mathrm{ii}')\)] \uline{the same condition as (ii), but with} \(u\) \uline{an} \(n\)\uline{-isomorphism of} \(\Dpx(\underline{C}_X)\).
\end{enumerate}

\uline{Proof.} a) For an integer \(i\) sufficiently large, there exists an \(i\)-quasi-isomorphism \(v_i:F'_i\longrightarrow E\), with \(F'_i\) strictly \(n\)-pseudo-coherent (take \(F'_i=0\)). Let us show that, for \(i>n\), there exists, after refining \(S\), an \((i-1)\)-quasi-isomorphism \(v_{i-1}:F'_{i-1}\longrightarrow E\), with \(F'_{i-1}\) strictly \(n\)-pseudo-coherent. By (1.4.4), it suffices to show that \(\Ho^{i-1}(M')\), where \(M'\) is the cone of \(v_i\), is of \(\underline{C}_0\)-finite type. Now, with \(M''\) denoting the cone of \(uv_i\), the hypothesis that \(u\) is a quasi-isomorphism and (1.4.2) imply that
\[
  \Ho^j(M')\qis \Ho^j(M'')\qquad\text{for every }j
\]
and \(\Ho^j(M')=0\) for \(j\geq i\).
%%%%%%%%%%  scan idx 110  |  folio 97  %%%%%%%%%%
\SGAPrintPage{97}%
The cone formula \(M''{}^j=F'_i{}^{\,j+1}\oplus F^j\) shows, on the other hand, that \(M''\) is strictly \((i-1)\)-pseudo-coherent. Thus, by (2.1.1), \(\Ho^{i-1}(M'')=\Ho^{i-1}(M')\) is of finite type, whence our assertion. By iteration, and after refining \(S\), we deduce an \(n\)-quasi-isomorphism \(v_n:F'_n\longrightarrow E\), where \(F'_n\) is strictly \(n\)-pseudo-coherent. But, by (1.4.5), we can extend \(v_n\) to a quasi-isomorphism \(v:F'\longrightarrow E\) without changing anything in degree \(\geq n\), hence in particular with \(F'\) strictly \(n\)-pseudo-coherent, which proves a).

b) The equivalence of conditions (i) and \((\mathrm{i}')\) follows at once from a) and the definition of the isomorphisms of \(\Dpx^-(\underline{C}_X)\). For \(E\in\Ob\Dpx^-(\underline{C}_S)\) strictly \(n\)-pseudo-coherent, there exists an exact sequence
\[
  0\longrightarrow E'\longrightarrow E\longrightarrow E''\longrightarrow 0,
\]
where \(E'\) is strictly perfect, and \(E''{}^i=0\) for \(i\geq n\); consequently \(E'\to E\) is an \(n\)-quasi-isomorphism, giving the implication \((\mathrm{i})\Rightarrow(\mathrm{ii})\). The implication \((\mathrm{ii})\Rightarrow(\mathrm{ii}')\) is trivial, so it remains to prove \((\mathrm{ii}')\Rightarrow(\mathrm{i})\). Let, then, \(u:E\longrightarrow F\) be an \(n\)-isomorphism of \(\Dpx^-(\underline{C}_S)\), with \(E\) strictly perfect; \(u\) is therefore defined by a diagram of \(\Cpx^-(\underline{C}_S)\)
\[
\begin{tikzcd}
& E' \arrow[dl,"s"'] \arrow[dr,"u'"]\\
E && F ,
\end{tikzcd}
\]
where \(s\) is a quasi-isomorphism and \(u'\) an \(n\)-quasi-isomorphism. By a), there exists, after refining \(S\), a quasi-isomorphism \(t:E''\longrightarrow E'\), with \(E''\) strictly \(n\)-pseudo-coherent. We then apply (1.4.5) to extend \(u't:E''\to F\) to a quasi-isomorphism \(E'''\longrightarrow F\) with \(E'''\) strictly \(n\)-pseudo-coherent, which proves \((\mathrm{ii}')\Rightarrow(\mathrm{i})\) and completes the proof of (2.2).

%%%%%%%%%%  scan idx 111  |  folio 98  %%%%%%%%%%
\SGAPrintPage{98}%
\sgaStableTarget{sga6.e01.def.2-3}\textbf{Definition 2.3.} \uline{Under the hypotheses of (2.2), we shall say that} \(F\in\Ob\Dpx(\underline{C}_S)\) \uline{is} \(\underline{n}\)-\(\underline{C}_0\)-pseudo-coherent \uline{when} \(F\) \uline{satisfies the equivalent conditions b) of (2.2). We shall say that} \(F\) \uline{is} \(\underline{C}_0\)-pseudo-coherent \uline{when} \(F\) \uline{is} \(n\)\uline{-}\(\underline{C}_0\)\uline{-pseudo-coherent for every} \(n\)\uline{.}

When this will cause no confusion, we shall sometimes allow ourselves to say \(n\)-pseudo-coherent, pseudo-coherent, instead of \(n\)-\(\underline{C}_0\)-pseudo-coherent, \(\underline{C}_0\)-pseudo-coherent.

\sgaStableTarget{sga6.e01.rem.2-3-1}\textbf{Remark 2.3.1.} Let \(F\in\Ob\Dpx(\underline{C}_S)\) be a pseudo-coherent complex. The set of \(X\) over \(S\) such that there exists a quasi-isomorphism \(E\longrightarrow F_X\) with \(E\) strictly pseudo-coherent (2.1) is in general not a refinement of \(S\). However, we shall see later (2.7. a) and Exp.II) important cases where it is.

The proof of (2.2 a)) gives the following result:

\sgaStableTarget{sga6.e01.prop.2-4}\textbf{Proposition 2.4.} \uline{Let} \(n,i\in\mathbb Z\)\uline{, and let} \(u:E\to F\) \uline{be an} \(i\)\uline{-quasi-isomorphism of} \(\Cpx(\underline{C}_S)\)\uline{, with} \(F\) \(n\)\uline{-pseudo-coherent and} \(E\) \uline{strictly} \(n\)\uline{-pseudo-coherent. Then, after refining} \(S\)\uline{, there exists a commutative triangle of} \(\Cpx(\underline{C}_S)\)
\[
\begin{tikzcd}
E \arrow[r,"v"] \arrow[d,"u"'] & E' \arrow[dl,"u'"]\\
F & ,
\end{tikzcd}
\]
\uline{where} \(E'\) \uline{is strictly} \(n\)\uline{-pseudo-coherent,} \(u'\) \uline{a quasi-isomorphism and} \(v\) \uline{a monomorphism inducing the identity in degree} \(\geq i\)\uline{.}

\sgaStableTarget{sga6.e01.prop.2-5}\textbf{Proposition 2.5.} a) \uline{When} \(F\in\Ob\Dpx(\underline{C}_S)\) \uline{is} \(n\)\uline{-pseudo-coherent (resp. pseudo-coherent),} \(F[m]\) \uline{is} \((n-m)\)\uline{-pseudo-coherent (resp. pseudo-coherent).}

%%%%%%%%%%  scan idx 112  |  folio 99  %%%%%%%%%%
\SGAPrintPage{99}%
b) \uline{Let}
\sgaStableTarget{sga6.e01.eq.unnamed.occ2}%
\[
\begin{tikzcd}
& G \arrow[dl,"+1"'] & \\
E \arrow[rr] && F \arrow[ul]
\end{tikzcd}
\tag{*}
\]
\uline{be a distinguished triangle of} \(\Dpx(\underline{C}_S)\)\uline{. When} \(E\) \uline{and} \(G\) \uline{are} \(n\)\uline{-pseudo-coherent (resp. pseudo-coherent), then} \(F\) \uline{is} \(n\)\uline{-pseudo-coherent (resp. pseudo-coherent).}

\uline{Proof.} a) Trivial from the definitions.

b) It of course suffices to prove the assertion without the ``resp.'' alternatives. The triangle \((*)\) can also be written
\[
\begin{tikzcd}
& F \arrow[dl,"+1"'] & \\
G[-1] \arrow[rr] && E \arrow[ul] & .
\end{tikzcd}
\]
By a), \(G[-1]\) is \((n+1)\)-pseudo-coherent. We are therefore reduced to proving that, when \(E\) is \((n+1)\)-pseudo-coherent and \(F\) \(n\)-pseudo-coherent, then \(G\) is \(n\)-pseudo-coherent. After localizing, we can assume that \(E\) (resp.\ \(F\)) is strictly \((n+1)\)-pseudo-coherent (resp.\ \(n\)-pseudo-coherent). The arrow \(E\xrightarrow{u}F\) of \(\Dpx^-(\underline{C}_S)\) is defined by a pair of arrows of \(\Cpx^-(\underline{C}_S)\)
\[
\begin{tikzcd}
& E' \arrow[dl,"s"'] \arrow[dr,"u'"]\\
E && F ,
\end{tikzcd}
\]
where \(s\) is a quasi-isomorphism. By (2.2 a)), there exists, after localizing, a quasi-isomorphism \(t:E''\to E'\), where \(E''\) is strictly \((n+1)\)-pseudo-coherent. We are therefore reduced to the case where \(u\) is defined by an
%%%%%%%%%%  scan idx 113  |  folio 100  %%%%%%%%%%
\SGAPrintPage{100}%
``actual'' arrow of \(\Cpx^-(\underline{C}_S)\). The assertion then follows from the definitions and the formula \(C(u)^i=E^{i+1}\oplus F^i\), where \(C(u)\) denotes the cone of \(u\).

\sgaStableTarget{sga6.e01.schol.2-6}\textbf{Scholium 2.6.} By ``rotation'' of the triangles, we see that pseudo-coherence properties of two of the objects of the triangle \((*)\) imply another such property of the third. The different cases can be summarized in the table below (the lowercase letters denote degrees of pseudo-coherence and each row corresponds to an implication whose conclusion is underlined):
\[
\begin{array}{ccc}
E&F&G\\[0.2em]
n+1&n&\underline n\\
n&\underline n&n\\
\underline n&n&n-1 .
\end{array}
\]

We shall denote by \(\Dpx(\underline{C})_{\underline{C}_0\text{-}n\text{-coh}}\) (resp.\ \(\Dpx(\underline{C})_{\underline{C}_0\text{-}\mathrm{coh}}\)) the strictly full fibered subcategory of \(\Dpx(\underline{C})\) whose fibre at each \(X\in\Ob \underline{S}\) consists of the objects of \(\Dpx(\underline{C}_X)\) that are \(n\)-pseudo-coherent (resp.\ pseudo-coherent) relative to \(\underline{C}_0\). By (2.5), the category \(\Dpx(\underline{C})_{\underline{C}_0\text{-}\mathrm{coh}}\) is a \uline{triangulated} sub-\(\underline{S}\)\uline{-category} of \(\Dpx(\underline{C})\) (1.1).

\sgaStableTarget{sga6.e01.prop.2-7}\var{\textbf{Proposition 2.7}}{\textbf{Proposition 2.7.}}{EN-VAR-084} \uline{Assume that} \(\underline{S}\) \uline{is a punctual site, so that} \(\underline{C}=\underline{C}_S\)\uline{,} \(\underline{C}_0=\underline{C}_{0S}\)\uline{. Then}:

a) \uline{When} \(F\in\Ob\Dpx(\underline{C})\) \uline{is pseudo-coherent, there exists a quasi-isomorphism} \(E\longrightarrow F\)\uline{, with} \(E\) \uline{strictly pseudo-coherent} (2.1).

b) \uline{Denote by} \(\Kpx^{-,\varnothing}(\underline{C}_0)\) \uline{the full triangulated subcategory of} \(\Kpx^-(\underline{C}_0)\) \uline{consisting of the acyclic complexes. The canonical functor}
\[
  \Kpx^-(\underline{C}_0)/\Kpx^{-,\varnothing}(\underline{C}_0)
  \longrightarrow \Dpx(\underline{C})
\]
%%%%%%%%%%  scan idx 114  |  folio 101  %%%%%%%%%%
\SGAPrintPage{101}%
\uline{is fully faithful and has essential image} \(\Dpx(\underline{C})_{\underline{C}_0\text{-}\mathrm{coh}}\).

\uline{Proof.} a) Let \(F\in\Ob\Dpx(\underline{C})_{\underline{C}_0\text{-}\mathrm{coh}}\), and suppose that an \(i\)-quasi-isomorphism (of \(\Cpx(\underline{C})\)) \(u_i:E_i\longrightarrow F\) has been constructed with \(E_i\) strictly pseudo-coherent (this is possible for \(i\) sufficiently large with \(E_i=0\)). Let \(G_i\) be the cone of \(u_i\). It follows from (1.4.2) and (2.1.1) that \(\Ho^{i-1}(G_i)\) is of \(\underline{C}_0\)-finite type. Thus, by (1.4.4), there exists a commutative triangle of \(\Cpx(\underline{C})\)
\[
\begin{tikzcd}
E_i \arrow[r,"v_{i-1}"] \arrow[d,"u_i"'] & E_{i-1} \arrow[dl,"u_{i-1}"]\\
F & ,
\end{tikzcd}
\]
where \(E_{i-1}\) is strictly pseudo-coherent, \(u_{i-1}\) an \((i-1)\)-quasi-isomorphism, and \(v_{i-1}\) a monomorphism inducing the identity in degree \(\geq i\). By passage to the limit (cf.\ proof of (1.4.5)), we deduce a quasi-isomorphism \(u:E\to F\) with \(E\) strictly pseudo-coherent.

b) Apply the criterion ([V] chap I \S\ 2 4.2 b)(i)) and a).

Observe that the proof of a) gives the following more precise result, which generalizes (1.4.5):

\sgaStableTarget{sga6.e01.prop.2-7-1}\textbf{Proposition 2.7.1.} \uline{Under the hypotheses of (2.7), let} \(u:E\to F\) \uline{be an} \(n\)\uline{-quasi-isomorphism of} \(\Cpx(\underline{C})\)\uline{, with} \(F\) \uline{pseudo-coherent and} \(E\) \uline{strictly pseudo-coherent. There then exists a commutative triangle of} \(\Cpx(\underline{C})\)
\[
\begin{tikzcd}
E \arrow[r,"v"] \arrow[d,"u"'] & E' \arrow[dl,"u'"]\\
F & ,
\end{tikzcd}
\]
\uline{where} \(E'\) \uline{is strictly pseudo-coherent,} \(u'\) \uline{a quasi-isomorphism, and} \(v\) \uline{a}
%%%%%%%%%%  scan idx 115  |  folio 102  %%%%%%%%%%
\SGAPrintPage{102}%
\uline{monomorphism inducing the identity in degree} \(\geq n\).

\sgaStableTarget{sga6.e01.def.2-8}\textbf{Definition 2.8.} \uline{Let} \(n\in\mathbb N\) \uline{and} \(F\in\Ob\underline{C}_S\)\uline{. We shall say that} \(F\) \uline{is of} \(\underline{n}\)-\(\underline{C}_0\)-finite presentation \uline{when the set of} \(X\) \uline{over} \(S\) \uline{such that there exists an exact sequence}
\[
  E^{-n}\longrightarrow\ldots\longrightarrow E^0\longrightarrow F_X\longrightarrow 0,
\]
\uline{with} \(E^i\in\Ob\underline{C}_{0X}\) \uline{for} \(-n\leq i\leq 0\)\uline{, is a refinement of} \(S\)\uline{. We shall say that} \(F\) \uline{is of} \(\infty\)-\(\underline{C}_0\)-finite presentation \uline{when} \(F\) \uline{is of} \(n\)\uline{-}\(\underline{C}_0\)\uline{-finite presentation for every} \(n\in\mathbb N\).

Of course, we shall sometimes simply say \(n\)-finite presentation when no ambiguity concerning \(\underline{C}_0\) is possible.

We shall say ``of finite presentation'' for ``of 1-finite presentation'', and ``of finite type'' for ``of 0-finite presentation'', the latter notion coinciding with the one already introduced in (1.2).

\sgaStableTarget{sga6.e01.prop.2-9}\textbf{Proposition 2.9.} \uline{Let} \(n\in\mathbb N\) \uline{and} \(F\in\Ob\underline{C}_S\)\uline{. Equivalent conditions}:

(i) \(F\) \uline{is of} \(n\)\uline{-finite presentation (resp.\ of} \(\infty\)\uline{-finite presentation)};

(ii) \(F\) \uline{is} \((-n)\)\uline{-pseudo-coherent (resp.\ pseudo-coherent)}.

\uline{In the case where} \(\underline{S}\) \uline{is a punctual site,} \(F\) \uline{is pseudo-coherent if and only if there exists an exact sequence}
\[
  \ldots\longrightarrow E^i\longrightarrow\ldots\longrightarrow E^0\longrightarrow F\longrightarrow 0,
\]
\uline{with} \(E^i\in\Ob\underline{C}_0\) \uline{for every} \(i\).

\uline{Proof.} For the equivalence \((\mathrm{i})\Leftrightarrow(\mathrm{ii})\), it clearly suffices to treat the case without the ``resp.'' alternatives. The implication \((\mathrm{i})\Rightarrow(\mathrm{ii})\) follows immediately from the definitions, and
%%%%%%%%%%  scan idx 116  |  folio 103  %%%%%%%%%%
\SGAPrintPage{103}%
the implication \((\mathrm{ii})\Rightarrow(\mathrm{i})\) follows at once from (2.4): the arrow \(0\to F\) is a 1-quasi-isomorphism; we deduce (after localizing) a quasi-isomorphism \(E\longrightarrow F\) with \(E\) strictly \(\var{n}{(-n)}{EN-VAR-085}\)-pseudo-coherent and \(E^i=0\) for \(i\geq 1\), i.e.\ an \(n\)-finite presentation of \(F\). The second assertion is proved in a similar way (use (2.7.1)).

\sgaStableTarget{sga6.e01.prop.2-10}\textbf{Proposition 2.10.} a) \uline{Let} \(m,n\in\mathbb Z\)\uline{, with} \(n\leq m\)\uline{. When} \(F\in\Ob\Dpx^-(\underline{C}_S)\) \uline{is} \(n\)\uline{-pseudo-coherent and acyclic in degree} \(\geq m\)\uline{, there locally exists a quasi-isomorphism} \(E\longrightarrow F\)\uline{, with} \(E\) \uline{strictly} \(n\)\uline{-pseudo-coherent and} \(E^i=0\) \uline{for} \(i\geq m\)\uline{. When} \(\underline{S}\) \uline{is punctual, and} \(F\) \uline{pseudo-coherent and acyclic in degree} \(\geq m\)\uline{, there exists a quasi-isomorphism} \(E\longrightarrow F\) \uline{with} \(E\) \uline{strictly pseudo-coherent and} \(E^i=0\) \uline{for} \(i\geq m\)\uline{.}

b) \uline{Let} \(F\in\Ob\Dpx^-(\underline{C}_S)\) \uline{be acyclic in degree} \(>n\)\uline{. For} \(F\) \uline{to be} \(n\)\uline{-pseudo-coherent, it is necessary and sufficient that} \(\Ho^n(F)\) \uline{be of finite type.}

\uline{Proof.} a) Since \(F\) is acyclic in degree \(\geq m\), the arrow \(0\to F\) is an \(m\)-quasi-isomorphism. The assertion follows from (2.4) and (2.7.1).

b) There exists a distinguished triangle of \(\Dpx^-(\underline{C}_S)\):
\[
\begin{tikzcd}
& F' \arrow[dl,"+1"'] & \\
F \arrow[rr] && \Ho^n(F)[-n] \arrow[ul] & ,
\end{tikzcd}
\]
(to see this, begin by truncating \(F\), replacing \(F^n\) by \(Z^n(F)\) and \(F^i\) by \(0\) for \(i>n\)). The exact cohomology sequence implies \(\Ho^i(F')=0\) for \(i\geq n-1\), so \(F'\) is \((n-1)\)-pseudo-coherent. We conclude using (2.5) and (2.6).

% Additional notation for Expose I continuation
\providecommand{\RHom}{\mathrm R\!\Hom}
\providecommand{\Ext}{\operatorname{Ext}}
\providecommand{\coh}{\operatorname{coh}}
\providecommand{\PreCoh}{\operatorname{precoh}}
\providecommand{\Tr}{\operatorname{Tr}}
\providecommand{\Qis}{\operatorname{Qis}}
\providecommand{\Rad}{\operatorname{rad}}
\providecommand{\Spec}{\operatorname{Spec}}
\providecommand{\End}{\operatorname{End}}
\providecommand{\Id}{\operatorname{Id}}
\providecommand{\varinjlim}{\mathop{\mathrm{lim}}\limits_{\longrightarrow}}
\providecommand{\varprojlim}{\mathop{\mathrm{lim}}\limits_{\longleftarrow}}
\providecommand{\qcoh}{\operatorname{qcoh}}
\providecommand{\amp}{\operatorname{amp}}

%%%%%%%%%%  scan idx 117  |  folio 104  %%%%%%%%%%
\SGAPrintPage{104}%
\sgaStableTarget{sga6.e01.expose.occ2}%
\sgaStableTarget{sga6.e01.prop.2-11}\textbf{Proposition 2.11.} \uline{Let} \(n\in\mathbb Z\)\uline{.}

a) \uline{Let} \(F\in\Ob\Cpx^-(\underline{C}_S)\)\uline{. When, for every} \(i\geq n\) \uline{(resp.\ for every} \(i\)\uline{),} \(F^i\) \uline{is of} \((i-n)\)\uline{-finite presentation (resp.\ of} \(\infty\)\uline{-finite presentation), then} \(F\) \uline{is} \(n\)\uline{-pseudo-coherent (resp.\ pseudo-coherent).}

b) \uline{Let} \(F\in\Ob\Dpx^-(\underline{C}_S)\)\uline{. When, for every} \(i\geq n\) \uline{(resp.\ for every} \(i\)\uline{),} \(\Ho^i(F)\) \uline{is of} \((i-n)\)\uline{-finite presentation (resp.\ of} \(\infty\)\uline{-finite presentation), then} \var{\(F\) \uline{is pseudo-coherent.}}{\(F\) \uline{is} \(n\)\uline{-pseudo-coherent (resp.\ pseudo-coherent).}}{EN-VAR-086}

\uline{Proof.} It suffices to treat the assertions without the ``resp.'' alternatives.

a) We argue by induction on \(N\) such that \(\var{F^i\neq0}{F^i=0}{EN-VAR-087}\) for \(i>n+N\). For \(N=0\), the assertion follows at once from (2.10 b)). Suppose \(N\geq1\). There is an exact sequence
\[
  0\longrightarrow F^{n+N}[-n-N]\longrightarrow F\longrightarrow F'\longrightarrow 0,
\]
By the induction hypothesis, \(F'\) is \(n\)-pseudo-coherent; on the other hand, \(F^{n+N}[-n-N]\) is \(n\)-pseudo-coherent by hypothesis (2.9 and 2.5.a)); thus, by (2.5.b)), \(F\) is \(n\)-pseudo-coherent.

b) We argue by induction on \(N\) such that \(\Ho^i(F)=0\) for \(i>n+N\). For \(N=0\), the assertion follows from (2.10 b)). Suppose \(N\geq1\). There is a distinguished triangle of \(\Dpx^-(\underline{C}_S)\) (cf.\ proof of (2.10.b))):
\[
\begin{tikzcd}
& F' \arrow[dl,"+1"'] & \\
F \arrow[rr] && \Ho^{n+N}(F)[-n-N] \arrow[ul] & .
\end{tikzcd}
\]
The exact cohomology sequence implies
\[
  \Ho^i(F')=0\quad\text{for }i\geq n+N-1,\quad\text{i.e. }i>(n-1)+(N-1),
\]
and
\[
  \Ho^i(F')\qis\Ho^{i+1}(F)\quad\text{for }i<n+N-1 .
\]
%%%%%%%%%%  scan idx 118  |  folio 105  %%%%%%%%%%
\SGAPrintPage{105}%
By the induction hypothesis (applied with \(n-1\) in place of \(n\)), \(F'\) is therefore \((n-1)\)-pseudo-coherent. On the other hand, \(\Ho^{n+N}(F)[-n-N]\) is \(n\)-pseudo-coherent (2.9 and 2.5.a)). Thus, by (2.6), \(F\) is \(n\)-pseudo-coherent.

\sgaStableTarget{sga6.e01.rem.2-11-1}\textbf{Remark 2.11.1.} When \(F\) is pseudo-coherent, the \(F^i\) (resp.\ \(\Ho^i(F)\)) are in general not pseudo-coherent, nor even of finite type. However, in the next paragraph, in the sorites on notions of coherence, we shall see important examples where the pseudo-coherence of \(F\) implies that of the \(\Ho^i(F)\).

\sgaStableTarget{sga6.e01.prop.2-12}\textbf{Proposition 2.12.} \uline{Let} \(F',F''\in\Ob\Dpx(\underline{C}_S)\)\uline{, and} \(F=F'+F''\)\uline{. For} \(F\) \uline{to be} \(n\)\uline{-pseudo-coherent (resp.\ pseudo-coherent), it is necessary and sufficient that} \(F'\) \uline{and} \(F''\) \uline{be} \(n\)\uline{-pseudo-coherent (resp.\ pseudo-coherent).}

\uline{Proof.} It suffices to prove the assertion without the ``resp.'' alternatives. Sufficiency is clear; let us prove necessity. By induction, we may suppose we have already proved that \(F'\) and \(F''\) are \((n+1)\)-pseudo-coherent. After localizing, we can even suppose that \(F'\) and \(F''\) are strictly \((n+1)\)-pseudo-coherent and that \(F=F'+F''\) in \(\Cpx^-(\underline{C}_S)\). Denote by \(F'_{>n}\) (resp.\ \(F'_{\leq n}\)) the complex obtained from \(F'\) by replacing \(F'^i\) by \(0\) for \(i\leq n\) (resp.\ \(i>n\)), and similarly \(F''_{>n}\) (resp.\ \(F''_{\leq n}\))\dots{} There is an exact sequence
\[
  0\longrightarrow F'_{>n}+F''_{>n}\longrightarrow F
  \longrightarrow F'_{\leq n}+F''_{\leq n}\longrightarrow 0 .
\]
\(F'_{>n}+F''_{>n}\) is strictly perfect and \(F\) is \(n\)-pseudo-coherent, so (2.6) \(F'_{\leq n}+F''_{\leq n}\) is \(n\)-pseudo-coherent. By (2.10 b)), this means that \(\Ho^n(F'_{\leq n}+F''_{\leq n})\) is of finite type. But we have
\[
  \Ho^n(F'_{\leq n}+F''_{\leq n})
  =\Ho^n(F'_{\leq n})+\Ho^n(F''_{\leq n}),
\]
%%%%%%%%%%  scan idx 119  |  folio 106  %%%%%%%%%%
\SGAPrintPage{106}%
so \(\Ho^n(F'_{\leq n})\) (resp.\ \(\Ho^n(F''_{\leq n})\)) is of finite type, so again, by (2.10.b)), \var{\(F'_{\leq n}\)) (resp.}{\(F'_{\leq n}\) (resp.}{EN-VAR-088}\ \(F''_{\leq n}\)) is \(n\)-pseudo-coherent. Applying (2.5) to the exact sequence
\[
  0\longrightarrow F'_{>n}\longrightarrow F'\longrightarrow F'_{\leq n}\longrightarrow 0,
\]
(resp.\dots), we deduce that \(F'\) (resp.\ \(F''\)) is \(n\)-pseudo-coherent, QED.

\sgaStableTarget{sga6.e01.prop.2-13}\textbf{Proposition 2.13.} \uline{Let} \((\underline{S}',\underline{C}',\underline{C}'_0)\) \uline{be a triple satisfying the same hypotheses} (2.0) \uline{as} \((\underline{S},\underline{C},\underline{C}_0)\)\uline{. Let} \(f:\underline{S}'\to \underline{S}\) \uline{be a morphism of sites} ([G] I 1.4)\uline{, corresponding to an underlying functor} \(f^{-1}:\underline{S}\to \underline{S}'\)\uline{, and let}
\(u:\Dpx^*(\underline{C})\longrightarrow \Dpx^*(\underline{C}')\) \((*=+,-,\varnothing)\)
\uline{be a cartesian functor over} \(f\)\uline{,} i.e.\ \uline{a cartesian} \(\underline{S}\)\uline{-functor} \(\Dpx^*(\underline{C}')\times_{\underline{S}'}\underline{S}\longleftarrow\Dpx^*(\underline{C})\)\uline{. Assume that} \(u\) \uline{is exact on each fibre and that there exist} \(a,b\in\mathbb Z\) \uline{such that the following conditions hold for every} \(S\in\Ob \underline{S}\):

(i) \uline{when} \(E\in\Ob\underline{C}_{0S}\)\uline{,} \(u(E)\) \uline{is} \(a\)\uline{-}\(\underline{C}'_0\)\uline{-pseudo-coherent (resp.} \(\underline{C}'_0\)\uline{-pseudo-coherent)}

(ii) \uline{when} \(F\in\Ob\Dpx^*(\underline{C}_S)\) \uline{is acyclic in degree} \(>0\)\uline{,} \(u(F)\) \uline{is acyclic in degree} \(>b\)\uline{.}

\uline{Then, when} \(F\in\Ob\Dpx^*(\underline{C}_S)\) \uline{is} \(m\)\uline{-pseudo-coherent rel.\ to} \(\underline{C}_0\) \uline{and acyclic in degree} \(>n\)\uline{,} \(u(F)\in\Ob\Dpx^*(\underline{C}'_{S'})\) \uline{(where} \(S'=f^{-1}(S)\)\uline{) is} \(\sup(m+b,n+a)\)\uline{-pseudo-coherent (rel.\ to} \(\underline{C}'_0\)\uline{) (resp.} \((m+b)\)\uline{-pseudo-coherent) and acyclic in degree} \(>n+b\)\uline{.}

\uline{In particular, in the case of the ``resp.'' alternative, when} \(F\in\Ob\Dpx^*(\underline{C}_S)\) \uline{is pseudo-coherent} (rel.\ to \(\underline{C}_0\)), \(u(F)\) \uline{is pseudo-coherent (rel.\ to} \(\underline{C}'_0\)\uline{).}

\uline{Proof.} We can restrict ourselves to the case without the ``resp.'' alternative. Let us first show that, for \(F\in\Ob\Dpx^*(\underline{C}_S)\) strictly perfect (2.1) with \(F^i=0\) for \(i>n\), \(u(F)\) is \((n+a)\)-\(\underline{C}'_0\)-pseudo-coherent. We argue by induction on the number \(N\) of indices \(i\) such that \(F^i\neq0\). For \(N=1\),
%%%%%%%%%%  scan idx 120  |  folio 107  %%%%%%%%%%
\SGAPrintPage{107}%
we have \(F=F^p[-p]\) for some \(p\leq n\), so \(u(F)\) is \((p+a)\)-\(\underline{C}'_0\)-pseudo-coherent, hence a fortiori \((n+a)\)-pseudo-coherent. For \(N>1\), there is an exact sequence
\[
  0\longrightarrow F'\longrightarrow F\longrightarrow F''\longrightarrow 0,
\]
where \(F'\) and \(F''\) are strictly perfect of length \(\leq N-1\) and satisfy \(F'^i=F''^i=0\) for \(i>n\). Applying (2.5.b)) to the distinguished triangle
\[
\begin{tikzcd}
& u(F'') \arrow[dl,"+1"'] & \\
u(F') \arrow[rr] && u(F) \arrow[ul] & ,
\end{tikzcd}
\]
we deduce from the induction hypothesis that \(u(F)\) is \((n+a)\)-pseudo-coherent.

Let us now pass to the general case: \(F\) \(m\)-\(\underline{C}_0\)-pseudo-coherent and acyclic in degree \(>n\). After localizing, we can, by (2.10.a)), suppose that \(F\) is strictly \(m\)-pseudo-coherent and satisfies \(F^i=0\) for \(i>n\). There is then an exact sequence
\[
  0\longrightarrow F'\longrightarrow F\longrightarrow F''\longrightarrow 0,
\]
where \(F''\) is zero in degree \(\geq m\), and \(F'\) strictly perfect and zero in degree \(>n\). In the distinguished triangle
\[
\begin{tikzcd}
& u(F'') \arrow[dl,"+1"'] & \\
u(F') \arrow[rr] && u(F) \arrow[ul] & ,
\end{tikzcd}
\]
\(u(F'')\) is acyclic in degree \(\geq m+b\), hence \((m+b)\)-pseudo-coherent, and \(u(F')\) is \((n+a)\)-pseudo-coherent by what we have seen. We conclude by (2.5.b)).

\sgaStableTarget{sga6.e01.cor.2-14}\textbf{Corollary 2.14.} \uline{Let} \(\underline{C}_1\) \uline{be a strictly full additive sub-}\(\underline{S}\)\uline{-category of} \(\underline{C}\) \uline{such that} \(\underline{C}_0\subset\underline{C}_1\subset\underline{C}\)\uline{. Assume that every object of} \(\underline{C}_1\) \uline{is} \(\underline{C}_0\)\uline{-pseudo-coherent.}

%%%%%%%%%%  scan idx 121  |  folio 108  %%%%%%%%%%
\SGAPrintPage{108}%
\uline{Then}:

(i) \(\underline{C}_1\) \uline{is locally quasi-liftable in} \(\underline{C}\) (1.2) \uline{and locally quasi-stable under kernels of epimorphisms} (1.2)\uline{.}

(ii) \uline{Let} \(S\in\Ob \underline{S}\) \uline{and} \(F\in\Ob\Dpx(\underline{C}_S)\)\uline{. For} \(F\) \uline{to be} \(n\)\uline{-pseudo-coherent (resp.\ pseudo-coherent) relative to} \(\underline{C}_1\)\uline{, it is necessary and sufficient that it be so relative to} \(\underline{C}_0\)\uline{.}

\uline{Proof.} (i) Let \(E\to F\) be an epimorphism of \(\underline{C}_S\), with \(F\in\Ob\underline{C}_{1S}\). Since \(F\) is \(\underline{C}_0\)-pseudo-coherent (hence a fortiori of \(\underline{C}_0\)-finite type), there exists (after localizing) an epimorphism \(E'\to F\) with \(E'\in\Ob\underline{C}_{0S}\). Using the quasi-liftability of \(\underline{C}_0\) in \(\underline{C}\), we can suppose (after localizing again) that this factors through \(E\), whence the quasi-liftability of \(\underline{C}_1\). On the other hand, when \(E\in\Ob\underline{C}_S\) admits a finite right resolution by objects of \(\underline{C}_1\), \(E\) is \(\underline{C}_0\)-pseudo-coherent by (2.11.a)), hence a fortiori of \(\underline{C}_1\)-finite type.

(ii) It suffices to prove necessity. For this we can apply, as we choose, (2.11.a)) or (2.13) with \(\underline{S}'=\underline{S}\), \(\underline{C}'=\underline{C}\), \(f=\text{Identity}\), \(u=\text{Identity}\), replacing \(\underline{C}'_0\) and \(\underline{C}_0\) respectively by \(\underline{C}_0\) \var{and \(\underline{C}_1\)).}{and \(\underline{C}_1\).}{EN-VAR-089}

\sgaStableTarget{sga6.e01.schol.2-14-1}\textbf{Scholium 2.14.1.} There thus exists a largest strictly full additive sub-\(\underline{S}\)-category of \(\underline{C}\) giving the same notion of \(n\)-pseudo-coherence (resp.\ pseudo-coherence) as \(\underline{C}_0\), namely the category formed by the \(\underline{C}_0\)-pseudo-coherent objects. The latter is locally stable under kernels of epimorphisms.

\sgaStableTarget{sga6.e01.ex.2-15}\textbf{Example 2.15.} Let \((\underline{S},A)\) be a ringed topos and \(\underline{C},\underline{C}_0\) the \(\underline{S}\)-categories considered in (1.3.d)). The fibered category \(\Dpx(\underline{C})\) will often be denoted by \(\Dpx(\underline{S},A)\), or \(\Dpx({}_A \underline{S})\) (to
%%%%%%%%%%  scan idx 122  |  folio 109  %%%%%%%%%%
\SGAPrintPage{109}%
suggest that these are left \(A\)-Modules), or even \(\Dpx(\underline{S})\) when there is no ambiguity concerning \(A\). We shall denote by \(S\) the final object of \(\underline{S}\) and put \(\Dpx(\underline{C}_S)=\Dpx({}_A S)\) (or \(\Dpx(S)\)). An object of \(\Dpx(\underline{S})\) will be called \(\underline{n}\)\uline{-pseudo-coherent} (resp.\ \uline{pseudo-coherent}) when it is \(n\)-\(\underline{C}_0\)-pseudo-coherent (resp.\ \(\underline{C}_0\)-pseudo-coherent) (likewise for the strict notions). The fibered category \(\Dpx(\underline{C})_{n\text{-}\underline{C}_0\text{-}\coh}\) (resp.\ \(\Dpx(\underline{C})_{\underline{C}_0\text{-}\coh}\)) (2.6) will be denoted by \(\Dpx({}_A \underline{S})_{n\text{-}\coh}\) (resp.\ \(\Dpx({}_A \underline{S})_{\coh}\)). Observe that, by (2.14), we obtain the same notion of pseudo-coherence on replacing \(\underline{C}_0\) by the \(\underline{S}\)-category \(\underline{C}_1\) defined by \(U\longmapsto\) (category of left \(A_U\)-Modules locally direct factors of free Modules of finite type). When an \(A\)-Module \(M\) is a direct factor of a free \(A\)-Module of finite type, it does not follow, in general, that \(M\) is locally free of finite type. However, we have:

\sgaStableTarget{sga6.e01.lem.2-15-1}\textbf{Lemma 2.15.1.} \uline{Assume that} \(\underline{S}\) \uline{satisfies one of the following two hypotheses:}

(i) \(\underline{S}\) \uline{``has enough points''} (SGA 4 IV 6.4) \uline{and, for every point} \(s\) \uline{of} \(\underline{S}\)\uline{, the projective} \(A_s\)\uline{-modules} (on the left) \uline{of finite type are free, this last condition being satisfied for example} (BOURBAKI Alg. Comm. chap.\ II \S~3 n\textsuperscript{o}~2 cor.~2 of prop.\ 5) \uline{when} \(A_s/\underline{\operatorname{radical}}(A_s)\) \var{\uline{is a field or when} \(\underline{\operatorname{radical}}(A_s)\) \uline{is nilpotent.}}{\uline{is a field.}}{EN-VAR-090}

(ii) \((\underline{S},A)\) \uline{is a ``locally ringed'' topos in the sense of} [1]\uline{, which means that} \(A\) \uline{is commutative and that, for every} \(X\in\Ob \underline{S}\) \uline{and every} \(f\in\Gamma(X,A)\)\uline{, we have} \(X=X_f\cup X_{1-f}\)\uline{, where} \(X_f\) \uline{(resp.} \(X_{1-f}\)\uline{) denotes the largest subobject of} \(X\) \uline{on which} \(f\) \uline{(resp.} \(1-f\)\uline{) is invertible.}

\uline{Then every direct factor of a free} \(A\)\uline{-module (on the left) of finite type is locally free of finite type.}

\uline{Proof.} Suppose (i). We first observe that, for two \(A\)-Modules \(E\) and \(F\), with \(E\) of finite presentation, the canonical arrow
%%%%%%%%%%  scan idx 123  |  folio 110  %%%%%%%%%%
\SGAPrintPage{110}%
\[
  \underline{\Hom}(E,F)_s\longrightarrow\Hom(E_s,F_s)
\]
is an isomorphism for every point \(s\) of \(\underline{S}\). It follows that, when \(E\) and \(F\) are both of finite presentation, and \(E_s\simeq F_s\) at a point \(s\) of \(\underline{S}\), there exists an \(s\)-pointed object \(U\) of \(\underline{S}\)\SGAAuthorFootnote{*}{i.e.\ such that \(s\) factors through \(\underline{S}_{/U}\).} such that \(E|U\simeq F|U\). The conclusion follows at once.

\var{Suppose (ii), Let}{Suppose (ii). Let}{EN-VAR-091} \(n\in\mathbb N\), and denote by \(\mathbb D^n\) the functor on the category of locally ringed topoi (belonging to a fixed universe\dots) defined by \(T\longmapsto\Gamma(T,\underline{O}_T)^n\). We know [1] that the functor \(\underline{\End}_{\mathbb D}(\mathbb D^n)\) is represented by \(M_n=\Spec\mathbb Z[T_{ij}]\ (1\leq i\leq n,\ 1\leq j\leq n)\), equipped with the endomorphism \var{\(U_n\)}{\(u_n\)}{EN-VAR-092} of \((\underline{O}_{M_n})^n\) defined by the matrix \((T_{ij})\). Now let \(p:A^n\longrightarrow A^n\) be a projector; let us show that \(p(A^n)\) is locally free of finite type. By what we have just said, there exists a morphism \(s:\underline{S}\longrightarrow M_n\) of locally ringed topoi such that \(s^*(u_n)=p\). We can interpret \(s\) as a point of \(M_n\) with values in \(\Gamma(S,A)\), i.e.\ a matrix of order \(n\) with coefficients in \(\Gamma(S,A)\); since this matrix is a projector, \(s\) factors through the closed subscheme \(i:P_n\to M_n\) corresponding to projectors as \(\underline{S}\xrightarrow{s'}P_n\xrightarrow{i}M_n\). Thus \(p=s'^*i^*(u_n)\), and we are done, because by the assertion under hypothesis (i), the image of \(i^*(u_n)\) is a locally free sheaf of finite type.

\sgaStableTarget{sga6.e01.cor.2-16}\textbf{Corollary 2.16.} \uline{Let} \(f:(\underline{S}',A')\to(\underline{S},A)\) \uline{be a morphism of ringed topoi (defined by a functor} \(f^{-1}:\underline{S}\to \underline{S}'\) \uline{and a morphism of sheaves of rings} \(f^{-1}(A)\to A'\)\uline{). Denote by} \(\Dpx({}_{A'}\underline{S}'_A)\) \uline{the derived category of the category of} \(A'\)\uline{-}\(A\)\uline{-biModules (on the left over} \(A'\)\uline{, on the right over} \(A\)\uline{), and by}
\[
  \otimes_A^{\mathrm L}:\Dpx^-(_{A'}\underline{S}'_A)\times_{\underline{S}'}\Dpx^-(_A \underline{S})
  \longrightarrow \Dpx^-(_{A'}\underline{S}')
\]
%%%%%%%%%%  scan idx 124  |  folio 111  %%%%%%%%%%
\SGAPrintPage{111}%
\uline{the derived functor of the tensor product} \((E',E)\longmapsto E'\otimes_{f^{-1}(A)} f^{-1}(E)\)\uline{.} \uline{Let} \(X\in\Ob \underline{S}\)\uline{,} \(\var{X'=u(X)}{X'=f^{-1}(X)}{EN-VAR-093}\)\uline{,} \(E'\in\Ob\Dpx^-(_{A'}X'_A)\)\uline{,} \(E\in\Ob\Dpx^-(_A X)\)\uline{. Then, when} \(E'\) \uline{is} \(m'\)\uline{-pseudo-coherent (as an object of} \(\Dpx^-(_{A'}\underline{S}')\)\uline{) and acyclic in degree} \(>n'\) \uline{and} \(E\) \uline{is} \(m\)\uline{-pseudo-coherent and acyclic in degree} \(>n\)\uline{,} \(E'\otimes_A^{\mathrm L}E\) \uline{is} \(\sup(m+n',m'+n)\)\uline{-pseudo-coherent and acyclic in degree} \(>n+n'\)\uline{.}

\uline{Proof.} Replacing \(\underline{S}\) by \(\underline{S}_{/X}\), we can suppose that \(X=S\), and we apply (2.13) with \(u=E'\otimes_A^{\mathrm L}:\Dpx^-(_A \underline{S})\to\Dpx^-(_{A'}\underline{S}')\). Let \(U\in\Ob \underline{S}\), \(U'=f^{-1}(U)\), and \(E\) be a complex of left \(A_U\)-Modules. When \(E\) is acyclic in degree \(>n\), \(E'\otimes_A^{\mathrm L}E\) is acyclic in degree \(>n+n'\) (definition of \(\otimes_A^{\mathrm L}\)); on the other hand, for \(E=(A_U)^p\), we have \(E'\otimes_A^{\mathrm L}E=(E'|U')^p\), which is \(m'\)-pseudo-coherent by hypothesis (and (2.12)). The conditions of (2.13) are therefore satisfied, with \(a=m'\), \(b=n'\), whence the conclusion.

In particular:

\sgaStableTarget{sga6.e01.cor.2-16-1}\textbf{Corollary 2.16.1.} \uline{Under the hypotheses of} (2.16), \uline{we have}:

a) \uline{The functor} \(\mathbb{L}f^*: \Dpx^-(_A \underline{S})\to \Dpx^-(_{A'}\underline{S}')\) \uline{(which is simply the functor} \(A'\otimes_A^{\mathrm L}\)\uline{) induces a functor}
\[
  \mathbb{L}f^*: \Dpx^-(_A \underline{S})_{\coh}\longrightarrow \Dpx^-(_{A'}\underline{S}')_{\coh}
\]
(\uline{notation of} (2.6)).

b) \uline{When} \(A\) \uline{is commutative, the functor} \(\otimes_A^{\mathrm L}:\Dpx^-(\underline{S})\times_{\underline{S}}\Dpx^-(\underline{S})\longrightarrow \Dpx^-(\underline{S})\) \uline{induces a functor}
\[
  \otimes_A^{\mathrm L}:\Dpx^-(\underline{S})_{\coh}\times_{\underline{S}}\Dpx^-(\underline{S})_{\coh}
  \longrightarrow \Dpx^-(\underline{S})_{\coh} .
\]

%%%%%%%%%%  scan idx 125  |  folio 112  %%%%%%%%%%
\SGAPrintPage{112}%
\sgaStableTarget{sga6.e01.exercise.2-16-2}\textbf{Exercise 2.16.2.} Let \(\underline{S}\) be a topos, and \(A\to A'\) a morphism of Rings of \(\underline{S}\). We make one of the following hypotheses:
\begin{enumerate}[label=(\roman*)]
\item \(A\to A'\) is surjective with kernel a nilpotent Ideal
\item the functor \(A'\otimes_A-\) on the category of left \(A\)-modules is exact and faithful.
\end{enumerate}
Show that \(E\in\Ob\Dpx^-(_A \underline{S})\) is \(n\)-pseudo-coherent (resp.\ pseudo-coherent) if and only if \(A'\otimes_A^{\mathrm L} E\in\Ob\Dpx^-(_{A'}\underline{S})\) is \(n\)-pseudo-coherent (resp.\ pseudo-coherent).

Hints: First show that a left \(A\)-Module \(L\) is of finite type if and only if \(A'\otimes_A L\) is (cf.\ Bourbaki, Alg. commutative, chap.\ I, \S\ 3, prop.\ 11 and chap.\ II, \S\ 3, prop.\ 4), then reduce to this case by truncation, using (2.5) and (2.10 b)).

\sgaStableTarget{sga6.e01.sec.3}\section*{3. Link with the classical notion of coherence}

\sgaStableTarget{sga6.e01.para.3-0}3.0. The hypotheses and notation are those of (2.0), except that \(\underline{C}_0\) is no longer assumed locally quasi-stable under kernels of epimorphisms.

\sgaStableTarget{sga6.e01.def.3-1}\textbf{Definition 3.1.} \uline{Let} \(F\in\Ob\underline{C}_S\)\uline{. We shall say that} \(F\) \uline{is} precoherent \uline{(relative to} \(\underline{C}_0\)\uline{)}\SGAAuthorFootnote{*}{We introduce this notion here as a technical intermediate step; the only notion to be retained is that of coherence.} \uline{when, for every} \(U\) \uline{over} \(S\) \uline{and every morphism} \(u:E\to F|U\) \uline{with} \(E\in\Ob\underline{C}_{0U}\)\uline{,} \(\Ker(u)\) \uline{is of} \(\underline{C}_0\)\uline{-finite type} (1.2)\uline{. We shall say that} \(F\) \uline{is} coherent \uline{(relative to} \(\underline{C}_0\)\uline{) when} \(F\) \uline{is precoherent and of finite type.}

\sgaStableTarget{sga6.e01.rem.3-2}\textbf{Remark 3.2.} Let \(F\in\Ob\underline{C}_S\). It is clear from the definitions that, when \(F\) is precoherent (resp.\ of finite type), every subobject (resp.\ quotient) of \(F\) is precoherent (resp.\ of finite type).

\sgaStableTarget{sga6.e01.lem.3-3}\var{\textbf{Lemma 3.3}}{\textbf{Lemma 3.3.}}{EN-VAR-094} a) \uline{Let}
\[
  0\longrightarrow F\longrightarrow G\longrightarrow H
\]
%%%%%%%%%%  scan idx 126  |  folio 113  %%%%%%%%%%
\SGAPrintPage{113}%
\uline{be an exact sequence of} \(\underline{C}_S\)\uline{. When} \(G\) \uline{is of finite type and} \(H\) \uline{precoherent,} \(F\) \uline{is of finite type.}

b) \uline{Let}
\[
  F\longrightarrow G\longrightarrow H\longrightarrow 0
\]
\uline{be an exact sequence of} \(\underline{C}_S\)\uline{. When} \(F\) \uline{is of finite type and} \(G\) \uline{precoherent,} \(H\) \uline{is precoherent.}

c) \uline{Let}
\[
  0\longrightarrow F\longrightarrow G\longrightarrow H\longrightarrow 0
\]
\uline{be an exact sequence of} \(\underline{C}_S\)\uline{. When} \(F\) \uline{and} \(H\) \uline{are precoherent (resp.\ of finite type),} \(G\) \uline{is too.}

\uline{Proof.} a) After localizing, we have an epimorphism \(G'\to G\) with \(G'\in\Ob\underline{C}_0\). Hence a commutative diagram with exact rows
\[
\begin{tikzcd}
0\arrow[r]&F'\arrow[r]\arrow[d]&G'\arrow[r]\arrow[d]&H\arrow[d,"\Id"]\\
0\arrow[r]&F\arrow[r]&G\arrow[r]&H .
\end{tikzcd}
\]
\(F'\) is of finite type because \(H\) is precoherent, but \(F'\to F\) is an epimorphism, so \(F\) is of finite type.

b) Let \(U\) be over \(S\), and \(H'\xrightarrow{f}H|U\) a morphism of \(\underline{C}_U\) with \(H'\in\Ob\underline{C}_{0U}\). Changing notation, we can suppose that \(U=S\). By the quasi-liftability of \(\underline{C}_0\) in \(\underline{C}\), there exists (after localizing) a commutative square
\[
\begin{tikzcd}
G' \arrow[r,"p"] \arrow[d] & H' \arrow[d,"f"]\\
G \arrow[r] & H,
\end{tikzcd}
\]
%%%%%%%%%%  scan idx 127  |  folio 114  %%%%%%%%%%
\SGAPrintPage{114}%
where the horizontal arrows are epimorphisms and \(G'\in\Ob\underline{C}_0\). Since \(\Ker(f)\) is a quotient of \(\Ker(fp)\), it suffices to prove that \(\Ker(fp)\) is of finite type, so we can suppose that \(H'=G'\). The arrows \(F\to G\) and \(H'\to G\) define \(g:F\oplus H'\to G\). Putting \(K=\Ker(G\to H)\), we therefore have a commutative diagram in which the rows are exact and \(q\) is an epimorphism:
\[
\begin{tikzcd}
0\arrow[r]&F\arrow[r]\arrow[d,"q"']&F\oplus H'\arrow[r]\arrow[d,"g"]&H'\arrow[r]\arrow[d,"f"]&0\\
0\arrow[r]&K\arrow[r]&G\arrow[r]&H\arrow[r]&0 .
\end{tikzcd}
\]
The snake diagram shows that \(\Ker(g)\to\Ker(f)\) is an epimorphism. But, since \(F\oplus H'\) is of finite type and \(G\) precoherent, \(\Ker(g)\) is of finite type by a), and we are done.

c) Let us first prove the assertion without the ``resp.'' alternative. Up to a change of notation, we must show that, for a morphism \(g:G'\to G\) of \(\underline{C}_S\) with \(G'\in\Ob\underline{C}_0\), \(\Ker(g)\) is of finite type. Denoting by \(H'\) the image of the composite arrow \(G'\to H\), we have a commutative diagram whose rows are exact
\[
\begin{tikzcd}
0\arrow[r]&F'\arrow[r]\arrow[d,"f"']&G'\arrow[r]\arrow[d,"g"]&H'\arrow[r]\arrow[d,"h"]&0\\
0\arrow[r]&F\arrow[r]&G\arrow[r]&H\arrow[r]&0 .
\end{tikzcd}
\]
Since \(h\) is a monomorphism, the snake diagram gives an isomorphism \(\Ker(f)\qis\Ker(g)\). But \(F'\) is of finite type by a), so (again by a)) \(\Ker(f)\) is of finite type, which proves the assertion without the ``resp.'' alternative. Let us pass to the ``resp.'' assertion (strictly speaking, we cannot invoke (2.5 b)) because we do not assume that \(\underline{C}_0\) is locally quasi-stable under kernels of epimorphisms). After localizing, there exist epimorphisms \(F'\xrightarrow{u}F\), \(H'\xrightarrow{v}H\), with
%%%%%%%%%%  scan idx 128  |  folio 115  %%%%%%%%%%
\SGAPrintPage{115}%
\(F'\in\Ob\underline{C}_0\), \(H'\in\Ob\underline{C}_0\). After localizing further, we can suppose that \(v\) factors through \(G\), whence a commutative diagram with exact rows
\[
\begin{tikzcd}
0\arrow[r]&F'\arrow[r]\arrow[d,"u"']&F'\oplus H'\arrow[r]\arrow[d,"w"]&H'\arrow[r]\arrow[d,dashed,"v"]&0\\
0\arrow[r]&F\arrow[r]&G\arrow[r]&H\arrow[r]&0 .
\end{tikzcd}
\]
Since \(u\) and \(v\) are epimorphisms, so is \(w\), and we are done. This completes the proof of (3.3).

We deduce immediately the

\sgaStableTarget{sga6.e01.cor.3-4}\textbf{Corollary 3.4.} \uline{The (strictly) full subcategory of} \(\underline{C}_S\) \uline{formed by the coherent objects is stable under finite projective (resp.\ inductive) limits. In particular, the kernel, cokernel and image of an arrow between two coherent objects are coherent. Let}
\[
  E^0\longrightarrow E^1\longrightarrow E^2\longrightarrow E^3\longrightarrow E^4
\]
\uline{be an exact sequence of} \(\underline{C}_S\)\uline{. When} \(E^0,E^1,E^3,E^4\) \uline{are coherent,} \(E^2\) \uline{is too.}

\sgaStableTarget{sga6.e01.cor.3-5}\textbf{Corollary 3.5.} \uline{Assume that the objects of} \(\underline{C}_0\) \uline{are} \(\underline{C}_0\)\uline{-coherent (which implies in particular that} \(\underline{C}_0\) \uline{is locally quasi-stable under kernels of epimorphisms,} i.e.\ \uline{that the conditions of} (2.0) \uline{hold}).

a) \uline{Equivalent conditions on} \(F\in\Ob\underline{C}_S\):
\begin{enumerate}[label=(\roman*)]
\item \(F\) \uline{is of finite presentation};
\item \(F\) \uline{is coherent};
\item \(F\) \uline{is of} \(\infty\)\uline{-finite presentation} (2.8).
\end{enumerate}

b) \uline{Let} \(F\in\Ob\Dpx^-(\underline{C}_S)\), \(n\in\mathbb Z\). \uline{We have}:
\[
\begin{split}
F\ n\text{-\uline{pseudo-coherent}}
\quad\Longleftrightarrow\quad&
\Ho^i(F)\text{ \uline{coherent for} }i\geq n+1,\\
&\text{\uline{and} }\Ho^n(F)\text{ \uline{of finite type}.}
\end{split}
\]
%%%%%%%%%%  scan idx 129  |  folio 116  %%%%%%%%%%
\SGAPrintPage{116}%
\uline{In particular, for} \(F\) \uline{to be pseudo-coherent, it is necessary and sufficient that} \(\Ho^n(F)\) \uline{be coherent for every} \(n\)\uline{.}

\uline{Proof.} Assertion a) and the implication \(\Rightarrow\) of b) follow trivially from (3.4). The implication \(\Leftarrow\) of b) follows at once from a) and (2.11.b)).

\sgaStableTarget{sga6.e01.ex.3-6}\textbf{Examples 3.6.} Let \((\underline{S},A)\) be a ringed topos, and \(\underline{C},\underline{C}_0\) the \(\underline{S}\)-categories of (1.3.d)) (cf.\ also (2.15)). For the objects of \(\underline{C}_0\) to be coherent, it is necessary and sufficient that \(A\) be coherent (i.e.\ precoherent). Here are examples where this holds:
\begin{itemize}[label=-]
\item \(A\) is constant with value a left noetherian ring (SGA 4 IX 2.1).
\item \(\underline{S}\) is the topos of sheaves of sets (for the Zariski topology) on a locally noetherian scheme \(S\)\SGAAuthorFootnote{*}{In accordance with the new conventions, we shall say ``scheme'' for ``prescheme'', and ``separated scheme'' for ``scheme''.}, and \(A=\underline{O}_S\).
\item \(\underline{S}\) is the topos of sheaves of sets on a complex analytic space \(S\), \(A=\underline{O}_S\).
\end{itemize}
On the other hand, the sheaf of continuous complex-valued functions on a topological space is in general not coherent.

Of course, when \((\underline{S},A)\) is defined by a ringed space, the notion of coherence considered here coincides with Serre's usual notion (EGA \(0_I\) 5.3.1).

\sgaStableTarget{sga6.e01.prop.3-7}\textbf{Proposition 3.7.} a) \uline{Let} \(f:(\underline{S}',A')\to(\underline{S},A)\) \uline{be a morphism of ringed topoi. Assume that} \(f\) \uline{is flat on the right,} i.e.\ \uline{that the functor} \(f^*=A'\otimes_A.\) \uline{is exact. For} \(F\in\Ob\Dpx(_A S)\)\uline{,} \(G\in\Ob\Dpx^+(_A S)\)\uline{, there is a functorial canonical arrow:}
\[
  f^*\mathbb{R}\underline{\Hom}(F,G)\longrightarrow\mathbb{R}\underline{\Hom}(f^*(F),f^*(G)).
\]
%%%%%%%%%%  scan idx 130  |  folio 117  %%%%%%%%%%
\SGAPrintPage{117}%
\uline{This is an isomorphism when} \(F\) \uline{is pseudo-coherent.}

b) \uline{Suppose that} \(A\) \uline{is commutative and coherent, and let} \(F\in\Ob\Dpx^-(S)\)\uline{,} \(G\in\Ob\Dpx^+(S)\)\uline{. If} \(\underline{\Ho}^i(F)\) and \(\underline{\Ho}^i(G)\) \uline{are coherent for every} \(i\)\uline{, the same is true of} \(\underline{\Ext}^i(F,G)\)\uline{.}

\uline{Proof.} a) To define the arrow (cf.\ SGA 4 XVII), we can suppose that \(G\) has injective components, and take the composite arrow
\[
  f^*\underline{\Hom}(F,G)\longrightarrow\underline{\Hom}(f^*F,f^*G)\longrightarrow\mathbb{R}\underline{\Hom}(f^*F,f^*G).
\]
Let us show that it is an isomorphism for \(F\) pseudo-coherent. For brevity, write \(T\) (resp.\ \(T'\)) for the functor \(f^*\mathbb{R}\underline{\Hom}(.,G)\) (resp.\ \(\mathbb{R}\underline{\Hom}(f^*.,f^*G)\)), and let \(a\in\mathbb Z\) be such that \(\Ho^i(G)=0\) for \(i<a\). It is clear that \(TA\qis T'A\), hence, by dévissage, \(TF\qis T'F\) for \(F\) strictly perfect (2.1). On the other hand, when \(F\) is acyclic in degree \(>n\), \(TF\) and \(T'F\) are acyclic in degree \(<a-n\). To verify that \(TF\qis T'F\) for \(F\) pseudo-coherent, we can, since the question is local, suppose that \(F\) is strictly \(n\)-pseudo-coherent, i.e.\ that there is a distinguished triangle
\[
\begin{tikzcd}[column sep=small]
& F'' \arrow[dl,"+1"'] & \\
F' \arrow[rr] & & F \arrow[ul]
\end{tikzcd}
\;,
\]
with \(F'\) strictly perfect and \(F''\) acyclic in degree \(>n-1\). We deduce a morphism of distinguished triangles
\[
\begin{tikzcd}[column sep=small]
& TF'' \arrow[dr] & \\
TF' \arrow[ur,"-1"'] & & TF \arrow[ll]
\end{tikzcd}
\qquad\longrightarrow\qquad
\begin{tikzcd}[column sep=small]
& T'F'' \arrow[dr] & \\
T'F' \arrow[ur,"-1"'] & & T'F \arrow[ll]
\end{tikzcd}
\;.
\]
Now \(TF'\qis T'F'\) since \(F'\) is strictly perfect, and \(TF''\) and \(T'F''\) are
%%%%%%%%%%  scan idx 131  |  folio 118  %%%%%%%%%%
\SGAPrintPage{118}%
acyclic in degree \(\leq a-n\); writing the induced morphism on the exact cohomology sequences of the preceding triangles, we deduce that \(TF\to T'F\) induces an isomorphism \(\Ho^i(TF)\qis\Ho^i(T'F)\) for \(i\leq a-n\). Since we can choose \(n\) arbitrarily negative, we are done.

b) The proof is left as an exercise to the reader: see, for example, ([H] II 3.3).

\sgaStableTarget{sga6.e01.rem.3-7-1}\textbf{Remark 3.7.1.} We shall see below (7.1.2) another case of ``compatibility'' of \(\mathbb{R}\underline{\Hom}\) with \(\mathbb{L}f^*\).

\sgaStableTarget{sga6.e01.cor.3-7-2}\textbf{Corollary 3.7.2.} (cf.\ EGA \(0_I\) 5.2.7). \uline{Let} \((\underline{S},A)\) \uline{be a ringed topos,} \(s\) \uline{a point of} \(\underline{S}\) \var{(SGAA VIII 7.8)}{(SGA 4 VIII 7.8)}{EN-VAR-095}\uline{,} \(E,F\in\Ob\D^b(_A S)_{\coh}\)\uline{. If} \(E_s\simeq F_s\)\uline{, there exists an} \(s\)\uline{-pointed object} \(U\) \uline{of} \(\underline{S}\) \uline{such that} \(E|U\simeq F|U\)\uline{.}

\uline{Proof.} Let \(\alpha:E_s\to F_s\) and \(\beta:F_s\to E_s\) be reciprocal isomorphisms. \var{By (2.22 a)),}{By (3.7 a)),}{EN-VAR-096} we have
\[
  \underline{\Hom}(E,F)_s\qis\Hom(E_s,F_s)
  \qquad
  (\text{resp. }\underline{\Hom}(F,E)_s\qis\Hom(F_s,E_s)).
\]
There are therefore a ``neighbourhood'' \(U\) of \(s\) and a section \(u\) (resp.\ \(v\)) of \(\underline{\Hom}(E,F)\) (resp.\ \(\underline{\Hom}(F,E)\)) over \(U\) such that \(u_s=\alpha\) (resp.\ \(v_s=\beta\)). The homomorphisms \(vu:E|U\to E|U\) and \(uv:F|U\to F|U\) induce the identity at \(s\), hence over a suitable \(V\) over \(U\), as required.

%%%%%%%%%%  scan idx 132  |  folio 119  %%%%%%%%%%
\SGAPrintPage{119}%
\sgaStableTarget{sga6.e01.sec.4}\section*{4. Perfect complexes}

\sgaStableTarget{sga6.e01.para.4-0}4.0. In this section, \(\underline{S}\) denotes a site, \(S\) an object of \(\underline{S}\), \(\underline{C}\) a flat abelian \(\underline{S}\)-category, and \(\underline{C}_0\) a strictly full additive sub-\(\underline{S}\)-category of \(\underline{C}\) (1.1). Unless otherwise mentioned, we assume that \(\underline{C}_0\) is locally liftable in \(\underline{C}\) and locally stable under kernels of epimorphisms (1.2).

\sgaStableTarget{sga6.e01.lem.4-1}\textbf{Lemma 4.1}\SGAAuthorFootnote{raisedstar}{This lemma was communicated to the editor by D. Ferrand.}. \uline{Let} \(u:E\to F\) \uline{be an arrow of} \(\Cpx(\underline{C}_S)\)\uline{. If} \(E\) \uline{is strictly} \(\underline{C}_0\)\uline{-perfect} (2.1) \uline{and} \(F\) \uline{is acyclic, then} \(u\) \uline{is locally homotopic to zero.}

\uline{Proof.} We must construct locally a homotopy \(k\) such that \(u=dk+kd\). Since \(E\) is strictly perfect, there exist \(a,b\in\mathbb Z\), with \(a\leq b\), such that \(E^i=0\) for \(i\notin[a,b]\). We can therefore take \(k^i=0\) for \(i\notin[a,b]\). We then construct \(k^i\) for \(i\geq a\) by descending induction on \(i\). Suppose that \(k^i:E^i\to F^{i-1}\) has been constructed so that \(u^i=k^{i+1}d_E^i+d_F^{i-1}k^i\) for \(i\geq n>a\), and let us show that, after localizing, there exists \(k^{n-1}:E^{n-1}\to F^{n-2}\) such that \(u^{n-1}=k^nd_E^{n-1}+d_F^{n-2}k^{n-1}\). By the induction hypothesis, the arrow
\[
  v^{n-1}=u^{n-1}-k^nd_E^{n-1}:E^{n-1}\to F^{n-1}
\]
factors through \(Z^{n-1}(F)\). But \(Z^{n-1}(F)=B^{n-1}(F)\) since \(F\) is acyclic. Since \(\underline{C}_0\) is locally liftable in \(\underline{C}\), it follows that, after localizing, there exists \(k^{n-1}:E^{n-1}\to F^{n-2}\) such that \(v^{n-1}=d_F^{n-2}k^{n-1}\), as required.

\sgaStableTarget{sga6.e01.cor.4-2}\textbf{Corollary 4.2.} \uline{Let}
\[
\begin{tikzcd}
& P \arrow[d,"v"] \\
E \arrow[r,"f"] & F
\end{tikzcd}
\]
%%%%%%%%%%  scan idx 133  |  folio 120  %%%%%%%%%%
\SGAPrintPage{120}%
\uline{be a diagram of} \(\Cpx(\underline{C}_S)\)\uline{, where} \(P\) \uline{is strictly} \(\underline{C}_0\)\uline{-perfect and} \(f\) \uline{a quasi-isomorphism. The set of} \(X\) \uline{over} \(S\) \uline{for which there exists an arrow} \(u:P_X\to E_X\) \uline{of} \(\Cpx(\underline{C}_X)\) \uline{such that the diagram}
\[
\begin{tikzcd}
& P_X \arrow[d,"v_X"] \arrow[dl,"u"']\\
E_X \arrow[r,"f_X"] & F_X
\end{tikzcd}
\]
\uline{is commutative in} \(\Kpx(\underline{C}_X)\) \uline{is a refinement of} \(S\)\uline{.}

\uline{Proof.} Let \(G\) be the cone of \(f\). We have the exact sequence
\[
  \longrightarrow \Hom_{\Kpx(\underline{C}_S)}(P,E)
  \longrightarrow \Hom_{\Kpx(\underline{C}_S)}(P,F)
  \longrightarrow \Hom_{\Kpx(\underline{C}_S)}(P,G)
  \longrightarrow .
\]
Since \(G\) is acyclic, the image of \(v\) in \(\Hom_{\Kpx(\underline{C}_S)}(P,G)\) is locally zero by (4.1). Thus, after localizing, there exists \(u:P\to E\) such that \(v=fu\) in \(\Kpx(\underline{C}_S)\), as required.

\sgaStableTarget{sga6.e01.cor.4-3}\textbf{Corollary 4.3.} \uline{Let} \(f:E\to F\) \uline{be a quasi-isomorphism of} \(\Cpx(\underline{C}_S)\)\uline{, with} \(F\) \uline{strictly perfect (relative to} \(\underline{C}_0\)\uline{). The set of} \(X\) \uline{over} \(S\) \uline{for which there exists a quasi-isomorphism} \(g:F_X\to E_X\) \uline{such that} \(f_Xg\) \uline{is homotopic to} \(\Id_{F_X}\) \uline{is a refinement of} \(S\)\uline{.}

\uline{Proof.} This is a special case of (4.2).

\sgaStableTarget{sga6.e01.schol.4-4}\textbf{Scholium 4.4.} Let \(F\in\Ob\Kpx(\underline{C}_S)\). We can define (for example using a cleavage of \(\underline{C}\) (SGA 1 VI 7.1)) a category \((\underline{\Qis/F})\) fibered over \(\underline{S}_{/S}\), whose fibre at \(X\) over \(S\) is the category \((\Qis/F_X)^\circ\) (opposite to the category of quasi-isomorphisms with target \(F_X\)). We can then express (4.3) by saying that,
%%%%%%%%%%  scan idx 134  |  folio 121  %%%%%%%%%%
\SGAPrintPage{121}%
when \(F\) is strictly perfect, the arrow \(\Id_F\) is locally cofinal in \((\underline{\Qis/F})\).

\sgaStableTarget{sga6.e01.notation.4-5}\textbf{Notation 4.5.} Let \(\underline{M}\) be a fibered \(\underline{S}\)-category. Let \(X\in\Ob \underline{S}\), and \(E,F\in\Ob\underline{M}_X\). We shall denote by \(\underline{\Hom}_{\underline{M}}(E,F)\) the sheaf on \(\underline{S}_{/X}\) associated to the presheaf of \(\underline{S}\)-morphisms from \(E\) to \(F\) in \(\underline{M}\) (for the definition of this presheaf, see [G] I 2.6; take care that, even when \(\underline{C}\) is a stack ([G]), the corresponding presheaf for \(\underline{M}=\Kpx(\underline{C})\) (or \(\Dpx(\underline{C})\)) is in general not a sheaf).

\sgaStableTarget{sga6.e01.cor.4-6}\textbf{Corollary 4.6.} \uline{Let} \(E,F\in\Ob\Cpx(\underline{C}_S)\)\uline{, with} \(E\) \uline{strictly perfect. The canonical arrow}
\[
  \underline{\Hom}_{\Kpx(\underline{C})}(E,F)
  \longrightarrow
  \underline{\Hom}_{\Dpx(\underline{C})}(E,F)
\]
\uline{is an isomorphism}.

\uline{Proof.} This is an immediate consequence of (4.4) and of the formula
\[
  \Hom_{\Dpx(\underline{C}_X)}(E_X,F_X)
  = \varinjlim_{(\Qis/E_X)^\circ}\Hom_{\Kpx(\underline{C}_X)}(\,.\,,F_X)
\]
for \(X\) over \(S\).

\sgaStableTarget{sga6.e01.def.4-7}\textbf{Definition 4.7.} \uline{Let} \(F\in\Ob\Dpx(\underline{C}_S)\) \uline{and let} \(a,b\in\mathbb Z\) \uline{be such that} \(a\leq b\)\uline{. We shall say that} \(F\) \uline{has} \(\underline{C}_0\)-perfect amplitude contained in the interval \([a,b]\), \uline{and we shall write}
\[
  \underline{C}_0\text{-}\mathrm{parf}\text{-}\amp(F)\subset[a,b],
\]
\uline{when the set of} \(X\) \uline{over} \(S\) \uline{for which there exists a quasi-isomorphism}
%%%%%%%%%%  scan idx 135  |  folio 122  %%%%%%%%%%
\SGAPrintPage{122}%
\(F'\to F\) \uline{of} \(\Cpx(\underline{C}_X)\) \uline{with} \(F'\) \uline{strictly} \(\underline{C}_0\)\uline{-perfect} (2.1)\uline{, and such that} \(F'^i=0\) \uline{for} \(i\notin[a,b]\)\uline{, is a refinement of} \(S\). (\uline{By} (4.3) \uline{this condition is equivalent to the one obtained by replacing the expression ``quasi-isomorphism of} \(\Cpx(\underline{C}_X)\)\uline{'' by ``isomorphism of} \(\Dpx(\underline{C}_X)\)\uline{''.}) \uline{We shall say that} \(F\) \uline{has} finite \(\underline{C}_0\)-perfect amplitude \uline{when there exist integers} \(a\) \uline{and} \(b\) \uline{such that} \(\underline{C}_0\)-\(\mathrm{parf}\)-\(\amp(F)\subset[a,b]\). \uline{We shall say that} \(F\) \uline{is} \(\underline{C}_0\)-perfect \uline{when} \(F\) \uline{locally has finite} \(\underline{C}_0\)\uline{-perfect amplitude}\SGAAuthorFootnote{raisedstar}{Of course, when no confusion can arise, we shall omit \(\underline{C}_0\) from the preceding expressions.}.

\sgaStableTarget{sga6.e01.ex.4-8}\textbf{Example 4.8.} Let \((\underline{S},A)\) be a ringed topos, \(\underline{C}\) the \(\underline{S}\)-category defined by \(U\longmapsto\) (category of left \(A_U\)-Modules), and \(\underline{C}_0\) the sub-\(\underline{S}\)-category defined by \(U\longmapsto\) (category of left \(A_U\)-Modules that are direct factors of free Modules of finite type). By (1.3.1), the pair \((\underline{C},\underline{C}_0)\) satisfies the hypotheses (4.0). An object \(F\) of \(\Dpx(_A \underline{S})\) (cf.\ (2.15)) will be said to have perfect amplitude contained in \([a,b]\) (resp.\ finite perfect amplitude, resp.\ to be perfect) when \(F\) has \(\underline{C}_0\)-perfect amplitude contained in \([a,b]\) (resp.\dots). When the direct factors of free Modules of finite type are locally free (cf.\ (2.15.1)), saying that \(F\) is perfect therefore means that \(F\) is locally isomorphic (in \(\Dpx(\underline{S})\)) to a complex with bounded degrees and components free of finite type.

\sgaStableTarget{sga6.e01.notation.4-9}\textbf{Notation 4.9.} We shall denote by \(\Dpx(\underline{C})_{\underline{C}_0\text{-}\mathrm{parf}}\) or \(\Parf(\underline{C},\underline{C}_0)\) the strictly full sub-\(\underline{S}\)-category of \(\Dpx(\underline{C})\) formed by the \(\underline{C}_0\)-perfect objects. In the case of Example (4.8), we shall write \(\Dpx(_A \underline{S})_{\mathrm{parf}}\), or \(\Parf(_A \underline{S})\), or simply \(\Parf(\underline{S})\) when there is no possible ambiguity about \(A\).

%%%%%%%%%%  scan idx 136  |  folio 123  %%%%%%%%%%
\SGAPrintPage{123}%
\sgaStableTarget{sga6.e01.prop.4-10}\var{\textbf{Proposition 4.10}}{\textbf{Proposition 4.10.}}{EN-VAR-097} \uline{The category} \(\underline{\Parf}(\underline{C},\underline{C}_0)\) \uline{is a triangulated sub-}\(\underline{S}\)\uline{-category of} \(\Dpx(\underline{C})\) (1.1.4)\uline{. The canonical functor}
\[
  \Kpx^b(\underline{C}_0)\longrightarrow\Parf(\underline{C},\underline{C}_0)
\]
(resp.\ \(\Tr(\Kpx^b(\underline{C}_0))\to\Tr(\Parf(\underline{C},\underline{C}_0))\) (1.1.3)) \uline{induces an} \(\underline{S}\)\uline{-equivalence on the associated stacks} ([G] II 2.2.4)\uline{, which means}:
\begin{enumerate}[label=(\roman*)]
\item \uline{for every} \var{\(X\in\Ob S\)}{\(X\in\Ob \underline{S}\)}{EN-VAR-098} \uline{and all} \(E,F\in\Ob\Kpx^b(\underline{C}_{0X})\) (resp.\ \(\Tr(\Kpx^b(\underline{C}_{0X}))\)), \uline{the canonical arrow}
\[
  \underline{\Hom}_{\Kpx^b(\underline{C}_0)}(E,F)
  \longrightarrow
  \underline{\Hom}_{\Dpx(\underline{C})}(E,F)
\]
(\uline{resp.}\ \(\underline{\Hom}_{\Tr(\Kpx^b(\underline{C}_0))}(E,F)\to\underline{\Hom}_{\Tr(\Dpx(\underline{C}))}(E,F)\)) \uline{is an isomorphism, and}
\item \uline{every object of} \(\underline{\Parf}(\underline{C},\underline{C}_0)\) (resp.\ \(\Tr(\Parf(\underline{C},\underline{C}_0))\)) \uline{is locally isomorphic in} \(\Dpx(\underline{C})\) \uline{(resp.}\ \var{\(\Tr(\Dpx(\underline{C}))\) \uline{to an object of}}{\(\Tr(\Dpx(\underline{C}))\)) \uline{to an object of}}{EN-VAR-099} \(\Kpx^b(\underline{C}_0)\) \var{\uline{(resp.}\ \(\Tr(\Kpx^b(\underline{C}_0))\).}{\uline{(resp.}\ \(\Tr(\Kpx^b(\underline{C}_0))\)).}{EN-VAR-100}
\end{enumerate}

\uline{Proof.} Assertion (i) without the ``resp.'' alternative is a special case of (4.6), and assertion (ii) without the ``resp.'' alternative is true by definition ((4.7) and (4.9)). It is also clear that \(\Parf(\underline{C},\underline{C}_0)\) is stable under translation of degrees. To prove that \(\Parf(\underline{C},\underline{C}_0)\) is a triangulated sub-\(\underline{S}\)-category of \(\Dpx(\underline{C})\), it therefore suffices to show that, when
\sgaStableTarget{sga6.e01.eq.unnamed.occ3}\[
\begin{tikzcd}
& G \arrow[dl,"+1"'] & \\
E \arrow[rr,"u"] && F \arrow[ul]
\end{tikzcd}
\tag{*}
\]
is a distinguished triangle of \(\Dpx(\underline{C})\) such that \(E\) and \(F\) are perfect, \(G\)
%%%%%%%%%%  scan idx 137  |  folio 124  %%%%%%%%%%
\SGAPrintPage{124}%
is perfect. After localizing, we can suppose that \(E\) and \(F\) are objects of \(\Kpx^b(\underline{C}_0)\). After localizing further, we can suppose, by assertion (i) without the ``resp.'' alternative, that \(u\) is an arrow of \(\Kpx^b(\underline{C}_0)\); then \(G\) is perfect because it is isomorphic to the cone of \(u\), which is strictly perfect. Thus \(\Parf(\underline{C},\underline{C}_0)\) is indeed a triangulated sub-\(\underline{S}\)-category of \(\Dpx(\underline{C})\). Moreover, the preceding argument also proves the ``resp.'' assertion (ii). It remains to prove the ``resp.'' assertion (i). Let \(X\in\Ob \underline{S}\), and let \(E=(E_1\xrightarrow{a_1}E_2\xrightarrow{a_2}E_3\xrightarrow[+1]{a_3}E_1)\), \(F=(F_1\xrightarrow{b_1}F_2\xrightarrow{b_2}F_3\xrightarrow[+1]{b_3}F_1)\) be distinguished triangles of \(\Kpx^b(\underline{C}_{0X})\). Let us show that the arrow
\sgaStableTarget{sga6.e01.eq.unnamed.occ4}\[
  \underline{\Hom}_{\Tr(\Kpx^b(\underline{C}_0))}(E,F)
  \longrightarrow
  \underline{\Hom}_{\Tr(\Dpx(\underline{C}))}(E,F)
  \tag{**}
\]
is an isomorphism. By assertion (i) without the ``resp.'' alternative, we already know that \((**)\) is injective. Let \((u_i):E\to F\) be a morphism of triangles in \(\Dpx(\underline{C})\). After localizing, we can suppose, by assertion (i) without the ``resp.'' alternative, that \(u_i\) (\(1\leq i\leq3\)) is an arrow of \(\Kpx^b(\underline{C}_0)\). The arrow \(u_{i+1}a_i-b_iu_i\), for \(1\leq i\leq3\) (with the convention \(u_4=u_1\)), is zero in \(\Dpx(\underline{C})\), hence locally zero in \(\Kpx^b(\underline{C}_0)\) by assertion (i) without the ``resp.'' alternative. In other words, \((u_i)\) locally defines a morphism of triangles in \(\Kpx^b(\underline{C}_0)\). This proves that \((**)\) is surjective and completes the proof of (4.10).

\sgaStableTarget{sga6.e01.rem.4-10-1}\textbf{Remark 4.10.1.} In (4.10), we can replace \(\Parf(\underline{C},\underline{C}_0)\) by \(\Dpx^*(\underline{C})_{\underline{C}_0\text{-}\mathrm{parf}}=\Dpx^*(\underline{C})\cap\Parf(\underline{C},\underline{C}_0)\) \((*=-\text{ or }b)\); nothing in the proof needs to be changed.

\sgaStableTarget{sga6.e01.complement.4-11}\textbf{Complement 4.11.} We can describe more precisely the triangulated structure of \(\Parf(\underline{C},\underline{C}_0)\) in terms of perfect amplitude (4.7). When \(E\in\Ob\Dpx(\underline{C})\) has perfect amplitude contained in \([a,b]\), \(E[n]\) has perfect amplitude contained in \([a-n,b-n]\).

%%%%%%%%%%  scan idx 138  |  folio 125  %%%%%%%%%%
\SGAPrintPage{125}%
Let
\[
\begin{tikzcd}
& G \arrow[dl,"+1"'] & \\
E \arrow[rr] && F \arrow[ul]
\end{tikzcd}
\]
be a distinguished triangle of \(\Dpx(\underline{C})\). We can draw up a table analogous to (2.6), in which the entries indicate perfect amplitudes and each row corresponds to an implication whose conclusion is underlined.
\[
\begin{array}{ccc}
E&F&G\\[0.3em]
{}[a+1,b+1]&{}[a,b]&\underline{[a,b]}\\
{}[a,b]&\underline{[a,b]}&{}[a,b]\\
\underline{[a,b]}&{}[a,b]&{}[a-1,b-1].
\end{array}
\]
From (4.10) we immediately deduce the

\sgaStableTarget{sga6.e01.cor.4-12}\textbf{Corollary 4.12.} \uline{Suppose that} \(\underline{S}\) \uline{is a punctual site. For} \(F\in\Ob\Dpx(\underline{C})\) \uline{to have perfect amplitude (rel.}\ to \(\underline{C}_0\)\uline{) contained in} \([a,b]\)\uline{, it is necessary and sufficient that there exist a quasi-isomorphism} \(F'\to F\)\uline{, with} \(F'\) \uline{strictly perfect and such that} \(F'^i=0\) \uline{for} \(i\notin[a,b]\)\uline{. The canonical functors}
\[
  \Kpx^b(\underline{C}_0)\longrightarrow\Parf(\underline{C},\underline{C}_0)
\]
\uline{and}
\[
  \Tr(\Kpx^b(\underline{C}_0))\longrightarrow\Tr(\Parf(\underline{C},\underline{C}_0))
\]
\uline{are equivalences of categories}.

\sgaStableTarget{sga6.e01.rem.4-12-1}\textbf{Remark 4.12.1.} When \(\underline{S}\) is a punctual site, the objects of \(\underline{C}_0\) are projective. When \(\underline{C}_0\) contains all projective objects of \(\underline{C}\), one says ``projective amplitude'' instead of ``\(\underline{C}_0\)-perfect amplitude''.

% Additional notation for Expose I continuation
\providecommand{\Cpx}{\mathrm C}
\providecommand{\Kpx}{\mathrm K}
\providecommand{\Dpx}{\mathrm D}
\providecommand{\Ho}{\mathrm H}
\providecommand{\RHom}{\mathrm R\!\operatorname{Hom}}
\renewcommand{\Ob}{\operatorname{ob}}
\providecommand{\qis}{\xrightarrow{\sim}}
\providecommand{\RHom}{\mathrm R\!\Hom}
\providecommand{\Ext}{\operatorname{Ext}}
\providecommand{\Hom}{\operatorname{Hom}}
\providecommand{\Tr}{\operatorname{Tr}}

\providecommand{\Add}{\operatorname{Add}}
\providecommand{\Cart}{\underline{\operatorname{Cart}}}
\providecommand{\perfamp}{\operatorname{parf.amp}}
\providecommand{\perfdim}{\operatorname{parf.dim}}
\providecommand{\toramp}{\operatorname{tor.amp}}
\providecommand{\torf}{\operatorname{torf}}
\providecommand{\injf}{\operatorname{injf}}
\providecommand{\rg}{\operatorname{rg}}
\providecommand{\is}{\mathrm{is}}
\providecommand{\Mod}{\operatorname{Mod}}
\sgaStableTarget{sga6.e01.expose.occ3}%
%%%%%%%%%%  scan idx 139  |  folio 126  %%%%%%%%%%
\SGAPrintPage{126}%
\sgaStableTarget{sga6.e01.rem.4-12-2}\textbf{Remark 4.12.2.} Let \(\underline{C}\) be an abelian category (which can be regarded as fibered over a punctual site), and let \(\underline{C}_0\) be a strictly full additive subcategory. Suppose that \(\underline{C}_0\) is stable under kernels of epimorphisms and \uline{quasi-liftable} in \(\underline{C}\). Then an acyclic object of \(\Kpx^b(\underline{C}_0)\) is in general not zero, so there is no hope that the functor
\var{\(\Kpx^b(C_0)\longrightarrow \Dpx(\underline{C})\)}{\(\Kpx^b(\underline{C}_0)\longrightarrow \Dpx(\underline{C})\)}{EN-VAR-101}
is faithful. However, when \(\Kpx^{b,\varnothing}(\underline{C}_0)\) denotes the thick subcategory of \(\Kpx^b(\underline{C}_0)\) formed by the acyclic objects (in \(\underline{C}\)), the canonical functor
\[
  \Kpx^b(\underline{C}_0)/\Kpx^{b,\varnothing}(\underline{C}_0)
  \longrightarrow \Dpx(\underline{C})
\]
is fully faithful.

Hints: We already know (2.7) that the functor
\[
  \Kpx^-(\underline{C}_0)/\Kpx^{-,\varnothing}(\underline{C}_0)
  \longrightarrow \Dpx^-(\underline{C})
\]
is fully faithful. Show, for example using the criterion ([V] 4.2.b)(ii)'), that the functor
\[\var{
  \Kpx^b(\underline{C}_0)\Kpx^{b,\varnothing}(\underline{C}_0)
  \longrightarrow
  \Kpx^-(\underline{C}_0)/\Kpx^{-,\varnothing}(\underline{C}_0)
}{
  \Kpx^b(\underline{C}_0)/\Kpx^{b,\varnothing}(\underline{C}_0)
  \longrightarrow
  \Kpx^-(\underline{C}_0)/\Kpx^{-,\varnothing}(\underline{C}_0)
}{EN-VAR-102}\]
is fully faithful.

\sgaStableTarget{sga6.e01.lem.4-13}\textbf{Lemma 4.13.} \uline{Let} \(F\in\Ob\Dpx(\underline{C}_S)\)\uline{,} \(a,b,c\in\mathbb Z\) \uline{be such that} \(a\leq b\) \uline{and} \(a\leq c\)\uline{. If} \(F\) \uline{has perfect amplitude contained in} \([a,c]\) (4.7) \uline{and} \(\Ho^i(F)=0\) \uline{for} \(i>b\)\uline{, then} \(F\) \uline{has perfect amplitude contained in} \([a,d]\)\uline{, where} \(d=\inf(b,c)\)\uline{.}

\uline{Proof.} We can suppose that \(b\leq c\), and that \(F\) is strictly perfect with \(F^i=0\) for \(i\notin[a,c]\). Then \(F\) is isomorphic to its truncation \(F'\) defined by
%%%%%%%%%%  scan idx 140  |  folio 127  %%%%%%%%%%
\SGAPrintPage{127}%
\(F'^i=F^i\) for \(i<b\), \(F'^b=Z^b(F)\), \(F'^i=0\) for \(i>b\). But \(Z^b(F)\) admits a finite right resolution by objects of \(\underline{C}_0\), so it is locally an object of \(\underline{C}_0\), whence the assertion.

\sgaStableTarget{sga6.e01.cor.4-14}\textbf{Corollary 4.14.} \uline{Let} \(F\in\Ob\underline{C}_S\) \uline{and} \(n\in\mathbb N\)\uline{. Equivalent conditions}:
\begin{enumerate}[label=(\roman*)]
\item \(\underline{\perfamp}(F)\subset[-n,0]\);
\item \(F\) \uline{locally admits a left resolution of length} \(n\) \uline{by objects of} \(\underline{C}_0\)\uline{, i.e.\ the set of} \(X\) \uline{over} \(S\) \uline{for which there exists an exact sequence}
\[
  0\longrightarrow L^{-n}\longrightarrow L^{-n+1}\longrightarrow \ldots
  \longrightarrow L^0\longrightarrow F_X\longrightarrow 0,
\]
\uline{with} \(L^i\in\Ob\underline{C}_{0X}\) \uline{for every} \(i\), \uline{is a refinement of} \(S\).
\end{enumerate}

\uline{Proof.} This is an immediate consequence of (4.7) and (4.13).

\sgaStableTarget{sga6.e01.def.4-14-1}\textbf{Definition 4.14.1.} \uline{We shall say that} \(F\) \uline{has perfect dimension} \(\leq n\)\uline{, and we shall write} \(\perfdim(F)\leq n\)\uline{, when} \(F\) \uline{satisfies the equivalent conditions of} (4.14)\uline{. We shall call the smallest integer} \(n\) \uline{such that} \(\underline{\perfdim}(F)\leq n\) \uline{the} perfect dimension \uline{of} \(F\) \uline{(when there is no such integer, we shall say that} \(F\) \uline{has infinite perfect dimension).}

When \(\underline{S}\) is the punctual site and \(\underline{C}_0\) is the full subcategory of \(\underline{C}\) formed by the projective objects, perfect dimension is the usual projective dimension.

\sgaStableTarget{sga6.e01.prop.4-15}\textbf{Proposition 4.15.} \uline{Let} \(a,b,n\in\mathbb Z\) \uline{(}\(a\leq b\)\uline{,} \(n\geq0\)\uline{), and let} \(F\in\Ob\Dpx(\underline{C})\) \uline{be such that} \(F^i=0\) \uline{(resp.}\ \(\Ho^i(F)=0\)\uline{) for} \(i\notin[a,b]\)\uline{. If, for every} \(i\in[a,b]\)\uline{,} \(F^i\) \uline{(resp.}\ \(\Ho^i(F)\)\uline{) has perfect dimension} \(\leq n+(i-a)\)\uline{, then} \(F\) \uline{has perfect amplitude contained in} \([a-n,b]\)\uline{. Consequently, if for every} \(i\in[a,b]\)\uline{,} \(F^i\) \uline{(resp.}\ \(\Ho^i(F)\)\uline{) is perfect,} \(F\) \uline{is too.}

%%%%%%%%%%  scan idx 141  |  folio 128  %%%%%%%%%%
\SGAPrintPage{128}%
\uline{Proof.} Proceed by induction on \(b\) (with \(a\) and \(n\) fixed), truncating \(F\) and applying (4.11) (cf.\ the proof of (2.11)).

\sgaStableTarget{sga6.e01.lem.4-16}\textbf{Lemma 4.16.} \uline{Let} \(a,b\in\mathbb Z\) \uline{(}\(a\leq b\)\uline{), and} \(F\in\Ob\Dpx(\underline{C})\) \uline{be such that} \(F^i=0\) \uline{for} \(i\notin[a,b]\)\uline{. Assume that} \(\underline{\perfamp}(F)\subset[a,b]\)\uline{, and that} \(F^i\in\Ob\underline{C}_0\) \uline{for} \(i>a\)\uline{. Then} \(F^a\) \uline{is locally an object of} \(\underline{C}_0\)\uline{.}

\uline{Proof.} We have an exact sequence (obtained by truncating \(F\))
\[
  0\longrightarrow F'\longrightarrow F\longrightarrow F^a[-a]\longrightarrow 0,
\]
with \(\perfamp(F')\subset[a+1,b+1]\). From (4.11) and (4.13) we deduce that \(\perfdim(F^a)=0\), i.e.\ that \(F^a\) is locally an object of \(\underline{C}_0\).

\sgaStableTarget{sga6.e01.prop.4-17}\textbf{Proposition 4.17.} \uline{Let} \(F',F''\in\Ob\Dpx(\underline{C}_S)\)\uline{, and} \(a,b\in\mathbb Z\)\uline{, with} \(a\leq b\)\uline{. Consider the conditions}:
\begin{enumerate}[label=(\roman*)]
\item \(\perfamp(F')\subset[a,b]\) and \(\perfamp(F'')\subset[a,b]\);
\item \(\perfamp(F'+F'')\subset[a,b]\).
\end{enumerate}
\uline{We have} (i) \(\Rightarrow\) (ii)\uline{, and, when} \(\underline{C}_0\) \uline{is locally stable under direct factors (}i.e.\ \uline{the relation} \(E'+E''\in\Ob\underline{C}_0\)\uline{, for} \(E',E''\in\Ob\underline{C}_X\) \uline{and} \(X\in\Ob \underline{S}\)\uline{, implies that locally} \(E'\) \uline{and} \(E''\) \uline{are objects of} \(\underline{C}_0\)\uline{), we have} (i) \(\Leftrightarrow\) (ii).

\uline{Proof.} The implication (i) \(\Rightarrow\) (ii) is a special case of (4.11). Let us prove the converse. Put \(F=F'+F''\). Since \(F\) is pseudo-coherent, \(F'\) and \(F''\) are too (2.12). On the other hand, the relation \(\Ho^i(F)=0\) for \(i\notin[a,b]\) implies \(\Ho^i(F')=\Ho^i(F'')=0\) for \(i\notin[a,b]\). By (2.10 a)), we can therefore suppose (after localizing) that \(F'^i\) (resp.\ \(F''^i\)) \(\in\Ob\underline{C}_0\) for \(i\geq a+1\) and \(F'^i=F''^i=0\) for \(i>b\), then, by truncating, that
%%%%%%%%%%  scan idx 142  |  folio 129  %%%%%%%%%%
\SGAPrintPage{129}%
\(F'^i=F''^i=0\) for \(i<a\). Since \var{\(\perfdim(F)\subset[a,b]\)}{\(\perfamp(F)\subset[a,b]\)}{EN-VAR-103}, we deduce from (4.16) that \(F^a\) is locally an object of \(\underline{C}_0\). But \(F^a=F'^a+F''^a\), so \(F'^a\) and \(F''^a\) are locally objects of \(\underline{C}_0\), since \(\underline{C}_0\) is locally stable under direct factors.

\sgaStableTarget{sga6.e01.prop.4-18}\textbf{Proposition 4.18.} \uline{Let} \((\underline{S}',\underline{C}',\underline{C}'_0)\) \uline{be a triple satisfying the same hypotheses} (4.0) \uline{as} \((\underline{S},\underline{C},\underline{C}_0)\)\uline{. Let} \(f:\underline{S}'\to \underline{S}\) \uline{be a morphism of sites (defined by} \(f^{-1}:\underline{S}\to \underline{S}'\)\uline{), and let}
\(u:\Dpx^*(\underline{C})\longrightarrow \Dpx^*(\underline{C}')\qquad (*=-,+,\varnothing)\)
\uline{be a functor over} \(f\) (cf.\ (2.13)). \uline{If} \(u\) \uline{takes objects of} \(\underline{C}_0\) \uline{to} \(\underline{C}'_0\)\uline{-perfect complexes, then} \(u\) \uline{defines a functor}
\(\Dpx^*(\underline{C})_{\underline{C}_0\text{-}\mathrm{parf}}\longrightarrow\Dpx^*(\underline{C}')_{\underline{C}'_0\text{-}\mathrm{parf}}\).

\uline{Proof.} Let \(F\in\Ob\Dpx^*(\underline{C})_{\underline{C}_0\text{-}\mathrm{parf}}\); let us show that \(u(F)\in\Ob\Dpx^*(\underline{C}')_{\underline{C}'_0\text{-}\mathrm{parf}}\). Since the question is local, we can suppose that \(F\) is strictly perfect. We need only use dévissage, taking into account that \(\Dpx^*(\underline{C}')_{\underline{C}'_0\text{-}\mathrm{parf}}\) is a triangulated subcategory of \var{\(\Dpx^*(\underline{C})\) (4.10).}{\(\Dpx^*(\underline{C}')\) (4.10).}{EN-VAR-104}

\sgaStableTarget{sga6.e01.rem.4-18-1}\textbf{Remark 4.18.1.} Suppose that there exist \(m,n\in\mathbb Z\) (\(m\leq n\)) such that \(F\in\Ob\underline{C}_0\) implies \(\underline{C}'_0\text{-}\perfamp(u(F))\subset[m,n]\). Then, for \(F\in\Ob\Dpx^*(\underline{C})\), we have
\[
  \underline{C}_0\text{-}\perfamp(F)
  \subset[a,b]
  \quad\Rightarrow\quad
  \underline{C}'_0\text{-}\perfamp(u(F))
  \subset[a+m,b+n].
\]
The proof is left as an exercise to the reader (use (4.11)).

\sgaStableTarget{sga6.e01.cor.4-19}\textbf{Corollary 4.19.} \uline{Under the hypotheses of} (2.16), \uline{if} \(E'\in\Ob\Dpx^-({}_{A'} \underline{S}'_A)\) \uline{is perfect as an object of} \(\Dpx^-({}_{A'} \underline{S}')\) \uline{and} \(E\in\Ob\Dpx^-({}_A \underline{S})\) \uline{is perfect, then}
\(E'\overset{\mathrm L}{\otimes}_A E\) \uline{is perfect}.

\uline{Proof.} We can suppose that \(E'\in\Ob\Dpx^-({}_{A'} S'_A)\) (\(S'\)= final object of \(\underline{S}'\)), and apply (4.18) to
\(u=E'\overset{\mathrm L}{\otimes}_A:\Dpx^-({}_A \underline{S})\longrightarrow \Dpx^-({}_{A'} \underline{S}')\).
In particular:

%%%%%%%%%%  scan idx 143  |  folio 130  %%%%%%%%%%
\SGAPrintPage{130}%
\sgaStableTarget{sga6.e01.cor.4-19-1}\textbf{Corollary 4.19.1.} a) \uline{The functor}
\(\mathbb L f^*:\Dpx^-({}_A \underline{S})\longrightarrow \Dpx^-({}_{A'} \underline{S}')\)
\uline{induces a functor}
\[
  \mathbb L f^*:\Dpx^-({}_A \underline{S})_{\mathrm{parf}}
  \longrightarrow
  \Dpx^-({}_{A'} \underline{S}')_{\mathrm{parf}}.
\]

b) \uline{If} \(A\) \uline{is commutative, the functor}
\(\overset{\mathrm L}{\otimes}_A:\Dpx^-(\underline{S})\times_{\underline{S}}\Dpx^-(\underline{S})\longrightarrow \Dpx^-(\underline{S})\)
\uline{induces a functor}
\[
  \overset{\mathrm L}{\otimes}_A:
  \Dpx^-(\underline{S})_{\mathrm{parf}}\times_{\underline{S}}\Dpx^-(\underline{S})_{\mathrm{parf}}
  \longrightarrow \Dpx^-(\underline{S})_{\mathrm{parf}}.
\]

\sgaStableTarget{sga6.e01.rem.4-19-2}\textbf{Remark 4.19.2.} We can make (4.19) more precise in terms of perfect amplitude:
\[
  \perfamp(E)\subset[a,b]\quad\text{and}\quad
  \perfamp(E')\subset[a',b']
  \quad\Rightarrow\quad
  \perfamp(E'\overset{\mathrm L}{\otimes}_A E)\subset[a+a',b+b'].
\]
In particular:
\[
  \perfamp(E)\subset[a,b]
  \quad\Rightarrow\quad
  \perfamp(\mathbb L f^*(E))\subset[a,b].
\]

\sgaStableTarget{sga6.e01.sec.5}\section*{5. Finite Tor-dimension and perfection}

\sgaStableTarget{sga6.e01.para.5-0}5.0. The reader will surely have asked: what does a pseudo-coherent complex \(E\) lack in order to be perfect? Is it enough, for example, that \(E\) have bounded cohomology? The answer is no. Indeed, take for \(\underline{C}\) the category of modules over a commutative noetherian ring \(A\), and for \(\underline{C}_0\) the full subcategory of projective \(A\)-modules of finite type. Any \(A\)-module of finite type \(E\) is pseudo-coherent and has bounded cohomology, but \(E\) is perfect only when \(E\) also has finite projective dimension (or, equivalently, finite tor-dimension). This example suggests that the difference between a pseudo-coherent complex and a perfect complex lies in a question of cohomological dimension. This is what we shall now examine
%%%%%%%%%%  scan idx 144  |  folio 131  %%%%%%%%%%
\SGAPrintPage{131}%
in the case of the category of Modules of a ringed topos (cf.\ (3.3)). We shall need some preliminary sorites on the notion of finite tor-dimension. We shall be brief, since this question already appears in the literature: cf.\ ([H] II 4) and (SGA 4 XVII 4.1.9.).

\sgaStableTarget{sga6.e01.prop.5-1}\textbf{Proposition 5.1.} \uline{Let} \((\underline{S},A)\) \uline{be a ringed topos,} \(X\in\Ob \underline{S}\)\uline{,} \(F\in\Ob\Dpx^-({}_A X)\)\uline{,} \(m,n\in\mathbb Z\) \uline{be such that} \(m\leq n\)\uline{. Equivalent conditions}:
\begin{enumerate}[label=(\roman*)]
\item \(F\) \uline{is isomorphic, in} \(\Dpx^-({}_A X)\)\uline{, to a complex} \(F'\) \uline{such that} \(F'^i\) \uline{is flat for every} \(i\) \uline{and zero for} \(i\notin[m,n]\);
\item \uline{for every right} \(A_X\)\uline{-Module} \(E\)\uline{, we have} \(\underline{\Ho}^i(E\overset{\mathrm L}{\otimes}_A F)=0\) \uline{for} \(i\notin[m,n]\).
\end{enumerate}

\uline{Proof.} The implication (i) \(\Rightarrow\) (ii) is trivial; the implication (ii) \(\Rightarrow\) (i) follows, by truncation, from the

\sgaStableTarget{sga6.e01.lem.5-1-1}\textbf{Lemma 5.1.1.} (cf.\ (4.16)). \uline{Let} \(F\in\Ob\Dpx^-({}_A X)\)\uline{. Assume that} \(F\) \uline{satisfies hypothesis} (ii) of (5.1)\uline{, that} \(F^i=0\) \uline{for} \(i\notin[m,n]\)\uline{, and that} \(F^i\) \uline{is flat for} \(i\ne m\)\uline{. Then} \(F^m\) \uline{is flat.}

\uline{Proof.} We have a distinguished triangle
\[
\begin{tikzcd}
& F^m[-m] \arrow[dl,"+1"'] & \\
F' \arrow[rr] && F \arrow[ul] & ,
\end{tikzcd}
\]
where \(F'\) has flat components that vanish outside the interval \([m+1,n]\) (suppose \(m<n\)). Apply the functor \(E\overset{\mathrm L}{\otimes}_A.\) for \(E\) an arbitrary right \(A_X\)-Module, and write the exact cohomology sequence.

%%%%%%%%%%  scan idx 145  |  folio 132  %%%%%%%%%%
\SGAPrintPage{132}%
\sgaStableTarget{sga6.e01.def.5-2}\textbf{Definition 5.2.} \uline{An object} \(F\) \uline{of} \(\Dpx^-({}_A \underline{S})\) \uline{satisfying the equivalent conditions of} (5.1) \uline{will be said to have} flat amplitude \uline{contained in the interval} \([m,n]\)\uline{, and we shall write}
\[
  \toramp(F)\subset[m,n].
\]
\uline{We shall say that} \(F\) \uline{has finite flat amplitude (or finite tor-dimension) when there exist integers} \(m\) \uline{and} \(n\) \uline{such that the preceding condition holds.}

\sgaStableTarget{sga6.e01.para.5-3}5.3. Let \(X\in\Ob \underline{S}\). We shall denote by \(\Dpx^-({}_A X)_{\torf}\) the full subcategory of \(\Dpx^-({}_A X)\) formed by the objects of finite tor-dimension. By criterion (ii) of (5.1), it is a triangulated subcategory of \(\Dpx^-({}_A X)\). As \(X\) varies, the categories \(\Dpx^-({}_A X)_{\torf}\) form a triangulated sub-\(\underline{S}\)-category (1.1.4) of \(\Dpx^-({}_A \underline{S})\), which we shall denote by \(\Dpx^-({}_A \underline{S})_{\torf}\). The functor \(\overset{\mathrm L}{\otimes}_A\) induces an exact \(\underline{S}\)-functor
\[
  \overset{\mathrm L}{\otimes}_A:
  \Dpx^b(\underline{S}_A)\times_{\underline{S}} \Dpx({}_A \underline{S})_{\torf}
  \longrightarrow \Dpx^b(\underline{S}_{\mathbb Z}).
\]

\sgaStableTarget{sga6.e01.para.5-4}5.4. For a left \(A\)-Module \(F\) to have flat amplitude contained in \([-n,0]\) (\(n\in\mathbb N\)), it is necessary and sufficient that \(F\) have tor-dimension \(\leq n\) in the usual sense, i.e.\ that \(F\) admit a left resolution of length \(n\) by flat \(A\)-Modules.

%%%%%%%%%%  scan idx 146  |  folio 133  %%%%%%%%%%
\SGAPrintPage{133}%
\sgaStableTarget{sga6.e01.para.5-5}5.5. Let \(X\in\Ob \underline{S}\), and \(m,n\in\mathbb Z\) (\(m\leq n\)). For \(F\in\Ob\Dpx^-({}_A X)\) to have flat amplitude contained in \([m,n]\), it is necessary and sufficient that there exist a covering family \((X_i\to X)\) such that \(F|X_i\) has flat amplitude contained in \([m,n]\). In particular, when \(\underline{S}_{/X}\) is of finite type (SGA 4 VI 1.7), and \(F\) is locally of finite tor-dimension, \(F\) has finite tor-dimension.

\sgaStableTarget{sga6.e01.prop.5-6}\textbf{Proposition 5.6.} \uline{Under the hypotheses of} (2.16)\uline{, if} \(E'\) \uline{has flat amplitude contained in} \([m',n']\) \uline{as an object of} \(\Dpx^-({}_{A'} \underline{S}')\) \uline{and} \(E\) \uline{has flat amplitude contained in} \([m,n]\)\uline{, then} \(E'\overset{\mathrm L}{\otimes}_A E\) \uline{has flat amplitude contained in} \([m+m',n+n']\)\uline{.}

\uline{Proof.} This is immediate: use form (5.1 (i)) of the notion of finite tor-dimension and the fact that, when \(E\) is a flat left \(A\)-Module, \(f^{-1}(E)\) is a flat left \(f^{-1}(A)\)-Module (SGA 4 IV App.\ II).

In particular:

\sgaStableTarget{sga6.e01.cor.5-6-1}\textbf{Corollary 5.6.1.} a) \uline{The functor} \(\mathbb L f^*:\Dpx^-({}_A \underline{S})\longrightarrow \Dpx^-({}_{A'} \underline{S}')\) \uline{induces a functor}
\[
  \mathbb L f^*:\Dpx({}_A \underline{S})_{\torf}\longrightarrow \Dpx({}_{A'} \underline{S}')_{\torf}.
\]

b) \uline{If} \(A\) \uline{is commutative, the functor} \(\overset{\mathrm L}{\otimes}_A:\Dpx^-(\underline{S})\times_{\underline{S}}\Dpx^-(\underline{S})\longrightarrow \Dpx^-(\underline{S})\) \uline{induces a functor}
\[
  \overset{\mathrm L}{\otimes}_A:
  \Dpx(\underline{S})_{\torf}\times_{\underline{S}}\Dpx(\underline{S})_{\torf}
  \longrightarrow \Dpx(\underline{S})_{\torf}.
\]

\sgaStableTarget{sga6.e01.rem.5-6-2}\textbf{Remark 5.6.2.} Let \(E\in\Ob\Dpx^-({}_A S)\) (\(S\)= final object of \(\underline{S}\)), and \(m,n\in\mathbb Z\) (\(m\leq n\)). When \(E\) has flat amplitude contained in \([m,n]\), so does \(E_s\), by (5.6), for every point \(s\) of \(\underline{S}\). The converse is true when \(\underline{S}\) has enough points (use form (5.1.(ii)) of the notion of flat
%%%%%%%%%%  scan idx 147  |  folio 134  %%%%%%%%%%
\SGAPrintPage{134}%
amplitude, and the fact that \(\overset{\mathrm L}{\otimes}_A\) commutes with passage to the fibres).

\sgaStableTarget{sga6.e01.exercise.5-7}\textbf{Exercises 5.7.} a) Suppose that \(A\) is commutative and \((\underline{S}',A')\) faithfully flat over \((\underline{S},A)\) (cf.\ (2.16.2)). For \(E\in\Ob\Dpx^-({}_A S)\)\SGAAuthorFootnote{1}{\(S\) = final object of \(\underline{S}\)} to have flat amplitude contained in \([m,n]\), it is necessary and sufficient that \(f^*(E)\) have that property.

b) Let \(F\in\Ob\Dpx^+({}_A S)\)\SGAAuthorFootMark{1}. Show the equivalence of the conditions:
\begin{enumerate}[label=(\roman*)]
\item \(F\) is isomorphic, in \(\Dpx^+({}_A S)\), to a complex \(F'\) such that \(F'^i\) is injective for every \(i\) and zero for \(i\notin[m,n]\);
\item for every left \(A_S\)-Module \(E\), we have \(\Ext^i(E,F)=0\) for \(i\notin[m,n]\).
\end{enumerate}
When \(F\) satisfies the preceding conditions, we say that \(F\) has injective amplitude contained in \([m,n]\). The full subcategory of \(\Dpx^+({}_A S)\) formed by the objects of finite injective amplitude is a triangulated subcategory, denoted by \(\Dpx^+({}_A S)_{\injf}\). Show that, when \(A\) is commutative, the functor \(\mathbb R\underline{\Hom}:\Dpx(S)^\circ\times\Dpx^+(S)\longrightarrow \Dpx(S)\) induces a functor
\[
  \mathbb R\underline{\Hom}:\Dpx(S)^\circ_{\torf}\times\Dpx(S)_{\injf}
  \longrightarrow \Dpx(S)_{\injf}.
\]

We can now answer the question of (5.0) and state the main result of this section.

\sgaStableTarget{sga6.e01.thm.5-8}\textbf{Theorem 5.8.} \uline{Let} \((\underline{S},A)\) \uline{be a ringed topos,} \(E\in\Ob\Dpx({}_A \underline{S})\)\uline{,} \(m,n\in\mathbb Z\)\uline{, with} \(m\leq n\)\uline{. Equivalent conditions}:
\begin{enumerate}[label=(\roman*)]
\item \(\perfamp(E)\subset[m,n]\) (4.8);
\item \(E\) \uline{is} \((m-1)\)\uline{-pseudo-coherent} (2.15) \uline{and we have} \(\toramp(E)\subset[m,n]\) (5.2).
\end{enumerate}
\uline{Since} ``\uline{perfect}'' \uline{obviously implies} ``\uline{pseudo-coherent}''\uline{, we deduce}:

%%%%%%%%%%  scan idx 148  |  folio 135  %%%%%%%%%%
\SGAPrintPage{135}%
\sgaStableTarget{sga6.e01.cor.5-8-1}\textbf{Corollary 5.8.1.} \uline{For} \(E\) \uline{to be perfect} (4.8)\uline{, it is necessary and sufficient that} \(E\) \uline{be pseudo-coherent} \var{(2,15)}{(2.15)}{EN-VAR-105} \uline{and locally of finite tor-dimension} (5.2)\uline{.}

For the proof, we shall need some lemmas.

\sgaStableTarget{sga6.e01.lem.5-8-2}\textbf{Lemma 5.8.2.} \uline{Let} \(B\) \uline{be a ring of} \(\underline{S}\)\uline{. If} \(E\) \uline{is a left} \(A\)\uline{-Module,} \(G\) \uline{a right} \(B\)\uline{-Module, and} \(F\) \uline{an} \(A\)\uline{-}\(B\)\uline{-Bimodule (on the left over} \(A\)\uline{, on the right over} \(B\)\uline{), there is a functorial canonical homomorphism}
\sgaStableTarget{sga6.e01.eq.5-8-2-1}\begin{equation}\tag{5.8.2.1}
  \underline{\Hom}_B(F,G)\otimes_A E
  \longrightarrow
  \underline{\Hom}_B(\underline{\Hom}_A(E,F),G).
\end{equation}
\uline{This is an isomorphism for} \(G\) \uline{injective and} \(E\) \uline{of finite presentation}.

\uline{Proof.} Let us first define the arrow. By the adjunction formula dear to Cartan, it amounts to defining a homomorphism of \(B\)-Modules:
\sgaStableTarget{sga6.e01.eq.unnamed.occ5}\begin{equation}\tag{*}
  \underline{\Hom}_B(F,G)\otimes_A E\otimes_{\mathbb Z_S}\underline{\Hom}_A(E,F)
  \longrightarrow G\text{\SGAAuthorFootMark{1}}.
\end{equation}
\SGAAuthorFootText{1}{we write \(S\) for the final object of \(\underline{S}\) and \(\mathbb Z_S\) for the constant sheaf with value \(\mathbb Z\).}
Now we have a canonical homomorphism of \(A\)-\(B\)-biModules
\sgaStableTarget{sga6.e01.eq.unnamed.occ6}\begin{equation}\tag{**}
  E\otimes_{\mathbb Z_S}\underline{\Hom}_A(E,F)
  \longrightarrow F
\end{equation}
(coming by adjunction from the identity arrow of \(\underline{\Hom}_A(E,F)\)). Likewise, we have a canonical homomorphism of \(B\)-Modules
\sgaStableTarget{sga6.e01.eq.unnamed.occ7}\begin{equation}\tag{***}
  \underline{\Hom}_B(F,G)\otimes_A F\longrightarrow G.
\end{equation}
We then define (*) as the composite of (**) and (***). To see that the arrow (5.8.2.1) is an isomorphism when \(G\) is injective and \(E\) of finite presentation, it suffices to observe that, when \(G\) is injective, the source and target of the arrow are right-exact functors in \(E\), and that the arrow induces the identity for \(E=A\).

%%%%%%%%%%  scan idx 149  |  folio 136  %%%%%%%%%%
\SGAPrintPage{136}%
Here is now a special case of (5.8).

\sgaStableTarget{sga6.e01.lem.5-8-3}\textbf{Lemma 5.8.3.} \uline{Let} \(E\) \uline{be a left} \(A\)-\uline{Module}. \uline{Equivalent conditions}:
\begin{enumerate}[label=(\roman*)]
\item \(E\) \uline{is locally a direct factor of a free Module of finite type};
\item \(E\) \uline{is flat and of finite presentation}.
\end{enumerate}

\uline{Proof} (after BOURBAKI Alg. Comm. chap.\ I, \S\ 2 Exer 15). It is clear that (i) implies (ii). Let us prove (ii) \(\Rightarrow\) (i). By (1.3.1), we must show that the functor \(\underline{\Hom}_A(E,.)\) is exact. Let \(F\longrightarrow F''\longrightarrow 0\) be an exact sequence of \(A\)-Modules (on the left). We must show that the sequence
\[
  \underline{\Hom}_A(E,F)\longrightarrow \underline{\Hom}_A(E,F'')\longrightarrow 0
\]
is exact. For this it suffices to prove that, for every injective \(\mathbb Z_S\)-Module \(I\), the sequence obtained by applying the functor \(\underline{\Hom}_{\mathbb Z}(.,I)\) is exact. Now we have a commutative diagram, in which the vertical arrows are the canonical arrows of \var{(3.17.1)}{(5.8.2.1)}{EN-VAR-106}:
\[
\begin{array}{ccccccccc}
0&\to&\underline{\Hom}_{\mathbb Z_S}(F'',I)\otimes_A E&\to&\underline{\Hom}_{\mathbb Z_S}(F,I)\otimes_A E\\
&&\downarrow&&\downarrow\\
0&\to&\underline{\Hom}_{\mathbb Z_S}(\underline{\Hom}_A(E,F''),I)&\to&\underline{\Hom}_{\mathbb Z_S}(\underline{\Hom}_A(E,F),I).
\end{array}
\]
The vertical arrows are isomorphisms by (5.8.2). The top row is exact because \(E\) is flat. Thus the bottom row is exact, as required.

\uline{Proof of} (5.8). The implication (i) \(\Rightarrow\) (ii) follows trivially from the definitions. Let us prove (ii) \(\Rightarrow\) (i). Since the question is local, we can suppose that \(E\) is strictly \((m-1)\)-pseudo-coherent. Since, on the other hand, \(E\) is acyclic in degree \(<m\), the
%%%%%%%%%%  scan idx 150  |  folio 137  %%%%%%%%%%
\SGAPrintPage{137}%
truncation arrow
\[
\begin{array}{ccccccc}
&\to&E^{m-1}&\to&E^m&\to&E^{m+1}\to\ldots\\
&&\downarrow&&\downarrow&&\downarrow\\
&\to&0&\to&E^m/B^m&\to&E^{m+1}\to\ldots
\end{array}
\]
is a quasi-isomorphism. Denote the truncated complex by \(E'\). By construction, \(E'\) has flat amplitude \(\subset[m,n]\), and \(E'^i\) is free of finite type for \(i>m\). Thus, by (5.1.1), \(E'^m\) is flat. On the other hand, we have a distinguished triangle
\[
\begin{tikzcd}
& E'^m[-m] \arrow[dl,"+1"'] & \\
E'' \arrow[rr] && E' \arrow[ul] & ,
\end{tikzcd}
\]
where \(E''\) is strictly perfect. By (2.6), \(E'^m[-m]\) is \((m-1)\)-pseudo-coherent, i.e.\ (2.9) \(E'^m\) is of finite presentation. By (5.8.3), we conclude that \(E'^m\) is locally a direct factor of a free Module of finite type, hence that \(\operatorname{parf.amp}(E)\subset[m,n]\).

\sgaStableTarget{sga6.e01.rem.5-8-4}\textbf{Remarks 5.8.4.} a) Using (2.5), we recover the fact that perfect complexes form a triangulated subcategory of \(\Dpx({}_A \underline{S})\).

b) A perfect complex, even with bounded cohomology, does not necessarily have finite tor-dimension. It does, however, when \(\underline{S}\) is of finite type (5.5). Denote by \(\Dpx({}_A \underline{S})_{\underline{\mathrm{parf}}}\) the full sub-\(\underline{S}\)-category of \(\Dpx^b({}_A \underline{S})\) formed by the perfect complexes of finite tor-dimension. It follows at once from (2.16.1 b)) and (5.3) that, when \(A\) is commutative, the functor \(\overset{\mathrm L}{\otimes}_A\) induces a functor
\[
  \overset{\mathrm L}{\otimes}_A:
  \Dpx(\underline{S})_{\underline{\mathrm{parf}}}\times_{\underline{S}} \Dpx^b(\underline{S})_{\mathrm{coh}}
  \longrightarrow \Dpx^b(\underline{S})_{\mathrm{coh}}.
\]
%%%%%%%%%%  scan idx 151  |  folio 138  %%%%%%%%%%
\SGAPrintPage{138}%
This formula is the starting point of the theory of intersections on schemes. Observe, on the other hand, that the (derived) tensor product of two pseudo-coherent complexes with bounded cohomology does not in general have bounded cohomology.

\sgaStableTarget{sga6.e01.cor.5-9}\textbf{Corollary 5.9.} \uline{Suppose that} \(\underline{S}=(\mathrm{Ens})\) \uline{(so} \(A\) ``\uline{is}'' \uline{an ordinary ring) and that every left} \(A\)\uline{-module has finite tor-dimension (this holds, for example, when} \(A\) \uline{has finite global cohomological dimension, see} (MULT VI 2.6) \uline{and} (EGA \(0_{\mathrm{IV}}\) 17.2.8)). \uline{Then we have}
\[
\begin{aligned}
\text{(i)}\quad & \Dpx^b({}_A \underline{S})=\Dpx({}_A \underline{S})_{\torf},\quad\text{\uline{and}}\\
\text{(ii)}\quad & \Dpx^b({}_A \underline{S})_{\mathrm{coh}}=\Dpx({}_A \underline{S})_{\underline{\mathrm{parf}}}.
\end{aligned}
\]
\uline{In other words, for} \(E\in\Ob\Dpx({}_A \underline{S})\) \uline{to have finite tor-dimension (resp.\ finite tor-dimension and be perfect), it is necessary and sufficient that} \(E\) \uline{have bounded cohomology (resp.\ bounded cohomology and be pseudo-coherent)}.

\uline{Proof}. By (5.8), it suffices to prove (i). Let \(E\in\Ob\Dpx^b({}_A \underline{S})\). We can suppose (extending \(E\) by \(0\)) that \(E\in\Ob\Dpx^b({}_A S)\), where \(S\) is the final object of \(\underline{S}\). Proceeding by induction on the number of nonzero cohomology objects of \(E\), we reduce by dévissage to the case where \(E\) is an ordinary \(A\)-Module, and we are done.

More generally, we can state:

\sgaStableTarget{sga6.e01.cor.5-10}\textbf{Corollary 5.10} \uline{Let} \(X\in\Ob \underline{S}\)\uline{. Suppose that} \(\underline{S}_{/X}\) \uline{has enough points} (SGA 4 IV 6.4.1.) \uline{and that, for every point} \(s\) \uline{of} \(\underline{S}_{/X}\)\uline{, every left} \(A_s\)\uline{-module has finite tor-dimension. Then we have}
\[
  \Dpx^b({}_A X)_{\mathrm{coh}}=\Dpx^b({}_A X)_{\mathrm{parf}}.
\]
%%%%%%%%%%  scan idx 152  |  folio 139  %%%%%%%%%%
\SGAPrintPage{139}%
In other words, for \(E\in\Ob\Dpx^b({}_A X)\) to be pseudo-coherent, it is necessary and sufficient that \(E\) be perfect.

\uline{Proof}. Let \(E\in\Ob\Dpx^b({}_A X)_{\mathrm{coh}}\) and \(s\) be a point of \(\underline{S}_{/X}\). We must show that \(E\) is perfect ``in a neighbourhood of \(s\)'', i.e.\ that there exists an \(s\)-pointed object \(U\) of \(\underline{S}\) over \(X\) such that \(E|U\) is perfect. Since the functor \(s^*\) is exact, by (2.16.1 a)) we have \(E_s\in\Ob\Dpx^b(A_s)_{\mathrm{coh}}\). Thus, by (5.9), \(E_s\) is perfect. The assertion will therefore follow from (3.7.2) and the following lemma:

\sgaStableTarget{sga6.e01.lem.5-10-1}\textbf{Lemma 5.10.1.} \uline{Let} \(s\) \uline{be a point of} \(\underline{S}\) \uline{and let} \(F\in\Ob\Dpx(A_s)_{\mathrm{parf}}\)\uline{. There exist an} \(s\)\uline{-pointed object} \(U\) \uline{of} \(\underline{S}\) \uline{and} \(F'\in\Ob\Dpx({}_A U)_{\mathrm{parf}}\) \uline{such that} \(F'_s=F\)\uline{.}

\uline{Proof}. This is evident when \(F\) is a free \(A_s\)-module of finite type. Now suppose that \(F\) is projective of finite type, hence the image of a projector \(p:L\longrightarrow L\), where \(L\) is free of finite type. There then exist a ``neighbourhood'' \(U\) of \(s\), a free \(A_U\)-Module \(L'\), and a homomorphism \(p':L'\to L'\) such that \(p'_s=p\). Since \({p'_s}^2=p'_s\), we can suppose, after shrinking \(U\), that \(p'^2=p'\). Thus \(F'=\Im(p')\) is a direct factor of \(L'\) and satisfies \(F'_s=F\). This generalizes immediately to the case where \(F\) is strictly perfect, i.e.\ a complex with bounded degrees made up of projective modules of finite type.

\sgaStableTarget{sga6.e01.ex.5-11}\textbf{Examples 5.11.} Corollary (5.10) applies in particular when \((\underline{S},A)\) is defined by a ringed space with \uline{regular} local rings, for example a regular scheme, or a complex analytic space smooth over \(\mathbb C\).

%%%%%%%%%%  scan idx 153  |  folio 140  %%%%%%%%%%
\SGAPrintPage{140}%
\sgaStableTarget{sga6.e01.sec.6}\section*{6. Rank of a perfect complex}

\sgaStableTarget{sga6.e01.para.6-0}6.0. Let \((X,\underline{O}_X)\) be a ringed space with local rings. It is well known that to every locally free sheaf of finite type \(L\) on an open set \(U\) of \(X\) we can attach a function \(\rg_U(L)\), locally constant on \(U\) and with integer values, called the \uline{rank} of \(L\), the family of functions \(\rg_U(L)\) being characterized by the following properties:
\begin{enumerate}[label=(\roman*)]
\item for an inclusion \(V\subset U\) of open sets, we have \(\rg_V(L|V)=\rg_U(L)|V\);
\item for every exact sequence of locally free \(\underline{O}_U\)-Modules of finite type
\[
  0\longrightarrow L'\longrightarrow L\longrightarrow L''\longrightarrow 0,
\]
we have \(\rg_U(L)=\rg_U(L')+\rg_U(L'')\);
\item \(\rg_X(\underline{O}_X)=1\).
\end{enumerate}
In this section, we propose to extend the notion of rank to perfect complexes.

\sgaStableTarget{sga6.e01.def.6-1}\textbf{Definition 6.1} (cf.\ SGA 5 VIII 2). \uline{Let} \(C\) \uline{be an additive (resp.\ triangulated) category and} \(F\) \uline{an abelian group. A function} \(f:\Ob(C)\longrightarrow F\) \uline{is called an} additive function from \(C\) to \(F\) \uline{when for all} \(L,L',L''\in\Ob(C)\) \uline{such that} \(L\simeq L'+L''\) \uline{(resp.\ there is a distinguished triangle}
\[
\begin{tikzcd}
& L'' \arrow[dl, dashed] & \\
L' \arrow[rr] && L \arrow[ul]
\end{tikzcd}
\]
) \uline{we have} \(f(L)=f(L')+f(L'')\).

\uline{The additive functions from} \(C\) \uline{to} \(F\) \uline{form an abelian group, denoted by} \(\Add(C,F)\).

%%%%%%%%%%  scan idx 154  |  folio 141  %%%%%%%%%%
\SGAPrintPage{141}%
\sgaStableTarget{sga6.e01.para.6-2}6.2. Let \(C\) be a triangulated category and \(f\) an additive function from \(C\) to an abelian group \(F\). Then \(f\) induces an additive function on the additive category \(C^{\mathrm{ad}}\) underlying \(C\), hence an inclusion \(\Add(C,F)\subset\Add(C^{\mathrm{ad}},F)\). In particular, we have \(f(0)=0\) and \(f(L)=f(M)\) when \(L\simeq M\). On the other hand, we have \(f(L[n])=(-1)^n f(L)\) for every \(L\in\Ob(C)\) and \(n\in\mathbb Z\).

\sgaStableTarget{sga6.e01.para.6-3}6.3. Let \(C\) be an additive (resp.\ triangulated) category. The group \(\Add(C,F)\) depends functorially on \(F\). The functor \(\Add(C,.)\) is representable (cf.\ SGA 5 VIII 2) by an abelian group denoted by \(k(C)\)\SGAAuthorFootnote{raisedstar}{we shall not use here the usual notation \(\Kpx(C)\), since \(\Kpx(C)\) already denotes, when \(C\) is additive, the category of complexes of \(C\) ``up to homotopy''.} endowed with an additive map \(\cl_{C}:\Ob(C)\longrightarrow k(C)\) (when \(C\) is additive, we take for \(k(C)\) the symmetrization of the monoid of isomorphism classes of objects of \(C\); when \(C\) is triangulated, we take for \(k(C)\) the quotient of the free abelian group generated by \(\Ob(C)\) by the relations identifying \(L\) with \(L'+L''\) whenever there is a distinguished triangle \(L'\longrightarrow L\longrightarrow L''\longrightarrow L'[1]\)).

The group \(k(C)\) depends functorially on \(C\) (with respect to additive (resp.\ exact) functors). When \(u,u':C\to D\) are isomorphic additive (resp.\ exact) functors, we have \(k(u)=k(u')\).

When \(C\) is a triangulated category, we have a canonical epimorphism \(k(C^{\mathrm{ad}})\longrightarrow k(C)\) (cf.\ (6.2)).

\sgaStableTarget{sga6.e01.lem.6-4}\textbf{Lemma 6.4.} \uline{Let} \(C\) \uline{be an additive category}. \uline{The canonical functor} \(C\longrightarrow\Kpx^b(C)\) \uline{induces an isomorphism} (\uline{functorial in} \(C\)) \(k(C)\longrightarrow k(\Kpx^b(C))\).

\uline{Proof}. This is an immediate consequence of the following formula (which follows trivially from (6.2)):
\sgaStableTarget{sga6.e01.eq.6-4-1}\begin{equation}\tag{6.4.1}
  f(L)=\sum (-1)^i f(L^i)
\end{equation}
%%%%%%%%%%  scan idx 155  |  folio 142  %%%%%%%%%%
\SGAPrintPage{142}%
for \(L\in\Ob\Kpx^b(C)\) and \(f\) an additive function on \(\Kpx^b(C)\).

\sgaStableTarget{sga6.e01.def.6-5}\textbf{Definition 6.5.} \uline{Let} \(\underline{S}\) \uline{be a category,} \(\underline{C}\) \uline{an additive (resp}.\ \uline{triangulated)} \(\underline{S}\)\uline{-category} (1.1.1)\uline{, and} \(F\) \uline{an abelian presheaf on} \(\underline{S}\). \uline{A family of functions} \(f_X:\Ob(\underline{C}_X)\longrightarrow F(X)\) \uline{indexed by} \(X\in\Ob \underline{S}\) \uline{is called an additive} \(\underline{S}\)\uline{-function from} \(\underline{C}\) \uline{to} \(F\) \uline{when the following conditions are satisfied}:
\begin{enumerate}[label=(\roman*)]
\item \var{\uline{for every} \(X\in\Ob\underline{C}_X\)}{\uline{for every} \(X\in\Ob \underline{S}\)}{EN-VAR-107}\uline{,} \(f_X\) \uline{is an additive function} (6.1);
\item \uline{for every arrow} \(u:X\to Y\) \uline{of} \(\underline{S}\)\uline{, the following diagram is commutative}
% SGA display variant EN-VAR-108
\SGACaptureAppState{EN-VAR-108}%
\ifnum\SGAReaderForm=1\relax\SGAAppMark{EN-VAR-108}\fi
\ifnum\SGAReaderForm=0\relax
\[
\begin{tikzcd}
\Ob\underline{C}_Y \arrow[r,"f_Y"] \arrow[d,"u^*"'] & F(Y) \arrow[d,"F(u)"]\\
\Ob\underline{C}_X \arrow[r,"f_X",dash] & F(X).
\end{tikzcd}
\]
\else
\[
\begin{tikzcd}
\Ob\underline{C}_Y \arrow[r,"f_Y"] \arrow[d,"u^*"'] & F(Y) \arrow[d,"F(u)"]\\
\Ob\underline{C}_X \arrow[r,"f_X"] & F(X).
\end{tikzcd}
\]
\fi
\end{enumerate}
\uline{The additive} \(\underline{S}\)\uline{-functions from} \(\underline{C}\) \uline{to} \(F\) \uline{form an abelian group, denoted by} \(\Add_{\underline{S}}(\underline{C},F)\).

\sgaStableTarget{sga6.e01.rem.6-6}\textbf{Remarks 6.6.} a) Denote by \(\underline{C}^{\is}\) the fibered \(\underline{S}\)-category whose fiber over \(X\in\Ob \underline{S}\) is the category with the same objects as \(\underline{C}_X\), but whose arrows are the isomorphisms of \(\underline{C}_X\). On the other hand, denote by \(\underline F\) the fibered \(\underline{S}\)-category whose fiber over \(X\in\Ob \underline{S}\) is the discrete category (i.e.\ the only arrows are the identity arrows) \(F(X)\), the inverse-image functor by \(u\in\mathrm{Fl}(\underline{S})\) being \(F(u)\). Then an additive \(\underline{S}\)-function from \(\underline{C}\) to \(F\) is precisely a cartesian \(\underline{S}\)-functor from \(\underline{C}^{\is}\) to \(\underline F\), inducing an additive function in the sense of (6.1) on each fiber. It will be convenient for us to interpret additivity as follows. To fix ideas, suppose that \(\underline{C}\) is triangulated. The category \(\Cart_{\underline{S}}(\underline{C}^{\is},\underline F)\) (resp.\ \(\Cart_{\underline{S}}(\Tr(\underline{C})^{\is},\underline F)\) (1.1.2)) of cartesian \(\underline{S}\)-functors from \(\underline{C}^{\is}\) (resp.\ \(\Tr(\underline{C})^{\is}\)) to \(\underline F\) is the discrete category associated to an abelian
%%%%%%%%%%  scan idx 156  |  folio 143  %%%%%%%%%%
\SGAPrintPage{143}%
group, and we have a homomorphism
\[
  \phi:\Cart_{\underline{S}}(\underline{C}^{\is},\underline F)
  \longrightarrow
  \Cart_{\underline{S}}(\Tr(\underline{C})^{\is},\underline F),
\]
defined by \(\phi(f)(T)=f(L)-f(L')-f(L'')\) \var{for \(f\in\Cart_{\underline{S}}(\underline{C}^{\is},F)\) and}{for \(f\in\Cart_{\underline{S}}(\underline{C}^{\is},\underline F)\) and}{EN-VAR-109} \(T=(L'\to L\to L''\to L'[1])\in\Ob\Tr(\underline{C})\). To say that \(f\in\Cart_{\underline{S}}(\underline{C}^{\is},\underline F)\) is additive means that \(\phi(f)=0\); in other words, we have a (canonical functorial) exact sequence
\sgaStableTarget{sga6.e01.eq.6-6-1}\begin{equation}\tag{6.6.1}
  0\longrightarrow \Add_{\underline{S}}(\underline{C},F)
  \longrightarrow \Cart_{\underline{S}}(\underline{C}^{\is},\underline F)
  \xrightarrow{\phi}
  \Cart_{\underline{S}}(\Tr(\underline{C})^{\is},\underline F).
\end{equation}

b) Let \(u:X\to Y\) be an arrow of \(\underline{S}\). The ``inverse-image by \(u\)'' functors, being all isomorphic to one another, define (6.3) a unique homomorphism \(k(Y)\longrightarrow k(X)\). Consequently, \(X\mapsto k(X)\) defines an abelian presheaf on \(\underline{S}\), denoted by \(k(\underline{C})\). An additive \(\underline{S}\)-function from \(\underline{C}\) to \(F\) can then be interpreted simply as a homomorphism of presheaves \(k(\underline{C})\longrightarrow F\); in other words, we have
\[
  \Add_{\underline{S}}(\underline{C},F)=\Hom(k(\underline{C}),F).
\]
The presheaf \(k(\underline{C})\) of course depends functorially on \(\underline{C}\) with respect to additive (resp.\ exact) \(\underline{S}\)-functors.

\sgaStableTarget{sga6.e01.prop.6-7}\textbf{Proposition 6.7.} \uline{Let} \(\underline{S}\) \uline{be a site,} \(\underline{C}\) \uline{a flat abelian} \(\underline{S}\)\uline{-category, and} \(\underline{C}_0\) \uline{a strictly full additive sub-}\(\underline{S}\)\uline{-category of} \(\underline{C}\) \uline{satisfying the hypotheses} (4.0). \uline{Then the morphism of presheaves}
\[
  k(\underline{C}_0)\longrightarrow k(\Parf(\underline{C},\underline{C}_0))\qquad(\text{cf. }(4.9)),
\]
\uline{induced by the inclusion of} \(\underline{C}_0\) \uline{in} \(\underline{\Parf}(\underline{C},\underline{C}_0)\)\uline{, defines an isomorphism on the associated sheaves. In other words, for every abelian sheaf} \(F\) \uline{on} \(\underline{S}\)\uline{,}
%%%%%%%%%%  scan idx 157  |  folio 144  %%%%%%%%%%
\SGAPrintPage{144}%
\uline{the restriction homomorphism}
\[
  \Add_{\underline{S}}(\Parf(\underline{C},\underline{C}_0),F)
  \longrightarrow
  \Add_{\underline{S}}(\underline{C}_0,F)
\]
\uline{is an isomorphism}.

\uline{Proof}. Let \(F\) be an abelian sheaf on \(\underline{S}\); let us show that the arrow
\[
  \Add_{\underline{S}}(\Parf(\underline{C},\underline{C}_0),F)
  \longrightarrow
  \Add_{\underline{S}}(\underline{C}_0,F)
\]
is an isomorphism. It factors as
\[
\Add_{\underline{S}}(\Parf(\underline{C},\underline{C}_0),F)
 \xrightarrow{a}
\Add_{\underline{S}}(\Kpx^b(\underline{C}_0),F)
 \xrightarrow{b}
\Add_{\underline{S}}(\underline{C}_0,F).
\]
Now, by (6.4), the arrow \(b\) is an isomorphism. We are therefore reduced to proving that \(a\) is an isomorphism.

Consider the following commutative diagram, in which the rows are exact (6.6.1), and the vertical arrows are the ``restriction'' arrows:
\[
\begin{array}{ccccccccc}
0&\to&\Add_{\underline{S}}(\Parf(\underline{C}),F)&\to&\Cart_{\underline{S}}(\Parf(\underline{C})^{\is},\underline F)&\to&\Cart_{\underline{S}}(\Tr(\Parf(\underline{C}))^{\is},\underline F)\\
&&\downarrow a&&\downarrow&&\downarrow\\
0&\to&\Add_{\underline{S}}(\Kpx^b(\underline{C}_0),F)&\to&\Cart_{\underline{S}}(\Kpx^b(\underline{C}_0)^{\is},\underline F)&\to&\Cart_{\underline{S}}(\Tr(\Kpx^b(\underline{C}_0))^{\is},\underline F).
\end{array}
\]
Since \(F\) is a sheaf, \(\underline F\) is a stack (G II 2.2.4); thus, by (4.10), the two vertical arrows on the right are isomorphisms.

The same is therefore true of \(a\), q.e.d.

\sgaStableTarget{sga6.e01.rem.6-7-1}\textbf{Remark 6.7.1.} (cf.\ 4.10.1). In (6.7), we can replace \(\Parf(\underline{C},\underline{C}_0)\) by \(\Dpx^*(\underline{C})_{\underline{C}_0\text{-}\mathrm{parf}}\) (\(*=-\) or \(b\)).

%%%%%%%%%%  scan idx 158  |  folio 145  %%%%%%%%%%
\SGAPrintPage{145}%
\sgaStableTarget{sga6.e01.para.6-8}6.8. Let \((\underline{S},A)\) be a ringed topos and \(\underline{C}\) and \(\underline{C}_0\) the fibered \(\underline{S}\)-categories considered in (4.8). We have homomorphisms of abelian presheaves
\[
  k(\underline{C}_0)
  \xrightarrow{\alpha}
  k(\Dpx^-({}_A \underline{S})_{\mathrm{parf}})
  \xrightarrow{\beta}
  k(\Parf({}_A \underline{S})),
\]
which, by (6.7), induce isomorphisms on the associated sheaves. The homomorphism \(\alpha\) is compatible with inverse images, i.e.\ when \(u:(\underline{S}',A')\longrightarrow(\underline{S},A)\) is a morphism of ringed topoi, we have a commutative diagram, in which the vertical arrows are induced by the functor \(\mathbb L u^*\) (cf.\ (4.19.1)):
\[
\begin{tikzcd}
 k(\underline{C}_0) \arrow[r] \arrow[d] & k(\Dpx^-({}_A \underline{S})_{\mathrm{parf}}) \arrow[d]\\
 k(\underline{C}'_0) \arrow[r] & k(\Dpx^-({}_{A'} \underline{S}')_{\mathrm{parf}})
\end{tikzcd}
\]
(\(\underline{C}'_0\) being defined analogously to \(\underline{C}_0\)). When \(A\) is commutative, the tensor product \(\overset{\mathrm L}{\otimes}\) equips \(k(\underline{C}_0)\) and \(k(\Dpx^-({}_A \underline{S})_{\mathrm{parf}})\) with structures of presheaves of rings (cf.\ (4.19.1)), and \(\alpha\) is a homomorphism of presheaves of rings.

\sgaStableTarget{sga6.e01.para.6-9}6.9. In the situation of (6.8), suppose that \((\underline{S},A)\) is locally ringed (cf.\ (2.15.1)). Then (2.15.1), for every \(X\in\Ob \underline{S}\), \(\underline{C}_{0X}\) is the category of locally free \(A_X\)-Modules of finite type. Denote by \(\underline{C}_1\) the fibered \(\underline{S}\)-category defined by \(X\longmapsto\) (the category of free \(A_X\)-Modules of finite type). It is clear that \(\underline{C}_1\) satisfies the hypotheses (4.0). An object of \(\Dpx(\underline{S})\) is perfect if and only if it is perfect relative to \(\underline{C}_1\). We have homomorphisms of abelian presheaves
\[
  k(\underline{C}_1)\longrightarrow k(\underline{C}_0)\longrightarrow
  k(\Dpx^-(\underline{S})_{\mathrm{parf}})\longrightarrow k(\Parf(\underline{S}))
\]
%%%%%%%%%%  scan idx 159  |  folio 146  %%%%%%%%%%
\SGAPrintPage{146}%
(the two on the left being homomorphisms of presheaves of rings), which, by (6.7), induce isomorphisms on the associated sheaves. We shall show that the associated sheaves are in fact canonically isomorphic to the constant sheaf \(\mathbb Z_{\underline{S}}\).

\sgaStableTarget{sga6.e01.lem.6-10}\textbf{Lemma 6.10.} \uline{Let} \((\underline{S},A)\) \uline{be a locally ringed topos,} \(U\) \uline{an object of} \(\underline{S}\) \uline{distinct from the empty object, and} \(p,q\in\mathbb N\)\uline{. If we have} \(A_U^p\simeq A_U^q\)\uline{, then} \(p=q\)\uline{.}

\uline{Proof}. By an argument similar to the one given in the proof of (2.15.1), we reduce to the case where \((\underline{S},A)\) is the locally ringed topos defined by a prescheme, and in this case the assertion is trivial (take a point in \(U\)).

\sgaStableTarget{sga6.e01.def.6-11}\textbf{Definition 6.11.} \uline{We call} rank, \uline{and denote by} \(\rg\), \uline{the additive} \(\underline{S}\)\uline{-function from} \(\underline{C}_1\) \uline{to} \(\mathbb Z_{\underline{S}}\) \uline{defined by}
\[
  \rg(L)=p\quad\text{if } L\simeq A_X^p,
  \qquad X\ne\varnothing,
\]
\[
  \rg(L)=0\quad\text{if } L\in\Ob\underline{C}_{1X},
  \qquad X=\varnothing.
\]
By (6.7), the rank function extends uniquely to an additive \(\underline{S}\)-function from \(\Parf(\underline{S})\) to \(\mathbb Z_{\underline{S}}\) (which we shall again call rank and denote by \(\rg\)).

\sgaStableTarget{sga6.e01.prop.6-12}\textbf{Proposition 6.12.} a) \uline{For} \(L,M\in\Ob\Dpx^-(X)_{\mathrm{parf}}\) (\(X\in\Ob \underline{S}\)), \uline{we have} \(\rg(L\overset{\mathrm L}{\otimes} M)=\rg(L)\,\rg(M)\).

b) \uline{Let} \(u:(\underline{S}',A')\longrightarrow(\underline{S},A)\) \uline{be a morphism of locally ringed topoi. For} \(M\in\Ob\Dpx^-(\underline{S})_{\mathrm{parf}}\), we have \(\rg(\mathbb L u^*(M))=u^*(\rg(M))\). \uline{In other words, the rank homomorphism}: \(k(\Dpx^-(\underline{S})_{\mathrm{parf}})\longrightarrow \mathbb Z_{\underline{S}}\) \uline{is a homomorphism of presheaves of rings, compatible with inverse images}.

% Additional notation for Exposé I
\providecommand{\Cpx}{\mathrm C}
\providecommand{\Kpx}{\mathrm K}
\providecommand{\Dpx}{\mathrm D}
\providecommand{\Ho}{\mathrm H}
\providecommand{\RHom}{\mathrm R\!\operatorname{Hom}}
\renewcommand{\Ob}{\operatorname{ob}}
\providecommand{\Hom}{\operatorname{Hom}}
\providecommand{\RHom}{\mathrm R\!\operatorname{Hom}}
\providecommand{\Ext}{\operatorname{Ext}}
\providecommand{\Tor}{\operatorname{Tor}}
\providecommand{\Tr}{\operatorname{Tr}}
\providecommand{\Add}{\operatorname{Add}}
\providecommand{\Cart}{\underline{\operatorname{Cart}}}
\providecommand{\Parf}{\operatorname{Parf}}
\providecommand{\coh}{\operatorname{coh}}
\providecommand{\amp}{\operatorname{amp}}
\providecommand{\cl}{\operatorname{cl}}
\providecommand{\rg}{\operatorname{rg}}
\providecommand{\Id}{\operatorname{Id}}
\providecommand{\Ker}{\operatorname{Ker}}
\renewcommand{\Im}{\operatorname{Im}}
\providecommand{\cH}{\mathcal H}
\providecommand{\Zs}{\mathbb Z_S}
\providecommand{\A}{\mathcal A}
\providecommand{\cC}{\mathcal C}
\providecommand{\cCzero}{\mathcal C_0}
\providecommand{\otL}{\mathop{\otimes}^{\mathrm L}}
\providecommand{\tamp}{\operatorname{tor.amp}}
\providecommand{\pamp}{\operatorname{parf.amp}}
\providecommand{\pdim}{\operatorname{parf.dim}}
\providecommand{\bon}{\mathrm b}
\providecommand{\parf}{\operatorname{parf}}
\providecommand{\parff}{\operatorname{parf}}
\providecommand{\ac}{\mathrm{ac}}


% Additional notation for Exposé I, completion pages 154--166
\providecommand{\End}{\operatorname{End}}
\providecommand{\tr}{\operatorname{tr}}

%%%%%%%%%%  scan idx 160  |  folio 147  %%%%%%%%%%
\SGAPrintPage{147}%
\uline{Proof}. It suffices to show that the homomorphism induced by rank on the sheaf associated to \(k(\Dpx^-(\underline{S})_{\mathrm{parf}})\) is a homomorphism of sheaves of rings compatible with inverse images. By (6.9), we can replace \(k(\Dpx^-(\underline{S})_{\mathrm{parf}})\) by \(k(\underline{C}_1)\). We are therefore reduced to checking a) and b) for \(L\) and \(M\) free of finite type, which is trivial.

\sgaStableTarget{sga6.e01.para.6-13}6.13. With \((\underline{S},A)\) still denoting a locally ringed topos, let \(X\in\Ob \underline{S}\) and \(\underline n\in H^0(X,\mathbb Z_{\underline{S}})\). We define an object of \(\Parf(X)\), denoted by \(A_X^{\underline n}\), as follows. When \(\underline n\) is constant with value an integer \(n\geq0\) (resp.\ \(n<0\)), we put \(A_X^{\underline n}=A_X^n\) (resp.\ \(A_X^{\underline n}=A_X^{-n}[1]\)); in the general case, we decompose \(X\) into disjoint pieces \(X_i\) such that \(\underline n|X_i\) is constant, and \(A_X^{\underline n}\) is defined as the unique object of \(\Parf(X)\) inducing \(A_{X_i}^{\underline n|X_i}\) on each \(X_i\). (More generally, we can define \(L^{\underline n}\) for \(L\in\Ob\Parf(X)\), by putting \(L^{\underline n}=A_X^{\underline n}\otimes L\).).
The map \(i_X:H^0(X,\mathbb Z_{\underline{S}})\longrightarrow k(\Parf(X))\), defined by \(i_X(\underline n)=\cl(A_X^{\underline n})\), is a ring homomorphism (such that \(i_X(1)=1\)), and defines, as \(X\) varies, a homomorphism of presheaves of rings
\[
  i:\mathbb Z_{\underline{S}}\longrightarrow k(\Parf(\underline{S})).
\]

\sgaStableTarget{sga6.e01.prop.6-14}\textbf{Proposition 6.14.} \uline{Under the hypotheses of} (6.13), \uline{we have}
\[
  \rg\, i=\Id(\mathbb Z_{\underline{S}})
\]
(\uline{in other words,} \(k(\Parf(\underline{S}))\)\uline{, considered as a presheaf of} \(\mathbb Z_{\underline{S}}\)\uline{-algebras by means of} \(i\)\uline{, is augmented toward} \(\mathbb Z_{\underline{S}}\) \uline{by rank}). \uline{Moreover, rank}
\[
  \rg:k(\Parf(\underline{S}))\longrightarrow\mathbb Z_{\underline{S}}
\]
\uline{defines an isomorphism on the associated sheaves}.

%%%%%%%%%%  scan idx 161  |  folio 148  %%%%%%%%%%
\SGAPrintPage{148}%
\uline{Proof}. The first assertion follows trivially from the relations \(\cl(L[1])=-\cl(L)\), \(\cl(L^n)=n\cl(L)\) for \(L\in\Ob\Parf(\underline{S})\). Let us prove the second. First observe that \(i\) factors in the evident way through \(k(\underline{C}_0)\). Since \(\rg\, i\) is the identity, we are reduced (by (6.9)) to proving that the homomorphism \(i\, \rg:k(\underline{C}_0)\longrightarrow k(\underline{C}_0)\) induces an isomorphism on the associated sheaves, i.e.\ that, for \(L\in\Ob\underline{C}_0\), \(i\, \rg(L)\) is locally isomorphic to \(L\), which is immediate.

We deduce, by (6.6.b)),

\sgaStableTarget{sga6.e01.cor.6-15}\textbf{Corollary 6.15.} \uline{Under the hypotheses of} (6.13), \uline{for every abelian sheaf} \(F\) \uline{on} \(\underline{S}\)\uline{, the group of additive} \(\underline{S}\)\uline{-functions from} \(\underline{\Parf}(\underline{S})\) \uline{to} \(F\) \uline{is canonically identified with} \(H^0(S,F)\) (\(S\) = \uline{final object of} \(\underline{S}\)), \uline{the rank function corresponding to} \(F=\mathbb Z_{\underline{S}}\) \uline{and to the unit section of} \(\mathbb Z_{\underline{S}}\).

\sgaStableTarget{sga6.e01.sec.7}\section*{7. Duality of perfect complexes}

\sgaStableTarget{sga6.e01.para.7-0}7.0. In this section, \((\underline{S},A)\) denotes a topos ringed in commutative rings; we denote by \(S\) the final object of \(\underline{S}\). We propose to generalize to perfect complexes the notions of duality well known for locally free sheaves of finite type.

\sgaStableTarget{sga6.e01.prop.7-1}\textbf{Proposition 7.1.} \uline{Let} \(E\in\Ob\Dpx(S)_{\mathrm{parf}}\), \(F\in\Ob\Dpx^+(S)\).

a) \uline{When} \(F\) \uline{has locally bounded cohomology, the same holds for} \(\mathbb{R}\underline{\Hom}(E,F)\). \uline{When} \(F\in\Ob\Dpx(S)\) \uline{is pseudo-coherent (resp.\ locally of finite tor-dimension, resp.\ perfect), the same holds for} \(\mathbb{R}\underline{\Hom}(E,F)\).

%%%%%%%%%%  scan idx 162  |  folio 149  %%%%%%%%%%
\SGAPrintPage{149}%
b) Suppose that \(E\) is moreover of finite tor-dimension (a condition satisfied when \(\underline{S}\) is of finite type). Then, when \(F\) has bounded cohomology (resp.\ finite tor-dimension), the same holds for \(\mathbb{R}\underline{\Hom}(E,F)\). More precisely, let \(a,b,m,n\) be integers such that \(\tamp(E)\subset[a,b]\) (or, equivalently, by (5.8), \(\pamp(E)\subset[a,b]\)) and \(\Ho^i(F)=0\) for \(i\notin[m,n]\) (resp.\ \(\tamp(F)\subset[m,n]\)). Then we have \(\Ho^i(\mathbb{R}\underline{\Hom}(E,F))=0\) for \(i\notin[m-b,n-a]\) (resp.\ \(\tamp(\mathbb{R}\underline{\Hom}(E,F))\subset[m-b,n-a]\)).

\uline{Proof}. Assertion a) is local, so we can suppose that \(E\) is strictly perfect, and we conclude by dévissage on \(E\) (reducing to \(E=A\)). For assertion b), we can suppose that \(E\) is strictly perfect with \(E^i=0\) for \(i\notin[a,b]\), and \(F^i=0\) for \(i\notin[m,n]\) (resp.\ \(F^i\) flat for every \(i\) and \(F^i=0\) for \(i\notin[m,n]\)). It then suffices to apply the following lemma:

\sgaStableTarget{sga6.e01.lem.7-1-1}\textbf{Lemma 7.1.1.} \uline{Let} \(E\in\Ob\Dpx^-(S)\), \(F\in\Ob\Dpx^+(S)\). \uline{When} \(E\) \uline{is strictly pseudo-coherent} (i.e.\ (2.1) \uline{with components direct factors of free Modules of finite type}), \uline{the canonical arrow of} \(\Dpx^+(S)\)
\[
  \underline{\Hom}^\bullet(E,F)\longrightarrow\mathbb{R}\underline{\Hom}(E,F)
\]
\uline{is an isomorphism}.

\uline{Proof}. The spectral sequence
\[
  E_2^{pq}=\underline{\Ho}^p(\underline{\Ext}^q(E,F)^\bullet)\Longrightarrow \underline{\Ext}^*(E,F),
\]
where
\[
  \underline{\Ext}^q(E,F)^i=\coprod_{i_2-i_1=i}\underline{\Ext}^q(E^{i_1},F^{i_2}),
\]
degenerates since \(\underline{\Ext}^q(E,F)^\bullet=0\) for every \(q>0\) (1.3.1). We deduce isomorphisms
%%%%%%%%%%  scan idx 163  |  folio 150  %%%%%%%%%%
\SGAPrintPage{150}%
\[
  E_2^{p0}=\underline{\Ho}^p(\underline{\Hom}^\bullet(E,F))
  \xrightarrow{\sim}\underline{\Ext}^p(E,F)=\underline{\Ho}^p(\mathbb{R}\underline{\Hom}(E,F)),
\]
which proves the assertion.

\sgaStableTarget{sga6.e01.prop.7-1-2}\textbf{Proposition 7.1.2.} \uline{Let} \(f:(\underline{S}',A')\longrightarrow(\underline{S},A)\) \uline{be a morphism of ringed topoi. Let} \(E\in\Ob\Dpx^-(S)\) \uline{and} \(F\in\Ob\Dpx^b(S)\) \uline{be such that} \(\mathbb{R}\underline{\Hom}(E,F)\in\Ob\Dpx^b(S)\) \uline{and} \(\mathbb{L}f^*(F)\in\Ob\Dpx^b(S')\). \uline{There exists a canonical functorial arrow}
\[
  \mathbb{L}f^*\mathbb{R}\underline{\Hom}(E,F)
  \longrightarrow
  \mathbb{R}\underline{\Hom}(\mathbb{L}f^*E,\mathbb{L}f^*F).
\]
\uline{It is an isomorphism when} \(E\) \uline{is perfect}.

\uline{Proof}. By the adjunction of the functors \(\mathbb{L}f^*\) and \(\mathbb{R}f_*\), it suffices to define an arrow
\[
  \mathbb{R}\underline{\Hom}(E,F)
  \longrightarrow
  \mathbb{R}f_*\mathbb{R}\underline{\Hom}(\mathbb{L}f^*E,\mathbb{L}f^*F).
\]
For this we use the adjunction arrow \(F\to\mathbb{R}f_*\mathbb{L}f^*F\) and the duality isomorphism
\[
  \mathbb{R}f_*\mathbb{R}\underline{\Hom}(\mathbb{L}f^*E,\mathbb{L}f^*F)
  \simeq
  \mathbb{R}\underline{\Hom}(E,\mathbb{R}f_*\mathbb{L}f^*F).
\]
To see that the arrow of (7.1.2) is an isomorphism for \(E\) perfect, we reduce trivially, by localization and dévissage, to the case \(E=A\), in which it is a tautology.

\sgaStableTarget{sga6.e01.prop.7-2}\textbf{Proposition 7.2.} \uline{There exists} (cf.\ SGA 5 I 1.6), \uline{for} \(E\in\Ob\Dpx(S)\) \uline{and} \(F\in\Ob\Dpx^+(S)\)\uline{, a canonical functorial arrow}
\[
  E\longrightarrow\mathbb{R}\underline{\Hom}(\mathbb{R}\underline{\Hom}(E,F),F).
\]
\uline{It is an isomorphism when} \(E\in\Ob\Dpx(S)_{\mathrm{parf}}\) and \(F=A\).

%%%%%%%%%%  scan idx 164  |  folio 151  %%%%%%%%%%
\SGAPrintPage{151}%
\uline{Proof}. We define the arrow by resolving \(F\) injectively and using the canonical arrow of \(\Kpx(S)\)
\[
  E\longrightarrow\underline{\Hom}^\bullet(\underline{\Hom}^\bullet(E,F),F).
\]
To see that the arrow in question is an isomorphism when \(E\) is perfect and \(F=A\), we reduce, by localization and dévissage, to \(E=A\), in which case it is evident.

\sgaStableTarget{sga6.e01.def.7-3}\textbf{Definition 7.3.} \uline{Let} \(E\in\Ob\Dpx(S)\). \uline{We shall put} \(\mathbb{R}\underline{\Hom}(E,A)=E^\vee\), \uline{and say that} \(E^\vee\) \uline{is the} dual \uline{of} \(E\).

By (7.1) and (7.2), when \(E\in\Ob\Dpx(S)\) is perfect (resp.\ perfect and of finite tor-dimension), \(E^\vee\) is also, and the canonical arrow \(E\to E^{\vee\vee}\) is an isomorphism. In other words, the functor \(\mathbb{R}\underline{\Hom}(\quad,A)\) defines an anti-auto-equivalence of the category \(\Dpx(S)_{\mathrm{parf}}\) (resp.\ \(\Dpx(S)_{\underline{\mathrm{parf}}}\) (5.8.4.b))) with itself.

\sgaStableTarget{sga6.e01.lem.7-4}\textbf{Lemma 7.4} (CARTAN isomorphism) (cf.\ ([H] II 5.15)). \uline{For} \(E\in\Ob\Dpx(S)\), \(F\in\Ob\Dpx^-(S)\), \(G\in\Ob\Dpx^+(S)\), \uline{there exist canonical functorial isomorphisms}:
\[
\begin{aligned}
\mathrm R\underline{\Hom}(E\otL F,G)&\xrightarrow{\sim}\mathrm R\underline{\Hom}(E,\mathrm R\underline{\Hom}(F,G)),\\
\RHom(E\otL F,G)&\xrightarrow{\sim}\RHom(E,\mathrm R\underline{\Hom}(F,G)),\\
\Hom(E\otL F,G)&\xrightarrow{\sim}\Hom(E,\mathrm R\underline{\Hom}(F,G)).
\end{aligned}
\]

%%%%%%%%%%  scan idx 165  |  folio 152  %%%%%%%%%%
\SGAPrintPage{152}%
\uline{Proof}: This is standard: resolve \(G\) injectively and \(F\) flatly.

\sgaStableTarget{sga6.e01.lem.7-5}\textbf{Lemma 7.5.} \uline{For} \(E\in\Ob\Dpx^-(S)\), \(F\in\Ob\Dpx^+(S)\), \uline{there exists a canonical functorial homomorphism}
\[
  E\otL\mathrm R\underline{\Hom}(E,F)\longrightarrow F
\]

\sgaStableTarget{sga6.e01.prop.7-6}\textbf{Proposition 7.6} (cf.\ ([H] II 5.13)). \uline{Let} \(E\in\Ob\Dpx(S)\), \(F\in\Ob\Dpx^+(S)\), \(G\in\Ob\Dpx^+(S)\). \uline{Under one of the following hypotheses}
\begin{enumerate}[label=\roman*)]
\item \(G\) \uline{of finite tor-dimension},
\item \(E\in\Ob\Dpx^-(S)\), \(F\) \uline{of finite tor-dimension and} \(G\in\Ob\Dpx^b(S)\),
\end{enumerate}
\uline{there exists a canonical functorial homomorphism}
\[
  \mathrm R\underline{\Hom}(E,F)\otL G
  \longrightarrow
  \mathrm R\underline{\Hom}(E,F\otL G).
\]
\uline{It is an isomorphism when} \(E\) \uline{is perfect}.

\uline{Proof}. Let us define the arrow. Under hypothesis (i), it suffices to resolve \(F\) injectively and \(G\) flatly. Under hypothesis (ii), it is equivalent, by (7.4), to define an arrow
\[
  E\otL\mathbb{R}\underline{\Hom}(E,F)\otL G\longrightarrow F\otL G.
\]
We apply the functor \(\otL G\) to the arrow of (7.5).

To see that the arrow is an isomorphism when \(E\) is perfect, we use the standard method (cf.\ the proof of (7.2)): we reduce, by localization and dévissage, to the case \(E=A\).

%%%%%%%%%%  scan idx 166  |  folio 153  %%%%%%%%%%
\SGAPrintPage{153}%
\sgaStableTarget{sga6.e01.cor.7-7}\textbf{Corollary 7.7.} \uline{There exists, for} \(E\in\Ob\Dpx(S)_{\underline{\mathrm{parf}}}\) and \(F\in\Ob\Dpx^+(S)\), \uline{a functorial isomorphism}
\[
  E^\vee\otL F\xrightarrow{\sim}\mathbb{R}\underline{\Hom}(E,F),
\]
\uline{with the notation of} (7.3).

\uline{Proof}. Apply (7.6), replacing \(F\) and \(G\) respectively by \(A\) and \(F\), and use the fact (7.1) that \(E^\vee\in\Ob\Dpx(S)_{\underline{\mathrm{parf}}}\).

\sgaStableTarget{sga6.e01.cor.7-8}\textbf{Corollary 7.8.} \uline{Let} \(E,F\in\Ob\Dpx(S)_{\underline{\mathrm{parf}}}\). \uline{There exists a natural arrow}
\[
  \mathbb{R}\underline{\Hom}(E,F)\otL\mathbb{R}\underline{\Hom}(F,E)\longrightarrow A,
\]
\uline{which is a} ``\uline{perfect duality}'', i.e.\ \uline{identifies, modulo} (7.4), \uline{each of the complexes} \(\mathbb{R}\underline{\Hom}(E,F)\), \(\mathbb{R}\underline{\Hom}(F,E)\) \uline{with the dual of the other} (\uline{in the sense of} (7.3)).

\uline{Proof}. By (7.5), we have natural pairings \(E^\vee\otL E\to A\), \(F^\vee\otL F\to A\), hence, by tensor product, a natural pairing
\[
  E^\vee\otL F\otL F^\vee\otL E\longrightarrow A,
\]
which, in view of (7.7), gives the stated pairing on the \(\mathbb{R}\underline{\Hom}\). Let us show that it is a perfect duality. We have the canonical isomorphisms
\[
\begin{aligned}
\mathbb{R}\underline{\Hom}(F^\vee\otL E,A)&\simeq\mathbb{R}\underline{\Hom}(F^\vee,E^\vee) &&(7.4)\\
&\simeq F^{\vee\vee}\otL E^\vee &&(7.6)\\
&\simeq F\otL E^\vee &&(7.2).
\end{aligned}
\]
Thus, after checking compatibility, we are done. We could also localize and use dévissage to reduce to \(E=F=A\).

%%%%%%%%%%  scan idx 167  |  folio 154  %%%%%%%%%%
\SGAPrintPage{154}%
\sgaStableTarget{sga6.e01.exercise.7-9}\textbf{Exercise 7.9.} Let \(E,F\in\Ob\Dpx(S)_{\underline{\mathrm{parf}}}\). Show that there exists a canonical isomorphism
\[
  E^\vee\otL F^\vee \xrightarrow{\sim} (E\otL F)^\vee .
\]

\sgaStableTarget{sga6.e01.sec.8}\section*{8. Traces and cup-products}

\sgaStableTarget{sga6.e01.para.8-0}8.0. In this section, we generalize to perfect complexes the notion of the trace of an endomorphism of a projective module of finite type (see SGA 5 III for applications). The hypotheses and notation are those of (7.0).

\sgaStableTarget{sga6.e01.para.8-1}8.1. Let \(E\in\Ob\Dpx(S)_{\underline{\mathrm{parf}}}\) (5.8.4.b)). There exists a natural arrow
\[
  \tau:\mathbb{R}\underline{\Hom}(E,E)\longrightarrow A,
\]
composed of the inverse isomorphism of (7.7), \(\mathbb{R}\underline{\Hom}(E,E)\xrightarrow{\sim}E^\vee\otL E\), and the canonical arrow (7.5) \(E^\vee\otL E\to A\).

Let \(U\in\Ob \underline{S}\). Applying the functor \(H^0(U,.)\) to \(\tau\), and using the fact that \(H^0(U,\mathbb{R}\underline{\Hom}(E,E))\) is identified with \(\Hom_{\Dpx(U)}(E|U,E|U)\), we deduce a homomorphism of \(H^0(U,A)\)-modules,
\[
  \Tr_{E|U}:\Hom_{\Dpx(U)}(E|U,E|U)\longrightarrow H^0(U,A),
\]
which is called the \uline{trace homomorphism}\SGAAuthorFootnote{raisedstar}{The finite tor-dimension hypothesis on \(E\) is in fact superfluous for the definition of \(\Tr_{E|U}\), see (8.5).}.

Before going further, let us give some properties of the trace that follow immediately from the definition.

%%%%%%%%%%  scan idx 168  |  folio 155  %%%%%%%%%%
\SGAPrintPage{155}%
\sgaStableTarget{sga6.e01.para.8-1-1}8.1.1. When \(E=A\), \(\Tr_{A|U}:\Hom(A|U,A|U)\longrightarrow H^0(U,A)\) is precisely the well-known isomorphism. When \(E\) is a direct factor of a free Module of finite type, a homomorphism from \(E\) to \(E\) over \(U\) is locally defined by
\[
  e\longmapsto \sum\langle e,x_i^\vee\rangle x_i,
\]
where \(x_i\) (resp.\ \(x_i^\vee\)) is a local section of \(E\) (resp.\ \(E^\vee\)); the corresponding trace is \(\sum\langle x_i,x_i^\vee\rangle\). We easily deduce that, when \(E\) and \(F\) are direct factors of free Modules of finite type, and \(u:(E+F)|U\longrightarrow(E+F)|U\) is defined by a matrix \(\begin{pmatrix}a&b\\ c&d\end{pmatrix}\), we have
\[
  \Tr_{(E+F)|U}(u)=\Tr_{E|U}(a)+\Tr_{F|U}(d).
\]

\sgaStableTarget{sga6.e01.para.8-1-2}8.1.2. Let \(E\) be a strictly perfect complex, \(U\in\Ob \underline{S}\), and \(u:E|U\longrightarrow E|U\) a homomorphism of complexes, with components \(u^i:E^i|U\to E^i|U\). We then have
\[
  \Tr_{E|U}(u)=\sum(-1)^i\Tr_{E^i|U}(u^i).
\]
Indeed, since the question is local, we can suppose that \(u\) is defined by a section of \(\underline{H}^0(E^\vee\otimes E)\), and even that this section is of the form \(\coprod_i\sum_j x_{ij}^\vee\otimes x_{ij}\), where \(x_{ij}\) (resp.\ \(x_{ij}^\vee\)) is a section of \(E^i\) (resp.\ \(E^{i\vee}\)). Making the arrow (7.5) \(E^\vee\otimes E\to A\) explicit, we find
\[
  \Tr_{E|U}(u)=\sum(-1)^i\langle x_{ij},x_{ij}^\vee\rangle
  =\sum(-1)^i\Tr_{E^i|U}(u^i),
\]
q.e.d.

\sgaStableTarget{sga6.e01.para.8-2}8.2. Let \(E,F\in\Ob\Dpx(S)_{\underline{\mathrm{parf}}}\). We have seen (7.8) that there exists a canonical homomorphism
\[
  \gamma:\mathbb{R}\underline{\Hom}(E,F)\otL\mathbb{R}\underline{\Hom}(F,E)\longrightarrow A.
\]
Let \(U\in\Ob \underline{S}\), and take \(u\in H^0(U,\mathbb{R}\underline{\Hom}(E,F))\) and \(v\in H^0(U,\mathbb{R}\underline{\Hom}(F,E))\).

%%%%%%%%%%  scan idx 169  |  folio 156  %%%%%%%%%%
\SGAPrintPage{156}%
We can interpret \(u\) (resp.\ \(v\)) as a homomorphism, in \(\Dpx(U)\), from \(A|U\) to \(\mathbb{R}\underline{\Hom}(E,F)|U\) (resp.\ \(\mathbb{R}\underline{\Hom}(F,E)|U\)), and form
\[
  u\otL v:A|U\longrightarrow(\mathbb{R}\underline{\Hom}(E,F)\otL\mathbb{R}\underline{\Hom}(F,E))|U.
\]
The map \((u,v)\mapsto\langle u,v\rangle=\gamma\circ(u\otL v)\) defines an \(H^0(U,A)\)-linear map,
\[
  \Hom_{\Dpx(U)}(E|U,F|U)\otimes_{H^0(U,A)}\Hom_{\Dpx(U)}(F|U,E|U)
  \longrightarrow H^0(U,A),
\]
which we shall call the \uline{cup-product}\SGAAuthorFootnote{raisedstar}{the same remark as for the definition of \(\Tr_{E|U}\)}.

The relation between trace and cup-product is expressed in

\sgaStableTarget{sga6.e01.prop.8-3}\textbf{Proposition 8.3.} \uline{In the situation of} (8.2), \uline{for} \(u\in\Hom(E|U,F|U)\) \uline{and} \(v\in\Hom(F|U,E|U)\), \uline{we have}
\[
  \Tr_{E|U}(vu)=\Tr_{F|U}(uv)=\langle u,v\rangle.
\]
\uline{In particular}, \uline{when} \(E=F\), \uline{we have}
\[
  \Tr_{E|U}(u)=\langle u,\Id\rangle=\langle\Id,u\rangle.
\]

\uline{Proof}. For simplicity, we shall suppose that \(U=S\). We therefore have two arrows \(u:A\longrightarrow E^\vee\otL F\) and \(v:A\longrightarrow F^\vee\otL E\), hence
\[
  u\otL v:A\longrightarrow E^\vee\otL F\otL F^\vee\otL E.
\]
On the other hand, we have \(vu:A\to E^\vee\otL E\) and \(uv:A\to F^\vee\otL F\). The proposition follows from the commutativity of the following diagrams (whose verification is ``purely mechanical''):
%%%%%%%%%%  scan idx 170  |  folio 157  %%%%%%%%%%
\SGAPrintPage{157}%
\[
\begin{tikzcd}[column sep=large,row sep=large]
 A \arrow[r,"u\otL v"] \arrow[dr,"vu"'] & E^\vee\otL F\otL F^\vee\otL E \arrow[d,"\Id\otL\tau_F"] \\
 & E^\vee\otL E
\end{tikzcd}
\qquad
\begin{tikzcd}[column sep=large,row sep=large]
 A \arrow[r,"u\otL v"] \arrow[dr,"uv"'] & E^\vee\otL F\otL F^\vee\otL E \arrow[d,"\tau_E\otL\Id"] \\
 & F^\vee\otL F
\end{tikzcd}
\]
% SGA display variant EN-VAR-110
\SGACaptureAppState{EN-VAR-110}%
\ifnum\SGAReaderForm=1\relax\SGAAppMark{EN-VAR-110}\fi
\ifnum\SGAReaderForm=0\relax
\[
\begin{tikzcd}[column sep=
large,row sep=large]
E^\vee\otL F\otL F^\vee\otL E \arrow[r,"\Id\otL\tau_E"] \arrow[d,"\Id\otimes\tau_F"'] & F^\vee\otL F \arrow[d,"\tau_F"]\\
E^\vee\otL E \arrow[r,"\tau_E"'] & A .
\end{tikzcd}
\]
\else
\[
\begin{tikzcd}[column sep=
large,row sep=large]
E^\vee\otL F\otL F^\vee\otL E \arrow[r,"\Id\otL\tau_E"] \arrow[d,"\Id\otL\tau_F"'] & F^\vee\otL F \arrow[d,"\tau_F"]\\
E^\vee\otL E \arrow[r,"\tau_E"'] & A .
\end{tikzcd}
\]
\fi

\sgaStableTarget{sga6.e01.cor.8-4}\textbf{Corollary 8.4.} \(E,F,U\) \uline{being given as in} (8.2), \uline{let} \(u\in\Hom(E|U,E|U)\), \(s:E|U\to F|U\) \uline{be an isomorphism}, \uline{so that} \(sus^{-1}\in\Hom(F|U,F|U)\). \uline{We then have}
\[
  \Tr_{F|U}(sus^{-1})=\Tr_{E|U}(u).
\]

\uline{Proof}. This is immediate.

\sgaStableTarget{sga6.e01.rem.8-5}\textbf{Remark 8.5.} When \(\underline{C}\) is a fibered \(\underline{S}\)-category, denote by \(E(\underline{C})\) the fibered \(\underline{S}\)-category whose fiber over \(U\in\Ob \underline{S}\) is the category with objects the pairs \((L,u)\), where \(L\in\Ob\underline{C}_U\) and \(u\in\operatorname{End}_{\underline{S}}(L)\), and with arrows from \((L,u)\) to \((M,v)\) the arrows \(f:L\to M\) such that \(fu=vf\). By (8.4), we can regard the trace homomorphism as an \(\underline{S}\)-functor from \(E(\Dpx(\underline{S})_{\underline{\parf}})^{\mathrm{is}}\) to \(\underline A\) (where the exponent \(\mathrm{is}\) means that the only arrows considered are the isomorphisms, and \(\underline A\) is the \(\underline{S}\)-category with discrete fibers defined by \(A\) (cf.\ (6.6.a))). Since every perfect complex is locally of finite tor-dimension, the forgetful \(\underline{S}\)-functor
\[
  E(\Dpx(\underline{S})_{\underline{\parf}})^{\mathrm{is}}
  \longrightarrow
  E(\Dpx(\underline{S})_{\parf})^{\mathrm{is}}
\]
induces an isomorphism on the associated stacks. Consequently, the trace
%%%%%%%%%%  scan idx 171  |  folio 158  %%%%%%%%%%
\SGAPrintPage{158}%
homomorphism extends uniquely to an \(\underline{S}\)-functor from \(E(\Dpx(\underline{S})_{\parf})^{\mathrm{is}}\) to \(\underline A\).

For \(E\in\Ob\Parf(U)\) (\(U\in\Ob \underline{S}\)), we again denote by
\[
  \Tr_E:\Hom(E,E)\longrightarrow H^0(U,A)
\]
the corresponding homomorphism, which is a homomorphism of \(H^0(U,A)\)-modules. We see similarly that the definition of the cup-product (8.2), as well as assertions (8.3) and (8.4), generalize to the case of perfect complexes (not necessarily of finite tor-dimension).

\sgaStableTarget{sga6.e01.prop.8-6}\textbf{Proposition 8.6.} \uline{Let} \(U\in\Ob \underline{S}\), \(E,F\in\Parf(U)\), \(u\in\Hom(E,E)\), \(v\in\Hom(F,F)\), \uline{so that} \(u\otL v\in\Hom(E\otL F,E\otL F)\). \uline{We have}
\[
  \Tr_{E\otL F}(u\otL v)=\Tr_E(u)\,\Tr_F(v).
\]

\uline{Proof}. Since the question is local, by (8.4) and (4.10) we can suppose that \(u\) and \(v\) are arrows of \(K^b(\underline{C}_0)\). Using (8.1.2), we reduce to the case where \(u\) and \(v\) are arrows of \(\underline{C}_0\). We conclude by (8.1.1).

%%%%%%%%%%  scan idx 172  |  folio 159  %%%%%%%%%%
\SGAPrintPage{159}%
\sgaStableTarget{sga6.e01.bibliography}\section*{BIBLIOGRAPHY}
\begin{description}[leftmargin=2.5cm,style=nextline]
\item[{[H]}]
Hartshorne, R. Residues and Duality, Lecture Notes no.\ 20, Springer, 1966.
\item[{[V]}] Verdier, J-L. Catégories dérivées, I.H.E.S., 1963.
\item[{[G]}] Giraud, J. Cohomologie non abélienne, to appear in North Holl. Pub. Comp.
\item[{SGA 4}] Cohomologie étale des schémas, by M. Artin, A. Grothendieck, and J-L. Verdier, to appear in North Holl. Pub. Comp.
\item[{SGA 5}] Cohomologie \(\ell\)-adique et fonctions L, by A. Grothendieck, to appear in North Holl. Pub. Comp.
\item[{[1]}] Hakim, M. Topos annelés et schémas relatifs, thesis (in preparation).
\item[{[2]}] Verdier, J-L. Les catégories dérivées des catégories abéliennes, to appear in North Holl. Pub. Comp.
\item[{[MULT]}] Cartan, H. and Eilenberg, S. Homological Algebra, Princeton, 1956.
\end{description}

\providecommand{\cS}{\mathcal S}
\providecommand{\cCz}{\mathcal C_0}
\providecommand{\cCone}{\mathcal C_1}
\providecommand{\cA}{\mathcal A}
\providecommand{\cB}{\mathcal B}
\providecommand{\cL}{\mathcal L}
\providecommand{\cI}{\mathcal I}
\providecommand{\cJ}{\mathcal J}
\providecommand{\Db}{\mathrm D}
\providecommand{\Kb}{\mathrm K}
\providecommand{\Dminus}{\mathrm D^{-}}
\providecommand{\Kminus}{\mathrm K^{-}}
\providecommand{\Kbd}{\mathrm K^{b}}
\providecommand{\Mod}{\operatorname{Mod}}
\providecommand{\Lib}{\operatorname{Lib}}
\providecommand{\Loclib}{\operatorname{Loclib}}
\providecommand{\Qcoh}{\operatorname{Qcoh}}
\providecommand{\Coh}{\operatorname{Coh}}
\providecommand{\Ext}{\operatorname{Ext}}
\providecommand{\Hom}{\operatorname{Hom}}
\renewcommand{\Ob}{\operatorname{ob}}
\providecommand{\im}{\operatorname{Im}}
\providecommand{\coker}{\operatorname{Coker}}
\providecommand{\coh}{\operatorname{coh}}
\providecommand{\parf}{\operatorname{parf}}
\providecommand{\amp}{\operatorname{amp}}
\providecommand{\RHom}{\operatorname{RHom}}

%%%%%%%%%%  scan idx 173  |  folio 160  %%%%%%%%%%
\SGAPrintPage{160}%
\sgaStableTarget{sga6.e02.expose}\section*{\underline{Exposé II}\\[3pt] \underline{\underline{EXISTENCE OF GLOBAL RESOLUTIONS}}}
\begin{center}
\underline{by L. ILLUSIE}
\end{center}

\sgaStableTarget{sga6.e02.sec.1}\section*{1. General globalization criteria}

\sgaStableTarget{sga6.e02.sec.1.0}1.0. In this section, we are given a site $\underline{S}$, a flat abelian $\underline{S}$-category $\underline{C}$ (I 1.1), and a strictly full additive sub-$\underline{S}$-category $\underline{C}_0$ (I 1.1). We assume that $\underline{C}_0$ is locally quasi-liftable in $\underline{C}$ (I 1.2) and locally quasi-stable under kernels of epimorphisms (I 1.2).

\sgaStableTarget{sga6.e02.prop.1-1}\textbf{Proposition 1.1.} \uline{Let} $S\in\Ob \underline{S}$, \uline{and suppose that}:
\begin{enumerate}[label=(\roman*)]
\item $\underline{C}_{0S}$ \uline{is quasi-liftable in} $\underline{C}_S$ (I 1.2) ($\underline{C}_S$ \uline{being regarded as a category fibered over a punctual site});
\item \uline{every object of} $\underline{C}_S$ \uline{which is of} $\underline{C}_0$\uline{-finite type} (I 1.2) (\uline{resp.\ of} $\underline{C}_0$\uline{-finite presentation} (I 2.8)) \uline{is of} $\underline{C}_{0S}$\uline{-finite type}.
\end{enumerate}
\uline{Then}:
\begin{enumerate}[label=\alph*)]
\item $\underline{C}_{0S}$ \uline{is quasi-stable under kernels of epimorphisms} (I 1.2) ($\underline{C}_{0S}$ \uline{being regarded as a category fibered over a punctual site}).
\item \uline{Let} $E\in\Ob \mathrm D^{-}(\underline{C}_S)$ \uline{and} $n\in\mathbb Z$. \uline{If} $E$ \uline{is} $n$-$\underline{C}_0$\uline{-pseudo-coherent} (I 2.3), $E$ \uline{is} $n$-$\underline{C}_{0S}$\uline{-pseudo-coherent} (\uline{resp.} $(n+1)$-$\underline{C}_{0S}$\uline{-pseudo-coherent}). \uline{In other words}, \uline{we have}
\[
\mathrm D^{-}(\underline{C}_S)_{n-\underline{C}_0\text{-coh}}
=
\mathrm D^{-}(\underline{C}_S)_{n-\underline{C}_{0S}\text{-coh}}
\]
(\uline{resp.}
\[
\mathrm D^{-}(\underline{C}_S)_{n-\underline{C}_0\text{-coh}}
\subset
\mathrm D^{-}(\underline{C}_S)_{(n+1)-\underline{C}_{0S}\text{-coh}}
\subset
\mathrm D^{-}(\underline{C}_S)_{(n+1)-\underline{C}_0\text{-coh}},
\]
\uline{inclusions of strictly full subcategories of} $\mathrm D(\underline{C}_S)$, \uline{the inclusion on the right being trivial, moreover}).
\item \uline{For} $E\in\Ob\mathrm D^{-}(\underline{C}_S)$ \uline{to be pseudo-coherent relative to} $\underline{C}_0$ (I 2.3), \uline{it is necessary and sufficient that it be so relative to} $\underline{C}_{0S}$; \uline{in other words}, \uline{we have}
\[
\mathrm D^{-}(\underline{C}_S)_{\underline{C}_0\text{-coh}}
=
\mathrm D^{-}(\underline{C}_S)_{\underline{C}_{0S}\text{-coh}}.
\]
\uline{Moreover}, \uline{the canonical functor}
\[
\mathrm K^{-}(\underline{C}_{0S})/\mathrm K^{-,\varnothing}(\underline{C}_{0S})
\longrightarrow \mathrm D(\underline{C}_S),
\]
\uline{where} $\mathrm K^{-,\varnothing}(\underline{C}_{0S})$ \uline{denotes the thick subcategory of} $\mathrm K^{-}(\underline{C}_{0S})$ \uline{formed by}
%%%%%%%%%%  scan idx 174  |  folio 161  %%%%%%%%%%
\SGAPrintPage{161}%
\uline{acyclic complexes}, \uline{is fully faithful and has as essential image} $\mathrm D^{-}(\underline{C}_S)_{\underline{C}_0\text{-coh}}$. If $\underline{C}_{0S}$ \uline{is liftable in} $\underline{C}_S$ (I 1.2) (\uline{i.e.\ if the objects of} $\underline{C}_{0S}$ \uline{are projective}), \uline{the category} $\mathrm K^{-,\varnothing}(\underline{C}_{0S})$ \uline{is reduced to zero objects}, \uline{and consequently the functor}
\[
\mathrm K^{-}(\underline{C}_{0S})\longrightarrow \mathrm D(\underline{C}_S)
\]
\uline{is fully faithful} (\uline{and has as essential image} $\mathrm D^{-}(\underline{C}_S)_{\underline{C}_0\text{-coh}}$).
\end{enumerate}

\uline{Proof.} a) Let
\[
0\longrightarrow L\longrightarrow L^0\longrightarrow\ldots\longrightarrow L^m\longrightarrow0
\]
be an exact sequence in $\underline{C}_S$, with $L^i\in\Ob\underline{C}_{0S}$ for every $i$; we have to show that $L$ is of $\underline{C}_{0S}$-finite type. Now, by (I 2.2), $L$ is $\underline{C}_0$-pseudo-coherent, hence (I 2.9) of $\underline{C}_0$-finite presentation, and the assertion follows from (ii).

b) Let us show, by descending induction on $i$, that $E$ is $i$-$\underline{C}_{0S}$-pseudo-coherent for $i\ge n$ (resp. $i\ge n+1$). We may already assume that $E$ is strictly $(n+1)$-$\underline{C}_{0S}$-pseudo-coherent (resp. $(n+2)$-$\underline{C}_{0S}$-pseudo-coherent) (I 2.1), and we have to prove that $E$ is $n$-$\underline{C}_{0S}$-pseudo-coherent (resp. $(n+1)$-$\underline{C}_{0S}$-pseudo-coherent). We therefore have an exact sequence
\[
0\longrightarrow E'\longrightarrow E\longrightarrow E''\longrightarrow0,
\]
where $E'$ is strictly $\underline{C}_{0S}$-perfect (I 2.1) and $E''{}^i=0$ for $i>n$ (resp. $i>n+1$). Replacing $E$ by $E''$, we may therefore assume (I 2.6) that $E^i=0$ for $i>n$ (resp. $i>n+1$). Then, by (I 2.10 a)), $H^n(E)$ (resp. $H^{n+1}(E)$) is of $\underline{C}_0$-finite type (resp. of $\underline{C}_0$-finite presentation). Thus, by (ii), $H^n(E)$ (resp. $H^{n+1}(E)$) is of $\underline{C}_{0S}$-finite type. Applying (I 2.10 b)) again, we deduce that $E$ is $n$-$\underline{C}_{0S}$-pseudo-coherent (resp. $(n+1)$-$\underline{C}_{0S}$-pseudo-coherent), q.e.d. (N.B.\ hypothesis (i) and assertion a) were used implicitly in the references to (I 2.6 and 2.10)).

c) The first assertion follows from b), the second from (I 2.7); the last is evident.

\sgaStableTarget{sga6.e02.prop.1-2}\textbf{Proposition 1.2.} \uline{Suppose that} $\underline{C}_0$ \uline{is locally liftable in} $\underline{C}$ (I 1.2) \uline{and} \uline{locally stable under kernels of epimorphisms} (I 1.2), \uline{and make the}
%%%%%%%%%%  scan idx 175  |  folio 162  %%%%%%%%%%
\SGAPrintPage{162}%
\uline{hypotheses} (1.1 (i) \uline{and} (ii)
\uline{in the resp.\ case). Denote by} $\underline{C}_{1S}$ \uline{the strictly full subcategory of} $\underline{C}_S$ \uline{formed by objects which are locally in} $\underline{C}_0$\uline{. Then:}
\begin{enumerate}[label=\alph*)]
\item $\underline{C}_{1S}$ \uline{is stable under kernels of epimorphisms} (I 1.2) \uline{and quasi-liftable in} $\underline{C}_S$ (I 1.2) \uline{(as a category fibered over a punctual site).}
\item \uline{Let} $E\in\Ob\mathrm D(\underline{C}_S)$ \uline{and} $[m,n]$ \uline{be an interval in} $\mathbb Z$\uline{. For} $E$ \uline{to have} $\underline{C}_0$\uline{-perfect amplitude contained in} $[m,n]$\uline{, it is necessary and sufficient that} $E$ \uline{be isomorphic, in} $\mathrm D(\underline{C}_S)$\uline{, to a complex} $F$ \uline{such that} $F^i=0$ \uline{for} $i\notin[m,n]$, $F^i\in\Ob\underline{C}_{0S}$ \uline{for} $i>m$\uline{, and} $F^m\in\Ob\underline{C}_{1S}$.
\item \uline{The canonical functor}
\[
\mathrm K^b(\underline{C}_{1S})/\mathrm K^{b,\varnothing}(\underline{C}_{1S})
\longrightarrow \mathrm D(\underline{C}_S)
\]
\uline{is fully faithful and has as essential image the full subcategory of} $\mathrm D(\underline{C}_S)$ \uline{formed by complexes of finite} $\underline{C}_0$\uline{-perfect amplitude.}
\item \uline{If} $\underline{C}_{1S}$ \uline{is not only quasi-liftable but liftable in} $\underline{C}_S$\uline{, one may replace in} b) \uline{the expression \textquotedblleft}$E$ \uline{is isomorphic to} $F$ \uline{in} $\mathrm D(\underline{C}_S)$\uline{\textquotedblright{} by \textquotedblleft{}there exists a quasi-isomorphism} $F\to E$\uline{\textquotedblright, and in} c) \uline{the category} $\mathrm K^{b,\varnothing}(\underline{C}_{1S})$ \uline{is reduced to zero objects.}
\end{enumerate}

\uline{Proof.} a) It is clear that $\underline{C}_{1S}$ is stable under kernels of epimorphisms. On the other hand, the objects of $\underline{C}_{1S}$ are $\underline{C}_0$-pseudo-coherent, hence $\underline{C}_{0S}$-pseudo-coherent by (1.1 c)); the second assertion therefore follows from (I 2.14) (where we take the punctual site for $\underline{S}$).

b) It suffices to prove necessity. We already know, by (1.1 c)), that $E$ is $\underline{C}_{0S}$-pseudo-coherent. Since $E$ is acyclic in degree $>n$, we may therefore assume (I 2.10), after replacing $E$ by an isomorphic complex in $\mathrm D^{-}(\underline{C}_S)$, that $E$ is strictly $(m+1)$-$\underline{C}_{0S}$-pseudo-coherent and that $E^i=0$ for $i>n$. Since $E$ is acyclic in degree $<m$, we may further assume, after replacing $E$ by its truncation
\[
0\longrightarrow E^m/B^m\longrightarrow E^{m+1}\longrightarrow\ldots\longrightarrow E^n\longrightarrow0,
\]
that $E^i=0$ for $i<m$. Then, by (I 4.16), $E^m$ is locally an object of $\underline{C}_0$, hence an object of $\underline{C}_{1S}$, which proves the assertion.

%%%%%%%%%%  scan idx 176  |  folio 163  %%%%%%%%%%
\SGAPrintPage{163}%
c)
In view of (1.2 a)), the assertion follows from (1.2 b)) and (I 4.12.2).

d) The fact that $\mathrm K^{b,\varnothing}(\underline{C}_{1S})$ is reduced to zero objects is evident. The other assertion follows from (I 4.6) (applied over the punctual site).

\sgaStableTarget{sga6.e02.sec.2}\section*{2. Application to certain categories of Modules}
\sgaStableTarget{sga6.e02.sec.2.0}\subsection*{2.0. Notation}
Let $(\underline{S},A)$ be a ringed topos. For $U\in\Ob \underline{S}$, we write $\Mod(U)$ for the category of left $A_U$-Modules, and $\Lib(U)$ (resp. $\Loclib(U)$) for the full subcategory of $\Mod(U)$ formed by free $A_U$-Modules (resp. Modules locally direct factors of free Modules) of finite type.

\sgaStableTarget{sga6.e02.sec.2.1}\subsection*{2.1. Lifting lemmas}

\sgaStableTarget{sga6.e02.lem.2-1-1}\textbf{Lemma 2.1.1} (cf. EGA $0_{\mathrm I}$ 5.2.3). \uline{Let} $(\underline{S},A)$ \uline{be a ringed topos, and} $S\in\Ob \underline{S}$\uline{. We assume that} \var{$\underline{S}$ \uline{is} quasi-compact}{$S$ \uline{is} quasi-compact}{EN-VAR-111}\uline{, i.e.\ that every covering family} $(S_i\to S)$ \uline{is dominated by a finite covering family.}
\begin{enumerate}[label=\alph*)]
\item \uline{Let} $F$ \uline{be a left} $A_S$\uline{-Module which is the inductive limit, over a filtered category} $\Lambda$\uline{, of} $A_S$\uline{-Modules} $F_\lambda$\uline{, and let}
\[
u=\varinjlim u_\lambda:F\longrightarrow G
\]
\uline{be an epimorphism, where} $G$ \uline{is an} $A_S$\uline{-Module of finite type. Then there exists} $\lambda\in\Lambda$ \uline{such that} $u_\lambda$ \uline{is an epimorphism.}
\item \uline{Let} $F$ \uline{be a left} $A_S$\uline{-Module} generated by its sections \uline{(i.e.\ such that there exists an epimorphism} $A_S^{(I)}\to F$ \uline{for a suitable index set} $I$\uline{). Then, if} $F$ \uline{is of finite type, there exists an epimorphism} $A_S^p\to F$ \uline{for} $p\in\mathbb N$.
\end{enumerate}
The proof of a) is analogous to that of (loc.\ cit.) and is left as an exercise to the reader, and b) follows immediately from a).

\sgaStableTarget{sga6.e02.lem.2-1-2}\textbf{Lemma 2.1.2.} \uline{Let} $(\underline{S},A)$ \uline{be a ringed topos,} $\underline{C}$ \uline{the fibered} $\underline{S}$\uline{-category of left} $A$\uline{-Modules,} $\underline{C}_0$ \uline{the fibered subcategory of free} $A$\uline{-Modules of finite type} (I 1.3 d)), $S\in\Ob \underline{S}$, \var{$C'_S$}{$\underline{C}'_S$}{EN-VAR-112} \uline{a strictly full additive subcategory of} $\underline{C}_S$ \uline{containing} $\underline{C}_{0S}$ \uline{and stable under kernels, cokernels and extensions. We assume that}
%%%%%%%%%%  scan idx 177  |  folio 164  %%%%%%%%%%
\SGAPrintPage{164}%
\[
H^1(S,F)=0\quad
\text{\uline{for every} }F\in\Ob\underline{C}'_S.
\]
\uline{Then:}
\begin{enumerate}[label=\alph*)]
\item \uline{Every exact sequence in} $\underline{C}'_S$:
\[
0\longrightarrow E\longrightarrow F\longrightarrow G\longrightarrow0,
\]
\uline{such that} $\underline{\Hom}(G,E)
\in\Ob\underline{C}'_S$ \uline{and such that} $G$ \uline{is locally a direct factor of a free Module of finite type, is split. The category} $\underline{C}_{0S}$ \uline{is liftable in} $\underline{C}'_S$ (I 1.2).
\item \uline{Denote by} $\Mod(A(S))$ \uline{the category of} $A(S)$\uline{-modules. The functor}
\[
\Gamma_S:\underline{C}'_S\longrightarrow\Mod(A(S))
\]
\uline{is exact. Moreover:}
\begin{enumerate}[label=(\roman*)]
\item \uline{If every object} $F$ \uline{of} $\underline{C}'_S$ \uline{admits a finite global presentation}
\[
A_S^p\longrightarrow A_S^q\longrightarrow F\longrightarrow0,
\]
\uline{the functor} $\Gamma_S:\underline{C}'_S\longrightarrow\Mod(A(S))$ \uline{is fully faithful and has as essential image the full subcategory of} $\Mod(A(S))$ \uline{formed by} $A(S)$\uline{-modules of finite presentation;}
\item \uline{If} $\underline{C}'_S$ \uline{contains the free} $A_S$\uline{-Modules, if every object} $F$ \uline{of} $\underline{C}'_S$ \uline{admits a global presentation}
\[
A_S^{(I)}\longrightarrow A_S^{(J)}\longrightarrow F\longrightarrow0,
\]
\uline{and if the functor} $\Gamma_S$ \uline{commutes with filtered inductive limits}\SGAAuthorFootnote{1}{a hypothesis satisfied when $\underline{S}$ is locally of finite type (SGA 4 VI 1.2) or when $\underline{S}$ is the topos of sheaves of sets on a compact space and $S$ the final object (T.F II 3.10.1).}\uline{, the functor} $\Gamma_S:\underline{C}'_S\longrightarrow\Mod(A(S))$ \uline{is an equivalence of categories.}
\end{enumerate}
\end{enumerate}

\uline{Proof.} a) Since $G$ is locally a direct factor of a free Module of finite type, the spectral sequence
\[
E_2^{pq}=H^p(S,\underline{\Ext}^q(G,E))\Longrightarrow\Ext^{*}(G,E)
\]
degenerates (I 1.3.1) and gives an isomorphism
\[
H^1(S,\underline{\Hom}(G,E))\simeq\Ext^1(G,E),
\]
%%%%%%%%%%  scan idx 178  |  folio 165  %%%%%%%%%%
\SGAPrintPage{165}%
whence the
first assertion. The second is a particular case, since, for $G\in\Ob\underline{C}_{0S}$, i.e.\ $G$ free and of finite type, and $E\in\Ob\underline{C}'_S$, we have $\underline{\Hom}(G,E)\in\Ob\underline{C}'_S$.

b) The first assertion is trivial. Let us now prove, in case (i) or (ii), the full faithfulness of $\Gamma_S:\underline{C}'_S\longrightarrow\Mod(A(S))$. Let $F,G\in\Ob\underline{C}'_S$, and let
\sgaStableTarget{sga6.e02.eq.unnamed}\[
A_S^{(I)}\longrightarrow A_S^{(J)}\longrightarrow F\longrightarrow0
\tag{$*$}
\]
be a global presentation of $F$ (with $I$ and $J$ finite in case (i)). Applying to $(*)$ the functor $\Gamma_S$, which is exact on $\underline{C}'_S$ and commutes, in case (ii), with filtered inductive limits, we obtain an exact sequence
\sgaStableTarget{sga6.e02.eq.unnamed.occ2}\[
A(S)^{(I)}\longrightarrow A(S)^{(J)}\longrightarrow\Gamma_S(F)\longrightarrow0.
\tag{$**$}
\]
On the other hand, applying to $(*)$ (resp. $(**)$) the left-exact functor $\operatorname{Hom}_{A_S}(.,G)$ (resp. $\var{\operatorname{Hom}_{A(S)}(.,\Gamma_S(F))}{\operatorname{Hom}_{A(S)}(.,\Gamma_S(G))}{EN-VAR-113}$), we obtain a commutative diagram with exact rows:
% SGA display variant EN-VAR-114
\SGACaptureAppState{EN-VAR-114}%
\ifnum\SGAReaderForm=1\relax\SGAAppMark{EN-VAR-114}\fi
\ifnum\SGAReaderForm=0\relax
\[
\begin{tikzcd}[column sep=small]
0 \arrow[r] & \operatorname{Hom}_{A_S}(F,G) \arrow[r] \arrow[d] &
\operatorname{Hom}_{A_S}(A_S^{(J)},G) \arrow[r] \arrow[d] &
\operatorname{Hom}_{A_S}(A_S^{(I)},G) \arrow[d] \\
0 \arrow[r] & \operatorname{Hom}_{A(S)}(F,G) \arrow[r] &
\operatorname{Hom}_{A(S)}(A(S)^{(J)},\Gamma_S(G)) \arrow[r] &
\operatorname{Hom}_{A(S)}(A(S)^{(I)},\Gamma_S(G)).
\end{tikzcd}
\]
\else
\[
\begin{tikzcd}[column sep=small]
0 \arrow[r] & \operatorname{Hom}_{A_S}(F,G) \arrow[r] \arrow[d] &
\operatorname{Hom}_{A_S}(A_S^{(J)},G) \arrow[r] \arrow[d] &
\operatorname{Hom}_{A_S}(A_S^{(I)},G) \arrow[d] \\
0 \arrow[r] & \operatorname{Hom}_{A(S)}(\Gamma_S(F),\Gamma_S(G)) \arrow[r] &
\operatorname{Hom}_{A(S)}(A(S)^{(J)},\Gamma_S(G)) \arrow[r] &
\operatorname{Hom}_{A(S)}(A(S)^{(I)},\Gamma_S(G)).
\end{tikzcd}
\]
\fi
The two vertical arrows on the right \var{are visibly isomorphisms the same}{are visibly isomorphisms; the same}{EN-VAR-115} is therefore true of the vertical arrow on the left, which proves the full faithfulness of $\Gamma_S:\underline{C}'_S\longrightarrow\Mod(A(S))$. The characterization of the essential image, in case (i) or (ii), is immediate and is left to the reader.

\sgaStableTarget{sga6.e02.lem.2-1-3}\textbf{Lemma 2.1.3.} \uline{Let} $A$ \uline{and} $B$ \uline{be abelian categories, and} $\phi:A\to B$ \uline{an exact functor.}
\begin{enumerate}[label=\alph*)]
\item \uline{For} $M\in\Ob A$\uline{, the following conditions are equivalent:}
\begin{enumerate}[label=(\roman*)]
\item \uline{for every} $L\in\Ob A$\uline{, the canonical arrow (induced by} $\phi$)
\[
\Ext^n_{A}(L,M)\longrightarrow\Ext^n_{B}(\phi(L),\phi(M))
\]
\uline{is an isomorphism for every} $n$;
\item \uline{for every} $L\in\Ob A$\uline{, the canonical arrow}
%%%%%%%%%%  scan idx 179  |  folio 166  %%%%%%%%%%
\SGAPrintPage{166}%
\[
\operatorname{Hom}_{A}(L,M)\longrightarrow\operatorname{Hom}_{B}(\phi(L),\phi(M))
\]
\uline{is an
isomorphism, and for every} $n\ge1$ \uline{and every} $x\in\Ext^n_{B}(\phi(L),\phi(M))$ \uline{there exists an epimorphism} $L'\to L$ \uline{in} $A$ \uline{which effaces} $x$.
\end{enumerate}
\item \uline{The following conditions are equivalent}:
\begin{enumerate}[label=(\roman*)]
\item \uline{for all} $L,M\in\Ob A$ \uline{and every} $n\in\mathbb Z$\uline{, the canonical arrow}
\[
\Ext^n_{A}(L,M)\longrightarrow\Ext^n_{B}(\phi(L),\phi(M))
\]
\uline{is an isomorphism};
\item \uline{the canonical functor} (\uline{induced by} $\phi$)
\[
\mathrm D^b(A)\longrightarrow\mathrm D^b(B)
\]
\uline{is fully faithful and has as essential image the full subcategory of} $\mathrm D^b(B)$ \uline{formed by complexes} $F$ \uline{such that} $H^i(F)$ \uline{is in the essential image of} $\phi$ \uline{for every} $i$.
\end{enumerate}
\end{enumerate}

\uline{Proof.} a) For the implication (i)$\Rightarrow$(ii), it suffices to show the following: given $n\ge1$, $L\in\Ob A$, and $y\in\Ext^n_{A}(L,M)$, there exists an epimorphism $L'\to L$ which effaces $y$. For this, one can use Yoneda's interpretation ([18] III 3.2): $y$ is the class of an exact sequence
\[
0\longrightarrow M\longrightarrow E_{n-1}\longrightarrow\ldots\longrightarrow E_0\longrightarrow L\longrightarrow0;
\]
the epimorphism $E_0\to L$ effaces $y$ ([18] III 3.2.5 and 3.2.8). Let us prove (ii)$\Rightarrow$(i). For brevity, put
\[
\Ext^n_{A}(.,M)=F^n,\qquad
\Ext^n_{B}(\phi(.),\phi(M))=G^n.
\]
The effaceability hypothesis means that, for $n\ge1$ and $L\in\Ob A$, we have
\[
G^n(L)=\varinjlim\operatorname{Coker}(G^{n-1}(L_0)\to G^{n-1}(L_1)),
\]
the inductive limit being taken over the set of exact sequences
%%%%%%%%%%  scan idx 180  |  folio 167  %%%%%%%%%%
\SGAPrintPage{167}%
$0\to L_1\to L_0\to L\to0$. From
what we saw above, we likewise have
\[
F^n(L)=\varinjlim\operatorname{Coker}(F^{n-1}(L_0)\to F^{n-1}(L_1)).
\]
since $\phi=\phi^0:F^0\to G^0$ is an isomorphism, it follows, by induction on $n$, that the morphism $\phi^n:F^n\to G^n$ induced by $\phi$ is an isomorphism for every $n$.

b) The equivalence (i)$\Leftrightarrow$(ii) follows from an immediate dévissage, left to the reader.

\sgaStableTarget{sga6.e02.rem.2-1-3-1}\textbf{Remarks 2.1.3.1.} There is of course a dual equivalence to a), which we leave the reader to state. On the other hand, when $A$ and $B$ have enough injectives, conditions (i) and (ii) of b) are equivalent to
\begin{enumerate}[label=(\roman*),start=3]
\item for $L\in\Ob\mathrm D^{-}(A)$ and $M\in\Ob\mathrm D^{+}(A)$, the canonical arrow
\[
\underline{\underline{R}}\!\operatorname{Hom}(L,M)\longrightarrow
\underline{\underline{R}}\!\operatorname{Hom}(\phi(L),\phi(M))
\]
is an isomorphism,
\end{enumerate}
as one sees immediately by a spectral-sequence argument.

\sgaStableTarget{sga6.e02.sec.2.2}\subsection*{2.2. Noetherian schemes and divisorial schemes}

\sgaStableTarget{sga6.e02.notation.2-2-1}\textbf{Notation 2.2.1.} Let $S$ be a scheme. We write $\Qcoh(S)$ (resp. $\Coh(S)$) for the (abelian) category of quasi-coherent (resp. coherent) $\underline{O}_S$-Modules. Let $[m,n]$ be an interval in $\mathbb Z$. We shall say that an object $E$ of $\mathrm D(\Qcoh(S))$ is $n$-pseudo-coherent (resp. pseudo-coherent, resp. of perfect amplitude contained in $[m,n]$, resp. perfect) when $E$ is $n$-pseudo-coherent (resp. pseudo-coherent, resp. of perfect amplitude contained in $[m,n]$, resp. perfect) relative to the category fibered over the Zariski site of $S$ defined by $U\mapsto\Loclib(U)$ ((I 2.3), (I 4.7) and (2.0)). We denote by $\mathrm D^{-}(\Qcoh(S))_{\coh}$ (resp. $\mathrm D(\Qcoh(S))_{\parf}$) the full triangulated subcategory of $\mathrm D(\Qcoh(S))$ formed by pseudo-coherent (resp. perfect) complexes.

\sgaStableTarget{sga6.e02.prop.2-2-2}\textbf{Proposition 2.2.2}\SGAAuthorFootnote{raisedstar}{This ancillary result will not be used in the remainder of the Seminar: the reader is advised to skip it on a first reading.}. \uline{Let} $S$ \uline{be a noetherian scheme. The canonical}
\uline{functor}
%%%%%%%%%%  scan idx 181  |  folio 168  %%%%%%%%%%
\SGAPrintPage{168}%
\[
\mathrm D^{-}(\Coh(S))\longrightarrow \mathrm D(\Qcoh(S))
\]
\uline{is fully faithful and has as essential image} $\var{\mathrm D^{-}(\Qcoh(S)_{\coh}}{\mathrm D^{-}(\Qcoh(S))_{\coh}}{EN-VAR-116}$.

\uline{Proof.} We apply (1.1), taking for $\underline{S}$ the Zariski site of $S$, and for $\underline{C}$ (resp. $\underline{C}_0$) the fibered category $U\mapsto\Qcoh(U)$ (resp. $\Coh(U)$). Since $\underline{O}_S$ is coherent, the objects of $\underline{C}_0$ are pseudo-coherent (I 3.5), so an object of $\mathrm D(\Qcoh(U))$ ($U$ = open subset of $S$) is pseudo-coherent (in the sense of (2.2.1)) precisely when it is so relative to $\underline{C}_0$ (I 2.14.1). It therefore suffices, by (1.1 c)), to show that conditions (i) and (ii) in the non-resp.\ case of (1.1) are satisfied. For (ii), this is trivial, since a quasi-coherent Module which is a quotient of a coherent Module is coherent, and the notion of coherence is local. For (i), we use (2.1.1 a)) and the fact that every quasi-coherent $\underline{O}_S$-Module is a filtered inductive limit of coherent $\underline{O}_S$-Modules (EGA I 9.4.9).

Admitting (3.5) and (3.7), which are independent of the present section, one easily deduces:

\sgaStableTarget{sga6.e02.cor.2-2-2-1}\textbf{Corollary 2.2.2.1.} \uline{Let} $S$ \uline{be a noetherian scheme (resp.\ noetherian and of finite dimension). The canonical functor}
\[
\mathrm D^b(\Coh(S))\longrightarrow\mathrm D(S)
\]
\[
(\text{resp. }\mathrm D^{-}(\Coh(S))\longrightarrow\mathrm D(S))
\]
\uline{is fully faithful and has as essential image} $\mathrm D^b(S)_{\coh}$ (resp. $\mathrm D^{-}(S)_{\coh}$) (I 2.15).

\sgaStableTarget{sga6.e02.prop.2-2-3}\textbf{Proposition 2.2.3} (cf. EGA II 4.5.2 and 4.5.5). \uline{Let} $S$ \uline{be a quasi-compact and quasi-separated scheme, and let} $(L_i)_{i\in I}$ \uline{be a family of invertible sheaves on} $S$\uline{. For every} $\underline{O}_S$\uline{-Module} $F$\uline{, and every pair} $(i,n)\in I\times\mathbb Z$\uline{, we denote by} $F_{(i,n)}$ \uline{the sub-}$\underline{O}_S$\uline{-Module of} $F\otimes L_i^{\otimes n}$ \uline{generated by the sections of} $F\otimes L_i^{\otimes n}$ \uline{over} $S$\uline{. The following conditions are equivalent}:
\begin{enumerate}
%%%%%%%%%%  scan idx 182  |  folio 169  %%%%%%%%%%
\item[(i)] \SGAPrintPage{169}%
\uline{The
open sets} $S_f$\uline{, for} $f\in\Gamma(S,L_i^{\otimes n})$\uline{,} $i\in I$ \uline{and} $n\in\mathbb N$\uline{, form a basis for the topology of} $S$.
\item[(i$'$)] \uline{Those of the} $S_f$ \uline{which are affine cover} $S$.
\item[(ii)] \uline{For every quasi-coherent} $\underline{O}_S$\uline{-Module} $F$ \uline{of finite type, there exists a family} $(r_i)_{i\in I}$ (resp. $(n_i)_{i\in I}$) \uline{of integers} $\ge0$ (resp. $>0$), \uline{and an epimorphism}
\[
\coprod_{i\in I}\underline{O}_S^{r_i}\otimes L_i^{\otimes-n_i}\longrightarrow F.
\]
\item[(iii)] \uline{For every} $F\in\Ob\Qcoh(S)$\uline{, we have}
\[
F=\sum_{(i,n)\in I\times\mathbb N^+}F_{(i,n)}\otimes L_i^{\otimes-n}
\]
(\uline{sum (not direct) of subsheaves of} $F$).
\item[(iii$'$)] \uline{The same condition as} (iii), \uline{but with} $F$ \var{\uline{a quasi-coherent Ideal of} $\underline{O}_S$}{\uline{a quasi-coherent Ideal of} $\underline{O}_S$.}{EN-VAR-117}
\end{enumerate}

\uline{Proof.} It consists of a simple paraphrase of (loc.\ cit.); we give it briefly for the reader's convenience.

(i)$\Rightarrow$(i$'$): For $x\in S$, choose an affine open subset $U$ containing $x$. An open subset $S_f$ such that $x\in S_f\subset U$, with $f\in\Gamma(S,L_i^{\otimes n})$ ($n\ge0$), is necessarily affine, as follows from

\sgaStableTarget{sga6.e02.lem.2-2-3-1}\textbf{Lemma 2.2.3.1.} (EGA II 5.5.8). \uline{Let} $X$ \uline{be a scheme,} $L$ \uline{an invertible} $\underline{O}_X$\var{\uline{-Module} $f\in\Gamma(X,L)$}{\uline{-Module,} $f\in\Gamma(X,L)$}{EN-VAR-118}\uline{. The canonical immersion} $i:X_f\to X$ \uline{is affine}.

\uline{Proof.} This is trivial: the question is local on $X$, so we may assume that $X$ is affine and $L\simeq\underline{O}_X$, and we are done.

(i$'$)$\Rightarrow$(ii): Take a finite covering of $S$ by affine open subsets $S_{f_{ij}}$, where $f_{ij}\in\Gamma(S,L_i^{\otimes m_{ij}})$; put $S_i=\bigcup_j S_{f_{ij}}$. Replacing the $f_{ij}$ by suitable powers, we may assume that the $m_{ij}$ corresponding to the same index $i$ are equal to an integer $m_i$. The sheaf
%%%%%%%%%%  scan idx 183  |  folio 170  %%%%%%%%%%
\SGAPrintPage{170}%
$F\,|\,S_{f_{ij}}$ is
generated by finitely many sections $h_{ijk}$ over $S_{f_{ij}}$. There exists (EGA I 9.3.1 and IV 1.7.5) an integer $q_i\ge0$ such that $h_{ijk}\otimes f_{ij}^{\otimes q_i}$ extends to a section of $F\otimes L_i^{\otimes q_im_i}$ over $S$. In other words, $F\otimes L_i^{\otimes q_im_i}$ is generated over $S_i$ by finitely many of its sections over $S$. Proceeding as in the proof of (EGA II 4.5.5), we deduce that there exists an integer $n_i$, which we may take to be $>0$, such that, for every $n\ge n_i$, $F\otimes L_i^{\otimes n}$ is generated over $S_i$ by finitely many of its sections over $S$. In particular, there exist an integer $r_i\ge0$ and a morphism
\[
\underline{O}_S^{r_i}\otimes L_i^{\otimes-n_i}\longrightarrow F
\]
which induces an epimorphism over $S_i$. The implication (i$'$)$\Rightarrow$(ii) follows immediately.

(ii)$\Rightarrow$(iii): Since every quasi-coherent sheaf on $S$ is an inductive limit of its quasi-coherent subsheaves of finite type (EGA I 9.4.9 and IV 1.7.7), we may restrict ourselves to proving (iii) for $F$ of finite type. There then exist, for $i\in I$, morphisms
\[
u_i:\underline{O}_S^{r_i}\otimes L_i^{\otimes-n_i}\longrightarrow F
\]
such that $\sum u_i$ is an epimorphism. Now we have
\[
\begin{aligned}
\operatorname{Im}(u_i)&=\operatorname{Im}\bigl(u_i\otimes\operatorname{Id}(L_i^{\otimes n_i})\otimes\operatorname{Id}(L_i^{\otimes-n_i})\bigr)\\
&=\operatorname{Im}\bigl(u_i\otimes\operatorname{Id}(L_i^{\otimes n_i})\bigr)\otimes L_i^{\otimes-n_i}
\end{aligned}
\]
But $\operatorname{Im}(u_i\otimes\operatorname{Id}(L_i^{\otimes n_i}))$ is a subsheaf of $F_{(i,n_i)}$, whence the assertion.

(iii)$\Rightarrow$(iii$'$): This is trivial.

(iii$'$)$\Rightarrow$(i): Let $x\in S$, $U$ be an open subset containing $x$, and $F$ a
%%%%%%%%%%  scan idx 184  |  folio 171  %%%%%%%%%%
\SGAPrintPage{171}%
quasi-coherent Ideal
defining $S-U$. By (iii$'$), there exist $(i,n)\in I\times\mathbb N^+$ and a section $f$ of $F\otimes L_i^{\otimes n}$ over $S$ such that $f(x)\ne0$. We have $x\in S_f\subset U$, which proves (i).

\sgaStableTarget{sga6.e02.def.2-2-4}\textbf{Definition 2.2.4.} \uline{Let} $S$ \uline{be a quasi-compact and quasi-separated scheme, and let} $(L_i)_{i\in I}$ \uline{be a family of invertible sheaves on} $S$\uline{. We say that the} $L_i$ \uline{form an} ample family\SGAAuthorFootnote{raisedstar}{This notion, in form (i$'$), was introduced by M.~BORELLI and S.~L.~KLEIMAN; the terminology \textquotedblleft ample family\textquotedblright{} is proposed by the editor.} \uline{when the equivalent conditions of} (2.2.3) \uline{hold. (In particular, a family consisting of a single sheaf} $L$ \uline{is ample precisely when} $L$ \uline{is ample in the usual sense} (EGA II 4.5.3)).

\sgaStableTarget{sga6.e02.def.2-2-5}\textbf{Definition 2.2.5.} \uline{A scheme} $S$ \uline{is said to be} divisorial \uline{when} $S$ \uline{is quasi-compact and quasi-separated and there exists an ample family} (2.2.4) \uline{of invertible} $\underline{O}_S$\uline{-Modules.}

\sgaStableTarget{sga6.e02.lem.2-2-6}\textbf{Lemma 2.2.6.} \uline{Let} $X$ \uline{be a locally noetherian scheme. Let} $U$ \uline{be an open subset of} $X$ \uline{such that the inclusion} $U\to X$ \uline{is affine (for example} $U$ \uline{affine and} $X$ \uline{separated). Then} $X-U$ \uline{is purely of codimension} $0$ \uline{or} $1$ \uline{(i.e.\ at every maximal point} $x$ \uline{of} $X-U$ \uline{we have} $\dim(\underline{O}_{X,x})=0$ or $1$\uline{), hence purely of codimension} $1$ \uline{when} $U$ \uline{contains the maximal points of} $X$.

\uline{Proof.} Put $Y=X-U$. Let $x$ be a maximal point of $Y$, and make the base change $X'=\operatorname{Spec}(\underline{O}_{X,x})\to X$: we obtain a cartesian diagram:
\[
\begin{tikzcd}
U' \arrow[r] \arrow[d] & X' \arrow[d] & Y' \arrow[l] \arrow[d] \\
U \arrow[r] & X & Y \arrow[l]
\end{tikzcd},
\]
where $\var{Y'=X'-U}{Y'=X'-U'}{EN-VAR-119}$. Since $x$ is a maximal point of $Y$, $Y'$ is precisely the closed point of $X'$. On the other hand, since $X'$ is affine, $U'$ is an affine open subset of $X'$. Let us show that the dimension $n$ of $X'$ is equal to $0$ or $1$. Indeed, we know (SGA 2 V 5.7) that $H^n_{Y'}(X',\underline{O}_{X'})\ne0$. But, for $n\ge2$, we have
\[
H^n_{Y'}(X',\underline{O}_{X'})\simeq H^{n-1}(U',\underline{O}_{U'}),
\]
whereas
\[
H^{n-1}(U',\underline{O}_{U'})=0,
\]
%%%%%%%%%%  scan idx 185  |  folio 172  %%%%%%%%%%
\SGAPrintPage{172}%
since $U'$ is affine and $n-1\ge1$, whence a contradiction.

\sgaStableTarget{sga6.e02.rem.2-2-6-1}\textbf{Remark 2.2.6.1.} Without the separation hypothesis on $X$, (2.2.6) for an affine open subset $U$ of $X$ would be false, as shown by the example of the affine plane with the origin doubled (EGA I 5.5.11).

\sgaStableTarget{sga6.e02.prop.2-2-7}\textbf{Proposition 2.2.7.} \uline{Let} $S$ \uline{be a noetherian, separated, locally factorial scheme (i.e.\ one whose local rings are factorial). Then} $S$ \uline{is divisorial} (2.2.5).

In particular (EGA IV 21.11.1):

\sgaStableTarget{sga6.e02.cor.2-2-7-1}\textbf{Corollary 2.2.7.1.} \uline{Every noetherian regular separated scheme is divisorial.}

\uline{Proof} of 2.2.7. We may assume that $S$ is integral. Let $(S_i)_{i\in I}$ be a finite covering of $S$ by nonempty affine open subsets. By (2.2.6), $S-S_i$ is purely of codimension $1$, hence, by (EGA IV 21.7.4), is the support of a positive divisor. In other words, for each $i\in I$, there exist an invertible sheaf $L_i$ and a section $f_i\in\Gamma(S,L_i)$ such that $S_i=S_{f_i}$. The family $(L_i)$ is then ample (2.2.4) by (2.2.3 (i$'$)). This completes the proof.

\sgaStableTarget{sga6.e02.rem.2-2-7-2}\textbf{Remark 2.2.7.2.} As in (2.2.6), the separation hypothesis on $S$ is essential. Indeed, again take for $S$ the affine plane with the origin doubled. Recall that $S$ is obtained by gluing two copies of the affine plane $\underline{\underline{E}}^2_k$ ($k$, a field) by the identity on the complement $U$ of the origin $O$; we therefore have a canonical projection
\[
p:S\longrightarrow\underline{\underline{E}}^2_k,
\]
which induces an isomorphism over $U$, while $p^{-1}(O)$ consists of two points $O_1$ and $O_2$. An invertible sheaf on $S$ is obtained by gluing two invertible sheaves on $\underline{\underline{E}}^2_k$ by an isomorphism over $U$; since every invertible sheaf on $\underline{\underline{E}}^2_k$ is trivial, an invertible sheaf on $S$ is therefore defined by an element $g\in\Gamma(U,\underline{O}_U^*)$; but $g$ is necessarily a constant, since the restriction $\Gamma(\underline{\underline{E}}^2_k,\underline{O}_{\underline{\underline{E}}^2_k})\to\Gamma(U,\underline{O}_U)$ is an isomorphism, $O$ being of depth $2$ in the plane; thus every
%%%%%%%%%%  scan idx 186  |  folio 173  %%%%%%%%%%
\SGAPrintPage{173}%
invertible sheaf on $S$ is trivial. On the other hand, it is immediate that $p$ induces an isomorphism
\[
p^*:\Gamma(\mathbb E^2_k,\underline{O}_{\mathbb E^2_k})\longrightarrow\Gamma(S,\underline{O}_S).
\]
Consequently, every global section of $\underline{O}_S$ which vanishes at $O_1$ also vanishes at $O_2$, so the $S_f$, for $f\in\Gamma(S,\underline{O}_S)$, do not form a basis for the topology of $S$.

\sgaStableTarget{sga6.e02.thm.2-2-8}\textbf{Lemma 2.2.8.} \uline{Let} $S$ \uline{be a divisorial scheme} (2.2.5), $(L_i)_{i\in I}$ \uline{an ample family of invertible} $\underline{O}_S$\uline{-Modules. Denote by} $\underline{L}(S)$ \uline{the strictly full additive subcategory of} $\Qcoh(S)$ (2.2.1) \uline{generated by the} $L_i^{\otimes-n}$ \uline{for} $(i,n)\in I\times\mathbb N^+$\uline{. Let} $E\in\Ob\mathrm D(\Qcoh(S))$.
\begin{enumerate}[label=\alph*)]
\item \uline{The categories} $\underline{L}(S)$ \uline{and} $\Loclib(S)$ (2.0) \uline{are quasi-liftable in} $\Qcoh(S)$ (I 1.2) \uline{and quasi-stable under kernels of epimorphisms} (I 1.2).
\item \uline{Let} $n\in\mathbb Z$\uline{. For} $E$ \uline{to be} $n$\uline{-pseudo-coherent (resp.\ pseudo-coherent)} (2.2.1), \uline{it is necessary and sufficient that} $E$ \uline{be} $n$\uline{-pseudo-coherent (resp.\ pseudo-coherent) relative to} $\underline{L}(S)$ \uline{or} $\Loclib(S)$ (I 2.3)\uline{. The canonical functors}
\[
\mathrm K^{-}(\underline{L}(S))/\mathrm K^{-,\varnothing}(\underline{L}(S))\longrightarrow \mathrm D(\Qcoh(S))
\]
\uline{and}
\[
\mathrm K^{-}(\Loclib(S))/\mathrm K^{-,\varnothing}(\Loclib(S))\longrightarrow \mathrm D(\Qcoh(S))
\]
\uline{are fully faithful and have as essential image} $\mathrm D^{-}(\Qcoh(S))_{\coh}$ (2.2.1).
\item \uline{Let} $[m,n]$ \uline{be an interval in} $\mathbb Z$\uline{. For} $E$ \uline{to have perfect amplitude contained in} $[m,n]$ (2.2.1)\uline{, it is necessary and sufficient that} $E$ \uline{be isomorphic, in} $\mathrm D(\Qcoh(S))$\uline{, to a complex} $F$ \uline{such that} $F^i=0$ \uline{for} $i\notin[m,n]$, $F^i\in\Ob \underline{L}(S)$ \uline{for} $i>m$, \uline{with} $F^m$ \uline{locally free of finite type. The canonical functor}
\[
\mathrm K^b(\Loclib(S))/\mathrm K^{b,\varnothing}(\Loclib(S))\longrightarrow \mathrm D(\Qcoh(S))
\]
%%%%%%%%%%  scan idx 187  |  folio 174  %%%%%%%%%%
\SGAPrintPage{174}%
is
fully faithful and has as essential image $\mathrm D(\Qcoh(S))_{\parf}$ (2.2.1).
\end{enumerate}

\uline{Proof.} Let $F\to G$ be an epimorphism in $\Qcoh(S)$, where $G$ is of finite type. Since $F$ is an inductive limit of its quasi-coherent subsheaves of finite type (EGA I 9.4.9), there exists, by (2.1.1 a)), a quasi-coherent subsheaf $F'$ of $F$ of finite type such that the composite $F'\to F\to G$ is an epimorphism. By (2.2.3 (ii)), there exists an epimorphism $E\to F'$, with $E\in\Ob \underline{L}(S)$, so the composite $E\to F'\to F$, composed with $F\to G$, is an epimorphism, which proves that the category of quasi-coherent $\underline{O}_S$-Modules of finite type is quasi-liftable in $\Qcoh(S)$; a fortiori, the same holds for the category $\Loclib(S)$ or $\underline{L}(S)$. This proves the first assertion of a). For the second, we apply (1.1), taking for $\underline{S}$ the Zariski site of $S$, and for $\underline{C}$ (resp. $\underline{C}_0$) the fibered category $U\mapsto\Qcoh(U)$ (resp. $\underline{L}(U)$) ($U$ denotes a variable open subset of $S$, and $\underline{L}(U)$ the strictly full additive subcategory of $\Qcoh(U)$ generated by the $(L_i^{\otimes-n}\,|\,U)$). Since every object of $\underline{L}$ is pseudo-coherent, the notions of $n$-pseudo-coherence (resp. pseudo-coherence) (2.2.1) in the fibered category $\mathrm D(\Qcoh)$ coincide with the same notions relative to $\underline{L}$ (I 2.14). Condition (ii) (in the non-resp.\ case) of (1.1) being satisfied by (2.2.3 (ii)), we deduce, by (1.1 a)), the second part of a) concerning $\underline{L}(S)$; the assertion concerning $\Loclib(S)$ is proved similarly. Assertions b) follow immediately from (1.1 b) and c)), taking into account the fact that, when $E\in\Ob\mathrm D(\Qcoh(S))$ is $n$-pseudo-coherent, $E$ is an object of $\mathrm D^{-}(\Qcoh(S))$, since $S$ is quasi-compact. An object of the fibered category $\Qcoh$ is locally in $\underline{L}$ precisely when it is locally free of finite type: the category $\underline{L}$ is therefore locally liftable in $\Qcoh$ and locally stable under kernels of epimorphisms. Consequently, \var{assertions c)) follow}{assertions c) follow}{EN-VAR-120} from (1.2) applied with $\underline{C}=\Qcoh$, $\underline{C}_0=\underline{L}$, $\underline{C}_{1S}=\Loclib(S)$. \var{This completes the proof}{This completes the proof.}{EN-VAR-121}

\sgaStableTarget{sga6.e02.prop.2-2-9}\textbf{Proposition 2.2.9.} \uline{The hypotheses and notation are those of} (2.2.8)\uline{, but we further assume that} $S$ \uline{is separated or noetherian.}
\begin{enumerate}[label=\alph*)]
\item \uline{The canonical functors}
\[
\mathrm K^{-,b}(\underline{L}(S))/\mathrm K^{-,\varnothing}(\underline{L}(S))\longrightarrow\mathrm D(S)
\]
and
\[
\mathrm K^{-,b}(\Loclib(S))/\mathrm K^{-,\varnothing}(\Loclib(S))\longrightarrow\mathrm D(S)
\]
%%%%%%%%%%  scan idx 188  |  folio 175  %%%%%%%%%%
\SGAPrintPage{175}%
\uline{are fully faithful and have as essential image} $\mathrm D^b(S)_{\coh}$ (I 2.15)\uline{. Moreover, when} $S$ \uline{has \textquotedblleft finite cohomological dimension\textquotedblright{} in the sense of} (3.7)\SGAAuthorFootnote{1}{Section~3 is independent of the rest of the exposé.} \uline{(for example, when the underlying space of} $S$ \uline{is noetherian and of finite Krull dimension), the functors}
\[
\mathrm K^{-}(\underline{L}(S))/\mathrm K^{-,\varnothing}(\underline{L}(S))\longrightarrow
\mathrm K^{-}(\Loclib(S))/\mathrm K^{-,\varnothing}(\Loclib(S))\longrightarrow
\mathrm D^{-}(S)_{\coh}
\]
\uline{are equivalences of categories.}
\item \uline{The canonical functor}
\[
\mathrm K^b(\Loclib(S))/\mathrm K^{b,\varnothing}(\Loclib(S))\longrightarrow\mathrm D(S)
\]
\uline{is fully faithful and has as essential image} $\mathrm D(S)_{\parf}$ (I 4.9)\uline{. For} $E\in\Ob\mathrm D(S)$ \uline{to have perfect amplitude contained in} $[m,n]$ (I 4.8)\uline{, it is necessary and sufficient that} $E$ \uline{be isomorphic, in} $\mathrm D(S)$\uline{, to a complex} $F$ \uline{such that} $F^i=0$ \uline{for} $i\notin[m,n]$, $\var{F^i\in\Ob \underline{L}(S))}{F^i\in\Ob \underline{L}(S)}{EN-VAR-122}$ \uline{for} $i>m$, \uline{with} $F^m$ \uline{locally free of finite type.}
\end{enumerate}

\uline{Proof.} In view of (3.5) and (3.7)\SGAAuthorFootMark{1}, this is an immediate corollary of (2.2.8): the only point to verify, which follows trivially from (3.5) (resp. (3.7)), is that, for an object $E$ of $\mathrm D^b(\Qcoh(S))$ (resp. $\mathrm D(\Qcoh(S))$ when $S$ has finite cohomological dimension in the sense of (3.7)), the notions of pseudo-coherence and perfection of (2.2.1) coincide with the usual notions (I 2.15 and 4.9) when $E$ is regarded as an object of $\mathrm D(S)$.

\sgaStableTarget{sga6.e02.rem.2-2-10}\textbf{Remark 2.2.10.} a) When $S$ is an affine scheme with ring $A$, the family consisting of the sheaf $\underline{O}_S$ is ample. For this family, we have $\underline{L}(S)=\Lib(S)$. It is also clear that the categories $\mathrm K^{-,\varnothing}(\Lib(S))$ and $\mathrm K^{-,\varnothing}(\Loclib(S))$ are reduced to zero objects. We therefore conclude from (2.2.9) that, writing $\Mod(A)$ for the category of $A$-modules and $\mathrm D(A)=\mathrm D(\Mod(A))$, the functor $\underline{\underline{R}}\Gamma_S:\mathrm D^+(S)\longrightarrow\mathrm D^+(A)$ induces a fully faithful functor $\mathrm D^b(S)_{\coh}\longrightarrow\mathrm D(A)$ (resp. $\mathrm D(S)_{\parf}\longrightarrow\mathrm D(A)$) whose essential image is
%%%%%%%%%%  scan idx 189  |  folio 176  %%%%%%%%%%
\SGAPrintPage{176}%
the
triangulated subcategory $\mathrm D^b(A)_{\coh}$ (resp. $\mathrm D(A)_{\parf}$) formed by complexes with bounded cohomology and pseudo-coherent (resp. perfect) relative to the subcategory $\Lib(A)$ (resp. $\Loclib(A)$) of free (resp. projective) modules of finite type. We have commutative diagrams of equivalences of categories
\[
\begin{tikzcd}
\mathrm K^{-,b}(\Lib(S)) \arrow[r] \arrow[d,"\Gamma_S"] &
\mathrm K^{-,b}(\Loclib(S)) \arrow[r] \arrow[d,"\Gamma_S"] &
\mathrm D^b(S)_{\coh} \arrow[d,"\mathbb R\Gamma_S"] \\
\mathrm K^{-,b}(\Lib(A)) \arrow[r] &
\mathrm K^{-,b}(\Loclib(A)) \arrow[r] &
\mathrm D^b(A)_{\coh},
\end{tikzcd}
\]
\[
\begin{tikzcd}
\mathrm K^b(\Loclib(S)) \arrow[r] \arrow[d,"\Gamma_S"] &
\mathrm D(S)_{\parf} \arrow[d,"\mathbb R\Gamma_S"] \\
\mathrm K^b(\Loclib(A)) \arrow[r] &
\mathrm D(A)_{\parf},
\end{tikzcd}
\]
(where, in the upper diagram, one may replace $b$ by $-$ when $S$ has finite cohomological dimension in the sense of (3.7)).

b) Let $S$ be a scheme. It follows from a) that a complex of $\underline{O}_S$-Modules which is pseudo-coherent and has bounded cohomology is locally isomorphic (in $\mathrm D(S)$) to a strictly pseudo-coherent complex, which was by no means evident a priori (cf. (I 2.3.1)).

\sgaStableTarget{sga6.e02.sec.2.3}\subsection*{2.3. Soft sheaves}

\sgaStableTarget{sga6.e02.para.2-3-1}2.3.1. Recall some definitions and results from [T.F]. Let $X$ be a paracompact space. A sheaf of sets $F$ on $X$ is said to be \uline{soft} when, for every closed subset $Y$ of $X$, every section of $F$ over $Y$ extends to a section of $F$ over $X$. Soft abelian sheaves are acyclic, i.e.\ when $F$ is soft, we have $H^i(X,F)=0$ for $i>0$ (T.F. II 4.4.3).

Let $A$ be a sheaf of rings on $X$. When $F$ is a soft $A$-Module, $F$ is generated by its sections (cf. 2.1.1): indeed, the homomorphism $A_X^{(I)}\to F$ defined by $I=H^0(X,F)$ is clearly an epimorphism.
%%%%%%%%%%  scan idx 190  |  folio 177  %%%%%%%%%%
\SGAPrintPage{177}%
On the other hand, when $A$ is soft, every $A$-Module is soft (T.F II 3.7.1). As examples of soft sheaves of rings, we mention the sheaf of continuous functions (real or complex) on a paracompact space, and the sheaf of $C^r$ functions on a manifold of class $C^r$.

\sgaStableTarget{sga6.e02.prop.2-3-2}\textbf{Proposition 2.3.2.} \uline{Let} $S$ \uline{be a compact space, equipped with a soft sheaf of rings} $A$\uline{. We write} $\Mod(S)$ \uline{for the category of left} $A$\uline{-Modules,} $\mathrm D(S)=\mathrm D(\Mod(S))$.
\begin{enumerate}[label=\alph*)]
\item \uline{The category} $\Lib(S)$ \uline{(resp.} $\underline{\operatorname{Loclib}}(S)$) (2.0) \uline{is liftable in} $\Mod(S)$ \uline{and quasi-stable} (resp. stable) \uline{under kernels of epimorphisms} (I 1.2)\SGAAuthorFootnote{1}{This assertion holds more generally for $S$ paracompact.}.
\item \uline{Let} $E\in\Ob\mathrm D(S)$ \uline{and} $n\in\mathbb Z$\uline{. For} $E$ \uline{to be} $n$\uline{-pseudo-coherent (resp.\ pseudo-coherent), it is necessary and sufficient that} $E$ \uline{be so relative to} $\Lib(S)$ \uline{or to} $\Loclib(S)$\uline{. The canonical functors}
\[
\mathrm K^{-}(\Lib(S))\longrightarrow \mathrm D(S)
\]
\uline{and}
\[
\mathrm K^{-}(\Loclib(S))\longrightarrow \mathrm D(S)
\]
\uline{are fully faithful and have as essential image} $\mathrm D^{-}(S)_{\coh}$ (I 2.15).
\item \uline{Let} $E\in\Ob\mathrm D(S)$\uline{, and} $[m,n]$ \uline{be an interval in} $\mathbb Z$\uline{. For} $E$ \uline{to have perfect amplitude contained in} $[m,n]$ \var{(4.8)}{(I 4.8)}{EN-VAR-123}, \uline{it is necessary and sufficient that there exist a quasi-isomorphism} $F\to E$\uline{, where} $F$ \uline{is a complex such that} $F^i=0$ \uline{for} $i\notin[m,n]$, $F^i\in\Ob\Lib(S)$ \uline{for} $i>m$ \uline{and} $F^m\in\Ob\Loclib(S)$\uline{. The canonical functor}
\[
\mathrm K^b(\Loclib(S))\longrightarrow\mathrm D(S)
\]
\uline{is fully faithful and has as essential image} $\mathrm D(S)_{\parf}$ \var{(4.9)}{(I 4.9)}{EN-VAR-124}.
\item \uline{Denote by} $\Mod(A(S))$ \uline{the category of left} $A(S)$\uline{-modules, and by} $\mathrm D(A(S))$ \uline{the corresponding derived category. The functor}
\[
H^0(S,.):\Mod(S)\longrightarrow\Mod(A(S))
\]
\uline{is exact and an equivalence of categories. Moreover, it
defines}
%%%%%%%%%%  scan idx 191  |  folio 178  %%%%%%%%%%
\SGAPrintPage{178}%
\uline{equivalences of categories:}
\[
\begin{gathered}
\mathrm D(S)\longrightarrow\mathrm D(A(S))\\
\mathrm D^{-}(S)_{\coh}\longrightarrow\mathrm D^{-}(A(S))_{\coh}\\
\mathrm D(S)_{\parf}\longrightarrow\mathrm D(A(S))_{\parf},
\end{gathered}
\]
\uline{the notions of pseudo-coherence (resp.\ perfection) in} $\mathrm D(A(S))$ \uline{being understood relative to the category of free (resp.\ projective) modules of finite type. It also defines an equivalence from the full subcategory of} $\mathrm D(S)$ \uline{formed by} $n$\uline{-pseudo-coherent objects (resp.\ objects of perfect amplitude contained in} $[m,n]$\uline{) onto the analogous full subcategory of} $\mathrm D(A(S))$.
\end{enumerate}

\uline{Proof.} The first assertion of a) follows from (2.1.2 a)) applied by taking for $\underline{S}$ the topos of sheaves of sets on $S$, and $\underline{C}'_S=\Mod(S)$ (in view of (2.3.1)). The second follows from (1.1 a)) (resp. (1.2 a))) with $\underline{C}=\Mod$ (the fibered category of $A$-Modules on variable open subsets), $\underline{C}_0=\Lib$ (resp. $\Loclib$), condition (ii) of (1.1) being verified by (2.3.1) and (2.1.1.b)). For assertion b), we apply (1.1 b),c)) (with $\underline{C}$ and $\underline{C}_0$ as before). \var{Assertion c) follows from (1.2.1 b),c),d))}{Assertion c) follows from (1.2 b),c),d))}{EN-VAR-126}, where $\underline{C}=\Mod$ and $\underline{C}_0=\Lib$. The first assertion of d) follows, by (2.3.1), from (2.1.2 b) (ii)), where we take $\underline{C}'_S=\underline{C}_S=\Mod(S)$. The other assertions of d) follow from the preceding one, in view of b) and c), and the fact that the functor $\Gamma_S$ sends \var{free Modules (resp.) direct factors of free Modules) of finite type}{free Modules (resp. direct factors of free Modules) of finite type}{EN-VAR-125} to free modules (resp. projective modules of finite type).

\sgaStableTarget{sga6.e02.sec.2.4}\subsection*{2.4. Stein spaces}

\sgaStableTarget{sga6.e02.def.2-4-1}\textbf{Definition 2.4.1.} \uline{Let} $\underline{S}$ \uline{be a locally ringed topos} (I 2.15.1 (ii))\uline{, with final object} $S$\uline{. We say that} $\underline{S}$ \uline{is} Stein \uline{when the structure sheaf} $\underline{O}_S$ \uline{is coherent} (I 3.1) \uline{and the following properties hold:}

\textquotedblleft\uline{theorem} A\textquotedblright{} \uline{: every coherent} $\underline{O}_S$\uline{-Module} (I 3.1) \uline{is generated by its sections,}

\textquotedblleft\uline{theorem} B\textquotedblright{} \uline{: for every coherent} $\underline{O}_S$\uline{-Module} $F$\uline{, we have} $H^i(S,F)=0$ \uline{for} $i>0$.

%%%%%%%%%%  scan idx 192  |  folio 179  %%%%%%%%%%
\SGAPrintPage{179}%
\sgaStableTarget{sga6.e02.rem.2-4-1-1}\textbf{Remark 2.4.1.1.} When $\underline{S}$ has \textquotedblleft enough closed points\textquotedblright\SGAAuthorFootnote{1}{i.e.\ enough fiber functors (SGA~4 IV 6.4.1) of the type $i_s:s\to\underline{S}$, where $i_s$ is the inclusion of a closed subtopos of $\underline{S}$.} $s$ such that the Ideal $\underline{m}_s$ defined by $s$\SGAAuthorFootnote{2}{i.e.\ formed by sections of $\underline{O}_S$ whose fiber at $s$ belongs to the maximal ideal $m_s$ of $\underline{O}_{S,s}$.} is coherent, the relation \textquotedblleft$H^1(S,F)=0$ for every coherent sheaf $F$\textquotedblright{} implies theorem A. Indeed, let $F$ be a coherent $\underline{O}_S$-Module and $s$ a closed point of $S$. We have the exact sequence
\[
0\longrightarrow\underline{m}_sF\longrightarrow F\longrightarrow F/\underline{m}_sF\longrightarrow0,
\]
where $F/\underline{m}_sF$ is a coherent sheaf concentrated at $s$. Applying the functor $H^0(S,.)$ and writing the cohomology exact sequence, we find that $H^0(S,F)\longrightarrow H^0(S,F/\underline{m}_sF)$ is an epimorphism. Since $H^0(S,F/\underline{m}_sF)\simeq F_s/m_sF_s$, where $m_s$ is the maximal ideal of $\underline{O}_{S,s}$, we conclude by Nakayama.

The editor does not know whether the preceding implication is true in general (when one assumes only that $\underline{O}_S$ is a coherent sheaf of local rings).

The editor also does not know whether the relation \textquotedblleft$H^1(S,F)=0$ for every coherent sheaf $F$\textquotedblright{} implies theorem B.

\sgaStableTarget{sga6.e02.ex.2-4-2}\textbf{Examples 2.4.2.} A complex analytic Stein space [3], a compact convex subset of $\mathbb C^n$ equipped with the sheaf induced by the sheaf of holomorphic functions [3], the spectrum of a noetherian ring, and an affine rigid analytic space [13] define Stein topoi.

\sgaStableTarget{sga6.e02.lem.2-4-3}\textbf{Lemma 2.4.3.} \uline{Let} $\underline{S}$ \uline{be a Stein topos} (2.4.1)\uline{. We assume that the final object} $S$ \uline{is quasi-compact} (2.1.1)\uline{, and we denote by} $\Coh(S)$ \uline{the abelian category of coherent} $\underline{O}_S$\uline{-Modules, i.e.\ of finite presentation. Then:}
\begin{enumerate}[label=\alph*)]
\item \uline{The category} $\Lib(S)$ \uline{(resp.} $\Loclib(S)$) (2.0) \uline{is liftable in} $\Coh(S)$ \uline{and quasi-stable (resp.\ stable) under kernels of epimorphisms} (I 1.2).
\item \uline{The canonical functors}
\[
\mathrm K^{-}(\Lib(S))\longrightarrow\mathrm K^{-}(\Loclib(S))\longrightarrow\mathrm D^{-}(\Coh(S))
\]
\uline{are equivalences of categories.}
%%%%%%%%%%  scan idx 193  |  folio 180  %%%%%%%%%%
\item \SGAPrintPage{180}%
\uline{Put} $A=\Gamma(S,\underline{O}_S)$\uline{, and denote by} $\Mod(A)$ \uline{the category of} $A$\uline{-modules, and by} $\mathrm D(A)$ \uline{the corresponding derived category. The functor}
\[
\Gamma_S:\Coh(S)\longrightarrow\Mod(A)
\]
\uline{is exact, fully faithful and has as essential image the strictly full additive subcategory of} $\Mod(A)$ \uline{formed by} $A$\uline{-modules of finite presentation. It induces a fully faithful functor}
\[
\mathrm D^{-}(\Coh(S))\longrightarrow\mathrm D(A),
\]
\uline{whose essential image is the full subcategory} $\mathrm D^{-}(A)_{\coh}$ \uline{formed by pseudo-coherent complexes (relative to the subcategory of free modules of finite type).}
\end{enumerate}

\uline{Proof.} For the first assertion of a), we apply (2.1.2 a)) with $\underline{C}'_S=\Coh(S)$, the hypothesis of (2.1.2) being satisfied by theorem B: with the notation of (2.1.2 a)), when $E$ is coherent and $G$ locally free of finite type, $\underline{\operatorname{Hom}}(G,E)$ is coherent, so the category $\Loclib(S)$ (and a fortiori $\Lib(S)$) is liftable in $\Coh(S)$. The second assertion of a) follows from (1.1 a)) (resp.(1.2 a)) ) applied with $\underline{C}_0=\Lib$ (resp. $\Loclib$)\SGAAuthorFootnote{1}{Lib and Loclib are locally stable under kernels of epimorphisms by (I 2.15.1).}, the resp.\ condition (ii) of (1.1) being satisfied by theorem A and (2.1.1 b)). Assertion b) follows from (1.1 c)) applied successively with $\underline{C}_0=\Lib$ and $\underline{C}_0=\Loclib$. The first assertion of c) follows from (2.1.2 b)(ii)) applied with $\underline{C}'_S=\Coh(S)$. To prove that the functor $\Gamma_S:\mathrm D^{-}(\Coh(S))\to\mathrm D(A)$ is fully faithful, it suffices, in view of b), to prove that the composite
\[
\mathrm K^{-}(\Lib(S))\longrightarrow\mathrm D^{-}(\Coh(S))\xrightarrow{\Gamma_S}\mathrm D(A)
\]
is so. But this factors as
\[
\mathrm K^{-}(\Lib(S))\xrightarrow{(*)}\mathrm K^{-}(\Lib(A))\xrightarrow{(**)}\mathrm D(A),
\]
%%%%%%%%%%  scan idx 194  |  folio 181  %%%%%%%%%%
\SGAPrintPage{181}%
where $\Lib(A)$
denotes the category of free $A$-modules of finite type and $(*)$ is induced by $\Gamma_S$. It is trivial that $(*)$ is an equivalence of categories. On the other hand, (I 2.7 b)) applied with $\underline{C}=\Mod(A)$, $\underline{C}_0=\Lib(A)$ (and taking account of the fact that $\mathrm K^{-,\varnothing}(\Lib(A))=0$) shows that $(**)$ is fully faithful and has as essential image $\mathrm D^{-}(A)_{\coh}$. This proves the last assertion of c) and completes the proof of (2.4.3).

\sgaStableTarget{sga6.e02.lem.2-4-4}\textbf{Lemma 2.4.4.} \uline{Under the hypotheses of} (2.4.3)\uline{, the canonical functor}
\[
\mathrm D^b(\Coh(S))\longrightarrow\mathrm D(S)
\]
\uline{is fully faithful and has as essential image} $\mathrm D^b(S)_{\coh}$ (I 2.15).

\uline{Proof} We apply (2.1.3) with $A=\Coh(S)$, $B=\Mod(S)$, $\phi:A\to B$ being the inclusion functor. Since this functor is fully faithful, it suffices to check, by (2.1.3 a)(ii)), that for fixed $M\in\Ob\Coh(S)$ and $n>0$ the functor $L\longmapsto \Ext^n_{\Mod(S)}(L,M)$ is coeffaceable in $L\in\Ob\Coh(S)$. Now, by theorem A and (2.1.1 b)), every coherent $\underline{O}_S$-Module is a quotient of a free Module of finite type, and, when $L$ is free of finite type, we have $\Ext^n_{\Mod(S)}(L,M)=0$ by theorem B, whence our assertion. It follows that the functor $\mathrm D^b(\Coh(S))\to\mathrm D(S)$ is fully faithful and has as essential image the full subcategory of $\mathrm D(S)$ formed by complexes with bounded and coherent cohomology, which is precisely $\mathrm D^b(S)_{\coh}$ since $\underline{O}_S$ is coherent (I 3.5. b)).

\sgaStableTarget{sga6.e02.prop.2-4-5}\textbf{Proposition 2.4.5.} \uline{With the hypotheses and notation of} (2.4.3)\uline{, we have:}
\begin{enumerate}[label=\alph*)]
\item \uline{The canonical functors}
\[
\mathrm K^{-,b}(\Lib(S))\longrightarrow\mathrm K^{-,b}(\Loclib(S))\longrightarrow\mathrm D^b(S)_{\coh}
\]
\uline{are equivalences of categories. Moreover, the functor} $\underline{\underline{R}}\Gamma_S:\mathrm D^+(S)\to\mathrm D^+(A)$ \uline{induces a fully faithful functor}
\[
\var{\underline{\underline{R}}\Gamma_S=\mathrm D^b(S)_{\coh}\longrightarrow\mathrm D^+(A)}{\underline{\underline{R}}\Gamma_S:\mathrm D^b(S)_{\coh}\longrightarrow\mathrm D^+(A)}{EN-VAR-127},
\]
%%%%%%%%%%  scan idx 195  |  folio 182  %%%%%%%%%%
\SGAPrintPage{182}%
\uline{whose essential image is} $\mathrm D^b(A)_{\coh}$ \uline{(the triangulated subcategory of} $\mathrm D(A)$ \uline{formed by complexes with bounded cohomology and pseudo-coherent relative to the category of free modules of finite type).}
\item \uline{Let} $E\in\Ob\mathrm D(S)$\uline{, and let} $[m,n]$ \uline{be an interval of} $\mathbb Z$\uline{. When} $E$ \uline{has perfect amplitude contained in} $[m,n]$ (I 4.8), $E$ \uline{is isomorphic to a complex} $F$ \uline{such that} $F^i=0$ \uline{for} $\var{i\in[m,n]}{i\notin[m,n]}{EN-VAR-128}$, $F^i\in\Ob\Lib(S)$ \uline{for} $i>m$ \uline{and} $F^m\in\Ob\Loclib(S)$.
\item \uline{The canonical functor}
\[
\mathrm K^b(\Loclib(S))\longrightarrow\mathrm D(S)_{\parf}\qquad(\text{I }4.9)
\]
\uline{is an equivalence of categories. Moreover, the functor} $\underline{\underline{R}}\Gamma_S:\mathrm D^+(S)\to\mathrm D^+(A)$ \uline{induces a fully faithful functor}
\[
\underline{\underline{R}}\Gamma_S:\mathrm D(S)_{\parf}\longrightarrow\mathrm D^+(A),
\]
\uline{whose essential image is} $\mathrm D(A)_{\parf}$ \uline{(the triangulated subcategory of} $\mathrm D(A)$ \uline{formed by complexes perfect relative to the category of projective modules of finite type).}
\end{enumerate}

\uline{Proof.} a) The functors
\[
\mathrm K^{-,b}(\Lib(S))\longrightarrow\mathrm K^{-,b}(\Loclib(S))\longrightarrow\mathrm D^b(\Coh(S))\xrightarrow{\Gamma_S}\mathrm D^b(A)_{\coh}
\]
are equivalences of categories by (2.4.3 b) and c)). Since the (fully faithful) functor $\Gamma_S:\mathrm D^b(\Coh(S))\to\mathrm D(A)$ factors as
\[
\mathrm D^b(\Coh(S))\longrightarrow\mathrm D^b(S)_{\coh}\xrightarrow{\underline{\underline{R}}\Gamma_S}\mathrm D(A),
\]
we conclude by (2.4.4).

b) Suppose that $\operatorname{parf.amp}(E)\subset[m,n]$. In particular, $E$ is pseudo-coherent; by a), after replacing $E$ by an isomorphic object, we may therefore assume that $E^i\in\Ob\Lib(S)$ for $i>m$ and $E^i=0$ for $i>n$.
%%%%%%%%%%  scan idx 196  |  folio 183  %%%%%%%%%%
\SGAPrintPage{183}%
After truncation, we may even assume that $E^i=0$ for $i<m$. By (I 4.16) applied with $\underline{C}=\Mod$, $\underline{C}_0=\Lib$, we then have $\var{E^m\in\Ob\Loclib(S))}{E^m\in\Ob\Loclib(S)}{EN-VAR-129}$, whence the assertion.

c) Hypothesis (1.1 (i))(resp. (1.1 (ii) resp.)) being satisfied for $\underline{C}_S=\Coh(S)$, $\underline{C}_0=\Lib$ by (2.4.3 a)) (resp. (2.1.1 b)) and theorem A), we know from (1.2 c) and d)) that the functor
\[
\mathrm K^b(\Loclib(S))\longrightarrow\mathrm D^b(\Coh(S))
\]
is fully faithful. Thus, by (2.4.4), the same holds for the functor
\[
\mathrm K^b(\Loclib(S))\longrightarrow\mathrm D(S),
\]
whose essential image is $\mathrm D(S)_{\parf}$ by b). By an argument analogous to that in the proof of (2.4.3 c)), we see that the functor $\var{\mathrm K^b(\Loclib(S))\to\mathrm D(S)}{\mathrm K^b(\Loclib(S))\to\mathrm D(A)}{EN-VAR-130}$ induced by $\Gamma_S$ is fully faithful and has as essential image $\var{\mathrm D(S)_{\parf}}{\mathrm D(A)_{\parf}}{EN-VAR-131}$, and we complete the proof as in a).

\sgaStableTarget{sga6.e02.rem.2-4-6}\textbf{Remark 2.4.6.} The editor does not know whether, in (2.4.4), one may replace $\mathrm D^b$ by $\mathrm D^{-}$, and consequently whether (2.4.5 a)) remains true when one replaces $\mathrm K^{-,b}$ by $\mathrm K^{-}$ and $\mathrm D^b$ by $\mathrm D^{-}$.

To finish, here is an interesting complement to (2.4.5). Let us first give a definition:

\sgaStableTarget{sga6.e02.def.2-4-7}\textbf{Definition 2.4.7.} \uline{A complex pseudo-analytic space is a ringed space locally isomorphic to a ringed space of the form} $(Y,\underline{O}_X|Y)$\uline{, where} $X$ \uline{is a complex analytic space and} $Y$ \uline{a closed subset of} $X$\uline{.

The structure sheaf of a complex pseudo-analytic space is therefore coherent.

A complex pseudo-analytic space is said to be} Stein \uline{when the corresponding topos is Stein} (2.4.1).

% Exposé II continuation macros
\providecommand{\Qcoh}{\operatorname{Qcoh}}
\providecommand{\Mod}{\operatorname{Mod}}
\providecommand{\Coh}{\operatorname{Coh}}
\renewcommand{\Ob}{\operatorname{ob}}
\providecommand{\Hom}{\operatorname{Hom}}
\providecommand{\Ext}{\operatorname{Ext}}
\providecommand{\Spec}{\operatorname{Spec}}
\providecommand{\qcoh}{\mathrm{qcoh}}
\providecommand{\parf}{\mathrm{parf}}
\providecommand{\coh}{\mathrm{coh}}
\providecommand{\Loclib}{\operatorname{Loclib}}
\providecommand{\Lib}{\operatorname{Lib}}
\providecommand{\im}{\operatorname{Im}}

%%%%%%%%%%  scan idx 197  |  folio 184  %%%%%%%%%%
\SGAPrintPage{184}%
\sgaStableTarget{sga6.e02.expose.occ2}%
\sgaStableTarget{sga6.e02.prop.2-4-8}\textbf{Proposition 2.4.8.}\SGAAuthorFootnote{1}{This result, together with its proof, was kindly communicated to the editor by A. Grothendieck; it was obtained independently by \var{J. Frish [5]}{J. Frisch [5]}{EN-VAR-132}.}. \uline{Let} $S$ \uline{be a compact complex pseudo-analytic Stein space} (2.4.7)\uline{. We assume that, for every open subset} $U$ \uline{of} $S$ \uline{and every reduced closed subspace} $X$ \uline{of} $U$\uline{, the normalization of} $X$ \uline{is locally connected}\SGAAuthorFootnote{2}{We shall admit that the notions of local Analytic Geometry in [10] extend to complex pseudo-analytic spaces; the reduction to the usual case is indeed immediate each time.} \uline{(a condition satisfied, for example, when} $S$ \uline{is locally defined by real analytic equations or inequalities in a space} $\underline{\underline{R}}^n$ [15]\uline{). Then the ring} $A=\Gamma(S,\underline{O}_S)$ \uline{is a noetherian ring.}

\uline{Proof.} We have to show that, for an $A$-module $M$ of finite presentation, the set of submodules (or, equivalently, the set of submodules of finite type) of $M$ satisfies the maximal condition. Now, when $N$ is a submodule of finite type of $M$, $M/N$ is of finite presentation by (I 2.6 and 2.9), so $N$ is of finite presentation by the equivalence (2.4.3 c)). Using again the equivalence (2.4.3 c)) between the category of coherent $\underline{O}_S$-Modules and the category of $A$-modules of finite presentation, and the compactness of $S$, we are reduced to proving the following lemma:

\sgaStableTarget{sga6.e02.lem.2-4-8-1}\textbf{Lemma 2.4.8.1.} \uline{Let} $S$ \uline{be a complex pseudo-analytic space satisfying the local connectedness condition of} (2.4.8)\uline{. Let} $F$ \uline{be a coherent} $\underline{O}_S$\uline{-Module,} $(E_n)$ \uline{an increasing sequence of coherent sub-}$\underline{O}_S$\uline{-Modules of} $F$, $s$ \uline{a point of} $S$\uline{. Then there exists an open neighbourhood} $U$ \uline{of} $s$ \uline{such that the sequence} $(E_n|U)$ \uline{is stationary.}

\uline{Proof.} Since the ring $\underline{O}_{S,s}$ is noetherian, there exists an integer $n_0$ such that $(E_n)_s=(E_{n_0})_s$ for $n\geq n_0$. Replacing the sequence $(E_n)$ by the sequence $(E_n/E_{n_0})$ for $n\geq n_0$, we may therefore assume that $(E_n)_s=0$ for every $n$. On the other hand, $F_s=M$ is an $\underline{O}_s$-module of finite type, so there exists (Bourbaki Alg. com. chap. IV \S1 th.1) a finite filtration
\[
 M=M_0\supset M_1\supset\ldots\supset M_q=0
\]
such that $M_i/M_{i+1}\simeq \underline{O}_s/p_i$, where $p_i$ is a prime ideal of $\underline{O}_s$. Since the category of $\underline{O}_s$-modules of finite type is equivalent to the category of germs at $s$ of coherent $\underline{O}_S$-Modules, we may suppose, after
%%%%%%%%%%  scan idx 198  |  folio 185  %%%%%%%%%%
\SGAPrintPage{185}%
localizing in a neighbourhood of $s$, that the $M_i$ (resp. $p_i$) are the fibers at $s$ of coherent sub-$\underline{O}_S$-Modules $F_i$ (resp. coherent Ideals $\underline{p}_i$) of $F$ (resp. $\underline{O}_S$) such that
\[
 F=F_0\supset F_1\supset\ldots\supset F_q=0
\]
and $F_i/F_{i+1}\simeq\underline{O}_S/\underline{p}_i$. Proving that the sequence $(E_n)$ is stationary in a neighbourhood of $s$ is evidently equivalent to proving that the sequences induced on each graded piece of the filtration $(F_i)$ of $F$ are stationary. Thus we may assume that $F=\underline{O}_S$, and that $\underline{O}_{S,s}$ is an integral domain, the $E_n$ forming an increasing sequence of Ideals of $\underline{O}_S$ such that $(E_n)_s=0$. Since $\underline{O}_s$ is an integral domain, hence in particular reduced, Oka's theorem ([10] IV \S3) further allows us to assume, after localizing at $s$, that $S$ is \uline{reduced}. Let us show that there exists an open neighbourhood $U$ of $s$ such that $E_n|U=0$ for sufficiently large $n$. For this, consider the normalization $\widetilde S$ of $S$ ([10] IV \S4). We know (loc. cit.) \var{that $S$ is the analytic spectrum}{that $\widetilde S$ is the analytic spectrum}{EN-VAR-133} ([10] II \S3) of a finite $\underline{O}_S$-Algebra $\widetilde{\underline{O}}_S$ whose fiber at each $t\in S$ is the integral closure of $\underline{O}_{S,t}$ (cf. note $(^2)$ to (2.4.8)). We have $E_n\subset E_n\cdot\widetilde{\underline{O}}_S$, so it suffices to prove that $E_n\cdot\widetilde{\underline{O}}_S=0$ for sufficiently large $n$. Now we have \var{(10 II \S3 Prop.8)}{([10] II \S3 Prop.8)}{EN-VAR-134}
\[
 E_n\cdot\widetilde{\underline{O}}_S = p_*(\widetilde{E_n\cdot\widetilde{\underline{O}}_S}),
\]
where $p:\widetilde S\longrightarrow S$ is the canonical projection, and $\widetilde{E_n\cdot\widetilde{\underline{O}}_S}$ the Ideal of $\underline{O}_{\widetilde S}$ canonically defined by $E_n\cdot\widetilde{\underline{O}}_S$. It therefore suffices to show that there exists an open neighbourhood $V$ of $s$ in $\widetilde S$ such that $\var{\widetilde{E_n\cdot\widetilde{\underline{O}}_S}|U=0}{\widetilde{E_n\cdot\widetilde{\underline{O}}_S}|V=0}{EN-VAR-135}$ for sufficiently large $n$, which reduces us to the case where $S$ is \uline{normal}, as we shall now assume. Since $(E_n)_s=0$, the closed immersion $i_n:S_n\longrightarrow S$ defined by $E_n$ is \uline{étale} at $s$. But, since $S$ is normal, the set of points where $i_n$ is étale is both open and closed (see the lemma below), so its image in $S$ is a subset of $S$ that is both open and closed. Thus the restriction of $E_n$ to the connected component of $s$ in $S$ is zero; this connected component is an open set, since $S$ is locally connected, and we are done. It therefore remains to establish the following
%%%%%%%%%%  scan idx 199  |  folio 186  %%%%%%%%%%
\SGAPrintPage{186}%
\sgaStableTarget{sga6.e02.lem.2-4-8-2}\textbf{Lemma 2.4.8.2.} (cf. EGA IV 18.10.3) \uline{Let} $Y$ \uline{be a normal pseudo-analytic space}\SGAAuthorFootnote{1}{It suffices that the local rings of $Y$ be integral domains.}\uline{, and} $f:X\longrightarrow Y$ \uline{an unramified morphism. The set of points of} $X$ \uline{where} $f$ \uline{is étale is both open and closed.}

\uline{Proof.} It is immediate that the set of points where $f$ is étale is open (since saying that $f$ is étale at a point means that $f$ is an isomorphism in a neighbourhood of that point ([8] VI)). It remains to show that the set of points where $f$ is not étale is open. Since an unramified morphism is locally an immersion ([8] VI), we may assume that $f$ is a closed immersion, defined by a coherent Ideal $I$. Let $x$ be a point of $X$ such that $I_x\ne0$. Let $\underline{p}$ be a prime ideal of $\underline{O}_{Y,x}$ containing $I_x$. We have $\dim(\underline{O}_{Y,x})_{\underline{p}}>0$, since otherwise we would have $\underline{p}=0$, as $\underline{O}_{Y,x}$ is an integral domain. \var{Thus $Y$ has codimension}{Thus $X$ has codimension}{EN-VAR-136} $>0$ at $x$, hence codimension $>0$ in a neighbourhood of $x$ ([10] III A Cor 1 to Lemma 3), and we are done.

\sgaStableTarget{sga6.e02.sec.3}\section*{3. Complements on quasi-coherent sheaves on schemes\protect\SGAAuthorFootnote{2}{This section is essentially due to J.L. Verdier.}}

\begin{samepage}
\sgaStableTarget{sga6.e02.para.3-0}3.0. In this section, we fix a \uline{quasi-compact and quasi-separated} scheme $S$. The canonical (fully faithful) inclusion functor
\[
 \phi:\Qcoh(S)\longrightarrow\Mod(S)\qquad\text{\SGAAuthorFootMark{3}}
\]\SGAAuthorFootText{3}{$\Mod(S)$ (resp. $\Qcoh(S)$) denotes the category of $\underline{O}_S$-Modules (resp. quasi-coherent $\underline{O}_S$-Modules).}
extends to a functor, still denoted by $\phi$, from $\mathrm D(\Qcoh(S))$ to $\mathrm D(S)=\mathrm D(\Mod(S))$, which we propose to study\SGAAuthorFootnote{4}{We shall omit $\phi$ in formulas where there is no risk of confusion.}.
\end{samepage}

\sgaStableTarget{sga6.e02.lem.3-1}\textbf{Lemma 3.1.} \uline{The category} $\Qcoh(S)$ \uline{is an abelian category satisfying} (AB 5) \uline{and for which the quasi-coherent sheaves of finite type form a \textquotedblleft small\textquotedblright{} family of generators. In particular} (Tohoku 1.10.1), $\Qcoh(S)$ \uline{has enough injectives.}
%%%%%%%%%%  scan idx 200  |  folio 187  %%%%%%%%%%
\SGAPrintPage{187}%
\uline{Proof.} The assertion concerning generators follows immediately from the fact that every quasi-coherent sheaf on $S$ is the filtered inductive limit of its quasi-coherent subsheaves of finite type (EGA I 9.4.9 and IV 1.7.7). The other assertions are evident.

\sgaStableTarget{sga6.e02.rem.3-1-1}\textbf{Remark 3.1.1.} The editor does not know whether $\Qcoh(S)$ still has enough injectives when one assumes only that $S$ is quasi-separated.

\sgaStableTarget{sga6.e02.lem.3-2}\textbf{Lemma 3.2.} \uline{The functor} $\phi:\Qcoh(S)\longrightarrow\Mod(S)$ \uline{has a right adjoint}
\[
 Q:\Mod(S)\longrightarrow\Qcoh(S),
\]
\uline{(called the} coherer).

\uline{Proof.} When $S$ is affine, the existence of $Q$ is immediate: take $Q=\widetilde{\Gamma}_S$ (EGA $\mathrm{Err}_{\mathrm I}$ 1.8.7). Now let $f:U\longrightarrow S$ be an affine open subset of $S$; it is clear that, for $E\in\Ob\Qcoh(S)$ and $F\in\Ob\Mod(U)$, we have a canonical isomorphism
\sgaStableTarget{sga6.e02.eq.1}\begin{equation}\tag{1}
 \Hom(E,f_*(F))\simeq \Hom(E,f_*Q(F)).
\end{equation}
Then let $(f_i:U_i\longrightarrow S)$ be a finite covering of $S$ by affine open subsets, and for each $(i,j)$ let $(U_{ijk}\longrightarrow U_{ij})$ be a finite covering of $U_{ij}=U_i\cap U_j$ by affine open subsets ($U_{ij}$ is quasi-compact since $S$ is quasi-separated). Let $F\in\Ob\Mod(S)$. We have an exact sequence
\sgaStableTarget{sga6.e02.eq.2}\begin{equation}\tag{2}
 0\longrightarrow F\longrightarrow \coprod_i f_{i*}(F|U_i)
 \longrightarrow \coprod_{ijk}f_{ijk*}(F|U_{ijk}),
\end{equation}
(where $f_{ijk}$ denotes the inclusion of $U_{ijk}$ into $S$). Note that, for an inclusion $g:V\longrightarrow U$ of affine open subsets of $S$, we have an evident natural arrow $Q(F|U)\longrightarrow g_*Q(F|V)$. Consequently, (2) gives an arrow
\sgaStableTarget{sga6.e02.eq.3}\begin{equation}\tag{3}
 \coprod_i f_{i*}Q(F|U_i)\longrightarrow
 \coprod_{ijk}f_{ijk*}Q(F|U_{ijk}),
\end{equation}
whose kernel we shall denote by $QF$. Let $E\in\Ob\Qcoh(S)$. Applying the functor $\Hom(E,.)$ to the sequence (2), we obtain, in view of (1), an
%%%%%%%%%%  scan idx 201  |  folio 188  %%%%%%%%%%
\SGAPrintPage{188}%
isomorphism functorial in $E$:
\[
 \Hom(E,F)\simeq\Hom(E,QF).
\]
This proves the existence of $Q$; uniqueness is evident, since it is an adjoint.

\sgaStableTarget{sga6.e02.lem.3-3}\textbf{Lemma 3.3.} \begin{enumerate}[label=\alph*)]
\item \uline{The coherer is an additive functor taking injectives to injectives.}
\item \uline{Write}
\[
 \underline{\underline{R}}Q:\mathrm D^+(S)\longrightarrow \mathrm D^+(\Qcoh(S))
\]
\uline{for the right derived functor of the coherer. For} $E\in\Ob\mathrm D(\Qcoh(S))$ and $F\in\Ob\mathrm D^+(S)$\uline{, there are canonical functorial isomorphisms}
\sgaStableTarget{sga6.e02.eq.i}\sgaStableTarget{sga6.e02.eq.ii}\begin{align}
 \underline{\underline{R}}\Hom_{\Mod}(E,F)&\simeq \underline{\underline{R}}\Hom_{\Qcoh}(E,\underline{\underline{R}}Q(F)) \qquad\text{\SGAAuthorFootMark{1}}\tag{i}\\
 \Hom_{\Mod}(E,F)&\simeq \Hom_{\Qcoh}(E,\underline{\underline{R}}Q(F)). \tag{ii}
\end{align}\SGAAuthorFootText{1}{The subscript $\Mod$ (resp. $\Qcoh$) is an abbreviation for $\mathrm D(S)$ (resp. $\mathrm D(\Qcoh(S))$).}
\end{enumerate}

\uline{Proof.} Assertion a) is an immediate consequence of the definition. For b), it suffices to establish the isomorphism (i), since (ii) will follow by passing to cohomology. For this, we may assume that $F$ has injective components; $QF$ then has injective components by a), and the isomorphism (i) comes simply from the definition of $Q$.

\sgaStableTarget{sga6.e02.lem.3-4}\textbf{Lemma 3.4.} \uline{Let} $f:X\longrightarrow Y$ \uline{be a morphism of quasi-compact and quasi-separated schemes. For} $E\in\Ob\Mod(X)$\uline{, there is a canonical functorial isomorphism}
\[
 Qf_*(E)\simeq f_*Q(E).
\]
%%%%%%%%%%  scan idx 202  |  folio 189  %%%%%%%%%%
\SGAPrintPage{189}%
\uline{Proof.} Since $f$ is quasi-compact and quasi-separated (EGA IV 1.2.3 and 1.2.4), $f_*$ takes quasi-coherent sheaves to quasi-coherent sheaves (III 1.4.10 and IV 1.7.21). The functors $Qf_*$ and $f_*Q$ from $\Mod(X)$ to $\Qcoh(Y)$ are both right adjoints to the functor $f^*:\Qcoh(Y)\longrightarrow\Mod(X)$, and hence canonically isomorphic.

\sgaStableTarget{sga6.e02.prop.3-5}\textbf{Proposition 3.5.} \uline{When} $S$ \uline{is noetherian, or quasi-compact and separated, the functor}
\begin{equation}\tag{3.0}
 \phi:\mathrm D^+(\Qcoh(S))\longrightarrow \mathrm D^+(S)
\end{equation}
\uline{is fully faithful and has as essential image the full subcategory} $\mathrm D^+(S)_{\qcoh}$ of $\mathrm D^+(S)$ \uline{formed by complexes with quasi-coherent cohomology.}

\uline{Proof.} We may restrict ourselves to proving the assertion in the case where $S$ is separated, the case where $S$ is noetherian being well known ([H] II 7.19). It is evidently equivalent to prove the following assertions:

\begin{enumerate}[label={(3.5.\arabic*)~:}]
\item \uline{for} $E\in\Ob\mathrm D^+(\Qcoh(S))$, \uline{the adjunction arrow}
\[
 E\longrightarrow \underline{\underline{R}}Q(E)
\]
\var{(\uline{deduced from} (3.3 b) (ii)) ) \uline{is an isomorphism}}{(\uline{deduced from} (3.3 b) (ii)) \uline{is an isomorphism}}{EN-VAR-137};
\item \uline{for} $E\in\Ob\mathrm D^+(S)_{\qcoh}$, \uline{the adjunction arrow}
\[
 \underline{\underline{R}}Q(E)\longrightarrow E
\]
\var{(\uline{deduced from} (3.3 b)(ii)) ) \uline{is an isomorphism}}{(\uline{deduced from} (3.3 b)(ii)) \uline{is an isomorphism}}{EN-VAR-138}.
\end{enumerate}

\uline{Proof of} (3.5.1). By an immediate dévissage\SGAAuthorFootnote{1}{Using, for example, the spectral sequence $R^pQ(H^q(E))\Longrightarrow R^{*}Q(E)$.}, we are reduced to proving that quasi-coherent sheaves are acyclic for the coherer, (i.e.\ that, for $E\in\Ob\Qcoh(S)$, we have $R^iQ(E)=0$ for $i>0$), and that, on the other hand, for $E\in\Ob\Qcoh(S)$, the adjunction arrow $E\longrightarrow QE$ is an isomorphism. When $S$ is affine, this follows from the formula $Q=\widetilde{\Gamma}_S$ (cf. proof of (3.2)). In the general case, cover $S$ by finitely many affine open subsets $U_i$. Since $S$ is separated, the inclusions $f_i:U_i\longrightarrow S$ are affine morphisms,
%%%%%%%%%%  scan idx 203  |  folio 190  %%%%%%%%%%
\SGAPrintPage{190}%
as are the inclusions $f_{i_0\ldots i_p}:U_{i_0\ldots i_p}=U_{i_0}\cap\ldots\cap U_{i_p}\longrightarrow S$. Let $E\in\Ob\Qcoh(S)$, and consider the resolution of $E$ defined by the covering $(U_i)$ (T.F. II 5.2):
\sgaStableTarget{sga6.e02.eq.unnamed.occ3}\begin{equation}\tag{*}
 0\longrightarrow E\longrightarrow \coprod f_{i*}f_i^*(E)\longrightarrow\ldots\longrightarrow
 \coprod f_{i_0\ldots i_p*}f_{i_0\ldots i_p}^*(E)\longrightarrow .
\end{equation}
It gives a spectral sequence
\sgaStableTarget{sga6.e02.eq.unnamed.occ4}\begin{equation}\tag{**}
 E_1^{pq}=\coprod_{i_0\ldots i_p}R^qQ\bigl(f_{i_0\ldots i_p*}f_{i_0\ldots i_p}^*(E)\bigr)
 \Longrightarrow R^{*}Q(E).
\end{equation}
Since $f_{i_0\ldots i_p}$ is affine, (3.4) gives:
\sgaStableTarget{sga6.e02.eq.unnamed.occ5}\begin{equation}\tag{***}
 R^qQ\bigl(f_{i_0\ldots i_p*}f_{i_0\ldots i_p}^*(E)\bigr)
 \simeq f_{i_0\ldots i_p*}R^qQ\bigl(f_{i_0\ldots i_p}^*(E)\bigr)
\end{equation}
for every $q$, whence in particular
\[
 E_1^{pq}=0\quad\text{for }q>0,
\]
by the affine case, already treated. From (**) we deduce isomorphisms
\[
 R^nQ(E)\simeq H^n(QE^\bullet)\quad\text{for }n\geq0,
\]
where $E^\bullet$ is the resolution of $E$ defined by the exact sequence (*). But (***) gives $QE^n\simeq E^n$ for $n\geq0$, whence finally
\[
 R^nQ(E)=0\quad\text{for }n>0,
\]
and $Q(E)\simeq E$, q.e.d.

\uline{Proof of} (3.5.2). The spectral sequence
\[
 R^pQ(H^q(E))\Longrightarrow R^{*}Q(E),
\]
%%%%%%%%%%  scan idx 204  |  folio 191  %%%%%%%%%%
\SGAPrintPage{191}%
immediately reduces us to the case where $E\in\Ob\Qcoh(S)$. Since the composite
\[
 E\longrightarrow \underline{\underline{R}}QE\longrightarrow E
\]
of the adjunction arrows is the identity, the assertion therefore follows from (3.5.1). This completes the proof of (3.5).

\sgaStableTarget{sga6.e02.rem.3-6}\textbf{Remark 3.6.} The conclusion of (3.5) would be false with the hypothesis \textquotedblleft separated or noetherian\textquotedblright{} omitted, as a neat counterexample of J.L. Verdier shows (see Appendix I to this exposé).

\sgaStableTarget{sga6.e02.prop.3-7}\textbf{Proposition 3.7.} \uline{Under the hypothesis of} (3.5)\uline{, suppose further that there exists an integer} $N$ \uline{such that, for every} $F\in\Ob\Mod(S)$ \uline{and every quasi-compact open subset} $U$ \uline{of} $S$,
\[
 H^i(U,F)=0\quad\text{for }i>N\qquad\text{\SGAAuthorFootMark{1}}\ .
\]
\SGAAuthorFootText{1}{By (Tohoku 3.6.5), this condition holds when the underlying space of $S$ is noetherian and has Krull dimension $\leq N$.}
\uline{Then:}
\begin{enumerate}[label=\alph*)]
\item \uline{The coherer has finite cohomological dimension.}
\item \uline{The functor}
\[
 \phi:\mathrm D(\Qcoh(S))\longrightarrow \mathrm D(S)
\]
\uline{is fully faithful and has as essential image the full subcategory} $\mathrm D(S)_{\qcoh}$ \uline{of} $\mathrm D(S)$ \uline{formed by complexes with quasi-coherent cohomology.}
\end{enumerate}

\uline{Proof of} a). When $S$ is affine, the formula $Q=\widetilde{\Gamma}_S$ (cf. proof of (3.2)) implies that $Q$ has cohomological dimension $\leq N$. We pass to the general case. Let us show that there exists an integer $a$ such that $R^qQ(E)=0$ for every $E\in\Ob\Mod(S)$ and every $q>a$. The hypothesis implies that every $E\in\Ob\Mod(S)$ admits a finite right resolution $E^\bullet$ of length $N$ such that the $E^i$ are acyclic, i.e.\ satisfy $H^q(U,E^i)=0$ for every open subset $U$ of $S$ and every $q>0$. It therefore suffices to prove that there exists an integer $b$ such that $R^qQ(E)=0$ for every acyclic $\underline{O}_S$-Module $E$ and every $q>b$. (We may then take $a=b+N$). Let $(U_i)$ be a finite covering of $S$ by $N'$ affine
%%%%%%%%%%  scan idx 205  |  folio 192  %%%%%%%%%%
\SGAPrintPage{192}%
open subsets. Let $E$ be an acyclic $\underline{O}_S$-Module, and consider the resolution $C^\bullet(E)$ of $E$ defined by the covering $(U_i)$ (T.F. II 5.2). There are two cases to examine:

1°) $S$ is \uline{separated}. The components of $C^\bullet(E)$ are then sums of terms of the form $f_*f^*(E)$, where $f:U\longrightarrow S$ is the inclusion of an affine open subset (hence an affine morphism, since $S$ is separated). We have
\[
 \underline{\underline{R}}(Qf_*)f^*E\simeq \underline{\underline{R}}Q\,\underline{\underline{R}}f_*f^*E\simeq \underline{\underline{R}}Q\,f_*f^*E,
\]
since $f^*E$ is acyclic. On the other hand, by (3.4), we have
\[
 \underline{\underline{R}}(Qf_*)f^*E\simeq \underline{\underline{R}}(f_*Q)f^*E\simeq f_*\underline{\underline{R}}Qf^*E,
\]
since $f$ is affine. We deduce that
\[
 R^nQ(f_*f^*E)\simeq f_*R^nQ(f^*E)
\quad\text{for every }n\geq0,
\]
whence, by the affine case,
\[
 R^nQ(f_*f^*E)=0\quad\text{for }n>N.
\]
Since $C^\bullet(E)$ has length $N'$, we finally have
\[
 R^nQ(E)=0\quad\text{for }n>N+N',
\]
and we are done.

2°) $S$ is \uline{noetherian}. The components of $C^\bullet(E)$ are sums of terms of the form $f_*f^*(E)$, where $f:U\longrightarrow S$ is the inclusion of a quasi-compact and \uline{separated} open subset. As in case 1°), we have
%%%%%%%%%%  scan idx 206  |  folio 193  %%%%%%%%%%
\SGAPrintPage{193}%
\begin{equation}\tag{*}
 \underline{\underline{R}}(Qf_*)f^*E\simeq \underline{\underline{R}}Q\,f_*f^*E\simeq \underline{\underline{R}}(f_*Q)f^*E.
\end{equation}
Since $U$ is noetherian, $\Qcoh(U)$ has enough objects that are injective in $\Mod(U)$ ([H] II 7.18), so the right derived functor of $f_*:\Qcoh(U)\longrightarrow\Qcoh(S)$ is \textquotedblleft induced\textquotedblright{} by $\underline{\underline{R}}f_*:\mathrm D^+(U)\longrightarrow\mathrm D^+(S)$ when we identify $\mathrm D^+(\Qcoh(U))$ (resp. $\mathrm D^+(\Qcoh(S))$) with a full subcategory of $\mathrm D^+(U)$ (resp. $\mathrm D^+(S)$) by (3.5). Consequently, we have
\[
 \underline{\underline{R}}(f_*Q)f^*E\simeq \underline{\underline{R}}f_*\underline{\underline{R}}Qf^*E.
\]
We deduce a spectral sequence
\[
 E_2^{pq}=R^pf_*R^qQf^*E\Longrightarrow R^{*}(f_*Q)f^*E.
\]
By case 1°), there exists an integer $n(U)$ depending only on $U$ such that $Q$ has cohomological dimension $\leq n(U)$ on $\Mod(U)$. Thus we have
\[
 E_2^{pq}=0\quad\text{for }q>n(U).
\]
But we have
\[
 E_2^{pq}=0\quad\text{for }p>N,
\]
because of the \textquotedblleft cohomological dimension\textquotedblright{} hypothesis on $S$. We therefore have
\[
 R^n(f_*Q)f^*E=0\quad\text{for }n>N+N'',
\]
where $N''$ is an integer greater than all the $n(U)$, with $U$ ranging over the set of intersections $U_{i_0}\cap\ldots\cap U_{i_p}$. Taking account of (*), we obtain
\[
 R^nQ(f_*f^*E)=0\quad\text{for }n>N+N'',
\]
whence finally
\[
 R^nQ(E)=0\quad\text{for }n>N+N'+N'',
\]
%%%%%%%%%%  scan idx 207  |  folio 194  %%%%%%%%%%
\SGAPrintPage{194}%
which completes the proof of a).

\uline{Proof of} b). The coherer, having finite cohomological dimension by a), admits ([H] I 5.3) a right derived functor
\[
 \underline{\underline{R}}Q:\mathrm D(S)\longrightarrow\mathrm D(\Qcoh(S))
\]
(whose restriction to $\mathrm D^+(S)$ coincides with the functor $\underline{\underline{R}}Q$ defined above). For $E\in\Ob\mathrm D(\Qcoh(S))$, there is a canonical functorial isomorphism
\sgaStableTarget{sga6.e02.eq.1.occ2}\begin{equation}\tag{1}
 E\xrightarrow{\sim}\underline{\underline{R}}Q(E),
\end{equation}
since, by (3.5.1), the $E^i$ are acyclic for $Q$ and satisfy $E^i\xrightarrow{\sim}Q(E^i)$. On the other hand, for $E\in\Ob\mathrm D(S)$, we have a canonical functorial homomorphism
\sgaStableTarget{sga6.e02.eq.2.occ2}\begin{equation}\tag{2}
 \underline{\underline{R}}Q(E)\longrightarrow E
\end{equation}
(coming, in the case where the $E^i$ are $Q$-acyclic, from the ordinary adjunction homomorphism between $\phi$ and $Q$). We check trivially that the homomorphisms (1) and (2) make $\underline{\underline{R}}Q:\mathrm D(S)\longrightarrow\mathrm D(\Qcoh(S))$ a right adjoint to $\phi:\mathrm D(\Qcoh(S))\longrightarrow\mathrm D(S)$. The fact that (1) is an isomorphism implies that $\phi$ is fully faithful. Finally, using the spectral sequence
\[
 R^pQ(H^q(E))\Longrightarrow R^{*}Q(E),
\]
(which converges since $Q$ has finite cohomological dimension), we see that (2) is an isomorphism for $E\in\Ob\mathrm D(S)_{\qcoh}$, so $\phi$ has as essential image $\mathrm D(S)_{\qcoh}$. This proves b) and completes the proof of (3.7).

%%%%%%%%%%  scan idx 208  |  folio 195  %%%%%%%%%%
\SGAPrintPage{195}%
\sgaStableTarget{sga6.e02.appendix-i-a-counterexample-of-verdier}\section*{APPENDIX I\\[3pt] A COUNTEREXAMPLE OF VERDIER\protect\SGAAuthorFootnote{1}{Apart from a few details of presentation, this appendix reproduces two letters from J.-L. Verdier (30/9/67 and 11/10/67) to the editor, in reply to a question from the latter.}}

\sgaStableTarget{sga6.e02.sec.0}\subsection*{0. Introduction}
For a scheme $S$, we denote by $\Mod(S)$ (resp. $\Qcoh(S)$) the category of $\underline{O}_S$-Modules (resp. quasi-coherent $\underline{O}_S$-Modules). The purpose of this appendix is to draw the reader's attention to certain pathological properties of the category $\mathrm D(\Qcoh(S))$, and in particular to prove the following facts:

\sgaStableTarget{sga6.e02.para.0-1}0.1. Let $X$ be an affine scheme with ring $A$, and let $I$ be an injective $A$-module. We know ([H] II 7.14) that, when $A$ is noetherian, the sheaf $\widetilde I$ is an injective $\underline{O}_X$-Module. When $A$ is not noetherian, $\widetilde I$ need not be injective, or even flasque, even when the underlying space of $X$ is noetherian and of finite Krull dimension. Moreover, the restriction of $\widetilde I$ to an open subset $U$ of $X$ is not in general an injective object of $\Qcoh(U)$.

\sgaStableTarget{sga6.e02.para.0-2}0.2. Let $f:X\longrightarrow Y$ be a morphism of quasi-compact and quasi-separated schemes. The functor $\underline{\underline{R}}f_*:\mathrm D^+(X)\longrightarrow\mathrm D^+(Y)$ does not in general induce on $\mathrm D^+(\Qcoh(X))$ (a full subcategory of $\mathrm D^+(X)$ by (II 3.5)) the right derived functor of $f_*:\Qcoh(X)\longrightarrow\Mod(Y)$. More precisely, an injective object of $\Qcoh(X)$ need not be acyclic for $f_*:\Mod(X)\longrightarrow\Mod(Y)$.

\sgaStableTarget{sga6.e02.para.0-3}0.3. Let $S$ be a quasi-compact and quasi-separated scheme. An object of $\Qcoh(S)$, even an injective one, need not be acyclic for the coherer (II 3.2). The functor $\phi:\mathrm D^b(\Qcoh(S))\longrightarrow\mathrm D^b(S)$ (II 3.0) is not in general fully faithful.

\sgaStableTarget{sga6.e02.sec.1.occ2}\subsection*{1. The counterexamples}
We first recall the following lemma (SGA 2 II 9):

\sgaStableTarget{sga6.e02.lem.1-1}\textbf{Lemma 1.1.} \uline{Let} $A$ \uline{be a ring,} $X=\Spec(A)$, $\underline{f}=(f_\alpha)$ \uline{a finite system of elements of} $A$, $Y$ \uline{the closed subset defined by} $\underline{f}(x)=0$\uline{. For} $i>0$\uline{, the following conditions are equivalent:}
\begin{enumerate}[label=(\roman*)]
\item \uline{for every injective} $A$\uline{-module} $I$\uline{, we have} $H^i_Y(X,\widetilde I)=0$;
\item \uline{the pro-object} $n\mapsto H_i(\underline{f}^n,A)$ \uline{is zero, i.e. for every} $n$\uline{, there exists} $n'\geq n$
%%%%%%%%%%  scan idx 209  |  folio 196  %%%%%%%%%%
\SGAPrintPage{196}%
\uline{such that}
\[
 H_i(\underline{f}^{n'},A)\longrightarrow H_i(\underline{f}^n,A)
\]
\uline{is zero}.
\end{enumerate}

\sgaStableTarget{sga6.e02.para.1-2}1.2 We shall use (1.1) to prove (0.1). Let $k$ be a field, $B=k[[x,y]]$, $\underline{f}$ the system $(x,y)$, $\var{M=\coprod_{n\ 0}B/(\underline{f}^n)}{M=\coprod_{n>0}B/(\underline{f}^n)}{EN-VAR-139}$, $A$ the $B$-algebra $D_B(M)$ (EGA $0_{\mathrm{IV}}$ 18.2.3). We have trivially
\[
 H_2^A(\underline{f}^n,A)\simeq H_2^B(\underline{f}^n,A).\qquad\text{\SGAAuthorFootMark{1}}
\]
\SGAAuthorFootText{1}{The superscript A or B indicates the ring over which the Koszul complex is computed.}
On the other hand, we have
\[
 H_2^B(\underline{f}^n,A)\simeq H_2^B(\underline{f}^n,B)\oplus H_2^B(\underline{f}^n,M)
\]
but $H_2^B(\underline{f}^n,B)=0$ since $\underline{f}$ is a regular sequence (EGA III 1.1.4 and $0_{\mathrm{IV}}$ 15.1.20), so we have
\[
 \var{H_2^B(\underline{f}^n,A)\simeq H_2^B(\underline{f}^n,M),}{H_2^B(\underline{f}^n,A)\simeq H_2^B(\underline{f}^n,M).}{EN-VAR-140}
\]
\var{Finally, we have (EGA II 1.1.3.3)}{Finally, we have (EGA III 1.1.3.3)}{EN-VAR-141}
\[
 H_2^B(\underline{f}^n,M)\simeq H_B^0(\underline{f}^n,M),
\]
where $H_B^0(\underline{f}^n,M)$ is canonically identified with the submodule ${}_{(\underline{f}^n)}M$ of $M$ annihilated by $(\underline{f}^n)$. Hence finally:
\sgaStableTarget{sga6.e02.eq.unnamed.occ6}\begin{equation}\tag{*}
 \var{H_2^A(\underline{f}^n,A\simeq {}_{(\underline{f}^n)}M.}{H_2^A(\underline{f}^n,A)\simeq {}_{(\underline{f}^n)}M.}{EN-VAR-142}
\end{equation}
Since the transition morphism $H_2^A(\underline{f}^{n'},A)\longrightarrow H_2^A(\underline{f}^n,A)$ is induced by multiplication by $\underline{f}^{\,n'-n}$ (EGA III 1.1.6), we check immediately, using (*), that \uline{the pro-object} $n\mapsto H_2^A(\underline{f}^n,A)$ \uline{is not zero}. From (1.1) we then deduce that there exists an injective $A$-module $I$ such that $H^2_Y(X,\widetilde I)\ne0$, and thus such that $\widetilde I$ is not flasque (SGA 2 I 2.12). We note in passing that the projection $\Spec(A)\longrightarrow\Spec(B)$ is a homeomorphism,
%%%%%%%%%%  scan idx 210  |  folio 197  %%%%%%%%%%
\SGAPrintPage{197}%
so the underlying space of $X$ is noetherian of dimension $2$. Let $I$ be an injective $A$-module such that $H^2_Y(X,\widetilde I)\ne0$, and denote by $U$ the open complement of $Y$ in $X$. We have (SGA 2 II 4)
\[
 H^2_Y(X,I)\simeq H^1(U,I)\ne0.
\]
Now, since $U$ is separated, the functor $\mathrm D^+(\Qcoh(U))\to\mathrm D^+(U)$ is fully faithful (II 3.5), so $\widetilde I|U$ is not injective in $\Qcoh(U)$. We have therefore proved assertion (0.1).

Denote by $Z$ the quasi-compact and quasi-separated scheme obtained by gluing two copies of $X$ along $U$, and by $j_1,j_2$ the two canonical open immersions of $X$ into $Z$, so that we have a cartesian square
\[
\begin{tikzcd}
 U \arrow[r,"i"] \arrow[d,"i"'] & X \arrow[d,"j_1"]\\
 X \arrow[r,"j_2"] & Z .
\end{tikzcd}
\]
We have
\[
 j_2^*R^1j_{1*}(\widetilde I)\simeq R^1i_*(i^*\widetilde I),
\]
whence
\[
 H^0(X,R^1j_{1*}(\widetilde I))\simeq H^1(U,\widetilde I)\ne0,
\]
so $R^1j_{1*}(\widetilde I)\ne0$, which proves (0.2).

On the other hand, by (II 3.4 and 3.5), we have
\[
 \underline{\underline{R}}Q\,\underline{\underline{R}}j_{1*}(\widetilde I)\simeq j_{1*}(\widetilde I).
\]
We deduce a spectral sequence whose exact sequence of low-degree terms gives the isomorphism
\[
 R^1j_{1*}(\widetilde I)\simeq R^2Q(j_{1*}\widetilde I).
\]
Since the sheaf $j_{1*}\widetilde I$ is injective in $\Qcoh(Z)$, we have proved the first
%%%%%%%%%%  scan idx 211  |  folio 198  %%%%%%%%%%
\SGAPrintPage{198}%
part of (0.3).

Finally, let $E\in\Ob\Qcoh(Z)$. We have the duality isomorphism
\[
 \underline{\underline{R}}\Hom_{\Mod}(E,\underline{\underline{R}}j_{1*}(\widetilde I))
 \simeq \underline{\underline{R}}\Hom_{\Mod}(j_1^*E,\widetilde I),
\]
but, since $X$ is affine, we have (II 3.5)
\[
 \underline{\underline{R}}\Hom_{\Mod}(j_1^*E,\widetilde I)
 \simeq \underline{\underline{R}}\Hom_{\Qcoh}(j_1^*E,I)
 \simeq \Hom(j_1^*E,\widetilde I).
\]
We deduce
\[
 \underline{\underline{R}}\Hom_{\Mod}(E,\underline{\underline{R}}j_{1*}(\widetilde I))
 \simeq \Hom(j_1^*E,\widetilde I),
\]
whence a spectral sequence whose exact sequence of low-degree terms gives the isomorphism
\[
 \Ext^2_{\Mod}(E,j_{1*}\widetilde I)
 \simeq \Hom(E,R^1j_{1*}\widetilde I).
\]
Consequently, the identity morphism of $E=R^1j_{1*}(\widetilde I)$ gives a nonzero element of $\Ext^2_{\Mod}(E,j_{1*}\widetilde I)$, so the functor $\mathrm D^b(\Qcoh(Z))\longrightarrow \mathrm D^b(Z)$ is not fully faithful, which completes the proof of (0.3).

% Original English continuation definitions, retained in their original order.
\providecommand{\Diff}{\operatorname{Diff}}
\providecommand{\Hom}{\operatorname{Hom}}
\providecommand{\im}{\operatorname{Im}}
\providecommand{\ind}{\operatorname{ind}}
\providecommand{\Cone}{\operatorname{Cone}}
\providecommand{\Loclib}{\operatorname{Loclib}}
\providecommand{\ob}{\operatorname{ob}}
\providecommand{\Parf}{\operatorname{Parf}}
\providecommand{\Mod}{\operatorname{Mod}}
\providecommand{\Ob}{\operatorname{Ob}}
\providecommand{\Coker}{\operatorname{Coker}}
\providecommand{\Ker}{\operatorname{Ker}}
\providecommand{\cH}{\mathcal H}
\providecommand{\cO}{\mathcal O}
\providecommand{\cP}{\mathcal P}
\providecommand{\cE}{\mathcal E}
\providecommand{\cF}{\mathcal F}
\providecommand{\cG}{\mathcal G}
\providecommand{\cD}{\mathcal D}
\providecommand{\cT}{\mathcal T}
\providecommand{\cL}{\mathcal L}
\providecommand{\RHom}{\mathrm R\!\operatorname{Hom}}
\providecommand{\RGamma}{\mathrm R\Gamma}
\providecommand{\dlim}{\operatorname*{lim}}
\providecommand{\plim}{\operatorname*{lim}}

% Reviewed French continuation definitions, retained verbatim.
% Additions for Expose II, Appendix II
\providecommand{\cO}{\mathcal O}
\providecommand{\cE}{\mathscr E}
\providecommand{\cF}{\mathscr F}
\providecommand{\cG}{\mathscr G}
\providecommand{\cH}{\mathcal H}
\providecommand{\cL}{\mathscr L}
\providecommand{\cP}{\mathscr P}
\providecommand{\cM}{\mathscr M}
\providecommand{\cD}{\mathcal D}
\providecommand{\R}{\mathbb R}
\providecommand{\C}{\mathbb C}
\providecommand{\N}{\mathbb N}
\providecommand{\Zell}{\mathbb Z_\ell}
\providecommand{\Qell}{\mathbb Q_\ell}
\providecommand{\C}{\mathbb C}
\providecommand{\Ch}{\operatorname{Ch}}
\providecommand{\Sym}{\operatorname{Sym}}
\providecommand{\Td}{\operatorname{Td}}
\providecommand{\Diff}{\operatorname{Diff}}
\providecommand{\Homop}{\operatorname{Hom}}
\providecommand{\Spec}{\operatorname{Spec}}
\providecommand{\pr}{\operatorname{pr}}
\providecommand{\ind}{\operatorname{ind}}
\providecommand{\soft}{\operatorname{soft}}
\providecommand{\Ban}{\operatorname{Ban}}
\providecommand{\Lop}{\operatorname{L}}
\providecommand{\im}{\operatorname{Im}}
\providecommand{\Parf}{\operatorname{Parf}}
\renewcommand{\id}{\operatorname{id}}

\newpage

%%%%%%%%%%  scan idx 212  |  folio 199  %%%%%%%%%%
\SGAPrintPage{199}%
\sgaStableTarget{sga6.e02.appendix-ii-definition-of-the-analytic-index-of}\section*{Appendix II\\ Definition of the analytic index of a relative elliptic complex}

In this appendix, we propose to show that the notion of a perfect complex leads to a simple and natural definition of the \textquotedblleft analytic index\textquotedblright{} of a continuous family of elliptic differential operators, or more generally of a \textquotedblleft relative elliptic complex\textquotedblright. For the initiated reader, let us say that the definition given here has the advantage over the classical definitions (see for example [2][17]) of requiring no deformation of the elliptic complex under consideration.

Sections 1 and 2 consist of reminders about the classical notion of a \textquotedblleft mixed structure\textquotedblright\SGAAuthorFootnote{1}{The informed reader is advised to skip them.}. We give, of course, only the minimum of definitions and results required to state and prove, in \S3, the central theorem we have in view, which says, a little more precisely, that the direct image, under the projection (assumed proper) onto the base (assumed locally compact), of a relative elliptic complex is a perfect complex. When the base $S$ is compact, the definition of the analytic index as an element of $K(S)$ (\textquotedblleft Grothendieck group\textquotedblright{} of complex vector bundles of finite rank on $S$) then follows from the globalization sorite of (II 2.3.2).

Throughout what follows, for a topological space $S$, we shall denote by $\underline{O}_S$ the sheaf of continuous real-valued functions on $S$.

\sgaStableTarget{sga6.e02.sec.1.occ3}\subsection*{1. The sorite of mixed structures}

In this section, we fix a topological space $S$.

\sgaStableTarget{sga6.e02.sec.1.1}\subsubsection*{1.1. Definition of mixed varieties}
Let $V$ be an open subset of $\underline{\underline{R}}^n$. For $r\in\underline{\underline{N}}$, we denote by $C^r(V)$ the Fréchet space of real-valued functions of class $C^r$ on $V$. We put
\[
  C^\infty(V)=\varprojlim_r C^r(V).
\]
Let $U$ be an open subset of $S$. For $r\in\underline{\underline{N}}\cup\{\infty\}$, we denote by ${}^r\underline{O}_{U\times V}$ the sheaf of rings on $U\times V$ defined by
%%%%%%%%%%  scan idx 213  |  folio 200  %%%%%%%%%%
\SGAPrintPage{200}%
\[
  \Gamma(U'\times V',{}^r\underline{O}_{U'\times V'})=C(U',C^r(V')),
\]
where $U'$ (resp. $V'$) is an open subset of $U$ (resp. $V$), and $C(U',C^r(V'))$ denotes the space of continuous maps from $U'$ to $C^r(V')$. There is a canonical injection
\[
  \pr_1^{-1}(\underline{O}_U)\longrightarrow{}^r\underline{O}_{U\times V},
\]
which makes $(U\times V,{}^r\underline{O}_{U\times V})$ a ringed space over $U$.

A ringed space $X$ over $S$ is called a \uline{mixed variety of class} $C^r$ ($r\in\underline{\underline{N}}\cup\{\infty\}$) when $X$ is locally isomorphic (as a ringed space over $S$) to a space of the form $(U\times V,{}^r\underline{O}_{U\times V})$. (When $S$ consists of one point, a mixed variety of class $C^r$ over $S$ is a variety of class $C^r$ in the usual sense).

Let $f:X\to S$ be a mixed variety of class $C^r$. For $s\in S$, the fiber $X_s=f^{-1}(s)$, equipped with the sheaf $i_s^{-1}(\underline{O}_X)/I_s$, where $i_s:X_s\to X$ is the injection and $I_s$ is the Ideal of sections that \textquotedblleft vanish\textquotedblright{} on $X_s$, is a variety of class $C^r$.

\sgaStableTarget{sga6.e02.sec.1.2}\subsubsection*{1.2. \textquotedblleft Relative\textquotedblright{} differential operators}
It is easy to make an infinitesimal study of mixed varieties on the model of Algebraic Geometry (EGA IV 16, SGA 3 VII). We shall limit ourselves to a brief review of the notions we shall need.

\sgaStableTarget{sga6.e02.para.1-2-1}1.2.1. Let $f:X\to S$ be a mixed variety of class $C^\infty$ (1.1). For every $n\in\underline{\underline{N}}$, we define the sheaf $P^n_{X/S}$ of \uline{principal parts of order} $n$ on $X$. In two different ways, corresponding to the two projections $\pr_1$ and $\pr_2$ from $X\times_S X$ onto $X$, it is an $\underline{O}_X$-Algebra augmented to $\underline{O}_X$, locally free of finite type as an $\underline{O}_X$-Module. When we do not specify the $\underline{O}_X$-Module structure considered on $P^n_{X/S}$, it will be understood to be the one coming from $\pr_1$. There are canonical isomorphisms
%%%%%%%%%%  scan idx 214  |  folio 201  %%%%%%%%%%
\SGAPrintPage{201}%
\[
  \begin{aligned}
  P^0_{X/S}&\simeq\underline{O}_X,\\
  P^1_{X/S}&\simeq\underline{O}_X\oplus\Omega^1_{X/S},
  \end{aligned}
\]
where $\Omega^1_{X/S}$ is the sheaf of relative differentials of $X$ over $S$, or \uline{relative cotangent sheaf}. We denote by $T_{X/S}$ the dual of $\Omega^1_{X/S}$, called the \uline{relative tangent sheaf}. The sheaf $T_{X/S}$ (resp. $\Omega^1_{X/S}$) is the sheaf of sections of a vector bundle of finite rank on $X$, called the \uline{relative} \uline{tangent bundle} (resp. \uline{cotangent bundle}), which is a mixed $S$-variety over $X$.

The $f^{-1}(\underline{O}_S)$-linear homomorphism (corresponding to the $\underline{O}_X$-Algebra structure defined by $\pr_2$)
\[
  d^n_{X/S}:\underline{O}_X\longrightarrow P^n_{X/S},
\]
which associates to a section of $\underline{O}_X$ its principal part of order $n$, is called the \uline{universal differential operator of order} $n$ on $X$ relative to $S$.

\sgaStableTarget{sga6.e02.para.1-2-2}1.2.2. More generally, for a locally free $\underline{O}_X$-Module $E$ of finite type, we can consider the sheaf $\var{P^n_{X/S}(E)=P^n_{X/S}\oplus_{\underline{O}_X}E}{P^n_{X/S}(E)=P^n_{X/S}\otimes_{\underline{O}_X}E}{EN-VAR-143}$ of principal parts of order $n$ of $E$ (the tensor product being defined by means of the second $\underline{O}_X$-Module structure on $P^n_{X/S}$), and we have a universal differential operator
\[
  d^n_{X/S}(E):E\longrightarrow P^n_{X/S}(E).
\]

Let $E,F$ be locally free $\underline{O}_X$-Modules of finite type. By definition, a differential operator from $E$ to $F$ of order $\le n$ \var{is an $f^{-1}(\underline{O}_X)$-linear homomorphism}{is an $f^{-1}(\underline{O}_S)$-linear homomorphism}{EN-VAR-144} from $E$ to $F$ that factors through $d^n_{X/S}(E)$ as an $\underline{O}_X$-linear homomorphism from $P^n_{X/S}(E)$ to $F$ (the factorization is then unique). In other words, the operator $d^n_{X/S}(E)$ gives a bijection
\sgaStableTarget{sga6.e02.eq.1-2-2-1}\begin{equation}\tag{1.2.2.1}
  \Homop_{\underline{O}_X}\bigl(P^n_{X/S}(E),F\bigr)
  \xrightarrow{\sim}
  \Diff^n(E,F),
\end{equation}
where $\Diff^n(E,F)$ denotes the set of differential operators from $E$
%%%%%%%%%%  scan idx 215  |  folio 202  %%%%%%%%%%
\SGAPrintPage{202}%
to $F$ of order $\le n$.

\sgaStableTarget{sga6.e02.para.1-2-3}1.2.3. We shall also have to consider differential operators with complex coefficients. Denote by $\underline{O}_S\otimes\underline{\underline{C}}$ (resp. $\underline{O}_X\otimes\underline{\underline{C}}$) the complexification of $\underline{O}_S$ (resp. $\underline{O}_X$). Let $E,F$ be locally free $(\underline{O}_X\otimes\underline{\underline{C}})$-Modules of finite type. By a differential operator from $E$ to $F$ of order $\le n$, we shall mean an $f^{-1}(\underline{O}_S\otimes\underline{\underline{C}})$-linear homomorphism from $E$ to $F$ that factors through $d^n_{X/S}(E)$ as an $(\underline{O}_X\otimes\underline{\underline{C}})$-linear homomorphism from $P^n_{X/S}(E)$ to $F$. Denoting the set of these operators by $\Diff^n_{\underline{\underline{C}}}(E,F)$, the operator $d^n_{X/S}(E)$ therefore defines a bijection
\sgaStableTarget{sga6.e02.eq.1-2-3-1}\begin{equation}\tag{1.2.3.1}
  \Homop_{\underline{O}_X\otimes\underline{\underline{C}}}\bigl(P^n_{X/S}(E),F\bigr)
  \xrightarrow{\sim}
  \Diff^n_{\underline{\underline{C}}}(E,F).
\end{equation}

\begin{samepage}
\sgaStableTarget{sga6.e02.para.1-2-4}1.2.4. For $n\ge1$, there is a canonical exact sequence
\[
  0\longrightarrow S^n(\Omega^1_{X/S})
  \longrightarrow P^n_{X/S}
  \longrightarrow P^{n-1}_{X/S}
  \longrightarrow 0.\text{\SGAAuthorFootMark{1}}
\]
\SGAAuthorFootText{1}{$S^n(\ )$ denotes the $n$th symmetric power.}%
\end{samepage}
Hence, tensoring over $\underline{O}_X$\SGAAuthorFootnote{2}{For the second $\underline{O}_X$-Module structure of $P^n_{X/S}$ and $P^{n-1}_{X/S}$.} with a locally free $\underline{O}_X$-Module (resp. $(\underline{O}_X\otimes\underline{\underline{C}})$-Module) $E$ of finite type, we obtain an exact sequence of $\underline{O}_X$ (resp. $(\underline{O}_X\otimes\underline{\underline{C}})$)-Modules:
\sgaStableTarget{sga6.e02.eq.1-2-4-1}\begin{equation}\tag{1.2.4.1}
  0\longrightarrow S^n(\Omega^1_{X/S})\otimes_{\underline{O}_X}E
  \longrightarrow P^n_{X/S}(E)
  \longrightarrow P^{n-1}_{X/S}(E)
  \longrightarrow 0.
\end{equation}
Let $F$ be a locally free $\underline{O}_X$ (resp. $(\underline{O}_X\otimes\underline{\underline{C}})$)-Module of finite type. Applying the functor $\Homop_{\underline{O}_X}(.,F)$ (resp. $\Homop_{\underline{O}_X\otimes\underline{\underline{C}}}(.,F)$) to (1.2.4.1), we deduce, taking account of \var{the identifications (1.2.2.1) and (1.2.4.1)}{the identifications (1.2.2.1) and (1.2.3.1)}{EN-VAR-145}, an exact sequence
\sgaStableTarget{sga6.e02.eq.1-2-4-2}\begin{equation}\tag{1.2.4.2}
\begin{gathered}
  0\longrightarrow \Diff^{n-1}(E,F)
  \longrightarrow \Diff^n(E,F)
  \xrightarrow{\sigma_n}
  \Homop_{\underline{O}_X}\bigl(S^n(\Omega^1_{X/S})\otimes_{\underline{O}_X}E,F\bigr)\\
  (\text{resp. } 0\longrightarrow \Diff^{n-1}_{\underline{\underline{C}}}(E,F)
  \longrightarrow \Diff^n_{\underline{\underline{C}}}(E,F)
  \longrightarrow
  \Homop_{\underline{O}_X\otimes\underline{\underline{C}}}\bigl(S^n(\Omega^1_{X/S})\otimes_{\underline{O}_X}E,F\bigr)).
\end{gathered}
\end{equation}
%%%%%%%%%%  scan idx 216  |  folio 203  %%%%%%%%%%
\SGAPrintPage{203}%
(When $X$ is paracompact and $\underline{O}_X$ is soft (cf. II 2.3.1), we can complete the exact sequence (1.2.4.2) by adding a zero on the right). The homomorphism $\sigma_n$ is called the \uline{symbol}; it associates to a differential operator from $E$ to $F$ of order $\le n$ a \textquotedblleft homogeneous polynomial map of degree $n$ from $\Omega^1_{X/S}$ to $\underline{\Homop}_{\underline{O}_X}(E,F)$ (resp. $\underline{\Homop}_{\underline{O}_X\otimes\underline{\underline{C}}}(E,F)$)\textquotedblright, which we may also regard, if we wish, as a homomorphism $p^*(E)\longrightarrow p^*(F)$ (where $p$ is the projection onto $X$ from the relative cotangent bundle (1.2.1)), satisfying a certain homogeneity condition.

Let $E,F,G$ be locally free $\underline{O}_X$-Modules of finite type, and $u\in\Diff^m(E,F)$, $v\in\Diff^n(F,G)$, whence $vu\in\Diff^{m+n}(E,G)$. We easily check the relation
\sgaStableTarget{sga6.e02.eq.1-2-4-3}\begin{equation}\tag{1.2.4.3}
  \sigma_{m+n}(vu)=\sigma_n(v)\sigma_m(u).
\end{equation}
(The same relation holds for differential operators with complex coefficients). In particular, when the composite of two differential operators is zero, the composite of their symbols is also zero.

\sgaStableTarget{sga6.e02.para.1-2-5}1.2.5. Let $E$ be a complex, with bounded degrees, of differential operators with complex coefficients on $X$ relative to $S$ (1.2.3), with $d^i:E^i\longrightarrow E^{i+1}$ of order $\le n_i$. By (1.2.4.3), we have
\[
  \sigma_{n_{i+1}}(d^{i+1})\,\sigma_{n_i}(d^i)=0.
\]
We call $E$ an \uline{elliptic complex} of order $\underline{n}=(n_i)$ when the complex $(p^*(E^i),\sigma_{n_i}(d^i))$ (where $p$ is the projection onto $X$ from the cotangent bundle) is acyclic outside the zero section.

When $S$ consists of one point, the preceding notion of an elliptic complex on $X$ coincides with the usual notion, studied in detail in [4] or [16].

\sgaStableTarget{sga6.e02.para.1-2-6}1.2.6. Let $s\in S$, and let $i_s$ be the inclusion of the fiber $X_s$ into $X$ \var{(considered as a monomorphism of ringed spaces) cf.~(1.1)).}{(considered as a monomorphism of ringed spaces) cf.~(1.1).}{EN-VAR-146} For $n\in\underline{\underline{N}}$, there is a canonical isomorphism

%%%%%%%%%%  scan idx 217  |  folio 204  %%%%%%%%%%
\SGAPrintPage{204}%
\[
  i_s^*(P^n_{X/S})\simeq P^n_{X_s},
\]
whence a restriction homomorphism
\[
  i_s^*:\Diff^n(E,F)
  \longrightarrow
  \Diff^n\bigl(i_s^*(E),i_s^*(F)\bigr),
\]
for locally free $\underline{O}_X$-Modules $E,F$ of finite type (and an analogous homomorphism for operators with complex coefficients). Let $E$ be a complex, with bounded degrees, of differential operators with complex coefficients on $X$; for $E$ to be elliptic (1.2.5), it is necessary and sufficient that, for every $s\in S$, $i_s^*(E)$ be an elliptic complex (verification left to the reader).

\sgaStableTarget{sga6.e02.sec.2.occ2}\subsection*{2. Rigidity lemmas}

\sgaStableTarget{sga6.e02.para.2-0}2.0. It is \textquotedblleft well known\textquotedblright{} that compact differentiable varieties are \textquotedblleft rigid\textquotedblright, i.e. admit no continuous deformation other than a trivial one. However, since this result will be essential to us in \S3, and references seem rather scarce, it may be useful to give precise statements and a few proofs.

We fix a topological space $S$. The mixed varieties considered will be assumed to be of class $C^r$, with $r\ge1$ (possibly $r=\infty$).

\sgaStableTarget{sga6.e02.para.2-1}\textbf{2.1. Notation.} Let $X,Y$ be mixed $S$-varieties. We denote by $\Homop(X,Y)$ the set of $S$-morphisms from $X$ to $Y$, and by $\underline{\Homop}(X,Y)$ the sheaf on $X$ of germs of $S$-morphisms from $X$ to $Y$ (defined by
\[
  \underline{\Homop}(X,Y)(U)=\Homop(U,Y)
\]
for $U$ open in $X$).

Let $s\in S$. We denote by $i_s$ the inclusion of $X_s$ into $X$. There is a canonical induction homomorphism
\sgaStableTarget{sga6.e02.eq.2-1-1}\begin{equation}\tag{2.1.1}
  i_s^{-1}\underline{\Homop}(X,Y)
  \longrightarrow
  \underline{\Homop}(X_s,Y_s),
\end{equation}
%%%%%%%%%%  scan idx 218  |  folio 205  %%%%%%%%%%
\SGAPrintPage{205}%
which is an epimorphism, as we see immediately. When $Y=S\times\underline{\underline{R}}$, the sheaf $\underline{\Homop}(X,Y)$ is canonically identified with $\underline{O}_X$, and the epimorphism (2.1.1) with the epimorphism
\[
  i_s^{-1}(\underline{O}_X)\longrightarrow i_s^*(\underline{O}_X)=\underline{O}_{X_s}
\]
considered in (1.1).

\sgaStableTarget{sga6.e02.prop.2-2}\textbf{Proposition}\SGAAuthorFootnote{1}{Communicated to the editor by A. Douady.} \textbf{2.2.} \uline{Let} $X,Y$ \uline{be mixed} $S$\uline{-varieties, and} $s\in S$\uline{. Suppose that the fiber} $X_s$ \uline{is paracompact and that} $s$ \uline{has a fundamental system of paracompact neighborhoods. Let} $A$ \uline{be a closed subset of} $X_s$\uline{, and} $f\in H^0(A,\underline{\Homop}(X_s,Y_s))$\uline{.
The subsheaf} $F$ \uline{of} $i_s^{-1}\underline{\Homop}(X,Y)$ \uline{formed by the germs \textquotedblleft that extend} $f$\uline{\textquotedblright{} (i.e. whose image under} (2.1.1) \uline{coincides with} $f$ \uline{over} $A$\uline{) is soft.}

\uline{Proof.} Let $X'$ be an open neighborhood of $x$ in $X$, isomorphic (as a mixed variety) to $S'\times U$, where $S'$ is a paracompact neighborhood of $s$ and $U$ a relatively compact open subset of $\underline{\underline{R}}^m$, and such that the image under $f$ of $X'_s\cap A$ is contained in an open subset of $Y$ isomorphic to $S'\times V$, where $V$ is a convex open subset of $\underline{\underline{R}}^n$ containing the origin. By (T.F. II 3.4.1), it suffices to show that every section of $F$ over a closed subset $A_0$ of $X_s$ contained in $X'_s$ extends to $X'_s$. Now $X'_s$ ($\simeq\{s\}\times U$) is relatively compact, so a closed subset of $X_s$ contained in $X'_s$ is compact. Identify $X'$ (resp. $Y'$) with $S'\times U$ (resp. $S'\times V$), put $A'=A\cap X'_s$, and let $A_0$ be a compact subset of $X'_s=\{s\}\times U$. Since $A'$ is closed in $\{s\}\times U$, which is paracompact, we have ([T.F.]~II 3.3.1)\SGAAuthorFootnote{1}{A closed subset of a paracompact space has a fundamental system of paracompact neighborhoods.}
\[
  H^0\bigl(A',\underline{\Homop}(X_s,Y_s)\bigr)
  =\varinjlim_{U'\supset A'}\Homop(U',Y'_s),
\]
over the open subsets $U'$ of $X'_s$ containing $A'$. Similarly, since $A_0$ is closed in $S'\times \overline U$, which is paracompact (a product of a paracompact space and a compact space), we have
\[
  H^0\bigl(A_0,i_s^{-1}\underline{\Homop}(X,Y)\bigr)
  =\mathop{\underrightarrow{\mathrm{Lim}}}\limits_{U_0\supset A_0}\Homop(U_0,Y'),
\]
over the open subsets $U_0$ of $X'$ containing $A_0$. The proposition will therefore be a
%%%%%%%%%%  scan idx 219  |  folio 206  %%%%%%%%%%
\SGAPrintPage{206}%
consequence of the following lemma:

\sgaStableTarget{sga6.e02.lem.2-2-1}\textbf{Lemma 2.2.1.} \uline{Let} $U$ \uline{be an open subset of} $\underline{\underline{R}}^m$, $V$ \uline{a convex open subset of} $\underline{\underline{R}}^n$ \uline{containing} $0$\uline{. Let} $A_0,A_1$ \uline{be closed subsets of} $\{s\}\times U$\uline{, let} $U_0$ (resp. $U_1$\uline{) be an open subset of} $S\times U$ (resp. $\{s\}\times U$\uline{) containing} $A_0$ (resp. $A_1$\uline{). Let} $f_1:U_1\to V$ \uline{be a morphism and let} $f_0:U_0\to S\times V$ \uline{be an} $S$\uline{-morphism inducing on} $U_0\cap(\{s\}\times U)$ \uline{a morphism defining the same germ as} $f_1$ \uline{in a neighborhood of} $A_1$\uline{. There then exists an open subset} $U'$ \uline{of} $S\times U$ \uline{containing} $s\times U$\uline{, and an} $S$\uline{-morphism} $\bar f:U'\longrightarrow S\times V$ \uline{such that}
\begin{itemize}
\item[(i)] \uline{the germ of} $\bar f$ \uline{along} $\{s\}\times U$ \uline{extends the germ defined by} $f_0$;
\item[(ii)] \uline{the morphism induced by} $\bar f$ \uline{on} $\{s\}\times U$ \uline{defines the same germ as} $f_1$ \uline{in a neighborhood of} $A_1$.
\end{itemize}

\uline{Proof.} Replacing $f_1$ by $af_1$, where $a$ is a $C^r$ function on $U$, equal to $1$ in a neighborhood of $A_1$ and with support in $U_1$, we may assume that $U_1=U$. Let $b$ be a $C^r$ function on $U$, equal to $1$ on $A_0$ and with support in $U_0\cap(\{s\}\times U)$. Then
\[
  U'=U_0\cup(S\times(U-\operatorname{Supp}b))
\]
is an open subset of $S\times U$ containing $\{s\}\times U$. The morphism $\bar f:U'\to S\times V$ defined by
\[
  \bar f(t,x)=\bigl(t,\ b(x)\pr_2f_0(t,x)+(1-b(x))f_1(x)\bigr)
\]
visibly satisfies conditions (i) and (ii).

% Author-claim audit ASTRA-II-AUTHOR-CLAIM-206 is recorded separately.
% The source statement and proof are retained in both English forms.
\sgaStableTarget{sga6.e02.cor.2-3}\textbf{Corollary 2.3.} \uline{Under the hypotheses of} (2.2)\uline{, the sheaves} $i_s^{-1}\underline{\Homop}(X,Y)$, $\underline{\Homop}(X_s,Y_s)$ \uline{are soft. In particular, the sheaves} $i_s^{-1}(\underline{O}_X)$, $\underline{O}_{X_s}$ \uline{are soft}\SGAAuthorFootnote{1}{For $\underline{O}_{X_s}$, this is \textquotedblleft very well known\textquotedblright. The editor does not know whether $\underline{O}_X$ is soft when the hypotheses of (2.2) hold for every $s\in S$ and $X$ is paracompact. Under those hypotheses, this would trivially imply the softness of $i_s^{-1}(\underline{O}_X)$ as well as (2.5) below.}\uline{. For every closed subset} $A$ \uline{of} $X_s$\uline{, the map}
\[
  H^0\bigl(X_s,i_s^{-1}\underline{\Homop}(X,Y)\bigr)
  \longrightarrow
  H^0\bigl(A,\underline{\Homop}(X_s,Y_s)\bigr)
\]
\uline{deduced from} (2.1.1) \uline{is surjective.}
%%%%%%%%%%  scan idx 220  |  folio 207  %%%%%%%%%%
\SGAPrintPage{207}%
\uline{Proof.} The softness of $i_s^{-1}\underline{\Homop}(X,Y)$ follows from (2.2) applied with $A=\varnothing$; taking $S=s$, we deduce that $\underline{\Homop}(X_s,Y_s)$ is soft. The last assertion follows from (2.2).

\sgaStableTarget{sga6.e02.cor.2-4}\textbf{Corollary 2.4.} \uline{Under the hypotheses of} (2.2)\uline{, let} $E$ \uline{be an} $\underline{O}_X$\uline{-Module. For every closed subset} $A$ \uline{of} $X_s$\uline{, the canonical homomorphism}
\[
  H^0\bigl(X_s,i_s^{-1}(E)\bigr)
  \longrightarrow
  H^0\bigl(A,i_s^*(E)\bigr),
\]
\uline{deduced from the epimorphism}
\[
  i_s^{-1}(E)\longrightarrow i_s^*(E)
  =i_s^{-1}(E)\otimes_{i_s^{-1}(\underline{O}_X)}\underline{O}_{X_s},
\]
\uline{is an epimorphism.}

\uline{Proof.} Since $i_s^{-1}(\underline{O}_X)$ is soft (2.3), the functor $H^0(A,.)$ is exact on the category of $i_s^{-1}(\underline{O}_X)$-Modules (T.F. II 3.7.1, 3.4.2, 4.4.3), so the homomorphism
\[
  H^0\bigl(A,i_s^{-1}(E)\bigr)
  \longrightarrow
  H^0\bigl(A,i_s^*(E)\bigr)
\]
is an epimorphism. Since $i_s^{-1}(E)$ is soft, the homomorphism
\[
  H^0\bigl(X_s,i_s^{-1}(E)\bigr)
  \longrightarrow
  H^0\bigl(A,i_s^{-1}(E)\bigr)
\]
is an epimorphism, and we are done.

\sgaStableTarget{sga6.e02.cor.2-5}\textbf{Corollary 2.5.} \uline{Let} $p:X\to S$ \uline{be a mixed variety, and} $s\in S$\uline{. Suppose that} $s$ \uline{has a fundamental system of paracompact neighborhoods} $S'$ \uline{such that the} $p^{-1}(S')$ \uline{form a fundamental system of paracompact neighborhoods of} $X_s$ \uline{(a condition satisfied, for example, when} $S$ \uline{is locally compact and} $p$ \uline{is a proper map of topological spaces). Then, for every} $\underline{O}_X$\uline{-Module} $E$\uline{, we have}
\[
  R^i p_*(E)_s=0 \quad \text{for } i>0.
\]
%%%%%%%%%%  scan idx 221  |  folio 208  %%%%%%%%%%
\SGAPrintPage{208}%
\uline{Proof.} By (T.F. II 4.17.1), we have
\[
  R^i p_*(E)_s=H^i\bigl(X_s,i_s^{-1}(E)\bigr).
\]
But, since $i_s^{-1}(\underline{O}_X)$ is soft (2.3), we have
\[
  H^i\bigl(X_s,i_s^{-1}(E)\bigr)=0\quad \text{for } i>0,
\]
whence the assertion.

\sgaStableTarget{sga6.e02.cor.2-6}\textbf{Corollary 2.6.} \uline{Let} $X,Y$ \uline{be mixed} $S$\uline{-varieties, and} $s\in S$\uline{. Suppose that} $X$ \uline{satisfies the hypotheses of} (2.5)\uline{. For every morphism} $f:X_s\to Y_s$\uline{, there exists an open neighborhood} $S'$ \uline{of} $s$ \uline{and a morphism} $\bar f:X|S'\longrightarrow Y|S'$ \uline{such that} $\bar f_s=f$.

\uline{Proof.} This follows immediately from (2.3) (with $A=X_s$) and the fact that, by the hypotheses on $X$, we have
\[
  H^0\bigl(X_s,i_s^{-1}\underline{\Homop}(X,Y)\bigr)
  =\operatorname*{Lim}_{S'\ni s}\Homop(X|S',Y)\qquad (S'\ \text{open})
\]
(T.F. II 3.3.1).

\sgaStableTarget{sga6.e02.prop.2-7}\textbf{Proposition 2.7.} \uline{Let} $X,Y$ \uline{be mixed} $S$\uline{-varieties,} $f:X\to Y$ \uline{a morphism,} $s\in S$, $x\in X_s$\uline{. If the tangent linear map of} $f_s:X_s\to Y_s$ \uline{is invertible at} $x$\uline{, then} $f$ \uline{is a local isomorphism at} $x$\uline{, i.e. there exists an open neighborhood} $U$ \uline{of} $x$ \uline{in} $X$ \uline{and an open subset} $V$ \uline{of} $Y$ \uline{such that} $f$ \uline{induces an isomorphism from} $U$ \uline{onto} $V$.

\uline{Proof.} This is an easy consequence of the implicit function theorem. The details are left to the reader.

\sgaStableTarget{sga6.e02.cor.2-8}\textbf{Corollary 2.8.} \uline{Let} $X,Y$ \uline{be mixed} $S$\uline{-varieties, and} $s\in S$\uline{. Suppose that} $s$ \uline{has a fundamental system of neighborhoods} $S'$ \uline{such that the} $X|S'$ \uline{(resp.} $Y|S'$\uline{) form a fundamental system of neighborhoods of} $X_s$ \uline{(resp.} $Y_s$\uline{) (a condition satisfied, for example, when} $S$ \uline{is locally compact and} $X$ \uline{(resp.} $Y$\uline{) is proper over} $S$\uline{). Let} $f:X\to Y$ \uline{be an} $S$\uline{-morphism such that} $f_s:X_s\to Y_s$ \uline{is an isomorphism. There then exists a neighborhood} $S'$

%%%%%%%%%%  scan idx 222  |  folio 209  %%%%%%%%%%
\SGAPrintPage{209}%
\uline{of} $s$ \uline{such that} $f|X|S':X|S'\longrightarrow Y|S'$
\uline{is an isomorphism.}

\uline{Proof.} By (2.7), there exists an open neighborhood $U$ of $X_s$ in $X$ such that $f|U$ is étale (i.e. a local isomorphism). We may assume that $U=X|S'$, where $S'$ is a neighborhood of $s$ in $S$; in other words, after restricting $S$, we may assume that $f$ is étale. We then easily check that the fiber product $Z=X\times_Y X$ is naturally equipped with the structure of a mixed $S$-variety (étale over $X$ for each projection), and that the diagonal morphism $\Delta:X\to Z$ is an open immersion\SGAAuthorFootnote{1}{Note the analogy with Algebraic Geometry (SGA 1 I 5.1).}. Moreover, it is immediate that $Z$ satisfies the same hypothesis as $X$ and $Y$, namely that the open subsets of the form $Z|S'$, for $S'$ an open neighborhood of $s$ in $S$, form a fundamental system of neighborhoods of $Z_s$. Since $\Delta_s:X_s\to Z_s$ is an isomorphism, it follows, after restricting $S$, that $\Delta$ is an isomorphism. Thus $f$ is an open immersion, hence an isomorphism after shrinking $S$ again, and we are done.

\sgaStableTarget{sga6.e02.cor.2-9}\textbf{Corollary 2.9.} \uline{Let} $X$ \uline{be a mixed} $S$\uline{-variety, and} $s\in S$\uline{. Suppose that the hypothesis of} (2.5) \uline{holds. There then exists a neighborhood} $S'$ \uline{of} $s$ \uline{in} $S$ \uline{such that} $X|S'$ \uline{is} $S'$\uline{-isomorphic to} $S'\times X_s$.

\uline{Proof.} Apply (2.6) and (2.8).

\sgaStableTarget{sga6.e02.cor.2-10}\textbf{Corollary 2.10.} \uline{Let} $s\in S$ \uline{and let} $V$ \uline{be a paracompact variety. Put} $X=S\times V$\uline{, and denote by} $i_s$ \uline{the injection of} $V$ \uline{into} $X$ \uline{defined by} $s$\uline{. Suppose that the hypothesis of} (2.5) \uline{holds. Let} $E$ \uline{be a locally free} $\underline{O}_X$ \uline{(resp.} $\underline{O}_X\otimes\underline{\underline{C}}$\uline{)-Module of finite type. There then exists a neighborhood} $S'$ \uline{of} $s$ \uline{in} $S$ \uline{such that} $E|X|S'$ \uline{is isomorphic to} $\pr_2^*i_s^*(E)$ ($\pr_2:S'\times V\to V$ \uline{being the second projection).}

\uline{Proof.} Put $F=\pr_2^*i_s^*(E)$. By (2.4), applied to the Module $\underline{\Homop}(E,F)$, the induction homomorphism
\[
  \Homop\bigl(i_s^{-1}(E),i_s^{-1}(F)\bigr)
  \longrightarrow
  \Homop\bigl(i_s^*(E),i_s^*(F)\bigr)
\]
%%%%%%%%%%  scan idx 223  |  folio 210  %%%%%%%%%%
\SGAPrintPage{210}%
is an epimorphism. On the other hand, we have (T.F. II 3.3.1):
\[
\begin{aligned}
  \Homop\bigl(i_s^{-1}(E),i_s^{-1}(F)\bigr)
  &\simeq H^0\bigl(X_s,i_s^{-1}\underline{\Homop}(E,F)\bigr)\\
  &=\varinjlim_{S'\ni s} H^0\bigl(X|S',\underline{\Homop}(E,F)\bigr)\qquad (S'\ \text{open})\\
  &\simeq \varinjlim_{S'\ni s}\Homop(E|X|S',F|X|S').
\end{aligned}
\]
Thus the canonical isomorphism $i_s^*(E)\simeq i_s^*(F)$ is induced by a homomorphism $f:E\to F$ (after shrinking $S$). Since $f_s$ is an isomorphism, $f$ is an isomorphism in a neighborhood of $X_s$ by Nakayama, hence over $X|S'$, where $S'$ is a suitable neighborhood of $s$, q.e.d.

\sgaStableTarget{sga6.e02.sec.3.occ2}\subsection*{3. The finiteness theorem}

\sgaStableTarget{sga6.e02.sec.3.0}\subsubsection*{3.0. Conventions and notation}
\sgaStableTarget{sga6.e02.para.3-0-1}3.0.1. Let $X$ and $Y$ be topological spaces. We denote by $C(X,Y)$ the set of continuous maps from $X$ to $Y$, and by $\underline C(X,Y)$ the sheaf on $X$ defined by $U\mapsto C(U,Y)$.

\sgaStableTarget{sga6.e02.para.3-0-2}3.0.2. Let $E,F$ be Banach spaces (real or complex). We denote by $\Lop(E,F)$ the Banach space of continuous linear maps from $E$ to $F$.

\sgaStableTarget{sga6.e02.para.3-0-3}3.0.3. The mixed varieties considered in this section will be assumed to be of class $C^\infty$. Let $X$ be a mixed $S$-variety; we shall often identify a locally free $\underline{O}_X$ (resp. $\underline{O}_X\otimes\underline{\underline{C}}$)-Module of finite type with the corresponding real (resp. complex) vector bundle.

\sgaStableTarget{sga6.e02.para.3-0-4}3.0.4. Let $X$ be a compact (ordinary) variety, and let $E$ be a vector bundle (real or complex) of finite rank on $X$. We denote by $C^\infty(X,E)$ the Fréchet space of sections (of class $C^\infty$) of $E$ on $X$. For $r\in\underline{\underline{Z}}$, we denote by $C^r_2(X,E)$ the Hilbert space of sections of $E$ that are $r$ times differentiable in $L^2$ ([16]~IX). For $s\ge r$, $C^s_2(X,E)$ maps into $C^r_2(X,E)$ by a continuous injective linear map with dense image, and we have
%%%%%%%%%%  scan idx 224  |  folio 211  %%%%%%%%%%
\SGAPrintPage{211}%
(loc.cit.)
\[
  C^\infty(X,E)=\varprojlim_r C^r_2(X,E).
\]

\sgaStableTarget{sga6.e02.para.3-0-5}3.0.5.
Let $S$ be a topological space. We shall need the notion of a \textquotedblleft Hilbert bundle\textquotedblright{} over $S$, for which we refer the reader to [12] or [14]. There is an additive functor from the category of Hilbert bundles (real resp. complex) over $S$ \var{to that of $\underline{O}_S$ rep.}{to that of $\underline{O}_S$ resp.}{EN-VAR-147} $\underline{O}_S\otimes\underline{\underline{C}}$-Modules, associating to each Hilbert bundle $E$ the sheaf $\underline E$ of continuous sections of $E$.

\sgaStableTarget{sga6.e02.sec.3.1}\subsubsection*{3.1. Auxiliary technical results}
\sgaStableTarget{sga6.e02.lem.3-1-1}\textbf{Lemma 3.1.1.} \uline{Let} $X,Y$ \uline{be topological spaces, with} $X$ \uline{locally compact. Denote by} $\underline C'(X,Y)$ \uline{the presheaf defined on the basis of relatively compact open subsets of} $X$ \uline{by} $\underline C'(X,Y)(U)=C(\bar U,Y)$ (3.0.1)\uline{. The canonical restriction morphism}
\[
  \underline C'(X,Y)\longrightarrow \underline C(X,Y)
\]
\uline{induces an isomorphism on the associated sheaves.}

\uline{Proof.} This is immediate.

\sgaStableTarget{sga6.e02.lem.3-1-2}\textbf{Lemma 3.1.2.} \uline{Let} $E,F$ \uline{be locally convex topological vector spaces (over} $\underline{\underline{C}}$\uline{, to fix ideas), and} $u:E\to F$ \uline{a continuous linear map (resp. continuous with dense image). Let} $X$ \uline{be a compact space. Equip} $C(X,E)$, $C(X,F)$ \uline{with the topology of uniform convergence. Then the linear map}
\[
  u_X:C(X,E)\longrightarrow C(X,F)
\]
\uline{induced by} $u$ \uline{is continuous (resp. continuous with dense image).}
%%%%%%%%%%  scan idx 225  |  folio 212  %%%%%%%%%%
\SGAPrintPage{212}%
\uline{Proof.} It is clear that $u_X$ is continuous. On the other hand (Bourbaki, Intégration, chap. III \S~2, lemmas 1 and 2), the subspace $C(X,\underline{\underline{C}})\otimes_{\underline{\underline{C}}}E$ of $C(X,E)$ (formed by the continuous maps from $X$ to $E$ whose values lie in a finite-dimensional vector subspace) is dense; similarly, $C(X,\underline{\underline{C}})\otimes_{\underline{\underline{C}}}F$ is dense in $C(X,F)$. Suppose $u$ is continuous with dense image. Then
\[
  u_X(C(X,\underline{\underline{C}})\otimes_{\underline{\underline{C}}}E)=C(X,\underline{\underline{C}})\otimes_{\underline{\underline{C}}}u(E)
\]
is dense in $C(X,\underline{\underline{C}})\otimes_{\underline{\underline{C}}}F$, whence, taking closures, we deduce that $u_X(C(X,E))$ is dense in $C(X,F)$, q.e.d.

\sgaStableTarget{sga6.e02.para.3-1-3}3.1.3. Let $V$ be a compact variety, $E$ and $F$ complex vector bundles\SGAAuthorFootnote{1}{To fix ideas (analogous results hold in the real case).} of finite rank on $V$, and $n\in\underline{\underline{N}}$. The space $\Diff^n_{\underline{\underline{C}}}(E,F)$, identified by (1.2.3.1) with
\[
  \Homop(P^n(E),F)\simeq C^\infty(V,\underline{\Homop}(P^n(E),F)),
\]
is naturally equipped with the structure of a Fréchet space. Let $u\in\Diff^n_{\underline{\underline{C}}}(E,F)$; we easily check that the linear map from $C^\infty(V,E)$ to $C^\infty(V,F)$ defined by $u$ is continuous. Denoting by $\Lop(C^\infty(V,E),C^\infty(V,F))$ the space of continuous linear maps from $C^\infty(V,E)$ to $C^\infty(V,F)$, we therefore have a linear map
\[
  \phi:\Diff^n_{\underline{\underline{C}}}(E,F)
  \longrightarrow \Lop(C^\infty(V,E),C^\infty(V,F)).
\]
On the other hand, for $r\in\underline{\underline{Z}}$, $u$ acts by extension of scalars ([16]~IX Th.~12) on the $C^r_2$ sections of $E$ and defines a linear map
\[
  \phi_r:\Diff^n_{\underline{\underline{C}}}(E,F)
  \longrightarrow \Lop(C^r_2(V,E),C^{r-n}_2(V,F)).
\]
We readily check that:
\begin{enumerate}[label=(\roman*)]
\item the $\phi_r$ are \uline{continuous} (the target is a Banach space
%%%%%%%%%%  scan idx 226  |  folio 213  %%%%%%%%%%
\SGAPrintPage{213}%
(3.0.2));
\item $\phi=\varprojlim\phi_r$ (i.e.
for $u\in\Diff^n_{\underline{\underline{C}}}(E,F)$, $\phi(u)$ is the limit of the projective system of the $\phi_r(u)$).
\end{enumerate}
(The preceding assertions follow easily from the \var{results of (loc.cit):}{results of (loc.cit.):}{EN-VAR-148} we reduce immediately to $n=0$; (ii) then follows directly from the definitions, while (i) follows (after localizing the problem by a partition of unity) from the formula giving the successive derivatives of a product).

\sgaStableTarget{sga6.e02.lem.3-1-4}\textbf{Lemma 3.1.4.} \uline{Let} $S$ \uline{be a topological space,} $V$ \uline{a compact variety,} $E$ \uline{and} $F$ \uline{complex vector bundles of finite rank on} $V$\uline{, and} $n\in\underline{\underline{N}}$\uline{. Denote by} $X$ \uline{the mixed} $S$\uline{-variety} $S\times V$\uline{. For} $u\in\Diff^n_{\underline{\underline{C}}}(\pr_2^*(E),\pr_2^*(F))$\uline{, the map} $s\mapsto i_s^*(u)$ \uline{from} $S$ \uline{to} $\Diff^n_{\underline{\underline{C}}}(E,F)$ \uline{is continuous. The linear map}
\[
  \Diff^n_{\underline{\underline{C}}}(\pr_2^*(E),\pr_2^*(F))
  \longrightarrow C(S,\Diff^n_{\underline{\underline{C}}}(E,F))
\]
\uline{thus defined is an isomorphism.}

\uline{Proof.} Using the identification (1.2.3.1), we may restrict to $n=0$. Putting $\underline{\Homop}(E,F)=G$, we are then reduced to showing that, for $u\in\Gamma(X,\pr_2^*(G))$, the map $s\mapsto u_s=i_s^*(u)$ from $S$ to $C^\infty(V,G)$ is continuous, and that the linear map
\[
  \Gamma(X,\pr_2^*(G))\longrightarrow C(S,C^\infty(V,G))
\]
so obtained is an isomorphism. But this follows immediately from the definition of mixed varieties (1.1).

\sgaStableTarget{sga6.e02.para.3-1-5}3.1.5. Under the hypotheses of (3.1.4), let
\[
  u\in\Diff^n(\pr_2^*(E),\pr_2^*(F)).\text{\SGAAuthorFootMark{1}}
\]
\SGAAuthorFootText{1}{We omit the subscript $\underline{\underline{C}}$ for brevity.}

%%%%%%%%%%  scan idx 227  |  folio 214  %%%%%%%%%%
\SGAPrintPage{214}%
Identifying $u$, by (3.1.4), with a continuous map from $S$ to $\Diff^n(E,F)$, we see that, for $U$ open in $S$, the morphism
\[
  \pr_{1*}(u)(U):\Gamma(X|U,E)\longrightarrow \Gamma(X|U,F)
\]
is canonically identified with the linear map
\[
  C(U,C^\infty(V,E))\longrightarrow C(U,C^\infty(V,F))
\]
deduced from
\begin{equation}\tag{3.1.3}
  \phi u:U\longrightarrow \Lop(C^\infty(V,E),C^\infty(V,F)).
\end{equation}
Suppose $S$ is locally compact. Then (3.1.1) the morphism
\[
  \pr_{1*}(u):\pr_{1*}\pr_2^*(E)
  \longrightarrow \pr_{1*}\pr_2^*(F)
\]
is the morphism of sheaves associated to the morphism of presheaves
\[
  u_\infty:\underline C'(S,C^\infty(V,E))
  \longrightarrow \underline C'(S,C^\infty(V,F))
\]
defined by $\phi u$. On the other hand, for $r\in\underline{\underline{Z}}$,
\[
  \phi_r u:S\longrightarrow \Lop(C^r_2(V,E),C^{r-n}_2(V,F))
\]
defines a morphism of presheaves
\[
  \var{u_r:\underline C'(S,C^r_2(V,E))
  \longrightarrow \underline C'(S,C^r_2(V,F)).}{u_r:\underline C'(S,C^r_2(V,E))
  \longrightarrow \underline C'(S,C^{r-n}_2(V,F)).}{EN-VAR-149}
\]
By (3.1.3 (ii)), the $u_r$ form a projective system whose limit is $u_\infty$. It will be convenient to interpret $\phi_r u$ as a morphism of trivial Hilbert bundles (3.0.5)
\[
  S\times C^r_2(V,E)\longrightarrow S\times C^{r-n}_2(V,F),
\]
%%%%%%%%%%  scan idx 228  |  folio 215  %%%%%%%%%%
\SGAPrintPage{215}%
inducing on the sheaves of sections the morphism of sheaves associated to $u_r$.

\sgaStableTarget{sga6.e02.sec.3.2}\subsubsection*{3.2. Definition of the index}
\sgaStableTarget{sga6.e02.thm.3-2-1}\textbf{Theorem 3.2.1.} \uline{Let} $S$ \uline{be a locally compact space,} $p:X\to S$ \uline{a proper mixed variety,} $E^\bullet$ \uline{an elliptic complex of order} $\underline{n}$ \uline{of complex vector bundles on} $X$ (1.2.5)\uline{. Then the complex} $\underline{\underline{R}}p_*(E^\bullet)$ \uline{(which, by} (2.5)\uline{, is isomorphic, to} $p_*(E^\bullet)$ \uline{in the derived category of that of} $\underline{O}_S\otimes\underline{\underline{C}}$\uline{-Modules) is perfect.}

\uline{Proof.} The question being local on $S$, we may assume (2.9) that $X=S\times V$, where $V$ is a compact variety, and $p=\pr_1$. Moreover, we may assume (2.10) that, for every $i$, $E^i=\pr_2^*(F^i)$, where $F^i$ is a complex vector bundle of finite rank on $V$. By (3.1.5), the complex $p_*(E)$ is associated to the complex of presheaves $\bigl(\underline C'(S,C^\infty(V,F^i)),d_\infty^i\bigr)$ (where the $d^i:E^i\longrightarrow E^{i+1}$ are the given differential operators), which we shall denote simply by $\underline C'(S,C^\infty(V,F^\bullet))$. Put, for each $i\in\underline{\underline{Z}}$,
\[
  m_i=\sum_{j<i} n_j
  \qquad (\text{we agree that } n_i=0 \text{ when } E^i=0).
\]
For each $r\in\underline{\underline{Z}}$, we have (3.1.5) a complex of presheaves $\bigl(\underline C'(S,C_2^{r-m_i}(V,F^i)),d_{r-m_i}^i\bigr)$, which we shall denote by $\underline C'(S,C_2^r(V,F^\bullet))$.
The $\underline C'(S,C_2^r(V,F^\bullet))$ form a projective system, and we have (3.1.5):
\sgaStableTarget{sga6.e02.eq.3-2-1-1}\begin{equation}\tag{3.2.1.1}
  \underline C'(S,C^\infty(V,F^\bullet))=\varprojlim_r \underline C'(S,C_2^r(V,F^\bullet)).
\end{equation}
Finally, we shall denote by $F_r^\bullet$ the complex of trivial Hilbert bundles over $S$ defined by $\phi_{r-m_i}d^i$, and by $\underline F_r^\bullet$ the corresponding complex of sheaves of sections, which is precisely (3.1.5) the complex of sheaves associated to $\underline C'(S,C_2^r(V,F^\bullet))$. With this notation, to prove the theorem it suffices to prove the following assertions:

(a) \uline{For} $r\in\underline{\underline{Z}}$\uline{, the canonical morphism}
%%%%%%%%%%  scan idx 229  |  folio 216  %%%%%%%%%%
\SGAPrintPage{216}%
\[
  \underline C'(S,C^\infty(V,F^\bullet))
  \longrightarrow \underline C'(S,C_2^r(V,F^\bullet))
\]
\uline{is a
quasi-isomorphism of complexes of presheaves, i.e.\ induces isomorphisms on the cohomology presheaves.}

(b) \uline{For} $r\in\underline{\underline{Z}}$\uline{, the complex} $\underline F_r^\bullet$ \uline{is perfect.}

\uline{Proof of} (b). The ellipticity of $E^\bullet$ implies ([16]~XI Th.6\SGAAuthorFootnote{1}{which readily generalizes to elliptic complexes.}) that, for every $s\in S$, the differential of $(F_r^\bullet)_s$ has closed image and the $H^i((F_r^\bullet)_s)$ are finite-dimensional vector spaces. In other words, to use the terminology of [11], $F_r^\bullet$ is a direct quasi-acyclic complex of Hilbert bundles. Consequently, by ([12]~2.3), $F_r^\bullet$ is locally homotopically equivalent to a bounded complex of vector bundles of finite rank. Since the \textquotedblleft sheaf of sections\textquotedblright{} functor is additive (3.0.5), $\underline F_r^\bullet$ is locally homotopically equivalent to a bounded complex of locally free sheaves of finite type, and hence is perfect.

For the proof of (a), we shall need the following intermediate assertion:

(a') \uline{For} $p\geq q$\uline{, the canonical morphism}
\[
  \underline F_p^\bullet\longrightarrow \underline F_q^\bullet
\]
\uline{is a quasi-isomorphism.}

\uline{Proof of} (a'). Denote by $M_{pq}^\bullet$ the complex of Hilbert bundles that is the cone of the morphism $F_p^\bullet\longrightarrow F_q^\bullet$. The ellipticity of $E$ implies that, for every $s\in S$, the morphism $(F_p^\bullet)_s\longrightarrow (F_q^\bullet)_s$ is a quasi-isomorphism ([16]~XI Cor.~2 of Th.~6). In other words, $M_{pq}^\bullet$ is a direct acyclic complex. It follows ([11]~2.3) that $M_{pq}^\bullet$ is locally homotopically trivial, and hence that $\underline M_{pq}^\bullet$, which is precisely the cone of $\underline F_p^\bullet\longrightarrow\underline F_q^\bullet$, is acyclic, q.e.d.

\uline{Proof of} (a). We must show that, for every relatively compact open subset $U$ of $S$, the morphism
%%%%%%%%%%  scan idx 230  |  folio 217  %%%%%%%%%%
\SGAPrintPage{217}%
\[
  C(\overline U,C^\infty(V,F^\bullet))
  \longrightarrow C(\overline U,C_2^r(V,F^\bullet))
\]
is a
quasi-isomorphism. Replacing $S$ by $\overline U$ (and $E^\bullet$ by the complex induced over $\overline U$), we are reduced to showing that
\[
  C(S,C^\infty(V,F^\bullet))
  \longrightarrow C(S,C_2^r(V,F^\bullet))
\]
is a quasi-isomorphism, with $S$ assumed \uline{compact}. By (3.2.1.1), we have
\[
  C(S,C^\infty(V,F^\bullet))=
  \varprojlim_r C(S,C_2^r(V,F^\bullet))
\]
For each $i$, $C(S,C_2^r(V,F^i))$ is a Banach space (for the topology of uniform convergence), and the transition morphisms
\[
  C(S,C_2^p(V,F^i))\longrightarrow C(S,C_2^q(V,F^i))
\]
are, by (3.0.4) and (3.1.2), continuous linear maps with dense image. On the other hand, applying to the quasi-isomorphisms (a') $\underline F_p^\bullet\longrightarrow\underline F_q^\bullet$ the functor $\Gamma_S$, which is exact since $\underline{O}_S$ is soft, we see that the transition morphisms
\[
  C(S,C_2^p(V,F^\bullet))\longrightarrow C(S,C_2^q(V,F^\bullet))
\]
are quasi-isomorphisms. Assertion (a) then follows from (EGA $0_{\mathrm{III}}$ 13.2.3, \textquotedblleft topological\textquotedblright{} variant). This completes the proof of (3.2.1).

\sgaStableTarget{sga6.e02.para.3-2-2}3.2.2. For a topological space $S$, we shall denote by $\Parf(S)$ the triangulated category of perfect complexes of $(\underline{O}_S\otimes\underline{\underline{C}})$-Modules, and by $K^\bullet(S)$ the corresponding \textquotedblleft Grothendieck group\textquotedblright{}\SGAAuthorFootnote{1}{see (I 6.3). For a general review of the functors $K$, see Exp.~IV.}.

Let $S$, $X$, $E^\bullet$ be as in (3.2.1); the \uline{analytic index} of $E^\bullet$ is the element
\[
  \operatorname{ind.an}(E^\bullet)=C\ell(\underline{\underline{R}}p_*(E^\bullet))\in K^\bullet(S).
\]
%%%%%%%%%%  scan idx 231  |  folio 218  %%%%%%%%%%
\SGAPrintPage{218}%
For \uline{compact} $S$, $\operatorname{ind.an}(E^\bullet)$ canonically defines an element in the usual group $K(S)$ of topologists. Indeed, in general, for every topological space $S$, denote by $\mathrm{Loclib}(S)$ the additive category of locally free $(\underline{O}_S\otimes\underline{\underline{C}})$-Modules of finite type, and by $K(S)$ the \textquotedblleft Grothendieck group\textquotedblright{} of the triangulated category $K^b(\mathrm{Loclib}(S))/K^{b,\varnothing}(\mathrm{Loclib}(S))$ (the latter is, moreover, identified with $K^b(\mathrm{Loclib}(S))$ when $S$ is paracompact (II 2.3.2 (a)), so that in this case (I 6.4) $K(S)$ is canonically isomorphic to the \textquotedblleft Grothendieck group\textquotedblright{} of the additive category $\mathrm{Loclib}(S)$). There is a natural homomorphism
\[
  K(S)\longrightarrow K^\bullet(S).
\]
But, by (II 2.3.2), this is an isomorphism for compact $S$, proving our assertion. More generally, $\operatorname{ind.an}(E^\bullet)$ defines an element in the completed group
\[
  \widehat K(S)=\varprojlim K(A),
\]
the projective limit being taken over the compact subsets $A$ of $S$.

\sgaStableTarget{sga6.e02.para.3-2-3}3.2.3. The writer does not know whether the finiteness theorem (3.2.1) remains valid when, instead of assuming $S$ locally compact, one assumes only that every point of $S$ has a fundamental system of paracompact neighborhoods. Moreover, the local compactness hypothesis on $S$ does not seem very natural. The space of \textquotedblleft universal local parameters\textquotedblright{} for elliptic complexes of fixed order is indeed an (infinite-dimensional) subspace of a Fréchet space. Let us make this precise. Let $V$ be a compact variety, $F=(F^i)_{i\in\underline{\underline{Z}}}$ a sequence of complex vector bundles of finite rank on $V$, all but finitely many zero, and $\underline n=(n_i)$ a sequence of integers $\geq 0$. Denote by $\operatorname{Comp.diff}^{\underline n}(F)$ the space of complexes of differential operators \textquotedblleft with base $F$\textquotedblright{} and of order $\underline n$, i.e.\ the closed subspace of $\prod\Diff^{n_i}(F^i,F^{i+1})$ defined by the equations $d^{i+1}d^i=0$. (the space $\prod\Diff^{n_i}(F^i,F^{i+1})$ is a Fréchet space). The subset
%%%%%%%%%%  scan idx 232  |  folio 219  %%%%%%%%%%
\SGAPrintPage{219}%
$\operatorname{Ell}^{\underline n}(F)$ of
$\operatorname{Comp.diff}^{\underline n}(F)$ formed by the elliptic complexes of order $\underline n$ is open; it is this space (which in general is not locally compact) that plays the role of a space of universal local parameters.

\sgaStableTarget{sga6.e02.para.3-2-4}3.2.4. Nevertheless, even should the generalization of (3.2.1) mentioned above prove impossible, we can always associate to $E^\bullet$, under the sole hypothesis that the fibration $p$ is proper and locally trivial, an element of $K^\bullet(S)$ deserving the name analytic index of $E^\bullet$. Indeed, suppose (for simplicity\SGAAuthorFootnote{1}{for the general case, use the notion of a \textquotedblleft weak Hilbert bundle\textquotedblright{} or continuous field of Hilbert spaces from [12].}) that we are in the \textquotedblleft local\textquotedblright{} situation where $X=S\times V$ ($V$, a compact variety), with $E^\bullet$ given by a continuous map from $S$ to $\operatorname{Ell}^{\underline n}(F)$ (3.2.3). Assertions (a') and (b) in the proof of (3.2.1) are in fact valid without any hypothesis on $S$: it is therefore reasonable to put, with the notation of (loc.cit.):
\[
  \operatorname{ind.an}(E^\bullet)=C\ell(\underline F_r^\bullet)\in K^\bullet(S),
\]
this element being independent of $r$.

\sgaStableTarget{sga6.e02.para.3-2-5}3.2.5. Suppose $S$ is paracompact, and denote by $H$ a complex Hilbert space of infinite Hilbert dimension, and by $I(H)$ the subspace of $L(H,H)$ formed by the operators \textquotedblleft with index\textquotedblright{}. Then the complexes $F_r^\bullet$ define [12] a canonical element in the group $[S,I(H)]$ of homotopy classes of continuous maps from $S$ to $I(H)$, but the writer does not know, for noncompact $S$ (even when $S$ is locally compact), \var{the relations between the groups $[S,I(S)]$ and}{the relations between the groups $[S,I(H)]$ and}{EN-VAR-150} $K^\bullet(S)$ (is there a natural arrow from one to the other? are they isomorphic?), and therefore a fortiori between the two indices thus defined.

For a generalization, under the hypotheses of (3.2.1), of the Atiyah--Singer index theorem, we refer to the recent work of Shih [17], or to work in preparation by M. F. Atiyah and I. M. Singer.
%%%%%%%%%%  scan idx 233  |  folio 220  %%%%%%%%%%
\SGAPrintPage{220}%
\sgaStableTarget{sga6.e02.bibliography}\subsection*{BIBLIOGRAPHY}
\begin{enumerate}[label={[\arabic*]}]
\item
Artin, M., Grothendieck, A., and Verdier, J-L., Cohomologie étale des schémas, to appear in North Holl. Pub. Comp., (cited as SGA 4).
\item Atiyah, M. and Segal, G., secret papers.
\item Cartan, H. Funzioni e varietà complesse, Faisceaux analytiques cohérents, Course at the Centro Internazionale Matematico Estivo, 1964, Roma.
\item Cartan, H., and Schwartz, L., Théorème d'Atiyah--Singer, Seminar 1963-64, Institut Henri Poincaré, Paris.
\item Frisch, J., \var{Fonctions analyliques sur un ensemble semi-analytique}{Fonctions analytiques sur un ensemble semi-analytique}{EN-VAR-151}, note in CRAS Paris, t. 260, p. 2974-2976 (1965).
\item Godement, R., Théorie des Faisceaux, Hermann (1958), (cited as T.F).
\item Grothendieck, A., Cohomologie locale des faisceaux cohérents et théorèmes de Lefschetz locaux et globaux, Séminaire de Géométrie Algébrique de l'IHES 1962, to appear in North Holl. Pub. Comp., (cited as SGA 2).
\item Grothendieck, A., Fondements de la Géométrie Analytique, Séminaire Cartan 1959-60, Institut Henri Poincaré, Paris.
\item Hartshorne, R., Residues and Duality, Lecture Notes, Springer n\textsuperscript{o}~20, cited as [H].
\item Houzel, C., Géométrie Analytique Locale, Séminaire Cartan 1959-60, Institut Henri Poincaré, Paris.
\item Illusie, L., Complexes de fibrés et définitions du foncteur $K$, Séminaire Shih Weishu, IHES 1964-65.
\item Illusie, L., Complexes quasi-acycliques directs de fibrés banachiques, note in CRAS Paris, t. 260, p. 6499-6502.
\item Kiehl, R., Theorem A und Theorem B in der nichtarchimedischen Funktionentheorie, Invent. math. 2, p. 256-273 (1967).
%%%%%%%%%%  scan idx 234  |  folio 221  %%%%%%%%%%
\SGAPrintPage{221}%
\item
Lang, S., Introduction to differentiable manifolds,
\item \L{}ojasiewicz, S., Ann. della Scuola Norm. Sup. Pisa, serie III, 18, fasc. IV, 1964.
\item Palais, R., Lectures at the Princeton Seminar 1963-64 on the index theorem.
\item Shih, W., Fiber cobordism and the index of a family of elliptic differential operators,
\item Verdier, J.-L., Les catégories dérivées des catégories abéliennes, to appear in North Holl. Pub. Comp.
\end{enumerate}

% Original English continuation definitions, retained in their original order.
\providecommand{\Ob}{\operatorname{Ob}}
\providecommand{\qcoh}{\mathrm{qcoh}}
\providecommand{\coh}{\mathrm{coh}}
\providecommand{\torf}{\mathrm{torf}}
\providecommand{\toramp}{\operatorname{tor.amp}}
\providecommand{\amp}{\operatorname{amp}}
\providecommand{\Id}{\operatorname{Id}}
\providecommand{\pr}{\operatorname{pr}}
\providecommand{\coker}{\operatorname{coker}}
\providecommand{\im}{\operatorname{Im}}
\providecommand{\Spec}{\operatorname{Spec}}

% Reviewed French continuation definitions, retained verbatim.
% Exposé III additions
\renewcommand{\Ob}{\operatorname{ob}}
\providecommand{\qcoh}{\mathrm{qcoh}}
\providecommand{\coh}{\mathrm{coh}}
\providecommand{\torf}{\mathrm{torf}}
\providecommand{\toramp}{\operatorname{tor.amp}}
\providecommand{\amp}{\operatorname{amp}}
\providecommand{\Id}{\operatorname{Id}}
\providecommand{\pr}{\operatorname{pr}}
\providecommand{\coker}{\operatorname{coker}}
\providecommand{\im}{\operatorname{Im}}
\providecommand{\Spec}{\operatorname{Spec}}

%%%%%%%%%%  scan idx 235  |  folio 222  %%%%%%%%%%
\SGAPrintPage{222}
\sgaStableTarget{sga6.e03.expose}\section*{\underline{EXPOSÉ III}\\[3pt] \underline{\underline{RELATIVE FINITENESS CONDITIONS}}}
\begin{center}
\underline{by L. ILLUSIE}
\end{center}

\sgaStableTarget{sga6.e03.sec.1}\section*{1. Relative pseudo-coherence}

\sgaStableTarget{sga6.e03.thm.1-1}\textbf{Theorem 1.1.} \uline{Let} \(n\in\mathbb Z\). \uline{For each locally finite type morphism of schemes} \(f:X\longrightarrow Y\), \uline{there exists one and only one full subcategory} \(D(f)_{n\text{-}\coh}\) \uline{of} \(D(X)\)\SGAAuthorFootnote{1}{\(D(X)\) = derived category of the category of \(\underline{O}_X\)-Modules.} \uline{such that the following conditions are satisfied}:
\begin{enumerate}[label=(\roman*)]
\item \uline{If} \(E\in\operatorname{Ob} D(f)_{n\text{-}\coh}\), \uline{then} \(i^*(E)\in\operatorname{Ob} D(fi)_{n\text{-}\coh}\) \uline{for every open immersion} \(i:U\to X\).

\item \uline{If} \(E\in\ob D(X)\) \uline{and there exists an open covering} \((i_\alpha:U_\alpha\to X)_{\alpha\in A}\) \uline{such that} \(i_\alpha^*(E)\in\ob D(fi_\alpha)_{n\text{-}\coh}\) \uline{for every} \(\alpha\in A\), \uline{then} \(E\in\ob D(f)_{n\text{-}\coh}\).

\item \uline{If}
\[
\begin{tikzcd}
X \arrow[rr,"f"] \arrow[dr,"h"'] & & Y \arrow[dl,"g"]\\
& Z &
\end{tikzcd}
\]
\uline{is a commutative triangle of locally finite type morphisms}, \uline{where} \(f\) \uline{is a closed immersion and} \(g\) \uline{a smooth morphism}, \uline{then we have}
\[
  E\in\operatorname{Ob} D(h)_{n\text{-}\coh}
  \quad\Longleftrightarrow\quad
  f_*(E)\in D(Y)_{n\text{-}\coh}\qquad\mbox{\SGAAuthorFootMark{2}}\ .
\]
\SGAAuthorFootText{2}{\(D(Y)_{n\text{-}\coh}\) = full subcategory of \(D(Y)\) formed by the \(n\)-pseudo-coherent objects (I~2.15).}
\end{enumerate}

\uline{Proof}. \uline{Uniqueness} is evident. Indeed, let \(f:X\to Y\) be a locally finite type morphism. Choose an open covering \((i_\alpha:X_\alpha\to X)_{\alpha\in A}\) such that, for each \(\alpha\in A\), \(fi_\alpha\) factors as
\[
  X_\alpha\xrightarrow{i'_\alpha}Y_\alpha\xrightarrow{f'_\alpha}Y,
\]
where \(i'_\alpha\) is a closed immersion and \(f'_\alpha\) a smooth morphism. It follows at once from (i), (ii), (iii) that the objects of \(D(f)_{n\text{-}\coh}\) are the objects \(E\) of \(D(X)\) such that
\[
  i'_{\alpha *}i_\alpha^*(E)\in \ob D(Y_\alpha)_{n\text{-}\coh};
\]
thus \(D(f)_{n\text{-}\coh}\)

%%%%%%%%%%  scan idx 236  |  folio 223  %%%%%%%%%%
\SGAPrintPage{223}
is indeed determined (since it is a full subcategory of \(D(X)\)). To prove existence, we shall need a few lemmas.

\sgaStableTarget{sga6.e03.lem.1-1-1}\textbf{Lemma 1.1.1.} \uline{Let} \(i:X\to Y\) \uline{be a closed immersion of schemes such that} \(i_*(\underline{O}_X)\) \uline{is pseudo-coherent}, \uline{and let} \(E\in\ob D(X)\). \uline{For} \(E\) \uline{to be \(n\)-pseudo-coherent}, \uline{it is necessary and sufficient that} \(i_*(E)\) \uline{be so}.

\uline{Proof}. The ``necessity'' follows from the general criterion for stability of pseudo-coherence under application of a functor (I~2.13). Conversely, suppose that \(i_*(E)\) is \(n\)-pseudo-coherent, and prove that the same is true of \(E\). Since \(i_*(E)\) has locally bounded-above cohomology, the same is true of \(E\); thus, the question being local, we may suppose that \(E\) is \(m\)-pseudo-coherent for sufficiently large \(m\). Let us show, by descending induction on \(m\), that \(E\) is \(m\)-pseudo-coherent for every \(m\geq n\). We may suppose it already proved that \(E\) is \((n+1)\)-pseudo-coherent. After localizing, we therefore have a distinguished triangle
\[
\begin{tikzcd}
& E'' \arrow[dl,"|" marking] & \\
E' \arrow[rr] & & E \arrow[ul]
\end{tikzcd}
\]
where \(E'\) is perfect and \(E''^i=0\) for \(i>n\). By (I~2.6), we are reduced to showing that \(E''\) is \(n\)-pseudo-coherent. Now, applying \(i_*\) to the preceding triangle, we obtain a distinguished triangle
\[
\begin{tikzcd}
& i_*(E'') \arrow[dl,"|" marking] & \\
i_*(E') \arrow[rr] & & i_*(E) \arrow[ul]
\end{tikzcd}\ .
\]
By the first part of the proof, \(i_*(E')\) is pseudo-coherent, so by (I~2.6), \(i_*(E'')\) is \(n\)-pseudo-coherent. Since \(i_*(E'')\) is zero in degree \(>n\), this means, by (I~2.10 b)), that \var{\(H^n(i_*(E'')\)}{\(H^n(i_*(E''))\)}{EN-VAR-152} is of finite type. Since \(H^n(i_*(E''))\simeq i_*H^n(E'')\), we are reduced to showing, by (I~2.10 b)), that if \(M\) is an \(\underline{O}_X\)-Module such that \(i_*(M)\) is of finite type, then \(M\) is of finite type, which is immediate (and \var{independent of the hypothesis on \(\bar i_*(\underline{O}_X)\)}{independent of the hypothesis on \(i_*(\underline{O}_X)\)}{EN-VAR-153}).

%%%%%%%%%%  scan idx 237  |  folio 224  %%%%%%%%%%
\SGAPrintPage{224}
\sgaStableTarget{sga6.e03.rem.1-1-2}\textbf{Remark 1.1.2.} The preceding lemma applies in particular when \(i\) is a \uline{regular immersion} (VII~1.4). Indeed, by means of \var{the Kozsul complex}{the Koszul complex}{EN-VAR-154}, we have locally a finite left resolution of \(i_*(\underline{O}_X)\) by free \(\underline{O}_Y\)-Modules of finite type, so \(i_*(\underline{O}_X)\) is perfect (I 4.8 and 4.14) (and therefore a fortiori pseudo-coherent).

\sgaStableTarget{sga6.e03.lem.1-1-3}\textbf{Lemma 1.1.3.} \uline{Let there be a Cartesian square of schemes}
\[
\begin{tikzcd}
X \arrow[d,"i"'] & X' \arrow[l,"f"'] \arrow[d,"i'"]\\
Y & Y'{,} \arrow[l,"g"']
\end{tikzcd}
\]
\uline{where} \(i\) \uline{is a closed immersion}, \uline{and let} \(E\) \uline{be a complex of} \(\underline{O}_X\)-\uline{Modules}. \uline{Then} the \uline{canonical base-change arrow} (EGA~III 1.4.15.4)
\[
  g^*i_*(E)\longrightarrow i'_*f^*(E)
\]
\uline{is an isomorphism}.

(ii) \uline{Put} \(E'=f^*(E)\). \uline{If} \(g\) \uline{is flat and} \(i_*(E)\) \(n\)-\uline{pseudo-coherent}, \uline{then} \(i'_*(E')\) \uline{is} \(n\)-\uline{pseudo-coherent}. \uline{Conversely}, \uline{if} \(g\) \uline{is faithfully flat and} \(i'_*(E')\) \(n\)-\uline{pseudo-coherent}, \uline{then} \(i_*(E)\) \uline{is} \(n\)-\uline{pseudo-coherent}.

\uline{Proof}. The assertion is checked trivially on the stalks.

%%%%%%%%%%  scan idx 238  |  folio 225  %%%%%%%%%%
\SGAPrintPage{225}
\sgaStableTarget{sga6.e03.lem.1-1-4}\textbf{Lemma 1.1.4.}
\uline{Let there be a commutative diagram of schemes}
\[
\begin{tikzcd}
X' \arrow[dr,"f'"'] & X \arrow[l,hook',"i'"'] \arrow[r,hook,"i''"] & X'' \arrow[dl,"f''"]\\
& Y{,} &
\end{tikzcd}
\]
\uline{where} \(i'\) (\uline{resp.} \(i''\)) \uline{is a closed immersion}, \uline{and} \(f'\) (\uline{resp.} \(f''\)) \uline{a smooth morphism}, \uline{and let} \(E\in\ob D(X)\). \uline{We then have}:
\[
  i'_*(E)\ n\text{-\uline{pseudo-coherent}}
  \quad\Longleftrightarrow\quad
  \var{i''_*\ n\text{-\uline{pseudo-coherent}}}{i''_*(E)\ n\text{-\uline{pseudo-coherent}}}{EN-VAR-155}.
\]

\uline{Proof}. Following a device due to Lichtenbaum, we send \(X\) diagonally into \(X'\times_YX''\), and are at once reduced to proving the assertion when one of the morphisms \(f'\), \(f''\) is the identity. Changing notation, we thus have a commutative triangle
\begin{equation}\tag{$*$}
\begin{tikzcd}
& Y \arrow[d,"f"]\\
X \arrow[r,hook,"j"'] \arrow[ur,hook,"i"] & Z{,}
\end{tikzcd}
\end{equation}
where \(i\) and \(j\) are closed immersions and \(f\) a smooth morphism, and we must show that ``\(i_*(E)\) \(n\)-pseudo-coherent'' is equivalent to ``\(j_*(E)\) \(n\)-pseudo-coherent''. Since the question is local in \(Y\) and in \(Z\) near a point of \(X\), we may suppose that \(X,Y,Z\) are affine and that we have a factorization
\[
\begin{tikzcd}
Y \arrow[r,hook,"i'"] \arrow[d,"f"'] & Z[T_1,\ldots,T_r]\arrow[dl,"f'"]\\
Z{,}&
\end{tikzcd}
\]
where \(i'\) is a closed immersion, and the oblique arrow the canonical projection. Then (VII 1.10) \(i'\) is a regular immersion, so (1.1.1) an object \(F\) of \(D(Y)\) is \(n\)-pseudo-coherent if and only if \(i'_*(F)\) is so.

%%%%%%%%%%  scan idx 239  |  folio 226  %%%%%%%%%%
\SGAPrintPage{226}
We may therefore suppose that \(Y=Z[T_1,\ldots,T_r]\), with \(f\) the canonical projection. The diagram \((*)\) then corresponds to a diagram of rings
\[
\begin{tikzcd}
& B=C[T_1,\ldots,T_r] \arrow[dl,"u"'] & \\
A & C \arrow[l,"v"] \arrow[u,"\varepsilon"'] & {,}
\end{tikzcd}
\]
where \(u\) and \(v\) are epimorphisms, and \(\varepsilon\) the canonical inclusion. For \(1\leq i\leq r\), put \(u(T_i)=x_i\), and let \(z_i\in C\) satisfy \(v(z_i)=x_i\). The retraction \(\sigma:B\longrightarrow C\) of \(\varepsilon\) defined by \(\sigma(T_i)=z_i\) satisfies \(v\sigma=u\). In geometric language, it corresponds to a section \(s\) of \(f\) such that \(i=sj\). By (VII 1.10), \(s\) is a regular immersion, so the assertion follows from (1.1.1).

\uline{Proof of} (1.1)(\uline{end}). Let \(f:X\to Y\) be a locally finite type morphism. There exists an open covering \((X_\alpha)\) of \(X\) such that each \(X_\alpha\) embeds by a closed immersion into a scheme \(Y_\alpha\) smooth over \(Y\). Putting \(\coprod X_\alpha=X'\), \(\coprod Y_\alpha=Y'\), we thus have a commutative diagram
% SGA display variant EN-VAR-156
\SGACaptureAppState{EN-VAR-156}%
\ifnum\SGAReaderForm=1\relax\SGAAppMark{EN-VAR-156}\fi
\ifnum\SGAReaderForm=0\relax
\begin{equation}\tag{$a'$}
\begin{tikzcd}
X' \arrow[r,"i''"] \arrow[d,"p'"'] & Y' \arrow[ddl,"f'"]\\
X \arrow[d,"f"']\\
Y
\end{tikzcd}
\end{equation}
\else
\begin{equation}\tag{$a'$}
\begin{tikzcd}
X' \arrow[r,"i'"] \arrow[d,"p'"'] & Y' \arrow[ddl,"f'"]\\
X \arrow[d,"f"']\\
Y
\end{tikzcd}
\end{equation}
\fi
where \(p'\) is a morphism that is the sum of a surjective family of open immersions (we shall say \uline{special} for short), \(i'\) a closed immersion, and \(f'\) a smooth morphism. We shall define \(D(f)_{n\text{-}\coh}\) as the full subcategory of \(D(X)\) formed by the objects \(E\) such that \(i'_*p'^*(E)\) is \(n\)-pseudo-coherent. Let us first show that this definition makes sense, i.e. that \(D(f)_{n\text{-}\coh}\) is independent of the choice of \(i'\), \(p'\), \(f'\). Thus, let

%%%%%%%%%%  scan idx 240  |  folio 227  %%%%%%%%%%
\SGAPrintPage{227}
\begin{equation}\tag{$a''$}
\begin{tikzcd}
X'' \arrow[r,"i''"] \arrow[d,"p''"'] & Y'' \arrow[ddl,"f''"]\\
X \arrow[d,"f"']\\
Y
\end{tikzcd}
\end{equation}
be a commutative diagram analogous to \((a')\) (\(p''\) a special morphism, \(i''\) a closed immersion, \(f''\) a smooth morphism). From \((a')\) and \((a'')\), we deduce a commutative diagram
\[
\begin{tikzcd}
X'\times_X X'' \arrow[rr,hook] \arrow[dr] \arrow[dd] & & Y'\times_Y Y'' \arrow[dr] \arrow[dd] & \\
& X' \arrow[rr,hook,"i'"] \arrow[dd] & & Y' \arrow[dd,"f'"]\\
X'' \arrow[rr,hook,"i''"] \arrow[dr] & & Y'' \arrow[dr,"f''"'] & \\
& X \arrow[rr,"f"] & & Y,
\end{tikzcd}
\]
where the horizontal arrows, except for \(f\), are closed immersions, and the arrows of the left (resp. right) Cartesian square are special morphisms (resp. smooth morphisms). We are therefore reduced to proving the following assertion:

(A): \uline{Let}
\[
\begin{tikzcd}
X_1 \arrow[r,hook,"i_1"] \arrow[d,"p"'] & Y_1 \arrow[d,"q"]\\
X_2 \arrow[r,hook,"i_2"] & Y_2
\end{tikzcd}
\]
\uline{be a commutative square of schemes}, \uline{where} \(p\) \uline{is a special morphism}, \(q\) \uline{a smooth morphism}, \(i_1\) \uline{and} \(i_2\) \uline{closed immersions}, \uline{and let} \(E\in\operatorname{ob} D(X_2)\). \uline{Then}, \uline{for} \(i_{2*}(E)\) \uline{to be} \(n\)-\uline{pseudo-coherent}, \uline{it is necessary and sufficient that} \(i_{1*}p^*(E)\) \uline{be so}.

%%%%%%%%%%  scan idx 241  |  folio 228  %%%%%%%%%%
\SGAPrintPage{228}
\uline{Proof of} (A). By definition, \(p\) is the sum of a family of open subsets \(U_\alpha\) of \(X_2\) that cover \(X_2\). Each \(U_\alpha\) is induced by an open subset \(V_\alpha\) of \(Y_2\), and the \(V_\alpha\) cover an open neighborhood \(X'_2\) of \(X_2\) in \(Y_2\). We therefore have a commutative diagram with a Cartesian square
\[
\begin{tikzcd}
X_1 \arrow[r,hook,"i'_1"] \arrow[d,"p"'] & X'_1 \arrow[d,"p'"]\\
X_2 \arrow[r,hook,"i'_2"] \arrow[rr,bend right,"i_2"'] & X'_2 \arrow[r,hook,"j"] & Y_2,
\end{tikzcd}
\]
where \(i'_1\) and \(i'_2\) are closed immersions, \(p\) and \(p'\) special morphisms, and \(j\) an open immersion. The condition ``\(i_{2*}(E)\) \(n\)-pseudo-coherent'' is clearly equivalent to ``\(i'_{2*}(E)\) \(n\)-pseudo-coherent''. On the other hand, it follows at once from (1.1.3) that the condition ``\(i'_{2*}(E)\) \(n\)-pseudo-coherent'' is equivalent to ``\(i'_{1*}p^*(E)\) \(n\)-pseudo-coherent''. Applying (1.1.4) to the factorizations \((jp')i'_1\) and \(qi_1\) of \(i_2p\), and to \(p^*(E)\in\Ob D(X_1)\), we finally conclude that ``\(i_{2*}(E)\) \(n\)-pseudo-coherent'' is equivalent to ``\(i_{1*}p^*(E)\) \(n\)-pseudo-coherent'', q.e.d.

We have thus shown that the category \(D(f)_{n\text{-}\coh}\) defined above by means of the diagram \((a')\) in fact depends only on \(f\). Moreover, it is trivial from the definition that it does satisfy conditions (i) to (iii) of (1.1). This therefore completes the proof of (1.1).

\sgaStableTarget{sga6.e03.def.1-2}\textbf{Definition 1.2.} \uline{Let} \(f:X\to Y\) \uline{be a locally finite type morphism}, \(E\in\Ob D(X)\), \(n\in\mathbb Z\). \uline{We shall say that} \(E\) \uline{is} \(n\)-pseudo-coherent relative to \(f\) \uline{if} \(E\in\Ob D(f)_{n\text{-}\coh}\) (cf. (1.1)). \uline{We shall say that} \(E\) \uline{is} pseudo-coherent (\uline{or equivalently} \((-\infty)\)-pseudo-coherent) relative to \(f\) \uline{if} \(E\) \uline{is} \(n\)-\uline{pseudo-coherent relative to} \(f\) \uline{for every} \(n\in\mathbb Z\). \uline{We shall write} \(D(f)_{\mathrm{coh}}\) \uline{for the full subcategory}

%%%%%%%%%%  scan idx 242  |  folio 229  %%%%%%%%%%
\SGAPrintPage{229}
\uline{of} \(D(X)\) \uline{formed by the objects pseudo-coherent relative to} \(f\). \uline{We shall say that} \(f\) \uline{is} \(n\)-\uline{pseudo-coherent} (\uline{resp.} \uline{pseudo-coherent}) \uline{if} \(\underline{O}_X\) \uline{is} \(n\)-\uline{pseudo-coherent} (\uline{resp.} \uline{pseudo-coherent}) \uline{relative to} \(f\). \uline{We shall sometimes say}, \uline{when this cannot lead to confusion}, ``\(n\)-\uline{pseudo-coherent} (\uline{resp.} \uline{pseudo-coherent}) \uline{relative to} \(Y\)'' \uline{instead of} ``\uline{relative to} \(f\)'', \uline{and shall also write} \(D(X)_{Y\text{-}n\text{-}\coh}\) (\uline{resp.} \(D(X)_{Y\text{-}\coh}\)) \uline{instead of} \(D(f)_{n\text{-}\coh}\) (\uline{resp.} \(D(f)_{\mathrm{coh}}\)). \uline{Finally, for} \(*=+,-,b\), \uline{we shall put}
\[
\begin{gathered}
  D^*(f)_{n\text{-}\coh}=D^*(X)\cap D(f)_{n\text{-}\coh},\\
  D^*(f)_{\mathrm{coh}}=D^*(X)\cap D(f)_{\mathrm{coh}}.
\end{gathered}
\]

\sgaStableTarget{sga6.e03.rem.1-2-1}\textbf{Remarks 1.2.1.}
\begin{enumerate}[label=\alph*)]
\item The categories \(D(f)_{n\text{-}\coh}\), \(D(f)_{\mathrm{coh}}\) are not in general contained in \(D^{-}(X)\). This is so, however, when \(X\) is quasi-compact, since an object of \(D(X)\) with locally bounded-above cohomology then has bounded-above cohomology.

\item The objects of \(D(f)_{\mathrm{coh}}\) have quasi-coherent cohomology. Indeed, let \(E\in\Ob D(f)_{\mathrm{coh}}\); the question being local, we may suppose that \(f\) is a closed immersion. Then \(f_*(E)\) is pseudo-coherent, hence has quasi-coherent cohomology, and the result follows from the fact that an \(\underline{O}_X\)-Module \(M\) is quasi-coherent if and only if \(f_*(M)\) is so.

\item When speaking of an \(\underline{O}_X\)-Module, we shall use the terminology ``of finite presentation rel. to \(f\)'' in preference to ``\((-1)\)-pseudo-coherent rel. to \(f\)''. For an \(\underline{O}_X\)-Module \(E\) to be of finite presentation rel. to \(f\), it is necessary and sufficient, in view of (I 2.9), that there exist locally a closed immersion \(i:X\to X'\), with \(X'\) smooth over \(Y\), such that \(i_*(E)\) is of finite presentation (in the ordinary sense). In particular, saying that \(f\) is

%%%%%%%%%%  scan idx 243  |  folio 230  %%%%%%%%%%
\SGAPrintPage{230}
\((-1)\)-pseudo-coherent is equivalent to saying that \(f\) is locally of finite presentation in the sense of (EGA IV 1.4).
\end{enumerate}

\sgaStableTarget{sga6.e03.ex.1-3}\textbf{Examples 1.3.} A smooth morphism and a regular immersion (1.1.2) are pseudo-coherent morphisms. The same holds for any locally finite type morphism \(f:X\to Y\) when \(Y\) is locally noetherian.

\sgaStableTarget{sga6.e03.prop.1-4}\textbf{Proposition 1.4.} \uline{Let} \(f:X\to Y\) \uline{be a locally finite type morphism}, \uline{and let}
\[
\begin{tikzcd}
& G \arrow[dl] & \\
E \arrow[rr] && F \arrow[ul]
\end{tikzcd}
\]
\uline{be a distinguished triangle in} \(D(X)\). \uline{We then have the following table}, \uline{where} \(n\in\mathbb Z\cup\{-\infty\}\) \uline{indicates a degree of pseudo-coherence relative to} \(f\), \uline{and each row represents an implication whose conclusion is underlined}:
\[
\begin{array}{ccc}
E&F&G\\
n+1&n&\underline n\\
n&\underline n&n\\
\underline n&n&n-1.
\end{array}
\]
\uline{In particular}, \(D(f)_{\mathrm{coh}}\) \uline{is a triangulated subcategory of} \(D(X)\).

\uline{Proof}. By properties (i) to (iii) of (1.1), we may suppose that \(f\) is a closed immersion. We then have a distinguished triangle
\[
\begin{tikzcd}
& f_*(G) \arrow[dl] & \\
f_*(E) \arrow[rr] && f_*(F) \arrow[ul]
\end{tikzcd},
\]
and the assertion follows from (I 2.6).

\sgaStableTarget{sga6.e03.def.1-5}\textbf{Definition 1.5.} \uline{Let} \(X\) \uline{and} \(Y\) \uline{be two} \(S\)-\uline{schemes}. \uline{We shall say that} \(X\) \uline{and} \(Y\) \uline{are tor-independent over} \(S\) \uline{if we have}
\[
  \var{\underline{\Tor}_i^S(\underline{O}_X,O_Y)}{\underline{\Tor}_i^S(\underline{O}_X,\underline{O}_Y)}{EN-VAR-157}=0\quad\text{\uline{for} }i\ne 0,
\]
\uline{where the} \(\underline{\Tor}_i^S\) \uline{are the local hyper-tor defined in} (EGA III 6.5).

%%%%%%%%%%  scan idx 244  |  folio 231  %%%%%%%%%%
\SGAPrintPage{231}
This is a condition of a local nature; when \(X,Y,S\) are affine with respective rings \(B,C,A\), it means that the canonical arrow
\[
  B\otimes_A^{L}C\longrightarrow B\otimes_A C
\]
is an isomorphism. It is therefore satisfied in particular when \(X\) or \(Y\) is flat over \(S\).

\sgaStableTarget{sga6.e03.lem.1-5-1}\textbf{Lemma 1.5.1.} \uline{Let there be Cartesian squares of schemes}
\[
\begin{tikzcd}
X \arrow[d] & X' \arrow[l] \arrow[d] & X'' \arrow[l] \arrow[d]\\
Y & Y' \arrow[l] & Y'' \arrow[l].
\end{tikzcd}
\]
If \(X\) and \(Y'\) are tor-independent over \(Y\), and if \(X'\) and \(Y''\) are tor-independent over \(Y'\), then \(X\) and \(Y''\) are tor-independent over \(Y\).

\uline{Proof}. An easy exercise on associativity of the tensor product.

\sgaStableTarget{sga6.e03.notation.1-6}\textbf{Notation 1.6.} Let \(X\) and \(Y\) be two \(S\)-schemes, \(E\in\operatorname{ob} D^{-}(X)\), \(F\in\operatorname{ob} D^{-}(Y)\). We shall put
\[
  E\otimes_S^{L}F
  =
  L\pr_1^*(E)\otimes^{L}L\pr_2^*(F),
\]
where \(\pr_1,\pr_2\) are the two projections from \(X\times_SY\).

\sgaStableTarget{sga6.e03.lem.1-7}\textbf{Lemma 1.7.} \uline{Let} \(S\) \uline{be a scheme}, \uline{and, for} \(i=1,2\), \(f_i:X_i\longrightarrow Y_i\) \uline{be a closed} \(S\)-\uline{immersion}. \uline{Suppose that} \(X_1\) \uline{and} \(X_2\), \uline{as well as} \(Y_1\) \uline{and} \(Y_2\), \uline{are tor-independent over} \(S\). \uline{Then}, \uline{for} \(E_i\in\operatorname{ob} D^{-}(X_i)_{\mathrm{qcoh}}\)\SGAAuthorFootnote{1}{\(D(X_i)_{\mathrm{qcoh}}\) denotes the full subcategory of \(D(X_i)\) formed by complexes with quasi-coherent cohomology.}, \uline{there exists a canonical isomorphism}
\[
  f_{1*}(E_1)\otimes_S^{L}f_{2*}(E_2)
  \xrightarrow{\sim}
  (f_1\times_Sf_2)_*(E_1\otimes_S^{L}E_2).
\]

\uline{Proof}. Put \(f_1\times_Sf_2=f\). By adjunction, one easily defines a canonical arrow

%%%%%%%%%%  scan idx 245  |  folio 232  %%%%%%%%%%
\SGAPrintPage{232}
\[
  f_{1*}(E_1)\otimes_S^{L}f_{2*}(E_2)
  \longrightarrow
  f_*(E_1\otimes_S^{L}E_2).
\]
Let us show that it is an isomorphism. The question being local, we may suppose that \(S,X_i,Y_i\) are affine with respective rings \(C,A_i,B_i\). By dévissage, we reduce to \(E_i=\underline{O}_{X_i}\). We must then check that
\[
  A_1\otimes_{B_1}^{L}B\otimes_{B_2}^{L}A_2
  \longrightarrow A
\]
is an isomorphism (we have put \(A=A_1\otimes_C A_2\), \(B=B_1\otimes_C B_2\)). But, taking account of the tor-independence hypotheses, this follows at once from associativity of the tensor product.

\sgaStableTarget{sga6.e03.prop.1-8}\textbf{Proposition 1.8.} \uline{Let} \(S\) \uline{be a scheme}, \uline{and}, \uline{for} \(i=1,2\), \(f_i:X_i\longrightarrow Y_i\) \uline{be a locally finite type} \(S\)-\uline{morphism}. \uline{Suppose} \(X_1\) \uline{and} \(X_2\), \uline{as well as} \(Y_1\) \uline{and} \(Y_2\), \uline{tor-independent over} \(S\). \uline{Let} \([a_i,b_i]\) \uline{be an interval of} \(\mathbb Z\), \uline{and} \(E_i\in\operatorname{ob} D^{-}(X_i)_{\mathrm{qcoh}}\) \uline{a complex that is} \(a_i\)-\uline{pseudo-coherent rel.} \uline{to} \(f_i\) \uline{and acyclic in degree} \(>b_i\). \uline{Then} \(E_1\otimes_S^{L}E_2\) \uline{is} \(a\)-\uline{pseudo-coherent rel.} \uline{to} \(f\) \uline{and acyclic in degree} \(>b\), \uline{with} \(f=f_1\times_Sf_2\), \(a=\operatorname{Sup}(a_1+b_2,a_2+b_1)\), \(b=b_1+b_2\).

\uline{Proof}. The question being local on \(X_1\times_SX_2\), we may suppose that \(f_i\) factors as
\[
  X_i\xrightarrow{f'_i}Y'_i\xrightarrow{p_i}Y_i,
\]
where \(f'_i\) is a closed immersion and \(p_i\) a smooth morphism. Putting \(f'=f'_1\times_Sf'_2\), we must show that \(f'_*(E_1\otimes_S^{L}E_2)\) is \(a\)-pseudo-coherent and acyclic in degree \(>b\). By (1.5.1), \(Y'_1\) and \(Y'_2\) are tor-independent over \(S\): after changing notation, we may therefore suppose that \(f_i=f'_i\), in other words that \(f_i\) is a closed immersion. By hypothesis, \(f_{i*}(E_i)\) is \(a_i\)-pseudo-coherent and acyclic in degree \(>b_i\), so (I 2.16) \(f_{1*}(E_1)\otimes_S^{L}f_{2*}(E_2)\) is \(a\)-pseudo-coherent and acyclic in degree \(>b\). We conclude by (1.7).

%%%%%%%%%%  scan idx 246  |  folio 233  %%%%%%%%%%
\SGAPrintPage{233}
\sgaStableTarget{sga6.e03.cor.1-9}\textbf{Corollary 1.9.} \uline{Under the hypotheses of} (1.8), \uline{the functor}
\[
  \otimes_S^{L}:D^{-}(X_1)\times D^{-}(X_2)\longrightarrow D^{-}(X)
\]
\uline{induces a functor}
\[
  \otimes_S^{L}:D^{-}(f_1)_{\mathrm{coh}}\times D^{-}(f_2)_{\mathrm{coh}}
  \longrightarrow D^{-}(f)_{\mathrm{coh}}\qquad(\text{cf. }(1.2)).
\]
\uline{In particular}, \uline{if} \(f_1\) \uline{and} \(f_2\) \uline{are pseudo-coherent} (1.2), \uline{the same holds for} \(f\).

\sgaStableTarget{sga6.e03.cor.1-10}\textbf{Corollary 1.10.} \uline{Let}
\[
\begin{tikzcd}
X \arrow[d,"f"'] & X' \arrow[l,"g"'] \arrow[d,"f'"]\\
Y & Y' \arrow[l,"h"']
\end{tikzcd}
\]
\uline{be a Cartesian square}, \uline{with} \(f\) \uline{locally of finite type}. \uline{Then}:
\begin{enumerate}[label=\alph*)]
\item \uline{Suppose} \(X\) \uline{and} \(Y'\) \uline{tor-independent over} \(Y\) (1.5). \uline{For} \(n\in\mathbb Z\cup\{-\infty\}\), \uline{the functor} \(Lg^*:D^{-}(X)\to D^{-}(X')\) \uline{induces a functor}
\[
  Lg^*:D^{-}(f)_{n\text{-}\coh}\longrightarrow D^{-}(f')_{n\text{-}\coh}.
\]
\uline{In particular}, \uline{if} \(f\) \uline{is} \(n\)-\uline{pseudo-coherent} (1.2), \uline{the same holds for} \(f'\).

\item \uline{Suppose} \(h\) \uline{faithfully flat}, \uline{let} \(n\in\mathbb Z\cup\{-\infty\}\) \uline{and} \(E\in\operatorname{ob} D^{-}(X)\) \uline{be a complex with quasi-coherent cohomology}. \uline{For} \(E\) \uline{to be} \(n\)-\uline{pseudo-coherent rel.} \uline{to} \(f\), \uline{it is necessary and sufficient that} \(g^*(E)\) \uline{be} \(n\)-\uline{pseudo-coherent} \uline{relative to} \(f'\). \uline{In particular}, \uline{for} \(f\) \uline{to be} \(n\)-\uline{pseudo-coherent}, \uline{it is necessary and sufficient that} \(f'\) \uline{be so}.
\end{enumerate}

\uline{Proof}. Taking account of (1.2.1 b)), all the assertions of (1.9) and (1.10) follow trivially from (1.8), except the ``sufficiency'' part of (1.10 b)). For this, the question being local on \(X\), we may suppose that \(f\) factors as a closed immersion followed by a smooth morphism. By (1.1.3), we are then reduced to the case where \(f\) (resp. \(f'\)) is the identity of \(X\) (resp. \(X'\)). The assertion then follows, by a standard dévissage (cf. I 2.16.2), from the fact that a quasi-coherent sheaf \(L\) on \(X\) is of finite type if and only if \(g^*L\) is of finite type.

%%%%%%%%%%  scan idx 247  |  folio 234  %%%%%%%%%%
\SGAPrintPage{234}
\sgaStableTarget{sga6.e03.prop.1-11}\textbf{Proposition 1.11.} \uline{Let}
% SGA display variant EN-VAR-158
\SGACaptureAppState{EN-VAR-158}%
\ifnum\SGAReaderForm=1\relax\SGAAppMark{EN-VAR-158}\fi
\ifnum\SGAReaderForm=0\relax
\[
\begin{tikzcd}
X \arrow[r,"f"] \arrow[dr,"f"'] & Y \arrow[d,"g"]\\
& Z
\end{tikzcd}
\]
\else
\[
\begin{tikzcd}
X \arrow[r,"f"] \arrow[dr,"h"'] & Y \arrow[d,"g"]\\
& Z
\end{tikzcd}
\]
\fi
\uline{be a commutative triangle of locally finite type morphisms}, \uline{let} \(E\in\operatorname{ob} D(X)\), \(F\in\operatorname{ob} D(Y)\), \uline{and} \([a,b]\), \([c,d]\) \uline{be intervals of} \(\mathbb Z\).
\begin{enumerate}[label=\alph*)]
\item \uline{If} \(E\) \uline{is} \(a\)-\uline{pseudo-coherent rel.} \uline{to} \(f\) \uline{and acyclic in degree} \(>b\) \uline{and if} \(F\) \uline{is} \(c\)-\uline{pseudo-coherent rel.} \uline{to} \(g\) \uline{and acyclic in degree} \(>d\), \uline{then} \(E\otimes_Y^{L}F\) (cf. (1.6)) \uline{is} \(m\)-\uline{pseudo-coherent rel.} \uline{to} \(h\) \uline{and acyclic in degree} \(>n\), \uline{where} \(m=\operatorname{Sup}(a+d,b+c)\), \uline{and} \(n=b+d\).

\item \uline{If} \(E\) \uline{is} \(a\)-\uline{pseudo-coherent rel.} \uline{to} \(h\) \uline{and acyclic in degree} \(>b\), \uline{and if} \(g\) \uline{is} \((a+1-b)\)-\uline{pseudo-coherent}, \uline{then} \(E\) \uline{is} \(a\)-\uline{pseudo-coherent rel.} \uline{to} \(f\).
\end{enumerate}

\uline{Proof}. The question being local on \(X\), we may suppose that we have a commutative diagram with a Cartesian square
\[
\begin{tikzcd}[column sep=large,row sep=large]
& & E^{r+s}_Z \arrow[dr] & \\
& E^r_Y \arrow[ur,hook,"j"] \arrow[d,"p"'] & & E^s_Z \arrow[d,"q"]\\
X \arrow[ur,hook,"i"] \arrow[r,"f"] & Y \arrow[urr,hook,"k"'] \arrow[rr,"g"] & & Z{,}
\end{tikzcd}
\]
where \(i,j,k\) are closed immersions, and \(p,q\) the canonical projections of the standard affine spaces. If \(F\) is \(c\)-pseudo-coherent rel. to \(Z\), then \(k_*(F)\) is \(c\)-pseudo-coherent by (1.1 (iii)), and \(j_*(p^*(F))\) is \(c\)-pseudo-coherent by (1.1.3). Similarly, if \(g\) is \(x\)-pseudo-coherent, then \(j_*(\underline{O}_{E^r_Y})\) is \(x\)-pseudo-coherent. By the definition of relative pseudo-coherence (1.2), we are therefore reduced to proving a) and b) when \(f\) and \(g\) are closed immersions.

%%%%%%%%%%  scan idx 248  |  folio 235  %%%%%%%%%%
\SGAPrintPage{235}
\uline{Proof of} a). By (1.1 (iii)), we must show the following: suppose \(f_*(E)\) is \(a\)-pseudo-coherent and acyclic in degree \(>b\), and \(g_*(F)\) is \(c\)-pseudo-coherent and acyclic in degree \(>d\); then \(h_*(E\otimes_Y^{L}F)\) is \(m\)-pseudo-coherent and acyclic in degree \(>n\). Now we have a canonical isomorphism
\[
  f_*(E)\otimes^{L}F\xrightarrow{\sim} f_*(E\otimes_Y^{L}F)
\]
(a special case of \var{the isomorphism (*) in the proof of (1.8)}{the isomorphism in (1.7)}{EN-VAR-159}). Since the functor \(g_*(.\otimes^{L}F)\) takes locally free sheaves of finite type to \(c\)-pseudo-coherent complexes, and complexes acyclic in degree \(>0\) to complexes acyclic in degree \(>d\), the assertion therefore follows from (I 2.13).

\uline{Proof of} b). We must prove the following assertion: suppose \(\var{f_*(E)}{h_*(E)}{EN-VAR-160}\) is \(a\)-pseudo-coherent and acyclic in degree \(>b\), and \(g_*(\underline{O}_Y)\) is \((a+1-b)\)-pseudo-coherent; then \(f_*(E)\) is \(a\)-pseudo-coherent. Since the assertion concerns \(f_*(E)\), we may suppose \(X=Y\) and \(f=\Id_X\). We then proceed as in the proof of (1.1.1), proving by descending induction that \(E\) is pseudo-coherent in degree \(\ge a\). We suppose it already proved that \(E\) is \((a+1)\)-pseudo-coherent, whence (I 2.10), after localizing, a distinguished triangle
\[
\begin{tikzcd}
& E'' \arrow[dl,hook] & \\
E' \arrow[rr] & & E \arrow[ul],
\end{tikzcd}
\]
where \(E'^i\) is locally free for \(i\in[a+1,b]\), \(E'^i=0\) for \(i\notin[a+1,b]\), and \(E''^i=0\) for \(i>a\). The hypothesis that \(g_*(\underline{O}_Y)\) is \((a+1-b)\)-pseudo-coherent then implies (I 2.11 a)) that \(g_*(E')\) is \((a+1)\)-pseudo-coherent. Since \(g_*(E)\) is \(a\)-pseudo-coherent, it follows from (I 2.6) that \(g_*(E'')\) is \(a\)-pseudo-coherent. We conclude by (I 2.10 b)) and the fact that an \(\underline{O}_Y\)-Module \(M\) is of finite type if and only if \(g_*(M)\) is so.

The following corollaries are deduced trivially from (1.11):

%%%%%%%%%%  scan idx 249  |  folio 236  %%%%%%%%%%
\SGAPrintPage{236}
\sgaStableTarget{sga6.e03.cor.1-12}\textbf{Corollary 1.12.} \uline{In the situation of} (1.11), \uline{if} \(g\) \uline{is pseudo-coherent}, \uline{we have}, \uline{for} \(r\in\mathbb Z\cup\{-\infty\}\),
\[
  D^{-}(f)_{r\text{-}\coh}=D^{-}(h)_{r\text{-}\coh};
\]
\uline{in particular} (\uline{taking} \(X=Y\)), \uline{the notions of} \(r\)-\uline{pseudo-coherence over} \(Y\), ``\uline{in the absolute sense}'' \uline{or rel.} \uline{to} \(g\), \uline{coincide}.

\sgaStableTarget{sga6.e03.cor.1-13}\textbf{Corollary 1.13.} \uline{In the situation of} (1.11), \uline{if} \(f\) \uline{is pseudo-coherent}, \uline{the functor} \(\underline{O}_X\otimes_Y^{\mathrm L}\cdot=Lf^*:D^{-}(Y)\longrightarrow D^{-}(X)\) \uline{induces}, \uline{for} \(r\in\mathbb Z\cup\{-\infty\}\), \uline{a functor}
\[
  Lf^*:D^{-}(g)_{r\text{-}\coh}\longrightarrow D^{-}(h)_{r\text{-}\coh}.
\]

\sgaStableTarget{sga6.e03.cor.1-14}\textbf{Corollary 1.14.} \uline{In the situation of} (1.11), \uline{if} \(f\) \uline{is} \(r\)-\uline{pseudo-coherent} \uline{and} \(g\) \uline{is} \(s\)-\uline{pseudo-coherent} (\(r,s\in\mathbb Z\cup\{-\infty\}\)), \uline{then} \(h\) \uline{is} \(\operatorname{Sup}(r,s)\)-\uline{pseudo-coherent}. \uline{If} \(h\) \uline{is} \(r\)-\uline{pseudo-coherent} \uline{and} \(g\) \uline{is} \((r+1)\)-\uline{pseudo-coherent}, \(f\) \uline{is} \(r\)-\uline{pseudo-coherent}.

\sgaStableTarget{sga6.e03.rem.1-14-1}\textbf{Remark 1.14.1.} In particular, for \(r=s=-1\), we recover (cf. (1.2.1 c))) assertions (ii) and (v) of (EGA IV 1.4.3).

\sgaStableTarget{sga6.e03.sec.2}\section*{2. The finiteness theorem.}

\sgaStableTarget{sga6.e03.conj.2-1}\textbf{Conjecture 2.1.} \uline{Let} \(S\) \uline{be a scheme}, \(f:X\to Y\) \uline{be a} proper \(S\)-\uline{morphism of schemes locally of finite type over} \(S\), \(E\in\Ob D^+(X)\). \uline{Then, if} \(E\) \uline{is pseudo-coherent rel.} \uline{to} \(S\) (1.2), \uline{the same holds for} \(Rf_*(E)\).

\sgaStableTarget{sga6.e03.thm.2-2}\textbf{Theorem 2.2.} \uline{The preceding conjecture is true in the following cases}:
\begin{enumerate}[label=(\roman*)]
\item \(S\) \uline{is locally noetherian};
\item \(f\) \uline{is projective}.
\end{enumerate}

%%%%%%%%%%  scan idx 250  |  folio 237  %%%%%%%%%%
\SGAPrintPage{237}
\uline{Proof in case} (i). Since \(S\) is locally noetherian, the projection of \(X\) (resp. \(Y\)) onto \(S\) is a pseudo-coherent morphism (1.3), and the notions of pseudo-coherence over \(X\) (resp. \(Y\)) and relative to \(S\) coincide (1.12). On the other hand, since \(X\) (resp. \(Y\)) is locally noetherian, hence \(\underline{O}_X\) (resp. \(\underline{O}_Y\)) coherent, an object \(M\) of \(D(X)\) (resp. \(D(Y)\)) is pseudo-coherent if and only if (I 3.5) the \(H^i(M)\) are coherent and \(M\) is locally in \(D^{-}\). The theorem therefore follows from the following one:

\sgaStableTarget{sga6.e03.thm.2-2-1}\textbf{Theorem 2.2.1.} \uline{Let} \(f:X\to Y\) \uline{be a proper morphism of locally noetherian schemes}, \uline{and let} \(E\in\Ob D^+(X)\). \uline{If the} \(H^i(E)\) \uline{are coherent}, \uline{the same holds for the} \(R^if_*(E)\). \uline{If, moreover,} \(E\) \uline{is locally in} \(D^{-}\), \uline{the same holds for} \(Rf_*(E)\).

\uline{Proof}. The first assertion follows from the spectral sequence
\begin{equation}\tag{$*$}
  R^pf_*(H^q(E)) \Longrightarrow R^*f_*(E)
\end{equation}
and the finiteness theorem (EGA III 3.2.1). For the second, we may suppose \(Y\) affine; \(X\) is then quasi-compact and separated. Suppose it covered by \(r\) affine open subsets. We know (EGA III 1.4.12) that, under these conditions, \(R^if_*(M)=0\) for every quasi-coherent \(\underline{O}_X\)-Module \(M\) and every \(i>r\). If \(E\) is locally in \(D^{-}\), \(E\) is in \(D^{-}\), since \(X\) is quasi-compact; say \(H^i(E)=0\) for \(i>N\). From the spectral sequence (*) we then deduce that \(R^if_*(E)=0\) for \(i>N+r\), whence the assertion.

\uline{Proof in case} (ii). The question being local on \(Y\), we may suppose \(S\) and \(Y\) affine. We then have a commutative diagram with a Cartesian square
\[
\begin{tikzcd}[column sep=large,row sep=large]
& & P^r_{Y'} \arrow[dr,"p'"] & \\
& P^r_Y \arrow[ur,hook,"j"] \arrow[d,"p"'] & & Y' \arrow[d,"g'"]\\
X \arrow[ur,hook,"i"] \arrow[r,"f"] & Y \arrow[urr,hook,"j'"] \arrow[rr,"g"] & & S{,}
\end{tikzcd}
\]

%%%%%%%%%%  scan idx 251  |  folio 238  %%%%%%%%%%
\SGAPrintPage{238}
where \(i,j,j'\) are closed immersions, \(g'\) a smooth affine morphism, and \(p,p'\) the projections of the standard projective spaces. Suppose \(E\) pseudo-coherent rel. to \(S\). By (1.1 (iii)), this means that \((ji)_*(E)\) is pseudo-coherent. We must show that \(Rf_*(E)\) is pseudo-coherent rel. to \(S\), i.e. (1.1 (iii)) that \(j'_*Rf_*(E)\) is pseudo-coherent. But we have
\[
  j'_*Rf_*(E)\simeq Rp'_*(ji)_*(E),
\]
so the assertion will follow from the following more precise result:

\sgaStableTarget{sga6.e03.thm.2-2-2}\textbf{Theorem 2.2.2.} \uline{Let} \(Y\) \uline{be an affine scheme}, \(f:X=P^r_Y\longrightarrow Y\) \uline{be the projection of a standard projective space}, \(E\in\Ob D^b(X)_{\mathrm{coh}}\), \(a\in\mathbb Z\). \uline{Then} \(Rf_*(E)\) \uline{is pseudo-coherent}. \uline{Moreover}, \uline{if} \(E\) \uline{is acyclic in degree} \(>a\), \(Rf_*(E)\) \uline{is acyclic in degree} \(>a+r\), \uline{and there exists an integer} \(N\) \uline{such that}, \uline{for every} \(n\geq N\), \(Rf_*(E(n))\) (\uline{where} \(E(n)=E\otimes\underline{O}_X(n)\)) \uline{is acyclic in degree} \(>a\).

\uline{Proof}. Since \(Y\) is affine, \(X\) is quasi-compact and separated, so (II 2.2.9 a)) \(E\) is isomorphic, in \(D(X)\), to a complex \(F\) bounded above in degree such that, for every \(i\), \(F^i\in\Ob \underline{L}(X)\), i.e. (loc. cit.) is a finite direct sum of sheaves of the type \(\underline{O}_X(-n)\) \((n>0)\). Moreover, if \(E\) is acyclic in degree \(>a\), we may choose \(F\) so that \(F^i=0\) for \(i>a\): this follows formally from the preceding assertion, taking account of (I 2.10 a) and II 2.2.8 a)). Replacing \(E\) by \(F\), we may therefore suppose that \(E^i\in\Ob \underline{L}(X)\) for every \(i\), and \(E^i=0\) for \(i>a\). Since \(X\) is covered by \((r+1)\) affine open subsets, we have, by (EGA III 1.4.10) and the hypercohomology spectral sequence, \(R^qf_*(E)=0\) for \var{\(i>a+r\)}{\(q>a+r\)}{EN-VAR-161}. Let \(t\) be an integer \(\leq a\), and put \(E'=(0\to E^t\to\ldots\to E^a\to0)\). We have an exact sequence
\[
  0\longrightarrow E'\longrightarrow E\longrightarrow E''\longrightarrow0,
\]

%%%%%%%%%%  scan idx 252  |  folio 239  %%%%%%%%%%
\SGAPrintPage{239}
where \(E''\) satisfies the same conditions as \(E\), but with \(a\) replaced by \(t-1\). Applying \(Rf_*\), we obtain a distinguished triangle
\begin{equation}\tag{$*$}
\begin{tikzcd}
& Rf_*(E'') \arrow[dl,hook] & \\
Rf_*(E') \arrow[rr] & & Rf_*(E) \arrow[ul] & .
\end{tikzcd}
\end{equation}
Since \(t\) can be chosen arbitrarily far in the negative direction, it suffices, for the first assertion, to show that \(Rf_*(E)\) is \((t+r)\)-pseudo-coherent. Now \(Rf_*(E'')\) is acyclic in degree \(\geq t+r\), hence \((t+r)\)-pseudo-coherent. Thus, by (I 2.6), the result will follow once we prove that \(Rf_*(E')\) is perfect. For this, we reduce by dévissage (I 4.10) to the case where \(E'\) is concentrated in degree \(0\) and equal to \(\underline{O}_X(-n)\), \(n>0\). Then, by (EGA III 2.1.16), \(R^rf_*(\underline{O}_X(-n))\) is locally free of finite type, and we have
\[
  Rf_*(\underline{O}_X(-n))\simeq R^rf_*(\underline{O}_X(-n))[-r].
\]
Thus \(Rf_*(\underline{O}_X(-n))\) is perfect, and we deduce the first assertion of (2.2.2). We have already seen that \(Rf_*(E)\) is acyclic in degree \(>a+r\). It remains to prove the last assertion. For \(t<-r+a\), \(Rf_*(E'')\) is acyclic in degree \(>a\), so, taking account of the triangle \((*)\), we are reduced to proving the result for \(E'\) in place of \(E\). Since \(E'\) is bounded in degree and its components are objects of \(\underline{L}(X)\), there exists an integer \(N\) such that, for \(n\geq N\), the components of \(E'(n)\) are finite direct sums of sheaves of the type \(\underline{O}_X(p)\), with \(p\geq0\). Then, by (EGA III 2.1.16) and the hypercohomology spectral sequence, \(Rf_*(E'(n))\) is acyclic in degree \(>a\) for every \(n\geq N\), which completes the proof of (2.2.2), and hence of (2.2).

We have a ``classical'' complement to (2.2 (ii)), which we shall give for later reference:

%%%%%%%%%%  scan idx 253  |  folio 240  %%%%%%%%%%
\SGAPrintPage{240}
\sgaStableTarget{sga6.e03.cor.2-3}\textbf{Corollary 2.3} (EGA III 2.2.1). \uline{Let} \(S\) \uline{be a quasi-compact scheme}, \(f:X\to Y\) \uline{be a projective} \(S\)-\uline{morphism of schemes of finite type over} \(S\), \(L=i^*\underline{O}_P(1)\) \uline{be the invertible sheaf defined by an embedding} \(i:X\to P=P(F)\), \uline{where} \(F\) \uline{is} \uline{a quasi-coherent} \(\underline{O}_Y\)-\uline{Module of finite type}. \uline{For} \(M\in\Ob D(X)\) \uline{and} \(n\in\mathbb Z\), \uline{put} \(M(n)=M\otimes^{\mathrm L} L^{\otimes n}\). \uline{Then}, \uline{if} \(E\in\Ob D^+(X)\) \uline{is pseudo-coherent rel.} \uline{to} \(S\) \uline{and} \uline{acyclic in degree} \(>a\) \((a\in\mathbb Z)\), \uline{there exists an integer} \(N\) \uline{such that}, \uline{for every} \(n\geq N\), \(Rf_*E(n)\) \uline{satisfies the same conditions}.

\uline{Proof}. We reduce trivially, as in the proof of (2.2 (ii)), to the situation of (2.2.2): \(Y=S\), \(f:X=P^r_Y\to Y\) the canonical projection, \(L=\underline{O}_X(1)\) the canonical sheaf, \(E\in\Ob D^b(X)\) a pseudo-coherent complex acyclic in degree \(>a\), and the result follows by (2.2.2).

\sgaStableTarget{sga6.e03.cor.2-4}\textbf{Corollary 2.4.} \uline{Let} \(S\) \uline{be a quasi-compact scheme}, \uline{and} \(f:X\to Y\) \uline{be a proper} \(S\)-\uline{morphism of schemes of finite type over} \(S\). \uline{If} \(S\) \uline{is noetherian or} \(f\) \uline{is projective}, \uline{the functor} \(Rf_*:D^+(X)\longrightarrow D^+(Y)\) \uline{induces a functor}
\[
  Rf_*:D^b(X)_{S\text{-}\coh}\longrightarrow D^b(Y)_{S\text{-}\coh}
\]

\uline{Proof}. A consequence of (2.2) and the fact that there is then (EGA III 1.4.12 and IV 1.7.21) an integer \(N\) such that \(R^if_*(E)=0\) for every \(i>N\) and every quasi-coherent \(\underline{O}_X\)-Module \(E\).

\sgaStableTarget{sga6.e03.cor.2-5}\textbf{Corollary 2.5.} \uline{Let} \(f:X\to Y\) \uline{be a proper and pseudo-coherent morphism} (1.2), \uline{and} \(E\in\Ob D^+(X)\). \uline{Suppose} \(Y\) \uline{locally noetherian or} \(f\) \uline{projective}. \uline{Then}, \uline{if} \(E\) \uline{is pseudo-coherent} (\uline{in the absolute sense}), \uline{the same holds for} \(Rf_*(E)\). \uline{If} \(Y\) \uline{is quasi-compact}, \uline{the functor} \(Rf_*\) \uline{induces a functor}
\[
  Rf_*:D^b(X)_{\mathrm{coh}}\longrightarrow D^b(Y)_{\mathrm{coh}}.
\]

%%%%%%%%%%  scan idx 254  |  folio 241  %%%%%%%%%%
\SGAPrintPage{241}%
\uline{Proof}. A consequence of \var{(2.2), (2.5) and (1.12)}{(2.2), (2.4) and (1.12)}{EN-VAR-162}.

\sgaStableTarget{sga6.e03.schol.2-6}\textbf{Scholium 2.6.} To prove conjecture (2.1) in the general case, it will no doubt be necessary to introduce new techniques. One may hope that these will one day help to prove the analogous conjecture in Analytic Geometry, for a proper morphism of analytic spaces over a locally ringed topos in suitably topologized \(\mathbb C\)-algebras, a conjecture that is at present ``proved'' only over a complex analytic space (Grauert's theorem [1]).\SGAAuthorFootnote{1}{(added in April 1971) see an article in preparation by R. Kiehl and J.L. Verdier, in which finiteness theorem 1 is proved by different, simpler techniques and under more general hypotheses.}

\sgaStableTarget{sga6.e03.sec.3}\section*{3. Relative finite Tor-dimension}

\sgaStableTarget{sga6.e03.def.3-1}\textbf{Definition 3.1.} \uline{Let} \(f:(X,\underline{O}_X)\longrightarrow(Y,\underline{O}_Y)\) \uline{be a morphism of ringed topoi}, \(E\in\Ob D^{-}(X)\), \([m,n]\) \uline{an interval of} \(\mathbb Z\). \uline{We say that} \(E\) \uline{is} of flat amplitude relative to \(f\) (or \(Y\), when no confusion is possible) contained in \([m,n]\), \uline{and write}
\[
  \toramp_f(E)\subset[m,n]\qquad(\text{or }\toramp_Y(E)\subset[m,n])
\]
\uline{if} \(E\), \uline{considered as a complex of} \(f^{-1}(\underline{O}_Y)\)-\uline{Modules}, \uline{is of flat amplitude contained} \uline{in} \([m,n]\) (I 5.2). \uline{We say that} \(f\) \uline{is of flat amplitude contained} \uline{in} \([-a,0]\) \((a\geq0)\), \uline{or of} tor-dimension \(\leq a\), \uline{if} \(\underline{O}_X\) \uline{is of flat amplitude} \uline{rel.} \uline{to} \(f\), \uline{contained in} \([-a,0]\). \uline{We say that} \(E\) \uline{is} of finite flat amplitude rel. to \(f\) \uline{if there is an interval} \([m,n]\) \uline{such that} \(\toramp_f(E)\subset[m,n]\).

\sgaStableTarget{sga6.e03.notation.3-2}\textbf{Notation 3.2.} We denote by \(D(f)_{\mathrm{torf}}\), or \(D(X)_{f\text{-}\torf}\), or \(D(X)_{Y\text{-}\torf}\) when no confusion is possible, the full subcategory of \(D^{-}(X)\) formed by the complexes of finite flat amplitude relative to \(f\). By (I 5.3), this is a triangulated subcategory of \(D^b(X)\).
%%%%%%%%%%  scan idx 255  |  folio 242  %%%%%%%%%%
\SGAPrintPage{242}%
\sgaStableTarget{sga6.e03.prop.3-3}\textbf{Proposition 3.3.} \uline{With the data of} (3.1), \uline{suppose that} \(X\) \uline{has enough points} \var{(SGA 4 XVII IV 6.4)}{(SGA 4 IV 6.4)}{EN-VAR-163}. \uline{The following conditions are then equivalent}:
\begin{enumerate}[label=(\roman*)]
\item \(\toramp_f(E)\subset[m,n]\);
\item \uline{for every point} \(x\) \uline{of} \(X\), \uline{we have}
\[
  \toramp_y(E_x)\subset[m,n],
\]
\uline{where} \(y\) \uline{is the image point of} \(x\) \uline{under} \(f\), \uline{ringed by} \(\underline{O}_{Y,y}\);
\item \uline{for every object} \(V\) \uline{of} \(Y\) \uline{and every} \(\underline{O}_V\)-\uline{Module} \(M\), \uline{we have}
\[
  H^i\left((E|f^{-1}(V))\otimes_V^{\mathrm L}M\right)=0
  \quad\text{\uline{for} }i\notin[m,n],
\]
\uline{where} \((E|f^{-1}(V))\otimes_V^{\mathrm L}M=(E|f^{-1}(V))\otimes^{\mathrm L}\mathrm Lf^*(M)\);
\item \uline{the same condition as} (iii), \uline{but with} \(M\) \uline{the cokernel of a homomorphism} \(\underline{O}_V^q\longrightarrow\underline{O}_V^p\), \uline{where} \(p,q\in\mathbb N\).
\end{enumerate}

\uline{Proof}. The equivalence (i)\(\Longleftrightarrow\)(ii) follows from (I 5.6.2), and the implications (ii)\(\Longrightarrow\)(iii)\(\Longrightarrow\)(iv) are trivial. Let us prove (iv)\(\Longrightarrow\)(ii). Let \(x\) be a point of \(X\), \(y=f(x)\); we must show that
\begin{equation}\tag{$*$}
  H^i(E_x\otimes_{\underline{O}_{Y,y}}^{\mathrm L}M)=0
\end{equation}
for every \(\underline{O}_{Y,y}\)-module \(M\) and every \(i\notin[m,n]\). Since every \(\underline{O}_{Y,y}\)-module is a filtered inductive limit of modules of finite presentation, we may restrict ourselves to checking the preceding relation for \(M\) of finite presentation. Now, when \(M\) is an \(\underline{O}_{Y,y}\)-module of finite presentation, the same argument as for ordinary ringed spaces (EGA \(0_{\mathrm I}\) 5.3.9) shows that there is a \(y\)-pointed object \(V\) of \(Y\) and an \(\underline{O}_V\)-Module \(P\), the cokernel of a homomorphism \(\underline{O}_V^q\longrightarrow\underline{O}_V^p\) \((p,q\in\mathbb N)\), such that \(P_y\simeq M\). We then have
\[
  \bigl((E|f^{-1}(V))\otimes_V^{\mathrm L}P\bigr)_x
  \simeq
  E_x\otimes_{\underline{O}_{Y,y}}^{\mathrm L}M,
\]
%%%%%%%%%%  scan idx 256  |  folio 243  %%%%%%%%%%
\SGAPrintPage{243}%
by the commutation of tensor product with passage to stalks, and hypothesis (iv) implies (*), q.e.d.

\sgaStableTarget{sga6.e03.rem.3-3-1}\textbf{Remark 3.3.1.} The writer does not know whether the equivalences (i)\(\Longleftrightarrow\)(iii)\(\Longleftrightarrow\)(iv) hold without the hypothesis that \(X\) has enough points.

\sgaStableTarget{sga6.e03.cor.3-4}\textbf{Corollary 3.4.} \uline{Let}
\[
\begin{tikzcd}[column sep=large]
X \arrow[rr,"f"] \arrow[dr,"h"'] & & Y \arrow[dl,"g"]\\
& Z &
\end{tikzcd}
\]
\uline{be a commutative triangle of ringed topoi}, \uline{where we suppose that} \(X\) \uline{has enough points}, \uline{and let} \(E\in\Ob D(X)\), \(F\in\Ob D(Y)\), \([a,b]\), \([c,d]\) \uline{be intervals of} \(\mathbb Z\). \uline{Then we have}:
\[
  \toramp_f(E)\subset[a,b]\quad\text{and}\quad\toramp_g(F)\subset[c,d]
\]
\[
  \Longrightarrow\toramp_h(E\otimes_Y^{\mathrm L}F)\subset[a+c,b+d].
\]

\uline{Proof}. By (3.3), we must show that the relations on the left imply
\[
  H^i\left(\bigl((E\otimes_Y^{\mathrm L}F)|h^{-1}(Z')\bigr)
  \otimes_{Z'}^{\mathrm L}M\right)=0
\]
for every \(i\notin[a+c,b+d]\) and every \(\underline{O}_{Z'}\)-Module \(M\), where \(Z'\) is an object of \(Z\). Now, putting \(Y'=g^{-1}(Z')\), \(X'=f^{-1}(Y')=h^{-1}(Z')\), we have
\[
  \bigl((E\otimes_Y^{\mathrm L}F)|X'\bigr)
  \otimes_{Z'}^{\mathrm L}M
  \simeq
  (E|X')\otimes_{Y'}^{\mathrm L}
  \bigl((F|Y')\otimes_{Z'}^{\mathrm L}M\bigr),
\]
and the assertion follows at once from the hypotheses.

In particular:

\sgaStableTarget{sga6.e03.cor.3-4-1}\textbf{Corollary 3.4.1.} \uline{Under the hypotheses of} (3.4),
\begin{enumerate}[label=\alph*)]
\item \uline{if} \(g\) \uline{is of tor-dimension} \(\leq n\), \uline{we have}
\[
  \toramp_f(E)\subset[a,b]\Longrightarrow\toramp_h(E)\subset[a-n,b];
\]
%%%%%%%%%%  scan idx 257  |  folio 244  %%%%%%%%%%
\SGAPrintPage{244}%
\item \uline{if} \(f\) \uline{is of tor-dimension} \(\leq m\), \uline{we have}
\[
  \toramp_g(F)\subset[c,d]\Longrightarrow\toramp_h(\mathrm Lf^*(F))\subset[c-m,d].
\]
\end{enumerate}

\sgaStableTarget{sga6.e03.rem.3-4-2}\textbf{Remark 3.4.2.} Let \(f:X\to Y\) be a morphism of finite tor-dimension and \(E\in\Ob D(X)\). Then, by (3.4.1 a)), when \(E\) is of finite flat amplitude (in the absolute sense), \(E\) is of finite flat amplitude rel. to \(f\). But the converse is false, even when \(f\) is flat, as the example of a singularity of a scheme over a field shows.

\sgaStableTarget{sga6.e03.prop.3-5}\textbf{Proposition 3.5.} \uline{Let} \(S\) \uline{be a scheme}, \uline{and}, \uline{for} \(i=1,2\), \(f_i:X_i\longrightarrow Y_i\) \uline{an} \(S\)-\uline{morphism}, \(E_i\in\Ob D^{-}(X_i)\), \([a_i,b_i]\) \uline{an interval of} \(\mathbb Z\). \uline{Suppose} \(X_1\) \uline{and} \(X_2\), \uline{as well as} \(Y_1\) \uline{and} \(Y_2\), \uline{are tor-independent over} \(S\) (1.5). \uline{Put} \(a=a_1+a_2\), \(b=b_1+b_2\), \(f=f_1\times_Sf_2:X\to Y\). \uline{Then}:
\begin{enumerate}[label=\alph*)]
\item \uline{If we have} \(\toramp_{f_1}(E_1)\subset[a_1,b_1]\) \uline{and if} \(E_2\) \uline{is acyclic outside} \([a_2,b_2]\), \(E_1\otimes_S^{\mathrm L}E_2\) (\uline{cf.} (1.6)) \uline{is acyclic outside} \([a,b]\).
\item \uline{If we have} \(\toramp_{f_i}(E_i)\subset[a_i,b_i]\) \uline{for} \(i=1,2\), \uline{we have}
\[
  \toramp_f(E_1\otimes_S^{\mathrm L}E_2)\subset[a,b].
\]
\end{enumerate}

\uline{Proof}. By (3.3), it suffices to make the verification ``on the stalks''. Let \(s\) be a point of \(S\), \(x\) a point of \(X\) over \(s\), \(y=f(x)\), \(x_i\) (resp. \(y_i\)) the canonical projection of \(x\) (resp. \(y\)) onto \(X_i\) (resp. \(Y_i\)). Put
\[
  (E_i)_{x_i}=M_i,\quad \underline{O}_{X_i,x_i}=A_i,\quad
  \underline{O}_{Y_i,y_i}=B_i,\quad \underline{O}_{S,s}=C,
\]
\[
  \underline{O}_{X,x}=A=A_1\otimes_C A_2,\qquad
  \underline{O}_{Y,y}=B=B_1\otimes_C B_2.
\]
We have
\begin{equation}\tag{$*$}
  (E_1\otimes_S^{\mathrm L}E_2)_x
  \simeq
  M_1\otimes_{A_1}^{\mathrm L}A\otimes_{A_2}^{\mathrm L}M_2.
\end{equation}
The tor-independence hypotheses imply, on the other hand:
\[
  A\simeq A_1\otimes_{B_1}^{\mathrm L}B\otimes_{B_2}^{\mathrm L}A_2.
\]
Consequently, the right-hand side of (*) is canonically isomorphic, in the
%%%%%%%%%%  scan idx 258  |  folio 245  %%%%%%%%%%
\SGAPrintPage{245}%
derived category \(D(B)\) of \(B\)-modules, to \(M_1\otimes_{B_1}^{\mathrm L}B\otimes_{B_2}^{\mathrm L}M_2\). The relations \(\toramp_{f_i}(E_i)\subset[a_i,b_i]\) \((i=1,2)\) imply \(\toramp_{B_i}(M_i)\subset[a_i,b_i]\), hence (I 5.6) \(\toramp_B(M_1\otimes_{B_1}^{\mathrm L}B\otimes_{B_2}^{\mathrm L}M_2)\subset[a,b]\), giving assertion b). On the other hand, when we have \(\toramp_{f_1}(E_1)\subset[a_1,b_1]\), \(M_1\) is isomorphic to a complex of flat \(B_1\)-modules, zero outside \([a_1,b_1]\), and consequently when \(E_2\) (hence \(M_2\)) is acyclic outside \([a_2,b_2]\), \(M_1\otimes_{B_1}^{\mathrm L}B\otimes_{B_2}^{\mathrm L}M_2\) is acyclic outside \([a,b]\), giving assertion a), which completes the proof.

\sgaStableTarget{sga6.e03.cor.3-5-1}\textbf{Corollary 3.5.1.} \uline{Under the hypotheses of} (3.5), \uline{if} \(f_1\) \uline{and} \(f_2\) \uline{are of finite tor-dimension} (3.1), \uline{the same holds for} \(f\).

\sgaStableTarget{sga6.e03.cor.3-5-2}\textbf{Corollary 3.5.2.} \uline{Let}
\[
\begin{tikzcd}
X \arrow[d,"f"'] & X' \arrow[l,"g"'] \arrow[d,"f'"]\\
Y & Y' \arrow[l,"h"']
\end{tikzcd}
\]
\uline{be a Cartesian square of schemes}, \uline{where} \(X\) \uline{and} \(Y'\) \uline{are supposed tor-independent over} \(Y\), \uline{and let} \(E\in\Ob D(X)\). \uline{Then we have}
\[
  \toramp_f(E)\subset[a,b]\Longrightarrow
  \toramp_{f'}(\mathrm Lg^*(E))\subset[a,b],
\]
\uline{and consequently the functor} \(\mathrm Lg^*:D^{-}(X)\to D^{-}(X')\) \uline{induces a functor}
\[
  \mathrm Lg^*:D(f)_{\mathrm{torf}}\longrightarrow D(f')_{\mathrm{torf}} .
\]


% Original English continuation definitions, retained in their original order.
\providecommand{\Ob}{\operatorname{Ob}}
\providecommand{\D}{\mathrm D}
\providecommand{\Ho}{\mathrm H}
\providecommand{\RHom}{\mathrm R\!\operatorname{Hom}}
\providecommand{\parfamp}{\operatorname{parf.amp}}
\providecommand{\tamp}{\operatorname{tor.amp}}
\providecommand{\Qcoh}{\operatorname{Qcoh}}
\providecommand{\Mod}{\operatorname{Mod}}
\providecommand{\Coh}{\operatorname{Coh}}
\providecommand{\Lib}{\operatorname{Lib}}
\providecommand{\Loclib}{\operatorname{Loclib}}
\providecommand{\qcoh}{\mathrm{qcoh}}
\providecommand{\coh}{\mathrm{coh}}
\providecommand{\parf}{\mathrm{parf}}
\providecommand{\torf}{\mathrm{torf}}
\providecommand{\id}{\operatorname{id}}
\providecommand{\Coker}{\operatorname{Coker}}
\providecommand{\rank}{\operatorname{rank}}
\providecommand{\rg}{\operatorname{rg}}
\providecommand{\Hom}{\operatorname{Hom}}
\providecommand{\Ext}{\operatorname{Ext}}
\providecommand{\Spec}{\operatorname{Spec}}
\providecommand{\Cone}{\operatorname{Cone}}
\providecommand{\pt}{\operatorname{Pt}}
\providecommand{\ind}{\operatorname{ind}}
\hbadness=10000
\hfuzz=20pt
\emergencystretch=3em

% Reviewed French continuation definitions, retained verbatim.
% Additions for Exposé III, pages 253--280
\renewcommand{\Ob}{\operatorname{ob}}
\providecommand{\D}{\mathrm D}
\providecommand{\Ho}{\mathrm H}
\providecommand{\RHom}{\mathrm R\!\operatorname{Hom}}
\providecommand{\parfamp}{\operatorname{parf.amp}}
\providecommand{\tamp}{\operatorname{tor.amp}}
\providecommand{\Qcoh}{\operatorname{Qcoh}}
\providecommand{\Mod}{\operatorname{Mod}}
\providecommand{\Coh}{\operatorname{Coh}}
\providecommand{\Lib}{\operatorname{Lib}}
\providecommand{\Loclib}{\operatorname{Loclib}}
\providecommand{\qcoh}{\mathrm{qcoh}}
\providecommand{\coh}{\mathrm{coh}}
\providecommand{\parf}{\mathrm{parf}}
\providecommand{\torf}{\mathrm{torf}}
\renewcommand{\id}{\operatorname{id}}
\providecommand{\uparf}{\underline{\mathrm{parf}}\rule[-1ex]{0pt}{0pt}}
\providecommand{\Coker}{\operatorname{Coker}}
\providecommand{\rank}{\operatorname{rank}}
\providecommand{\rg}{\operatorname{rg}}
\providecommand{\Hom}{\operatorname{Hom}}
\providecommand{\Ext}{\operatorname{Ext}}
\providecommand{\Spec}{\operatorname{Spec}}
\providecommand{\Cone}{\operatorname{Cone}}
\providecommand{\pt}{\operatorname{Pt}}
\providecommand{\ind}{\operatorname{ind}}
\hbadness=10000
\hfuzz=20pt
\emergencystretch=3em
%%%%%%%%%%  scan idx 259  |  folio 246  %%%%%%%%%%
\SGAPrintPage{246}%
\sgaStableTarget{sga6.e03.prop.3-6}\textbf{Proposition 3.6.} \uline{Let}
\[
\begin{tikzcd}
X \arrow[r,hook,"i"] \arrow[d,"f"'] & X' \arrow[dl,"f'"]\\
Y &
\end{tikzcd}
\]
\uline{be a commutative triangle of schemes}, \uline{where} \(i\) \uline{is a closed immersion and} \(f'\) \uline{a smooth morphism}, \uline{and let} \(E\in\Ob D(X)\). \uline{For} \(E\) \uline{to be locally of finite flat amplitude rel.} \uline{to} \(f\), \uline{it is necessary and sufficient that} \(i_*(E)\) \uline{be locally of finite flat amplitude} (\uline{in the absolute sense}). \uline{More precisely}, \uline{if} \(f'\) \uline{is of relative dimension} \(\leq d\), \uline{we have}
\begin{align*}
\text{(i)}\quad &\tamp(i_*(E))\subset[a,b] \Longrightarrow \tamp_f(E)\subset[a,b],\\
\text{(ii)}\quad &\tamp_f(E)\subset[a,b] \Longrightarrow \tamp(i_*(E))\subset[a-d,b].
\end{align*}
%%%%%%%%%%  scan idx 260  |  folio 247  %%%%%%%%%%
\SGAPrintPage{247}%
\uline{Proof}. For an interval \([m,n]\) of \(\mathbb Z\), we evidently have
\[
 \tamp_f(E)\subset[m,n]
 \Longleftrightarrow
 \tamp_{f'}(i_*(E))\subset[m,n],
\]
so that we may suppose \(X=X'\), \(i=\Id_X\), \(f=f'\). Relation (i) follows trivially from (3.4.1 a)). Let us prove (ii). We have a Cartesian square
\[
\begin{tikzcd}
X \arrow[d,"f"'] & X\times_Y X \arrow[l,"\operatorname{pr}_2"'] \arrow[d,"\operatorname{pr}_1"]\\
Y & X \arrow[l,"f"']
\end{tikzcd},
\]
and the diagonal morphism \(s:X\longrightarrow X\times_YX\) is a section of \(\operatorname{pr}_1\) and \(\operatorname{pr}_2\). Suppose \(f\) is smooth of relative dimension \(d\). Then (VII 1.10) \(s\) is a regular immersion of codimension \(\leq d\), hence is of tor-dimension \(\leq d\) (as is seen using the Koszul complex). By (3.5.2), we have
\[
 \tamp_f(E)\subset[a,b]
 \Longrightarrow
 \tamp_{\operatorname{pr}_1}(\operatorname{pr}_2^*(E))\subset[a,b],
\]
and by (3.4.1 b)) we have
\[
 \tamp_{\operatorname{pr}_1}(\operatorname{pr}_2^*(E))\subset[a,b]
 \Longrightarrow
 \tamp(\mathrm Ls^*\operatorname{pr}_2^*(E))\subset[a-d,b],
\]
giving implication (ii), since \(\mathrm Ls^*\operatorname{pr}_2^*\simeq\Id\).

\sgaStableTarget{sga6.e03.rem.3-6-1}\textbf{Remark 3.6.1.} Statement (3.6) would be false with only the assumption that \(f'\) is flat; see (3.4.2).

\sgaStableTarget{sga6.e03.prop.3-7}\textbf{Proposition 3.7.} (``\uline{Projection formula}'')\SGAAuthorFootnote{1}{cf. (H II 5.6).}. \uline{Let} \(S\) \uline{be a quasi-compact scheme}, \(f:X\to Y\) \uline{a quasi-compact and quasi-separated} \(S\)-\uline{morphism of quasi-compact schemes over} \(S\), \(E\in\Ob D^b(X)_{\qcoh}\), \(M\in\Ob D^b(S)_{\qcoh}\)\SGAAuthorFootnote{2}{The subscript ``qcoh'' indicates the full subcategory formed by complexes with quasi-coherent cohomology.}. \uline{If} \(M\) \uline{is} \uline{of finite flat amplitude}, \uline{or if} \(E\) \uline{is of finite flat amplitude rel.} \uline{to} \(S\), \uline{there}
%%%%%%%%%%  scan idx 261  |  folio 248  %%%%%%%%%%
\SGAPrintPage{248}%
\uline{is a canonical functorial isomorphism in} \(D(Y)_{\qcoh}\)
\[
 Rf_*(E)\otimes_S^{\mathrm L}M
 \xrightarrow{\sim}
 Rf_*(E\otimes_S^{\mathrm L}M)\qquad(\text{cf. }(1.6)).
\]

\uline{Proof}. First note that there is (EGA III 1.4.12 and IV 1.7.21) an integer \(N\) such that \(R^if_*(F)=0\) for every \(i>N\) and every quasi-coherent \(\underline{O}_X\)-Module \(F\). Consequently the functor \(Rf_*\) sends \(D^b(X)_{\qcoh}\) into \(D^b(Y)\) (and even \(D^b(Y)_{\qcoh}\) (EGA III 1.4.10 and IV 1.7.21)). \var{Suppose \(E\) (resp. \(M\)) is of finite flat amplitude rel. to \(f\) (resp. finite flat amplitude).}{Suppose \(E\) (resp. \(M\)) is of finite flat amplitude rel. to \(S\) (resp. finite flat amplitude).}{EN-VAR-164} The adjunction arrow \(Lf^*Rf_*(E)\longrightarrow E\) defines, by tensoring with \(M\), a canonical arrow of \(D^-(X)\)
\[
 Lf^*(Rf_*(E)\otimes_S^{\mathrm L}M)
 \longrightarrow E\otimes_S^{\mathrm L}M,
\]
where \(E\otimes_S^{\mathrm L}M\in\Ob D^b(X)\). By adjunction we deduce a canonical arrow of \(D^-(Y)\)
\sgaStableTarget{sga6.e03.eq.unnamed}\[
 Rf_*(E)\otimes_S^{\mathrm L}M
 \longrightarrow Rf_*(E\otimes_S^{\mathrm L}M).\tag{*}
\]
Its target and source moreover have quasi-coherent cohomology, since \(Rf_*\) sends \(D^b(X)_{\qcoh}\) into \(D^b(Y)_{\qcoh}\) (EGA III 1.4.10 and IV 1.7.21), and \(\otimes_S^{\mathrm L}\) sends \(D^-(X)_{\qcoh}\times D^-(S)_{\qcoh}\) into \(D^-(X)_{\qcoh}\) (likewise with \(Y\) in place of \(X\)). Let us prove that \((*)\) is an isomorphism. The question being local on \(Y\), we may suppose \(S\) and \(Y\) affine. By dévissage, in the first case we immediately reduce to the case where \(M\) is an ordinary quasi-coherent \(\underline{O}_S\)-Module. Noting that the functors \(Rf_*(E)\otimes_S^{\mathrm L}.\) and \(Rf_*(E\otimes_S^{\mathrm L}.)\) are ``way-out to the left'' on the category \(D^b(S)_{\qcoh}\) ([H] I 7), we then reduce, by taking a left resolution of \(M\) by free Modules and truncating, to the case where \(M\) is free. Since the \(R^if_*\) commute with inductive limits (\(f\) being quasi-compact and quasi-separated), we may even suppose \(M\) free of finite type, and then \(M=\underline{O}_S\). The arrow \((*)\) is then tautologically an isomorphism. In the second case, by dévissage we reduce to \(M=P\), where \(P\) is a flat \(A\)-module
%%%%%%%%%%  scan idx 262  |  folio 249  %%%%%%%%%%
\SGAPrintPage{249}%
\((S=\Spec(A))\). Putting \(A'=D_A(P)\) (EGA \(0_{\mathrm{IV}}\) 18.2.3), \(S'=\Spec(A')\), everything reduces to proving that \(Rf_*(E)\) commutes with the flat base change \(S'\to S\), i.e. that the canonical arrow
\[
 q^*Rf_*(E)\longrightarrow Rf'_*p^*(E)
\]
is an isomorphism (we have put \(f_{(S')}=f'\) and denoted by \(p\) (resp. \(q\)) the canonical projection \(X_{(S')}\longrightarrow X\) (resp. \(Y_{(S')}\longrightarrow Y\))). But this follows from (EGA III 1.4.15 and IV 1.7.21). This completes the proof of (3.7).

\sgaStableTarget{sga6.e03.cor.3-7-1}\textbf{Corollary 3.7.1.} \uline{Let} \(S\) \uline{be a scheme}, \(f:X\to Y\) \uline{a quasi-compact and quasi-separated} \(S\)-\uline{morphism of quasi-compact schemes over} \(S\), \(E\in\Ob D^b(X)_{\qcoh}\). \uline{If} \(E\) \uline{is locally of finite flat amplitude rel.} \uline{to} \(S\), \uline{the same holds for} \(Rf_*(E)\). \uline{More precisely}, \uline{if} \(S\) \uline{is quasi-compact}, \uline{and if} \(N\) \uline{is an integer such that} \(R^if_*(F)=0\) \uline{for every} \(i>N\) \uline{and every quasi-coherent} \(\underline{O}_X\)-\uline{Module} \(F\) (EGA III 1.4.12 \uline{and} IV 1.7.21), \uline{then we have}
\[
 \tamp_S(E)\subset[a,b]
 \Longrightarrow
 \tamp_S(Rf_*(E))
 \subset[a,b+N].
\]

\uline{Proof}. We may evidently restrict ourselves to proving the last assertion. By (3.3 (i)\(\Longleftrightarrow\)(iv)), and taking into account that \(Rf_*\) ``commutes with restriction to an open subset of \(S\)'', we must prove that the hypothesis \(\tamp_S(E)\subset[a,b]\) implies
\[
 H^i(Rf_*(E)\otimes_S^{\mathrm L}M)=0
\]
for every \(i\notin[a,b+N]\) and every quasi-coherent \(\underline{O}_S\)-Module \(M\). Now the preceding hypothesis means that \(H^i(E\otimes_S^{\mathrm L}M)=0\) for every \(i\notin[a,b]\) and every quasi-coherent \(\underline{O}_S\)-Module \(M\), and (3.7) gives the result.
%%%%%%%%%%  scan idx 263  |  folio 250  %%%%%%%%%%
\SGAPrintPage{250}%
\sgaStableTarget{sga6.e03.cor.3-7-2}\textbf{Corollary 3.7.2.} \uline{Let} \(f:X\to Y\) \uline{be a quasi-compact}, \uline{quasi-separated morphism}, \uline{locally of finite tor-dimension} (3.1), \uline{and let} \(E\in\Ob D^b(X)_{\qcoh}\). \uline{If} \(E\) \uline{is locally of finite flat amplitude} (\uline{in the absolute sense}), \uline{the same holds for} \(Rf_*(E)\).

\uline{Proof}. Combine (3.4.1 a)) and (3.7.1).

\sgaStableTarget{sga6.e03.sec.4}\section*{4. Relative perfection}

\sgaStableTarget{sga6.e03.def.4-1}\textbf{Definition 4.1.} \uline{Let} \(f:X\to Y\) \uline{be a locally finite type morphism of schemes}, \(E\in\Ob D(X)\), \([a,b]\) \uline{an interval of} \(\mathbb Z\). \uline{We say that} \(E\) \uline{is} of perfect amplitude \uline{relative to} \(f\) (\uline{or} \(Y\) \uline{when there is no ambiguity}), contained in \([a,b]\), \uline{and write}
\[
 \parfamp_f(E)\subset[a,b]\qquad(\text{or }\parfamp_Y(E)\subset[a,b])
\]
\uline{if}, \uline{relative to} \(f\), \(E\) \uline{is pseudo-coherent} (1.2) \uline{and of flat amplitude contained in} \([a,b]\) (3.1). \uline{We say that} \(f\) is of perfect amplitude contained in \([-n,0]\), \uline{or of} perfect dimension \(\leq n\) \((n\geq0)\), \uline{if} \(\underline{O}_X\) \uline{is of perfect amplitude contained in} \([-n,0]\). \uline{We say that} \(E\) \uline{is} of finite perfect amplitude relative to \(f\) \uline{if there is an interval} \([a,b]\) \uline{such that} \(\parfamp_f(E)\subset[a,b]\). \uline{We say that} \(E\) \uline{is} perfect rel. to \(f\) \uline{if} \(E\) \uline{is locally of finite perfect amplitude rel.} \uline{to} \(f\). \uline{We say that} \(f\) is perfect \uline{if} \(f\) \uline{is} \uline{locally of finite perfect amplitude}.

\sgaStableTarget{sga6.e03.ex.4-1-1}\textbf{Example 4.1.1.} A regular immersion of codimension \(d\) (VII 1.4) is of perfect dimension \(\leq d\) (as is seen from the Koszul complex). A locally complete-intersection morphism (i.e. locally a composite
%%%%%%%%%%  scan idx 264  |  folio 251  %%%%%%%%%%
\SGAPrintPage{251}%
of a regular immersion and a smooth morphism) \var{is a perfect morphism}{is a perfect morphism.}{EN-VAR-165}

\sgaStableTarget{sga6.e03.notation.4-2}\textbf{Notation 4.2.} Let \(f:X\to Y\) be a locally finite type morphism of schemes. We denote by \(D(f)_{\parf}\), or \(D(X)_{f\text{-}\parf}\), or \(D(X)_{Y\text{-}\parf}\) when no confusion is to be feared (resp. \(D(f)_{\uparf}\), \(D(X)_{f\text{-}\uparf}\), \(D(X)_{Y\text{-}\uparf}\)) the full subcategory of \(D(X)\) formed by the complexes perfect (resp. of finite perfect amplitude) rel. to \(f\). By (I 5.3) and (1.4), \(D(f)_{\parf}\) and \(D(f)_{\uparf}\) are triangulated subcategories of \(D(X)\). When \(X\) is quasi-compact, we have \(D(f)_{\parf}=D(f)_{\uparf}\).

\sgaStableTarget{sga6.e03.prop.4-3}\textbf{Proposition 4.3} (cf. I 5.8). \uline{Let} \(f:X\to Y\) \uline{be a locally finite type morphism of schemes}, \(E\in\Ob D(X)\), \([a,b]\) \uline{an interval of} \(\mathbb Z\), \(x\) \uline{a point of} \(X\), \(y=f(x)\). \uline{The following conditions are equivalent}:
\begin{enumerate}[label=(\roman*)]
\item \(\parfamp_f(E)\subset[a,b]\) \uline{in a neighborhood of} \(x\):
\item \(E\) \uline{is}, \uline{in a neighborhood of} \(x\), \((a-1)\)-\uline{pseudo-coherent rel.} \uline{to} \(f\) \uline{and acyclic outside} \([a,b]\), \uline{and we have} \(\tamp_{\underline{O}_{Y,y}}(E_x)\subset[a,b]\).
\end{enumerate}

\uline{Proof}. The implication (i) \(\Longrightarrow\) (ii) is trivial; let us prove (ii) \(\Longrightarrow\) (i). The question being local in a neighborhood of \(x\), we immediately reduce, by the definitions, to the case where \(f\) is smooth. Then \(E\) is \((a-1)\)-pseudo-coherent in a neighborhood of \(x\); hence, after localizing, we may suppose \(E\) strictly \((a-1)\)-pseudo-coherent (I 2.1) and zero in degree \(>b\), and after truncating (i.e. replacing \(E\) by the complex \((0\to E^a/B^a\to E^{a+1}\to\ldots E^b\to 0)\)), we may suppose that \(E^i\) is free of finite type for \(i>a\), \(E^a\) of finite presentation, and \(E^i=0\) for \(i\notin[a,b]\). By (I 5.1.1), \(E^a_x\) is flat rel. to \(\underline{O}_{Y,y}\), hence (EGA IV 11.3.1) \(E^a\) is \(f\)-flat in a neighborhood of \(x\). Write \(E^a=F\) for brevity; we must show that \(F\) is pseudo-coherent (in a neighborhood of \(x\)). After embedding \(X\) in a standard affine space over \(Y\), we may
%%%%%%%%%%  scan idx 265  |  folio 252  %%%%%%%%%%
\SGAPrintPage{252}%
suppose that \(X=Y[t_1,\dots,t_r]\), with \(Y\) affine and \(f\) the canonical projection. Then \(Y\) is a projective limit of noetherian affine schemes \(Y_\alpha\). Denote by \(f_\alpha:X_\alpha=Y_\alpha[t_1,\dots,t_r]\to Y_\alpha\) the canonical projection. There is (EGA IV 8.5.2 and 11.2.6) an index \(\alpha\) and an \(\underline{O}_{X_\alpha}\)-Module \(F_\alpha\) of finite presentation and \(f_\alpha\)-flat such that \(g_\alpha^*(F_\alpha)\simeq F\), where \(g_\alpha:X\to X_\alpha\) is the canonical morphism. Since \(X_\alpha\) is noetherian, \(F_\alpha\) is coherent, hence pseudo-coherent rel. to \(f_\alpha\). Since, on the other hand, \(F_\alpha\) is \(f_\alpha\)-flat and \(f_\alpha\) is flat, we have \(\var{g_\alpha^*(F)\simeq F_\alpha}{g_\alpha^*(F_\alpha)\simeq F_\alpha}{EN-VAR-166}\otimes^{\mathrm L}_{Y_\alpha}\underline{O}_Y\), and it follows from (1.10 a)) that \(F\) is pseudo-coherent rel. to \(Y\), hence pseudo-coherent, q.e.d.

\sgaStableTarget{sga6.e03.cor.4-3-1}\textbf{Corollary 4.3.1.} \uline{Let} \(f:X\to Y\) \uline{be a locally finite type morphism}, \uline{and} \(E\) \uline{an} \(\underline{O}_X\)-\uline{Module}. \uline{For} \(\parfamp_f(E)\subset[0,0]\) \uline{to hold, it is necessary and sufficient that} \(E\) \uline{be}, \uline{rel.} \uline{to} \(f\), \uline{flat and of finite presentation} (1.2.1 c)). \uline{In particular}, \uline{for} \(f\) \uline{to have perfect dimension zero} (4.1), \uline{it is necessary and sufficient that} \(f\) \uline{be flat and locally of finite presentation}.

\uline{Proof}. Apply (4.3 (ii) \(\Longrightarrow\) (i)) with \(a=b=0\).

\sgaStableTarget{sga6.e03.cor.4-3-2}\textbf{Corollary 4.3.2.} \uline{Let} \(f:X\to Y\) \uline{be a locally finite type morphism}, \uline{and} \(E\in\Ob D(X)\) \uline{a complex of} \(\underline{O}_X\)-\uline{Modules flat and of finite presentation} \uline{rel.} \uline{to} \(f\) (1.2.1 c)), \uline{and} \([a,b]\) \uline{an interval of} \(\mathbb Z\). \uline{If} \(E^i=0\) \uline{for} \(i\notin[a,b]\), \uline{we have} \(\parfamp_f(E)\subset[a,b]\). \uline{If} \(E\) \uline{is bounded above in degree}, \(E\) \uline{is pseudo-coherent rel.} \uline{to} \(f\).

\uline{Proof}. A consequence of (4.3.1) and (I 4.15 and 2.11).

\sgaStableTarget{sga6.e03.prop.4-4}\textbf{Proposition 4.4.} \uline{Let there be a commutative triangle of schemes}
\[
\begin{tikzcd}
X \arrow[r,hook,"i"] \arrow[d,"f"'] & Y \arrow[dl,"g"]\\
S &
\end{tikzcd},
\]
%%%%%%%%%%  scan idx 266  |  folio 253  %%%%%%%%%%
\SGAPrintPage{253}%
\uline{where} \(i\) \uline{is a closed immersion and} \(g\) \uline{a smooth morphism}, \uline{and let} \(E\in\Ob D(X)\). \uline{For} \(E\) \uline{to be perfect rel.} \uline{to} \(f\), \uline{it is necessary and sufficient that} \(i_*(E)\) \uline{be perfect} (\uline{in the absolute sense} (I 4.8)). \uline{More precisely}, \uline{suppose} \(g\) \uline{is of relative dimension} \(\leq d\), \uline{and let} \([a,b]\) \uline{be an interval of} \(\mathbb Z\). \uline{Then we have}:
\begin{enumerate}[label=(\roman*)]
\item \(\var{\parfamp(i_*(E)\subset[a,b]\Longrightarrow\parfamp_f(E)\subset[a,b]}{\parfamp(i_*(E))\subset[a,b]\Longrightarrow\parfamp_f(E)\subset[a,b]}{EN-VAR-167}\);
\item \(\parfamp_f(E)\subset[a,b]\Longrightarrow\parfamp(i_*(E))\subset[a-d,b]\).
\end{enumerate}

\uline{Proof}. An immediate consequence of (3.6).

\sgaStableTarget{sga6.e03.prop.4-5}\textbf{Proposition 4.5.} \uline{Let}
\[
\begin{tikzcd}
X \arrow[rr,"f"] \arrow[dr,"h"'] && Y \arrow[dl,"g"]\\
& S &
\end{tikzcd}
\]
\uline{be a commutative triangle of locally finite type morphisms of schemes}, \uline{and let} \(E\in\Ob D^-(X)\), \(F\in\Ob D^-(Y)\), \([a,b]\), \([c,d]\) \uline{be intervals of} \(\mathbb Z\). \uline{Then we have}:
\[
 \parfamp_f(E)\subset[a,b]\text{ and }\parfamp_g(F)\subset[c,d]
 \Longrightarrow
 \parfamp_h(E\otimes_Y^{\mathrm L}F)
 \subset[a+c,b+d].
\]
\uline{Consequently}, \uline{if} \(E\) \uline{is perfect rel.} \uline{to} \(f\) \uline{and} \(F\) \uline{perfect rel.} \uline{to} \(g\), \(E\otimes_Y^{\mathrm L}F\) \uline{is perfect rel.} \uline{to} \(h\).

\uline{Proof}. Combine (1.11) and (3.4).

\sgaStableTarget{sga6.e03.cor.4-5-1}\textbf{Corollary 4.5.1.} \uline{Under the hypotheses of} (4.5):
\begin{enumerate}[label=\alph*)]
\item \uline{if} \(g\) \uline{has perfect dimension} \(\leq n\), \uline{we have}
\[
 \parfamp_f(E)
 \subset[a,b]
 \Longrightarrow
 \parfamp_h(E)
 \subset[a-n,b];
\]
\uline{consequently}, \uline{if} \(g\) \uline{is perfect}, \(E\) \uline{perfect rel.} \uline{to} \(f\) \uline{implies} \(E\) \uline{perfect rel.} \uline{to} \(h\);
\item \uline{if} \(f\) \uline{has perfect dimension} \(\leq m\), \uline{we have}
\[
 \parfamp_g(F)
 \subset[c,d]
 \Longrightarrow
 \parfamp_h(\mathrm Lf^*(F))
 \subset[c-m,d];
\]
%%%%%%%%%%  scan idx 267  |  folio 254  %%%%%%%%%%
\SGAPrintPage{254}%
\uline{consequently}, \uline{if} \(f\) \uline{is perfect}, \(F\) \uline{perfect rel.} \uline{to} \(g\) \uline{implies} \(\mathrm Lf^*(F)\) \uline{perfect rel.} \uline{to} \(h\).
\end{enumerate}

\sgaStableTarget{sga6.e03.rem.4-5-2}\textbf{Remark 4.5.2.} With the preceding notation, when \(g\) is perfect, \(E\) perfect rel. to \(h\) does not necessarily imply \(E\) perfect rel. to \(f\), as the ever-recurring example of a singular point of a scheme over a field shows. We have the following criterion:

\sgaStableTarget{sga6.e03.prop.4-6}\textbf{Proposition 4.6} (\uline{criterion of perfection by fibers})\SGAAuthorFootnote{1}{Analogous to the ``criterion of flatness by fibers'' (EGA IV 11.3.10).}. \uline{Let}
\[
\begin{tikzcd}
X \arrow[rr,"f"] \arrow[dr,"h"'] && Y \arrow[dl,"g"]\\
& S &
\end{tikzcd}
\]
\uline{be a commutative triangle of morphisms of schemes}. \uline{We suppose} \(f\), \(h\) \uline{locally of finite type}, \uline{and} \(g\) \uline{locally of finite presentation and} flat. \uline{Let} \(x\in X\), \(y=f(x)\), \(s=h(x)=g(y)\). \uline{Finally}, \uline{let} \(E\in\Ob D^-(X)\), \uline{and let} \([a,b]\) \uline{be an interval of} \(\mathbb Z\). \uline{We denote by} \(f_s:X_s\to Y_s\) \uline{the fiber of} \(f\) \uline{at} \(s\), \uline{i.e.} \uline{the morphism obtained from} \(f\) \uline{by the base change} \(\Spec(k(s))\to S\). \uline{Then the following conditions are equivalent}:
\begin{enumerate}[label=(\roman*)]
\item \(\parfamp_h(E)\subset[a,b]\) \uline{in a neighborhood of} \(x\), \uline{and}
\[
 \tamp_{\underline{O}_{Y_s,y}}(E_x\otimes_{\underline{O}_{Y,y}}^{\mathrm L}\underline{O}_{Y_s,y})
 \subset[a,b];
\]
\item \(\parfamp_f(E)\subset[a,b]\) \uline{in a neighborhood of} \(x\).
\end{enumerate}

\uline{Proof}. (ii) \(\Longrightarrow\) (i). We already know (4.3.1) that \(g\) has perfect dimension zero. Thus (4.5.1 a)), condition (ii) implies \(\parfamp_h(E)\subset[a,b]\) in a neighborhood of \(x\). \var{It also implies. by 4.3),}{It also implies (by 4.3),}{EN-VAR-168} \(\tamp_{\underline{O}_{Y,y}}(E_x)\subset[a,b]\), hence (I 5.1 (i)), \(\tamp_{\underline{O}_{Y_s,y}}(E_x\otimes_{\underline{O}_{Y,y}}^{\mathrm L}\underline{O}_{Y_s,y})\subset[a,b]\), which
%%%%%%%%%%  scan idx 268  |  folio 255  %%%%%%%%%%
\SGAPrintPage{255}%
completes the proof of (ii) \(\Longrightarrow\) (i).

(i) \(\Longrightarrow\) (ii). The question being local in a neighborhood of \(x\), we may suppose that \(f\) factors as \(X\overset{i}{\hookrightarrow}X'\overset{f'}{\longrightarrow}Y\), where \(i\) is a closed immersion and \(f'\) a smooth morphism. It is immediate that conditions (i) and (ii) are unchanged when we replace \(E\) by \(i_*(E)\), \(f\) by \(f'\) (and \(h\) by \(gf'\)), so that we may suppose \(f\) smooth. Since \(g\) is perfect (4.3.1), \(h\) is perfect, hence (1.12) \(E\) is pseudo-coherent. After localizing, by (II 2.2.10 b)) and (I 2.10 a)), we may suppose that \(E\) is strictly pseudo-coherent (I 2.1) and zero in degree \(>b\), and that the truncation arrow
\begin{equation}
 E\longrightarrow(0\longrightarrow E^a/B^a\longrightarrow E^{a+1}\longrightarrow\ldots\longrightarrow E^b\longrightarrow0)\tag*{$t)$}
\end{equation}
is a quasi-isomorphism (hypothesis (i) indeed implies that \(E\) \var{is, in a neighborhood of \(s\), acyclic outside}{is, in a neighborhood of \(x\), acyclic outside}{EN-VAR-169} \([a,b]\)). We must prove that \(E^a/B^a\) is \(f\)-flat in a neighborhood of \(x\). Denote by \(i:X_s\longrightarrow X\) the canonical arrow. We have:
\begin{equation}
 E_x\otimes_{\underline{O}_{Y,y}}^{\mathrm L}\underline{O}_{Y_s,y}\simeq i^*(E)_x.\notag
\end{equation}
Since \(h\) is flat, hypothesis (i) implies, by (3.5.2), \(\tamp_{h_s}(i^*(E))\subset[a,b]\) in a neighborhood of \(x\). \var{Thus, after localizing,. we may suppose}{Thus, after localizing, we may suppose}{EN-VAR-170} that the truncation arrow
\begin{equation}
 i^*(E)\longrightarrow(0\longrightarrow i^*(E^a)/B^a(i^*(E))\longrightarrow\ldots\longrightarrow i^*(E^b)\longrightarrow0)\tag*{$i^*(t)$}
\end{equation}
is a quasi-isomorphism. The hypothesis \(\tamp_{\underline{O}_{Y_s,y}}i^*(E)_x\subset[a,b]\) then implies, by (I 5.1.1), that \(i^*(E^a/B^a)_x\) is a flat \(\underline{O}_{Y_s,y}\)-module. On the other hand, since the \var{arrow \((*)\)}{arrow \((t)\)}{EN-VAR-171} is a quasi-isomorphism, the hypothesis \(\tamp_h(E)\subset[a,b]\) implies, again by (I 5.1.1),
%%%%%%%%%%  scan idx 269  |  folio 256  %%%%%%%%%%
\SGAPrintPage{256}%
that \(E^a/B^a\) is \(h\)-flat. By the criterion of flatness by fibers (EGA IV 11.3.10), we conclude that \(E^a/B^a\) is \(f\)-flat in a neighborhood of \(x\), q.e.d.

\sgaStableTarget{sga6.e03.cor.4-6-1}\textbf{Corollary 4.6.1.} \uline{Let} \(f:X\to S\) \uline{be a flat morphism locally of finite presentation}, \(x\in X\), \(s=f(x)\), \(E\in\Ob D^-(X)\), \([a,b]\) \uline{an interval of} \(\mathbb Z\). \uline{Denote by} \(i_s:X_s=X\times_S\Spec(k(s))\to X\) \uline{the canonical projection}, \uline{and put} \(E_{X_s}=\mathrm Li_s^*(E)\). \uline{Then we have}
\[
\begin{gathered}
 \parfamp_f(E)\subset[a,b]\text{ \uline{in a neighborhood of} }x\text{ \uline{and} }\tamp(E_{X_s})_x\subset[a,b]\\
 \Longleftrightarrow \parfamp(E)\subset[a,b]\text{ \uline{in a neighborhood of} }x.
\end{gathered}
\]
\uline{Consequently}, \(E\) \uline{is perfect if and only if} \(E\) \uline{is perfect rel.} \uline{to} \(f\) \uline{and} \(E_{X_s}\) \uline{is perfect for every} \(s\in S\).

\uline{Proof}. Apply (4.6) with \(X=Y\), \(f=\Id_X\).

\sgaStableTarget{sga6.e03.rem.4-6-2}\textbf{Remark 4.6.2.} From (4.6.1) and (I 5.10) it follows that, when \(f\) is locally of finite presentation, flat and has regular schemes as fibers, \(E\) perfect rel. to \(f\) is equivalent to \(E\) perfect (this equivalence had already been obtained (4.4) for \(f\) smooth).

\sgaStableTarget{sga6.e03.prop.4-7}\textbf{Proposition 4.7.} \uline{Let} \(S\) \uline{be a scheme}, \uline{and for} \(i=1,2\), \(f_i:X_i\to Y_i\) \uline{an} \(S\)-\uline{morphism locally of finite type}, \(E_i\in\Ob D(X_i)\), \([a_i,b_i]\) \uline{an interval of} \(\mathbb Z\). \uline{Suppose} \(X_1\) \uline{and} \(X_2\), \uline{as well as} \(Y_1\) \uline{and} \(Y_2\), \uline{are tor-independent over} \(S\) (1.5). \uline{Put} \(a=a_1+a_2\), \(b=b_1+b_2\), \(f=f_1\times_Sf_2\). \uline{Then we have}:
\begin{enumerate}[label=\alph*)]
\item \uline{If} \(\parfamp_{f_1}(E_1)\subset[a_1,b_1]\), \uline{and if} \(E_2\) \uline{is pseudo-coherent rel.} \uline{to} \(f_2\) \uline{and acyclic in degree outside} \([a_2,b_2]\), \(E_1\otimes_S^{\mathrm L}E_2\) (\uline{cf.} (1.6)) \uline{is pseudo-coherent rel.} \uline{to} \(f\) \uline{and acyclic in degree outside} \([a,b]\).
%%%%%%%%%%  scan idx 270  |  folio 257  %%%%%%%%%%
\SGAPrintPage{257}%
\item \(\parfamp_{f_i}(E_i)\subset[a_i,b_i]\) \uline{for} \(i=1,2\)
\[
 \Longrightarrow\parfamp_f(E_1\otimes_S^{\mathrm L}E_2)\subset[a,b].
\]
\end{enumerate}

\uline{Proof}. Combine (1.9) and (3.5).

\sgaStableTarget{sga6.e03.cor.4-7-1}\textbf{Corollary 4.7.1.} \var{\uline{Under the hypotheses of} (4.6)}{\uline{Under the hypotheses of} (4.7)}{EN-VAR-172}, \uline{if}, \uline{for} \(i=1,2\), \(f_i\) \uline{has perfect dimension} \(\leq n_i\), \(f\) \uline{has perfect dimension} \(\leq n_1+n_2\). \uline{Consequently}, \uline{if} \(f_1\) \uline{and} \(f_2\) \uline{are perfect}, \uline{the same holds for} \(f\).

\sgaStableTarget{sga6.e03.cor.4-7-2}\textbf{Corollary 4.7.2.} \uline{Let}
\[
\begin{tikzcd}
X \arrow[d,"f"'] & X' \arrow[l,"g"'] \arrow[d,"f'"]\\
Y & Y' \arrow[l,"h"']
\end{tikzcd}
\]
\uline{be a Cartesian square of schemes}, \uline{where} \(f\) \uline{is locally of finite type and} \(X\) \uline{and} \(Y'\) \uline{are supposed tor-independent over} \(Y\), \uline{and let} \(E\in\Ob D(X)\), \uline{and} \([a,b]\) \uline{be an interval of} \(\mathbb Z\). \uline{We have}:
\[
 \parfamp_f(E)\subset[a,b]\Longrightarrow\parfamp_{f'}(\mathrm Lg^*(E))\subset[a,b],
\]
\uline{and consequently the functor} \(\mathrm Lg^*:D^-(X)\to D^-(X')\) \uline{sends} \(D(f)_{\parf}\) (resp. \(D(f)_{\uparf}\)) into \(D(f')_{\parf}\) (resp. \(D(f')_{\uparf}\)) (4.2).

\sgaStableTarget{sga6.e03.prop.4-8}\textbf{Proposition 4.8.} \uline{Let} \(S\) \uline{be a scheme}, \(f:X\to Y\) \uline{a proper} \(S\)-\uline{morphism of} \(S\)-\uline{schemes locally of finite type}, \(E\in\Ob D^b(X)\). \uline{Then}, \uline{subject to the validity of} (2.1) \uline{for} \(f\) (\uline{hence in any case when} \(S\) \uline{is locally noetherian} \uline{or when} \(f\) \uline{is projective}):
\begin{enumerate}[label=\alph*)]
\item \uline{If} \(E\) \uline{is perfect rel.} \uline{to} \(S\), \uline{the same holds for} \(Rf_*(E)\).
\item \uline{Suppose} \(S\) \uline{quasi-compact}, \(X\) \uline{and} \(Y\) \uline{of finite type over} \(S\); \uline{let} \([a,b]\) \uline{be an interval of} \(\mathbb Z\), \uline{and} \(N\) \uline{an integer such that} \(R^if_*(M)=0\) \uline{for every} \(i>N\) \uline{and}
%%%%%%%%%%  scan idx 271  |  folio 258  %%%%%%%%%%
\SGAPrintPage{258}%
\uline{every} \(\underline{O}_X\)-\uline{Module} \(M\) \uline{that is quasi-coherent} (\uline{such an integer exists by} (EGA III 1.4.12 \uline{and} IV 1.7.21))\,; \uline{we then have}
\[
 \parfamp_S(E)\subset[a,b]\Longrightarrow\parfamp_S(Rf_*(E))\subset[a,b+N].
\]
\end{enumerate}
\uline{Consequently}, \uline{the functor} \(Rf_*\) \uline{sends} \(D(X)_{S\text{-}\uparf}\) \uline{into} \(D(Y)_{S\text{-}\uparf}\)\SGAAuthorFootnote{1}{moreover, \(D(X)_{S\text{-}\uparf}=D(X)_{S\text{-}\parf}\) (and likewise with \(Y\)), since \(X\) and \(Y\) are quasi-compact (cf. (4.2)).}.

\uline{Proof}. Assertion a) follows from b), and b) is obtained by combining (2.1) and (3.7.1).

\sgaStableTarget{sga6.e03.cor.4-8-1}\textbf{Corollary 4.8.1.} \uline{Let} \(f:X\to Y\) \uline{be a proper and perfect morphism}, \uline{and} \(E\in\Ob D^+(X)\). \uline{Then}, \uline{subject to the validity of} (2.1) \uline{for} \(f\), \uline{if} \(E\) \uline{is perfect}, \uline{the same holds for} \(Rf_*(E)\)\,; \uline{moreover}, \uline{if} \(Y\) \uline{is quasi-compact}, \uline{the functor} \(Rf_*\) \uline{sends} \(D(X)_{\uparf}\) \uline{into} \(D(Y)_{\uparf}\)\SGAAuthorFootnote{2}{We have \(D(X)_{\uparf}=D(X)_{\parf}\) (and likewise with \(Y\)).}.

\uline{Proof}. Indeed, when \(E\) is perfect, \(E\) is perfect rel. to \(f\) by (4.5.1 a)).

\sgaStableTarget{sga6.e03.prop.4-9}\textbf{Proposition 4.9}\SGAAuthorFootnote{3}{This result and its corollaries will not be used in the remainder of this Seminar, but play an important role in the Lefschetz--Verdier formula for coherent sheaves; see SGA 5 Exp.\ III.}. \uline{Let} \(S\) \uline{be a noetherian scheme}, \(X\) \uline{and} \(Y\) \uline{be} \(S\)-\uline{schemes of finite type}, \(f:X\to Y\) \uline{a compactifiable} \(S\)-\uline{morphism} (\uline{i.e.} \uline{a composite of an open immersion and a proper morphism}), \(E\in\Ob D(X)\), \(F\in\Ob D^+(Y)\). \uline{If} \(E\) \uline{is perfect rel.} \uline{to} \(f\), \uline{and} \(F\) \uline{perfect rel.} \uline{to} \(S\), \(\mathrm R\underline{\Hom}(E,f^!(F))\) \uline{is perfect rel.} \uline{to} \(S\)\SGAAuthorFootnote{4}{For the definition of the functor \(f^!\), see [H] App.\ n\textsuperscript{o}~4.}.

\uline{Proof}. The question being local on \(X\), we may suppose that \(f\) factors as \(X\xrightarrow{i}X'\xrightarrow{f'}Y\), where \(i\) is a closed immersion and \(f'\) a smooth morphism. When \(E\) is perfect rel. to \(f\), \(E\) is perfect rel. to \(i\) (4.4), and when \(F\) is perfect rel. to \(S\), \(f'^!(F)\) is perfect rel. to \(S\) by (4.5.1 b)), since \(f'^!(F)\) is locally isomorphic to \(f'^*(F)\) up to a shift. Since \(f^!\simeq i^!f'^!\), we therefore reduce to proving the assertion when \(f\) is a closed immersion. Thus suppose \(E\) perfect rel. to \(f\), i.e. \(f_*(E)\) perfect, and \(F\) perfect rel. to \(S\). We must show that \(f_*\mathrm R\underline{\Hom}(E,f^!F)\) is
%%%%%%%%%%  scan idx 272  |  folio 259  %%%%%%%%%%
\SGAPrintPage{259}%
perfect rel. to \(S\). But by the duality theorem ([H] III 6.7), we have
\[
 f_*\mathrm R\underline{\Hom}(E,f^!F)\simeq\mathrm R\underline{\Hom}(f_*(E),F).
\]
Since \(f_*(E)\) is perfect, we have, by (I 7.7),
\[
 \mathrm R\underline{\Hom}(f_*(E),F)\simeq(f_*(E))^{\vee}\otimes^{\mathrm L}F,
\]
and (4.5) (where \(X=Y\) and \(f=\Id_X\)) gives the result.

\sgaStableTarget{sga6.e03.cor.4-9-1}\textbf{Corollary 4.9.1.} \var{\uline{Under the hypotheses of} (4.8)}{\uline{Under the hypotheses of} (4.9)}{EN-VAR-173}, \uline{if} \(f\) \uline{is perfect}, \(f^!\) \uline{sends} \(D(Y)_{S\text{-}\parf}\) \uline{into} \(D(X)_{S\text{-}\parf}\).

\uline{Proof}. Take \(E=\underline{O}_X\).

\sgaStableTarget{sga6.e03.cor.4-9-2}\textbf{Corollary 4.9.2.} \uline{Let} \(S\) \uline{be a noetherian scheme}, \(f:X\to S\) \uline{a compactifiable morphism}. \uline{Put} \(f^!(\underline{O}_S)=K_X\), \(\mathrm R\underline{\Hom}(.,K_X)=\underline{D}_X\). \uline{For} \(E\in\Ob D(f)_{\parf}\), \uline{we have} \(\underline{D}_X(E)\in\Ob D(f)_{\parf}\), \uline{and the canonical arrow} (I 7.2)
\[
 E\longrightarrow \underline{D}_X\underline{D}_X(E)
\]
\uline{is an isomorphism}.

\uline{Proof}. \var{For the first assertion, apply (4.8)}{For the first assertion, apply (4.9)}{EN-VAR-174} with \(Y=S\), \(F=\underline{O}_S\). For the second, we easily reduce to the case where \(f\) is a closed immersion. It then suffices to prove that the arrow obtained from the canonical arrow by applying \(f_*\) is an isomorphism. By global duality ([H] III 6.7), we have
\[
 f_*\underline{D}_X\simeq\underline{D}_Sf_*.
\]
We therefore reduce to showing that the canonical arrow
\[
 f_*(E)\longrightarrow \underline{D}_S\underline{D}_Sf_*(E)
\]
is an isomorphism for \(E\) perfect rel. to \(f\), i.e. \(f_*(E)\) perfect, which follows from ordinary biduality for perfect complexes (I 7.2).

\sgaStableTarget{sga6.e03.sec.5}\section*{5. Applications: exchange and semi-continuity theorems\SGAAuthorFootnote{1}{The results of this section will not be used in the remainder of the Seminar.}}

\sgaStableTarget{sga6.e03.lem.5-1}\textbf{Lemma 5.1.} \uline{Let} \((X,\underline{O}_X)\) \uline{be a ringed topos} (\uline{commutative}), \uline{denote by} \(\Mod(X)\) \uline{the category of} \(\underline{O}_X\)-\uline{Modules}, \(D(X)=D(\Mod(X))\). \uline{Let} \(E\in\Ob D^-(X)\), \uline{and} \(p\in\mathbb Z\). \uline{Consider the following conditions}:
%%%%%%%%%%  scan idx 273  |  folio 260  %%%%%%%%%%
\SGAPrintPage{260}%
\begin{enumerate}[label=\alph*)]
\item \uline{The functor} \(M\mapsto H^p(E\otimes^{\mathrm L}M)\) \uline{from} \(\Mod(X)\) \uline{to} \(\Mod(X)\) \uline{is right exact}.
\item[a')] \uline{The canonical arrow}
\[
 H^p(E)\otimes M\longrightarrow H^p(E\otimes^{\mathrm L}M)
\]
\uline{is an isomorphism for every} \(M\in\Ob\Mod(X)\).
\item[b)] \uline{The functor} \(M\mapsto H^{p+1}(E\otimes^{\mathrm L}M)\) \uline{from} \(\Mod(X)\) \uline{to} \(\Mod(X)\) \uline{is left exact}.
\item[b')] \uline{The canonical arrow defined below}
\begin{equation}
 H^{p+1}(E\otimes^{\mathrm L}M)\longrightarrow\underline{\Hom}(H^{-p-1}(E^{\vee}),M)\tag{$*$}
\end{equation}
\uline{is an isomorphism for every} \(M\in\Ob\Mod(X)\).
\end{enumerate}
\uline{Then we have}:
\[
 \mathrm a)\Longleftrightarrow\mathrm a')\Longleftrightarrow\mathrm b)\Longleftarrow\mathrm b').
\]
\uline{Moreover}, \uline{if} \(E\) \uline{is pseudo-coherent}, \uline{the conditions} a), a'), b), b') \uline{are equivalent}.

\uline{Definition of} \((*)\): The canonical arrow (I 7.5) \(E^{\vee}\otimes^{\mathrm L}E\longrightarrow\underline{O}_X\), where \(E^{\vee}=\mathrm R\underline{\Hom}(E,\underline{O}_X)\), gives, by tensoring with \(M\in\Ob\Mod(X)\), an arrow \(E^{\vee}\otimes^{\mathrm L}E\otimes^{\mathrm L}M\longrightarrow M\); hence, by the adjunction formula dear to Cartan (I 7.4), a canonical arrow
\begin{equation}
 E\otimes^{\mathrm L}M\longrightarrow\mathrm R\underline{\Hom}(E^{\vee},M).\tag{$**$}
\end{equation}
Composing the arrow induced by \((**)\) on \(H^{p+1}\) with the canonical arrow
\begin{equation}
 \underline{\Ext}^{p+1}(E^{\vee},M)\longrightarrow\underline{\Hom}(H^{-p-1}(E^{\vee}),M),\tag{$***$}
\end{equation}
we obtain the arrow \((*)\).

\uline{Proof}. The implications
\[
 \mathrm a')\Longrightarrow\mathrm a)\Longleftrightarrow\mathrm b)\Longleftarrow\mathrm b')
\]
are trivial. Let us prove \(\mathrm a)\Longrightarrow\mathrm a')\). Let \(M\) be an \(\underline{O}_X\)-Module. There is an exact sequence \(M^1\longrightarrow M^0\longrightarrow M\longrightarrow0\), where \(M^0\) and \(M^1\) are direct sums of sheaves
%%%%%%%%%%  scan idx 274  |  folio 261  %%%%%%%%%%
\SGAPrintPage{261}%
of the type \(\underline{O}_{U,X}\), where \(U\) is an object of \(X\). Using the right exactness and commutation with direct sums of the functors \(H^p(E)\otimes\,.\,\) and \(H^p(E\otimes^{\mathrm L}\,.\,)\), we reduce to proving that the arrow of a') is an isomorphism for \(M=\underline{O}_{U,X}\), \(U\in\Ob X\). Now \(H^p(E)\otimes\underline{O}_{U,X}\) (resp. \(H^p(E\otimes^{\mathrm L}\underline{O}_{U,X})\)) identifies canonically (SGA 4 IV 2.13) with \(i_!(H^p(E)|U)\) (resp. \(H^p(i_!(E|U))\)), where \(i_!\) is extension by \(0\) outside \(U\), and the result follows because \(i_!\) is exact.

Suppose \(E\) pseudo-coherent, and let us prove \(\mathrm b)\Longrightarrow\mathrm b')\). It clearly suffices to show that \((*)\) is an isomorphism for \(M\) injective. Since the arrow \((***)\) is an isomorphism for \(M\) injective, the assertion will therefore follow from the following one:

\((5.1.1)\): \uline{the arrow} \((**)\) \uline{is an isomorphism if} \(E\) \uline{is pseudo-coherent and} \(M\) \uline{injective}.

\uline{Proof of} \((5.1.1)\). Denote by \(T\) (resp. \(T'\)) the functor from \(D^-(X)\) to \(D(X)\) defined by \(F\mapsto F\otimes^{\mathrm L}M\) (resp. \(\mathrm R\underline{\Hom}(F^{\vee},M)\)). Let \(n\in\mathbb Z\); if \(H^i(F)=0\) for \(i>n\), then \(H^iT(F)=0\) for \(i>n\), \(H^i(F^{\vee})=0\) for \(i<-n\), and \(H^iT'(F)=0\) for \(i>n\) if \(M\) is injective. Since, on the other hand, the arrow \((**)\) commutes with restriction to an open, and an injective restricted to an open is still injective, we may, after localizing and truncating, suppose \(E\) \var{strictly perfect (2.1)}{strictly perfect (I 2.1)}{EN-VAR-175}. By dévissage we then reduce to \(E=\underline{O}_X\), and the result follows.

\sgaStableTarget{sga6.e03.rem.5-2}\textbf{Remarks 5.2.} (i) Condition a) (resp. a'), b), b')) is local in nature: if \((X_\alpha)\) is an open covering of the final object of \(X\), and if, for each \(\alpha\), condition a) (resp\ldots) holds for \(X_\alpha\) and \(E\,|\,X_\alpha\), then a) (resp\ldots) holds.

(ii) If \(x\) is a point of \(X\), denote by \(\mathrm a_x)\) (resp. \(\mathrm a'_x)\), etc.) the condi%
%%%%%%%%%%  scan idx 275  |  folio 262  %%%%%%%%%%
\SGAPrintPage{262}%
tion a) (resp. a'), etc.) relative to \(E_x\) and \(\Mod(\underline{O}_{X,x})\). Then, if \(X\) has enough points, condition a) (resp. a'), b)) is equivalent to the condition ``for every point \(x\) of \(X\), \(\mathrm a_x)\) (resp. \(\mathrm a'_x)\), \(\mathrm b_x))\)'', and is equivalent to the condition ``for every point \(x\) of \(X\), \(\mathrm b'_x)\)'' if \(E\) is pseudo-coherent; moreover, condition a') is equivalent to the analogous condition in which ``for every \(M\in\Ob\Mod(X)\)'' is replaced by ``for every \(\underline{O}_X\)-Module that is the cokernel of a homomorphism \(\underline{O}_{U,X}^p\longrightarrow\underline{O}_{U,X}^q\), where \(U\) is an object of \(X\)''.

(iii) The criterion (EGA III 7.4.2) extends word for word to our situation: if \(E\) has flat components, condition a) simply means that \(\Coker(E^p\longrightarrow E^{p+1})\) is flat, hence locally a direct factor of a free Module of finite type when \(E^p\) and \(E^{p+1}\) are also supposed of finite presentation (I 5.8.3). In particular, we deduce that, if \(E\) is perfect, then \(E\) satisfies condition a) relative to the integer \(p\) if and only if \(E^{\vee}\) satisfies condition a) relative to the integer \(-p-1\).

(iv) Let \(X\) be a \uline{quasi-separated} scheme, \(\underline{O}_X\) the structure sheaf, \(E\in\Ob D^-(X)\) a complex with locally bounded and quasi-coherent cohomology. Denote by \(\mathrm a_{\qcoh})\) (resp. \(\mathrm a'_{\qcoh})\), etc.) the \var{condition analogous to a) (resp. a'), etc). in which}{condition analogous to a) (resp. a'), etc.) in which}{EN-VAR-176} \(\Mod(X)\) is replaced by the category \(\Qcoh(X)\) of quasi-coherent \(\underline{O}_X\)-Modules. Then we have
\[
 \mathrm a)\Longleftrightarrow\mathrm a_{\qcoh})\Longleftrightarrow\mathrm a'_{\qcoh})\Longleftrightarrow\mathrm b_{\qcoh})\Longleftarrow\mathrm b'_{\qcoh}),
\]
and the preceding conditions are equivalent if \(E\) is pseudo-coherent.

Indeed, since \(X\) is quasi-separated, every quasi-coherent sheaf on an affine open of \(X\) is induced by a quasi-coherent sheaf on \(X\) (EGA I 9.4.2 and IV 1.2.7). It follows that condition \(\mathrm a_{\qcoh})\) (resp. \(\mathrm a'_{\qcoh})\), \(\mathrm b_{\qcoh})\), \(\mathrm b'_{\qcoh}))\) holds if and only if it holds over every affine open. To prove the implications above, we may therefore suppose \(X\) affine. Then (II 3.5) \(E\) is isomorphic in \(D(X)\) to a complex of quasi-coherent Modules, and the implications
\[
 \mathrm a_{\qcoh})\Longleftrightarrow\mathrm a'_{\qcoh})\Longleftrightarrow\mathrm b_{\qcoh})\Longleftarrow\mathrm b'_{\qcoh})
\]
%%%%%%%%%%  scan idx 276  |  folio 263  %%%%%%%%%%
\SGAPrintPage{263}%
follow from (5.1) applied to the ``one-point'' topos ringed by \(\Gamma(X,\underline{O}_X)\); the same holds, by (II 2.2.10 a)), for the implication \(\mathrm b_{\qcoh})\Longrightarrow\mathrm b'_{\qcoh})\) when \(E\) is pseudo-coherent. Finally, the equivalence \(\mathrm a)\Longleftrightarrow\mathrm a_{\qcoh})\), or, equivalently, \(\mathrm a')\Longleftrightarrow\mathrm a'_{\qcoh})\) follows from the last remark of (ii) (and the fact that quasi-coherent sheaves ``extend'', since \(X\) is quasi-separated).

\sgaStableTarget{sga6.e03.prop.5-3}\textbf{Proposition 5.3.} (\uline{exchange property}) (EGA III 7.7.5 II). \uline{Let} \(f:X\to Y\) \uline{be a morphism of quasi-compact and quasi-separated schemes}, \(E\in\Ob D^b(X)\) \uline{a complex with quasi-coherent cohomology and finite flat amplitude rel.} \uline{to} \(f\), \uline{and} \(p\in\mathbb Z\). \uline{Consider the following conditions}:
\begin{enumerate}[label=\alph*)]
\item[a)] \uline{The functor} \(M\mapsto R^pf_*(E\otimes_Y^{\mathrm L}M)\) \uline{from} \(\Qcoh(Y)\)\SGAAuthorFootnote{1}{\(\Qcoh(Y)\) = category of quasi-coherent \(\underline{O}_Y\)-Modules.} \uline{to} \(\Qcoh(Y)\)\SGAAuthorFootnote{2}{by (3.7).} \uline{is right exact}.
\item[a')] \uline{The canonical arrow}
\[
 R^pf_*(E)\otimes M\longrightarrow R^pf_*(E\otimes_Y^{\mathrm L}M)
\]
(\uline{the composite of the canonical arrow} \(R^pf_*(E)\otimes M\longrightarrow H^p(Rf_*(E)\otimes^{\mathrm L}M)\) \uline{and the isomorphism} (3.7) \(H^p(Rf_*(E)\otimes^{\mathrm L}M)\xrightarrow{\sim}R^pf_*(E\otimes_Y^{\mathrm L}M)\)) \uline{is an isomorphism for every} \(M\in\Ob\Qcoh(Y)\).
\item[b)] \uline{The functor} \(M\mapsto R^{p+1}f_*(E\otimes_Y^{\mathrm L}M)\) \uline{from} \(\Qcoh(Y)\) \uline{to} \(\Qcoh(Y)\) \uline{is left exact}.
\item[b')] \uline{The canonical arrow}
\[
 R^{p+1}f_*(E\otimes_Y^{\mathrm L}M)\longrightarrow\underline{\mathrm{Hom}}(H^{-p-1}(Rf_*(E)^{\vee}),M)
\]
(\uline{the composite of the isomorphism} (3.7) \(R^{p+1}f_*(E\otimes_Y^{\mathrm L}M)\simeq H^{p+1}(Rf_*(E)\otimes^{\mathrm L}M)\) \uline{and the canonical arrow} (5.1 (\(*\))) \(H^{p+1}(Rf_*(E)\otimes^{\mathrm L}M)\longrightarrow\underline{\mathrm{Hom}}(H^{-p-1}(Rf_*(E)^{\vee}),M)\)) \uline{is an isomorphism for every} \(M\in\Ob\Qcoh(Y)\).
%%%%%%%%%%  scan idx 277  |  folio 264  %%%%%%%%%%
\SGAPrintPage{264}%
\item[c)]
\uline{For every base change} \(Y'\to Y\), \uline{the canonical arrow}
\[
 R^pf_*(E)\otimes_Y\underline{O}_{Y'}\longrightarrow H^p(Rf_*(E)\otimes_Y^{\mathrm L}\underline{O}_{Y'})
\]
\uline{is an isomorphism}.
\end{enumerate}
\uline{Then we have}:
\[
 \mathrm a)\Longleftrightarrow\mathrm a')\Longleftrightarrow\mathrm b)\Longleftrightarrow\mathrm c)\Longleftarrow\mathrm b').
\]
\uline{Moreover}, \uline{if} \(f\) \uline{is proper}, \(E\) \uline{pseudo-coherent rel.} \uline{to} \(f\) (\uline{hence perfect rel.} \uline{to} \(f\)), \uline{and subject to the validity of} (2.1) \uline{for} \(f\) (\uline{hence in particular if} \(f\) \uline{is projective or if} \(Y\) \uline{is noetherian}), \uline{the preceding conditions are equivalent}.

\uline{Proof}. The implications \(\mathrm a)\Longleftrightarrow\mathrm a')\Longleftrightarrow\mathrm b)\Longleftarrow\mathrm b')\) follow at once from (3.7), (5.1) and (5.2 (iv)), as does the implication \(\mathrm b')\Longleftarrow\mathrm b)\) (for \(f\) proper and \(E\) pseudo-coherent rel. to \(f\)) by (2.1). It remains to prove that \(\mathrm a)\Longleftrightarrow\mathrm c)\). This will follow from the first assertion of the following lemma:

\sgaStableTarget{sga6.e03.lem.5-3-1}\textbf{Lemma 5.3.1.} \uline{Let} \(X\) \uline{be a quasi-separated scheme}, \uline{and let} \(E\in\Ob D^{-}(X)\) \uline{be a complex with locally bounded and quasi-coherent cohomology}. \uline{Then condition} (5.1 a)) \uline{is equivalent to the condition}
\begin{enumerate}[label=\alph*)]
\item[c)] \uline{For every scheme} \(X'\) \uline{over} \(X\), \uline{the canonical arrow}
\[
 H^p(E)\otimes_X\underline{O}_{X'}\longrightarrow H^p(E\otimes_X^{\mathrm L}\underline{O}_{X'})
\]
\uline{is an isomorphism}.
\end{enumerate}
\uline{Moreover}, \uline{if} \(E\) \uline{is pseudo-coherent}, \uline{condition} (5.1 a)) \uline{implies that}, \uline{for every morphism} \(X'\to X\), \uline{the canonical arrow}
\sgaStableTarget{sga6.e03.eq.5-3-1-1}\begin{equation}
 H^{-p-1}(E^{\vee})\otimes_X\underline{O}_{X'}\longrightarrow H^{-p-1}((E\otimes_X^{\mathrm L}\underline{O}_{X'})^{\vee}),\tag{5.3.1.1}
\end{equation}
%%%%%%%%%%  scan idx 278  |  folio 265  %%%%%%%%%%
\SGAPrintPage{265}%
\uline{the composite of} \(H^{-p-1}(E^{\vee})\otimes_X\underline{O}_{X'}\longrightarrow H^{-p-1}(E^{\vee}\otimes_X^{\mathrm L}\underline{O}_{X'})\) \uline{and the canonical arrow} (I 7.1.2) \(H^{-p-1}(E^{\vee}\otimes_X^{\mathrm L}\underline{O}_{X'})\longrightarrow H^{-p-1}((E\otimes_X^{\mathrm L}\underline{O}_{X'})^{\vee})\), \uline{is an isomorphism}.

\uline{Proof}. \(\mathrm a)\Longrightarrow\mathrm c)\). The question being local on \(X'\), we may suppose \(X\) (resp. \(X'\)) affine with ring \(A\) (resp. \(A'\)). By (II 3.5), we may also suppose that \(E=\widetilde F\), where \(F\) is a complex of \(A\)-modules. By hypothesis (condition a')), the arrow
\[
 H^p(F)\otimes_A A'\longrightarrow H^p(F\otimes_A^{\mathrm L}A')
\]
is an isomorphism, and the result follows since \(H^p(E)\otimes_X\underline{O}_{X'}\) (resp. \(H^p(E\otimes_X^{\mathrm L}\underline{O}_{X'})\)) identifies canonically with the sheaf on \(X'\) defined by the \(A'\)-module \(H^p(F)\otimes_A A'\) (resp. \(H^p(F\otimes_A^{\mathrm L}A')\)).

\(\mathrm c)\Longrightarrow\mathrm a)\). In view of (5.2 (iv)), it suffices to prove \(\mathrm c)\Longrightarrow\mathrm a'_{\qcoh})\). The question being local (since \(X\) is quasi-separated), we may suppose \(X\) affine with ring \(A\), and \(E=\widetilde F\), where \(F\) is a complex of \(A\)-modules. Let \(M=\widetilde N\) be a quasi-coherent \(\underline{O}_X\)-Module. We must show that the arrow
\[
 H^p(F)\otimes_A N\longrightarrow H^p(F\otimes_A^{\mathrm L}N)
\]
is an isomorphism. But, putting \(A'=D_A(N)\) (EGA \(\mathrm O_{\mathrm{IV}}\) 18.2.3), \(X'=\Spec(A')\), this follows from c) applied to the canonical morphism \(X'\to X\).

Let us prove the last assertion. The question being local on \(X'\), we may, by (II 2.2.10 b)), suppose \(E\) strictly pseudo-coherent. Moreover, since the assertion involves only a part of \(E\) of bounded length, we may, after truncating, suppose \(E\) strictly perfect. Then \((E\otimes_X^{\mathrm L}\underline{O}_{X'})^{\vee}\) identifies canonically (I 7.1.2) with \(E^{\vee}\otimes_X^{\mathrm L}\underline{O}_{X'}\), and the arrow (5.3.1.1) may be rewritten as the canonical arrow
\sgaStableTarget{sga6.e03.eq.unnamed.occ2}\begin{equation}
 H^{-p-1}(E^{\vee})\otimes_X\underline{O}_{X'}\longrightarrow H^{-p-1}(E^{\vee}\otimes_X^{\mathrm L}\underline{O}_{X'}).\tag{$*$}
\end{equation}
%%%%%%%%%%  scan idx 279  |  folio 266  %%%%%%%%%%
\SGAPrintPage{266}%
If \(E\) satisfies condition (5.1 a)) for the integer \(p\), \(E^{\vee}\) satisfies the same condition for the integer \(-p-1\) by (5.2 (iii)), and the arrow \((*)\) is an isomorphism by the first assertion of (5.3.1). This completes the proof of (5.3.1), and hence of (5.3).

\sgaStableTarget{sga6.e03.rem.5-3-2}\textbf{Remark 5.3.2.} Proposition (5.3) applies in particular when \(E\) is a complex of bounded degrees of quasi-coherent \(\underline{O}_X\)-Modules that are flat and of finite presentation rel. to \(f\): indeed, \(E\) then has finite perfect amplitude rel. to \(f\) by (4.3.2).

%%%%%%%%%%  scan idx 280  |  folio 267  %%%%%%%%%%
\SGAPrintPage{267}%
\sgaStableTarget{sga6.e03.def.5-4}\textbf{Definition 5.4.} \uline{Let} \(X\) \uline{be a topos}, \uline{and denote by} \(\pt(X)\) \uline{the set of its points} (SGA 4 IV 6.1). \uline{Let} \(M\) \uline{be an ordered set}, \(f:\pt(X)\to M\) \uline{a map}, \uline{and} \(x\in\pt(X)\). \uline{We say that} \(f\) \uline{is upper semi-continuous} (\uline{resp.} \uline{lower semi-continuous}, \uline{resp.} \uline{locally constant}) \uline{at} \(x\) \uline{if there exists} \uline{an} \(x\)-\uline{pointed object} \(U\) \uline{of} \(X\) \uline{such that for every point} \(y\) \uline{of} \(U\) \uline{we have} \(f(y)\leq f(x)\) (\uline{resp.} \(f(y)\geq f(x)\), \uline{resp.} \(f(y)=f(x)\)). \uline{We say that} \(f\) \uline{is upper semi-continuous} (\uline{resp.} \uline{lower semi-continuous}, \uline{resp.} \uline{locally constant}) \uline{if it is so at every point}\SGAAuthorFootnote{1}{When \(X\) comes from a topological space, we recover the usual notions.}.

\sgaStableTarget{sga6.e03.ex.5-4-1}\textbf{Examples 5.4.1.} a) Denote by \(M_X\) the constant sheaf with value \(M\) on \(X\), and let \(s\in H^0(X,M)\). The map \(x\mapsto s(x)=x^*(s)\) from \(\pt(X)\) to \(M\) is locally constant. In particular, let \(E\) be a perfect complex on \(X\), supposed locally ringed: the rank of \(E\) (I 6.11) defines a locally constant function on \(\pt(X)\) with values in \(\mathbb Z\).

b) Suppose \(X\) locally ringed, and let \(P\) be an \(\underline{O}_X\)-Module of finite type. For \(x\in\pt(X)\), denote by \(\underline{O}_{X,x}\) the local stalk ring of \(\underline{O}_X\) at \(x\), and by \(k(x)\) its residue field. Then, for \(x\in\pt(X)\), \(P_x\otimes_{\underline{O}_{X,x}}k(x)\) has finite rank over \(k(x)\), and the function \(x\mapsto\operatorname{rg}_{k(x)}(P\otimes_{\underline{O}_{X,x}}k(x))\) is upper semi-continuous. Indeed, if, for \(x\in\pt(X)\), \(P_x\otimes_{\underline{O}_{X,x}}k(x)\) has rank \(n\), there is an epimorphism \(\underline{O}_{X,x}^n\to P_x\) by Nakayama, and, since \(P\) is of finite type, there is an \(x\)-pointed object \(U\) of \(X\) and an epimorphism \(\underline{O}_U^n\to P\,|\,U\), giving the assertion.

\sgaStableTarget{sga6.e03.lem.5-5}\textbf{Lemma 5.5.} \uline{Let} \((X,\underline{O}_X)\) \uline{be a locally ringed topos}, \(E\in\Ob D^{-}(X)\), \uline{and} \(p\in\mathbb Z\). \uline{Then we have}, \uline{with the notation of} (5.4.1 b)):
\begin{enumerate}[label=(\roman*)]
%%%%%%%%%%  scan idx 281  |  folio 268  %%%%%%%%%%
\item\SGAPrintPage{268}%
\uline{If} \(E\) \uline{is perfect}, \uline{the} \(H^i(E_x\otimes_{\underline{O}_{X,x}}^{\mathrm L}k(x))\), \uline{for} \(x\in\pt(X)\), \uline{are vector spaces of finite rank over} \(k(x)\), \uline{all but finitely many zero}, \uline{and the function}
\[
 x\longmapsto\sum(-1)^i\operatorname{rg}_{k(x)}H^i(E_x\otimes_{\underline{O}_{X,x}}^{\mathrm L}k(x))
\]
\uline{is locally constant on} \(\pt(X)\).
\item \uline{If} \(E\) \uline{is} \((p-1)\)-\uline{pseudo-coherent}, \(H^p(E_x\otimes_{\underline{O}_{X,x}}^{\mathrm L}k(x))\) \uline{is}, \uline{for} \(x\in\pt(X)\), \uline{a vector space of finite rank over} \(k(x)\), \uline{and the function}
\[
 x\longmapsto\operatorname{rg}_{k(x)}H^p(E_x\otimes_{\underline{O}_{X,x}}^{\mathrm L}k(x))
\]
\uline{is upper semi-continuous on} \(\pt(X)\). \uline{If moreover the functor} \(M\mapsto H^p(E\otimes^{\mathrm L}M)\) \uline{from} \(\Mod(X)\) \uline{to} \(\Mod(X)\) \uline{is exact}, \(H^p(E)\) \uline{is locally free of finite type and the preceding function is locally constant}.
\item \uline{Suppose} \(E\) \(p\)-\uline{pseudo-coherent} (\uline{resp.} \((p-1)\)-\uline{pseudo-coherent}), \uline{and let} \(x\in\pt(X)\). \uline{If the functor} \(M\mapsto H^p(E_x\otimes_{\underline{O}_{X,x}}^{\mathrm L}M)\) \uline{is right exact} (\uline{resp.} \uline{left exact}) \uline{on the category} \(\Mod(x)\) \uline{of} \(\underline{O}_{X,x}\)-\uline{Modules}, \uline{there is an} \(x\)-\uline{pointed open} \(U\) \uline{of} \(X\) \uline{such that the functor} \(M\mapsto H^p(E\,|\,U\otimes^{\mathrm L}M)\) \uline{is right exact} (\uline{resp.} \uline{left exact}) \uline{on} \(\Mod(U)\).
\end{enumerate}

\uline{Proof}. (i) By (I 4.19.1), \(E_x\otimes_{\underline{O}_{X,x}}^{\mathrm L}k(x)\) is a perfect complex of \(k(x)\)-modules, so (I 5.9 and 3.5) the \(H^i(E_x\otimes_{\underline{O}_{X,x}}^{\mathrm L}k(x))\) have finite rank over \(k(x)\) and all but finitely many are zero. We have
\begin{equation}
 \operatorname{rg}(E_x\otimes_{\underline{O}_{X,x}}^{\mathrm L}k(x))=\sum(-1)^i\operatorname{rg}_{k(x)}H^i(E_x\otimes_{\underline{O}_{X,x}}^{\mathrm L}k(x)),\qquad\mbox{\SGAAuthorFootMark{1}}\tag{5.5.1}
\end{equation}
\SGAAuthorFootText{1}{\(\operatorname{rg}(\ )\) denotes the rank of a perfect complex (I 6.11).}%
as is seen by an immediate dévissage using the additivity of the rank function. In view of (5.4.1 a)), assertion (i) therefore follows from
%%%%%%%%%%  scan idx 282  |  folio 269  %%%%%%%%%%
\SGAPrintPage{269}%
the compatibility of rank with inverse images (I 6.12).

(ii) After truncating and shifting, we may, for the first assertion, suppose \(E\) strictly perfect, \(p=0\), and \(E^i=0\) for \(i\notin[-1,+1]\). For \(x\in\pt(X)\), put \(E_x\otimes_{\underline{O}_{X,x}}^{\mathrm L}k(x)=E_{k(x)}\). By (5.5.1) we have
\[
\begin{gathered}
 \operatorname{rg}_{k(x)}H^0(E_{k(x)})=\operatorname{rg}_{k(x)}H^{-1}(E_{k(x)})+\operatorname{rg}_{k(x)}H^1(E_{k(x)})\\
 +\operatorname{rg}(E_x),
\end{gathered}
\]
so it suffices to prove that the function \(x\mapsto\operatorname{rg}_{k(x)}H^{-1}(E_{k(x)})\) \var{(resp. \(\operatorname{rg}_{k(x)}H^1(E_{k(x)})\) is upper semi-continuous}{(resp. \(\operatorname{rg}_{k(x)}H^1(E_{k(x)})\)) is upper semi-continuous}{EN-VAR-177} (since the function \(x\mapsto\var{\operatorname{rg}(E_x)=\operatorname{rg}(E_{k(X)})}{\operatorname{rg}(E_x)=\operatorname{rg}(E_{k(x)})}{EN-VAR-178}\) is locally constant). For the function \(x\mapsto\operatorname{rg}_{k(x)}H^1(E_{k(x)})\), this is an immediate consequence of (5.4.1 b))\SGAAuthorFootnote{*}{taking into account that \(\var{H^1(E_{k(x)}\simeq H^1(E)_{k(x)}}{H^1(E_{k(x)})\simeq H^1(E)_{k(x)}}{EN-VAR-180}\).}. For the function \(x\mapsto\operatorname{rg}_{k(x)}H^{-1}(E_{k(x)})\), we may suppose \(E^1=0\), i.e. replace \(E\) by \(E'=(0\to E^{-1}\to E^0\to0)\). We already know that the function \(x\mapsto\operatorname{rg}_{k(x)}H^0(E'_{k(x)})\) is upper semi-continuous, so the same holds for the function \(x\mapsto\operatorname{rg}_{k(x)}H^{-1}(E'_{k(x)})=\operatorname{rg}_{k(x)}H^{-1}(E_{k(x)})\), since the function \(x\mapsto\operatorname{rg}_{k(x)}H^0(E'_{k(x)})-\operatorname{rg}_{k(x)}H^{-1}(E'_{k(x)})=\operatorname{rg}(E'_x)\) is locally constant. Let us now prove the second assertion. After truncating, we may suppose \(E\) strictly perfect and \(E^i=0\) for \(i<p-1\). The \(\underline{O}_X\)-Module \(Z'^i(E)=\Coker(E^{i-1}\to E^i)\) is therefore of finite presentation. We have an exact sequence:
\begin{equation}
 0\to H^p(E)\to Z'^p(E)\to E^{p+1}\to Z'^{p+1}(E)\to0.\tag{$*$}
\end{equation}
Since the functor \(M\mapsto H^p(E\otimes^{\mathrm L}M)\) is right and left exact, \(Z'^{p+1}(E)\) and \var{\(Z'^p(E)\)) are flat}{\(Z'^p(E)\) are flat}{EN-VAR-179} by (5.2 (iii)), hence (I 5.8.3) locally free of finite type. The exactness of (*) then implies that \(H^p(E)\) is locally free of finite type, and consequently that the function \(x\mapsto\operatorname{rg}_{k(x)}H^p(E_{k(x)})=\operatorname{rg}(H^p(E)_x)\)
%%%%%%%%%%  scan idx 283  |  folio 270  %%%%%%%%%%
\SGAPrintPage{270}%
is locally constant, which completes the proof of (ii).

(iii) We may suppose \(E\) strictly perfect and \(E^i=0\) for \(i<p\) (resp. \(p-1\)). Saying that the functor \(M\mapsto H^p(E_x\otimes_{\underline{O}_{X,x}}^{\mathrm L}M)\) is right exact (resp. left exact) on \(\Mod(x)\) is equivalent, by (5.2 (iii)), to saying that \(Z'^{p+1}(E_x)\) (resp. \(Z'^p(E_x)\)) (with the notation of (ii)) is flat, hence free of finite type (since it is of finite presentation). Since \(Z'^{p+1}(E_x)\simeq Z'^{p+1}(E)_x\) (resp. \(Z'^p(E_x)\simeq Z'^p(E)_x\)), the usual argument (cf. EGA \(\mathrm O_I\) 5.2.7) shows that there is an \(x\)-pointed object \(U\) of \(X\) such that \(Z'^{p+1}(E\,|\,U)\) (resp. \(Z'^p(E\,|\,U)\)) is free of finite type, and the result follows.

\sgaStableTarget{sga6.e03.cor.5-6}\textbf{Corollary 5.6} (EGA III 7.6.9). \uline{Let} \(X\) \uline{be a scheme}, \(p\in\mathbb Z\), and \(E\in\Ob D^{-}(X)\) \uline{a} \((p-1)\)-\uline{pseudo-coherent complex}. \uline{Then}:
\begin{enumerate}[label=(\roman*)]
\item \uline{For} \(x\in X\), \(H^p(E\otimes_X^{\mathrm L}k(x))\) \uline{has finite rank over} \(k(x)\) \uline{and the function}
\[
 x\longmapsto\operatorname{rg}_{k(x)}H^p(E\otimes_X^{\mathrm L}k(x))
\]
\uline{is constructible and upper semi-continuous on} \(X\).
\item \uline{Suppose} \(X\) \uline{quasi-separated and} \(E\) \uline{with quasi-coherent and locally bounded cohomology}. \uline{If the functor} \(M\mapsto H^p(E\otimes^{\mathrm L}M)\) \uline{is exact on the category of quasi-coherent} \(\underline{O}_X\)-\uline{Modules}, \uline{the function} \(x\mapsto\operatorname{rg}_{k(x)}H^p(E\otimes_X^{\mathrm L}k(x))\) \uline{is locally constant and} \(H^p(E)\) \uline{is locally free of finite type}. \uline{Conversely}, \uline{if} \(X\) \uline{is locally noetherian and reduced and the preceding function is locally constant}, \uline{the functor} \(M\mapsto H^p(E\otimes^{\mathrm L}M)\) \uline{is exact on the category of quasi-coherent} \(\underline{O}_X\)-\uline{Modules}.
\end{enumerate}

\uline{Proof}. If \(X\) is integral with generic point \(x\), the functor \(M\mapsto H^p(E_x\otimes_{\underline{O}_{X,x}}^{\mathrm L}M)\) is exact on the category of \(\underline{O}_{X,x}\)-modules (since \(\underline{O}_{X,x}=k(x)\)), so (5.5 (iii)) there is an open neighborhood \(U\) of \(x\) such that the functor \(M\mapsto H^p(E\,|\,U\otimes^{\mathrm L}M)\) is exact on \(\Mod(U)\), and, on \(U\), the function
%%%%%%%%%%  scan idx 284  |  folio 271  %%%%%%%%%%
\SGAPrintPage{271}%
\(\var{y\mapsto\operatorname{rg}_{k(x)}H^p(E\otimes_X^{\mathrm L}k(y))}{y\mapsto\operatorname{rg}_{k(y)}H^p(E\otimes_X^{\mathrm L}k(y))}{EN-VAR-181}\) is locally constant by (5.5 (ii)). The constructibility assertion of (i) follows at once. The other assertions follow from (5.5 (ii) and (iii)) and (5.2 (iv)), except for the converse of (ii), for which it suffices, after reducing to \(E\) strictly perfect (by localization and truncation), to apply (EGA III 7.6.2 and 7.6.5) to the functor \(T_{-p}=H^p(E\otimes\,.\,)\).

\sgaStableTarget{sga6.e03.prop.5-7}\textbf{Proposition 5.7.} (EGA III 7.7.5 \uline{and} 7.9.4). \uline{Let} \(f:X\to Y\) \uline{be a proper morphism of quasi-compact schemes}, \(E\in\Ob D^b(f)_{\uparf}\) (4.2), \uline{and} \(p\in\mathbb Z\). \uline{Then}, \uline{subject to the validity of} (2.1), \uline{hence for example if} \(f\) \uline{is projective or if} \(Y\) \uline{is noetherian}:
\begin{enumerate}[label=(\roman*)]
\item \uline{For} \(y\in Y\), \uline{the} \var{\(H^i(Rf_*(E)\otimes_Y^{\mathrm L}k(y)\) \uline{are vector spaces of finite rank over}}{\(H^i(Rf_*(E)\otimes_Y^{\mathrm L}k(y))\) \uline{are vector spaces of finite rank over}}{EN-VAR-182} \(k(y)\), \uline{all but finitely many zero}, \uline{and the function}
\[
 \var{y\longmapsto\sum(-1)^i\operatorname{rg}_{k(y)}H^i(Rf_*(E)\otimes_Y^{\mathrm L}k(y)}{y\longmapsto\sum(-1)^i\operatorname{rg}_{k(y)}H^i(Rf_*(E)\otimes_Y^{\mathrm L}k(y))}{EN-VAR-183}
\]
\uline{is locally constant on} \(Y\).
\item \uline{The function}
\[
 \var{y\longmapsto\operatorname{rg}_{k(y)}H^p(Rf_*(E)\otimes_Y^{\mathrm L}k(y)).}{y\longmapsto\operatorname{rg}_{k(y)}H^p(Rf_*(E)\otimes_Y^{\mathrm L}k(y))}{EN-VAR-184}
\]
\uline{is constructible and upper semi-continuous on} \(Y\).
\item \uline{Suppose} \(Y\) \uline{quasi-separated}. \uline{Then}, \uline{if the functor} \(M\mapsto R^pf_*(E\otimes_Y^{\mathrm L}M)\) \uline{is exact on the category of quasi-coherent} \(\underline{O}_Y\)-\uline{Modules}, \(R^pf_*(E)\) \uline{is a locally free} \(\underline{O}_Y\)-\uline{Module of finite type}, \uline{and the function of} (ii) \uline{is locally constant}. \uline{Conversely}, \uline{if} \(Y\) \uline{is noetherian and reduced}, \uline{and if the function of} (ii) \uline{is locally constant}, \uline{the functor} \(M\mapsto R^pf_*(E\otimes_Y^{\mathrm L}M)\) \uline{is exact on the category of quasi-coherent} \(\underline{O}_Y\)-\uline{Modules}.
\end{enumerate}

%%%%%%%%%%  scan idx 285  |  folio 272  %%%%%%%%%%
\SGAPrintPage{272}%
\uline{Proof}. By (4.8), \(Rf_*(E)\) is a perfect complex.

Assertion (i) (resp. (ii)) therefore follows from (5.5 (i)) (resp. (5.6 (i)) ). Finally, if \(M\) is a quasi-coherent \(\underline{O}_Y\)-Module, the canonical arrow
\[
 Rf_*(E)\otimes_Y^{\mathrm L}M\longrightarrow Rf_*(E\otimes_Y^{\mathrm L}M)
\]
is an isomorphism by (3.7), so assertion (iii) follows from (5.6 (ii)).

\sgaStableTarget{sga6.e03.rem.5-8}\textbf{Remarks 5.8.}
\begin{enumerate}[label=\alph*)]
\item In the situation of (5.7), denote by \(i_y:X_y\to X\) the inclusion of the fiber of \(X\) over a point \(y\) of \(Y\). One can show \var{(cf. (EGA III 6 and (IV 3.1.0)}{(cf. (EGA III 6) and (IV 3.1.0))}{EN-VAR-185} that, for \(f\) flat (resp. \(E\) bounded with quasi-coherent and \(Y\)-flat components), we have the Künneth formula:
\[
 H^i(Rf_*(E)\otimes_Y^{\mathrm L}k(y))\xrightarrow{\sim}H^i(X_y,E_y),
\]
where \(E_y=Li_y^*(E)\) (resp. \(i_y^*(E)\)).
%%%%%%%%%%  scan idx 286  |  folio 273  %%%%%%%%%%
\SGAPrintPage{273}%
\item One can develop for complex analytic spaces a body of results analogous to that of this section (exchange, semi-continuity, continuity properties)\SGAAuthorFootnote{1}{It was moreover in this setting that these properties were first formulated (see for example Grauert [1]), before being transposed to Algebraic Geometry in (EGA III 7).}. The essential ingredients are: 1) the finiteness theorem for proper morphisms (``Grauert's theorem''); 2) an analogue of the projection formula (3.7), which may be stated as follows: let \(f:X\to Y\) be a morphism of analytic spaces, where \(X\) is supposed of bounded dimension, and let \(E\in\Ob D^-(X)\), \(M\in\Ob D^-(Y)_{\coh}\) (``coh'' means ``with coherent cohomology''). The canonical arrow
\[
 \var{Rf_*(E)\otimes^{\mathrm L}M)\longrightarrow Rf_*(E\otimes_Y^{\mathrm L}M)}{Rf_*(E)\otimes^{\mathrm L}M\longrightarrow Rf_*(E\otimes_Y^{\mathrm L}M)}{EN-VAR-186}
\]
is an isomorphism (the hypothesis on \(X\) ensures (T.F II 5.12) that \(f_*\) has finite cohomological dimension, hence that \(Rf_*\) is defined on \(D(X)\)).

Unlike statement 1), statement 2) is essentially trivial, and follows from a dévissage analogous to that of the proof of (3.7); we moreover have greater freedom since \(Rf_*\) is defined on \(D(X)\), and in particular no hypothesis on \(E\) is needed (apart from \(E\in\Ob D^-(X)\)).
\end{enumerate}

\sgaStableTarget{sga6.e03.bibliography}\subsection*{BIBLIOGRAPHY}
\begin{enumerate}[label={[\arabic*]}]
\item[{[1]}] Grauert, H., Ein Theorem der Analytischen Garbentheorie \ldots, Pub. IHES, n\textsuperscript{o}5.
\item[{[H]}] Hartshorne, R., Residues and Duality, Lecture Notes n\textsuperscript{o}20, Springer, 1966.
\item[{[T.F]}] Godement, R., Théorie des faisceaux, Hermann, 1958.
\end{enumerate}

% SGA_NATIVE_CONTINUATION_273
\providecommand{\Ob}{\operatorname{Ob}}
\providecommand{\Mod}{\operatorname{Mod}}
\providecommand{\Coh}{\operatorname{Coh}}
\providecommand{\Loclib}{\operatorname{Loclib}}
\providecommand{\coh}{\mathrm{coh}}
\providecommand{\parf}{\mathrm{parf}}
\providecommand{\naif}{\mathrm{naif}}
\providecommand{\loc}{\mathrm{loc}}
\providecommand{\fin}{\mathrm{fin}}
\providecommand{\clA}{\operatorname{cl}}
\providecommand{\Rg}{\operatorname{rg}}
\providecommand{\varinjlim}{\varinjlim}
\providecommand{\varprojlim}{\varprojlim}
\providecommand{\id}{\operatorname{id}}
\providecommand{\Th}{\operatorname{Th}}
\providecommand{\Hom}{\operatorname{Hom}}
\providecommand{\Dcat}{\mathrm D}
\providecommand{\Kcat}{\mathrm K}
\providecommand{\Kdot}{K_{\bullet}}
\providecommand{\Kprime}{K'}
\providecommand{\naif}{\mathrm{na\ddot{\i}f}}


% SGA_REVIEWED_FRENCH_CONTINUATION_273
% Exposé IV macros
\renewcommand{\Ob}{\operatorname{ob}}
\providecommand{\Mod}{\operatorname{Mod}}
\providecommand{\Coh}{\operatorname{Coh}}
\providecommand{\Loclib}{\operatorname{Loclib}}
\providecommand{\coh}{\mathrm{coh}}
\providecommand{\parf}{\mathrm{parf}}
\providecommand{\naif}{\text{naïf}}
\providecommand{\loc}{\mathrm{loc}}
\providecommand{\fin}{\mathrm{fin}}
\providecommand{\clA}{\operatorname{cl}}
\providecommand{\Rg}{\operatorname{rg}}
\providecommand{\varinjlim}{\varinjlim}
\providecommand{\varprojlim}{\varprojlim}
\renewcommand{\id}{\operatorname{id}}
\providecommand{\Th}{\operatorname{Th}}
\providecommand{\Hom}{\operatorname{Hom}}
\providecommand{\Dcat}{\mathrm D}
\providecommand{\Kcat}{\mathrm K}
\providecommand{\Kdot}{K_{\bullet}}
\providecommand{\Kprime}{K'}
\providecommand{\naif}{\text{naïf}}



% SGA_SOURCE_MACRO_RESTORATION_273: native provide must not shadow the reviewed accented label.
\renewcommand{\naif}{\text{naïf}}

%%%%%%%%%%  scan idx 287  |  folio 274  %%%%%%%%%%
\SGAPrintPage{274}%
\sgaStableTarget{sga6.e04.expose}\section*{\underline{EXPOSÉ IV}\\[3pt] \underline{\underline{GROTHENDIECK GROUPS OF RINGED TOPOI}}}
% The reviewed Exposé IV source renews this local abbreviation with one rule.
\renewcommand{\parf}{\underline{\mathrm{parf}}}
\begin{center}
\underline{by L. ILLUSIE}
\end{center}

\sgaStableTarget{sga6.e04.sec.1}\section*{1. Reminders and generalities on Grothendieck groups}

\sgaStableTarget{sga6.e04.para.1-1}1.1. Let \(C\) be a \uline{triangulated} category. Recall (I 6.3 and SGA 5 VIII 2) that a map \(f\) from \(\Ob C\) to an abelian group \(G\) is called \uline{additive} if \(f(E)=f(E')+f(E'')\) for every distinguished triangle \(E'\longrightarrow E\longrightarrow E''\longrightarrow E'[1]\). The additive maps from \(\Ob C\) to \(G\) form an abelian group, depending functorially on \(G\). The functor thus obtained is represented by an abelian group \(k(C)\) and a universal additive map \(\cl_C:\Ob C\longrightarrow k(C)\) (denoted by \(\cl\) when there is no risk of confusion).

The group \(k(C)\) depends functorially on \(C\) with respect to exact functors. If \(C,C'\) are triangulated categories, two isomorphic exact functors from \(C\) to \(C'\) induce the same homomorphism \(k(C)\to k(C')\); in particular, if \(u:C\to C'\) is an exact functor that is an equivalence of categories, \(k(u):k(C)\to k(C')\) is an isomorphism.

\sgaStableTarget{sga6.e04.lem.1-2}\textbf{Lemma 1.2.} \uline{Let} \(C\) \uline{be a triangulated category, and let} \(L\in\Ob C\).
\begin{enumerate}[label=\alph*)]
\item \uline{For every} \(n\in\mathbb Z\)\uline{, we have}
\[
 \cl(L[n])=(-1)^n\cl(L).
\]
\uline{If} \(L',L''\) \uline{are objects of} \(C\) \uline{such that} \(L\simeq L'\oplus L''\)\uline{, then} \(\cl(L)=\cl(L')+\cl(L'')\).
\item \uline{Suppose there exists} \(n\in\mathbb Z\) \uline{such that} \(\coprod_{i\geq0} L[2ni]\) \uline{is representable. Then} \(\cl(L)=0\).
\end{enumerate}

\uline{Proof.} Assertion a) follows directly from the definitions (cf. SGA 5 VIII 2). For b), it suffices to observe that we have

%%%%%%%%%%  scan idx 288  |  folio 275  %%%%%%%%%%
\SGAPrintPage{275}%
\[
 \coprod_{i\geq0} L[2ni]\simeq
 L\oplus\left(\coprod_{i\geq0}L[2ni]\right)[2n],
\]
and to apply a).

\sgaStableTarget{sga6.e04.lem.1-3}\textbf{Lemma 1.3.} \uline{Let} \(C\) \uline{be a triangulated category, and} \(A\) \uline{a thick subcategory} \(([\mathrm V]\ \mathrm I\ 2.1)\)\uline{. The inclusion and quotient functors define an exact sequence}
\[
 k(A)\longrightarrow k(C)\longrightarrow k(C/A)\longrightarrow 0.
\]

\uline{Proof.} See (SGA 5 VIII 3.1).

\sgaStableTarget{sga6.e04.lem.1-4}\textbf{Lemma 1.4.} \uline{Let} \(A\) \uline{be an additive (resp. abelian) category, and} \(f\) \uline{an additive function} (1.1) \uline{on} \(\Ob K^b(A)\) \uline{(resp.} \(\Ob D^b(A)\)\uline{). For} \(E\in\Ob K^b(A)\) \uline{(resp.} \(D^b(A)\)\uline{) we have}
\[
 f(E)=\sum(-1)^i f(E^i)
\]
\uline{(resp.}
\[
 f(E)=\sum(-1)^i f(H^i(E))
\]).

\uline{Proof.} Left as an exercise to the reader.

\sgaStableTarget{sga6.e04.para.1-5}1.5. Let \(C\) be an \uline{abelian} category and \(A\) a full subcategory of \(C\), stable under finite sums. A map \(f\) from \(\Ob A\) to an abelian group \(G\) is called \uline{additive} if \(\sum(-1)^i f(E^i)=0\) for every acyclic complex \(E\) of \(C\), with bounded degrees and components belonging to \(A\). When \(A\) is stable under kernels of epimorphisms (I 1.2), this amounts to saying that \(f(L)=f(L')+f(L'')\) for every exact sequence \(0\longrightarrow L'\longrightarrow L\longrightarrow L''\longrightarrow0\) of \(C\) whose terms belong to \(A\). The additive maps from \(\Ob A\) to \(G\) form an abelian group depending functorially on \(G\). Denote by \(K^{b,\emptyset}(A)\) the thick subcategory of \(K^b(A)\) formed by the acyclic complexes. It follows easily from (1.3) and (1.4) that the map
\sgaStableTarget{sga6.e04.eq.1-5-1}\begin{equation}\tag{1.5.1}
 \cl_A:\Ob A\longrightarrow k\bigl(K^b(A)/K^{b,\emptyset}(A)\bigr)
\end{equation}

%%%%%%%%%%  scan idx 289  |  folio 276  %%%%%%%%%%
\SGAPrintPage{276}%
induced by \(\cl_{K^b(A)/K^{b,\emptyset}(A)}\) (1.1) is a universal additive map, i.e. the pair \(\left(k\bigl(K^b(A)/K^{b,\emptyset}(A)\bigr),\cl_A\right)\) represents the functor assigning to each abelian group \(G\) the abelian group of additive maps from \(\Ob A\) to \(G\). We shall put
\sgaStableTarget{sga6.e04.eq.1-5-2}\begin{equation}\tag{1.5.2}
 k\bigl(K^b(A)/K^{b,\emptyset}(A)\bigr)=k(A)
\end{equation}
(we shall write \(k(A,C)\) instead of \(k(A)\) when we wish to avoid any ambiguity about the ambient abelian category)\SGAAuthorFootnote{1}{The notation \(k(A)\) was introduced in (I 6.3) to denote the group \(k(K^b(A))\), which is not in general isomorphic to the group defined by (1.5.2). We therefore change convention here: in the present exposé, the notation \(k(A)\) will denote exclusively the group defined by (1.5.2).}.
In particular, for \(A=C\), we have
\sgaStableTarget{sga6.e04.eq.1-5-3}\begin{equation}\tag{1.5.3}
 k(D^b(C))=k(C).
\end{equation}

\sgaStableTarget{sga6.e04.lem.1-6}\textbf{Lemma 1.6.} \uline{Let} \(C\) \uline{be an abelian category, and} \(A\) \uline{a thick subcategory (i.e. a full subcategory stable under kernels, cokernels, and extensions). Denote by} \(D_A(C)\) \uline{the full subcategory of} \(D(C)\) \uline{formed by the objects} \(E\) \uline{such that} \(H^i(E)\in\Ob A\) \uline{for every} \(i\)\uline{. Then} \(D_A(C)\) \uline{is a triangulated subcategory} \(([\mathrm V]\ \mathrm I\ 1.2.3)\) \uline{of} \(D(C)\)\uline{, and the canonical functor} \(D^b(A)\longrightarrow D_A^b(C)=D^b(C)\cap D_A(C)\) \uline{induces an isomorphism}
\[
 k(A)\xrightarrow{\sim} k(D_A^b(C)),
\]
\uline{where} \(k(A)\) \uline{is defined by} (1.5.2).

\uline{Proof.} The first assertion is an immediate consequence of the long exact cohomology sequence. For the second, consider the map \(f:\Ob D_A^b(C)\longrightarrow k(A)\) defined by
\[
 f(E)=\sum(-1)^i\cl_A(H^i(E)),
\]
where \(\cl_A\) is the map (1.5.1). By the long exact cohomology sequence, \(f\) is additive (1.1), hence defines a homomorphism
\[
 \varphi:k(D_A^b(C))\longrightarrow k(A).
\]

%%%%%%%%%%  scan idx 290  |  folio 277  %%%%%%%%%%
\SGAPrintPage{277}%
It follows at once from (1.4) that \(\varphi\) is inverse to the canonical homomorphism \(k(A)\longrightarrow k(D_A^b(C))\).

\sgaStableTarget{sga6.e04.rem.1-6-1}\textbf{Remark 1.6.1.} In the situation of (1.6), the canonical functor \(D^b(A)\longrightarrow D_A^b(C)\) is not in general an equivalence of categories, as shown by the example where \(S\) is a scheme, \(C\) the category of \(\underline{O}_S\)-Modules, and \(A\) the category of quasi-coherent \(\underline{O}_S\)-Modules (II Ap. I 0.3). See, however, the criterion (II 2.1.3).

\sgaStableTarget{sga6.e04.def.1-7}\textbf{Definition 1.7.} \uline{Let} \(C\) \uline{be a triangulated (resp. abelian) category, and} \(C'\) \uline{a subset of} \(\Ob C\)\uline{. We shall say that} \(C'\) \uline{is an} exact \uline{subset (cf.} EGA III 3.1.1\uline{) if, for every distinguished triangle} \(E'\to E\to E''\to E'[1]\) \uline{(resp. every exact sequence} \(0\to E'\to E\to E''\to0\)\uline{) such that two of the terms} \(E',E,E''\) \uline{belong to} \(C'\)\uline{, the third also belongs to it. We shall say that} \(C'\) \uline{is} dense \uline{in} \(\Ob C\) \uline{if} \(\Ob C\) \uline{is the smallest exact subset of} \(\Ob C\) \uline{containing} \(C'\).

\sgaStableTarget{sga6.e04.lem.1-8}\textbf{Lemma 1.8.} \uline{Let} \(C\) \uline{be a triangulated (resp. abelian) category, and} \(C'\) \uline{a dense subset of} \(\Ob C\)\uline{. Then the elements} \(\cl_C(E)\)\uline{, for} \(E\in C'\)\uline{, generate} \(k(C)\).

\uline{Proof.} Put \(C'=C_0\), and, for \(i\geq1\), define \(C_i\) inductively as the set of objects of \(C\) appearing in a distinguished triangle (resp. a short exact sequence) two of whose terms belong to \(C_{i-1}\). The \(C_i\) form an increasing filtration of \(\Ob C\), which is exhaustive by hypothesis. For \(i\geq0\), denote by \(k(C)_i\) the subgroup of \(k(C)\) generated by the elements \(\cl_C(E)\) for \(E\in C_i\). It is immediate from the definitions that, for every \(i\geq0\), \(k(C)_i=k(C)_{i+1}\), hence \(k(C)_i=k(C)_0\). Now let \(x\in k(C)\);

%%%%%%%%%%  scan idx 291  |  folio 278  %%%%%%%%%%
\SGAPrintPage{278}%
as is easily seen, we can write \(x=\cl(E')-\cl(E'')\) for \(E',E''\in\Ob C\) (in the first case, we can even arrange that \(E''=0\)). Since the filtration \((C_i)\) is exhaustive, there is an integer \(i\geq0\) such that \(E'\) and \(E''\) belong to \(C_i\); then \(x\in k(C)_i=k(C)_0\), as required.

\sgaStableTarget{sga6.e04.exercise.1-9}\textbf{Exercise 1.9.} Let \(C\) be an abelian category, and \(A\) a full subcategory stable under finite sums, subobjects, and quotients, such that \(\Ob A\) is dense in \(\Ob C\). Then the homomorphism \(k(A,C)\longrightarrow k(C)\) defined by the additive function \(E\mapsto\cl_C(E)\) on \(\Ob A\) (cf. (1.5)) is an isomorphism.

\uline{Hint}: Put \(A=A_0\), and define \(A_i\), for \(i\geq1\), as the full subcategory of \(C\) formed by the objects of \(C\) that are extensions of two objects of \(C_{i-1}\). Prove that \(A_i\) satisfies the same hypotheses as \(A\), then show that the canonical homomorphism \(\varinjlim k(A_i,C)\longrightarrow k(C)\) is an isomorphism (the desired conclusion follows at once).

\sgaStableTarget{sga6.e04.sec.2}\section*{2. The functors \(K_{\bullet}\) and \(K^{\bullet}\) of a ringed topos}

\sgaStableTarget{sga6.e04.para.2-1}2.1. Let \((X,A)\) be a ringed topos. Recall the following notation:

\(D^b({}_A X)_{\coh}\) = category of pseudo-coherent complexes with bounded cohomology of left \(A\)-Modules (I 2.15).

\(D({}_A X)_{\uparf}\) = category of complexes of left \(A\)-Modules of finite perfect amplitude (I 5.8.4).

\(\Loclib({}_A X)\) = category of left \(A\)-Modules locally direct summands of free Modules of finite type (II 2.0).

\(\Coh({}_A X)\) = category of coherent left \(A\)-Modules (relative to the category of free \(A\)-Modules of finite type) (I 3.1).

%%%%%%%%%%  scan idx 292  |  folio 279  %%%%%%%%%%
\SGAPrintPage{279}%
The last two categories are full subcategories of the category \(\Mod({}_A X)\) of left \(A\)-Modules, and the first two are full triangulated subcategories of \(D({}_A X)\) (\(= D(\Mod({}_A X))\)).

\sgaStableTarget{sga6.e04.def.2-2}\textbf{Definition 2.2.} \uline{Let} \((X,A)\) \uline{be a ringed topos. Unless otherwise stated, we shall put}
\[
 \begin{gathered}
  K_\bullet({}_A X)=k(D^b({}_A X)_{\coh}),\\
  K^{\bullet}({}_A X)=k(D({}_A X)_{\uparf}),
 \end{gathered}
\]
\uline{where} \(k(\ )\) \uline{is defined in} (1.1)\uline{, and}
\[
 \begin{gathered}
  K_\bullet({}_A X)_{\naif}=k(\Coh({}_A X)),\\
  K^{\bullet}({}_A X)_{\naif}=k(\Loclib({}_A X)),
 \end{gathered}
\]
\uline{where} \(k(\ )\) \uline{is defined by} (1.5.2)\uline{, the ambient abelian category being} \(\Mod({}_A X)\)\SGAAuthorFootnote{1}{We shall allow ourselves to omit the index \(A\) when no confusion can result. Sometimes we shall also drop the index ``naïf'', but shall indicate this explicitly.}\uline{. We shall say that} \(K_\bullet({}_A X)\) \uline{(resp.} \(K^{\bullet}({}_A X)\)\uline{) is the Grothendieck group of pseudo-coherent (resp. perfect) complexes of left} \(A\)\uline{-Modules, and that} \(K_\bullet({}_A X)_{\naif}\) \uline{(resp.} \(K^{\bullet}({}_A X)_{\naif}\)\uline{) is the Grothendieck group of coherent (resp. locally free) left} \(A\)\uline{-Modules.}

\sgaStableTarget{sga6.e04.rem.2-3}\textbf{Remarks 2.3.}
\begin{enumerate}[label=\alph*)]
\item One can define variants of \(K_\bullet(\ )\) and \(K^{\bullet}(\ )\) by putting
\[
 K_\bullet^{\loc}({}_A X)=k(D^{\loc b}({}_A X)_{\coh})
\]
(where ``\(\loc b\)'' means ``with locally bounded cohomology''), and
\[
 K^{\bullet}_{\loc}({}_A X)=k(D({}_A X)_{\mathrm{parf}}).
\]
These groups coincide with those defined in (2.2) when the final object of \(X\) is quasi-compact, but in general are less suited to intersection calculations.

%%%%%%%%%%  scan idx 293  |  folio 280  %%%%%%%%%%
\item\SGAPrintPage{280}%
The group \(K_\bullet({}_A X)_{\naif}\) is practically of interest only when \(A\) is coherent (otherwise there may well be very few coherent sheaves \ldots).
\item Considering the \(k\) of triangulated categories that are ``too large'' is of no interest. Thus, from (1.2 b)) one easily deduces the relations
\[
 k(D({}_A X))=k(D^-({}_A X))=k(D^+({}_A X))=k(D^b({}_A X))=k(D^-({}_A X)_{\coh})=0
\]
(for the last equality, observe that \(L\in\Ob D^-({}_A X)_{\coh}\) implies \(\coprod_{i\geq0}L[2i]\in\Ob D^-({}_A X)_{\coh}\)).
\end{enumerate}

\sgaStableTarget{sga6.e04.para.2-4}2.4. \uline{Case of a coherent sheaf of rings}. Let \((X,A)\) be a ringed topos, with \(A\) supposed \uline{coherent} as a left \(A\)-Module (I 3.1 and 3.6). Every coherent left \(A\)-Module is a pseudo-coherent complex (I 3.5), so we have a fully faithful functor
\sgaStableTarget{sga6.e04.eq.2-4-1}\begin{equation}\tag{2.4.1}
 \Coh({}_A X)\hookrightarrow D^b({}_A X)_{\coh}.
\end{equation}
Now (I 3.4) \(\Coh({}_A X)\) is a thick subcategory of \(\Mod({}_A X)\), and, by (I 3.5), \(D^b({}_A X)_{\coh}\) is the full subcategory of \(D^b({}_A X)\) formed by the objects \(L\) such that \(H^i(L)\in\Ob\Coh({}_A X)\) for every \(i\). Consequently, it follows from (1.6) that \uline{the homomorphism}
\sgaStableTarget{sga6.e04.eq.2-4-2}\begin{equation}\tag{2.4.2}
 K_\bullet({}_A X)_{\naif}\longrightarrow K_\bullet({}_A X)
\end{equation}
\uline{defined by} (2.4.1) \uline{is an isomorphism}.

\sgaStableTarget{sga6.e04.para.2-5}2.5. \uline{The homomorphism} \(K^{\bullet}\to K_\bullet\). Let \((X,A)\) be a ringed topos. The inclusion functor \(D({}_A X)_{\uparf}\hookrightarrow D^b({}_A X)_{\coh}\) defines a homomorphism
\sgaStableTarget{sga6.e04.eq.2-5-1}\begin{equation}\tag{2.5.1}
 K^{\bullet}({}_A X)\longrightarrow K_\bullet({}_A X).
\end{equation}
This is in general neither injective nor surjective. However, make the following hypotheses:
\begin{enumerate}[label=(\roman*)]

%%%%%%%%%%  scan idx 294  |  folio 281  %%%%%%%%%%
\item\SGAPrintPage{281}%
the final object of \(X\) is quasi-compact;
\item \(X\) has enough points, and, for every point \(x\) of \(X\), every left \(A_x\)-module has finite tor-dimension.
\end{enumerate}
Then, by (I 5.10), \(D({}_A X)_{\parf}=D^b({}_A X)_{\coh}\), and the homomorphism (2.5.1) is an isomorphism.

Hypotheses (i) and (ii) are satisfied, for example, by a compact space ringed by regular local rings.

\sgaStableTarget{sga6.e04.para.2-6}2.6. \uline{Pairings between} \(K^{\bullet}\) \uline{and} \(K^{\bullet}\) (\uline{noncommutative case}). Let
\[
 f:(X',A')\longrightarrow (X,A)
\]
be a morphism of ringed topoi. Denote by \(\Mod({}_{A'} X'_A)\) the category of bi-Modules, left over \(A'\) and right over \(f^{-1}(A)\), and by \(\Loclib_{A'}({}_{A'} X'_A)\) the full subcategory formed by the objects belonging to \(\Loclib({}_{A'} X')\). By analogy with (2.2), we shall put
\[
 K^{\bullet}_{\naif}({}_{A'} X'_A)_{A'}=k(\Loclib_{A'}({}_{A'} X'_A)).
\]
The tensor product
\[
 \otimes_A:\Mod({}_{A'} X'_A)\times\Mod({}_A X)\longrightarrow \Mod({}_{A'} X')
\]
\[
 (E',E)\longmapsto E'\otimes_{f^{-1}(A)}f^{-1}(E)
\]
induces a functor
\sgaStableTarget{sga6.e04.eq.2-6-1}\begin{equation}\tag{2.6.1}
 \Loclib_{A'}({}_{A'} X'_A)\times\Loclib({}_A X)\longrightarrow\Loclib({}_{A'} X').
\end{equation}
The map \((E',E)\mapsto\cl(E'\otimes_A E)\) deduced from (2.6.1) is biadditive in the sense of (1.5), and therefore defines a homomorphism
\sgaStableTarget{sga6.e04.eq.2-6-2}\begin{equation}\tag{2.6.2}
 K^{\bullet}_{\naif}({}_{A'} X'_A)_{A'}\otimes K^{\bullet}_{\naif}({}_A X)\longrightarrow K^{\bullet}_{\naif}({}_{A'} X').
\end{equation}
Denote by \(D({}_{A'} X'_A)\) the derived category of \(\Mod({}_{A'} X'_A)\), and by \(D({}_{A'} X'_A)_{A'\text{-}\parf}\) the full subcategory formed by the complexes of finite perfect amplitude rel. to \(A'\) (i.e. as objects of \(D({}_{A'} X')\)), and put, as in (2.2),
\[
 \var{K^{\bullet}({}_{A'} X'_A)_{A'}=k(D({}_{A'} X'_A)_{A'\text{-}\parf}.}{K^{\bullet}({}_{A'} X'_A)_{A'}=k(D({}_{A'} X'_A)_{A'\text{-}\parf}).}{EN-VAR-187}
\]

%%%%%%%%%%  scan idx 295  |  folio 282  %%%%%%%%%%
\SGAPrintPage{282}%
The inclusion functors plainly define homomorphisms
\sgaStableTarget{sga6.e04.eq.2-6-3-1}\sgaStableTarget{sga6.e04.eq.2-6-3-2}\sgaStableTarget{sga6.e04.eq.2-6-3-3}\begin{align}
 K^{\bullet}_{\naif}({}_{A'} X'_A)_{A'}&\longrightarrow K^{\bullet}({}_{A'} X'_A)_{A'}, \tag{2.6.3.1}\\
 K^{\bullet}_{\naif}({}_A X)&\longrightarrow K^{\bullet}({}_A X), \tag{2.6.3.2}\\
 K^{\bullet}_{\naif}({}_{A'} X')&\longrightarrow K^{\bullet}({}_{A'} X'). \tag{2.6.3.3}
\end{align}
On the other hand, we know (I 4.18, 4.19) that the derived functor of \(\otimes_A\),
\[
 \otimes_A^{\mathrm L}:D^-({}_{A'} X'_A)\times D^-({}_A X)\longrightarrow D^-({}_{A'} X'),
\]
induces an (exact) functor
\sgaStableTarget{sga6.e04.eq.2-6-4}\begin{equation}\tag{2.6.4}
 D({}_{A'} X'_A)_{A'\text{-}\parf}\times D({}_A X)_{\parf}
 \longrightarrow D({}_{A'} X')_{\parf}.
\end{equation}
The latter defines a homomorphism
\sgaStableTarget{sga6.e04.eq.2-6-5}\begin{equation}\tag{2.6.5}
 K^{\bullet}({}_{A'} X'_A)_{A'}\otimes K^{\bullet}({}_A X)\longrightarrow K^{\bullet}({}_{A'} X').
\end{equation}
Since the functor (2.6.1) is induced by (2.6.4), the homomorphisms (2.6.2) and (2.6.5) are compatible with the homomorphisms (2.6.3.1), (2.6.3.2), and (2.6.3.3); in other words, we have a commutative diagram
\sgaStableTarget{sga6.e04.eq.2-6-6}\[
\begin{tikzcd}[column sep=tiny]
K^{\bullet}_{\naif}({}_{A'} X'_A)_{A'}\otimes K^{\bullet}_{\naif}({}_A X)\arrow[r]\arrow[d]
&K^{\bullet}_{\naif}({}_{A'} X')\arrow[d]\\
K^{\bullet}({}_{A'} X'_A)_{A'}\otimes K^{\bullet}({}_A X)\arrow[r]
&K^{\bullet}({}_{A'} X'),
\end{tikzcd}
\tag{2.6.6}
\]
where the upper and lower horizontal arrows are respectively (2.6.2) and (2.6.5), and the vertical arrows are defined by (2.6.3.1), (2.6.3.2), and (2.6.3.3).

%%%%%%%%%%  scan idx 296  |  folio 283  %%%%%%%%%%
\SGAPrintPage{283}%
\sgaStableTarget{sga6.e04.para.2-7}2.7. \uline{Contravariant law and ring structure on} \(K^{\bullet}\) We shall spell out some important special cases of the pairings (2.6) (their direct study is moreover immediate).
\begin{enumerate}[label=\alph*)]
\item Let \(f:(X',A')\longrightarrow (X,A)\) be a morphism of ringed topoi. The functor \(Lf^*:D^-({}_A X)\longrightarrow D^-({}_{A'} X')\) sends (I 4.19.1) \(D({}_A X)_{\parf}\) to \(D({}_{A'} X')_{\parf}\) (resp. \(\Loclib({}_A X)\) to \(\Loclib({}_{A'} X')\)), and defines a homomorphism
\sgaStableTarget{sga6.e04.eq.2-7-1}\begin{equation}\tag{2.7.1}
 f^*:K^{\bullet}({}_A X)\longrightarrow K^{\bullet}({}_{A'} X')
\end{equation}
(resp.
\sgaStableTarget{sga6.e04.eq.2-7-2}\begin{equation}\tag{2.7.2}
 f^*:K^{\bullet}({}_A X)_{\naif}\longrightarrow K^{\bullet}({}_{A'} X')_{\naif}
\end{equation}
), which can also be interpreted as multiplication by \(\cl(A')\) \var{in (2.6.2) (resp. (2.6.5))}{in (2.6.5) (resp. (2.6.2))}{EN-VAR-188}. We have a commutative diagram (a special case of (2.6.6))
\[
\begin{tikzcd}
K^{\bullet}({}_A X)_{\naif}\arrow[r,"f^*"]\arrow[d]
&K^{\bullet}({}_{A'} X')_{\naif}\arrow[d]\\
K^{\bullet}({}_A X)\arrow[r,"f^*"]
&K^{\bullet}({}_{A'} X'),
\end{tikzcd}
\]
where the left and right vertical arrows are respectively (2.6.3.2) and (2.6.3.3).

Let \(g:(X,A)\to(X'',A'')\) be a morphism of ringed topoi. The transitivity of inverse-image functors implies that the homomorphisms of type (2.7.1) (resp. 2.7.2) satisfy the relation
\sgaStableTarget{sga6.e04.eq.2-7-3}\begin{equation}\tag{2.7.3}
 (gf)^*=f^*g^*.
\end{equation}
Thus \(K^{\bullet}({}_A X)\) (resp. \(K^{\bullet}({}_A X)_{\naif}\)) \uline{is a contravariant functor of} \((X,A)\).

\item Let \((X,\underline{O}_X)\) be a commutatively ringed topos. By the associativity and commutativity of the tensor product, \uline{the homomorphism} (2.6.5) (resp. (2.6.2)) (where we put \(A=\underline{O}_X\), \((X',A')=(X,A)\), \(f=\Id\)) \uline{endows} \(K^{\bullet}(X)\) (resp. \(K^{\bullet}(X)_{\naif}\)) \uline{with the structure of a commutative ring with unit}
%%%%%%%%%%  scan idx 297  |  folio 284  %%%%%%%%%%
\SGAPrintPage{284}%
\((1=\cl(\underline{O}_X))\), \uline{and the canonical homomorphism} (2.6.3.2)
\[
 K^{\bullet}(X)_{\naif}\longrightarrow K^{\bullet}(X)
\]
\uline{is a homomorphism of unital rings}.

Let \(f:(X,\underline{O}_X)\to(Y,\underline{O}_Y)\) be a morphism of topoi ringed by commutative rings. Then compatibility of the tensor product with inverse images shows that \uline{the homomorphism}
\[
 f^*:K^{\bullet}(Y)\longrightarrow K^{\bullet}(X)
\]
(resp.
\[
 f^*:K^{\bullet}(Y)_{\naif}\longrightarrow K^{\bullet}(X)_{\naif}
\]
) \uline{is a homomorphism of unital rings}.
\end{enumerate}

\sgaStableTarget{sga6.e04.para.2-8}2.8. \uline{Augmentation of} \(K^{\bullet}\) \uline{to} \(H^0(\ ,\mathbb Z)\). Let \((X,\underline{O}_X)\) be a locally ringed topos (I 2.15.1). Recall that we have defined (I 6.11 to 6.14) a homomorphism of (unital) rings
\[
 i:H^0(X,\mathbb Z)\longrightarrow K^{\bullet}(X)_{\naif},
\]
thus giving \(K^{\bullet}(X)_{\naif}\) (and consequently also \(K^{\bullet}(X)\) via the canonical homomorphism \var{\(K^{\bullet}(X)_{\naif}\to K^{\bullet}(X)\) the structure}{\(K^{\bullet}(X)_{\naif}\to K^{\bullet}(X)\)) the structure}{EN-VAR-189} of an algebra (with unit) over \(H^0(X,\mathbb Z)\). Moreover (loc.cit.), the algebra \(K^{\bullet}(X)_{\naif}\) (resp. \(K^{\bullet}(X)\)) is augmented to \(H^0(X,\mathbb Z)\) by the \uline{rank}, and the canonical homomorphism \(K^{\bullet}(X)_{\naif}\to K^{\bullet}(X)\) is a homomorphism of augmented algebras.

Let \(f:(X,\underline{O}_X)\to(Y,\underline{O}_Y)\) be a morphism of locally ringed topoi. The homomorphisms \(f^*:K^{\bullet}(Y)\longrightarrow K^{\bullet}(X)\) and \(f^*:K^{\bullet}(Y)_{\naif}\longrightarrow K^{\bullet}(X)_{\naif}\) are compatible (I 6.12, 6.13) with the preceding augmented algebra structures; this is expressed by the commutativity of the diagram (where \(\Rg\) denotes rank):

%%%%%%%%%%  scan idx 298  |  folio 285  %%%%%%%%%%
\SGAPrintPage{285}%
\[
\begin{tikzcd}[column sep=tiny]
H^0(Y,\mathbb Z)\arrow[r,"i"]\arrow[d]
&K^{\bullet}(Y)_{\naif}\arrow[r]\arrow[d]
&K^{\bullet}(Y)\arrow[r,"\Rg"]\arrow[d]
&H^0(Y,\mathbb Z)\arrow[d]\\
H^0(X,\mathbb Z)\arrow[r,"i"]
&K^{\bullet}(X)_{\naif}\arrow[r]
&K^{\bullet}(X)\arrow[r,"\Rg"]
&H^0(X,\mathbb Z).
\end{tikzcd}
\]

\textsuperscript{*}We shall see later in the Seminar \var{(V 3.9.1 and VI\qquad)}{(V 3.9.1 and VI)}{EN-VAR-190} that \(K^{\bullet}(X)_{\naif}\) can be given a richer structure than that of an algebra augmented to \(H^0(X,\mathbb Z)\), namely that of a \(\lambda\)-algebra augmented to \(H^0(X,\mathbb Z)\) regarded as a binomial ring (V 2.8), this structure being compatible with inverse-image homomorphisms. \var{Deligne [5]}{Deligne [1]}{EN-VAR-191} recently proved that \(K^{\bullet}(X)\) can be given an analogous structure (at least when \(X\) has characteristic \(p\geq0\)) and that the canonical homomorphism \(K^{\bullet}(X)_{\naif}\longrightarrow K^{\bullet}(X)\) is a homomorphism of augmented \(\lambda\)-algebras, thus answering a conjecture of Grothendieck (XIV 1) in the affirmative.\textsubscript{*}

\sgaStableTarget{sga6.e04.para.2-9}2.9. \uline{Comparison of} \(K^{\bullet}\) \uline{and} \(K^{\bullet}_{\naif}\). Let \((X,A)\) be a ringed topos. In general there is no reason for the canonical homomorphism
\[
 K^{\bullet}(X)_{\naif}\longrightarrow K^{\bullet}(X)
\]
to be injective or surjective. However, the ``criteria for the existence of global resolutions'' of Exposé II show that it is an isomorphism in the following cases:
\begin{enumerate}[label=(\roman*)]
\item \(X\) is the Zariski topos of a divisorial scheme, separated or noetherian (and \(A=\underline{O}_X\)) (II 2.2.9);
\item \(X\) is the topos defined by a compact space and \(A\) is a soft sheaf (II 2.3.2);
\item \(X\) is a Stein topos whose final object is quasi-compact and \(A=\underline{O}_X\) (II 2.4.5).
\end{enumerate}

%%%%%%%%%%  scan idx 299  |  folio 286  %%%%%%%%%%
\SGAPrintPage{286}%
\sgaStableTarget{sga6.e04.para.2-10}2.10. \uline{Pairing between} \(K^{\bullet}\) \uline{and} \(K_{\bullet}\). Let \((X,\underline{O}_X)\) be a commutatively ringed topos. We know (I 5.8.4 b)) that there is a pairing
\[
 \otimes_A^{\mathrm L}:D(X)_{\uparf}\times D^b(X)_{\coh}
 \longrightarrow D^b(X)_{\coh}.
\]
This defines a homomorphism
\sgaStableTarget{sga6.e04.eq.2-10-1}\begin{equation}\tag{2.10.1}
 K^{\bullet}(X)\otimes K_{\bullet}(X)\longrightarrow K_{\bullet}(X).
\end{equation}
\var{Thanks,to the associativity}{Thanks to the associativity}{EN-VAR-192} of the tensor product, we see that (2.10.1) \uline{endows} \(K_{\bullet}(X)\) \uline{with the structure of a} \(K^{\bullet}(X)\)-\uline{module}. More generally, if \(f:(X,\underline{O}_X)\longrightarrow(Y,\underline{O}_Y)\) is a morphism of commutatively ringed topoi, we can regard \(K_{\bullet}(X)\) as a \(K^{\bullet}(Y)\)-module via the homomorphism \(f^*:K^{\bullet}(Y)\longrightarrow K^{\bullet}(X)\).

\sgaStableTarget{sga6.e04.para.2-11}2.11. \uline{Covariant law on} \(K_{\bullet}\) \uline{and projection formula}. Let \((X,A)\) be a ringed topos. Unlike \(K^{\bullet}({}_A X)\), the group \(K_\bullet({}_A X)\) has no natural variance with respect to \((X,A)\). However, in Algebraic Geometry and Analytic Geometry, \(K_{\bullet}(X)\) tends to be \uline{covariant} in \(X\) for \uline{proper} morphisms. Let us make this precise:

\sgaStableTarget{sga6.e04.para.2-11-1}2.11.1. Let \(f:X\to Y\) be a proper and pseudo-coherent morphism (III 1.2) of quasi-compact schemes. Then, subject to the validity of the finiteness conjecture for \(f\) (III 2.1), hence in particular (III 2.2) when \(Y\) is noetherian (in which case the pseudo-coherence hypothesis on \(f\) is automatically satisfied (III 1.3)) or when \(f\) is projective, the functor \(Rf_*\) induces, by (III 2.6), a functor
\[
 Rf_*:D^b(X)_{\coh}\longrightarrow D^b(Y)_{\coh}.
\]
This defines a homomorphism
\sgaStableTarget{sga6.e04.eq.2-11-1-1}\begin{equation}\tag{2.11.1.1}
 f_*:K_{\bullet}(X)\longrightarrow K_{\bullet}(Y)
\end{equation}

%%%%%%%%%%  scan idx 300  |  folio 287  %%%%%%%%%%
\SGAPrintPage{287}%
This homomorphism is \(K^{\bullet}(Y)\)-linear (cf. (2.10)); in other words, for \(y\in K^{\bullet}(Y)\) and \(x\in K_{\bullet}(X)\), we have the so-called \uline{projection formula}
\sgaStableTarget{sga6.e04.eq.2-11-1-2}\begin{equation}\tag{2.11.1.2}
 f_*(yx)=y f_*(x).
\end{equation}
Indeed, if \(x=\cl(E)\) for \(E\in\Ob D^b(X)_{\coh}\) and \(y=\cl(M)\) for \(M\in\Ob D(Y)_{\uparf}\), we have (III 3.7)
\[
 Rf_*(E)\otimes^{\mathrm L}M\simeq
 Rf_*(E\otimes_Y^{\mathrm L}M).
\]
Finally, if \(g:Y\to Z\) is a morphism satisfying the same hypotheses as \(f\) and for which the finiteness conjecture (III 2.1) holds, the transitivity of direct-image functors gives, for the homomorphisms (2.11.1.1), the formula
\sgaStableTarget{sga6.e04.eq.2-11-1-3}\begin{equation}\tag{2.11.1.3}
 g_*f_*=(gf)_*.
\end{equation}

\sgaStableTarget{sga6.e04.para.2-11-2}2.11.2. Let \(f:X\to Y\) be a proper morphism of finite-dimensional analytic spaces. Then, by Grauert's finiteness theorem [2], \(Rf_*\) sends \(D^b(X)_{\coh}\) to \(D^b(Y)_{\coh}\), and therefore defines a homomorphism
\sgaStableTarget{sga6.e04.eq.2-11-2-1}\begin{equation}\tag{2.11.2.1}
 f_*:K_{\bullet}(X)\longrightarrow K_{\bullet}(Y).
\end{equation}
Like its algebraic analogue (2.11.1.1), this is \(K^{\bullet}(Y)\)-linear, i.e. formula (2.11.1.2) is valid, as follows from the projection formula at the level of derived categories \var{(III 5.9 b))}{(III 5.8 b))}{EN-VAR-199}. We likewise have the transitivity formula (2.11.1.3).

\sgaStableTarget{sga6.e04.para.2-11-3}2.11.3. Once we leave the algebraic or analytic setting, this covariance of \(K_{\bullet}\) for proper morphisms disappears. Thus, let \(f:X\to Y\) be a proper morphism of topological spaces ringed by the sheaf of complex continuous functions. Then, for \(E\in\Ob D^b(X)_{\coh}\), in general we do not have \(Rf_*(E)\in\Ob D^b(Y)_{\coh}\), as is already evident for \(E=\underline{O}_X\) and \(Y\) reduced to a point. In fact, definition (2.2) of \(K_{\bullet}\)

%%%%%%%%%%  scan idx 301  |  folio 288  %%%%%%%%%%
\SGAPrintPage{288}%
appears especially suited to the needs of Algebraic Geometry and Analytic Geometry, and a good definition of the \(K_{\bullet}\) of a topological space remains to be found. It seems moreover that we shall have to leave the setting of \(\underline{O}_X\)-linear complexes and define, for a morphism \(f:X\to Y\) (of topological spaces, resp. of mixed varieties (II App. II), etc.), a suitable class of complexes of topological vector sheaves on \(X\), which should include, for example, when \(f\) is the structural morphism of a mixed variety, the elliptic complexes on \(X\) relative to \(Y\). This is suggested by the finiteness theorem (II App. II), as well as by a similar theorem (where \(Y\) is supposed reduced to a point) proved by Grothendieck in [3].

\sgaStableTarget{sga6.e04.para.2-12}2.12. \uline{Unusual variances of} \(K^{\bullet}\) \uline{and} \(K_{\bullet}\). Let \(f:(X,A)\longrightarrow(Y,B)\) be a morphism of finite tor-dimension of ringed topoi (III 3.1). The functor \(Lf^*:D^-({}_B Y)\longrightarrow D^-({}_A X)\) sends \(D^b({}_B Y)\) to \(D^b({}_A X)\), hence (I 2.16.1) \(D^b({}_B Y)_{\coh}\) to \(D^b({}_A X)_{\coh}\), and therefore defines a homomorphism
\sgaStableTarget{sga6.e04.eq.2-12-1}\begin{equation}\tag{2.12.1}
 f^*:K_\bullet({}_B Y)\longrightarrow K_\bullet({}_A X).
\end{equation}
If \(A\) and \(B\) are commutative, commutation of inverse image with the tensor product shows that the homomorphism (2.12.1) is compatible with the module structures of \(K_{\bullet}(Y)\) and \(K_{\bullet}(X)\) over \(K^{\bullet}(Y)\) and \(K^{\bullet}(X)\) respectively (2.10); in other words, for \(a\in K^{\bullet}(Y)\) and \(y\in K_{\bullet}(Y)\), we have the formula
\sgaStableTarget{sga6.e04.eq.2-12-2}\begin{equation}\tag{2.12.2}
 f^*(ay)=a f^*(y).
\end{equation}
If \(g:(Y,B)\to(Z,C)\) is a morphism of finite tor-dimension of ringed topoi, transitivity of inverse-image functors implies that the homomorphisms (2.12.1) satisfy the formula:
\sgaStableTarget{sga6.e04.eq.2-12-1-1}\begin{equation}\tag{2.12.1.1}
 f^*g^*=(gf)^*.
\end{equation}

%%%%%%%%%%  scan idx 302  |  folio 289  %%%%%%%%%%
\SGAPrintPage{289}%
Suppose now that we are in situation (2.11.1) (resp. (2.11.2)), with \(f\) of finite tor-dimension. Then \var{(III 4.8.1 and 5.9 b))}{(III 4.8.1 and 5.8 b))}{EN-VAR-200} the functor \(Rf_*\) sends \(D(X)_{\uparf}\) to \(D(Y)_{\uparf}\), and therefore defines a homomorphism
\sgaStableTarget{sga6.e04.eq.2-12-3}\begin{equation}\tag{2.12.3}
 f_*:K^{\bullet}(X)\longrightarrow K^{\bullet}(Y).
\end{equation}
This is not in general a ring homomorphism. However, for \(y\in K_{\bullet}(Y)\) \var{(resp. \(K^{\bullet}(Y)\) and}{(resp. \(K^{\bullet}(Y)\)) and}{EN-VAR-193} \(x\in K^{\bullet}(X)\), we have the formula
\sgaStableTarget{sga6.e04.eq.2-12-4}\begin{equation}\tag{2.12.4}
 f_*(f^*(y)x)=y f_*(x),
\end{equation}
where \(f^*\) is the homomorphism (2.12.1) (resp. (2.7.1)). The proof is analogous to that of (2.11.1.2). If \(g:Y\to Z\) is a morphism satisfying the same hypotheses as \(f\), the homomorphisms (2.12.3) satisfy the transitivity formula:
\sgaStableTarget{sga6.e04.eq.2-12-3-1}\begin{equation}\tag{2.12.3.1}
 g_*f_*=(gf)_*.
\end{equation}

The homomorphism (2.12.3) is one of the essential ingredients of the Riemann--Roch theorem, in both Algebraic Geometry and Analytic Geometry (for indications in the analytic case, see (Exposé 0, 5)).

Suppose that \(X\) and \(Y\) are separated and divisorial schemes (II 2.2.5) (as is the case, for example, when \(X\) and \(Y\) possess ample invertible sheaves). Then (2.9 (i)) \(K^{\bullet}(X)\) (resp. \(K^{\bullet}(Y)\)) is identified with \(K^{\bullet}(X)_{\naif}\) (resp. \(K^{\bullet}(Y)_{\naif}\)), and (2.12.3) therefore defines a homomorphism
\sgaStableTarget{sga6.e04.eq.2-12-3.occ2}\begin{equation}\tag{2.12.3'}
 f_*:K^{\bullet}(X)_{\naif}\longrightarrow K^{\bullet}(Y)_{\naif}.
\end{equation}
As already explained in (Exposé 0, 3.1), the definition of (2.12.3') was by no means evident a priori. Indeed, from the standpoint of our Seminar, it is one of the main motivations for the sorites of Exposé I through Exposé IV.

%%%%%%%%%%  scan idx 303  |  folio 290  %%%%%%%%%%
\SGAPrintPage{290}%
\sgaStableTarget{sga6.e04.sec.3}\section*{3. Supplements on the Grothendieck groups of schemes\protect\SGAAuthorFootnote{1}{The results of this section will not be used in Exposés V through VIII.}}

\sgaStableTarget{sga6.e04.prop.3-1-0}\textbf{Proposition 3.1.0.} Let
\[
\begin{tikzcd}
X \arrow[d,"f"'] & X' \arrow[l,"g"']\arrow[d,"f'"]\\
Y & Y' \arrow[l,"h"']
\end{tikzcd}
\]
be a Cartesian square of schemes, where \(Y\) is quasi-compact and \(f\) quasi-compact and quasi-separated, and let \(E\in\Ob D^b(X)\) be a complex with quasi-coherent cohomology. Suppose \(X\) and \(Y'\) tor-independent over \(Y\) (III 1.5). Suppose moreover either \(h\) of finite tor-dimension or \(E\) of finite flat amplitude rel. to \(f\). Under these conditions, the base-change arrow (cf. SGA 4 XVII 4)
\[
 Lh^*\,Rf_*E\longrightarrow Rf'_*\,Lg^*E
\]
is defined and is an isomorphism.

\uline{Proof}: To define the arrow, we use the trivial duality formula (SGA 4 XVII 2.3) for \(h\), which says that there is a canonical functorial isomorphism
\begin{equation}\tag{$*$}
 \Hom(Lh^*L,M)\xrightarrow{\sim}\Hom(L,Rh_*M)
\end{equation}
for \(L\in\Ob D^-(Y)\), \(M\in\Ob D^+(Y')\). Using (III 3.5.2) and (III 3.7.1), we easily verify that \(Lg^*E\in\Ob D^b(X')\) and \(Rf_*E\in\Ob D^-(Y)\), so, applying (*) with \(L=Rf_*E\) and \(M=Rf'_*Lg^*E\), we are reduced to defining an arrow
\[
 Rf_*E\longrightarrow Rf_*Rg_*Lg^*E.
\]
For this we take the image under \(Rf_*\) of the adjunction arrow
\[
 E\longrightarrow Rg_*Lg^*E,
\]
which comes from the trivial duality formula for \(g\). The base-change arrow

%%%%%%%%%%  scan idx 304  |  folio 291  %%%%%%%%%%
\SGAPrintPage{291}%
is therefore defined; it remains to show that it is an isomorphism. For this, we may first suppose \(Y\) and \(Y'\) affine. Then, by the Leray spectral sequence for an open covering of \(X\), we reduce to the case where \(X\) and \(X'\) are affine. We therefore have a cocartesian square of rings
\[
\begin{tikzcd}
B\arrow[r] & B'\\
A\arrow[u]\arrow[r] & A'\arrow[u],
\end{tikzcd}
\]
and we may suppose that \(E\) is defined by a bounded-above complex of \(B\)-modules. The tor-independence of \(A'\) and \(B\) over \(A\), which is written
\[
 B\otimes_A^{\mathrm L}A'\xrightarrow{\sim}B',
\]
then implies that in \(D(A')\) we have
\[
 E\otimes_A^{\mathrm L}A'\xrightarrow{\sim}E\otimes_B^{\mathrm L}B',
\]
whence the proposition.

Translating (3.1.0) into the language of Grothendieck groups, we obtain:

\sgaStableTarget{sga6.e04.prop.3-1-1}\textbf{Proposition 3.1.1.} \uline{Let}
\[
\begin{tikzcd}
X \arrow[d,"f"'] & X' \arrow[l,"g"']\arrow[d,"f'"]\\
Y & Y' \arrow[l,"h"']
\end{tikzcd}
\]

%%%%%%%%%%  scan idx 305  |  folio 292  %%%%%%%%%%
\SGAPrintPage{292}%
\uline{be a Cartesian square of schemes}. \uline{Suppose} \(Y\) \uline{and} \(Y'\) \uline{quasi-compact}, \(X\) \uline{and} \(Y'\) \uline{tor-independent over} \(Y\) (III 1.5), \(f\) \uline{proper}, \uline{and pseudo-coherent}, \uline{and the finiteness conjecture} (III 2.1) \uline{satisfied for} \(f\) \uline{and} \(f'\) (\uline{for example} \(Y\) \uline{and} \(Y'\) \uline{noetherian}, \uline{or} \(f\) \uline{projective}). \uline{Then}, \uline{if} \(f\) \uline{is perfect} (\uline{resp}. \(h\) \uline{of finite tor-dimension}), \uline{we have}, \uline{for} \(x\in K^{\bullet}(X)\) (\uline{resp}. \(K_{\bullet}(X)\)),
\[
 h^* f_*(x)=f'_* g^*(x),
\]
\uline{where} \(h^*,g^*\) \uline{are defined by} (2.7.1) (\uline{resp}. (2.12.1)), \uline{and} \(f_*,f'_*\) \uline{by} (2.12.3) (\uline{resp}. (2.11.1.1)) (\uline{note that by} (III 4.7.1) \(f'\) \uline{is perfect if} \(f\) \uline{is}, \uline{and that by} (III 3.5.1) \(g\) \uline{has finite tor-dimension if} \(h\) \uline{does}).

\sgaStableTarget{sga6.e04.para.3-2}3.2. \uline{Passage to the limit for} \(K^{\bullet}_{\naif}\) \uline{and} \(\var{K}{K_{\bullet}}{EN-VAR-194}\) Let \((X_i,f_{ij})_{i\in I}\) be a filtered projective system of quasi-compact and quasi-separated schemes, with affine transition morphisms (cf. EGA IV 8.2.2). Put \(X=\varprojlim X_i\), and denote by \(f_i:X\to X_i\) the canonical homomorphism.

\sgaStableTarget{sga6.e04.para.3-2-1}3.2.1. The homomorphisms (2.7.2) \(f_i^*:K^{\bullet}(X_i)_{\naif}\longrightarrow K^{\bullet}(X)_{\naif}\) define, by the transitivity formula (2.7.3), a homomorphism
\sgaStableTarget{sga6.e04.eq.3-2-1-1}\begin{equation}\tag{3.2.1.1}
 \varinjlim K^{\bullet}(X_i)_{\naif}\longrightarrow K^{\bullet}(X)_{\naif}.
\end{equation}
I claim that \uline{the homomorphism} (3.2.1.1) \uline{is an isomorphism}. Indeed, let \(E\) be a locally free \(\underline O_X\)-Module of finite type. By (EGA IV 8.5.2 and 8.5.5), there is an index \(i\) and a locally free \(\underline O_{X_i}\)-Module of finite type \(E_i\) such that \(E\simeq f_i^*(E_i)\). By (EGA IV 8.5.2.5), the image of \(\cl(E_i)\) in \(\varinjlim K^{\bullet}(X_i)_{\naif}\) is independent of the choice of \(E_i\); denote it by \(u(E)\). By (EGA IV 8.5.2 and 8.5.6), the map
\(\var{u:\Ob\Loclib(X)\longrightarrow \lim K^{\bullet}(X_i)_{\naif}}{u:\Ob\Loclib(X)\longrightarrow \varinjlim K^{\bullet}(X_i)_{\naif}}{EN-VAR-195}\) thus defined is additive (1.5);

%%%%%%%%%%  scan idx 306  |  folio 293  %%%%%%%%%%
\SGAPrintPage{293}%
it therefore defines a homomorphism \(K^{\bullet}(X)_{\naif}\longrightarrow \varinjlim K^{\bullet}(X_i)_{\naif}\), which is immediately seen, by (EGA IV 8.5.2), to be inverse to (3.2.1.1), proving our assertion.

\sgaStableTarget{sga6.e04.para.3-2-2}3.2.2. In the same way, one defines a canonical homomorphism
\sgaStableTarget{sga6.e04.eq.3-2-2-1}\begin{equation}\tag{3.2.2.1}
 \varinjlim K^{\bullet}(X_i)\longrightarrow K^{\bullet}(X).
\end{equation}
Using Deligne's localization technique (SGA 4 XVII), which implies that the category of perfect complexes on \(X\) is equivalent to the category \(\mathop{\underrightarrow{\mathrm{Lim}}}\) of the analogous categories \(\Parf(X_i)\), one again finds that (3.2.2.1) is an isomorphism.

\sgaStableTarget{sga6.e04.para.3-2-3}3.2.3. Suppose the \(X_i\) and \(X\) \uline{noetherian}, and the morphisms \(f_{ij}\) \uline{flat} (the \(f_i\) then are as well). Then the homomorphisms (2.12.1)
\(f_i^*:K_{\bullet}(X_i)\longrightarrow K_{\bullet}(X)\) define an isomorphism
% SGA display variant EN-VAR-196
\SGACaptureAppState{EN-VAR-196}%
\ifnum\SGAReaderForm=1\relax\SGAAppMark{EN-VAR-196}\fi
\ifnum\SGAReaderForm=0\relax
\begin{equation}\tag{3.2.2.1}
 \varinjlim K_{\bullet}(X_i)\xrightarrow{\sim}K_{\bullet}(X).
\end{equation}
\else
\begin{equation}\tag{3.2.3.1} \varinjlim K_{\bullet}(X_i)\xrightarrow{\sim}K_{\bullet}(X).
\end{equation}
\fi
The proof is analogous to that of (3.2.1), using the fact (2.4) that \(K_{\bullet}(X)\simeq K_{\bullet}(X)_{\naif}\) is computed in terms of coherent sheaves on \(X\).

Since the module structures of the \(K_{\bullet}(X_i)\) over the \(K^{\bullet}(X_i)_{\naif}\) are compatible with the inverse images \(f_{ij}^*\) by (2.7 a)) and (2.12.2), we find, by passing to the limit, that \(\varinjlim K_{\bullet}(X_i)\) has a module structure over \(\varinjlim K^{\bullet}(X_i)_{\naif}\); this is compatible, \var{via the isomorphisms (3.2.1.1) and (3.2.2.1) with}{via the isomorphisms (3.2.1.1) and (3.2.3.1) with}{EN-VAR-197} the module structure of \(K_{\bullet}(X)\) over \(K^{\bullet}(X)_{\naif}\). In other words, we have a commutative diagram, where \var{the vertical arrows are defined by (3.2.1.1) and (3.2.2.1)}{the vertical arrows are defined by (3.2.1.1) and (3.2.3.1)}{EN-VAR-198}:
\[
\begin{tikzcd}[column sep=tiny]
\varinjlim K^{\bullet}(X_i)_{\naif}\otimes \varinjlim K_{\bullet}(X_i)
\arrow[r]\arrow[d,"\simeq"]
&\varinjlim K_{\bullet}(X_i)\arrow[d,"\simeq"]\\
K^{\bullet}(X)_{\naif}\otimes K_{\bullet}(X)\arrow[r]
&K_{\bullet}(X).
\end{tikzcd}
\]

%%%%%%%%%%  scan idx 307  |  folio 294  %%%%%%%%%%
\SGAPrintPage{294}%
\sgaStableTarget{sga6.e04.para.3-3}3.3. \uline{Relative Grothendieck groups}

\sgaStableTarget{sga6.e04.para.3-3-1}3.3.1. Let \(f:X\longrightarrow Y\) be a locally finite type morphism of schemes. We shall put:
\[
\begin{gathered}
 K^{\bullet}(f)=k(D(f)_{\parf}),\\
 K_{\bullet}(f)=k(D^b(f)_{\coh}),
\end{gathered}
\]
where \(D^b(f)_{\coh}\) denotes the category of complexes pseudo-coherent rel. to \(f\) and with bounded cohomology (III 1.2), and \(D(f)_{\parf}\) the category of complexes of finite perfect amplitude rel. to \(f\) (III 4.2). When no ambiguity is possible, we shall sometimes write \(K^{\bullet}(X/Y)\) (resp. \(K_{\bullet}(X/Y)\)) instead of \(K^{\bullet}(f)\) (resp. \(K_{\bullet}(f)\)).

The groups \(K^{\bullet}(f)\) and \(K_{\bullet}(f)\) do not possess the rich structure of their ``absolute'' analogues \(K^{\bullet}(X)\) and \(K_{\bullet}(X)\), since the tensor product does not in general define a homomorphism \(K^{\bullet}(f)\otimes K^{\bullet}(f)\longrightarrow K^{\bullet}(f)\) (resp. \(K^{\bullet}(f)\otimes K_{\bullet}(f)\longrightarrow K_{\bullet}(f)\)). On the other hand, they satisfy interesting variance properties, which we shall review.

\sgaStableTarget{sga6.e04.para.3-3-2}3.3.2. Let
\[
\begin{tikzcd}
X\arrow[rr,"f"]\arrow[dr,"h"'] &&Y\arrow[dl,"g"]\\
&S&
\end{tikzcd}
\]
be a commutative triangle of locally finite type morphisms of schemes.
\begin{enumerate}[label=3.3.2.\arabic*., leftmargin=4em]
\item If \(g\) is pseudo-coherent, we have \(K_{\bullet}(f)=K_{\bullet}(h)\) (III 1.12). If \(g\) is smooth and \(X\) quasi-compact, we have \(K^{\bullet}(f)=K^{\bullet}(h)\) (III 4.4.1).

\item By (III 1.11 and 4.5), the functor \(\otimes_Y^{\mathrm L}\) defines pairings
\[
 D(f)_{\parf}\times D^b(g)_{\coh}\longrightarrow D^b(h)_{\coh},
\]
\[
 D(f)_{\parf}\times D(g)_{\parf}\longrightarrow D(h)_{\parf},
\]
which induce homomorphisms
\[
 K^{\bullet}(f)\otimes K_{\bullet}(g)\longrightarrow K_{\bullet}(h),
\]
\[
 K^{\bullet}(f)\otimes K^{\bullet}(g)\longrightarrow K^{\bullet}(h).
\]

%%%%%%%%%%  scan idx 308  |  folio 295  %%%%%%%%%%
\SGAPrintPage{295}%
We deduce in particular (by putting \(X=Y\), \(f=\Id\)) that \(K_{\bullet}(g)\) is a \(K^{\bullet}(Y)\)-module. On the other hand, if \(f\) is perfect, multiplication by \(\cl(\underline O_X)\in K^{\bullet}(f)\) defines \(K^{\bullet}(Y)\)-linear homomorphisms
\[
 f^*:K_{\bullet}(Y/S)\longrightarrow K_{\bullet}(X/S),
\]
and
\[
 f^*:K^{\bullet}(Y/S)\longrightarrow K^{\bullet}(X/S).
\]

\item Suppose \(S\) quasi-compact, \(X\) and \(Y\) of finite type over \(S\), and \(f\) proper. Then, subject to the validity of the finiteness conjecture (III 2.1), hence in particular when \(Y\) is noetherian or \(f\) is projective, the functor \(Rf_*\) sends \(D^b(h)_{\coh}\) to \(D^b(g)_{\coh}\), and therefore defines a homomorphism
\[
 f_*:K_{\bullet}(X/S)\longrightarrow K_{\bullet}(Y/S).
\]
If moreover \(f\) is perfect, \(Rf_*\) sends \(D(h)_{\parf}\) to \(D(g)_{\parf}\) (III 4.8), and defines a homomorphism
\[
 f_*:K^{\bullet}(X/S)\longrightarrow K^{\bullet}(Y/S).
\]
From the projection formula (III 3.7) we deduce the formula
\[
 f_*(a x)=a f_*(x)
\]
for \(f\) proper, \(a\in K^{\bullet}(S)\), \(\var{x\in K_{\bullet}(X)}{x\in K_{\bullet}(X/S)}{EN-VAR-201}\) (resp. \(f\) proper and perfect, \(a\in K_{\bullet}(S)\), \(x\in K^{\bullet}(X/S)\)).

\item Suppose \(S\) noetherian, \(X\) and \(Y\) of finite type over \(S\), and \(f\) compactifiable and perfect. Then (III 4.9.1) the functor \(f^!\) sends \(D(g)_{\parf}\) to \(D(h)_{\parf}\), and defines a homomorphism
\[
 f^!:K^{\bullet}(Y/S)\longrightarrow K^{\bullet}(X/S).
\]
\end{enumerate}

\sgaStableTarget{sga6.e04.para.3-3-3}3.3.3. Under the hypotheses of (III 4.7), the functor \(\otimes_S^{\mathrm L}\) defines pairings
\[
 D(f_1)_{\parf}\times D^b(f_2)_{\coh}\longrightarrow D^b(f)_{\coh},
\]
\[
 D(f_1)_{\parf}\times D(f_2)_{\parf}\longrightarrow D(f)_{\parf},
\]
which induce homomorphisms
\[
 K^{\bullet}(f_1)\otimes K_{\bullet}(f_2)\longrightarrow K_{\bullet}(f),
\]
\[
 K^{\bullet}(f_1)\otimes K^{\bullet}(f_2)\longrightarrow K^{\bullet}(f).
\]

%%%%%%%%%%  scan idx 309  |  folio 296  %%%%%%%%%%
\SGAPrintPage{296}%
(Of course, for \(X_i=Y_i=S\), we recover the \(K^{\bullet}(S)\)-module structure of \(K_{\bullet}(S)\).)

If \(X\) is flat over \(S\), we deduce, for \(n\geq1\), homomorphisms
\[
 K^{\bullet}(X/S)^{\otimes n}\longrightarrow K^{\bullet}\bigl((X/S)^n\bigr),
\]
which should in a sense replace the absent ring structure on \(K^{\bullet}(X/S)\).

\sgaStableTarget{sga6.e04.bibliography}\subsection*{BIBLIOGRAPHY}
\begin{enumerate}[label={[\arabic*]}]
\item Deligne, P., papiers secrets.
\item Grauert, H., Ein Theorem der Analytischen Garbentheorie \dots, Pub. Math. IHES, n\textsuperscript{o}5 (1960).
\item Grothendieck, A., Théorèmes de finitude pour la cohomologie des faisceaux, Bull. Soc. Math. France, 84, 1956, p.1-7.
\item[{[V]}] Verdier, J.-L., Catégories dérivées, IHES (1964), miméographié.
\end{enumerate}


% SGA_NATIVE_CONTINUATION_296
\providecommand{\Chat}{\widehat C}
\providecommand{\Ghat}{\widehat G}
\providecommand{\Gg}{\widehat G_g}
\providecommand{\Ga}{\widehat G_a}
\providecommand{\Grhat}{\widehat{\operatorname{Gr}}}
\providecommand{\Abfil}{\mathrm{Ab\,fil}}
\providecommand{\Kalgfil}{K\text{-}\mathrm{Alg.fil.}}
\providecommand{\Hom}{\operatorname{Hom}}
\providecommand{\Sym}{\mathfrak S}
\providecommand{\id}{\operatorname{Id}}
\providecommand{\Log}{\operatorname{Log}}
\providecommand{\Toddc}{\operatorname{Todd}}
\providecommand{\Ob}{\operatorname{Ob}}
\providecommand{\cG}{\mathcal G}
\providecommand{\cCzero}{\mathcal C_0}
\providecommand{\Fil}{\operatorname{Fil}}



% SGA_REVIEWED_FRENCH_CONTINUATION_296
% Expose V macros
\providecommand{\Chat}{\widehat C}
\providecommand{\Ghat}{\widehat G}
\providecommand{\Gg}{\widehat G_g}
\providecommand{\Ga}{\widehat G_a}
\providecommand{\Grhat}{\widehat{\operatorname{Gr}}}
\providecommand{\Abfil}{\underline{\mathrm{Ab}}\,\underline{\mathrm{fil}}}
\providecommand{\Kalgfil}{K\text{-}\mathrm{Alg.fil.}}
\providecommand{\Hom}{\operatorname{Hom}}
\providecommand{\Sym}{\mathfrak S}
\renewcommand{\id}{\operatorname{Id}}
\providecommand{\Log}{\operatorname{Log}}
\providecommand{\Toddc}{\operatorname{Todd}}
\renewcommand{\Ob}{\operatorname{Ob}}
\providecommand{\cG}{\mathcal G}
\providecommand{\cCzero}{\mathcal C_0}
\providecommand{\Fil}{\operatorname{Fil}}



% SGA_SOURCE_MACRO_RESTORATION_296: native provide must not shadow the two reviewed word rules.
\renewcommand{\Abfil}{\underline{\mathrm{Ab}}\,\underline{\mathrm{fil}}}

% SGA_END_REVIEWED_EN_BODY
\sgaStableTarget{sga6.e05.expose}\section*{Exposé V. Generalities on \(\lambda\)-rings}
\begin{center}
by P. Berthelot
\end{center}

The theory of \(\lambda\)-rings is useful whenever one uses rings such as the ring of classes of vector bundles on a ringed site, the representation ring of a group, and so on, for which the exterior-power operations define maps \(\lambda^i\). We shall use it abundantly later in the Seminar. It is therefore necessary first to develop the general formalism. This exposé is accordingly entirely independent of the rest of the Seminar, and may be useful to users of \(\lambda\)-rings outside algebraic geometry.

The first paragraph essentially gives the theory of Hirzebruch's ``multiplicative sequences of polynomials'' (cf. [2]). Besides its generality, the functorial viewpoint adopted here has the advantage of considerably reducing the calculations. The following paragraphs will show how, to every augmented \(\lambda\)-ring, one can associate a ``Chern ring'' which permits the development of a theory of Chern classes. This is the theory which will allow us to state the Riemann--Roch theorem in the general form of \hyperlink{sga6.e08.expose}{Exposé VIII}.

\sgaStableTarget{sga6.e05.sec.1}\section*{1. Universal polynomials}

\sgaStableTarget{sga6.e05.para.1-1}\textbf{1.1.} Let \(K\) be a commutative ring with unit, fixed throughout this paragraph. Let \(\mathcal C\) denote the category of filtered commutative \(K\)-algebras with unit, whose filtration is decreasing and such that the group of elements of filtration zero is the whole algebra. We define on \(\mathcal C\) a functor with values in the category of abelian groups as follows. If \(A\in\operatorname{Ob}(\mathcal C)\), associate to \(A\) the set of formal series \(\psi\in A[[t]]\) satisfying the following conditions: if
\[
  \psi=a_0+a_1t+\cdots+a_nt^n+\cdots ,
\]
then
\sgaStableTarget{sga6.e05.eq.1-1-1}\begin{equation}\tag{1.1.1}
  a_0=1,\qquad \forall i\geq 1,\quad a_i\in \Fil^i A .
\end{equation}
This set will be denoted \(\Ghat_K(A)\), or \(\Ghat(A)\) when no confusion is possible. It is clearly an abelian group for multiplication of formal series, and this defines a functor
\[
  \Ghat_K:K\text{-Alg.fil.}\longrightarrow \mathrm{Ab}.
\]
If \(A\) has the gross filtration, then \(\Ghat(A)\) is the group of formal series with augmentation \(1\), which will also be denoted
\[
  1+A[[t]]^+.
\]

\sgaStableTarget{sga6.e05.lem.1-2}\textbf{Lemma 1.2.} The functor \(\Ghat_K\) is representable.

Let \((X_i)_{i\in\mathbb N,\ i\geq 1}\) be a family of indeterminates. Assign to \(X_i\) the weight \(i\), and filter the ring \(K[X_i]_{i\geq 1}\) by weights with respect to the \(X_i\). One then obtains a bijection
\[
  \Ghat_K(A)\xrightarrow{\sim}
  \Hom_{K\text{-Alg.fil.}}\bigl(K[X_i]_{i\geq 1},A\bigr)
\]
by associating to every series
\[
  \psi=1+a_1t+\cdots+a_nt^n+\cdots\in\Ghat_K(A)
\]
the morphism \(f_\psi\) determined by \(f_\psi(X_i)=a_i\). This homomorphism is compatible with the filtrations, and one thus obtains an isomorphism of functors.

Put
\[
  \varphi=1+X_1t+\cdots+X_nt^n+\cdots .
\]
For every filtered \(K\)-algebra \(A\) and every series \(\psi\in\Ghat(A)\), there exists a unique homomorphism
\[
  f_\psi\in K[X_i]_{i\geq 1}\longrightarrow A
\]
such that
\[
  \Ghat(f_\psi)(\varphi)=\psi .
\]

\sgaStableTarget{sga6.e05.para.1-3}\textbf{1.3.} For every \(A\in\operatorname{Ob}(\mathcal C)\) we define a filtration on \(\Ghat(A)\) as follows: a series
\[
  \psi=1+a_1t+\cdots
\]
is said to be of filtration \(\geq k\) if \(a_i=0\) for \(1\leq i<k\). Let \(F^k\Ghat(A)\) be the subgroup so obtained. This is a functor on \(\mathcal C\), and from now on we shall regard \(\Ghat\) as a functor with values in the category \(\Abfil\) of filtered abelian groups by means of this filtration, with \(F^0\Ghat=\Ghat\).

Let \(H\) be a functor on \(\mathcal C\) with values in \(\Abfil\). We shall suppose that \(H\) satisfies the following two conditions:

\[
\begin{array}{ll}
\text{i)} & \forall A\in\operatorname{Ob}(\mathcal C),\ H(A)\text{ is separated and complete},\\[0.3em]
\text{ii)} & \text{if }A\in\operatorname{Ob}(\mathcal C)\text{ is a graded }K\text{-algebra, with the filtration associated}\\
&\text{to its grading, and if a finite group }G\text{ acts on }A\text{ by graded morphisms, then}\\
& H(A^G)\simeq H(A)^G .
\end{array}
\]
Here \(H(A)^G\) denotes the invariants under \(G\).

In fact we shall need these conditions only when \(A\) is a polynomial ring and \(G\) is a finite product of symmetric groups acting by permutation of the indeterminates.

\sgaStableTarget{sga6.e05.ex.unnamed}\textbf{Examples.}
\[
\begin{array}{ll}
1)& H(A)=\Ghat(A),\\
\sgaStableTarget{sga6.e05.ex.1-3-2}2)& H(A)=1+A[[t]]^+.
\end{array}
\]
The latter functor will be denoted \(\Gg\), with \(g\) standing for ``gross''. Finally,
\[
3)\quad H(A)=\Ga(A),
\]
where \(\Ga(A)\) denotes the group, for addition, of formal series
\[
  a=a_0+a_1t+\cdots
\]
such that \(a_0=0\) and \(a_i\in\Fil^i A\) for \(i\geq 1\), endowed with the filtration analogous to that on \(\Ghat\).

\sgaStableTarget{sga6.e05.para.1-4}\textbf{1.4.} For every \(n\), denote by \(f_n\) the endomorphism of \(K[X_1,\ldots,X_n]\) defined by
\[
  f_n(X_i)=X_i\quad (i<n),\qquad f_n(X_n)=0.
\]
It factors through the inclusion of \(K[X_1,\ldots,X_{n-1}]\) in \(K[X_1,\ldots,X_n]\). Hence the homomorphism
\[
  H(K[X_1,\ldots,X_{n-1}])\longrightarrow H(K[X_1,\ldots,X_n])
\]
is injective. This will allow us to identify the first group with a subgroup of the second; similarly, we shall regard \(H(K[X_1,\ldots,X_n])\) as contained in \(H(K[X_i]_{i\geq 1})\).

One should note that \(f_1\) is none other than the augmentation of \(K[X]\).

\sgaStableTarget{sga6.e05.para.1-5}\textbf{1.5.} Embed \(K[X_1,\ldots,X_n]\) in \(K[T_1,\ldots,T_n]\) by sending \(X_i\) to the \(i\)-th elementary symmetric function in the \(T_j\). Let the symmetric group \(\Sym_n\) act by permutation of the \(T_j\). Using the theorem on symmetric polynomials and (1.3.ii), one sees that the induced homomorphism
\[
  H(K[X_1,\ldots,X_n])\longrightarrow H(K[T_1,\ldots,T_n])
\]
is injective, which again permits us to identify the former with a subgroup of the latter:
\[
  H(K[X_1,\ldots,X_n])\simeq H(K[T_1,\ldots,T_n])^{\Sym_n}
  \subset H(K[T_1,\ldots,T_n]).
\]
Let
\[
  u_i^n:K[X]\longrightarrow K[T_1,\ldots,T_n]
\]
be the homomorphism defined by \(u_i^n(X)=T_i\), and put
\[
  u_H^n=\sum_{i=1}^n H(u_i^n):H(K[X])\longrightarrow H(K[T_1,\ldots,T_n]).
\]
The symmetric group permutes the morphisms \(H(u_i^n)\), and so leaves \(u_H^n\) invariant; thus \(u_H^n\) takes its values in
\[
  H(K[T_1,\ldots,T_n])^{\Sym_n}
   \simeq H(K[X_1,\ldots,X_n])
\]
by (1.3.ii). Notice, finally, that if \(H=\Ghat\), then
\[
  u_H^n(1+Xt)=\prod_{i=1}^n(1+T_it)=1+X_1t+\cdots+X_nt^n.
\]

\sgaStableTarget{sga6.e05.thm.1-6}\textbf{Theorem 1.6.} The map from \(\Hom_{\Abfil}(\Ghat,H)\) to \(H(K[X])\) which, to a morphism of functors, associates its value on the element \(1+Xt\) of \(\Ghat(K[X])\), is injective. Its image is the subgroup of \(H(K[X])\) formed by the elements \(\theta\) such that, with the notation of \hyperlink{sga6.e05.para.1-4}{1.4} for the \(f_n\),
\[
\begin{array}{ll}
\text{i)}& H(f_1)(\theta)=0,\\[0.3em]
\text{ii)}& \forall n>1,\quad
 H(f_n)\bigl(u_H^n(\theta)\bigr)\equiv u_H^n(\theta)
 \pmod{F^nH(K[X_1,\ldots,X_n])}.
\end{array}
\]

First prove injectivity. By \hyperlink{sga6.e05.lem.1-2}{Lemma 1.2} it is enough to see that a morphism of functors \(\xi\) which is zero on \(1+Xt\) is zero on
\[
  \varphi=1+X_1t+\cdots+X_nt^n+\cdots .
\]
The commutative diagram
\sgaStableTarget{sga6.e05.eq.1-6-1}\[
\begin{tikzcd}
\Ghat(K[X]) \arrow[r,"u^n_{\Ghat}"] \arrow[d,"\xi"'] &
\Ghat(K[X_1,\ldots,X_n]) \arrow[d,"\xi"]\\
H(K[X]) \arrow[r,"u_H^n"'] &
H(K[X_1,\ldots,X_n])
\end{tikzcd}
\tag{1.6.1}
\]
shows that, for every \(n\),
\[
  \xi(1+X_1t+\cdots+X_nt^n)=0.
\]
But
\[
  \varphi\equiv 1+X_1t+\cdots+X_nt^n
  \pmod{F^{n+1}\Ghat},
\]
hence
\[
  \xi(\varphi)\equiv \xi(1+X_1t+\cdots+X_nt^n)
  \pmod{F^{n+1}H};
\]
therefore \(\xi(\varphi)\equiv 0\pmod{F^{n+1}H}\) for every \(n\). Since \(H(K[X_i]_{i\geq 1})\) is separated and complete, \(\xi(\varphi)=0\).

Now let \(\xi\) be any element of \(\Hom(\Ghat,H)\), and set
\[
  \theta=\xi(1+Xt).
\]
We show that \(\theta\) satisfies conditions i) and ii). First,
\[
  H(f_1)(\theta)=H(f_1)\xi(1+Xt)
  =\xi\,\Ghat(f_1)(1+Xt)=\xi(1)=0.
\]
Second,
\[
  u_H^n(\theta)=\xi(1+\cdots+X_nt^n)
\]
by \hyperlink{sga6.e05.eq.1-6-1}{(1.6.1)}, while
\[
  1+\cdots+X_nt^n\equiv 1+\cdots+X_{n-1}t^{\,n-1}
  \pmod{F^n\Ghat},
\]
and
\[
  1+\cdots+X_{n-1}t^{\,n-1}=\Ghat(f_n)(1+\cdots+X_nt^n).
\]
Since \(H\circ \Ghat(f_n)=H(f_n)\circ \xi\), and since \(\xi\) is a morphism of filtered functors, condition ii) follows.

It remains to show that if \(\theta\in H(K[X])\) satisfies i) and ii), then there exists \(\xi\in\Hom(\Ghat,H)\) with
\[
  \theta=\xi(1+Xt).
\]
First determine the value of \(\xi\) on the element \(1+\cdots+X_nt^n\) of \(\Ghat(K[X_1,\ldots,X_n])\). In view of \hyperlink{sga6.e05.eq.1-6-1}{(1.6.1)}, it must be, by definition,
\[
  \xi(1+\cdots+X_nt^n)
  =
  \sum_{i=1}^n \xi(1+T_it),
\]
where the \(\xi(1+T_it)\) are known by functoriality from the morphisms \(u_i^n\).

One then glues the \(\xi(1+\cdots+X_nt^n)\), for variable \(n\), by means of the commutative diagram
\sgaStableTarget{sga6.e05.eq.1-6-2}\[
\begin{tikzcd}
K[X_1,\ldots,X_n] \arrow[r,"f_n"] \arrow[d,"\sigma_n"'] &
K[X_1,\ldots,X_n] \arrow[d,"\sigma'_n"]\\
K[T_1,\ldots,T_n] \arrow[r,"f'_n"'] &
K[T_1,\ldots,T_n],
\end{tikzcd}
\tag{1.6.2}
\]
where \(f'_n\) is determined by \(f'_n(T_i)=T_i\) if \(i<n\), and \(f'_n(T_n)=0\); where \(\sigma_n\) sends \(X_i\) to the \(i\)-th elementary symmetric function in \(T_1,\ldots,T_n\); and where \(\sigma'_n\) sends \(X_i\) to the \(i\)-th elementary symmetric function in \(T_1,\ldots,T_{n-1}\) for \(i<n\), and sends \(X_n\) to \(0\).

Indeed one obtains from this
\[
\begin{aligned}
H(f_n)\bigl(\xi(1+\cdots+X_nt^n)\bigr)
&=H(f'_n)\left(\sum_{i=1}^n \xi(1+T_it)\right)\\
&=\sum_{i=1}^{n-1}\xi(1+T_it)
   =\xi(1+\cdots+X_{n-1}t^{\,n-1}),
\end{aligned}
\]
inside \(H(K[T_1,\ldots,T_{n-1}])\) by the identifications of 1.4. By 1.6.ii the elements \(\xi(1+\cdots+X_nt^n)\) define an element \(\psi\) of \(H(K[X_i]_{i\geq 1})\), since this is a separated complete group. Put
\[
  \xi(\varphi)=\psi .
\]

Finally, if \(A\in\operatorname{Ob}(\mathcal C)\) and \(x\in\Ghat(A)\), define
\[
  \xi_A(x)=H(f_x)(\psi),
\]
where \(f_x:K[X_i]_{i\geq 1}\to A\) is the homomorphism defined in 1.2. Since \(\xi\) has already been defined on the elements \(1+\cdots+X_nt^n=\varphi_n\), one must verify that the two definitions give the same result, i.e.
\[
  H(f_{\varphi_n})(\psi)=\xi(1+\cdots+X_nt^n).
\]
But for every \(p\) one has
\[
  \psi\equiv \xi(1+\cdots+X_pt^p)
  \pmod{F^{p+1}H(K[X_i]_{i\geq 1})}.
\]
Thus, for \(p\geq n\),
\[
  H(f_{\varphi_n})(\psi)
  \equiv
  H(f_{\varphi_n})\bigl(\xi(1+\cdots+X_pt^p)\bigr)
  \pmod{F^{p+1}H(K[X_i]_{i\geq 1})}.
\]
Using the identifications of \hyperlink{sga6.e05.para.1-4}{1.4} and the gluing of the \(\xi(\varphi_m)\) for variable \(m\), one obtains
\[
  H(f_{\varphi_n})(\xi(\varphi_p))=\xi(\varphi_n).
\]
Hence finally
\[
  \forall p\geq n,\quad
  H(f_{\varphi_n})(\psi)\equiv \xi(1+\cdots+X_nt^n)
  \pmod{F^pH(K[X_1,\ldots,X_n])}.
\]
Since \(H\) is separated, the required equality follows.

We have thus defined a system of maps
\[
  \xi_A:\Ghat(A)\longrightarrow H(A).
\]
To verify that they constitute a morphism of functors with values in groups, one must prove first that, for every morphism \(h:A\to A'\), the square
\[
\begin{tikzcd}
\Ghat(A) \arrow[r,"\Ghat(h)"] \arrow[d,"\xi_A"'] &
\Ghat(A') \arrow[d,"\xi_{A'}"]\\
H(A) \arrow[r,"H(h)"'] & H(A')
\end{tikzcd}
\]
is commutative. If \(x\in\Ghat(A)\), then \(h\circ f_x=f_{\Ghat(h)(x)}\), whence
\[
  H(h)(\xi_A(x))=H(h)H(f_x)(\psi)
  =H(f_{\Ghat(h)(x)})(\psi)=\xi_{A'}(\Ghat(h)(x)).
\]
It remains only to show that the morphism of functors so constructed is compatible with the structures of filtered groups.

Suppose \(x\in F^n\Ghat(A)\). To say that \(x\) is of filtration \(>n\) is to say that
\(f_x:K[X_i]_{i\geq 1}\to A\) factors through the endomorphism \(g_n\) of \(K[X_i]_{i\geq 1}\) defined by
\[
  g_n(X_i)=0\quad (1\leq i\leq n-1),\qquad
  g_n(X_i)=X_i\quad (i\geq n).
\]
To show that \(\xi\) is compatible with filtrations, it is enough to show that
\[
  H(g_n)(\xi(\varphi))\in F^nH.
\]
But by construction
\[
  \xi(\varphi)\equiv \xi(1+\cdots+X_{n-1}t^{\,n-1})
  \pmod{F^nH}.
\]
Therefore
\[
  H(g_n)(\xi(\varphi))
  =H(g_n)\bigl(\xi(1+\cdots+X_{n-1}t^{\,n-1})\bigr)
  \equiv 0\pmod{F^nH}.
\]

To verify the additivity of \(\xi\), it suffices to do this for the elements of finite degree of \(\Ghat(A)\), and then pass to the limit. Let \(X_1,X_2\) be two elements of degree \(n\) of \(\Ghat(A)\), and let \(X\) be their product. Consider the commutative diagram
\[
\begin{tikzcd}[column sep=tiny]
K[X_1,\ldots,X_{2n}] \arrow[rr,"f_X"] \arrow[d] && A\\
\bigl(K[T_1,\ldots,T_{2n}]\bigr)^{\Sym_n\times\Sym_n}
\arrow[urr,"f'_X"'] &&
\end{tikzcd}
\]
where the first factor of \(\Sym_n\times\Sym_n\) permutes \(T_1,\ldots,T_n\) and the second permutes \(T_{n+1},\ldots,T_{2n}\). The ring \(K[X_1,\ldots,X_{2n}]\) is embedded into the ring of invariants by sending \(X_j\) to the \(j\)-th elementary symmetric function of the \(T_i\); \(f'\) sends the \(i\)-th elementary symmetric function of \(T_1,\ldots,T_n\), respectively of \(T_{n+1},\ldots,T_{2n}\), to the \(i\)-th coefficient of \(X_1\), respectively \(X_2\). Then
\[
\begin{aligned}
\xi(X)
&=H(f'_X)\left(\xi\left(\prod_{i=1}^{2n}(1+T_it)\right)\right),\\
\xi(X_1)
&=H(f'_X)\left(\xi\left(\prod_{i=1}^{n}(1+T_it)\right)\right),
\end{aligned}
\]
and similarly for \(X_2\), using \(i=n+1,\ldots,2n\). Since \(H(f'_X)\) is a homomorphism of groups, we are reduced to the case
\[
  X=\prod_{i=1}^{2n}(1+T_it),
\]
and in this case the additivity of \(\xi\) follows from the definition of the \(\xi(1+X_1t+\cdots+X_k t^k)\).

\sgaStableTarget{sga6.e05.para.1-7}\textbf{1.7.} One can generalize the preceding considerations to the case of a product of functors \(\Ghat\). The functor \((\Ghat)^k\) is representable by the ring
\[
  K[X_i^{(j)}]_{j=1,\ldots,k;\ i\geq 1},
\]
where \(X_i^{(j)}\) has weight \(i\). We denote by \(\varphi^{(j)}\) the series
\[
  1+X_1^{(j)}t+\cdots+X_n^{(j)}t^n+\cdots.
\]
The element \((\varphi^{(j)})_j\) of
\[
  \bigl(\Ghat(K[X_i^{(j)}])\bigr)^k
\]
has the universal property analogous to that of \(\varphi\) (cf. \hyperlink{sga6.e05.lem.1-2}{1.2}).

One can introduce the morphisms \(f_n^{(j)}\) and \(u_H^{n,(j)}\) as above, by repeating the construction for each family \((X_i^{(j)})_{i\geq 1}\) separately. One obtains a result analogous to \hyperlink{sga6.e05.thm.1-6}{Theorem 1.6} by considering morphisms of functors from \((\Ghat)^k\) to \(H\) which are \(\mathbb Z\)-multilinear and compatible with the filtrations in each of the \(k\) arguments separately. They are determined by their value on
\[
  (1+X^{(j)}t)_{j=1,\ldots,k}
  \in \bigl(\Ghat(K[X^{(j)}]_j)\bigr)^k.
\]
An element
\[
  \theta\in H(K[X^{(j)}]_j)
\]
determines such a morphism of functors if and only if it satisfies
\[
\begin{array}{ll}
\text{i)} & \forall j,\quad H(f_1^{(j)})(\theta)=0,\\[0.3em]
\text{ii)}& \forall n,j,\quad
 H(f_n^{(j)})\bigl(u_H^{n,(j)}(\theta)\bigr)
 \equiv u_H^{n,(j)}(\theta)\pmod{F^nH}.
\end{array}
\]
The reader will easily convince himself that, apart from notational complications, it is enough to copy the proof of 1.6.

\sgaStableTarget{sga6.e05.para.1-8}\textbf{1.8.} Let \(\mathcal C_0\) be the category of \(K\)-algebras, and let
\[
  i:\mathcal C_0\longrightarrow\mathcal C
\]
be the functor which assigns to a \(K\)-algebra \(A\) the \(K\)-algebra \(A\) endowed with the gross filtration. Let
\[
  g:\mathcal C\longrightarrow\mathcal C_0
\]
be the forgetful functor which forgets the filtration. These functors satisfy the following relations:
\[
  g\circ i=\id_{\mathcal C_0},\qquad
  \Ghat\circ i\circ g=\Gg,\qquad
  \Gg\circ i\circ g=\Gg
\]
(cf. \hyperlink{sga6.e05.para.1-3}{1.3}, \hyperlink{sga6.e05.ex.1-3-2}{example 2)}.

\sgaStableTarget{sga6.e05.lem.1-9}\textbf{Lemma 1.9.} Let \(H:\mathcal C_0\to\Abfil\) be a functor on \(\mathcal C_0\). Then the canonical homomorphism
\[
  \Hom_{\Abfil}(\Ghat,H\circ g)
  \longrightarrow
  \Hom_{\Abfil}(\Ghat\circ i,H)
\]
is an isomorphism.

To see that it is injective, one uses the sequence of canonical homomorphisms
\[
 \Hom(\Ghat,H\circ g)\to
 \Hom(\Ghat\circ i,H)\to
 \Hom(\Gg,H\circ g)\to
 \Hom(\Ghat,H\circ g),
\]
the last arrow coming from the fact that \(\Ghat\) is a subfunctor of \(\Gg\). To verify that the composite is the identity, observe that a morphism of functors
\[
  \xi:\Ghat\longrightarrow H\circ g
\]
is determined by its values on objects of \(\mathcal C\) whose filtration is gross. Indeed, if \(A\in\operatorname{Ob}(\mathcal C)\), and if \(A_g\) is \(A\) endowed with the gross filtration, the canonical morphism \(A\to A_g\) defines a commutative diagram
\[
\begin{tikzcd}
\Ghat(A) \arrow[r,"\xi_A"] \arrow[d] &
(H\circ g)(A) \arrow[d]\\
\Ghat(A_g) \arrow[r,"\xi_{A_g}"'] &
(H\circ g)(A_g),
\end{tikzcd}
\]
which shows that \(\xi_A\) is the restriction of \(\xi_{A_g}\) to \(\Ghat(A)\).

Next consider the composite homomorphism
\[
  \Hom(\Ghat\circ i,H)\longrightarrow
  \Hom(\Gg,H\circ g)\longrightarrow
  \Hom(\Ghat,H\circ g)\longrightarrow
  \Hom(\Ghat\circ i,H).
\]
It is clear that this is the identity, whence the asserted result.

\sgaStableTarget{sga6.e05.thm.1-10}\textbf{Theorem 1.10.} Let \(H\) be a functor on the category \(\mathcal C_0\) of \(K\)-algebras, with values in the category of separated complete filtered abelian groups, and commuting with invariants under a finite group. Then the map from
\[
  \Hom_{\Abfil}(\Ghat\circ i,H)
\]
to \(H(K[X])\) which, to a morphism of functors, associates its value on the element \(1+Xt\) of \(\Ghat\circ i(K[X])\), is injective. Its image is the subgroup of \(H(K[X])\) formed by the elements \(\theta\) such that, with notation analogous to that of \hyperlink{sga6.e05.para.1-4}{1.4} and \hyperlink{sga6.e05.para.1-5}{1.5},
\[
\begin{array}{ll}
\text{i)}& H(f_1)(\theta)=0,\\[0.3em]
\text{ii)}& \forall n>1,\quad
 H(f_n)\bigl(u_H^n(\theta)\bigr)\equiv u_H^n(\theta)
 \pmod{F^nH(K[X_1,\ldots,X_n])}.
\end{array}
\]
The theorem follows immediately from \hyperlink{sga6.e05.thm.1-6}{Theorem 1.6} and \hyperlink{sga6.e05.lem.1-9}{Lemma 1.9}, since \(H\circ g(K[X])\) does not depend on the chosen filtration on \(K[X]\).

The analogous theorem obtained by replacing \(\Ghat\) by \((\Ghat)^k\) is also true (cf. \hyperlink{sga6.e05.para.1-7}{1.7}).

When there is no risk of confusion, we write the functor \(\Ghat\circ i\) simply as \(\Ghat\). This is legitimate because \(i\) identifies \(\mathcal C_0\) with a subcategory of \(\mathcal C\).

\sgaStableTarget{sga6.e05.para.1-11}\textbf{1.11.} Let \(\mathcal C_g\) denote the category of graded \(K\)-algebras, and let
\[
  j:\mathcal C_g\longrightarrow\mathcal C
\]
be the functor which assigns to a graded \(K\)-algebra \(B\) the \(K\)-algebra \(B\) endowed with the filtration canonically defined by its grading. The associated graded functor
\[
  \operatorname{Gr}:\mathcal C\longrightarrow \mathcal C_g
\]
will be denoted \(\operatorname{Gr}\). There is an isomorphism of functors
\[
  \varepsilon:\id_{\mathcal C_g}\xrightarrow{\sim}\operatorname{Gr}\circ j.
\]

If \(B\in\operatorname{Ob}(\mathcal C_g)\), denote by
\[
  1+\widehat{B}^{+}
\]
the multiplicative group of elements of
\[
  \prod_{i=0}^{\infty}B^i
\]
having augmentation \(1\). With the filtration analogous to that on \(\Ghat\), one obtains a functor
\[
  \mathcal C_g\longrightarrow\Abfil.
\]
Furthermore, for every \(A\in\operatorname{Ob}(\mathcal C)\), one defines a homomorphism of filtered groups
\[
  \Ghat(A)\longrightarrow 1+\widehat{\operatorname{Gr}A}^{+}
\]
by sending the series \(1+\cdots+a_nt^n+\cdots\) to the series \(1+\cdots+\bar a_n+\cdots\), where \(\bar a_n\) is the class of \(a_n\) in \(\operatorname{Gr}^nA\). Its kernel is the group \(\Ghat'(A)\) of formal series such that
\[
  a_0=1,\qquad \forall i\geq 1,\quad a_i\in\Fil^{i+1}A.
\]
Thus one has a functorial isomorphism in \(A\):
\[
  \Ghat(A)/\Ghat'(A)\xrightarrow{\sim} 1+\widehat{\operatorname{Gr}A}^{+}.
\]

\sgaStableTarget{sga6.e05.lem.1-12}\textbf{Lemma 1.12.} Let \(H:\mathcal C_g\to\Abfil\) be a functor on \(\mathcal C_g\). Then the canonical homomorphism
\[
  \rho:\Hom_{\Abfil}(1+\widehat{\,\cdot\,}^{+},H)
  \longrightarrow
  \Hom_{\Abfil}(\Ghat,H\circ\operatorname{Gr})
\]
is an isomorphism.

The homomorphism \(\rho\) is defined as the composite
\[
 \Hom(1+\widehat{\,\cdot\,}^{+},H)
 \longrightarrow
 \Hom(1+\widehat{\operatorname{Gr}(\,\cdot\,)}^{+},H\circ\operatorname{Gr})
 \longrightarrow
 \Hom(\Ghat,H\circ\operatorname{Gr}),
\]
the last arrow being defined by the homomorphism
\[
  \Ghat\longrightarrow 1+\widehat{\operatorname{Gr}(\,\cdot\,)}^{+}.
\]
We construct an inverse morphism
\[
  \rho':\Hom(\Ghat,H\circ\operatorname{Gr})
  \longrightarrow
  \Hom(1+\widehat{\,\cdot\,}^{+},H).
\]
For this, we show that every morphism of functors from \(\Ghat\) to \(H\circ\operatorname{Gr}\) factors through the canonical morphism
\[
  \Ghat\longrightarrow 1+\widehat{\operatorname{Gr}(\,\cdot\,)}^{+}.
\]
Equivalently, it suffices to show that it vanishes on \(\Ghat'\), by 1.11. Let \(A\in\operatorname{Ob}(\mathcal C)\) and let \(x\in\Ghat'(A)\). Define a homomorphism of filtered \(K\)-algebras
\[
  K[X_i]_{i\geq 1}\longrightarrow A
\]
by sending \(X_1\) to \(0\) and \(X_i\) to the \((i-1)\)-st coefficient of \(x\) for \(i>1\). Using this homomorphism, it is enough to verify that the morphism of functors considered vanishes on the element
\[
  1+X_2t+\cdots+X_{n+1}t^n+\cdots .
\]
Use the endomorphism \(f\) of \(K[X_i]_{i\geq 1}\) defined by
\[
  f(X_i)=X_{i+1}\quad\text{for all }i,
\]
and the commutative diagram
\[
\begin{tikzcd}
\Ghat(K[X_i]) \arrow[r] \arrow[d,"\Ghat(f)"'] &
(H\circ\operatorname{Gr})(K[X_i]) \arrow[d,"H(\operatorname{Gr}(f))"]\\
\Ghat(K[X_i]) \arrow[r] &
(H\circ\operatorname{Gr})(K[X_i]).
\end{tikzcd}
\]
Now \(\operatorname{Gr}(f)=\operatorname{Gr}(f_1)\), where \(f_1\) is the augmentation homomorphism of \(K[X_i]_{i\geq 1}\). Since the latter annihilates \(\varphi\), the analogous diagram written for \(f_1\) shows that \(H(\operatorname{Gr}(f))\) annihilates the image of \(\varphi\) in \(H\circ\operatorname{Gr}(K[X_i]_{i\geq 1})\), and the result follows.

Passing to the quotient, one associates to every morphism of functors
\[
  \Ghat\longrightarrow H\circ\operatorname{Gr}
\]
a morphism of functors
\[
  1+\widehat{\operatorname{Gr}}^{+}\longrightarrow H\circ\operatorname{Gr}.
\]
By composing with \(j\) and using the isomorphism \(\varepsilon\), one obtains a morphism
\[
  1+\widehat{\,\cdot\,}^{+}\longrightarrow H,
\]
which defines the desired homomorphism \(\rho'\). It is immediate to verify that \(\rho\) and \(\rho'\) are inverses of one another.

\sgaStableTarget{sga6.e05.thm.1-13}\textbf{Theorem 1.13.} Let \(H\) be a functor on \(\mathcal C_g\), with values in the category of separated complete filtered abelian groups, and commuting with invariants under a finite group. Then the map from
\[
  \Hom_{\Abfil}(1+\widehat{\,\cdot\,}^{+},H)
\]
to \(H(K[X])\), which to a morphism of functors associates its value on the element \(1+X\) of \(1+\widehat{K[X]}^{+}\), is injective. Its image is the subgroup of \(H(K[X])\) formed by the elements \(\theta\) such that, with notation as in \hyperlink{sga6.e05.para.1-4}{1.4} and \hyperlink{sga6.e05.para.1-5}{1.5},
\[
\begin{array}{ll}
\text{i)}&H(f_1)(\theta)=0,\\[0.3em]
\text{ii)}&\forall n>1,\quad
H(f_n)\bigl(u_H^n(\theta)\bigr)\equiv u_H^n(\theta)
 \pmod{F^nH(K[X_1,\ldots,X_n])}.
\end{array}
\]
This is immediate from \hyperlink{sga6.e05.thm.1-6}{1.6} and 1.12. There is an analogous statement replacing \(\Ghat\) by \((\Ghat)^k\).

\sgaStableTarget{sga6.e05.rem.1-14}\textbf{Remark 1.14.} If \(H\) is a functor on \(\mathcal C_g\), with values in abelian groups, commuting with invariants under a finite group, and satisfying the condition
\[
  \forall B\in\operatorname{Ob}(\mathcal C_g),\qquad
  H(B)\simeq \varprojlim_n H(B/B^{\geq n}),
  \qquad B^{\geq n}=\prod_{i\geq n}B^i,
\]
then condition ii) of \hyperlink{sga6.e05.thm.1-13}{1.13} is always satisfied if one takes for the filtration on \(H\) the filtration by the kernels of
\[
  H(B)\longrightarrow H(B/B^{\geq n}),
\]
which is separated and complete. Indeed, if \(p_n\) is the canonical homomorphism
\[
  K[X_1,\ldots,X_n]\longrightarrow K[X_1,\ldots,X_n]/K[X_1,\ldots,X_n]^{\geq n},
\]
then
\[
  p_n\circ f_n=p_n,
\]
and hence
\[
  H(p_n)\bigl[H(f_n)\circ u_H^n(\theta)-u_H^n(\theta)\bigr]=0.
\]
This says precisely that
\[
  H(f_n)\circ u_H^n(\theta)-u_H^n(\theta)
  \in F^nH(K[X_1,\ldots,X_n]).
\]
Thus, in order to give a morphism of functors from \(1+\widehat{\,\cdot\,}^{+}\) to \(H\), it is enough to give an element of \(H(K[X])\) satisfying only condition i) of 1.13.

\sgaStableTarget{sga6.e05.para.1-15}\textbf{1.15.} Let \(H:\mathcal C\to\Abfil\) (respectively \(\mathcal C_0\to\Abfil\), \(\mathcal C_g\to\Abfil\)) be a functor satisfying the conditions of \hyperlink{sga6.e05.para.1-3}{1.3} (respectively \hyperlink{sga6.e05.thm.1-10}{1.10}, \hyperlink{sga6.e05.thm.1-13}{1.13}), and let
\[
  \xi:(\Ghat)^k\to H
\]
(respectively \((\Ghat\circ i)^k\to H\), \((1+\widehat{\,\cdot\,}^{+})^k\to H\)) be a morphism of functors. We have seen in \hyperlink{sga6.e05.para.1-7}{1.7} that \(\xi\) is determined by its value on the element
\[
  (\varphi^{(j)})_{j=1,\ldots,k}
\]
of
\[
  \bigl(\Ghat(K[X_i^{(j)}]_{j=1,\ldots,k;\ i\geq 1})\bigr)^k
\]
(respectively similarly). Suppose \(H\) is one of the functors
\[
  \Ghat,\quad \Gg,\quad \Ga,\quad 1+\widehat{\operatorname{Gr}}^{+},\quad\text{etc.}
\]
Then the element \(\xi((\varphi^{(j)}))\) is a formal series with coefficients in
\[
  K[X_i^{(j)}]_{i\geq 1,\ j=1,\ldots,k}.
\]
The \(n\)-th coefficient of this series is called the \(n\)-th universal polynomial of \(\xi\). The knowledge of these polynomials is equivalent to the knowledge of \(\xi\).

According to the proof of \hyperlink{sga6.e05.thm.1-6}{1.6}, these polynomials are themselves determined by the element
\[
  \xi(1+X^{(1)}t,\ldots,1+X^{(k)}t).
\]
The calculation is performed by the method of the proof: one computes the elements
\[
  \xi(1+X_1^{(1)}t+\cdots+X_n^{(1)}t^n,\ldots,
        1+X_1^{(k)}t+\cdots+X_n^{(k)}t^n),
\]
by decomposing
\[
  1+X_1^{(1)}t+\cdots+X_n^{(1)}t^n
  =\prod_{j=1}^n(1+T_j^{(1)}t)
\]
and by using the additivity of \(\xi\). The polynomials thus obtained are symmetric with respect to the \(T_j^{(i)}\), for fixed \(i\) and variable \(j\), and can be expressed as functions of the \(X_j^{(i)}\). Gluing by means of the filtration of \(H\), one thus determines the universal polynomials of \(\xi\). If
\[
  (\psi^{(1)},\ldots,\psi^{(k)})\in(\Ghat(A))^k,
\]
the calculation of \(\xi(\psi^{(1)},\ldots,\psi^{(k)})\) is made by substituting the coefficients of the \(\psi^{(i)}\) for the \(X_j^{(i)}\) in the universal polynomials.

Suppose that there exists an integer \(m\) such that the coefficient of \(t^n\) in
\[
  \theta=\xi(1+X^{(1)}t,\ldots,1+X^{(k)}t)
\]
is a polynomial whose isobaric components have weights between \(n\) and \(m\). Then the \(k\)-th universal polynomial of \(\xi\) has isobaric components whose weights range between \(k\) and \(km\). This follows immediately from the calculation just described and from the theorem on symmetric functions.

In particular, if \(H=\Ghat\), then the \(n\)-th universal polynomial of \(\xi\) has no component of weight \(<n\). Now consider the canonical morphism
\[
  H\longrightarrow 1+\widehat{\operatorname{Gr}}^{+}.
\]
Composing it with \(\xi\) gives a morphism
\[
  (\Ghat)^k\longrightarrow 1+\widehat{\operatorname{Gr}}^{+},
\]
which, by the proof of \hyperlink{sga6.e05.lem.1-12}{1.12}, factors through
\[
  (1+\widehat{\operatorname{Gr}}^{+})^k
\]
and defines a morphism
\[
  \bar\xi:(1+\widehat{\,\cdot\,}^{+})^k\longrightarrow 1+\widehat{\,\cdot\,}^{+}.
\]
One sees at once that the \(n\)-th universal polynomial of \(\bar\xi\) is the same as that of
\[
  (\Ghat)^k\longrightarrow 1+\widehat{\operatorname{Gr}}^{+},
\]
that is, the isobaric component of weight \(n\) of the \(n\)-th universal polynomial of \(\xi\).

The same remarks apply to the case \(H=\Ga\). Then one may define a morphism of functors
\[
  H\longrightarrow \widehat{\operatorname{Gr}}^{+}
\]
analogously to \hyperlink{sga6.e05.para.1-11}{1.11}, by putting
\[
  \widehat B^{+}=\prod_{i\geq 1}B^i
\]
for a graded algebra \(B\); one repeats the same reasoning and reaches the conclusions of the preceding paragraph.

\sgaStableTarget{sga6.e05.para.1-16}\textbf{1.16.} We finish with an example. Suppose \(K=\mathbb Q\). Define an endomorphism of the functor \(1+\widehat{\,\cdot\,}^{+}\), denoted \(\Toddc\) and called the Todd character, by putting
\[
  \Toddc(1+X)=\frac{X}{1-e^{-X}}.
\]
The conditions i) and ii) of \hyperlink{sga6.e05.thm.1-13}{Theorem 1.13} are immediately verified, which shows that an endomorphism of \(1+\widehat{\,\cdot\,}^{+}\) has indeed been defined.

\sgaStableTarget{sga6.e05.sec.2}\section*{2. Definition of \(\lambda\)-rings; examples}

In this paragraph the base ring of paragraph 1 will be \(\mathbb Z\), except in 2.9. We place ourselves in the category \(\mathcal C_0\) of \hyperlink{sga6.e05.para.1-8}{1.8}, and write \(\Ghat\) for the restriction of \(\Ghat\) to \(\mathcal C_0\), as in 1.10.

\sgaStableTarget{sga6.e05.def.2-1}\textbf{Definition 2.1.} A \emph{pre-\(\lambda\)-ring} is a commutative ring \(K\) with unit, endowed with a homomorphism of abelian groups
\[
  \lambda_t:K\longrightarrow 1+K[[t]]^+
\]
such that, for every \(x\in K\),
\[
  \lambda_t(x)=1+xt+\cdots.
\]
We shall write
\[
  \lambda_t(x)=\sum_{i\geq 0}\lambda^i(x)t^i.
\]
Thus giving a pre-\(\lambda\)-ring structure on a ring \(K\) is equivalent to giving a family of maps
\[
  \lambda^i:K\longrightarrow K\qquad (i\in\mathbb N)
\]
such that
\[
\begin{array}{ll}
\text{i)}&\forall x\in K,\quad \lambda^0(x)=1,\\[0.3em]
\text{ii)}&\forall x\in K,\quad \lambda^1(x)=x,\\[0.3em]
\text{iii)}&\forall x,y\in K,\ \forall n\in\mathbb N,\quad
\lambda^n(x+y)=\displaystyle\sum_{p+q=n}\lambda^p(x)\lambda^q(y).
\end{array}
\]
If \(K,K'\) are two pre-\(\lambda\)-rings and \(f:K\to K'\) is a ring homomorphism, we say that \(f\) is a homomorphism of pre-\(\lambda\)-rings, or a \(\lambda\)-homomorphism, if the diagram
\[
\begin{tikzcd}
K \arrow[r,"\lambda_t"] \arrow[d,"f"'] &
1+K[[t]]^+ \arrow[d,"\Ghat(f)"]\\
K' \arrow[r,"\lambda_t"'] &
1+K'[[t]]^+
\end{tikzcd}
\]
is commutative. This is equivalent to the relations
\[
  \forall x\in K,\ \forall i\in\mathbb N,\qquad
  \lambda^i(f(x))=f(\lambda^i(x)).
\]

\sgaStableTarget{sga6.e05.ex.2-2}\textbf{Examples 2.2.}
\begin{enumerate}[label=\alph*)]
\item One can give all pre-\(\lambda\)-ring structures on \(\mathbb Z\). It is enough to give \(\lambda_t(1)\), which is then a formal series with integral coefficients subject to the sole condition that it begin with \(1+t\).

\item Let \((T,A)\) be a commutatively ringed topos, and let \(K\) be the ``Grothendieck group'' of \(A\)-modules locally isomorphic to a direct factor of a free \(A\)-module of finite type. We know that \(K\) can be given a ring structure by the tensor product (\hyperlink{sga6.e04.expose}{Exposé IV}), and we now endow it with a pre-\(\lambda\)-ring structure as follows. For every \(A\)-module \(E\) locally a direct factor of a free module of finite type, let \(\lambda^i(E)\) denote the class in \(K\) of the \(i\)-th exterior power of \(E\), which is again locally a direct factor of a free module of finite type. This defines a function from the category of such \(A\)-modules to the abelian group \(1+K[[t]]^+\). It remains to show that this function is additive, which will imply that it factors through \(K\) and hence gives the desired structure.
\end{enumerate}

\sgaStableTarget{sga6.e05.lem.2-2-1}\textbf{Lemma 2.2.1.} Let
\[
  0\longrightarrow E' \xrightarrow{u} E \xrightarrow{v} E''\longrightarrow 0
\]
be an exact sequence of \(A\)-modules which is locally split, for example with \(E''\) locally a direct factor of a free module of finite type. There exists on \(\Lambda^n(E)\) a finite filtration and isomorphisms
\[
  \Lambda^{n-p}(E')\otimes_A\Lambda^p(E'')
  \xrightarrow{\sim}
  \operatorname{Gr}^p\Lambda^n(E).
\]
Take for \(F^p\Lambda^n(E)\) the image of the canonical morphism
\[
  \Lambda^{n-p}(E')\otimes_A\Lambda^p(E)\longrightarrow \Lambda^n(E)
  \quad (0\leq p\leq n).
\]
If \(p>n\), put \(F^p\Lambda^n(E)=\Lambda^n(E)\), and if \(p<0\), put \(F^p\Lambda^n(E)=0\). This does indeed define a finite filtration on \(\Lambda^n(E)\).

Define a homomorphism
\[
  f_p:\Lambda^{n-p}(E')\otimes_A\Lambda^p(E'')
  \longrightarrow
  F^p\Lambda^n(E)/F^{p-1}\Lambda^n(E)
\]
by
\[
  f_p(x'_1\wedge\cdots\wedge x'_{n-p}\otimes x''_1\wedge\cdots\wedge x''_p)
  =
  \operatorname{cl}\bigl(u(x'_1)\wedge\cdots\wedge u(x'_{n-p})\wedge
  x_1\wedge\cdots\wedge x_p\bigr),
\]
where the \(x_i\) are lifts in \(E\) of the \(x''_i\) in \(E''\). This visibly does not depend on the choice of the \(x_i\). If the exact sequence is locally split, then \(f_p\) is an isomorphism, since locally
\[
  \Lambda^n(E)\simeq\bigoplus_{p+q=n}
  \Lambda^p(E')\otimes_A\Lambda^q(E'').
\]

When \(E',E,E''\) are locally direct factors of free modules of finite type, the \(\operatorname{Gr}^p\Lambda^n(E)\) are locally direct factors of free modules of finite type, and hence so are the \(F^p\Lambda^n(E)\). Thus one may write in \(K\)
\[
  \lambda^n(E)=\sum_{p=0}^{n}\operatorname{cl}(\operatorname{Gr}^p\Lambda^n(E))
  =
  \sum_{p=0}^{n}\lambda^{n-p}(E')\lambda^p(E''),
\]
which gives the required additivity.

Let \((T,A)\) and \((T',A')\) be two commutatively ringed topoi, and let
\[
  f:(T,A)\longrightarrow (T',A')
\]
be a morphism of ringed topoi. We know (\hyperlink{sga6.e04.para.2-7}{Exposé IV, 2.7)} that \(f\) defines a morphism of rings on the corresponding \(K\)'s, namely
\[
  f^*:K^*_{\mathrm{naif}}(T',A')\longrightarrow K^*_{\mathrm{naif}}(T,A).
\]
Since inverse image commutes with exterior powers, the corresponding morphism \(f^*\) on \(K\) is a \(\lambda\)-homomorphism for the structure just defined.

We point out two special cases:
\begin{itemize}
\item If \(X\) is a scheme, then \(K^*_{\mathrm{naif}}(X)\) is thus endowed with a pre-\(\lambda\)-ring structure. This is the case which will be used later in the Seminar, and it may be regarded as the model for the general theory of \(\lambda\)-rings.
\item If \(A\) is a ring, \(X=\operatorname{Spec}A\), and \(G\) is a group acting trivially on \(A\), then, taking for \((T,A)\) the topos of sheaves of sets on \(X\) on which \(G\) acts, the ring \(K\) thus obtained is the representation ring of \(G\) in the projective \(A\)-modules of finite type.
\end{itemize}

\sgaStableTarget{sga6.e05.para.2-3}\textbf{2.3.} Let \(K\) be a commutative ring with unit. We propose to endow the abelian group
\[
  \Ghat(K)=1+K[[t]]^+
\]
with a pre-\(\lambda\)-ring structure, functorially in \(K\). To this end, we apply the results of \hyperlink{sga6.e05.thm.1-10}{1.10} to the functor \(\Ghat\).

First define multiplication by giving a \(\mathbb Z\)-bilinear morphism of functors
\[
  \Ghat\times\Ghat\longrightarrow\Ghat
\]
through the element \(1+XYt\) of \(\Ghat(\mathbb Z[X,Y])\). This element satisfies 1.7.i; to see that it satisfies 1.7.ii, one must show that in
\[
  \prod_{i=1}^{n}(1+T_iYt),
\]
when expressed by means of the elementary symmetric functions \(X_j\) in the \(T_i\), the variable \(X_n\) does not occur in the coefficients of \(t^k\) for \(k<n\). This is evident by degree considerations.

The law so defined will be denoted by the sign \(\circ\). Thus, by definition,
\sgaStableTarget{sga6.e05.eq.2-3-1}\begin{equation}\tag{2.3.1}
  (1+Xt)\circ(1+Yt)=1+XYt .
\end{equation}
It is distributive by construction, and it is associative and commutative by the diagrams
\[
\begin{tikzcd}[column sep=tiny]
\Ghat\times\Ghat\times\Ghat \arrow[r] \arrow[d] &
\Ghat\times\Ghat \arrow[d]\\
\Ghat\times\Ghat \arrow[r] & \Ghat
\end{tikzcd}
\qquad
\begin{tikzcd}[column sep=tiny]
\Ghat\times\Ghat \arrow[r] \arrow[d] &
\Ghat\times\Ghat \arrow[d]\\
\Ghat \arrow[r,"\id"'] & \Ghat,
\end{tikzcd}
\]
where the meaning of the arrows is evident, and where it is enough to verify commutativity on the elements \(1+Xt,1+Yt,1+Zt\), respectively \(1+Xt,1+Xt\), of \(\Ghat(\mathbb Z[X,Y,Z])\), respectively \(\Ghat(\mathbb Z[X,Y])\), by \hyperlink{sga6.e05.thm.1-10}{1.10}; this verification is immediate.

By \hyperlink{sga6.e05.para.1-8}{1.8}, the product under this law of two series is calculated by substituting their coefficients into the universal polynomials of the law. The \(n\)-th universal polynomial will be denoted \(P_n\). It is a polynomial in
\[
  \mathbb Z[X_i,Y_j]_{1\leq i,j\leq n}
\]
which is isobaric of weight \(n\) with respect to each of the two groups of variables (see \hyperlink{sga6.e05.para.1-15}{1.15}), and symmetric with respect to these two groups. For example,
\[
  P_2(X_1,X_2;Y_1,Y_2)=X_2Y_1^2+Y_2X_1^2-2X_2Y_2 .
\]
One also has the following relation: if \(K\in\operatorname{Ob}(\mathcal C_0)\), \(\psi\in\Ghat(K)\), and \(a\in K\), then
\sgaStableTarget{sga6.e05.eq.2-3-2}\begin{equation}\tag{2.3.2}
  \psi(t)\circ(1+at)=\psi(at).
\end{equation}
As usual, by the standard functorial argument this reduces to the case \(K=\mathbb Z[X,Y]\), \(\psi=1+Xt\), \(a=Y\), where it is nothing other than \hyperlink{sga6.e05.eq.2-3-1}{(2.3.1)}.

Applying \hyperlink{sga6.e05.eq.2-3-2}{(2.3.2)} to \(\Ghat(K)=1+K[[t]]^+\) with \(a=1\), one sees that \(1+t\) is the unit for this multiplication. To avoid confusion, we sometimes denote it by \(1\), or by \(1_{\circ}\).

The functor \(\Ghat\) can now be considered as a functor from \(\mathcal C_0\) to the category of commutative rings with unit, and hence can be composed with itself. To avoid confusion, choose different indeterminates for each of the two factors in the product of \(\Ghat\) with itself; we indicate the indeterminate by an exponent when necessary:
\[
  \Ghat^t,\quad \Ghat^u,\quad\text{etc.}
\]
For \(A\in\operatorname{Ob}(\mathcal C_0)\) one has
\[
  \Ghat^u\circ \Ghat^t(A)=1+\{1+A[[t]]^+\}[[u]]^+.
\]
To define the operations \(\lambda^i\) on \(1+K[[t]]^+\), we are therefore going to give a morphism of functors
\[
  \Ghat\longrightarrow \Ghat^u\circ\Ghat^t.
\]
For this, one must first define a filtration of the functor \(\Ghat^u\circ\Ghat^t\) which permits the gluing of 1.6.

First note that \(F^k\Ghat\) is an ideal for the law \(\circ\). For this follows, for the product
\[
  \left(1+\sum_{i\geq k}X_i t^i\right)\circ(1+Yt),
\]
by reducing by functoriality to the case \(A=\mathbb Z[X_i,Y]_{i\geq 1}\), and using \hyperlink{sga6.e05.eq.2-3-2}{(2.3.2)}. Therefore \(F^n(\Ghat^u\circ\Ghat^t)\) is defined to be the set of series of the form
\[
  1+\psi_1u+\cdots+\psi_iu^i+\cdots,\qquad
  \psi_i\in F^{[n/i]}\Ghat,
\]
where \([n/i]\) denotes the integral part of \(n/i\).

The element defining the morphism \(\lambda_u\) is
\[
  1+\{1+Xt\}u.
\]
The reader will verify as an exercise that the conditions of \hyperlink{sga6.e05.thm.1-6}{1.6} are indeed fulfilled; in fact the filtration was chosen ad hoc.

Again, the operations \(\lambda^i\) are calculated by substitution from
\[
  \lambda_u(1+\cdots+X_nt^n+\cdots),
\]
and so are determined by universal polynomials. Denote by
\[
  Q_{i,j}(X_1,\ldots,X_{ij})
\]
the coefficient of \(t^j\) in
\[
  \lambda^i(1+\cdots+X_nt^n+\cdots).
\]
It is an isobaric polynomial of weight \(ij\) with respect to the variables \(X_1,\ldots,X_{ij}\) (by a verification analogous to that of \hyperlink{sga6.e05.para.1-15}{1.15}). It is calculated by symmetric functions from the defining element
\sgaStableTarget{sga6.e05.eq.2-3-3}\begin{equation}\tag{2.3.3}
  \lambda_u(1+Xt)=1+\{1+Xt\}u .
\end{equation}
For example,
\[
  Q_{2,2}(X_1,X_2,X_3,X_4)=X_1X_3-X_4.
\]

Finally, we indicate the origin of formulas \hyperlink{sga6.e05.eq.2-3-1}{(2.3.1)} and \hyperlink{sga6.e05.eq.2-3-3}{(2.3.3)}. Consider the \hyperlink{sga6.e05.ex.2-2}{example 2.2}.b, where the operations \(\lambda^i\) are defined by exterior powers of \(A\)-modules. We want
\[
  \lambda_t:K\longrightarrow \Ghat(K)
\]
to be a \(\lambda\)-homomorphism. Thus on \(\Ghat\) one must define laws such that, for all \(A\)-modules \(E,F\),
\[
  \lambda_t(EF)=\lambda_t(E)\circ\lambda_t(F),\qquad
  \lambda_u(\lambda_t(E))=\Ghat^u_{\lambda_t}\bigl(\lambda_t(\lambda_u(E))\bigr).
\]
The relations \hyperlink{sga6.e05.eq.2-3-1}{(2.3.1)} and \hyperlink{sga6.e05.eq.2-3-3}{(2.3.3)} are precisely the preceding relations applied to the case where \(E\) and \(F\) are invertible modules.

We shall see in \hyperlink{sga6.e06.expose}{Exposé VI} to what extent the desired property is verified.

\sgaStableTarget{sga6.e05.def.2-4}\textbf{Definition 2.4.} A pre-\(\lambda\)-ring \(K\) will be called a \emph{\(\lambda\)-ring} if the homomorphism of abelian groups
\[
  \lambda_t:K\longrightarrow 1+K[[t]]^+
\]
is a \(\lambda\)-homomorphism for the pre-\(\lambda\)-ring structure on the second member defined in 2.3. Thus the following relations hold for all \(x,y\in K\):
\sgaStableTarget{sga6.e05.eq.2-4-1}\begin{align}
  \lambda_t(1)&=1+t,&
  \lambda_t(xy)&=\lambda_t(x)\circ\lambda_t(y),&
  \lambda_u(\lambda_t(x))&=\Ghat(\lambda_t)\bigl(\lambda_u(x)\bigr).
  \tag{2.4.1}
\end{align}
In terms of the maps \(\lambda^i\), the last two relations are equivalent, for \(x,y\in K\) and \(i,j\in\mathbb N\), to
\sgaStableTarget{sga6.e05.eq.2-4-2}\sgaStableTarget{sga6.e05.eq.2-4-3}\begin{align}
  \lambda^i(xy)
  &=P_i(\lambda^1(x),\ldots,\lambda^i(x);
       \lambda^1(y),\ldots,\lambda^i(y)), \tag{2.4.2}\\
  \lambda^j(\lambda^i(x))
  &=Q_{i,j}(\lambda^1(x),\ldots,\lambda^{ij}(x)). \tag{2.4.3}
\end{align}

We note that here we depart from the usual terminology (cf. [1]), where what we call a pre-\(\lambda\)-ring was called a \(\lambda\)-ring, and our \(\lambda\)-rings were called ``special'' \(\lambda\)-rings. The reason is that only ``special'' \(\lambda\)-rings are useful in practice.

\sgaStableTarget{sga6.e05.ex.2-5}\textbf{Examples 2.5.}
\begin{enumerate}[label=\alph*)]
\item We have determined all pre-\(\lambda\)-ring structures on \(\mathbb Z\). It follows from \hyperlink{sga6.e05.eq.2-4-1}{(2.4.1)} that only the structure defined by
\[
  \lambda_t(1)=1+t
\]
is capable of being a \(\lambda\)-ring structure. One then has
\[
  \lambda_t(n)=(1+t)^n,
\]
hence
\[
  \forall n,n'\in\mathbb Z,\quad
  \lambda_t(n)\circ\lambda_t(n')=(1+t)^n\circ(1+t)^{n'}
  =(1+t)^{nn'}=\lambda_t(nn'),
\]
and
\[
  \lambda_u((1+t)^n)
  =
  (\lambda_u(1+t))^n
  =
  (1+\{1+t\}u)^n
  =
  \Ghat(\lambda_t)((1+u)^n)
  =
  \Ghat(\lambda_t)(\lambda_u(n)).
\]
Thus this is indeed a \(\lambda\)-ring structure. If \(\binom{X}{k}\) denotes the polynomial with rational coefficients
\[
  \frac{1}{k!}X(X-1)\cdots(X-k+1),
\]
then
\sgaStableTarget{sga6.e05.eq.2-5-1}\begin{equation}\tag{2.5.1}
  \lambda^i(n)=\binom{n}{i}.
\end{equation}

\item For every ring \(K\), the pre-\(\lambda\)-ring \(1+K[[t]]^+\) is a \(\lambda\)-ring. First,
\[
  \lambda_u(1)=\lambda_u(1+t)=1+\{1+t\}u=1+1\cdot u.
\]
The other relations are verified by the commutativity of the diagrams
\[
\begin{tikzcd}[column sep=tiny]
\Ghat^t\times\Ghat^t \arrow[r,"\lambda_u\times\lambda_u"] \arrow[d,"\circ"'] &
(\Ghat^u\circ\Ghat^t)\times(\Ghat^u\circ\Ghat^t) \arrow[d,"\circ"]\\
\Ghat^t \arrow[r,"\lambda_u"'] &
\Ghat^u\circ\Ghat^t,
\end{tikzcd}
\qquad
\begin{tikzcd}[column sep=tiny]
\Ghat^t \arrow[r,"\lambda_v"] \arrow[d,"\lambda_u"'] &
\Ghat^v\circ\Ghat^t \arrow[d,"\Ghat^v(\lambda_u)"]\\
\Ghat^u\circ\Ghat^t \arrow[r,"\lambda_v"'] &
\Ghat^v\circ\Ghat^u\circ\Ghat^t,
\end{tikzcd}
\]
which are checked by \hyperlink{sga6.e05.thm.1-10}{1.10} on the elements \(1+Xt\) and \(1+Yt\). For the first,
\[
\begin{aligned}
\lambda_u(1+Xt)\circ\lambda_u(1+Yt)
&=(1+\{Xt\}u)\circ(1+\{1+Yt\}u)\\
&=1+\{(1+Xt)\circ(1+Yt)\}u\\
&=1+\{1+XYt\}u\\
&=\lambda_u((1+Xt)\circ(1+Yt)).
\end{aligned}
\]
For the second,
\[
  \lambda_v(\lambda_u(1+Xt))
  =
  1+\{1+\{1+Xt\}u\}v
  =
  \Ghat^v(\lambda_u)\lambda_v(1+Xt).
\]

\item The geometric examples of \hyperlink{sga6.e05.ex.2-2}{2.2} will be treated later.
\end{enumerate}

\sgaStableTarget{sga6.e05.para.2-6}\textbf{2.6.} One can generalize \hyperlink{sga6.e05.ex.2-5}{Example 2.5}.a as follows. Let \(K\) be a commutative ring with unit such that:
\[
\begin{array}{ll}
\text{i)}& K\text{ is torsion-free},\\
\text{ii)}& \forall x\in K,\ \forall n\in\mathbb N,\quad \binom{x}{n}\in K .
\end{array}
\]
Then, for all \(x\in K\) and all \(n\in\mathbb N\), put
\[
  \lambda^n(x)=\binom{x}{n},
\]
that is,
\sgaStableTarget{sga6.e05.eq.2-6-1}\begin{equation}\tag{2.6.1}
  \lambda_t(x)=\sum_{n=0}^{\infty}\binom{x}{n}t^n
  =(1+t)^x=\exp(x\Log(1+t)),
\end{equation}
where the last identity is considered in \(K\otimes_{\mathbb Z}\mathbb Q[[t]]\). By multiplicativity of the exponential, this defines a pre-\(\lambda\)-ring structure on \(K\).

To see that it is a \(\lambda\)-ring structure, one must prove the relations \hyperlink{sga6.e05.eq.2-4-2}{(2.4.2)}, and it is enough to prove them in \(K\otimes\mathbb Q\), by i). One is then reduced by functoriality to the case \(K=\mathbb Q[X,Y]\) and to the elements \(X,Y\). In view of the definition of the \(\lambda^i\), the two sides of the equalities \hyperlink{sga6.e05.eq.2-4-2}{(2.4.2)} are polynomials in \(X,Y\), with rational coefficients. It is therefore enough, in order to prove their equality, to show that they take the same value for all integral values of \(X,Y\). But the relations obtained by substituting integral values in \(X,Y\) are precisely those expressing that the pre-\(\lambda\)-ring \(\mathbb Z\) of 2.5.a is a \(\lambda\)-ring. They are therefore satisfied.

\sgaStableTarget{sga6.e05.def.2-7}\textbf{Definition 2.7.} A \emph{binomial ring} is a ring satisfying the conditions 2.6.i and ii), endowed with the \(\lambda\)-ring structure defined by \hyperlink{sga6.e05.eq.2-6-1}{(2.6.1)}.

\sgaStableTarget{sga6.e05.ex.2-8}\textbf{Example 2.8.} Let \(X\) be a site, \(K\) a binomial ring, and \(\mathcal K\) the sheaf associated to the constant presheaf with value \(K\) on \(X\). Then \(H^0(X,\mathcal K)\) is a binomial ring. Indeed, the conditions \hyperlink{sga6.e05.para.2-6}{2.6} are preserved under passage to an inductive limit. In particular, \(H^0(X,\mathbb Z)\) is a binomial ring.

\sgaStableTarget{sga6.e05.para.2-9}\textbf{2.9.} Let \(K\) be a binomial ring and \(A\) a \(K\)-algebra. For every series
\[
  \varphi\in\Ghat(A)
\]
and every \(a\in K\), one can define a series \(\varphi^a\) in the usual way. Writing \(\varphi=1+\psi\), put
\sgaStableTarget{sga6.e05.eq.2-9-1}\begin{equation}\tag{2.9.1}
  \varphi^a=(1+\psi)^a
  =
  \sum_{n=0}^{\infty}\binom{a}{n}\psi^n
  \left(=\exp(a\Log\varphi)\right),
\end{equation}
where the last equality can be written only modulo torsion. These formulas are obtained by substitution from the analogous formulas for \(\varphi=1+t\). Therefore the following relations, the classical properties of the exponential, remain valid for \(a,b\in K\) and \(\varphi,\psi\in\Ghat(A)\):
\sgaStableTarget{sga6.e05.eq.2-9-2}\begin{equation}\tag{2.9.2}
  \varphi^a\varphi^b=\varphi^{a+b},\qquad
  \varphi^{ab}=(\varphi^a)^b,\qquad
  \varphi^a\psi^a=(\varphi\psi)^a .
\end{equation}



\providecommand{\Ghat}{\widehat G}
\providecommand{\Gg}{\widehat G_g}
\providecommand{\Gcirc}{\widehat G_{\circ}}
\providecommand{\Gnabla}{\widehat G_{\nabla}}
\providecommand{\Grhat}{\widehat{\operatorname{Gr}}}
\providecommand{\Log}{\operatorname{Log}}
\providecommand{\Fil}{\operatorname{Fil}}
\providecommand{\eps}{\varepsilon}
\providecommand{\id}{\operatorname{id}}
\providecommand{\Sym}{\operatorname{Sym}}
\providecommand{\Mon}{\operatorname{Mon}}


\sgaStableTarget{sga6.e05.continuation-of-2-9}\subsection*{Continuation of 2.9}

We still suppose that \(K\) is binomial. Multiplication, for the law \(\circ\), by \(\lambda_t(a)\), where \(a\in K\), defines a morphism of functors
\[
  \Ghat_K\longrightarrow \Ghat_K
\]
which is characterized by its value on the element \(1+Xt\) of \(G(K[X])\). Thus
\[
  \lambda_t(a)\circ(1+Xt)=\lambda_{Xt}(a),
\]
and hence
\sgaStableTarget{sga6.e05.eq.2-9-3}\begin{equation}\tag{2.9.3}
  \lambda_t(a)\circ(1+Xt)
  =(1+Xt)^a
  =\exp(a\Log(1+Xt))
  =\sum_{n=0}^{\infty}\binom{a}{n}X^nt^n .
\end{equation}
It follows that
\sgaStableTarget{sga6.e05.eq.2-9-4}\begin{equation}\tag{2.9.4}
  \lambda_t(a)\circ\varphi=\varphi^a
  \quad(=\exp(a\Log\varphi)).
\end{equation}
Indeed, the two members are additive functors which coincide on \(1+Xt\). From \hyperlink{sga6.e05.eq.2-9-3}{(2.9.3)} and \hyperlink{sga6.e05.para.1-15}{1.15} one obtains that the universal polynomial of degree \(n\) defining multiplication by \(\lambda_t(a)\) is isobaric of weight \(n\). We denote it by \(R_{n,a}\).

Finally, if \(A\) is a \(K\)-\(\lambda\)-algebra and \(b\in A\), then
\sgaStableTarget{sga6.e05.eq.2-9-5}\begin{equation}\tag{2.9.5}
  \lambda_t(ab)=(\lambda_t(b))^a .
\end{equation}
This follows from \(\lambda_t(ab)=\lambda_t(a)\circ\lambda_t(b)\) and from \hyperlink{sga6.e05.eq.2-9-4}{(2.9.4)}.

\sgaStableTarget{sga6.e05.sec.3}\section*{3. The operations \(\gamma\)}

We still work over the category \(\mathscr C_0\).

\sgaStableTarget{sga6.e05.para.3-1}\textbf{3.1.} Let \(K\) be a \(\lambda\)-ring. We shall define from the operations \(\lambda^i\) new operations denoted \(\gamma^i\). They are the operations which are adapted to Chern classes. If \(E\) is a vector bundle of rank \(d\), the top Chern class of \(E\) is represented, in \(K\)-theoretic language, by the alternating sum of the exterior powers. To obtain the \(n\)-th class, one first replaces \(E\) by \(E\oplus\mathbf1^{\,n-d}\), so that the same top-class recipe applies. This motivates the formula below.

\sgaStableTarget{sga6.e05.prop.3-2}\textbf{Proposition 3.2.} Let \(K\) be a \(\lambda\)-ring.
\begin{enumerate}[label=\roman*)]
\item For every \(x\in K\) and every \(n\in\mathbb N\),
\[
  \lambda^n(x+n-1)=(-1)^n\sum_{i=0}^n(-1)^i\lambda^i(x+n).
\]
\item If
\[
  \gamma^n(x)=\lambda^n(x+n-1),\qquad
  \gamma_t(x)=\sum_{n\ge0}\gamma^n(x)t^n,
\]
then
\[
  \gamma_t(x)=\lambda_{t/(1-t)}(x),
  \qquad\text{hence}\qquad
  \lambda_t(x)=\gamma_{t/(1+t)}(x).
\]
\end{enumerate}
Thus knowledge of the operations \(\gamma\) is equivalent to knowledge of the operations \(\lambda\).

\emph{Proof.} For i),
\[
  \lambda^n(x+n-1)=\sum_{i=0}^{n}\lambda^i(x+n)\lambda^{n-i}(-1).
\]
Since \(\lambda_t(1)=1+t\), one has \(\lambda_t(-1)=1/(1+t)\), and hence \(\lambda^k(-1)=(-1)^k\), giving the desired relation. For ii),
\[
  \lambda_{t/(1-t)}(x)=\sum_{n\ge0}\lambda^n(x)t^n(1-t)^{-n},
\]
and
\[
  t^n(1-t)^{-n}=\sum_{p\ge n}\binom{p-1}{p-n}t^p.
\]
The coefficient of \(t^p\) is therefore
\[
  \sum_{n=0}^{p}\lambda^n(x)\binom{p-1}{p-n}
  =\sum_{n=0}^{p}\lambda^n(x)\lambda^{p-n}(p-1)
  =\lambda^p(x+p-1).
\]

\sgaStableTarget{sga6.e05.para.3-3}\textbf{3.3.} If \(f:K\to K'\) is a homomorphism of \(\lambda\)-rings, then \(f\) is a \(\lambda\)-homomorphism if and only if, for every \(i\in\mathbb N\),
\[
  f(\gamma^i(x))=\gamma^i(f(x)),
\]
or equivalently
\[
  \gamma_t(f(x))=\Ghat(f)(\gamma_t(x)).
\]

\sgaStableTarget{sga6.e05.para.3-4}\textbf{3.4.} From \hyperlink{sga6.e05.prop.3-2}{3.2} one obtains
\sgaStableTarget{sga6.e05.eq.3-4-1}\begin{equation}\tag{3.4.1}
  \gamma^0(x)=1,
  \qquad \gamma^1(x)=x,
  \qquad \gamma^n(x+y)=\sum_{p+q=n}\gamma^p(x)\gamma^q(y).
\end{equation}
Thus \(\gamma_t\) defines a new pre-\(\lambda\)-ring structure. It is not, in general, a \(\lambda\)-ring structure. Let \(\rho_t\) denote the automorphism of the functor \(\Ghat\) induced by the change of variable
\[
  t\longmapsto \frac{t}{1-t}.
\]
Then \(\gamma_t=\rho_t\lambda_t\). Transporting the \(\lambda\)-ring functor structure of \(\Ghat\) by \(\rho_t\) gives a second structure on \(\Ghat\), for which \(\gamma_t:K\to\Ghat(K)\) is a \(\lambda\)-homomorphism whenever \(K\) is a \(\lambda\)-ring.

\sgaStableTarget{sga6.e05.para.3-5}\textbf{3.5.} The multiplication in this transported structure is the composite
\[
  \Ghat\times\Ghat\xrightarrow{\rho_t^{-1}\times\rho_t^{-1}}
  \Ghat\times\Ghat\xrightarrow{\circ}\Ghat\xrightarrow{\rho_t}\Ghat .
\]
It will be denoted \(\nabla\). It is determined by
\sgaStableTarget{sga6.e05.eq.3-5-1}\begin{align}
 (1+Xt)\nabla(1+Yt)
 &=\rho_t\!\left(\rho_t^{-1}(1+Xt)\circ\rho_t^{-1}(1+Yt)\right) \notag\\
 &=\rho_t\!\left(\frac{1+(X+1)t}{1+t}\circ
                 \frac{1+(Y+1)t}{1+t}\right) \notag\\
 &=\rho_t\!\left(\frac{(1+((X+1)(Y+1))t)(1+t)}
                      {(1+(X+1)t)(1+(Y+1)t)}\right) \notag\\
 &=\frac{1+(1+X+Y+XY)t}{(1+Xt)(1+Yt)} . \tag{3.5.1}
\end{align}
Thus, for \(x,y\in K\),
\sgaStableTarget{sga6.e05.eq.3-5-2}\begin{equation}\tag{3.5.2}
  \gamma_t(xy)=\gamma_t(x)\nabla\gamma_t(y).
\end{equation}
The universal polynomial of degree \(n\) for this law will be denoted \(P'_n\). Its isobaric components have weights between \(n\) and \(2n\) in each of the two groups of variables, and
\sgaStableTarget{sga6.e05.eq.3-5-3}\begin{equation}\tag{3.5.3}
  \gamma^n(xy)
  =P'_n(\gamma^1(x),\ldots,\gamma^n(x),
       \gamma^1(y),\ldots,\gamma^n(y)).
\end{equation}
The unit for \(\nabla\) is \(\gamma_t(1)=1/(1-t)\).

\sgaStableTarget{sga6.e05.para.3-6}\textbf{3.6.} The operations \(\lambda'\) of the transported structure are defined by
\sgaStableTarget{sga6.e05.eq.3-6-1}\[
\begin{tikzcd}[column sep=tiny]
\Ghat \\[-0.6em]
1+K[[t]]^+ \arrow[r,"\lambda_u"] \arrow[d,"\rho_t"'] &
1+\{1+K[[t]]^+\}[[u]]^+ \arrow[d,"\Ghat(\rho_t)"]\\
1+K[[t]]^+ \arrow[r,"\lambda'_u"'] &
1+\{1+K[[t]]^+\}[[u]]^+ .
\end{tikzcd}
\tag{3.6.1}
\]
If \(\gamma'\) denotes the \(\gamma\)-operations attached to these operations \(\lambda'\), then
\sgaStableTarget{sga6.e05.eq.3-6-2}\begin{equation}\tag{3.6.2}
  \gamma'_u(1+Xt)=\frac1{1-u}+\{1+Xt\}u .
\end{equation}
Let \(Q'_{i,j}\) be the coefficient of \(t^j\) in \(\gamma'^i(1+X_1t+X_2t^2+\cdots)\). It is a polynomial in \(\mathbb Z[X_1,\ldots,X_{ij}]\), whose isobaric components have weights from \(j\) to \(ij\). For a \(\lambda\)-ring \(K\) and \(x\in K\),
\sgaStableTarget{sga6.e05.eq.3-6-3}\begin{equation}\tag{3.6.3}
  \gamma_t(\lambda^i(x))=\lambda'^i(\gamma_t(x)),
  \qquad
  \gamma_t(\gamma^i(x))=\gamma'^i(\gamma_t(x)),
\end{equation}
and therefore
\sgaStableTarget{sga6.e05.eq.3-6-4}\begin{equation}\tag{3.6.4}
  \gamma^j(\gamma^i(x))
  =Q'_{i,j}(\gamma^1(x),\ldots,\gamma^{ij}(x)).
\end{equation}
The transported pre-\(\lambda\)-ring structure on \(\Ghat\) is itself a \(\lambda\)-ring structure.

\sgaStableTarget{sga6.e05.para.3-7}\textbf{3.7.} We denote by \(\Gcirc\) the functor \(\Ghat\) with the structure of §2, and by \(\Gnabla\) the same functor with the structure of §3. Giving a \(\lambda\)-ring structure on a ring \(A\) is equivalent either to giving a ring homomorphism \(A\to\Gcirc(A)\) satisfying the usual diagram for \(\lambda_t\), or to giving a ring homomorphism \(A\to\Gnabla(A)\) satisfying the corresponding diagram for \(\gamma_t\). Explicitly, the two diagrams are
\sgaStableTarget{sga6.e05.eq.3-7-1}\[
\begin{tikzcd}[column sep=tiny]
A \arrow[r,"\lambda_u"] \arrow[d,"\lambda_t"'] & \Gcirc^u(A) \arrow[d,"\Gcirc^u(\lambda_t)"]\\
\Gcirc^t(A) \arrow[r,"\lambda_u"'] & \Gcirc^u(\Gcirc^t(A))
\end{tikzcd}
\tag{3.7.1}
\]
and
\sgaStableTarget{sga6.e05.eq.3-7-2}\[
\begin{tikzcd}[column sep=tiny]
A \arrow[r,"\gamma_u"] \arrow[d,"\gamma_t"'] & \Gnabla^u(A) \arrow[d,"\Gnabla^u(\gamma_t)"]\\
\Gnabla^t(A) \arrow[r,"\gamma'_u"'] & \Gnabla^u(\Gnabla^t(A)).
\end{tikzcd}
\tag{3.7.2}
\]

\sgaStableTarget{sga6.e05.para.3-8}\textbf{3.8.} Let \(K\) be a binomial ring. From \hyperlink{sga6.e05.eq.2-6-1}{(2.6.1)}, for \(x\in K\),
\sgaStableTarget{sga6.e05.eq.3-8-1}\begin{equation}\tag{3.8.1}
  \gamma_t(x)=\left(\frac1{1-t}\right)^x
  =\exp(-x\Log(1-t)).
\end{equation}
For \(a\in K\), multiplication by \(\gamma_t(a)\) for the law \(\nabla\) is characterized by its value on \(1+Xt\):
\sgaStableTarget{sga6.e05.eq.3-8-2}\begin{equation}\tag{3.8.2}
  \gamma_t(a)\nabla(1+Xt)
  =(1+Xt)^a
  =\exp(a\Log(1+Xt))
  =\sum_{n=0}^{\infty}\binom{a}{n}X^nt^n .
\end{equation}
Thus the universal polynomial of degree \(n\) defining this multiplication is again \(R_{n,a}\).

\sgaStableTarget{sga6.e05.def.3-9}\textbf{Definition 3.9.} Let \(K\) be a \(\lambda\)-ring. An \emph{augmented \(K\)-\(\lambda\)-algebra} is a \(K\)-\(\lambda\)-algebra \(A\), together with a \(\lambda\)-homomorphism of \(K\)-algebras
\[
  \epsilon:A\longrightarrow K.
\]
Morphisms are \(\lambda\)-homomorphisms compatible with the augmentations.

\sgaStableTarget{sga6.e05.ex.3-9-1}\textbf{Example 3.9.1.} Let \((X,A)\) be a locally ringed topos. The Grothendieck group \(K^\circ(X)\) of locally free finite-type \(A\)-modules has a natural pre-\(\lambda\)-ring structure by exterior powers. There is a natural \(\lambda\)-homomorphism from \(H^0(X,\mathbb Z)\) to \(K^\circ(X)\): a section \(a\) of the constant sheaf \(\mathbb Z\) partitions the final object into opens \(U_n\) on which \(a=n\), and is sent to the class of the locally free module which is \(A^n\) on \(U_n\) for \(n\ge0\), minus the analogous class on the negative pieces. The rank map is the augmentation. This is the augmented structure used below.

\sgaStableTarget{sga6.e05.para.3-10}\textbf{3.10.} Let \(A\) be an augmented \(K\)-\(\lambda\)-algebra. For \(n>0\), define \(\Fil^n A\) to be the ideal generated by products
\sgaStableTarget{sga6.e05.eq.3-10-1}\begin{equation}\tag{3.10.1}
  \gamma^{i_1}(x_1)\cdots\gamma^{i_k}(x_k),
  \qquad i_1+\cdots+i_k\ge n,
  \qquad \epsilon(x_i)=0.
\end{equation}
If \(A\) is generated as an abelian group by \((x_\alpha)\), then \(\Fil^nA\) is generated over \(K\) by
\sgaStableTarget{sga6.e05.eq.3-10-2}\begin{equation}\tag{3.10.2}
  \gamma^{i_1}(x_{\alpha_1}-\epsilon(x_{\alpha_1}))\cdots
  \gamma^{i_k}(x_{\alpha_k}-\epsilon(x_{\alpha_k})),
  \qquad i_1+\cdots+i_k\ge n.
\end{equation}
The filtration is decreasing and multiplicative, and
\[
  \Fil^0A=A,
  \qquad \Fil^1A=\Ker\epsilon.
\]
It is functorial for homomorphisms of augmented \(K\)-\(\lambda\)-algebras.

\sgaStableTarget{sga6.e05.prop.3-11}\textbf{Proposition 3.11.} Let \(A\) be a \(\lambda\)-ring and let \(y_1,\ldots,y_n\) be elements of \(A\) such that \(\gamma^j(y_i)=0\) for \(j\ge2\). Put \(x=\sum_i y_i\). Then \(\gamma^k(x-n)=0\) for \(k>n\), and, for \(1\le k\le n\), \(\gamma^k(x-n)\) is the \(k\)-th elementary symmetric function of the \(y_i-1\). This follows from
\[
  \gamma_t(x-n)=\prod_{i=1}^{n}\gamma_t(y_i-1).
\]

\sgaStableTarget{sga6.e05.sec.4}\section*{4. \(\lambda\)-rings generated by generators and relations}

\sgaStableTarget{sga6.e05.para.4-1}\textbf{4.1.} Let \(K\) be a \(\lambda\)-ring and let \((X_i)_{i\in I}\) be indeterminates. Put
\[
  A=K[X_i,{}^nX_i]_{i\in I,\ n\ge2}.
\]
Define
\sgaStableTarget{sga6.e05.eq.4-1-1}\begin{equation}\tag{4.1.1}
  \lambda_t(X_i)=1+X_it+\sum_{n\ge2}{}^nX_i t^n,
  \qquad
  \lambda_t({}^nX_i)=\lambda^n(\lambda_t(X_i)).
\end{equation}
This gives a \(\lambda\)-ring structure on \(A\) extending that of \(K\).

\sgaStableTarget{sga6.e05.prop.4-2}\textbf{Proposition 4.2.} For every \(K\)-\(\lambda\)-algebra \(B\) and every family \((b_i)_{i\in I}\) in \(B\), there exists a unique \(K\)-\(\lambda\)-homomorphism \(A\to B\) carrying \(X_i\) to \(b_i\). Necessarily it carries \({}^nX_i\) to \(\lambda^n(b_i)\), and this condition is sufficient.

\sgaStableTarget{sga6.e05.para.4-3}\textbf{4.3.} The same construction may be made using \(\gamma\)'s. Let
\[
  A'=K[X_i,{}^\gamma_nX_i]_{i\in I,\ n\ge2},
\]
and define
\sgaStableTarget{sga6.e05.eq.4-3-1}\begin{equation}\tag{4.3.1}
  \gamma_t(X_i)=1+X_it+\sum_{n\ge2}{}^\gamma_nX_i t^n,
  \qquad
  \gamma_t({}^\gamma_nX_i)=\gamma'^n(\gamma_t(X_i)).
\end{equation}
This gives a \(\lambda\)-ring satisfying the same universal property, hence canonically isomorphic to \(A\).

\sgaStableTarget{sga6.e05.def.4-4}\textbf{Definition 4.4.} The preceding \(\lambda\)-ring is the \emph{free \(\lambda\)-ring} over \(K\) generated by the family \((X_i)\).

\sgaStableTarget{sga6.e05.def.4-5}\textbf{Definition 4.5.} An ideal \(J\) of a \(\lambda\)-ring \(A\) is a \emph{\(\lambda\)-ideal} if \(\lambda^n(x)\in J\) for every \(x\in J\) and \(n>1\). Kernels of \(\lambda\)-homomorphisms are \(\lambda\)-ideals; conversely, if \(J\) is a \(\lambda\)-ideal, \(A/J\) receives a unique \(\lambda\)-ring structure. The intersection of \(\lambda\)-ideals is a \(\lambda\)-ideal. Thus every family of elements has a smallest containing \(\lambda\)-ideal, generated by the elements and all their iterated \(\lambda^n\)'s.

\sgaStableTarget{sga6.e05.def.4-6}\textbf{Definition 4.6.} The quotient of the free \(\lambda\)-ring by the \(\lambda\)-ideal generated by prescribed relations \((P_i)\) is called the \(\lambda\)-ring generated by the generators \((X_i)\) subject to the relations \((P_i)\). Equivalently, one may formulate the same construction in terms of \(\gamma\)-operations.

\sgaStableTarget{sga6.e05.prop.4-7}\textbf{Proposition 4.7.} This quotient has the expected universal property: every family \((b_i)\) in a \(K\)-\(\lambda\)-algebra satisfying the relations \((P_i)\) determines a unique \(K\)-\(\lambda\)-homomorphism from the quotient to that algebra.

\sgaStableTarget{sga6.e05.prop.4-8}\textbf{Proposition 4.8.} Let \(n\ge0\). In the free \(\lambda\)-ring over \(K\) on one generator \(X\), the ideal generated by the \(\lambda^kX\), \(k>n\), is a \(\lambda\)-ideal. Indeed, after quotienting, one has only to check that if \(\varphi\in\Ghat(C)\) has no terms of degree \(>n\), then \(\lambda^k(\varphi)=0\) for \(k>n\). By the splitting principle, it is enough to write
\[
  \varphi=\prod_{i=1}^{n}(1+T_it),
\]
and then
\[
  \lambda_u(\varphi)=\prod_{i=1}^{n}\bigl(1+\{1+T_it\}u\bigr),
\]
which has degree \(n\) in \(u\).

\sgaStableTarget{sga6.e05.cor.4-9}\textbf{Corollary 4.9.} The \(\lambda\)-ring generated over \(K\) by one element \(X\), subject to \(\lambda^k(X)=0\) for \(k>n\), is
\[
  K[X,\lambda^2X,\ldots,\lambda^nX].
\]
More generally, for a family \((X_i)\) with bounds \((n_i)\), the quotient subject to \(\lambda^k(X_i)=0\) for \(k>n_i\) is the polynomial ring generated by \(X_i,\lambda^2X_i,\\ldots,\lambda^{n_i}X_i\).

\sgaStableTarget{sga6.e05.para.4-10}\textbf{4.10.} Let \(K\) be binomial and let \(A\) be an augmented \(K\)-\(\lambda\)-algebra. Put
\[
  B=A[X,\gamma^2X,\gamma^3X,\ldots],
\]
the free \(A\)-\(\lambda\)-algebra on \(X\), written in \(\gamma\)-coordinates. Extend the augmentation by \(\epsilon(X)=0\), hence \(\epsilon(\gamma^iX)=0\) for \(i>0\). Let \(M_k\) be the set of monomials in the variables \(\gamma^jX\), counted with weight \(j\), of total weight \(k\), and set \(M_0=\{1\}\).

\sgaStableTarget{sga6.e05.lem.4-11}\textbf{Lemma 4.11.} If \(x\in\Fil^pA\), then every coefficient of \(\gamma^n(x+X)\) which is homogeneous of weight \(k\) in the \(\gamma^jX\) belongs to \(\Fil^{p+n-k}A\).

\sgaStableTarget{sga6.e05.prop.4-12}\textbf{Proposition 4.12.} For every \(k\ge0\),
\sgaStableTarget{sga6.e05.eq.4-12-1}\begin{equation}\tag{4.12.1}
  \Fil^kB=
  \Fil^kA+\Fil^{k-1}A\,M_1+\cdots+\Fil^1A\,M_{k-1}
  +\Fil^0A\,M_k+B\,M_{k+1}.
\end{equation}
The inclusion from right to left follows from \hyperlink{sga6.e05.lem.4-11}{Lemma 4.11}. The reverse inclusion is obtained by writing the generators of \(\Fil^kB\) as products of \(\gamma\)-operations and reducing, term by term, to the displayed weights.

\sgaStableTarget{sga6.e05.rem.4-13}\textbf{Remarks 4.13.} The proof gives the same assertion for a free \(\lambda\)-algebra on any finite or infinite family of generators, with the evident multi-weight. It is also unchanged when one uses the \(\lambda\)-coordinates, since \(\lambda\) and \(\gamma\) coordinates differ by triangular universal formulas.

\sgaStableTarget{sga6.e05.cor.4-14}\textbf{Corollary 4.14.} Passing to the associated graded ring, one has
\sgaStableTarget{sga6.e05.eq.4-14-1}\begin{equation}\tag{4.14.1}
  \Gr^kB\simeq
  \Gr^kA\oplus \Gr^{k-1}A\,M_1\oplus\cdots\oplus\Gr^0A\,M_k.
\end{equation}

\sgaStableTarget{sga6.e05.cor.4-15}\textbf{Corollary 4.15.} If \(B'=K[X,\gamma^2X,\ldots,\gamma^nX]\) is the \(K\)-\(\lambda\)-algebra generated by one element satisfying \(\gamma^k(X)=0\) for \(k>n\), then the \(\lambda\)-ring filtration of \(B'\) is exactly the weight filtration in the variables \(\gamma^jX\).

\sgaStableTarget{sga6.e05.sec.5}\section*{5. The elements \(\lambda^p(N,x)\)}

This section will be used only in \hyperlink{sga6.e07.expose}{Exposé VII}, §4, in the study of regular immersions; it may be omitted on first reading.

\sgaStableTarget{sga6.e05.para.5-1}\textbf{5.1.} Let \(d\ge1\). Let \(A\) be the \(\mathbb Z\)-\(\lambda\)-algebra generated by \(N\) and \(x\), subject to the relations \(\lambda^i(N)=0\) for \(i>d\). Thus
\[
  A=\mathbb Z[N,\ldots,\lambda^dN;
       x,\ldots,\lambda^jx,\ldots].
\]
Put
\sgaStableTarget{sga6.e05.eq.5-1-1}\begin{equation}\tag{5.1.1}
  \lambda_{-1}(N)=\sum_{i=0}^{d}(-1)^i\lambda^i(N).
\end{equation}

\sgaStableTarget{sga6.e05.lem.5-2}\textbf{Lemma 5.2.} For every \(p\ge1\), the element \(\lambda^p(x\lambda_{-1}(N))\) is divisible by \(\lambda_{-1}(N)\).

It is enough, by isobaricity, to prove this for \(\lambda^p(\lambda_{-1}(N))\). Let \(B\) be the \(\lambda\)-ring generated by \(N_1,\ldots,N_d\), with \(\lambda^i(N_j)=0\) for \(i\ge2\). The map \(N\mapsto\sum_iN_i\) identifies \(A\) with the symmetric polynomials in the \(N_i\), and
\[
  \lambda_{-1}(N)=\prod_{i=1}^{d}(1-N_i).
\]
If some \(N_i=1\), then \(\lambda_t(\lambda_{-1}(N))=\lambda_t(0)=1\), so every positive coefficient vanishes. Thus each \(\lambda^p(\lambda_{-1}(N))\) is divisible by every \(1-N_i\), hence by their product.

\sgaStableTarget{sga6.e05.def.5-3}\textbf{Definition 5.3.} Since \(A\) is integral, define \(\lambda^p(N,x)\) by
\sgaStableTarget{sga6.e05.eq.5-3-1}\begin{equation}\tag{5.3.1}
  \lambda^p(N,x)\lambda_{-1}(N)=\lambda^p(x\lambda_{-1}(N)),
  \qquad p\ge1.
\end{equation}
Thus \(\lambda^p(N,x)\) is a polynomial in the \(\lambda^i(x)\) and \(\lambda^j(N)\), isobaric of weight \(p\) in the \(\lambda^i(x)\). If a \(\lambda\)-ring \(B\) contains an element \(N'\) with \(\lambda^k(N')=0\) for \(k>d\), then \(\lambda^p(N',x')\) is defined by substitution; it is independent of the chosen bound \(d\).

\sgaStableTarget{sga6.e05.para.5-4}\textbf{5.4.} In the \(\lambda\)-ring generated over \(\mathbb Z\) by \(N,x,y\), subject to \(\lambda^k(N)=0\) for \(k>d\), formula \hyperlink{sga6.e05.eq.5-3-1}{(5.3.1)} gives
\begin{align*}
 \lambda^p(N,x+y)\lambda_{-1}(N)
 &=\lambda^p\bigl((x+y)\lambda_{-1}(N)\bigr)\\
 &=\sum_{i+j=p}
   \lambda^i\bigl(x\lambda_{-1}(N)\bigr)
   \lambda^j\bigl(y\lambda_{-1}(N)\bigr)\\
 &=\sum_{\substack{i+j=p\\ i,j\ne0}}
   \lambda^i(N,x)\lambda^j(N,y)\bigl(\lambda_{-1}(N)\bigr)^2
   +\bigl(\lambda^p(N,x)+\lambda^p(N,y)\bigr)\lambda_{-1}(N).
\end{align*}
Since the ring under consideration is an integral domain, this yields the relation, valid by substitution in every \(\lambda\)-ring,
\sgaStableTarget{sga6.e05.eq.5-4-1}\begin{equation}\tag{5.4.1}
 \lambda^p(N,x+y)
 =\sum_{\substack{i+j=p\\ i,j\ne0}}
   \lambda^i(N,x)\lambda^j(N,y)\lambda_{-1}(N)
   +\lambda^p(N,x)+\lambda^p(N,y).
\end{equation}
Now consider the \(\lambda\)-ring generated over \(\mathbb Z\) by \(N,N',x\), subject to \(\lambda^k(N)=0\) for \(k>d\) and \(\lambda^k(N')=0\) for \(k>d'\). Then \(\lambda^k(N+N')=0\) for \(k>d+d'\), and
\begin{align*}
 \lambda^p(N+N',x)\lambda_{-1}(N)\lambda_{-1}(N')
 &=\lambda^p\bigl(x\lambda_{-1}(N)\lambda_{-1}(N')\bigr)\\
 &=\lambda^p\bigl(N',x\lambda_{-1}(N)\bigr)\lambda_{-1}(N').
\end{align*}
Here we used \(\lambda_{-1}(N+N')=\lambda_{-1}(N)\lambda_{-1}(N')\). Cancelling in the integral domain gives the relation, again valid by substitution,
\sgaStableTarget{sga6.e05.eq.5-4-2}\begin{equation}\tag{5.4.2}
  \lambda^p(N+N',x)\lambda_{-1}(N)
  =\lambda^p\bigl(N',x\lambda_{-1}(N)\bigr).
\end{equation}

\sgaStableTarget{sga6.e05.para.5-5}\textbf{5.5.} Let \(A\) be a \(\lambda\)-ring and let \(N\in A\) satisfy \(\lambda^k(N)=0\) for \(k>d\). Define on the underlying abelian group of \(A\) a new multiplication by
\[
  x\cdot_N y=xy\lambda_{-1}(N).
\]
After adjoining a unit, one obtains a ring \(\mathbb Z\times A\) with
\[
  (n,x)(m,y)=(nm,ny+mx+x\cdot_N y).
\]
The formulas
\sgaStableTarget{sga6.e05.eq.5-5-1}\begin{equation}\tag{5.5.1}
  \lambda_t(n,x)=(1+t)^n\left(1+
  \sum_{p\ge1}\lambda^p(N,x)t^p\right)
\end{equation}
define a pre-\(\lambda\)-ring structure. In fact, when \(A\) is a \(\lambda\)-ring, this is a \(\lambda\)-ring structure; it is the one transported from the original structure by multiplication by \(\lambda_{-1}(N)\). We denote this augmented \(\lambda\)-ring by \(\widehat A_N\). Its \(\gamma\)-operations satisfy
\sgaStableTarget{sga6.e05.eq.5-5-2}\begin{equation}\tag{5.5.2}
  \gamma^p(N,x)\lambda_{-1}(N)=\gamma^p(x\lambda_{-1}(N)).
\end{equation}

\sgaStableTarget{sga6.e05.para.5-6}\textbf{5.6.} Suppose that, in addition, \(L\) is an element such that \(\lambda^k(L)=0\) for \(k>1\). Then \(\lambda^p(N+L,x)\) is the coefficient of \(t^p\) in the series
\sgaStableTarget{sga6.e05.eq.5-6-1}\begin{equation}\tag{5.6.1}
  \frac{\displaystyle\sum_{k\ge1}\lambda^k(N,x)
        (1+L+\cdots+L^{k-1})t^k}
       {\displaystyle 1+\sum_{k\ge1}\lambda^k(N,x)
        \lambda_{-1}(N)L^k t^k}.
\end{equation}
This is obtained by applying \hyperlink{sga6.e05.eq.5-4-2}{(5.4.2)} to \(N+L\) and using \(\lambda_{-1}(L)=1-L\).

\sgaStableTarget{sga6.e05.prop.5-7}\textbf{Proposition 5.7.} Let \(K\) be a binomial ring, \(A\) an augmented \(K\)-\(\lambda\)-algebra, and \(N\in A\) such that \(\lambda^k(N)=0\) for \(k>d\). Then, for every \(x\in A\) and every \(p\ge1\),
\sgaStableTarget{sga6.e05.eq.5-7-1}\begin{equation}\tag{5.7.1}
  \gamma^p(N,x)\in\Fil^{p-d}A.
\end{equation}
Indeed, let \(A_0\) be the \(K\)-\(\lambda\)-algebra generated by \(N'_0,X'_0\), subject to \(\gamma^k(N'_0)=0\) for \(k>d\), and augmented by \(\epsilon(N'_0)=\epsilon(X'_0)=0\). This is a polynomial algebra over \(K\) in the variables \(\gamma^i(N'_0)\), \(1\le i\le d\), and \(\gamma^j(X'_0)\), \(j\ge1\), and its \(\lambda\)-filtration is the weight filtration, with weight \(i\) assigned to \(\gamma^i(N'_0)\) and \(\gamma^i(X'_0)\) (\hyperlink{sga6.e05.prop.4-12}{4.12} and \hyperlink{sga6.e05.cor.4-15}{4.15}). Put
\[
  N_0=N'_0+d,
  \qquad X_0=X'_0+\epsilon(x).
\]
For \(N_0,X_0\), the assertion is equivalent to
\[
  \gamma^p(N_0,X_0)\gamma^d(N'_0)\in\Fil^pA_0.
\]
Since \(\gamma^d(N'_0)=(-1)^d\lambda_{-1}(N_0)\), this becomes
\[
  (-1)^d\gamma^p\bigl(X_0\gamma^d(N'_0)\bigr)\in\Fil^pA_0,
\]
by \hyperlink{sga6.e05.eq.5-5-2}{(5.5.4)}. This holds because \(\gamma^d(N'_0)\) has zero augmentation. By \hyperlink{sga6.e05.prop.4-7}{4.7} there is a \(K\)-\(\lambda\)-homomorphism \(A_0\to A\), compatible with augmentations and sending \(N'_0\) to \(N-d\) and \(X'_0\) to \(x-\epsilon(x)\), hence sending \(N_0\) to \(N\) and \(X_0\) to \(x\). Compatibility with the filtrations gives the result.

\sgaStableTarget{sga6.e05.sec.6}\section*{6. Chern ring: opening}

Throughout this section the base ring is a fixed binomial ring \(K\). We use the functor \(G\) of \hyperlink{sga6.e05.para.1-1}{1.1}, and return to the notation of §1: \(G\) denotes the functor there defined and \(G_i\) its restriction to \(\mathscr C_0\).

\sgaStableTarget{sga6.e05.para.6-1}\textbf{6.1.} In \hyperlink{sga6.e05.para.3-5}{3.5} and \hyperlink{sga6.e05.para.3-6}{3.6} we defined morphisms of functors
\[
  G_i\times G_i\xrightarrow{\star}G_i,
  \qquad
  G_i^\circ\xrightarrow{\gamma'}G_i^\circ\circ G_S,
  \qquad
  G_i\xrightarrow{a}G_i,
\]
where the last map is multiplication, for the law \(\nabla\), by \(\gamma_t(a)\) with \(a\in K\). By \hyperlink{sga6.e05.lem.1-9}{1.9} these give morphisms
\[
  G\times G\xrightarrow{\star}G,
  \qquad
  G\xrightarrow{\gamma'}G\circ G_S,
  \qquad
  G\xrightarrow{a}G.
\]
The isobaricity of the universal polynomials involved implies that the first and third images lie in \(G\), while the second lies in the subgroup of formal series in \(u\) whose non-zero terms have the prescribed degrees.



\providecommand{\Ghat}{\widehat G}
\providecommand{\Aqhat}{\widehat A_{\mathbb Q}}
\providecommand{\Ch}{\operatorname{Ch}}
\providecommand{\Log}{\operatorname{Log}}
\providecommand{\id}{\operatorname{id}}
\providecommand{\gr}{\operatorname{gr}}
\providecommand{\Hom}{\operatorname{Hom}}
\providecommand{\epsi}{\varepsilon}
\providecommand{\clgr}{\operatorname{cl}}
\providecommand{\Fil}{\operatorname{Fil}}
\providecommand{\Gr}{\operatorname{Gr}}
\providecommand{\Sym}{\operatorname{Sym}}


\sgaStableTarget{sga6.e05.expose.occ2}\section*{Exposé V, 6. Chern ring: completion}

\sgaStableTarget{sga6.e05.def.6-2}\textbf{Definition 6.2.} The functor defined in \hyperlink{sga6.e05.para.6-1}{6.1} on the category of graded \(K\)-algebras, with values in the category of \(K\)-\(\lambda\)-algebras, is called the \emph{Chern ring}; its value on \(A\in\operatorname{Ob}(\mathscr C)\) is denoted \(\Ch(A)\), or, if there is a risk of confusion, \(\Ch_K(A)\).

From formulas \hyperlink{sga6.e05.para.6-1}{(6.1.1)}, \hyperlink{sga6.e05.para.6-1}{(6.1.2)}, \hyperlink{sga6.e05.para.6-1}{(6.1.4)}, and \hyperlink{sga6.e05.eq.3-6-2}{(3.6.2)}, one deduces the relations
\sgaStableTarget{sga6.e05.eq.6-2-1}\begin{equation}\tag{6.2.1}
 (1,1+X)(1,1+Y)=(1,1+X+Y),
 \qquad
 \lambda^i(1,1+X)=0\quad\text{for }i\ge2,
\end{equation}
the relations which served in [1] to define the Chern ring. Conversely, from these formulas one can recover the relations \hyperlink{sga6.e05.para.6-1}{(6.1.1)} and \hyperlink{sga6.e05.para.6-1}{(6.1.4)}.

Let us indicate why the Chern ring is introduced. We place ourselves in the situation of \hyperlink{sga6.e05.ex.3-9-1}{Example 3.9}, and write \(A\) for the graded ring associated, by the filtration defined in \hyperlink{sga6.e05.para.3-10}{3.10}, with the \(K^\bullet\) of a commutatively locally ringed topos; for simplicity suppose that \(K^\bullet\) is a \(\lambda\)-ring. The augmentation ring is \(K\). We shall see below that for every element of an augmented \(\lambda\)-ring one can define Chern classes with values in the associated graded ring \hyperlink{sga6.e05.para.6-7}{(6.7)}. To every locally free sheaf of finite type one can therefore associate an element of \(\Ch(A)\), whose component on the factor \(K\) is the rank of the sheaf and whose component on the factor \(1+\widehat A^+\) is the sum of the Chern classes. Formulas \hyperlink{sga6.e05.eq.6-2-1}{(6.2.1)} are then the usual formulas giving the Chern classes and the rank of the tensor product of two invertible sheaves and of the exterior powers of an invertible sheaf; in other words, the map just defined from \(K^\bullet\) to \(\Ch(A)\) is a \(\lambda\)-homomorphism.

\sgaStableTarget{sga6.e05.para.6-3}\textbf{6.3.} For every graded \(K\)-algebra \(A\), denote by \(A_{\mathbb Q}\) the \(K\)-algebra \(A\otimes\mathbb Q\), graded in the evident way. By \hyperlink{sga6.e05.thm.1-13}{1.13} one defines a functorial homomorphism
\[
  1+\widehat A^+\xrightarrow{\ \ch\ }\widehat A_{\mathbb Q}^{+}
\]
by giving its value on \(1+X\):
\sgaStableTarget{sga6.e05.eq.6-3-1}\begin{equation}\tag{6.3.1}
  \ch(1+X)=\exp(X)-1.
\end{equation}
This homomorphism is compatible with the multiplication defined in \hyperlink{sga6.e05.para.6-1}{6.1} on \(1+\widehat A^+\). By \hyperlink{sga6.e05.thm.1-13}{1.13} it is enough to show this for the elements \(1+X\), \(1+Y\) of the universal ring \(1+K[X,Y]^+\). One has
\[
\begin{aligned}
 \ch\bigl((1+X)*(1+Y)\bigr)
 &= \ch\!\left(\frac{1+X+Y}{(1+X)(1+Y)}\right)\\
 &= \ch(1+X+Y)-\ch(1+X)-\ch(1+Y)\\
 &= \exp(X+Y)-\exp(X)-\exp(Y)+1\\
 &= \ch(1+X)\,\ch(1+Y).
\end{aligned}
\]

Let \(\eta\) be the endomorphism of \(\widehat A_{\mathbb Q}^{+}\), functorial in \(A\), defined by
\[
 \eta\!\left(\sum_{k\ge1} a_k t^k\right)
 =
 \sum_{k\ge1}\frac{(-1)^{k-1}}{(k-1)!}\,a_k t^k.
\]
Then, for every \(\varphi\in1+\widehat A^+\),
\sgaStableTarget{sga6.e05.eq.6-3-2}\begin{equation}\tag{6.3.2}
  \ch(\varphi)=\eta(\Log\varphi).
\end{equation}
Indeed, \(\ch\) and \(\eta\circ\Log\) are two functorial homomorphisms of abelian groups, and it is enough for them to agree on \(1+X\), which is evident.

It follows that \(\ch\) commutes with the operations of \(K\): again it is enough to show this for \(1+X\), and for \(a\in K\) one has
\[
\begin{aligned}
 \ch((1+X)^a)
 &=\ch(\exp(a\Log(1+X)))\\
 &=\eta(a\Log(1+X))\\
 &=a\,\eta(\Log(1+X))\\
 &=a\,\ch(1+X).
\end{aligned}
\]
Thus one can extend \(\ch\) in the evident way to a homomorphism of \(K\)-algebras
\[
  K\times(1+\widehat A^+)\longrightarrow K\oplus\widehat A_{\mathbb Q}^{+},
\]
namely
\sgaStableTarget{sga6.e05.eq.6-3-3}\begin{equation}\tag{6.3.3}
  \ch(a,\varphi)=a+\eta(\Log\varphi).
\end{equation}

\sgaStableTarget{sga6.e05.def.6-4}\textbf{Definition 6.4.} The homomorphism
\[
  \ch:\Ch A\longrightarrow K\oplus\widehat A_{\mathbb Q}^{+}
\]
just defined is called the \emph{Chern character}.

If \(A\) is of characteristic zero, that is, is a \(\mathbb Q\)-algebra, the Chern character is an isomorphism. Indeed, if \(\psi\in K\oplus\widehat A^+\), one may write \(\psi=a+\psi_1\), where \(\psi_1\) is of positive degree, and then an inverse morphism is defined by
\[
  \psi\longmapsto (a,\exp(\eta^{-1}(\psi_1))).
\]

\sgaStableTarget{sga6.e05.prop.6-5}\textbf{Proposition 6.5.} Let \(A\) be a graded \(K\)-algebra. On \(\Ch(A)\) consider the filtration defined as follows, the degree filtration:
\sgaStableTarget{sga6.e05.eq.6-5-1}\begin{equation}\tag{6.5.1}
  (\Ch(A))_0=\Ch(A),\qquad
  (\Ch(A))_n=0\times(1+\widehat A^{\ge n})\quad(n\ge1),
\end{equation}
where \(\widehat A^{\ge n}=\prod_{i\ge n}A^i\). The following congruences hold.

\begin{enumerate}[label=\roman*)]
\item If \(n\ge1\) and \(\varphi\in(\Ch A)_n\), writing \(\varphi=1+a_n+\cdots\), then
\sgaStableTarget{sga6.e05.eq.6-5-2}\begin{equation}\tag{6.5.2}
  \ch(\varphi)=\frac{(-1)^{n-1}}{(n-1)!}\,a_n+\cdots .
\end{equation}

\item If \(\varphi\in(\Ch A)_m\), \(\psi\in(\Ch A)_n\), and if
\(\varphi=1+a_m+\cdots\), \(\psi=1+b_n+\cdots\), then
\sgaStableTarget{sga6.e05.eq.6-5-3}\begin{equation}\tag{6.5.3}
  \varphi*\psi
  =
  1-\frac{(m+n-1)!}{(m-1)!(n-1)!}\,a_m b_n+\cdots .
\end{equation}
Consequently the degree filtration is a ring filtration.
\end{enumerate}

For (i), the Chern character is by construction a morphism of filtered functors; from the congruence \(\varphi\equiv1+a_n\pmod{(\Ch A)_{n+1}}\) one obtains
\[
  \ch(\varphi)=\ch(1+a_n)\pmod{\widehat A^{\ge n+1}},
\]
and, since \(\ch(1+a_n)=\eta(\Log(1+a_n))\), this is
\[
  \frac{(-1)^{n-1}}{(n-1)!}a_n+\cdots .
\]
For (ii), by substitution it is enough to prove \hyperlink{sga6.e05.eq.6-5-3}{(6.5.3)} for the torsion-free universal polynomial algebra
\[
 A=K[X_i,Y_j]_{i,j\in\mathbb N}
\]
and for \(\varphi=1+\sum_{k\ge m}X_k\), \(\psi=1+\sum_{k\ge n}Y_k\). Since \(K\) is torsion-free in the universal case, one can embed this algebra in its tensor product with \(\mathbb Q\), so that \(\ch\) is injective. Put \(\varphi*\psi=1+Z_p+\cdots\), where \(Z_p\) is the first non-zero term different from \(1\). The multiplicativity of \(\ch\), together with \hyperlink{sga6.e05.eq.6-5-2}{(6.5.2)}, gives
\[
  1+\frac{(-1)^{p-1}}{(p-1)!}Z_p+\cdots
  =
  \left(1+\frac{(-1)^{m-1}}{(m-1)!}X_m+\cdots\right)
  \left(1+\frac{(-1)^{n-1}}{(n-1)!}Y_n+\cdots\right),
\]
which gives the stated value of \(Z_p\) by comparing the terms of lowest degree.

\sgaStableTarget{sga6.e05.para.6-6}\textbf{6.6.} The \(\lambda\)-ring \(\Ch A\) is augmented over \(K\): one defines a \(\lambda\)-homomorphism \(\Ch A\to K\) by \((a,\varphi)\mapsto a\). It is therefore equipped with the canonical filtration of an augmented \(\lambda\)-ring. We compare it with the filtration just defined.

\sgaStableTarget{sga6.e05.prop.6-6-1}\textbf{Proposition 6.6.1.} Let \(x\in\Ch A\) be an element of the form \(x=(n,1+\cdots+a_n)\). Then:
\[
\begin{array}{ll}
\text{i)}& \lambda^i(x)=0\quad\text{for }i>n,\\[0.3em]
\text{ii)}& \lambda^n(x)=(1,1+a_1),\\[0.3em]
\text{iii)}& \displaystyle\sum_{i=0}^n(-1)^i\lambda^i(x)
       =(0,1-(n-1)!a_n+\cdots),
\end{array}
\]
where the unwritten terms have degree \(>n\).

By functoriality, it is enough to prove this in the case where
\(A=K[X_1,\ldots,X_n]\), filtered by weight with respect to the \(X_i\), \(X_i\) having weight \(i\), and for the element \((n,1+X_1+\cdots+X_n)\). Embed \(A\) in \(A'=K[T_1,\ldots,T_n]\) by the symmetric functions. Then
\[
 (n,1+X_1+\cdots+X_n)
 =
 \left(n,\prod_{i=1}^n(1+T_i)\right)
 =
 \sum_{i=1}^n(1,1+T_i).
\]
Hence, by \hyperlink{sga6.e05.eq.6-2-1}{(6.2.1)},
\[
  \lambda_t(x)=\prod_{i=1}^n\bigl(1+(1,1+T_i)t\bigr),
\]
so that
\[
  \lambda^k(x)=0\quad(k>n),\qquad
  \lambda^n(x)=\prod_{i=1}^n(1,1+T_i)
   =\left(1,1+\sum_{i=1}^nT_i\right)=(1,1+X_1).
\]
Finally,
\[
  \sum_{i=1}^n(-1)^i\lambda^i(x)=\lambda_{-1}(x)
  =\prod_{i=1}^n\bigl(1-(1,1+T_i)\bigr)
  =(-1)^n\prod_{i=1}^n(0,1+T_i).
\]
By \hyperlink{sga6.e05.eq.6-5-3}{(6.5.3)},
\[
  \prod_{i=1}^n(0,1+T_i)
  =
  \left(0,1+(-1)^{n-1}(n-1)!\prod_{i=1}^nT_i+\cdots\right),
\]
hence
\[
  \lambda_{-1}(x)=(0,1-(n-1)!X_n+\cdots).
\]

\sgaStableTarget{sga6.e05.cor.6-6-2}\textbf{Corollary 6.6.2.} Relation \hyperlink{sga6.e05.prop.6-6-1}{6.6.1(iii)} is true for every element \(X\) of augmentation \(n\).

For this we shall use the following lemma.

\sgaStableTarget{sga6.e05.lem.6-6-3}\textbf{Lemma 6.6.3.} For every \(n\), \((\Ch A)_n\) is a \(\lambda\)-ideal of \(\Ch A\).

We may suppose \(n>1\). Let \(x=(0,\varphi)\in(\Ch A)_n\), and write
\[
  \lambda^i(0,\varphi)=(0,1+Q^i_{i,1}+\cdots+Q^i_{i,j}+\cdots),
\]
where the \(Q^i_{i,j}\) are the isobaric components of weight \(j\) of the \(Q^i_i\) of \hyperlink{sga6.e05.para.3-6}{(3.6)}. For every \(j\), the polynomial \(Q^i_{i,j}\) is a polynomial in the first \(j\) coefficients of \(\varphi\). If \(j<n\), these coefficients are zero; hence the \(Q^i_{i,j}\) are zero, and so the \(\lambda^i(x)\) belong to \((\Ch A)_n\). Thus \((\Ch A)_n\) is a \(\lambda\)-ideal.

Let \(x=(n,1+a_1+\cdots+a_n+\cdots)\) be an element of augmentation \(n\) in \(\Ch A\). Then
\[
  x\equiv (n,1+\cdots+a_n)\pmod{(\Ch A)_{n+1}},
\]
and therefore
\[
  \lambda_{-1}(x)\equiv \lambda_{-1}(n,1+\cdots+a_n)
  =(0,1-(n-1)!a_n+\cdots),
\]
which proves the corollary.

\sgaStableTarget{sga6.e05.cor.6-6-4}\textbf{Corollary 6.6.4.} Let \(\epsilon\) be the augmentation homomorphism of \(\Ch A\). Then, for every \(x\in\Ch A\), one has
\[
  \gamma^n(x-\epsilon(x))
  =
  (0,1+(-1)^{n-1}(n-1)!a_n+\cdots).
\]
By definition,
\[
  \gamma^n(x-\epsilon(x))
  =
  (-1)^n\sum_{i=0}^n(-1)^i\lambda^i(x+n-\epsilon(x)).
\]
Since \(x+n-\epsilon(x)\) has augmentation \(n\), \hyperlink{sga6.e05.cor.6-6-2}{Corollary 6.6.2} applies to it, giving the result.

\sgaStableTarget{sga6.e05.cor.6-6-5}\textbf{Corollary 6.6.5.} With the notation of \hyperlink{sga6.e05.para.3-10}{3.10},
\[
  \Fil^n\Ch A\subset(\Ch A)_n.
\]
This follows from the preceding corollary, which gives \(\gamma^n(x-\epsilon(x))\in(\Ch A)_n\), and from the fact that the two filtrations are ring filtrations and that the elements \(\gamma^n(x-\epsilon(x))\), for \(n\) and \(x\) variable, generate the \(\lambda\)-ring filtration.

\sgaStableTarget{sga6.e05.cor.6-6-7}\textbf{Corollary 6.6.7.} Let \(A\) be a graded \(K\)-algebra of characteristic zero, bounded by degree. Then the two filtrations defined on \(\Ch A\) are identical.

Suppose that \(\Fil^p A=0\), and let \(x=(0,1+a_n+\cdots+a_p)\) be an element of \((\Ch A)_n\). Put
\[
  x'=\left(0,1+\frac{(-1)^{n-1}}{(n-1)!}a_n+a_{n+1}+\cdots+a_p\right).
\]
Then \(\gamma^n(x')=(0,1+a_n+\cdots+a'_p)\), since \(x'\) has zero augmentation. Hence \(x-\gamma^n(x')\in(\Ch A)_{n+1}\). Iterating the operation \(p-n\) times, one reaches \((\Ch A)_{p+1}=0\), and therefore \(x\) is written as an element of \(\Fil^n\Ch A\). In view of \hyperlink{sga6.e05.cor.6-6-5}{(6.6.5)}, the two filtrations are identical.

\sgaStableTarget{sga6.e05.para.6-7}\textbf{6.7.} Now suppose that \(A\) is a \(\lambda\)-ring augmented over \(K\), and let \(\epsilon\) be the augmentation homomorphism. Then \(A\) is equipped with its canonical filtration of augmented \(\lambda\)-ring, and one can associate to \(\Gr A\) its Chern ring.

If \(x\in A\), then \(\gamma_t(x-\epsilon(x))\in\Ghat(A)\) by definition of the filtration of an augmented \(\lambda\)-ring. We write \(c(x)\) for the image of \(\gamma_t(x-\epsilon(x))\) in
\[
  \Ghat(A)/\Ghat'(A)=1+\widehat{\Gr A}^{+}.
\]
It is called the \emph{total Chern class} of \(x\). We write \(c^i(x)\), and call it the \(i\)-th Chern class of \(x\), for the \(i\)-th coefficient of \(c(x)\). Thus
\sgaStableTarget{sga6.e05.eq.6-7-1}\begin{equation}\tag{6.7.1}
  c^i(x)=\cl_{\Gr^iA}\!\left(\gamma^i(x-\epsilon(x))\right),
  \qquad
  c(x)=1+c^1(x)+\cdots .
\end{equation}
Put
\sgaStableTarget{sga6.e05.eq.6-7-2}\begin{equation}\tag{6.7.2}
  \widetilde c(x)=(\epsilon(x),c(x)).
\end{equation}
This is called the \emph{completed Chern class} of \(x\).

\sgaStableTarget{sga6.e05.prop.6-8}\textbf{Proposition 6.8.} The completed Chern class
\[
  \widetilde c:A\longrightarrow\Ch(\Gr A)
\]
is a \(K\)-\(\lambda\)-homomorphism of augmented \(K\)-\(\lambda\)-algebras.

It is clear that the homomorphism so defined is additive. To verify that it is multiplicative, write
\[
\begin{aligned}
\gamma_t(xy-\epsilon(xy))
&=\gamma_t\bigl(\epsilon(y)(x-\epsilon(x))\bigr)
  \gamma_t\bigl(\epsilon(x)(y-\epsilon(y))\bigr)
  \gamma_t\bigl((x-\epsilon(x))(y-\epsilon(y))\bigr)\\
&=\bigl(\gamma_t(x-\epsilon(x))\bigr)^{\epsilon(y)}
  \bigl(\gamma_t(y-\epsilon(y))\bigr)^{\epsilon(x)}
  \bigl(\gamma_t(x-\epsilon(x))\circ\gamma_t(y-\epsilon(y))\bigr),
\end{aligned}
\]
which gives the result by \hyperlink{sga6.e05.para.6-1}{(6.1.2)} after passage to the quotient. If \(a\in K\), then \(\widetilde c(a)=(a,1)\), so \(\widetilde c\) is a homomorphism of \(K\)-algebras; moreover it is compatible with augmentations.

It remains to see that it is a \(\lambda\)-homomorphism. One must prove the commutativity of the diagram
\[
\begin{tikzcd}[column sep=large,row sep=large]
A \arrow[r,"\gamma_t"] \arrow[d,"\widetilde c"'] &
\Ghat(A) \arrow[d,"\Ghat(\widetilde c)"]\\
\Ch A \arrow[r,"\gamma_t"'] &
\Ghat(\Ch A).
\end{tikzcd}
\]
By additivity, it is enough to prove it for the elements of \(K\), where it is evident, and for the elements of zero augmentation in \(A\), where it follows, by the definition of \(c\) and \hyperlink{sga6.e05.para.6-1}{(6.1.4)}, from diagram \hyperlink{sga6.e05.eq.3-7-2}{(3.7.2)}.

\sgaStableTarget{sga6.e05.prop.6-9}\textbf{Proposition 6.9.} Let \(A\) be an augmented \(K\)-\(\lambda\)-algebra and let \(z\in\Fil^mA\). Then
\[
  \gamma^m(z)=(-1)^{m-1}(m-1)!z\pmod{\Fil^{m+1}A},
\]
that is,
\[
  c^m(z)=\cl_{\Gr^mA}\!\left((-1)^{m-1}(m-1)!z\right).
\]

One must see that these two elements give the same image in \(\Gr^mA\), that is, that the \(m\)-th Chern class of \(z\) is the image of \((-1)^{m-1}(m-1)!z\).

One may write \(z\) as a sum
\[
  \sum a\,\gamma^{i_1}(x_1)\cdots\gamma^{i_k}(x_k),
\]
with \(a\in A\), \(\epsilon(x_i)=0\) for all \(i\), and \(i_1+\cdots+i_k\ge m\). We then have
\[
  \widetilde c(z)=
  \sum \widetilde c(a)\,
        \widetilde c(\gamma^{i_1}(x_1))\cdots
        \widetilde c(\gamma^{i_k}(x_k)).
\]
Since \(\widetilde c\) is a \(\lambda\)-homomorphism, one may write
\[
  \widetilde c(\gamma^i(x_j))=\gamma^i(\widetilde c(x_j)).
\]
Taking account of the definition of \(c\) and of \hyperlink{sga6.e05.cor.6-6-4}{(6.6.4)}, one obtains
\[
  \widetilde c(\gamma^{i_j}(x_j))
  =
  \left(0,1+(-1)^{i_j-1}(i_j-1)!\,\gamma^{i_j}(x_j)+\cdots\right).
\]
By \hyperlink{sga6.e05.eq.6-5-3}{(6.5.3)}, and since \(\widetilde c(a)=(a,1)\), one finally gets
\[
\begin{aligned}
\widetilde c(z)
&=\sum
\left(0,\left(1+
(-1)^{i_1+\cdots+i_k-1}
(i_1+\cdots+i_k-1)!
\,\gamma^{i_1}(x_1)\cdots\gamma^{i_k}(x_k)
+\cdots\right)^a\right)\\
&=
\left(0,1+
\sum_{\substack{(x_1,\ldots,x_k)\\ i_1+\cdots+i_k=m}}
(-1)^{m-1}(m-1)!a\,\gamma^{i_1}(x_1)\cdots\gamma^{i_k}(x_k)
+\cdots\right)\\
&=
(0,1+(-1)^{m-1}(m-1)!z+\cdots),
\end{aligned}
\]
which proves the assertion.

\sgaStableTarget{sga6.e05.prop.6-10}\textbf{Proposition 6.10.} Let \(K\) be a binomial ring, \(R\) an augmented \(K\)-\(\lambda\)-algebra, \(N\in R\) an element such that \(\lambda^k(N)=0\) for \(k>d\), and \(x\in\Fil^iR\). Then
\[
  \gamma^{d+i}(N,x)-(-1)^{d+i-1}(d+i-1)!x\in\Fil^{i+1}R.
\]

One can find elements \(z_{\alpha,j}\in R\) and \(a_\alpha\in K\) such that
\[
  x=\sum_\alpha a_\alpha\,
       \gamma^{i_1}(x_{\alpha,1})\cdots
       \gamma^{i_k}(x_{\alpha,k}),
\]
with \(i_1+\cdots+i_k\ge i\) and \(\epsilon(x_{\alpha,j})=0\) for every \(\alpha\) and every \(j\). Let \(R_0\) be the \(K\)-\(\lambda\)-algebra generated by \(N'_0\) and \(X_{\alpha,j}\), subject to the relations
\[
  \gamma^k(N'_0)=0\quad\text{for }k>d.
\]
Define an augmentation homomorphism \(R_0\to K\) by \(\epsilon(N'_0)=\epsilon(X_{\alpha,j})=0\). By \hyperlink{sga6.e05.cor.4-9}{(4.9.1)}, \(R_0\) is the ring of polynomials with coefficients in \(K\) in the indeterminates \(\gamma^k(N'_0)\), \(1\le k\le d\), and \(\gamma^p(X_{\alpha,j})\), \(p\ge1\); the \(\lambda\)-filtration of \(R_0\) is its weight filtration, where \(\gamma^k(N'_0)\) and \(\gamma^k(X_{\alpha,j})\) have weight \(k\) (\hyperlink{sga6.e05.prop.4-12}{4.12} and \hyperlink{sga6.e05.cor.4-15}{4.15}).

Put
\[
  X_0=\sum_\alpha a_\alpha\,
       \gamma^{i_1}(X_{\alpha,1})\cdots
       \gamma^{i_k}(X_{\alpha,k}).
\]
We shall prove \hyperlink{sga6.e05.prop.6-10}{6.10} first for \(X_0\) and \(N'_0=N_0+d\), which do satisfy the hypotheses of 6.10. In view of the preceding remark about the filtration of \(R_0\), it is enough to show that
\[
  \gamma^{d+i}(N_0,X_0)\gamma^d(N'_0)
  -
  (-1)^{d+i-1}(d+i-1)!X_0\gamma^d(N'_0)
  \in\Fil^{d+i+1}R_0 .
\]
But \(\lambda_{-1}(N_0)=(-1)^d\gamma^d(N'_0)\). Thus this relation becomes
\[
  (-1)^d\gamma^{d+i}(X_0,(-1)^d\gamma^d(N'_0))
  -
  (-1)^{d+i-1}(d+i-1)!X_0\gamma^d(N'_0)
  \in\Fil^{d+i+1}R_0 .
\]
Putting \(z=(-1)^dX_0\gamma^d(N'_0)\), we have \(z\in\Fil^{d+i}R_0\), and the statement is reduced to
\[
  \gamma^{d+i}(z)-(-1)^{d+i-1}(d+i-1)!z\in\Fil^{d+i+1}R_0,
\]
which is just 6.9.

To deduce \hyperlink{sga6.e05.prop.6-10}{6.10} from this particular case, consider the \(K\)-\(\lambda\)-homomorphism \(R_0\to R\) which sends \(N'_0\) to \(N-d\) and \(X_{\alpha,j}\) to \(x_{\alpha,j}\), for every \(\alpha\) and \(j\). This homomorphism is compatible with the \(\lambda\)-filtrations and sends \(N_0\) to \(N\) and \(X_0\) to \(x\). Hence the result follows.

\sgaStableTarget{sga6.e05.prop.6-11}\textbf{Proposition 6.11.} Suppose that \(K\) is of characteristic zero, and let \(n\ge0\) be an integer. Then the functor \(\Ch\) is an equivalence from the category of graded \(K\)-algebras \(A\) such that \(A_0=K\) and \(A_k=0\) for \(k>n\), to the category of augmented \(K\)-\(\lambda\)-algebras \(\Lambda\) such that \(\Fil^{n+1}\Lambda=0\). A quasi-inverse functor is given by \(\Lambda\mapsto\Gr\Lambda\).

Let \(A\) be a graded \(K\)-algebra such that \(A_k=0\) for \(k>n\), equipped with the filtration associated to its grading. Then \((\Ch A)_{n+1}=0\), and hence \(\Fil^{n+1}\Ch A=0\) by \hyperlink{sga6.e05.cor.6-6-5}{(6.6.5)}. Conversely, if \(\Lambda\) is an augmented \(K\)-\(\lambda\)-algebra such that \(\Fil^{n+1}\Lambda=0\), then \(\Gr^0\Lambda=K\) and \(\Gr^k\Lambda=0\) for \(k>n\).

We shall define isomorphisms of functors
\[
  A\longrightarrow \Gr(\Ch A),
  \qquad
  \Lambda\longrightarrow\Ch(\Gr\Lambda).
\]
Since \(A\) is of finite grading, the associated gradeds of the \(\lambda\)-ring filtration on \(\Ch A\) and of the degree filtration are the same by \hyperlink{sga6.e05.cor.6-6-7}{(6.6.7)}. Define
\[
  u:A\longrightarrow \Gr_{\deg}(\Ch A)
\]
to be the identity on \(A_0\), and, for \(i\ge1\) and \(a_i\in A_i\),
\[
  u(a_i)=
  \cl_{\Gr^i}\left(0,1+(-1)^{i-1}(i-1)!a_i\right),
\]
then extend by additivity. By relation \hyperlink{sga6.e05.eq.6-5-3}{(6.5.3)} this is a homomorphism of graded rings, and it is easy to see that it is compatible with the operations of \(K\). The inverse morphism \(u^{-1}\) may be defined as follows: for \(i\ge1\) and \(\bar\varphi\in\Gr^i(\Ch A)\), the \(i\)-th coefficient \(a_i\) of \(\varphi\) is independent of the lift of \(\bar\varphi\) in \((\Ch A)_i\), and one puts
\[
  u^{-1}(\bar\varphi)=\frac{(-1)^{i-1}}{(i-1)!}a_i.
\]
Thus \(u\) is an isomorphism.

The isomorphism \(\Lambda\to\Ch(\Gr\Lambda)\) is defined by the completed Chern class \(\widetilde c\). To verify that it really defines an isomorphism, it is enough to verify that the homomorphism obtained by passing to associated graded rings is one, the filtrations being discrete. One gets
\[
  \Gr\Lambda
   \xrightarrow{\Gr(\widetilde c)}
  \Gr_{\lambda}(\Ch(\Gr\Lambda))
   \simeq
  \Gr_{\deg}(\Ch(\Gr\Lambda))
   \xrightarrow{u^{-1}}
  \Gr\Lambda .
\]
It remains to show that the composite is the identity. Since the homomorphisms are homomorphisms of graded \(K\)-algebras, it is enough, by the definition of the \(\lambda\)-ring filtration of \(\Lambda\), to show that it is the identity on the \(\gamma^i(x)\), where \(\epsilon(x)=0\). One then obtains
\[
  \Gr(\widetilde c)\bigl(\gamma^i(x)\bigr)
  =
  \cl_{\Gr^i}\bigl(\gamma^i(c(x))\bigr).
\]
But \(\gamma^i(\widetilde c(x))=(0,1+(-1)^{i-1}(i-1)!c^i(x)+\cdots)\), by \hyperlink{sga6.e05.cor.6-6-4}{(6.6.4)}. Therefore
\[
  \cl_{\Gr^i}\bigl(\gamma^i(\widetilde c(x))\bigr)=u(c^i(x)),
\]
which shows that \(u^{-1}\circ\Gr(\widetilde c)=\id\).

\sgaStableTarget{sga6.e05.cor.6-12}\textbf{Corollary 6.12.} If \(K\) is of characteristic zero, and if \(\Lambda\) is a \(K\)-\(\lambda\)-algebra with discrete filtration, then the completed Chern class
\[
  \widetilde c:\Lambda\longrightarrow \Ch(\Gr\Lambda)
\]
is an isomorphism.

The proof is contained in that of 6.11.

\sgaStableTarget{sga6.e05.def.6-13}\textbf{Definition 6.13.} Let \(A\) be a graded \(K\)-algebra and \(\Lambda\) an augmented \(K\)-\(\lambda\)-algebra. A \emph{Chern theory for \(\Lambda\) with values in \(A\)} is a \(\lambda\)-homomorphism
\[
  \Lambda\longrightarrow \Ch(A)
\]
compatible with augmentations.

Thus one defines a functor on the category of graded \(K\)-algebras, with values in the category of sets, by associating with a graded \(K\)-algebra \(A\) the set of Chern theories of \(\Lambda\) with values in \(A\), namely
\[
  \Hom_\lambda(\Lambda,\Ch(A)).
\]

\sgaStableTarget{sga6.e05.prop.6-14}\textbf{Proposition 6.14.} The functor defined in \hyperlink{sga6.e05.def.6-13}{6.13} is representable. If \(K\) is of characteristic zero and \(\Lambda\) has discrete filtration, it is represented by the graded ring associated with \(\Lambda\).

First prove the first assertion. Consider the graded \(K\)-algebra generated by generators \(c^i(x)\), with \(i\ge1\), \(x\in\Fil^1\Lambda\), \(c^i(x)\) of degree \(i\), and take its quotient by the following relations:
\[
\begin{array}{ll}
\text{i)}&
 c^k(x+y)=\displaystyle\sum_{i+j=k}c^i(x)c^j(y),
 \quad k\ge1,\ x,y\in\Fil^1\Lambda,\\[0.8em]
\text{ii)}&
 c^k(xy)=P''_k(c^i(x),c^j(y))
 \quad(\text{cf. }\hyperlink{sga6.e05.para.6-1}{6.1}),
 \quad k\ge1,\ x,y\in\Fil^1\Lambda,\\[0.4em]
\text{iii)}&
 c^k(ax)=R_{k,a}(c^i(x))
 \quad(\text{cf. }\hyperlink{sga6.e05.para.6-1}{6.1}),
 \quad k\ge1,\ a\in K,\ x\in\Fil^1\Lambda,\\[0.4em]
\text{iv)}&
 c^k(\gamma^j(x))=Q''_{j,k}(c^i(x))
 \quad(\text{cf. }\hyperlink{sga6.e05.para.6-1}{6.1}),
 \quad j,k\ge1,\ x\in\Fil^1\Lambda.
\end{array}
\]
These relations are isobaric, so the quotient is a graded \(K\)-algebra \(\Omega\). Define a morphism
\[
  \widetilde c_0:\Lambda\longrightarrow \Ch(\Omega)
\]
by
\[
  x\longmapsto
  \bigl(\epsilon(x),1+c^1(x-\epsilon(x))+\cdots\bigr).
\]
The imposed relations show that this is a Chern theory for \(\Lambda\) with values in \(\Omega\). Conversely, if \(f\) is a Chern theory with values in a graded \(K\)-algebra \(A\), then there is a unique homomorphism \(\Omega\to A\) through which \(f\) factors through \(\Ch(\Omega)\): it sends, for every \(i\ge1\) and every \(x\in\Fil^1\Lambda\), the element \(c^i(x)\) to the \(i\)-th coefficient of \(f(x)\). This is well defined because \(f\) is a \(\lambda\)-homomorphism.

Now suppose that \(K\) is of characteristic zero and \(\Lambda\) has discrete filtration. Let \(n\) be such that \(\Fil^n\Lambda=0\). Then the graded algebra \(\Omega\) is zero in degrees \(>n\). To show this, it is enough to prove that, for
\[
  x_1,\ldots,x_p\in\Fil^1\Lambda
  \quad\text{and}\quad
  k_1+\cdots+k_p\ge n,
\]
one has
\[
  c^{k_1}(x_1)\cdots c^{k_p}(x_p)=0.
\]
By hypothesis,
\[
  \gamma^{k_1}(x_1)\cdots\gamma^{k_p}(x_p)=0.
\]
Applying the \(\lambda\)-homomorphism \(\widetilde c_0\), one gets
\[
  \gamma^{k_1}(\widetilde c_0(x_1))\cdots
  \gamma^{k_p}(\widetilde c_0(x_p))=(0,1).
\]
By \hyperlink{sga6.e05.cor.6-6-4}{(6.6.4)} and \hyperlink{sga6.e05.eq.6-5-3}{(6.5.3)}, this implies
\[
  (k_1+\cdots+k_p)!\,
  c^{k_1}(x_1)\cdots c^{k_p}(x_p)=0.
\]
Since \(K\) is of characteristic zero, the desired result follows.

The universal property of \(\Omega\), applied to the canonical homomorphism \(\widetilde c\), gives a commutative diagram
\[
\begin{tikzcd}[column sep=large,row sep=large]
\Lambda \arrow[r,"\widetilde c"] \arrow[d,"\widetilde c_0"'] &
\Ch(\Gr\Lambda)\\
\Ch(\Omega) \arrow[ur] &
\end{tikzcd}
\]
By \hyperlink{sga6.e05.prop.6-11}{Proposition 6.11}, the homomorphism \(\Omega\to\Gr\Lambda\) is an isomorphism if and only if the homomorphism
\[
  \Ch(\Omega)\longrightarrow \Ch(\Gr\Lambda)
\]
is one. Since it is visibly surjective, the latter is also surjective; it remains to show that
\(\Ch(\Omega)\to\Ch(\Gr\Lambda)\) is injective. By functoriality of \(\Ch\), one obtains the diagram
\[
\begin{tikzcd}[column sep=large,row sep=large]
\Lambda \arrow[r,"\widetilde c_0"] \arrow[d,"\widetilde c"'] &
\Ch(\Omega) \arrow[d,"\widetilde c"]\\
\Ch(\Gr\Lambda) \arrow[r,"\Ch(\widetilde c_0)"'] &
\Ch\bigl(\Gr_\Lambda(\Ch(\Omega))\bigr).
\end{tikzcd}
\]
Since \(\Omega\) has bounded degrees, the right vertical arrow is an isomorphism by \hyperlink{sga6.e05.cor.6-12}{(6.12)}. Using \hyperlink{sga6.e05.prop.6-11}{6.11}, one therefore obtains a morphism
\[
  \Gr\Lambda\longrightarrow\Omega
\]
whose two composites with \(\Omega\to\Gr\Lambda\) are the identity, by the universal property of \(\Omega\).

\sgaStableTarget{sga6.e05.sec.7}\section*{7. Appendix: the Adams operations \(\psi^k\)}

We follow here a presentation due to J.-P. Serre.

\sgaStableTarget{sga6.e05.para.7-1}\textbf{7.1.} Let \(K\) be a pre-\(\lambda\)-ring. For \(x\in K\), write
\[
  \frac{d}{dt}(\lambda_t(x))
\]
for the derivative with respect to \(t\) of \(\lambda_t(x)\). Define elements of \(K\), denoted \(\psi^k(x)\), by putting
\sgaStableTarget{sga6.e05.eq.7-1-1}\begin{equation}\tag{7.1.1}
  \frac{d}{dt}(\lambda_t(x))/\lambda_t(x)
  =
  \sum_{k=1}^{\infty}(-1)^{k-1}\psi^k(x)t^{k-1}.
\end{equation}
The maps \(\psi^k:K\to K\) are additive. Indeed,
\[
\begin{aligned}
 \frac{d}{dt}(\lambda_t(x+y))/\lambda_t(x+y)
 &=
 \frac{d}{dt}(\lambda_t(x)\lambda_t(y))/\lambda_t(x)\lambda_t(y)\\
 &=
 \frac{d}{dt}(\lambda_t(x))/\lambda_t(x)
 +\frac{d}{dt}(\lambda_t(y))/\lambda_t(y),
\end{aligned}
\]
and therefore, for every \(k\),
\[
  \psi^k(x+y)=\psi^k(x)+\psi^k(y).
\]
By \hyperlink{sga6.e05.eq.7-1-1}{(7.1.1)}, the \(\psi^k\) are isobaric polynomials of weight \(k\) in the \(\lambda^i\), \(i=1,\ldots,k\). For example,
\[
  \psi^1(x)=\lambda^1(x)=x,
  \qquad
  \psi^2(x)=x^2-2\lambda^2(x).
\]

\sgaStableTarget{sga6.e05.para.7-2}\textbf{7.2.} Let \(K\) be a ring of characteristic zero. Suppose given additive maps \(\psi^k:K\to K\) such that \(\psi^1=\id_K\). Define a pre-\(\lambda\)-structure on \(K\) by
\sgaStableTarget{sga6.e05.eq.7-2-1}\begin{equation}\tag{7.2.1}
  \lambda_t(x)=
  \exp\left(\sum_{k=1}^{\infty}(-1)^{k-1}\psi^k(x)t^k/k\right).
\end{equation}
Multiplicativity follows from the additivity of the \(\psi^k\). Thus there is an equivalence between giving a pre-\(\lambda\)-structure on \(K\) and giving the corresponding operations \(\psi^k\).

\sgaStableTarget{sga6.e05.para.7-3}\textbf{7.3.} It is clear from formula \hyperlink{sga6.e05.eq.7-1-1}{(7.1.1)} that the \(\psi^k\) commute with \(\lambda\)-homomorphisms. If \(K,K'\) are pre-\(\lambda\)-rings of characteristic zero, it is necessary and sufficient for a homomorphism \(f:K\to K'\) to be a \(\lambda\)-homomorphism that it commute with the \(\psi^k\). This follows from \hyperlink{sga6.e05.para.7-2}{7.2}, since the \(\lambda^i\) are expressed as polynomials in the \(\psi^k\).

\sgaStableTarget{sga6.e05.ex.7-4}\textbf{Examples 7.4.} i) For \(K=\mathbb Z\), with its \(\lambda\)-ring structure, one has
\[
  \frac{d}{dt}(\lambda_t(n))/\lambda_t(n)
  =
  n(1+t)^{n-1}/(1+t)^n
  =
  n\sum_{k=1}^{\infty}(-1)^{k-1}t^{k-1}.
\]
Consequently the \(\psi^k\) are all equal to the identity map of \(\mathbb Z\).

ii) Work in the category \(\mathscr C_0\) of \(\mathbb Z\)-algebras \hyperlink{sga6.e05.para.1-8}{(1.8)}. By \hyperlink{sga6.e05.para.7-3}{7.3}, the operations \(\psi^k\) define additive endomorphisms of the functor \(G\) on \(\lambda\)-rings. By \hyperlink{sga6.e05.thm.1-10}{1.10} they are characterized by their value on \(1+Xt\). One has
\[
  \frac{d}{dt}(\lambda_u(1+Xt))/\lambda_u(1+Xt)
  =
  \frac{1+Xt}{1+\{1+Xt\}u}
  =
  \sum_{k=1}^{\infty}(-1)^{k-1}\{1+Xt\}^{\circ k}u^{k-1}.
\]
Thus
\sgaStableTarget{sga6.e05.eq.7-4-1}\begin{equation}\tag{7.4.1}
  \psi^k(1+Xt)=1+X^k t.
\end{equation}

\sgaStableTarget{sga6.e05.prop.7-5}\textbf{Proposition 7.5.} i) If \(K\) is a \(\lambda\)-ring, the operations \(\psi^k\) satisfy the identities
\sgaStableTarget{sga6.e05.eq.7-5-1}\begin{equation}\tag{7.5.1}
  \psi^k(1)=1,\qquad
  \psi^k(xy)=\psi^k(x)\psi^k(y),\qquad
  \psi^k\circ\psi^{k'}=\psi^{kk'}.
\end{equation}
ii) Conversely, if the \(\psi^k\) satisfy these identities and if \(K\) is of characteristic zero, then \(K\) is a \(\lambda\)-ring.

For (i), if \(K\) is a \(\lambda\)-ring, \(\lambda_t(1)=1+t\), hence
\[
  \frac{d}{dt}(\lambda_t(1))/\lambda_t(1)
  =
  \frac1{1+t}
  =
  \sum_{k=1}^{\infty}(-1)^{k-1}t^{k-1},
\]
and therefore \(\psi^k(1)=1\).

To prove the second relation of \hyperlink{sga6.e05.eq.7-5-1}{(7.5.1)}, consider the diagram
\sgaStableTarget{sga6.e05.eq.7-5-2}\begin{equation}\tag{7.5.2}
\begin{tikzcd}[column sep=large,row sep=large]
K\times K \arrow[r,"\lambda_t\times\lambda_t"] \arrow[d] &
\widehat G(K)\times\widehat G(K) \arrow[r,"g^k\times g^k"] \arrow[d,"\circ"] &
K\times K \arrow[d]\\
K \arrow[r,"\lambda_t"'] &
\widehat G(K) \arrow[r,"g^k"'] &
K ,
\end{tikzcd}
\end{equation}
where the vertical arrows are defined by multiplication, and the maps \(g^k\) as follows: the logarithmic derivative is a group homomorphism from \(G(K)\) to \(K[[t]]\), and the map which to a series assigns its \((k-1)\)-st coefficient, multiplied by \((-1)^{k-1}\), is an additive homomorphism from \(K[[t]]\) to \(K\); \(g^k\) is their composite. Clearly \(g^k\circ\lambda_t=\psi^k\) by \hyperlink{sga6.e05.eq.7-1-1}{definition (7.1.1)}.

The second of formulas \hyperlink{sga6.e05.eq.7-5-1}{(7.5.1)} therefore expresses the commutativity of the diagram just written. If \(K\) is a \(\lambda\)-ring, the left square is commutative, and it is enough to verify the commutativity of the right square. This is a diagram of functor homomorphisms on \(\mathscr C_0\), with source \(\widehat G\times\widehat G\). By \hyperlink{sga6.e05.thm.1-10}{1.10} it is enough to verify commutativity on \((1+Xt,1+Yt)\). Then
\[
  g^k(1+Xt)=X^k,\qquad
  g^k(1+Yt)=Y^k,
\]
and
\[
  g^k\bigl((1+Xt)\circ(1+Yt)\bigr)
  =
  g^k(1+XYt)
  =
  (XY)^k
  =
  g^k(1+Xt)g^k(1+Yt),
\]
which gives the result.

To prove the last relation, first observe that the diagram
\sgaStableTarget{sga6.e05.eq.7-5-3}\begin{equation}\tag{7.5.3}
\begin{tikzcd}[column sep=large,row sep=large]
\widehat G_t(K) \arrow[r,"\lambda_u"] \arrow[d,"g^k"'] &
\widehat G_u\widehat G_t(K) \arrow[d,"g^k_u"]\\
K \arrow[r,"\lambda_u"'] &
\widehat G^u(K)
\end{tikzcd}
\end{equation}
is commutative. As above, it is enough to see this for \(1+Xt\), and one has
\[
  g^{k'}_u\lambda_u(1+Xt)
  =
  g^{k'}_u(1+\{1+Xt\}u)
  =
  (1+Xt)^{\circ k'}
  =
  1+X^{k'}t,
\]
and then
\[
  g^k(1+X^{k'}t)=X^{kk'}=g^k(1+Xt).
\]
Consequently
\[
  \psi^k\circ\psi^{k'}
  =
  g^k\circ g^{k'}_u\circ\lambda_u\circ\lambda_t
  =
  g^k\circ\lambda_t\circ\psi^{k'}.
\]
Moreover
\[
  \psi^k\circ\psi^{k'}
  =
  g^k\circ\lambda_t\circ\psi^{k'}
  =
  \psi^k\circ\lambda_t\circ\psi^{k'}.
\]
If \(K\) is a \(\lambda\)-ring, the diagram
\sgaStableTarget{sga6.e05.eq.7-5-4}\begin{equation}\tag{7.5.4}
\begin{tikzcd}[column sep=large,row sep=large]
K \arrow[r,"\lambda_t"] \arrow[d,"\psi^{k'}"'] &
\widehat G(K) \arrow[d,"\psi^{k'}"]\\
K \arrow[r,"\lambda_t"'] &
\widehat G(K)
\end{tikzcd}
\end{equation}
is commutative by \hyperlink{sga6.e05.para.7-3}{7.3}, since \(\lambda_t\) is a \(\lambda\)-homomorphism. The desired relations follow at once.

For (ii), if \(K\) is of characteristic zero, the logarithmic derivative is an isomorphism from \(\widehat G(K)\) onto \(K[[t]]\), since one can integrate. Therefore two elements of \(\widehat G(K)\) are equal if and only if, for every \(k\), they have the same image by \(g^k\).

To show that \(\lambda_t\) is a ring homomorphism, it is enough to show that the left square in \hyperlink{sga6.e05.eq.7-5-2}{(7.5.2)} is commutative. By the preceding remark, it is enough to show that, for every \(k\), diagram \hyperlink{sga6.e05.eq.7-5-2}{(7.5.2)} is commutative; the right square always is. But this is exactly the relation
\[
  \psi^k(xy)=\psi^k(x)\psi^k(y).
\]
To show that \(\lambda_t\) is a \(\lambda\)-homomorphism, it is enough, by \hyperlink{sga6.e05.para.7-3}{7.3}, to show that diagram \hyperlink{sga6.e05.eq.7-5-4}{(7.5.4)} is commutative for every \(k'\). By the preceding remark, this is equivalent to showing, for all \(k,k'\), that
\[
  g^k\circ\psi^{k'}\circ\lambda_t
  =
  g^k\circ\lambda_t\circ\psi^{k'},
\]
that is, the relations \(\psi^{kk'}=\psi^k\circ\psi^{k'}\). Finally, the relation \(\lambda_t(1)=1+t\) is obtained by integration from the relations \(\psi^k(1)=1\).

\sgaStableTarget{sga6.e05.bibliography}\subsection*{Bibliography}
\begin{enumerate}[label={[\arabic*]}]
\item A. Grothendieck, \emph{La théorie des classes de Chern}, Bull. Soc. math. France, 86, 1958, p. 137 à 154.
\item F. Hirzebruch, \emph{Neue topologische Methoden in der algebraischen Geometrie}.
\end{enumerate}



\providecommand{\Sym}{\operatorname{Sym}}
\providecommand{\Proj}{\operatorname{Proj}}
\providecommand{\Projz}{\operatorname{Proj}_0}
\providecommand{\Spec}{\operatorname{Spec}}
\providecommand{\Comp}{\operatorname{Comp}}
\providecommand{\Dplus}{D_{+}}
\providecommand{\id}{\operatorname{id}}
\providecommand{\cG}{\mathcal G}
\providecommand{\cA}{\mathcal A}
\providecommand{\cB}{\mathcal B}
\providecommand{\cL}{\mathscr L}
\providecommand{\cM}{\mathscr M}
\providecommand{\cN}{\mathscr N}
\providecommand{\cE}{\mathscr E}
\providecommand{\cK}{\mathscr K}
\providecommand{\cP}{\mathscr P}
\providecommand{\cQ}{\mathscr Q}
\providecommand{\OO}{\mathcal O}
\providecommand{\rk}{\operatorname{rk}}
\providecommand{\colim}{\operatorname*{colim}}


\sgaStableTarget{sga6.e06.expose}\section*{Exposé VI. The \(K^\bullet\) of a projective bundle: computations and consequences}
\begin{center}
Pierre Berthelot
\end{center}

\sgaStableTarget{sga6.e06.sec.1}\section*{1. Computation of \(K^\bullet\) of a projective bundle: the case of locally free sheaves of finite type}

In this section \(K^\bullet(X)\) denotes the Grothendieck group of locally free sheaves of finite type on a scheme \(X\). The aim is to prove the following structure theorem.

\sgaStableTarget{sga6.e06.thm.1-1}\textbf{Theorem 1.1.} Let \(S\) be a quasi-compact scheme, let \(E\) be a locally free \(\mathcal O_S\)-module of rank \(r+1\), and let
\[
  X=\mathbf P(E)
\]
be the associated projective bundle. The structural morphism \(f:X\to S\) gives \(K^\bullet(X)\) the structure of a \(K^\bullet(S)\)-algebra. If \(\xi\in K^\bullet(X)\) denotes the class of \(\mathcal O_X(1)\), then \(K^\bullet(X)\) is a free \(K^\bullet(S)\)-module with basis
\[
  1,\xi,\ldots,\xi^r .
\]

\sgaStableTarget{sga6.e06.sec.1.2}\subsection*{1.2. Plan of the proof}
The proof is divided into three parts.

First, one proves that the powers of \(\xi\) generate \(K^\bullet(X)\) as a \(K^\bullet(S)\)-module. For this it is enough to establish the following facts. If \(F\) is locally free of finite type on \(X\), one can find a graded \(\Sym(E)\)-module \(M\), of finite presentation and flat over \(\mathcal O_S\), such that \(F=\Proj_0(M)\). Such an \(M\) admits a finite graded resolution
\[
  0\longrightarrow L_{r+1}\longrightarrow \cdots \longrightarrow L_0\longrightarrow M\longrightarrow 0
\]
by graded \(\Sym(E)\)-modules locally free of finite type. Finally, each \(L_i\) has a finite filtration whose quotients are of the form
\[
  N\otimes_{\mathcal O_S}\Sym(E)(n),
\]
with \(N\) locally free of finite type on \(S\). Applying \(\Proj_0\) then gives the required generation by the powers of \(\xi\).

Second, one proves that \(1,\xi,\ldots,\xi^{r+1}\) are linearly dependent over \(K^\bullet(S)\). This is obtained from the canonical epimorphism \(f^*E\to\mathcal O_X(1)\).

Third, one proves that \(1,\xi,\ldots,\xi^r\) are linearly independent over \(K^\bullet(S)\).

Let \(S\) be a scheme, let \(B=\bigoplus_{n\geq 0}B_n\) be a graded \(\mathcal O_S\)-algebra, and let \(M\) be a graded \(B\)-module. If \(N\) is a graded \(B\)-module, \(N(n)\) denotes the graded module defined by
\[
  N(n)_k=N_{n+k}.
\]
A graded \(B\)-module is said to be free, respectively locally free, respectively locally free of finite type, if locally on \(S\) it is a direct sum, respectively locally a direct sum, respectively locally a finite direct sum, of modules \(B(n)\). It is said to be of finite presentation if, locally on \(S\), it admits a presentation by free graded modules of finite type.

For every integer \(n\) there is a canonical homomorphism of graded \(B\)-modules
\sgaStableTarget{sga6.e06.eq.1-2-1}\begin{equation}\tag{1.2.1}
  M_n\otimes_{\mathcal O_S}B(-n)\longrightarrow M,
\end{equation}
defined by \(x\otimes a\mapsto ax\). A free graded module of finite type admits a decomposition
\sgaStableTarget{sga6.e06.eq.1-2-2}\begin{equation}\tag{1.2.2}
  M\simeq \bigoplus_{n=n_0}^{n_1} B(-n)^{k_n}.
\end{equation}

\sgaStableTarget{sga6.e06.prop.1-2}\textbf{Proposition 1.2.} The integers \(k_n\) occurring in a decomposition \hyperlink{sga6.e06.eq.1-2-2}{(1.2.2)} are independent of the chosen decomposition.

Indeed, after restricting to an open set where \(B_0=\mathcal O_S\), one compares the first non-zero homogeneous component. This determines \(k_{n_0}\); passing to the quotient and arguing by induction gives all the \(k_n\).

\sgaStableTarget{sga6.e06.prop.1-3}\textbf{Proposition 1.3.} Let \(S\) be quasi-compact, let \(E\) be a locally free \(\mathcal O_S\)-module of finite type, let \(X=\mathbf P(E)\), and let \(F\) be an \(\mathcal O_X\)-module of finite presentation which is flat over \(S\). Then there exists an integer \(n_0\) such that, for all \(n\geq n_0\),
\[
\begin{array}{ll}
\text{\rm (i)}& R^i f_*(F(n))=0\quad (i\geq 1),\\[2mm]
\text{\rm (ii)}& f_*(F(n))\text{ is locally free of finite type on }S,\\[2mm]
\text{\rm (iii)}& \displaystyle M=\bigoplus_{n\geq n_0} f_*(F(n))
\text{ is a graded }\Sym(E)\text{-module of finite presentation.}
\end{array}
\]

The statement is local on \(S\). Since \(S\) is quasi-compact, one is reduced to the affine case. If \(S\) is noetherian, (i) is Serre's theorem, (ii) follows from EGA III, 7.9.10, and (iii) from the usual comparison between a coherent sheaf on a projective bundle and its associated graded module.

In the general affine case one uses passage to the limit. Write \(S=\Spec(A)\) as a filtered projective limit of noetherian affine schemes \(S_\lambda=\Spec(A_\lambda)\), with affine transition morphisms. The locally free module \(E\), the scheme \(X\), and the sheaf \(F\) descend, after increasing \(\lambda\), to corresponding data \(E_\lambda,X_\lambda,F_\lambda\), with \(F_\lambda\) flat over \(S_\lambda\). The noetherian case applies to \(F_\lambda\), and the theorem of change of base gives the assertions after pullback to \(S\).

In this reduction one uses the Cartesian square
\[
\begin{tikzcd}[column sep=large,row sep=large]
X \arrow[r,"g"] \arrow[d,"f"'] & X_0 \arrow[d,"f_0"]\\
S \arrow[r,"g_0"'] & S_0 .
\end{tikzcd}
\]
Here \(S_0\) is noetherian, \(X_0=\mathbf P(E_0)\), and \(F=g^*(F_0)\). The previous paragraph proves the assertion for \(F_0\), and EGA III, 7.7.5 and 7.8.5 give the required base-change identifications for \(F\).

\sgaStableTarget{sga6.e06.cor.1-4}\textbf{Corollary 1.4.} Under the hypotheses of \hyperlink{sga6.e06.prop.1-3}{Proposition 1.3} there exists a graded \(\Sym(E)\)-module \(M\), of finite presentation and flat over \(\mathcal O_S\), such that
\[
  F=\Proj_0(M).
\]

Choose an integer \(n_0\) for which the properties of \hyperlink{sga6.e06.prop.1-3}{Proposition 1.3} hold. For
\[
 M=\bigoplus_{n\ge n_0}f_*(F(n)),
\]
we have seen that \(M\) is of finite presentation; it is flat over \(S\), since the \(f_*(F(n))\), for \(n\ge n_0\), are locally free on \(S\). Moreover, \(F=\Proj_0(M)\) by EGA II 3.4. We now turn to part b) of the proof.

\sgaStableTarget{sga6.e06.lem.1-5-1}\textbf{Lemma 1.5.1.} Let \(A\) be a ring, let \(B\) be a positively graded \(A\)-algebra with \(B_0=A\), and let \(M\) be a graded \(B\)-module. Suppose that \(M\), regarded as an ungraded \(B\)-module, is projective of finite type. Let \(\widetilde B\) be the sheaf of algebras on \(\Spec(A)\) defined by \(B\), and let \(\widetilde M\) be the associated graded \(\widetilde B\)-module. Then, locally on \(\Spec(A)\), \(\widetilde M\) is free of finite type over \(\widetilde B\) as a graded module \hyperlink{sga6.e06.prop.1-2}{(1.2)}.

Indeed, since \(M\) is of finite type, its degrees are bounded below. Further, \(M\otimes_B A\) is projective of finite type over \(A\), and therefore has only finitely many non-zero homogeneous components, each projective of finite type over \(A\). After localizing on \(\Spec(A)\), one may thus suppose that \(M\otimes_B A\) is free of finite type over \(A\). Finally, projectivity of \(M\) gives
\[
 \Tor_1^B(M,A)=0.
\]
These properties imply that \(M\) is free of finite type over \(B\) as a graded module (Bourbaki, \emph{Algèbre}, ch.~II, \S~11, prop.~7), and hence \(\widetilde M\) is free of finite type over \(\widetilde B\), as required.

\sgaStableTarget{sga6.e06.lem.1-5-2}\textbf{Lemma 1.5.2.} Let \(A\) be a ring, let \(B\) be an \(A\)-algebra of finite presentation, and let \(M\) be a \(B\)-module of finite presentation. Suppose that \(B\) and \(M\) are flat over \(A\), and that, for every residue field \(k\) of \(A\), the \(B\otimes_A k\)-module \(M\otimes_A k\) has projective dimension at most \(n\) (as is the case if \(B=A[t_0,\ldots,t_r]\)). Then \(M\) has a resolution of length \(n\) by projective \(B\)-modules of finite type:
\[
  0\longrightarrow L_n\longrightarrow\cdots\longrightarrow L_0\longrightarrow M\longrightarrow 0.
\]
This is the local criterion of EGA IV, 12.3.2.

\sgaStableTarget{sga6.e06.prop.1-6}\textbf{Proposition 1.6.} Let \(S\) be quasi-compact, let \(E\) be locally free of rank \(r+1\), let \(B=\Sym(E)\) with its usual grading, and let \(M\) be a graded \(B\)-module of finite presentation, flat over \(S\). Then \(M\) admits a graded resolution
\[
  0\longrightarrow L_{r+1}\longrightarrow\cdots\longrightarrow L_0\longrightarrow M\longrightarrow 0,
\]
where the \(L_i\) are graded \(B\)-modules locally free of finite type.

First, \(M\) admits a graded epimorphism \(L\to M\) with \(L\) locally free of finite type: choose an integer \(N\) such that \(M\) is generated in degrees \(\leq N\), and take
\[
  L=\bigoplus_{i\leq N} M_i\otimes_{\mathcal O_S} B(-i),
\]
with the morphism \hyperlink{sga6.e06.eq.1-2-1}{(1.2.1)}. Repeating the construction for successive kernels gives a global graded resolution. After \(r+1\) steps, \hyperlink{sga6.e06.lem.1-5-2}{Lemma 1.5.2} shows locally on affine opens that the final kernel is projective over \(B\), and \hyperlink{sga6.e06.lem.1-5-1}{Lemma 1.5.1} then shows that it is locally free as a graded \(B\)-module; truncation gives the asserted resolution.

\sgaStableTarget{sga6.e06.prop.1-7}\textbf{Proposition 1.7.} Let \(S\) be quasi-compact, let \(E\) be locally free of finite type, let \(B=\Sym(E)\), and let \(L\) be a graded \(B\)-module locally free of finite type. Then \(L\) admits a finite decreasing filtration by graded submodules \(L^{(m)}\) such that
\[
\begin{array}{ll}
\text{\rm (i)}& L^{(m)}=L\quad\text{for }m\leq m_0,\\
\text{\rm (ii)}& L^{(m)}=0\quad\text{for }m\geq m_1,\\
\text{\rm (iii)}& \operatorname{gr}^m L
  \simeq N_m\otimes_{\mathcal O_S}B(m),
\end{array}
\]
where \(N_m\) is locally free of finite type on \(S\).

The question is local on \(S\). Since \(S\) is quasi-compact, the degrees of homogeneous bases that occur locally lie between two fixed integers. One argues by induction on the length of this interval. The first step is the canonical map from the lowest non-zero homogeneous component,
\[
  L_{n_0}\otimes_{\mathcal O_S}B(-n_0)\longrightarrow L,
\]
which locally identifies the corresponding direct summand. The quotient has local homogeneous bases in higher degrees, and induction applies.

\sgaStableTarget{sga6.e06.sec.1.8}\subsection*{1.8. Generation by the powers of \(\xi\)}
Let \(F\) be locally free of finite type on \(X\). By \hyperlink{sga6.e06.cor.1-4}{Corollary 1.4} and \hyperlink{sga6.e06.prop.1-6}{Proposition 1.6}, \(F\) comes from a graded module \(M\) admitting a finite graded resolution by \(L_i\)'s. Applying \(\Proj_0\) gives an exact sequence of locally free sheaves on \(X\), and therefore in \(K^\bullet(X)\)
\[
  \cl(F)=\sum_i(-1)^i\cl(\Proj_0(L_i)).
\]
Using the filtration of \hyperlink{sga6.e06.prop.1-7}{Proposition 1.7}, each term is a sum of classes
\[
  \cl\bigl(f^*N_{ij}(j)\bigr)=f^*\cl(N_{ij})\,\xi^j .
\]
Thus \(K^\bullet(X)\) is generated as a \(K^\bullet(S)\)-module by the powers of \(\xi\).

\sgaStableTarget{sga6.e06.sec.1.9}\subsection*{1.9. The relation}
For linear dependence, use the canonical epimorphism \(f^*E\to\mathcal O_X(1)\). Its kernel is locally free of rank \(r\). Hence
\[
  \lambda^{r+1}(f^*E-\xi)=0.
\]
Since \(\lambda_t(-\xi)=1/(1+\xi t)\), one obtains the relation
\sgaStableTarget{sga6.e06.eq.1-9-1}\begin{equation}\tag{1.9.1}
  \sum_{i=0}^{r+1}(-1)^i\,\lambda^i(E)\,\xi^{r+1-i}=0,
\end{equation}
equivalently the characteristic equation of \(\xi\) associated with the locally free module \(E\).

\sgaStableTarget{sga6.e06.sec.1.10}\subsection*{1.10--1.11. Linear independence}
The proof of the linear-independence assertion is due to J.-P. Jouanolou.

\sgaStableTarget{sga6.e06.lem.1-10}\textbf{Lemma 1.10.} With the hypotheses of \hyperlink{sga6.e06.thm.1-1}{Theorem 1.1}, let \(F\) and \(G\) be locally free sheaves of finite type on \(X\), and suppose that they have the same class in \(K^\bullet(X)\). Then there is an integer \(n_0\) such that, for all \(n\geq n_0\),
\[
\begin{array}{ll}
\text{\rm (i)}& f_*(F(n))\text{ and }f_*(G(n))\text{ are locally free of finite type on }S,\\
\text{\rm (ii)}& \cl f_*(F(n))=\cl f_*(G(n))\text{ in }K^\bullet(S).
\end{array}
\]

An equality of classes in \(K^\bullet(X)\) is a finite integral combination of relations coming from short exact sequences of locally free sheaves. Twisting by \(\mathcal O_X(n)\) and applying \(f_*\), for \(n\) large enough all higher direct images vanish and all direct images are locally free of finite type; the same relations then hold in \(K^\bullet(S)\).

Suppose that
\sgaStableTarget{sga6.e06.eq.1-11-1}\begin{equation}\tag{1.11.1}
  f^*(x_0)+f^*(x_1)\xi+\cdots+f^*(x_r)\xi^r=0 .
\end{equation}
Writing each \(x_i\) as a difference of classes of locally free sheaves and adding suitable summands, \hyperlink{sga6.e06.eq.1-11-1}{(1.11.1)} is converted into an equality \(\cl(F)=\cl(G)\) for locally free sheaves of finite type on \(X\). By \hyperlink{sga6.e06.lem.1-10}{Lemma 1.10}, for all sufficiently large \(n\) one has
\sgaStableTarget{sga6.e06.eq.1-11-2}\begin{equation}\tag{1.11.2}
  \sum_{i=0}^{r} x_i\,\sigma_{n+i}(E)=0
  \quad\text{in }K^\bullet(S),
\end{equation}
where \(\sigma_p(E)\) denotes the class of \(\Sym^p(E)\), and \(\sigma_p(E)=0\) for \(p<0\).

The Koszul complex of the symmetric algebra gives, for \(p\geq 1\),
\sgaStableTarget{sga6.e06.eq.1-11-3}\begin{equation}\tag{1.11.3}
  \sum_{i=0}^{r+1}(-1)^i\lambda^i(E)\,\sigma_{p-i}(E)=0 .
\end{equation}
Using this relation, one descends by induction from large \(n\) in \hyperlink{sga6.e06.eq.1-11-2}{(1.11.2)}. It follows that \hyperlink{sga6.e06.eq.1-11-2}{(1.11.2)} holds for all \(n\geq -r\). Taking \(n=-r\) gives \(x_r=0\); repeating the argument gives successively all \(x_i=0\). This proves the linear independence and hence \hyperlink{sga6.e06.thm.1-1}{Theorem 1.1}.

\sgaStableTarget{sga6.e06.rem.1-12}\textbf{Remark 1.12.} Let
\[
  \sigma_t(E)=\sum_{i\geq 0}\sigma_i(E)t^i .
\]
The relation \hyperlink{sga6.e06.eq.1-11-3}{(1.11.3)} is equivalent to
\[
  \lambda_t(E)\,\sigma_{-t}(E)=1.
\]
It follows in particular that \(\sigma_t\) is additive and extends to a homomorphism
\[
  K^\bullet(S)\longrightarrow 1+K^\bullet(S)[[t]]_+ .
\]

\sgaStableTarget{sga6.e06.rem.1-13}\textbf{Remark 1.13.} Any sequence of \(r+1\) consecutive powers of \(\xi\) is a basis of \(K^\bullet(X)\) over \(K^\bullet(S)\). If
\[
  x=1-\xi^{-1}
\]
is the class of a hyperplane section, then \(1,x,\ldots,x^r\) is also a basis. The dependence relation becomes
\sgaStableTarget{sga6.e06.eq.1-13-1}\begin{equation}\tag{1.13.1}
  \sum_{k=0}^{r+1}\lambda^{r+1-k}(\check E-r-1)\,x^k=0.
\end{equation}
For \(E=\mathcal O_S^{r+1}\), so that \(X=\mathbf P^r_S\), this reduces to
\[
  x^{r+1}=0.
\]
Using relative schemes in Hakim's sense, \hyperlink{sga6.e06.thm.1-1}{Theorem 1.1} also makes sense when \(S\) is a locally ringed topos; if the final object of the topos is quasi-compact, the proof transports without essential change.

\sgaStableTarget{sga6.e06.sec.2}\section*{2. Computation of \(K^\bullet\) of a projective bundle: the case of perfect complexes}

The present section will not be used in the rest of the Seminar. In this section the notation \(K^\bullet(X)\) denotes the Grothendieck group of perfect complexes on \(X\) (\hyperlink{sga6.e04.def.2-2}{IV 2.2}). We propose to prove the following theorem, analogous to \hyperlink{sga6.e06.thm.1-1}{1.1}:

\sgaStableTarget{sga6.e06.thm.2-1}\textbf{Theorem 2.1.} Let \(S\) be quasi-compact, let \(E\) be a locally free \(\mathcal O_S\)-module of rank \(r+1\), let \(X=\mathbf P(E)\) be the associated projective bundle; let \(\xi\) be the class in \(K^\bullet(X)\) of the complex equal to \(\mathcal O_X(1)\) in degree zero, and zero in the other degrees. Then \(K^\bullet(X)\), endowed with its natural \(K^\bullet(S)\)-algebra structure, is a free \(K^\bullet(S)\)-module with basis
\[
  1,\xi,\ldots,\xi^r .
\]

\sgaStableTarget{sga6.e06.sec.2.2}\subsection*{2.2. Notation}
Fix the notation: \(B\) will denote the symmetric algebra of \(E\); write \(\mathcal A\) for the category of \(\mathcal O_X\)-modules, \(\mathcal G\) for the category of graded \(B\)-modules with graded morphisms as morphisms, \(\mathcal G_0\) for the full subcategory of \(\mathcal G\) whose objects are the \(B\)-modules locally free of finite type as graded \(B\)-modules, and \(\mathcal G'\) for the category of \(\mathcal O_S\)-modules. For every \(n\), and every complex \(L^\bullet\) of \(D^+(\mathcal A)\), we shall write \(L^\bullet(n)\) for the complex
\[
  L^\bullet(n)=L^\bullet\otimes_{\mathcal O_X}\mathcal O_X(n).
\]

\sgaStableTarget{sga6.e06.sec.2.3}\subsection*{2.3. The functor \(\Proj_0\)}
We shall need to apply the functor \(\Proj_0\) to arbitrary graded \(B\)-modules, and not only to quasi-coherent graded \(B\)-modules as in the construction of EGA II 3.2. We shall therefore indicate briefly how this construction can be made in the general case.

\sgaStableTarget{sga6.e06.para.2-3-1}\textbf{2.3.1.} First let \(S\) be a scheme, let \(A\) be a quasi-coherent \(\mathcal O_S\)-algebra, let \(Y=\Spec(A)\), and let \(f:Y\to S\) be the structural morphism. From the isomorphism
\[
  A\simeq f_*\mathcal O_Y
\]
one obtains a canonical homomorphism
\sgaStableTarget{sga6.e06.eq.2-3-1-1}\begin{equation}\tag{2.3.1.1}
  f^{-1}(A)\longrightarrow \mathcal O_Y,
\end{equation}
where \(f^{-1}\) denotes inverse image in the sense of sheaves of sets; in other words, one has a morphism of ringed spaces
\sgaStableTarget{sga6.e06.eq.2-3-1-2}\begin{equation}\tag{2.3.1.2}
  \bar f:(Y,\mathcal O_Y)\longrightarrow (S,A).
\end{equation}
It gives an inverse-image functor \(\bar f^{\,*}\) from the category of \(A\)-modules to the category of \(\mathcal O_Y\)-modules, left adjoint to the direct-image functor \(\bar f_*\) from the category of \(\mathcal O_Y\)-modules to the category of \(f_*(\mathcal O_Y)\)-modules. Thus
\sgaStableTarget{sga6.e06.eq.2-3-1-3}\begin{equation}\tag{2.3.1.3}
  \bar f^{\,*}(M)=f^{-1}(M)\otimes_{f^{-1}A}\mathcal O_Y
\end{equation}
for every \(A\)-module \(M\).

Since the homomorphism \hyperlink{sga6.e06.eq.2-3-1-1}{(2.3.1.1)} is flat (in each fiber it corresponds to a localization), the functor \(\bar f^{\,*}\) just defined is exact (cf. [4], p. 164). If \(M\) is a quasi-coherent \(A\)-module, this construction gives the module \(\widetilde M\) defined in EGA II 1.4.3. Indeed, the isomorphism
\[
  M\simeq f_*(\widetilde M)=\bar f_*(\widetilde M)
\]
gives, by adjunction, a homomorphism
\[
  \bar f^{\,*}(M)\longrightarrow \widetilde M
\]
and one checks immediately on each fiber that it is an isomorphism.

\sgaStableTarget{sga6.e06.para.2-3-2}\textbf{2.3.2.} Now let \(B\) be a quasi-coherent graded \(\mathcal O_S\)-algebra, let \(M\) be a graded \(B\)-module, and let \(X=\Proj(B)\). For every affine open \(U\) of \(S\), and every section \(s\) of \(B\) over \(U\), homogeneous of degree \(d>0\), write \(M_{(s),U}\) for the homogeneous component of degree \(0\) of the module of fractions \(M_s\); thus
\[
  M_{(s),U}=
  \bigl((M|_U)\otimes_{B|_U}(B|_U)_s\bigr)_0 .
\]
If \(V\) is an affine open of \(U\), then \(M_{(s),V}\) is the restriction to \(V\) of \(M_{(s),U}\).

Consider the base of opens of \(X\) formed by the affine opens
\[
  D_+(U,s)\simeq \Spec\bigl((B|_U)_{(s)}\bigr)
\]
Write \(f\) for the structural morphism \(X\to S\), and \(f_{U,s}\) for the affine morphism
\[
  f_{U,s}:D_+(U,s)\longrightarrow U .
\]
Put
\[
  \Proj_0(M)_{U,s}
  =
  \bar f_{U,s}^{\,*}\bigl(M_{(s),U}\bigr).
\]

With the notation of \hyperlink{sga6.e06.para.2-3-1}{2.3.1}, we shall show that the \(\Proj_0(M)_{U,s}\) glue as \(U\) and \(s\) vary, thereby defining an \(\mathcal O_X\)-module denoted \(\Proj_0(M)\). The gluing for variable \(U\) and fixed \(s\) being evident, one is reduced (cf. EGA II 2.5.2) to showing that, if \(t\) is a section of \(B\) on \(U\), homogeneous of degree \(e>0\), then the restriction of \(\Proj_0(M)_{U,s}\) to \(D_+(U,st)\) is canonically isomorphic to \(\Proj_0(M)_{U,st}\). Consider
\[
\begin{tikzcd}[column sep=large,row sep=large]
D_+(st) \arrow[r,"i"] \arrow[dr,"f_{U,st}"'] &
D_+(s) \arrow[d,"f_{U,s}"]\\
& U
\end{tikzcd}
\]
It is a matter of showing that the modules
\[
  i^*\bar f_{U,s}^{\,*}\bigl(M_{(s),U}\bigr)
  \quad\text{and}\quad
  \bar f_{U,st}^{\,*}\bigl(M_{(st),U}\bigr)
\]
are canonically isomorphic. On the one hand, one has the canonical isomorphisms
\[
\begin{aligned}
i^*\bar f_{U,s}^{\,*}\bigl(M_{(s),U}\bigr)
&\simeq
i^*\Bigl(f_{U,s}^{-1}\bigl(M_{(s),U}\bigr)
 \otimes_{f_{U,s}^{-1}((B|_U)_{(s)})}\mathcal O_{D_+(s)}\Bigr)\\
&\simeq
i^{-1}\Bigl(f_{U,s}^{-1}\bigl(M_{(s),U}\bigr)
 \otimes_{f_{U,s}^{-1}((B|_U)_{(s)})}\mathcal O_{D_+(s)}\Bigr)
 \otimes_{i^{-1}(\mathcal O_{D_+(s)})}\mathcal O_{D_+(st)}\\
&\simeq
f_{U,st}^{-1}\bigl(M_{(s),U}\bigr)
 \otimes_{f_{U,st}^{-1}((B|_U)_{(s)})}\mathcal O_{D_+(st)} .
\end{aligned}
\]
On the other hand,
\[
  \bar f_{U,st}^{\,*}\bigl(M_{(st),U}\bigr)
  \simeq
  f_{U,st}^{-1}\bigl(M_{(st),U}\bigr)
  \otimes_{f_{U,st}^{-1}((B|_U)_{(st)})}\mathcal O_{D_+(st)} .
\]
Finally, it follows from EGA II 2.2.2, translated into sheaf language, that one has the canonical isomorphisms
\[
  M_{(st),U}\simeq
  \bigl(M_{(s),U}\bigr)_{t^d/s^e}
  \simeq
  M_{(s),U}\otimes_{(B|_U)_{(s)}}(B|_U)_{(st)},
\]
and hence the desired isomorphism.

It is clear that, if \(M\) is quasi-coherent, the module \(\Proj_0(M)\) thus constructed is indeed the one defined in EGA II.

\sgaStableTarget{sga6.e06.para.2-3-3}\textbf{2.3.3.} It follows from \hyperlink{sga6.e06.para.2-3-1}{2.3.1} and \hyperlink{sga6.e06.para.2-3-2}{2.3.2} that the functor \(\Proj_0\) is exact. Returning to the notation of \hyperlink{sga6.e06.sec.2.2}{2.2}, by passage to derived categories one therefore obtains a functor, still denoted \(\Proj_0\):
\[
  \Proj_0:D^+(\mathcal G)\longrightarrow D^+(\mathcal A).
\]
It is also clear from the preceding constructions that the functor \(\Proj_0\), on the category of graded \(B\)-modules, is compatible with tensor product; the same remains true of the functor obtained on the derived categories.

\sgaStableTarget{sga6.e06.sec.2.4}\subsection*{2.4. The functor \(F_k\)}
For an \(\mathcal O_X\)-module \(L\) and an integer \(k\), set
\sgaStableTarget{sga6.e06.eq.2-4-1}\begin{equation}\tag{2.4.1}
  F_k(L)=\bigoplus_{n\geq k} f_*(L(n)).
\end{equation}
The functor \(F_k\) is left exact and therefore has a right derived functor
\sgaStableTarget{sga6.e06.eq.2-4-2}\begin{equation}\tag{2.4.2}
  RF_k:D^+(\mathcal A)\longrightarrow D^+(\mathcal G).
\end{equation}
If \(L\) is locally free of finite type, then for \(k\) sufficiently large it is \(F_k\)-acyclic, by the vanishing theorem used above in 1.3.

For every \(k\) there is a functorial morphism
\sgaStableTarget{sga6.e06.eq.2-5-1}\begin{equation}\tag{2.5.1}
  \beta:\Proj_0(F_k(F))\longrightarrow F,
\end{equation}
constructed on each \(D_+(U,s)\) by the usual localization formula: a fraction represented by a homogeneous section of sufficiently high degree is sent to the corresponding local section of \(F\). In the quasi-coherent case this is the canonical morphism of EGA. Passing to derived categories gives
\sgaStableTarget{sga6.e06.eq.2-5-2}\begin{equation}\tag{2.5.2}
  \Proj_0\circ RF_k \longrightarrow \operatorname{id}_{D^+(\mathcal A)}.
\end{equation}

\sgaStableTarget{sga6.e06.prop.2-6}\textbf{Proposition 2.6.} Let \(S\) be quasi-compact, let \(E\) be locally free of finite type, let \(X=\mathbf P(E)\), and let \(L^\bullet\) be a perfect complex on \(X\). Then, for \(k\) sufficiently large,
\[
\begin{array}{ll}
\text{\rm (i)}& RF_k(L^\bullet)\text{ is }\mathcal G_0\text{-perfect},\\[1mm]
\text{\rm (ii)}& \Proj_0 RF_k(L^\bullet)\longrightarrow L^\bullet
\text{ is an isomorphism in }D^+(\mathcal A).
\end{array}
\]

The assertion is local on \(S\). One may therefore assume \(S\) affine. Since \(\mathcal O_X(1)\) is ample, the perfect complex \(L^\bullet\) is locally, and then globally after refining, represented by a bounded complex whose terms are locally free of finite type. For \(k\) sufficiently large all those finitely many terms are \(F_k\)-acyclic, and the graded modules \(F_k(L^i)\) are \(\mathcal G_0\)-perfect. This gives (i), and the morphism \hyperlink{sga6.e06.eq.2-5-2}{(2.5.2)} is then an isomorphism term by term, giving (ii).

\sgaStableTarget{sga6.e06.sec.2.7}\subsection*{2.7. Decomposition of a \(\mathcal G_0\)-perfect complex}
On \(\mathcal G\), the functor which associates to a graded \(B\)-module its \(n\)-th homogeneous component is exact with values in \(\mathcal G'\); it defines a functor
\sgaStableTarget{sga6.e06.eq.2-7-1}\begin{equation}\tag{2.7.1}
  \Comp_n:D^+(\mathcal G)\longrightarrow D^+(\mathcal G').
\end{equation}
The homomorphism \hyperlink{sga6.e06.eq.1-2-1}{(1.2.1)} gives a morphism of functors
\sgaStableTarget{sga6.e06.eq.2-7-2}\begin{equation}\tag{2.7.2}
  \Comp_n\otimes_{\mathcal O_S}B(-n)\longrightarrow
  \operatorname{id}_{D^+(\mathcal G)}.
\end{equation}

Let \(M^\bullet\) be a \(\mathcal G_0\)-perfect complex. We shall show that there exists a finite sequence of distinguished triangles
\[
\begin{array}{ccccccccccc}
& P_1^\bullet &&&& P_2^\bullet &&&& P_k^\bullet &\\
\swarrow && \nwarrow && \swarrow && \nwarrow && \swarrow && \nwarrow\\
N_1^\bullet & \longrightarrow & M^\bullet, & \qquad &
N_2^\bullet & \longrightarrow & P_1^\bullet, & \cdots, &
N_k^\bullet & \longrightarrow & P_{k-1}^\bullet,
\end{array}
\]
where each \(P_k^\bullet\) at the last stage is acyclic, the terms \(N_i^\bullet\) are of the form
\[
  N_i^\bullet=K_i^\bullet\otimes_{\mathcal O_S}B(d_i),
\]
and the \(K_i^\bullet\) are perfect complexes on \(S\) in the ordinary sense.

Indeed, because \(M^\bullet\) is \(\mathcal G_0\)-perfect and \(S\) is quasi-compact, there exists a finite covering of \(S\) and an integer \(d_0\) such that, on every open of the covering, \(M^\bullet\) is represented by a bounded complex of locally free \(B\)-modules of finite type, with no homogeneous generators in degrees \(<d_0\). Taking \(d_0\) minimal, the morphism \hyperlink{sga6.e06.eq.2-7-2}{(2.7.2)} extracts the part of the lowest degree. Its cone has all homogeneous generators in strictly higher degrees. Repeating the process, and using the finiteness of the degree interval from \hyperlink{sga6.e06.prop.1-7}{Proposition 1.7}, one obtains the desired finite sequence of triangles.

\sgaStableTarget{sga6.e06.sec.2.8}\subsection*{2.8. Proof of Theorem 2.1}
Let \(L^\bullet\) be a perfect complex on \(X\). For \(k\) sufficiently large put
\[
  M^\bullet=RF_k(L^\bullet).
\]
By \hyperlink{sga6.e06.prop.2-6}{Proposition 2.6}, \(M^\bullet\) is \(\mathcal G_0\)-perfect and \(\Proj_0(M^\bullet)\simeq L^\bullet\). Applying the decomposition of \hyperlink{sga6.e06.sec.2.7}{2.7} and then \(\Proj_0\), one obtains in \(K^\bullet(X)\)
\[
  \cl(L^\bullet)
  =
  \sum_i \cl\bigl(f^*K_i^\bullet(d_i)\bigr)
  =
  \sum_i f^*\cl(K_i^\bullet)\,\xi^{d_i}.
\]
Thus the powers of \(\xi\) generate \(K^\bullet(X)\) as a \(K^\bullet(S)\)-module.

The dependence relation is the same relation \hyperlink{sga6.e06.eq.1-9-1}{(1.9.1)}, now interpreted in the Grothendieck group of perfect complexes. For linear independence one uses the \(K^\bullet(S)\)-linear homomorphism
\[
  f_*:K^\bullet(X)\longrightarrow K^\bullet(S).
\]
For the structure sheaf and negative twists of the projective bundle one has
\[
  f_*(1)=1,\qquad f_*(\xi^n)=0\quad(-r\leq n<0).
\]
Multiplying a relation among \(1,\xi,\ldots,\xi^r\) by suitable negative powers of \(\xi\) and applying \(f_*\) successively gives all its coefficients equal to zero. This completes the proof of \hyperlink{sga6.e06.thm.2-1}{Theorem 2.1}.



\providecommand{\Sym}{\operatorname{Sym}}
\providecommand{\GL}{\operatorname{GL}}
\providecommand{\Grass}{\operatorname{Grass}}
\providecommand{\gr}{\operatorname{gr}}
\providecommand{\ctop}{c_{\mathrm{top}}}
\providecommand{\Rep}{\operatorname{Rep}}
\providecommand{\RZ}{R_{\mathbb Z}}
\providecommand{\Dflag}{D}
\providecommand{\Gmn}{G}


This completes the induction, for if \(0\le s\le r\) and if \(x_j=0\) for \(s<j\le r\), then from a relation
\[
  \sum_{i=0}^{s} f^*(x_i)\xi^i=0
\]
after multiplication by \(\xi^{r-s}\) one gets
\[
  \sum_{i=0}^{s} f^*(x_i)\xi^{i+r-s}=0,
\]
and hence \(x_s=0\) by what precedes.

\sgaStableTarget{sga6.e06.sec.3}\section*{3. Consequences of the structure theorem for the \(K^\bullet\) of a projective bundle}

\sgaStableTarget{sga6.e06.prop.3-1}\textbf{Proposition 3.1.} Let \(S\) be a quasi-compact scheme, let \(E_1,\ldots,E_n\) be \(\mathcal O_S\)-modules locally free of finite type, and let \(D(E_i)\) be the flag scheme of \(E_i\). Put
\[
  X=D(E_1)\times_S\cdots\times_S D(E_n).
\]
Then the canonical homomorphism
\[
  K^\bullet(S)\longrightarrow K^\bullet(X)
\]
is injective, where \(K^\bullet\) is defined either from \(\mathcal O_S\)-modules locally free of finite type or from perfect complexes.

One verifies this immediately by using the universal property of flag schemes. The scheme \(X\) can be constructed as follows: one considers a sequence of \(S\)-schemes \(X_1,\ldots,X_n\), defined by induction by taking \(X_k\) to be the flag scheme of the inverse image of \(E_k\) on \(X_{k-1}\); then \(X=X_n\). Thus, to prove the proposition, it is enough to prove it when \(n=1\), the general case following by induction. We may therefore suppose that \(X\) is the flag scheme of an \(\mathcal O_S\)-module \(E\), locally free of finite type.

After decomposing \(S\) into finitely many open-and-closed subschemes, we may suppose that \(E\) has constant rank \(r\). The scheme \(X\) may then be constructed as follows. Put \(E=F_0\), \(S=X_0\), and define successively schemes \(X_1,\ldots,X_r\) and, for each \(i\), an \(\mathcal O_{X_i}\)-module locally free of finite type \(F_i\), by
\[
  X_i=\mathbb P(F_{i-1}),\qquad
  F_i=\Ker\bigl(f_i^*(F_{i-1})\longrightarrow \mathcal O_{X_i}(1)\bigr),
\]
where \(f_i:X_i\to X_{i-1}\) is the structural morphism. It follows that it is enough to prove the proposition for a projective bundle. This is precisely the content of \hyperlink{sga6.e06.thm.1-1}{1.1}, respectively of \hyperlink{sga6.e06.thm.2-1}{2.1}, according to which definition of \(K^\bullet\) is being used.

\sgaStableTarget{sga6.e06.thm.3-2}\textbf{Theorem 3.2.} Let \(S\) be a quasi-compact scheme. The pre-\(\lambda\)-ring \(K^\bullet(S)\), defined from \(\mathcal O_S\)-modules locally free of finite type, is a \(\lambda\)-ring.

The homomorphism
\[
  \lambda_t:K^\bullet(S)\longrightarrow 1+K^\bullet(S)[[t]]^+
\]
being additive, it is enough, in order to verify that it is a \(\lambda\)-homomorphism, to check the universal relations on elements of the form \(\cl(E)\), where \(E\) is an \(\mathcal O_S\)-module locally free of finite type. Thus one has to show that for every pair \(E,F\) of such modules, with classes again denoted \(E,F\), the relations \hyperlink{sga6.e05.eq.2-4-2}{V 2.4.2} are satisfied:
\[
  \lambda^i(EF)=P_i\bigl(\lambda^1(E),\ldots,\lambda^i(E);\lambda^1(F),\ldots,\lambda^i(F)\bigr),
\]
\[
  \lambda^j(\lambda^i(E))=Q_{i,j}\bigl(\lambda^1(E),\ldots,\lambda^{ij}(E)\bigr).
\]
Since \(S\) is quasi-compact, it admits a finite covering by disjoint opens such that on each of them \(E\) and \(F\) have constant rank. Since the \(K^\bullet\) of a finite disjoint union is the product of the corresponding \(K^\bullet\)'s, it is enough to check the relations on the members of this covering; hence we may suppose that \(E\) and \(F\) have constant ranks on \(S\).

Put then
\[
  X=D(E)\times_S D(F).
\]
By \hyperlink{sga6.e06.prop.3-1}{Proposition 3.1} the \(\lambda\)-homomorphism \(K^\bullet(S)\to K^\bullet(X)\) is injective. Consequently, in order to prove the relations \hyperlink{sga6.e05.eq.2-4-2}{V 2.4.2}, we may work in \(K^\bullet(X)\). But in \(K^\bullet(X)\) the elements \(E\) and \(F\) split as sums of classes of invertible sheaves. By additivity it is therefore enough to prove the same relations for elements \(\xi,\eta\) which are classes of invertible sheaves. Written in the form \hyperlink{sga6.e05.eq.2-4-1}{V 2.4.1}, the relations are
\[
  1+\xi\eta t=(1+\xi t)\circ(1+\eta t),
  \qquad
  \lambda_u(1+\xi t)=1+\{1+\xi t\}u.
\]
These are just the relations \hyperlink{sga6.e05.eq.2-3-1}{V 2.3.1} and \hyperlink{sga6.e05.eq.2-3-3}{V 2.3.3} defining the \(\lambda\)-ring structure on \(1+K^\bullet(S)[[t]]^+\).

Let us also record the following more general result.

\sgaStableTarget{sga6.e06.thm.3-3}\textbf{Theorem 3.3.} Let \((T,A)\) be a commutatively ringed topos, and let \(K\) be the Grothendieck group of \(A\)-modules locally free of bounded rank on \(T\). Then \(K\), endowed with its canonical pre-\(\lambda\)-structure, is a \(\lambda\)-ring.

To prove that \(K\) is a \(\lambda\)-ring, it is necessary to prove the relations \hyperlink{sga6.e05.eq.2-4-1}{V 2.4.1}; by additivity, it is enough to prove them for classes of \(A\)-modules \(E,F\) locally free of bounded rank. Since the ranks of \(E\) and \(F\) are bounded, one may decompose \(T\) as a finite disjoint sum of opens over which \(E\) and \(F\) have constant rank. Using the compatibility of \(K\) with finite disjoint sums, and writing the relations \hyperlink{sga6.e05.eq.2-4-1}{V 2.4.1} in the form \hyperlink{sga6.e05.eq.2-4-2}{V 2.4.2}, we are reduced to the case of modules of constant rank.

Let \(E\) be an \(A\)-module locally free of finite rank \(n\). There is a torsor \(P\) under \(\GL(n,A)\) such that
\[
  E=P\times_{\GL(n,A)} A^n,
\]
where \(A^n\) is considered as a \(\GL(n,A)\)-module by the standard representation. One can then define a \(\lambda\)-homomorphism from the representation ring \(R_{\mathbb Z}(\GL(n))\) to \(K\) by associating to a representation \(\rho\) of \(\GL(n)\) in a free \(\mathbb Z\)-module of rank \(p\) the class of
\[
  P\times_{\GL(n,A)} A^p.
\]
Under this homomorphism, the class of \(E\) is the image of the standard representation. Thus any polynomial relation among the \(\lambda^i(E)\)'s follows from the corresponding relation for the standard representation in \(R_{\mathbb Z}(\GL(n))\).

Similarly, for two locally free modules \(E,F\) of ranks \(n,p\), one uses a torsor under \(\GL(n,A)\times\GL(p,A)\), with \(\GL(n,A)\) acting through the standard representation and \(\GL(p,A)\) through the trivial one on the first factor, and conversely for \(F\). We are reduced to proving that, for every group scheme \(G\) over \(\Spec\mathbb Z\) which is a finite product of linear group schemes, the pre-\(\lambda\)-ring \(R_{\mathbb Z}(G)\) is a \(\lambda\)-ring. This may be obtained by the explicit computation of \(R_{\mathbb Z}(G)\) in [5]; alternatively, one can reason as in \hyperlink{sga6.e06.thm.3-2}{3.2}, using 1.14.

\sgaStableTarget{sga6.e06.sec.4}\section*{4. Computation of the \(K^\bullet\) of a flag bundle}

\sgaStableTarget{sga6.e06.sec.4.1}\subsection*{4.1}
Let \(S\) be a scheme and let \(E\) be an \(\mathcal O_S\)-module locally free of rank \(n\). Let
\[
  \pi=(p_1,\ldots,p_k),\qquad p_i\geq 1,
  \qquad \sum_{i=1}^k p_i=n.
\]
The flag functor of type \(\pi\) of \(E\) is the functor on the category of \(S\)-schemes, with values in sets, which to an \(S\)-scheme \(f:X\to S\) associates the set of chains of quotients of \(f^*E\)
\[
  (L_{k-1},\ldots,L_1),
\]
where, for each \(i\), \(L_i\) is a quotient of \(L_{i+1}\), locally free of rank \(p_1+\cdots+p_i\). This functor is representable by an \(S\)-scheme \(D_\pi(E)\), endowed with a canonical sequence of quotients of the inverse image of \(E\), locally free of ranks \(p_1,p_1+p_2,\ldots\). If \(\pi=(1,n-1)\), then \(D_\pi(E)=\mathbb P(E)\); if \(\pi=(r,n-r)\), then \(D_\pi(E)=G_{r,n}(E)\), the Grassmannian of type \((r,n-r)\) of \(E\). If all \(p_i\)'s are equal to \(1\), one obtains the usual flag bundle of \(E\), denoted \(D(E)\).

\sgaStableTarget{sga6.e06.lem.4-2}\textbf{Lemma 4.2.} Let \(S\) be a scheme, let \(E\) be an \(\mathcal O_S\)-module locally free of rank \(n\), and let
\[
  \pi=(p_1,\ldots,p_i,\ldots,p_k),\qquad p_1+\cdots+p_k=n.
\]
Let \(p',p''\geq 1\) be such that \(p'+p''=p_i\), and put
\[
  \pi'=(p_1,\ldots,p_{i-1},p',p'',p_{i+1},\ldots,p_k),
  \qquad \pi''=(p',p'').
\]
Let \(L_1,\ldots,L_{k-1}\) be the canonical quotients of the inverse image of \(E\) on \(D_\pi(E)\), and let \(F\) be the kernel of \(L_i\to L_{i-1}\). Then there is a canonical \(S\)-isomorphism
\[
  D_{\pi'}(E)\simeq D_{\pi''}(F).
\]
The proof, immediate from the functorial definition of flag schemes, is left to the reader.

\sgaStableTarget{sga6.e06.lem.4-3}\textbf{Lemma 4.3.} Let \(A\) be a commutative ring with unit, let \(n>0\), and let \((c_i)_{1\leq i\leq n}\) be elements of \(A\). Let \(C\) be the \(A\)-algebra generated by \(\xi_1,\ldots,\xi_n\), subject to the relations
\[
  \sigma_i(\xi_1,\ldots,\xi_n)=c_i,
  \qquad i=1,\ldots,n,
\]
where \(\sigma_i\) denotes the elementary symmetric functions. Let \(p_1,\ldots,p_k\) be integers \(\geq 1\) whose sum is \(n\). For \(1\leq j\leq k\) put
\[
  \xi^{(j)}=(\xi^{(j)}_1,\ldots,\xi^{(j)}_{p_j}),
  \qquad
  \xi_i^{(j)}=\xi_{\sum_{\alpha<j}p_\alpha+i}.
\]
Put
\[
  c_i^{(j)}=\sigma_i(\xi^{(j)}_1,\ldots,\xi^{(j)}_{p_j}),
  \qquad 1\leq j\leq k,
  \quad 1\leq i\leq p_j,
\]
and let \(B\) be the subalgebra of \(C\) generated by the \(c_i^{(j)}\)'s. Then:

\begin{enumerate}[label=\textup{\roman*)}]
\item The elements
\[
  (\xi^{(1)})^{\alpha(1)}\cdots(\xi^{(k)})^{\alpha(k)},
  \qquad 0\leq \alpha_i^{(j)}\leq p_j-i,
\]
form a basis of \(C\), considered as a \(B\)-module.
\item The generators \(c_i^{(j)}\) of the \(A\)-algebra \(B\) satisfy the \(n\) isobaric relations deduced from
\[
  1+\sum_{i=1}^n c_i
  =
  \prod_{j=1}^{k}\left(1+\sum_{i=1}^{p_j} c_i^{(j)}\right),
\]
where \(c_i\) and \(c_i^{(j)}\) are assigned weight \(i\); these \(n\) relations generate the ideal of relations among the \(c_i^{(j)}\)'s.
\end{enumerate}

\sgaStableTarget{sga6.e06.cor.4-4}\textbf{Corollary 4.4.} The elements of \(C\) of the form
\[
  \xi_1^{\alpha_1}\cdots \xi_{n-1}^{\alpha_{n-1}},
  \qquad 0\leq \alpha_i\leq n-i,
\]
form a basis of \(C\) as an \(A\)-module.

For the proof of the lemma and corollary we refer to [2], Exposé 4, pages 4.19, from which they are borrowed.

\sgaStableTarget{sga6.e06.sec.4.5}\subsection*{4.5}
Let \(S\) be quasi-compact, let \(E\) be an \(\mathcal O_S\)-module locally free of rank \(n\), let \(\pi=(p_1,\ldots,p_k)\) be a sequence of integers \(\geq 1\) with sum \(n\), and put \(X=D_\pi(E)\). The scheme \(X\) is equipped with canonical quotients \(L_i\). Put \(L_0=0\), \(L_k=f^*E\), where \(f:X\to S\) is the structural morphism. For \(1\leq i\leq k\), let \(F_i\) be the kernel of \(L_i\to L_{i-1}\). Let \(K^\bullet(X)\) be the Grothendieck group of locally free finite-type sheaves on \(X\), respectively of perfect complexes. Denote by \(F_j\) the class of \(F_j\), and by \(\lambda_i^{(j)}\) the class of \(\lambda^i(F_j)\). In the case where \(K^\bullet(X)\) is defined by perfect complexes, one gives the same evident meaning to the series \(\lambda_t(E)\) and \(\lambda_t(F_j)\), where \(E\) is the class of \(E\). Finally,
\[
  \sum_{j=1}^{k}F_j=E.
\]

\sgaStableTarget{sga6.e06.thm.4-6}\textbf{Theorem 4.6.} With the hypotheses and notations of \hyperlink{sga6.e06.sec.4.5}{4.5}, the \(K^\bullet(S)\)-algebra \(K^\bullet(X)\) is generated by the \(\lambda_i^{(j)}\), and the ideal of relations among the \(\lambda_i^{(j)}\)'s is generated by the relations deduced from
\[
  \lambda_t(E)=\prod_{j=1}^{k}\lambda_t(F_j).
\]
Equivalently, \(K^\bullet(X)\) is the \(K^\bullet(S)\)-\(\lambda\)-algebra generated by \(F_1,\ldots,F_k\), subject to the relations
\[
  \text{(i)}\quad \lambda^p(F_i)=0\quad(p>p_i),
  \qquad
  \text{(ii)}\quad \sum_i F_i=E.
\]

We first treat the case where all \(p_i\)'s are equal to \(1\). Then the theorem reads as follows.

\sgaStableTarget{sga6.e06.cor.4-7}\textbf{Corollary 4.7.} Let \(S\) be quasi-compact, let \(E\) be an \(\mathcal O_S\)-module locally free of rank \(n\), let \(X=D(E)\) be the flag bundle of \(E\), and let \(\xi_1,\ldots,\xi_n\) be the classes of the invertible sheaves \(F_1,\ldots,F_n\). Then \(K^\bullet(X)\) is the \(K^\bullet(S)\)-algebra generated by the \(\xi_i\)'s, subject to the relations
\[
  \sigma_i(\xi_1,\ldots,\xi_n)=\lambda^i(E).
\]

We argue by induction on \(n\). For \(n=1\) the assertion is trivial. Suppose it known for \(n-1\). Consider \(\mathbb P(E)\), and on \(\mathbb P(E)\) the exact sequence
\[
  0\longrightarrow F\longrightarrow g^*E\longrightarrow \mathcal O_{\mathbb P(E)}(1)\longrightarrow 0,
\]
where \(g\) is the structural morphism of \(\mathbb P(E)\). By \hyperlink{sga6.e06.lem.4-2}{Lemma 4.2}, \(X\) identifies with the usual flag bundle \(D(F)\) over \(\mathbb P(E)\), and the canonical sheaves correspond under this identification. The diagram
\[
\begin{tikzcd}[column sep=large,row sep=large]
D(E)=D(F) \arrow[r,"h"] \arrow[dr,"f"'] & \mathbb P(E) \arrow[d,"g"]\\
& S
\end{tikzcd}
\]
shows that \(K^\bullet(S)\to K^\bullet(X)\) factors through \(K^\bullet(\mathbb P(E))\). By the induction hypothesis applied to \(D(F)\), and by \hyperlink{sga6.e06.cor.4-4}{Corollary 4.4}, \(K^\bullet(X)\) has over \(K^\bullet(\mathbb P(E))\) a basis formed by
\[
  \xi_2^{\alpha_2}\cdots \xi_n^{\alpha_{n-1}},
  \qquad 0\leq \alpha_i\leq n-i.
\]
Moreover \(K^\bullet(\mathbb P(E))\) has over \(K^\bullet(S)\) the basis \(1,\xi_1,\ldots,\xi_1^{n-1}\). Hence \(K^\bullet(X)\) has over \(K^\bullet(S)\) the basis
\[
  \xi_1^{\alpha_1}\cdots \xi_{n-1}^{\alpha_{n-1}},
  \qquad 0\leq \alpha_i\leq n-i.
\]
In addition, the additivity of \(\lambda_t\) gives \(\prod_i(1+\xi_i t)=\lambda_t(E)\), so the displayed relations are satisfied. If \(C\) is the \(K^\bullet(S)\)-algebra generated by \(\xi_i\)'s and subjected to these relations, then there is a surjective homomorphism \(C\to K^\bullet(X)\). By \hyperlink{sga6.e06.cor.4-4}{Corollary 4.4} this homomorphism sends a basis of \(C\) to a basis of \(K^\bullet(X)\), hence it is an isomorphism.

For a general flag bundle of type \(\pi\), the usual flag bundle \(D(E)\) is obtained from \(D_\pi(E)\) by constructing successively the flag bundles of the \(F_j\)'s. It follows that \(K^\bullet(D(E))\), considered as a module over \(K^\bullet(X)\), has a basis formed by the elements appearing in 4.3.i), the \(\xi_i^{(j)}\)'s being the classes of the sheaves analogous to the \(F_j\)'s in the splitting of the inverse images of the \(F_j\)'s on \(D(E)\). Thus \(K^\bullet(X)\) is a subring of \(K^\bullet(D(E))\). Moreover, additivity gives
\[
  \lambda_i^{(j)}=\sigma_i(\xi_1^{(j)},\ldots,\xi_{p_j}^{(j)}),
\]
and therefore \(K^\bullet(X)\) contains the sub-\(K^\bullet(S)\)-algebra of \(K^\bullet(D(E))\) generated by the \(\lambda_i^{(j)}\)'s. If \(B\) is this subalgebra, then by 4.3.i) the same family of elements is a basis both for the \(B\)-module structure and for the \(K^\bullet(X)\)-module structure. Hence \(K^\bullet(X)=B\), proving generation. The ideal of relations among the \(\lambda_i^{(j)}\)'s is then given by 4.3.ii), taking into account the structure of \(K^\bullet(D(E))\) over \(K^\bullet(S)\) determined by \hyperlink{sga6.e06.cor.4-7}{4.7}, and the fact that the \(\xi_i^{(j)}\)'s are precisely the classes involved in the splitting of the inverse image of \(E\) on \(D(E)\). This proves the first form of 4.6.

For the second form, one observes that the elements \(F_i\) of the \(K^\bullet(S)\)-\(\lambda\)-algebra \(K^\bullet(X)\) satisfy (i) and (ii). If \(A\) is the \(K^\bullet(S)\)-\(\lambda\)-algebra generated by elements \(F_i'\) satisfying (i) and (ii), there is therefore a \(K^\bullet(S)\)-\(\lambda\)-homomorphism \(A\to K^\bullet(X)\) sending \(F_i'\) to \(F_i\). By \hyperlink{sga6.e06.thm.4-6}{4.6} this homomorphism is surjective. Since the \(\lambda\)-ideal generated by (i) and (ii) in \(K^\bullet(S)[F_1,\ldots,F_k]\) contains the usual ideal generated by the first form of the relations, the homomorphism is also injective. Thus it is an isomorphism.

\sgaStableTarget{sga6.e06.cor.4-8}\textbf{Corollary 4.8.} With the hypotheses and notations of \hyperlink{sga6.e06.sec.4.5}{4.5}, put
\[
  c_i^{(j)}=\gamma^i(F_j-p_j),
  \qquad 1\leq i\leq p_j.
\]
Then the \(K^\bullet(S)\)-algebra \(K^\bullet(X)\) is generated by the \(c_i^{(j)}\)'s, and the ideal of relations among them is generated by the relations deduced from
\[
  \gamma_t(E-n)=\prod_{j=1}^{k}\gamma_t(F_j-p_j).
\]
This is merely another formulation of 4.6. Indeed, the \(c_i^{(j)}\)'s may be expressed as functions of the \(\lambda_i^{(j)}\)'s and conversely. Moreover the relation \(\lambda_t(E)=\prod\lambda_t(F_j)\) is equivalent to \(\gamma_t(E-n)=\prod\gamma_t(F_j-p_j)\).

\sgaStableTarget{sga6.e06.sec.4.9}\subsection*{4.9}
Let \(S\) be a scheme. Denote by \(G_S^{n,r}\) the Grassmannian of type \((n,r-n)\) of \(\mathcal O_S^r\), with \(r\ge n\). For variable \(r\), the \(G_S^{n,r}\) form an inductive system, the morphism
\[
  G_S^{n,r}\longrightarrow G_S^{n,r+1}
\]
being defined by considering \(\mathcal O_{G_S^{n,r}}^r\) as the quotient of \(\mathcal O_{G_S^{n,r}}^{r+1}\) by the \((r+1)\)-st factor. In this system of morphisms the canonical quotients therefore correspond by inverse image. The inductive system so defined will be denoted \(G_S^{n,\infty}\).

We define its \(K^\bullet\) by the formula
\sgaStableTarget{sga6.e06.eq.4-9-1}\begin{equation}\tag{4.9.1}
  K^\bullet(G_S^{n,\infty})
  = \varprojlim_r K^\bullet(G_S^{n,r}) .
\end{equation}
Let \(F_1,F_2\) be the canonical sheaves on \(G_S^{n,r}\) \hyperlink{sga6.e06.sec.4.5}{(4.5)}, and \(F_1,F_2\) their classes in \(K^\bullet(G_S^{n,r})\). Put
\[
  c_i=\gamma^i(F_1-n).
\]
When \(r\) varies, the \(c_i\) correspond under the transition morphisms and consequently define elements, still denoted \(c_i\), of \(K^\bullet(G_S^{n,\infty})\).

\sgaStableTarget{sga6.e06.prop.4-10}\textbf{Proposition 4.10.} Let \(S\) be a quasi-compact scheme. Then, with the notation of \hyperlink{sga6.e06.sec.4.9}{4.9}, one has an isomorphism
\[
  K^\bullet(G_S^{n,\infty})\simeq K^\bullet(S)[[c_1,\ldots,c_n]] .
\]

By \hyperlink{sga6.e06.cor.4-8}{4.8}, \(K^\bullet(G_S^{n,r})\) is the \(K^\bullet(S)\)-algebra generated by
\[
  c_i=\gamma^i(F_1-n),\quad 1\le i\le n,
\]
and by
\[
  \gamma^j(F_2-r+n),\quad 1\le j\le r-n,
\]
subject to the relations deduced from
\[
  \gamma_t(F_1-n)\cdot\gamma_t(F_2-r+n)
  = \gamma_t(r-r)=1,
\]
or equivalently
\[
  \gamma_t(F_2-r+n)=1/\gamma_t(F_1-n).
\]
This relation makes it possible to express the \(\gamma^j(F_2-r+n)\) in terms of the \(c_i\), and gives a series of isobaric relations among the \(c_i\) (where \(c_i\) is assumed to have weight \(i\)) by writing that the coefficients of \(1/\gamma_t(F_1-n)\) vanish in degree \(>r-n\); moreover, these relations generate the ideal of relations among the \(c_i\). We may therefore write
\[
  K^\bullet(G_S^{n,r})=K^\bullet(S)[c_1,\ldots,c_n]/I_r,
\]
where \(I_r\) is a homogeneous ideal of \(K^\bullet(S)[c_1,\ldots,c_n]\), of degree \(>r-n\). It then follows from the \hyperlink{sga6.e06.eq.4-9-1}{definition (4.9.1)} that \(K^\bullet(G_S^{n,\infty})\) is the separated completion of \(K^\bullet(S)[c_1,\ldots,c_n]\) for the topology defined by the \(I_r\). This topology is finer than the topology defined by weight; we shall show that it is equivalent to it, which will imply the proposition.

Put \(A=K^\bullet(S)[c_1,\ldots,c_n]/I_r\), and let \(C\) be the \(A\)-algebra generated by generators \(x_1,\ldots,x_n\) subject to the relations
\[
  \sigma_j(x_1,\ldots,x_n)=c_j,
\]
where the \(\sigma_j\) denote the elementary symmetric functions. By \hyperlink{sga6.e06.cor.4-4}{4.4}, the canonical homomorphism \(A\to C\) is injective. We obtain
\[
  \gamma_t(F_1-n)
  =1+c_1t+\cdots+c_nt^n
  =\prod_{i=1}^n(1+x_it).
\]
Since \(1/\gamma_t(F_1-n)\) is a polynomial, so is \(1/(1+x_it)\), for every \(i\). Hence the \(x_i\) are nilpotent, and there exists an integer \(N\) such that every polynomial of degree \(>N\) in the \(x_i\), with coefficients in \(K^\bullet(S)\), is zero. Consequently every polynomial of \(A\) of weight \(>N\) is zero, i.e. every polynomial of \(K^\bullet(S)[c_1,\ldots,c_n]\) of weight \(>N\) belongs to \(I_r\). It follows that the two topologies under consideration are equivalent, completing the proof.

\sgaStableTarget{sga6.e06.sec.5}\section*{5. Applications to projective bundles; study of \(f_*\)}

\sgaStableTarget{sga6.e06.sec.5.1}\subsection*{5.1}
Let \(S\) be quasi-compact, let \(E\) be locally free of rank \(r+1\), and put \(X=\mathbb P(E)\), with structural morphism \(f:X\to S\). On \(X\) one has the canonical exact sequence
\[
  0\longrightarrow F\longrightarrow f^*E\longrightarrow \mathcal O_X(1)\longrightarrow 0.
\]
Let \(E,F,\xi\) denote the corresponding classes in \(K^\bullet\), and put \(x=\xi-1\).

\sgaStableTarget{sga6.e06.sec.5.2}\subsection*{5.2}
For a locally free sheaf \(G\) on \(X\), the higher direct images \(R^if_*G\) need not be locally free of finite type on \(S\). Thus the direct image \(f_*:K^\bullet(X)\to K^\bullet(S)\) is not defined here by alternating sums of higher direct images. \hyperlink{sga6.e06.thm.1-1}{Theorem 1.1} gives a substitute: every \(y\in K^\bullet(X)\) is uniquely expressible in the form
\[
  y=a_0+a_1\xi+\cdots+a_r\xi^r,
  \qquad a_i\in K^\bullet(S),
\]
and we put
\[
  f_*(y)=a_0.
\]
This additive homomorphism satisfies the projection formula. Moreover,
\[
  f_*(1)=1,
  \qquad
  f_*(\xi^n)=0\quad(-r\leq n<0),
  \qquad
  f_*(\xi^n)=s_n\quad(n\geq 0),
\]
where \(s_n\) is the \(n\)-th complete symmetric function in the roots of \(E\). The last relations follow by induction from the dependence relation and the Koszul relation.

\sgaStableTarget{sga6.e06.prop.5-3}\textbf{Proposition 5.3.} For every \(k\geq 0\),
\[
  \Fil^k K^\bullet(X)
  =
  \Fil^k K^\bullet(S)
  +\Fil^{k-1}K^\bullet(S)x+\cdots+\Fil^{k-r}K^\bullet(S)x^r,
\]
as a direct sum of \(K^\bullet(S)\)-modules, where the filtration is the canonical filtration of augmented \(\lambda\)-rings.

Here \(x\) has augmentation zero and
\[
  \gamma_t(x)=\lambda_{t/(1-t)}(1-\xi^{-1})=\frac{1}{1-xt},
\]
so that \(\gamma^i(x)=x^i\). The result follows from the computation of the filtration for free \(\lambda\)-rings in \hyperlink{sga6.e05.prop.4-12}{V, 4.12} and \hyperlink{sga6.e05.rem.4-13}{V, 4.13}, together with the relation \hyperlink{sga6.e06.eq.1-9-1}{(1.9.1)}. Directness follows from the independence of \(1,x,\ldots,x^r\).

\sgaStableTarget{sga6.e06.cor.5-4}\textbf{Corollary 5.4.} The filtration induced on \(K^\bullet(S)\) by \(K^\bullet(X)\) is the original filtration of \(K^\bullet(S)\).

\sgaStableTarget{sga6.e06.cor.5-5}\textbf{Corollary 5.5.} If \(X\) is a product over \(S\) of flag bundles of locally free sheaves, then
\[
  \Gr K^\bullet(S)\longrightarrow \Gr K^\bullet(X)
\]
is injective.

\sgaStableTarget{sga6.e06.cor.5-6}\textbf{Corollary 5.6.} With the notation of the projective bundle above, \(\Gr K^\bullet(X)\) is a free \(\Gr K^\bullet(S)\)-module with basis \(1,x,\ldots,x^r\), and
\[
  \Gr^k K^\bullet(X)
  =\Gr^kK^\bullet(S)\oplus\Gr^{k-1}K^\bullet(S)x\oplus\cdots\oplus\Gr^{k-r}K^\bullet(S)x^r.
\]
The elements \(1,x,\ldots,x^{r+1}\) satisfy
\[
  \sum_{i=0}^{r+1}(-1)^i c_i(E)x^{r+1-i}=0,
\]
where \(c_i(E)\) is the \(i\)-th Chern class of \(E\).

\sgaStableTarget{sga6.e06.cor.5-7}\textbf{Corollary 5.7.} For a flag bundle as in \hyperlink{sga6.e06.sec.4.5}{4.5}, the \(\Gr K^\bullet(S)\)-algebra \(\Gr K^\bullet(X)\) is generated by the classes \(\xi_i^{(\alpha)}\) of \(\xi_i^{(\alpha)}-1\), with relations generated by
\[
  c_t(E)=\prod_{i,\alpha}c_t(\xi_i^{(\alpha)}),
\]
where \(c_t\) denotes the total Chern class.

\sgaStableTarget{sga6.e06.cor.5-8}\textbf{Corollary 5.8.} For every \(k\),
\[
  f_*(\Fil^k K^\bullet(X))\subset \Fil^k K^\bullet(S).
\]
This follows immediately from \hyperlink{sga6.e06.prop.5-3}{Proposition 5.3} and the projection formula.

\sgaStableTarget{sga6.e06.prop.5-9}\textbf{Proposition 5.9.} With the notation of \hyperlink{sga6.e06.sec.5.1}{5.1}, for every \(y\in K^\bullet(X)\), one has
\[
  f_*(\lambda_{-1}(F)y)=0,
  \qquad
  f_*(\gamma^*(F)y)=0.
\]
In particular,
\[
  f_*(\lambda_{-1}(F))=1.
\]
If
\[
  \xi_t:1+K^\bullet(X)[[t]]^+\longrightarrow 1+K^\bullet(S)[[t]]^+
\]
is defined coefficientwise by
\[
  \xi_t\left(\sum a_i t^i\right)=\sum f_*(a_i)t^i,
\]
then, although it is not in general a homomorphism of groups, it satisfies
\[
  \xi_t(mf^*y)=\xi_t(m)y.
\]
Using \(E=F+\xi\) and \(\lambda_t(E)=\lambda_t(F)(1+\xi t)\), one obtains the asserted identities.

\sgaStableTarget{sga6.e06.prop.5-10}\textbf{Proposition 5.10.} Let \(z\in K^\bullet(X)\). The condition
\[
  z(1-\xi)=0
\]
is equivalent to the existence of an element \(y\in K^\bullet(S)\) such that
\[
  \text{(i)}\quad z=y\lambda_{-1}(F),
  \qquad
  \text{(ii)}\quad y\lambda_{-1}(E)=0.
\]
Sufficiency follows from
\[
  \lambda_{-1}(E)=\lambda_{-1}(F)\lambda_{-1}(\xi)=\lambda_{-1}(F)(1-\xi).
\]
Conversely, write \(z=\sum_{i=0}^r a_i\xi^i\). Combining the dependence relation among the powers of \(\xi\) with \(z(1-\xi)=0\), and using the independence of \(1,\ldots,\xi^r\), expresses \(z\) as a multiple of \(\lambda_{-1}(F)\). The top relation gives \(y\lambda_{-1}(E)=0\).

\sgaStableTarget{sga6.e06.prop.5-11}\textbf{Proposition 5.11.} Let \(z\in \Gr K^\bullet(X)\). The condition \(zx=0\) is equivalent to the existence of \(y\in \Gr K^\bullet(S)\) such that
\[
  \text{(i)}\quad z=y\,c_{\mathrm{top}}(F),
  \qquad
  \text{(ii)}\quad y\,c_{\mathrm{top}}(E)=0.
\]
The proof is the same as that of \hyperlink{sga6.e06.prop.5-10}{5.10}, after passing to associated graded rings and using 5.6.

\sgaStableTarget{sga6.e06.sec.6}\section*{6. Study of the filtration of \(K^\bullet(X)\) when \(X\) has an ample sheaf}

Let \(X\) be a scheme possessing an ample sheaf. Then the two definitions of \(K^\bullet(X)\), by locally free finite-type sheaves and by perfect complexes, coincide. Let
\[
  \varepsilon:K^\bullet(X)\longrightarrow H^0(X,\mathbb Z)
\]
be the augmentation.

\sgaStableTarget{sga6.e06.prop.6-1}\textbf{Proposition 6.1.} If \(X\) is quasi-compact and has an ample sheaf, then \(\Fil^1 K^\bullet(X)\) is a nilideal.

It is enough to show that \(E-\varepsilon(E)\) is nilpotent, where \(E\) is the class of a locally free sheaf. By decomposing \(X\) into finitely many disjoint opens and then passing to the flag bundle of \(E\), one reduces to the case where \(E\) is the class of an invertible sheaf. Let \(L\) be the class of an ample invertible sheaf. For \(n\) sufficiently large, \(E^\vee(n)\) is generated by global sections, giving a surjection \((E(-n))^m\to \mathcal O_X\) with locally free kernel. Thus in \(K^\bullet(X)\)
\[
  mEL^{-n}=1+F.
\]
The corresponding \(\gamma\)-relation shows that \(1-EL^{-n}\), and therefore \(L^n-E\), is nilpotent. Since \(E-1=(L^n-1)-(L^n-E)\), and since \(L^n-1\) is also nilpotent for \(n\) large, the assertion follows.

\sgaStableTarget{sga6.e06.prop.6-2}\textbf{Proposition 6.2.} Let \(R\) be an augmented sub-\(\lambda\)-algebra of \(K^\bullet(X)\), generated by the classes of finitely many sheaves. Then the \(\lambda\)-ring filtration of \(R\) is discrete and \(\Fil^1R\) is nilpotent. Moreover, \(R\) is a finite module over its image in \(H^0(X,\mathbb Z)\).

After reducing to constant ranks, \(R\) is generated by finitely many elements \(x_i=E_i-\varepsilon(E_i)\) and their \(\gamma\)-operations. \hyperlink{sga6.e06.prop.6-1}{Proposition 6.1} shows that all these elements are nilpotent. Hence the augmentation ideal of \(R\) is nilpotent; only finitely many monomials in the \(x_i\)'s and the \(\gamma^j(x_i)\)'s remain. This gives both discreteness and finiteness.

\sgaStableTarget{sga6.e06.cor.6-3}\textbf{Corollary 6.3.} Let \(A\) be a commutative ring. An element \(y\) of \(K^\bullet(X)\otimes_{\mathbb Z}A\) is invertible if and only if its augmentation is invertible. In particular, the class of every nowhere-zero locally free sheaf of finite type becomes invertible in \(K^\bullet(X)\otimes_{\mathbb Z}A\).

Indeed, the augmentation ideal is nil.

\sgaStableTarget{sga6.e06.cor.6-4}\textbf{Corollary 6.4.} Let \(X\) be quasi-compact with an ample sheaf, and let \(f:X'\to X\) be proper and perfect. Assume either that \(f\) is projective or that \(X\) is noetherian, so that the finiteness theorem applies. If there is an element \(x'\in K^\bullet(X')\otimes A\) such that the augmentation of \(f_*(x')\) is invertible, then
\[
  f^*:K^\bullet(X)\otimes A\longrightarrow K^\bullet(X')\otimes A
\]
is injective.
Indeed, by the projection formula,
\[
  f_*(f^*y\cdot x')=y\cdot f_*(x'),
\]
and \(f_*(x')\) is invertible by 6.3.

\sgaStableTarget{sga6.e06.sec.6.5}\subsection*{6.5. The codimension filtration}
Assume now that \(X\) is noetherian of dimension \(d\), and has an ample sheaf. Define a filtration \(K^\bullet(X)^i\) as follows: an element \(x\in K^\bullet(X)\) lies in \(K^\bullet(X)^i\) if, for every closed subscheme \(Y\subset X\), it can be represented by a perfect complex \(M^\bullet(Y)\) such that
\begin{enumerate}[label=\textup{\alph*)}]
\item \(\cl(M^\bullet)=x\);
\item \(\codim(\Supp H^j(M^\bullet)\cap Y,Y)\geq i+j\), and \(H^j(M^\bullet)=0\) for \(j\gg 0\).
\end{enumerate}
These \(K^\bullet(X)^i\)'s are additive subgroups.

\sgaStableTarget{sga6.e06.prop.6-6}\textbf{Proposition 6.6.} The filtration of \hyperlink{sga6.e06.sec.6.5}{6.5} is a ring filtration, and \(K^\bullet(X)^{d+1}=0\).

If \(x\) and \(x'\) are represented by perfect complexes \(M^\bullet\) and \(M'{}^\bullet\) satisfying the above codimension estimates, then \(xx'\) is represented by \(M^\bullet\otimes M'{}^\bullet\). The support of \(H^n(M^\bullet\otimes M'{}^\bullet)\) is contained in the union of the intersections of supports with total degree \(n\), and the codimension estimate follows from EGA \(0_{\mathrm{IV}}\), 14.2.2.3. Finally \(K^\bullet(X)^{d+1}=0\), because a support in codimension greater than \(d\) is empty.

\sgaStableTarget{sga6.e06.lem.6-7}\textbf{Lemma 6.7.} Let \(X=\Spec A\) be affine noetherian, let \(x_1,\ldots,x_n\) be finitely many points of \(X\), and let \(E\) be invertible. There exists a section \(s\in \Gamma(X,E)\) which is a basis of \(E_{x_i}\) for every \(i\).

The proof uses closed specializations \(x_i'\) with distinct maximal ideals \(\mathfrak m_i\), the isomorphisms \(M/\mathfrak m_iM\simeq A/\mathfrak m_i\), for \(M=\Gamma(X,E)\), and Nakayama's lemma.

\sgaStableTarget{sga6.e06.lem.6-8}\textbf{Lemma 6.8.} Let \(X\) be noetherian with an ample sheaf \(L\), let \(E\) be invertible, and let \(x_1,\ldots,x_n\) be points of \(X\). There exists an integer \(N\) such that, for every \(k\geq1\), \(E(kN)\) has a global section which is a basis at every \(x_i\).

Choose a power of the ample sheaf and a section whose non-vanishing locus is affine and contains all the \(x_i\)'s; apply \hyperlink{sga6.e06.lem.6-7}{Lemma 6.7} there, and extend after multiplying by a further power.

\sgaStableTarget{sga6.e06.thm.6-9}\textbf{Theorem 6.9.} Let \(X\) be a noetherian scheme of dimension \(d\) with an ample sheaf. Then
\[
  \Fil^{d+1}K^\bullet(X)=0,
\]
and therefore \((\Fil^1K^\bullet(X))^{d+1}=0\).

First, if \(E\) and \(E'\) are invertible, with classes still denoted \(E,E'\), then \(E-E'\in K^\bullet(X)^1\). For a closed \(Y\subset X\), choose one point in each irreducible component. \hyperlink{sga6.e06.lem.6-8}{Lemma 6.8} gives morphisms from suitable negative powers of the ample sheaf to \(E\) and \(E'\) which are isomorphisms at those points. The two-term complex built from these morphisms has class \(E-E'\), and is acyclic at the chosen generic points of \(Y\); hence its cohomology supports have codimension at least \(1\) in \(Y\).

Now take \(x\in \Fil^{d+1}K^\bullet(X)\). It is enough to treat products
\[
  \gamma^{i_1}(E_1-F_1)\cdots\gamma^{i_k}(E_k-F_k),
  \qquad \sum i_j\geq d+1,
\]
where the \(E_i,F_i\) are represented by locally free sheaves of the same constant rank. Pull back to the product of the flag bundles of these sheaves. The pullback is injective on \(K^\bullet\), and each class splits into sums of invertible classes. For invertible classes \(\xi,\eta\),
\[
  \gamma_t(\xi-\eta)
  =\frac{\gamma_t(\xi-1)}{\gamma_t(\eta-1)}
  =\frac{1+(\xi-1)t}{1+(\eta-1)t},
\]
so
\[
  \gamma^i(\xi-\eta)=(-1)^{i-1}(\xi-\eta)(\eta-1)^{i-1}\in K^\bullet(X')^i.
\]
Thus the pullback of \(x\) belongs to \(K^\bullet(X')^{d+1}\). Multiplication by a basis element of \(K^\bullet(X')\) over \(K^\bullet(X)\) of relative codimension \(r\) gives an element of \(K^\bullet(X')^{d+1+r}=0\), since \(\dim X'=d+r\). The basis property and injectivity then give \(x=0\).

\sgaStableTarget{sga6.e06.problem.6-10}\textbf{Problem 6.10.} Under the hypotheses of \hyperlink{sga6.e06.thm.6-9}{Theorem 6.9}, is it true that for every \(k\) one has
\[
  K^\bullet(X)^k=\Fil^kK^\bullet(X)?
\]

\sgaStableTarget{sga6.e06.bibliography-for-expose-vi}\subsection*{Bibliography for Exposé VI}
\begin{enumerate}[label={\textup{[\arabic*]}}]
\item H. Bass, \emph{K-Theory and stable algebra}, Publications Mathématiques de l'IHÉS, no. 22.
\item C. Chevalley, Séminaire 1958: \emph{Anneaux de Chow et applications}.
\item M. Hakim, \emph{Topos localement annelés et schémas relatifs}, thesis in preparation.
\item R. Hartshorne, \emph{Residues and duality}, Lecture Notes in Mathematics 20, Springer-Verlag.
\item J.-P. Serre, \emph{Le groupe de Grothendieck des schémas en groupes réductifs}, Publications Mathématiques de l'IHÉS, no. 34.
\end{enumerate}



\providecommand{\Kosz}{\mathrm K}
\providecommand{\Proj}{\operatorname{Proj}}
\providecommand{\Sym}{\operatorname{Sym}}
\providecommand{\gr}{\operatorname{gr}}
\providecommand{\Tor}{\operatorname{Tor}}
\providecommand{\Spec}{\operatorname{Spec}}
\providecommand{\Imf}{\operatorname{Im}}
\providecommand{\Ker}{\operatorname{Ker}}
\providecommand{\Id}{\operatorname{Id}}
\providecommand{\rad}{\operatorname{rad}}
\providecommand{\cO}{\mathcal O}
\providecommand{\cS}{\mathcal S}
\providecommand{\cJ}{\mathcal J}
\providecommand{\cN}{\mathcal N}
\providecommand{\bfP}{\mathbf P}
\providecommand{\bfA}{\mathbf A}


%%%%%%%%%%  scan idx 429  |  folio 416  %%%%%%%%%%
\SGAPrintPage{416}%
% Exposé VII: back to the volume's Sym and Proj
\renewcommand{\Sym}{\operatorname{Sym}}
\renewcommand{\Projz}{\operatorname{Proj}_0}
\par\noindent\textbf{Disjoint English VII rendering probe, printed416--465. This probe does not extend the assembled prefix beyond296.}\par
\sgaStableTarget{sga6.e07.expose}\section*{Exposé VII\\[3pt] REGULAR IMMERSIONS AND COMPUTATION OF \texorpdfstring{$K^{\bullet}$}{K0} OF A BLOWN-UP SCHEME}
\begin{center}
\uline{by Pierre BERTHELOT}
\end{center}

\sgaStableTarget{sga6.e07.sec.1}\section*{1. Generalities on regular immersions}

As the classical notion of regular immersion (EGA IV 16.9.2) proves inadequate in the non-noetherian case, we shall define a weaker notion, which nevertheless gives the desired properties and is equivalent to the usual definition in the noetherian case.

\sgaStableTarget{sga6.e07.def.1-1}\textbf{Definition 1.1.} \uline{Let} $A$ \uline{be a commutative ring}, $E$ \uline{a projective} $A$-\uline{module of finite type}, $u:E\longrightarrow A$ \uline{an} $A$-\uline{linear homomorphism}. \uline{We say that} $u$ \uline{is regular if the Koszul complex} $K_\bullet(u)$ \uline{defined by} $u$ \uline{is acyclic in dimension} $>0$, \uline{and hence is a resolution of} $B=A/u(E)$.

Recall that the Koszul complex associated with $u$ is defined as follows: the graded $A$-module $K_\bullet(u)$ is the exterior algebra $\Lambda(E)$, and the boundary operator is interior multiplication by $u$ regarded as an element of the dual $\check E$: it is an antiderivation of degree $-1$ of $\Lambda(E)$, equal to $u$ on $E$.

\sgaStableTarget{sga6.e07.lem.1-2}\textbf{Lemma 1.2.} \uline{Let} $u:E\longrightarrow A$ \uline{be a regular homomorphism}. \uline{Let} $A'$ \uline{be an} $A$-\uline{algebra}, $u':E'\longrightarrow A'$ \uline{the homomorphism deduced from} $u$ \uline{by base change}; \uline{suppose that} $u'$ \uline{is regular}. \uline{Put} $J=u(E)$, $J'=JA'=u(E')$. \uline{Then we have the following relations for every integer} $n\geq0$:
\sgaStableTarget{sga6.e07.eq.1-2-1}\sgaStableTarget{sga6.e07.eq.1-2-2}\sgaStableTarget{sga6.e07.eq.1-2-3}\sgaStableTarget{sga6.e07.eq.1-2-4}\begin{align}
 J^n/J^{n+1}\otimes_AA'&\xrightarrow{\sim}J'^n/J'^{n+1}, \tag{1.2.1}\\
 \Tor_i^A(J^n/J^{n+1},A')&=0\quad\text{\uline{for} }i\geq1; \tag{1.2.2}\\
 J^n\otimes_AA'&\xrightarrow{\sim}J'^n, \tag{1.2.3}\\
 \Tor_i^A(J^n,A')&=0\quad\text{\uline{for} }i\geq1. \tag{1.2.4}
\end{align}

a) First suppose that the $J^n/J^{n+1}$ are projective over $B=A/J$; since $J$ is of finite type, they are locally free of finite type over $B$.

%%%%%%%%%%  scan idx 430  |  folio 417  %%%%%%%%%%
\SGAPrintPage{417}%
The relations to be proved can be checked pointwise on $\Spec(A)$, so we are reduced to the case where $A$ is local, with $u(E)\subset\rad(A)$. The $J^n/J^{n+1}$ are then free over $B$, and the relations $\Tor_i^A(J^n/J^{n+1},A')=0$ for $i\geq1$ are equivalent to the relations $\Tor_i^A(B,A')=0$ for $i\geq1$. Now we can compute the $\Tor_i^A$ using the projective resolution $K_\bullet(u)$ of $B$; they are the homology groups of the complex $K_\bullet(u)\otimes_AA'=K_\bullet(u')$. Since $u'$ is assumed regular, this last complex is acyclic in degrees $>0$, whence the relations $\Tor_i^A(J^n/J^{n+1},A')=0$, for $i\geq1$.

We deduce the relations $\Tor_i^A(A/J^n,A')=0$ for $i\geq1$ by induction on $n$ using the exact Tor sequence associated with the exact sequence:
\[
 0\longrightarrow J^n/J^{n+1}\longrightarrow A/J^{n+1}\longrightarrow A/J^n\longrightarrow0.
\]
Since for every $k$ we have $J^kA'=J'^k$, this exact sequence and $\Tor_1^A(A/J^n,A')=0$ give the relations $J^n/J^{n+1}\otimes_AA'\xrightarrow{\sim}J'^n/J'^{n+1}$. From the exact sequence:
\[
 0\longrightarrow J^n\longrightarrow A\longrightarrow A/J^n\longrightarrow0,
\]
and $\Tor_1^A(A/J^n,A')=0$, we deduce the relations $J^n\otimes_AA'\simeq J'^n$. Finally, from the exact Tor sequence associated with the exact sequence:
\[
 0\longrightarrow J^{n+1}\longrightarrow J^n\longrightarrow J^n/J^{n+1}\longrightarrow0,
\]
we deduce the relations $\Tor_i^A(J^n,A')=0$ for $i\geq1$ by induction on $n$, the case $n=0$ being evident.

b) We shall now show that the $J^n/J^{n+1}$ are projective over $B=A/J$. Since the question is local on $\Spec(A)$, we may suppose $E$ free, say $E=(A)^d$, so that $u$ is defined by a sequence $f_1,\ldots,f_d$ of elements of $A$. Put $A_0=\mathbb Z[T_1,\ldots,T_d]$ and let $A_0\longrightarrow A$ be the homomorphism that sends

%%%%%%%%%%  scan idx 431  |  folio 418  %%%%%%%%%%
\SGAPrintPage{418}%
$T_i$ to $f_i$; denote by $u_0:E_0=(A_0)^d\longrightarrow A_0$ the homomorphism defined by the sequence $(T_1,\ldots,T_d)$, so that $u$ is deduced from $u_0$ by base change from $A_0$ to $A$. Since the sequence $(T_1,\ldots,T_d)$ is regular (in the classical sense), the homomorphism $u_0$ is regular, and the hypotheses of the lemma hold for the base change $A_0\longrightarrow A$. Now the $J_0^n/J_0^{n+1}$ are free of finite type over $\mathbb Z$, so we can apply the results of a). Consequently $J^n/J^{n+1}\xleftarrow{\sim}J_0^n/J_0^{n+1}\otimes_{A_0}A$, and $J^n/J^{n+1}$ is free of finite type.

\sgaStableTarget{sga6.e07.prop.1-3}\textbf{Proposition 1.3.} \uline{Let} $A$ \uline{be a ring}, $E$ \uline{a projective} $A$-\uline{module of finite type}, $u:E\longrightarrow A$ \uline{a regular homomorphism}, $J=u(E)$. \uline{Then the ideal} $J$ \uline{is quasi-regular} (EGA IV 16.9.7), \uline{that is} (EGA IV 16.9.8):
\begin{enumerate}[label=\roman*)]
\item $J$ \uline{is of finite type}.
\item $J/J^2$ \uline{is a projective} $A/J$-\uline{module of finite type};
\item \uline{the canonical homomorphism}
\end{enumerate}
\sgaStableTarget{sga6.e07.eq.1-3-1}\begin{equation}\tag{1.3.1}
 \Sym_{A/J}(J/J^2)\longrightarrow \gr_J(A)=\bigoplus_{n\geq0}J^n/J^{n+1},
\end{equation}
\uline{is an isomorphism}.

We saw in the proof of 1.2 b) that $J/J^2$ is projective of finite type. Resume the notation of 1.2 b) (since the question is local on $\Spec(A)$, we again suppose $E$ free of finite type). Then the canonical homomorphism (1.3.1) is deduced from the analogous homomorphism for $J_0$ by the base change $A_0\longrightarrow A$, by the relations 1.2. The latter is an isomorphism, so the same is true of (1.3.1).

\sgaStableTarget{sga6.e07.def.1-4}\textbf{Definition 1.4.} \uline{Let} $(X,\underline{A})$ \uline{be a commutatively locally ringed site}, $\underline{E}$ \uline{an} $\underline{A}$-\uline{Module locally free of finite type}, $u:\underline{E}\longrightarrow \underline{A}$ \uline{an} $\underline{A}$-\uline{linear homomorphism}. \uline{We say that} $u$ \uline{is} regular \uline{if the Koszul complex} $\underline{K}_\bullet(u)$ \uline{defined by} $u$ \uline{is acyclic in degrees} $>0$. \uline{Let} $\underline{J}$ \uline{be an Ideal of} $\underline{A}$; \uline{we say that} $\underline{J}$ \uline{is}

%%%%%%%%%%  scan idx 432  |  folio 419  %%%%%%%%%%
\SGAPrintPage{419}%
\uline{regular}, \uline{if}, \uline{locally on} $X$, \uline{there exist a free} $\underline{A}$-\uline{Module of finite type} $\underline{E}$, \uline{and a surjective regular homomorphism}: $\underline{E}\longrightarrow \underline{J}$.

\uline{Let} $i:Y\longrightarrow X$ \uline{be an immersion of preschemes}, \uline{and let} $U$ \uline{be an open subset of} $X$ \uline{such that} $i(Y)\subset U$ \uline{and such that} $i$ \uline{is a closed immersion of} $Y$ \uline{into} $U$. \uline{We say that} $i$ \uline{is regular if the ideal of} $\underline{O}_U$ \uline{defining} $i(Y)$ \uline{is regular} (\uline{a condition independent of} $U$). \uline{If} $y\in Y$, \uline{we say that} $i$ \uline{is regular at} $y$ \uline{if there exists a neighborhood} $V$ \uline{of} $i(y)$ \uline{such that the immersion} $i^{-1}(V)\longrightarrow V$ \uline{is regular}. \uline{For an immersion to be regular}, \uline{it is therefore necessary and sufficient that it be regular in a neighborhood of each point}.

\uline{This definition conflicts with} EGA IV 16.9.2, \uline{and replaces that definition}. However, in the noetherian case the two definitions are equivalent (EGA IV 19.5.1). An immersion that can locally be defined by a regular sequence is regular; a regular immersion is quasi-regular (1.3).

\textbf{1.4.1.} Let $X$ be a scheme, $\underline{J}$ an Ideal of finite type of $\underline{O}_X$, $x$ a point of $\Supp(\underline{O}_X/\underline{J})$. We shall say that a family of sections $f_1,\ldots,f_d$ of $\underline{J}$ on a neighborhood of $x$ forms a \uline{minimum system of generators of} $\underline{J}$ \uline{in a neighborhood of} $x$ if the $(f_i)_x$ form a family of generators of $\underline{J}_x$ of minimum cardinality, i.e. their images in $\underline{J}_x\otimes k(x)$ form a basis of this vector space over $k(x)$. Note that the $f_i$ generate $\underline{J}$ on a neighborhood of $x$.

\textbf{1.4.2.} To say that an immersion is regular is also to say that it is defined in some open subset of $X$ by an ideal of finite type $\underline{J}$, and that for every point $x\in\Supp(\underline{O}_X/\underline{J})$ and every minimum system of generators of $\underline{J}$ in a neighborhood of $x$, the Koszul complex associated with this sequence is acyclic in a neighborhood of $x$: indeed, let $f_1,\ldots,f_d$ be a minimum system of

%%%%%%%%%%  scan idx 433  |  folio 420  %%%%%%%%%%
\SGAPrintPage{420}%
generators of $\underline{J}$ in a neighborhood of $x$ such that the associated Koszul complex is acyclic (one exists when the immersion is assumed regular); let $g_1,\ldots,g_d$ be another minimum system of generators of $\underline{J}$; there exists a neighborhood of $x$ on which $\underline{J}$ \var{can be generated either the $f_i$ or by the $g_j$}{can be generated either by the $f_i$ or by the $g_j$}{EN-CONTRIB-VII-001}, and consequently there exists on this neighborhood a matrix with coefficients in $\underline{O}_X$ expressing the $g_j$ in terms of the $f_i$. Now the images of the $f_i$ (resp. $g_j$) in $\underline{J}/\underline{J}^2$ form a basis at the point $x$, hence in a neighborhood of $x$, and consequently this matrix is invertible in a neighborhood of $x$. It therefore defines an isomorphism from the Koszul complex defined by the $g_j$ onto the Koszul complex defined by the $f_i$, which shows that the Koszul complex defined by the $g_j$ is acyclic.

\textbf{1.4.3.} Assuming only that $\underline{J}$ is of finite type and $\underline{J}/\underline{J}^2$ is locally free of finite type, note that every minimum system of generators of $\underline{J}$ in a neighborhood of a point $y$ of $Y$ gives a basis of $\underline{J}/\underline{J}^2$ in a neighborhood of $y$. Conversely, given a basis of $\underline{J}/\underline{J}^2$ in a neighborhood of $y$, we can lift it in a neighborhood of $y$ to a family of sections of $\underline{J}$, which generate $\underline{J}_y$ by Nakayama, hence $\underline{J}$ \var{in a neighborhood of $\ddot y$ since}{in a neighborhood of $y$ since}{EN-CONTRIB-VII-002} $\underline{J}$ is of finite type; moreover these generators have minimum cardinality, since they give a basis of $\underline{J}/\underline{J}^2$.

\sgaStableTarget{sga6.e07.prop.1-5}\textbf{Proposition 1.5.} (\uline{Compare} EGA IV 19.1.5 (ii)). \uline{Let} $i:Y\longrightarrow X$ \uline{be an immersion}. \uline{If} $i$ \uline{is regular}, \uline{the same is true of every immersion deduced from} $i$ \uline{by flat base change}. \uline{For} $i$ \uline{to be regular}, \uline{it is necessary and sufficient that it be so locally for the faithfully flat quasi-compact topology}.

Stability under flat base change is evident.

Conversely, let $i'$ be the immersion deduced from $i$ by a

%%%%%%%%%%  scan idx 434  |  folio 421  %%%%%%%%%%
\SGAPrintPage{421}%
faithfully flat quasi-compact base change. If $i'$ is regular, then the Ideal $\underline{J}$ defining $i$ is of finite type, and $\underline{J}/\underline{J}^2$ is locally free of finite type, since these properties descend under faithfully flat quasi-compact morphisms. If $f_1,\ldots,f_d$ is a sequence of sections of $\underline{J}$ lifting a basis of $\underline{J}/\underline{J}^2$, the Koszul complex associated with the inverse images of the $f_i$ is acyclic in degrees $>0$ by 1.4; since it is the inverse image of the Koszul complex associated with the $f_i$, the latter is acyclic in degrees $>0$, and $i$ is regular.

\sgaStableTarget{sga6.e07.def.1-6}\textbf{Definition 1.6.} \uline{Let} $i:Y\longrightarrow X$ \uline{be an immersion defined in an open subset} $U$ \uline{of} $X$ \uline{by an Ideal} $\underline{J}$ \uline{of} $\underline{O}_X$; \uline{the} $\underline{O}_Y$-\uline{Module} $\underline{J}/\underline{J}^2$ (\uline{independent of} $U$) \uline{is called the} conormal sheaf of $Y$ in $X$, \uline{and is denoted by} $\underline{N}_{X/Y}$. \uline{If} $i$ \uline{is regular}, $\underline{J}/\underline{J}^2$ \uline{is locally free of finite type}, \uline{and its rank is called the} codimension of $Y$ in $X$, \uline{or the} \uline{codimension of} $i$.

Given two closed immersions: $Z\xrightarrow{j}Y\xrightarrow{i}X$, the conormal sheaves form an exact sequence:
\sgaStableTarget{sga6.e07.eq.1-6-1}\begin{equation}\tag{1.6.1}
 j^*(\underline{N}_{X/Y})\longrightarrow \underline{N}_{X/Z}\longrightarrow \underline{N}_{Y/Z}\longrightarrow0.
\end{equation}
Indeed, let $\underline{I}$ be the ideal of $Y$ in $X$, $\underline{J}$ the ideal of $Z$ in $X$, $\underline{K}$ the ideal of $Z$ in $Y$. We have the exact sequence:
\[
0\longrightarrow \underline{I}\longrightarrow \underline{J}\longrightarrow \underline{K}\longrightarrow0;
\]
whence, by tensoring with $\underline{O}_X/\underline{I}=\underline{O}_Y$, the exact sequence:
\[
 \underline{I}/\underline{I}^2\longrightarrow \underline{J}\otimes_{\underline{O}_X}\underline{O}_Y\longrightarrow \underline{K}\otimes_{\underline{O}_X}\underline{O}_Y\longrightarrow0;
\]
tensoring with $\underline{O}_Z=\underline{O}_X/\underline{J}=\underline{O}_Y/\underline{K}$, we obtain the exact sequence:
\[
 j^*(\underline{I}/\underline{I}^2)\longrightarrow \underline{J}\otimes_{\underline{O}_X}\underline{O}_Z\longrightarrow \underline{K}\otimes_{\underline{O}_Y}\underline{O}_Z\longrightarrow0,
\]
which is the required sequence.
%%%%%%%%%%  scan idx 435  |  folio 422  %%%%%%%%%%
\SGAPrintPage{422}%
\sgaStableTarget{sga6.e07.prop.1-7}\textbf{Proposition 1.7.} (\uline{Compare} EGA IV 19.1.5 (iv)). \uline{Let} $i:Y\longrightarrow X$, \uline{and} $j:Z\longrightarrow Y$ \uline{be two immersions}:
\begin{enumerate}[label=\roman*)]
\item \uline{If} $i$ \uline{and} $j$ \uline{are regular}, $i\circ j$ \uline{is regular}, \uline{and the sequence}:
\[
0\longrightarrow j^*(\underline{N}_{X/Y})\longrightarrow \underline{N}_{X/Z}\longrightarrow \underline{N}_{Y/Z}\longrightarrow0
\]
\uline{is exact}; \uline{in particular}, $\codim(i\circ j)=\codim(i)+\codim(j)$.
\item \uline{If} $i\circ j$ \uline{and} $i$ \uline{are regular}, \uline{and if the sequence}:
\[
0\longrightarrow j^*(\underline{N}_{X/Y})\longrightarrow \underline{N}_{X/Z}\longrightarrow \underline{N}_{Y/Z}\longrightarrow0,
\]
\uline{is exact and locally split}, $j$ \uline{is regular};
\item \uline{If} $X$ \uline{is noetherian}, \uline{and if} $i\circ j$ \uline{and} $j$ \uline{are regular}, $i$ \uline{is regular at the points of} $Z$.
\end{enumerate}

Since the question is local on $X$, we may suppose $X$ affine and the immersions closed; let $X=\Spec(A)$, $Y=\Spec(B)$, $Z=\Spec(C)$, with $B=A/I$, $C=A/J=B/K$.

i) After localizing, we may suppose $I/I^2$ free over $B$, and $K/K^2$ free over $C$ (1.3). Let $f_1,\ldots,f_n$ be elements of $A$ lifting a basis of $I/I^2$, $g_1,\ldots,g_m$ elements of $A$ lifting a basis of $K/K^2$; then, after further localization, we may suppose that $f_1,\ldots,f_n$ generate $I$, that $f_1,\ldots,f_n,g_1,\ldots,g_m$ generate $J$, and that $\bar g_1,\ldots,\bar g_m$ generate $K$, where $\bar g_j$ is the image of $g_j$ in $B$ (1.4); moreover, $f_1,\ldots,f_n$ form a minimum system of generators of $I$, and $\bar g_1,\ldots,\bar g_m$ a minimum system of generators of $K$.

Consider the diagram:
\sgaStableTarget{sga6.e07.eq.1-7-1}\begin{equation}\tag{1.7.1}
\begin{tikzcd}[column sep=large,row sep=large]
A_0=\mathbb Z[T_1,\ldots,T_{n+m}] \arrow[r] \arrow[d] & A \arrow[d]\\
B_0=\mathbb Z[T_1,\ldots,T_m] \arrow[r] \arrow[d] & B \arrow[d]\\
C_0=\mathbb Z \arrow[r] & C,
\end{tikzcd}
\end{equation}

%%%%%%%%%%  scan idx 436  |  folio 423  %%%%%%%%%%
\SGAPrintPage{423}%
where the horizontal arrows are defined by sending $T_i$ to $f_i$ for $1\leq i\leq n$, and $T_{n+j}$ to $g_j$ for $1\leq j\leq m$, where the homomorphism $A_0\longrightarrow B_0$ is defined by sending $T_i$ to $0$ for $1\leq i\leq n$, and $T_{n+j}$ to $T_j$ for $1\leq j\leq m$, and the homomorphism $B_0\longrightarrow C_0$ by setting the $T_j$, $1\leq j\leq m$, to zero. We then have:
\[
 B=A\otimes_{A_0}B_0,
 \qquad
 C=A\otimes_{A_0}C_0=B\otimes_{B_0}C_0.
\]
The situation under consideration is therefore deduced from the situation $(A_0,B_0,C_0)$ by the base change $A_0\longrightarrow A$. The corresponding immersions $i_0$, $j_0$ and $i_0\circ j_0$ are regular; consequently, the homology groups of the Koszul complexes associated with the sequences $(f_1,\ldots,f_n)$, $(f_1,\ldots,f_n,g_1,\ldots,g_m)$ and $(\bar g_1,\ldots,\bar g_m)$ are respectively the $\Tor_i^{A_0}(B_0,A)$, $\Tor_i^{A_0}(C_0,A)$, \var{$\Tor_i^{A_0}(C_0,B)$}{$\Tor_i^{B_0}(C_0,B)$}{EN-CONTRIB-VII-003}. Consequently:
\sgaStableTarget{sga6.e07.eq.1-7-2}\begin{equation}\tag{1.7.2}
\begin{aligned}
 i\ \text{regular} &\Longleftrightarrow \forall i\geq1,
 \Tor_i^{A_0}(B_0,A)=0;\\
 i\circ j\ \text{regular} &\Longleftrightarrow \forall i\geq1,
 \Tor_i^{A_0}(C_0,A)=0;\\
 j\ \text{regular} &\Longleftrightarrow \forall i\geq1,
 \Tor_i^{B_0}(C_0,B)=0.
\end{aligned}
\end{equation}

We then use the transitivity spectral sequence ([1], XVI 5)
\sgaStableTarget{sga6.e07.eq.1-7-3}\begin{equation}\tag{1.7.3}
E^2_{pq}=\Tor_p^{B_0}\bigl(C_0,\Tor_q^{A_0}(B_0,A)\bigr)
\Longrightarrow \Tor_n^{A_0}(C_0,A)..
\end{equation}

By the hypotheses of i), $E^2_{pq}=0$ for $p$ or $q\ne0$. It therefore degenerates, and $i\circ j$ is regular by (1.7.2).

By (1.6.1), the sequence:
\[
 j^*(\underline{N}_{X/Y})\longrightarrow \underline{N}_{X/Z}\longrightarrow \underline{N}_{Y/Z}\longrightarrow0
\]
is exact. Since $i\circ j$ is regular, the three conormal sheaves are locally free of finite type; moreover, their rank at a point is the

%%%%%%%%%%  scan idx 437  |  folio 424  %%%%%%%%%%
\SGAPrintPage{424}%
number of elements in a minimum system of generators of the corresponding ideal; by the preceding proof, we have
\[
 \rang(\underline{N}_{X/Z})=\rang(\underline{N}_{X/Y})+\rang(\underline{N}_{Y/Z}).
\]
It follows that the kernel of $j^*(\underline{N}_{X/Y})\longrightarrow \underline{N}_{X/Z}$ is locally free of rank zero, and hence that \var{the sequence (1.8.1) is exact.}{the sequence (1.6.1) is exact.}{EN-CONTRIB-VII-004}

ii) We may suppose $J/J^2$ free over $C$, and $I/I^2$ free over $B$; moreover, the hypothesis on the sequence of conormal sheaves allows us to suppose that $I/I^2\otimes_A A/J$ is a direct summand of $J/J^2$. We can therefore find elements $f_1,\ldots,f_n$ of $A$ lifting a basis of $I/I^2$, and elements $g_1,\ldots,g_m$ of $A$ such that $f_1,\ldots,f_n,g_1,\ldots,g_m$ lift a basis of $J/J^2$, the images of $g_1,\ldots,g_m$ generating $K/K^2$; after further localization, we may suppose that $f_1,\ldots,f_n$ form a minimum system of generators of $I$, that $f_1,\ldots,f_n,g_1,\ldots,g_m$ form a minimum system of generators of $J$, and that $\bar g_1,\ldots,\bar g_m$ generate $K$. We can then construct the diagram (1.7.1) and use the spectral sequence (1.7.3).

By (1.7.2), since $i$ is regular, we have $E^2_{pq}=0$ for $q\geq1$; we therefore obtain edge-isomorphisms: $E^2_{p,0}\simeq\Tor_p^{A_0}(C_0,A)$. Since $i\circ j$ is regular, this implies by (1.7.2) that $E^2_{p,0}=0$ for $p\geq1$, and hence that $j$ is regular, again by (1.7.2).

iii) This last property is proved in EGA IV 19.1.5 (iv), since in the noetherian case the definition in EGA IV coincides with the one we adopt here.

\sgaStableTarget{sga6.e07.prop.1-8}\textbf{Proposition 1.8.} \uline{Let} \(X\) \uline{be a scheme}, \(\underline{E}\) \uline{an} \(\underline{O}_X\)-\uline{Module locally free of finite type}, \(u:\underline{E}\to\underline{O}_X\) \uline{a regular homomorphism}, \(\underline{J}=u(\underline{E})\). \uline{Denote by} \(Y\)

%%%%%%%%%%  scan idx 438  |  folio 425  %%%%%%%%%%
\SGAPrintPage{425}%
\uline{the closed subscheme of} \(X\) \uline{defined by} \(\underline{J}\), \(X'\) \uline{the scheme obtained by blowing up} \(Y\), \(P\) \uline{the projective bundle associated with} \(\underline{E}\), \(p\) \uline{the structural morphism} \(P\to X\). \uline{On} \(P\), \uline{consider the canonical exact sequence}:
\sgaStableTarget{sga6.e07.eq.1-8-1}\[
0\to \underline{F}\to p^*(\underline{E})\to\underline{O}_P(1)\to0,
\]
\uline{and the composite morphism}:
\[
 v:\underline{F}\longrightarrow p^*(\underline{E})\xrightarrow{p^*(u)}\underline{O}_P .
\]

\uline{Then}:
\begin{enumerate}[label=\roman*)]
\item \uline{The formation of} \(X'\) \uline{commutes with every base change} \(f:T\to X\) \uline{such that} \(f^*(u)\) \uline{is regular}.
\item \uline{The closed immersion} \(j:X'\to P\) \uline{defined by the surjective homomorphism} \var{\(\underline{\underline{S}}(\underline{E})\longrightarrow\bigoplus_{n\ge0}\underline{J}^n\)}{\(\mathbf{S}(\underline{E})\longrightarrow\bigoplus_{n\ge0}\underline{J}^n\)}{EN-CONTRIB-VII-005} \uline{defined by} \(u\) \uline{in degree} \(1\) \uline{is a regular immersion}, \uline{its ideal being the image of the regular morphism} \(v\); \uline{in particular, if} \(\operatorname{codim}(i)\) \uline{is constant}, \(\operatorname{codim}(j)=\operatorname{codim}(i)-1\).
\item \var{\uline{The} \(\underline{\underline{S}}(\underline{E})\)-\uline{Module}}{\uline{The} \(\mathbf{S}(\underline{E})\)-\uline{Module}}{EN-CONTRIB-VII-006} \(\bigoplus_{n\ge0}\underline{J}^n\) \var{\uline{is perfect as a graded} \(\underline{\underline{S}}(\underline{E})\)-\uline{Module}}{\uline{is perfect as a graded} \(\mathbf{S}(\underline{E})\)-\uline{Module}}{EN-CONTRIB-VII-007}, \uline{i}.\uline{e}. \uline{locally there exists a finite}, \uline{graded resolution of} \(\bigoplus_{n\ge0}\underline{J}^n\) \var{\uline{by graded} \(\underline{\underline{S}}(\underline{E})\)-\uline{Modules locally free of finite type}}{\uline{by graded} \(\mathbf{S}(\underline{E})\)-\uline{Modules locally free of finite type}}{EN-CONTRIB-VII-008} (VI 1.2).
\end{enumerate}

i) Since \(X'=\Proj(\bigoplus_{n\ge0}\underline{J}^n)\), the result is an immediate consequence of the relations (1.2.3) and of EGA II.

ii) Since the stated properties are local on \(X\), we may suppose \(X\) affine, and moreover \(\underline{E}\) free over \(X\). Put \(X=\Spec(A)\).

First we give the proof in the standard case \(A=\mathbb Z[T_1,\ldots,T_d]\), with \(E\) endowed with a basis \(e_1,\ldots,e_d\) such that \(u(e_i)=T_i\). Then \(P=\Proj\bigl(\mathbb Z[T_1,\ldots,T_d,X_1,\ldots,X_d]\bigr)\), the grading being taken with respect to the \(X_i\). To show that \(v\) is a regular morphism whose image is the Ideal of \(j\), it suffices to do so locally on \(P\), hence on the

%%%%%%%%%%  scan idx 439  |  folio 426  %%%%%%%%%%
\SGAPrintPage{426}%
\(D^+(X_i)\); let us work, for example, on \(D^+(X_1)\). Then \(p^*(\underline{E})\) is defined by the module \(\mathbb Z[T_1,\ldots,T_d,X_2/X_1,\ldots,X_d/X_1]\otimes E\), and \(\underline{O}_P(1)\) by \(\mathbb Z[T_1,\ldots,T_d,X_2/X_1,\ldots,X_d/X_1]\), the canonical morphism sending \(e_i\) to \(X_i/X_1\) for every \(i\). Then \(F\) has as a basis the elements \(e_i-e_1.X_i/X_1\) for \(i>1\).

We have \(v(e_i-e_1.X_i/X_1)=T_i-T_1.X_i/X_1\); put \(T'_i=T_i-T_1X_i/X_1\).
Since the sequence of the \(T_i\) is regular in \(\mathbb Z[T_1,\ldots,T_d,X_2/X_1,\ldots,X_d/X_1]\), the same is true of the sequence \(T_1,T'_2,\ldots,T'_d\) deduced from it by an invertible matrix, and hence also of the subsequence \(T'_2,\ldots,T'_d\) (we are in a noetherian case). Consequently, \(v\) is a regular morphism.

We must therefore show that the \(T'_i\) generate the Ideal of \(j\). In \(\bigoplus_{n\ge0}J^n\), for \(m\le p\), we shall denote by \(x^{(m)}\) the element \(x\) of \(J^p\), considered as having degree \(m\). Over \(D^+(X_1)\), \(X'=\Spec\left(\left(\bigoplus_{n\ge0}J^n\right)_{(T_1^{(1)})}\right)\), and \(j\) is defined by the morphism sending \(X_i\) to \(T_i^{(1)}\). Under this morphism, the image of \(T'_i\) is \(T_i^{(0)}-T_1^{(0)}.T_i^{(1)}/T_1^{(1)}=\bigl((T_iT_1)^{(1)}-(T_1T_i)^{(1)}\bigr)/T_1^{(1)}=0\). Thus the \(T'_i\) belong to the Ideal of \(j\).

\var{Since every element of \(\mathbb Z[T_1,\ldots,T_d,X_2/X_1,\ldots,X_d,X_1]\)}{Since every element of \(\mathbb Z[T_1,\ldots,T_d,X_2/X_1,\ldots,X_d/X_1]\)}{EN-CONTRIB-VII-009} can be expressed as a polynomial in \(T_1,T'_2,\ldots,T'_d,X_2/X_1,\ldots,X_d/X_1\), to show that the \(T'_i\) generate the Ideal of \(j\) it suffices to show that no nonzero element of \(\mathbb Z[T_1,X_2/X_1,\ldots,X_d/X_1]\) belongs to the Ideal of \(j\), which amounts to showing that the homomorphism
\[
\mathbb Z[T_1,X_1,\ldots,X_d]\longrightarrow\bigoplus_{n\ge0}J^n
\]
is injective, since \(T_1^{(1)}\) is not a zero divisor in \(\bigoplus_{n\ge0}J^n\). Grading \(\mathbb Z[T_1,X_1,\ldots,X_d]\) by assigning degree \(0\) to \(T_1\) and degree \(1\) to the \(X_i\), the kernel of the homomorphism under consideration is a graded ideal. Now it is clear that
%%%%%%%%%%  scan idx 440  |  folio 427  %%%%%%%%%%
\SGAPrintPage{427}%
this kernel is contained in that of the composite homomorphism:
\[
\mathbb Z[T_1,X_1,\ldots,X_d]\longrightarrow
\bigoplus_{n\ge0}J^n\longrightarrow A,
\]
the last homomorphism being defined in degree \(n\) by the inclusion of \(J^n\) in \(A\). Now the kernel of the composite homomorphism is the principal ideal generated by \((X_1-T_1)\), which, for reasons of integrality, contains no nonzero homogeneous elements. Thus the first homomorphism is injective, completing the proof in this case.

Return to the general case where \(A\) is arbitrary; we may still suppose \(\underline{E}\) free over \(X\). Attach the index \(0\) to the special case just treated; then, by part i) of the proposition, the situation under consideration is obtained by the base change \(A_0\to A\) from the situation studied above, the homomorphism \(A_0\to A\) being defined by taking a basis of \(E\), its image \(f_1,\ldots,f_d\) in \(J\), and sending \(T_i\) to \(f_i\). Put \(B_0=\left(\bigoplus_{n\ge0}J_0^n\right)_{(T_1^{(1)})}\), \(B=\left(\bigoplus_{n\ge0}J^n\right)_{(f_1^{(1)})}\); then \(B=B_0\otimes_{A_0}A\); the Koszul complex defined by \(v:\underline{F}\to\underline{O}_P\) is deduced by base change from the Koszul complex defined by \(v_0\); since the latter is a projective resolution of \(B_0\) (working on a suitable open subset), the homology groups of \(K_\bullet(v)\) on this open subset are the \(\Tor_i^{A_0}(A,B_0)\). These are the degree \(0\) components of the \(\Tor_i^{A_0}\left(A,\left(\bigoplus_{n\ge0}J_0^n\right)_{T_1^{(1)}}\right)=\left(\bigoplus_{n\ge0}\Tor_i^{A_0}(A,J_0^n)\right)_{T_1^{(1)}}\); the \(\Tor_i^{A_0}(A,J_0^n)\) vanish for \(i\ge1\) by the relations 1.2, so \(K_\bullet(v)\) is a resolution of \(B\), which shows that \(j\) is regular.

iii) Since the property is local on \(X\), we may again suppose \(X\) affine and \(\underline{E}\) free of finite type. Resume the notation of ii).

First consider the standard case of ii). The ring \(A_0\) is then noetherian and of finite cohomological dimension, and the algebra \var{\(\underline{\underline{S}}(E_0)\)}{\(\mathbf{S}(E_0)\)}{EN-CONTRIB-VII-010}

%%%%%%%%%%  scan idx 441  |  folio 428  %%%%%%%%%%
\SGAPrintPage{428}%
likewise. We can then construct step by step a partial graded resolution of \(\bigoplus_{n\ge0}J_0^n\) \var{by graded free \(\underline{\underline{S}}(E_0)\)-modules}{by graded free \(\mathbf{S}(E_0)\)-modules}{EN-CONTRIB-VII-011} of finite type, of length greater than the cohomological dimension of \(A_0\); \var{the last kernel is then a graded \(\underline{\underline{S}}(E_0)\)-module}{the last kernel is then a graded \(\mathbf{S}(E_0)\)-module}{EN-CONTRIB-VII-012}, projective of finite type \var{as an ungraded \(\underline{\underline{S}}(E_0)\)-module}{as an ungraded \(\mathbf{S}(E_0)\)-module}{EN-CONTRIB-VII-013}, hence locally free of finite type over \(\Spec(A_0)\) \var{as a graded \(\underline{\underline{S}}(\underline{E}_0)\)-module}{as a graded \(\mathbf{S}(\underline{E}_0)\)-module}{EN-CONTRIB-VII-014} by VI 1.5.1. We obtain the announced finite resolution by replacing the last module of the partial resolution by the corresponding kernel.

Tensoring the resolution thus constructed with \(A\), we obtain \var{a graded complex of graded \(\underline{\underline{S}}(E)\)-modules}{a graded complex of graded \(\mathbf{S}(E)\)-modules}{EN-CONTRIB-VII-015} locally free of finite type, and must show that it is a resolution of \(\bigoplus_{n\ge0}J^n\). The \(i\)-th homology Module of the resulting complex \var{is a graded \(\underline{\underline{S}}(E)\)-Module whose}{is a graded \(\mathbf{S}(E)\)-Module whose}{EN-CONTRIB-VII-016} \(n\)-th component is \(\Tor_i^{A_0}(A,J_0^n)\), which vanishes for \(i\ge1\) by (1.2.4). Consequently, by (1.2.3), we have constructed a resolution of the required type.

\sgaStableTarget{sga6.e07.prop.1-9}\textbf{Proposition 1.9.} \uline{Let} \(i:Y\to X\) \uline{be a regular closed immersion}, \(X'\) \uline{the} \(X\)-\uline{scheme obtained by blowing up} \(Y\), \uline{and} \(f:X'\to X\) \uline{the structural morphism}. \uline{Then} \(i\) \uline{and} \(f\) \uline{are perfect morphisms}. (III 4.1.).

Since the property is local on \(X\), we may suppose given an \(\underline{O}_X\)-Module locally free of finite type \(\underline{E}\), and a regular homomorphism \(u:\underline{E}\to\underline{O}_X\) whose image is the ideal \(\underline{J}\) of \(Y\). The Koszul complex \(\underline{K}_\bullet(u)\) is then a finite resolution of \(\underline{O}_Y\) by \(\underline{O}_X\)-Modules locally free of finite type; consequently, \(\underline{O}_Y\) is a perfect \(\underline{O}_X\)-Module. It follows that every regular immersion is a perfect morphism. By 1.8, \(f\) can be written as the composite of a regular immersion \var{\(X'\to\underline{\underline{P}}(\underline{E})\)}{\(X'\to\bfP(\underline{E})\)}{EN-CONTRIB-VII-017} and the structural morphism \var{\(\underline{\underline{P}}(\underline{E})\to X\)}{\(\bfP(\underline{E})\to X\)}{EN-CONTRIB-VII-018}, which is smooth and hence perfect; it is therefore

%%%%%%%%%%  scan idx 442  |  folio 429  %%%%%%%%%%
\SGAPrintPage{429}%
the composite of two perfect morphisms, and consequently a perfect morphism.

\sgaStableTarget{sga6.e07.prop.1-10}\textbf{Proposition 1.10.} \uline{Let} \(X\), \(Y\) \uline{be two schemes smooth over} \(S\). \uline{Then every} \(S\)-\uline{immersion of} \(X\) \uline{into} \(Y\) \uline{is regular}. \uline{In particular}, \uline{every section of a smooth morphism is a regular immersion}.

The proposition follows immediately from EGA IV 17.12.1, since a regular immersion in the sense of EGA IV is regular in the sense of 1.4.

\sgaStableTarget{sga6.e07.sec.2}\section*{2. Computations on regular immersions}

\sgaStableTarget{sga6.e07.para.2-1}\subsection*{2.1.}
Let \(A\) be a ring, and \(B\) and \(C\) two \(A\)-algebras. We propose to endow \(\bigoplus_{i\ge0}\Tor_i^A(B,C)\) with a graded \(A\)-algebra structure, and hence to define homomorphisms:
\begin{equation}\tag{2.1.1}
 \Tor_i^A(B,C)\otimes_A\Tor_j^A(B,C)\longrightarrow \Tor_{i+j}^A(B,C).
\end{equation}

Let \(P_\bullet\) be a resolution of \(B\) by projective \(A\)-modules; the \(\Tor_i^A(B,C)\) are the homology groups of the complex \(P_\bullet\otimes_A C\). We know how to define canonical homomorphisms ([1], IV 6):
\begin{equation}\tag{2.1.2}
 H_i(P_\bullet\otimes_A C)\otimes_A H_j(P_\bullet\otimes_A C)
 \longrightarrow H_{i+j}\bigl((P_\bullet\otimes_A C)\otimes_A(P_\bullet\otimes_A C)\bigr).
\end{equation}

Furthermore, \(P_\bullet\otimes_A P_\bullet\) is a complex of projective \(A\)-modules, in degrees \(\geq0\), and \(H_0(P_\bullet\otimes_A P_\bullet)=B\otimes_A B\); moreover, \(P_\bullet\) is a resolution of \(B\). Consequently, the augmentation morphism \(B\otimes_A B\to B\) extends to a morphism of complexes: \(P_\bullet\otimes_A P_\bullet\to P_\bullet\) and therefore defines, using the homomorphism \(C\otimes_A C\to C\) deduced from the algebra structure of \(C\), homomorphisms:
\begin{equation}\tag{2.1.3}
 H_{i+j}\bigl((P_\bullet\otimes_A C)\otimes_A(P_\bullet\otimes_A C)\bigr)
 \longrightarrow H_{i+j}(P_\bullet\otimes_A C).
\end{equation}

%%%%%%%%%%  scan idx 443  |  folio 430  %%%%%%%%%%
\SGAPrintPage{430}%
Composing (2.1.2) and (2.1.3), we obtain the homomorphisms (2.1.1).

The algebra structure defined by the homomorphisms (2.1.1) is associative (cf. [1] XI 2). Furthermore, the \(A\)-module \(H_i(P_\bullet\otimes_A C)\) has a standard \(B\otimes_A C\)-module structure, with \(C\) acting in the evident way, and \(B\) by functoriality; likewise, the \(H_j\bigl((P_\bullet\otimes_A C)\otimes_A(P_\bullet\otimes_A C)\bigr)\) carry two \(B\otimes_A C\)-module structures. The homomorphisms (2.1.2) and (2.1.3) are clearly compatible with these two structures: for (2.1.2), by functoriality (working in the category of \(A\)-modules), and for (2.1.3), by the classical result on morphisms of an acyclic complex. We have therefore endowed \(\bigoplus_{i\ge0}\Tor_i^A(B,C)\) with an associative \(B\otimes_A C\)-algebra structure.

\subsection*{2.2.}
Suppose that \(B\) admits a resolution by projective \(A\)-modules \(P_\bullet\) such that \(P_\bullet=\bigoplus_{i\ge0}P_i\) is a graded \(A\)-algebra, with the differential satisfying:
\[
\begin{array}{ll}
\text{a)}& P_0\longrightarrow B\text{ is a ring homomorphism~;}\SGAAuthorFootMark{1}\\[0.4ex]
\text{b)}& \forall x\in P_i,\ \forall y\in P_j,\quad d(xy)=d(x)y+(-1)^i x\,d(y).
\end{array}
\]
\SGAAuthorFootText{1}{One can show that such a resolution always exists.}
Then the ``product'' maps define commutative diagrams for \(i,j\ge0\):
\sgaStableTarget{sga6.e07.eq.2-2-1}\begin{equation}\notag
\begin{tikzcd}[column sep=large,row sep=large]
P_i\otimes_A P_j \arrow[r] \arrow[d] & (P_{i-1}\otimes_A P_j)\oplus(P_i\otimes_A P_{j-1}) \arrow[d]\\
P_{i+j} \arrow[r] & P_{i+j-1},
\end{tikzcd}
\end{equation}
and consequently a morphism of complexes \(P_\bullet\otimes_A P_\bullet\to P_\bullet\) extending the augmentation homomorphism: \(B\otimes_A B\to B\). This shows that the algebra structure of \(\bigoplus_{i\ge0}\Tor_i^A(B,C)\) can be computed from that of \(P_\bullet\). In particular, if \(P_\bullet\) is an alternating \(A\)-algebra (i.e. such that for every

%%%%%%%%%%  scan idx 444  |  folio 431  %%%%%%%%%%
\SGAPrintPage{431}%
\(x\in P_\bullet\), we have \(x^2=0\)), then the same is true of \(\bigoplus_{i\ge0}\Tor_i^A(B,C)\).

This last remark applies when \(B\) is the quotient of \(A\) by a regular ideal, since we can then take for \(P_\bullet\) the Koszul complex associated with a regular homomorphism defining the ideal of \(B\).

\subsection*{2.3.}
The \(B\otimes_A C\)-algebra structure defined in 2.1 on \(\bigoplus_{i\ge0}\Tor_i^A(B,C)\) is functorial in \(B\) and \(C\). We can therefore globalize the preceding constructions: let \(X\) be a scheme, \(X'\), \(X''\) two \(X\)-schemes; we can then define the morphisms (2.1.1) locally on \(X'\times_X X''\); they glue by functoriality and define on \(\bigoplus_{i\ge0}\underline{\Tor}_i^{\underline{O}_X}(\underline{O}_{X'},\underline{O}_{X''})\) an associative graded \(\underline{O}_{X'\times_X X''}\)-Algebra structure.

\sgaStableTarget{sga6.e07.lem.2-4}\textbf{Lemma 2.4.} \uline{Let} \(X',X''\) \uline{be two} \(X\)-\uline{schemes}, \uline{and} \(\underline{M}\) \uline{a flat} \(\underline{O}_{X'}\)-\uline{Module}. \uline{Then for every} \(i\ge0\), \uline{we have a canonical isomorphism}:
\sgaStableTarget{sga6.e07.eq.2-4-1}\begin{equation}\notag
\underline{\Tor}_i^{\underline{O}_X}(\underline{M},\underline{O}_{X''})\simeq
\underline{M}\otimes_{\underline{O}_{X'}}\underline{\Tor}_i^{\underline{O}_X}(\underline{O}_{X'},\underline{O}_{X''}).
\end{equation}

It suffices to define a functorial isomorphism in the affine case; we can then define it globally by gluing in the general case.

Let \(X=\Spec(A)\), \(X'=\Spec(A')\), \(X''=\Spec(A'')\); \(\underline{M}\) is defined by a flat \(A'\)-module \var{\(M\), Let \(P_\bullet\)}{\(M\). Let \(P_\bullet\)}{EN-CONTRIB-VII-019} be a projective resolution of \(A''\) over \(A\); the \(\Tor_i^A(A',A'')\) are the homology groups of \(A'\otimes_A P_\bullet\). Since \(M\) is flat over \(A'\), we have a canonical isomorphism: \(M\otimes_{A'}\Tor_i^A(A',A'')\xrightarrow{\sim}H_i\bigl(M\otimes_{A'}(A'\otimes_A P_\bullet)\bigr)\). Furthermore, the canonical isomorphism: \(M\otimes_{A'}(A'\otimes_A P_\bullet)\xrightarrow{\sim}M\otimes_A P_\bullet\) gives a canonical isomorphism: \(H_i\bigl(M\otimes_{A'}(A'\otimes_A P_\bullet)\bigr)\xrightarrow{\sim}\Tor_i^A(M,A'')\), whence the result.

\sgaStableTarget{sga6.e07.prop.2-5}\textbf{Proposition 2.5.} \uline{Let} \(X\) \uline{be a scheme}, \(Y\) \uline{and} \(Z\) \uline{two closed subschemes of} \(X\), \(T=Y\times_X Z\). \uline{Suppose that for every} \(x\in T\), \uline{there exists an open neighborhood}
%%%%%%%%%%  scan idx 445  |  folio 432  %%%%%%%%%%
\SGAPrintPage{432}%
\(U\) \uline{of} \(x\), \uline{and a sequence of sections} \(f_1,\ldots,f_p,g_1,\ldots,g_r,h_1,\ldots,h_s\) \uline{of} \(\underline{O}_U\) \uline{such that the Koszul complex associated with} \(f_1,\ldots,f_p,g_1,\ldots,g_r\) \uline{is a resolution of} \(\underline{O}_{Y\cap U}\), \uline{the Koszul complex associated with} \(f_1,\ldots,f_p,h_1,\ldots,h_s\) \uline{a resolution of} \(\underline{O}_{Z\cap U}\), \uline{and the Koszul complex associated with} \(f_1,\ldots,f_p,g_1,\ldots,g_r,h_1,\ldots,h_s\) \uline{a resolution of} \(\underline{O}_{T\cap U}\) (\var{\uline{conditions satisfied in particular in the noetherian case when the sequence} \(f_1,\ldots,f_p,g_1,\ldots g_r,h_1,\ldots,h_s\)}{\uline{conditions satisfied in particular in the noetherian case when the sequence} \(f_1,\ldots,f_p,g_1,\ldots,g_r,h_1,\ldots,h_s\)}{EN-CONTRIB-VII-020} \uline{is regular and} \(Y=V(f_1,\ldots,f_p,g_1,\ldots,g_r)\), \(Z=V(f_1,\ldots,f_p,h_1,\ldots,h_s)\)). \uline{Let} \(\underline{I},\underline{J}\) \uline{be the ideals of} \(Y,Z\). \uline{We then have the following isomorphisms}:
\[
\begin{array}{ll}
\text{i)}& \forall i\ge0,\quad
\underline{\Tor}_i^{\underline{O}_X}(\underline{O}_Y,\underline{O}_Z)\xrightarrow{\sim}
\bigwedge^i\underline{\Tor}_1^{\underline{O}_X}(\underline{O}_Y,\underline{O}_Z);\\[0.8ex]
\text{ii)}&
\underline{\Tor}_1^{\underline{O}_X}(\underline{O}_Y,\underline{O}_Z)\simeq
\underline{I}\cap \underline{J}/\underline{I}.\underline{J}\simeq
\Ker\bigl(i^*(\underline{N}_{Y/X})\times j^*(\underline{N}_{Z/X})\to \underline{N}_{T/X}\bigr),
\end{array}
\]
\uline{where we denote by} \(i\) \uline{the immersion of} \(T\) \uline{into} \(Y\) \uline{and by} \(j\) \uline{the immersion of} \(T\) \uline{into} \(Z\).

\uline{In particular}, \uline{for every regular closed immersion} \(Y\to X\), \uline{we have}
\[
\begin{array}{ll}
\text{i')}& \forall i\ge0,\quad
\underline{\Tor}_i^{\underline{O}_X}(\underline{O}_Y,\underline{O}_Y)\xrightarrow{\sim}
\bigwedge^i\underline{\Tor}_1^{\underline{O}_X}(\underline{O}_Y,\underline{O}_Y);\\[0.8ex]
\text{ii')}&
\underline{\Tor}_1^{\underline{O}_X}(\underline{O}_Y,\underline{O}_Y)\simeq \underline{N}_{Y/X}.
\end{array}
\]

i) We saw in 2.3 that \(\bigoplus_{i\ge0}\underline{\Tor}_i^{\underline{O}_X}(\underline{O}_Y,\underline{O}_Z)\) could be endowed with an \(\underline{O}_T\)-Algebra structure; by 2.2, this algebra is alternating, since the immersion of \(Y\) (or \(Z\)) into \(X\) is regular. There therefore exists one and only one algebra homomorphism:
\sgaStableTarget{sga6.e07.eq.2-5-1}\begin{equation}\tag{2.5.1}
\bigwedge \underline{\Tor}_1^{\underline{O}_X}(\underline{O}_Y,\underline{O}_Z)\longrightarrow
\bigoplus_{i\ge0}\underline{\Tor}_i^{\underline{O}_X}(\underline{O}_Y,\underline{O}_Z),
\end{equation}
extending the identity in degree \(1\).

To see that it is an isomorphism, we can check locally. Thus we work in a neighborhood of a point of \(T\) and consider sections \(f_1,\ldots,f_p,g_1,\ldots,g_r,h_1,\ldots,h_s\) of \(\underline{O}_X\) on this neighborhood, having the properties stated in 2.5. Put \(\underline{E}=(\underline{O}_X)^p\), \(\underline{E}'=(\underline{O}_X)^r\), and let \(u:\underline{E}\to\underline{O}_X\) be the homomorphism defined by the \(f_i\), \(v:\underline{E}'\to\underline{O}_X\) the homomorphism defined by the \(g_j\), and \(\underline{K}_\bullet(u)\) and \(\underline{K}_\bullet(v)\) the associated

%%%%%%%%%%  scan idx 446  |  folio 433  %%%%%%%%%%
\SGAPrintPage{433}%
Koszul complexes. The Koszul complex associated with the sequence \(f_1,\ldots,f_p,g_1,\ldots,g_r\) is isomorphic to the tensor-product complex \(\underline{K}_\bullet(u)\otimes_{\underline{O}_X}\underline{K}_\bullet(v)\). Since it is a resolution of \(\underline{O}_Y\) by free \(\underline{O}_X\)-Modules of finite type, the \(\underline{\Tor}_i^{\underline{O}_X}(\underline{O}_Y,\underline{O}_Z)\) are the homology sheaves of the complex \(\underline{K}_\bullet(u)\otimes_{\underline{O}_X}\underline{K}_\bullet(v)\otimes_{\underline{O}_X}\underline{O}_Z\). We can regard it as the tensor product of the complexes \(\underline{K}_\bullet(u)\) and \(\underline{K}_\bullet(v)\otimes_{\underline{O}_X}\underline{O}_Z\), and compute its homology using the spectral sequences associated with a bicomplex (EGA \(0_{\mathrm{III}}\) (11.3.5)), which abut to the homology of the associated simple complex.

Consider the spectral sequence such that \({}'E^2_{pq}=H^I_p\bigl(H^{II}_q(K_{\bullet\bullet})\bigr)\), in the notation of EGA \(0_{\mathrm{III}}\) (11.3.5), with the first index referring to the complex \(\underline{K}_\bullet(u)\) and the second to the complex \(\underline{K}_\bullet(v)\otimes_{\underline{O}_X}\underline{O}_Z\). Now the complex \(\underline{K}_\bullet(v)\otimes_{\underline{O}_X}\underline{O}_Z\) is the Koszul complex over \(\underline{O}_Z\) associated with the sequence of sections \(\overline g_1,\ldots,\overline g_r\), where the \(\overline g_i\) are the images of the \(g_i\) in \(\underline{O}_Z\), and this complex is a resolution of \(\underline{O}_T\) over \(\underline{O}_Z\) since the immersions \(Z\to X\), and \(T\to Z\to X\) are regular (cf. the proof of 1.7 ii). Since \(\bigwedge^p\underline{E}\) is free over \(\underline{O}_X\), the complex \(\bigwedge^p\underline{E}\otimes_{\underline{O}_X}(\underline{K}_\bullet(v)\otimes_{\underline{O}_X}\underline{O}_Z)\) is a resolution of \(\bigwedge^p\underline{E}\otimes_{\underline{O}_X}\underline{O}_T\). Consequently the spectral sequence degenerates, and the \(\underline{\Tor}_i^{\underline{O}_X}(\underline{O}_Y,\underline{O}_Z)\) are the homology sheaves of the complex \(\underline{K}_\bullet(u)\otimes_{\underline{O}_X}\underline{O}_T\). The differential of \(\underline{K}_\bullet(u)\) is expressed in terms of the \(f_i\), which annihilate \(\underline{O}_T\); consequently the differential of \(\underline{K}_\bullet(u)\otimes_{\underline{O}_X}\underline{O}_T\) is zero, and we therefore obtain the isomorphisms:
\sgaStableTarget{sga6.e07.eq.2-5-2}\begin{equation}\tag{2.5.2}
\underline{\Tor}_i^{\underline{O}_X}(\underline{O}_Y,\underline{O}_Z)
\xrightarrow{\sim}
\underline{O}_T\otimes_{\underline{O}_X}\bigwedge^i\underline{E}
\longrightarrow
\bigwedge^i(\underline{O}_T)^p.
\end{equation}

In particular, for \(i=1\) we find:
% VII probe display variant EN-CONTRIB-VII-021
\SGACaptureAppState{EN-CONTRIB-VII-021}%
\ifnum\SGAReaderForm=1\relax\SGAAppMark{EN-CONTRIB-VII-021}\fi
\ifnum\SGAReaderForm=0\relax
\[
\underline{\Tor}_1^{\underline{O}_X}(\underline{O}_Y,\underline{O}_Z)\xrightarrow{\sim}(\underline{O}_T)^n,
\]
\else
\[
\underline{\Tor}_1^{\underline{O}_X}(\underline{O}_Y,\underline{O}_Z)\xrightarrow{\sim}(\underline{O}_T)^p,
\]
\fi
whence for every \(i\):
\sgaStableTarget{sga6.e07.eq.2-5-3}% VII probe display variant EN-CONTRIB-VII-022
\SGACaptureAppState{EN-CONTRIB-VII-022}%
\ifnum\SGAReaderForm=1\relax\SGAAppMark{EN-CONTRIB-VII-022}\fi
\ifnum\SGAReaderForm=0\relax
\begin{equation}\tag{2.5.3}
\underline{\Tor}_i^{\underline{O}_X}(\underline{O}_Y,\underline{O}_Z)
\xrightarrow{\sim}
\bigwedge^i\underline{\Tor}_i^{\underline{O}_X}(\underline{O}_Y,\underline{O}_Z).
\end{equation}
\else
\begin{equation}\tag{2.5.3}
\underline{\Tor}_i^{\underline{O}_X}(\underline{O}_Y,\underline{O}_Z)
\xrightarrow{\sim}
\bigwedge^i\underline{\Tor}_1^{\underline{O}_X}(\underline{O}_Y,\underline{O}_Z).
\end{equation}
\fi

%%%%%%%%%%  scan idx 447  |  folio 434  %%%%%%%%%%
\SGAPrintPage{434}%
Now it follows from 2.2 that the isomorphism:
% VII probe display variant EN-CONTRIB-VII-023
\SGACaptureAppState{EN-CONTRIB-VII-023}%
\ifnum\SGAReaderForm=1\relax\SGAAppMark{EN-CONTRIB-VII-023}\fi
\ifnum\SGAReaderForm=0\relax
\[
\bigwedge(\underline{O}_T)^n\xrightarrow{\sim}
\bigoplus_{i\ge0}\underline{\Tor}_i^{\underline{O}_X}(\underline{O}_Y,\underline{O}_Z)
\]
\else
\[
\bigwedge(\underline{O}_T)^p\xrightarrow{\sim}
\bigoplus_{i\ge0}\underline{\Tor}_i^{\underline{O}_X}(\underline{O}_Y,\underline{O}_Z)
\]
\fi
defined by the isomorphisms (2.5.2) is compatible with the algebra structures of the two sides; hence the same is true of the isomorphism:
% VII probe display variant EN-CONTRIB-VII-024
\SGACaptureAppState{EN-CONTRIB-VII-024}%
\ifnum\SGAReaderForm=1\relax\SGAAppMark{EN-CONTRIB-VII-024}\fi
\ifnum\SGAReaderForm=0\relax
\[
\bigwedge\underline{\Tor}_i^{\underline{O}_X}(\underline{O}_Y,\underline{O}_Z)\xrightarrow{\sim}
\bigoplus_{i\ge0}\underline{\Tor}_i^{\underline{O}_X}(\underline{O}_Y,\underline{O}_Z)
\]
\else
\[
\bigwedge\underline{\Tor}_1^{\underline{O}_X}(\underline{O}_Y,\underline{O}_Z)\xrightarrow{\sim}
\bigoplus_{i\ge0}\underline{\Tor}_i^{\underline{O}_X}(\underline{O}_Y,\underline{O}_Z)
\]
\fi
defined by the isomorphisms (2.5.3). By uniqueness, the latter is therefore the restriction of the homomorphism (2.5.1) that was defined globally. The homomorphism (2.5.1) is thus an isomorphism.

ii) Consider the exact sequence:
\[
0\longrightarrow \underline{I}\longrightarrow \underline{O}_X\longrightarrow \underline{O}_Y\longrightarrow0,
\]
and the resulting exact \(\underline{\Tor}\) sequence:
\[
0\longrightarrow\underline{\Tor}_1^{\underline{O}_X}(\underline{O}_Y,\underline{O}_Z)
\longrightarrow \underline{I}\otimes_{\underline{O}_X}\underline{O}_Z
\longrightarrow \underline{O}_Z.
\]
It follows that \(\underline{\Tor}_1^{\underline{O}_X}(\underline{O}_Y,\underline{O}_Z)\) is the kernel of the homomorphism:
\[
\underline{I}\otimes_{\underline{O}_X}(\underline{O}_X/\underline{J})\longrightarrow \underline{O}_X/J,
\]
that is, \(\underline{I}\cap \underline{J}/\underline{I}.\underline{J}\).

To prove the last relation, we shall show that there is an exact sequence:
\[
\underline{0}\longrightarrow \underline{I}\cap \underline{J}/\underline{I}.\underline{J}
\longrightarrow i^*(\underline{N}_{Y/X})\times j^*(\underline{N}_{Z/X})
\longrightarrow \underline{N}_{T/X}\longrightarrow0.
\]
The homomorphisms \(i^*(\underline{N}_{Y/X})\to \underline{N}_{T/X}\) and \(j^*(\underline{N}_{Z/X})\to \underline{N}_{T/X}\) were defined in (1.6.1); here we take their difference. Furthermore, \(\underline{I}\cap \underline{J}\) maps naturally into \(\underline{I}/\underline{I}^2\), hence into \(i^*(\underline{I}/\underline{I}^2)\), and likewise into \(j^*(\underline{J}/\underline{J}^2)\), and these homomorphisms vanish on \(\underline{I}.\underline{J}\), since \(\underline{I},\underline{J}\subset \underline{K}\) (where \(\underline{K}\) denotes the ideal of \(T\)), defining the arrow on the left.

%%%%%%%%%%  scan idx 448  |  folio 435  %%%%%%%%%%
\SGAPrintPage{435}%
To check exactness of the sequence, we can work locally on \(T\) and use the sections defined in the statement of 2.5. We then have a basis of \(\underline{N}_{T/X}\) formed by the images of the \(f_i,g_j,h_k\), a basis of \(i^*(\underline{N}_{Y/X})\) formed by the images of the \(f_i,g_j\), and a basis of \(j^*(\underline{N}_{Z/X})\) formed by the images of the \(f_i,h_k\). It is then clear that \(i^*(\underline{N}_{Y/X})\times j^*(\underline{N}_{Z/X})\to \underline{N}_{T/X}\) is surjective, and that its kernel consists of the elements whose coordinates on the \((\overline g_j,0)\) and the \((0,\overline h_k)\) vanish, and whose coordinates on \((\overline f_i,0)\) and \((0,\overline f_i)\) are equal for every \(i\). Consequently the sequence is also exact in the middle. The three \(\underline{O}_T\)-Modules involved are locally free of finite type; since \(\underline{\Tor}_1^{\underline{O}_X}(\underline{O}_Y,\underline{O}_Z)\) has rank \(p\) by the proof of i), we see immediately that the sequence is exact.

\sgaStableTarget{sga6.e07.para.2-6}\subsection*{2.6.}
We now suppose that \(X\) is a scheme having an ample sheaf. Every closed subscheme \(Y\) of \(X\) then has an ample sheaf. We know that in this case ``global resolutions exist'', and that the \(K^{\bullet}\) formed respectively with locally free sheaves of finite type and with perfect complexes are identical (II, 2.2.9). If moreover the immersion \(i:Y\to X\) is regular, it is a perfect morphism, and we can define an additive homomorphism \(i_*:K^{\bullet}(Y)\longrightarrow K^{\bullet}(X)\) (IV 2.12) by associating with the class of a perfect complex \(\underline{L}^\bullet\) the class of the complex \var{\(\underline{\underline{R}}i_*(\underline{L}^\bullet)=i_*(\underline{L}^\bullet)\)}{\(\mathbb Ri_*(\underline{L}^\bullet)=i_*(\underline{L}^\bullet)\)}{EN-CONTRIB-VII-025} (\(i_*\) being an exact functor).

\textbf{Proposition 2.7.} \uline{Let} \(X\) \uline{be a scheme having an ample sheaf}, \(i:Y\to X\) \uline{a regular closed immersion}, \(\underline{N}\) \uline{the conormal sheaf of} \(Y\) \uline{in} \(X\). \uline{Then}, \uline{for every} \(x\in K^{\bullet}(Y)\), \uline{we have the relation}:
\sgaStableTarget{sga6.e07.eq.2-7-1}\begin{equation}\notag
i^*(i_*(x))=x.\lambda_{-1}(N),
\end{equation}
\uline{where} \(N\) \uline{is the class of} \(\underline{N}\) \uline{in} \(K^{\bullet}(Y)\), \uline{and} \(\lambda_{-1}(N)=\sum_{i\ge0}(-1)^i\lambda^i(N)\).

%%%%%%%%%%  scan idx 449  |  folio 436  %%%%%%%%%%
\SGAPrintPage{436}%
Since both sides of the relation to be proved are additive, we may suppose that \(x\) is the class of an \(\underline{O}_Y\)-Module locally free of finite type \(\underline{M}\). The element \(i^*(i_*(x))\) is then \var{the class of the complex \(\underline{\underline{L}} i^*(\underline{\underline{R}}i_*(\underline{M}))\)}{the class of the complex \(\mathbb L i^*(\mathbb Ri_*(\underline{M}))\)}{EN-CONTRIB-VII-026}, obtained by taking \var{a projective resolution of the complex \(\underline{\underline{R}}i_*(\underline{M})=i_*(\underline{M})\)}{a projective resolution of the complex \(\mathbb Ri_*(\underline{M})=i_*(\underline{M})\)}{EN-CONTRIB-VII-027} and tensoring it with \(\underline{O}_Y\). The homology sheaves of this complex are therefore the \(\underline{\Tor}_i^{\underline{O}_X}(\underline{M},\underline{O}_Y)\). By 2.4, we have the isomorphism:
\[
\underline{\Tor}_i^{\underline{O}_X}(\underline{M},\underline{O}_Y)
\xrightarrow{\sim}
\underline{M}\otimes_{\underline{O}_Y}\underline{\Tor}_i^{\underline{O}_X}(\underline{O}_Y,\underline{O}_Y).
\]
By 2.5, we have the isomorphism:
\[
\underline{\Tor}_i^{\underline{O}_X}(\underline{O}_Y,\underline{O}_Y)
\xrightarrow{\sim}
\bigwedge^i\underline{N}.
\]

The \(\underline{\Tor}_i^{\underline{O}_X}(\underline{M},\underline{O}_Y)\) are therefore \(\underline{O}_Y\)-Modules locally free of finite type, and the class in \(K^{\bullet}(Y)\) of the complex \var{\(\underline{\underline{L}} i^*(\underline{\underline{R}}i_*(\underline{M}))\) is the alternating sum}{\(\mathbb L i^*(\mathbb Ri_*(\underline{M}))\) is the alternating sum}{EN-CONTRIB-VII-028} of the classes of its homology sheaves, namely:
\[
i^*(i_*(x))=\sum_{i\ge0}(-1)^i
\operatorname{cl}\bigl(\underline{M}\otimes_{\underline{O}_Y}\underline{\Tor}_i^{\underline{O}_X}(\underline{O}_Y,\underline{O}_Y)\bigr)
=\sum_{i\ge0}(-1)^i x.\lambda^i(N).
\]

\sgaStableTarget{sga6.e07.cor.2-8}\textbf{Corollary 2.8.} \uline{Under the hypotheses of} 2.7, \uline{let} \(x,y\in K^{\bullet}(Y)\); \uline{we have the relation}:
\sgaStableTarget{sga6.e07.eq.2-8-1}\begin{equation}\notag
i_*(xy\lambda_{-1}(N))=i_*(x).i_*(y).
\end{equation}

By 2.7, \(y\lambda_{-1}(N)=i^*(i_*(y))\).

Applying the projection formula, we therefore find:
\[
i_*(xy\lambda_{-1}(N))
=i_*(x.i^*(i_*(y)))
=i_*(x).i_*(y).
\]

\sgaStableTarget{sga6.e07.sec.3}\section*{3. Computation of \texorpdfstring{$K^{\bullet}$}{K} of a blown-up scheme}

\sgaStableTarget{sga6.e07.para.3-1}\textbf{3.1.} In this section, \(i:Y\to X\) denotes a regular closed immersion. Denote by \(X'\) the scheme obtained by blowing up \(Y\), and put \(Y'=Y\times_X X'\).
%%%%%%%%%%  scan idx 450  |  folio 437  %%%%%%%%%%
\SGAPrintPage{437}%
We know (EGA II 8.1.8) that \(Y'\) is a closed subscheme of \(X'\) \var{defined by an ideal of \(\underline{O}_X\), canonically}{defined by an ideal of \(\underline{O}_{X'}\), canonically}{EN-CONTRIB-VII-029} isomorphic to \(\underline{O}_{X'}(1)\); since the immersion \var{\(j:Y'\to X\)}{\(j:Y'\to X'\)}{EN-CONTRIB-VII-030} is defined by an invertible ideal, it is a regular closed immersion of codimension \(1\). Denote by \(\underline{J}\) the ideal of \(\underline{O}_X\) defining \(Y\), and by \(\underline{J}'\) the ideal of \(\underline{O}_{X'}\) defining \(Y'\); thus \(\underline{J}'\simeq \underline{J}.\underline{O}_{X'}\simeq\underline{O}_{X'}(1)\). Finally, denote by \(f:X'\to X\), and \(g:Y'\to Y\) the structural morphisms. We therefore have the commutative diagram:
\sgaStableTarget{sga6.e07.eq.3-1-1}\begin{equation}\tag{3.1.1}
\begin{tikzcd}[column sep=large,row sep=large]
Y' \arrow[r,"j"] \arrow[d,"g"'] & X' \arrow[d,"f"]\\
Y \arrow[r,"i"'] & X.
\end{tikzcd}
\end{equation}

By definition, \(X'=\Proj(\underline{S})\), where \(\underline{S}=\bigoplus_{n\ge0}\underline{J}^n\). Consequently, we have:
\[
Y'=\Proj\bigl(i^*(\underline{S})\bigr)
=\Proj\left(\left(\bigoplus_{n\ge0}\underline{J}^n\right)\otimes_{\underline{O}_X}(\underline{O}_X/\underline{J})\right)
=\Proj\left(\bigoplus_{n\ge0}\underline{J}^n/\underline{J}^{n+1}\right).
\]
Since \(i\) is regular, \(\bigoplus_{n\ge0}\underline{J}^n/\underline{J}^{n+1}\simeq \underline{\Sym}_{\underline{O}_Y}(\underline{J}/\underline{J}^2)\) (1.3). Denote by \(\underline{N}\) the conormal sheaf \(\underline{J}/\underline{J}^2\); thus:
% VII probe display variant EN-CONTRIB-VII-031
\SGACaptureAppState{EN-CONTRIB-VII-031}%
\ifnum\SGAReaderForm=1\relax\SGAAppMark{EN-CONTRIB-VII-031}\fi
\ifnum\SGAReaderForm=0\relax
\[
Y'=\underline{\underline{P}}(\underline{N}).
\]
\else
\[
Y'=\bfP(\underline{N}).
\]
\fi

On \(Y'\), we can write the exact sequence:
\sgaStableTarget{sga6.e07.eq.3-1-2}\begin{equation}\tag{3.1.2}
0\to \underline{F}\to g^*(\underline{N})\to\underline{O}_{Y'}(1)\to0.
\end{equation}

Observe that \(\underline{O}_{Y'}(1)=j^*(\underline{O}_{X'}(1))=\underline{J}'/\underline{J}'^2\); it is therefore the conormal sheaf of the immersion \(j\), and we shall denote it by \(\underline{L}\).

\sgaStableTarget{sga6.e07.lem.3-2}\textbf{Lemma 3.2.} \uline{With the hypotheses and notation of} 3.1, \uline{we have the following isomorphisms}:
\sgaStableTarget{sga6.e07.eq.3-2-i}\sgaStableTarget{sga6.e07.eq.3-2-ii}\[
\begin{array}{ll}
\text{i)}&
\forall i\ge0,\quad
\underline{\Tor}_i^{\underline{O}_X}(\underline{O}_{X'},\underline{O}_Y)
\simeq
\bigwedge^i\underline{\Tor}_1^{\underline{O}_X}(\underline{O}_{X'},\underline{O}_Y);\\[4pt]
\text{ii)}&
\underline{\Tor}_1^{\underline{O}_X}(\underline{O}_{X'},\underline{O}_Y)\simeq \underline{F}.
\end{array}
\]

%%%%%%%%%%  scan idx 451  |  folio 438  %%%%%%%%%%
\SGAPrintPage{438}%
i) Since \(\underline{O}_Y\) can locally be resolved by a Koszul complex, the graded algebra \(\bigoplus_{i\ge0}\underline{\Tor}_i^{\underline{O}_X}(\underline{O}_{X'},\underline{O}_Y)\) is alternating (2.2). There therefore exists one and only one algebra homomorphism:
\sgaStableTarget{sga6.e07.eq.3-2-1}\begin{equation}\tag{3.2.1}
\bigwedge \underline{\Tor}_1^{\underline{O}_X}(\underline{O}_{X'},\underline{O}_Y)
\longrightarrow
\bigoplus_{i\ge0}\underline{\Tor}_i^{\underline{O}_X}(\underline{O}_{X'},\underline{O}_Y),
\end{equation}
extending the inclusion of \(\underline{\Tor}_1^{\underline{O}_X}(\underline{O}_{X'},\underline{O}_Y)\) in
\(\bigoplus_{i\ge0}\underline{\Tor}_i^{\underline{O}_X}(\underline{O}_{X'},\underline{O}_Y)\). It suffices to check locally that it is an isomorphism.

We can therefore work on an affine open subset \(U=\Spec(A)\) of \(X\), with \(\underline{J}\) defined by an ideal \(J\) of \(A\), and choose a regular homomorphism \(u:E\to A\) defining \(J\), with \(E=(A)^n\). Let \(f_1,\ldots,f_d\) be the images of a basis of \(E\).

In \(\bigoplus_{n\ge0}J^n\), for \(m\le p\), we shall denote by \(x^{(m)}\) the element \(x\) of \(J^p\) considered as having degree \(m\). We can restrict to the open subset \(V\) of \(X'\) where, for example, the element \(f_1^{(1)}\) is invertible.

The \(\underline{\Tor}_i^{\underline{O}_X}(\underline{O}_{X'},\underline{O}_Y)\) are the homology sheaves of the complex \(K_\bullet(u)\otimes_{\underline{O}_X}\underline{O}_{X'}\). To study this complex on \(V\), make the following change of basis in \(E'=E\otimes_A\Gamma(V,\underline{O}_{X'})\): if \(e_1,\ldots,e_d\) is the given basis of \(E\), put:
\[
e'_1=e_1;\qquad \text{for }i\ne1,\quad
e'_i=e_i-(f_i^{(1)}/f_1^{(1)})e_1.
\]
Put \(A'=\Gamma(V,\underline{O}_{X'})\); the differential of the complex \(K_\bullet(u)\otimes_A A'\) is given by
\[
d(e'_1)=f_1^{(0)};\qquad \text{for }i\ne1,\quad d(e'_i)=0.
\]
Taking as a basis of \(\bigwedge^k(E')\) the elements
\(e'_{i_1}\wedge\cdots\wedge e'_{i_k}\) with \var{\(i_1<..<i_k\)}{\(i_1<\cdots<i_k\)}{EN-CONTRIB-VII-032}, we therefore find:
\[
\begin{array}{ll}
d(a.e'_{i_1}\wedge\cdots\wedge e'_{i_k})=0
&\text{if }i_1\ne1;\\[3pt]
d(a.e'_{i_1}\wedge\cdots\wedge e'_{i_k})=
af_1^{(0)}.e'_{i_2}\wedge\cdots\wedge e'_{i_k}
&\text{if }i_1=1.
\end{array}
\]

%%%%%%%%%%  scan idx 452  |  folio 439  %%%%%%%%%%
\SGAPrintPage{439}%
Denoting by \(E''\) the submodule of \(E'\) generated by \(e'_2,\ldots,e'_d\), we therefore see that the image of \(d:\bigwedge^kE'\longrightarrow\bigwedge^{k-1}E'\) is \(f_1^{(0)}\bigwedge^{k-1}E''\). Furthermore, \(f_1^{(0)}\) is regular in \(A'\), since on \(V\) the invertible ideal \(\underline{J}'\) is defined by the equation \(f_1^{(0)}\). The kernel of \(d:\bigwedge^kE'\longrightarrow\bigwedge^{k-1}E'\) is therefore \(\bigwedge^kE''\). Consequently, the \(k\)-th homology group of \(K_\bullet(u)\otimes_AA'\) is \(\bigwedge^kE''/f_1^{(0)}\bigwedge^kE''\); we thus obtain on \(V\) isomorphisms:
\sgaStableTarget{sga6.e07.eq.3-2-2}\begin{equation}\tag{3.2.2}
\underline{\Tor}_i^{\underline{O}_X}(\underline{O}_{X'},\underline{O}_Y)
\xrightarrow{\sim}
(\underline{O}_{X'}/f_1^{(0)}\underline{O}_{X'})\otimes_{\underline{O}_{X'}}
\bigwedge^i(\underline{O}_{X'})^{d-1}
\xrightarrow{\sim}
\bigwedge^i(\underline{O}_{Y'})^{d-1}.
\end{equation}
In particular, for \(i=1\) we find:
\[
\underline{\Tor}_1^{\underline{O}_X}(\underline{O}_{X'},\underline{O}_Y)
\xrightarrow{\sim}
(\underline{O}_{Y'})^{d-1}.
\]
Consequently, for every \(i\) we have:
\sgaStableTarget{sga6.e07.eq.3-2-3}\begin{equation}\tag{3.2.3}
\underline{\Tor}_i^{\underline{O}_X}(\underline{O}_{X'},\underline{O}_Y)
\xrightarrow{\sim}
\bigwedge^i\underline{\Tor}_1^{\underline{O}_X}(\underline{O}_{X'},\underline{O}_Y)
\end{equation}

By 2.2, the isomorphism:
\[
\bigwedge(\underline{O}_{Y'})^{d-1}
\xrightarrow{\sim}
\bigoplus_{i\ge0}\underline{\Tor}_i^{\underline{O}_X}(\underline{O}_{X'},\underline{O}_Y)
\]
defined by the isomorphisms (3.2.2) is an algebra isomorphism. Hence the same is true of the isomorphism:
\[
\bigwedge\underline{\Tor}_1^{\underline{O}_X}(\underline{O}_{X'},\underline{O}_Y)
\xrightarrow{\sim}
\bigoplus_{i\ge0}\underline{\Tor}_i^{\underline{O}_X}(\underline{O}_{X'},\underline{O}_Y)
\]
defined by the isomorphisms (3.2.3). By uniqueness, it is then the restriction to \(V\) of the homomorphism (3.2.1), and the latter is therefore an isomorphism.

ii) Consider the exact sequence:
\[
0\longrightarrow \underline{J}\longrightarrow\underline{O}_X\longrightarrow\underline{O}_Y\longrightarrow0,
\]

%%%%%%%%%%  scan idx 453  |  folio 440  %%%%%%%%%%
\SGAPrintPage{440}%
and the exact \(\underline{\Tor}\) sequence obtained by tensoring with \(\underline{O}_{X'}\):
\[
0\longrightarrow
\Tor_1^{\underline{O}_X}(\underline{O}_{X'},\underline{O}_Y)
\longrightarrow
\underline{J}\otimes_{\underline{O}_X}\underline{O}_{X'}
\longrightarrow
\underline{O}_{X'}
\longrightarrow
\underline{O}_{Y'}
\longrightarrow0.
\]
The image of the morphism: \(\underline{J}\otimes_{\underline{O}_X}\underline{O}_{X'}\to\underline{O}_{X'}\) is
\(\underline{J}.\underline{O}_{X'}=\underline{O}_{X'}(1)\). Consequently, we obtain the exact sequence:
\[
0\longrightarrow
\underline{\Tor}_1^{\underline{O}_X}(\underline{O}_{X'},\underline{O}_Y)
\longrightarrow
\underline{J}\otimes_{\underline{O}_X}\underline{O}_{X'}
\longrightarrow
\underline{O}_{X'}(1)
\longrightarrow0.
\]

Then tensor it with \(\underline{O}_{Y'}\); since
\(\underline{\Tor}_1^{\underline{O}_X}(\underline{O}_{X'},\underline{O}_Y)\) is an \(\underline{O}_{Y'}\)-Module, and
\(\underline{O}_{X'}(1)\) is invertible, we obtain the exact sequence:
\sgaStableTarget{sga6.e07.eq.3-2-4}\[
0\longrightarrow
\underline{\Tor}_1^{\underline{O}_X}(\underline{O}_{X'},\underline{O}_Y)
\longrightarrow
g^*(\underline{N})
\longrightarrow
\underline{O}_{Y'}(1)
\longrightarrow0.
\]
One readily checks that the morphism \(g^*(\underline{N})\to\underline{O}_{Y'}(1)\) occurring in this sequence is the canonical morphism, so that this sequence is identified with (3.1.2). By definition, we therefore have:
\[
\underline{\Tor}_1^{\underline{O}_X}(\underline{O}_{X'},\underline{O}_Y)=\underline{F}.
\]

\sgaStableTarget{sga6.e07.para.3-3}\textbf{3.3.} We now suppose that \(X\) is a scheme having an ample sheaf. Then the same is true of \(X'\), which is projective over \(X\), and of \(Y\) and \(Y'\). Furthermore, \(i\) and \(j\) are perfect morphisms, since they are regular immersions; \(g\) is smooth and hence perfect, since \var{\(Y'=\underline{\underline{P}}(\underline{N})\)}{\(Y'=\bfP(\underline{N})\)}{EN-CONTRIB-VII-033}; finally, \(f\) is perfect by 1.9. Since these morphisms are projective, we can define the additive homomorphisms (IV
\[
i_*:K^{\bullet}(Y)\to K^{\bullet}(X);\quad
j_*:K^{\bullet}(Y')\to K^{\bullet}(X');\quad
f_*:K^{\bullet}(X')\to K^{\bullet}(X);\quad
g_*:K^{\bullet}(Y')\to K^{\bullet}(Y).
\]

%%%%%%%%%%  scan idx 454  |  folio 441  %%%%%%%%%%
\SGAPrintPage{441}%
\sgaStableTarget{sga6.e07.prop.3-4}\textbf{Proposition 3.4.} \uline{With the hypotheses and notation of} 3.1 \uline{and} 3.3, \uline{for every} \(x\in K^{\bullet}(Y)\) \uline{we have the relation}:
\sgaStableTarget{sga6.e07.eq.3-4-1}\[
f^*(i_*(x))=j_*\bigl(g^*(x).\lambda_{-1}(F)\bigr).
\]

Since both sides of the relation to be proved are additive, we may suppose that \(x\) is the class of an \(\underline{O}_Y\)-Module locally free of finite type \(\underline{M}\). The element \(f^*(i_*(x))\) is then \var{the class of the complex \(\underline{\underline{L}}f^*(\underline{\underline{R}}i_*(\underline{M}))\)}{the class of the complex \(\mathbb Lf^*(\mathbb Ri_*(\underline{M}))\)}{EN-CONTRIB-VII-034}, obtained by taking a projective resolution of \(i_*(\underline{M})\) and applying the functor \var{\(f_*\)}{\(f^*\)}{EN-CONTRIB-VII-035} to it. The homology sheaves of the resulting complex are the \(\underline{\Tor}_i^{\underline{O}_X}(\underline{M},\underline{O}_{X'})\). By 2.4 and 3.2, we have:
\[
\underline{\Tor}_i^{\underline{O}_X}(\underline{M},\underline{O}_{X'})
\xrightarrow{\sim}
\underline{M}\otimes_{\underline{O}_Y}\underline{\Tor}_i^{\underline{O}_X}(\underline{O}_{X'},\underline{O}_Y)
\xrightarrow{\sim}
\underline{M}\otimes_{\underline{O}_Y}\bigwedge^i\underline{F}.
\]

It follows that these homology sheaves are locally free of finite type over \(\underline{O}_{Y'}\), hence perfect over \(\underline{O}_{X'}\); consequently, the class in \(K^{\bullet}(X')\) of the complex \var{\(\underline{\underline{L}}f^*(\underline{\underline{R}}i_*(\underline{M}))\) is the alternating sum}{\(\mathbb Lf^*(\mathbb Ri_*(\underline{M}))\) is the alternating sum}{EN-CONTRIB-VII-036} of the classes of its homology sheaves, namely:
\[
f^*(i_*(x))
=\sum_{i\ge0}(-1)^i\operatorname{cl}
\bigl(j_*(\underline{M}\otimes_{\underline{O}_Y}\bigwedge^i\underline{F})\bigr)
=\sum_{i\ge0}(-1)^i
j_*\bigl(g^*(x)\lambda^i(F)\bigr).
\]

\sgaStableTarget{sga6.e07.lem.3-5}\textbf{Lemma 3.5.} \uline{With the hypotheses and notation of} 3.1, \uline{we have the following isomorphisms}:
\begin{enumerate}[label=\roman*)]
\item \uline{for every} \(n\ge0\), \uline{and every} \(i\ge1\), \(R^if_*\bigl(\underline{O}_{X'}(n)\bigr)=0\);
\item \uline{for every} \(n\ge0\), \(f_*\bigl(\underline{O}_{X'}(n)\bigr)=\underline{J}^n\).
\end{enumerate}

We may clearly suppose \(X\) affine, hence quasi-compact.

i) First recall the following lemma, \var{which was proved in exp.}{which was proved in exp. III.}{EN-CONTRIB-VII-037}
%%%%%%%%%%  scan idx 455  |  folio 442  %%%%%%%%%%
\SGAPrintPage{442}%
\textbf{Lemma.} \uline{Let} \(X'\xrightarrow{f}X\) \uline{be a projective and perfect morphism}, \uline{with} \(X\) \uline{quasi-compact}, \uline{and let} \(\underline{L}^\bullet\) \uline{be a perfect complex over} \(\underline{O}_{X'}\), \uline{acyclic in degrees} \(>0\). \uline{Then there exists an integer} \(n_0\) \uline{such that for} \(n\ge n_0\) \uline{the complex} \var{\(\underline{\underline{R}}f_*\bigl(\underline{L}^\bullet\otimes_{\underline{O}_{X'}}\underline{O}_{X'}(n)\bigr)\)}{\(\mathbb Rf_*\bigl(\underline{L}^\bullet\otimes_{\underline{O}_{X'}}\underline{O}_{X'}(n)\bigr)\)}{EN-CONTRIB-VII-038} \uline{is acyclic in degrees} \(>0\).

Apply this lemma to the complex reduced to \(\underline{O}_{X'}\); there is therefore an integer \(n_0\) such that for \(n\ge n_0\), we have \(R^if_*(\underline{O}_{X'}(n))=0\) when \(i\ge1\). We then descend by induction to \(n=0\). For this, consider the exact sequence:
\[
0\longrightarrow
\underline{O}_{X'}(n+1)
\longrightarrow
\underline{O}_{X'}(n)
\longrightarrow
\underline{O}_{Y'}(n)
\longrightarrow0,
\]
and apply the functor \(f_*\) to it. We have \var{\(\underline{\underline{R}}f_*(\underline{O}_{Y'}(n))=i_*\circ\underline{\underline{R}}g_*(\underline{O}_{Y'}(n))\)}{\(\mathbb Rf_*(\underline{O}_{Y'}(n))=i_*\circ\mathbb Rg_*(\underline{O}_{Y'}(n))\)}{EN-CONTRIB-VII-039}. Since \var{\(Y'=\underline{\underline{P}}(\underline{N})\)}{\(Y'=\bfP(\underline{N})\)}{EN-CONTRIB-VII-040}, \(\underline{O}_{Y'}(n)\) is \(g_*\)-acyclic for \(n\ge0\). Thus for every \(n\ge0\), and every \(i\ge1\):
\[
R^if_*(\underline{O}_{Y'}(n))=0,
\]
and consequently we have the exact sequences:
\[
R^if_*(\underline{O}_{X'}(n+1))
\longrightarrow
R^if_*(\underline{O}_{X'}(n))
\longrightarrow0,
\]
which complete the induction.

ii) We must show that the canonical morphism:
\[
\bigoplus_{n\ge0}\underline{J}^n
\longrightarrow
\bigoplus_{n\ge0}f_*\bigl(\underline{O}_{X'}(n)\bigr)
\]
is an isomorphism. Since the property is local, we may suppose that \(\underline{J}\) is defined by a sequence \(f_1,\ldots,f_d\) of sections of \(\underline{O}_X\) such that the corresponding homomorphism \((\underline{O}_X)^d\longrightarrow\underline{O}_X\) is regular. Put: \(\underline{S}=\bigoplus_{n\ge0}\underline{J}^n\); we define a surjective algebra homomorphism: \(\underline{O}_X[T_1,\ldots,T_d]\xrightarrow{v}\underline{S}\), by sending

%%%%%%%%%%  scan idx 456  |  folio 443  %%%%%%%%%%
\SGAPrintPage{443}%
\(T_i\) to \(f_i^{(1)}\). This homomorphism defines a closed immersion \var{\(X'\longrightarrow\underline{\underline{P}}_X^{d-1}\)}{\(X'\longrightarrow\bfP_X^{d-1}\)}{EN-CONTRIB-VII-041}, regular by 1.8 ii).

By 1.8 iii), \(\underline{S}\) is perfect as a graded \(\underline{O}_X[T_1,\ldots,T_d]\)-Module. In particular, \(\underline{S}\) is finitely presented as a graded Module. Let therefore:
\[
\underline{L}''\longrightarrow \underline{L}'\longrightarrow \underline{S}\longrightarrow0
\]
be a finite graded presentation of \(\underline{S}\) over \(\underline{O}_X[T_1,\ldots,T_d]\). \var{On \(\underline{\underline{P}}_X^{d-1}\) we obtain}{On \(\bfP_X^{d-1}\) we obtain}{EN-CONTRIB-VII-042} the exact sequence:
\[
\underline{M}''\longrightarrow \underline{M}'\longrightarrow\underline{O}_{X'}\longrightarrow0,
\]
where \(\underline{M}''=\underline{\Proj}(\underline{L}'')\), \(\underline{M}'=\underline{\Proj}(\underline{L}')\), \var{so these are free of finite type}{so these are locally free of finite type}{EN-CONTRIB-VII-043} \var{over \(\underline{\underline{P}}_X^{d-1}\). We can therefore}{over \(\bfP_X^{d-1}\). We can therefore}{EN-CONTRIB-VII-044} write, for every \(n\), the commutative diagram:
\[
\begin{tikzcd}[column sep=large,row sep=large]
\underline{L}''_n \arrow[r] \arrow[d,"\sim"'] &
\underline{L}'_n \arrow[r] \arrow[d,"\sim"'] &
\underline{J}^n \arrow[r] \arrow[d] &
0\\
p_*(\underline{M}''(n)) \arrow[r] &
p_*(\underline{M}'(n)) \arrow[r] &
p_*(\underline{O}_{X'}(n)) \arrow[r] &
0,
\end{tikzcd}
\]
where \(p\) is the structural morphism: \var{\(\underline{\underline{P}}_X^{d-1}\to X\)}{\(\bfP_X^{d-1}\to X\)}{EN-CONTRIB-VII-045}, and \(\underline{L}''_n\), \(\underline{L}'_n\) are the degree \(n\) components of \(\underline{L}''\), \(\underline{L}'\). The upper row is exact by hypothesis; the first two vertical arrows are isomorphisms by the classical results for the direct images of \(\underline{O}_P(n)\), \(P\) being a projective bundle (EGA III, 2.1.15); finally, the last row is exact when \(n\) is sufficiently large, since the kernel of \(\underline{M}'\to\underline{O}_{X'}\) is a perfect \(\underline{O}_P\)-Module, and we can apply the lemma recalled above, which shows that this kernel, tensored with \(\underline{O}_P(n)\) for \(n\) sufficiently large, is \(p_*\)-acyclic; likewise for the kernel of \(\underline{M}''\to \underline{M}'\). It follows that for \(n\) sufficiently large, \(\underline{J}^n\longrightarrow p_*(\underline{O}_{X'}(n))\) is an isomorphism. Since \(p_*(\underline{O}_{X'}(n))=f_*(\underline{O}_{X'}(n))\), we therefore have

%%%%%%%%%%  scan idx 457  |  folio 444  %%%%%%%%%%
\SGAPrintPage{444}%
the result for \(n\ge n_0\).

It remains to descend to \(n=0\). We proceed by descending induction, using the commutative diagram:
\[
\begin{tikzcd}[column sep=large,row sep=large]
0 \arrow[r] \arrow[d] &
\underline{J}^{n+1} \arrow[r] \arrow[d,"\sim"'] &
\underline{J}^n \arrow[r] \arrow[d] &
\underline{J}^n/\underline{J}^{n+1} \arrow[r] \arrow[d,"\sim"'] &
0 \arrow[d]\\
0 \arrow[r] &
f_*(\underline{O}_{X'}(n+1)) \arrow[r] &
f_*(\underline{O}_{X'}(n)) \arrow[r] &
f_*(\underline{O}_{Y'}(n)) \arrow[r] &
0.
\end{tikzcd}
\]
The bottom row is exact because \(R^1f_*(\underline{O}_{X'}(n+1))=0\) for every \(n\ge0\), by i), already proved. Furthermore, we have:
% VII probe display variant EN-CONTRIB-VII-046
\SGACaptureAppState{EN-CONTRIB-VII-046}%
\ifnum\SGAReaderForm=1\relax\SGAAppMark{EN-CONTRIB-VII-046}\fi
\ifnum\SGAReaderForm=0\relax
\[
f_*(\underline{O}_{Y'}(n))
=i_*\bigl(g_*(\underline{O}_{Y'}(n))\bigr)
=i_*(\underline{\underline{S}}(\underline{N})_n)
=\underline{J}^n/\underline{J}^{n+1}.
\]
\else
\[
f_*(\underline{O}_{Y'}(n))
=i_*\bigl(g_*(\underline{O}_{Y'}(n))\bigr)
=i_*(\mathbf S(\underline{N})_n)
=\underline{J}^n/\underline{J}^{n+1}.
\]
\fi
The arrow on the right is therefore an isomorphism; the one on the left is an isomorphism by the induction hypothesis; \var{hence the middle one is also an isomorphism}{hence the middle one is also an isomorphism.}{EN-CONTRIB-VII-047}

\sgaStableTarget{sga6.e07.prop.3-6}\textbf{Proposition 3.6.} \uline{With the hypotheses and notation of} 3.1 \uline{and} 3.3, \uline{we have the relation}:
\sgaStableTarget{sga6.e07.eq.3-6-1}\[
f_*\circ f^*=\Id_{K^{\bullet}(X)}.
\]

Let \(x\in K^{\bullet}(X)\); by the projection formula, we can write:
\[
f_*(f^*(x))=f_*(1.f^*(x))=x.f_*(1).
\]
It therefore suffices to show that \(f_*(1)=1\).

Now 3.5 shows that the cohomology of \(\mathbb Rf_*(\underline{O}_{X'})\) vanishes in degrees \(>0\), and equals \(\underline{O}_X\) in degree \(0\). The proposition follows at once.

\sgaStableTarget{sga6.e07.thm.3-7}\textbf{Theorem 3.7.} \uline{Let} \(i:Y\to X\) \uline{be a regular closed immersion}, \(X'\) \uline{the scheme obtained by blowing up} \(Y\), \uline{and} \(Y'=X'\times_XY\), \uline{so that we have the commutative diagram}

%%%%%%%%%%  scan idx 458  |  folio 445  %%%%%%%%%%
\SGAPrintPage{445}%
\[
\begin{tikzcd}[column sep=large,row sep=large]
Y' \arrow[r,"j"] \arrow[d,"g"'] & X' \arrow[d,"f"]\\
Y \arrow[r,"i"] & X.
\end{tikzcd}
\]
\uline{Suppose that} \(X\) \uline{is quasi-compact and has an ample sheaf}. \uline{Then}, \uline{we have the exact sequence}:
\sgaStableTarget{sga6.e07.eq.3-7-0}\[
0\longrightarrow K^{\bullet}(Y)\xrightarrow{u}K^{\bullet}(Y')\times K^{\bullet}(X)
\xrightarrow{v}K^{\bullet}(X'),
\]
\uline{where we put}, \uline{for} \(y\in K^{\bullet}(Y)\):
\sgaStableTarget{sga6.e07.eq.3-7-u}\[
u(y)=\bigl(g^*(y).\lambda_{-1}(F),-i_*(y)\bigr),
\]
\uline{and for} \(y'\in K^{\bullet}(Y')\), \(x\in K^{\bullet}(X)\):
\sgaStableTarget{sga6.e07.eq.3-7-v}\[
v(y',x)=j_*(y')+f^*(x),
\]
\uline{with the notation of} 3.1 \uline{and} 3.3. \uline{Moreover}, \(u\) \uline{has a left inverse} \(u'\), \uline{defined by} \(u'(y',x)=g_*(y')\).

\uline{If in addition} \(X\) \uline{is noetherian}, \uline{and} \(Y\) \uline{regular}, \uline{then} \(v\) \uline{is surjective}, \uline{so we have the exact sequence}:
\sgaStableTarget{sga6.e07.eq.3-7-2}\[
0\longrightarrow K^{\bullet}(Y)\xrightarrow{u}K^{\bullet}(Y')\times K^{\bullet}(X)
\xrightarrow{v}K^{\bullet}(X')\longrightarrow0.
\]

First, \(u\) has the left inverse \(u'\), and is a fortiori injective. Indeed, we have:
\[
g_*\bigl(g^*(y).\lambda_{-1}(F)\bigr)
=y.g_*(\lambda_{-1}(F))
\]
by the projection formula. Furthermore, \(g_*(\lambda_{-1}(F))=1\) \var{(VI 5.9) which gives the result.}{(VI 5.9), which gives the result.}{EN-CONTRIB-VII-048}

Next, \(v\circ u=0\). Indeed, for \(y\in K^{\bullet}(Y)\) we have:
% VII probe display variant EN-CONTRIB-VII-049
\SGACaptureAppState{EN-CONTRIB-VII-049}%
\ifnum\SGAReaderForm=1\relax\SGAAppMark{EN-CONTRIB-VII-049}\fi
\ifnum\SGAReaderForm=0\relax
\[
v\circ u(y)
=v\bigl(g^*(y)\lambda_{-1}(F),-i_*(y)\bigr)
=j_*\bigl(g^*(y)\lambda_{-1}(F)\bigr)-f^*(i_*(x))=0,
\]
\else
\[
v\circ u(y)
=v\bigl(g^*(y)\lambda_{-1}(F),-i_*(y)\bigr)
=j_*\bigl(g^*(y)\lambda_{-1}(F)\bigr)-f^*(i_*(y))=0,
\]
\fi
by 3.4.

It therefore remains to show that \(\Ker(v)\subset \Im(u)\). Let therefore \(x\in K^{\bullet}(X)\), \(y'\in K^{\bullet}(Y')\), be such that:

%%%%%%%%%%  scan idx 459  |  folio 446  %%%%%%%%%%
\SGAPrintPage{446}%
\[
j_*(y')+f^*(x)=0.
\]
We must find \(y\in K^{\bullet}(Y)\) such that:
\sgaStableTarget{sga6.e07.eq.3-7-1}\begin{equation}\tag{3.7.1}
x=-i_*(y);\qquad y'=g^*(y).\lambda_{-1}(F).
\end{equation}
Since \(u\) has a left inverse \(u'\), we necessarily have
\[
y=u'(y',x)=g_*(y').
\]
Thus we must show that \(y\) satisfies the relations (3.7.1).

To show that \(x=-i_*(y)\), start with the relation \(f^*(x)=-j_*(y')\) and apply \(f_*\) to it. We therefore obtain: \(f_*(f^*(x))=-f_*(j_*(y'))=-i_*(g_*(y'))=-i_*(y)\) by the definition of \(y\). By Proposition 3.6, however, \(f_*(f^*(x))=x\); hence the required relation.

To complete the proof, we must show that
\[
y'=g^*(g_*(y'))\lambda_{-1}(F).
\]
Put \(z=y'-g^*(g_*(y'))\lambda_{-1}(F)\). We obtain: \(g_*(z)=g_*(y')-g_*(y')\,g_*(\lambda_{-1}(F))=0\). On the other hand, by hypothesis we have: \var{\(j_*(y')=f^*(x)\)}{\(j_*(y')=-f^*(x)\)}{EN-CONTRIB-VII-050}, whence, by the preceding calculation, \(j_*(y')=f^*(i_*(y))\), and by Proposition 3.4: \(j_*(y')=j_*\bigl(g^*(y)\lambda_{-1}(F)\bigr)\); consequently, we also have \(j_*(z)=0\). Now \(\Ker(j_*)\cap\Ker(g_*)=0\). Indeed, if \(j_*(z)=0\), then \(j^*(j_*(z))=0\); but \(j^*(j_*(z))=z(1-L)\) (Prop. 2.7), and for \(z(1-L)=0\) to hold there must exist \(z'\in K^{\bullet}(Y)\) such that \(z=z'\lambda_{-1}(F)\) (VI,5.10); applying \(g_*\) gives \(z'=g_*(z)=0\), hence \(z=0\).

Suppose now that \(X\) is noetherian and \(Y\) regular. Then \(Y'\) is also regular, since it is a projective bundle over \(Y\), and for every point \(y'\) of \(Y'\), the local ring \(\underline{O}_{X',y'}\) of \(X'\) at \(y'\) is regular, since \(\underline{O}_{Y',y'}\) is the quotient of \(\underline{O}_{X',y'}\) by a regular element (see EGA \(0_{\mathrm{IV}}\) 17.1.8).
%%%%%%%%%%  scan idx 460  |  folio 447  %%%%%%%%%%
\SGAPrintPage{447}%
We know that for every regular scheme having an ample sheaf, the Grothendieck groups formed with locally free sheaves of finite type and with coherent sheaves are identical. In particular, this is the case for \(Y'\).

Let \(\underline{E}\) be a coherent \(\underline{O}_{X'}\)-Module supported on \(Y'\); then \(\underline{E}\) is a perfect \(\underline{O}_{X'}\)-Module. Indeed, it is perfect at the points of \(Y'\), since \(X'\) is regular at these points, and at the points of \(X'-Y'\) since it vanishes there.

If \(E\) is the class of \(\underline{E}\) in \(K^{\bullet}(X')\), then \(E\) belongs to the image of \(j_*\). Indeed, we can filter \(\underline{E}\) by the \(\underline{J}'^n.\underline{E}\). This filtration is finite, since, \(X'\) being noetherian, there is an integer \(n_0\) such that \(\underline{J}'^{\,n_0}.\underline{E}=0\); the quotients of this filtration are then \(\underline{O}_{Y'}\)-Modules, so their classes in \(K^{\bullet}(X')\) belong to the image of \(j_*\). Since \(E\) is the sum of the classes of these quotients, \(E\) also belongs to this image.

Let now \(x\in K^{\bullet}(X')\). We can find a perfect complex \(\underline{M}^\bullet\) whose class is \(x\). Consider the distinguished triangle:
% VII probe display variant EN-CONTRIB-VII-051
\SGACaptureAppState{EN-CONTRIB-VII-051}%
\ifnum\SGAReaderForm=1\relax\SGAAppMark{EN-CONTRIB-VII-051}\fi
\ifnum\SGAReaderForm=0\relax
\[
\begin{tikzcd}[column sep=large,row sep=large]
& \underline{N}^\bullet \arrow[dl] &\\
\underline{M}^\bullet \arrow[rr] && \underline{\underline{L}}f^*(\underline{\underline{R}}f_*(\underline{M}^\bullet)) \arrow[ul] ,
\end{tikzcd}
\]
\else
\[
\begin{tikzcd}[column sep=large,row sep=large]
& \underline{N}^\bullet \arrow[dr] &\\
\underline{M}^\bullet \arrow[ur] && \mathbb Lf^*(\mathbb Rf_*(\underline{M}^\bullet)) \arrow[ll] ,
\end{tikzcd}
\]
\fi
where the horizontal arrow is the canonical homomorphism. This arrow is an isomorphism on the complement of \(Y'\), and consequently the cohomology sheaves of \(\underline{N}^\bullet\) are coherent sheaves supported on \(Y'\). By the preceding argument, these are perfect \(\underline{O}_{X'}\)-Modules, and their classes belong to the image of \(j_*\). Thus we have:
\[
x-f^*(f_*(x))\in\operatorname{Im}(j_*)\,;
\]
and consequently the homomorphism \(v\) is surjective.



% Macros for Exposé VII continuation
\providecommand{\Q}{\mathbb Q}
\providecommand{\bfP}{\mathbf P}
\providecommand{\gr}{\operatorname{gr}}
\providecommand{\Id}{\operatorname{Id}}
\providecommand{\Fil}{\operatorname{Fil}}
\providecommand{\Imf}{\operatorname{Im}}
\providecommand{\rank}{\operatorname{rank}}
\providecommand{\codim}{\operatorname{codim}}
\providecommand{\cJ}{\mathcal J}
\providecommand{\cI}{\mathcal I}
\providecommand{\cN}{\mathcal N}
\providecommand{\cE}{\mathcal E}
\providecommand{\cO}{\mathcal O}



%%%%%%%%%%  scan idx 461  |  folio 448  %%%%%%%%%%
\SGAPrintPage{448}%
\sgaStableTarget{sga6.e07.remark-3-8}\subsection*{Remark 3.8.}

When \(X\) is noetherian and \(Y\) regular, 3.7 determines the abelian group \(K^{\bullet}(X')\) entirely from the \(\lambda\)-rings \(K^{\bullet}(X)\), \(K^{\bullet}(Y)\), the maps \(i_*,i^*\) between these rings, and the element \(N\) of \(K^{\bullet}(Y)\): indeed, \(K^{\bullet}(Y')\) can be computed from \(K^{\bullet}(Y)\) and \(N\) by (VI 1.1). We shall see in the next section that the multiplicative structure of \(K^{\bullet}(X')\) is determined by the same data, as is the \(\lambda\)-structure after tensoring with \(\mathbb Q\) (this restriction is probably unnecessary: see 4.5).

\sgaStableTarget{sga6.e07.sec.4}\section*{4. Regular immersions and the filtration of \(K^{\bullet}\)}

Throughout this section we keep the notation and hypotheses of 3.1.

\sgaStableTarget{sga6.e07.lem.4-1}\textbf{Lemma 4.1.} \uline{Suppose that there exists on} \(Y\) \uline{a locally free sheaf of finite type} \(\underline{N}'\), \uline{such that}, \uline{if} \(N,N'\) \uline{are the classes in} \(K^{\bullet}(Y)\) \uline{of} \(\underline{N},\underline{N}'\), \uline{we have} \(N=N'+2\). \uline{Then the following congruence holds}:
\[
  \lambda_{-1}(F)\equiv 0 \ \mathrm{mod.}(1-L).
\]

Let \(d\) be the rank of \(\underline{N}\) (which we may suppose constant); the rank of \(\underline{F}\) is then \(d-1\). By definition:
\[
\begin{aligned}
\lambda_{-1}(F)
 &= \sum_{i\ge0}(-1)^i\lambda^i(F)\\
 &= (-1)^{d-1}\lambda^{d-1}(F-1)\\
 &= (-1)^{d-1}\lambda^{d-1}(F+L-1-L)\\
 &= (-1)^{d-1}\lambda^{d-1}(N-2+1-L).
\end{aligned}
\]
Now \(\lambda_t(1-L)=\lambda_t(1)/\lambda_t(L)=(1+t)/(1+Lt)\). Thus for every \(i\ge1\), \(\lambda^i(1-L)\) is divisible by \(1-L\). Consequently:
\[
\begin{aligned}
\lambda_{-1}(F)
  &\equiv (-1)^{d-1}\lambda^{d-1}(N-2) \qquad\mathrm{mod.}(1-L)\\
  &\equiv(-1)^{d-1}\lambda^{d-1}(N') \qquad\mathrm{mod.}(1-L).
\end{aligned}
\]

%%%%%%%%%%  scan idx 462  |  folio 449  %%%%%%%%%%
\SGAPrintPage{449}%
Now \(\underline{N}'\) has rank \(d-2\), so \(\lambda^{d-1}(N')=0\), which gives the result.

\sgaStableTarget{sga6.e07.sec.4.2}\subsection*{4.2.}

Let \(i:Y\to X\) be a regular immersion. Let \(r\) be an integer \(\ge1\), and \var{\(s:X\longrightarrow \underline{\underline{P}}_X^r\)}{\(s:X\longrightarrow \mathbf P_X^r\)}{EN-CONTRIB-VII-052} the section defined by the homomorphism: \((\underline{O}_X)^{r+1}\to\underline{O}_X\), passage to the quotient by the subsheaf generated by the last \(r\) coordinates. If \var{\(\underline{\underline{P}}_X^r\)}{\(\mathbf P_X^r\)}{EN-CONTRIB-VII-053} is defined by the graded \(\underline{O}_X\)-Algebra \(\underline{O}_X[T_0,\ldots,T_r]\), the ideal \(\underline{I}\) defining the section \(s\) is the ideal generated by the sections \(T_1,\ldots,T_r\), so \(\underline{I}/\underline{I}^2\) is isomorphic to \((\underline{O}_X[T_0])^r\), and hence the corresponding ideal of \var{\(\underline{O}_p\)}{\(\underline{O}_P\)}{EN-CONTRIB-VII-054} is isomorphic to \((\underline{O}_X)^r\).

Since the section \(s\) is clearly a regular immersion, the same holds for \(s\circ i\), and we can write the corresponding exact sequence of conormal sheaves (1.7). We thus obtain the following equality in \(K^{\bullet}(Y)\), writing \(\underline{N}\) for the conormal sheaf of \(i\), and \(\underline{N}_1\) for the conormal sheaf of \(s\circ i\):
\[
  N_1=N+r.
\]
Thus for every regular immersion \(i\), composing it with a section: \var{\(X\to\underline{\underline{P}}_X^2\)}{\(X\to\mathbf P_X^2\)}{EN-CONTRIB-VII-055}, or: \var{\(X\to\underline{\underline{P}}_{\underline{\underline{P}}_X^1}^{1}\)}{\(X\to\mathbf P_{\mathbf P_X^1}^{1}\)}{EN-CONTRIB-VII-056}, gives a regular immersion satisfying the property of Lemma 4.1.

\sgaStableTarget{sga6.e07.thm.4-3}\textbf{Theorem 4.3.} \uline{Let} \(i:Y\to X\) \uline{be a regular closed immersion}, \(N\) \uline{the class in} \(K^{\bullet}(Y)\) \uline{of the conormal sheaf of} \(Y\) \uline{in} \(X\). \uline{Suppose that} \(X\) \uline{is quasi-compact and has an ample sheaf}. \uline{Then}, \uline{with the notation of} 3.3, \uline{we have for every} \(y\in K^{\bullet}(Y)\), \uline{and every} \(p\ge1\):
\sgaStableTarget{sga6.e07.eq.4-3-1}\begin{equation}\tag{4.3.1}
 i_*\bigl(\lambda^p(N,y)\bigr)\equiv \lambda^p(i_*(y))
 \qquad\text{\uline{mod}. \uline{torsion}},
\end{equation}
\(\lambda^p(N,y)\) \uline{being the element defined in} (V 5.3).

i) First suppose that the hypothesis of 4.1 is satisfied; in this case we shall show that \(i_*\bigl(\lambda^p(N,y)\bigr)=\lambda^p(i_*(y))\). Since \(f_*\circ f^*=\Id\), this

%%%%%%%%%%  scan idx 463  |  folio 450  %%%%%%%%%%
\SGAPrintPage{450}%
relation is equivalent to the relation:
\[
  f^*\bigl(i_*(\lambda^p(N,y))\bigr)=
  f^*\bigl(\lambda^p(i_*(y))\bigr).
\]
Now we have:
\sgaStableTarget{sga6.e07.eq.3-4}\[
\begin{aligned}
 f^*\bigl(\lambda^p(i_*(y))\bigr)
 &=\lambda^p\bigl(f^*(i_*(y))\bigr)\\
 &=\lambda^p\bigl(j_*(g^*(y)\lambda_{-1}(F))\bigr)
\qquad(3.4)
\end{aligned}
\]
On the other hand:
\sgaStableTarget{sga6.e07.eq.3-4.occ2}\[
  f^*\bigl(i_*(\lambda^p(N,y))\bigr)
  =
  j_*\bigl(g^*(\lambda^p(N,y))\lambda_{-1}(F)\bigr)\ (3.4).
\]
Since \(g^*(N)=F+L\), we can apply formula (V (5.4.2)):
\[
  g^*(\lambda^p(N,y))\lambda_{-1}(F)
  =
  \lambda^p(g^*(N),g^*(y))\lambda_{-1}(F)
  =
  \lambda^p(L,g^*(y)\lambda_{-1}(F)).
\]
Putting: \(z=g^*(y)\lambda_{-1}(F)\), we are therefore reduced to proving the relation:
\[
  j_*\bigl(\lambda^p(L,z)\bigr)=\lambda^p(j_*(z)),
\]
that is, the analogous relation for the immersion \(j\) and the element \(z\). By hypothesis, \(\lambda_{-1}(F)\) is a multiple of \((1-L)\); we can find \(z'\) such that \(z=z'(1-L)\). Then \(z=j^*(u')\) with \(u'=j_*(z')\) (2.7); since \(L=j^*(J')\), we finally have to verify:
\[
  j_*\bigl(j^*(\lambda^p(J',u'))\bigr)=
  \lambda^p\bigl(j_*(j^*(u'))\bigr).
\]
The exact sequence
\[
  0\longrightarrow \underline{J}'\longrightarrow \underline{O}_{X'}
  \longrightarrow \underline{O}_{Y'}\longrightarrow 0
\]
gives the relation:
\sgaStableTarget{sga6.e07.eq.4-3-2}\begin{equation}\tag{4.3.2}
  j_*(1)=1-J'.
\end{equation}
Thus, applying the projection formula, the relation to be verified finally becomes:

%%%%%%%%%%  scan idx 464  |  folio 451  %%%%%%%%%%
\SGAPrintPage{451}%
\[
  \lambda^p(J',u')(1-J')=\lambda^p(u'(1-J')).
\]
Since \(1-J'=\lambda_{-1}(J')\), this is the formula defining \(\lambda^p(J',u')\).

ii) We turn to the general case. Compose \(i\) with a section: \var{\(X\to\underline{\underline{P}}_X^1\)}{\(X\to\mathbf P_X^1\)}{EN-CONTRIB-VII-057}, then with a section: \var{\(\underline{\underline{P}}_X^1\to\underline{\underline{P}}_{\underline{\underline{P}}_X^1}^1\)}{\(\mathbf P_X^1\to\mathbf P_{\mathbf P_X^1}^1\)}{EN-CONTRIB-VII-058}, so that the composite immersion satisfies the condition of 4.1. (see 4.2). Proposition 4.3 is then true for the composite immersion by the preceding argument; it therefore remains to ``descend'' to \(X\). We descend first to \var{\(\underline{\underline{P}}_X^1\)}{\(\mathbf P_X^1\)}{EN-CONTRIB-VII-059}, then to \(X\), so that the proof consists in applying the following lemma twice (due to J.P. Jouanolou):

\sgaStableTarget{sga6.e07.lem.unnamed}\textbf{Lemma}: \uline{Under the conditions of} 4.3, \uline{let} \var{\(s:X\to\underline{\underline{P}}_X^1\)}{\(s:X\to\mathbf P_X^1\)}{EN-CONTRIB-VII-060} \uline{be the section defined in} 4.2. \uline{If Proposition} 4.3 \uline{is true for} \(s\circ i\), \uline{it is true for} \(i\).

Supposing 4.3 true for \(s\circ i\), we have for every \(p\ge1\), and every \(y\in K^{\bullet}(Y)\):
\sgaStableTarget{sga6.e07.eq.4-3-3}\begin{equation}\tag{4.3.3}
 (s\circ i)_*\bigl(\lambda^p(N+1,y)\bigr)
 \equiv
 \lambda^p\bigl((s\circ i)_*(y)\bigr)
 \quad\text{mod.\ torsion}.
\end{equation}

The section \(s\) is defined by an ideal isomorphic to \var{\(\underline{O}_p(-1)\)}{\(\underline{O}_{\mathbf P}(-1)\)}{EN-CONTRIB-VII-061} (4.2); write \(\xi\) for its class in \var{\(K^{\bullet}(\underline{\underline{P}}_X^1)\)}{\(K^{\bullet}(\mathbf P_X^1)\)}{EN-CONTRIB-VII-062}. We then have the following relation for every \(x\in K^{\bullet}(X)\):
\[
  s_*(x)=p^*(x)(1-\xi),
\]
where \(p\) is the structural morphism \var{\(\underline{\underline{P}}_X^1\to X\)}{\(\mathbf P_X^1\to X\)}{EN-CONTRIB-VII-063}. Indeed, since \(p\circ s=\Id_X\), we have \(x=s^*(p^*(x))\), whence \(s_*(x)=s_*(s^*(p^*(x)))=p^*(x)s_*(1)\); since \(s_*(1)=1-\xi\), we have the required relation:

Relation (4.3.3) then becomes:
\[
\begin{aligned}
(s\circ i)_*\bigl(\lambda^p(N+1,y)\bigr)
&\equiv \lambda^p\bigl(p^*(i_*(y))(1-\xi)\bigr)\quad\text{mod.\ torsion}\\
&\equiv \lambda^p(\xi,p^*(i_*(y)))(1-\xi)\quad\text{mod.\ torsion}.
\end{aligned}
\]
%%%%%%%%%%  scan idx 465  |  folio 452  %%%%%%%%%%
\SGAPrintPage{452}%
Now \(\lambda^p(\xi,p^*(i_*(y)))-\lambda^p(1,p^*(i_*(y)))\) is divisible by \(1-\xi\): to see this, we can work in the universal \(\lambda\)-ring used to define the \(\lambda^p(\xi,X)\) (V 5); then \(\lambda^p(\xi,X)-\lambda^p(1,X)\) is a polynomial in \(\xi\), vanishing for \(\xi=1\), and hence divisible by \(\xi-1\). Furthermore, the dependence equation of \(\xi\) over \(K^{\bullet}(X)\) is \((1-\xi)^2=0\) (VI (1.13.1)). We therefore obtain:
\[
\begin{aligned}
(s\circ i)_*\bigl(\lambda^p(N+1,y)\bigr)
&\equiv \lambda^p(1,p^*(i_*(y)))(1-\xi) \quad\text{mod.\ torsion}\\
&\equiv p^*(\lambda^p(1,i_*(y)))(1-\xi) \quad\text{mod.\ torsion}.
\end{aligned}
\]
Applying \(p_*\), we obtain:
\sgaStableTarget{sga6.e07.eq.4-3-4}\begin{equation}\tag{4.3.4}
  i_*\bigl(\lambda^p(N+1,y)\bigr)
  \equiv \lambda^p(1,i_*(y)) \quad\text{mod.\ torsion},
\end{equation}
since \(p_*(\xi)=0\).

By the formula in (V 5.6), applied with \(L=1\), \(\lambda^p(N+1,y)\) is the coefficient of \(t^p\) in the series:
\[
  \frac{\sum_{k\ge1}\lambda^k(N,y)kt^k}
       {1+\sum_{k\ge1}\lambda^k(N,y)\lambda_{-1}(N)t^k}.
\]

Similarly, the same formula applied with \(N=0\) and \(L=1\) shows that \(\lambda^p(1,y)\) is the coefficient of \(t^p\) in the series:
\[
  \frac{\sum_{k\ge1}\lambda^k(y)kt^k}
       {\sum_{k\ge0}\lambda^k(y)t^k}.
\]

The congruences (4.3.4) can therefore be written:
\[
i_*\left(
  \frac{\sum_{k\ge1}k\lambda^k(N,y)t^k}
       {1+\sum_{k\ge1}\lambda^k(N,y)\lambda_{-1}(N)t^k}
\right)
\equiv
  \frac{\sum_{k\ge1}k\lambda^k(i_*(y))t^k}
       {\sum_{k\ge0}\lambda^k(i_*(y))t^k}
       \quad\text{mod.\ torsion}
\]

%%%%%%%%%%  scan idx 466  |  folio 453  %%%%%%%%%%
\SGAPrintPage{453}%
Applying formula 2.8, we obtain:
\[
  \frac{\sum_{k\ge0}k\,i_*(\lambda^k(N,y))t^k}
       {1+\sum_{k\ge1} i_*(\lambda^k(N,y))t^k}
  \equiv
  \frac{\sum_{k\ge0}k\lambda^k(i_*(y))t^k}
       {\sum_{k\ge0}\lambda^k(i_*(y))t^k}
       \quad\text{mod.\ torsion}
\]
Putting:
\[
  u(t)=1+\sum_{k\ge1}i_*(\lambda^k(N,y))t^k,\qquad
  v(t)=\sum_{k\ge0}\lambda^k(i_*(y))t^k,
\]
we can also write this last relation as:
\[
  \frac{d}{dt}\!\left(\frac{u(t)}{v(t)}\right)\equiv0
  \quad\text{mod.\ torsion}.
\]
Since the constant term of \(u(t)/v(t)\) is \(1\), and we are writing congruences modulo torsion, we can integrate and write the congruence \(u(t)\equiv v(t)\). This gives the stated congruences.

\sgaStableTarget{sga6.e07.cor.4-4}\textbf{Corollary 4.4.} \uline{With the notation and hypotheses of} 4.3, \uline{we have for every} \(p\ge1\), \uline{and every} \(y\in K^{\bullet}(Y)\):
% VII probe display variant EN-CONTRIB-VII-064
\SGACaptureAppState{EN-CONTRIB-VII-064}%
\ifnum\SGAReaderForm=1\relax\SGAAppMark{EN-CONTRIB-VII-064}\fi
\ifnum\SGAReaderForm=0\relax
\[
  i_*\bigl(\gamma^p(N,y)\bigr)\equiv \gamma^p(i'_*(y))
  \quad\text{mod.\ torsion}.
\]
\else
\[
  i_*\bigl(\gamma^p(N,y)\bigr)\equiv \gamma^p(i_*(y))
  \quad\text{mod.\ torsion}.
\]
\fi

The corollary \var{follows immediately from 4.3, and using}{follows immediately from 4.3, using}{EN-CONTRIB-VII-065} the expression of the \(\gamma^p\) as linear combinations with integer coefficients of the \(\lambda^p\).

\sgaStableTarget{sga6.e07.rem.4-5}\textbf{Remark 4.5.} The proof of 4.3 itself suggests that (4.3.1) should hold not only modulo torsion, but also as an equality. In [RRR], the equality was proved in characteristic \(0\) by means of a geometric characterization of the \(\lambda^p(N,x)\). An extension of this method to all characteristics would probably be possible if one knew how to define exterior-power operations on perfect complexes (cf. XIV).

%%%%%%%%%%  scan idx 467  |  folio 454  %%%%%%%%%%
\SGAPrintPage{454}%
\sgaStableTarget{sga6.e07.thm.4-6}\textbf{Theorem 4.6.}
\uline{Let} \(X\) \uline{be a quasi-compact scheme having an ample sheaf}, \(d\) \uline{an integer} \(>0\) \uline{and} \(i:Y\to X\) \uline{a regular immersion of codimension} \(d\). \uline{Then}, \uline{for every} \(k\), \uline{we have}:
\[
 i_*\bigl((\Fil^kK^{\bullet}(Y))\otimes_{\mathbb Z}\mathbb Q\bigr)
 \subset
 (\Fil^{k+d}K^{\bullet}(X))\otimes_{\mathbb Z}\mathbb Q.
\]

First observe that \(i_*(K^{\bullet}(Y))\subset \Fil^1K^{\bullet}(X)\). Indeed, \(X-Y\) is dense in \(X\): otherwise, \(Y\) would contain an affine open subset, and on this open subset the ideal \(\underline{J}\) defining \(Y\) would be a nil ideal, hence a locally nilpotent ideal since it is of finite type. This is impossible, however, since the \(\underline{J}^n/\underline{J}^{n+1}\) are locally free \(\underline{O}_Y\)-Modules of rank \(\ne0\). The augmentation in \(K^{\bullet}(X)\) of an element of \(i_*(K^{\bullet}(Y))\) is a locally constant function on \(X\); since it vanishes on \(X-Y\), it vanishes on all of \(X\).

For every abelian group \(A\), we put \(A_{\mathbb Q}=A\otimes_{\mathbb Z}\mathbb Q\). Let \(x\in\Fil^kK^{\bullet}(Y)_{\mathbb Q}\); we can write \(x\) in the form:
\[
  x=\sum a\,\gamma^{i_1}(x_1)\cdots\gamma^{i_n}(x_n),
\]
with \(\epsilon(x_i)=0\) for every \(i\), \(a\in\mathbb Q\), and \(i_1+\cdots+i_n\ge k\). The \(x_i\) can be expressed as linear combinations with integer coefficients of the classes \(E_\alpha\) of finitely many locally free \(\underline{O}_Y\)-Modules of finite type \(\underline{E}_\alpha\). Let \(R\) be the sub-\(\mathbb Q\)-\(\lambda\)-algebra of \(K^{\bullet}(Y)_{\mathbb Q}\) generated by the \(E_\alpha\) and by \(N\). We then have:
\[
  x\in\Fil^kR.
\]
Furthermore, the \(\lambda\)-ring filtration of \(R\) is discrete (VI 6.2); let \(m\) satisfy \(\Fil^mR=0\). Finally, \var{by decomposing \(X\) into a finite sum}{by decomposing \(Y\) into a finite sum}{EN-CONTRIB-VII-066} of disjoint open subsets and carrying out the proof on each open subset (since \(K^{\bullet}\) commutes with finite direct sums), we may suppose that the \(\underline{E}_\alpha\) have constant rank \(r_\alpha\). The augmentation ring of \(R\) is then \(\mathbb Q\).

%%%%%%%%%%  scan idx 468  |  folio 455  %%%%%%%%%%
\SGAPrintPage{455}%
We shall show that for every \(i\), we have:
\[
  i_*(\Fil^iR)\subset \Fil^{i+d}(K^{\bullet}(X)_{\mathbb Q}).
\]
This will imply in particular:
% VII probe display variant EN-CONTRIB-VII-067
\SGACaptureAppState{EN-CONTRIB-VII-067}%
\ifnum\SGAReaderForm=1\relax\SGAAppMark{EN-CONTRIB-VII-067}\fi
\ifnum\SGAReaderForm=0\relax
\[
  i_*(x)\in\Fil^{k+d}K^{\bullet}(X),
\]
\else
\[
  i_*(x)\in\Fil^{k+d}K^{\bullet}(X)_{\mathbb Q},
\]
\fi
and hence prove the theorem.

The proof is by descending induction on \(i\), the case \(i=m\) being trivial. Let \(y\in\Fil^iR\), and put, for every \(j\):
\[
  a_j=(-1)^{j-1}(j-1)!.
\]
Then by V 6.10,
\[
  \gamma^{d+i}(N,y)-a_{d+i}y\in\Fil^{i+1}R.
\]

We can therefore apply the induction hypothesis, supposing that:
\[
  i_*(\Fil^{i+1}R)\subset \Fil^{i+d+1}K^{\bullet}(X)_{\mathbb Q}.
\]
Consequently, we have:
\[
  i_*\bigl(\gamma^{d+i}(N,y)-a_{d+i}y\bigr)
  \in\Fil^{i+d+1}K^{\bullet}(X)_{\mathbb Q}.
\]
By 4.4, this implies:
\[
  \gamma^{d+i}\bigl(i_*(y)\bigr)-a_{d+i}i_*(y)
  \in\Fil^{i+d+1}K^{\bullet}(X)_{\mathbb Q}.
\]
Since \(i_*(y)\) has zero augmentation, \var{\(\gamma^{d+i}\bigl(i_*(y)\bigr)\in\Fil^{d+i}K^{\bullet}(X)\)}{\(\gamma^{d+i}\bigl(i_*(y)\bigr)\in\Fil^{d+i}K^{\bullet}(X)_{\mathbb Q}\)}{EN-CONTRIB-VII-068}.
Since we are calculating modulo torsion, we therefore find:
\[
  i_*(y)\in\Fil^{i+d}K^{\bullet}(X)_{\mathbb Q}.
\]

\sgaStableTarget{sga6.e07.rem.4-6-1}\textbf{Remark 4.6.1.} We see that, even if 4.3 were true without tensoring with \(\mathbb Q\), the proof above works only modulo torsion. For a counterexample, see XIV 4.6.

%%%%%%%%%%  scan idx 469  |  folio 456  %%%%%%%%%%
\SGAPrintPage{456}%
\sgaStableTarget{sga6.e07.sec.4.7}\subsection*{4.7.}

We return to the situation of 3.7, where we had the commutative diagram:
\[
\begin{tikzcd}[column sep=large,row sep=large]
Y' \arrow[r,"j"] \arrow[d,"g"'] & X' \arrow[d,"f"]\\
Y \arrow[r,"i"] & X,
\end{tikzcd}
\]
where \(i\) is a regular immersion, \(X'\) the \(X\)-scheme obtained by blowing up \(Y\), and \(Y'=X'\times_XY\). Suppose that \(X\) is noetherian with an ample sheaf, and \(Y\) regular. We saw in 3.7 that we then have the exact sequence:
\[
0\longrightarrow K^{\bullet}(Y)\xrightarrow{u}
K^{\bullet}(Y')\times K^{\bullet}(X)\xrightarrow{v}
K^{\bullet}(X')\longrightarrow0.
\]
We shall show that this exact sequence describes the whole situation formed by the \(K^{\bullet}\) of \(X,Y,X',Y'\) from the data of the \(\lambda\)-rings \(K^{\bullet}(X)\), \(K^{\bullet}(Y)\), the maps \(i^*\) and \(i_*\), and the element \(N\) of \(K^{\bullet}(Y)\). (at least up to torsion, for the time being).

Let therefore \(K,K'\) be two \(\lambda\)-rings, \(i_*:K\to K'\) an additive homomorphism, \(i^*:K'\to K\) a \(\lambda\)-homomorphism, \(d\) an integer, and \(N\) an element of \(K\) such that \(\lambda^i(N)=0\) for \(i>d\) and \(\lambda^d(N)\) is invertible. Suppose moreover that the following properties hold:
\begin{enumerate}[label=\roman*)]
\item \(i_*\) is \(K'\)-linear (``projection formula'');
\item \(i^*\circ i_*\) is the homomorphism ``multiplication by \(\lambda_{-1}(N)\)'';
\item for every \(n\ge1\), and every \(y\in K\), \(i_*\bigl(\lambda^n(N,y)\bigr)=\lambda^n(i_*(y))\) \var{(relation (4.3.1)}{(relation (4.3.1)).}{EN-CONTRIB-VII-069}
\end{enumerate}

\var{First we use VI.I.1 to construct}{First we use VI.1.1 to construct}{EN-CONTRIB-VII-070} a \(K\)-\(\lambda\)-algebra \(K_1\) which plays the role of \(K^{\bullet}(Y')\). Let \(L\) be an indeterminate, and \(K[L]\) the \(K\)-\(\lambda\)-algebra generated by the generator \(L\) subject to the relations \var{\({}^i(L)=0\)}{\(\lambda^i(L)=0\)}{EN-CONTRIB-VII-071} for \(i\ge2\) (V 4.9). Put:
\[
  P(L)=\lambda^d(N-L)=\sum_{i=0}^{d}(-1)^i\lambda^{d-i}(N)L^i.
\]
%%%%%%%%%%  scan idx 470  |  folio 457  %%%%%%%%%%
\SGAPrintPage{457}%
The ideal generated by \(P\) is a \(\lambda\)-ideal. To see this, first observe that the relations \(\lambda^i(N)=0\) for \(i>d\) imply that the \(\lambda^j(N-L)\) for \(j\ge d\) are multiples of \(\lambda^d(N-L)\). Working in \var{\(1+\widehat{K/(P)[[t]]}^+\)}{\(1+\widehat{K[L]/(P)[[t]]}^+\)}{EN-CONTRIB-VII-072}, the image of \(\lambda_t(N-L)\) is a polynomial of degree \(<d\), and a universal argument then shows that its \(\lambda^i\) for \(i\ge d\) equal \(1\) (cf. V 4.8). Since in \var{\(1+\widehat{K[[t]]}^+\)}{\(1+\widehat{K[L][[t]]}^+\)}{EN-CONTRIB-VII-073} we have \(\lambda_t(\lambda^d(N-L))=\lambda^d(\lambda_t(N-L))\), we deduce that the \(\lambda^i(\lambda^d(N-L))\) for \(i\ge1\) are multiples of \(P\). By the relations V 2.4.2, this implies that the ideal generated by \(P\) is a \(\lambda\)-ideal.

Put then \var{\(K_1=K/(P)\)}{\(K_1=K[L]/(P)\)}{EN-CONTRIB-VII-074}. We denote the natural homomorphism \(K\to K_1\) by \(g^*\); it is a \(\lambda\)-homomorphism. Furthermore, \(K_1\) is a free \(K\)-module with basis the elements \(1,L,\ldots,L^{d-1}\). The relation \(\lambda^d(N)=\lambda^{d-1}(N-L).L\) and the invertibility hypothesis on \(\lambda^d(N)\) imply that \(L\) is invertible. Consequently, \(K_1\) has as a basis over \(K\) the elements \(1,L^{-1},\ldots,L^{-d+1}\); we then define a map \(g_*:K_1\to K\) by assigning to every element \(y'=y_0+y_1L^{-1}+\ldots+y_{d-1}L^{-d+1}\) of \(K_1\) the element \(y_0\) of \(K\). Clearly the map thus defined is an additive, \(K\)-linear homomorphism, satisfying \(g_*\circ g^*=\Id\).

Next we endow \(K_1\times K'\) with a \(\lambda\)-ring structure, putting:
\sgaStableTarget{sga6.e07.eq.4-7-1}\sgaStableTarget{sga6.e07.eq.4-7-2}\begin{align}
 (y'_1,x_1)*(y'_2,x_2)
   &=(y'_1y'_2(1-L)+x_1y'_2+x_2y'_1,\; x_1x_2); \tag{4.7.1}\\
 \lambda^n_*(y',x)
   &=\left(\sum_{p=1}^{n}\lambda^{\,n-p}(x)\lambda^p(L,y'),\;
     \lambda^n(x)\right). \tag{4.7.2}
\end{align}
In these formulas, \(K_1\) is regarded as a \(K'\)-algebra through the homomorphism \(g^*\circ i^*\).
Next, for \(y_1,y_2\in K\) and \(n\ge1\), put:
\sgaStableTarget{sga6.e07.eq.4-7-3}\sgaStableTarget{sga6.e07.eq.4-7-4}\begin{align}
 y_1*y_2&=-\,y_1y_2\lambda_{-1}(N); \tag{4.7.3}\\
 \lambda^n_*(y_1)&=-\,\lambda^n(N,-y_1); \tag{4.7.4}
\end{align}

%%%%%%%%%%  scan idx 471  |  folio 458  %%%%%%%%%%
\SGAPrintPage{458}%
and define a homomorphism:
\[
  u:K\longrightarrow K_1\times K'
\]
by putting:
\[
  u(y)=\bigl(g^*(y).\lambda_{-1}(F),-i_*(y)\bigr),
\]
with \(F=N-L\).

\sgaStableTarget{sga6.e07.prop.4-7-5}\textbf{Proposition 4.7.5.} \uline{The operations defined by} (4.7.1) \uline{and} (4.7.2) \uline{endow} \(K_1\times K'\) \uline{with a} \(\lambda\)-\uline{ring structure}, \uline{and} \(u\) \uline{is an injective homomorphism identifying} \(K\), \uline{equipped with the operations defined by} (4.7.3) \uline{and} (4.7.4), \uline{with a} \(\lambda\)-\uline{ideal of} \(K_1\times K'\).

It is immediate to verify that \(K_1\times K'\) is a \(\lambda\)-ring: one can refer, for example, to V 5.5, where an entirely analogous construction is given. The second part of the proposition is also proved without difficulty using the preceding relations, the relation \(i_*(y_1y_2\lambda_{-1}(N))=i_*(y_1).i_*(y_2)\) which follows immediately from the hypotheses, the relation \(g_*(\lambda_{-1}(F))=1\) (see VI 5.9), and the relations of V 5.

Put then:
\[
  K'_1=(K_1\times K')/K
\]
and denote by \(v\) the canonical homomorphism \(K_1\times K'\to K'_1\). Regard \(K'_1\) as a \(\lambda\)-ring through the quotient structure.

Next define additive homomorphisms:
\[
\begin{array}{llll}
j_*:K_1\to K'_1; & j^*:K'_1\to K_1; \\
f_*:K'_1\to K'; & f^*:K'\to K'_1 .
\end{array}
\]

The homomorphism \(j_*\) is the composite \(K_1\to K_1\times K'\to K'_1\). Likewise, the homomorphism \(f^*\) is the composite \(K'\to K_1\times K'\to K'_1\). We then have:
\[
  v=j_*+f^*.
\]

%%%%%%%%%%  scan idx 472  |  folio 459  %%%%%%%%%%
\SGAPrintPage{459}%
Next define a homomorphism \(K_1\times K'\to K_1\) by assigning to \((y',x)\) the element \(y'(1-L)+g^*\circ i^*(x)\); this homomorphism vanishes on \(K\), and consequently defines a homomorphism \(j^*:K'_1\longrightarrow K_1\). Finally, define a homomorphism \(K_1\times K'\to K'\) by assigning to \((y',x)\) the element \(i_*\circ g_*(y')+x\); this homomorphism vanishes on \(K\) and therefore defines a homomorphism \(f_*:K'_1\longrightarrow K'\).

\sgaStableTarget{sga6.e07.prop.4-7-6}\textbf{Proposition 4.7.6.}
\begin{enumerate}[label=\roman*)]
\item \(j^*\) \uline{and} \(f^*\) \uline{are} \(\lambda\)-\uline{homomorphisms};
\item \(j_*\) \uline{is} \(K'_1\)-\uline{linear and} \(f_*\) \uline{is} \(K'\)-\uline{linear};
\item \(j^*\circ j_*\) \uline{is the homomorphism} ``\uline{multiplication by} \(1-L\)'';
\item \uline{for every} \(y'\in K_1\), \uline{and every} \(p\ge1\), \(j_*\bigl(\lambda^p(L,y')\bigr)=\lambda^p(j_*(y'))\);
\item \(f_*\circ f^*=\Id\).
\end{enumerate}

We need only set the machinery in motion.

Returning to the geometric situation recalled above, we see that \(K^{\bullet}(X')\) is constructed from \(K^{\bullet}(Y)\), \(K^{\bullet}(X)\), the maps \(i^*\) and \(i_*\), and the element \(N\) of \(K^{\bullet}(Y)\) by the preceding method; the \(\lambda\)-operations on \(K^{\bullet}(X')\), however, agree with the \(\lambda\)-operations defined above on \(K'_1\) only up to torsion, as long as relation (4.3.1) is proved only modulo torsion.

\sgaStableTarget{sga6.e07.thm.4-8}\textbf{Theorem 4.8.} \uline{Let} \(i:Y\to X\) \uline{be a regular closed immersion of constant codimension} \(d\), \(X'\) \uline{the scheme obtained by blowing up} \(Y\), \uline{and} \(Y'=X'\times_XY\), \uline{so that we have the commutative diagram}:
\[
\begin{tikzcd}[column sep=large,row sep=large]
Y' \arrow[r,"j"] \arrow[d,"g"'] & X' \arrow[d,"f"]\\
Y \arrow[r,"i"] & X.
\end{tikzcd}
\]
\uline{Suppose that} \(X\) \uline{is quasi-compact and has an ample sheaf}. \uline{Then for every} \(k\) \uline{we have the exact sequence}:

%%%%%%%%%%  scan idx 473  |  folio 460  %%%%%%%%%%
\SGAPrintPage{460}%
\[
0\longrightarrow
\operatorname{Gr}^{k-d}K^{\bullet}(Y)_{\mathbb Q}
\xrightarrow{u^{\gr}}
\operatorname{Gr}^{k-1}K^{\bullet}(Y')_{\mathbb Q}
\times
\operatorname{Gr}^{k}K^{\bullet}(X)_{\mathbb Q}
\xrightarrow{v^{\gr}}
\operatorname{Gr}^{k}K^{\bullet}(X')_{\mathbb Q}
\]
\uline{where we put}, \uline{for} \(y\in \operatorname{Gr}^{k-d}K^{\bullet}(Y)_{\mathbb Q}\):
\[
  u^{\gr}(y)=
  \bigl(g^{\gr}(y).(-1)^{d-1}c^{d-1}(F),-i_{\gr}(y)\bigr),
\]
\uline{and for} \(y'\in \operatorname{Gr}^{k-1}K^{\bullet}(Y')_{\mathbb Q}\) \uline{and} \(x\in \operatorname{Gr}^{k}K^{\bullet}(X)_{\mathbb Q}\):
\[
  v^{\gr}(y',x)=j_{\gr}(y')+f^{\gr}(x),
\]
\uline{with the notation of} 3.1 \uline{and} 3.3, \uline{and} \(c^{d-1}(F)\) \uline{denoting the} \((d-1)\)-\uline{st Chern class of} \(F\) (V 6.7). \uline{Moreover}, \(u^{\gr}\) \uline{has a left inverse} \(u'^{\gr}\), \uline{defined by} \(u'^{\gr}(y',x)=g_{\gr}(y')\).

\uline{If in addition} \(X\) \uline{is noetherian and} \(Y\) \uline{regular}, \uline{we have the exact sequence}:
\[
0\longrightarrow
\operatorname{Gr}^{k-d}K^{\bullet}(Y)_{\mathbb Q}
\xrightarrow{u^{\gr}}
\operatorname{Gr}^{k-1}K^{\bullet}(Y')_{\mathbb Q}
\times
\operatorname{Gr}^{k}K^{\bullet}(X)_{\mathbb Q}
\xrightarrow{v^{\gr}}
\operatorname{Gr}^{k}K^{\bullet}(X')_{\mathbb Q}
\longrightarrow0
\]

\uline{Proof.} Observe that, when \(Y\) is regular, one could give a purely algebraic proof in terms of the theory of \(\lambda\)-rings, using the description of the situation in 4.7. We shall instead follow the proof of 3.7.

We know that the ``inverse-image'' homomorphisms \(f^*,g^*\), being \(\lambda\)-homomorphisms, are compatible with the \(\lambda\)-ring filtrations, and consequently define homomorphisms on the associated gradeds; hence the homomorphisms:
\[
  f^{\gr}:\operatorname{Gr}^kK^{\bullet}(X)\to\operatorname{Gr}^kK^{\bullet}(X');
  \qquad
  g^{\gr}:\operatorname{Gr}^{k-d}K^{\bullet}(Y)\to\operatorname{Gr}^{k-d}K^{\bullet}(Y'),
\]
and the analogous homomorphisms obtained after tensoring with \(\mathbb Q\). Furthermore, 4.6 shows that after tensoring with \(\mathbb Q\), the homomorphisms \(i_*,j_*\) pass to the associated gradeds, and therefore define:
\[
  i_{\gr}:\operatorname{Gr}^{k-d}K^{\bullet}(Y)_{\mathbb Q}
  \to \operatorname{Gr}^{k}K^{\bullet}(X)_{\mathbb Q};
  \qquad
  j_{\gr}:\operatorname{Gr}^{k-1}K^{\bullet}(Y')_{\mathbb Q}
  \to \operatorname{Gr}^{k}K^{\bullet}(X')_{\mathbb Q}.
\]

%%%%%%%%%%  scan idx 474  |  folio 461  %%%%%%%%%%
\SGAPrintPage{461}%
Finally, observe that \(\lambda_{-1}(F)=(-1)^{d-1}\gamma^{d-1}(F)\), so the morphisms \(u^{\gr}\) and \(v^{\gr}\) are obtained by passage to the quotient from the morphisms \(u\) and \(v\) of 3.7. Thus we have \(v^{\gr}\circ u^{\gr}=0\).

By VI 5.8, \(g_*(\Fil^{k-1}K^{\bullet}(Y'))\subset \Fil^{k-d}K^{\bullet}(Y)\) for every \(k\), since \var{\(Y'=\underline{\underline{P}}(\underline{N})\)}{\(Y'=\mathbf P(\underline{N})\)}{EN-CONTRIB-VII-075}, with \(\underline{N}\) of rank \(d\). We therefore obtain a homomorphism:
\[
  g_{\gr}:\operatorname{Gr}^{k-1}K^{\bullet}(Y')
  \to\operatorname{Gr}^{k-d}K^{\bullet}(Y),
\]
and an analogous homomorphism after tensoring with \(\mathbb Q\). The homomorphism \(u'^{\gr}\) is therefore obtained from \(u'\) (3.7) by passage to the associated graded. Since \(u'\) is a left inverse of \(u\), \(u'^{\gr}\) is a left inverse of \(u^{\gr}\), proving at the same time the injectivity of \(u^{\gr}\).

To show that \(\Ker(v^{\gr})=\Im(u^{\gr})\), we must show that for \(y'\in\Fil^{k-1}K^{\bullet}(Y')_{\mathbb Q}\), and \(x\in\Fil^kK^{\bullet}(X)_{\mathbb Q}\), satisfying:
\sgaStableTarget{sga6.e07.eq.4-6-1}% VII probe display variant EN-CONTRIB-VII-076
\SGACaptureAppState{EN-CONTRIB-VII-076}%
\ifnum\SGAReaderForm=1\relax\SGAAppMark{EN-CONTRIB-VII-076}\fi
\ifnum\SGAReaderForm=0\relax
\begin{equation}\tag{4.6.1}
  j_*(y')+f^*(x)\in\Fil^{k+1}K^{\bullet}(X')_{\mathbb Q},
\end{equation}
\else
\begin{equation}\tag{4.8.1}
  j_*(y')+f^*(x)\in\Fil^{k+1}K^{\bullet}(X')_{\mathbb Q},
\end{equation}
\fi
there exists \(y\in\Fil^{k-d}K^{\bullet}(Y)_{\mathbb Q}\) such that:
\[
\begin{aligned}
\text{i)}\quad &x+i_*(y)\in\Fil^{k+1}K^{\bullet}(X)_{\mathbb Q},\\
\text{ii)}\quad &y'-g^*(y)\lambda_{-1}(F)\in\Fil^{k}K^{\bullet}(Y')_{\mathbb Q}.
\end{aligned}
\]

Put:
\[
  y=g_*(y').
\]
\var{Then by VI 5.8, \(y\in\Fil^{k}K^{\bullet}(Y)\)}{Then by VI 5.8, \(y\in\Fil^{k-d}K^{\bullet}(Y)_{\mathbb Q}\)}{EN-CONTRIB-VII-077}. Next observe that \(f_*\) preserves the filtration (using results of Exposé VIII, which causes no difficulty, since 4.8 is not used to prove this fact). Indeed, it follows from 1.8 ii) that \(f\) is a complete-intersection morphism of virtual relative dimension zero (VIII 4.2). Since \(f\) is projective, the result follows from VIII 3.2. Applying \(f_*\) to the relation
%%%%%%%%%%  scan idx 475  |  folio 462  %%%%%%%%%%
\SGAPrintPage{462}%
(4.8.1), we obtain:
\[
  f_*(j_*(y'))+f_*(f^*(x))\in\Fil^{k+1}K^{\bullet}(X)_{\mathbb Q},
\]
or, by 3.6:
\[
  i_*(y)+x\in\Fil^{k+1}K^{\bullet}(X)_{\mathbb Q}.
\]

To prove relation ii), we must show that:
\[
  y'-g^*(g_*(y'))\lambda_{-1}(F)\in\Fil^kK^{\bullet}(Y')_{\mathbb Q}.
\]
Put \(z=y'-g^*(g_*(y'))\lambda_{-1}(F)\). We have:
\[
\begin{aligned}
j_*(z)
  &=j_*(y')-j_*\bigl(g^*(g_*(y'))\lambda_{-1}(F)\bigr)\\
  &=j_*(y')-f^*(i_*(y))\qquad\text{by 3.4.}
\end{aligned}
\]
Now relation i) gives:
\[
  f^*(i_*(y))+f^*(x)\in\Fil^{k+1}K^{\bullet}(X')_{\mathbb Q}.
\]
\var{Taking (4.6.1) into account, we therefore obtain:}{Taking (4.8.1) into account, we therefore obtain:}{EN-CONTRIB-VII-078}
\[
  j_*(z)\in\Fil^{k+1}K^{\bullet}(X')_{\mathbb Q},
\]
whence:
\[
  j^*(j_*(z))=z(1-L)\in\Fil^{k+1}K^{\bullet}(Y')_{\mathbb Q}.
\]

Let \(\bar z\) be the image of \(z\) in \(\operatorname{Gr}^{k-1}K^{\bullet}(Y')_{\mathbb Q}\), and \(\overline{(1-L)}\) that of \((1-L)\) in \(\operatorname{Gr}^{1}K^{\bullet}(Y')_{\mathbb Q}\). We therefore have: \(\overline{z(1-L)}=0\). We can then apply (VI 5.11) (which remains valid after tensoring with \(\mathbb Q\), since its proof uses only the structure of \(K^{\bullet}(Y')\) as a free \(K^{\bullet}(Y)\)-algebra), and hence find \(z'\in K^{\bullet}(Y)_{\mathbb Q}\) such that:
\[
  z-g^*(z')\lambda_{-1}(F)\in\Fil^kK^{\bullet}(Y')_{\mathbb Q}.
\]
Applying \(g_*\), and using \(g_*(z)=0\), we find:
\[
  z'\in\Fil^{k+1-d}K^{\bullet}(Y)_{\mathbb Q},
\]
whence, since \(\lambda_{-1}(F)\in\Fil^{d-1}K^{\bullet}(Y')_{\mathbb Q}\):
\[
  z\in\Fil^kK^{\bullet}(Y')_{\mathbb Q},
\]
which completes the proof.

%%%%%%%%%%  scan idx 476  |  folio 463  %%%%%%%%%%
\SGAPrintPage{463}%
Suppose now that \(X\) is noetherian and \(Y\) regular, which implies that \(v\) is surjective (3.7). We saw in the proof of 3.7 that for every \(x'\in K^{\bullet}(X')\) there is \(y'\in K^{\bullet}(Y')\) such that:
\[
  x'=f^*(f_*(x'))+j_*(y').
\]
The same property clearly remains true after tensoring with \(\mathbb Q\). Consequently, for \(x'\in\Fil^kK^{\bullet}(X')_{\mathbb Q}\) there is \((y',x)\in K^{\bullet}(Y')_{\mathbb Q}\times K^{\bullet}(X)_{\mathbb Q}\) such that:
\[
  x'=v(y',x),\qquad x\in\Fil^kK^{\bullet}(X)_{\mathbb Q}.
\]
Indeed, it suffices to take \(x=f_*(x')\).

Filter the ring \(K^{\bullet}(Y')_{\mathbb Q}\times K^{\bullet}(X)_{\mathbb Q}\) defined in 4.7 by putting:
\[
  \Fil^k\bigl(K^{\bullet}(Y')_{\mathbb Q}\times K^{\bullet}(X)_{\mathbb Q}\bigr)
  =
  \Fil^{k-1}K^{\bullet}(Y')_{\mathbb Q}\times \Fil^kK^{\bullet}(X)_{\mathbb Q}.
\]
It follows immediately from (4.7.1) that this is a ring filtration. Thus, to show that for every \(k\) the homomorphism:
\[
  \Fil^k\bigl(K^{\bullet}(Y')_{\mathbb Q}\times K^{\bullet}(X)_{\mathbb Q}\bigr)
  \xrightarrow{v}
  \Fil^kK^{\bullet}(X')_{\mathbb Q}
\]
is surjective, it suffices to show that for \(x'= \gamma^k(x'_1)\), with \(\epsilon(x'_1)=0\), there is \((y',x)\in\Fil^k(K^{\bullet}(Y')_{\mathbb Q}\times K^{\bullet}(X)_{\mathbb Q})\) such that \(x'=v(y',x)\).

By the preceding argument, there is \((y'_1,x_1)\in K^{\bullet}(Y')_{\mathbb Q}\times K^{\bullet}(X)_{\mathbb Q}\) such that \(x'_1=v(y'_1,x_1)\) and \(x_1\in\Fil^1K^{\bullet}(X)_{\mathbb Q}\). By 4.7, we then have:
\[
  x'=v\bigl(\gamma^k_*(y'_1,x_1)\bigr).
\]
Now:
\[
  \gamma^k_*(y'_1,x_1)
  =
  \left(\sum_{p=1}^{k}\gamma^{k-p}(x_1)\gamma^p(L,y'_1),\;
  \gamma^k(x_1)\right)
\]
(an immediate consequence of (4.7.2)). Since \(x_1\) has zero augmentation, \(\gamma^k(x_1)\in\Fil^kK^{\bullet}(X)_{\mathbb Q}\), and, using V 5.7,
\[
  \sum_{p=1}^{k}\gamma^{k-p}(x_1)\gamma^p(L,y'_1)
  \in\Fil^{k-1}K^{\bullet}(Y')_{\mathbb Q},
\]
which gives the result.

%%%%%%%%%%  scan idx 477  |  folio 464  %%%%%%%%%%
\SGAPrintPage{464}%
\sgaStableTarget{sga6.e07.rem.4-9}\textbf{Remark 4.9.} When \(X\) is noetherian and \(Y\) regular, proceeding as in 4.7 we can describe the graded ring structure of \(\operatorname{Gr}^*K^{\bullet}(X')_{\mathbb Q}\) in terms of the graded rings \(\operatorname{Gr}^*K^{\bullet}(X)_{\mathbb Q}\), \(\operatorname{Gr}^*K^{\bullet}(Y)_{\mathbb Q}\), the additive homomorphisms \(i^{\gr}\) and \(i_{\gr}\) induced by \(i^*\) and \(i_*\), and the element \(c^{\bullet}(N)\) of \(\operatorname{Gr}^*K^{\bullet}(Y)\). Likewise, the same data describe the homomorphisms \(f^{\gr},f_{\gr},j^{\gr},j_{\gr}\) defined from \(f^*,f_*,j^*,j_*\) by passage to the associated gradeds.

\sgaStableTarget{sga6.e07.sec.4.10}\subsection*{4.10.}

Let \(X\) be a \uline{regular} scheme having an ample sheaf. We know (IV\qquad) that the Grothendieck groups formed with locally free \(\underline{O}_X\)-Modules of finite type and with coherent \(\underline{O}_X\)-Modules are isomorphic. We use this fact to define a ``topological'' filtration on \(K^{\bullet}(X)\), taking \(\Fil^i_{\mathrm{top}}K^{\bullet}(X)\) to be the subgroup generated by the classes of coherent sheaves whose support has codimension \(\ge i\); it is easy to show that \(\Fil^i_{\mathrm{top}}K^{\bullet}(X)\) is generated by the classes of the sheaves \(\underline{O}_Z\), where \(Z\) ranges over the integral closed subschemes of \(X\) of codimension \(\ge i\)(X\qquad).

\sgaStableTarget{sga6.e07.prop.4-11}\textbf{Proposition 4.11.} \uline{Let} \(X\) \uline{be a regular scheme having an ample sheaf}. \uline{Then} \uline{for every} \(k\ge0\):
\[
  \Fil_{\mathrm{top}}^kK^{\bullet}(X)_{\mathbb Q}
  \subset
  \Fil_\lambda^kK^{\bullet}(X)_{\mathbb Q}.
\]
\((\mathrm{N.B.}\) The reverse inclusion, valid without tensoring with \(\mathbb Q\), will be proved in X\qquad).

By 4.10, it suffices to show that for an integral closed subscheme \(Z\) of \(X\), of codimension \(k\), we have:
\[
  \operatorname{cl}(\underline{O}_Z)\in\Fil^kK^{\bullet}(X)_{\mathbb Q}.
\]

For this, we proceed by noetherian induction on the integral closed subschemes of \(X\). Suppose therefore that the property is proved for the integral closed subschemes of \(Z\), and let us prove it for \(Z\). Let \(\underline{J}\) be the Ideal of \(\underline{O}_X\)

%%%%%%%%%%  scan idx 478  |  folio 465  %%%%%%%%%%
\SGAPrintPage{465}%
defining \(Z\), and let \(z\) be the generic point of \(Z\). Since \(Z\) is integral, \(\underline{J}_z\) is the maximal ideal of \(\underline{O}_{X,z}\), and since \(X\) is regular, \(\underline{J}_z\) can be generated by an \(\underline{O}_{X,z}\)-regular sequence. We can therefore find a sequence of sections of \(\underline{J}\) in a neighborhood of \(z\) forming an \(\underline{O}_X\)-regular sequence. Consequently, there is an open neighborhood \(U\) of \(z\) such that \(Z'=X-U\) is contained in \(Z\), and \(\underline{J}|U\) is a regular ideal of \(\underline{O}_U\). The immersion \(Z\cap U\xrightarrow{i}U\) is then a regular immersion, with codimension \(\dim(\underline{O}_{X,z})=\codim(Z,X)\) (EGA IV 5.1.2). Furthermore, we have \var{the exact sequence (IV\qquad:}{the exact sequence (IV\qquad):}{EN-CONTRIB-VII-079}
\[
  K^{\bullet}(Z')\xrightarrow{j_*}K^{\bullet}(X)
  \xrightarrow{u^*}K^{\bullet}(U)\to0,
\]
where \(j\) is the immersion \(Z'\to X\) and \(u\) the immersion \(U\to X\); tensoring with \(\mathbb Q\) gives an analogous sequence. Now \(u^*(\underline{O}_Z)=i_*(\underline{O}_{Z\cap U})\), so by 4.6 the image \(u^*(1_Z)\) of the class of \(\underline{O}_Z\) under \(u^*\) belongs to \(\Fil^kK^{\bullet}(U)_{\mathbb Q}\). Since \(u^*\) is a surjective \(\lambda\)-homomorphism, it induces a \uline{surjection} \(\Fil^kK^{\bullet}(X)\longrightarrow \Fil^kK^{\bullet}(U)\), so \(u^*(1_Z)=u^*(x)\) with \(x\in\Fil^kK^{\bullet}(X)\); there is therefore \(y\in j_*(K^{\bullet}(Z'))\) such that:
\[
  \operatorname{cl}(\underline{O}_Z)-y\in\Fil^kK^{\bullet}(X)_{\mathbb Q}.
\]
Now \(Z'\subset Z\), and since \(z\notin Z'\), \(\codim(Z',X)>k\). \var{Consequently, \(x\) can be expressed}{Consequently, \(y\) can be expressed}{EN-CONTRIB-VII-080} as a combination of classes of sheaves \(\underline{O}_{Z''_\alpha}\), where the \(Z''_\alpha\) are integral closed subschemes of \(Z\), distinct from \(Z\). By the induction hypothesis, \var{\(x\in\Fil^{k+1}K^{\bullet}(X)_{\mathbb Q}\)}{\(y\in\Fil^{k+1}K^{\bullet}(X)_{\mathbb Q}\)}{EN-CONTRIB-VII-081}, which completes the proof.

\sgaStableTarget{sga6.e07.bibliography}\subsection*{BIBLIOGRAPHY}
\begin{enumerate}[label={[\arabic*]}]
\item Cartan--Eilenberg, Homological Algebra.
\end{enumerate}



% Expose VIII macro additions
\providecommand{\etale}{\'etale}
\providecommand{\Ch}{\operatorname{Ch}}


\sgaStableTarget{sga6.e08.expose}\section*{Expos\'e VIII. The Riemann--Roch theorem}
\begin{center}
by Pierre Berthelot
\end{center}

\sgaStableTarget{sga6.e08.sec.1}\subsection*{1. Complete-intersection morphisms}

\sgaStableTarget{sga6.e08.def.1-1}\textbf{Definition 1.1.} Let $f:X\to Y$ be a morphism of schemes. One says that $f$ is a \emph{complete-intersection morphism} if, for every point $x$ of $X$, there exist an open neighbourhood $U$ of $x$ and a smooth $Y$-scheme $V$ such that the restriction of $f$ to $U$ factors as
\[
\begin{tikzcd}[column sep=tiny]
U \arrow[r,"i"] \arrow[dr,"f"'] & V \arrow[d,"g"]\\
& Y,
\end{tikzcd}
\]
where $i$ is a regular immersion (\hyperlink{sga6.e07.def.1-4}{VII 1.4}). After replacing $V$ by an open subset, one may suppose that $i$ is a regular closed immersion.

In this definition one may replace $V$ by a vector bundle $Y[T_1,\ldots,T_n]$. Indeed, since the notion is local, one may suppose that $V$ is affine and that its image in $Y$ is contained in an affine open $W$ of $Y$. There then exists an integer $n$ and a closed $W$-immersion
\[
V\longrightarrow W[T_1,\ldots,T_n],
\]
hence a $Y$-immersion $V\to Y[T_1,\ldots,T_n]$, which is regular by \hyperlink{sga6.e07.prop.1-10}{VII 1.10}. Since the composite of two regular immersions is regular, one obtains a regular $Y$-immersion
\[
U\longrightarrow Y[T_1,\ldots,T_n]
\]
factorizing $f$.

\sgaStableTarget{sga6.e08.prop.1-2}\textbf{Proposition 1.2.} Let $f:X\to Y$ be a complete-intersection morphism. Then every $Y$-immersion $i:X\to Z$ of $X$ in a smooth $Y$-scheme is regular.

We first prove the following corollary.

\sgaStableTarget{sga6.e08.cor.1-3}\textbf{Corollary 1.3.} Let
\[
\begin{tikzcd}[column sep=large,row sep=large]
& V' \arrow[d,"p"]\\
X \arrow[ur,"j"] \arrow[r,"i"'] & V
\end{tikzcd}
\]
be a commutative diagram in which $i$ and $j$ are immersions, and $p$ is a smooth morphism. Then $i$ is a regular immersion if and only if $j$ is a regular immersion.

(i) Suppose $i$ regular. Put $X'=V'\times_V X$. The datum of $j$ defines a section $j'$ of $X'$ over $X$, such that $j=i'j'$:
\[
\begin{tikzcd}[column sep=large,row sep=large]
X' \arrow[r,"i'"] \arrow[d,"p'"'] & V' \arrow[d,"p"]\\
X \arrow[r,"i"'] \arrow[ur,"j"] & V .
\end{tikzcd}
\]
Now $j'$ is a regular immersion, being a section of a smooth morphism (\hyperlink{sga6.e07.prop.1-10}{VII 1.10}); $i'$ is a regular immersion because it is obtained from $i$ by the smooth, hence flat, base change (\hyperlink{sga6.e07.prop.1-5}{VII 1.5}). Therefore $j$ is a regular immersion (\hyperlink{sga6.e07.prop.1-7}{VII 1.7}).

(ii) Suppose $j$ regular and $p$ \etale. Then $j'$ is an open immersion (EGA IV 17.9.3), hence an isomorphism of $X$ onto $j'(X)$; put $U=j'(X)$. Let $i''$ be the restriction of $i'$ to $U$, and let $U'$ be an open subset of $V'$ such that $U=U'\cap X'$. Then
\[
i''=j\circ(j')^{-1},
\]
so that $i''$ is a regular immersion of $U$ into $U'$. But the morphism $U'\to V$, being smooth, is flat and locally of finite presentation, hence open. Consequently it is faithfully flat and quasi-compact from $U'$ onto an open neighbourhood of $X$ in $V$. It follows from \hyperlink{sga6.e07.prop.1-5}{VII 1.5} that $i$ is a regular immersion.

(iii) Suppose $j$ regular, and suppose that $V'$ is a vector bundle over $V$. Since the property to be proved is local on $V$, one may suppose $V$ affine, $i$ closed, and $V'=V[T_1,\ldots,T_n]$. Put $V=\Spec(A)$, $X=\Spec(B)$. One has the commutative diagram
\[
\begin{tikzcd}[column sep=large,row sep=large]
& A[T_1,\ldots,T_n] \arrow[d,"v"]\\
A \arrow[r] \arrow[ur,"u"] & B,
\end{tikzcd}
\]
where the morphisms $u$ and $v$ are surjective. Choose elements $t_i\in A$, $1\leq i\leq n$, such that $u(t_i)=v(T_i)$ for every $i$. Define a homomorphism
\[
w:A[T_1,\ldots,T_n]\longrightarrow A,
\qquad w(T_i)=t_i.
\]
This gives a closed immersion $s:V\to V'$ which is a section of $p$, and such that $j=s\circ i$. Moreover, being a section of a smooth morphism, $s$ is a regular immersion (\hyperlink{sga6.e07.prop.1-10}{VII 1.10}). Since $s$ and $s\circ i$ are regular, it will follow that $i$ is regular if one proves that the sequence
\[
0\longrightarrow i^*(K/K^2)\longrightarrow J/J^2\longrightarrow I/I^2\longrightarrow 0
\]
is exact and locally split; here $I,J,K$ denote respectively the ideals defining $i,j,s$ (\hyperlink{sga6.e07.prop.1-7}{VII 1.7}). It is enough to show that $i^*(K/K^2)$ is a direct factor of $J/J^2$. But $j'$ is obtained from $s$ by the base change $i':X'\to V'$. Moreover $j'$ and $s$ are regular immersions (\hyperlink{sga6.e07.prop.1-10}{VII 1.10}). Thus, by \hyperlink{sga6.e07.lem.1-2}{VII 1.2}, the ideal defining $j'$ is $i'^*(K)$, and the conormal sheaf of $j'$ is $i'^*(K/K^2)$. The exact sequence of conormal sheaves attached to the decomposition $j=i'j'$ therefore gives in particular a homomorphism
\[
J/J^2\longrightarrow i^*(K/K^2).
\]
It is immediate to verify that the composite
\[
i^*(K/K^2)\longrightarrow J/J^2\longrightarrow i^*(K/K^2)
\]
is the identity. In affine terms, this corresponds to the decomposition of $J$ as the direct sum of an ideal of $A$ and the ideal generated by the $T_i-t_i$. This gives the required splitting.

(iv) Finally suppose $j$ regular and $p$ smooth. For every $x\in X$, one can find an open neighbourhood $W$ of $i(x)$, an open neighbourhood $W'$ of $j(x)$ lying over $W$, and a commutative diagram
\[
\begin{tikzcd}[column sep=large,row sep=large]
W' \arrow[r,"q"] \arrow[d,"p"'] & W[T_1,\ldots,T_n] \arrow[d,"p_0"]\\
W \arrow[r,equal] & W,
\end{tikzcd}
\]
where $p_0$ is the structural morphism, and $q$ is an \etale morphism (EGA IV 17.11.4). One may then suppose $V=W$ and $V'=W'$. It is easy to see that $qj$ is an immersion, by introducing the fiber product $X[T_1,\ldots,T_n]$. Applying (ii) and (iii), one deduces that $qj$ is regular, and then that $i$ is regular.

We now prove \hyperlink{sga6.e08.prop.1-2}{Proposition 1.2}. Since regularity is local on $Z$, we may work in a neighbourhood of a point $x\in X$. There is then an open neighbourhood $U$ of $x$ in $X$, a smooth $Y$-scheme $V$, and a regular $Y$-immersion $i':U\to V$. We have a commutative diagram
\[
\begin{tikzcd}[column sep=large,row sep=large]
U \arrow[r,"i'"] \arrow[d,"i"'] & V \arrow[d]\\
Z \arrow[r] & Y,
\end{tikzcd}
\]
where the unlabeled vertical arrows are the structural morphisms. From it one deduces a closed immersion
\[
i'':U\longrightarrow Z\times_Y V,
\]
and a commutative diagram
\[
\begin{tikzcd}[column sep=large,row sep=large]
& U \arrow[dl,"i"'] \arrow[dr,"i'"] &\\
Z & Z\times_Y V \arrow[l,"p"'] \arrow[r,"p'"] & V.
\end{tikzcd}
\]
Since $i'$ is regular and $p,p'$ are smooth, the corollary implies that $i''$ and $i$ are regular.

The following proposition will not be used in the sequel of the Seminar.

\sgaStableTarget{sga6.e08.prop.1-4}\textbf{Proposition 1.4.} Let $f:X\to Y$ be a flat morphism locally of finite presentation. For $f$ to be a complete-intersection morphism in the sense of \hyperlink{sga6.e08.def.1-1}{1.1}, it is necessary and sufficient that it be a complete-intersection morphism in the sense of EGA IV 19.3.6.

Suppose that $f$ is a complete-intersection morphism in the sense of \hyperlink{sga6.e08.def.1-1}{1.1}, and let $x\in X$. Then there are an open neighbourhood $U$ of $x$ and a factorization
\[
\begin{tikzcd}[column sep=large,row sep=large]
U \arrow[r,"i"] \arrow[dr,"f"'] & V \arrow[d,"g"]\\
& Y,
\end{tikzcd}
\]
where $i$ is a regular immersion and $g$ is smooth. Let $y=f(x)$. Then the immersion $i_y:U_y\to V_y$ is regular at $x$. Indeed, after possibly restricting $V$, one may suppose that there is a finite-type free $\mathcal O_V$-module $E$ and a regular homomorphism $u:E\to\mathcal O_V$ such that the Koszul complex associated to $u$ is a resolution of $\mathcal O_U$. Let $u_y$ be the homomorphism induced on the fiber over $y$, and let $K_\bullet(u_y)=K_\bullet(u)_y$ be the associated Koszul complex. Since $g$ is smooth, hence flat, $K_\bullet(u)$ is a $Y$-flat resolution of $\mathcal O_U$. Consequently the homology groups of $K_\bullet(u)_y$ are the groups
\[
\Tor_i^{\mathcal O_Y}(\mathcal O_U,k(y)),
\]
where $k(y)$ is the residue field of $\mathcal O_{Y,y}$. Since $f$ is flat, these groups are zero; hence $i_y$ is regular at $x$. Since $g$ is smooth, the local ring of $V_y$ at $x$ is regular. The local ring of $U_y$ at $y$ is therefore the quotient of a regular ring by a regular ideal, hence a complete-intersection ring (EGA IV 19.3.2). It follows that $f$ is a complete intersection at $x$ in the sense of EGA IV 19.3.6.

Conversely, suppose that $f$ is a complete-intersection morphism in the sense of EGA IV 19.3.6, and let $x\in X$. Let $U$ be an affine neighbourhood of $x$ lying over an affine neighbourhood $U'$ of $y=f(x)$. One can then find a smooth $U'$-scheme $V$, and a closed immersion $U\to V$ (for instance by taking for $V$ a vector bundle over $U'$). We shall show that $i$ is regular. By EGA IV 19.3.7, $i$ is transversally regular at $x$, since $V$ is smooth over $Y$, hence $V_y$ is regular at $x$, and $U$ is flat over $Y$. It follows from EGA IV 19.2.4, a)3c), that $i$ is regular at $x$. Therefore $f$ is a complete-intersection morphism in the sense of 1.1.

\sgaStableTarget{sga6.e08.prop.1-5}\textbf{Proposition 1.5.} Let $f:X\to Y$ and $g:Y\to Z$ be two locally complete-intersection morphisms. Then $gf$ is locally complete-intersection.

Let $x\in X$. One can find open neighbourhoods $U,V$ of $x,f(x)$ and regular immersions $i,j$ such that one has factorizations
\[
\begin{tikzcd}[column sep=large,row sep=large]
U \arrow[r,"i"] \arrow[dr,"f"'] & V[T_1,\ldots,T_n] \arrow[d]\\
& V
\end{tikzcd}
\qquad
\begin{tikzcd}[column sep=large,row sep=large]
V \arrow[r,"j"] \arrow[dr,"g"'] & Z[T'_1,\ldots,T'_m] \arrow[d]\\
& Z.
\end{tikzcd}
\]
This diagram can be completed by the commutative square
\[
\begin{tikzcd}[column sep=large,row sep=large]
V[T_1,\ldots,T_n] \arrow[r,"j'"] \arrow[d] & Z[T'_1,\ldots,T'_m,T_1,\ldots,T_n] \arrow[d]\\
V \arrow[r,"j"'] & Z[T'_1,\ldots,T'_m],
\end{tikzcd}
\]
where $j'$ is obtained from $j$ by base change. This base change is flat; hence $j'$ is a regular immersion, and so $j'i$ is a regular immersion of $U$ in a scheme smooth over $W$, where $W$ is the relevant open of $Z$. This proves the assertion.

\sgaStableTarget{sga6.e08.prop.1-6}\textbf{Proposition 1.6.} Let $f:X\to Y$ be a morphism and let $g:Y'\to Y$ be a flat base change. If $f$ is complete-intersection, then the morphism
\[
f':X\times_Y Y'\longrightarrow Y'
\]
deduced from $f$ is complete-intersection. The converse is true if, moreover, $g$ is quasi-compact and surjective.

One can find an open covering of $X$ by open subsets $U_i$ such that, for every $i$, the restriction of $f$ to $U_i$ factors as
\[
U_i\longrightarrow V_i\longrightarrow Y,
\]
where $U_i\to V_i$ is a regular immersion and $V_i\to Y$ is smooth. The $U_i\times_Y Y'$ form an open covering of $X\times_Y Y'$, the immersions
\[
U_i\times_Y Y'\longrightarrow V_i\times_Y Y'
\]
are regular by \hyperlink{sga6.e07.prop.1-5}{VII 1.5}, and the morphisms $V_i\times_Y Y'\to Y'$ are smooth, giving the first assertion. For the converse, EGA IV 2.7.1 (iv) first implies that $f$ is locally of finite presentation; this reduces us easily, using \hyperlink{sga6.e08.prop.1-2}{1.2}, to the case where $f$ is a closed immersion. In that case the assertion is just \hyperlink{sga6.e07.prop.1-5}{VII 1.5}.

\sgaStableTarget{sga6.e08.prop.1-7}\textbf{Proposition 1.7.} Let $f:X\to Y$ be a complete-intersection morphism. Then $f$ is a perfect morphism.

Since the property is local on $X$, one may work in a neighbourhood of a point $x\in X$ and suppose that there is an open subset $U$ containing $x$ and a factorization
\[
\begin{tikzcd}[column sep=large,row sep=large]
U \arrow[r,"i"] \arrow[dr,"f"'] & V \arrow[d,"g"]\\
& Y,
\end{tikzcd}
\]
where $i$ is a regular immersion and $g$ is smooth. Since $g$ is smooth, it is a perfect morphism (\hyperlink{sga6.e03.def.4-1}{III 4.1}). By \hyperlink{sga6.e07.prop.1-9}{VII 1.9}, $i$ is also a perfect morphism. Since the composite of two perfect morphisms is perfect (\hyperlink{sga6.e03.prop.4-5}{III 4.5}), the proposition follows.

\sgaStableTarget{sga6.e08.prop.1-8}\textbf{Proposition 1.8.} Let $f:X\to Y$ be a complete-intersection morphism, let $x$ be a point of $X$, and let
\[
\begin{tikzcd}[column sep=large,row sep=large]
U \arrow[r,"i"] \arrow[dr,"f"'] & V \arrow[d,"g"]\\
& Y
\end{tikzcd}
\]
be a factorization in which $U$ is an open neighbourhood of $x$, $i$ is a regular immersion, and $g$ is smooth. Let $J$ be the ideal of $\mathcal O_V$ defining $U$. Then the integer
\[
d_x=\operatorname{rank}\bigl((\Omega^1_{V/Y})_{i(x)}\bigr)-\operatorname{rank}\bigl((J/J^2)_x\bigr)
\]
is independent of the chosen factorization of $f$.

Suppose that the restriction of $f$ to $U$ factors through two smooth $Y$-schemes $V$ and $V'$, and put $V''=V\times_Y V'$. Then $V''$ is smooth over $Y$, and also over $V$ and $V'$. Moreover, the immersions $X\to V$ and $X\to V'$ define an immersion $X\to V''$, which is regular by 1.2. It is therefore enough to compare the integers $d_x$ attached to two factorizations through smooth $Y$-schemes $V$ and $V'$ such that there exists a commutative diagram over $Y$
\[
\begin{tikzcd}[column sep=large,row sep=large]
& V' \arrow[d,"p"]\\
X \arrow[ur,"j"] \arrow[r,"i"'] & V,
\end{tikzcd}
\]
where $p$ is smooth. We may also suppose that $\Omega^1_{V/Y}$, $\Omega^1_{V'/Y}$, and $\Omega^1_{V'/V}$ have constant rank, as do the conormal sheaves of $i$ and $j$.

From the commutative diagram, all morphisms being smooth,
\[
\begin{tikzcd}[column sep=large,row sep=large]
V' \arrow[r,"p"] \arrow[dr] & V \arrow[d]\\
& Y,
\end{tikzcd}
\]
one deduces
\sgaStableTarget{sga6.e08.eq.1-8-1}\begin{equation}\tag{1.8.1}
\operatorname{rank}(\Omega^1_{V'/Y})=
\operatorname{rank}(\Omega^1_{V/Y})+
\operatorname{rank}(\Omega^1_{V'/V}).
\end{equation}
Let $X'=V'\times_V X$; then the rank of $\Omega^1_{V'/V}$ is equal to that of $\Omega^1_{X'/X}$. Further, $j$ defines a section $j':X\to X'$, which is a regular immersion of codimension equal to the rank of $\Omega^1_{X'/X}$ (EGA IV 17.10.4). If $i'$ is the immersion obtained from $i$ by the base change $V'\to V$, then $j=i'j'$, and consequently (\hyperlink{sga6.e07.prop.1-7}{VII 1.7})
\sgaStableTarget{sga6.e08.eq.1-8-2}\begin{equation}\tag{1.8.2}
\operatorname{rank}(N_{V'/X})=
\operatorname{rank}(N_{V/X})+
\operatorname{rank}(\Omega^1_{V'/V}).
\end{equation}
Comparing \hyperlink{sga6.e08.eq.1-8-1}{(1.8.1)} and \hyperlink{sga6.e08.eq.1-8-2}{(1.8.2)} gives the proposition.

\sgaStableTarget{sga6.e08.def.1-9}\textbf{Definition 1.9.} The integer $d_x$ of \hyperlink{sga6.e08.prop.1-8}{Proposition 1.8} is called the \emph{relative virtual dimension} of $X$ over $Y$ (or of $f$) at the point $x$.

The relative virtual dimension of $X$ over $Y$ is therefore a locally constant function on $X$, with integral values, positive or negative. If $f$ is a smooth morphism of relative dimension $n$, its relative virtual dimension is $n$. If $f$ is a regular immersion of codimension $n$, its relative virtual dimension is $-n$.

\sgaStableTarget{sga6.e08.prop.1-10}\textbf{Proposition 1.10.} Let $f:X\to Y$ and $g:Y\to Z$ be two complete-intersection morphisms, and let $x\in X$. Then
\[
\operatorname{dim.rel.virt.}_x(gf)=
\operatorname{dim.rel.virt.}_x(f)+
\operatorname{dim.rel.virt.}_{f(x)}(g).
\]
This follows immediately from \hyperlink{sga6.e08.def.1-9}{1.9} and from the proof of 1.5.

\sgaStableTarget{sga6.e08.sec.2}\subsection*{2. The relative cotangent complex}

\sgaStableTarget{sga6.e08.para.2-1}\textbf{2.1.} We give here a definition of the cotangent complex which, without having the greatest possible generality (see \hyperlink{sga6.e00.sec.4.4}{Expos\'e 0, 4.4)}, will be sufficient for the needs of the Riemann--Roch theorem. Of course the definition given in \hyperlink{sga6.e00.expose}{Expos\'e 0} coincides, in the case considered there, with the definition given here.

Let $f:X\to Y$ be a smoothable morphism, that is, a morphism admitting a factorization
\[
\begin{tikzcd}[column sep=large,row sep=large]
X \arrow[r,"i"] \arrow[dr,"f"'] & V \arrow[d,"g"]\\
& Y,
\end{tikzcd}
\]
where $g$ is smooth and $i$ is an immersion, which may be supposed closed. Let $J$ be the ideal of $X$ in $V$. The canonical derivation
\[
d:\mathcal O_V\longrightarrow\Omega^1_{V/Y}
\]
defines a homomorphism
\[
J\longrightarrow\mathcal O_V\xrightarrow{d}\Omega^1_{V/Y}\longrightarrow
\Omega^1_{V/Y}\otimes_{\mathcal O_V}\mathcal O_X
\]
which vanishes on $J^2$, and hence defines a homomorphism, still denoted $d$,
\sgaStableTarget{sga6.e08.eq.2-1-1}\begin{equation}\tag{2.1.1}
J/J^2\xrightarrow{d} i^*(\Omega^1_{V/Y}).
\end{equation}
One checks easily that this homomorphism is $\mathcal O_X$-linear. One then defines a chain complex $L^{X/Y}_\bullet$ by putting
\[
L^{X/Y}_0=i^*(\Omega^1_{V/Y}),\qquad
L^{X/Y}_1=J/J^2,
\qquad
L^{X/Y}_i=0\quad(i\neq 0,1),
\]
and taking as homomorphism $L^{X/Y}_1\to L^{X/Y}_0$ the homomorphism \hyperlink{sga6.e08.eq.2-1-1}{(2.1.1)}.

\sgaStableTarget{sga6.e08.prop.2-2}\textbf{Proposition 2.2.} The complex $L^{X/Y}_\bullet$ is independent, up to canonical isomorphism in the derived category $D(X)$, of the factorization chosen through a smooth $Y$-scheme.

Let $V,V'$ be two smooth $Y$-schemes through which $f$ factors. Put $V''=V\times_Y V'$. Then $V''$ is smooth over $Y$, and also over $V$ and $V'$. Moreover the immersions $X\to V$ and $X\to V'$ define an immersion $X\to V''$ which factors $f$. It is therefore enough to compare the complexes $L^{X/Y}_\bullet$ attached to two factorizations through smooth $Y$-schemes $V$ and $V'$ for which there is a commutative diagram over $Y$
\[
\begin{tikzcd}[column sep=large,row sep=large]
& V' \arrow[d,"p"]\\
X \arrow[ur,"j"] \arrow[r,"i"'] & V,
\end{tikzcd}
\]
where $i$ and $j$ are closed immersions, and $p$ is smooth. Functoriality then gives homomorphisms
\[
I/I^2\longrightarrow J/J^2,
\qquad
p^*(\Omega^1_{V/Y})\longrightarrow\Omega^1_{V'/Y},
\]
hence a diagram
\sgaStableTarget{sga6.e08.eq.2-2-1}\begin{equation}\tag{2.2.1}
\begin{tikzcd}[column sep=large,row sep=large]
I/I^2 \arrow[r] \arrow[d,"d"'] & J/J^2 \arrow[d,"d"]\\
i^*(\Omega^1_{V/Y}) \arrow[r] & j^*(\Omega^1_{V'/Y}).
\end{tikzcd}
\end{equation}
Since $V$ and $V'$ are smooth over $Y$, one obtains the exact sequence
\[
0\longrightarrow p^*(\Omega^1_{V/Y})\longrightarrow\Omega^1_{V'/Y}
\longrightarrow\Omega^1_{V'/V}\longrightarrow0.
\]
Since the modules of differentials considered are locally free on $\mathcal O_V$, applying the functor $j^*$ gives the exact sequence
\sgaStableTarget{sga6.e08.eq.2-2-2}\begin{equation}\tag{2.2.2}
0\longrightarrow i^*(\Omega^1_{V/Y})\longrightarrow j^*(\Omega^1_{V'/Y})
\longrightarrow j^*(\Omega^1_{V'/V})\longrightarrow0.
\end{equation}
Furthermore, $j$ defines a section $j'$ of $X'=V'\times_V X$:
\[
\begin{tikzcd}[column sep=large,row sep=large]
X' \arrow[r,"i'"] \arrow[d,"p'"'] & V' \arrow[d,"p"]\\
X \arrow[r,"i"'] \arrow[ur,"j"] & V .
\end{tikzcd}
\]
Let $I,J,K$ be the ideals defining $i,j,j'$. Then $j=i'j'$, and the immersion $j'$ is regular; moreover $i'$ is defined by the ideal $p^*(I)$, because $p$ is flat, and its conormal sheaf is therefore
\[
i'^*(p^*(I))=p'^*(I/I^2).
\]
One obtains the exact sequence
\sgaStableTarget{sga6.e08.eq.2-2-3}\begin{equation}\tag{2.2.3}
0\longrightarrow I/I^2\longrightarrow J/J^2\longrightarrow K/K^2\longrightarrow0.
\end{equation}
We shall show: (a) the diagram \hyperlink{sga6.e08.eq.2-2-1}{(2.2.1)} is commutative; this will give a homomorphism of complexes $\theta$ from the complex $L^{X/Y}_\bullet$ defined by $V$ to the complex $L^{X/Y}_\bullet$ defined by $V'$; (b) the homomorphism
\[
K/K^2\longrightarrow j^*(\Omega^1_{V'/V})
\]
defined by \hyperlink{sga6.e08.eq.2-2-2}{(2.2.2)}, \hyperlink{sga6.e08.eq.2-2-3}{(2.2.3)}, and \hyperlink{sga6.e08.eq.2-2-1}{(2.2.1)} is an isomorphism. This will imply that the mapping cylinder of $\theta$ is acyclic, hence that $\theta$ is a quasi-isomorphism, and therefore defines an isomorphism in the derived category.

The properties (a) and (b) are local on $X$. One may therefore suppose $X,Y,V,V'$ affine. Put $Y=\Spec(A)$, $V=\Spec(B)$, $V'=\Spec(B')$, $X=\Spec(C)$, and keep the same letters for the corresponding homomorphisms and ideals. One has the diagram
\[
\begin{tikzcd}[column sep=large,row sep=large]
B\otimes_A B \arrow[r,"p\otimes p"] \arrow[d,"\Delta"'] & B'\otimes_A B' \arrow[d,"\Delta"] & B'\otimes_B B' \arrow[l] \arrow[d]\\
B \arrow[r,"p"] \arrow[dr,"i"'] & B' \arrow[d,"j"] & B'\otimes_B C = B'/IB' \arrow[dl,"p'" ]\\
{} & C & {}
\end{tikzcd}
\]
where the left and right displayed squares are denoted (1) and (2), respectively.
The homomorphism $I/I^2\to J/J^2$ is deduced from the homomorphism $IB'\to J$ and coincides with that deduced from $p:I\to J$; the homomorphism
\[
i^*(\Omega^1_{V/Y})\longrightarrow j^*(\Omega^1_{V'/Y})
\]
is deduced from the commutative square (1). Hence property (a) is evident.

The homomorphism $K/K^2\to j^*(\Omega^1_{V'/V})$ considered in (b) may be defined as follows. Let $\overline x\in K/K^2$, and let $x$ be an element of $J$ lifting $\overline x$ (such an element exists because $K=J/IB'$). To $\overline x$ one associates the image of $x$ by the composite
\[
J\xrightarrow{d}\Omega^1_{B'/A}\otimes_{B'}C
\longrightarrow\Omega^1_{B'/B}\otimes_{B'}C.
\]
On the other hand, let $H$ be the augmentation ideal of
\[
B'\otimes_B B'\longrightarrow B'.
\]
The commutative square (2) gives a homomorphism $H\to K$, hence a homomorphism
\[
H/H^2\otimes_{B'}C\longrightarrow K/K^2.
\]
But $H/H^2=\Omega^1_{B'/B}$. It remains only to check that the homomorphisms thus constructed are inverse to each other, which is straightforward.

\sgaStableTarget{sga6.e08.def.2-3}\textbf{Definition 2.3.} Let $f:X\to Y$ be a smoothable morphism. The complex $L^{X/Y}_\bullet$, defined up to canonical isomorphism in the derived category $D(X)$, is called the \emph{relative cotangent complex} of $X$ over $Y$ (or of $f$).

\sgaStableTarget{sga6.e08.prop.2-4}\textbf{Proposition 2.4.} Let $f:X\to Y$ be a smoothable complete-intersection morphism. Then $L^{X/Y}_\bullet$ is perfect. More precisely, for every factorization of $f$ as the composite of an immersion and a smooth morphism, the complex $L^{X/Y}_\bullet$ defined in \hyperlink{sga6.e08.para.2-1}{2.1} is strictly perfect.

Let $f=g i$, where $g$ is smooth and $i$ is an immersion. By \hyperlink{sga6.e08.prop.1-2}{1.2}, $i$ is a regular immersion. If $J$ is the ideal defining $i$, then $J/J^2$ is a locally free finite-type $\mathcal O_X$-module. Likewise, since $g:V\to Y$ is smooth, $i^*(\Omega^1_{V/Y})$ is locally free of finite type. This proves 2.4.

\sgaStableTarget{sga6.e08.cor.2-5}\textbf{Corollary 2.5.} Let $f:X\to Y$ be a smoothable complete-intersection morphism, and let $f=gi$ be a factorization of $f$ through a smooth $Y$-scheme $V$. Let $J$ be the ideal of $\mathcal O_V$ defining $i$. If $K^{\bullet}(X)$ denotes the Grothendieck group of locally free finite-type $\mathcal O_X$-modules (respectively of perfect complexes on $X$), the element $T_f\in K^{\bullet}(X)$ defined by
\[
T_f=\operatorname{cl}(i^*\Omega^1_{V/Y})-\operatorname{cl}(J/J^2)
\]
respectively
\[
T_f=\operatorname{cl}(L^{X/Y}_\bullet)
\]
does not depend on the chosen factorization of $f$. The value at a point $x$ of $X$ of the augmentation of $T_f$ is precisely the relative virtual dimension of $X$ over $Y$ at $x$ \hyperlink{sga6.e08.def.1-9}{(1.9)}.

The element $T_f$ is well defined in both cases by \hyperlink{sga6.e08.prop.2-4}{2.4}; its class in $K^{\bullet}(X)$ is independent of the factorization, by \hyperlink{sga6.e08.prop.2-2}{2.2} in the second case and by the proof of \hyperlink{sga6.e08.prop.2-2}{2.2} in the first. The last assertion is trivial.

\sgaStableTarget{sga6.e08.prop.2-6}\textbf{Proposition 2.6.} Let $f:X\to Y$ and $g:Y\to Z$ be two smoothable complete-intersection morphisms, and suppose in addition that the following condition is satisfied:

(G) There exists a smooth $Y$-scheme $V$ factorizing $f$ such that, for every closed immersion $Y\to V'$ into a smooth $Z$-scheme, one can find a smooth $V'$-scheme $V''$ with
\[
V \simeq V''\times_{V'}Y.
\]
This condition is in particular satisfied if $f$ is quasi-projective and if $Y$ has an ample sheaf, by taking for $V$ some $\mathbb P^r_Y$.

Then $gf$ is complete-intersection and smoothable, and there is a distinguished triangle
\[
\begin{tikzcd}[column sep=large,row sep=large]
& L^{X/Y}_\bullet \arrow[dl] \\
f^*L^{Y/Z}_\bullet \arrow[rr] && L^{X/Z}_\bullet \arrow[ul]
\end{tikzcd}
\]
in the derived category.

Let $V'$ be a smooth $Z$-scheme factorizing $g$. Consider the commutative diagram
\[
\begin{tikzcd}[column sep=large,row sep=large]
X \arrow[r,"i"] \arrow[dr,"f"'] & V \arrow[r,"j'"] \arrow[d,"p"] & V'' \arrow[d,"p''"]\\
& Y \arrow[r,"j"] \arrow[dr,"g"'] & V' \arrow[d,"p'"]\\
& & Z .
\end{tikzcd}
\]
Then the immersion $j'$ is regular, since it is obtained from $j$ by a flat base change. Hence $j'i$ is a regular immersion factorizing $gf$, which is therefore smoothable and complete-intersection. Moreover, one has the exact sequence of conormal sheaves
\[
0\longrightarrow i^*(J'/J'^2)\longrightarrow K/K^2\longrightarrow I/I^2\longrightarrow0,
\]
where $I,J,J',K$ denote the sheaves defining respectively the immersions $i,j,j',j'i$. Since $p''$ is flat, $J'=p''^*(J)$, and the sequence can be written
\sgaStableTarget{sga6.e08.eq.2-6-1}\begin{equation}\tag{2.6.1}
0\longrightarrow f^*(J/J^2)\longrightarrow K/K^2\longrightarrow I/I^2\longrightarrow0.
\end{equation}
Furthermore, since $V'$ and $V''$ are smooth over $Z$, there is an exact sequence
\[
0\longrightarrow p''^*(\Omega^1_{V'/Z})\longrightarrow\Omega^1_{V''/Z}\longrightarrow\Omega^1_{V''/V'}\longrightarrow0.
\]
Applying $(j'i)^*$, one obtains an exact sequence
\[
0\longrightarrow (p''ji)^*(\Omega^1_{V'/Z})\longrightarrow (j'i)^*(\Omega^1_{V''/Z})\longrightarrow (j'i)^*(\Omega^1_{V''/V'})\longrightarrow0.
\]
Since
\[
j'^*(\Omega^1_{V''/V'})=\Omega^1_{V/Y}
\]
by hypothesis (G), one obtains
\sgaStableTarget{sga6.e08.eq.2-6-2}\begin{equation}\tag{2.6.2}
0\longrightarrow f^*(j^*\Omega^1_{V'/Z})\longrightarrow (j'i)^*(\Omega^1_{V''/Z})\longrightarrow i^*(\Omega^1_{V/Y})\longrightarrow0.
\end{equation}
Combining \hyperlink{sga6.e08.eq.2-6-1}{(2.6.1)} and \hyperlink{sga6.e08.eq.2-6-2}{(2.6.2)}, one is reduced to checking that the diagram
\[
\begin{tikzcd}[column sep=large,row sep=large]
0 \arrow[r] & f^*(J/J^2) \arrow[r] \arrow[d,"f^*(d)"'] & K/K^2 \arrow[r] \arrow[d,"d"] & I/I^2 \arrow[r] \arrow[d,"d"] & 0\\
0 \arrow[r] & f^*(j^*\Omega^1_{V'/Z}) \arrow[r] & (j'i)^*\Omega^1_{V''/Z} \arrow[r] & i^*\Omega^1_{V/Y} \arrow[r] & 0
\end{tikzcd}
\]
is commutative. This is verified locally without difficulty by returning to the definitions of the homomorphisms involved. The announced distinguished triangle follows from the resulting exact sequence of complexes.

\sgaStableTarget{sga6.e08.cor.2-7}\textbf{Corollary 2.7.} With the notation and hypotheses of \hyperlink{sga6.e08.prop.2-6}{2.6}, one has in $K^{\bullet}(X)$ the equality
\[
T_{gf}=T_f+f^*(T_g).
\]

\sgaStableTarget{sga6.e08.sec.3}\subsection*{3. Riemann--Roch theorem: statement}

\sgaStableTarget{sga6.e08.para.3-1}\textbf{3.1.} In no. 3, $Y$ is a quasi-compact scheme having an ample sheaf.

Let $f:X\to Y$ be a projective complete-intersection morphism \hyperlink{sga6.e08.def.1-1}{(1.1)}; suppose moreover that the relative virtual dimension of $f$ \hyperlink{sga6.e08.def.1-9}{(1.9)} is constant on $X$. Since $f$ is projective, there is a closed immersion $X\to P$, where $P$ is a projective bundle over $Y$; since $Y$ has an ample sheaf, one may suppose that $P=\mathbb P^r_Y$. Thus there is a factorization
\[
\begin{tikzcd}[column sep=large,row sep=large]
X \arrow[r,"i"] \arrow[dr,"f"'] & \mathbb P^r_Y \arrow[d,"g"]\\
& Y.
\end{tikzcd}
\]
Since $f$ is complete-intersection, the immersion $i$ is regular \hyperlink{sga6.e08.prop.1-2}{(1.2)}.

By \hyperlink{sga6.e08.prop.1-7}{1.7}, $f$ is a perfect morphism. Moreover $Y$ has an ample sheaf, and therefore so do $X$ and $P$. One can therefore define an additive homomorphism
\[
f_*:K^{\bullet}(X)\longrightarrow K^{\bullet}(Y)
\]
(\hyperlink{sga6.e04.para.2-12}{IV 2.12}), the groups $K^{\bullet}$ defined in terms of locally free sheaves of finite type, or of perfect complexes, being identical. One can also define additive homomorphisms
\[
i_*:K^{\bullet}(X)\longrightarrow K^{\bullet}(P),
\qquad
 g_*:K^{\bullet}(P)\longrightarrow K^{\bullet}(Y),
\]
and one has $f_*=g_*i_*$.

Let $d$ be the relative virtual dimension of $X$ over $Y$. The codimension of $i$ is then $r-d$. By \hyperlink{sga6.e07.thm.4-6}{VII 4.6}, the homomorphism $i_*$ gives, after tensoring with $\mathbb Q$, inclusions
\[
i_*\bigl(\Fil^kK^{\bullet}(X)_{\mathbb Q}\bigr)
\subset
\Fil^{k+r-d}K^{\bullet}(P)_{\mathbb Q}.
\]
By \hyperlink{sga6.e06.cor.5-8}{VI 5.8}, the homomorphism $g_*$ gives inclusions
\[
g_*\bigl(\Fil^{k'}K^{\bullet}(P)\bigr)
\subset
\Fil^{k'-r}K^{\bullet}(Y).
\]
Composing, one obtains the following.

\sgaStableTarget{sga6.e08.prop.3-2}\textbf{Proposition 3.2.} Let $f:X\to Y$ be a projective complete-intersection morphism whose relative virtual dimension $d$ is constant. Suppose that $Y$ is quasi-compact and has an ample sheaf. Then for every $k\in\mathbb Z$,
\[
f_*\bigl(\Fil^kK^{\bullet}(X)_{\mathbb Q}\bigr)
\subset
\Fil^{k-d}K^{\bullet}(Y)_{\mathbb Q}.
\]
Consequently one obtains an additive homomorphism
\[
f^{\mathrm{gr}}_*:\Gr K^{\bullet}(X)_{\mathbb Q}\longrightarrow
\Gr K^{\bullet}(Y)_{\mathbb Q}
\]
sending $\Gr^kK^{\bullet}(X)_{\mathbb Q}$ into $\Gr^{k-d}K^{\bullet}(Y)_{\mathbb Q}$.

\sgaStableTarget{sga6.e08.para.3-3}\textbf{3.3.} Recall that in \hyperlink{sga6.e05.para.6-1}{V 6.1} we defined a functor
\[
A\longmapsto \operatorname{Ch}(A)
\]
from the category of graded $K$-algebras to the category of augmented $K$-$\lambda$-algebras, where $K$ is a binomial ring (\hyperlink{sga6.e05.def.2-7}{V 2.7}). If $A$ is a graded $K$-algebra, then
\[
\operatorname{Ch}(A)=K\times(1+\widehat A^+).
\]
For every augmented $K$-$\lambda$-algebra $\Lambda$ we defined a $K$-$\lambda$-homomorphism, called the completed Chern class (\hyperlink{sga6.e05.para.6-7}{V 6.7}),
\[
\widetilde c=(\varepsilon,c):\Lambda\longrightarrow\operatorname{Ch}(\Gr_{\Lambda}).
\]
Finally, we defined a ring homomorphism, functorial in $A$, called the Chern character (\hyperlink{sga6.e05.def.6-4}{V 6.4}),
\[
\ch:\operatorname{Ch}(A)\longrightarrow K\oplus\widehat A^+_{\mathbb Q},
\]
and a group endomorphism, called the Todd character (\hyperlink{sga6.e05.para.1-16}{V 1.16}),
\[
\Todd:1+\widehat A^+\longrightarrow 1+\widehat A^+.
\]
We shall apply these definitions to the augmented $\lambda$-rings $K^{\bullet}(X)$ and $K^{\bullet}(Y)$; the augmentation rings $H^0(X,\mathbb Z)$ and $H^0(Y,\mathbb Z)$ are binomial (\hyperlink{sga6.e05.ex.2-8}{V 2.8}). To simplify notation, for every scheme $Z$ and every $z\in K^{\bullet}(Z)$ we write
\[
\ch(z)=\ch(\widetilde c(z)),
\qquad
\Todd(z)=\Todd(c(z)).
\]
For every graded algebra $A$ we regard the group $1+\widehat A^+$ as contained in $\widehat A$. Finally, if $f:A\to B$ is additive and compatible with the gradings up to a shift, we shall still write $f$ for the homomorphism $\widehat A\to\widehat B$ deduced from it.

\sgaStableTarget{sga6.e08.para.3-4}\textbf{3.4.} Let $X$ be a scheme. For every locally free finite-type $\mathcal O_X$-module $E$, denote by $\check E$ the class in $K^{\bullet}(X)$ of the dual $\check E$ of $E$. It is clear that the function $E\mapsto\check E$ is additive on the category of locally free finite-type $\mathcal O_X$-modules, and therefore defines an endomorphism of $K^{\bullet}(X)$, denoted $x\mapsto\check x$.

\sgaStableTarget{sga6.e08.prop.3-5}\textbf{Proposition 3.5.} Let $X$ be a quasi-compact scheme. Then, for all $x\in K^{\bullet}(X)$ and all $i\geq0$, one has
\[
c^i(\check x)=(-1)^i c^i(x).
\]
Equivalently,
\[
c_t(\check x)=c_{-t}(x),
\qquad
c_t(x)=\sum_{i\geq0}c^i(x)t^i.
\]
By additivity one may suppose that $x$ is the class of a locally free finite-type $\mathcal O_X$-module $E$. Moreover, after replacing $X$ by the flag scheme of $E$, which gives an injection on the groups $\Gr K^{\bullet}$ (\hyperlink{sga6.e06.cor.5-5}{VI 5.5}), one may suppose that $E$ is an invertible sheaf. Then $\check E=E^{-1}$, so $\check x=x^{-1}$.

The $c^i(\check x)$ and $c^j(x)$ are zero for $i,j\geq2$. It is enough to prove
\[
c^1(x^{-1})=-c^1(x).
\]
Now $c^1(x)$ is the class in $\Gr^1K^{\bullet}(X)$ of $x-1$, while $c^1(x^{-1})$ is that of $x^{-1}-1$. The equality
\[
x^{-1}-1=1-x+(x-1)(1-x^{-1})
\]
gives
\[
x^{-1}-1=1-x\pmod{\Fil^2K^{\bullet}(X)},
\]
and the result follows.

\sgaStableTarget{sga6.e08.thm.3-6}\textbf{Theorem 3.6 (Riemann--Roch).} Let $f:X\to Y$ be a projective complete-intersection morphism \hyperlink{sga6.e08.def.1-1}{(1.1)}, whose relative virtual dimension $d$ is constant \hyperlink{sga6.e08.def.1-9}{(1.9)}. Suppose that $Y$ is quasi-compact and admits an ample sheaf. With the notation of \hyperlink{sga6.e08.para.3-3}{3.3}, for every $x\in K^{\bullet}(X)$ one has
\[
\ch(f_*(x))=f^{\mathrm{gr}}_*\bigl(\ch(x)\Todd(\check T_f)\bigr),
\]
where $\check T_f$ denotes the dual element of the class of the relative cotangent complex $L^{X/Y}_\bullet$ (\hyperlink{sga6.e08.cor.2-5}{2.5} and \hyperlink{sga6.e08.def.2-3}{2.3}).

In other words, the diagram
\[
\begin{tikzcd}[column sep=huge,row sep=large]
K^{\bullet}(X) \arrow[r,"\ch\circ\widetilde c"] \arrow[d,"f_*"'] & \Gr K^{\bullet}(X)_{\mathbb Q} \arrow[d,"x\mapsto f^{\mathrm{gr}}_*(x\Todd(\check T_f))"]\\
K^{\bullet}(Y) \arrow[r,"\ch\circ\widetilde c"'] & \Gr K^{\bullet}(Y)_{\mathbb Q}
\end{tikzcd}
\]
is commutative.



\sgaStableTarget{sga6.e08.continuation-of-the-proof-of-theorem-3-6}\subsection*{Continuation of the proof of Theorem 3.6}

The proof of \hyperlink{sga6.e08.thm.3-6}{Theorem 3.6} will occupy the end of this exposé. Let us give some preliminaries.

\sgaStableTarget{sga6.e08.prop.3-7}\textbf{Proposition 3.7.} For arbitrary
\[
x\in \Gr K^{\bullet}(X)_{\mathbb Q},\qquad
y\in \Gr K^{\bullet}(Y)_{\mathbb Q},
\]
one has
\[
 f^{\mathrm{gr}}_*\bigl(x\cdot f^{\mathrm{gr}}(y)\bigr)
 =
 f^{\mathrm{gr}}_*(x)\,y .
\]

By additivity, one may suppose \(x\) and \(y\) homogeneous; say
\[
x\in \Gr^iK^{\bullet}(X)_{\mathbb Q},\qquad
y\in \Gr^jK^{\bullet}(Y)_{\mathbb Q}.
\]
Let
\[
x'\in \Fil^iK^{\bullet}(X)_{\mathbb Q},\qquad
y'\in \Fil^jK^{\bullet}(Y)_{\mathbb Q}
\]
be elements lifting \(x\) and \(y\). Then the projection formula gives
\[
 f_*\bigl(x'\cdot f^*(y')\bigr)=f_*(x')\,y'.
\]
By \hyperlink{sga6.e08.prop.3-2}{3.2}, this equality is written in
\(\Fil^{i+j-d}K^{\bullet}(Y)_{\mathbb Q}\). One deduces from it an equality in
\(\Gr^{i+j-d}K^{\bullet}(Y)_{\mathbb Q}\), and one verifies immediately that this is the required equality.

\sgaStableTarget{sga6.e08.lem.3-8}\textbf{Lemma 3.8.} Let \(f:X\to Y\) and \(g:Y\to Z\) be two morphisms satisfying the hypotheses of 3.6. Then \(g\circ f\) also satisfies them. If \(x\) is an element of \(K^{\bullet}(X)\), and if \hyperlink{sga6.e08.thm.3-6}{3.6} is verified by \(f\) for the element \(x\), and by \(g\) for the element \(f_*(x)\), then it is verified by \(g\circ f\) for the element \(x\).

By hypothesis, one has
\sgaStableTarget{sga6.e08.eq.3-8-1}\begin{equation}\tag{3.8.1}
 \ch(f_*(x))=
 f^{\mathrm{gr}}_*\bigl(\ch(x)\Todd(\check T_f)\bigr),
\end{equation}
and
\sgaStableTarget{sga6.e08.eq.3-8-2}\begin{equation}\tag{3.8.2}
 \ch(g_*(f_*(x)))=
 g^{\mathrm{gr}}_*\bigl(\ch(f_*(x))\Todd(\check T_g)\bigr).
\end{equation}
Combining \hyperlink{sga6.e08.eq.3-8-1}{(3.8.1)} and \hyperlink{sga6.e08.eq.3-8-2}{(3.8.2)}, one obtains
\[
\ch(g_*(f_*(x)))=
g^{\mathrm{gr}}_*\!\left(
 f^{\mathrm{gr}}_*\bigl(\ch(x)\Todd(\check T_f)\bigr)\cdot \Todd(\check T_g)
\right).
\]
Applying \hyperlink{sga6.e08.prop.3-7}{3.7}, one finds
\[
\ch(g_*(f_*(x)))=
g^{\mathrm{gr}}_*\!\left(
 f^{\mathrm{gr}}_*\bigl(\ch(x)\Todd(\check T_f)\cdot
 f^{\mathrm{gr}}(\Todd(\check T_g))\bigr)
\right).
\]
The functoriality of \(c\) and of \(\Todd\) implies
\[
 f^{\mathrm{gr}}(\Todd(\check T_g))=\Todd(f^*(\check T_g)).
\]
By \hyperlink{sga6.e08.cor.2-7}{2.7}, one has
\[
 \check T_{g\circ f}=\check T_f+f^*(\check T_g).
\]
The multiplicativity of the Todd character therefore gives
\[
 \Todd(\check T_f)\,f^{\mathrm{gr}}(\Todd(\check T_g))
 =
 \Todd(\check T_{g\circ f}).
\]
Thus finally
\[
 \ch(g_*(f_*(x)))=
 g^{\mathrm{gr}}_* f^{\mathrm{gr}}_*\bigl(\ch(x)\Todd(\check T_{g\circ f})\bigr),
\]
which is the announced result.

\sgaStableTarget{sga6.e08.para.3-9}\textbf{3.9.} If one applies \hyperlink{sga6.e08.lem.3-8}{3.8} to the factorization of \(f\) given in \hyperlink{sga6.e08.para.3-1}{3.1}, one sees that, to prove the Riemann--Roch theorem, it is enough to prove it in the case of a regular immersion of constant codimension, and in the case of the structural morphism
\[
 \mathbb P^r_Y\longrightarrow Y .
\]
Moreover, one may suppose that the immersion \(i\) satisfies the condition of \hyperlink{sga6.e07.lem.4-1}{VII 4.1}. Indeed, let \(s\) be the section of
\[
 \mathbb P^2_{\mathbb P^r_Y}
\]
introduced in \hyperlink{sga6.e07.sec.4.2}{VII 4.2}. Then \(s\circ i\) is a regular immersion satisfying the condition of \hyperlink{sga6.e07.lem.4-1}{VII 4.1} (\hyperlink{sga6.e07.sec.4.2}{VII 4.2}), and by \hyperlink{sga6.e08.lem.3-8}{3.8} it is enough to prove \hyperlink{sga6.e08.thm.3-6}{3.6} for
\[
s\circ i,\qquad
\mathbb P^2_{\mathbb P^r_Y}\longrightarrow\mathbb P^r_Y,\qquad
\mathbb P^r_Y\longrightarrow Y .
\]

We shall therefore prove \hyperlink{sga6.e08.thm.3-6}{3.6} by proving the following two particular cases:
\begin{enumerate}[label=\roman*)]
\item \(f\) is a regular immersion of constant codimension, satisfying the condition of \hyperlink{sga6.e07.lem.4-1}{VII 4.1};
\item \(f\) is the structural morphism \(\mathbb P^r_Y\to Y\).
\end{enumerate}

\sgaStableTarget{sga6.e08.sec.4}\section*{4. Riemann--Roch theorem: case of a regular closed immersion}

\sgaStableTarget{sga6.e08.para.4-1}\textbf{4.1.} Let \(i:Y\to X\) be a regular closed immersion. The Riemann--Roch theorem for \(i\) is essentially relation \hyperlink{sga6.e07.cor.4-4}{VII 4.4}, which makes it possible to compute the Chern classes of \(i_*(x)\) from the \(\gamma^p(N,x)\). More precisely, if one introduces the \(\lambda\)-ring \(K^{\bullet}(Y)_{\mathbb Q N}\) defined in \hyperlink{sga6.e05.para.5-5}{V 5.5}, then the homomorphism
\[
K^{\bullet}(Y)_{\mathbb Q N}\longrightarrow K^{\bullet}(X)_{\mathbb Q}
\]
deduced from \(i_*\) is a \(\lambda\)-homomorphism by \hyperlink{sga6.e07.cor.2-8}{VII 2.8} and \hyperlink{sga6.e07.thm.4-3}{VII 4.3}. The Riemann--Roch theorem is then reduced to an analogous formula concerning the additive homomorphism
\[
K^{\bullet}(Y)\longrightarrow K^{\bullet}(Y)_{\mathbb Q N}
\]
defined by \(y\mapsto(0,y)\), that is, to a formal relation valid for every \(\lambda\)-ring. We shall not develop this proof here, in order to avoid introducing still more formalism on \(\lambda\)-rings (cf. \(\RRR\), I 1.5), and we shall resume the method of [1], which used blow-up schemes. This is the method already employed for the proof of \hyperlink{sga6.e07.thm.4-3}{VII 4.3}, and also the method at the end of \(\RRR\).

As indicated in \hyperlink{sga6.e08.para.3-9}{3.9}, one may restrict to the case of a regular immersion of constant codimension \(d\) satisfying the condition of \hyperlink{sga6.e07.lem.4-1}{VII 4.1}; the case of a general regular immersion is then deduced from 3.8. In this paragraph we resume the notation of \hyperlink{sga6.e07.para.3-1}{VII 3.1}. Thus one has the Cartesian square
\[
\begin{tikzcd}[column sep=large,row sep=large]
Y' \arrow[r,"j"] \arrow[d,"g"'] & X' \arrow[d,"f"]\\
Y \arrow[r,"i"'] & X ,
\end{tikzcd}
\]
where \(X'\) is the \(X\)-scheme obtained by blowing up \(Y\). We suppose moreover that \(X\) is quasi-compact and has an ample sheaf.

\sgaStableTarget{sga6.e08.prop.4-2}\textbf{Proposition 4.2.} The morphism \(f\) is projective, a complete-intersection morphism, and of virtual relative dimension zero. In addition one has the relation
\[
 f^{\mathrm{gr}}_* f^{\mathrm{gr}}=\operatorname{Id}_{\Gr K^{\bullet}(X)_{\mathbb Q}} .
\]
The first assertion follows immediately from \hyperlink{sga6.e07.prop.1-8}{VII 1.8}. The second relation is a consequence of the relation
\[
 f_*f^*=\operatorname{Id}_{K^{\bullet}(X)}
\]
proved in \hyperlink{sga6.e07.prop.3-6}{VII 3.6}.

\sgaStableTarget{sga6.e08.prop.4-3}\textbf{Proposition 4.3.} For every \(y\in K^{\bullet}(Y)_{\mathbb Q}\), one has the relation
\[
 f^{\mathrm{gr}}_* i^{\mathrm{gr}}_*(y)
 =
 j^{\mathrm{gr}}_*\bigl(g^{\mathrm{gr}}(y)\,c^{d-1}(\check F)\bigr),
\]
where \(c^{d-1}\) denotes the \((d-1)\)-st Chern class.

This relation is contained in \hyperlink{sga6.e07.thm.4-8}{VII 4.8}, for it is nothing other than
\[
 v^{\mathrm{gr}}\circ u^{\mathrm{gr}}=0 .
\]

\sgaStableTarget{sga6.e08.prop.4-4}\textbf{Proposition 4.4.} Let \(X\) be a quasi-compact scheme. If \(E\) is a locally free sheaf of finite type on \(X\), of class \(E\) in \(K^{\bullet}(X)\), and of rank \(n\), then
\[
 \ch(\lambda_{-1}(E))=c^n(\check E)\Todd(\check E)^{-1}.
\]

The two sides of the relation to be proved are multiplicative. This is clear for \(\ch\circ\lambda_{-1}\) and for \(\Todd\); it is also true for \(c^n\), which is the class of \((-1)^nE\lambda_{-1}(E)\). Consequently, by passing to the flag bundle of \(E\), one may suppose that \(E\) is a sum of classes of invertible sheaves, and reduce to proving the formula for an invertible sheaf. Let \(\xi\) be the class of such a sheaf. Then
\[
 \ch(\lambda_{-1}(\xi))=\ch(1-\xi)=1-\ch(\xi).
\]
Since \(\widetilde c(\xi)=(1,1+c^1(\xi))\), the definition of \(\ch\) gives
\[
 \ch(\lambda_{-1}(\xi))=1-\exp(c^1(\xi)).
\]
Moreover, \(\xi\) has augmentation \(1\); the Chern class which intervenes in formula \hyperlink{sga6.e08.prop.4-4}{4.4} relative to \(\xi\) is therefore \(c^1\). Taking into account the definition of \(\Todd\), and \hyperlink{sga6.e08.prop.3-5}{3.5}, the relation to be proved is written
\[
 1-\exp(c^1(\xi))
 =
 -c^1(\xi)\left(\frac{-c^1(\xi)}{1-\exp(c^1(\xi))}\right)^{-1},
\]
and is therefore verified.

\sgaStableTarget{sga6.e08.rem.4-4-1}\textbf{Remark 4.4.1.} It is also possible to deduce \hyperlink{sga6.e08.prop.4-4}{4.4} from a corresponding purely algebraic statement; cf. \(\RRR\), \hyperlink{sga6.e01.sec.1.3}{I 1.3}.

\sgaStableTarget{sga6.e08.para.4-5}\textbf{4.5.} We shall now show that, in order to prove the Riemann--Roch theorem in the case considered in this paragraph, one may reduce to the case where \(Y\) is a divisor of \(X\), and where the element to which the formula is applied is the inverse image of an element of \(K^{\bullet}(X)\).

By \hyperlink{sga6.e08.prop.3-2}{3.2}, the formula to be proved,
\[
 \ch(i_*(x))=
 i^{\mathrm{gr}}_*\bigl(\ch(x)\Todd(\check T_i)\bigr),
\]
is equivalent to the formula
\[
 f^{\mathrm{gr}}\bigl(\ch(i_*(x))\bigr)
 =
 f^{\mathrm{gr}}i^{\mathrm{gr}}_*\bigl(\ch(x)\Todd(\check T_i)\bigr).
\]
Taking account of the functoriality of \(\ch\), and of \hyperlink{sga6.e08.prop.4-3}{4.3}, it may also be written
\[
 \ch\bigl(f^*(i_*(x))\bigr)
 =
 j^{\mathrm{gr}}_*\bigl(\ch(g^*(x))\Todd(g^*(\check T_i))\,c^{d-1}(\check F)\bigr),
\]
or, by \hyperlink{sga6.e07.prop.3-4}{VII 3.4},
\[
 \ch\bigl(j_*(g^*(x)\lambda_{-1}(F))\bigr)
 =
 j^{\mathrm{gr}}_*\bigl(\ch(g^*(y))\Todd(g^*(\check T_i))\,c^{d-1}(\check F)\bigr).
\]

By definition,
\[
 T_i=-c^1(J/J^2)=-N,\qquad
 T_j=-c^1(J'/J'^2)=-L.
\]
Moreover, the exact sequence
\[
0\longrightarrow F\longrightarrow g^*(N)\longrightarrow \mathcal O_{Y'}(1)\longrightarrow0
\]
gives the relation
\[
 g^*(T_i)=-(F+L),
\]
hence
\[
 \Todd(g^*(\check T_i))=\Todd(\check T_j)\Todd(\check F)^{-1}.
\]
Taking account of \hyperlink{sga6.e08.prop.4-4}{4.4}, the relation to be proved is therefore written
\[
 \ch\bigl(j_*(g^*(x)\lambda_{-1}(F))\bigr)
 =
 j^{\mathrm{gr}}_*\bigl(\ch(g^*(x)\lambda_{-1}(F))\Todd(\check T_j)\bigr).
\]
Putting
\[
z=g^*(x)\lambda_{-1}(F),
\]
this relation is nothing but the Riemann--Roch theorem applied to the regular immersion of codimension \(1\), namely \(j\), and to the element \(z\in K^{\bullet}(Y')\).

Furthermore, since \(i\) was supposed to satisfy the hypothesis of \hyperlink{sga6.e07.lem.4-1}{VII 4.1}, \(\lambda_{-1}(F)\) is divisible by \(1-L\). Thus one can find \(z'\) such that \(z=z'(1-L)\), or still
\[
 z=j^*(j_*(z'))
\]
(\hyperlink{sga6.e07.para.2-6}{Exp. VII 2.7)}. We are therefore reduced to the case where \(Y\) is a divisor, and to an element \(x\) for which there exists \(y\), with \(x=i^*(y)\).

\sgaStableTarget{sga6.e08.para.4-6}\textbf{4.6.} It remains to show the relation
\[
 \ch\bigl(i_*(i^*(y))\bigr)
 =
 i^{\mathrm{gr}}_*\bigl(\ch(i^*(y))\Todd(\check T_i)\bigr).
\]
The projection formula gives
\[
 i_*(i^*(y))=y\,i_*(1),
\]
and from the exact sequence
\[
0\longrightarrow J\longrightarrow\mathcal O_X\longrightarrow\mathcal O_Y\longrightarrow0,
\]
where \(J\) is here invertible, one obtains
\[
 i_*(1)=1-J.
\]
Thus
\[
 \ch(i_*(i^*(y)))=\ch(y)\,\ch(1-J).
\]
Since \(J\) is invertible, one has
\[
 \widetilde c(J)=(1,1+c^1(J)),\qquad \ch(J)=\exp(c^1(J)).
\]
Consequently
\[
 \ch(i_*(i^*(y)))=\ch(y)(1-\exp(c^1(J))).
\]

On the other hand, the functoriality of \(\ch\) and the projection formula \hyperlink{sga6.e08.prop.3-7}{3.7} imply
\[
 i^{\mathrm{gr}}_*\bigl(\ch(i^*(y))\Todd(\check T_i)\bigr)
 =
 \ch(y)\,i^{\mathrm{gr}}_*(\Todd(\check T_i)).
\]
Now
\[
T_i=-c^1(J/J^2)=-i^*(J).
\]
Hence
\[
 i^{\mathrm{gr}}_*(\Todd(\check T_i))
 =
 i^{\mathrm{gr}}_*\,i^{\mathrm{gr}}(\Todd(-\check J))
 =
 i^{\mathrm{gr}}_*(1)\Todd(-\check J).
\]
Since \(i_*(1)=1-J\), \(i^{\mathrm{gr}}_*(1)\) is the image in \(\Gr^1K^{\bullet}(X)\) of \(1-J\), hence it is \(-c^1(J)\). Moreover
\[
 \Todd(-\check J)=\Todd(J)^{-1}
 =
 \frac{1-\exp(c^1(J))}{-c^1(J)}.
\]
Thus finally
\[
 i^{\mathrm{gr}}_*\bigl(\ch(i^*(y))\Todd(\check T_i)\bigr)
 =
 \ch(y)(1-\exp(c^1(J))),
\]
which proves the relation to be shown.

\sgaStableTarget{sga6.e08.sec.5}\section*{5. Riemann--Roch theorem: case of the structural morphism of a projective bundle}

\sgaStableTarget{sga6.e08.para.5-1}\textbf{5.1.} By \hyperlink{sga6.e08.para.3-9}{3.9}, in order to prove the Riemann--Roch theorem, it remains to prove it for the structural morphism
\[
f:\mathbb P^r_Y\longrightarrow Y,
\]
where \(Y\) is a quasi-compact scheme admitting an ample sheaf. Put
\[
P=\mathbb P^r_Y,\qquad \xi=c^1(\mathcal O_P(1)),\qquad E=(\mathcal O_Y)^{r+1}.
\]

\sgaStableTarget{sga6.e08.prop.5-2}\textbf{Proposition 5.2.} On \(P\) one has the exact sequence
\[
0\longrightarrow\Omega^1_{P/Y}\longrightarrow f^*(E)(-1)\longrightarrow\mathcal O_P\longrightarrow0,
\]
i.e. one has an isomorphism
\[
\Omega^1_{P/Y}\simeq F(-1),
\qquad
F=\Ker\bigl(f^*(E)\longrightarrow\mathcal O_P(1)\bigr).
\]

Consider the vector bundle over \(Y\) defined by \(E\), and denote by \(Z\) the complement of the zero section of this bundle. One has the commutative diagram
\[
\begin{tikzcd}[column sep=large,row sep=large]
P \arrow[d,"f"'] & Z \arrow[l,"h"'] \arrow[dl,"g"]\\
Y &
\end{tikzcd}
\]
The morphism \(h\) identifies \(Z\) with the affine \(P\)-scheme defined by the sheaf of algebras
\[
B=\bigoplus_{n\in\mathbb Z}\mathcal O_P(n);
\]
\(B\) is a graded \(\mathcal O_P\)-algebra.

Since \(f\), \(g\), and \(h\) are smooth morphisms, one may write the exact sequence of sheaves of relative differentials. Since \(Z\) is affine over \(P\), one may write this sequence in terms of \(B\)-modules; thus one obtains
\sgaStableTarget{sga6.e08.eq.5-2-1}\begin{equation}\tag{5.2.1}
0\longrightarrow
 \Omega^1_{P/Y}\otimes_{\mathcal O_P}B
\longrightarrow
 \Omega^1_{B/Y}
\longrightarrow
 \Omega^1_{B/\mathcal O_P}
\longrightarrow0 .
\end{equation}
This exact sequence is an exact sequence of graded \(B\)-modules. Indeed, it follows from the commutative diagram
\[
\begin{tikzcd}[column sep=large,row sep=large]
S(E) \arrow[rr] && \displaystyle\bigoplus_{n\in\mathbb Z}S(E)(n)\simeq S(E)[T,T^{-1}]\\
& \mathcal O_X \arrow[ul] \arrow[ur] &
\end{tikzcd}
\]
where the arrows are algebra homomorphisms, compatible with the grading of \(S(E)[T,T^{-1}]\) (grading defined by the degree in the indeterminate \(T\)), the two other algebras being supposed to be reduced to degree \(0\).

The module of differentials \(\Omega^1_{B/\mathcal O_P}\) is isomorphic to
\[
 B(-1)\otimes_{\mathcal O_P}\mathcal O_P(1),
\]
for one may define an \(S(E)\)-derivation
\[
S(E)[T,T^{-1}]
\longrightarrow
S(E)[T,T^{-1}]\otimes_{S(E)}S(E)\,dT
\]
satisfying the universal property which defines \(\Omega^1_{B/\mathcal O_P}\), by setting
\[
 d(aT^n)=naT^{n-1}\otimes dT;
\]
moreover, \(d\) is compatible with the gradings if the degree \(1\) is attributed to the factor \(S(E)\,dT\).

Let \(V\) be the vector bundle defined by \(E\) over \(Y\). Then \(\Omega^1_{Z/Y}\) is the restriction to \(Z\) of \(\Omega^1_{V/Y}\). But \(V\) is the affine \(Y\)-scheme defined by \(S(E)\), so \(\Omega^1_{V/Y}\) is defined by
\[
S(E)\otimes_{\mathcal O_Y}E
\]
(since \(E\) is locally free of finite type), and therefore \(\Omega^1_{V/Y}\) is equal to \(g^*(E)\). Consequently
\[
\Omega^1_{B/Y}=f^*(E)\otimes_{\mathcal O_P}B.
\]
To see its grading, one looks locally, and one sees at once that \(f^*(E)\) must be considered as being of degree \(1\). Taking the terms of degree \(0\) in the exact sequence \hyperlink{sga6.e08.eq.5-2-1}{(5.2.1)}, one obtains the announced exact sequence.

\sgaStableTarget{sga6.e08.rem.5-2-1}\textbf{Remark 5.2.1.} Using the general results on the differential study of Hilbert functors ([2], 5), one directly obtains an isomorphism
\[
\Omega^1_{P/Y}\xrightarrow{\sim}F(-1)
\]
from the functorial characterization given of \(g^*(\Omega^1_{P/Y})\) for every \(Y\)-morphism \(g:Y'\to P\).

\sgaStableTarget{sga6.e08.para.5-3}\textbf{5.3.} By \hyperlink{sga6.e08.prop.5-2}{5.2}, one has
\[
 T_f=c_1(\Omega^1_{P/Y})=f^*(E)\xi^{-1}-1.
\]
Consequently
\[
 \widetilde c(T_f)=\widetilde c(E)\,\widetilde c(\xi^{-1})-\widetilde c(1).
\]
Since \(E=(\mathcal O_Y)^{r+1}\), one has \(E=r+1\), and hence
\[
 \widetilde c(E)=(r+1,1).
\]
Put
\[
 x=c^1_{\Gr^1K^{\bullet}(P)}(\xi-1)
 =
 c^1_{\Gr^1K^{\bullet}(P)}(1-\xi^{-1})
 =
 c^1(\xi).
\]
Since \(\xi^{-1}=\check \xi\), one has
\[
 \widetilde c(\xi^{-1})=(1,1-x).
\]
By \hyperlink{sga6.e05.eq.6-2-1}{V (6.2.1)}, one deduces
\[
 \widetilde c(T_f)=(r,(1-x)^{r+1}).
\]
Taking \hyperlink{sga6.e08.prop.3-5}{3.5} into account, one finally obtains
\[
 \widetilde c(\check T_f)=(r,(1+x)^{r+1}).
\]
By the definition of \(\Todd\), one therefore has
\sgaStableTarget{sga6.e08.eq.5-3-1}\begin{equation}\tag{5.3.1}
 \Todd(\check T_f)=
 \left(\frac{x}{1-\exp(-x)}\right)^{r+1}.
\end{equation}

\sgaStableTarget{sga6.e08.para.5-4}\textbf{5.4.} Using the basis of \(K^{\bullet}(P)\) over \(K^{\bullet}(Y)\) formed by the elements \(\xi^k\), with \(-r\leq k\leq0\), one sees by additivity that it is enough to prove the theorem for elements of the form
\[
 f^*(a)\xi^k .
\]
Then
\[
 \ch\bigl(f_*(f^*(a)\xi^k)\bigr)
 =
 \ch(a f_*(\xi^k))
 =
 \ch(a)\,\ch(f_*(\xi^k)).
\]
On the other hand,
\[
 f^{\mathrm{gr}}_*\bigl(\ch(f^*(a)\xi^k)\Todd(\check T_f)\bigr)
 =
 f^{\mathrm{gr}}_*\bigl(f^{\mathrm{gr}}(\ch(a))\ch(\xi^k)\Todd(\check T_f)\bigr)
 =
 \ch(a)\,f^{\mathrm{gr}}_*\bigl(\ch(\xi^k)\Todd(\check T_f)\bigr).
\]
Thus it is enough to prove the theorem for the \(\xi^k\), with \(-r\leq k\leq0\). By \hyperlink{sga6.e06.sec.5.2}{VI (5.2.3)}, one has
\[
 f_*(\xi^k)=0\quad\text{for }-r\leq k<0,\qquad f_*(1)=1.
\]
Consequently
\sgaStableTarget{sga6.e08.eq.5-4-1}\begin{equation}\tag{5.4.1}
 \ch(f_*(\xi^k))=0\quad\text{if }-r\leq k<0,\qquad
 \ch(f_*(1))=1.
\end{equation}

We now calculate the second member of the Riemann--Roch formula. For every \(k\), one has
\[
 \widetilde c(\xi^k)=(\widetilde c(\xi))^k=(1,1+x)^k=(1,1+kx)
\]
by \hyperlink{sga6.e05.eq.6-2-1}{V (6.2.1)}. Hence
\sgaStableTarget{sga6.e08.eq.5-4-2}\begin{equation}\tag{5.4.2}
 \ch(\xi^k)=\exp(kx).
\end{equation}
The second member of the Riemann--Roch formula is therefore
\sgaStableTarget{sga6.e08.eq.5-4-3}\begin{equation}\tag{5.4.3}
 f^{\mathrm{gr}}_*\left(\exp(kx)\,
 \frac{x^{r+1}}{(1-\exp(-x))^{r+1}}\right).
\end{equation}
Since \(x\) is the class in \(\Gr^1K^{\bullet}(P)\) of \(1-\xi^{-1}\), \(f^{\mathrm{gr}}_*(x^n)\) is, for \(0\leq n\leq r\), the class in
\(\Gr^{n-r}K^{\bullet}(Y)\) of \(f_*((1-\xi^{-1})^n)=1\). Thus
\sgaStableTarget{sga6.e08.eq.5-4-4}\begin{equation}\tag{5.4.4}
 f^{\mathrm{gr}}_*(x^n)=0\quad\text{if }0\leq n<r,\qquad
 f^{\mathrm{gr}}_*(x^r)=1.
\end{equation}
Finally recall that the elements
\[
1,x,\ldots,x^{r+1}
\]
are linearly dependent over \(\Gr K^{\bullet}(Y)\) (\hyperlink{sga6.e06.cor.5-6}{VI 5.6}), and satisfy the equation
\sgaStableTarget{sga6.e08.eq.5-4-5}\begin{equation}\tag{5.4.5}
 x^{r+1}=0,
\end{equation}
and that \(1,\ldots,x^r\) are linearly independent over \(\Gr K^{\bullet}(Y)\).

One can therefore write the series
\[
\varphi_k=\exp(kx)\,
\frac{x^{r+1}}{(1-\exp(-x))^{r+1}}
\]
in a unique way as a linear combination, with coefficients in
\(\Gr K^{\bullet}(Y)_{\mathbb Q}\), of the elements \(1,\ldots,x^r\). The image under \(f^{\mathrm{gr}}_*\) of \(\varphi_k\) is then, by \hyperlink{sga6.e08.eq.5-4-4}{(5.4.4)} and \hyperlink{sga6.e08.prop.3-7}{3.7}, its coefficient of \(x^r\). Taking account of \hyperlink{sga6.e08.eq.5-4-1}{(5.4.1)}, the Riemann--Roch theorem for \(f\) is therefore equivalent to
\sgaStableTarget{sga6.e08.eq.rr}\[
\left\{
\begin{array}{ll}
\text{for }-r\leq k<0,\text{ the coefficient of }x^r\text{ in }\varphi_k\text{ is }0,\\
\text{for }k=0,\text{ the coefficient of }x^r\text{ in }\varphi_k\text{ is }1.
\end{array}
\right.
\tag{RR}
\]
This coefficient is the residue of the differential form
\[
 \frac{\exp(kx)\,dx}{(1-\exp(-x))^{r+1}}.
\]
Putting \(y=1-\exp(-x)\), one is reduced to calculating the residue of
\[
 \frac{(1-y)^{-k-1}\,dy}{y^{r+1}},
\]
that is, the coefficient of \(y^r\) in \((1-y)^{-k-1}\). For \(-r\leq k<0\), it is zero, and for \(k=0\), it is \(1\).

This completes the proof of the Riemann--Roch theorem.

\sgaStableTarget{sga6.e08.bibliography}\subsection*{Bibliography}
\begin{enumerate}[label={[\arabic*]}]
\item A. Borel and J.-P. Serre, \emph{Le théorème de Riemann--Roch}, Bull. Soc. Math. de France, t. 86, 1958.
\item A. Grothendieck, \emph{Techniques de construction et théorèmes d'existence en Géométrie Algébrique}, Sém. Bourbaki, 1960--61.
\end{enumerate}



\sgaStableTarget{sga6.e09.expose}\section*{Exposé IX. Some computations of $K_\bullet$-groups}
\begin{center}
by P. Berthelot
\end{center}

This exposé corresponds to no oral exposé and will not be used in the sequel of the seminar. It contains a certain number of computations of the groups $K_\bullet$ for various fibrations: vector bundles, projective bundles, torsors under a linear group, and so on. Throughout the exposé, the base scheme will be assumed noetherian, and the notation $K_\bullet(X)$ will denote the Grothendieck group of coherent sheaves on a scheme $X$. Recall that, if $X$ is noetherian, the homomorphism from the $K$ of the category of coherent $\mathcal O_X$-modules to the $K$ of the category of pseudo-coherent complexes on $X$ is an isomorphism (\hyperlink{sga6.e04.para.2-4}{IV 2.4}).

The results set out here are the analogues of those one obtains for the Chow ring, and which are to be found in Grothendieck's exposés in the Chevalley seminar of 1958. The question also arises of determining to what extent statements of the same type can be obtained for the graded group associated with $K_\bullet$ for the filtration defined by the dimension of the support (cf. X). It does not seem that the answer is known, even for \hyperlink{sga6.e09.prop.1-1}{Proposition 1.1}, the homotopy exact sequence, which is one of the properties most used in the computations which follow.

\sgaStableTarget{sga6.e09.sec.1}\subsection*{1. Vector bundles}

\sgaStableTarget{sga6.e09.prop.1-1}\textbf{Proposition 1.1 (homotopy exact sequence).} Let $X$ be a noetherian scheme, let $X'$ be a closed subscheme of $X$, let $U$ be the open complement of $X'$ in $X$, let $i$ be the immersion of $X'$ in $X$, and let $j$ be the immersion of $U$ in $X$. Then the sequence
\[
 K_\bullet(X')\xrightarrow{i_*}K_\bullet(X)\xrightarrow{j^*}K_\bullet(U)\longrightarrow 0
\]
is exact.

If $F'$ is a coherent $\mathcal O_{X'}$-module, $i_*(F')$ is an $\mathcal O_X$-module whose support is contained in $X'$; consequently its restriction to $U$ is zero, and $j^*i_*(\cl(F'))=0$. Since the elements of the form $\cl(F')$ generate $K_\bullet(X')$, $j^*i_*=0$. We therefore obtain a homomorphism
\[
 \rho:K_\bullet(X)/\operatorname{Im}K_\bullet(X')\longrightarrow K_\bullet(U).
\]
We shall define an inverse homomorphism to $\rho$. Let $G$ be a coherent $\mathcal O_U$-module. Since $X$ is noetherian, there exists a coherent $\mathcal O_X$-module $F$ whose restriction to $U$ is isomorphic to $G$. If $F'$ is a second coherent $\mathcal O_X$-module whose restriction to $U$ is isomorphic to $G$, then in $K_\bullet(X)$ one has
\[
 \cl(F)-\cl(F')\in \operatorname{Im}K_\bullet(X').
\]
To see this, consider $G$ as a subsheaf of $j^*(F\oplus F')$ by means of the diagonal map. Then $G$ extends to a subsheaf $F''$ of $F\oplus F'$, and one is reduced to proving that
\[
 \cl(F)-\cl(F'')\in \operatorname{Im}K_\bullet(X')
\]
and likewise for $F'$. Let $H$ and $K$ be the kernel and cokernel of the composite homomorphism
\[
 F''\longrightarrow F\oplus F'\longrightarrow F.
\]
This homomorphism restricts on $U$ to the identity of $G$, and hence $H$ and $K$ have support contained in $X'$. If $I$ is the ideal of $X'$ in $X$, it follows that there exist integers $m$ and $n$ such that $I^mH=0$ and $I^nK=0$ (EGA I 5.3.4). Thus one may dévissser $H$ and $K$ into extensions of modules annihilated by $I$, hence of $\mathcal O_{X'}$-modules, and
\[
 \cl(H)\in\operatorname{Im}K_\bullet(X'),\qquad \cl(K)\in\operatorname{Im}K_\bullet(X'),
\]
which proves the assertion. Thus, to every coherent $\mathcal O_U$-module, assigning the class of one of its extensions to $X$ defines a function on the category of coherent $\mathcal O_U$-modules with values in $K_\bullet(X)/\operatorname{Im}K_\bullet(X')$.

This function is additive. Indeed, let
\[
0\longrightarrow G'\longrightarrow G\longrightarrow G''\longrightarrow 0
\]
be an exact sequence of coherent $\mathcal O_U$-modules. Let $F$ be an extension of $G$ to $X$. There exists a submodule $F'$ of $F$ extending $G'$, and $F/F'$ extends $G''$. Hence the function factors through $K_\bullet(U)$, defining the desired homomorphism $K_\bullet(U)\to K_\bullet(X)/\operatorname{Im}K_\bullet(X')$. One verifies at once that it is inverse to $\rho$.

\sgaStableTarget{sga6.e09.prop.1-2}\textbf{Proposition 1.2.} Let $X$ and $S$ be noetherian schemes, and let $f:X\to S$ be a flat morphism. Let $(\xi_i)_{i\in I}$ be a family of elements of $K^\circ(X)$, and consider the homomorphism
\[
 \varphi:K_\bullet(S)^{(I)}\longrightarrow K_\bullet(X)
\]
defined by
\[
 \varphi((x_i)_{i\in I})=\sum_{i\in I}\xi_i\,f^*(x_i),
\]
where $f^*$ is the inverse-image homomorphism from $K_\bullet(S)$ to $K_\bullet(X)$ defined in \hyperlink{sga6.e04.para.2-12}{IV 2.12}, and where $K_\bullet(X)$ is considered as a $K^\circ(X)$-module by \hyperlink{sga6.e04.para.2-10}{IV 2.10}. Suppose that, for every point $s\in S$, the images of the $\xi_i$ in $K_\bullet(X_s)$ by the homomorphism
\[
 K^\circ(X)\longrightarrow K^\circ(X_s)=K_\bullet(X_s)
\]
generate $K_\bullet(X_s)$ as an abelian group. Then $\varphi$ is surjective.

Let $S'$ be a closed subscheme of $S$, and let $U$ be the open complement of $S'$ in $S$. Write $X_U$ for the restriction of $X$ over $U$, and $X_{S'}$ for the restriction of $X$ over $S'$. We again denote by $\xi_i$ the images of the $\xi_i$ in $K^\circ(X_{S'})$ and $K^\circ(X_U)$, and write
\[
 \varphi_{S'}:K_\bullet(S')^{(I)}\to K_\bullet(X_{S'}),\qquad
 \varphi_U:K_\bullet(U)^{(I)}\to K_\bullet(X_U)
\]
for the homomorphisms defined in the same way as $\varphi$.

Consider the diagram
\sgaStableTarget{sga6.e09.eq.1-2-1}\[
\begin{tikzcd}[column sep=tiny]
K_\bullet(S')^{(I)} \arrow[r,"i_*"] \arrow[d,"\varphi_{S'}"'] &
K_\bullet(S)^{(I)} \arrow[r,"j^*"] \arrow[d,"\varphi"] &
K_\bullet(U)^{(I)} \arrow[r] \arrow[d,"\varphi_U"] & 0\\
K_\bullet(X_{S'}) \arrow[r,"u_*"] & K_\bullet(X) \arrow[r,"v^*"] & K_\bullet(X_U) \arrow[r] & 0,
\end{tikzcd}\tag{1.2.1}
\]
where $i$ and $j$ denote the immersions of $S'$ and $U$ into $S$, and $u$ and $v$ the immersions of $X_{S'}$ and $X_U$ into $X$. We show that, if $\varphi_{S'}$ and $\varphi_U$ are surjective, then so is $\varphi$.

The two horizontal rows are exact by 1.1. To prove that $\varphi$ is surjective, it is therefore enough to prove that the two squares commute. For $x_i\in K_\bullet(S')$, one has
\[
 \varphi(i_*(x_i))=\sum_i \xi_i f^*i_*(x_i),
\]
whereas
\[
 u_*\varphi_{S'}((x_i))=u_*\left(\sum_i\xi_i u^*f^*(x_i)\right)=\sum_i \xi_i u_*(u^*f^*(x_i)),
\]
the last equality following from the projection formula \hyperlink{sga6.e04.eq.2-11-1-2}{IV 2.11.1.2}. The equality $\varphi i_*=u_*\varphi_{S'}$ therefore follows from $f^*i_*=u_*u^*f^*$, which is true by \hyperlink{sga6.e04.prop.3-1-1}{IV 3.1.1} since $i$ is proper and $f$ is flat. Similarly, for $(x_i)\in K_\bullet(S)^{(I)}$,
\[
 \varphi_U(j^*(x_i))=\sum_i\xi_i f_U^*j^*(x_i),
\]
and
\[
 v^*\varphi((x_i))=v^*\left(\sum_i\xi_i f^*(x_i)\right)=\sum_i\xi_i v^*f^*(x_i)=\sum_i\xi_i f_U^*j^*(x_i),
\]
the middle equality being \hyperlink{sga6.e04.eq.2-12-2}{IV 2.12.2}.

This result permits one to prove \hyperlink{sga6.e09.prop.1-2}{1.2} by noetherian induction on the set of closed subschemes of $S$. If one assumes that, for every closed subscheme $S'$ of $S$ distinct from $S$, $\varphi_{S'}$ is surjective, the preceding argument shows that one may reduce to proving \hyperlink{sga6.e09.prop.1-2}{1.2} for an irreducible affine base scheme, after possibly replacing $S$ by an irreducible affine open. Moreover, the diagram
\[
\begin{tikzcd}
K_\bullet(S_{\mathrm{red}})^{(I)} \arrow[r] \arrow[d,"\varphi_{S_{\mathrm{red}}}"'] & K_\bullet(S)^{(I)} \arrow[d,"\varphi"]\\
K_\bullet(X_{S_{\mathrm{red}}}) \arrow[r] & K_\bullet(X)
\end{tikzcd}
\]
where $K_\bullet(X_{S_{\mathrm{red}}})\to K_\bullet(X)$ is surjective by \hyperlink{sga6.e09.prop.1-1}{1.1}, shows that one may also assume $S$ reduced.

We are thus on an integral affine base scheme. Put $S=\Spec(A)$, and let $K$ denote the fraction field of $A$. If $s$ is the generic point of $S$, then $\mathcal O_{S,s}=K$. Let $y\in K_\bullet(X)$. By hypothesis, $K_\bullet(X_s)$ is generated by the images of the $\xi_i$; if $y_s$ is the image of $y$ in $K_\bullet(X_s)$, then one can find integers $n_i$ such that
\[
 y_s=\sum_i \xi_i n_i.
\]
On the other hand $K=\varinjlim_f A_f$, where $f$ runs through the elements of $A$. Hence $\Spec(K)=\varprojlim_f \Spec(A_f)$ and
\[
 X_s=X\times_S\Spec(K)=\varprojlim_f X\times_S\Spec(A_f)
\]
by EGA IV 8.2.3 and 8.2.5. Since the schemes considered are noetherian and the transition morphisms affine and flat, one has, by \hyperlink{sga6.e04.para.3-2-3}{IV 3.2.3},
\[
 K_\bullet(X_s)=\varinjlim_f K_\bullet(X\times_S\Spec(A_f)).
\]
As inverse-image homomorphisms on the $K_\bullet$ are linear with respect to $K^\circ$ (\hyperlink{sga6.e04.eq.2-12-2}{IV 2.12.2}), there exists a non-empty open $U=\Spec(A_f)$ of $S$ such that in $K_\bullet(X_U)$ one has
\[
 y_U=\sum_i \xi_i n_i.
\]
Let $S'$ be a closed subscheme of $S$ with support $S-U$. Then $\varphi_{S'}$ is surjective by the induction hypothesis, and consequently diagram \hyperlink{sga6.e09.eq.1-2-1}{(1.2.1)} shows that
\[
 y-
\sum_i\xi_i n_i\in \operatorname{Im}(\varphi).
\]
Since $n_i=f^*(n_i)$, $\varphi$ is surjective.

\sgaStableTarget{sga6.e09.cor.1-3}\textbf{Corollary 1.3.} Let $f:X\to S$ be a flat morphism of noetherian schemes. Suppose that, for every point $s$ of $S$, the homomorphism
\[
 K^\circ(X)\longrightarrow K_\bullet(X_s)
\]
is surjective. Then the homomorphism deduced from $f^*$,
\[
  K_\bullet(S)\otimes_{K^\circ(S)}K^\circ(X)\longrightarrow K_\bullet(X),
\]
is surjective.

Take for the family $(\xi_i)$ the family of all elements of $K^\circ(X)$. The hypotheses of \hyperlink{sga6.e09.prop.1-2}{1.2} are then satisfied, and hence
\[
 \varphi:K_\bullet(S)^{(K^\circ(X))}\longrightarrow K_\bullet(X)
\]
is surjective. By definition of $\varphi$, this is equivalent to the surjectivity of the displayed tensor-product homomorphism.

\sgaStableTarget{sga6.e09.cor.1-4}\textbf{Corollary 1.4.} Let $f:X\to S$ be a flat morphism of noetherian schemes. Suppose that, for every point $s$ of $S$, the canonical homomorphism $\mathbb Z\to K_\bullet(X_s)$ is an isomorphism. Then
\[
 f^*:K_\bullet(S)\longrightarrow K_\bullet(X)
\]
is surjective. If, moreover, $f$ has a section $g$ of finite Tor-dimension, then $f^*$ is an isomorphism.

Take for the family $(\xi_i)$ the family reduced to the single element $1\in K^\circ(X)$. Then $\varphi=f^*$, and the conditions of \hyperlink{sga6.e09.prop.1-2}{1.2} being fulfilled, $f^*$ is surjective. If $f$ admits a section $g$ of finite Tor-dimension, one may define a homomorphism
\[
 g^*:K_\bullet(X)\longrightarrow K_\bullet(S)
\]
(\hyperlink{sga6.e04.para.2-12}{IV 2.12}). Since $g^*f^*=(fg)^*=\operatorname{Id}$ by \hyperlink{sga6.e04.eq.2-12-1-1}{IV 2.12.1.1}, $f^*$ is injective, hence an isomorphism.

\sgaStableTarget{sga6.e09.rem.1-5}\textbf{Remark 1.5.} For $g$ to be of finite Tor-dimension, it is necessary and sufficient that the fibres of $X$ be regular at the points of $g(S)$, i.e. that $f$ be regular at the points of $g(S)$, or smooth at these points if one assumes $f$ of finite presentation.

These properties being local on $X$ and $S$, one may suppose them affine: $S=\Spec(A)$, $X=\Spec(B)$. Then
\[
 \operatorname{tor.dim}_B(A)=\operatorname{dim.proj}_B(A)=\sup_{s\in S}\operatorname{dim.proj}_{B\otimes_A k(s)}(k(s)),
\]
by EGA IV 12.3.2, where the finite-presentation hypothesis on $B$ may be replaced by noetherian hypotheses. But $k(s)$ is isomorphic to the residue field of $\mathcal O_{X_s,g(s)}$, and by EGA $0_{\mathrm{IV}}$ 17.2.5 one has
\[
 \operatorname{dim.proj}_{B\otimes_A k(s)}(k(s))=
 \operatorname{dim.proj}_{\mathcal O_{X_s,g(s)}}(k(s)).
\]
It follows from EGA $0_{\mathrm{IV}}$ 17.2.11 that the latter is finite if and only if $\mathcal O_{X_s,g(s)}$ has finite cohomological dimension, i.e. is regular.

\sgaStableTarget{sga6.e09.prop.1-6}\textbf{Proposition 1.6.} Let $S$ be a noetherian scheme, let $E$ be a locally free $\mathcal O_S$-module of finite type, let $X$ be the vector bundle associated with $E$, and let $f$ be the structural morphism of $X$. Then
\[
 f^*:K_\bullet(S)\longrightarrow K_\bullet(X)
\]
is an isomorphism.

Since $f$ is smooth and admits a section, namely the zero section, it is enough, to prove \hyperlink{sga6.e09.prop.1-6}{1.6}, to show that for every $s\in S$ one has $K_\bullet(X_s)\simeq \mathbb Z$ \hyperlink{sga6.e09.cor.1-4}{(1.4)}. Thus one is reduced to the case where the base scheme is a field. In this case $E$ is a direct sum of rank-one submodules, which permits one to decompose $f$ into a sequence of rank-one fibrations, and hence to reduce to the case where $X$ has rank one, namely $X=\Spec(k[T])$. The proposition then follows from the next lemma.

\sgaStableTarget{sga6.e09.lem.1-7}\textbf{Lemma 1.7.} Let $A$ be a principal ideal domain and let $X=\Spec(A)$. Then the canonical homomorphism
\[
 \mathbb Z\longrightarrow K_\bullet(X)
\]
is an isomorphism.

Since $X$ is regular, the homomorphism $K^\circ(X)\to K_\bullet(X)$ is an isomorphism. Hence $K_\bullet(X)$ is generated by the classes of locally free $\mathcal O_X$-modules of finite type, i.e. by the classes of projective $A$-modules of finite type. Since $A$ is principal, every projective module of finite type is free, which gives surjectivity. Injectivity follows from the fact that $K^\circ(X)$ is $\mathbb Z$-augmented.

\sgaStableTarget{sga6.e09.prop.1-8}\textbf{Proposition 1.8.} Let $S$ be a noetherian scheme, let $E$ be a locally free $\mathcal O_S$-module of rank $r$, let $X$ be the vector bundle associated with $E$, let $f$ be the structural morphism of $X$, and let $g$ be the zero section of $X$. Then the composite homomorphism
\[
 K_\bullet(S)\xrightarrow{g_*}K_\bullet(X)\xrightarrow{g^*}K_\bullet(S)
\]
where $g^*$ is defined by \hyperlink{sga6.e04.eq.2-12-1}{IV 2.12.1} and \hyperlink{sga6.e04.para.1-5}{1.5}, is multiplication by the element $\lambda_{-1}(E)$ of $K^\circ(S)$, writing $E$ for the class of $E$ in $K^\circ(S)$ and putting
\[
 \lambda_{-1}(E)=\sum_{i=0}^r(-1)^i\lambda^i(E).
\]

This statement is the analogue, for $K_\bullet$, of \hyperlink{sga6.e07.para.2-6}{VII 2.7}. Let $F$ be a coherent $\mathcal O_S$-module, of class $F$ in $K_\bullet(S)$. Then $g_*(F)$ is the class in $K_\bullet(X)$ of the complex $Rg_*(F)=g_*(F)$, hence of $F$ considered as an $\mathcal O_X$-module through $g$. To compute $g^*(g_*(F))$, one must take the class of
\[
 Lg^*(F)=F\otimes^L_{\mathcal O_X}\mathcal O_S.
\]
For this, consider the Koszul complex
\[
0\longrightarrow \Lambda^r(E)\otimes_{\mathcal O_S}S(E)\longrightarrow
\Lambda^{r-1}(E)\otimes_{\mathcal O_S}S(E)\longrightarrow \cdots\longrightarrow
E\otimes_{\mathcal O_S}S(E)\longrightarrow S(E)\longrightarrow 0,
\]
where $S(E)$ is the symmetric algebra of $E$. This is a resolution of $\mathcal O_S$ over $S(E)$ (\hyperlink{sga6.e06.sec.1.10}{VI 1.11}). The associated complex on $X$, denoted $L_\bullet$, has $k$-th term $\Lambda^k(f^*E)$ and is therefore a resolution of $\mathcal O_S$ by locally free $\mathcal O_X$-modules of finite type. Thus
\[
 L_\bullet\otimes_{\mathcal O_X}F=\mathcal O_S\otimes_{\mathcal O_X}F,
\]
and hence
\[
\begin{aligned}
 g^*(g_*(F))
 &=\cl(L_\bullet\otimes_{\mathcal O_X}F)
   =\sum_{i=0}^r(-1)^i\cl(L_i\otimes_{\mathcal O_X}F)\\
 &=\sum_{i=0}^r(-1)^i\cl(\Lambda^i(E)\otimes_{\mathcal O_S}F)
   =F\lambda_{-1}(E).
\end{aligned}
\]
The proposition follows by additivity.

\sgaStableTarget{sga6.e09.sec.2}\subsection*{2. Principal bundles under split tori}

\sgaStableTarget{sga6.e09.para.2-1}\textbf{2.1.} Let $S$ be a noetherian scheme, and let $T$ be a torus over $S$. Suppose $T$ is split; thus one has an isomorphism $T\simeq D_S(M)$, where $M$ is a free abelian group of rank $r$, or equivalently $T\simeq \mathbb G_{m,S}^r$. Let $P$ be a torsor under $T$. Then $P$ is affine over $S$, defined by a graded $\mathcal O_S$-algebra $A$ of type $M$. If
\[
 A=\bigoplus_{m\in M}A_m,
\]
then $A_m$ is an invertible $\mathcal O_S$-module for every $m$, and, for $m,m'\in M$, the homomorphism
\[
 A_m\otimes_{\mathcal O_S}A_{m'}\longrightarrow A_{m+m'}
\]
defined by the multiplication of $A$ is an isomorphism (SGA 3 VIII 4.1). One therefore defines a multiplicative homomorphism
\[
 \eta:M\longrightarrow K^\circ(S)
\]
by assigning to $m\in M$ the class of $A_m$ in $K^\circ(S)$.

\sgaStableTarget{sga6.e09.prop.2-2}\textbf{Proposition 2.2.} With the notation and hypotheses of \hyperlink{sga6.e09.para.2-1}{2.1}, let $f:P\to S$ be the structural morphism of $P$. Then the homomorphism
\[
 f^*:K_\bullet(S)\longrightarrow K_\bullet(P)
\]
is surjective, and its kernel is the subgroup of $K_\bullet(S)$ generated by the subgroups $(\eta(m)-1)K_\bullet(S)$ for $m\in M$, or, equivalently, for $m$ running through a system of generators of $M$.

First suppose $r=1$, i.e. $T=\mathbb G_{m,S}$. Then $M=\mathbb Z$, and for every $m\in M$ one has
\[
 A_m\simeq A_1^{\otimes m}.
\]
If $A_1=L$, then $A=\bigoplus_{n\in\mathbb Z}L^{\otimes n}$, which shows that $P$ identifies with the complement of the zero section in the bundle $X=\mathbb V(L)$. Considering $S$ as a closed subscheme of $X$ by the zero section, one has, by \hyperlink{sga6.e09.prop.1-1}{1.1}, the exact sequence
\[
 K_\bullet(S)\xrightarrow{g_*}K_\bullet(X)\longrightarrow K_\bullet(P)\longrightarrow 0,
\]
where $g:S\to X$ is the zero section of $X$. Moreover $f^*$ factors as
\[
 K_\bullet(S)\xrightarrow{f'^*}K_\bullet(X)\longrightarrow K_\bullet(P),
\]
where $f'$ is the structural morphism of $X$. By \hyperlink{sga6.e09.prop.1-6}{1.6}, $f'^*$ is an isomorphism, inverse to $g^*$. Thus we have the exact sequence
\[
 K_\bullet(S)\xrightarrow{g^*g_*}K_\bullet(S)\xrightarrow{f^*}K_\bullet(P)\longrightarrow 0,
\]
and by \hyperlink{sga6.e09.prop.1-8}{1.8} one obtains
\[
 K_\bullet(P)\simeq K_\bullet(S)/\lambda_{-1}(L)K_\bullet(S)=K_\bullet(S)/(1-L)K_\bullet(S),
\]
where $L$ is the class of $L$ in $K^\circ(S)$. Since, for every $n\in\mathbb Z$, $L^n-1$ is a multiple of $1-L$, this proves \hyperlink{sga6.e09.prop.2-2}{2.2} when $P$ has rank one.

Continue the proof by induction on the rank $r$ of $P$. Suppose the result proved for torsors under split tori of rank $<r$. Write $M=M_1\oplus M_2$, with $M_1\simeq\mathbb Z^{r-1}$ and $M_2\simeq\mathbb Z$. Put $A'=\bigoplus_{m\in M_1}A_m$ and $A''=\bigoplus_{m\in M_2}A_m$, so that $A=A'\otimes_{\mathcal O_S}A''$.
Let $P_1=\Spec(A')$. Then $P_1$ is a torsor under $D_S(M_1)$, and $P$ is, over $P_1$, a torsor under the group $\mathbb G_{m,P_1}$ with algebra the inverse image of $A''$ on $P_1$. By the preceding rank-one case, $K_\bullet(P)$ is the quotient of $K_\bullet(P_1)$ by the subgroup generated by the $(\eta(m_2)-1)K_\bullet(P_1)$ for $m_2\in M_2$. By the induction hypothesis, $K_\bullet(P_1)$ is the quotient of $K_\bullet(S)$ by the subgroup generated by the $(\eta(m_1)-1)K_\bullet(S)$ for $m_1\in M_1$. Taking account of the relation
\[
 \eta(m_1+m_2)-1=\eta(m_1)\eta(m_2)-1
 =\eta(m_1)(\eta(m_2)-1)+\eta(m_1)-1,
\]
we conclude that $K_\bullet(P)$ is the quotient of $K_\bullet(S)$ by the subgroup generated by the $(\eta(m)-1)K_\bullet(S)$ for $m\in M$, or equivalently for $m$ running through a system of generators of $M$.

\sgaStableTarget{sga6.e09.cor.2-3}\textbf{Corollary 2.3.} With the notation and hypotheses of \hyperlink{sga6.e09.para.2-1}{2.1}, suppose $S$ regular. Then one has a canonical isomorphism
\[
 K^\circ(P)\simeq K^\circ(S)/\bigl((\eta(m)-1)_{m\in M}\bigr),
\]
where $((\eta(m)-1)_{m\in M})$ is the ideal of $K^\circ(S)$ generated by the $\eta(m)-1$, for $m$ running through a system of generators of $M$.

If $S$ is regular, the canonical homomorphism $K^\circ(S)\to K_\bullet(S)$ is an isomorphism (\hyperlink{sga6.e04.para.2-5}{IV 2.5}). Likewise, $P$ is then regular, and $K^\circ(P)\to K_\bullet(P)$ is an isomorphism. The corollary therefore follows immediately from 2.2.

\sgaStableTarget{sga6.e09.sec.3}\subsection*{3. Projective bundles and flag bundles}

\sgaStableTarget{sga6.e09.prop.3-1}\textbf{Proposition 3.1.} Let $S$ be a noetherian scheme, let $E$ be a locally free $\mathcal O_S$-module of rank $n$, let $\pi=(p_1,\ldots,p_k)$ be a sequence of integers such that $\sum p_i=n$, and let $X=D_\pi(E)$ be the scheme of flags of type $\pi$ in $E$. Then the canonical homomorphism
\[
 K_\bullet(S)\otimes_{K^\circ(S)}K^\circ(X)\longrightarrow K_\bullet(X)
\]
is an isomorphism.

We resume the notation of \hyperlink{sga6.e06.sec.4.5}{VI 4.5}. Thus $L_i$ denotes the canonical quotients of the inverse image of $E$ on $X$, $F_j$ denotes the kernel of $L_j\to L_{j-1}$, and $\lambda_i^{(j)}$ denotes the class in $K^\circ(X)$ of $\Lambda^i(F_j)$. Then the $\lambda_i^{(j)}$ generate $K^\circ(X)$ as a $K^\circ(S)$-algebra (\hyperlink{sga6.e06.thm.4-6}{VI 4.6}). Likewise, the images of the $\lambda_i^{(j)}$ in $K^\circ(X_s)$ generate $K^\circ(X_s)$ as an algebra over $K^\circ(k(s))=\mathbb Z$ for every point $s$ of $S$. Since $X_s$ is a flag scheme over a field, it is regular, and hence $K^\circ(X_s)\to K_\bullet(X_s)$ is an isomorphism. It follows that $K^\circ(X)\to K_\bullet(X_s)$ is surjective. By \hyperlink{sga6.e09.cor.1-3}{1.3}, this implies that
\[
 K_\bullet(S)\otimes_{K^\circ(S)}K^\circ(X)\longrightarrow K_\bullet(X)
\]
with $K^\circ(X)$ in the second factor is surjective.

To prove injectivity, one may first reduce to the case of the complete flag bundle $D(E)$. Indeed, suppose that
\[
 K_\bullet(S)\otimes_{K^\circ(S)}K^\circ(D(E))\longrightarrow K_\bullet(D(E))
\]
is injective. Since $K_\bullet(S)\to K_\bullet(D(E))$ factors through $K_\bullet(X)$, the corresponding homomorphism factors as
\[
 K_\bullet(S)\otimes_{K^\circ(S)}K^\circ(X)\otimes_{K^\circ(X)}K^\circ(D(E))
 \longrightarrow
 K_\bullet(X)\otimes_{K^\circ(X)}K^\circ(D(E))
 \longrightarrow K_\bullet(D(E)).
\]
The first homomorphism is injective. Since, by \hyperlink{sga6.e06.lem.4-3}{VI 4.3}, \hyperlink{sga6.e06.thm.4-6}{4.6} and \hyperlink{sga6.e06.cor.4-7}{4.7}, $K^\circ(D(E))$ is free over $K^\circ(X)$, this implies injectivity for
\[
 K_\bullet(S)\otimes_{K^\circ(S)}K^\circ(X)\to K_\bullet(X).
\]
Thus assume $X=D(E)$. One proceeds by induction on the rank of $E$. The case $n=1$ is trivial. Suppose the property proved for $n-1$. Introduce $P=\mathbb P(E)$, and let $F$ be the kernel of the canonical homomorphism $p^*(E)\to\mathcal O_P(1)$, where $p$ is the structural morphism of $P$. Then $X$ identifies with the scheme of flags of $F$ over $P$, and by the induction hypothesis it is enough to prove \hyperlink{sga6.e09.prop.3-1}{3.1} for $P$. Taking \hyperlink{sga6.e06.thm.1-1}{VI 1.1} into account, the result to prove may therefore be stated as the following corollary.

\sgaStableTarget{sga6.e09.cor.3-2}\textbf{Corollary 3.2.} Let $S$ be a noetherian scheme, let $E$ be a locally free $\mathcal O_S$-module of rank $n$, let $X=\mathbb P(E)$ be the associated projective bundle, let $f$ be the structural morphism of $X$, and let $\xi$ be the class in $K^\circ(X)$ of the sheaf $\mathcal O_X(1)$. Then the homomorphism
\[
 \varphi:(K_\bullet(S))^n\longrightarrow K_\bullet(X)
\]
defined by
\[
 \varphi((x_i))=\sum_{i=0}^{n-1}\xi^{-i}f^*(x_i)
\]
is an isomorphism.

We have already seen that this homomorphism is surjective, since the elements $\xi^i$, $i=0,\ldots,-n+1$, form a basis of $K^\circ(X)$ over $K^\circ(S)$. To see that it is injective, use the homomorphism $f_*$, defined because $f$ is projective. By \hyperlink{sga6.e04.eq.2-12-4}{IV 2.12.4},
\[
 f_*\left(\sum_{i=0}^{n-1}\xi^{-i}f^*(x_i)\right)
 =\sum_{i=0}^{n-1}x_i f_*(\xi^{-i})=x_0,
\]
the last equality following from the fact that $f_*(\xi^i)=0$ for $-n+1\leq i<0$ (\hyperlink{sga6.e06.sec.2.8}{VI 2.8}). If $\varphi((x_i))=0$, what precedes shows that $x_0=0$. Applying the same reasoning to the elements $\xi^i\varphi((x_i))$, for $i=1,\ldots,n-1$, one obtains by induction that all the $x_i$ are zero. Thus $\varphi$ is injective.

\sgaStableTarget{sga6.e09.para.3-3}\textbf{3.3.} Let $S$ still be a noetherian scheme, and let $E$ be a locally free $\mathcal O_S$-module of rank $n$. Write $V=\mathbb V(E)$ for the vector bundle associated with $E$, contravariantly in $E$ (SGA 3 I 4.6), considered as a group scheme over $S$. Let $P$ be a torsor under $V$; we wish to compute $K_\bullet(P)$.

Recall how one embeds $P$ in a projective bundle over $S$. Torsors under $V$ correspond to extensions $F$ of $E$ by $\mathcal O_S$. Indeed, if one has an extension
\sgaStableTarget{sga6.e09.eq.3-3-1}\[
0\longrightarrow\mathcal O_S\xrightarrow{s}F\xrightarrow{u}E\longrightarrow 0,
\tag{3.3.1}
\]
one obtains a torsor under $V$ by considering the homomorphism $\mathbb V(F)\to\mathbb V(\mathcal O_S)$ defined by $s$, the section of $\mathbb V(\mathcal O_S)=\Spec(\mathcal O_S[T])$ defined by $T=1$, and taking
\[
 P=\mathbb V(F)\times_{\mathbb V(\mathcal O_S)}S,
\]
the operations of $V$ on $P$ being induced by those of $V$ on $\mathbb V(F)$.

Conversely, let $P$ be a torsor under $V$. Denote by $\underline P$ the sheaf of germs of homomorphisms from $S$ to $P$. It is a principal homogeneous sheaf under the sheaf of germs of homomorphisms from $S$ to $V$, which is identified with the dual $\check E$ of $E$. In the classical way this yields an extension of $\mathcal O_S$ by $\check E$, and by passing to duals one obtains \hyperlink{sga6.e09.eq.3-3-1}{(3.3.1)}.

Thus let $P$ be a torsor under $V$, corresponding to an extension $F$ of $E$ by $\mathcal O_S$. The algebra $A(P)$ of $P$ over $\mathcal O_S$ is
\[
 S(F)\otimes_{S(\mathcal O_S)}\mathcal O_S,
\]
that is,
\[
 S(F)/(s-1)S(F),
\]
which identifies with the ring of homogeneous fractions of degree $0$, $S(F)_{(s)}$. Consequently $P$ identifies with the open $D_+(s)$ of $\mathbb P(F)$. The complement of $P$ in $\mathbb P(F)$ is then $\operatorname{Proj}(S(F)/sS(F))$, which is identified, by \hyperlink{sga6.e09.eq.3-3-1}{(3.3.1)}, with $\mathbb P(E)$. Thus one has immersions
\[
\begin{tikzcd}[column sep=large,row sep=small]
\mathbb P(E) \arrow[r,hook] \arrow[dr] & \mathbb P(F) & P \arrow[l,hook'] \arrow[dl]\\
& S &
\end{tikzcd}
\]
where $\mathbb P(E)$ is closed in $\mathbb P(F)$, and $P$ is open and complementary to $\mathbb P(E)$.

\sgaStableTarget{sga6.e09.prop.3-4}\textbf{Proposition 3.4.} Let $S$ be a noetherian scheme, let $E$ be a locally free $\mathcal O_S$-module of rank $n$, let $V=\mathbb V(E)$ be the associated vector bundle, let $P$ be a torsor under $V$, and let $f$ be the structural morphism of $P$. Then the canonical homomorphism
\[
 f^*:K_\bullet(S)\longrightarrow K_\bullet(P)
\]
is an isomorphism.

With the identifications of \hyperlink{sga6.e09.para.3-3}{3.3}, one has the exact sequence
\sgaStableTarget{sga6.e09.eq.3-4-1}\[
 K_\bullet(\mathbb P(E))\xrightarrow{i_*}K_\bullet(\mathbb P(F))\longrightarrow K_\bullet(P)\longrightarrow 0,
\tag{3.4.1}
\]
where $i$ denotes the closed immersion $\mathbb P(E)\to\mathbb P(F)$ defined by $u$. By \hyperlink{sga6.e09.cor.3-2}{3.2},
\[
 K_\bullet(\mathbb P(E))=K_\bullet(S)\otimes_{K^\circ(S)}K^\circ(\mathbb P(E)),
\qquad
 K_\bullet(\mathbb P(F))=K_\bullet(S)\otimes_{K^\circ(S)}K^\circ(\mathbb P(F)).
\]
Since $\mathbb P(E)$ and $\mathbb P(F)$ are smooth, $i$ is a regular immersion, and hence we have a homomorphism
\[
 i_*:K^\circ(\mathbb P(E))\longrightarrow K^\circ(\mathbb P(F)).
\]
The homomorphism between the $K_\bullet$ groups is then obtained by tensoring the homomorphism between the $K^\circ$ groups by $K_\bullet(S)$. It is enough, to see this, to show that for $a\in K_\bullet(S)$ and $b\in K^\circ(\mathbb P(E))$ one has
\[
 i_*(b\,h^*(a))=i_*(b)g^*(a),
\]
where $g$ and $h$ are the structural morphisms of $\mathbb P(F)$ and $\mathbb P(E)$. Since $h=gi$,
\[
 i_*(b\,h^*(a))=i_*(b\,i^*g^*(a))=g^*(a)i_*(b)
\]
by \hyperlink{sga6.e04.eq.2-12-4}{IV 2.12.4}.

From \hyperlink{sga6.e09.eq.3-4-1}{(3.4.1)} one therefore gets
\[
 K_\bullet(P)\simeq K_\bullet(S)\otimes_{K^\circ(S)}
 \bigl(K^\circ(\mathbb P(F))/\operatorname{Im}K^\circ(\mathbb P(E))\bigr).
\]
Let $\xi_F$ and $\xi_E$ be the classes of the canonical invertible sheaves on $\mathbb P(F)$ and $\mathbb P(E)$. By \hyperlink{sga6.e06.thm.1-1}{VI 1.1}, $K^\circ(\mathbb P(F))$ (respectively $K^\circ(\mathbb P(E))$) has as basis over $K^\circ(S)$ the elements $\xi_F^k$ for $k=0,\ldots,-n$ (respectively $\xi_E^k$ for $k=0,\ldots,-n+1$). The projection formula gives
\[
 i_*(\xi_E^k)=i_*(i^*(\xi_F^k))=\xi_F^k i_*(1).
\]
If $I$ is the ideal of $\mathbb P(E)$ in $\mathbb P(F)$, then $I$ is invertible and $i_*(1)=1-I$. Moreover $I$ is defined by the ideal $J$ of $S(F)$, kernel of $S(F)\to S(E)$, which is generated by a submodule of $F$ isomorphic to $\mathcal O_S$. Hence
\[
 I\simeq \mathcal O_{\mathbb P(F)}(-1).
\]
The image of $K^\circ(\mathbb P(E))$ is thus the $K^\circ(S)$-submodule generated by the $\xi_F^k-\xi_F^{k-1}$, for $k=0,\ldots,-n+1$. Therefore
\[
 K^\circ(S)\longrightarrow K^\circ(\mathbb P(F))/\operatorname{Im}K^\circ(\mathbb P(E))
\]
is an isomorphism, and consequently so is $K_\bullet(S)\to K_\bullet(P)$.

\sgaStableTarget{sga6.e09.sec.4}\subsection*{4. Principal bundles under the groups $\mathrm{GL}(n)_S$}

\sgaStableTarget{sga6.e09.prop.4-1}\textbf{Proposition 4.1.} Let $S$ be a noetherian scheme, let $G=\mathrm{GL}(n)_S$, let $P$ be a torsor under $G$, and let $E$ be the locally free $\mathcal O_S$-module of rank $n$ corresponding to $P$. Denote by $E$ the class of $E$ in $K^\circ(S)$, and let $J$ be the $\lambda$-ideal generated by $E-n$. If $f$ is the structural morphism of $P$, then $f^*$, defined because $f$ is flat, is surjective and defines an isomorphism
\[
 K_\bullet(S)/J K_\bullet(S)\xrightarrow{\sim}K_\bullet(P).
\]

Let $E_0=(\mathcal O_S)^n$, with evident basis $e_1,\ldots,e_n$, and let $E_0^{(i)}$ be the submodule of $E_0$ generated by $e_1,\ldots,e_i$. The subfunctor of $G$ which to every $S$-scheme $S'$ assigns the subgroup of automorphisms of $E_{0,S'}$ leaving invariant the submodules $E_{0,S'}^{(i)}$ for all $i$ is represented by a group subscheme $B$ of $G$, a Borel subgroup of $G$. The quotient sheaf $G/B$ for the fpqc topology is represented by the flag scheme of $E_0$, denoted $X_0$. Let $L_0^{(i)}$ be the canonical quotients of the inverse image of $E_0$ on $X_0$, let $G_0^{(i)}$ be the kernel of $E_{0,X_0}\to L_0^{(i)}$, and put
\[
 F_0^{(i)}=G_0^{(i-1)}/G_0^{(i)}\simeq \ker(L_0^{(i)}\to L_0^{(i-1)}).
\]

Moreover, $B$ is the semi-direct product of its unipotent radical $U$ and of its maximal subtorus $T$ leaving the $e_i$ invariant. Let $Y_0=G/U$. Then $Y_0$ is a torsor over $X_0$ under the group $T_{X_0}$. If $S'$ is an $S$-scheme, the datum of an $S'$-point of $Y_0$ is equivalent to the datum, on some faithfully flat quasi-compact $S''$ over $S'$, of an automorphism of $E_{0,S''}$ modulo the unipotents of $U_{S''}$; that is, of a flag of $E_{0,S''}$ with free submodules, and, for each $i$, of a class of free elements $f_i^\alpha$ on $\mathcal O_{S''}$ such that $f_i^\alpha-f_i^\beta\in G_{0,S''}^{(i)}$. Equivalently, it is the datum of a flag of $E_{0,S''}$ and, for every $i$, of a basis of $F_{0,S''}^{(i)}$, this datum being subject to the gluing condition on $S''\times_{S'}S''$. By fpqc descent, this amounts to the same datum on $S'$. Thus $Y_0$ is the affine $X_0$-scheme defined by the algebra
\[
 A_0=\bigoplus_{k_i\in\mathbb Z}(F_0^{(1)})^{\otimes k_1}\otimes\cdots\otimes(F_0^{(n)})^{\otimes k_n}.
\]

The group $G$ operates on $E_0$, on $X_0$ provided with the sheaves $L_0^{(i)}$, and on $Y_0$. Taking the product of these objects with $P$ modulo $G$, one obtains the sheaf $E$ associated with $P$, the scheme $X=D(E)=P/B$ of flags of $E$, and the scheme $Y=P/U$. We denote the canonical sheaves on $X$ by the same letters as their analogues on $X_0$. Thus one has the diagram
\[
\begin{array}{ccccc}
&& Y=P/U &&\\[-0.2em]
& \nearrow && \searrow &\\[-0.2em]
P &&&& X=P/B=D(E)\\[-0.2em]
& \searrow && \swarrow &\\[-0.2em]
&& S &&
\end{array}
\]
We note that $P$ is a torsor over $Y$ with structural group $U_Y$, and that $Y$ is a torsor over $X$ with structural group $T_X$, defined by the $\mathcal O_X$-algebra
\[
 A=\bigoplus_{k_i\in\mathbb Z}(F^{(1)})^{\otimes k_1}\otimes\cdots\otimes(F^{(n)})^{\otimes k_n}.
\]

The homomorphism $K_\bullet(S)\to K_\bullet(P)$ therefore factors as
\[
 K_\bullet(S)\longrightarrow K_\bullet(X)\longrightarrow K_\bullet(Y)\longrightarrow K_\bullet(P).
\]
By \hyperlink{sga6.e09.prop.3-1}{3.1}, $K_\bullet(X)=K_\bullet(S)\otimes_{K^\circ(S)}K^\circ(X)$. Let $\xi_i$ be the class of $F^{(i)}$ in $K^\circ(X)$. By \hyperlink{sga6.e09.prop.2-2}{2.2}, $K_\bullet(Y)$ is the quotient of $K_\bullet(X)$ by the submodule generated by the $(\xi_i-1)K_\bullet(X)$. If $I$ is the ideal of $K^\circ(X)$ generated by the $\xi_i-1$, then
\[
 K_\bullet(Y)=K_\bullet(S)\otimes_{K^\circ(S)}(K^\circ(X)/I).
\]
By \hyperlink{sga6.e06.cor.4-8}{VI 4.8}, $K^\circ(X)$ is the $K^\circ(S)$-algebra generated by the $\xi_i-1$, subject to the relations deduced from
\[
 \gamma_t(E-n)=\prod_{i=1}^n\gamma_t(\xi_i-1).
\]
Therefore $K^\circ(X)/I=K^\circ(S)/J$, where $J$ is the ideal generated by the $\gamma^j(E-n)$ for $j=1,\ldots,n$, hence the $\lambda$-ideal generated by $E-n$. Thus
\[
 K_\bullet(Y)\simeq K_\bullet(S)/J K_\bullet(S).
\]
The demonstration is completed by the following lemma.

\sgaStableTarget{sga6.e09.lem.4-2}\textbf{Lemma 4.2.} Let $S$ be a noetherian scheme, let $B$ be a parabolic subgroup of a reductive $S$-group, and let $P$ be a torsor under the unipotent radical $U$ of $B$. Then one has an isomorphism
\[
 K_\bullet(S)\xrightarrow{\sim}K_\bullet(P).
\]

Indeed, there exists a finite sequence of subgroup schemes of $U$,
\[
 U_0=U\supset U_1\supset U_2\supset\cdots\supset U_p=e,
\]
closed in $B$, finite because $B$ is noetherian, and such that, for every $i$, $U_i/U_{i+1}$ is isomorphic to an $S$-vector group (SGA 3 XXVI 2.1). Consequently $P\to S$ factors as a sequence of torsors under vector groups, and the result follows from 3.4.

\sgaStableTarget{sga6.e09.cor.4-3}\textbf{Corollary 4.3.} Let $S$ be a noetherian scheme and let $G=\mathrm{GL}(n)_S$. Then the canonical homomorphism
\[
 K_\bullet(S)\longrightarrow K_\bullet(G)
\]
is an isomorphism.

This follows immediately from \hyperlink{sga6.e09.prop.4-1}{4.1}, applied to $P=G$, since then $E=\mathcal O_S^n$.

\sgaStableTarget{sga6.e09.cor.4-4}\textbf{Corollary 4.4.} Under the hypotheses of \hyperlink{sga6.e09.prop.4-1}{4.1}, suppose $S$ regular. Then one has an isomorphism
\[
 K^\circ(S)/J\xrightarrow{\sim}K^\circ(P).
\]

If $S$ is regular, $G$ is also regular, and $P$ as well. Thus $K^\circ(S)\simeq K_\bullet(S)$ and $K^\circ(P)\simeq K_\bullet(P)$, whence the result follows from 4.1.

\clearpage

\providecommand{\Filt}{}\renewcommand{\Filt}{\operatorname{Filt}}
\providecommand{\Coh}{}\renewcommand{\Coh}{\operatorname{Coh}}
\providecommand{\Supp}{}\renewcommand{\Supp}{\operatorname{Supp}}
\providecommand{\longg}{}\renewcommand{\longg}{\operatorname{length}}
\providecommand{\cl}{}\renewcommand{\cl}{\operatorname{cl}}
\providecommand{\Ker}{}\renewcommand{\Ker}{\operatorname{Ker}}
\providecommand{\codim}{}\renewcommand{\codim}{\operatorname{codim}}
\providecommand{\degtr}{}\renewcommand{\degtr}{\operatorname{deg.tr}}
\providecommand{\Pic}{}\renewcommand{\Pic}{\operatorname{Pic}}
\providecommand{\Spec}{}\renewcommand{\Spec}{\operatorname{Spec}}
\providecommand{\Gr}{}\renewcommand{\Gr}{\operatorname{Gr}}
\providecommand{\Gro}{}\renewcommand{\Gro}{\operatorname{Gr}^{\circ}}
\providecommand{\Grb}{}\renewcommand{\Grb}{\operatorname{Gr}_{\bullet}}
\providecommand{\fpqc}{}\renewcommand{\fpqc}{\mathrm{fpqc}}
\providecommand{\Sch}{}\renewcommand{\Sch}{\operatorname{Sch}}
\providecommand{\Ob}{}\renewcommand{\Ob}{\operatorname{Ob}}
\providecommand{\Tor}{}\renewcommand{\Tor}{\operatorname{Tor}}
\providecommand{\red}{}\renewcommand{\red}{\mathrm{red}}
\providecommand{\Mod}{}\renewcommand{\Mod}{\operatorname{Mod}}
\providecommand{\Aut}{}\renewcommand{\Aut}{\operatorname{Aut}}
\providecommand{\GL}{}\renewcommand{\GL}{\operatorname{GL}}
\providecommand{\SL}{}\renewcommand{\SL}{\operatorname{SL}}
\providecommand{\Gm}{}\renewcommand{\Gm}{\mathbb G_{\mathrm m}}
\providecommand{\Der}{}\renewcommand{\Der}{\operatorname{Der}}
\providecommand{\id}{}\renewcommand{\id}{\operatorname{id}}
\providecommand{\Ch}{}\renewcommand{\Ch}{\operatorname{Ch}}
\providecommand{\cO}{}\renewcommand{\cO}{\mathcal O}
\providecommand{\C}{}\renewcommand{\C}{\mathbb C}
\providecommand{\Q}{}\renewcommand{\Q}{\mathbb Q}
\providecommand{\Z}{}\renewcommand{\Z}{\mathbb Z}
\providecommand{\Zell}{}\renewcommand{\Zell}{\mathbf Z_\ell}
\providecommand{\Qell}{}\renewcommand{\Qell}{\mathbf Q_\ell}
\providecommand{\ch}{}\renewcommand{\ch}{\operatorname{ch}}
\providecommand{\Todd}{}\renewcommand{\Todd}{\operatorname{Todd}}
\providecommand{\Gtop}{}\renewcommand{\Gtop}{\operatorname{Gr}_{\mathrm{top}}}
\providecommand{\Fil}{}\renewcommand{\Fil}{\operatorname{Fil}}
\providecommand{\Qcoh}{}\renewcommand{\Qcoh}{\operatorname{Qcoh}}
\providecommand{\Parf}{}\renewcommand{\Parf}{\operatorname{Parf}}
\providecommand{\Lib}{}\renewcommand{\Lib}{\operatorname{Lib}}
\providecommand{\Loclib}{}\renewcommand{\Loclib}{\operatorname{Loclib}}
\providecommand{\Proj}{}\renewcommand{\Proj}{\operatorname{Proj}}
\providecommand{\Hom}{}\renewcommand{\Hom}{\operatorname{Hom}}
\providecommand{\Diff}{}\renewcommand{\Diff}{\operatorname{Diff}}
\providecommand{\Dcat}{}\renewcommand{\Dcat}{\mathrm D}
\providecommand{\Tr}{}\renewcommand{\Tr}{\operatorname{Tr}}
\providecommand{\toramp}{}\renewcommand{\toramp}{\operatorname{tor.amp}}
\providecommand{\parfamp}{}\renewcommand{\parfamp}{\operatorname{par.amp}}
\providecommand{\coh}{}\renewcommand{\coh}{\operatorname{coh}}
\providecommand{\qcoh}{}\renewcommand{\qcoh}{\operatorname{qcoh}}
\providecommand{\parf}{}\renewcommand{\parf}{\operatorname{parf}}
\providecommand{\torf}{}\renewcommand{\torf}{\operatorname{torf}}
\providecommand{\amp}{}\renewcommand{\amp}{\operatorname{amp}}
\providecommand{\num}{}\renewcommand{\num}{\operatorname{num}}
\providecommand{\naif}{}\renewcommand{\naif}{\operatorname{naif}}
\providecommand{\alg}{}\renewcommand{\alg}{\operatorname{alg}}

\hbadness=1000
\hfuzz=0.1pt
\allowdisplaybreaks

\sgaStableTarget{sga6.e10.expose}\section*{Expos\'e X. Formalism of Intersections on Proper Algebraic Schemes}
\begin{center}
by O. Jussila\\
(with an appendix by A. Grothendieck)
\end{center}

\sgaStableTarget{sga6.e10.sec.1}\section*{1. Compatibility of the filtrations with the composition law}
\[
K^\circ(X)\times K_\bullet(X)\longrightarrow K_\bullet(X).
\]

\sgaStableTarget{sga6.e10.sec.1.1}\subsection*{1.1}
First let $X$ be a noetherian scheme, and let $K_\bullet(X)$ be the group of
classes of coherent $\mathcal O_X$-modules. For every coherent
$\mathcal O_X$-module $F$, we denote its class in $K_\bullet(X)$ by
$\cl_\bullet(F)$.

\sgaStableTarget{sga6.e10.def.1-1-1}\textbf{Definition 1.1.1.}
For every integer $n\in\mathbb Z$, let $\Filt_n(X)$ denote the subgroup of
$K_\bullet(X)$ generated by the classes $\cl_\bullet(F)$ of coherent
$\mathcal O_X$-modules $F$ such that
\[
\dim(\Supp(F))\leq n.
\]
This defines an increasing filtration $(\Filt_n(X))_{n\in\mathbb Z}$ of
$K_\bullet(X)$ such that $\Filt_n(X)=0$ for $n<0$.

\sgaStableTarget{sga6.e10.prop.1-1-2}\textbf{Proposition 1.1.2.}
Let $X$ be a noetherian scheme, let $F$ be a coherent $\mathcal O_X$-module
such that $\dim(\Supp(F))\leq n$, and, for $n\in\mathbb Z$, let $X(n)$ be the
subset
\[
\{x\in X\,;\,\dim(\overline{\{x\}})=n\}
\]
of $X$. For every $x\in X$, let $Z(x)$ denote the integral closed subscheme
associated with $\overline{\{x\}}$. Then the length of the
$\mathcal O_{X,x}$-module $F_x$ is finite for every $x\in X(n)$ and is
nonzero for only finitely many $x\in X(n)$, and one has
\sgaStableTarget{sga6.e10.eq.1-1-3}\begin{equation}\tag{1.1.3}
\cl_\bullet(F)=
\sum_{x\in X(n)}\longg(F_x)\cl_\bullet(\mathcal O_{Z(x)})+w,
\end{equation}
where $w\in\Filt_{n-1}(X)$.

By immediate induction on $n$, one obtains:

\sgaStableTarget{sga6.e10.cor.1-1-4}\textbf{Corollary 1.1.4.}
For every noetherian scheme $X$ and every integer $n$, $\Filt_n(X)$ is the
subgroup of $K_\bullet(X)$ generated by the classes
$\cl_\bullet(\mathcal O_Z)$, where $Z$ ranges over the integral closed
subschemes of $X$ of dimension at most $n$.

\emph{Proof of 1.1.2.}
Let $X'=\Supp(F)$. Since $X'$ is a noetherian space of dimension at most $n$,
the elements of $X'\cap X(n)$ are maximal in $X'$ and hence finite in number;
moreover, $\longg(F_x)$ is finite for every $x\in X'\cap X(n)$. This proves
the first assertion.

Let $\mathcal K=\Coh(X)$ be the category of coherent $\mathcal O_X$-modules,
and let $\mathcal K'$ be the subset of $\operatorname{Ob}(\mathcal K)$ formed
by the coherent $\mathcal O_X$-modules concentrated on $X'$ that satisfy
condition \hyperlink{sga6.e10.eq.1-1-3}{(1.1.3)}. Plainly $\mathcal K'$ contains the sheaves
$\mathcal O_{Z(x)}$ for $x\in X'$. It is therefore enough to show that
$\mathcal K'$ is an ``exact'' subset of $\operatorname{Ob}(\mathcal K)$,
for then the devissage lemma (EGA III 3.1.2) gives $F\in\mathcal K'$.

Thus let
\sgaStableTarget{sga6.e10.eq.1-1-5}\begin{equation}\tag{1.1.5}
0\longrightarrow G'\longrightarrow G\longrightarrow G''\longrightarrow0
\end{equation}
be an exact sequence of $\mathcal O_X$-modules having two of its terms in
$\mathcal K'$. The sheaves $G'$, $G$, and $G''$ are then all concentrated on
$X'$. Hence the lengths $\ell'_x$, $\ell_x$, and $\ell''_x$ of the
$\mathcal O_{X,x}$-modules $G'_x$, $G_x$, and $G''_x$ are finite for every
$x\in X(n)$, and
\sgaStableTarget{sga6.e10.eq.1-1-6}\begin{equation}\tag{1.1.6}
\ell_x=\ell'_x+\ell''_x\qquad(x\in X(n)).
\end{equation}
Assertion \hyperlink{sga6.e10.eq.1-1-3}{(1.1.3)} for all the terms of \hyperlink{sga6.e10.eq.1-1-5}{(1.1.5)} now follows immediately from
the hypothesis and the equation
\sgaStableTarget{sga6.e10.eq.1-1-7}\begin{equation}\tag{1.1.7}
\cl_\bullet(G)=\cl_\bullet(G')+\cl_\bullet(G'').
\end{equation}

It is known that $K_\bullet$ is a covariant functor on the category of locally
noetherian schemes with respect to proper morphisms. The following proposition
asserts that our filtration is functorial.

\sgaStableTarget{sga6.e10.prop.1-1-8}\textbf{Proposition 1.1.8.}
Let $f:X'\to X$ be a proper morphism of locally noetherian schemes, and let
\[
f_*:K_\bullet(X')\longrightarrow K_\bullet(X)
\]
be the morphism induced by $f$. Then
\sgaStableTarget{sga6.e10.eq.1-1-9}\begin{equation}\tag{1.1.9}
f_*(\Filt_n(X'))\subset\Filt_n(X)
\end{equation}
for every integer $n$.

\emph{Proof.}
If $F$ is a coherent $\mathcal O_{X'}$-module, then by definition (\hyperlink{sga6.e04.para.2-11}{IV 2.11}),
\[
f_*(\cl_\bullet(F))=
\sum_{i\in\mathbb N}(-1)^i\cl_\bullet(R^if_*(F)).
\]
The $\mathcal O_X$-modules $R^if_*(F)$ are concentrated on
$f(\Supp(F))$, and by EGA IV 5.4.1(ii),
\[
\dim(\Supp(F))\geq\dim(f(\Supp(F))),
\]
which proves the assertion.

\sgaStableTarget{sga6.e10.sec.1.2}\subsection*{1.2}
Let now $X$ be an arbitrary scheme, and consider the group $K^\circ(X)$ of
classes of locally free $\mathcal O_X$-modules (henceforth understood to be of
finite type). It is an augmented pre-$\lambda$-algebra over the ring
$H^0(X,\mathbb Z)$ (\hyperlink{sga6.e05.ex.3-9-1}{V 3.9.1}). The map
$f\rightsquigarrow f\cdot1\in K^\circ(X)$ allows us henceforth to identify
$H^0(X,\mathbb Z)$ with a subring of $K^\circ(X)$. If $E$ is a locally free
$\mathcal O_X$-module, we denote its class in $K^\circ(X)$ by
$\cl^\circ(E)$. The operators $\lambda^k$ ($k\in\mathbb N$) are defined by
\sgaStableTarget{sga6.e10.eq.1-2-1}\begin{equation}\tag{1.2.1}
\lambda^k(\cl^\circ(E))=\cl^\circ(\wedge_{\mathcal O}^kE)
\end{equation}
and by
\sgaStableTarget{sga6.e10.eq.1-2-2}\begin{equation}\tag{1.2.2}
\lambda^k(x+y)=
\sum_{\substack{i+j=k\\ i,j\in\mathbb N}}\lambda^i(x)\lambda^j(y).
\end{equation}

(\hyperlink{sga6.e05.ex.2-2}{V 2.2}). Recall that in the $\lambda$-filtration
$(\Filt^i(X))_{i\in\mathbb Z}$ of $K^\circ(X)$, one has
$\Filt^i(X)=K^\circ(X)$ for $i\leq0$, while for $i\geq1$, $\Filt^i(X)$ is
the ideal generated by the products
\sgaStableTarget{sga6.e10.eq.1-2-3}\begin{equation}\tag{1.2.3}
P=\gamma^{i_1}(x_1)\cdots\gamma^{i_k}(x_k)
\end{equation}
with $1\leq i_1,\ldots,i_k$; $i_1+\cdots+i_k\geq i$; and
$x_1,\ldots,x_k\in I=\Ker(\varepsilon)$, where $\varepsilon$ is the
augmentation homomorphism (cf. \hyperlink{sga6.e05.para.3-10}{Expos\'e V 3.10)}. It was proved in \hyperlink{sga6.e05.para.3-10}{V 3.10} that
$\Filt^i(X)$ is even the $H^0(X,\mathbb Z)$-submodule, hence, when $X$ is
connected, the abelian subgroup generated by the products \hyperlink{sga6.e10.eq.1-2-3}{(1.2.3)}, where one
may moreover suppose that the elements $x_p$ ($p=1,\ldots,k$) have the form
\sgaStableTarget{sga6.e10.eq.1-2-4}\begin{equation}\tag{1.2.4}
x_p=\cl^\circ(E_p)-n_p,
\end{equation}
with $E_p$ a locally free $\mathcal O_X$-module of rank
$n_p\in H^0(X,\mathbb Z)$.

\sgaStableTarget{sga6.e10.sec.1.3}\subsection*{1.3}
Recall (\hyperlink{sga6.e04.para.2-10}{IV 2.10}) that the bifunctor
\[
(E,F)\rightsquigarrow E\otimes_{\mathcal O_X}F,
\]
where $E$ is a locally free $\mathcal O_X$-module and $F$ is a coherent
$\mathcal O_X$-module, induces a composition law
\begin{equation}\tag{\sgaStableTargetInline{sga6.e10.eq.1-3-1}1.3.1}
K^\circ(X)\times K_\bullet(X)\longrightarrow K_\bullet(X)
\end{equation}
that makes $K_\bullet(X)$ a unital $K^\circ(X)$-module. The key result of the
present expos\'e is the following.

\sgaStableTarget{sga6.e10.thm.1-3-2}\textbf{Theorem 1.3.2.}
Let $X$ be a noetherian scheme. For all $i,j\in\mathbb Z$, the composition law
\hyperlink{sga6.e10.eq.1-3-1}{(1.3.1)} induces by restriction a composition law
\sgaStableTarget{sga6.e10.eq.1-3-3}\begin{equation}\tag{1.3.3}
\Filt^i(X)\times\Filt_j(X)\longrightarrow\Filt_{j-i}(X).
\end{equation}

\sgaStableTarget{sga6.e10.cor.1-3-3}\textbf{Corollary 1.3.3.}
If $X$ is a noetherian scheme of dimension at most $j$, then the canonical
homomorphism
\[
\Theta_X:K^\circ(X)\longrightarrow K_\bullet(X),
\]
defined by $\Theta_X(x)=x\cdot\cl_\bullet(\mathcal O_X)$
($x\in K^\circ(X)$), induces for every $i\in\mathbb Z$ a homomorphism
\[
\Filt^i(X)\longrightarrow\Filt_{j-i}(X).
\]

\emph{Proof.}
First let us show that \hyperlink{sga6.e10.cor.1-3-3}{(1.3.3)}, for $i+j\leq k\in\mathbb Z$ and every integral
closed subscheme of dimension at most $j$, implies \hyperlink{sga6.e10.thm.1-3-2}{Theorem 1.3.2} for
$i+j\leq k$. By \hyperlink{sga6.e10.cor.1-1-4}{(1.1.4)}, it is enough to show that, for every
$x\in\Filt^i(X)$ and every integral closed subscheme $Z$ of dimension at most
$j$, one has
\[
x\cl_\bullet(\mathcal O_Z)\in\Filt_{j-i}(X).
\]
To avoid confusion, let $1_Z$ denote the class of $\mathcal O_Z$ in
$K_\bullet(Z)$. If $\alpha:Z\to X$ is the canonical immersion, the projection
formula (IV \ldots) gives
\[
x\cl_\bullet(\mathcal O_Z)
=x\alpha_*(1_Z)
=\alpha_*(\alpha^*(x)1_Z).
\]
Since $\alpha^*(x)\in\Filt^i(Z)$, \hyperlink{sga6.e10.prop.1-1-8}{(1.1.8)} and \hyperlink{sga6.e10.cor.1-3-3}{(1.3.3)}, applied with $X=Z$
and $i+j\leq k$, give
\[
\alpha_*(\alpha^*(x)1_Z)\in\Filt_{j-i}(X),
\]
which proves the assertion.

\hyperlink{sga6.e10.thm.1-3-2}{Theorem 1.3.2} is trivial for $i+j\leq0$. Suppose by induction that it has been
proved for $i+j<k$ and for every noetherian scheme $X$. By the preceding
remark, it remains to prove that, for every integral noetherian scheme $X$ of
dimension $j$ and every product \hyperlink{sga6.e10.eq.1-2-3}{(1.2.3)} satisfying \hyperlink{sga6.e10.eq.1-2-4}{(1.2.4)}, one has
\sgaStableTarget{sga6.e10.eq.1-3-4}\begin{equation}\tag{1.3.4}
P\cl_\bullet(\mathcal O_X)\in\Filt_{j-i}(X).
\end{equation}
If $k>1$ in \hyperlink{sga6.e10.eq.1-2-3}{(1.2.3)}, the assertion follows from the induction hypothesis and
the associativity of the law \hyperlink{sga6.e10.eq.1-3-1}{(1.3.1)}. We may therefore suppose, since $X$ is
connected, that
\sgaStableTarget{sga6.e10.eq.1-3-5}\begin{equation}\tag{1.3.5}
x=P=\gamma^m(\cl^\circ(E)-n),
\end{equation}
where $m\geq i$ and $E$ is a locally free $\mathcal O_X$-module of rank
$n>0$. We distinguish two cases.

I. Suppose that there exist invertible $\mathcal O_X$-modules
$L_1,\ldots,L_n$ such that
\[
\cl^\circ(E)=\cl^\circ(L_1)+\cdots+\cl^\circ(L_n)\in K^\circ(X).
\]

It follows from (V) that $x$ is the $m$-th elementary symmetric function in
the $\cl^\circ(L_r)-1$ ($r=1,\ldots,n$), and hence is a homogeneous
polynomial of degree $m$ with integral coefficients in the variables
$(\cl^\circ(L_r)-1)$ ($r=1,\ldots,n$). Applying again the induction
hypothesis and the associativity of the law \hyperlink{sga6.e10.eq.1-3-1}{(1.3.1)}, we are reduced to proving
that, for every invertible $\mathcal O_X$-module $L$, one has
\sgaStableTarget{sga6.e10.eq.1-3-6}\begin{equation}\tag{1.3.6}
(\cl^\circ(L)-1)\cl_\bullet(\mathcal O_X)\in\Filt_{j-1}(X).
\end{equation}
Since $X$ is integral, $L$ and $\mathcal O_X$ are isomorphic to submodules of
the constant $\mathcal O_X$-module $R_X$ of rational functions on $X$. The
$\mathcal O_X$-module sum $F=\mathcal O_X+L$ is coherent, and one obtains
\begin{align*}
(\cl^\circ(L)-1)\cl_\bullet(\mathcal O_X)
&=\cl_\bullet(L)-\cl_\bullet(\mathcal O_X)\\
&=(\cl_\bullet(F)-\cl_\bullet(\mathcal O_X))
  -(\cl_\bullet(F)-\cl_\bullet(L))\\
&=\cl_\bullet(F/\mathcal O_X)-\cl_\bullet(F/L).
\end{align*}
But the fibers at the generic point of the sheaves $F/\mathcal O_X$ and
$F/L$ vanish; their supports in $X$ are therefore nowhere dense (since $X$
is irreducible), and hence have dimension at most $j-1$. Assertion \hyperlink{sga6.e10.eq.1-3-6}{(1.3.6)}
follows immediately.

II. It remains to prove the general case. Introduce the flag bundle
\[
g:D\longrightarrow X
\]
of $E$ over $X$, on which there exist invertible $\mathcal O_D$-modules
$L_1,\ldots,L_n$ such that
\sgaStableTarget{sga6.e10.eq.1-3-7}\begin{equation}\tag{1.3.7}
g^*(\cl^\circ(E))=\cl^\circ(L_1)+\cdots+\cl^\circ(L_n)\in K^\circ(D).
\end{equation}
Let $e$ be the generic point of $X$, and let $k(e)$ be its residue field.
There exists a rational point $u$ of the generic fiber $D_e$ over $k(e)$,

because these points are in one-to-one correspondence with flags of type
$(1,\ldots,1)$ ($n$ terms) in the vector space $E_e$ of dimension $n>0$ over
$k(e)$. These rational points also correspond to germs of sections of $g$ in
a neighborhood of $e$, so there exists such a section $s:U\to D$ over an
open subset $U$ of $X$. Let $X'=\overline{s(U)}$, endowed with its reduced
induced structure, and let $f:X'\to X$ be the restriction of $g$ to $X'$.
The morphism $f$ is projective and, since it is an isomorphism over $s(U)$,
it is birational.

Thus $\dim(X')\leq\dim(X)$ (EGA IV 5.6.6.1). On the other hand,
$f(X')=g(X')$ is closed and contains $U$, so $f$ is surjective and
$\dim(X')\geq\dim(X)$ (EGA IV 5.4.1(ii)). The hypotheses of case I are
therefore satisfied for $X'$ and $x'=f^*(x)\in\Filt^i(X')$, and one obtains
\sgaStableTarget{sga6.e10.eq.1-3-7.occ2}\begin{equation}\tag{1.3.7'}
f^*(x)\cl_\bullet(\mathcal O_{X'})\in\Filt_{j-i}(X').
\end{equation}
Consequently, by the projection formula and \hyperlink{sga6.e10.prop.1-1-8}{(1.1.8)},
\sgaStableTarget{sga6.e10.eq.1-3-8}\begin{equation}\tag{1.3.8}
x f_*(\cl_\bullet(\mathcal O_{X'}))\in\Filt_{j-i}(X).
\end{equation}
On the other hand, the canonical homomorphism
\[
\mathcal O_X\longrightarrow f_*(\mathcal O_{X'})
\]
is an isomorphism over the dense open subset $U$, and the
$\mathcal O_X$-modules $R^if_*(\mathcal O_{X'})$ ($i\geq1$) are concentrated
on the nowhere-dense closed subset $X-U$. It follows that
\sgaStableTarget{sga6.e10.eq.1-3-9}\begin{equation}\tag{1.3.9}
f_*(\cl_\bullet(\mathcal O_{X'}))=\cl_\bullet(\mathcal O_X)+w,
\end{equation}

where $w\in\Filt_{j-1}(X)$. By the induction hypothesis,
$x\cdot w\in\Filt_{j-i-1}(X)$; hence, by \hyperlink{sga6.e10.eq.1-3-8}{(1.3.8)} and \hyperlink{sga6.e10.eq.1-3-9}{(1.3.9)},
\[
x\cl_\bullet(\mathcal O_X)\in\Filt_{j-i}(X).
\]
This proves the assertion.

\sgaStableTarget{sga6.e10.rem.1-4}\textbf{Remark 1.4.}
Let $X$ be a locally noetherian scheme. Observe that, in this section, we
have used only the following properties of the ``dimension function''
\[
D(x)=D_Y(x)=\dim(\overline{\{x\}})
\]
on proper $X$-schemes $Y$:
\begin{enumerate}[label=(\roman*)]
\item For every proper $X$-scheme $Y$ and every $x\in Y$, one has
  $D_Y(x)\in\mathbb Z$.
\item If $y\ne x$ is a specialization of $x$, then $D(y)<D(x)$.
\item For every $X$-morphism $f:Y_1\to Y_2$ of proper $X$-schemes and every
  $y\in Y_1$, one has
  \[
  D_{Y_2}(f(y))\leq D_{Y_1}(y).
  \]
\item Equality holds in (iii) if $f$ is an immersion, or if $Y_1$ is
  irreducible, $f$ is birational and surjective, and $y$ is the generic point
  of $Y_1$.
\end{enumerate}
Thus all the results of this section, in particular \hyperlink{sga6.e10.thm.1-3-2}{(1.3.2)} and \hyperlink{sga6.e10.cor.1-3-3}{(1.3.3)},
remain valid unchanged if ``dimension'' is replaced by a function $D$
satisfying conditions (i)--(iv) above and the filtration of $K_\bullet(X)$ is
defined using such a function $D$.

\sgaStableTarget{sga6.e10.ex.1-5}\textbf{Example 1.5.}
Let $S$ be a locally noetherian universally catenary scheme (for example, a
subscheme of a Cohen--Macaulay scheme). Put
\sgaStableTarget{sga6.e10.eq.1-5-1}\begin{equation}\tag{1.5.1}
D_S(x)=n-\codim(\overline{\{x\}},S)
      =n-\dim\mathcal O_{S,x}\qquad(x\in S),
\end{equation}
and, for every proper $S$-scheme $s_Y:Y\to S$ and every $y\in Y$, also put
\sgaStableTarget{sga6.e10.eq.1-5-2}\begin{equation}\tag{1.5.2}
D_Y(y)=D_S(s_Y(y))+\degtr_{k(s_Y(y))}k(y).
\end{equation}
If $f:Y_1\to Y_2$ is an $S$-morphism of proper $S$-schemes, the additivity of
transcendence degrees and \hyperlink{sga6.e10.eq.1-5-2}{(1.5.2)} immediately give
\sgaStableTarget{sga6.e10.eq.1-5-3}\begin{equation}\tag{1.5.3}
D_{Y_1}(y)=D_{Y_2}(f(y))+\degtr_{k(f(y))}k(y)
\qquad(y\in Y_1).
\end{equation}
We show that the functions $D_Y$ so defined satisfy conditions (i)--(iv) of
1.4. By \hyperlink{sga6.e10.eq.1-5-3}{(1.5.3)}, condition (ii) is the only nontrivial one. In view of (iv),
we may plainly suppose that $Y=\overline{\{x\}}$. Factor the structural
morphism $s_Y$ through its image $s_Y(Y)=Z\hookrightarrow S$, say
$s_Y=g\circ f$:
\[
\begin{tikzcd}
Y \arrow[r,"f"] & Z \arrow[r,hook,"g"] & S,
\end{tikzcd}
\]
where $Y$ and $Z$ are irreducible, $f$ is proper and surjective, and $g$ is a
closed immersion. Put $x'=s_Y(x)$ and $y'=s_Y(y)$. From \hyperlink{sga6.e10.eq.1-5-2}{(1.5.2)} and \hyperlink{sga6.e10.eq.1-5-1}{(1.5.1)}
one obtains
\[
\begin{split}
D_Y(x)-D_Y(y)={}&\dim\mathcal O_{S,y'}-\dim\mathcal O_{S,x'}\\
&+\degtr_{k(x')}k(x)-\degtr_{k(y')}k(y).
\end{split}
\]
Since $S$ is catenary, one has

\[
\dim\mathcal O_{S,y'}-\dim\mathcal O_{S,x'}
=\dim\mathcal O_{Z,y'}-\dim\mathcal O_{Z,x'}
\]
(EGA $0_{\mathrm{IV}}$ 14.3.2.1), whence
\sgaStableTarget{sga6.e10.eq.1-5-4}\begin{equation}\tag{1.5.4}
\begin{split}
D_Y(x)-D_Y(y)={}&
[\dim\mathcal O_{Z,y'}-\degtr_{k(y')}k(y)]\\
&-[\dim\mathcal O_{Z,x'}-\degtr_{k(x')}k(x)].
\end{split}
\end{equation}
Since $Z$ is universally catenary, EGA IV 5.6.5.1 applies and gives
\[
C_1=\dim\mathcal O_{Y,y}-e,
\qquad
C_2=\dim\mathcal O_{Y,x}-e,
\]
where $C_1$ and $C_2$ are the two terms on the right of \hyperlink{sga6.e10.eq.1-5-4}{(1.5.4)}. It follows
immediately from \hyperlink{sga6.e10.eq.1-5-4}{(1.5.4)} that
\[
D_Y(x)-D_Y(y)
=\dim\mathcal O_{Y,y}-\dim\mathcal O_{Y,x}>0,
\]
as required.

Thus \hyperlink{sga6.e10.thm.1-3-2}{(1.3.2)} and \hyperlink{sga6.e10.cor.1-3-3}{(1.3.3)} remain valid for a noetherian universally catenary
scheme $X$ of dimension $n$ if the filtration of $K_\bullet(X)$ is defined by
means of \hyperlink{sga6.e10.eq.1-5-1}{(1.5.1)}, according to the codimension of supports. This filtration
coincides with that of \hyperlink{sga6.e10.def.1-1-1}{(1.1.1)} when, in addition, $X$ is biequidimensional.

\sgaStableTarget{sga6.e10.sec.1.6}\subsection*{1.6}
One may ask whether it is always possible, on a noetherian scheme $X$ of
finite dimension, to define a function $D$ possessing property (i) of \hyperlink{sga6.e10.rem.1-4}{1.4}
and also the following property, stronger than (ii):
\begin{quote}
(ii bis) If $x\in X$ and $y\ne x$ is an immediate specialization of $x$,
then $D(x)=D(y)+1$.
\end{quote}

For this it is plainly necessary that $X$ be catenary, and sufficient that
$X$ be biequidimensional (EGA $0_{\mathrm{IV}}$ 14.3.3). There nevertheless
exist noetherian schemes of finite dimension that are catenary and
equidimensional but admit no such function $D$. To see this, modify the
example of EGA IV 5.6.11 as follows. Take two copies $Y_1$ and $Y_2$ of the
spectrum of the ring $E$ of loc. cit., and let $y_i$ and $z_i$ ($i=1,2$) be
the closed points of $Y_i$ corresponding respectively to the maximal ideals
$\mathfrak m$ and $\mathfrak m'$ of $E$, of heights $2$ and $1$. Recall that
the residue fields of $\mathfrak m$ and $\mathfrak m'$ are isomorphic. One
can therefore construct a scheme $X$ by gluing $Y_1$ and $Y_2$, identifying
$y_1$ with $z_2$ and $y_2$ with $z_1$. The schemes $Y_1$ and $Y_2$ may then
be identified with the two irreducible components of $X$; inside $Z$,
$Y_1\cap Y_2=\{y_1,y_2\}$ has dimension zero. Thus $X$ is a noetherian,
catenary, equidimensional scheme of dimension $2$. If (ii bis) held on $X$,
one would have
\[
D(y_1)=D(z_1)-1=D(y_2)-1,
\]
and, on the other hand,
\[
D(y_1)=D(z_2)=D(y_2)+1,
\]
which is a contradiction.

\sgaStableTarget{sga6.e10.sec.2}\section*{2. Snapper polynomials}

\sgaStableTarget{sga6.e10.sec.2.1}\subsection*{2.1}
We first recall the definition and a few elementary properties of polynomial
maps.

\sgaStableTarget{sga6.e10.def.2-1-1}\textbf{Definition 2.1.1.}
Let $f$ be a map from an abelian group $A$ to an abelian group $B$. We say
that $f$ is a polynomial map if, for every sequence $x_1,\ldots,x_r$
($r\geq1$) of elements of $A$, there is a family
\[
\bigl(y(x_1,\ldots,x_r;k_1,\ldots,k_r)\bigr)_{k_1,\ldots,k_r\geq0}
\]
of elements of $B$, almost all zero, satisfying the following condition: for
every sequence $n_1,\ldots,n_r$ of integers, one has
\[
f(n_1x_1+\cdots+n_rx_r)=
\sum_{k_1,\ldots,k_r\geq0}
n_1^{k_1}\cdots n_r^{k_r}
y(x_1,\ldots,x_r;k_1,\ldots,k_r).
\]
We say that $f$ has degree $\leq m$ if, for every integer $r\geq1$ and all
$x_1,\ldots,x_r\in A$, one has
\[
y(x_1,\ldots,x_r;k_1,\ldots,k_r)=0
\]
whenever $k_1+\cdots+k_r>m$.

\sgaStableTarget{sga6.e10.lem.2-1-2}\textbf{Lemma 2.1.2.}
Every homomorphism of abelian groups is a polynomial map of degree $\leq1$.
If $h=f\circ g$, where $f$ is a polynomial map of degree $p$ and $g$ is a
polynomial map of degree $q$, then $h$ is a polynomial map of degree
$\leq pq$.

The proof is trivial.

\sgaStableTarget{sga6.e10.lem.2-1-3}\textbf{Lemma 2.1.3.}
Let
\[
0\longrightarrow G'\longrightarrow G\xrightarrow{j}G''\longrightarrow0
\]
be an exact sequence of abelian groups such that $G'$ is $p$-divisible for
some prime $p$ and $G''$ is a torsion group. Then every polynomial map
$f:G\to\mathbb Q$ that factors through $\mathbb Z$ is constant.

\emph{Proof.}
It is enough to show that $f(x)=f(0)$ for every $x\in G$.

(a) Suppose first that $x$ is a torsion element, say $mx=0$
($m\in\mathbb Z$, $m\ne0$). Then the polynomial
$P_x(n)=f(nx)-f(0)$ has infinitely many roots, namely all integers divisible
by $m$. Thus $P_x$ is identically zero, and $f(x)=f(0)$.

(b) Suppose now that $x$ is an element of infinite order in $G$, divisible by every power
$p^k$ ($k\in\mathbb Z$). Consider the polynomials $P_k$ defined by
\[
P_k(n)=f(nx/p^k)=\sum_{i\in\mathbb N}n^iy(x/p^k;i)\in\mathbb Z
\qquad(n\in\mathbb Z;\ y(x/p^k;i)\in\mathbb Q).
\]
For every $n\in\mathbb Z$, one has the identity
\[
P_k(n)=P_{k+1}(pn),
\]
that is,
\[
\sum_{i\in\mathbb N}n^iy(x/p^k;i)=
\sum_{i\in\mathbb N}p^in^iy(x/p^{k+1};i).
\]
Comparing coefficients gives
\[
y(x/p^k;i)=p^iy(x/p^{k+1};i)\qquad(i\in\mathbb N).
\]

Every polynomial $P:\mathbb Q\to\mathbb Q$ is continuous for the topology
induced by $\mathbb R$. Thus, for $k$ sufficiently large,
\[
\lvert P_x(n/p^k)-P_x(0)\rvert<1.
\]
But, by definition, $P_x(n/p^k)$ is an integer for all
$n,k\in\mathbb Z$. It follows that $P_x(n/p^k)-P_x(0)=0$ for sufficiently
large $k$. Hence $P_x$ is constant and, in particular,
$P_x(1)=P_x(0)$, i.e. $f(x)=f(0)$.

The two cases above show that the restriction of $f$ to $G'$ is constant.
We now show that $f$ is constant on each coset of $G'$ in $G$. For every
$b\in G$, let $g_b:G'\to\mathbb Q$ be defined by
\[
g_b(y)=f(b+y)\qquad(y\in G').
\]
The map $g_b$ is a polynomial map from $G'$ to $\mathbb Q$ that factors
through $\mathbb Z$, and is therefore constant by the preceding argument.
Thus $f$ is constant on the cosets of $G'$ in $G$. Consequently one may
write $f=h\circ j$, where $h:G''\to\mathbb Q$ is a polynomial map that
factors through $\mathbb Z$. Since $G''$ is a torsion group, $h$ is constant.
Therefore $f$ is constant.

\sgaStableTarget{sga6.e10.lem.2-1-4}\textbf{Lemma 2.1.4.}
Let $A$ be a unital commutative ring and $I$ a nilpotent ideal of $A$. Let
$P$ be the multiplicative subgroup of $A^*$ consisting of the invertible elements
of $A$ of the form $x=1+y$ with $y\in I$, and let $i:P\to A$ be the natural
injection.

Then the composite map
\[
\alpha:P\xrightarrow{i}A\xrightarrow{j}A\otimes_{\mathbb Z}\mathbb Q,
\]
where $j(x)=x\otimes1$, is a polynomial map. If $I^{m+1}=0$, then $\alpha$
has degree $\leq m$.

\emph{Proof.}
Let $x_1,\ldots,x_r$ be a sequence of elements of $P$, with
$x_i=1+y_i$ ($i=1,\ldots,r$). The binomial formula gives, for every
$n_1,\ldots,n_r\in\mathbb Z$,
\sgaStableTarget{sga6.e10.eq.2-1-5}\begin{equation}\tag{2.1.5}
x_1^{n_1}\cdots x_r^{n_r}=
\sum_{k_1,\ldots,k_r\geq0}
\binom{n_1}{k_1}\cdots\binom{n_r}{k_r}
y_1^{k_1}\cdots y_r^{k_r},
\end{equation}
where the sum on the right is finite because the elements
$y_1,\ldots,y_r$ are nilpotent. The coefficients
$\binom{n_1}{k_1}\cdots\binom{n_r}{k_r}$ are polynomials with rational
coefficients in $n_1,\ldots,n_r$, of degree $k_1+\cdots+k_r$; hence
$\alpha$ is indeed a polynomial map. If $I^{m+1}=0$, then
$y_1^{k_1}\cdots y_r^{k_r}=0$ whenever $k_1+\cdots+k_r>m$. Thus $\alpha$
has degree $\leq m$.

\sgaStableTarget{sga6.e10.sec.2.2}\subsection*{2.2}
Now let $X$ be a noetherian scheme and let $y\in\Filt_j(X)$ \hyperlink{sga6.e10.def.1-1-1}{(1.1.1)}. The
assignment $x\mapsto xy$ defines a homomorphism of $K^\circ(X)$-modules
\[
\Theta_y:K^\circ(X)\longrightarrow K_\bullet(X),
\]
which, by \hyperlink{sga6.e10.thm.1-3-2}{(1.3.2)}, factors through the canonical epimorphism
\[
j:K^\circ(X)\longrightarrow
K^\circ(X)/\Filt^{j+1}(X)=A.
\]

Let $\beta_y:A\to K_\bullet(X)$ be the homomorphism so defined. If $I$
denotes the ideal
\[
I=\Filt^1(X)/\Filt^{j+1}(X)
\]
of $A$, then $I^{j+1}=0$. Keeping the notation of \hyperlink{sga6.e10.lem.2-1-4}{Lemma 2.1.4}, one obtains a
homomorphism of abelian groups
\[
\gamma:\Pic(X)\longrightarrow P
\]
that sends the class $\cl_P(L)$ of an invertible $\mathcal O_X$-module $L$
in $\Pic(X)$ to the element $j(\cl^\circ(L))$ of $P$.

Lemmas \hyperlink{sga6.e10.lem.2-1-2}{2.1.2} and \hyperlink{sga6.e10.lem.2-1-4}{2.1.4} immediately give:

\sgaStableTarget{sga6.e10.prop.2-2-1}\textbf{Proposition 2.2.1.}
For every noetherian scheme $X$ and every $y\in\Filt_j(X)$
($j\in\mathbb Z$), the composite map
\[
\Pic(X)\longrightarrow K_\bullet(X)\otimes_{\mathbb Z}\mathbb Q,
\]
defined by
\[
\cl_P(L)\longmapsto(\cl^\circ(L)y)\otimes1,
\]
is a polynomial map of degree $\leq j$.

\sgaStableTarget{sga6.e10.sec.2.3}\subsection*{2.3}
If $X$ is a proper scheme over a field $k$ and
$g:X\to\Spec(k)$ is the structural morphism, then the homomorphism
\[
g_*:K_\bullet(X)\longrightarrow K_\bullet(\Spec(k))\simeq\mathbb Z
\]
sends the class $\cl_\bullet(F)$ of every coherent $\mathcal O_X$-module
$F$ to the Euler--Poincar\'e characteristic $\chi(F)\in\mathbb Z$ of $F$
on $X$. \hyperlink{sga6.e10.prop.2-2-1}{Proposition 2.2.1} immediately gives:

\sgaStableTarget{sga6.e10.prop.2-3-1}\textbf{Proposition 2.3.1.}
For every scheme $X$ proper over $k$ and every $y\in\Filt_j(X)$, the map
\[
S_y:\Pic(X)\longrightarrow\mathbb Z\hookrightarrow\mathbb Q,
\]
defined by
\[
\cl_P(L)\longmapsto
g_*(\cl^\circ(L)y)=\chi(\cl^\circ(L)y),
\]
is a polynomial map of degree $\leq j$.

When $L_1,\ldots,L_n$ are invertible $\mathcal O_X$-modules and $F$ is a
coherent $\mathcal O_X$-module such that $\dim(\Supp(F))=j$, \hyperlink{sga6.e10.prop.2-3-1}{Proposition
2.3.1} implies in particular that the Euler--Poincar\'e characteristic
\sgaStableTarget{sga6.e10.eq.2-3-2}\begin{equation}\tag{2.3.2}
\chi(L_1^{\otimes m_1}\otimes\cdots\otimes
L_n^{\otimes m_n}\otimes F)
\end{equation}
is a numerical polynomial in $m_1,\ldots,m_n\in\mathbb Z$, of degree
$\leq\dim(\Supp(F))$, sometimes called the Snapper polynomial of $F$ and
the sequence $L_1,\ldots,L_n$.

\sgaStableTarget{sga6.e10.sec.2.4}\subsection*{2.4}
Recall [2] that, for the $\fpqc$ topology on $\Sch/k$, the Picard functor
$\Pic_{X/k}$ is the sheaf of groups associated with the presheaf
\[
T\longmapsto\Pic(X\times_kT)
\qquad(T\in\Ob(\Sch/k)).
\]
When $k$ is algebraically closed, the canonical homomorphism
\[
\Pic(X)\longrightarrow\Pic_{X/k}(\Spec(k))
\]
is an isomorphism.

When $X$ is proper over $k$, it is known [4] that the Picard functor is
representable by a group scheme, locally of finite type over $k$, also denoted
$\Pic_{X/k}$. The neutral component $\Pic^0_{X/k}$ of $\Pic_{X/k}$ is a
connected commutative algebraic group over $k$.

\sgaStableTarget{sga6.e10.def.2-4-1}\textbf{Definition 2.4.1.}
Let $X$ be a scheme proper over a field $k$, and let $L$ be an invertible
$\mathcal O_X$-module (respectively, let $x\in\Pic(X)$). We say that $L$
(respectively, $x$) is $\tau$-equivalent to zero if there is an integer
$n\ne0$ such that $L^{\otimes n}$ (respectively, $nx$) is algebraically
equivalent to zero. The elements of $\Pic(X)$ algebraically equivalent
(respectively, $\tau$-equivalent) to zero form a subgroup $\Pic^0(X)$
(respectively, $\Pic^\tau(X)$) of $\Pic(X)$.

\sgaStableTarget{sga6.e10.lem.2-4-3}\textbf{Lemma 2.4.3.}
For every proper scheme $X$ over an algebraically closed field $k$, the group
$\Pic^0(X)$ is divisible by every integer $n$ prime to the characteristic of
$k$.

\emph{Proof.}
Since $k$ is algebraically closed, one has
\[
\Pic^0(X)\simeq\Pic^0_{X/k}(k).
\]
Since $\Pic^0_{X/k}$ is connected and of finite type, the $n$-th-power
endomorphism is surjective when $n$ is prime to the characteristic of $k$
(SGA 3 XV 1.3 and VI 1.3.2). Since $k$ is algebraically closed, the
corresponding endomorphism of $\Pic^0_{X/k}(k)$ is surjective.

It follows from \hyperlink{sga6.e10.lem.2-4-3}{2.4.3} that, in the exact sequence
\[
0\longrightarrow\Pic^0(X)\longrightarrow\Pic^\tau(X)
\longrightarrow T\longrightarrow0,
\]
$\Pic^0(X)$ is $p$-divisible for some prime $p$ and $T$ is a torsion group.
Thus \hyperlink{sga6.e10.lem.2-1-3}{Lemma 2.1.3} applies to the restriction of the polynomial map \hyperlink{sga6.e10.prop.2-3-1}{(2.3.1)} to
$\Pic^\tau(X)$. One obtains:

\sgaStableTarget{sga6.e10.prop.2-4-4}\textbf{Proposition 2.4.4.}
The polynomial map $S_y$ of \hyperlink{sga6.e10.prop.2-3-1}{(2.3.1)} is constant on the cosets of
$\Pic^\tau(X)$ in $\Pic(X)$.

\emph{Proof.}
If $K$ is an algebraically closed extension of $k$, one has a commutative
diagram
\[
\begin{tikzcd}[column sep=large,row sep=large]
\Pic(X) \arrow[r,"S_y"] \arrow[d] &
K_\bullet(\Spec(k)) \arrow[r] &
\mathbb Z \arrow[r,hook] & \mathbb Q\\
\Pic(X_K) \arrow[r,"\overline S_y"'] &
K_\bullet(\Spec(K)) \arrow[ur] &&
\end{tikzcd}
\]
where $\overline S_y$ is the polynomial map on $\Pic(X_K)$ corresponding to
$S_y$. Hence one may suppose that $k$ is algebraically closed. The
assertion then follows immediately from Lemmas \hyperlink{sga6.e10.lem.2-1-3}{2.1.3} and 2.4.3.

\sgaStableTarget{sga6.e10.rem.2-4-5}\textbf{Remark 2.4.5.}
In \hyperlink{sga6.e13.expose}{Expos\'e XIII} this result will be proved without using representability of
the functor $\Pic_{X/k}$.

\sgaStableTarget{sga6.e10.sec.3}\section*{3. Projection formulas for the associated graded groups}

\sgaStableTarget{sga6.e10.sec.3.1}\subsection*{3.1}
Let $X$ be a noetherian scheme. Denote by $\Gro(X)$ the graded group associated
with the \mbox{$\lambda$-filtration} of $K^\circ(X)$, and by $\Grb(X)$ the graded
group associated with the filtration \hyperlink{sga6.e10.def.1-1-1}{1.1.1} of $K_\bullet(X)$. By \hyperlink{sga6.e10.thm.1-3-2}{Theorem
1.3.2}, the composition law \hyperlink{sga6.e10.eq.1-3-1}{(1.3.1)}
\[
K^\circ(X)\times K_\bullet(X)\longrightarrow K_\bullet(X)
\]
induces, on passing to the associated graded groups, composition laws
\sgaStableTarget{sga6.e10.eq.3-1-1}\begin{equation}\tag{3.1.1}
\Gr^i(X)\times\Gr_j(X)\longrightarrow\Gr_{j-i}(X)
\qquad(i,j\in\mathbb Z).
\end{equation}

Thus $\Grb(X)$ becomes a module over the ring $\Gro(X)$, and even a graded
module over the graded ring $\Gro(X)$, provided that the sign of the grading
of $\Grb(X)$ is changed by putting
\[
\Gr^i_\bullet(X)=\Gr_{-i}(X)\qquad(i\in\mathbb Z).
\]

Now let
\[
f:X'\longrightarrow X
\]
be a proper morphism of noetherian schemes. Then $f$ induces a homomorphism
\[
f^*:\Gro(X)\longrightarrow\Gro(X')
\]
of graded rings (\hyperlink{sga6.e05.expose}{Expos\'e V}), and, by \hyperlink{sga6.e10.prop.1-1-8}{Proposition 1.1.8}, a homomorphism
\[
f_*:\Grb(X')\longrightarrow\Grb(X)
\]
of graded abelian groups. Passing the projection formula (\hyperlink{sga6.e04.para.2-11}{IV 2.11}) and
\hyperlink{sga6.e10.eq.3-1-1}{(3.1.1)} to the associated graded groups gives the graded projection formula
\sgaStableTarget{sga6.e10.eq.3-1-2}\begin{equation}\tag{3.1.2}
x\cdot f_*(y)=f_*\bigl(f^*(x)\cdot y\bigr)
\qquad(x\in\Gro(X),\ y\in\Grb(X')).
\end{equation}

\sgaStableTarget{sga6.e10.sec.3.2}\subsection*{3.2}
Let $f:X'\to X$ be a projective complete-intersection morphism (\hyperlink{sga6.e08.expose}{Expos\'e
VIII}), of virtual dimension $r$, between noetherian schemes, with $X$ carrying
an ample invertible $\mathcal O_X$-module. In \hyperlink{sga6.e08.expose}{Expos\'e VIII} it was proved that
$f$ induces a homomorphism
\sgaStableTarget{sga6.e10.eq.3-2-1}\begin{equation}\tag{3.2.1}
f_*:\Filt^k(X')_{\mathbb Q}\longrightarrow
\Filt^{k-r}(X)_{\mathbb Q}\qquad(k\in\mathbb Z),
\end{equation}
which moreover satisfies the projection formula (\hyperlink{sga6.e04.expose}{Expos\'e IV}):
\sgaStableTarget{sga6.e10.eq.3-2-2}\begin{equation}\tag{3.2.2}
f_*\bigl(f^*(x)\cdot y\bigr)=x\cdot f_*(y)
\qquad(x\in K^\circ(X)_{\mathbb Q},\ y\in K^\circ(X')_{\mathbb Q}).
\end{equation}

Passing to the associated graded groups in \hyperlink{sga6.e10.eq.3-2-1}{(3.2.1)} and \hyperlink{sga6.e10.eq.3-2-2}{(3.2.2)}, one obtains a
homomorphism of degree $-r$:
\sgaStableTarget{sga6.e10.eq.3-2-3}\begin{equation}\tag{3.2.3}
f_*:\Gro(X')_{\mathbb Q}\longrightarrow\Gro(X)_{\mathbb Q},
\end{equation}
and, by \hyperlink{sga6.e10.thm.1-3-2}{Theorem 1.3.2}, a graded projection formula
\[
f_*\bigl(f^*(x)\cdot y\bigr)=x\cdot f_*(y)
\qquad(x\in\Gro(X)_{\mathbb Q},\ y\in\Gro(X')_{\mathbb Q}).
\]

\sgaStableTarget{sga6.e10.sec.3.3}\subsection*{3.3}
Finally, let $f:X'\to X$ be a morphism of finite Tor-dimension between
noetherian schemes. It is known (\hyperlink{sga6.e04.expose}{Expos\'e IV}) that $f$ induces a homomorphism
\[
f^*:K_\bullet(X)\longrightarrow K_\bullet(X')
\]
such that, for every coherent $\mathcal O_X$-module $F$,
\[
f^*(\cl_\bullet(F))=
\sum_i(-1)^i\cl_\bullet\bigl(
\Tor_i^{\mathcal O_X}(F,\mathcal O_{X'})\bigr).
\]
The $\mathcal O_{X'}$-modules
$\Tor_i^{\mathcal O_X}(F,\mathcal O_{X'})$ are supported on
$f^{-1}(\Supp(F))$. If the fibers of $f$ are everywhere of dimension
$\leq d$, then (EGA IV 5.6.7)
\[
\dim\bigl(f^{-1}(\Supp(F))\bigr)\leq\dim(\Supp(F))+d.
\]
It follows that $f$ induces, for every $k\in\mathbb Z$, a homomorphism
\[
f^*:\Filt_k(X)\longrightarrow\Filt_{k+d}(X'),
\]
hence a homomorphism of degree $d$ of graded abelian groups.

Thus one has
\sgaStableTarget{sga6.e10.eq.3-3-1}\begin{equation}\tag{3.3.1}
f^*:\Grb(X)\longrightarrow\Grb(X').
\end{equation}

\sgaStableTarget{sga6.e10.sec.3.4}\subsection*{3.4}
Suppose finally that the morphisms \hyperlink{sga6.e10.eq.3-2-3}{(3.2.3)} and \hyperlink{sga6.e10.eq.3-3-1}{(3.3.1)} are both defined for a
morphism $f:X'\to X$ and that $r=d$. This is the case, for example, if $X$
is quasi-compact with an ample invertible $\mathcal O_X$-module and if
$X'=\mathbf P_X(E)$, where $E$ is a locally free $\mathcal O_X$-module of
rank $r+1$. The projection formula in the derived category then gives, for
every $x\in\Filt^i(X')_{\mathbb Q}$ and $y\in\Filt_j(X)_{\mathbb Q}$,
\[
f_*\bigl(x\cdot f^*(y)\bigr)=f_*(x)\cdot y
\in\Filt_{j+r-i}(X)_{\mathbb Q}.
\]
When $X$ is noetherian, \hyperlink{sga6.e10.thm.1-3-2}{Theorem 1.3.2} yields the corresponding graded formula.

\sgaStableTarget{sga6.e10.sec.4}\section*{4. Intersection numbers}

\sgaStableTarget{sga6.e10.sec.4.1}\subsection*{4.1}
Let $X$ be a proper scheme over a field $k$. Composing the product \hyperlink{sga6.e10.eq.3-1-1}{(3.1.1)}
\[
\Gr^i(X)\times\Gr_i(X)\longrightarrow\Gr_0(X)=\Filt_0(X)
\qquad(i\in\mathbb Z)
\]
with the Euler--Poincar\'e characteristic
\[
\chi:\Filt_0(X)\longrightarrow\mathbb Z
\]
gives a pairing called the \emph{intersection number}:
\sgaStableTarget{sga6.e10.eq.4-1-1}\begin{equation}\tag{4.1.1}
\Gr^i(X)\times\Gr_i(X)\longrightarrow\mathbb Z.
\end{equation}

We denote this pairing by
\[
(x,y)\longmapsto\langle x,y\rangle\in\mathbb Z
\qquad(i\in\mathbb Z,\ x\in\Gr^i(X),\ y\in\Gr_i(X)).
\]

\sgaStableTarget{sga6.e10.ex.4-2}\textbf{Example 4.2.}
Suppose that $X$ has dimension $n$. Let $y\in\Gr_n(X)$, for example
$y=\cl_\bullet(\mathcal O_Z)$, where $Z$ is an integral closed subscheme of
$X$ of dimension $n$. Let $P$ also be a homogeneous polynomial of weight $n$,
\[
P\in\mathbb Z[T_1,\ldots,T_n],
\]
where $T_i$ is assigned weight $i$ ($i=1,\ldots,n$); for example, $P$ may
contain a monomial $T_{i_1}\cdots T_{i_r}$ with
$i_1+\cdots+i_r=n$. Substituting for each variable $T_i$ the Chern class
\[
c_i(x)=\cl^i\bigl(\gamma^i(x-\varepsilon(x))\bigr)\in\Gr^i(X)
\qquad(x\in K^\circ(X))
\]
gives an element
\[
P(c_1(x),\ldots,c_n(x))\in\Gr^n(X).
\]
The integer
\[
P_y(x)=\bigl\langle P(c_1(x),\ldots,c_n(x)),y\bigr\rangle
\]
is called the \emph{Chern number} of $x$ relative to $y$. When
$y=\cl_\bullet(\mathcal O_X)$, the reference to $y$ is omitted from this
terminology.

\sgaStableTarget{sga6.e10.sec.4.3}\subsection*{4.3}
Without any hypothesis on the dimension of $X$, consider invertible
$\mathcal O_X$-modules $\mathcal L_1,\ldots,\mathcal L_d$ and a coherent
$\mathcal O_X$-module $\mathcal F$ such that
$\dim(\Supp(\mathcal F))\leq d$. Put
\[
L_i=\cl^\circ(\mathcal L_i)\quad(i=1,\ldots,d),
\qquad F=\cl_\bullet(\mathcal F).
\]

\sgaStableTarget{sga6.e10.prop.4-3-1}\textbf{Proposition 4.3.1.}
With the preceding notation, for every sequence of nonnegative integers
$\alpha_1,\ldots,\alpha_d$ such that
$\sum_{i=1}^d\alpha_i=d$, the coefficient of
\[
\binom{n_1}{\alpha_1}\cdots\binom{n_d}{\alpha_d}
\]
in the Snapper polynomial
\[
\chi(\mathcal L_1^{\otimes n_1}\otimes\cdots\otimes
\mathcal L_d^{\otimes n_d}\otimes\mathcal F)
\]
is equal to
\sgaStableTarget{sga6.e10.eq.4-3-2}\begin{equation}\tag{4.3.2}
\bigl\langle c_1(L_1)^{\alpha_1}\cdots
c_1(L_d)^{\alpha_d},\lvert F\rvert\bigr\rangle,
\end{equation}
where $\lvert F\rvert$ is the class of $F$ in $\Gr_d(X)$. Thus, in the ordinary
expansion of this polynomial in monomials in the $n_i$ ($i=1,\ldots,d$),
the coefficient of $n_1^{\alpha_1}\cdots n_d^{\alpha_d}$ is
\[
  \frac{1}{\alpha_1!\cdots\alpha_d!}
\bigl\langle c_1(L_1)^{\alpha_1}\cdots
c_1(L_d)^{\alpha_d},\lvert F\rvert\bigr\rangle.
\]

In particular, one obtains immediately:

\sgaStableTarget{sga6.e10.cor.4-3-3}\textbf{Corollary 4.3.3.}
The coefficient of the monomial $n_1n_2\cdots n_d$ in the Snapper polynomial
\[
\chi(\mathcal L_1^{\otimes n_1}\otimes\cdots\otimes
\mathcal L_d^{\otimes n_d}\otimes\mathcal F)
\]
is equal to
\sgaStableTarget{sga6.e10.eq.4-3-4}\begin{equation}\tag{4.3.4}
\bigl\langle c_1(L_1)\cdots c_1(L_d),\lvert F\rvert\bigr\rangle
=\bigl(c_1(L_1)\cdots c_1(L_d),\lvert F\rvert\bigr).
\end{equation}

\emph{Proof.}
Put $L_i=1+M_i$ for $i=1,\ldots,d$. By the definitions, the integer \hyperlink{sga6.e10.eq.4-3-2}{(4.3.2)}
is equal to
\[
\chi(M_1^{\alpha_1}\cdots M_d^{\alpha_d}F).
\]
Substituting
\[
A=K^\circ(X)/\Filt^{d+1}(X),
\qquad I=\Filt^1(X)/\Filt^{d+1}(X)
\]
in \hyperlink{sga6.e10.lem.2-1-4}{Lemma 2.1.4}, formula \hyperlink{sga6.e10.eq.2-1-5}{(2.1.5)} gives, since
$M_i\in\Filt^1(X)$ for every $i=1,\ldots,d$, the congruence
\sgaStableTarget{sga6.e10.eq.4-3-5}\begin{equation}\tag{4.3.5}
L_1^{n_1}\cdots L_d^{n_d}\equiv
\sum_{\substack{\alpha_1,\ldots,\alpha_d\geq0\\
\alpha_1+\cdots+\alpha_d\leq d}}
\binom{n_1}{\alpha_1}\cdots\binom{n_d}{\alpha_d}
M_1^{\alpha_1}\cdots M_d^{\alpha_d}
\pmod{\Filt^{d+1}(X)}.
\end{equation}
Hence the coefficient of
$\binom{n_1}{\alpha_1}\cdots\binom{n_d}{\alpha_d}$ in
$\chi(\mathcal L_1^{\otimes n_1}\otimes\cdots\otimes
\mathcal L_d^{\otimes n_d}\otimes\mathcal F)$ is
$\chi(M_1^{\alpha_1}\cdots M_d^{\alpha_d}F)$, as required.

\sgaStableTarget{sga6.e10.sec.4.4}\subsection*{4.4}
The integer \hyperlink{sga6.e10.eq.4-3-4}{(4.3.4)} is sometimes called the \emph{intersection number} of the
invertible $\mathcal O_X$-modules $\mathcal L_1,\ldots,\mathcal L_d$ with
$\mathcal F$, and is denoted by
\[
(\mathcal L_1,\ldots,\mathcal L_d;\mathcal F)
\qquad(\emph{cf.}~[3]).
\]
Expression \hyperlink{sga6.e10.eq.4-3-4}{(4.3.4)} immediately gives the following properties.

\begin{enumerate}[label=\arabic*)]
\item The integer \hyperlink{sga6.e10.eq.4-3-4}{(4.3.4)} is a symmetric $d$-linear function of
  the classes $L_1,\ldots,L_d\in\Pic(X)$, is zero if
  $\dim(\Supp(\mathcal F))<d$, and is additive in $\mathcal F$, that is, in
  its class $\cl_\bullet(\mathcal F)=F$.

\item By \hyperlink{sga6.e10.prop.1-1-2}{Proposition 1.1.2}, one obtains
\sgaStableTarget{sga6.e10.eq.4-4-1}\begin{equation}\tag{4.4.1}
(\mathcal L_1,\ldots,\mathcal L_d;\mathcal F)=
\sum_{x\in X(d)}
\operatorname{length}_{\mathcal O_{X,x}}(\mathcal F_x)
(\mathcal L_1,\ldots,\mathcal L_d;\mathcal O_{Z(x)}).
\end{equation}

\item Let $Y$ be a closed subscheme of $X$, and let $\mathcal G$ be a coherent
  $\mathcal O_Y$-module such that $\mathcal F=f_*(\mathcal G)$, where
  $f:Y\to X$ is the
  canonical immersion. Then
\[
(f^*\mathcal L_1,\ldots,f^*\mathcal L_d;\mathcal G)
=(\mathcal L_1,\ldots,\mathcal L_d;\mathcal F).
\]
More generally, the projection formula gives
\sgaStableTarget{sga6.e10.eq.4-4-2}\begin{equation}\tag{4.4.2}
\chi(f^*(x)\cdot y)
=\chi\bigl(f_*(f^*(x)\cdot y)\bigr)
=\chi(x\cdot f_*(y))
\qquad(x\in K^\circ(X),\ y\in K_\bullet(Y))
\end{equation}
for every $k$-morphism $f:Y\to X$ of proper $k$-schemes.

\item Suppose that, in the morphism $f:Y\to X$ above, $X$ and $Y$ are
  irreducible of dimension $\leq d$. If $f$ has degree $n$, then

\sgaStableTarget{sga6.e10.eq.4-4-3}\begin{equation}\tag{4.4.3}
(f^*\mathcal L_1,\ldots,f^*\mathcal L_d;\mathcal O_Y)
=n(\mathcal L_1,\ldots,\mathcal L_d;\mathcal O_X).
\end{equation}
Recall the definition of the degree $n$ of $f$ [3]: it is zero unless
\[
\dim(X)=\dim(Y)=\dim(f(Y)).
\]
In the latter case, one puts
\sgaStableTarget{sga6.e10.eq.4-4-4}\begin{equation}\tag{4.4.4}
n=
\frac{\operatorname{length}_{\mathcal O_{X,x}}(\mathcal O_{Y,y})}
{\operatorname{length}_{\mathcal O_{X,x}}(\mathcal O_{X,x})},
\end{equation}
where $x$ is the generic point of $X$ and $y$ is the generic point of $Y$.
\end{enumerate}

In the proof of \hyperlink{sga6.e10.eq.4-4-3}{(4.4.3)}, the case $\dim(X)=\dim(Y)$ is the only nontrivial
one. If $\dim(f(Y))<d$, then
\sgaStableTarget{sga6.e10.eq.4-4-5}\begin{equation}\tag{4.4.5}
R^if_*(\mathcal O_Y)=0
\end{equation}
for $i>0$ in a neighborhood of $x$. Condition \hyperlink{sga6.e10.eq.4-4-5}{(4.4.5)} also holds when
$\dim(f(Y))=d$, since the generic fiber then has dimension zero (EGA III
4.2.2 and EGA IV 5.5.6). Proposition 1.1.2 therefore gives the congruences
\[
f_*(\cl_\bullet(\mathcal O_Y))
\equiv\cl_\bullet(f_*\mathcal O_Y)
\equiv\operatorname{length}_{\mathcal O_{X,x}}(\mathcal O_{Y,y})
\cl_\bullet(\mathcal O_{X_{\red}})
\pmod{\Filt_{d-1}(X)}.
\]
By \hyperlink{sga6.e10.eq.4-4-2}{(4.4.2)}, it follows that
\[
(f^*\mathcal L_1,\ldots,f^*\mathcal L_d;\mathcal O_Y)
=\operatorname{length}_{\mathcal O_{X,x}}(\mathcal O_{Y,y})
(\mathcal L_1,\ldots,\mathcal L_d;\mathcal O_{X_{\red}}).
\]
On the other hand, item 2 gives
\[
(\mathcal L_1,\ldots,\mathcal L_d;\mathcal O_X)
=\operatorname{length}_{\mathcal O_{X,x}}(\mathcal O_{X,x})
(\mathcal L_1,\ldots,\mathcal L_d;\mathcal O_{X_{\red}}),
\]
which proves \hyperlink{sga6.e10.eq.4-4-3}{(4.4.3)}.

\sgaStableTarget{sga6.e10.sec.4.5}\subsection*{4.5}
The intersection product $\langle\, ,\,\rangle$ defines an equivalence
relation called \emph{numerical equivalence}.

\sgaStableTarget{sga6.e10.def.4-5-1}\textbf{Definition 4.5.1.}
Two elements $x$ and $y$ of $\Gr^i(X)$ ($i\in\mathbb Z$), respectively of
$\Gr_i(X)$, are said to be \emph{numerically equivalent} if they define the
same homomorphism
\[
\Gr_i(X)\longrightarrow\mathbb Z,
\qquad\text{respectively}\qquad
\Gr^i(X)\longrightarrow\mathbb Z.
\]
Two invertible $\mathcal O_X$-modules are said to be numerically equivalent
if their first Chern classes in $\Gr^1(X)$ are numerically equivalent.

\sgaStableTarget{sga6.e10.prop.4-5-2}\textbf{Proposition 4.5.2.}
An invertible $\mathcal O_X$-module $\mathcal L$ is numerically equivalent to
zero if and only if
\sgaStableTarget{sga6.e10.eq.4-5-3}\begin{equation}\tag{4.5.3}
\chi(\mathcal L^{\otimes n}\otimes_{\mathcal O_X}\mathcal O_Z)
=\chi(\mathcal O_Z)
\end{equation}
for every $n\in\mathbb Z$ and every integral closed curve $Z\subset X$.

This follows from Propositions \hyperlink{sga6.e10.prop.1-1-2}{1.1.2} and 4.3.1.

Let $\Pic^N(X)$ denote the subgroup of $\Pic(X)$ consisting of the classes of
invertible $\mathcal O_X$-modules numerically equivalent to zero.
\hyperlink{sga6.e10.prop.2-4-4}{Proposition 2.4.4} gives:

\sgaStableTarget{sga6.e10.cor.4-5-4}\textbf{Corollary 4.5.4.}
\[
\Pic^\tau(X)\subset\Pic^N(X).
\]
It will be proved in \hyperlink{sga6.e13.expose}{Expos\'e XIII} that in fact
\[
\Pic^N(X)=\Pic^\tau(X).
\]

\sgaStableTarget{sga6.e10.sec.5}\section*{5. The isomorphism $\Pic(X)\simeq\Gr^1(X)$}

\sgaStableTarget{sga6.e10.sec.5.1}\subsection*{5.1}
Let $X$ be an arbitrary ringed topos. The ring $K^\circ(X)$ and the
filtration $(\Filt^i(X))_{i\in\mathbb Z}$ can be defined exactly as for
schemes. Likewise, define $\Pic(X)$ to be the group of isomorphism classes of
invertible $\mathcal O_X$-modules. Let

\sgaStableTarget{sga6.e10.eq.5-1-1}\begin{equation}\tag{5.1.1}
i:\Pic(X)\longrightarrow K^\circ(X)
\end{equation}
be the map sending the class $\cl_P(\mathcal L)\in\Pic(X)$ of an invertible
$\mathcal O_X$-module $\mathcal L$ to
$\cl^\circ(\mathcal L)\in K^\circ(X)$. This map is plainly a homomorphism
from $\Pic(X)$ to the multiplicative group $K^\circ(X)^*$ of units of
$K^\circ(X)$. Henceforth the group law of $\Pic(X)$ is written
multiplicatively.

The map $i$ has a left inverse defined by the determinant functor. If
$\mathcal E$ is a locally free $\mathcal O_X$-module of rank $n$, put
\[
\det(\mathcal E)=\bigwedge\nolimits^n\mathcal E.
\]
If
\sgaStableTarget{sga6.e10.eq.5-1-2}\begin{equation}\tag{5.1.2}
0\longrightarrow\mathcal E'\longrightarrow\mathcal E
\longrightarrow\mathcal E''\longrightarrow0
\end{equation}
is an exact sequence of locally free $\mathcal O_X$-modules, there is a
canonical isomorphism
\[
\det(\mathcal E)\xrightarrow{\sim}
\det(\mathcal E')\otimes_{\mathcal O_X}\det(\mathcal E''),
\]
arising from \hyperlink{sga6.e10.eq.6-2-9}{(6.2.9)} and the fact that \hyperlink{sga6.e10.eq.5-1-2}{(5.1.2)} is locally split. One therefore
obtains a homomorphism of abelian groups
\sgaStableTarget{sga6.e10.eq.5-1-3}\begin{equation}\tag{5.1.3}
\det:K^\circ(X)\longrightarrow\Pic(X)
\end{equation}
which is the required left inverse of $i$. In particular, $i$ is injective,
and henceforth $\Pic(X)$ is identified with its image in $K^\circ(X)^*$.

\sgaStableTarget{sga6.e10.sec.5.2}\subsection*{5.2}
Define another map
\sgaStableTarget{sga6.e10.eq.5-2-1}\begin{equation}\tag{5.2.1}
\delta:\Pic(X)\longrightarrow\Filt^1(X)
\end{equation}

by putting, for every $L\in\Pic(X)$,
\[
\delta(L)=1-L^{-1}.
\]

\sgaStableTarget{sga6.e10.rem.5-2-2}\textbf{Remarks 5.2.2.}
\begin{enumerate}[label=\alph*)]
\item Suppose that $X$ is the locally ringed topos of Zariski sheaves on a
scheme $S$. If $D$ is an effective Cartier divisor on $S$ and
\[
L=\cl^\circ(\mathcal O_S(D)),
\]
then
\sgaStableTarget{sga6.e10.eq.5-2-3}\begin{equation}\tag{5.2.3}
\Theta_S(\delta(L))=\cl_\bullet(\mathcal O_D),
\end{equation}
where $\Theta_S:K^\circ(X)\to K_\bullet(S)$ is the canonical homomorphism of
1.2.3. This follows immediately from the exact sequence
\[
0\longrightarrow\mathcal O_S(-D)\longrightarrow\mathcal O_S
\longrightarrow\mathcal O_D\longrightarrow0.
\]

\item Under the hypotheses of \hyperlink{sga6.e10.sec.5.1}{5.1} and \hyperlink{sga6.e10.sec.5.2}{5.2}, one has, for every
$L\in\Pic(X)$,
\sgaStableTarget{sga6.e10.eq.5-2-4}\begin{equation}\tag{5.2.4}
\delta(L)\equiv L-1\pmod{\Filt^2(X)},
\end{equation}
which justifies the notation $c_1(L)$ introduced below. Indeed,
\[
(L-1)-(1-L^{-1})=(L-1)(1-L^{-1})\in\Filt^2(X).
\]
\end{enumerate}

\sgaStableTarget{sga6.e10.sec.5.3}\subsection*{5.3}
Composing $\delta$ with the canonical epimorphism
\[
\Filt^1(X)\longrightarrow\Gr^1(X)
\]
gives the \emph{first Chern class} on $\Pic(X)$:
\sgaStableTarget{sga6.e10.eq.5-3-1}\begin{equation}\tag{5.3.1}
c_1:\Pic(X)\longrightarrow\Gr^1(X).
\end{equation}

\sgaStableTarget{sga6.e10.thm.5-3-2}\textbf{Theorem 5.3.2.}
The map $c_1$ in \hyperlink{sga6.e10.eq.5-3-1}{(5.3.1)} is an injective homomorphism, and the homomorphism
\[
\det:\Filt^1(X)\longrightarrow\Pic(X)
\]
induces on passing to the quotient a left inverse
\[
\alpha:\Gr^1(X)\longrightarrow\Pic(X)
\]
of $c_1$. If every section $n$ of the constant sheaf $\mathbb Z_X$ is
bounded---for example, if $X$ is quasi-compact, or if the final object $e$ of
$X$ is a finite disjoint sum of connected objects---then $c_1$ is an
isomorphism and $\alpha$ is its inverse.

\emph{Proof.}
The map $c_1$ is a homomorphism, since for all $L,L'\in\Pic(X)$ one has
\[
(1-(LL')^{-1})-(1-L^{-1})-(1-{L'}^{-1})
=-(1-L^{-1})(1-{L'}^{-1})\in\Filt^2(X).
\]
For the moment, admit that
\[
\det(\Filt^2(X))=\{1\}\subset\Pic(X),
\]
so that $\alpha$ is well defined. Then
\[
\alpha\circ c_1=\id_{\Pic(X)},
\]
because, by \hyperlink{sga6.e10.rem.5-2-2}{5.2.2} b),
\[
\det(1-L^{-1})=\det(L^{-1})^{-1}=L
\]
for every $L\in\Pic(X)$. Next we show that
\[
c_1\circ\alpha=\id_{\Gr^1(X)}
\]
when every section $n$ of $\mathbb Z_X$ is bounded. Since $\Filt^1(X)$ is
the abelian group generated by the elements $\cl^\circ(\mathcal E)-n$, where
$\mathcal E$ is a locally free $\mathcal O_X$-module of rank $n$, it is
enough, in view of \hyperlink{sga6.e10.rem.5-2-2}{5.2.2} b), to show that
\sgaStableTarget{sga6.e10.eq.5-3-3}\begin{equation}\tag{5.3.3}
\det(\cl^\circ(\mathcal E))-\cl^\circ(\mathcal E)+n-1\in\Filt^2(X).
\end{equation}
Since $n$ is bounded, the final object $e$ of $X$ decomposes as a finite
disjoint sum $\coprod_{i=1}^p e_i$ of subobjects such that $n$ is constant on
each subtopos $X/e_i$. Since the functor
\[
Y\longmapsto K^\circ(X/Y)\qquad(Y\in\Ob(X))
\]
takes finite disjoint sums to finite direct products, one may suppose that
$n$ is constant on $X$, say $n\in\mathbb Z$. Put
$E=\cl^\circ(\mathcal E)$. Then (cf. \hyperlink{sga6.e05.expose}{Expos\'e V}):
\[
\det(E)=\lambda^n(E)=\gamma^n((E-n)+1)
=\sum_{0\leq i\leq n}\gamma^i(E-n).
\]
This proves \hyperlink{sga6.e10.eq.5-3-3}{(5.3.3)}, since its left-hand side is
\[
\sum_{2\leq i\leq n}\gamma^i(E-n).
\]

It remains to prove the following lemma.

\sgaStableTarget{sga6.e10.lem.5-3-4}\textbf{Lemma 5.3.4.}
For every ringed topos $X$ and every $x\in\Filt^2(X)$,
\[
\det(x)=1.
\]

\emph{Proof.}
Let $\varepsilon:K^\circ(X)\to\Gamma(\mathbb Z_X)$ be the augmentation
homomorphism. Formula \hyperlink{sga6.e10.eq.6-2-8}{(6.2.8)}, together with \hyperlink{sga6.e10.para.6-5-3}{6.5.3} below, gives by linearity
\sgaStableTarget{sga6.e10.eq.5-3-5}\begin{equation}\tag{5.3.5}
\det(xy)=\det(x)^{\varepsilon(y)}\det(y)^{\varepsilon(x)}
\qquad(x,y\in K^\circ(X)).
\end{equation}

It is therefore enough to show that, for every $k>1$,
\[
\det\bigl(\gamma^k(E-n)\bigr)=1,
\]
where, as above, $E=\cl^\circ(\mathcal E)$ and
$n=\varepsilon(E)$. For every $x\in K^\circ(X)$,
\[
\gamma^k(x)=\lambda^k(x+k-1),
\]
and hence
\[
\begin{aligned}
\det\bigl(\gamma^k(E-n)\bigr)
&=\det\bigl(\lambda^k(E-n+k-1)\bigr)\\
&=\det\left(
\sum_{p+q=k}\binom{-n+k-1}{p}\lambda^q(E)
\right).
\end{aligned}
\]
But \hyperlink{sga6.e10.eq.6-2-10}{(6.2.10)} and \hyperlink{sga6.e10.para.6-5-3}{6.5.3} give
\[
\det\bigl(\lambda^q(E)\bigr)
=\det(E)^{\binom{n-1}{q-1}}
\qquad(q>0).
\]
It follows that
\[
\det\bigl(\gamma^k(E-n)\bigr)=\det(E)^r,
\]
where
\[
r=\sum_{\substack{p+q=k\\q\geq1}}
\binom{-n+k-1}{p}\binom{n-1}{q-1}.
\]
Now $r$ is a section of $\mathbb Z_X$, everywhere equal to the coefficient of
$t^{k-1}$ in $(1+t)^{k-2}$, and is therefore zero. This proves the lemma.

\sgaStableTarget{sga6.e10.sec.6}\section*{6. Appendix: Calculation of Determinants of Locally Free Sheaves}

\sgaStableTarget{sga6.e10.sec.6.1}\subsection*{6.1}
First let $X=(\mathcal C,\mathcal O)$ be an arbitrary ringed topos. For every
positive section $n\in\Gamma(X,\mathbb Z_X)$ of the constant sheaf and every
$\mathcal O$-module $M\in\Mod_{\mathcal O}(X)$, one can define the
$\mathcal O$-modules

\sgaStableTarget{sga6.e10.eq.6-1-1}\begin{equation}\tag{6.1.1}
M^n,\qquad \bigotimes\nolimits^n M,\qquad
\bigwedge\nolimits^n M,\quad\text{etc.}
\end{equation}
The functors $(\ )^n$, $\bigotimes^n$, and $\bigwedge^n$ carry free modules,
respectively locally free modules, to free modules, respectively locally free
modules.

More generally, let $n_1,\ldots,n_r$ be sections of $\mathbb Z_X$ and let
$T$ be an $r$-variable multifunctor, covariant for definiteness,
\[
T:\Mod_{\mathcal O}^{n_1}(X)\times\cdots\times
\Mod_{\mathcal O}^{n_r}(X)\longrightarrow\Mod_{\mathcal O}(X),
\]
where $\Mod_{\mathcal O}^n(X)$ denotes the category fibered over $X$ of
locally free $\mathcal O$-modules of rank $n$. Put
\[
M=T(\mathcal O^{n_1},\ldots,\mathcal O^{n_r})
\]
and
\[
G=\GL(n_1)\times\cdots\times\GL(n_r),
\qquad
\GL(n_i)=\Aut_{\mathcal O}(\mathcal O^{n_i}).
\]
It follows immediately that $T$ defines a representation
\sgaStableTarget{sga6.e10.eq.6-1-2}\begin{equation}\tag{6.1.2}
\alpha_T:G\longrightarrow\Aut_{\mathcal O}(M).
\end{equation}
If $M$ is itself locally free of rank $m$, then for every composite functor
$T'\circ T$, where
\[
T':\Mod_{\mathcal O}^m(X)\longrightarrow\Mod_{\mathcal O}(X),
\]
one has
\sgaStableTarget{sga6.e10.eq.6-1-3}\begin{equation}\tag{6.1.3}
\alpha_{T'\circ T}=\alpha_{T'}\circ\alpha_T.
\end{equation}

Conversely, from a representation
$\alpha=\alpha_T:G\to\Aut_{\mathcal O}(M)$ one recovers
$T(E_1,\ldots,E_r)$

for every sequence $E_1,\ldots,E_r$ of locally free $\mathcal O$-modules of
ranks $n_1,\ldots,n_r$, by gluing the free module
$T(\mathcal O^{n_1},\ldots,\mathcal O^{n_r})$ on local trivializing charts
through $\alpha$. More precisely, the sequence $E_1,\ldots,E_r$ determines a
right $G$-torsor $P=P_{E_1,\ldots,E_r}$ and canonical isomorphisms
\sgaStableTarget{sga6.e10.eq.6-1-3.occ2}\begin{equation}\tag{6.1.3}
E_i\xrightarrow{\sim}P\mathbin{\wedge^G}\mathcal O^{n_i}
\qquad(i=1,\ldots,r),
\end{equation}
where $\wedge^G$ denotes the contracted product and $G$ acts on
$\mathcal O^{n_i}$ through the $i$-th projection $G\to\GL(n_i)$. Every
representation
\[
\alpha:G\longrightarrow\Aut_{\mathcal O}(M)
\]
therefore defines a multifunctor $T=T_\alpha$ such that
\sgaStableTarget{sga6.e10.eq.6-1-4}\begin{equation}\tag{6.1.4}
T_\alpha(E_1,\ldots,E_r)
=P_{E_1,\ldots,E_r}\mathbin{\wedge^G}M.
\end{equation}
Moreover, for every composite representation
$\alpha=\alpha_1\circ\alpha_2$, one has
\sgaStableTarget{sga6.e10.eq.6-1-5}\begin{equation}\tag{6.1.5}
T_\alpha=T_{\alpha_1}\circ T_{\alpha_2}.
\end{equation}

\sgaStableTarget{sga6.e10.sec.6.2}\subsection*{6.2}
For every locally free $\mathcal O$-module $E$ of rank $n$, write
\[
\det(E)=\bigwedge\nolimits^n E.
\]
With the notation of \hyperlink{sga6.e10.sec.6.1}{6.1}, the functor $\det(T(E_1,\ldots,E_r))$ is
determined, up to canonical isomorphism, by the corresponding representation
\sgaStableTarget{sga6.e10.eq.6-2-1}\begin{equation}\tag{6.2.1}
\det\circ\alpha_T:G\longrightarrow\Gm=\GL(1).
\end{equation}

Suppose now that $X$ is the ringed topos of sheaves on the ringed site
$(\Sch/S)$ endowed with the \'etale topology. Then
$\det(T(E_1,\ldots,E_r))$, or equivalently the representation
$\det\circ\alpha_T$, can be calculated by means of the following proposition.

\sgaStableTarget{sga6.e10.prop.6-2-2}\textbf{Proposition 6.2.2.}
For every positive section $n$ of $\mathbb Z$, the \'etale sheaf of groups
$\Der(\GL(n))$ of commutators of $\GL(n)$ is represented by
\[
\SL(n)=\Ker\bigl(\det:\GL(n)\longrightarrow\Gm\bigr).
\]

\sgaStableTarget{sga6.e10.cor.6-2-3}\textbf{Corollary 6.2.3.}
Every representation
\sgaStableTarget{sga6.e10.eq.6-2-4}\begin{equation}\tag{6.2.4}
\alpha:G=\GL(n_1)\times\cdots\times\GL(n_r)\longrightarrow\Gm
\end{equation}
is of the form
\sgaStableTarget{sga6.e10.eq.6-2-5}\begin{equation}\tag{6.2.5}
\alpha(Y)(s)=\det(s_1)^{m_1}\cdots\det(s_r)^{m_r},
\end{equation}
where $Y\in\Ob(\Sch/S)$, $s=(s_1,\ldots,s_r)\in G(Y)$, and
$m_1,\ldots,m_r$ are sections of $\mathbb Z$ uniquely determined by
$\alpha$.

\sgaStableTarget{sga6.e10.cor.6-2-6}\textbf{Corollary 6.2.6.}
With the notation of \hyperlink{sga6.e10.sec.6.1}{6.1},
\sgaStableTarget{sga6.e10.eq.6-2-6}\begin{equation}\tag{6.2.6}
\det\bigl(T_\alpha(E_1,\ldots,E_r)\bigr)
=\det(E_1)^{\otimes m_1}\otimes\cdots\otimes
\det(E_r)^{\otimes m_r},
\end{equation}
where $m_1,\ldots,m_r$ are the sections of $\mathbb Z$ uniquely determined
by $\alpha$.

\sgaStableTarget{sga6.e10.rem.6-2-7}\textbf{Remark 6.2.7.}
The exponents $m_1,\ldots,m_r$ in \hyperlink{sga6.e10.eq.6-2-5}{(6.2.5)} and \hyperlink{sga6.e10.eq.6-2-6}{(6.2.6)} can be determined by
successively substituting in \hyperlink{sga6.e10.eq.6-2-5}{(6.2.5)}
\[
s_i=\lambda\,\id_{\mathcal O_Y^{n_i}}\quad(\lambda\in\Gm(Y)),
\qquad
s_j=\id_{\mathcal O_Y^{n_j}}\quad(j\ne i),
\]
for $i=1,\ldots,r$. Then
\[
\det(s_i)=\lambda^{n_i},
\qquad
\alpha(Y)(s)=\lambda^{n_i m_i}.
\]
For example, \hyperlink{sga6.e10.eq.6-2-6}{(6.2.6)} gives
\sgaStableTarget{sga6.e10.eq.6-2-8}\sgaStableTarget{sga6.e10.eq.6-2-9}\sgaStableTarget{sga6.e10.eq.6-2-10}\begin{align}
\det(E_1\otimes E_2)
&=\det(E_1)^{\otimes n_2}\otimes\det(E_2)^{\otimes n_1},
\tag{6.2.8}\\
\det(E_1\oplus E_2)
&=\det(E_1)\otimes\det(E_2),
\tag{6.2.9}\\
\det\bigl(\bigwedge\nolimits^m E_1\bigr)
&=\det(E_1)^{\otimes\binom{n_1-1}{m-1}}.
\tag{6.2.10}
\end{align}

\sgaStableTarget{sga6.e10.sec.6.3}\subsection*{6.3}
Let us first show how \hyperlink{sga6.e10.cor.6-2-3}{6.2.3} follows from 6.2.2. It is plainly enough to prove
\hyperlink{sga6.e10.eq.6-2-5}{(6.2.5)} for
\[
s_{(i)}=(1,\ldots,s_i,\ldots,1)\qquad(i=1,\ldots,r).
\]
Thus one may suppose that $r=1$ and $G=\GL(n_1)$. By \hyperlink{sga6.e10.prop.6-2-2}{6.2.2}, the
representation $\alpha$ factors through $\GL(n_1)/\SL(n_1)$; equivalently,
there is a commutative diagram
\[
\begin{tikzcd}[column sep=large,row sep=large]
\GL(n_1) \arrow[r,"\alpha"] \arrow[d,"j"'] & \Gm\\
\GL(n_1)/\SL(n_1) \arrow[ur,"\beta"'] &
\end{tikzcd}
\]
in which $j$ is the canonical epimorphism and $\beta$ is a representation
of $\GL(n_1)/\SL(n_1)$ in $\Gm$. The homomorphism
\[
\det:\GL(n_1)\longrightarrow\Gm
\]
has a section. Consequently, the monomorphism
\[
\GL(n_1)/\SL(n_1)\longrightarrow\Gm
\]
induced by $\det$ is an isomorphism. On the other hand, it is known
(SGA 3 VIII 1.3) that every representation of $\Gm$ in itself is of the
form
\[
s\longmapsto s^m,
\]
where $m$ is a section of $\mathbb Z$. The assertion follows.

\hyperlink{sga6.e10.cor.6-2-6}{Corollary 6.2.6} is simply a reformulation of 6.2.3. Notice that \hyperlink{sga6.e10.cor.6-2-3}{6.2.3} also
follows immediately from SGA 3 XXII 6.2.1, but we have wished to give an
elementary proof here by means of 6.2.2.

\sgaStableTarget{sga6.e10.sec.6.4}\subsection*{6.4. Proof of 6.2.2}
It is enough to prove that, for every $Y\in\Ob(\Sch/S)$, every $y\in Y$,
and every $s\in\SL(n)(Y)$, there is an \'etale neighborhood $Y'\to Y$ of
$y$ such that the restriction of $s$ to $Y'$ is a product of commutators
in $\GL(n)(Y')$. Since the question is local, one may suppose that $Y$ is
affine with ring $A$ and that the section $n$ is constant on $Y$. The case
$n=1$ is trivial; hence assume $n\geq2$. Let
$\mathfrak p$ be the prime ideal corresponding to $y$, and write
$s=(s_{ij})_{i,j=1,\ldots,n}$. For every $\lambda\in A$ and every
$h,k\in\{1,\ldots,n\}$ with $h\ne k$, let $c_{h,k,\lambda}$ be the matrix
$C=(c_{ij})$ defined by
\[
c_{hk}=\lambda,\qquad
c_{ij}=\delta_{ij}\quad\text{for }(i,j)\ne(h,k),
\]
where $\delta_{ij}$ is the Kronecker symbol.

Thus $c_{h,k,\lambda}=I+\lambda E_{hk}$. For every matrix $a$, left
multiplication by $c_{h,k,\lambda}$ adds $\lambda$ times row $k$ of $a$ to
row $h$, whereas right multiplication adds $\lambda$ times column $h$ of
$a$ to column $k$. In particular,
\[
\det(c_{h,k,\lambda})=1,
\]
and
\sgaStableTarget{sga6.e10.eq.6-4-1}\begin{equation}\tag{6.4.1}
c_{h,k,\lambda}c_{h,k,\mu}=c_{h,k,\lambda+\mu}
\end{equation}
for all $\lambda,\mu\in A$. Hence one obtains a representation
\[
c_{h,k}:\mathbb G_{\mathrm a}\longrightarrow\SL(n)
\qquad(h,k=1,\ldots,n,\ h\ne k).
\]
No invertibility assumption is needed in the row-reduction step. We reduce
to invertible parameters before the commutator step.

The proof has two steps. First, we show that $s$ is a product of matrices
$c_{h,k,\lambda}$ in a sufficiently small Zariski neighborhood of $y$.
Then we show that these matrices are products of commutators in $\GL(n)$ in
a sufficiently small \'etale neighborhood of $y$.

For the first step, one must transform $s$ into the identity matrix
$I=(\delta_{ij})$ in a Zariski neighborhood of $y$ by multiplying on the
left and right by suitable matrices $c_{h,k,\lambda}$. We first make
$s_{11}=1$. If $s_{11}\ne1$ and $s_{21}\in\mathfrak p$, then, since
$\det(s)=1\notin\mathfrak p$, there is an index $i$ for which
$s_{i1}\notin\mathfrak p$. Replacing $s$ by $c_{2,i,1}s$ replaces
$s_{21}$ by $s_{21}+s_{i1}$, which is not in $\mathfrak p$. Thus, after
localizing, one may suppose that $s_{21}$ is invertible.

If $s_{11}-1\in\mathfrak p$, replacing $s$ by $c_{1,2,1}s$ replaces
$s_{11}-1$ by
\[
(s_{11}+s_{21})-1\notin\mathfrak p.
\]
After a further localization, one may therefore suppose that $s_{11}-1$ is
invertible. Then
\[
\lambda=\frac{1-s_{11}}{s_{21}}
\]
is invertible, and left multiplication by $c_{1,2,\lambda}$ makes the new
entry $s_{11}$ equal to $1$.

Assume now that $s_{11}=1$. Left multiplication by
$c_{i,1,-s_{i1}}$ for $i>1$ clears the remaining entries in the first
column, and right multiplication by $c_{1,i,-s_{1i}}$ clears the remaining
entries in the first row. Repeating this procedure on the lower-right block
transforms $s$, in a Zariski neighborhood $U$ of $y$, into a matrix
$s'=(s'_{ij})$ satisfying
\[
s'_{ij}=\delta_{ij}\qquad\text{whenever }i<n\text{ or }j<n.
\]
By construction, $s'=c_1sc_2$ on $U$, where $c_1$ and $c_2$ are products
of matrices $c_{h,k,\lambda}$. Hence
\[
\det(s')=\det(s)=1,
\]
so $s'_{nn}=1$ and therefore $s'=I$. Since
$c_{h,k,\lambda}^{-1}=c_{h,k,-\lambda}$, this proves that $s$ is itself a
product of elementary matrices of this form.

After shrinking about $y$, a parameter $a\notin\mathfrak p$ is a unit. If
$a\in\mathfrak p$, then $a=1+(a-1)$, and both $1$ and $a-1$ are units near
$y$. By \hyperlink{sga6.e10.eq.6-4-1}{(6.4.1)},
\[
c_{h,k,a}=c_{h,k,1}c_{h,k,a-1}
\]
Thus it remains only to treat invertible parameters.

It remains to prove that these elementary matrices are products of
commutators in an \'etale neighborhood of $y$. For invertible $\lambda$ and
$\mu$, the two matrices

$c_{h,k,\lambda}$ and $c_{h,k,\mu}$ are conjugate. More precisely, there is
an invertible matrix $a=(a_{ij})$ such that
\[
a\,c_{h,k,\lambda}=c_{h,k,\mu}\,a.
\]
Indeed, the diagonal matrix defined by
\[
a_{hh}=\lambda^{-1},\qquad a_{kk}=\mu^{-1},\qquad
a_{ii}=1\quad(i\ne h,k)
\]
has this property; recall that $\lambda$ and $\mu$ are invertible and that
$h\ne k$. Thus, for every homomorphism
\[
\beta:\GL(n)(U)\longrightarrow B
\]
to an abelian group $B$, written additively, one has
\[
\beta(c_{h,k,\lambda})=\beta(c_{h,k,\mu}).
\]
If $\lambda+\mu$ is invertible, this equality and \hyperlink{sga6.e10.eq.6-4-1}{(6.4.1)} imply
\[
\beta(c_{h,k,\lambda})=0.
\]
It follows that $c_{h,k,\lambda}$ belongs to the commutator subgroup. If the
residue field $k(y)$ is not the prime field $\mathbb F_2=\{0,1\}$, then,
for every invertible $\lambda\in A$, there are a Zariski neighborhood
$U=\Spec(B)$ of $y$ and an invertible $\mu\in B$ such that
$\lambda+\mu$ is invertible in $B$. Hence $c_{h,k,\lambda}$ is a product
of commutators on $U$.

If $k(y)=\mathbb F_2$, then
\[
V=\Spec\bigl(A[T]/(T^2+T+1)\bigr)
\]
is an \'etale neighborhood of $y$ whose residue field is different from
$\mathbb F_2$. This completes the proof.

\sgaStableTarget{sga6.e10.sec.6.5}\subsection*{6.5. Remarks}
\sgaStableTarget{sga6.e10.para.6-5-1}\textbf{6.5.1.}
If no residue field of $S$ is isomorphic to $\mathbb F_2$, the preceding
proof shows that \hyperlink{sga6.e10.prop.6-2-2}{6.2.2} remains valid for the ringed topos of Zariski sheaves
on $S$. By contrast, already for $S=\Spec(\mathbb F_2)$ one has
\[
\SL(2,\mathbb F_2)=\GL(2,\mathbb F_2),
\]
which is different from the commutator subgroup of
$\GL(2,\mathbb F_2)$. The latter has order $3$, whereas
$\SL(2,\mathbb F_2)$ has order $6$ (cf. [1]).

\sgaStableTarget{sga6.e10.para.6-5-2}\textbf{6.5.2.}
On an arbitrary locally ringed topos $X$, it is not true in general that
every representation $\Gm\to\Gm$ is given by a

power. For example, on the topos $(\mathrm{Ens})$ ringed by the field
$\mathbb Q$, the representation $\Gm\to\Gm$ defined by
\[
r\longmapsto |r|\qquad(r\in\mathbb Q^*)
\]
is not a power. In particular, \hyperlink{sga6.e10.cor.6-2-3}{6.2.3} does not remain valid for an arbitrary
locally ringed topos $X$ and an arbitrary representation
$\alpha:G\to\Gm$.

\sgaStableTarget{sga6.e10.para.6-5-3}\textbf{6.5.3.}
On the other hand, \hyperlink{sga6.e10.cor.6-2-3}{6.2.3} remains valid, on an arbitrary ringed topos
$X=(\mathcal C,\mathcal O)$, for the representations associated with
formulas \hyperlink{sga6.e10.eq.6-2-8}{(6.2.8)}, \hyperlink{sga6.e10.eq.6-2-9}{(6.2.9)}, and \hyperlink{sga6.e10.eq.6-2-10}{(6.2.10)}, because these representations are
\emph{algebraic} in the following sense.

Let $e$ be the final object of $X$, put
$S=\Spec(\mathcal O(e))$, and consider the functor
\[
u:\mathcal C\longrightarrow(\Sch/S)
\]
which sends $Y\in\Ob(\mathcal C)$ to the affine $S$-scheme
$\Spec(\mathcal O(Y))$. For every positive integer $n$, the presheaf
$\GL(n)_X$ on $\mathcal C$ is obtained, by composition with $u$, from the
corresponding presheaf $\GL(n)_S$ on $(\Sch/S)$. A representation on $X$
\sgaStableTarget{sga6.e10.eq.6-5-4}\begin{equation}\tag{6.5.4}
\alpha:\GL(n)\longrightarrow\Gm
\end{equation}
is said to be \emph{algebraic} if it is induced via $u$ by a representation
\sgaStableTarget{sga6.e10.eq.6-5-5}\begin{equation}\tag{6.5.5}
\alpha_S:\GL(n)_S\longrightarrow\mathbb G_{\mathrm m,S}
\end{equation}
of group presheaves on $(\Sch/S)$; that is, if
\[
\alpha(s)=\alpha_S(s)
\]
for every $Y\in\Ob(\mathcal C)$ and every
$s\in\GL(n)(Y)=\GL(n)_S(u(Y))$. The representation $\alpha$ in \hyperlink{sga6.e10.eq.6-5-4}{(6.5.4)} is
of course completely determined by $\alpha_S$. It follows immediately that
\hyperlink{sga6.e10.cor.6-2-3}{6.2.3} remains valid in an arbitrary ringed topos whenever $\alpha$ is
algebraic. In particular, formulas \hyperlink{sga6.e10.eq.6-2-8}{(6.2.8)}, \hyperlink{sga6.e10.eq.6-2-9}{(6.2.9)}, and \hyperlink{sga6.e10.eq.6-2-10}{(6.2.10)} remain
valid on every ringed topos.

\sgaStableTarget{sga6.e10.sec.7}\section*{7. Appendix: Specialization in Intersection Theory}
\begin{center}
by A. Grothendieck
\end{center}

\sgaStableTarget{sga6.e10.lem.7-1}\textbf{Lemma 7.1.}
Let $i:Y\to X$ be a regular closed immersion (\hyperlink{sga6.e07.def.1-4}{VII 1.4}) of bounded
codimension. Then $i$ is a perfect morphism (\hyperlink{sga6.e07.prop.1-9}{VII 1.9}), so the inverse-image
homomorphism (\hyperlink{sga6.e04.para.2-12}{IV 2.12})
\[
i^*:K_\bullet(X)\longrightarrow K_\bullet(Y)
\]
is defined. Let $N$ be the conormal sheaf of $i$. It is a locally free
sheaf of finite type on $Y$ (\hyperlink{sga6.e07.def.1-6}{VII 1.6}), and hence one can form
\[
\lambda_{-1}(N)\in K^\circ(Y).
\]
Consider also the direct-image homomorphism (\hyperlink{sga6.e04.para.2-11}{IV 2.11})
\[
i_*:K_\bullet(Y)\longrightarrow K_\bullet(X).
\]
Then
\[
i^*i_*(x)=x\,\lambda_{-1}(N)\qquad(x\in K_\bullet(Y)),
\]
where the product on the right is taken with respect to the
$K^\circ(Y)$-module structure on $K_\bullet(Y)$ defined in \hyperlink{sga6.e04.para.2-10}{IV 2.10}.

The proof is, essentially without change, the same as that of the
corresponding formula for $K^\circ(X)$ and $K^\circ(Y)$ established in
\hyperlink{sga6.e07.para.2-6}{VII 2.7}.

\sgaStableTarget{sga6.e10.cor.7-2}\textbf{Corollary 7.2.}
Let $f:X\to S$ be flat, where $S$ is a trait, that is, the spectrum of a
discrete valuation ring $A$. Let $s$ be the closed point of $S$, let $t$ be
its generic point, and let $i:X_s\to X$ be the canonical inclusion. Then
$i$ has finite Tor-dimension, so (\hyperlink{sga6.e04.para.2-12}{IV 2.12}) the homomorphism
\[
i^*:K_\bullet(X)\longrightarrow K_\bullet(X_s)
\]
is defined.

On the other hand, (\hyperlink{sga6.e04.para.2-11}{IV 2.11}) gives a homomorphism
\[
i_*:K_\bullet(X_s)\longrightarrow K_\bullet(X).
\]
The composite
\[
i^*i_*:K_\bullet(X_s)\longrightarrow K_\bullet(X)
\longrightarrow K_\bullet(X_s)
\]
is zero.

Indeed, $i:X_s\to X$ is a regular immersion of codimension $1$. Its
conormal sheaf $N$ is free of rank $1$, because $X_s$ is defined by the
global equation $f^*(g)=0$, where $g$ is a uniformizer of $A$. Thus
$\lambda_{-1}(N)=0$, and the assertion follows from 7.1.

For the rest of this appendix, suppose that $X$ is noetherian and flat over
$S$. Apply the exact homotopy sequence (\hyperlink{sga6.e09.prop.1-1}{IX 1.1}) to $X$, with the open
subscheme $X_t$ and complementary closed subscheme $X_s$:
\[
K_\bullet(X_s)\xrightarrow{i_*}K_\bullet(X)
\xrightarrow{j^*}K_\bullet(X_t)\longrightarrow0,
\]
where $j:X_t\to X$ is the open immersion. \hyperlink{sga6.e10.cor.7-2}{Corollary 7.2} gives:

\sgaStableTarget{sga6.e10.prop.7-3}\textbf{Proposition 7.3.}
With the notation of \hyperlink{sga6.e10.cor.7-2}{7.2}, and assuming $X$ noetherian, there is a unique
group homomorphism
\sgaStableTarget{sga6.e10.eq.7-3-1}\begin{equation}\tag{7.3.1}
\sigma_\bullet:K_\bullet(X_t)\longrightarrow K_\bullet(X_s),
\end{equation}
called the \emph{specialization homomorphism}, for which the diagram
\sgaStableTarget{sga6.e10.eq.7-3-2}\begin{equation}\tag{7.3.2}
\begin{tikzcd}[column sep=large,row sep=large]
K_\bullet(X) \arrow[r,"j^*"] \arrow[d,"i^*"'] &
K_\bullet(X_t) \arrow[dl,"\sigma_\bullet"]\\
K_\bullet(X_s) &
\end{tikzcd}
\end{equation}
commutes.

Regard the two sides of \hyperlink{sga6.e10.eq.7-3-1}{(7.3.1)} as $K^\circ(X)$-modules via the ring
homomorphisms
\[
K^\circ(X)\xrightarrow{j^*}K^\circ(X_t),
\qquad
K^\circ(X)\xrightarrow{i^*}K^\circ(X_s)
\qquad\text{(\hyperlink{sga6.e04.para.2-7}{IV 2.7} b)}
\]
and via the $K^\circ(X_t)$- and $K^\circ(X_s)$-module structures on
$K_\bullet(X_t)$ and $K_\bullet(X_s)$, respectively (\hyperlink{sga6.e04.para.2-10}{IV 2.10}). Then
$\sigma_\bullet$ is a homomorphism of $K^\circ(X)$-modules.

This last assertion follows because
$i^*:K_\bullet(X)\to K_\bullet(X_s)$ is $K^\circ(X)$-linear, as is the
epimorphism $j^*:K_\bullet(X)\to K_\bullet(X_t)$, by \hyperlink{sga6.e04.eq.2-12-2}{IV 2.12.2}.

We now review other properties of the specialization homomorphism.

\sgaStableTarget{sga6.e10.sec.7.4}\subsection*{7.4. Compatibility with filtrations}
Endow the groups $K_\bullet$ with the filtration $\Filt_\bullet K_\bullet$
by the dimension of supports \hyperlink{sga6.e10.def.1-1-1}{(1.1.1)}. If $f$ is of finite type, and more
generally if $X$ is catenary, then the specialization homomorphism is
compatible with these filtrations and consequently induces a homomorphism,
again called specialization,
\sgaStableTarget{sga6.e10.eq.7-4-1}\begin{equation}\tag{7.4.1}
\sigma_\bullet:\Gr_\bullet(X_t)\longrightarrow\Gr_\bullet(X_s).
\end{equation}

To prove this, by \hyperlink{sga6.e10.cor.1-1-4}{1.1.4} one is reduced to showing that, if $Y_t$ is an
integral closed subscheme of dimension $d$ of $X_t$, then
$\sigma_\bullet(\cl_\bullet(\mathcal O_{Y_t}))$ has filtration at most
$d$. Let $Y$ be the schematic closure of $Y_t$ in $X$. Then $Y$ is flat
over $S$ (EGA IV 2.8.1), and one readily obtains
\[
i^*(\cl_\bullet(\mathcal O_Y))
=\cl_\bullet(\mathcal O_{Y_s}).
\]
It follows that
\[
\sigma_\bullet(\cl_\bullet(\mathcal O_{Y_t}))
=\cl_\bullet(\mathcal O_{Y_s}).
\]
Since $X$ is catenary, and hence $Y$ is catenary, one checks that
\[
\dim Y_s\leq\dim Y_t.
\]
When $f$ is of finite type, this follows immediately, after reducing to the
case in which $X$ is integral, from EGA IV 7.1.13. This proves the assertion.

\sgaStableTarget{sga6.e10.sec.7.5}\subsection*{7.5. ``Conservation of number''}
Suppose now that $f$ is proper. Then $X_t$ and $X_s$ are proper over
$k(t)$ and $k(s)$, respectively, and one has the homomorphisms defined in
\hyperlink{sga6.e10.sec.2.3}{2.3},
\[
\chi_t:K_\bullet(X_t)\longrightarrow\mathbb Z,
\qquad
\chi_s:K_\bullet(X_s)\longrightarrow\mathbb Z.
\]
For every $x\in K_\bullet(X_t)$ one has
\sgaStableTarget{sga6.e10.eq.7-5-1}\begin{equation}\tag{7.5.1}
\chi_s\bigl(\sigma_\bullet(x)\bigr)=\chi_t(x).
\end{equation}
This is the \emph{conservation-of-number theorem}; equivalently, the diagram
\sgaStableTarget{sga6.e10.eq.7-5-2}\begin{equation}\tag{7.5.2}
\begin{tikzcd}[column sep=large,row sep=large]
K_\bullet(X_t) \arrow[r,"\sigma_\bullet"] \arrow[dr,"\chi_t"'] &
K_\bullet(X_s) \arrow[d,"\chi_s"]\\
& \mathbb Z
\end{tikzcd}
\end{equation}
commutes.

To prove \hyperlink{sga6.e10.eq.7-5-1}{(7.5.1)}, one is again reduced to the case
$x=\cl_\bullet(\mathcal O_{Y_t})$, where $Y_t$ is an integral closed
subscheme of $X_t$. With the notation of the proof of \hyperlink{sga6.e10.sec.7.4}{7.4}, it remains to
prove
\[
\chi(Y_s)=\chi(Y_t),
\]
which is well known (EGA III 7.9.4, or Expos\'e III 5.7).

\sgaStableTarget{sga6.e10.sec.7.6}\subsection*{7.6}
It is immediate that
\sgaStableTarget{sga6.e10.eq.7-6-1}\begin{equation}\tag{7.6.1}
\sigma_\bullet\bigl(\cl_\bullet(\mathcal O_{X_t})\bigr)
=\cl_\bullet(\mathcal O_{X_s}).
\end{equation}
Since $\sigma_\bullet$ is $K^\circ(X)$-linear, the following diagram
\sgaStableTarget{sga6.e10.eq.7-6-2}\begin{equation}\tag{7.6.2}
\begin{tikzcd}[column sep=large,row sep=large]
& K^\circ(X) \arrow[dl,"j^*"'] \arrow[dr,"i^*"] &\\
K^\circ(X_t) \arrow[d] && K^\circ(X_s) \arrow[d]\\
K_\bullet(X_t) \arrow[rr,"\sigma_\bullet"'] && K_\bullet(X_s)
\end{tikzcd}
\end{equation}
commutes.

The vertical arrows are the canonical homomorphisms
$K^\circ\to K_\bullet$ (\hyperlink{sga6.e04.eq.2-5-1}{IV 2.5.1}).

\sgaStableTarget{sga6.e10.sec.7.7}\subsection*{7.7. Specialization for $K^\bullet$ and $\Gr^\bullet$}
When $X_t$ and $X_s$ are regular (for example, when $f$ is smooth), the two
preceding vertical arrows are isomorphisms (\hyperlink{sga6.e04.para.2-5}{IV 2.5}). The specialization
homomorphism can then be interpreted as a homomorphism
\sgaStableTarget{sga6.e10.eq.7-7-1}\begin{equation}\tag{7.7.1}
\sigma^\bullet:K^\bullet(X_t)\longrightarrow K^\bullet(X_s),
\end{equation}
which is equivalently the unique homomorphism making the diagram
\sgaStableTarget{sga6.e10.eq.7-7-2}\begin{equation}\tag{7.7.2}
\begin{tikzcd}[column sep=large,row sep=large]
K^\bullet(X) \arrow[r,"j^*"] \arrow[d,"i^*"'] &
K^\bullet(X_t) \arrow[dl,"\sigma^\bullet"]\\
K^\bullet(X_s) &
\end{tikzcd}
\end{equation}
commute.

In this diagram $i^*$ and $j^*$ are ring homomorphisms (\hyperlink{sga6.e04.para.2-7}{IV 2.7} b), and
$j^*$ is surjective. Hence \hyperlink{sga6.e10.eq.7-7-1}{(7.7.1)} is a ring homomorphism. For the same
reason, if $X$ is separated, it is even a homomorphism of $\lambda$-rings.
Indeed, a regular separated noetherian scheme is divisorial (\hyperlink{sga6.e02.cor.2-2-7-1}{II 2.2.7.1}),
so its $K^\bullet$ is canonically identified with its ``naive''
$K^\bullet$ (\hyperlink{sga6.e04.para.2-9}{IV 2.9(ii)}), which carries a natural $\lambda$-ring structure
(\hyperlink{sga6.e06.thm.3-2}{VI 3.2}). Moreover, the regularity of $X_t$ and $X_s$ implies that $X$ is
regular (EGA $0_{\mathrm{IV}}$ 17.1.8), and
\[
i^*:K^\bullet(X)\longrightarrow K^\bullet(X_s),
\qquad
j^*:K^\bullet(X)\longrightarrow K^\bullet(X_t)
\]
are $\lambda$-homomorphisms. Therefore the same is true of \hyperlink{sga6.e10.eq.7-7-1}{(7.7.1)}, which
is obtained from $i^*$ by passage to the quotient.

Since $X$ is regular, hence normal, its irreducible components are its
connected components, and they dominate $S$. Consequently the homomorphism
\[
H^0(X,\mathbb Z)\longrightarrow H^0(X_t,\mathbb Z)
\]
induced by $X_t\to X$ is bijective. The homomorphism induced by
$X_s\to X$ can therefore also be regarded as a specialization homomorphism
\sgaStableTarget{sga6.e10.eq.7-7-3}\begin{equation}\tag{7.7.3}
H^0(X_t,\mathbb Z)\longrightarrow H^0(X_s,\mathbb Z).
\end{equation}
It follows immediately from the definitions that \hyperlink{sga6.e10.eq.7-7-1}{(7.7.1)} is compatible,
via \hyperlink{sga6.e10.eq.7-7-3}{(7.7.3)}, with the augmented-algebra structures of its source and target
over $H^0(X_t,\mathbb Z)$ and $H^0(X_s,\mathbb Z)$, respectively, defined
in \hyperlink{sga6.e04.para.2-8}{IV 2.8}.

Still assuming that $X$ is separated, \hyperlink{sga6.e10.eq.7-7-1}{(7.7.1)} induces a homomorphism of
graded rings
\sgaStableTarget{sga6.e10.eq.7-7-4}\begin{equation}\tag{7.7.4}
\sigma^\bullet:\Gr^\bullet(X_t)\longrightarrow\Gr^\bullet(X_s).
\end{equation}
Using the Chern-class construction of V 6, one obtains the commutative
diagram
\[
\begin{tikzcd}[column sep=large,row sep=large]
K^\bullet(X_t) \arrow[r,"c_{X_t}"] \arrow[d] &
\Ch\,\Gr^\bullet(X_t) \arrow[d]\\
K^\bullet(X_s) \arrow[r,"c_{X_s}"'] &
\Ch\,\Gr^\bullet(X_s).
\end{tikzcd}
\]

\sgaStableTarget{sga6.e10.sec.7.8}\subsection*{7.8. Specialization for $\Pic$}
Retain the hypotheses of \hyperlink{sga6.e10.sec.7.7}{7.7}: $X$ is separated, and $X_t$ and $X_s$ are
regular. Using the canonical isomorphism \hyperlink{sga6.e10.thm.5-3-2}{(5.3.2)}
\[
\Gr^1\simeq\Pic,
\]
the homomorphism \hyperlink{sga6.e10.eq.7-7-4}{(7.7.4)} gives a canonical homomorphism
\sgaStableTarget{sga6.e10.eq.7-8-1}\begin{equation}\tag{7.8.1}
\Pic(X_t)\longrightarrow\Pic(X_s).
\end{equation}

An equivalent description is as follows. Since $X$ is regular, its local
rings are factorial by the Auslander--Buchsbaum theorem (EGA IV 21.11.1).
Thus every invertible sheaf $L_t$ on $X_t$ extends to an invertible sheaf
$L$ on $X$ (EGA IV 21.6.11). The image of
$\cl(L_t)\in\Pic(X_t)$ under (7.8.1) is
$\cl(L_s)\in\Pic(X_s)$, where $L_s=L|_{X_s}$.

In fact, assume only that the local rings of $X$ are factorial and that
$X_s$ is locally integral, without assuming regularity or separatedness.
Then $\cl(L_s)$ is independent of the chosen extension $L$ of $L_t$ to an
invertible sheaf on $X$, and this gives \hyperlink{sga6.e10.eq.7-8-1}{(7.8.1)} directly. To prove the
independence, it is enough to show that
\[
L_t\simeq\mathcal O_{X_t}\quad\Longrightarrow\quad
L_s\simeq\mathcal O_{X_s}.
\]
The hypothesis means that $L\simeq\mathcal O_X(D)$ for a divisor $D$ on
$X$ whose support is contained in $X_s$ (EGA IV 21.2.11). One reduces at
once to the case in which $X_s$ is integral. Then $D=nX_s$, and
$D=\operatorname{div}(g^n)$, where $g$ is a uniformizer of $A$. Hence
$\mathcal O_X(D)\simeq\mathcal O_X$, and a fortiori
$L_s\simeq\mathcal O_{X_s}$.

\sgaStableTarget{sga6.e10.sec.7.9}\subsection*{7.9. Specialization and numerical equivalence}
Suppose that $f$ is proper and that $X_s$ and $X_t$ are regular. Consider
the two specialization homomorphisms \hyperlink{sga6.e10.eq.7-4-1}{(7.4.1)} and \hyperlink{sga6.e10.eq.7-7-4}{(7.7.4)}:
\sgaStableTarget{sga6.e10.eq.7-9-1}\begin{equation}\tag{7.9.1}
\begin{gathered}
\sigma_\bullet:\Gr_\bullet(X_t)\longrightarrow\Gr_\bullet(X_s),\\
\sigma^\bullet:\Gr^\bullet(X_t)\longrightarrow\Gr^\bullet(X_s).
\end{gathered}
\end{equation}
They are induced by the homomorphism
\sgaStableTarget{sga6.e10.eq.7-9-2}\begin{equation}\tag{7.9.2}
\sigma^\bullet:K^\bullet(X_t)\longrightarrow K^\bullet(X_s).
\end{equation}

Since \hyperlink{sga6.e10.eq.7-9-2}{(7.9.2)} is compatible with the ring structures, formula \hyperlink{sga6.e10.eq.7-5-1}{(7.5.1)}
gives
\sgaStableTarget{sga6.e10.eq.7-9-3}\begin{equation}\tag{7.9.3}
\left\langle\sigma^\bullet(y_t),\sigma_\bullet(x_t)\right\rangle
=\left\langle y_t,x_t\right\rangle,
\qquad
x_t\in\Gr_\bullet(X_t),\quad y_t\in\Gr^\bullet(X_t),
\end{equation}
where the pairings between $\Gr^\bullet$ and $\Gr_\bullet$ are those of
4.1. It follows formally that, if $x\in\Gr_\bullet(X_t)$ (respectively,
$x\in\Gr^\bullet(X_t)$) and $\sigma_\bullet(x)$ (respectively,
$\sigma^\bullet(x)$) is numerically equivalent to zero \hyperlink{sga6.e10.def.4-5-1}{(4.5.1)}, then $x$
is numerically equivalent to zero. The converse holds if \hyperlink{sga6.e10.eq.7-9-2}{(7.9.2)}, or,
equivalently, either of the homomorphisms in \hyperlink{sga6.e10.eq.7-9-1}{(7.9.1)} obtained from it by passage to
the associated graded groups, is surjective modulo torsion. This condition
occurs only in exceptional cases.

\sgaStableTarget{sga6.e10.para.7-9-4}\textbf{7.9.4.}
It is plausible, at least when $f$ is proper and smooth, that the converse
just considered always holds: if $x$ is numerically equivalent to zero,
then so is $\sigma_\bullet(x)$ (respectively, $\sigma^\bullet(x)$). This
would imply that the homomorphisms \hyperlink{sga6.e10.eq.7-9-1}{(7.9.1)} induce monomorphisms on the
corresponding groups modulo numerical equivalence,
\[
\Gr_\bullet^{\mathrm{num}}
\qquad\text{and}\qquad
\Gr_{\mathrm{num}}^\bullet.
\]
At least the preceding result shows that
$\Gr_\bullet^{\mathrm{num}}(X_t)$ is isomorphic to a quotient of a subgroup
of $\Gr_\bullet^{\mathrm{num}}(X_s)$; in particular, its rank is at most
that of $\Gr_\bullet^{\mathrm{num}}(X_s)$. The same holds for
$\Gr_{\mathrm{num}}^\bullet$, and also for $K_{\mathrm{num}}^\bullet$.

\sgaStableTarget{sga6.e10.sec.7.10}\subsection*{7.10. Compatibility with direct and inverse images}
Let
\[
g:X\longrightarrow Y
\]
be a morphism of noetherian schemes flat over $S$. Denote by
$\sigma_{X,\bullet}$ and $\sigma_{Y,\bullet}$ the specialization
homomorphisms for $X$ and $Y$, and by
$i_X,j_X$ (respectively, $i_Y,j_Y$)

the inclusions of the special and generic fibers into the ambient schemes.
Write
\[
g_t:X_t\longrightarrow Y_t,
\qquad
g_s:X_s\longrightarrow Y_s
\]
for the induced morphisms. They give the Cartesian squares
\sgaStableTarget{sga6.e10.eq.7-10-0}\begin{equation}\tag{7.10.0}
\begin{tikzcd}[column sep=large,row sep=large]
X_t \arrow[r,"g_t"] \arrow[d,"j_X"'] & Y_t \arrow[d,"j_Y"'] &
X_s \arrow[r,"g_s"] \arrow[d,"i_X"'] & Y_s \arrow[d,"i_Y"']\\
X \arrow[r,"g"] & Y & X \arrow[r,"g"] & Y.
\end{tikzcd}
\end{equation}

First suppose that $g$ is proper. Then the diagram
\sgaStableTarget{sga6.e10.eq.7-10-1}\begin{equation}\tag{7.10.1}
\begin{tikzcd}[column sep=large,row sep=large]
K_\bullet(X_t) \arrow[r,"g_{t*}"] \arrow[d,"\sigma_{X,\bullet}"'] &
K_\bullet(Y_t) \arrow[d,"\sigma_{Y,\bullet}"]\\
K_\bullet(X_s) \arrow[r,"g_{s*}"'] & K_\bullet(Y_s)
\end{tikzcd}
\end{equation}
commutes. This generalizes \hyperlink{sga6.e10.sec.7.5}{7.5}: in the present notation take $Y=S$
(equivalently, take the scheme denoted by $X$ in \hyperlink{sga6.e10.eq.7-3-2}{(7.3.2)} to be $S$); the
specialization homomorphism for $S$ is the identity on $\mathbb Z$. If $X$
and $Y$ satisfy the conditions of \hyperlink{sga6.e10.sec.7.7}{7.7}, that is, if their fibers are
regular, then $g_t$ and $g_s$ automatically have finite Tor-dimension, and
\hyperlink{sga6.e10.eq.7-10-1}{(7.10.1)} may be read as a commutative square whose four terms are
$K^\bullet$-groups. We leave this reformulation to the reader.

Using the definitions of $\sigma_{X,\bullet}$ and $\sigma_{Y,\bullet}$,
the commutativity of \hyperlink{sga6.e10.eq.7-10-1}{(7.10.1)} reduces to that of the two front faces of the
prism
\sgaStableTarget{sga6.e10.eq.unnamed}\[
\begin{tikzcd}[column sep=large,row sep=large]
K_\bullet(X_t) \arrow[rr,"g_{t*}"] && K_\bullet(Y_t)\\
& K_\bullet(X) \arrow[ul,"j_X^*" description]
  \arrow[dl,"i_X^*" description] \arrow[rr,"g_*"] &&
K_\bullet(Y) \arrow[ul,"j_Y^*" description]
  \arrow[dl,"i_Y^*" description]\\
K_\bullet(X_s) \arrow[rr,"g_{s*}"'] && K_\bullet(Y_s).
\end{tikzcd}
\tag{*}
\]

These two faces are the base-change squares obtained from \hyperlink{sga6.e10.eq.7-10-0}{(7.10.0)}, so IV 3
applies. It remains only to verify that they correspond to the pairs
$(X,Y_t)$ and $(X,Y_s)$ of schemes over $Y$, and that these pairs are
Tor-independent over $Y$. This is clear for the first pair because
$Y_t\to Y$ is flat. It is also immediate for the second, because
$Y_s\to Y$ is a regular immersion of codimension $1$ and remains so after
base change by $X\to Y$.

Now drop the properness assumption and instead suppose that $g$ has finite
Tor-dimension. Then $g_t$ and $g_s$ also have finite Tor-dimension. This is
immediate for $g_t$; for $g_s$ it follows from the Tor-independence of
$X$ and $Y_s$ over $Y$ just noted. Hence one can form the square
\sgaStableTarget{sga6.e10.eq.7-10-2}\begin{equation}\tag{7.10.2}
\begin{tikzcd}[column sep=large,row sep=large]
K_\bullet(Y_t) \arrow[r,"g_t^*"] \arrow[d,"\sigma_{Y,\bullet}"'] &
K_\bullet(X_t) \arrow[d,"\sigma_{X,\bullet}"]\\
K_\bullet(Y_s) \arrow[r,"g_s^*"'] & K_\bullet(X_s),
\end{tikzcd}
\end{equation}
whose horizontal arrows are those defined in \hyperlink{sga6.e04.para.2-12}{IV 2.12}. If $X$ and $Y$
satisfy the conditions of \hyperlink{sga6.e10.sec.7.7}{7.7}, that is, if their fibers are regular, then
$Y$ is regular and $g$ automatically has finite Tor-dimension; the square
may then be interpreted as a commutative square of $K^\bullet$-rings.

To prove the commutativity of \hyperlink{sga6.e10.eq.7-10-2}{(7.10.2)}, it is again enough to prove that of
the two front faces of the prism analogous to $(*)$, with the direct-image
homomorphisms $g_{t*},g_*,g_{s*}$ replaced by the inverse-image
homomorphisms $g_t^*,g^*,g_s^*$. Their commutativity follows

immediately from that of the corresponding squares in \hyperlink{sga6.e10.eq.7-10-0}{(7.10.0)}, using the
transitivity of inverse images (\hyperlink{sga6.e04.para.2-12}{IV 2.12}).

\sgaStableTarget{sga6.e10.rem.7-10-3}\textbf{Remark 7.10.3.}
When $g$ is proper, the commutativity of \hyperlink{sga6.e10.eq.7-10-1}{(7.10.1)} immediately implies that
of the following diagram, obtained by passage to the associated graded
groups:
\sgaStableTarget{sga6.e10.eq.7-10-1bis}\begin{equation}\tag{7.10.1 bis}
\begin{tikzcd}[column sep=large,row sep=large]
\Gr_\bullet(X_t) \arrow[r,"g_{t*}"]
  \arrow[d,"\sigma_{X,\bullet}"'] &
\Gr_\bullet(Y_t) \arrow[d,"\sigma_{Y,\bullet}"]\\
\Gr_\bullet(X_s) \arrow[r,"g_{s*}"'] & \Gr_\bullet(Y_s).
\end{tikzcd}
\end{equation}
Likewise, if $X$ and $Y$ are separated and have regular fibers, the
commutativity of \hyperlink{sga6.e10.eq.7-10-2}{(7.10.2)} implies that of
\sgaStableTarget{sga6.e10.eq.7-10-2bis}\begin{equation}\tag{7.10.2 bis}
\begin{tikzcd}[column sep=large,row sep=large]
\Gr^\bullet(Y_t) \arrow[r,"g_t^*"] \arrow[d,"\sigma_Y^\bullet"'] &
\Gr^\bullet(X_t) \arrow[d,"\sigma_X^\bullet"]\\
\Gr^\bullet(Y_s) \arrow[r,"g_s^*"'] & \Gr^\bullet(X_s).
\end{tikzcd}
\end{equation}
If $X$ and $Y$ are moreover quasi-projective and smooth over $S$, then
$g_t^*$ and $g_s^*$ are compatible with the filtrations by codimension of
supports. Thus one obtains a commutative diagram of the same form as
(\hyperlink{sga6.e10.eq.7-10-2}{7.10.2} bis), with $\Gr^\bullet$ replaced by
$\Gr_{\mathrm{top}}^\bullet$ (\hyperlink{sga6.e00.sec.2.1}{Expos\'e 0, 2.1)}.

Finally, if $X$ and $Y$ are quasi-projective and separated over $S$ and have
regular fibers, then $g_t$ and $g_s$ are automatically complete-intersection
morphisms (\hyperlink{sga6.e08.def.1-1}{VIII 1.1}). Using \hyperlink{sga6.e08.prop.3-2}{VIII 3.2}, one obtains a commutative diagram
analogous to (\hyperlink{sga6.e10.eq.7-10-1}{7.10.1} bis), with $\Gr_\bullet$ replaced by
$\Gr^\bullet(-)_{\mathbb Q}$.

\sgaStableTarget{sga6.e10.para.7-11}\textbf{7.11.}
For a scheme $T$ over an algebraically closed field $k$, one can define the
subgroup of $K^\bullet(T)$ consisting of elements algebraically equivalent to
zero as the subgroup of elements of the form
\sgaStableTarget{sga6.e10.eq.7-11-1}\begin{equation}\tag{7.11.1}
E(x_1)-E(x_0),
\end{equation}
where $x_0$ and $x_1$ are two $k$-rational points of a proper, regular,
connected algebraic curve $C$ over $k$, where
$E\in K^\bullet(T\times_k C)$, and where, for every $k$-rational point $x$ of
$C$, one sets
\[
E(x)=i_x^*(E)\in K^\bullet(T),
\qquad\text{where}\qquad
i_x:T\longrightarrow T\times_k C
\]
is the morphism $t\mapsto(t,x)$ defined by $x$. That the elements of the form
\hyperlink{sga6.e10.eq.7-11-1}{(7.11.1)} constitute a subgroup of $K^\bullet(T)$ follows by a standard
argument, using a product $C\times_k C'$, the two points
$(x_1,x'_1)$ and $(x_0,x'_0)$ of that product, and the fact that one can
always find an irreducible algebraic curve through any two points of an
irreducible quasi-projective algebraic scheme over $k$ (here
$C\times_k C'$). (A reader
unwilling to admit this fact may define algebraic equivalence by the group
generated by the elements of the form \hyperlink{sga6.e10.eq.7-11-1}{(7.11.1)}.)

The elements algebraically equivalent to zero form an ideal of
$K^\bullet(T)$, as one sees at once. This ideal is moreover stable under the
operations $\lambda^i$ whenever they are defined, that is, when one works
with the ``naive'' $K^\bullet$; it is also contained in the augmentation
ideal of $K^\bullet(T)$, as is immediate. One can therefore define a quotient
ring
\sgaStableTarget{sga6.e10.eq.7-11-2}\begin{equation}\tag{7.11.2}
K^\bullet_{\mathrm{alg}}(T),
\end{equation}
which is even a $\lambda$-ring when one works with naive $K^\bullet$-groups.

In the same way one can define a subgroup of elements algebraically
equivalent to zero in $K_\bullet(T)$, taking the element $E$ above in
$K_\bullet(T\times_k C)$ and observing that

$i_x:T\to T\times_k C$ is a perfect morphism: it is obtained, by flat base
change along $T\times_k C\to C$, from the morphism $j_x:x\to C$, which is
perfect because $C$ is regular. One thereby obtains a quotient group
\sgaStableTarget{sga6.e10.eq.7-11-3}\begin{equation}\tag{7.11.3}
K_\bullet^{\mathrm{alg}}(T)
\end{equation}
of $K_\bullet(T)$. This group is naturally a module over
$K^\bullet_{\mathrm{alg}}(T)$, since the definitions show at once that the
product of an element of $K^\bullet(T)$ with an element of $K_\bullet(T)$ is
algebraically equivalent to zero if either factor is, using (\hyperlink{sga6.e04.eq.2-12-2}{IV 2.12.2}) for
$i_x$.

When one works with naive $K^\bullet$, define the ideal in $\Gr^\bullet(T)$
of elements algebraically equivalent to zero to be the ideal generated by the
$c^i(x)$, where $i\geq1$ and $x\in K^\bullet(T)$ is algebraically equivalent
to zero. The quotient graded ring
\sgaStableTarget{sga6.e10.eq.7-11-4}\begin{equation}\tag{7.11.4}
\Gr^\bullet_{\mathrm{alg}}(T)
\end{equation}
is then precisely the Chern ring associated with the $\lambda$-ring
$K^\bullet_{\mathrm{alg}}(T)$.

Finally, when $T$ is locally noetherian, define the graded subgroup of
$\Gr_\bullet(T)$ consisting of elements algebraically equivalent to zero by
requiring its component of dimension $i$ to be the image in $\Gr_i(T)$ of
the subgroup of $\operatorname{Fil}_iK_\bullet(T)$ formed by the elements algebraically
equivalent to zero in $K_\bullet(T)$. It follows at once that the quotient
graded group
\sgaStableTarget{sga6.e10.eq.7-11-5}\begin{equation}\tag{7.11.5}
\Gr_\bullet^{\mathrm{alg}}(T)
\end{equation}
is a module over the ring $\Gr^\bullet_{\mathrm{alg}}(T)$, by passage to the
quotient from the module structure of $\Gr_\bullet(T)$ over
$\Gr^\bullet(T)$ \hyperlink{sga6.e10.sec.3.1}{(3.1)}.

\sgaStableTarget{sga6.e10.para.7-11-6}\textbf{7.11.6.}
The preceding definitions would still make sense without assuming $k$
algebraically closed, provided that $C$ above is taken geometrically
connected over $k$. It is not clear, however, that these definitions would
be very useful. In particular, it does not appear that, for the group
$\Pic(T)$, one would

recover the classical ``geometric'' notion of algebraic equivalence
considered in 2.4.1. We therefore prefer to say that an element of
$K^\bullet(T)$ is algebraically equivalent to zero if its image in
$K^\bullet(T\otimes_k\bar k)$ is so, where $\bar k$ is an algebraic closure
of $k$, and to proceed in the same way when defining algebraic equivalence
in $K_\bullet(T)$, $\Gr^\bullet(T)$, and $\Gr_\bullet(T)$.

\sgaStableTarget{sga6.e10.para.7-12}\textbf{7.12.}
With these definitions in place, I claim that the \emph{specialization homomorphism}
\hyperlink{sga6.e10.eq.7-3-1}{(7.3.1)}
\emph{carries elements algebraically equivalent to zero to elements
algebraically equivalent to zero}, and consequently induces, by passage to
the quotient, a homomorphism
\sgaStableTarget{sga6.e10.eq.7-12-1}\begin{equation}\tag{7.12.1}
K_\bullet^{\mathrm{alg}}(X_t)
\longrightarrow K_\bullet^{\mathrm{alg}}(X_s).
\end{equation}

To see this, one reduces immediately to the case in which $A$ is complete
and $k(s)$ is algebraically closed, using the straightforward compatibility
made explicit in 7.15. It is enough to show that $\sigma_\bullet$ carries an
element of the form \hyperlink{sga6.e10.eq.7-11-1}{(7.11.1)} in $K_\bullet(X_t)$ to an element algebraically
equivalent to zero in $K_\bullet(X_s)$. Here $T=X_t$, and accordingly we
write $C_t,E_t$ in place of $C,E$ above. Since $C_t$ is projective over
$k(t)$, EGA IV 2.8.5 gives a projective flat scheme $C$ over $S$ whose
generic fiber is $C_t$. Using Abhyankar's resolution of singularities for
arithmetic surfaces [6], one sees that $C$ can even be chosen regular.
(N.B. Flatness is preserved under resolution, by EGA IV 14.3.8.) Since $C$ is
proper over $S$, the $k(t)$-rational points $x_{0t}$ and $x_{1t}$ of $C_t$
extend to sections $x_0,x_1$ of $C$ over $S$. Since $C$ is regular, it is
well known that $C$ is smooth over $S$ at the points of $x_0(S)$ and
$x_1(S)$ (EGA IV 17.12.1(c)$\Rightarrow$(a), and 19.1.1). Consequently,

the immersions $x_i:S\to C$ $(i=0,1)$ are regular and remain regular
immersions $x_i:s\to C_s$ after base change by $s\to S$. Note that, by the
exact homotopy sequence (\hyperlink{sga6.e09.prop.1-1}{IX 1.1}), the element
\[
E_t\in K_\bullet(X_t\times_{k(t)}C_t)
\]
comes from an element $E\in K_\bullet(X\times_SC)$. Thus the element
$E_t(x_{1t})-E_t(x_{0t})$ of $K_\bullet(X_t)$ comes from the element
$E(x_1)-E(x_0)$ of $K_\bullet(X)$, and its image in $K_\bullet(X_s)$ under
\hyperlink{sga6.e10.eq.7-3-1}{(7.3.1)} is therefore
\sgaStableTarget{sga6.e10.eq.unnamed.occ2}\[
i^*\bigl(E(x_1)-E(x_0)\bigr).
\tag{*}
\]

For a section $x$ of $C$ over $S$, of course, we set
\[
E(x)=i_x^*(E),
\qquad\text{where}\qquad
i_x:X\longrightarrow X\times_SC
\]
is defined as above. It may be viewed as obtained from $x$ by the flat base
change $X\times_SC\to C$; it is therefore a regular immersion and, a
fortiori, a perfect morphism, so that the expression $E(x)$ is well defined.
For the same reason $E_s(x_s)$ is defined. By transitivity of inverse images
(\hyperlink{sga6.e04.para.2-12}{IV 2.12}) in the following commutative diagram of perfect morphisms,
\[
\begin{tikzcd}[column sep=large,row sep=large]
X \arrow[r,"i_x"] & X\times_SC\\
X_s \arrow[u,"i"] \arrow[r,"i_{x_s}"] &
X_s\times_{k(s)}C_s \arrow[u]
\end{tikzcd}
\;,
\]
the expression $(*)$ is equal to
\sgaStableTarget{sga6.e10.eq.unnamed.occ3}\[
E_s(x_{1s})-E_s(x_{0s}).
\tag{**}
\]
Because $k(s)$ is algebraically closed, the reduced irreducible components
of $C_s$ are geometrically integral over $k(s)$, and their intersection
points are rational over $k(s)$. It follows readily that $(**)$ is
algebraically equivalent to zero in $K_\bullet(X_s)$, using the connectedness

of $C_s$, which follows from the geometric connectedness of $C_t$.

From \hyperlink{sga6.e10.eq.7-12-1}{(7.12.1)}, using \hyperlink{sga6.e10.sec.7.4}{7.4}, one obtains a homomorphism
\sgaStableTarget{sga6.e10.eq.7-12-2}\begin{equation}\tag{7.12.2}
\Gr_\bullet^{\mathrm{alg}}(X_t)
\longrightarrow \Gr_\bullet^{\mathrm{alg}}(X_s).
\end{equation}
If $X_s$ and $X_t$ are moreover regular, then \hyperlink{sga6.e10.eq.7-12-1}{(7.12.1)} can also be
interpreted as a ring homomorphism
\sgaStableTarget{sga6.e10.eq.7-12-3}\begin{equation}\tag{7.12.3}
K^\bullet_{\mathrm{alg}}(X_t)
\longrightarrow K^\bullet_{\mathrm{alg}}(X_s),
\end{equation}
which, if $X$ is also separated, is a homomorphism of $\lambda$-rings and
therefore induces a homomorphism
\sgaStableTarget{sga6.e10.eq.7-12-4}\begin{equation}\tag{7.12.4}
\Gr^\bullet_{\mathrm{alg}}(X_t)
\longrightarrow \Gr^\bullet_{\mathrm{alg}}(X_s).
\end{equation}

\sgaStableTarget{sga6.e10.para.7-12-5}\textbf{7.12.5.}
Let $T$ be a smooth proper algebraic scheme over a field $k$. It is easy to
verify that an element $x\in K^\bullet(T)$ for which there is an integer
$n\geq1$ such that $nx$ is algebraically equivalent to zero is numerically
equivalent to zero. It is plausible that the converse is always true, at
least when $T$ is projective. In considering this question, one may plainly
assume that $k$ is algebraically closed and replace $K^\bullet(T)$ by
$\Gr^\bullet(T)$ or $\Gr_\bullet(T)$ without changing the conjecture. One
obtains an equivalent conjecture by working with the Chow ring $A(T)$
(cf. \hyperlink{sga6.e14.para.4-1}{XIV 4.1}), namely the more or less classical conjecture that, if $Z$ is an
algebraic cycle on $T$ numerically equivalent to zero, then there is an
integer $n\geq1$ such that $nZ$ is algebraically equivalent to zero. If this
conjecture holds for $X_t$, the compatibility of \hyperlink{sga6.e10.eq.7-3-1}{(7.3.1)} with algebraic
equivalence would imply that this specialization homomorphism, when
$f:X\to S$ is projective and smooth, is also compatible with numerical
equivalence, as conjectured in 7.9.

\sgaStableTarget{sga6.e10.sec.7.13}\subsection*{7.13. Relations with $\ell$-adic cohomology}
Recall (SGA 5 VII) that, for every locally free module of finite type $E$ on
a scheme $T$, and every prime number $\ell$ distinct from the residue
characteristics of $T$, one defines $\ell$-adic Chern classes (for $i\geq0$)
\[
c^i(E)\in H^{2i}(T,\mathbb Z_\ell(i)),
\]
satisfying the additivity formula
\[
c(E)=c(E')c(E'')
\]
when $E$ is an extension of a locally free module $E'$ by a locally free
module $E''$. It follows from this additivity that the total Chern class
$c(E)=\sum_i c^i(E)$ arises from a homomorphism
\sgaStableTarget{sga6.e10.eq.7-13-0}\begin{equation}\tag{7.13.0}
c:K^\bullet(T)\longrightarrow
\widehat G\bigl(H^{2*}(T,\mathbb Z_\ell(*))\bigr)
=1+\prod_{i\geq1}H^{2i}(T,\mathbb Z_\ell(i)),
\end{equation}
where the target is a group under the product explained in V 1, and where
$K^\bullet(T)$ is the ``naive'' $K^\bullet$. One obtains canonical maps
\sgaStableTarget{sga6.e10.eq.7-13-1}\begin{equation}\tag{7.13.1}
c^i:K^\bullet(T)\longrightarrow H^{2i}(T,\mathbb Z_\ell(i)).
\end{equation}

Now let $X$ be a proper flat scheme over $S$, with regular fibers, and suppose
that $S$ is henselian. Assume that $\ell$ is different from the residue
characteristic of $S$. One can then form the diagram
\sgaStableTarget{sga6.e10.eq.7-13-2}\begin{equation}\tag{7.13.2}
\begin{tikzcd}[column sep=large,row sep=large]
K^\bullet(X_t) \arrow[r,"c^i"] \arrow[d,"\sigma"'] &
H^{2i}(X_t,\mathbb Z_\ell(i))\\
K^\bullet(X_s) \arrow[r,"c^i"'] &
H^{2i}(X_s,\mathbb Z_\ell(i)) \arrow[u,"\sigma^{-1}"']
\end{tikzcd}
\;,
\end{equation}
where the horizontal arrows are the Chern-class maps for $X_t$ and $X_s$,
respectively, and where $\sigma^{-1}$ is the composite homomorphism
\[
H(X_s)\xleftarrow[\sim]{i^*}H(X)\xrightarrow{j^*}H(X_t)\,;
\]
the arrows displayed here are restriction homomorphisms, and one uses the
fact that $i^*$ is an isomorphism by the proper base change theorem
(SGA 4 XII 5.5). I claim that the square (7.13.2) is commutative. The proof is
essentially the same as the proof of the commutativity of \hyperlink{sga6.e10.eq.7-10-2}{(7.10.2)}: it is
enough to use the compatibility of the formation of Chern classes with
inverse images for the morphisms $j:X_t\to X$ and $i:X_s\to X$.

\sgaStableTarget{sga6.e10.para.7-13-3}\textbf{7.13.3.}
The square \hyperlink{sga6.e10.eq.7-13-2}{(7.13.2)} is obtained by passage to the projective limit from the
analogous squares in which cohomology with coefficients in the $\ell$-adic
sheaf $\mathbb Z_\ell(i)$ is replaced by cohomology with coefficients in the
torsion sheaves $\mu_n^{\otimes i}$, where $n$ runs through the powers of
$\ell$. Thus the commutativity of \hyperlink{sga6.e10.eq.7-13-2}{(7.13.2)} merely expresses the
commutativity of the squares involving Chern classes modulo $n$, for every
power $n$ of $\ell$.

Choose an algebraic closure $\bar K$ of $K=k(t)$, and hence an algebraic
closure $\bar k$ of $k=k(s)$, namely the residue field of the integral closure
of $A$ in $\bar K$. For every finite subextension $K_\alpha$ of $\bar K$, one
can consider the analogous commutative diagram, for fixed $n$, in the
situation obtained from $f:X\to S$ by the base change $S_\alpha\to S$, where
$S_\alpha$ is the normalization of $S$ in $K_\alpha$. By Krull--Akizuki and
the henselian hypothesis on $S$, the scheme $S_\alpha$ is indeed a trait.
Here one must moreover assume that the fibers of $X$ over $S$ are
geometrically regular, that is, that $X$ is smooth over $S$, so that this
hypothesis is preserved by base change. As $\alpha$ varies, one obtains an
inductive system of commutative squares. Passing to the inductive limit and
using SGA 4 VII 5.8 and IV 3.2, one obtains a commutative square (cf. \hyperlink{sga6.e10.sec.7.16}{7.16}):
\[
\begin{tikzcd}[column sep=large,row sep=large]
K^\bullet(X_{\bar t}) \arrow[r,"c^i"] \arrow[d,"\sigma"'] &
H^{2i}(X_{\bar t},\mu_n^{\otimes i}) \arrow[d,"\sigma'"]\\
K^\bullet(X_{\bar s}) \arrow[r,"c^i"'] &
H^{2i}(X_{\bar s},\mu_n^{\otimes i})
\end{tikzcd}
\;,
\]
where $X_{\bar t}$ and $X_{\bar s}$ are the geometric fibers of $X$.
Moreover, the specialization homomorphism $\sigma^{-1}$ for cohomology is
now known to be an isomorphism (SGA 4 XVI 2.2). Hence the commutativity of the
preceding square can also be expressed as the commutativity of the square
obtained by reversing that arrow. Passing to the projective limit over $n$ in
the resulting commutative squares gives a commutative square
\sgaStableTarget{sga6.e10.eq.7-13-4}\begin{equation}\tag{7.13.4}
\begin{tikzcd}[column sep=large,row sep=large]
K^\bullet(X_{\bar t}) \arrow[r] \arrow[d,"\sigma"'] &
H^{2i}(X_{\bar t},\mathbb Z_\ell(i)) \arrow[d,"\sim"']\\
K^\bullet(X_{\bar s}) \arrow[r] &
H^{2i}(X_{\bar s},\mathbb Z_\ell(i))
\end{tikzcd}
\;.
\end{equation}

\sgaStableTarget{sga6.e10.para.7-13-5}\textbf{7.13.5.}
Another way to express this commutativity is to say that, if
$x_t\in K^\bullet(X_t)$ and $x_s$ is its image in $K^\bullet(X_s)$, then
$c^i(x_{\bar t})$ and $c^i(x_{\bar s})$ are the values at the geometric
points $\bar t$ and $\bar s$, respectively, of a section of
$R^{2i}f_*(\mathbb Z_\ell(i))$. This section is easily interpreted: take an
element $x\in K^\bullet(X)$ inducing $x_t$. Then
$c^i(x)\in H^{2i}(X,\mathbb Z_\ell(i))$ defines the required section of
$R^{2i}f_*(\mathbb Z_\ell(i))$.

\sgaStableTarget{sga6.e10.para.7-13-6}\textbf{7.13.6.}
Let $T$ be a scheme. By passage to the associated graded groups, the
homomorphism \hyperlink{sga6.e10.eq.7-13-0}{(7.13.0)} induces homomorphisms
\[
\psi_i:\Gr^i(T)\longrightarrow H^{2i}(T,\mathbb Z_\ell(i)),
\]
the image under $\psi_i$ of an element $x\in\Filt^iK^\bullet(T)$ being, by
definition, $c^i(x)$. (The compatibility of \hyperlink{sga6.e10.eq.7-13-0}{(7.13.0)} with the
$\lambda$-filtration on
$K^\bullet(T)$ and the degree filtration on the target follows readily from
\hyperlink{sga6.e05.cor.6-6-4}{V 6.6.4} and from the fact that the $c^i$ define a
Chern theory on $K^\bullet(T)$ in the sense of \hyperlink{sga6.e05.def.6-13}{V 6.13}, when $K^\bullet(T)$ is
regarded as a $\lambda$-algebra augmented over the binomial ring
$H^0(T,\mathbb Z)$.) One thereby obtains maps
\sgaStableTarget{sga6.e10.eq.7-13-7}\begin{equation}\tag{7.13.7}
\operatorname{cl}:\Gr^i(T)\longrightarrow H^{2i}(T,\mathbb Q_\ell(i))
\stackrel{\mathrm{dfn}}{=}
H^{2i}(T,\mathbb Z_\ell(i))\otimes_{\mathbb Z_\ell}\mathbb Q_\ell,
\end{equation}
by the formula
\sgaStableTarget{sga6.e10.eq.7-13-8}\begin{equation}\tag{7.13.8}
\operatorname{cl}(x^{(i)})=
\frac{1}{(-1)^{i-1}(i-1)!}\,\psi_i(x^{(i)}),
\end{equation}
compare \hyperlink{sga6.e14.para.5-1}{XIV 5.1} and 5.2. The commutativity of \hyperlink{sga6.e10.eq.7-13-2}{(7.13.2)} and \hyperlink{sga6.e10.eq.7-13-4}{(7.13.4)}
immediately implies that of the analogous diagrams
\sgaStableTarget{sga6.e10.eq.7-13-9}\begin{equation}\tag{7.13.9}
\begin{tikzcd}[column sep=large,row sep=large]
\Gr^i(X_t)_{\mathbb Q} \arrow[r,"\operatorname{cl}"]
  \arrow[d,"\sigma_{\mathbb Q}"'] &
H^{2i}(X_t,\mathbb Q_\ell(i))\\
\Gr^i(X_s)_{\mathbb Q} \arrow[r,"\operatorname{cl}"'] &
H^{2i}(X_s,\mathbb Q_\ell(i)) \arrow[u,"\sigma'"]
\end{tikzcd}
\;,
\end{equation}
and
\sgaStableTarget{sga6.e10.eq.7-13-10}\begin{equation}\tag{7.13.10}
\begin{tikzcd}[column sep=large,row sep=large]
\Gr^i(X_{\bar t})_{\mathbb Q} \arrow[r,"\operatorname{cl}"] \arrow[d] &
H^{2i}(X_{\bar t},\mathbb Q_\ell(i)) \arrow[d,"\sim"]\\
\Gr^i(X_{\bar s})_{\mathbb Q} \arrow[r,"\operatorname{cl}"'] &
H^{2i}(X_{\bar s},\mathbb Q_\ell(i))
\end{tikzcd}
\;.
\end{equation}

\sgaStableTarget{sga6.e10.rem.7-13-11}\textbf{Remark 7.13.11.}
The commutativity of the last two diagrams can be expressed by saying that
specialization of algebraic cycles is compatible with formation of the
associated $\ell$-adic cohomology classes. This statement is not very
satisfactory in the form just given, because we had to replace the
coefficients $\mathbb Z_\ell(i)$ by $\mathbb Q_\ell(i)$, that is, to neglect
torsion. To obtain a statement that does not neglect torsion, it seems that
one will be forced to work with the Chow ring, since it is on that ring that
the homomorphism $\operatorname{cl}$ is naturally defined with values in
$H^{2*}(-,\mathbb Z_\ell(*))$; compare \hyperlink{sga6.e14.para.4-1}{XIV 4.1}. Unfortunately, no theory of
specialization for the Chow ring has yet been developed (cf. \hyperlink{sga6.e10.sec.7.14}{7.14}) that would
allow one to write commutative squares analogous to \hyperlink{sga6.e10.eq.7-13-9}{(7.13.9)} and \hyperlink{sga6.e10.eq.7-13-10}{(7.13.10)} in
terms of Chow rings and cohomology with coefficients in
$\mathbb Z_\ell(*)$. All one can say when $f:X\to S$ is smooth and proper is
that there is a commutative square
\sgaStableTarget{sga6.e10.eq.7-13-12}\begin{equation}\tag{7.13.12}
\begin{tikzcd}[column sep=large,row sep=large]
Z^i(X_t) \arrow[r,"\operatorname{cl}"] \arrow[d,"\sigma"'] &
H^{2i}(X_t,\mathbb Z_\ell(i))\\
Z^i(X_s) \arrow[r,"\operatorname{cl}"'] &
H^{2i}(X_s,\mathbb Z_\ell(i)) \arrow[u]
\end{tikzcd}
\;,
\end{equation}
analogous to \hyperlink{sga6.e10.eq.7-13-9}{(7.13.9)}, together with a corresponding square analogous to
\hyperlink{sga6.e10.eq.7-13-10}{(7.13.10)}, which we leave to the reader to write down. Here $Z^i$ denotes the
group of algebraic cycles of codimension $i$, modulo no relation at all; the
horizontal arrows are the associated-cohomology-class homomorphisms defined
in SGA 5 IV; and $\sigma$ is the specialization homomorphism for algebraic
cycles, associating with a cycle $Z_t$ the cycle $Z_s$, where $Z$ is the
schematic closure of $Z_t$ in $X$ (cf. \hyperlink{sga6.e10.sec.7.14}{7.14}). This commutativity follows from
the following fact, proved in SGA 5 IV. If $Z$ is a closed subscheme of $X$
of codimension at least $i$ on every fiber, one can associate with it a
cohomology class
\[
\operatorname{cl}(Z)\in H^{2i}(X,\mathbb Z_\ell(i))
\]
whose formation commutes with every base change and whose restriction to
$X_t$ (respectively, $X_s$) is therefore
$\operatorname{cl}(Z_t)$ (respectively, $\operatorname{cl}(Z_s)$). Of course,
once it is known that the homomorphism $\sigma$ in \hyperlink{sga6.e10.eq.7-13-12}{(7.13.12)} carries cycles
rationally equivalent to zero to cycles rationally equivalent to zero, with
$X$ projective over $S$,
one can replace the $Z^i$ by the $A^i$ in this diagram and obtain a diagram
of the desired kind in terms of the Chow rings of $X_t$ and $X_s$.

\sgaStableTarget{sga6.e10.rem.7-13-13}\textbf{Remark 7.13.13.}
Let $T$ be a smooth projective scheme over an algebraically closed field. It
follows immediately from the formalism developed in SGA 5 IV that, if $Z$ is
a cycle of codimension $i$ on $T$ whose image in
$H^{2i}(T,\mathbb Q_\ell(i))$ is zero, then $Z$ is numerically equivalent to
zero. One of the standard conjectures on algebraic cycles asserts that the
converse is also true. By \hyperlink{sga6.e07.prop.4-11}{VII 4.11}, one obtains an equivalent conjecture by
replacing the group of algebraic cycles with $\Gr^\bullet(T)_{\mathbb Q}$,
using the \hyperlink{sga6.e10.eq.7-13-8}{definition (7.13.8)}, or with $K^\bullet(T)$, by stating that an
element of $K^\bullet(T)$ is numerically equivalent to zero if and only if
its augmentation and its Chern classes in the groups
$H^{2i}(T,\mathbb Q_\ell(i))$ vanish.

It follows immediately from the commutativity of \hyperlink{sga6.e10.eq.7-13-10}{(7.13.10)} that, if
$X_{\bar t}$ satisfies the preceding conjecture, then the specialization
homomorphism
\[
\Gr^\bullet(X_{\bar t})\longrightarrow\Gr^\bullet(X_{\bar s})
\]
(and, a fortiori, $\Gr^\bullet(X_t)\to\Gr^\bullet(X_s)$) is compatible with
numerical equivalence, as conjectured in 7.9. The conjecture used here
(numerical equivalence implies $\mathbb Q_\ell$-cohomological equivalence) is
weaker than the conjecture considered in 7.12.5.

\sgaStableTarget{sga6.e10.sec.7.14}\subsection*{7.14. Specialization homomorphism for groups of cycles and cycle classes}
If $T$ is a noetherian scheme, denote by $Z_i(T)$ (respectively, $Z^i(T)$)
the group of algebraic cycles on $T$ that are purely of dimension $i$
(respectively, purely of codimension $i$). When $f:X\to S$ is of finite type,
one defines
a specialization homomorphism
\sgaStableTarget{sga6.e10.eq.7-14-1}\begin{equation}\tag{7.14.1}
\sigma_i:Z_i(X_t)\longrightarrow Z_i(X_s)
\end{equation}
as follows. For every prime cycle $Z_t$ of $X_t$, purely of dimension $i$ and
identified with an integral subscheme of $X_t$, define $\sigma_i(Z_t)$ to be
the cycle associated with the closed subscheme $Z_s$ of $X_s$, where $Z$ is
the schematic closure of $Z_t$ in $X$. This $Z$ is a closed subscheme of $X$
flat over $S$ (EGA IV 2.8), with generic fiber equal to the original $Z_t$,
and it satisfies $\dim Z_s=\dim Z_t$ (EGA IV 7.1.13). When $X_s$ is locally
equidimensional, that is, when all irreducible components passing through a
given point have the same dimension, the specialization homomorphism just
defined immediately induces homomorphisms
\sgaStableTarget{sga6.e10.eq.7-14-1bis}\begin{equation}\tag{7.14.1 bis}
\sigma^i:Z^i(X_t)\longrightarrow Z^i(X_s).
\end{equation}

Now suppose that $f$ is quasi-projective, so that rational equivalence of
cycles is available on $X_t$ and $X_s$ (cf. [5], where the theory is
developed over an algebraically closed field). It seems very plausible that
\hyperlink{sga6.e10.eq.7-14-1}{(7.14.1)} carries cycles rationally equivalent to zero to cycles rationally
equivalent to zero and therefore induces a corresponding homomorphism on the
``Chow'' groups of cycle classes,
\sgaStableTarget{sga6.e10.eq.7-14-2}\begin{equation}\tag{7.14.2}
\sigma_i:A_i(X_t)\longrightarrow A_i(X_s),
\end{equation}
or, when $X_s$ is locally equidimensional,
\sgaStableTarget{sga6.e10.eq.7-14-2bis}\begin{equation}\tag{7.14.2 bis}
\sigma^i:A^i(X_t)\longrightarrow A^i(X_s).
\end{equation}

One might try to prove this by a method paralleling that of \hyperlink{sga6.e10.para.7-12}{7.12}, but here
resolution of singularities and reduction
to an algebraically closed residue field would be unnecessary, since
$C=\mathbb P^1_S$ is available as a
``parameter curve'' along which to deform cycles. A difficulty nevertheless
arises. With the notation of \hyperlink{sga6.e10.para.7-12}{7.12}, $E$ must now denote a cycle on
$\mathbb P^1_S\times_SX$ for which
$E_t(x_{1t})-E_t(x_{0t})$ is ``defined,'' but it does not follow that the
cycle $E_s(x_{1s})-E_s(x_{0s})$ is also defined. As was to be expected, one
is thus reduced to formulating a statement of the kind of Chow's moving
lemma in order to define \hyperlink{sga6.e10.eq.7-14-2}{(7.14.2)}.
Observe, however, that the existence of \hyperlink{sga6.e10.eq.7-4-1}{(7.4.1)} shows that \hyperlink{sga6.e10.eq.7-14-1}{(7.14.1)} is
compatible with the equivalence relation on cycles defined by the surjective
homomorphism
\[
Z_\bullet(T)\longrightarrow\Gr_\bullet(T),
\]
which is only slightly coarser than linear equivalence. Moreover, if $T$ is
quasi-projective and smooth over a field, this relation differs from linear
equivalence only by torsion (\hyperlink{sga6.e14.para.4-1}{XIV 4.1}, \hyperlink{sga6.e14.para.4-2}{4.2}). In that case one therefore obtains
a homomorphism analogous to \hyperlink{sga6.e10.eq.7-14-2}{(7.14.2)}, or better to (\hyperlink{sga6.e10.eq.7-14-2}{7.14.2} bis), with both
sides tensored by $\mathbb Q$.

The most natural way to define \hyperlink{sga6.e10.eq.7-14-2}{(7.14.2)} directly would seem to be to proceed
as in 7.3. This would require a theory of Chow groups $A_i$ (or of the Chow
ring in the regular case) for schemes more general than schemes
quasi-projective over a field, broad enough to include $X$, together with a
contravariant law for these groups making it possible to define the
restriction homomorphism to $X_s$,
\[
A_i(X)\longrightarrow A_i(X_s).
\]
Compare XIV 8. Once such a sufficiently general theory of Chow groups and
rings is available, it should be possible to transcribe into Chow theory all
the results stated above for the specialization homomorphism in $K$-theory.

\sgaStableTarget{sga6.e10.sec.7.15}\subsection*{7.15. Varying $S$}
Let $S'$ be a second trait, and let
\[
S'\longrightarrow S
\]
be a morphism of traits, that is, a dominant local morphism, carrying the
closed point to the closed point and the generic point to the generic point.
Denote the closed point (respectively, the generic point) of $S'$ by $s'$
(respectively, $t'$), and set $X'=X\times_SS'$. Suppose that $X'$ is
noetherian, as is the case, for example, when $X$ is of finite type over $S$.
There is a commutative diagram
\[
\begin{tikzcd}[column sep=large,row sep=large]
X & X_t \arrow[l,hook,"j_X"'] & & \\
& & X' \arrow[ull] & X'_{t'} \arrow[l,hook,"j_{X'}"'] \arrow[ull] \\
X_s \arrow[uu,hook,"i_X"] & & & \\
& & X'_{s'} \arrow[uu,hook,"i_{X'}"'] \arrow[ull] &
\end{tikzcd}
\;,
\]
in which the oblique arrows are flat morphisms. One can therefore take the
inverse-image homomorphisms on $K_\bullet$ for every arrow of the diagram
(\hyperlink{sga6.e04.para.2-12}{IV 2.12}). Applying \hyperlink{sga6.e10.prop.7-3}{Definition 7.3}, one finds that the diagram
\sgaStableTarget{sga6.e10.eq.7-15-1}\begin{equation}\tag{7.15.1}
\begin{tikzcd}[column sep=large,row sep=large]
K_\bullet(X_t) \arrow[r,"\sigma_X"] \arrow[d] &
K_\bullet(X_s) \arrow[d]\\
K_\bullet(X'_{t'}) \arrow[r,"\sigma_{X'}"'] &
K_\bullet(X'_{s'})
\end{tikzcd}
\end{equation}
is commutative, where the vertical arrows are the inverse-image
homomorphisms. One immediately obtains an analogous commutative diagram with
$\Gr_\bullet$ in place of $K_\bullet$. When $X_s$ and $X_t$ are regular, one
may likewise write $K^\bullet$ and $\Gr^\bullet$ in place of $K_\bullet$ and
$\Gr_\bullet$.


\sgaStableTarget{sga6.e10.sec.7.16}\subsection*{7.16. Passage to geometric fibers}
Let $\bar K$ be an algebraic closure of $K=k(t)$, expressed as the inductive
limit of its finite subextensions $K_i$, and let
$\bar A=\varinjlim_i\bar A_i$ be the normalization of $A$ in $\bar K$, the
inductive limit of its normalizations $\bar A_i$ in the $K_i$. Choose a
maximal ideal $\bar{\mathfrak m}$ of $\bar A$; it induces a maximal ideal
$\mathfrak m_i$ on each $\bar A_i$. Set
$B_i=(\bar A_i)_{\mathfrak m_i}$. By Krull--Akizuki, $B_i$ is again a
discrete valuation ring. Denote its spectrum by $S_i$ and its closed and
generic points by $s_i$ and $t_i$, respectively. The residue field
$\bar k=\bar A/\bar{\mathfrak m}$ is an algebraic closure of $k=k(s)$ and is
the inductive limit of the residue fields
$k_i=\bar A_i/\mathfrak m_i=k(s_i)$ of the $B_i$. Set
$\bar t=\Spec(\bar K)$ and $\bar s=\Spec(\bar k)$.

Assume now that $X$ is of finite type over $S$. Then $X_{\bar t}$ and
$X_{\bar s}$ are of finite type over $\bar t$ and $\bar s$, respectively,
and hence are noetherian. By \hyperlink{sga6.e04.para.3-2-3}{IV 3.2.3},
\[
K_\bullet(X_{\bar t})\xleftarrow{\sim}
  \varinjlim_i K_\bullet(X_{t_i}),
\qquad
K_\bullet(X_{\bar s})\xleftarrow{\sim}
  \varinjlim_i K_\bullet(X_{s_i}).
\]
On the other hand, \hyperlink{sga6.e10.sec.7.15}{7.15} gives a homomorphism from the inductive system of the
$K_\bullet(X_{t_i})$ to that of the $K_\bullet(X_{s_i})$. Passing to the
limit gives a canonical homomorphism
\sgaStableTarget{sga6.e10.eq.7-16-1}\begin{equation}\tag{7.16.1}
K_\bullet(X_{\bar t})\longrightarrow K_\bullet(X_{\bar s}),
\end{equation}
again, as expected, called the \emph{specialization homomorphism}. By
\hyperlink{sga6.e10.sec.7.4}{7.4}, it induces a homomorphism
$\Gr_\bullet(X_{\bar t})\to\Gr_\bullet(X_{\bar s})$, using the fact (whose
verification is left to the reader) that the functor $\Gr_\bullet$, like
$K_\bullet$, commutes with the limits occurring here. Similarly, when $f$
is smooth, one may interpret \hyperlink{sga6.e10.eq.7-16-1}{(7.16.1)} as a homomorphism on $K^\bullet$; it is
then a ring homomorphism, and even a $\lambda$-ring homomorphism if $X$ is
separated. In that case, it therefore induces a homomorphism on
$\Gr^\bullet$.

\sgaStableTarget{sga6.e10.para.7-16-2}\textbf{7.16.2.}
We shall not review here, for geometric fibers, the extension of all the
properties of specialization homomorphisms proved in the preceding sections
for the ``arithmetic'' fibers; they all follow immediately from the latter
by passage to the limit. Let us mention only the following interesting
consequence of 7.9.4. When $f$ is smooth and proper,
$\Gr_\bullet^{\mathrm{num}}(X_{\bar t})$ is isomorphic to a quotient of a
subgroup of $\Gr_\bullet^{\mathrm{num}}(X_{\bar s})$; in particular, its
rank is bounded above by that of the latter group. This is therefore a
semi-continuity result for the rank of the group of cycles modulo numerical
equivalence.

\sgaStableTarget{sga6.e10.sec.7.17}\subsection*{7.17. Case of an arbitrary base}
Let
\[
f:X\longrightarrow S
\]
be a flat and quasi-compact morphism of noetherian preschemes, and let $s,t$
be two points of $S$ such that $s$ is a specialization of $t$. We wish to
relate the groups $K_\bullet$, etc., for the two fibers $X_t$ and $X_s$, as
well as for the corresponding geometric fibers. For this, note (EGA II
7.1.9) that there exists a morphism
\[
g:S'\longrightarrow S,
\]
where $S'$ is a trait and $g$ carries the closed point $s'$ to $s$ and the
generic point $t'$ to $t$; one may even assume that
$k(t)\xrightarrow{\sim}k(t')$. Suppose that $X'=X\times_SS'$ is
noetherian. This is the case if $f$ is of finite type or if $S'$ is
essentially of finite type over $S$ (EGA IV 1.3.8), as one may assume when
$S$ itself is universally Japanese (EGA $0_{\mathrm{IV}}$ 23.1.1). The
specialization homomorphism for $X'$ then defines a homomorphism
\sgaStableTarget{sga6.e10.eq.7-17-1}\begin{equation}\tag{7.17.1}
\sigma_\bullet:K_\bullet(X_t)\xrightarrow{\sim}K_\bullet(X_{t'})
  \longrightarrow K_\bullet(X_{s'}).
\end{equation}
Thus one obtains a specialization homomorphism on $K_\bullet(X_t)$, but with
values not in $K_\bullet(X_s)$, but in $K_\bullet(X_{s'})$, where $X_{s'}$
is obtained from $X_s$ by an extension $k(s)\to k(s')$ of the base fields
(which may be assumed to be of finite type if $S$ is universally Japanese).
Passing to geometric fibers as explained in \hyperlink{sga6.e10.sec.7.16}{7.16}, one likewise obtains a
homomorphism
\sgaStableTarget{sga6.e10.eq.7-17-2}\begin{equation}\tag{7.17.2}
\sigma_\bullet:K_\bullet(X_{\bar t})
  \longrightarrow K_\bullet(X_{\overline{s'}}).
\end{equation}
Here again, the ``geometric'' specialization homomorphism takes its values
not in $K_\bullet(X_{\bar s})$, but in
$K_\bullet(X_{\overline{s'}})$, where $X_{\overline{s'}}$ is obtained from
$X_{\bar s}$ by the extension of algebraically closed fields
$k(\bar s)\to k(\overline{s'})$.

\sgaStableTarget{sga6.e10.para.7-17-3}\textbf{7.17.3.}
In general, the homomorphism
$K_\bullet(X_{\bar s})\to K_\bullet(X_{\overline{s'}})$ is not an
isomorphism, because of the possible presence of ``continuous'' components
in $K_\bullet(X_{\bar s})$. It is, however, easy to verify that it is an
isomorphism modulo algebraic equivalence \hyperlink{sga6.e10.para.7-11}{(7.11)}, and also modulo numerical
equivalence \hyperlink{sga6.e10.def.4-5-1}{(4.5.1)} if $f$ is smooth and proper. Consequently, using \hyperlink{sga6.e10.para.7-12}{7.12},
one finds that \hyperlink{sga6.e10.eq.7-17-2}{(7.17.2)} induces a specialization homomorphism
\sgaStableTarget{sga6.e10.eq.7-17-3-1}\begin{equation}\tag{7.17.3.1}
K_\bullet^{\mathrm{alg}}(X_{\bar t})
  \longrightarrow K_\bullet^{\mathrm{alg}}(X_{\bar s});
\end{equation}
and, if $f$ is proper and smooth, provided that the specialization
homomorphism \hyperlink{sga6.e10.eq.7-17-2}{(7.17.2)} is compatible with numerical equivalence (cf. \hyperlink{sga6.e10.para.7-9-4}{7.9.4}),
one likewise obtains a homomorphism
\sgaStableTarget{sga6.e10.eq.7-17-3-2}\begin{equation}\tag{7.17.3.2}
K_\bullet^{\mathrm{num}}(X_{\bar t})
  \longrightarrow K_\bullet^{\mathrm{num}}(X_{\bar s}).
\end{equation}
We leave to the reader the task of developing, if needed, the corresponding
constructions for $\Gr_\bullet$, $K^\bullet$, and $\Gr^\bullet$, and the
properties of these specialization homomorphisms that follow immediately
from those developed in the preceding sections. Let us mention only
that, by \hyperlink{sga6.e10.para.7-16-2}{7.16.2}, $\Gr_\bullet^{\mathrm{num}}(X_{\bar t})$ is isomorphic to a
quotient of a subgroup of $\Gr_\bullet^{\mathrm{num}}(X_{\bar s})$ when $f$
is proper and smooth; in particular, its rank is again bounded above by the
rank of the latter group.

\sgaStableTarget{sga6.e10.para.7-17-4}\textbf{7.17.4.}
When $S$ is regular at $s$ and $t$ is a maximal point, one can construct
$S'$ by
blowing up $s$ in the localization of $S$ at $s$, and then localizing at the
generic point of the fiber of $s$ in the blow-up. One then obtains a
\emph{pure extension} $k(s')$ of $k(s)$. Now it is not difficult to show
that, if $k'$ is a pure extension of a field $k$, then for every noetherian
scheme $T$ over $k$ the canonical homomorphism
\[
K_\bullet(T)\longrightarrow K_\bullet(T\otimes_k k')
\]
is an isomorphism: represent $k'$ as the inductive limit of the coordinate
rings of nonempty affine open subsets of schemes of the form
$\Spec k[T_1,\ldots,T_n]$, and use \hyperlink{sga6.e09.prop.1-6}{IX 1.6}. It follows that \hyperlink{sga6.e10.eq.7-17-1}{(7.17.1)} may be
viewed as a canonical homomorphism
\sgaStableTarget{sga6.e10.eq.7-17-4-1}\begin{equation}\tag{7.17.4.1}
K_\bullet(X_t)\longrightarrow K_\bullet(X_s).
\end{equation}
One should note, however, that this does not yield a homomorphism
$K_\bullet(X_{\bar t})\to K_\bullet(X_{\bar s})$: after passing to a finite
extension $K_i$ of $K=k(t)$ and normalizing $S$ in $K_i$, the regularity
hypothesis on $S$ is lost.

Another way to define \hyperlink{sga6.e10.eq.7-17-4-1}{(7.17.4.1)} is to reduce to the case in which $S$ is
local with closed point $s$ (which is immediate), and to consider an
increasing sequence of regular closed subschemes $S_i$ of $S$, with generic
points $t_i$, such that
\[
\dim S_i=i,\qquad 0\leq i\leq n=\dim S.
\]
For every $i$, $0\leq i\leq n-1$, the local ring
$\mathcal O_{S_{i+1},t_i}$ is a discrete valuation ring with residue field
$k(t_i)$ and fraction field $k(t_{i+1})$. There is therefore a
specialization homomorphism
\[
K_\bullet(X_{t_{i+1}})\longrightarrow K_\bullet(X_{t_i}),
\]
and the composite of these homomorphisms for $0\leq i\leq n-1$ defines
\hyperlink{sga6.e10.eq.7-17-4-1}{(7.17.4.1)}. By an argument analogous to that of \hyperlink{sga6.e10.prop.7-3}{7.3}, one can show that this
is the same homomorphism as the one defined above by means of a blow-up. The
advantage of the second method is that one even has homomorphisms of the
form \hyperlink{sga6.e10.eq.7-16-1}{(7.16.1)}:
\[
K_\bullet(X_{\bar t_{i+1}})\longrightarrow
K_\bullet(X_{\bar t_i}),
\]
with suitable coherent choices of algebraic closures $k(\bar t_i)$ of the
$k(t_i)$. Their composite defines a homomorphism
\sgaStableTarget{sga6.e10.eq.7-17-4-2}\begin{equation}\tag{7.17.4.2}
K_\bullet(X_{\bar t})\longrightarrow K_\bullet(X_{\bar s}),
\end{equation}
which probably depends essentially on the choice of the $S_i$, that is, of
the $t_i$.

\sgaStableTarget{sga6.e10.para.7-17-5}\textbf{7.17.5.}
The construction of a specialization homomorphism \hyperlink{sga6.e10.eq.7-17-4-2}{(7.17.4.2)} can be
generalized to the case in which $S$ is no longer assumed regular at $s$
and $t$ is no longer assumed to be a maximal point, provided that resolution
of singularities is available for the integral closed subscheme $T$ of $S$
with generic point $t$, or merely for its localization at $s$. This
hypothesis holds, in particular, if $\mathcal O_{S,s}$ is excellent of
characteristic zero, by Hironaka. Indeed, one can then find a morphism
\[
S'\longrightarrow\Spec(\mathcal O_{T,s})
\]
with $S'$ regular. Taking a closed point of the fiber $S'_s$, one can apply
the preceding construction to $X'=X\times_SS'$ over $S'$, its generic point
$t'$, and its closed point $s'$. This yields a homomorphism of the form
\hyperlink{sga6.e10.eq.7-17-4-2}{(7.17.4.2)}, since $k(s')$ is a finite extension of $k(s)$ and therefore has
the same algebraic closure.

\sgaStableTarget{sga6.e10.para.7-17-6}\textbf{7.17.6.}
In all the constructions considered in \hyperlink{sga6.e10.sec.7.17}{7.17}, one has reduced the
construction of a specialization homomorphism to the
case in which the base is a trait. It follows that the properties of the
specialization homomorphism reviewed in the preceding sections for a trait
as base extend automatically to the case of a more general base. Of course,
the methods of definition considered apply equally well to $\Gr_\bullet$,
to $K^\bullet$ when $f$ is smooth, and to $\Gr^\bullet$ when $f$ is also
separated.

\sgaStableTarget{sga6.e10.para.7-17-7}\textbf{7.17.7.}
Finally, note that the definition of a homomorphism \hyperlink{sga6.e10.eq.7-17-3-1}{(7.17.3.1)} or
\hyperlink{sga6.e10.eq.7-17-3-2}{(7.17.3.2)}, or of a homomorphism \hyperlink{sga6.e10.eq.7-17-4-2}{(7.17.4.2)}, depends essentially on certain
arbitrary choices: a suitable trait over $S$ in \hyperlink{sga6.e10.para.7-17-3}{7.17.3}, a sequence $(S_i)$
in \hyperlink{sga6.e10.para.7-17-4}{7.17.4}, and, in addition, the choice of a resolution of singularities in
7.17.5. The same is true of the variants for $\Gr_\bullet$, etc. Obvious
examples can be constructed with $s$, say, an ordinary double point of an
algebraic curve and $X$ projective and smooth over $S$ (for example, an
abelian scheme of relative dimension $2$), in which the specialization
homomorphism on $\Gr^1=\Pic$, and even on
$\Gr^1_{\mathrm{alg}}=\mathrm{NS}$ (the Néron--Severi group), actually
depends on the trait $S'$ chosen over $S$, or, what amounts to the same thing
here, on the choice of a point above $s$ in the normalization of $S$.

\sgaStableTarget{sga6.e10.sec.7.18}\subsection*{7.18. Appendix to the Appendix: base-field extension for certain
functors with values in commutative groups}
We spell out here a technical point raised in 7.17.4.

\sgaStableTarget{sga6.e10.para.7-18-1}\textbf{7.18.1.}
Let $k$ be a field, let $C$ be the category of integral $k$-algebras and
flat homomorphisms of such algebras, and let $F$ be a functor on $C$ with
values in the category $(\mathrm{Ens})$ of sets. Suppose that $F$ commutes
with filtered inductive limits. Given an extension $K$ of $k$, we seek
criteria
for
\[
F(k)\longrightarrow F(K)
\]
to be injective (respectively, surjective).

\sgaStableTarget{sga6.e10.para.7-18-2}\textbf{7.18.2.}
Suppose that one of the following two conditions holds:
\begin{enumerate}[label=\alph*)]
\item The field $k$ is separably closed, and $K/k$ is a separable extension.
\item The field $k$ is infinite, and $K$ is a pure extension of $k$.
\end{enumerate}
Then the map $F(k)\to F(K)$ is injective.

Indeed, $K$ is the filtered inductive limit of its finite-type
$k$-subalgebras, with flat transition morphisms. Using this fact and the
commutation of $F$ with filtered inductive limits, one is reduced in case
a) to proving the following: if $A$ is an integral $k$-algebra of finite
type whose fraction field $K$ is separable over $k$, then
$F(k)\to F(A)$ is injective. The hypotheses on $k$ and $K$ imply that there
is a $k$-homomorphism $A\to k$; the composite
\[
F(k)\longrightarrow F(A)\longrightarrow F(k)
\]
is therefore the identity, and the assertion follows. In case b), passage
to the limit similarly reduces the question to an integral finite-type
$k$-algebra $A$ whose fraction field is a pure extension of $k$. Since $k$
is infinite, the hypothesis again implies the existence of a
$k$-homomorphism $A\to k$, and one concludes as above.

\sgaStableTarget{sga6.e10.para.7-18-3}\textbf{7.18.3.}
Suppose now that $F$ takes values in the category $(\mathrm{Ab})$ of
abelian groups, and suppose that for every finite extension $k'$ of $k$ a
norm homomorphism
\[
N_{k'/k}:F(k')\longrightarrow F(k)
\]
has been defined, such that the composite
\[
F(k)\longrightarrow F(k')\xrightarrow{N_{k'/k}}F(k)
\]
is multiplication by $[k':k]$. Under these conditions:
\begin{enumerate}[label=(\roman*)]
\item for every extension $K$ of $k$, the kernel of
  $F(k)\to F(K)$ is a torsion group;
\item for every pure extension $K$ of $k$, the homomorphism
  $F(k)\to F(K)$ is injective. (Note that $k$ is not assumed infinite.)
\end{enumerate}

For (i), it is enough to prove that every element $x$ in the kernel of
$F(k)\to F(A)$, where $A$ is an integral $k$-algebra of finite type, is a
torsion element. Using a $k$-homomorphism $A\to k'$, where $k'$ is a finite
$k$-algebra, one reduces to the case in which $A=k'$ is a finite extension
of $k$. The norm property above then gives $[k':k]x=0$, proving that $x$ is
torsion. One even has $x=0$ when the greatest common divisor of the degrees
of the finite extensions $k'$ of $k$ that are quotients of $A$ is $1$. This
is the case when the fraction field $K$ of $A$ is a pure extension of $k$,
as is easily checked (only the case in which $k$ is finite requires a new
verification). This proves (ii).

\sgaStableTarget{sga6.e10.para.7-18-4}\textbf{7.18.4.}
Suppose that, for every nonempty affine open subset $U$ of the affine space
$\Spec k[T_1,\ldots,T_n]$, if $A$ is the ring of $U$, the homomorphism
$F(k)\to F(A)$ is surjective (compare \hyperlink{sga6.e09.prop.1-6}{IX 1.6}). Using the commutation of $F$
with filtered inductive limits, one then finds that
$F(k)\to F(K)$ is surjective for every pure extension $K$ of $k$. If, in
addition, $k$ is infinite, or if norm homomorphisms can be defined as in
\hyperlink{sga6.e10.para.7-18-3}{7.18.3}, this homomorphism is therefore bijective.

\sgaStableTarget{sga6.e10.para.7-18-5}\textbf{7.18.5.}
As an example, let $X$ be a scheme of finite type over $k$, and set
\[
F(A)=K_\bullet(X\otimes_k A),
\]
which is a functor on the category $C$ by \hyperlink{sga6.e04.para.2-12}{IV 2.12}. By \hyperlink{sga6.e04.para.3-2-3}{IV 3.2.3}, it commutes
with filtered inductive limits. Moreover, if $k'$ is a finite extension of
$k$, the projection
\[
p:X\otimes_k k'\longrightarrow X
\]
is finite, flat, and of finite presentation, and therefore defines a direct
image homomorphism
\[
p_*=N_{k'/k}:K_\bullet(X\otimes_k k')\longrightarrow K_\bullet(X).
\]
The condition on the norm homomorphisms in \hyperlink{sga6.e10.para.7-18-3}{7.18.3} is immediately verified.
Finally, if $A$ is the affine algebra of an open subset $U$ of the affine
space $\Spec k[T_1,\ldots,T_n]$, then $F(k)\to F(A)$ identifies with the map
\[
K_\bullet(X)\longrightarrow K_\bullet(X\times_k U),
\]
which is surjective for $X$ noetherian, by \hyperlink{sga6.e09.prop.1-6}{IX 1.6}. The results proved above
therefore give:
\begin{enumerate}[label=(\roman*)]
\item If $K/k$ is a pure extension of $k$, then
  $K_\bullet(X)\to K_\bullet(X_K)$ is injective, and is even bijective if
  $X$ is noetherian.
\item If $k$ is separably closed and $K/k$ is a separable extension, then
  $K_\bullet(X)\to K_\bullet(X_K)$ is injective.
\item For every extension $K$ of $k$, the kernel of
  $K_\bullet(X)\to K_\bullet(X_K)$ is a torsion group.
\end{enumerate}
These results remain valid if one assumes only that $X$ is noetherian,
provided one restricts to extensions $K$ of $k$ of finite type (so that
$X_K$ is still noetherian). One obtains essentially the same results for
$K^\bullet$ in place of the functor $K_\bullet$, with the sole difference
that in (i) bijectivity cannot be asserted in general, because there is no
homotopy theorem for $K^\bullet$ in place of $K_\bullet$. To conclude
bijectivity, one would have to assume that $X$ is noetherian and regular;
then the groups $K^\bullet$ and $K_\bullet$ agree for $X$ and, respectively,
for $X_K$.

One obtains essentially the same results for $\Gr^\bullet$ in place of
$K_\bullet$ or $K^\bullet$, with the same reservation concerning the
bijectivity assertion in (i). If one restricts to $\Gr^1=\Pic$, that
assertion is valid when $X$ is noetherian and normal: in that case, for every
open subset $U$ of an affine space over $k$, one sees at once that
\[
\Pic(X)\longrightarrow\Pic(X\times_k U)
\]
is surjective (EGA $\mathrm{Err}_{\mathrm{IV}}$ 53, 21.4.13).

\sgaStableTarget{sga6.e10.bibliography}\subsection*{Bibliography}
\begin{enumerate}[label={[\arabic*]}]
\item Dickson, E. \emph{Linear groups}, Dover, New York, 1958.
\item Grothendieck, A. TDTE, V, VI; S\'eminaire Bourbaki 232, 236.
\item Kleiman, S. \emph{Towards a numerical theory of ampleness}, Ann. of
  Math. 84 (1966), pp. 293--344.
\item Murre, J. \emph{On contravariant functors from the category of
  preschemes over a field into the category of abelian groups}, Publications
  Math\'ematiques de l'IHES 23 (1964).
\item Chevalley, C. \emph{Les classes d'\'equivalence rationnelles, I et II},
  S\'eminaire \`a l'ENS 1958, expos\'es 2 et 3.
\item S. Abhyankar, \emph{Resolution of singularities of arithmetical
  surfaces}, in: \emph{Arithmetical Algebraic Geometry}, Purdue University,
  Harper--Row, New York, 1965.
\end{enumerate}

\providecommand{\Hom}{\operatorname{Hom}}
\providecommand{\Isom}{\operatorname{Isom}}
\providecommand{\Coker}{\operatorname{Coker}}
\providecommand{\Ass}{\operatorname{Ass}}
\providecommand{\sgacircled}[1]{\tikz[baseline=(char.base)]\node[shape=circle,draw,inner sep=.35pt](char){\mbox{\scriptsize #1}};}
\newpage
\sgaStableTarget{sga6.e12.expose}\section*{Expos\'e XII. A relative representability theorem for the Picard functor}
\begin{center}
by M. Raynaud (written up by S. Kleiman)
\end{center}

\sgaStableTarget{sga6.e12.warning}\subsection*{Warning}

The purpose of \hyperlink{sga6.e12.expose}{Expos\'es XII} and XIII is to prove the nine finiteness theorems
for the Picard functor stated by A. Grothendieck in the comments
(pp. C 07--C 11) to his \emph{Fondements de la G\'eom\'etrie Alg\'ebrique}
(cited as FGA). In the first expos\'e, the first finiteness theorem is proved
in the following form, which is more precise than the statement in FGA. Let
$X$ and $Y$ be proper schemes over an integral noetherian base $S$, and let
$f:X\to Y$ be a surjective $S$-morphism. Then there is a nonempty open subset
$V$ of $S$ over which the induced morphism
$f^*:\Pic_{Y/S}\to\Pic_{X/S}$ is representable by a quasi-affine morphism of
finite presentation. In the second expos\'e, the remaining assertions are
proved.

The two expos\'es differ notably in their proof techniques. In the first,
Oort's d\'evissage, non-flat descent, and the representability of unramified
functors are used in an essential way. In the second, one first reduces to
the case in which the schemes $X/S$ under consideration are projective and
flat with geometrically integral fibers, or, when necessary, normal or
smooth, always by applying the first finiteness theorem to a suitable
morphism $X'\to X$. One then uses the tool of $(b)$-sheaves, forged by
I. Castelnuovo and sharpened and used by T. Matsusaka, D. Mumford, Y. Nakai,
and S. Kleiman. In particular, this tool makes it possible to dispense with
Chow coordinates, the Frobenius morphism, and the Lefschetz theorems used in
Grothendieck's original proofs.

The present \hyperlink{sga6.e12.expose}{Expos\'e XII} was presented orally by M. Raynaud following
Grothendieck; the editor wishes to thank him for communicating his personal
notes.

\sgaStableTarget{sga6.e12.sec.1}\section*{1. Statement of the main theorem and applications}

\sgaStableTarget{sga6.e12.thm.1-1}\textbf{Theorem 1.1.} Let $S$ be a noetherian scheme, let $X$ and $Y$ be two
proper $S$-schemes, and let $f:X\to Y$ be a morphism. Suppose that $S$ is
integral and that $f$ is surjective. Then there is a nonempty open subset $V$
of $S$ such that the morphism
\[
  f^*:\Pic_{Y/S}|_V\longrightarrow\Pic_{X/S}|_V
\]
between the Picard functors (localized for the fppf topology) is representable
by a quasi-affine morphism of finite presentation.

\sgaStableTarget{sga6.e12.cor.1-2}\textbf{Corollary 1.2.} Let $X$ be a proper scheme over a noetherian scheme
$S$, and suppose that $S$ is integral. Then there is a nonempty open subset
$V$ of $S$ such that the Picard functor $\Pic_{X/S}|_V$ is representable by a
$V$-group scheme that is a sum of quasi-projective open subschemes of finite
type, and hence, \emph{a fortiori}, is separated over $V$.

Indeed, recall the following fact (FGA no.~232). Let $S$ be a scheme and let
$X$ be a projective $S$-scheme, of finite presentation and integral over $S$
(that is, flat with geometrically integral fibers). Then the Picard functor is
represented by an $S$-group scheme $\Pic_{X/S}$, an increasing union of
quasi-projective open subschemes of finite presentation. Consideration of the
``Hilbert polynomials'' decomposes $\Pic_{X/S}$ here as a sum of open
subschemes $\Pic^q_{X/S}$, $q\in\mathbb Q[t]$. In the following expos\'e it is
proved, under the additional hypothesis that $S$ is noetherian, that the
$\Pic^q_{X/S}$ are of finite type over $S$. Thus \hyperlink{sga6.e12.cor.1-2}{1.2} is known when $X$ is in
addition projective and integral over $S$; the general case will be deduced
using the following lemma.

\sgaStableTarget{sga6.e12.lem.1-3}\textbf{Lemma 1.3.} Let $X$ be of finite type over an integral noetherian
scheme $S$. Then there are a nonempty open subset $V$ of $S$, an integral
scheme $S'$, and a finite
flat morphism $S'\to V$ such that every irreducible component $X'_i$ of
$X_{S'}$, endowed with the induced reduced structure, is integral over $S'$.

Indeed, there is a finite extension $L$ of the function field $k(S)$ such
that every irreducible component of $(X_L)_{\red}$ is geometrically integral
(EGA IV 4.6.8). Let $S_1$ be an integral $S$-scheme such that $k(S_1)=L$ and
such that there are a nonempty open subset $V_1$ of $S$ and a finite flat
morphism $h:S_1\to V_1$. There is a nonempty open subset $S_2$ of $S_1$ such
that every irreducible component of $X_{S_2}$, endowed with the induced
reduced structure, is integral over $S_2$ (EGA IV 9.7.7). Set
\[
  V=V_1\setminus h(S_1\setminus S_2),\qquad S'=h^{-1}(V).
\]

We now prove 1.2. Let $h:S'\to V_1$ and the $X'_i$ be as in 1.3. For each
$i$, let $f_i:Z_i\to X'_i$ be a Chow presentation, with $Z_i$ integral and
$f_i$ birational. After shrinking $V$ and $S'$, one may suppose that $Z_i$
is integral over $S'$ (EGA IV 9.7.7). Let
\[
  f:Z=\coprod_i Z_i\longrightarrow X'=X_{S'}
\]
be the canonical morphism. Then
$\Pic_{Z/S'}=\prod_i\Pic_{Z_i/S'}$ is representable by a group scheme that is
a sum of quasi-projective open subschemes of finite type. Finally, \hyperlink{sga6.e12.thm.1-1}{1.1} says
that, after shrinking $V$ and $S'$, the morphism
\[
  f^*:\Pic_{X'/S'}\longrightarrow\Pic_{Z/S'}
\]
is representable by a quasi-affine morphism. Consequently,
$\Pic_{X'/S'}$ is representable by a group scheme that is a sum of
quasi-projective open subschemes of finite type; by finite flat descent
(SGA 1 VIII 7.6 and EGA II 6.6.1), the same is true of
$\Pic_{X/S}|_V$.

\sgaStableTarget{sga6.e12.lem.1-4}\textbf{Lemma 1.4.} Let $k$ be a field, let $P$ and $Q$ be $k$-group schemes
locally of finite type, and let $f:P\to Q$ be a quasi-affine morphism. Then
$f$ is affine.

Indeed, let $R=\Ker(f)$. Then $R$ is a quasi-affine $k$-group scheme of
finite type; by SGA 3 $\mathrm{VI}_{\mathrm B}$ 11.11, $R$ is affine.
Thus $P\to P/R$, which makes $P$ an $R$-torsor over $P/R$, is affine.
Moreover, by SGA 3 I 4.2, $P/R\to Q$ is a closed immersion, hence affine;
therefore the composite $P\to P/R\to Q$ is affine.

\sgaStableTarget{sga6.e12.cor.1-5}\textbf{Corollary 1.5.} Let $k$ be a field.
\begin{enumerate}[label=\alph*)]
\item If $X$ is a proper $k$-scheme, then $\Pic_{X/k}$ is representable
  (J. P. Murre).
\item If $f:X\to Y$ is a surjective morphism between two proper $k$-schemes,
  then the induced morphism
  \[
    f^*:\Pic_{Y/k}\longrightarrow\Pic_{X/k}
  \]
  is affine.
\end{enumerate}

\sgaStableTarget{sga6.e12.cor.1-6}\textbf{Corollary 1.6.} Let $S$ be an integral noetherian scheme, and let
$f:X\to Y$ be a surjective $S$-morphism between two proper $S$-schemes.
\begin{enumerate}[label=\alph*)]
\item If $Y$ is normal over $S$ (that is, flat with geometrically normal
  fibers), then, after shrinking $S$,
  $f^*:\Pic_{Y/S}\to\Pic_{X/S}$ is quasi-finite.
\item If $\Pic_{Y/S}$ is essentially proper over $S$ (that is, if it
  satisfies the valuative criterion of properness; for example, this holds
  if $Y$ is smooth over $S$), then, after shrinking $S$,
  $f^*:\Pic_{Y/S}\to\Pic_{X/S}$ is finite (and, in particular, affine).
\end{enumerate}
Indeed, in (a), for every $s\in S$, $R_s=\Ker(f_s^*)$ is an affine algebraic
group; its identity component is a proper algebraic group over $k(s)$
(FGA 236, Theorem 2.1 (iii)), so $R_s$ is finite. In (b), $f^*$ is both
proper and quasi-affine, hence finite.

\sgaStableTarget{sga6.e12.rem.1-7}\textbf{Remark 1.7.} It is not known whether \hyperlink{sga6.e12.thm.1-1}{1.1} remains valid when
``quasi-affine'' is replaced by ``affine,'' as \hyperlink{sga6.e12.cor.1-5}{1.5} and \hyperlink{sga6.e12.cor.1-6}{1.6} might suggest.

\sgaStableTarget{sga6.e12.sec.2}\section*{2. First reductions}

\sgaStableTarget{sga6.e12.para.2-1}\textbf{2.1.} Apply \hyperlink{sga6.e12.lem.1-3}{1.3} to $X/S$ and $Y/S$. One obtains a nonempty open
subset $V$ of $S$, an integral scheme $S'$, and a finite flat morphism
$S'\to V$ such that every irreducible component $X'_i$ of $X_{S'}$ and
$Y'_j$ of $Y_{S'}$, endowed with the induced reduced structure, is integral
over $S'$. By fpqc descent (EGA IV 2.7.1 (xiv) and SGA 1 VIII 7.9), one may
replace $S$ by $S'$. It then follows from the commutative diagram
\[
\begin{tikzcd}
X_{\red} \arrow[r,"f_{\red}"] \arrow[d,"i"'] & Y_{\red} \arrow[d,"j"]\\
X \arrow[r,"f"'] & Y
\end{tikzcd}
\]
which implies the commutative diagram
\[
\begin{tikzcd}
\Pic_{X_{\red}/S} & \Pic_{Y_{\red}/S} \arrow[l]\\
\Pic_{X/S} \arrow[u] & \Pic_{Y/S} \arrow[l] \arrow[u]
\end{tikzcd}
\]
and from elementary sorites concerning morphisms of functors representable
by quasi-affine morphisms (EGA II 5.1.10 (ii), (v)) that it is enough to prove
the following two assertions:
\begin{enumerate}[label=\Roman*.]
\item Theorem 1.1 when $f$ is an immersion (necessarily nilpotent).
\item Theorem 1.1 when $X$ and $Y$ are reduced and every irreducible component
  $X_i$ of $X$ and $Y_i$ of $Y$, endowed with the induced reduced structure,
  is integral over $S$.
\end{enumerate}

\sgaStableTarget{sga6.e12.para.2-2}\textbf{2.2.} Consider assertion II. For every $i$, let $j=\varphi(i)$ be an
index such that $f(X_i)\subset Y_j$, and put
$\Sigma_j=\coprod_{\varphi(i)=j}X_i$. Then one has the two diagrams
\[
\begin{tikzcd}
Y & \coprod_jY_j \arrow[l]\\
X \arrow[u] & \coprod_j\Sigma_j \arrow[l] \arrow[u]
\end{tikzcd}
\]
and
\[
\begin{tikzcd}[column sep=large,row sep=large]
\Pic_{Y/S} \arrow[r,"\sgacircled{1}"] \arrow[d,"\sgacircled{2}"'] &
\Pic_{\coprod_jY_j/S}=\prod_j\Pic_{Y_j/S}
  \arrow[d,"\sgacircled{4}"]\\
\Pic_{X/S} \arrow[r,"\sgacircled{3}"'] &
\Pic_{\coprod_j\Sigma_j/S}=\prod_j\Pic_{\Sigma_j/S}.
\end{tikzcd}
\]
To establish II, it is enough to see that, after shrinking $S$,
$\sgacircled{1}$, $\sgacircled{3}$, and $\sgacircled{4}$ are quasi-affine.

Consider $\sgacircled{4}$ first. Replacing $Y$ by $Y_j$ and $X$ by
$\Sigma_j$, one is reduced to the case in which $Y$ is integral over $S$.
Since $f$ is surjective, one of the $X_i$, say $X_0$, dominates $Y$, and one
has the two sequences
\[
  X_0\hookrightarrow\coprod_iX_i\longrightarrow Y
\]
and
\[
  \Pic_{X_0/S}\longleftarrow\Pic_{\coprod_iX_i/S}
  \longleftarrow\Pic_{Y/S}.
\]
Thus, if the $\Pic_{X_i/S}$, and hence also
$\Pic_{\coprod_iX_i/S}$, are known to be representable and the composite
morphism is quasi-affine, then $\sgacircled{4}$ is quasi-affine. In view of
the proof of \hyperlink{sga6.e12.cor.1-2}{1.2}, proving $\sgacircled{4}$ is therefore reduced to proving
the following assertion:

\medskip
\noindent
$\mathrm{II}'$. \hyperlink{sga6.e12.thm.1-1}{Theorem 1.1} when $X$ and $Y$ are integral over $S$.
\medskip

For $\sgacircled{1}$ and $\sgacircled{3}$, one must treat the case of a
morphism $\coprod_iX_i\to X$, where $X$ is reduced and the $X_i$ are the
  irreducible components of $X$ and are integral over $S$. One may suppose that
  $X$ has connected fibers. Indeed, let $\eta$ be a geometric generic point of
  $S$.
Then the irreducible components of $X_\eta$ are the $(X_i)_\eta$, and, after
shrinking $S$, one may suppose that $X_i\cap X_j=\varnothing$ if and only if
$(X_i)_\eta\cap(X_j)_\eta=\varnothing$. Let
$Y_k=\bigcup_{i\in I(k)}X_i$ be the connected components of $X$. Then $f$
decomposes as
\[
  f=\coprod_k f_k:
  \coprod_k\left(\coprod_{i\in I(k)}X_i\right)
  \longrightarrow\coprod_kY_k,
\]
and the assertion follows. Finally, an immediate induction reduces us to the
case of a morphism $X_1\amalg X_2\to X$. After shrinking $S$,
$\sgacircled{1}$ and $\sgacircled{3}$ then follow from EGA III 7.8.6 and the
following proposition.

\sgaStableTarget{sga6.e12.prop.2-3}\textbf{Proposition 2.3.} Let $S$ be a scheme, let $g:Y\to S$ be a proper
morphism, and let $Y_1,Y_2$ be two closed subschemes defined by the ideals
$I_1,I_2$. Suppose that:
\begin{enumerate}[label=\alph*)]
\item $Y,Y_1,Y_2,Y_1\cap Y_2$ are faithfully flat and of finite presentation
  over $S$;
\item $g_*\mathcal O_Y=\mathcal O_S$ and
  $(g_i)_*\mathcal O_{Y_i}=\mathcal O_S$ universally, where
  $g_i=g|_{Y_i}$, $i=1,2$;
\item $I_1\cap I_2=0$.
\end{enumerate}
Then
\[
  \Pic_{Y/S}\longrightarrow\Pic_{(Y_1\amalg Y_2)/S}
\]
is representable by an affine morphism of finite presentation.

First note that the hypotheses are stable under every base change. For (c),
this follows from the two exact sequences
\[
  0\longrightarrow I_1+I_2\longrightarrow\mathcal O_Y
  \longrightarrow\mathcal O_{Y_1\cap Y_2}\longrightarrow0
\]
and
\[
  0\longrightarrow I_1\cap I_2\longrightarrow I_1\oplus I_2
  \longrightarrow I_1+I_2\longrightarrow0.
\]
The first shows that $I_1+I_2$ is flat over $S$, and the second then shows
that $I_1\cap I_2$ is stable under base change.

The assertion is local for the fppf topology. After the base change
$Y_1\cap Y_2\to S$, one may therefore suppose that $Y_1\cap Y_2$ has a
section $\alpha$, which defines sections $\alpha_i$ of $Y_i$ and $\beta$ of
$Y$. After a suitable base change, one is reduced to showing that every point
$\mu:S\to\Pic_{(Y_1\amalg Y_2)/S}$ has as inverse image a subfunctor of
$\Pic_{Y/S}$ that is representable by an affine scheme over $S$. Now $\mu$
is defined by an invertible sheaf $L_i$ on $Y_i$, furnished with a
rigidification
\[
  u_i:\alpha_i^*L_i\xrightarrow{\sim}\mathcal O_S,
  \qquad i=1,2.
\]
Likewise, a point $\lambda:S\to\Pic_{Y/S}$ corresponds to an invertible
sheaf $L$ on $Y$ furnished with a rigidification
$u:\beta^*L\xrightarrow{\sim}\mathcal O_S$. The image of $\lambda$ in
$\Pic_{(Y_1\amalg Y_2)/S}$ is $\mu$ if and only if there are isomorphisms of
rigidified sheaves $f_i^*L\simeq L_i$, where $f_i:Y_i\hookrightarrow Y$ is
the inclusion, for $i=1,2$. These isomorphisms define an isomorphism
\[
  \varphi:L_1|_{Y_1\cap Y_2}\xrightarrow{\sim}
  L_2|_{Y_1\cap Y_2}
\]
such that $u_1=u_2\circ\varphi$.

Conversely, if $\varphi$ is such an isomorphism, there is a rigidified
invertible sheaf $L$ on $Y$ obtained by gluing along $\varphi$. Explicitly,
if $f_{12}:Y_1\cap Y_2\hookrightarrow Y$ denotes the common restriction of
$f_1$ and $f_2$, then $L$ is the kernel of the canonical homomorphism
\[
  (f_1)_*L_1\oplus(f_2)_*L_2
  \longrightarrow
  (f_{12})_*\bigl(L_2|_{Y_1\cap Y_2}\bigr)
\]
defined by $\varphi$. This $L$ then corresponds to a point $\lambda$ of
$\Pic_{Y/S}$ mapping to $\mu$.

Finally, let $I$ be the functor that assigns, to an $S$-scheme $T$, the set
of isomorphisms
\[
  \varphi:(L_1)_T|_{(Y_1\cap Y_2)_T}\xrightarrow{\sim}
  (L_2)_T|_{(Y_1\cap Y_2)_T}
\]
such that $(u_1)_T=(u_2)_T\circ\varphi$. We have just established an
isomorphism between the inverse image of $\mu$ in $\Pic_{Y/S}$ and the
functor $I$. It is therefore enough to apply the following lemma (with
$X=Y_1\cap Y_2$).

\sgaStableTarget{sga6.e12.lem.2-4}\textbf{Lemma 2.4.} Let $X$ be a flat, proper scheme of finite presentation
over $S$. Let $L,M$ be two invertible sheaves on $X$, and let $\Isom(L,M)$ be
the functor of isomorphisms $L_T\to M_T$, for $T/S$. Then:
\begin{enumerate}[label=\roman*)]
\item $\Isom(L,M)$ is representable by an affine $S$-scheme of finite
  presentation.
\item If, moreover, $L$ and $M$ are furnished with rigidifications (relative
  to a
given section of $X$), the subfunctor of $\Isom(L,M)$ formed by the
isomorphisms compatible with these rigidifications is representable by a
closed subscheme of $\Isom(L,M)$.
\end{enumerate}

Indeed, (ii) is clear. For (i), recall that the functors
$T\mapsto\operatorname{Hom}(L_T,M_T)$ and
$T\mapsto\operatorname{Hom}(M_T,L_T)$ are representable by
vector bundles $\mathbb V(E)$ and $\mathbb V(F)$ (EGA III 7.7.8). Thus, if
$u\in\operatorname{Hom}(L_T,M_T)$ and
$v\in\operatorname{Hom}(M_T,L_T)$, the conditions
$u\circ v=1$ and $v\circ u=1$ define a closed subscheme $Z$ of
$\mathbb V(E)\times_S\mathbb V(F)$, and it is clear that
$Z\simeq\Isom(L,M)$.

\sgaStableTarget{sga6.e12.para.2-5}\textbf{2.5.} Finally, let us reduce assertion $\mathrm{II}'$ to the following
assertion:

\medskip
\noindent
$\mathrm{II}^*$. \hyperlink{sga6.e12.thm.1-1}{Theorem 1.1} when $X$ and $Y$ are integral over $S$ and the
sequence
\sgaStableTarget{sga6.e12.eq.unnamed}\begin{equation}\tag{*}
  \mathcal O_Y\longrightarrow f_*\mathcal O_X
  \rightrightarrows g_*\mathcal O_{X\times_YX}
\end{equation}
(where $g:X\times_YX\to Y$ is the structural morphism) is exact.
\medskip

For this, it is enough to prove that there is a factorization of $f$
\[
  X\xrightarrow{f_0}Y_0\xrightarrow{f_1}Y_1
  \longrightarrow\cdots\xrightarrow{f_n}Y_n=Y
\]
such that every $f_i$ satisfies (*) and, after possibly restricting $S$,
every $Y_i$ is integral over $S$.

For this purpose, put $Y_0=\operatorname{Spec}(f_*\mathcal O_X)$, with $f_0$ arising from
the Stein factorization of $f$ (EGA III 4.3.3), so that
$f_{0*}\mathcal O_X\xleftarrow{\sim}\mathcal O_{Y_0}$; hence $f_0$ does
satisfy (*). To factor $Y_0\to Y$ in the same way, one uses the following
lemma.

\sgaStableTarget{sga6.e12.lem.2-6}\textbf{Lemma 2.6.} Let $Y$ be a noetherian scheme, and let $f:X\to Y$ be a
finite epimorphism (that is, $\mathcal O_Y\to f_*\mathcal O_X$ is
injective). Then $f$ factors into a finite sequence of finite morphisms
\[
  X\xrightarrow{\sim}Y_0\xrightarrow{f_1}Y_1
  \longrightarrow\cdots\longrightarrow Y_n=Y
\]
that are effective epimorphisms (that is, for each $i$ the sequence

\[
  \mathcal O_{Y_i}\longrightarrow(f_i)_*\mathcal O_{Y_{i-1}}
  \rightrightarrows
  (g_i)_*\bigl(
    \mathcal O_{Y_{i-1}}\otimes_{\mathcal O_{Y_i}}
    \mathcal O_{Y_{i-1}}
  \bigr)
\]
is exact).

Put $A=\mathcal O_Y$, $B_0=f_*\mathcal O_X$, and define recursively
\[
  B_i=\Ker\bigl(
    B_{i-1}\rightrightarrows B_{i-1}\otimes_A B_{i-1}
  \bigr).
\]
One checks at once that
\[
  B_{i-1}\otimes_A B_{i-1}
  \xrightarrow{\sim}
  B_{i-1}\otimes_{B_i}B_{i-1}.
\]
It remains to prove that $B_i=A$ for all sufficiently large $i$. Indeed,
  consider the decreasing sequence of $A$-modules $M_i=B_i/A$, and put
$S_i=\operatorname{Supp}(M_i)$. The decreasing sequence of closed subsets $S_i$ of $Y$
is stationary; let $S_\infty$ be its stationary value. We show that
$S_\infty=\varnothing$. Otherwise, let $s$ be a maximal point of
$S_\infty$. After replacing $Y$ by $\operatorname{Spec}\mathcal O_{Y,s}$, we may suppose
that $Y$ is local with closed point $s$. By the choice of $s$, $M_i$ is
artinian for all sufficiently large $i$, and therefore the sequence $(M_i)$
is stationary. Thus $B_i=B_{i+1}$ for some $i$. It follows that
$\operatorname{Spec}(B_i)\to Y$ is a monomorphism, hence a closed immersion (Nakayama).
Since $A\subset B_i$, one has $A=B_i$, contrary to $s\in S_\infty$.

\sgaStableTarget{sga6.e12.sec.3}\section*{3. Proof of I: Oort's d\'evissage}

\sgaStableTarget{sga6.e12.prop.3-1}\textbf{Proposition 3.1.} (Oort, \emph{Bull. soc. math. France}, vol.~90
(1962)). Let $S$ be a locally noetherian scheme, let $f:X\to S$ be a proper
morphism, and let $X'\subset X$ be a closed subscheme. Suppose that:
\begin{enumerate}[label=\roman*)]
\item $X'$ is defined by an ideal $I$ contained in the nilradical $N$ of
  $X$, and $NI=0$.
\item $X'$ is flat over $S$ and cohomologically flat over $S$ in dimension
  $0$ (EGA III 7.8.1).
\item $I$ is flat over $S$ and cohomologically flat over $S$ in dimensions
  $0$ and $1$, so that (by EGA III 7.4.6) there are a coherent
  $\mathcal O_S$-module $Q$ and, functorially in the coherent
  $\mathcal O_S$-module $M$, an isomorphism

\[
  R^2f_*\bigl(I\otimes f^*M\bigr)
  \simeq\operatorname{Hom}_{\mathcal O_S}(Q,M).
\]
\item The cokernel $E$ of
  $f_*\mathcal O_{X'}\to R^1f_*(I)$ is a locally free
  $\mathcal O_S$-module.
\item The homomorphism
  \[
    (f_*\mathcal O_X)_{\red}
    \longrightarrow(f_*\mathcal O_{X'})_{\red}
  \]
  is surjective.
\end{enumerate}
Then the canonical morphism
\[
  \Pic_{X/S}\longrightarrow\Pic_{X'/S}
\]
is representable by an affine morphism of finite presentation. More
precisely, there is a canonical exact sequence of fppf sheaves
\[
  0\longrightarrow\mathbb V(E^\vee)\longrightarrow\Pic_{X/S}
  \longrightarrow\Pic_{X'/S}\xrightarrow{u}\mathbb V(Q),
\]
and, if $P=\Ker(u)$, the monomorphism
$P\hookrightarrow\Pic_{X'/S}$ is representable by a closed immersion and
$\Pic_{X/S}$ is a torsor over $P$ under the group $\mathbb V(E^\vee)$
(that is, the sequence of fppf sheaves
\[
  0\longrightarrow\mathbb V(E^\vee)\longrightarrow\Pic_{X/S}
  \longrightarrow P\longrightarrow0
\]
is exact).

For an $S$-scheme $T$, consider the following two diagrams with exact rows:
\[
\begin{tikzcd}
0 \arrow[r] & I_T \arrow[r] & \mathcal O_{X_T} \arrow[r] &
  \mathcal O_{X'_T}^{*} \arrow[r] & 0\\
0 \arrow[r] & I_T \arrow[r,"\alpha"] \arrow[u,equal] &
  \mathcal O_{X_T}^{*} \arrow[r] \arrow[u] &
  \mathcal O_{X'_T}^{*} \arrow[r] \arrow[u] & 0,
\end{tikzcd}
\qquad \alpha(i)=1+i,
\]
and
\[
\begin{tikzcd}[column sep=tiny]
(f_T)_*\mathcal O_{X'_T} \arrow[r,"\delta_T"] &
  R^1(f_T)_*I_T \arrow[d,equal]\\
(f_T)_*\mathcal O_{X'_T}^{*} \arrow[r,"\delta_T^*"'] \arrow[u] &
  R^1(f_T)_*I_T \arrow[r] &
  R^1(f_T)_*\mathcal O_{X_T}^{*} \arrow[r] &
  R^1(f_T)_*\mathcal O_{X'_T}^{*} \arrow[r] &
  R^2(f_T)_*I_T.
\end{tikzcd}
\]
Note carefully that the left-hand square in this diagram is \emph{not
necessarily commutative}. Nevertheless,
\[
  \sgaStableTarget{sga6.e12.para.3-2}\textbf{3.2.}\qquad \operatorname{im}\delta_T^*
  =\operatorname{im}\delta_T.
\]

Indeed, we may plainly suppose that $S$ and $T$ are affine. Choose an affine
covering $\{V_i\}$ of $X_T$. View $\mathcal O_{X'_T}$ as its direct image
under the closed immersion $X'_T\hookrightarrow X_T$. Let
$g\in\Gamma(X_T,\mathcal O_{X'_T}^{*})$, and let $g_i$ be a lift of
$g|_{V_i}$ in $\Gamma(V_i,\mathcal O_{X_T}^{*})$. Then $\delta_T^*(g)$ is
the class of the $1$-cocycle
\[
  \delta_{ij}^*=g_i g_j^{-1}-1,
\]
whereas $\delta_T(g)$ is the class of
\[
  \delta_{ij}=g_i-g_j=g_j\delta_{ij}^*=g\delta_{ij}^*.
\]
Consequently,
\[
  \delta_T(g)=g\,\delta_T^*(g),
\]
taking into account the action of
$\Gamma(X_T,\mathcal O_{X'_T})$ on $I_T$, and hence on $H^1(X_T,I_T)$.

By (ii),
\[
  \Gamma(X_T,\mathcal O_{X'_T})
  =\Gamma(X,\mathcal O_{X'})\otimes_S\Gamma(T,\mathcal O_T),
\]
and hence $g=\sum g_\beta\otimes t_\beta$, with
$g_\beta\in\Gamma(X,\mathcal O_{X'})$ and
$t_\beta\in\Gamma(T,\mathcal O_T)$. By (v), the map
\[
  (f_*\mathcal O_X)_{\red}\longrightarrow
  (f_*\mathcal O_{X'})_{\red}
\]
is surjective. Thus
$g_\beta=v(h_\beta)+n_\beta$, where
$h_\beta\in\Gamma(X,\mathcal O_X)$,
$v:\Gamma(X,\mathcal O_X)\to\Gamma(X,\mathcal O_{X'})$, and
$n_\beta\in\operatorname{nilrad}\bigl(\Gamma(X,\mathcal O_{X'})\bigr)$.

By (i), $(n_\beta\otimes t_\beta)I_T=0$. Therefore
$\delta_T(g)=h\,\delta_T^*(g)$, where
\[
  h=\sum h_\beta\otimes t_\beta
  \in\Gamma(X_T,\mathcal O_{X_T}).
\]
It follows that
\[
  \delta_T^*(g)=h^{-1}\delta_T(g)=\delta_T(h^{-1}g),
\]
and hence $\operatorname{im}\delta_T^*\subset
\operatorname{im}\delta_T$.

Finally, let
$g=\sum g_\beta\otimes t_\beta\in\Gamma(X_T,\mathcal O_{X'_T})$.
Write $g_\beta=v(h_\beta)+n_\beta$, with
$n_\beta\in\operatorname{nilrad}\bigl(\Gamma(X,\mathcal O_{X'})\bigr)$. Then
\[
\begin{split}
  \delta_T(g)
  &=\delta_T\Bigl(\sum n_\beta\otimes t_\beta\Bigr)\\
  &=\delta_T\Bigl(1+\sum n_\beta\otimes t_\beta\Bigr)\\
  &=\delta_T^*\Bigl(1+\sum n_\beta\otimes t_\beta\Bigr),
\end{split}
\]
whence $\operatorname{im}\delta_T\subset
\operatorname{im}\delta_T^*$.

By \hyperlink{sga6.e12.para.3-2}{3.2}, (iii), and (iv), there is an exact sequence, functorial in $T$,
\[
  0\longrightarrow\mathbb V(E^\vee)(T)
  \longrightarrow R^1(f_T)_*\mathcal O_{X_T}^{*}
  \longrightarrow R^1(f_T)_*\mathcal O_{X'_T}^{*}
  \longrightarrow\mathbb V(Q)(T),
\]
which, on passing to the associated fppf sheaves, gives an exact sequence
\[
  0\longrightarrow\mathbb V(E^\vee)\longrightarrow\Pic_{X/S}
  \longrightarrow\Pic_{X'/S}\xrightarrow{u}\mathbb V(Q).
\]
If $P=\Ker(u)$, the monomorphism $P\hookrightarrow\Pic_{X'/S}$ is obtained
by base change from the unit section $S\hookrightarrow\mathbb V(Q)$.
Since $\mathbb V(Q)$ is representable and separated, it follows that
$P\hookrightarrow\Pic_{X'/S}$ is representable by a

closed immersion.

Finally,
\[
  0\longrightarrow\mathbb V(E^\vee)\longrightarrow\Pic_{X/S}
  \xrightarrow{w}P\longrightarrow0
\]
is a torsor for the fppf topology. A fortiori, by fpqc descent from a
``splitting'' of the torsor, $w$ is representable by a faithfully flat
affine morphism of finite presentation (see Oort, \emph{Commutative group
schemes}, Proposition 17.4).

\sgaStableTarget{sga6.e12.cor.3-3}\textbf{Corollary 3.3.} Under the hypotheses of \hyperlink{sga6.e12.prop.3-1}{3.1}:
\begin{enumerate}[label=\roman*)]
\item If
  \[
    \Pic_{(\mathrm{zar})(X'/S)}\longrightarrow
    \Pic_{(\mathrm{fppf})(X'/S)}
  \]
  is an isomorphism, then
  \[
    \Pic_{(\mathrm{zar})(X/S)}\longrightarrow
    \Pic_{(\mathrm{fppf})(X/S)}
  \]
  is an isomorphism.
\item If the $S$-scheme $T$ is affine and
  \[
    \Pic(X'_T)\longrightarrow\Pic_{(\mathrm{zar})(X'/S)}^{(T)}
  \]
  is surjective, then
  \[
    \Pic(X_T)\longrightarrow\Pic_{(\mathrm{zar})(X/S)}^{(T)}
  \]
  is surjective.
\end{enumerate}

Indeed, (i) follows by applying the five lemma to the following diagram of
sheaves for the Zariski topology:
\[
\begin{tikzcd}[column sep=small]
0 \arrow[r] \arrow[d,"S"] &
  \mathbb V(E^\vee) \arrow[r] \arrow[d,"S"] &
  \Pic_{(\mathrm{zar})(X/S)} \arrow[r] \arrow[d] &
  \Pic_{(\mathrm{zar})(X'/S)} \arrow[r] \arrow[d,"S"] &
  \mathbb V(Q) \arrow[d,"S"]\\
0 \arrow[r] & \mathbb V(E^\vee) \arrow[r] &
  \Pic_{(\mathrm{fppf})(X/S)} \arrow[r] &
  \Pic_{(\mathrm{fppf})(X'/S)} \arrow[r] & \mathbb V(Q).
\end{tikzcd}
\]

Assertion (ii) follows from the following diagram with exact rows:
\[
\begin{tikzcd}[column sep=small]
H^1(X_T,I_T) \arrow[r] \arrow[d] &
  H^1(X_T,\mathcal O_{X_T}^{*}) \arrow[r] \arrow[d] &
  H^1(X_T,\mathcal O_{X'_T}^{*}) \arrow[r] \arrow[d] &
  H^2(X_T,I_T) \arrow[d]\\
H^0(T,R^1(f_T)_*I_T) \arrow[r] &
  H^0(T,R^1(f_T)_*\mathcal O_{X_T}^{*}) \arrow[r] &
  H^0(T,R^1(f_T)_*\mathcal O_{X'_T}^{*}) \arrow[r] &
  H^0(T,R^2(f_T)_*I_T),
\end{tikzcd}
\]
because the outside vertical arrows are isomorphisms when $T$ is affine.

\sgaStableTarget{sga6.e12.cor.3-4}\textbf{Corollary 3.4.} Let $k$ be a field, and let $X$ be a proper
$k$-scheme. Suppose that $\Pic_{X_{\red}/k}$ possesses a universal
invertible sheaf. (For example, descent theory shows that this is the case
if every connected component $X'_i$ of $X_{\red}$ has a rational point; in
particular, if $k$ is algebraically closed, or more generally if the greatest
common divisor of the degrees of the $0$-cycles on every $X'_i$ is $1$.)
Then, over every affine open subset $V$ of $\Pic_{X/k}$, there is a universal
invertible sheaf on $X\times_kV$.

Finally, we prove assertion I in the following more precise form.

\sgaStableTarget{sga6.e12.prop.3-5}\textbf{Proposition 3.5.} Let $S$ be an integral noetherian scheme, and let
$f:X\hookrightarrow Y$ be a nilpotent immersion of proper $S$-schemes.
Then there is a nonempty open subset $V$ of $S$ such that the morphism
\[
  f^*:\Pic_{Y/S}|_V\longrightarrow\Pic_{X/S}|_V
\]
is representable by an affine morphism of finite presentation. In fact,
$N=\Ker(f^*)$ is an affine group scheme with unipotent fibers (in the sense
of SGA 3 XVII).

Indeed, $X$ is defined by a nilpotent ideal $I$. The ideal $I$ has a
filtration by ideals
\[
  I_0=I\supset I_1\supset\cdots\supset I_n=0
\]
such that $NI_{i-1}\subset I_i$ for $1\leq i\leq n$, where $N$ is the
nilradical of $\mathcal O_Y$ (for example, $I_i=N^iI$). Let $X_i$ be the
closed subscheme of $Y$ defined by $I_i$. By induction it is enough to study
the immersions $X_{i-1}\hookrightarrow X_i$. We may therefore suppose that
$NI=0$; this is condition (i) of \hyperlink{sga6.e12.prop.3-1}{3.1}, and we are reduced finally to the case
in which the other conditions (ii)--(v) are satisfied as well. By the generic
flatness theorem (EGA IV 6.9.1), applied in the notation of \hyperlink{sga6.e12.prop.3-1}{3.1} to
$\mathcal O_{X'}$, the $R^if_*(I)$, and $E$, and taking account of the
K\"unneth relations of EGA III 6.9.9, after restricting $S$ we may suppose
that conditions (ii), (iii), and (iv) hold.

There is a finite extension $K$ of $k(S)$ such that
$(f_K)_*\mathcal O_{X'_K}$ has all its residue field extensions trivial.
If $S'$ is an integral scheme, finite and flat over an open subset $V$ of
$S$, with function field $k(S')=K$, there is an open subset $V'$ of $S'$
such that
\[
  \bigl[(f_{S'})_*\mathcal O_{X'_{S'}}\bigr]_{\red}
  \simeq
  \bigl[(f_{S'})_*\mathcal O_{X_{S'}}\bigr]_{\red}
\]
over $V'$. After restricting $S$ and $S'$ further, we may suppose that
$S=V$ and $S'=V'$. Since the assertion to be proved is local for the fpqc
topology, we may replace $S$ by $S'$. Finally, apply 3.1.

\sgaStableTarget{sga6.e12.rem.3-6}\textbf{Remark 3.6.} It follows in the same way, taking account of the exact
sequences
\[
  \Pic_{X_i/S}\longrightarrow\Pic_{X_{i-1}/S}
  \longrightarrow\mathbb V(Q_i),
\]
that, for every $s\in S$,
\[
  \operatorname{coker}\bigl(
    \Pic^0_{Y_s/s}\longrightarrow\Pic^0_{X_s/s}
  \bigr)
\]
is a unipotent group.

\sgaStableTarget{sga6.e12.sec.4}\section*{4. Proof of $\mathrm{II}^*$: the most delicate part of the proof}

Recall assertion $\mathrm{II}^*$: \hyperlink{sga6.e12.thm.1-1}{Theorem 1.1} when $X$ and $Y$ are integral
over $S$ and the sequence
\[
  \mathcal O_Y\longrightarrow f_*\mathcal O_X
  \rightrightarrows g_*\mathcal O_{X\times_YX}
\]
is exact. Applying the generic flatness theorem, EGA IV 6.9.1, and EGA III
7.8.8, we are reduced to proving the following proposition.

\sgaStableTarget{sga6.e12.prop.4-1}\textbf{Proposition 4.1.} Let $S$ be an integral noetherian scheme, let
$\lambda:X\to S$ and $\mu:Y\to S$ be proper morphisms of finite
presentation, let $f:X\to Y$ be a surjective $S$-morphism, and let
$g:X\times_YX\to Y$ be the canonical morphism. Suppose that:
\begin{enumerate}[label=\arabic*)]
\item $X$, $Y$, and $X\times_YX$ are flat over $S$.
\item $\lambda_*\mathcal O_X=\mathcal O_S$ and
  $\mu_*\mathcal O_Y=\mathcal O_S$ universally.

\item For all $i,j\geq0$, the modules $R^if_*\mathcal O_X$ and
  $R^jg_*\mathcal O_{X\times_YX}$ are flat over $S$.
\item The sequence
  \sgaStableTarget{sga6.e12.eq.unnamed.occ2}\begin{equation}\tag{*}
    \mathcal O_Y\longrightarrow f_*\mathcal O_X
    \mathrel{\substack{\xrightarrow{\pi_1}\\[-0.7ex]
                       \xrightarrow[\pi_2]{}}}
    g_*\mathcal O_{X\times_YX}
  \end{equation}
  is exact, and $\operatorname{coker}(\pi_1-\pi_2)$ is an
  $\mathcal O_Y$-module flat over $S$.
\end{enumerate}
Then
\[
  f^*:\Pic_{Y/S}\longrightarrow\Pic_{X/S}
\]
is representable by a quasi-affine morphism of finite presentation.

First note that (1)--(4) are stable under every base change $T\to S$. Only
the stability of (3) and (4) is not evident. By the K\"unneth formula (EGA III
6.9.9), (1) and (3) imply that the formation of $R^if_*\mathcal O_X$ and
$R^ig_*\mathcal O_{X\times_YX}$ commutes with base change $T\to S$; this
immediately implies the stability of condition (3). Moreover, the formation
of the cokernel of
\[
  \pi_1-\pi_2:f_*\mathcal O_X
  \longrightarrow g_*\mathcal O_{X\times_YX}
\]
also commutes with base change. Since it is flat over $S$ by (4),
$\operatorname{im}(\pi_1-\pi_2)$ is flat over $S$ and commutes with base
change. Hence $\Ker(\pi_1-\pi_2)$ commutes with base change $T\to S$ and is
therefore universally equal to $\mathcal O_Y$.

The stability of the hypotheses implies two things:
\begin{enumerate}[label=\alph*.]
\item Since the assertion to be proved is local for the fppf topology, by
  descent we may suppose that $X$ has a section $\alpha$; then $Y$ has the
  section $\beta=f\circ\alpha$.
\item To prove the assertion, it is enough to consider an $S$-point of
  $\Pic_{X/S}$, that is, by (a) and (2), an invertible sheaf $L$ on $X$
  endowed with a rigidification $u$ relative to $\alpha$, and to show that
  the fppf sheaf $Z$, the inverse image of this point in $\Pic_{Y/S}$, is
  representable by a quasi-affine $S$-scheme of finite presentation.
\end{enumerate}

\sgaStableTarget{sga6.e12.lem.4-2}\textbf{Lemma 4.2.} Under the hypotheses of \hyperlink{sga6.e12.prop.4-1}{4.1}, $f$ is a descent morphism
for the category of locally free sheaves; that is, for every pair $L,M$ of
locally free sheaves on $Y$, the diagram
\[
\begin{tikzcd}[column sep=tiny]
\operatorname{Hom}_Y(L,M) \arrow[r] &
  \operatorname{Hom}_X(f^*L,f^*M)
    \arrow[r,shift left=.7ex] \arrow[r,shift right=.7ex] &
  \operatorname{Hom}_{X\times_YX}(g^*L,g^*M)
\end{tikzcd}
\]
is exact.

Indeed, if $N=\operatorname{Hom}_{\mathcal O_Y}(L,M)$, one must show that the sequence
\[
  \Gamma(Y,N)\longrightarrow\Gamma(X,f^*N)
  \rightrightarrows\Gamma(X\times_YX,g^*N)
\]
is exact, or equivalently (by EGA $0_{\mathrm I}$ 5.4.10.1) that the sequence
\[
  \Gamma(Y,N)
  \longrightarrow
  \Gamma\bigl(Y,f_*\mathcal O_X\otimes_{\mathcal O_Y}N\bigr)
  \rightrightarrows
  \Gamma\bigl(Y,g_*\mathcal O_{X\times_YX}
    \otimes_{\mathcal O_Y}N\bigr)
\]
is exact. But the sequence (*)
\[
  \mathcal O_Y\longrightarrow f_*\mathcal O_X
  \rightrightarrows g_*\mathcal O_{X\times_YX}
\]
is exact by hypothesis. Applying the left exact functor
$\Gamma(N\otimes -)$ gives the lemma, since $N$ is flat over $Y$.

Let $L$ be an invertible sheaf on $X$, endowed with a rigidification $u$
relative to $\alpha$. By \hyperlink{sga6.e12.lem.4-2}{4.2} there is an equivalence between the category
$C$ of pairs $(M,\varphi)$ consisting of an invertible sheaf $M$ on $Y$ and
an isomorphism $\varphi:f^*M\xrightarrow{\sim}L$, and the category
$\mathrm{Deff}(S)$ of effective descent data on $L$ relative to $f$.
If $f^*M\simeq L$, then, since $L$ is rigidified relative to $\alpha$, $M$
is canonically endowed with a rigidification relative to $\beta$.
Consequently, $C$ is equivalent to the discrete category defined by the set
of classes of $\beta$-rigidified invertible sheaves $(M,v)$ on $Y$ such that
\[
  f^*(M,v)\simeq(L,u).
\]
But this is precisely the set $Z(S)$ of $S$-points of the inverse image in
$\Pic_{Y/S}$ of the $S$-point of $\Pic_{X/S}$ defined by $(L,u)$. Moreover,
this description of $Z(S)$ is functorial in $S$. Hence $Z$ identifies with
$\mathrm{Deff}$, and we may forget the rigidifications from now on.


\providecommand{\Isom}{\operatorname{Isom}}
\providecommand{\Coker}{\operatorname{Coker}}
\providecommand{\Ker}{\operatorname{Ker}}
\providecommand{\Supp}{\operatorname{Supp}}
\providecommand{\Ass}{\operatorname{Ass}}

There is a canonical monomorphism from \(\mathrm{Deff}\) into the functor
\(\mathrm{Rec}\) of gluing data on \(L\). By definition,
\(\mathrm{Rec}\) is the functor
\(\Isom(p_1^*L,p_2^*L)\), where the
\(p_i:X\times_YX\to X\) are the projections; by \hyperlink{sga6.e12.lem.2-4}{2.4} this functor is
represented by an affine scheme \(I\) of finite presentation over \(S\). It
therefore remains to represent the monomorphism
\[
  i:\mathrm{Deff}\hookrightarrow \Isom(p_1^*L,p_2^*L)
\]
by a morphism of finite presentation, which will necessarily be quasi-affine
(by EGA IV 8.11.2).

Let us note in passing that, if desired, one may restrict
\(\Isom(p_1^*L,p_2^*L)\) by imposing the cocycle condition, but this is not
needed below.

After the base change \(I\to S\), we are thus reduced to the following
problem: if \(L\) is an invertible sheaf on \(X\), endowed with a gluing datum
\[
  \varphi:p_1^*L\xrightarrow{\sim}p_2^*L,
\]
represent the subfunctor of \(S\) on which this datum becomes effective. We
have the following lemma.

\sgaStableTarget{sga6.e12.lem.4-3}\textbf{Lemma 4.3.} Under the hypotheses of \hyperlink{sga6.e12.prop.4-1}{4.1}, let \(L\) be an invertible
sheaf on \(X\), endowed with a gluing datum
\[
  \varphi:p_1^*L\xrightarrow{\sim}p_2^*L.
\]
This datum is effective if and only if the following five conditions hold.
\begin{enumerate}[label=\arabic*)]
\item For every \(i\geq0\), \(R^if_*(L)\) is flat over \(S\).
\item For every \(j\geq0\), \(R^jg_*(p_1^*L)\) is flat over \(S\).
\item The cokernel of
\[
  \pi_1-\pi_2:f_*(L)\longrightarrow
  g_*(p_1^*L)\simeq g_*(p_2^*L)
\]
is flat over \(S\), where the latter isomorphism is \(g_*(\varphi)\).
\item The kernel \(N\) of
\(\pi_1-\pi_2:f_*(L)\to g_*(p_1^*L)\) is invertible on \(Y\).
\item The composite morphism
\[
  f^*N\longrightarrow f^*f_*L\longrightarrow L
\]
is an isomorphism.
\end{enumerate}

\newpage
Indeed, these conditions are sufficient, since (4) and (5) alone imply
effectivity. Conversely, let \(M\) be invertible on \(Y\), and suppose that
there is an isomorphism \(f^*M\simeq L\) compatible with the gluing datum. We
may cover \(Y\) by open subsets \(V\) over which \(M\) is isomorphic to
\(\mathcal O_Y\). The descent datum on \(L\), restricted over
\(f^{-1}(V)\), is therefore isomorphic to the trivial descent datum on
\(\mathcal O_X|_{f^{-1}(V)}\). Finally, conditions (1)--(5) are local on
\(Y\), and the hypotheses of \hyperlink{sga6.e12.prop.4-1}{4.1} show that they hold for
\(L=\mathcal O_X\); hence they hold for \(L\).

It remains to represent conditions (1)--(5), one after another, by
monomorphisms of finite presentation. For (1) and (2), we apply the following
lemma.

\sgaStableTarget{sga6.e12.lem.4-4}\textbf{Lemma 4.4.} Let \(S\) be a scheme, let
\(\lambda:X\to S\) and \(\mu:Y\to S\) be proper morphisms of finite
presentation, and let \(f:X\to Y\) be an \(S\)-morphism. Let \(F\) be a
sheaf on \(X\), of finite presentation and flat over \(S\). For an integer
\(n\), let \(S_n\) be the subfunctor of \(S\) on which \(R^if_*(F)\) is
flat over \(S\) for every \(i\geq n\); more precisely, for a morphism
\(\phi:T\to S\),
\[
  S_n(T)=
  \begin{cases}
    \{\phi\},&\text{if \(R^i(f_T)_*(F_T)\) is flat for every \(i\geq n\),}\\
    \varnothing,&\text{otherwise.}
  \end{cases}
\]
Then \(S_n\hookrightarrow S\) is representable by a surjective monomorphism
of finite presentation; moreover, for \(i\geq n\), the formation of
\(R^i(f_{S_n})_*(F_{S_n})\) commutes with base changes \(T\to S_n\). In
addition, \(S_n\) coincides with the functor \(S'_n\) on which the functors
\[
  M\longmapsto
  R^i(f_V)_*\bigl(F_V\otimes\lambda_V^*M\bigr),
  \qquad i\geq n,
\]
are exact, where \(M\) ranges over the quasi-coherent
\(\mathcal O_V\)-modules, for every affine open subset \(V\) of \(S\).

Indeed, the assertion is local on \(S\), so we may suppose that \(S\) is
quasi-compact. The relative dimension of \(X/Y\) is then bounded, say by
\(r\), and, for \(n>r\),
\[
  R^nf_*(F_V\otimes\lambda_V^*M)=0.
\]
The assertion therefore holds for \(n>r\).

\newpage
Proceeding by descending induction, suppose it is known for \(n+1\). After
replacing \(S\) by \(S_{n+1}=S'_{n+1}\), we may suppose that
\(R^if_*(F)\) is flat over \(S\) for \(i\geq n+1\). Since \(F\) is flat
over \(S\), EGA III 6.9.8 then shows that the formation of \(R^nf_*(F)\)
commutes with every base change. Consequently, the representability of
\(S_n\) follows from the next lemma (see J. P. Murre,
\emph{Representation of unramified functors}, S\'eminaire Bourbaki, May
1965, p.~294-11, Theorem~2).

\sgaStableTarget{sga6.e12.lem.4-5}\textbf{Lemma 4.5.} Let \(X\to S\) be a proper morphism of finite
presentation, and let \(F\) be a sheaf of finite presentation on \(X\). Then
the subfunctor \(S_F\) of \(S\) on which \(F\) is flat over \(S\) is
representable by a surjective monomorphism of finite presentation (which, of
course, is not an immersion in general).

To represent condition (3) of \hyperlink{sga6.e12.lem.4-3}{4.3}, note that, by (1) and (2), the formation
of \(f_*(L)\) and \(g_*(p_1^*L)\) commutes with the base changes
\(I\to S\); hence so does the formation of
\(\Coker(\pi_1-\pi_2)\). We then apply 4.5.

To represent (4), note that, by (3), and because \(f_*(L)\) and
\(g_*(p_1^*L)\) are flat and commute with base change,
\[
  N=\Ker(\pi_1-\pi_2)
\]
also commutes with base change. We then apply the following lemma.

\sgaStableTarget{sga6.e12.lem.4-6}\textbf{Lemma 4.6.} Let \(\lambda:X\to S\) be a proper morphism of finite
presentation, and let \(u:F\to G\) be a homomorphism of
\(\mathcal O_X\)-modules of finite presentation, with \(G\) flat over
\(S\). Then the subfunctor \(S_1\) of \(S\) on which \(u\) is an
isomorphism is representable by an open subscheme of \(S\), of finite
presentation.

As usual, we reduce to the case where \(S\) is noetherian, using EGA IV 8
and 11.2.6; this relieves us of the need to consider finite presentation.
First let us represent the functor on which \(u\) is an epimorphism,

\newpage
that is, on which \(\Coker u\) is zero. Since the formation of
\(\Coker u\) commutes with base change, this functor is represented by the
open subset
\[
  S-\lambda\bigl(\Supp(\Coker u)\bigr)
\]
of \(S\). We may therefore suppose that \(u\) is already an epimorphism.
Using now the hypothesis that \(G\) is flat over \(S\), we find that the
formation of \(\Ker u\) commutes with base change; consequently, the functor
under consideration is represented by the open subset
\[
  S-\lambda\bigl(\Supp(\Ker u)\bigr)
\]
of \(S\).

\newpage
\sgaStableTarget{sga6.e13.expose}\section*{Expos\'e XIII. Finiteness Theorems for the Picard Functor}
\begin{center}
by S. Kleiman
\end{center}

\sgaStableTarget{sga6.e13.sec.1}\subsection*{1. The \texorpdfstring{\((b)\)}{(b)}-sheaves}

Let \(k\) be an algebraically closed field and \(X\) a projective
\(k\)-scheme, endowed with an ample invertible sheaf \(\mathcal O_X(1)\).

\sgaStableTarget{sga6.e13.def.1-1}\textbf{Definition 1.1 (Mumford).} Let \(m\in\mathbb Z\). A coherent sheaf
\(F\) on \(X\) is said to be \emph{\(m\)-regular} if, at every point of the
support of \(F\), \(\mathcal O_X(1)\) is generated by its global sections
and if, for every \(q\geq1\),
\[
  H^q(F(m-q))=0.
\]

\sgaStableTarget{sga6.e13.lem.1-2}\textbf{Lemma 1.2.} Let
\[
  0\longrightarrow F(-1)\longrightarrow F\longrightarrow G\longrightarrow0
\]
be an exact sequence of coherent sheaves on \(X\).\footnote{Here, as below
in analogous situations, it is understood that \(F(-1)\to F\) is defined
by tensoring with a given section of \(\mathcal O_X(1)\).} If \(F\) is
\(m\)-regular, then so is \(G\).

Indeed, we have the corresponding exact sequence
\[
  H^q(F(m-q))\longrightarrow H^q(G(m-q))
  \longrightarrow H^{q+1}(F(m-q-1)).
\]

\sgaStableTarget{sga6.e13.prop.1-3}\textbf{Proposition 1.3.} If a sheaf \(F\) on \(X\) is \(m\)-regular,
then, for every \(n\geq m\),
\begin{enumerate}[label=(\roman*)]
\item \(F\) is \(n\)-regular.
\item The map
\[
  H^0(F(n))\otimes H^0(\mathcal O_X(1))
  \longrightarrow H^0(F(n+1))
\]
is surjective.
\item \(F(n)\) is generated by \(H^0(F(n))\).
\end{enumerate}

\newpage
The proof is by induction on the dimension \(s\) of the support of \(F\); if
\(s=-1\), the assertion is trivial. Choose a section
\(\sigma\in H^0(\mathcal O_X(1))\) which vanishes at no point of the finite
set \(\Ass(F)\). Then \(\sigma\) defines an exact sequence
\[
  0\longrightarrow F(-1)\longrightarrow F\longrightarrow G\longrightarrow0,
\]
and \(G\) is \(m\)-regular by \hyperlink{sga6.e13.lem.1-2}{1.2}, with support of dimension \(s-1\).

Consider the corresponding exact sequence
\[
  H^q(F(n-q-1))\longrightarrow H^q(F(n-q))
  \longrightarrow H^q(G(n-q)).
\]
By the induction hypothesis, \(H^q(G(n-q))=0\) for \(n\geq m\), so the
vanishing of \(H^q(F(n-q-1))\) implies that of \(H^q(F(n-q))\). But
\(H^q(F(m-q))=0\), which proves (i).

To prove (ii), we chase through the following commutative diagram:
\[
\begin{tikzcd}[column sep=1.7em,row sep=2.4em]
  & & H^0(F(n))\otimes H^0(\mathcal O_X(1))
      \arrow[rr] \arrow[d]
  & & H^0(G(n))\otimes H^0(\mathcal O_X(1))
      \arrow[r] \arrow[d]
  & 0 \\
  0 \arrow[r]
  & H^0(F(n)) \arrow[ur] \arrow[r]
  & H^0(F(n+1)) \arrow[rr]
  & & H^0(G(n+1)) \arrow[d] \\
  & & & & 0
\end{tikzcd}
\]
Here the oblique arrow is defined by
\(\varphi\mapsto\varphi\otimes\sigma\); the horizontal rows are exact, the
first by the \(n\)-regularity of \(F\) and by (i), already proved, and the
second column is exact by the induction hypothesis.

Finally, (iii) follows from the commutative diagram
\[
\begin{tikzcd}[column sep=large,row sep=large]
H^0(F(n))_X\otimes H^0(\mathcal O_X(1))_X
  \arrow[r,"\alpha_n"]
  \arrow[d]
& H^0(F(n)(1))_X
  \arrow[d,"\beta_{n+1}"] \\
{[H^0(F(n))_X](1)}
  \arrow[r,"\beta_n\otimes\mathrm{id}"]
& F(n)(1).
\end{tikzcd}
\]

\newpage
Here, if \(M\) is a vector space over \(k\), we write \(M_X\) for
\(M\otimes_k\mathcal O_X\). For \(n\geq m\), \(\alpha_n\) is surjective
by (ii); therefore, if \(\beta_{n+1}\) is surjective, so is \(\beta_n\). By
Serre's theorem, \(\beta_{n+1}\) is surjective for \(n\gg0\); this proves
(iii).

\sgaStableTarget{sga6.e13.prop.1-4}\textbf{Proposition 1.4.} Let
\[
  0\longrightarrow F(-1)\longrightarrow F\longrightarrow G\longrightarrow0
\]
be an exact sequence on \(X\). If \(G\) is \(m\)-regular, then:
\begin{enumerate}[label=(\roman*)]
\item \(H^q(F(n))=0\) for \(q\geq2\) and \(n\geq m-q\).
\item \(h^1(F(n-1))\geq h^1(F(n))\) for \(n\geq m-1\).
\item \(H^1(F(n))=0\) for
\(n\geq(m-1)+h^1(F(m-1))\).
\end{enumerate}
In particular, \(F\) is \([m+h^1(F(m-1))]\)-regular.

Indeed, in the exact sequence
\[
  H^{q-1}(G(n))\longrightarrow H^q(F(n-1))
  \longrightarrow H^q(F(n))\longrightarrow H^q(G(n)),
\]
the two outer groups are zero for \(q\geq2\) and
\(n\geq m-(q-1)\), by \hyperlink{sga6.e13.prop.1-3}{1.3(i)}. Thus
\[
  H^q(F(m-q))\xrightarrow{\sim}H^q(F(m-q+1))
  \xrightarrow{\sim}H^q(F(m-q+2))\xrightarrow{\sim}\cdots;
\]
but, by Serre's theorem, \(H^q(F(n))=0\) for \(n\gg0\), proving (i).

For \(n\geq m-1\), consider the exact sequence
\[
  0\to H^0(F(n-1))\to H^0(F(n))
  \xrightarrow{\alpha_n}H^0(G(n))\to H^1(F(n-1))
  \to H^1(F(n))\to0.
\]
It implies (ii). Now suppose that \(\alpha_n\) is surjective, and consider
the commutative diagram
\[
\begin{tikzcd}[column sep=large,row sep=large]
H^0(F(n))\otimes H^0(\mathcal O_X(1))
  \arrow[r,"\alpha_n\otimes\mathrm{id}"]
  \arrow[d]
& H^0(G(n))\otimes H^0(\mathcal O_X(1))
  \arrow[r]
  \arrow[d,"\beta_n"]
& 0 \\
H^0(F(n+1))
  \arrow[r,"\alpha_{n+1}"]
& H^0(G(n+1)).
\end{tikzcd}
\]

\newpage
If \(n\geq m\), then \(\beta_n\) is surjective by \hyperlink{sga6.e13.prop.1-3}{1.3(ii)}, and hence
\(\alpha_{n+1}\) is surjective. Consequently,
\[
  H^1(F(n-1))\xrightarrow{\sim}H^1(F(n))
  \xrightarrow{\sim}\cdots\xrightarrow{\sim}0.
\]
On the other hand, if \(\alpha_n\) is not surjective, then
\(h^1(F(n-1))>h^1(F(n))\). Thus the function
\(n\longmapsto h^1(F(n))\) decreases \emph{strictly}, starting from
\(h^1(F(m-1))\), until it becomes zero; this proves (iii).

\sgaStableTarget{sga6.e13.def.1-5}\textbf{Definition 1.5.} Let \(F\) be a coherent sheaf on \(X\), let
\(r\geq\dim\Supp(F)\), and let \((b)=(b_0,\ldots,b_r)\) be a sequence of
\(r+1\) integers. Then \(F\) is said to be a \emph{\((b)\)-sheaf} if it
satisfies either of the following two equivalent sets of conditions.

\noindent\sgaStableTarget{sga6.e13.para.1-5-1}\textbf{(1.5.1)}
\begin{enumerate}[label=(\roman*)]
\item At every point of the support of \(F\), \(\mathcal O_X(1)\) is
generated by \(H^0(\mathcal O_X(1))\).
\item \(h^0(F(-1))\leq b_0\).
\item If \(r\geq1\), there is a section
\(\sigma\in H^0(\mathcal O_X(1))\) which defines an exact sequence
\[
  0\longrightarrow F(-1)\longrightarrow F\longrightarrow G\longrightarrow0
\]
such that \(G\) is a \((b_1,\ldots,b_r)\)-sheaf.
\end{enumerate}

\noindent\sgaStableTarget{sga6.e13.para.1-5-2}\textbf{(1.5.2)} There is an \(F\)-regular sequence of \(r\)
sections \(\sigma_1,\ldots,\sigma_r\in H^0(\mathcal O_X(1))\) such that,
for \(i=0,\ldots,r\), if \(F_i\) denotes the restriction of \(F\) to the
zero scheme of \(\sigma_1,\ldots,\sigma_i\), then
\[
  h^0(F_i(-1))\leq b_i.
\]

\sgaStableTarget{sga6.e13.prop.1-6}\textbf{Proposition 1.6.} Let \(F\) be a \((b)\)-sheaf on \(X\). Then:
\begin{enumerate}[label=(\roman*)]
\item For every sufficiently general sequence \(\sigma\) of \(r\) sections
\(\sigma_1,\ldots,\sigma_r\in H^0(\mathcal O_X(1))\), if
\(F_{\sigma,i}\) denotes the restriction of \(F\) to the zero scheme of
\(\sigma_1,\ldots,\sigma_i\), then
\(h^0(F_{\sigma,i}(-1))\leq b_i\), for \(0\leq i\leq r\).
\item Every subsheaf \(G\) of \(F\) is a \((b)\)-sheaf.
\end{enumerate}

\newpage
Indeed, let \(S'=\mathbb A_k^N\) be the affine space whose \(k\)-points
correspond to the sequences \(\sigma\), and let \(T\) be the open subset of
\(S'\) corresponding to the \(F\)-regular sequences. For fixed \(i\), the
sheaves \(F_{\sigma,i}(-1)\) corresponding to the \(k\)-points of \(T\) are
contained in a flat family over \(T\), and \(T\) is nonempty by hypothesis.
Thus (i) follows from upper semicontinuity of the function
\[
  \sigma\longmapsto h^0(F_{\sigma,i}(-1))
\]
(EGA III 7.7.5). Finally, since every sufficiently general sequence is
\(G\)-regular and is such that \(G_{\sigma,i}\) is a subsheaf of
\(F_{\sigma,i}\), for every \(i\), (ii) follows from (i).

\sgaStableTarget{sga6.e13.lem.1-7}\textbf{Lemma 1.7.} Let
\[
  0\longrightarrow F(-1)\longrightarrow F\longrightarrow G\longrightarrow0
\]
be an exact sequence on \(X\). If the Hilbert polynomial of \(F\) is written
in the form
\[
  \chi(F(n))=\sum_{i=0}^{r}a_i\binom{n+i}{i},
\]
then
\[
  \chi(G(n))=\sum_{i=0}^{r-1}a_{i+1}\binom{n+i}{i}.
\]
Indeed,
\[
\begin{aligned}
  \chi(G(n))
  &=\chi(F(n))-\chi(F(n-1))\\
  &=\sum_{i=0}^{r}a_i
    \left[\binom{n+i}{i}-\binom{n-1+i}{i}\right]\\
  &=\sum_{i=1}^{r}a_i\binom{n+i-1}{i-1},
\end{aligned}
\]
which proves the assertion.

\sgaStableTarget{sga6.e13.prop.1-8}\textbf{Proposition 1.8.} Let \(F\) be a \((b)\)-sheaf on \(X\), let
\(s=\dim\Supp(F)\), and let
\[
  \chi(F(n))=\sum_{i=0}^{s}a_i\binom{n+i}{i}
\]
be the Hilbert polynomial of \(F\). Then:
\begin{enumerate}[label=(\roman*)]
\item For \(n\geq-1\),
\[
  h^0(F(n))\leq\sum_{i=0}^{s}b_i\binom{n+i}{i}.
\]
\item \(a_s\leq b_s\), and \(F\) is also a
\((b_0,\ldots,b_{s-1},a_s)\)-sheaf.
\end{enumerate}
Indeed, we argue by induction on \(s\). If \(s=0\), the assertion is clear,
since \(a_0=h^0(F)=h^0(F(-1))\leq b_0\). If \(s\geq1\), we have an exact
sequence
\[
  0\longrightarrow F(-1)\longrightarrow F\longrightarrow G\longrightarrow0
\]

\newpage
with \(G\) a \((b_1,\ldots,b_r)\)-sheaf and
\(\dim\Supp(G)=s-1\). Then
\[
  h^0(F(n))-h^0(F(n-1))\leq h^0(G(n)),
\]
whereas
\[
  h^0(G(n))\leq
  \sum_{i=0}^{s-1}b_{i+1}\binom{n+i}{i}
\]
by induction on \(s\), and \(h^0(F(-1))\leq b_0\). Assertion (i) follows
by induction on \(n\geq-1\). Finally, by \hyperlink{sga6.e13.lem.1-7}{1.7} and induction on \(s\),
\(a_s\leq b_s\), and \(G\) is a
\((b_1,\ldots,b_{s-1},a_s)\)-sheaf; this proves (ii).

\sgaStableTarget{sga6.e13.def.1-9}\textbf{Definition 1.9.} The following polynomials with nonnegative rational
coefficients, defined recursively, are called the \emph{\((b)\)-polynomials}:
\[
\begin{cases}
P_{-1}=0,\\
P_r(X_0,\ldots,X_r)=P_{r-1}(X_1,\ldots,X_r)
+\displaystyle\sum_{i=0}^{r}X_i
\binom{P_{r-1}(X_1,\ldots,X_r)-1+i}{i}.
\end{cases}
\]

\sgaStableTarget{sga6.e13.rem.1-10}\textbf{Remark 1.10.} For the \((b)\)-polynomials, induction on \(r\geq t\)
immediately gives
\[
  P_r(X_0,\ldots,X_t,0,\ldots,0)=P_t(X_0,\ldots,X_t).
\]

\sgaStableTarget{sga6.e13.thm.1-11}\textbf{Theorem 1.11.} \emph{(The main theorem on \((b)\)-sheaves.)} Let
\(F\) be a \((b)\)-sheaf on \(X\), and let
\[
  \chi(F(n))=\sum_{i=0}^{r}a_i\binom{n+i}{i}
\]
be the Hilbert polynomial of \(F\). Let
\((c)=(c_0,\ldots,c_r)\) be a sequence of integers satisfying
\((c)\geq(b)-(a)\), and let \(m=P_r(c)\) be the integer given by the
\(r\)-th \((b)\)-polynomial \hyperlink{sga6.e13.def.1-9}{(1.9)}. Then \(m\geq0\), and \(F\) is
\(m\)-regular. In particular, if \(s=\dim\Supp(F)\), then \(F\) is
\[
  P_{s-1}(c_0,\ldots,c_{s-1})\text{-regular}.
\]

\newpage
Indeed, we argue by induction on \(r\). If \(r=0\), the assertion is clear,
since \(m=c_0\geq0\) and \(F\) is \(n\)-regular for every \(n\). If \(r\geq1\),
there is an exact sequence
\[
  0\longrightarrow F(-1)\longrightarrow F\longrightarrow G\longrightarrow0
\]
in which \(G\) is a \((b_1,\ldots,b_r)\)-sheaf. By \hyperlink{sga6.e13.lem.1-7}{1.7} and the induction
hypothesis, if \(n=P_{r-1}(c_1,\ldots,c_r)\), then \(n\geq0\) and \(G\) is
\(n\)-regular. Hence, by \hyperlink{sga6.e13.prop.1-4}{1.4}, \(F\) is
\([n+h^1(F(n-1))]\)-regular, and \(h^q(F(n-1))=0\) for \(q\geq2\).
Therefore
\[
  h^1(F(n-1))=h^0(F(n-1))-\chi(F(n-1)),
\]
which, by \hyperlink{sga6.e13.prop.1-8}{1.8(i)}, is at most
\[
  \sum_{i=0}^{r}(b_i-a_i)\binom{n-1+i}{i},
\]
since \(b_i\geq0\).
Since \(c_i\geq b_i-a_i\), this quantity is at most
\[
  \sum_{i=0}^{r}c_i\binom{n-1+i}{i}.
\]
Finally, by \hyperlink{sga6.e13.prop.1-3}{1.3(i)}, \(F\) is
\[
  \left[n+\sum_{i=0}^{r}c_i\binom{n-1+i}{i}\right]\text{-regular}.
\]
This proves the first assertion; the second follows from \hyperlink{sga6.e13.prop.1-8}{1.8(ii)} and 1.10.

\sgaStableTarget{sga6.e13.para.1-12}\textbf{1.12.} Let \(S\) be a noetherian scheme, and let \(X\) be an
\(S\)-scheme of finite type. Let \(\mathcal F\) be a family of classes of
coherent sheaves on the fibers of \(X/S\). In other words, for every point
\(s\in S\) and every extension \(K\) of \(k(s)\), coherent sheaves \(F_K\)
are specified on the algebraic scheme \(X_K\); and \(F_K\) and
\(F'_{K'}\) are declared to lie in the same class if there are
\(k(s)\)-homomorphisms from \(K\) and \(K'\) into a common extension
\(K''\) of \(k(s)\) such that
\[
  F_{K''}=F_K\otimes_KK''
  \quad\text{and}\quad
  F'_{K''}=F'_{K'}\otimes_{K'}K''
\]
are isomorphic on \(X_{K''}\). The family \(\mathcal F\) is said to be
\emph{bounded by the coherent sheaf} \(F\) \emph{on}
\(X_T=X\times_ST\), where \(T\) is of finite type over \(S\), if
\(\mathcal F\) is contained in the family of classes of the sheaves
\(F_{k(t)}\), with \(t\in T\). The family \(\mathcal F\) is said to be
\emph{bounded} if such \(T\) and \(F\) exist.

Suppose in addition that \(X/S\) is projective and endowed with a
(relatively) ample invertible sheaf \(\mathcal O_X(1)\). Then
\(\mathcal F\) is called a \((b)\)-\emph{family}, for a sequence of integers
\((b)=(b_0,\ldots,b_r)\), if every class in \(\mathcal F\) is represented
by an \(F_K\), with \(K\) algebraically closed, which is a
\((b)\)-sheaf.

\sgaStableTarget{sga6.e13.thm.1-13}\textbf{Theorem 1.13.} \emph{Let \(S\) be a noetherian scheme, and let
\(X\) be a projective \(S\)-scheme endowed with an ample sheaf
\(\mathcal O_X(1)\) such that the induced sheaves \(\mathcal O_{X_s}(1)\),
for \(s\in S\), are generated by their \(H^0(\mathcal O_{X_s}(1))\). Let
\(\mathcal F\) be a family of classes of coherent sheaves on the fibers of
\(X/S\). Then the following conditions are equivalent:}
\begin{enumerate}[label=(\roman*)]
\item \emph{\(\mathcal F\) is bounded. Moreover, if all the
\(F_K\in\mathcal F\) are locally free of rank \(p\), then \(\mathcal F\)
may be assumed to be bounded by a sheaf \(F\) on some \(X_T\), with \(F\)
locally free of rank \(p\).}
\item \emph{The set of Hilbert polynomials \(\chi(F_K(n))\), for
\(F_K\in\mathcal F\), is finite, and there exists a sequence of integers
\((b)\) such that \(\mathcal F\) is a \((b)\)-family.}
\item \emph{The set of Hilbert polynomials \(\chi(F_K(n))\), for
\(F_K\in\mathcal F\), is finite, and there exists an integer \(m\) such
that all the \(F_K\in\mathcal F\) are \(m\)-regular.}
\item \emph{The set of Hilbert polynomials \(\chi(F_K(n))\), for
\(F_K\in\mathcal F\), is finite, and \(\mathcal F\) is contained in the
family of classes of quotients of sheaves of the form \(E_K\), where \(E\)
is a coherent sheaf on some \(X_T\). Moreover, one may assume that \(T=S\)
and that \(E\) has the form \(\mathcal O_X(-m)^{\oplus M}\).}
\item \emph{\(\mathcal F\) is contained in the family of classes of
cokernels of homomorphisms of the form \(E'_K\to E_K\), where \(E'\) and
\(E\) are coherent sheaves on some \(X_T\). Moreover, one may assume that
\(T=S\) and that \(E'\) and \(E\) have the forms
\(\mathcal O_X(-m')^{\oplus M'}\) and \(\mathcal O_X(-m)^{\oplus M}\),
respectively.}
\end{enumerate}

Indeed, suppose that \(\mathcal F\) is bounded by a sheaf \(F\) on some
\(X_T\). After partitioning \(T\) into pieces and applying the generic
flatness theorem (EGA IV 6.9.1), one may assume that \(F\) is flat over
\(T\). Then,
the number of polynomials \(\chi(F_t(n))\) is at most the number of connected
components of \(T\) (EGA III 7.9.5). It is also an elementary lemma that,
if \(t\in T\) and \(F_t\) is locally free of rank \(p\), then \(F\) is
locally free of rank \(p\) over an open neighborhood of \(t\) in \(T\)
(see EGA IV 11.3.10, with \(X=Y\) and \(f=\mathrm{id}\)). The hypotheses
and EGA III 7.7.6 show, moreover, that after partitioning \(T\) still more
finely one can find a sequence \(\sigma\) of sections
\[
\sigma_1,\ldots,\sigma_r\in H^0(\mathcal O_X(1))
\]
that is \(F\)-regular. Thus (ii) follows from the upper semicontinuity of
the function
\[
t\longmapsto h^0(F_{t,i}(-1)),
\]
where, for \(t\in T\) and \(0\le i\le r\), \(F_{t,i}\) denotes the
restriction of \(F_t\) to the zero scheme of
\(\sigma_1,\ldots,\sigma_i\) (EGA III 7.7.5).

The implication (ii) \(\Rightarrow\) (iii) follows immediately from 1.11.
The implication (iii) \(\Rightarrow\) (iv) follows from \hyperlink{sga6.e13.prop.1-3}{1.3 (iii)}, on
taking
\[
M=\max\bigl(\chi(F_K(m))\bigr),
\qquad E=\mathcal O_X(-m)^{\oplus M}.
\]

Suppose that \(\mathcal F\) satisfies (iv). There are exact sequences
\[
0\longrightarrow F'_K\longrightarrow E_K\longrightarrow F_K
\longrightarrow0,
\]
and the set of polynomials \(\chi(F_K(n))\) is finite. By hypothesis, the
family of classes of the \(E_K\) is bounded. Hence, by the already proved
implication (i) \(\Rightarrow\) (ii), the set of polynomials
\(\chi(E_K(n))\) is finite, and there exists a sequence of integers \((b)\)
such that all the \(E_K\), with \(K\) algebraically closed, are
\((b)\)-sheaves. It follows that the set of polynomials
\(\chi(F'_K(n))\) is finite and, by \hyperlink{sga6.e13.prop.1-6}{1.6}, all the \(F'_K\) are
\((b)\)-sheaves. Applying the already proved implication (ii)
\(\Rightarrow\) (iv) to the family of classes of the \(F'_K\), one obtains
(v).

Finally, suppose that \(\mathcal F\) satisfies (v), and let us prove that
\(\mathcal F\) is bounded. One may assume that \(E\) (respectively \(E'\))
has the form \(\mathcal O_X(-m)^{\oplus M}\). Indeed,
by (i) \(\Rightarrow\) (iv), take a surjection
\[
L=\mathcal O_{X_T}(-m)^{\oplus M}\longrightarrow E.
\]
After partitioning \(T\) into pieces so that \(E\) becomes flat over \(T\),
formation of the exact sequence
\[
0\longrightarrow I\xrightarrow{u}L\longrightarrow E\longrightarrow0
\]
commutes with passage to the fibers. By (i) \(\Rightarrow\) (iii) and
(i) \(\Rightarrow\) (iv), choose \(m_1\) sufficiently large that all the
groups
\[
\operatorname{Ext}^1\bigl(\mathcal O_{X_K}(-m_1),I_K\bigr)
=H^1(I_K(m_1))
\]
vanish and that there exist surjections
\[
L_1=\mathcal O_{X_T}(-m_1)^{\oplus M_1}\longrightarrow E',
\qquad
\alpha:L_2=\mathcal O_{X_T}(-m_1)^{\oplus M_2}\longrightarrow I.
\]
The maps
\[
\operatorname{Hom}(L_{1,K},L_K)\longrightarrow
\operatorname{Hom}(L_{1,K},E_K)
\]
are then surjective. Let \(\beta:E'_K\to E_K\) be a homomorphism, and let
\(\gamma\) be the composite
\[
L_{1,K}\longrightarrow E'_K\xrightarrow{\beta}E_K.
\]
Then \(\gamma\) comes from a homomorphism
\(\delta:L_{1,K}\to L_K\), and
\[
(\delta,u_K\alpha):L_{1,K}\oplus L_{2,K}\longrightarrow L_K
\]
and \(\beta:E'_K\to E_K\) have the same cokernel. Consequently,
\(\mathcal F\) is contained in the family of classes of cokernels of
homomorphisms of the form
\[
\mathcal O_{X_K}(-m_1)^{\oplus(M_1+M_2)}
\longrightarrow
\mathcal O_{X_K}(-m)^{\oplus M}.
\]

Still after partitioning \(T\) into pieces, one may assume in addition
that \(X_T\) (and hence
\(\mathcal H\!om_{\mathcal O_{X_T}}(E',E)\)) and all the
\[
R^q(f_T)_*\mathcal H\!om_{\mathcal O_{X_T}}(E',E),
\]
where \(f_T:X_T\to T\) is the structural morphism, are flat over \(T\),
so that formation of
\[
G=(f_T)_*\mathcal H\!om_{\mathcal O_{X_T}}(E',E)
\]
commutes with every base change \(T'\to T\) (EGA III 6.9.9.2). Let
\[
R=\mathbb V(G^*).
\]
Then \(X_R\) canonically carries a \emph{universal} exact sequence
\[
E'_R\longrightarrow E_R\longrightarrow F\longrightarrow0,
\]
and \(F\) bounds \(\mathcal F\).

\sgaStableTarget{sga6.e13.sec.2}\subsection*{2. Several technical lemmas}

\sgaStableTarget{sga6.e13.para.2-1}\textbf{2.1.} Let \(k\) be a field and \(X\) a proper \(k\)-scheme. Let
\(L_1,\ldots,L_s\) be invertible sheaves on \(X\), let \(c_1(L_i)\) be the
Chern class of \(L_i\), and let \(Y\) be a closed subscheme of \(X\) of
dimension \(\le s\). Recall that the notation
\[
\bigl\langle c_1(L_1)\cdots c_1(L_s)\cdot Y\bigr\rangle
\]
denotes the intersection number, which may be calculated as the
coefficient of the monomial \(n_1\cdots n_s\) in the Snapper polynomial
\[
\chi\!\left(
L_1^{\otimes n_1}\otimes\cdots\otimes L_s^{\otimes n_s}
\otimes\mathcal O_Y
\right)
\]
(\hyperlink{sga6.e10.prop.4-3-1}{X 4.3.1}).

\sgaStableTarget{sga6.e13.para.2-2}\textbf{2.2.} From now on, suppose that \(k\) is algebraically closed,
that \(X\) is projective and endowed with a very ample invertible sheaf
\(H=\mathcal O_X(1)\), and that \(X\) is integral of dimension \(r\).

Let \(L\) be an invertible sheaf on \(X\). The integer
\[
\bigl\langle c_1(L)\cdot c_1(H)^{r-1}\bigr\rangle
\]
is called the \emph{degree} of \(L\) (relative to \(H\)) and is denoted
\(d(L)\).

In what follows we consider an invertible sheaf \(L\) and the Hilbert
polynomials
\[
\chi(L(n))=\sum_{i=0}^{r}a_i\binom{n+i}{i},
\qquad
\chi(\mathcal O_X(n))=\sum_{i=0}^{r}e_i\binom{n+i}{i}.
\]

\sgaStableTarget{sga6.e13.lem.2-3}\textbf{Lemma 2.3.} \emph{One has:}
\begin{enumerate}[label=(\roman*)]
\item \emph{\(a_r=e_r=d(H)\).}
\item \emph{If \(r\ge1\), then \(a_{r-1}-e_{r-1}=d(L)\).}
\end{enumerate}

Indeed, (i) and (ii) follow immediately by setting \(m=0,1\) in the
following polynomial:
\[
\chi\bigl(L^{\otimes m}\otimes H^{\otimes n}\bigr)
=\sum_{i+j\le r}f_{ij}
\binom{m+i}{i}\binom{n+j}{j}.
\]

\sgaStableTarget{sga6.e13.lem.2-4}\textbf{Lemma 2.4.} \emph{If \(d(L)\le0\) and \(h^0(L)\ne0\), then
\(L\simeq\mathcal O_X\); hence \(d(L)=0\) and \(h^0(L)=1\).}

Indeed, since \(X\) is integral, there is a divisor \(Y\ge0\) such that
\(L\simeq\mathcal O_X(Y)\). If \(Y>0\), one would have
\[
0\ge
\bigl\langle c_1(L)\cdot c_1(H)^{r-1}\bigr\rangle
=\bigl\langle c_1(H)^{r-1}\cdot Y\bigr\rangle
\ge1.
\]

\sgaStableTarget{sga6.e13.lem.2-5}\textbf{Lemma 2.5.} \emph{If \(d(L)\le0\), then \(L\) is a
\((0,\ldots,0,a_r)\)-sheaf (with \(r-1\) zeros).}

Indeed, if \(r=0\), then \(L\) is plainly an \(a_0\)-sheaf. If
\(r\ge1\), then by \hyperlink{sga6.e13.lem.2-4}{2.4} one has \(h^0(L(-1))=0\), since
\(d(L(-1))<0\). Let \(Y\) be a hyperplane section of \(X\), and let
\[
0\longrightarrow L(-1)\longrightarrow L\longrightarrow L_Y
\longrightarrow0
\]
be an exact sequence defined by \(Y\). If \(r=1\), then by \hyperlink{sga6.e13.lem.1-7}{1.7} one has
\(h^0(L_Y(-1))=a_1\); hence \(L\) is a \((0,a_1)\)-sheaf. If \(r\ge2\),
one may take \(Y\) integral by Bertini's theorem and may thus assume
inductively that \(L_Y\) is a \((0,\ldots,0,a_r)\)-sheaf; the assertion
follows.

\sgaStableTarget{sga6.e13.lem.2-6}\textbf{Lemma 2.6.} \emph{If \(r=2\) and \(d(L)\le0\), then:}
\begin{enumerate}[label=(\roman*)]
\item \emph{For every \(n\),
\[
h^1(L(n))\le a_1(a_1+1)a_2-a_0.
\]}
\item \emph{\(a_0\le a_1(a_1+1)\).}
\end{enumerate}

Let \(Y\) be an integral hyperplane section of \(X\). By \hyperlink{sga6.e13.lem.1-7}{1.7} and \hyperlink{sga6.e13.lem.2-6}{2.6},
\(L_Y\) is a \((0,a_2)\)-sheaf with Hilbert polynomial
\[
\chi(L_Y(n))=a_1+a_2\binom{n+1}{1}.
\]
Thus, by \hyperlink{sga6.e13.thm.1-11}{1.11}, \(-a_1\ge0\), and \(L_Y\) is
\([-a_1]\)-regular.

Consider the following exact sequence:
\[
H^0(L_Y(n))\longrightarrow H^1(L(n-1))
\longrightarrow H^1(L(n))\longrightarrow H^1(L_Y(n)).
\]
For \(n\ge-a_1-1\), one has
\(h^1(L(-a_1-2))\ge h^1(L(n))\), since \(h^1(L_Y(n))=0\) by \hyperlink{sga6.e13.prop.1-3}{1.3 (i)}.
For \(n\le-1\), one has
\(h^1(L(n-1))\le h^1(L(-1))\), since \(h^0(L_Y(n))=0\) by 2.4.
Finally, for every \(n\),
\(h^0(L_Y(n-1))\le h^0(L_Y(n))\), since there is an injection
\(L_Y(n-1)\to L_Y(n)\), and
\[
h^0(L_Y(-a_1-1))=a_2\binom{-a_1}{1}+a_1,
\]
since \(h^1(L_Y(-a_1-1))=0\). Hence, for \(n\le-a_1-2\),
\[
h^1(L(n))-h^1(L(n+1))\le -a_2a_1+a_1.
\]
It follows that, for every \(n\),
\sgaStableTarget{sga6.e13.eq.2-6-1}\begin{equation}
h^1(L(n))\le
h^1(L(-a_1-2))+(-a_1-1)(-a_2a_1+a_1).
\tag{2.6.1}
\end{equation}

Since \(L_Y\) is \([-a_1]\)-regular, one has
\(h^2(L(-a_1-2))=0\) by \hyperlink{sga6.e13.prop.1-4}{1.4}; hence
\[
h^1(L(-a_1-2))
=h^0(L(-a_1-2))-\chi(L(-a_1-2)).
\]
But
\[
h^0(L(-a_1-2))\le a_2\binom{-a_1}{2}
\]
by \hyperlink{sga6.e13.prop.1-8}{1.8 (i)}, since \(L\) is a \((0,0,a_2)\)-sheaf by 2.5. Therefore
\sgaStableTarget{sga6.e13.eq.2-6-2}\begin{equation}
0\le h^1(L(-a_1-2))\le -a_1(-a_1-1)-a_0.
\tag{2.6.2}
\end{equation}
Finally, \hyperlink{sga6.e13.eq.2-6-2}{2.6.2} gives (ii), and \hyperlink{sga6.e13.eq.2-6-1}{2.6.1} together with \hyperlink{sga6.e13.eq.2-6-2}{2.6.2} gives (i).

\sgaStableTarget{sga6.e13.lem.2-8}\textbf{Lemma 2.8.} \emph{Suppose that \(r=2\) and \(d(L)\le0\). Put
\(\delta=d(L)\) and \(\beta=\delta(e_2-1)+e_1\). Then:}

\[
a_0\ge
\bigl\langle c_1(L)^2\bigr\rangle+2e_0
+\sum_{i=1}^{2}(2e_i-a_i)\binom{\delta-1+i}{i}
+\beta(\beta+1).
\]

Indeed, write
\[
\chi\bigl(L^{\otimes m}(\delta-1)\bigr)
=\alpha_2m^2+\alpha_1m+\alpha_0.
\]
Comparison of polynomials gives
\[
2\alpha_2=\bigl\langle c_1(L)^2\bigr\rangle,
\qquad
\alpha_0=\sum_{i=0}^{2}e_i\binom{\delta-1+i}{i},
\]
and
\[
\alpha_2+\alpha_1+\alpha_0
=\sum_{i=0}^{2}a_i\binom{\delta-1+i}{i}.
\]
Write
\[
\chi\bigl(L^{-1}(\delta)(n)\bigr)
=\sum_{i=0}^{2}\beta_i\binom{n+i}{i}.
\]
Then \(\beta_0=\alpha_2-\alpha_1+\alpha_0\), and by \hyperlink{sga6.e13.lem.2-3}{2.3},
\[
\beta_1=d\bigl(L^{-1}(\delta)\bigr)+e_1=\beta.
\]
Thus
\[
a_0=2\alpha_2+2\alpha_0
-\sum_{i=1}^{2}a_i\binom{\delta-1+i}{i}-\beta_0.
\]
Since \(d(L^{-1}(\delta))\le0\), one has
\(\beta_0\le\beta(\beta+1)\) by \hyperlink{sga6.e13.lem.2-6}{2.6 (ii)}; whence the assertion.

\sgaStableTarget{sga6.e13.lem.2-9}\textbf{Lemma 2.9.} \emph{Suppose that \(r\ge2\), that \(d(L)\le0\),
and that \(X\) is normal. Then:}
\begin{enumerate}[label=(\roman*)]
\item \emph{\(h^1(L(-n))=0\) for
\(n\ge2+h_2(a_{r-2},a_{r-1},a_r)\);}
\item \emph{\(h^1(L(n))\le h_r(a_{r-2},a_{r-1},a_r)\) for every
\(n\);}
\item \emph{\(|a_i|\le g_{r-i}(a_{r-2},a_{r-1},a_r)\) for
\(i=0,\ldots,r\), where
\(h_2(X_0,X_1,X_2),h_3(X_0,X_1,X_2),\ldots\) and
\(g_0(X_0,X_1,X_2),g_1(X_0,X_1,X_2),\ldots\) are universal
polynomials with coefficients in \(\mathbb Q\), independent of both
\(k\) and \(X\).}
\end{enumerate}

Indeed, first consider the case \(r=2\). Put
\[
g_i(X_0,X_1,X_2)=X_{2-i}^2\quad(i=0,1,2),
\qquad
h_2(X_0,X_1,X_2)=X_1(X_1+1)X_2-X_0.
\]
Then (iii) is clear, and (ii) is \hyperlink{sga6.e13.lem.2-6}{2.6 (i)}.

Let \(X\subset P=\mathbb P_k^N\), and put
\[
\omega=
\mathcal E\!xt_{\mathcal O_P}^{N-2}
\bigl(\mathcal O_X,\Omega^N_{P/k}\bigr).
\]
Then
\[
h^1(L(-n))=h^1(\omega\otimes L^{-1}(n)),
\]
and \(\omega\) is invertible outside the finite set \(Z\) of singular
points of \(X\) (see FGA 149 \S~6). Let \(Y\) be a normal hyperplane
section of \(X\) not containing \(Z\). Then
\(\omega_Y(1)\simeq\Omega^1_{Y/k}\), and there is an exact sequence
\[
0\longrightarrow\omega\otimes L^{-1}(-1)
\longrightarrow\omega\otimes L^{-1}
\longrightarrow\Omega^1_{Y/k}\otimes L_Y^{-1}(-1)
\longrightarrow0.
\]
Now
\[
h^1\bigl(\Omega^1_{Y/k}\otimes L_Y^{-1}(n)\bigr)
=h^0(L_Y(-n))
\]
for every \(n\), and \(h^0(L_Y(-n))=0\) for \(n\ge1\) by 2.5. Thus
\(\Omega^1_{Y/k}\otimes L_Y^{-1}(-1)\) is 3-regular and, by \hyperlink{sga6.e13.prop.1-4}{1.4},
\(\omega\otimes L^{-1}\) is
\[
\bigl[3+h^1(\omega\otimes L^{-1}(2))\bigr]\text{-regular}.
\]
Consequently, \(h^1(L(-n))=0\) for
\(n\ge2+h^1(L(-2))\); this proves (i).

Now suppose that \(r\ge3\), and argue by induction on \(r\). Let \(Y\)
be a normal hyperplane section of \(X\), and consider the exact sequence
\sgaStableTarget{sga6.e13.eq.2-9-1}\begin{equation}
H^0(L_Y(-n))\longrightarrow H^1(L(-n-1))
\longrightarrow H^1(L(-n))\longrightarrow H^1(L_Y(-n)).
\tag{2.9.1}
\end{equation}
By \hyperlink{sga6.e13.lem.1-7}{1.7},
\[
\chi(L_Y(n))=
\sum_{i=0}^{r-1}a_{i+1}\binom{n+i}{i}.
\]
Thus, for \(n\ge2+h_2(a_{r-2},a_{r-1},a_r)\), the induction
hypothesis gives \(h^1(L_Y(-n))=0\); and since \(n\ge2\) by (ii) for
\(r=2\), \hyperlink{sga6.e13.lem.2-4}{2.4} gives \(h^0(L_Y(-n))=0\). Hence
\[
h^1(L(-n-1))=h^1(L(-n)).
\]
But for \(n\gg0\), one has \(h^1(L(-n))=0\) (see FGA 149 \S~6); hence
(i).

As induction hypothesis, suppose that the polynomials \(g_i\) and \(h_i\)
have been defined for \(i\le r-1\). For every \(n\), \hyperlink{sga6.e13.eq.2-9-1}{2.9.1} gives
\[
h^1(L(n))\le h^1(L(n-1))
+h_{r-1}(a_{r-2},a_{r-1},a_r).
\]
Therefore, by (i)
already proved, it follows that for every \(m\ge-2\),
\sgaStableTarget{sga6.e13.eq.2-9-2}\begin{equation}
h^1(L(n))\le
\bigl[2+h_2(a_{r-2},a_{r-1},a_r)+m\bigr]
h_{r-1}(a_{r-2},a_{r-1},a_r)
\tag{2.9.2}
\end{equation}
for every \(n\le m\).

By \hyperlink{sga6.e13.lem.2-5}{2.5}, \(L_Y\) is a \((0,\ldots,0,a_r)\)-sheaf. Put
\[
c_i=g_{r-i}(a_{r-2},a_{r-1},a_r)
\quad(i=1,\ldots,r),
\qquad
m=P_{r-2}(c_1,\ldots,c_{r-1}).
\]
By \hyperlink{sga6.e13.thm.1-11}{1.11}, \(m\ge0\), and \(L_Y\) is \(m\)-regular. Hence, by \hyperlink{sga6.e13.prop.1-4}{1.4 (ii)},
\[
h^1(L(m-2))\ge h^1(L(n))
\qquad\text{for every }n\ge m-1.
\]
Thus, if polynomials are defined by
\[
M=P_{r-2}(g_{r-1},\ldots,g_1),
\qquad
h_r=(h_2+M)h_{r-1},
\]
then (ii) follows from 2.9.2.

Finally, since \(L_Y\) is \(m\)-regular, one has \(h^q(L(m))=0\) for
\(q\ge2\) by \hyperlink{sga6.e13.prop.1-4}{1.4 (i)}. Hence
\[
\sum_{i=0}^{r}a_i\binom{m+i}{i}
=h^0(L(m))-h^1(L(m)).
\]
But, by \hyperlink{sga6.e13.prop.1-8}{1.8 (i)},
\[
h^0(L(m))\le a_r\binom{m+r}{r}.
\]
Therefore, if one defines
\[
g_r=
\sum_{i=1}^{r}g_{r-i}\binom{M+i}{i}
+a_r\binom{M+r}{r}+h_r,
\]
then (iii) follows.

\sgaStableTarget{sga6.e13.rem.2-10}\textbf{Remark 2.10.} For fixed \(m\), consider an equality of
polynomials in \(n\):
\[
\sum_{i=0}^{r}a_{m,i}\binom{n+i}{i}
=\sum_{j=0}^{r}a_j\binom{m+n+j}{j}.
\]
One readily finds the formula
\[
a_{m,i}=
\sum_{j=0}^{r-i}a_{j+i}\binom{m-1+j}{j}
\]
by induction on \(r\), reducing \(a_{m,i}\) to \(a_{m,0}\), and then
setting \(n=-1\).

Similarly, one finds a polynomial with rational coefficients
\[
P_i^{(m,r)}(X_r,\ldots,X_i),
\]
increasing in \(X_i\), such that in the equality
\[
\sum_{i=0}^{r}a_i^{(m)}\binom{n+i}{i}
=\sum_{j=0}^{r}a_j\binom{mn+j}{j},
\]
one has
\[
a_i^{(m)}=P_i^{(m,r)}(a_r,\ldots,a_i).
\]

\sgaStableTarget{sga6.e13.lem.2-11}\textbf{Lemma 2.11.} \emph{Let \(S\) be a noetherian scheme and \(X\)
a projective \(S\)-scheme endowed with a very ample sheaf
\(H=\mathcal O_X(1)\). Suppose that the fibers of \(X/S\) are
geometrically integral and all have the same dimension \(r\). Let
\(\mathcal L\) be a family of classes of invertible sheaves \(L_K\) on the
fibers of \(X/S\). For \(\mathcal L\) to be bounded, it is necessary and
sufficient that all the coefficients \(a_i\) in the Hilbert polynomials}
\[
\chi(L_K(n))=\sum_{i=0}^{r}a_i\binom{n+i}{i}
\]
\emph{of the \(L_K\in\mathcal L\) remain bounded. Moreover, if the
fibers of \(X/S\) are also geometrically normal, then it is enough that
the three coefficients \(a_r,a_{r-1},a_{r-2}\) remain bounded.}

Indeed, let
\[
\delta=\sup\{d(L_K)\}.
\]
If \(\delta>0\), one may replace \(\mathcal L\) by the family of classes
of the \(L_K(-\delta)\), by \hyperlink{sga6.e13.lem.2-3}{2.3} and \hyperlink{sga6.e13.rem.2-10}{2.10}; thus one may assume
\(\delta\le0\). Then \(\mathcal L\) is a \((b)\)-family by \hyperlink{sga6.e13.lem.2-5}{2.5}, and is
therefore bounded if all the coefficients \(a_i\) remain bounded, by
1.13. If the fibers of \(X/S\) are geometrically normal, it is enough
for this purpose that the three coefficients \(a_r,a_{r-1},a_{r-2}\)
remain bounded, by 2.9.

\sgaStableTarget{sga6.e13.sec.3}\subsection*{3. General finiteness theorems}

\sgaStableTarget{sga6.e13.lem.3-1}\textbf{Lemma 3.1.} \emph{Let \(S\) be a scheme and \(X\) a
quasi-compact and quasi-separated \(S\)-scheme. Let
\(\underline{\Pic}_{X/S}\) be the Picard functor, localized for the fppf
topology.}\footnote{Cf. SGA 3 IV 6.3 and FGA no.~232 for the definition of
\(\underline{\Pic}_{X/S}\); here, as indicated, the fpqc topology should be
replaced by the fppf topology.}
\emph{Then, for every filtered inverse system \((Z_\alpha)\) of
quasi-compact and quasi-separated \(S\)-schemes with affine transition
morphisms, the canonical map}
\sgaStableTarget{sga6.e13.eq.3-1-1}\begin{equation}
\varinjlim_\alpha\underline{\Pic}_{X/S}(Z_\alpha)
\longrightarrow
\underline{\Pic}_{X/S}\!\left(\varprojlim_\alpha Z_\alpha\right)
\tag{3.1.1}
\end{equation}
\emph{is bijective.}

\emph{Consequently} (EGA IV 8.14.2), \emph{if the Picard functor is
represented by an \(S\)-scheme \(\Pic_{X/S}\), then \(\Pic_{X/S}\) is
locally of finite presentation.}

Indeed, an element \(\lambda_\alpha\) of
\(\underline{\Pic}_{X/S}(Z_\alpha)\) is represented by an invertible sheaf
\(L'_\alpha\) on \(X_{Z'_\alpha}\), where
\(Z'_\alpha\to Z_\alpha\) is an fppf morphism. Put
\[
Z=\varprojlim_\alpha Z_\alpha.
\]
The image \(\lambda\) of \(\lambda_\alpha\) in
\(\underline{\Pic}_{X/S}(Z)\) is represented by the invertible sheaf
\(L'\) on \(X_{Z'}\) obtained by inverse image from \(L'_\alpha\), where
\[
Z'=Z'_\alpha\times_{Z_\alpha}Z.
\]

If \(\lambda\) is the neutral element of
\(\underline{\Pic}_{X/S}(Z)\), there exists an fppf morphism
\(Z''\to Z'\) such that the inverse image of \(L'\) on \(X_{Z''}\) is
isomorphic to \(\mathcal O_{X_{Z''}}\). By EGA IV 8.8.2, 8.10.5 (vi),
11.2.6 (ii), 8.5.2 (i), and 8.5.2.4, there exist
\(\beta\ge\alpha\) and an fppf morphism
\[
Z''_\beta\longrightarrow Z'_\beta,
\qquad
Z'_\beta=Z'_\alpha\times_{Z_\alpha}Z_\beta,
\]
such that the inverse image of \(L'_\alpha\) on \(X_{Z''_\beta}\) is
isomorphic to \(\mathcal O_{X_{Z''_\beta}}\). The image of
\(\lambda_\alpha\) is therefore already the neutral element in
\(\underline{\Pic}_{X/S}(Z_\beta)\). Thus \hyperlink{sga6.e13.eq.3-1-1}{3.1.1} is injective.

The surjectivity of \hyperlink{sga6.e13.eq.3-1-1}{3.1.1} is proved similarly.

\sgaStableTarget{sga6.e13.prop.3-2}\textbf{Proposition 3.2.} \emph{Let \(S\) be a noetherian scheme and
\(X\) a proper \(S\)-scheme. Suppose that the Picard functor is
represented by an \(S\)-scheme \(\Pic_{X/S}\). Then:}
\begin{enumerate}[label=(\roman*)]
\item \emph{\(\Pic_{X/S}\) is locally of finite type over \(S\).}
\item \emph{Subsets \(\Lambda\) of \(\Pic_{X/S}\) correspond bijectively
to families \(\mathscr L\) of classes of invertible sheaves on the fibers
of \(X/S\).}
\item \emph{A subset \(\Lambda\) of \(\Pic_{X/S}\) is quasi-compact if
and only if the corresponding family \(\mathscr L\) is bounded.}
\end{enumerate}

Indeed, (i) follows from 3.1. Assertion (ii) follows from the following
fact: if \(s\) is a point of \(S\) and \(K\) is an algebraically closed
extension of \(k(s)\), then the \(K\)-points of \(\Pic_{X/S}\) correspond
bijectively to the isomorphism classes of invertible sheaves on \(X_K\).

For (iii), suppose that the subset \(\Lambda\) of \(\Pic_{X/S}\) is
quasi-compact. Then \(\Lambda\) is contained in an open subscheme \(U\)
of finite type, and the inclusion \(U\to\Pic_{X/S}\) corresponds to an
invertible sheaf \(L\) on some \(X_T\), where \(T\to U\) is a surjective
(flat) morphism of finite type. Finally \(T\to S\) is of finite type, so
\(L\) bounds \(\mathscr L\).

Conversely, suppose that \(\mathscr L\) is bounded by a coherent sheaf
\(L\) on some \(X_T\), with \(T\) of finite type over \(S\). One may
assume that \(L\) is invertible \hyperlink{sga6.e13.thm.1-13}{(1.13 (i))}, so that \(L\) defines a
morphism
\[
f:T\longrightarrow\Pic_{X/S}.
\]
The image of \(f\) contains \(\Lambda\) and is quasi-compact, hence
noetherian by (i). It follows that \(\Lambda\) is quasi-compact, proving
(iii).

By \hyperlink{sga6.e13.prop.3-2}{3.2}, the usual definition of morphisms of finite type between two
Picard schemes generalizes as follows.

\sgaStableTarget{sga6.e13.def.3-3}\textbf{Definition 3.3.} \emph{Let \(X\) and \(Y\) be proper over a
noetherian scheme \(S\). A morphism between the Picard functors}
\[
\begin{tikzcd}
\underline{\Pic}_{Y/S}\arrow[r]&\underline{\Pic}_{X/S}
\end{tikzcd}
\]



\emph{is said to be of finite type if every bounded family of classes of
invertible sheaves on the fibers of \(X/S\) has as inverse image a bounded
family of classes of invertible sheaves on the fibers of \(Y/S\).}

\sgaStableTarget{sga6.e13.lem.3-4}\textbf{Lemma 3.4.} \emph{Let \(X,Y,X',Y'\) be proper over a noetherian
scheme \(S\), and let \(T\to S\) be surjective and of finite type. Then:}
\begin{enumerate}[label=(\roman*)]
\item \emph{A morphism
\(\Phi:\underline{\Pic}_{Y/S}\to\underline{\Pic}_{X/S}\)
is of finite type if and only if
\(\Phi_T:\underline{\Pic}_{Y_T/T}\to\underline{\Pic}_{X_T/T}\)
is of finite type.}
\item \emph{In a commutative diagram}
\[
\begin{tikzcd}
\underline{\Pic}_{Y'/S}\arrow[r,"\Phi'"]&
\underline{\Pic}_{X'/S}\\
\underline{\Pic}_{Y/S}\arrow[r,"\Phi"']\arrow[u,"\Phi_Y"]&
\underline{\Pic}_{X/S}\arrow[u,"\Phi_X"']
\end{tikzcd}
\]
\emph{in which \(\Phi'\) and \(\Phi_Y\) are of finite type, \(\Phi\) is
of finite type.}
\end{enumerate}

Suppose that \(\Phi_T\) is of finite type. Let \(\mathscr L\) be a family
on the fibers of \(X/S\), bounded by an invertible sheaf \(L\) on some
\(X_R\). The family \(\mathscr L_T=\mathscr L\times_S T\) on the fibers of
\(X_T/T\) is bounded by the inverse image of \(L\) on
\(X_{R\times_S T}\). If \(\Phi_T^{-1}(\mathscr L_T)\) is bounded by an
invertible sheaf \(M\) on some \(Y_Q\), then \(M\) also bounds
\(\Phi^{-1}(\mathscr L)\). The converse in (i), and assertion (ii), are
proved similarly.

\sgaStableTarget{sga6.e13.thm.3-5}\textbf{Theorem 3.5.} \emph{Let \(S\) be a noetherian scheme, let \(X\)
and \(Y\) be proper \(S\)-schemes, and let \(f:X\to Y\) be a morphism. If
\(f\) is surjective, then the induced morphism between the Picard
functors}
\[
f^*:\underline{\Pic}_{Y/S}\longrightarrow\underline{\Pic}_{X/S}
\]

\emph{is of finite type.}

By \hyperlink{sga6.e13.lem.3-4}{3.4 (i)}, one may assume that \(S\) is reduced and irreducible. There
is a nonempty open subset \(U\subset S\) such that
\(\underline{\Pic}_{Y/S}\) and \(\underline{\Pic}_{X/S}\) are
representable and the restriction of \(f^*\) is quasi-affine (\hyperlink{sga6.e12.thm.1-1}{XII 1.1}),
hence of finite type by \hyperlink{sga6.e13.prop.3-2}{3.2 (i)}. Give \(T=S-U\) its induced reduced
structure. By noetherian induction, one may suppose that
\[
\underline{\Pic}_{Y_T/T}\longrightarrow\underline{\Pic}_{X_T/T}
\]
induced by \(f^*\) is of finite type. The assertion follows from \hyperlink{sga6.e13.lem.3-4}{3.4 (i)},
applied to \(U\amalg T\to S\).

\sgaStableTarget{sga6.e13.thm.3-6}\textbf{Theorem 3.6.} \emph{Let \(S\) be a noetherian scheme, let \(X\)
be a proper \(S\)-scheme, and let \(\underline{\Pic}_{X/S}\) be its
Picard functor. For every nonzero integer \(n\), the \(n\)-th power
homomorphism}
\[
\Phi_n:\underline{\Pic}_{X/S}\longrightarrow\underline{\Pic}_{X/S}
\]
\emph{is a morphism of finite type.}

Let \(f:X'\to X\) be a Chow presentation: \(X'\) is projective over \(S\),
and \(f\) is surjective. If the theorem is known for \(X'\), then \hyperlink{sga6.e13.thm.3-5}{3.5} and
\hyperlink{sga6.e13.lem.3-4}{3.4 (ii)} imply it for \(X\). We may suppose \(X/S\) projective.

There are a nonempty open \(U\subset S\), an integral \(S'\), and a
finite surjective morphism \(S'\to U\), such that every irreducible
component \(X'_i\) of \(X_{S'}\), with its induced reduced structure, is
integral over \(S'\) (\hyperlink{sga6.e12.lem.1-3}{XII 1.3}). Give \(T=S-U\) its induced reduced
structure. By noetherian induction, the theorem may be assumed for
\(X_T/T\). If it is assumed for all \(X_i/S'\), it follows for
\(\coprod_iX_i/S'\) by the

following easy lemma, then for \(X_{S'}/S'\) by \hyperlink{sga6.e13.thm.3-5}{3.5} and \hyperlink{sga6.e13.lem.3-4}{3.4 (ii)}, and
finally for \(X/S\) by \hyperlink{sga6.e13.lem.3-4}{3.4 (i)}.

\sgaStableTarget{sga6.e13.lem.3-7}\textbf{Lemma 3.7.} \emph{Let \(X=\coprod_iX_i\) be of finite type over a
noetherian scheme \(S\), and let \(\mathcal F\) be a family of classes of
coherent sheaves on the fibers of \(X/S\). Then \(\mathcal F\) is bounded
if and only if, for every \(i\), the induced family \(\mathcal F_i\) on
the fibers of \(X_i/S\) is bounded.}

In the proof of \hyperlink{sga6.e13.thm.3-6}{3.6} we reduced to \(X\) projective and integral over \(S\);
endow \(X/S\) with a very ample sheaf \(\mathcal O_X(1)\). Let
\(\mathscr L\) be bounded on the fibers of \(X/S\), and let \(L\) on some
\(X_T\), with \(T\) of finite type over \(S\), bound it. After
partitioning \(T\), suppose \(L\) flat over \(T\); after taking connected
components, suppose \(T\) connected. Then
\[
P(q,m)=\chi\bigl(L_t^{\otimes q}(m)\bigr)
\]
is independent of \(t\in T\) (EGA III 7.9.4).

Put \(\mathscr L_n=\Phi_n^{-1}(\mathscr L)\). It consists of classes of
invertible sheaves \(M_K\) whose \(n\)-th power class belongs to
\(\mathscr L\). Write
\[
\chi\bigl(M_K^{\otimes p}(m)\bigr)
=\sum_{i=0}^{r}a_i(p)\binom{m+i}{i},
\]
where \(a_i(p)\) has degree at most \(r-i\) in \(p\). Since
\(\chi(M_K^{\otimes qn}(m))=P(q,m)\), every \(a_i(qn)\), and hence every
polynomial \(a_i(p)\), is independent of \(M_K\). Thus \(\mathscr L_n\)
has one Hilbert polynomial \(\chi(M_K(m))\), and is bounded by 2.11.

\sgaStableTarget{sga6.e13.thm.3-8}\textbf{Theorem 3.8.} \emph{Let \(S\) be noetherian, let \(Y\) be a
projective \(S\)-scheme with a relatively ample invertible sheaf
\(\mathcal O_Y(1)\), let \(X\) be the zero scheme of a section \(\varphi\)
of \(\mathcal O_Y(1)\), and let \(f:X\to Y\) be the inclusion. Assume}

\emph{that for every \(s\in S\) and every irreducible component \(Y_s^i\)
of \(Y_s\), the irreducible components of
\(X_s\cap Y_s^i=V(\varphi|_{Y_s^i})\) have dimension at least \(2\)
(as happens when the irreducible components of the fibers of \(Y\) have
dimension at least \(3\)). Then}
\[
f^*:\underline{\Pic}_{Y/S}\longrightarrow\underline{\Pic}_{X/S}
\]
\emph{is of finite type. Moreover, if \(S\) is integral, there is a
nonempty open \(U\subset S\) such that the restricted Picard functors
are representable and}
\[
\underline{\Pic}_{Y/S}|_U\longrightarrow\underline{\Pic}_{X/S}|_U
\]
\emph{is quasi-affine (and of finite type).}

For the first assertion, noetherian induction permits replacing \(S\) by
any nonempty open. One reduces to \(S\) integral and \(Y\) integral over
\(S\), replacing \(X\) by \(X\times_Y Y'\) when \(Y\) is replaced by
\(Y'\). One then reduces to \(Y\) normal over \(S\), using the following
lemma in place of \hyperlink{sga6.e12.lem.1-3}{XII 1.3}.

\sgaStableTarget{sga6.e13.lem.3-9}\textbf{Lemma 3.9.} \emph{Let \(S\) be integral noetherian and \(X\) an
\(S\)-scheme integral over \(S\). There are a nonempty open \(U\subset S\),
an integral scheme \(S'\) finite faithfully flat over \(U\), an
\(S'\)-scheme \(X'\) normal and integral over \(S'\), and a finite
surjective \(S'\)-morphism \(X'\to X_{S'}\).}

Let \(\eta\) be the generic point of \(S\), and let \(K\) be an algebraic
closure of \(k(\eta)\). The scheme \(\operatorname{Spec}(K)\) is the
limit of the filtered inverse system \((S_\alpha)\) of integral affine
schemes with finite faithfully flat morphisms
\(S_\alpha\to U_\alpha\), where \(U_\alpha\) is a nonempty open of \(S\).
The normalization of \(X_K\) descends to an \(S_\alpha\)-scheme \(X'\) for
some \(\alpha\), and the finite surjective map \(X'_K\to X_K\) descends
to a finite surjective \(S_\alpha\)-morphism \(X'\to X_{S_\alpha}\)
(EGA IV 8.8.2, 8.10.5). Since \(X'_K\) is normal and integral, there is
an open \(S'\subset S_\alpha\) such that \(X'_{S'}\) is normal and
integral over \(S'\) (EGA IV 9.9.5).

\sgaStableTarget{sga6.e13.rem.3-10}\textbf{Remark 3.10.} Retain \hyperlink{sga6.e13.lem.3-9}{3.9} and suppose there is a resolving
morphism \(X'_K\to X_K\) (EGA IV 7.9.1), as happens when
\(\dim(X_K)\le2\), by Abhyankar. A similar argument gives a finite
faithfully flat \(S'\to U\), a scheme \(X'\) smooth over \(S'\), and a
proper surjective birational morphism \(X'\to X_{S'}\).

In the proof of \hyperlink{sga6.e13.thm.3-8}{3.8} we have reduced to \(S\) integral and \(Y\) normal and
integral over \(S\). If \(\varphi=0\), then \(X=Y\). Otherwise, after
restricting \(S\), suppose that \(\varphi_s\ne0\) for every \(s\in S\).
Thus \(\varphi_s\) is \(\mathcal O_{Y_s}\)-regular, and the dimension
\(r\) of \(Y_s\) is at least \(3\).

Let \(\mathscr L\) be bounded on the fibers of \(X/S\), and let \(M_K\)
be invertible on a geometric fiber \(Y_K\), with
\[
\chi(M_K(n))=\sum_{i=0}^{r}a_i\binom{n+i}{i}.
\]
The Hilbert polynomial of \(f^*M_K\) is
\[
\sum_{i=0}^{r-1}a_{i+1}\binom{n+i}{i}.
\]
As the class of \(M_K\) varies through \((f^*)^{-1}(\mathscr L)\), that
is, as

the class of \(f^*M_K\) varies through \(\mathscr L\), the coefficients
\(a_r,a_{r-1},a_{r-2}\) remain bounded by 1.13.

Let \(H=\mathcal O_Y(m)\) be a very ample multiple of
\(\mathcal O_Y(1)\). In the Hilbert polynomials relative to \(H\),
\[
\chi(M_K(mn))=\sum_{i=0}^{r}a_i^{(m)}\binom{n+i}{i},
\]
the coefficients \(a_r^{(m)},a_{r-1}^{(m)},a_{r-2}^{(m)}\) remain bounded
by 2.10. Consequently \((f^*)^{-1}(\mathscr L)\) is bounded by \hyperlink{sga6.e13.lem.2-11}{2.11},
proving the first part of 3.8.

For the second part, generic representability and separatedness of the
Picard functors were established in \hyperlink{sga6.e12.cor.1-2}{XII 1.2}. An argument like XII 2,
using \hyperlink{sga6.e13.lem.3-9}{3.9} and \hyperlink{sga6.e12.thm.1-1}{XII 1.1}, reduces to \(Y\) normal over \(S\). It is enough
for \(\underline{\Pic}_{Y/S}\to\underline{\Pic}_{X/S}\) to be separated
and quasi-finite (EGA IV 8.12.8), which can be checked on geometric
fibers and follows from:

\sgaStableTarget{sga6.e13.lem.3-11}\textbf{Lemma 3.11.} \emph{Let \(k\) be algebraically closed, let \(Y\)
be normal and projective with an ample sheaf \(\mathcal O_Y(1)\), and let
\(X\) be the zero scheme of a section \(\varphi\). Suppose
\(\dim(Y)\ge3\) (respectively, \(\dim(Y)\ge2\)). Then}
\[
\Pic_{Y/k}\longrightarrow\Pic_{X/k}
\quad
\text{(respectively, }\Pic^\tau_{Y/k}\longrightarrow\Pic^\tau_{X/k}
\text{; see \hyperlink{sga6.e13.para.4-1}{4.1})}
\]
\emph{is of finite type and has a finite kernel \(N\), in general
nonreduced.}

Assume \(\varphi\ne0\); it is regular because \(Y\) is integral. The map
is of finite type by \hyperlink{sga6.e13.thm.3-8}{3.8} (respectively,\hyperlink{sga6.e13.thm.4-7}{ 4.7 (iii)}). Its kernel \(N\) is
an algebraic group contained in \(\Pic^\tau_{Y/k}\), hence is the same

in both cases. Since \(\Pic^\tau_{Y/k}\) is proper because \(Y\) is
normal (FGA no.~236, 2.13), and \(N\) is closed (SGA 3 VI\(_B\) 1.9.2),
it is enough to show that \(N\) contains no subgroup of type
\(\mu_\ell\;(=\mathbb Z/\ell\mathbb Z)\), with
\((\ell,\operatorname{char}(k))=1\).

Consider Kummer theory, arising from
\[
0\longrightarrow\mu_\ell\longrightarrow\mathbb G_m
\xrightarrow{x\mapsto x^\ell}\mathbb G_m\longrightarrow0
\]
on \(Y\) (respectively, on \(X\)):
\[
\begin{tikzcd}[column sep=1.6em,row sep=1.2em]
&&H^1(Y,\mu_\ell)\arrow[r]\arrow[d]&
H^1(Y,\mathbb G_m)\arrow[r,"\ell"]\arrow[d]&
H^1(Y,\mathbb G_m)\arrow[d]\\
H^0(X,\mathbb G_m)\arrow[r,"\ell"]&
H^0(X,\mathbb G_m)\arrow[r]&
H^1(X,\mu_\ell)\arrow[r]&
H^1(X,\mathbb G_m)\arrow[r,"\ell"]&
H^1(X,\mathbb G_m).
\end{tikzcd}
\]
By Hilbert's Theorem 90, \(H^1(Y,\mathbb G_m)=\Pic_{Y/k}(k)\)
(respectively, \(H^1(X,\mathbb G_m)=\Pic_{X/k}(k)\)). Since \(X\) is
proper and integral, \(H^0(X,\mathbb G_m)=k^*\), so
\(H^1(X,\mu_\ell)\to H^1(X,\mathbb G_m)\) is injective.

It remains to show that
\[
H^1(Y,\mathbb Z/\ell\mathbb Z)\longrightarrow
H^1(X,\mathbb Z/\ell\mathbb Z)
\]
is injective. This follows from \(\pi_1(X,\xi)\to\pi_1(Y,\xi)\)
surjective; equivalently (SGA 1 V 6.9), every connected \'etale cover
\(\widetilde Y\) of \(Y\) with group \(\mathbb Z/\ell\mathbb Z\) has
connected restriction \(\widetilde X=X\times_Y\widetilde Y\). This is a
consequence of Bertini and Zariski connectedness (SGA 1 X 2.10).

\sgaStableTarget{sga6.e13.rem.3-12}\textbf{Remark 3.12.} Retain the hypotheses of 3.11.
\begin{enumerate}[label=(\roman*)]
\item The kernel \(N\) is finite \emph{unipotent} (SGA 3 XVII 1.1),
since it contains no subgroup of multiplicative type \(\mu_n\), for any
\(n\) (SGA 3 XVII 4.6.1 (vi)). This is proved similarly, using instead

of Kummer theory the following fact of Grothendieck (SGA 1 XI, last
page): if \(X\) is proper over \(k\), with \(H^0(X,\mathcal O_X)=k\),
then there is an isomorphism of \(k\)-group schemes
\[
\underline{\operatorname{Hom}}_{\mathrm{grp}}
(\mu_n,\underline{\Pic}_{X/k})
\simeq \underline{H^1(X,\mathbb Z/n\mathbb Z)}.
\]
Consequently, when \(k\) is algebraically closed, its \(k\)-points give
\[
\operatorname{Hom}_{\mathrm{grp}}(\mu_n,\underline{\Pic}_{X/k})
\simeq H^1(X,\mathbb Z/n\mathbb Z).
\]
\item If \(k\) has characteristic zero, then \(N=0\), since every finite
unipotent group is zero (SGA 3 XVII 1.7). This approximately recovers
Kodaira's vanishing theorem in Mumford's form [2].
\end{enumerate}

\sgaStableTarget{sga6.e13.thm.3-13}\textbf{Theorem 3.13.} \emph{Let \(S\) be noetherian and \(X\) a
projective \(S\)-scheme with a relatively ample invertible sheaf
\(H=\mathcal O_X(1)\). Suppose the geometric fibers of \(X/S\) are
integral of dimension \(r\). For a family \(\mathscr L\) of classes of
invertible sheaves on the fibers, the following are equivalent:}
\begin{enumerate}[label=(\roman*)]
\item \emph{\(\mathscr L\) is bounded.}
\item \emph{In
\(\chi(L_K(n))=\sum_{i=0}^{r}a_i\binom{n+i}{i}\),
the coefficient \(a_{r-1}\) remains bounded and \(a_{r-2}\) remains
bounded below.}
\item \emph{The degree
\[
d(L_K)=\langle c_1(L_K)c_1(H_K)^{r-1}\rangle
\]
remains bounded and
\(\langle c_1(L_K)^2c_1(H_K)^{r-2}\rangle\) remains bounded below.}
\item \emph{There is an integer \(N\) such that
\[
H^{-N}<\mathscr L<H^N;
\]
that is, \(H_K^{\otimes N}\otimes L_K\) and
\(H_K^{\otimes N}\otimes L_K^{-1}\) are ample for every
\(L_K\in\mathscr L\).}
\end{enumerate}

Suppose \(\mathscr L\) is bounded by \(L\) on some \(X_T\). Then (iv)
holds by EGA II 4.6.12. Taking \(L\) flat over \(T\), the polynomial
\(\chi(L_t^{\otimes m}(n))\) depends only on the connected component of
\(T\) (EGA III 7.9.4). Hence (ii) holds, and \hyperlink{sga6.e13.para.2-1}{2.1} gives (iii).

The families of classes of \(H_K\) and \(\mathcal O_{X_K}\) are bounded,
so they satisfy (iii) and (ii), respectively. If \(\mathscr L\) satisfies
(iv), then
\[
\langle[Nc_1(H_K)+c_1(L_K)]c_1(H_K)^{r-1}\rangle>0,
\]
so \(-Nd(H_K)<d(L_K)<Nd(H_K)\). Moreover,
\[
\langle[Nc_1(H_K)+c_1(L_K)]^2c_1(H_K)^{r-2}\rangle>0,
\]
whence
\[
\langle c_1(L_K)^2c_1(H_K)^{r-2}\rangle
>-2Nd(L_K)-N^2d(H_K).
\]
Thus (iii) holds.

Finally, suppose (ii), respectively (iii), holds. By \hyperlink{sga6.e13.rem.2-10}{2.10} and \hyperlink{sga6.e13.lem.2-3}{2.3},
replace \(H\) by a multiple and assume it very ample. Partition \(S\) so
there is a linear section \(Y\) of \(X\), integral over \(S\), with
two-dimensional fibers. By \hyperlink{sga6.e13.thm.3-8}{3.8}, replace \(X,\mathscr L\) by
\(Y,\mathscr L|_Y\); \hyperlink{sga6.e13.lem.1-7}{1.7}, respectively the projection formula, preserves
the relevant condition. Thus \(r=2\). Put
\[
\delta=\sup\{d(L_K)\}.
\]
If \(\delta>0\), replace \(\mathscr L\) by the classes of
\(L_K(-\delta)\), using \hyperlink{sga6.e13.lem.2-3}{2.3} and \hyperlink{sga6.e13.rem.2-10}{2.10}, so that \(\delta\le0\).

If \(a_1\) remains bounded and \(a_0\) bounded below in
\(\sum_{i=0}^{2}a_i\binom{n+i}{i}\), all three coefficients remain
bounded by \hyperlink{sga6.e13.lem.2-3}{2.3} and 2.7. Thus \(\mathscr L\) is bounded by \hyperlink{sga6.e13.lem.2-11}{2.11}, so (ii)
implies

(i). On the other hand, (iii) implies (ii) by \hyperlink{sga6.e13.lem.2-3}{2.3} and 2.8.

\sgaStableTarget{sga6.e13.sec.4}\subsection*{4. Finiteness theorems for \(\Pic^\tau_{X/S}\)}

\sgaStableTarget{sga6.e13.para.4-1}\textbf{4.1.} Let \(S\) be a scheme and \(G\) a commutative group functor
over \(S\). Denote by \(G^0\), respectively \(G^\tau\), the subgroup
functor whose \(T\)-points \(f:T\to G\) satisfy: for every algebraically
closed field \(K\) and \(K\)-point \(t\) of \(T\), there are a connected
\(K\)-scheme \(Z\), two \(K\)-points \(z,z_0\), and \(g:Z\to G\), with
\(g(z)=f(t)\) and \(g(z_0)=0\) in \(G(K)\); respectively, there is a
nonzero \(n\) such that \(nf(t)\in G^0(K)\).

Formation of \(G^0\), respectively \(G^\tau\), commutes with base change,
and \(G\to G'\) induces \(G^0\to G'^0\), respectively
\(G^\tau\to G'^\tau\). If \(S\) is a field and \(G\) is represented by a
group scheme locally of finite type, \(G^0\) is represented by the
connected component of the identity, an open and closed finite-type
geometrically irreducible subscheme (SGA 3 VI\(_A\) 2.4).

\sgaStableTarget{sga6.e13.lem.4-2}\textbf{Lemma 4.2.} \emph{Let \(\Phi:G\to G'\) be a morphism of
commutative group functors over \(S\). Suppose that for every
algebraically closed \(K\) over \(S\), \(G_K,G'_K\) are represented by
group schemes locally of finite type and \(\Phi_K\) by a finite-type
(equivalently quasi-compact) morphism. Then}
\[
\Phi^{-1}(G'^\tau)=G^\tau.
\]

The inclusion \(\Phi^{-1}(G'^\tau)\supset G^\tau\) has already been
noted. Conversely, \(g\in G(K)\) belongs to \(G^\tau(K)\) if and only if
the set of multiples \(ng\), \(n\in\mathbb N\), is quasi-compact. Since
\(\Phi_K\) is quasi-compact, it is enough that the set
\(\Phi(ng)=n\Phi(g)\) have this property. This proves the lemma.

\sgaStableTarget{sga6.e13.lem.4-3}\textbf{Lemma 4.3.} \emph{Let \(X,Y\) be proper over \(S\), and let
\(f:X\to Y\) be surjective (respectively, the inclusion of the zero
scheme of a section \(\varphi\) of an ample \(\mathcal O_Y(1)\)). Then}
\[
(f^*)^{-1}\underline{\Pic}^\tau_{X/S}
=\underline{\Pic}^\tau_{Y/S}.
\]
This follows formally from \hyperlink{sga6.e13.lem.4-2}{4.2}, \hyperlink{sga6.e12.cor.1-2}{XII 1.2}, \hyperlink{sga6.e13.lem.3-1}{3.1}, and \hyperlink{sga6.e13.thm.3-5}{3.5} (respectively,
\hyperlink{sga6.e13.thm.3-8}{3.8}).

\sgaStableTarget{sga6.e13.para.4-4}\textbf{4.4.} Let \(k\) be algebraically closed, \(X\) proper over \(k\),
and \(L\) invertible on \(X\). The sheaf \(L\) is \emph{algebraically
equivalent} to zero if there are a connected finite-type \(k\)-scheme
\(T\), an invertible sheaf \(M\) on \(X_T=X\times_kT\), and two
\(k\)-points \(t,t_0\), such that \(L\simeq M_t\) and
\(\mathcal O_X\simeq M_{t_0}\). It is \emph{\(\tau\)-equivalent} to zero
if some nonzero tensor power is algebraically equivalent to zero.
Equivalently, the point \(\lambda\) corresponding to \(L\) belongs to
\(\underline{\Pic}^0_{X/k}(k)\), respectively
\(\underline{\Pic}^\tau_{X/k}(k)\).

Recall that \(L\) is \emph{numerically equivalent} to zero if
\[
\langle c_1(L)\cdot Y\rangle=0
\]
for every closed integral curve \(Y\subset X\) (see \hyperlink{sga6.e13.para.2-1}{2.1} or \hyperlink{sga6.e10.sec.4.5}{X 4.5}).

\sgaStableTarget{sga6.e13.prop.4-5}\textbf{Proposition 4.5 (Functorial behavior).} \emph{Let
\(f:X'\to X\) be a morphism between proper \(k\)-schemes, and let \(L\)
be invertible on \(X\). Then:}
\begin{enumerate}[label=(\roman*)]
\item \emph{If \(L\) is \(\tau\)-equivalent to zero (respectively,
numerically equivalent to zero), then so is \(f^*L\).}
\item \emph{Conversely, if \(f\) is surjective and \(f^*L\) is
\(\tau\)-equivalent to zero (respectively, numerically equivalent to
zero), then so is \(L\).}
\end{enumerate}

The assertions for \(\tau\)-equivalence follow from \hyperlink{sga6.e13.para.4-1}{4.1} and 4.3. Let
\(Y'\) be a closed integral curve in \(X'\), and let \(Y=f(Y')\), with
its induced reduced structure. If \(\dim(Y)=1\), the projection formula
(\hyperlink{sga6.e10.eq.4-4-4}{X 4.4.4}) gives
\sgaStableTarget{sga6.e13.eq.4-5-1}\begin{equation}
\langle c_1(f^*L)\cdot Y'\rangle
=[k(Y'):k(Y)]\,\langle c_1(L)\cdot Y\rangle.
\tag{4.5.1}
\end{equation}
If \(\dim(Y)=0\), then \((f^*L)_{Y'}\simeq\mathcal O_{Y'}\). This proves
(i) for numerical equivalence.

Conversely, for a closed integral curve \(Y\subset X\), choose a point
\(y'\in f^{-1}(y)\) algebraic over the function field \(k(y)\) of the
generic point \(y\), and let \(Y'\) be its closure. Then \(f(Y')=Y\), and
\hyperlink{sga6.e13.eq.4-5-1}{4.5.1} proves (ii).

\sgaStableTarget{sga6.e13.thm.4-6}\textbf{Theorem 4.6 (The numerical characterization of
\(\Pic^\tau_{X/k}\)).} \emph{Let \(k\) be algebraically closed, \(X\)
proper over \(k\), and \(L\) invertible on \(X\). The following are
equivalent:}
\begin{enumerate}[label=(\roman*)]
\item \emph{\(L\) is \(\tau\)-equivalent to zero.}
\item \emph{For every coherent \(F\) on \(X\),}

\[
\chi(F\otimes L)=\chi(F).
\]
\item \emph{\(L\) is numerically equivalent to zero.}
\end{enumerate}
\emph{If \(X/k\) is projective with an ample
\(H=\mathcal O_X(1)\), these are also equivalent to:}
\begin{enumerate}[label=(\roman*),start=4]
\item \emph{\(L^{\otimes m}(1)\) is ample for every integer \(m\).}
\item \emph{If \(X\) is irreducible, then
\[
\chi(L^{\otimes m}(n))=\chi(\mathcal O_X(n))
\]
for every \(m,n\).}
\item \emph{If \(X\) is irreducible of dimension \(r\ge2\), then
\[
\langle c_1(L)c_1(H)^{r-1}\rangle=0,\qquad
\langle c_1(L)^2c_1(H)^{r-2}\rangle=0.
\]}
\end{enumerate}

For (i) \(\Rightarrow\) (ii), devissage (EGA III 3.1) and \hyperlink{sga6.e13.prop.4-5}{4.5} reduce to
\(X\) integral and \(F=\mathcal O_X\). It remains to show that
\(\chi(L^{\otimes m})\) is constant. Replace \(L\) by a nonzero tensor
power and suppose it algebraically equivalent to zero. If \(M\) on
\(X\times_kT\), with \(T\) connected of finite type, deforms \(L\) to
\(\mathcal O_X\), then \(\chi(M_t^{\otimes m})\) is independent of
\(t\in T\) (EGA III 7.9.5).

For a closed integral curve \(Y\), take \(F=\mathcal O_Y\) in (ii).
The Riemann formula
\[
\chi(L_Y)=\langle c_1(L)\cdot Y\rangle+\chi(\mathcal O_Y)
\]
shows that (ii) implies (iii).

Let \(f:X'\to X\) be a Chow presentation. If (iii) implies (i) for
\(X'\), then \hyperlink{sga6.e13.prop.4-5}{4.5} gives it for \(X\). Hence

assume \(X\) projective with an ample \(H\). Passing through the finite
surjection \(\coprod_iX_i\to X\), where the \(X_i\) are the reduced
irreducible components, reduces (iii) \(\Rightarrow\) (i),(iv) to \(X\)
integral (\hyperlink{sga6.e13.prop.4-5}{4.5}, EGA III 2.6.2).

For irreducible \(X\), (ii) implies (v), (iii) implies (vi), and \hyperlink{sga6.e13.para.2-1}{2.1}
gives (v) implies (vi); condition (vi) is unchanged on
\(X_{\mathrm{red}}\). It remains to prove (vi) implies (i),(iv), and
(iv) implies (i), for \(X\) integral. Under (vi), respectively (iv), the
family of classes \(L^{\otimes m}\), \(m\in\mathbb N\), satisfies \hyperlink{sga6.e13.thm.3-13}{3.13
(iii)}, respectively (iv), and is bounded. Hence some invertible
\(\mathcal M\) on \(X_T\), with \(T\) of finite type over \(k\), bounds
it. Choose distinct \(p,q\) such that \(L^{\otimes p}\) and
\(L^{\otimes q}\) correspond to the same connected component of \(T\).
Then \(L^{\otimes(p-q)}\) is algebraically equivalent to zero, proving
(i).

Condition (vi) also gives an \(N\) such that
\(L^{\otimes m}\otimes H^{\otimes N}\) is ample for every \(m\).
Therefore \((L^{\otimes m}\otimes H)^{\otimes N}\) is ample, so
\(L^{\otimes m}\otimes H\) is ample. Thus (vi) implies (iv).

\sgaStableTarget{sga6.e13.thm.4-7}\textbf{Theorem 4.7.} \emph{Let \(S\) be noetherian and \(X\) proper over
\(S\). Then:}
\begin{enumerate}[label=(\roman*)]
\item \emph{\(\underline{\Pic}^{\tau}_{X/S}\to
\underline{\Pic}_{X/S}\) is representable by an open immersion.}
\item \emph{If \(X\) is projective and integral over \(S\), this map}

\emph{is also representable by a closed immersion.}
\item \emph{\(\underline{\Pic}^{\tau}_{X/S}\) is of finite type over
\(S\); that is, the family of classes \(\tau\)-equivalent to zero on the
fibers is bounded.}
\end{enumerate}

For (i), respectively (ii), take an arbitrary
\(T\to\underline{\Pic}_{X/S}\) and show that the inverse image of
\(\underline{\Pic}^{\tau}_{X/S}\) is an open, respectively open and
closed, subscheme of \(T\). Locally, \(T=\operatorname{Spec}(B)\) over
\(\operatorname{Spec}(A)\subset S\), and \(T\) is the limit of
\(T_\alpha=\operatorname{Spec}(B_\alpha)\), for the finite-type
\(A\)-subalgebras \(B_\alpha\subset B\). By \hyperlink{sga6.e13.lem.3-1}{3.1} the morphism descends to
some \(T_\alpha\), so assume \(T\) of finite type over \(S\).

The morphism corresponds to an fppf \(R\to T\) and an invertible sheaf
\(L\) on \(X_R\). By fppf descent replace \(T\) by \(R\), then base-change
along \(R\to S\). Compatibility with base change \hyperlink{sga6.e13.para.4-1}{(4.1)} reduces to a
section \(S\to\underline{\Pic}_{X/S}\) defined by \(L\) on \(X\).

Let \(S_0\) be the set of \(s\in S\) such that \(L_s\) is
\(\tau\)-equivalent to zero. It must be open, and under (ii) also closed.
Since \(S\) is noetherian, for \(s\in S_0\) it suffices that
\begin{enumerate}[label=(\alph*)]
\item every generalization of \(s\) lies in \(S_0\);
\item every specialization sufficiently near \(s\) lies in \(S_0\);
\item under (ii), every specialization of \(s\) lies in \(S_0\).
\end{enumerate}

Let \(f:X'\to X\) be a Chow presentation. If (a),(b),(c), and (iii) hold
for \(X'\), then \hyperlink{sga6.e13.lem.4-3}{4.3} and \hyperlink{sga6.e13.thm.3-5}{3.5} give them for \(X\); assume \(X\) projective
over \(S\).

Let \(H\) be ample for \(X/S\). If
\((L^{\otimes m}\otimes H)_s\) is ample, then it remains ample at every
generalization of \(s\) (EGA III 4.7.1). The equivalence \hyperlink{sga6.e13.thm.4-6}{4.6
(i)} \(\Longleftrightarrow\) (iv) gives (a).

For (b),(c), and (iii), reduce to \(S\) integral, with \(s\) generic in
(b),(c); for (iii), apply \hyperlink{sga6.e13.lem.3-4}{3.4 (i)} to
\(\coprod_iS_i\to S\), where \(S_i\) are the reduced irreducible
components.

There are an integral \(S'\) and an fppf \(S'\to S\) such that every
reduced irreducible component \(X'_i\) of \(X_{S'}\) is integral over
\(S'\) (\hyperlink{sga6.e12.lem.1-3}{XII 1.3}). It is enough to prove (b),(iii) for \(X_{S'}/S'\), by
descent and by noetherian induction with \hyperlink{sga6.e13.lem.3-4}{3.4 (i)}. The surjection
\(\coprod_iX'_i\to X_{S'}\), together with \hyperlink{sga6.e13.lem.4-3}{4.3}, \hyperlink{sga6.e13.thm.3-5}{3.5}, and \hyperlink{sga6.e13.lem.3-7}{3.7}, reduces
(b),(iii) to (c),(iii) for each \(X'_i/S'\).

Thus we are reduced to proving (c) and (iii), assuming that \(S\) is
integral with generic point \(s\), and that \(X\) is projective and
integral over \(S\). For a very ample \(\mathcal O_X(1)\) for \(X/S\),
the values \(\chi(L_t^{\otimes m}(n))\) and
\(\chi(\mathcal O_{X_t}(n))\), for every \(m,n\), are independent of
\(t\in S\) (EGA III 7.9.4). Hence (c) follows from \hyperlink{sga6.e13.thm.4-6}{4.6
(i)} \(\Longleftrightarrow\) (v), which also says that the family of
classes of invertible sheaves \(\tau\)-equivalent to zero on the fibers
of \(X/S\) has only one

Hilbert polynomial. Thus (iii) follows from 2.11.

\sgaStableTarget{sga6.e13.rem.4-8}\textbf{Remark 4.8.} Let \(X\) be proper and integral over a noetherian
scheme \(S\). Without the projectivity hypothesis on \(X/S\), it is
unknown whether
\(\underline{\Pic}^{\tau}_{X/S}\to\underline{\Pic}_{X/S}\) remains a
\emph{closed} immersion, except when \(X\) is normal over \(S\) and
\(\underline{\Pic}_{X/S}\) is \emph{representable} (FGA 236, 2.3).

\sgaStableTarget{sga6.e13.sec.5}\subsection*{5. Finiteness theorems of ``N\'eron--S\'everi'' type}

\sgaStableTarget{sga6.e13.thm.5-1}\textbf{Theorem 5.1.} \emph{Let \(X\) be proper over noetherian \(S\).
Then the N\'eron--S\'everi groups
\(\underline{\Pic}_{X_K/K}(K)/\underline{\Pic}^0_{X_K/K}(K)\) of the
geometric fibers \(X_K/K\) of \(X/S\) are finitely generated, and their
ranks and the orders of their torsion subgroups
\(\underline{\Pic}^\tau_{X_K/K}(K)/\underline{\Pic}^0_{X_K/K}(K)\)
remain bounded.}

By noetherian induction, replace \(S\) by an arbitrary nonempty open of
\(S_{\mathrm{red}}\), and suppose
\(\underline{\Pic}^{\tau}_{X/S}\) represented by a finite-type
\(S\)-scheme (\hyperlink{sga6.e12.cor.1-2}{XII 1.2},\hyperlink{sga6.e13.thm.4-7}{ 4.7 (iii)}). The order of
\(\underline{\Pic}^\tau_{X_K/K}(K)/\underline{\Pic}^0_{X_K/K}(K)\)
is the number of connected components of
\((\underline{\Pic}^{\tau}_{X/S})_K\), constant over a nonempty open of
\(S\) (EGA IV 9.7.9).

It remains to bound the rank of
\(\underline{\Pic}_{X_K/K}(K)/\underline{\Pic}^\tau_{X_K/K}(K)\).
Reduce as in \hyperlink{sga6.e13.thm.4-7}{4.7 (iii)} to
\(S\) integral and \(X\) projective and integral. Let \(r\) be the common
fiber dimension of \(X/S\). For \(r=0\), respectively \(1\), the rank is \(0\),
respectively \(1\). For \(r\ge2\), reduce to \(r=2\) and \(X\) smooth.
After restricting \(S\), take a linear section \(Y\) of \(X\), still

integral over \(S\), with two-dimensional fibers, and replace \(X\) by
\(Y\) using 4.3. Then use \hyperlink{sga6.e13.lem.4-3}{4.3} and \hyperlink{sga6.e13.rem.3-10}{3.10}, together with resolution of
surface singularities, and replace \(S\) by a flat cover of a nonempty
open, to make \(X\) smooth over \(S\).

By \hyperlink{sga6.e13.thm.4-6}{4.6 (i)} \(\Longleftrightarrow\) (iii) and finiteness and specialization
for \'etale cohomology (SGA 4 XVI 2.2), Theorem 5.1 follows from:

\sgaStableTarget{sga6.e13.prop.5-2}\textbf{Proposition 5.2.} \emph{Let \(k\) be algebraically closed and
\(X\) proper smooth over \(k\). The group \(C^p_{\mathrm{num}}(X)\) of
codimension-\(p\) algebraic cycles modulo numerical equivalence is free
of rank at most \(b^{2p}\), the \(2p\)-th Betti number.}

Disregarding the action of roots of unity, \'etale cohomology
\(H^*(X,\mathbb Q_\ell)=_{\mathrm{dfn}}\allowbreak
\left[\varprojlim_\nu H^*(X,\mathbb Z/\ell^\nu\mathbb Z)\right]
\allowbreak\otimes_{\mathbb Z_\ell}\mathbb Q_\ell\), where
\((\ell,\operatorname{char}(k))=1\), is a Weil cohomology: it satisfies
Poincar\'e duality and the K\"unneth formula, and receives a functorial,
nonzero homomorphism from the ring of algebraic cycles modulo rational
equivalence \(C^*_{\mathrm{rat}}(X)\) (see \hyperlink{sga6.e10.sec.7.13}{X 7.13}).

Choose \(a_1,\ldots,a_m\in C^{r-p}_{\mathrm{rat}}(X)\), \(r=\dim X\),
whose images form a basis for the \(\mathbb Q_\ell\)-subspace of
\(H^{2(r-p)}(X,\mathbb Q_\ell)\) spanned by algebraic cycles. Define
\[
\alpha:C^p_{\mathrm{rat}}(X)\longrightarrow\mathbb Z^m.
\]
It is defined by
\(\alpha(b)=(\langle b\cdot a_1\rangle,\ldots,\langle b,a_m\rangle)\),
where \(\langle\ \rangle\) denotes the intersection number. Its image
identifies with \(C^p_{\mathrm{num}}(X)\).

\sgaStableTarget{sga6.e13.rem.5-3}\textbf{Remark 5.3.} If \(X\) is a proper smooth surface over
algebraically closed \(k\), then its second Betti number is
\(b^2(X)=c_2(X)-2+2b^1(X)\), where
\(b^1(X)=2\dim\Pic^0_{X/k}\), and \(c_2(X)\), the second Chern
class, is the Euler-Poincar\'e characteristic, equivalently the
self-intersection of the diagonal in \(X\times X\). Thus \(b^2(X)\) has
a cohomology-independent interpretation. In characteristic \(>0\), the
inequality \(\operatorname{rank}C^1_{\mathrm{num}}(X)\le b^2(X)\) in
this form is due, before \'etale cohomology, to Igusa, who used deep
results on the tame fundamental group of an algebraic curve minus some
points.

\sgaStableTarget{sga6.e13.thm.5-4}\textbf{Theorem 5.4.} \emph{Let \(k\) be algebraically closed and \(X\)
proper over \(k\). There are finitely many closed integral curves
\(Y_1,\ldots,Y_p\) in \(X\) such that a subset \(\Lambda\) of
\(\Pic_{X/k}\) is quasi-compact if and only if the intersection numbers
\(\langle\lambda\cdot Y_i\rangle\), for \(\lambda\in\Lambda\), remain
bounded. Moreover, \(p\) may be taken to be the rank of the
N\'eron--S\'everi group \(\Pic_{X/k}(k)/\Pic^0_{X/k}(k)\).}

Indeed, \(\Pic^\tau_{X/k}\) is an open quasi-compact subgroup of
\(\Pic_{X/k}\) \hyperlink{sga6.e13.thm.4-7}{(4.7)}. Thus \(\Lambda\) is quasi-compact if and only if
the corresponding set of points in
\(\bigl[\Pic_{X/k}(k)/\Pic^\tau_{X/k}(k)\bigr]
\otimes_{\mathbb Z}\mathbb Q\) is finite. By \hyperlink{sga6.e13.thm.4-6}{4.6 (i)}
\(\Longleftrightarrow\) (iii),
\(\tau\)-equivalence equals numerical equivalence. Choose curves whose
images form a basis of the dual vector space; the assertion follows.



\sgaStableTarget{sga6.e13.sec.6}\subsection*{6. Appendix: Study of $(b)$-sheaves on $P=\mathbf P^N_k$.}

\sgaStableTarget{sga6.e13.para.6-1}\textbf{6.1.} Let $k$ be an algebraically closed field and let $P=\mathbf P^N_k$. Let $F$ be a coherent sheaf on $P$, with Hilbert polynomial
\[
  \chi(F(n))=\sum_{i=0}^{r}a_i\binom{n+i}{i},
\]
and let
\[
  (c)=(c_0,\ldots,c_r)
\]
be a sequence of integers with $(c)\geq(a)$. Put
\[
  a_{m,i}=\sum_{j=0}^{r-i}a_{j+i}\binom{m-1+j}{j}
  \quad\text{and}\quad
  c_{m,i}=\sum_{j=0}^{r-i}c_{j+i}\binom{m-1+j}{j}.
\]
Finally, for every integer $q\geq1$, denote by $N_q$ the largest coherent subsheaf of $F$ whose support has dimension at most $q-1$, and put $F_q=F/N_q$.

\sgaStableTarget{sga6.e13.lem.6-2}\textbf{Lemma 6.2.} \emph{Under the conditions of \hyperlink{sga6.e13.para.6-1}{6.1}, consider an exact sequence
\[
  0\longrightarrow I\longrightarrow\cO_P^{\oplus M}
  \longrightarrow F(m)\longrightarrow0
\]
with $m\geq0$. Then $I$ is $p$-regular, where
\(p=P_r(c_{m,0},\ldots,c_{m,r})\), and $P_r$ is the $r$-th
$(b)$-polynomial of 1.9.}

Indeed, $\cO_P$ is plainly a $(0,\ldots,0,1)$-sheaf ($N+1$ terms), so
$I$ is a $(0,\ldots,0,M)$-sheaf by 1.6. By \hyperlink{sga6.e13.rem.2-10}{2.10},
\(\chi(I(n))=M\binom{n+N}{N}
-\sum_{i=0}^{r}a_{m,i}\binom{n+i}{i}\). The assertion follows from
\hyperlink{sga6.e13.thm.1-11}{1.11} and 1.10.

\sgaStableTarget{sga6.e13.prop.6-3}\textbf{Proposition 6.3.} \emph{Under the conditions of \hyperlink{sga6.e13.para.6-1}{6.1}, suppose that $F$ is a $(b)$-sheaf. Then, by means of universal polynomials in the $a_i$, the $b_i$, and $N$, one can calculate integers $m,M,m_1,M_1$ for which there is an exact sequence
\[
  \cO_P(-m_1)^{\oplus M_1}\longrightarrow
  \cO_P(-m)^{\oplus M}\longrightarrow F\longrightarrow0.
\]
Moreover, $N$ enters only into the calculation of $M_1$.}

Indeed, take $m=P_{r-1}(c_0,\ldots,c_{r-1})$, with $c_i=b_i-a_i$,
and $M=\sum_i a_i\binom{m+i}{i}$, so that, by \hyperlink{sga6.e13.thm.1-11}{1.11} and \hyperlink{sga6.e13.prop.1-3}{1.3}, there is
an exact sequence
\[
  0\longrightarrow I\longrightarrow\cO_P^{\oplus M}
  \longrightarrow F(m)\longrightarrow0.
\]
Likewise, taking account of \hyperlink{sga6.e13.lem.6-2}{6.2}, take $m_1=m+p$, where
$p=P_r(a_{m,0},\ldots,a_{m,r})$, and
$M_1=M\binom{p+N}{N}-\sum_i a_i\binom{m_1+i}{i}
\quad(=\chi(I(p)))$.

\sgaStableTarget{sga6.e13.thm.6-4}\textbf{Theorem 6.4.} \emph{Under the conditions of \hyperlink{sga6.e13.para.6-1}{6.1}, suppose that
$F$ is a quotient of $\cO_P(-m)^{\oplus M}$, with $m\geq0$. Put
$b_i=P_{r-i}(c_{m,i},\ldots,c_{m,r})$ for $i=0,\ldots,r$, where $P_s$
is the $s$-th $(b)$-polynomial of \hyperlink{sga6.e13.def.1-9}{1.9}; put $b=b_1+m-1$, and put
$B=\sum_{i=1}^{r}a_i\binom{b+i}{i}$. Then:}
\begin{enumerate}[label=(\roman*)]
\item \emph{$F$ is $b$-regular.}
\item \emph{$-B\leq a_0=h^0(F(b))-B$}
\item \emph{$F$ is a $(b_0,\ldots,b_r)$-sheaf.}
\end{enumerate}

Indeed, choose a section $\sigma$ of $\cO_P(1)$ for which one obtains the following commutative diagram with exact rows and columns:
\[
\begin{tikzcd}[column sep=large,row sep=large]
&0 \arrow[d]&0 \arrow[d]&0 \arrow[d]&\\
0 \arrow[r]&I(-1) \arrow[r] \arrow[d]&I \arrow[r] \arrow[d]&J \arrow[r] \arrow[d]&0\\
0 \arrow[r]&\cO_P(-1)^{\oplus M} \arrow[r] \arrow[d]&\cO_P^{\oplus M} \arrow[r] \arrow[d]&\cO_Y^{\oplus M} \arrow[r] \arrow[d]&0\\
0 \arrow[r]&F(m-1) \arrow[r] \arrow[d]&F(m) \arrow[r] \arrow[d]&G(m) \arrow[r] \arrow[d]&0\\
&0&0&0&
\end{tikzcd}
\]

Since $\chi(G(n))=\sum_{i=0}^{r-1}a_{i+1}\binom{n+i}{i}$ by \hyperlink{sga6.e13.lem.1-7}{1.7},
$J$ is $b_1$-regular by \hyperlink{sga6.e13.lem.6-2}{6.2}; hence, by \hyperlink{sga6.e13.prop.1-4}{1.4(i)}, $H^q(I(p))=0$ for
$q\geq2$ and $p\geq b_1-q$. The exact sequence
$H^q\bigl(\cO_P^{\oplus M}(p)\bigr)\longrightarrow
H^q(F(m+p))\longrightarrow H^{q+1}(I(p))$ then shows that $F(m)$ is
$(b_1-1)$-regular; this proves (i).

Assertion (i) implies
$0\leq h^0(F(b))=\chi(F(b))
=\sum_{i=0}^{r}a_i\binom{b+i}{i}$, which proves (ii). By \hyperlink{sga6.e13.rem.2-10}{2.10} and \hyperlink{sga6.e13.def.1-9}{1.9},
\[
  \chi(F(b))
  =\sum_{i=0}^{r}a_{m,i}\binom{b_1-1+i}{i}
  \leq b_1+\sum_{i=0}^{r}c_{m,i}\binom{b_1-1+i}{i}
  =b_0.
\]
Thus $h^0(F(-1))\leq h^0(F(b))\leq b_0$. By induction on $r$, one may suppose that $G$ is a $(b_1,\ldots,b_r)$-sheaf, so (iii) follows from 1.6.2.

\sgaStableTarget{sga6.e13.lem.6-5}\textbf{Lemma 6.5.} \emph{Under the conditions of \hyperlink{sga6.e13.para.6-1}{6.1}, suppose that $F$
has no associated prime cycle of dimension $0$.}
\begin{enumerate}[label=(\roman*)]
\item \emph{If $h^0(F)\geq1$, then $h^0(F(-1))\leq h^0(F)-1$.}
\item \emph{$H^0(F(-n))=0$ for $n\geq h^0(F)$.}
\item \emph{Let $0\to F(-1)\to F\to G\to0$ be exact, and let
$n_0\geq0$ satisfy $H^0(G(-n_0))=0$. Then $H^0(F(-n_0))=0$ and
$h^0(F)\leq n_0\cdot h^0(G)$.}
\end{enumerate}

For (i), let $0\ne\sigma\in H^0(F)$. The section $\sigma$ defines a nonzero subsheaf $F'=\cO_X\cdot\sigma$ of $F$, so $s=\dim\Supp(F')\geq1$. Consider the commutative diagram with exact rows and columns
\[
\begin{tikzcd}[column sep=large,row sep=large]
&0 \arrow[d]&0 \arrow[d]&0 \arrow[d]&\\
0 \arrow[r]&F'(-1) \arrow[r] \arrow[d]&F' \arrow[r] \arrow[d]&G' \arrow[r] \arrow[d]&0\\
0 \arrow[r]&F(-1) \arrow[r]&F \arrow[r]&G \arrow[r]&0.
\end{tikzcd}
\]

Since $s\geq1$, one has $G'\ne0$; hence $F'$ is not contained in $F(-1)$. In other words, the section $\sigma$ is not in the image of the injection $H^0(F(-1))\to H^0(F)$, which proves (i).

Assertion (ii) follows immediately from (i), bearing in mind that
$h^0(F)=0$ plainly implies $h^0(F(-1))=0$, since
$h^0(F(-1))\leq h^0(F)$. In (iii) one obtains
$0\leq h^0(F(n))-h^0(F(n-1))\leq h^0(G(n))$.
Thus, for $n\geq n_0$, $h^0(F(-n_0))=h^0(F(-n))$, since $h^0(G(-n_0))=0$; hence $h^0(F(-n_0))=0$ by (ii). Consequently, for $p\geq1$,
$h^0(F(p-n_0))
\leq h^0(G(1-n_0))+\cdots+h^0(G(p-n_0))
\leq p\cdot h^0(G(p-n_0))$,
which gives (iii).

\sgaStableTarget{sga6.e13.prop.6-6}\textbf{Proposition 6.6.} \emph{Under the conditions of \hyperlink{sga6.e13.para.6-1}{6.1}, suppose
that $F$ is a $(b_0,\ldots,b_r)$-sheaf. Then $F_q$ is a
$\bigl((b_{q-1})^q,\ldots,(b_{q-1})^2,b_{q-1},b_q,\ldots,b_r\bigr)$-sheaf.}

Indeed, let $\sigma$ be a ``general'' section of $\cO_X(1)$. It defines a commutative diagram with exact rows and columns:
\[
\begin{tikzcd}[column sep=large,row sep=large]
&0 \arrow[d]&0 \arrow[d]&0 \arrow[d]&\\
0 \arrow[r]&N_q(-1) \arrow[r] \arrow[d]&N_q \arrow[r] \arrow[d]&G' \arrow[r] \arrow[d]&0\\
0 \arrow[r]&F(-1) \arrow[r] \arrow[d]&F \arrow[r] \arrow[d]&G \arrow[r] \arrow[d]&0\\
0 \arrow[r]&F_q(-1) \arrow[r] \arrow[d]&F_q \arrow[r] \arrow[d]&G'' \arrow[r] \arrow[d]&0\\
&0&0&0&
\end{tikzcd}
\]
where $G$ is a $(b_1,\ldots,b_r)$-sheaf by \hyperlink{sga6.e13.prop.1-6}{1.6(i)}, and where
$G''=G_{q-1}$ (a depth question).

Suppose first that $q=1$. Then $\dim\Supp(N_q)=0$, so
$H^1(N_q(-1))=0$ and $h^0(F_q(-1))\leq h^0(F(-1))\leq b_0$.
Moreover, $G'=0$, and hence $G\simeq G''$. Thus $F_1$ is indeed a $(b_0,\ldots,b_r)$-sheaf. In addition,\hyperlink{sga6.e13.lem.6-5}{ 6.5(ii)} shows that
$H^0(F_1(-b_0))=0$. If $q\geq2$, argue by induction on $q$. One may
suppose that $G_{q-1}$ is a
$\bigl((b_{q-1})^{q-1},\ldots,b_{q-1},\ldots,b_r\bigr)$-sheaf and that
$H^0(G_{q-1}(-b_{q-1}))=0$. \hyperlink{sga6.e13.lem.6-5}{Lemma 6.5} then gives
$H^0(F_q(-b_{q-1}))=0$ and $h^0(F_q(-1))\leq(b_{q-1})^q$,
which proves the assertion.

\sgaStableTarget{sga6.e13.thm.6-7}\textbf{Theorem 6.7.} \emph{There are two sequences
$\{A_i(X_0,\ldots,X_i;Y)\}$ and
$\{A_i^{(q)}(X_0,\ldots,X_q;Y)\}$ of polynomials having the following
properties. Assume the conditions of \hyperlink{sga6.e13.para.6-1}{6.1} and suppose that $F$ is a
quotient of $\cO_P(-m)^{\oplus M}$, with $m\geq0$.}
\begin{enumerate}[label=(\roman*)]
\item \emph{If $F$ is a $(b_0,\ldots,b_r)$-sheaf, then, for
$i=0,\ldots,r$,}
\[
  |a_i|\leq A_{r-i}(b_i,\ldots,b_r;m).
\]
\item \emph{If
$\chi(F_q(n))=\sum_{i=0}^{r}a_i^{(q)}\binom{n+i}{i}$ is the Hilbert
polynomial of $F_q$, then, for $i=0,\ldots,q-1$,}
\[
  |a_i^{(q)}|\leq A_{r-i}^{(r-q)}(c_{q-1},\ldots,c_r;m).
\]
\end{enumerate}
\emph{(Notice that $a_{q-1}^{(q)}\leq a_{q-1}$ and
$a_q^{(q)}=a_q,\ldots,a_r^{(q)}=a_r$, since
$\dim\Supp(N_q)\leq q-1$.)}

To prove (i), use induction on $r$ and apply \hyperlink{sga6.e13.thm.6-4}{6.4} to an exact sequence
$0\longrightarrow F(-1)\longrightarrow F\longrightarrow G\longrightarrow0$
in which $G$ is a $(b_1,\ldots,b_r)$-sheaf. By induction on $r$, one may
suppose that $A_0,\ldots,A_{r-1}$ have been defined, taking \hyperlink{sga6.e13.lem.1-7}{1.7} into
account. Then \hyperlink{sga6.e13.thm.6-4}{6.4(ii)}, together with \hyperlink{sga6.e13.prop.1-8}{1.8}, displays the definition of
$A_r$. Finally, (ii) follows formally from \hyperlink{sga6.e13.thm.6-4}{6.4(iii)}, \hyperlink{sga6.e13.prop.6-6}{6.6}, and assertion
(i), applied to $F_q$ in place of $F$.

\sgaStableTarget{sga6.e13.cor.6-8}\textbf{Corollary 6.8. (Grothendieck).} \emph{Let $X$ be projective over
a noetherian scheme $S$, let $\cO_X(1)$ be very ample for $X/S$, and let
$\mathcal J$ be a family of classes of coherent sheaves on the fibers of
$X/S$. Suppose that:}
\begin{enumerate}[label=(\alph*)]
\item \emph{There is a coherent sheaf $E$ on $X$ such that $\mathcal J$
is contained in the family of classes of quotients of sheaves of the
form $E_K$.}
\item[$(b_q)$] \emph{In the Hilbert polynomials $\chi(F_K(n))$, for
$F_K\in\mathcal J$, the coefficients in degrees at least $q-1$ remain
bounded above.}
\end{enumerate}
\emph{Then the $(F_K)_q$ (as $F_K$ ranges over $\mathcal J$) form a
bounded family; moreover, the coefficients of $\chi(F_K(n))$ in degrees
at least $q-2$ remain bounded below.}

Indeed, one may plainly suppose that $X=\mathbf P^N_S$ and
$E=\cO_X(-m)^{\oplus M}$. The first assertion follows formally from\hyperlink{sga6.e13.thm.6-7}{
6.7(ii)} and 1.13. The second follows by induction, using an exact sequence
$0\longrightarrow F(-1)\longrightarrow F\longrightarrow G\longrightarrow0$
and \hyperlink{sga6.e13.thm.6-4}{6.4(ii)}.

\sgaStableTarget{sga6.e13.def.6-9}\textbf{Definition 6.9.} \emph{Let $k$ be a field. We call a
\emph{special positive $k$-cycle of dimension $r$} any projective
$k$-scheme $X$, equipped with a very ample sheaf $\cO_X(1)$, which is a
union of closed subschemes $X_j$ of dimension $r$, each $X_j$ being
obtained by base change $\Spec(k)\to\Spec(k_j)$ as the inverse image of
an integral $k_j$-scheme $X'_j$ (equipped with $\cO_{X'_j}(1)$). We call
the \emph{degree} of $X$ the coefficient $a_r$ in the Hilbert polynomial
$\chi(\cO_X(n))=\sum_i a_i\binom{n+i}{i}$.}

\sgaStableTarget{sga6.e13.lem.6-10}\textbf{Lemma 6.10.} \emph{Let $k$ be an algebraically closed field, and let $X$ be a special positive $k$-cycle of dimension $r$ and degree $d$. Then $\cO_X$ is a $(0,\ldots,0,d)$-sheaf.}

Indeed, view $X$ as a subscheme of $\mathbf P^N_k$. After replacing $k$ by the algebraic closure of a purely transcendental extension, which is permissible by \hyperlink{sga6.e13.prop.1-6}{1.6(i)}, one may suppose that $k$ contains $r(N+1)$ elements algebraically independent over all the $k_j$. The $r$ successive ``general'' hyperplane sections of $X$ defined by these elements are then always special positive $k$-cycles of degree $d$.

Arguing by induction on $r$, it remains only to observe that, when $r=0$,
one has $h^0(\cO_X(-1))=d$, while, when $r\geq1$, one has
$h^0(\cO_X(-1))=0$. For the latter point, the injection
$\cO_X\longrightarrow\coprod_j\cO_{X_j}$ allows one to suppose that $X$
itself comes from an integral scheme $X'$. But then
$h^0(\cO_X(-1))=h^0(\cO_{X'}(-1))=0$, because $H^0(\cO_{X'})$ is a
field.

\sgaStableTarget{sga6.e13.cor.6-11}\textbf{Corollary 6.11.} (i) \emph{Let $r,d$ be two integers, put
$c_i=A_{r-i}(0,\ldots,0,d;0)$ for $i=0,\ldots,r$ \hyperlink{sga6.e13.thm.6-7}{(6.7(i))}, and put
$p=P_r(c_0,\ldots,c_r)$ \hyperlink{sga6.e13.def.1-9}{(1.9)}. Let $k$ be an algebraically closed field,
and let $X$ be a special positive $k$-cycle of dimension at most $r$ and
degree at most $d$, with Hilbert polynomial
$\chi(\cO_X(n))=\sum_{i=0}^{r}a_i\binom{n+i}{i}$. Then, for
$i=0,\ldots,r$,}
\[
  |a_i|\leq c_i.
\]
\emph{Moreover, $X$ embeds as a subscheme of $\mathbf P^N_k$, with
$N=d(r+1)-1$, defined by (at most) $\binom{N+p}{N}$ equations of degree
$p$.}

(ii) (Chow) \emph{Let $S$ be a noetherian scheme. As $K$ ranges over
the algebraically closed extensions of the $k(s)$, for $s\in S$, the
special positive $K$-cycles of bounded dimension and degree form a
``bounded family'' (either with the $X$ regarded as subschemes of a fixed
$\mathbf P^N_S$, or in an evident abstract sense).}

Indeed, in (i), the first assertion follows from \hyperlink{sga6.e13.thm.6-7}{6.7(i)}; the second
follows from \hyperlink{sga6.e13.lem.6-10}{6.10} and \hyperlink{sga6.e13.prop.1-8}{1.8(i)}, and from \hyperlink{sga6.e13.lem.6-2}{6.2}, applied to
$0\longrightarrow I\longrightarrow\cO_P\longrightarrow\cO_X\longrightarrow0$,
$P=\mathbf P^N_k$, together with 1.3. Finally, (ii) follows from (i) and
1.13.

\sgaStableTarget{sga6.e13.rem.6-12}\textbf{Remark 6.12.} Let $k$ be an algebraically closed field and let
$X$ be a projective $k$-scheme, purely of dimension $r$, with no embedded
associated prime cycle. If arbitrary nilpotent elements are allowed in
$\cO_X$, the coefficients in the Hilbert polynomial
$\chi(\cO_X(n))=\sum_{i=0}^{r}a_i\binom{n+i}{i}$ can no longer be bounded
in terms of $a_r$ alone.

For example, let $Z$ be a smooth curve of degree $d$ in $\mathbf P^3_k$.
For every $n\gg0$, there is a smooth surface $Y$ of degree $n$ containing
$Z$. Then $Z$ is a Cartier divisor on $Y$; put
$\cO_{X_n}=\cO_Y/\cO_Y(-2Z)$. There is an exact sequence
$0\longrightarrow\cO_Y(-Z)\longrightarrow
\cO_Y\longrightarrow\cO_Z\longrightarrow0$ and the formula
$\langle Z^2\rangle_Y=2p_a(Z)-2-(n-4)d$. It follows that
$a_0=\chi(\cO_{X_n}(-1))\longrightarrow\infty$ as $n\longrightarrow
\infty$, although $a_1=2d$.

\sgaStableTarget{sga6.e13.cor.6-13}\textbf{Corollary 6.13.} \emph{Under the conditions of \hyperlink{sga6.e13.para.6-1}{6.1}, consider an exact sequence
\[
  \cO_P(-m_1)^{\oplus M_1}\longrightarrow
  \cO_P(-m)^{\oplus M}\xrightarrow{\alpha}F\longrightarrow0
\]
with $m_1\geq m\geq0$. Then, for $i=0,\ldots,N$,
\[
  |a_i|\leq A_{N-i}(0,\ldots,0,M;m_1)
  +M\binom{m}{N-i}.
\]}

Indeed, let $I=\Ker(\alpha)$. Plainly, $I$ is a
$(0,\ldots,0,M)$-sheaf, and
$\chi(I(n))=\allowbreak\sum_{i=0}^{N}
\bigl[M\tbinom{-m-1+N-i}{N-i}-a_i\bigr]\allowbreak\tbinom{n+i}{i}$.
The assertion follows from \hyperlink{sga6.e13.thm.6-7}{6.7(i)} and the formula
$\tbinom{-m-1+Q}{Q}=(-1)^Q\tbinom{m}{Q}$.

\sgaStableTarget{sga6.e13.cor.6-14}\textbf{Corollary 6.14. (I. Hermann).} \emph{For every $N\geq0$, there
is a polynomial $R_N(m)$ such that, for every field $k_0$ and every ideal
$I$ of the polynomial ring $k_0[T_1,\ldots,T_N]$ generated by elements
of degree at most $m$, the radical $\sqrt I$ is generated by elements of
degree at most $R_N(m)$.}

Indeed, let $k$ be the algebraic closure of $k_0$. After introducing an
auxiliary indeterminate, let $Y$ (respectively $X$) be the closed
subscheme of $P=\mathbf P^N_k$ defined by the homogenized ideal of $I$
(respectively of $\sqrt I$). In the Hilbert polynomial
$\chi(\cO_Y(n))=\sum_{i=0}^{N}b_i\binom{n+i}{i}$, one has
$|b_i|\leq c_i$, where $c_i=A_{N-i}(0,\ldots,0,1;m)$, by 6.13.

Let $X=\bigcup X^q$ be the decomposition by dimension. Then
$\deg(X^q)\leq e_q$, where $e_q=P_{N-q}(c_q,\ldots,c_N)$.
Indeed, after cutting $X$ (respectively $Y$) by a ``general'' linear subspace of codimension $q$, one may suppose that $q=0$. Then
$\deg(X^q)=h^0(\cO_{X^q})$ and
$h^0(\cO_{X^q})\leq h^0(\cO_Y)$, and \hyperlink{sga6.e13.thm.6-4}{6.4(iii)} gives
$h^0(\cO_Y)\leq e_q$.

Consequently, $\cO_{X^q}$ is a $(0,\ldots,0,e_q)$-sheaf by 6.10. It
follows from the injection $\cO_X\longrightarrow\prod_q\cO_{X^q}$ that
$\cO_X$ is an $(e_0,\ldots,e_N)$-sheaf. Hence, in the Hilbert polynomial
$\chi(\cO_X(n))=\sum_{i=0}^{N}a_i\binom{n+i}{i}$, one has
$|a_i|\leq f_i$, where $f_i=A_{N-i}(e_i,\ldots,e_N;0)$, by 6.7.
Finally, the ideal $\sqrt I$ of $X$ is $R_N(m)$-regular, where
$R_N(m)=P_N(f_0,\ldots,f_N)$, and the assertion follows from 1.3.

\sgaStableTarget{sga6.e13.sec.7}\subsection*{7. Appendix: The Hodge index theorem}

Here we recover and complete some earlier results by a direct argument.

\sgaStableTarget{sga6.e13.thm.7-1}\textbf{Theorem 7.1.} \emph{Let $X$ be an irreducible surface, proper
over an algebraically closed field. Suppose that there is an invertible
sheaf $H$ on $X$ such that $\langle c_1(H)^2\rangle>0$ (for example,
$H$ ample). Then an invertible sheaf $L$ on $X$ satisfying
\[
  \langle c_1(L)\cdot c_1(H)\rangle=0,
  \qquad
  \langle c_1(L)^2\rangle\geq0,
\]
is numerically equivalent to $0$. Consequently, by \hyperlink{sga6.e13.thm.4-6}{4.6}, the intersection
form on $A^1_{\mathrm{num}}(X)_{\Q}
=\bigl[\Pic(X)/\Pic^\tau(X)\bigr]\otimes\Q$ is nondegenerate of type
$(1,\rho-1)$.}

\sgaStableTarget{sga6.e13.para.7-1-1}\textbf{7.1.1.} First suppose that $X$ is reduced (hence integral), that
$H$ is ample, that $\langle c_1(L)\cdot c_1(H)\rangle=0$, and that
$\langle c_1(L)^2\rangle>0$. Arguing by contradiction, choose $m>0$
such that $H_1=L\otimes H^{\otimes m}$ is ample, and choose $n$ such
that, on putting $N=L^{\otimes n}\otimes H^{-1}$, one has
$\langle c_1(N)\cdot c_1(H_1)\rangle
=n\langle c_1(L)^2\rangle-m\langle c_1(H)^2\rangle>0$. One also has
$\langle c_1(N)\cdot c_1(H)\rangle=-\langle c_1(H)^2\rangle<0$ and
$\langle c_1(N)^2\rangle
=n^2\langle c_1(L)^2\rangle+\langle c_1(H)^2\rangle>0$, contradicting
the following lemma.

\sgaStableTarget{sga6.e13.lem.7-1-2}\textbf{Lemma 7.1.2.} \emph{Let $X$ be an integral surface, projective over an algebraically closed field. For an invertible sheaf $N$ on $X$ such that $\langle c_1(N)^2\rangle>0$, the following conditions are equivalent:
\begin{enumerate}[label=(\roman*)]
\item for every ample sheaf $H$, one has $\langle c_1(N)\cdot c_1(H)\rangle>0$;
\item[(i$'$)] there is an ample sheaf $H$ such that $\langle c_1(N)\cdot c_1(H)\rangle>0$;
\item there is an integer $n>0$ such that $H^0(X,N^{\otimes n})\ne0$.
\end{enumerate}}

Indeed, (ii) implies that there is an effective divisor $D$ such that
$N^{\otimes n}\simeq\cO_X(D)$, and $\langle c_1(N)^2\rangle>0$ implies
$D>0$; hence (i).

To see that (i$'$) implies (ii), one may plainly suppose that
$H=\cO_X(1)$ is very ample. There is then a closed integral locally
principal curve $Y$ such that $H=\cO_X(Y)$. By the Riemann theorem,
(i$'$) implies $H^1\bigl(N_Y^{\otimes n}(p)\bigr)=0$ for $p\geq0$ and
$n\geq n_0$, where $n_0=2h^1(\cO_Y)-1$.
Since $H^2(N^{\otimes n}(p))=0$ for $p\gg0$, the exact sequence
\[
  H^1\bigl(N_Y^{\otimes n}(p)\bigr)\longrightarrow
  H^2\bigl(N^{\otimes n}(p-1)\bigr)\longrightarrow
  H^2\bigl(N^{\otimes n}(p)\bigr)\longrightarrow0
\]
gives $H^2(N^{\otimes n})=0$ for $n\geq n_0$. Thus, for $n\gg0$,
\[
  h^0(N^{\otimes n})
  =\chi(N^{\otimes n})+h^1(N^{\otimes n})
  \geq\chi(N^{\otimes n})
  =\frac12\langle c_1(N)^2\rangle n^2+\cdots>0,
\]
which completes the proof of the lemma.

\sgaStableTarget{sga6.e13.para.7-1-3}\textbf{7.1.3.} Now suppose in \hyperlink{sga6.e13.thm.7-1}{7.1} that there is an invertible sheaf $M$
on $X$ such that $\langle c_1(L)\cdot c_1(M)\rangle\ne0$, and argue by
contradiction. Replace $X$ by the underlying reduced scheme of a Chow
presentation of $X$, which is permissible by the projection formula \hyperlink{sga6.e10.eq.4-4-4}{X
4.4.4}, and let $G$ be ample on $X$. Choose $m>0$ such that
$H_1=G^{\otimes m}\otimes M$ is ample and
$\langle c_1(L)\cdot c_1(H_1)\rangle\ne0$. Choose $p,q$, with $q\ne0$,
such that, on putting $L_1=L^{\otimes p}\otimes H^{\otimes q}$, one has
$\langle c_1(L_1)\cdot c_1(H_1)\rangle
=p\langle c_1(L)\cdot c_1(H_1)\rangle
+q\langle c_1(H)\cdot c_1(H_1)\rangle=0$. Then one also has
$\langle c_1(L_1)^2\rangle=p^2\langle c_1(L)^2\rangle
+q^2\langle c_1(H)^2\rangle>0$. This is the case \hyperlink{sga6.e13.para.7-1-1}{7.1.1}, with $L=L_1$
and $H=H_1$.

\sgaStableTarget{sga6.e13.para.7-1-4}\textbf{7.1.4.} Finally, consider the general case of 7.1. Still arguing by contradiction, suppose that $L$ is not numerically equivalent to $0$. Replace $X$ by the normalization of $X_{\red}$, so that $X$ has

only finitely many singular points; this reduction is justified by the projection formula \hyperlink{sga6.e10.eq.4-4-4}{X 4.4.4} and by \hyperlink{sga6.e13.prop.4-5}{4.5(ii)}. There is therefore a closed integral curve $Y$ on $X$ such that
$\langle c_1(L)\cdot Y\rangle\ne0$, and $Y$ is locally principal
outside a finite set $Z$. Let $f:X'\to X$ be the blow-up of the ideal
$I\subset\cO_X$ defining $Y$. The ideal $\mathcal M=I\cO_{X'}$ is
invertible and defines a subscheme of $X'$ which is the union of two
subschemes $Y'$ and $Y''$ such that $f(Y')=Y$ and $f(Y'')\subset Z$.
By the projection formula \hyperlink{sga6.e10.eq.4-4-4}{X 4.4.4},
$\langle c_1(f^*L)\cdot c_1(\mathcal M)\rangle
=-\langle c_1(f^*L)\cdot Y'\rangle
-\langle c_1(f^*L)\cdot Y''\rangle
=-\langle c_1(L)\cdot Y\rangle-0\ne0$. Thus one may replace $X,L,H$
by $X',f^*L,f^*H$. This reduces the problem to the preceding case and
completes the proof of 7.1.

\sgaStableTarget{sga6.e13.rem.7-2}\textbf{Remark 7.2.} It is not known whether an irreducible proper surface over an algebraically closed field which possesses an invertible sheaf $H$ satisfying $\langle c_1(H)^2\rangle>0$ is necessarily projective. Nor is it known whether, for every irreducible proper surface, the intersection form on $A^1_{\mathrm{num}}(X)_{\Q}$ is nondegenerate. It follows from \hyperlink{sga6.e13.thm.7-1}{7.1} that, when it is degenerate, it is necessarily negative (non-definite).

\sgaStableTarget{sga6.e13.cor.7-3}\textbf{Corollary 7.3.} \emph{Under the conditions of \hyperlink{sga6.e13.thm.7-1}{7.1}, let $M$ be
an invertible sheaf on $X$ which is not numerically proportional to $H$,
and put $d=\langle c_1(M)\cdot c_1(H)\rangle$ and
$h=\langle c_1(H)^2\rangle$. Then
$\langle c_1(M)^2\rangle<d^2/h$.}

Indeed, let $L=M^{\otimes h}\otimes H^{\otimes(-d)}$. Since
$\langle c_1(L)\cdot c_1(H)\rangle=0$, \hyperlink{sga6.e13.thm.7-1}{Theorem 7.1} gives
$\langle c_1(L)^2\rangle<0$, which proves the assertion.

\sgaStableTarget{sga6.e13.cor.7-4}\textbf{Corollary 7.4.} \emph{Let $X$ be an irreducible scheme of
dimension $r\geq2$, projective over an algebraically closed field, and
let $H$ be an ample sheaf on $X$. Then:}
\begin{enumerate}[label=(\roman*)]
\item \emph{An invertible sheaf $L$ on $X$ is numerically equivalent to
$0$ if and only if}
\[
  \langle c_1(L)\cdot c_1(H)^{r-1}\rangle=0,
  \qquad
  \langle c_1(L)^2\cdot c_1(H)^{r-2}\rangle\geq0;
\]
\item \emph{For every invertible sheaf $M$ on $X$ which is not
numerically equivalent to $H$,}
\[
  \langle c_1(M)^2\cdot c_1(H)^{r-2}\rangle
  <\frac{\langle c_1(M)\cdot c_1(H)^{r-1}\rangle^2}
  {\langle c_1(H)^r\rangle}.
\]
\end{enumerate}

Assertion (ii) follows immediately from (i), by the argument of 7.3. To
prove (i), let $Y$ be a closed integral curve in $X$; we wish to prove
that $\langle c_1(L)\cdot Y\rangle=0$. Embed $X$ in $\mathbf P^N_k$ by
means of $H^{\otimes n}$, with $n>0$. Let $I=\bigoplus I_\nu$ be the
homogeneous ideal of $Y$, and let $\nu_0$ be an integer such that
$I_{\nu_0}$ generates $I_\nu$ for $\nu\geq\nu_0$. Then the
hypersurfaces $Z$ of degree $\nu_0$ containing $Y$ have no common points
outside $Y$. By induction on $r=\dim X$ and Bertini's theorem, one finds
$r-2$ hypersurfaces $Z_1,\ldots,Z_{r-2}$ of degree $\nu_0$ containing
$Y$ such that $X'=Z_1\cap\cdots\cap Z_{r-2}\cap X$ is irreducible.
Since $\langle c_1(L)\cdot Y\rangle
=\langle c_1(L|_{X'})\cdot Y\rangle$, it is enough to prove that
$L|_{X'}$ is numerically equivalent to $0$ on $X'$. By the projection
formula, $\langle c_1(L|_{X'})\cdot c_1(H|_{X'})\rangle=0$ and
$\langle c_1(L|_{X'})^2\rangle\geq0$, so the assertion follows from 7.1.

\sgaStableTarget{sga6.e13.cor.7-5}\textbf{Corollary 7.5.} \emph{Let $k$ be an algebraically closed field,
let $Y$ be a projective $k$-scheme equipped with an ample sheaf $H$, and
let $X$ be the zero scheme of a section $\varphi$ of $H$. Suppose that,
for every irreducible component $Y^i$ of $Y$, the irreducible components
of $X\cap Y^i=V(\varphi|_{Y^i})$
have dimension at least $2$. Then the morphism
$A^1_{\mathrm{num}}(Y)\longrightarrow A^1_{\mathrm{num}}(X)$ is
injective.}

Indeed, one may plainly suppose that $Y$ is integral, and the assertion follows from \hyperlink{sga6.e13.cor.7-4}{7.4(i)} by the projection formula.

\sgaStableTarget{sga6.e13.bibliography}\subsection*{Bibliography}
\begin{enumerate}[label={[\arabic*]}]
\item S. Abhyankar, ``Resolution of singularities of arithmetical surfaces,'' in \emph{Arithmetical Algebraic Geometry}, Purdue University, Harper \& Row, New York, 1965.
\item D. Mumford, ``Pathologies III,'' \emph{Amer. J. Math.}, 1967, pp. 94--103. \emph{Amer. J. Math.}, 1967, pp. 94--103.
\item A. Grothendieck and J. Dieudonn\'e, \emph{EGA, \'El\'ements de G\'eom\'etrie Alg\'ebrique}, Publ. Math. Inst. des Hautes \'Etudes Sci., 1960.
\item A. Grothendieck, \emph{FGA, Fondements de la g\'eom\'etrie alg\'ebrique}, extracts from the S\'eminaire Bourbaki 1957--1962.
\item A. Grothendieck et al., \emph{SGA, S\'eminaire de g\'eom\'etrie alg\'ebrique}, 1960--66, I.H.E.S.
\item O. Jussila, ``Formalisme des intersections sur les sch\'emas alg\'ebriques propres,'' Expos\'e X in this seminar.
\item M. Raynaud, ``Un th\'eor\`eme de repr\'esentabilit\'e relative sur le foncteur de Picard,'' Expos\'e XII in this seminar.
\end{enumerate}

\sgaStableTarget{sga6.e14.expose}\section*{Expos\'e XIV. Open problems in intersection theory}
\begin{center}
by A. Grothendieck
\end{center}

In the present expos\'e we review some of the problems connected with questions encountered in this seminar which deserve to be studied. We leave aside the problems connected with Serre's \emph{local intersection theory}~[14], which we have not had to use in the seminar, and refer in particular to \emph{loc. cit.} for the two fundamental conjectures which remain open in that direction (\emph{loc. cit.}, Chapter V, \S4). We also pass over the problems, very important for arithmetic applications, connected with the \emph{structure of groups of cycle classes}, especially on a smooth projective algebraic scheme. For discussion of these problems, we refer to S. Kleiman, ``Algebraic cycles and the Weil conjectures,'' in \emph{Dix expos\'es sur la cohomologie des sch\'emas} (Masson--North Holland Publishing Co.). In no. 4 of the present expos\'e we also give some supplements to the seminar, concerning in particular the \emph{Chow ring}, which had not found a place in the preceding expos\'es.

\sgaStableTarget{sga6.e14.sec.1}\section*{1. The operations $\Lambda^i$ in the derived category $D(X)$\footnote{Cf. the note on p. 23.}}

Let $X$ be a locally ringed topos, and let $K^\bullet(X)$ be the ring of classes of perfect complexes on $X$ (\hyperlink{sga6.e04.expose}{Expos\'e IV}). One would like to define on $K^\bullet(X)$ a $\lambda$-ring structure (\hyperlink{sga6.e05.expose}{Expos\'e V}), functorial in $X$, such that the natural homomorphism
\[
  K^\circ_{\mathrm{naive}}(X)\longrightarrow K^\bullet(X)
\]
(where the first member is the ring of classes of locally free modules on $X$, equipped with the $\lambda$-structure defined in \hyperlink{sga6.e05.expose}{Expos\'e V}) is a $\lambda$-homomorphism. For this purpose, one would have to define functors, in general nonadditive, of ``exterior power''
\[
  \Lambda^i:D^-(X)\longrightarrow D^-(X),
\]
transforming pseudo-coherent (respectively perfect) complexes into pseudo-coherent (respectively perfect) complexes and having the formal properties with respect to exact triangles needed to define a $\lambda$-structure on $K^\bullet(X)$ by the formula
\sgaStableTarget{sga6.e14.eq.1-1}\begin{equation}\tag{1.1}
  \lambda^i\bigl(\cl(L^\bullet)\bigr)
  =\cl\bigl(\Lambda^iL^\bullet\bigr).
\end{equation}

When $X$ is of characteristic zero, that is, when its structure sheaf is a sheaf of $\Q$-algebras, the problem is solved simply by putting
\sgaStableTarget{sga6.e14.eq.1-2}\begin{equation}\tag{1.2}
  \Lambda^iL^\bullet
  =\left({}^{\otimes i}L^\bullet\right)^{\mathrm{alt}},
\end{equation}
where the $i$-th tensor power, in the sense of derived categories, of
$L^\bullet$ is regarded as an object of $D^-(\mathcal O_i)$, with
$\mathcal O_i=\mathcal O_X[\mathfrak S_i]$ the algebra of the symmetric
group $\mathfrak S_i$ over $\mathcal O_X$, and where the superscript
\(\mathrm{alt}\) denotes passage to the ``alternating part under the
action of $\mathfrak S_i$,'' an operation which makes sense because of
the characteristic-zero hypothesis. A formula of this type was already
used in [8] (cited below as [RRR]) in the first proof of the
Riemann--Roch theorem.

Without a characteristic hypothesis, when $L^\bullet$ is zero in degrees $>0$, Dold--Puppe~[6] gives a canonical homotopy from $L^\bullet$ to the complex associated with a semisimplicial module $L^{\prime\bullet}$ canonically attached to $L^\bullet$. For $L^\bullet$ flat, say, a fairly natural candidate for $\Lambda^iL^\bullet$ would be the complex associated with the semisimplicial module $\Lambda^iL^{\prime\bullet}$, whose degree-$n$ component is $\Lambda^iL'_n$. But to extend the definition to complexes of $D^-(X)$ which are no longer assumed to be zero in degrees $>0$, one would have to study the behavior of the preceding operation with respect to translation $L^\bullet\mapsto L^\bullet[1]$, or \emph{suspension} in Dold--Puppe terminology~[6]. Unfortunately, there does not seem to be an isomorphism $\Lambda^i\bigl(L^\bullet[1]\bigr)\simeq\Lambda^i(L^\bullet)[i]$ which would make the desired extension trivial.

\sgaStableTarget{sga6.e14.sec.2}\section*{2. The Riemann--Roch formula without projectivity hypotheses}

\sgaStableTarget{sga6.e14.para.2-1}\textbf{2.1.} Once one has a natural $\lambda$-structure on the $K^\bullet(X)$, it is possible to formulate the Riemann--Roch theorem, in the form contemplated in \hyperlink{sga6.e08.expose}{Expos\'e VIII}, for any proper complete-intersection morphism $f:X\to Y$, and even, more generally, for a relative scheme $X$ over a locally ringed topos $Y$, with $X$ still subject to the corresponding relative conditions. The Todd class of $f$ is defined by means of the class in $K^\bullet(X)$ of the relative cotangent complex $T_f$ constructed in \hyperlink{sga6.e00.sec.4.4}{Expos\'e 0, 4.4}; that complex is perfect here by the local hypothesis on $f$. It seems clear that a proof of the Riemann--Roch formula in this general case will require essentially new ideas, and it is possible that these same ideas could provide the key

to a proof of the analogous formula in the complex analytic (or rigid-analytic!) case. Recall that, in the complex analytic case, only the case of a locally free sheaf $F$ on a compact nonsingular variety $X$, over a one-point $Y$, is currently known, by the Atiyah--Singer index formula; the case in which $F$ is only assumed coherent still seems to escape presently known techniques.

\sgaStableTarget{sga6.e14.para.2-2}\textbf{2.2.} As an essential ingredient of the Riemann--Roch formula,
recall also the \emph{finiteness problem} raised in \hyperlink{sga6.e03.expose}{Expos\'e III}: if
$f:X\to Y$ is a proper morphism and $L^\bullet$ is a complex on $X$
which is pseudo-coherent relative to $Y$, is $Rf_*(L^\bullet)$
pseudo-coherent? A general proof of this conjecture will probably use
algebraic ideas which will yield a proof of Grauert's ``finiteness
theorem''~[7] under appreciably more general conditions than those of
\emph{loc. cit.}, working from the outset over topologically ringed
spaces (or topoi) in $\mathbb C$-algebras $Y$ more general than analytic
spaces. This would include the cases in which $Y$ is a locally compact
space with the sheaf of continuous complex-valued functions, a
$C^m$-manifold with the sheaf of complex-valued $C^m$-functions, a real
analytic space, and so forth.

\sgaStableTarget{sga6.e14.para.2-3}\textbf{2.3.} Finally, in these questions the properness hypothesis on $f$ should be replaced by the hypothesis that the supports of the cohomology sheaves $H^i(L^\bullet)$ of the complexes under consideration on $X$ are proper over $Y$.

\sgaStableTarget{sga6.e14.sec.3}\section*{3. A Riemann--Roch formula ``without denominators'' for an immersion}

\sgaStableTarget{sga6.e14.para.3-1}\textbf{3.1.} Let $f:Y\to X$ be a regular closed immersion (\hyperlink{sga6.e07.expose}{Expos\'e VII}), with conormal sheaf $\underline N$. For simplicity, suppose that $X$ admits an ample invertible sheaf. The Riemann--Roch formula for $y\in K^\bullet(Y)$,
\[
  \ch_X(f_*(y))
  =f_*\bigl(\ch_Y(y)\Todd(-\check{\underline N})\bigr),
\]
together with the general formula (\hyperlink{sga6.e07.expose}{Expos\'e VII})
\[
  f_*(y')f_*(y'')
  =f_*\bigl(y'y''c_d(\underline N)\bigr),
  \qquad y',y''\in\Gr^\bullet(Y)_{\Q},
\]
allows one, assuming $f$ has codimension $d$, to express the Chern classes of $f_*(y)$ in $\Gr^\bullet(X)_{\Q}$ in the form

\sgaStableTarget{sga6.e14.eq.3-1}\begin{equation}\tag{3.1}
  c_i(f_*(y))
  =f_*\bigl(P_i((c_\alpha(y))_\alpha,(c_\beta(\underline N))_\beta)\bigr),
\end{equation}
where $P_i$ is a universal polynomial in the $c_\alpha(y)$ and the $c_\beta(\underline N)$, a polynomial which a priori has rational coefficients. In fact, one finds in [RRR] that these polynomials have \emph{integer coefficients}. This raises the question of the validity of a Riemann--Roch formula ``without denominators.''

In fact, such a formula for $c_i(f_*(y))$ cannot in general be sought in $\Gr^\bullet(X)$ itself, since the homomorphism $f_*:K^\bullet(Y)\longrightarrow K^\bullet(X)$ is not itself compatible with the $\lambda$-filtrations, modulo a shift by $d$ \hyperlink{sga6.e14.para.4-6}{(4.6)}, but is compatible only up to torsion. Thus it is not possible to define a homomorphism $f_*:\Gr^\bullet(Y)\longrightarrow\Gr^\bullet(X)$ inducing the natural homomorphism $f_*:\Gr^\bullet(Y)_{\Q}\longrightarrow\Gr^\bullet(X)_{\Q}$ and making it possible to interpret \hyperlink{sga6.e14.eq.3-1}{(3.1)} as a formula in $\Gr^\bullet(X)$ rather than only in $\Gr^\bullet(X)_{\Q}$.

When $X$ is regular, however, so that $K^\bullet(X)\simeq K_\bullet(X)$ (\hyperlink{sga6.e04.para.2-5}{IV 2.5}), one can equip $K^\bullet(X)$ with the \emph{topological filtration} by codimension of supports, inherited from the corresponding filtration of $K_\bullet(X)$. It is plausible that this filtration is compatible with the ring structure. This is proved at least in [RRR] when $X$ is smooth over a field $k$, using the theory of the Chow ring [4]. When this is so, the associated graded
\[
  \Gtop^\bullet(X)
\]
is a \emph{graded ring}; and, if $Y$ is also regular, one has precisely what is needed to obtain a natural homomorphism of degree $d$ of graded groups
\[
  f_*=f_*^{\mathrm{top}}:\Gtop^\bullet(Y)
  \longrightarrow\Gtop^\bullet(X).
\]
Since the $\lambda$-filtration of $K^\bullet(X)$ is finer than the topological filtration by codimension just considered (\hyperlink{sga6.e10.cor.1-3-3}{X 1.3.3}), one obtains a canonical homomorphism of graded groups (and even of graded rings),
\[
  \Gr^\bullet(X)\longrightarrow\Gtop^\bullet(X),
\]
which is, moreover, an isomorphism modulo torsion (Expos\'e VIII). This makes it possible, as in [RRR], to consider a theory of Chern classes on $X$ with values in $\Gtop^\bullet(X)$, and likewise on $Y$, and thus to interpret \hyperlink{sga6.e14.eq.3-1}{(3.1)} as a formula in $\Gtop^\bullet(X)$. In this form, the formula is proved in [RRR] when $X$ is of \emph{characteristic zero}, at least if $X$ is smooth over a field (an assumption used only to ensure that the topological filtration of $K^\bullet(X)$ is compatible with the ring structure). The characteristic-zero hypothesis is used essential%
ly to have formulas of the type \hyperlink{sga6.e14.eq.1-1}{(1.1)}, \hyperlink{sga6.e14.eq.1-2}{(1.2)}, and essentially the same proof should apply without a characteristic hypothesis once the question raised in no. 1 has been answered affirmatively.

\sgaStableTarget{sga6.e14.para.3-2}\textbf{3.2.} When $X$ is smooth over a field, one also has the Chow ring
\[
  A(X),
\]
and a theory of Chern classes on $X$ with values in $A(X)$~[9], finer than the theory with values in $\Gtop^\bullet(X)$ just introduced. The latter Chern classes are the images of the Chow-style Chern classes under the canonical surjective homomorphism (cf. \hyperlink{sga6.e14.para.4-1}{4.1})
\[
  \varphi:A(X)\longrightarrow\Gtop^\bullet(X).
\]
This homomorphism is an isomorphism modulo torsion, but not an isomorphism in general \hyperlink{sga6.e14.para.4-7}{(4.7)}. That said, if $X$ and $Y$ are smooth, the inclusion $f:Y\to X$ still defines a homomorphism of degree $d$ of graded groups
\[
  f_*:A(Y)\longrightarrow A(X),
\]
which makes it possible to interpret \hyperlink{sga6.e14.eq.3-1}{(3.1)} as a formula in $A^i(X)$.
This formula reduces to $0=0$ for $i<d$ and is therefore of interest
only for $i\geq d$. Even for $i=d=2$, this formula has not been proved at
present (cf. \hyperlink{sga6.e14.para.4-3}{4.3}), nor has the less precise formula deduced from it by
replacing rational equivalence of algebraic cycles by algebraic
equivalence. Unlike the analogous formula in $\Gtop^\bullet(X)$, there
does not seem to be any convincing reason to suppose that the formula
without denominators should be true in $A(X)$. A good way to test it,
given our almost total ignorance of the nature of the groups
$A^i(X)$\footnote{The formula, as well as the variant mentioned in \hyperlink{sga6.e14.para.3-3}{3.3},
has since been proved by P. Jouanolou (Nov. 1969). Cf.} for $i>2$, seems
to be to test the formulas that can be deduced from it in groups
$A^1(Z)=\Pic(Z)$, for example by applying to both sides of \hyperlink{sga6.e14.eq.3-1}{(3.1)}
$g_*:A(X)\longrightarrow A(Z)$,
where $g:X\to Z$ is a proper morphism to a smooth quasi-projective $Z$ of dimension $\dim X-i+1$.

\sgaStableTarget{sga6.e14.para.3-3}\textbf{3.3.} Finally, to give a meaning to \hyperlink{sga6.e14.eq.3-1}{(3.1)}, one may also use the cohomological theory of Chern classes (SGA 5 VII), whenever one has \emph{Gysin homomorphisms}
\[
  f_*:H^p(Y,\mu_n^{\otimes i})
  \longrightarrow H^{p+2d}(X,\mu_n^{\otimes(i+d)}).
\]
This holds at least when $Y$ and $X$ are both smooth over a scheme $S$, by the
theorem of relative cohomological purity (SGA 4 XVI 3.7); conjecturally, it
also holds when $Y$ and $X$ are both regular, by the theorem of absolute
cohomological purity, proved at least for excellent $X$ of characteristic
zero (SGA 4 XIX 3.2).

When $X$ and $Y$ are smooth over the field of complex numbers, one may also
use the integral Chern class on the locally compact space $X(\C)$ of complex
points of $X$. Formula \hyperlink{sga6.e14.eq.3-1}{(3.1)} was proved in this case by Atiyah and
Hirzebruch~[1], who more generally treat embeddings of complex analytic
varieties (working with naive $K'$). It follows, by the Lefschetz principle
and the general results of SGA 4, that the $\ell$-adic formula (3.1) holds
whenever $Y$ and $X$ are smooth over an algebraically closed field of
characteristic zero. This makes the cohomological formula \hyperlink{sga6.e14.eq.3-1}{(3.1)} extremely
plausible. Nevertheless, it should be noted that, at present, it has not been
proved even when $X$ and $Y$ are smooth over a field $k$ of characteristic
zero, if $k$ is not algebraically closed.

\noindent\sgaStableTarget{sga6.e14.para.3-4}\textbf{3.4.}
One may also work with the cohomological theory of Chern classes, in the de
Rham or Hodge style~[13]. The question of the validity of \hyperlink{sga6.e14.eq.3-1}{(3.1)} in this
setting arises only when $X$ is not of characteristic zero, since one knows
in any case (no.~6) that the formula in question is true modulo torsion. For
smooth quasi-projective schemes over a field $k$, the de Rham or Hodge
version of \hyperlink{sga6.e14.eq.3-1}{(3.1)} would moreover follow from the version using cycle classes
modulo algebraic equivalence. Thus the de Rham or Hodge cohomological version
provides another accessible way to test the validity of the cycle-class
version of \hyperlink{sga6.e14.eq.3-1}{(3.1)}.

\noindent\sgaStableTarget{sga6.e14.para.3-5}\textbf{3.5.}
The proof~[1] of the transcendental cohomological version of \hyperlink{sga6.e14.eq.3-1}{(3.1)} is
essentially transcendental in nature and uses tubular neighborhoods of $Y$
in $X$. This suggests that a proof of a formula of type \hyperlink{sga6.e14.eq.3-1}{(3.1)} may be tied to
the systematic consideration of $K^\bullet$-groups and cohomology groups with
support in $Y$, which might provide an adequate technical means of
\emph{localization around $Y$}.

\sgaStableTarget{sga6.e14.sec.4}\section*{4. Relations between $K^\bullet(X)$ and the Chow ring $A(X)$}

\noindent\sgaStableTarget{sga6.e14.para.4-1}\textbf{4.1.}
Suppose for simplicity that $X$ is a smooth quasi-projective scheme over a
field $k$ (see no.~8 for a more general case), so that the Chow ring $A(X)$
of cycle classes modulo rational equivalence is available~[4]. Two variants
of this ring, defined under appreciably more general conditions, are
$\operatorname{Gr}^\bullet(X)$, the graded ring associated with the
$\lambda$-filtration of $K^\bullet(X)$, and $\Gtop^\bullet(X)$, the graded
ring associated with the filtration of
$K^\bullet(X)\simeq K_\bullet(X)$ by codimension of supports. There are
theories of Chern classes with values in each of these three rings (see [9]
for the case of $A(X)$, and this seminar for
$\operatorname{Gr}^\bullet(X)$, which maps into $\Gtop^\bullet(X)$). There is
also an $\ell$-adic cohomological theory (SGA 5 VII). Thus, when $\ell$ is
prime to the characteristic of $k$,
\[
 c_i(E)\in H^{2i}(X,\Zell(i)),
\]
where $E$ is a vector bundle on $X$, or more generally any element of
$K^\bullet(X)$. These theories are related by the following canonical
homomorphisms, which are compatible with the multiplicative structures:
\sgaStableTarget{sga6.e14.eq.4-1}\begin{equation*}\tag{4.1}
\begin{tikzcd}[column sep=large,row sep=large]
 H^{2*}(X,\Zell(*)) & {} \\
 A(X) \arrow[u,"\gamma"] \arrow[r,"\varphi"] & \Gtop^\bullet(X) \\
 {} & \operatorname{Gr}^\bullet(X) \arrow[u,"j"']
\end{tikzcd}
\end{equation*}
where
\[
 H^{2*}(X,\Zell(*))=\bigoplus_i H^{2i}(X,\Zell(i)).
\]

The homomorphism $j$ comes from the fact that the $\lambda$-filtration of
$K^\bullet(X)$ is finer than the topological filtration by codimension of
support (\hyperlink{sga6.e10.expose}{Expose X}), and it is an isomorphism modulo torsion (Expose VII
4.11). The homomorphism $\varphi$, defined for example in [RRR], associates
to the class of the codimension-$i$ cycle carried by a closed subscheme
$Y\subset X$ the element $\cl^i(\mathcal O_Y)$ of
\[
 \Filt_{\mathrm{top}}^i K^\bullet(X)/
 \Filt_{\mathrm{top}}^{i+1}K^\bullet(X).
\]
It is surjective, and we shall see in \hyperlink{sga6.e14.para.4-2}{4.2} that it is an isomorphism modulo
torsion. Hence, modulo torsion, the three rings $A(X)$,
$\Gtop^\bullet(X)$, and $\operatorname{Gr}^\bullet(X)$ are canonically
isomorphic; consequently, the Riemann--Roch formulas expressed using any
one of these rings after tensoring with $\Q$ are equivalent.

The graded ring $\operatorname{Gr}^\bullet(X)$, used in this seminar chiefly
as a substitute for the Chow ring, has the considerable advantage of being
defined for every locally ringed topos, and in particular for every scheme,
whether regular or not. Its disadvantage \hyperlink{sga6.e14.para.4-6}{(4.6)} with respect to $A(X)$ and
$\Gtop^\bullet(X)$ is that it is not covariant as it stands for proper
morphisms, even between smooth quasi-projective schemes over a field; it
acquires that covariance only after tensoring with $\Q$.

Another advantage of $A(X)$ over its competitors $\Gtop^\bullet(X)$ and
$\operatorname{Gr}^\bullet(X)$ is the homomorphism
\[
 A(X)\longrightarrow H^{2*}(X,\Zell(*)),
\]
which permits cohomological Chern classes to be interpreted as images of
Chern classes with values in $A(X)$. In general this homomorphism does not
factor through $\Gtop^\bullet(X)$ \hyperlink{sga6.e14.para.4-7}{(4.7)}, and it is at least doubtful that
there is in general a homomorphism
\[
 \operatorname{Gr}^\bullet(X)\longrightarrow H^{2*}(X,\Zell(*))
\]
carrying Chern classes to Chern classes (see no.~5). These difficulties
disappear after tensoring with $\Q$, and one obtains a homomorphism
\sgaStableTarget{sga6.e14.eq.4-2}\begin{equation*}\tag{4.2}
 \Gtop^\bullet(X)\longrightarrow H^{2*}(X,\Qell(*))
 =H^{2*}(X,\Zell(*))\otimes_{\Zell}\Qell
\end{equation*}
that carries Chern classes to Chern classes.

\noindent\sgaStableTarget{sga6.e14.para.4-2}\textbf{4.2.}
To show that the surjective homomorphism $\varphi$ in \hyperlink{sga6.e14.eq.4-1}{(4.1)} is an
isomorphism modulo torsion, one uses the Chern homomorphism obtained from the
theory of Chern classes~[9],
\[
 K^\bullet(X)\longrightarrow \Ch(A(X)),
\]
where the target is the Chern ring of $A(X)$ as defined in \hyperlink{sga6.e05.expose}{Expose V}. This
homomorphism is compatible with the topological filtration of $K^\bullet(X)$
and the natural filtration of $\Ch(A(X))$, and therefore induces a
homomorphism of graded groups
\sgaStableTarget{sga6.e14.eq.4-3}\begin{equation*}\tag{4.3}
 \psi:\Gtop^\bullet(X)\longrightarrow A(X).
\end{equation*}
One proves the formulas
\sgaStableTarget{sga6.e14.eq.4-4}\begin{equation*}\tag{4.4}
\left\{
\begin{aligned}
 \psi\varphi(x)&=(-1)^{i-1}(i-1)!\,x
   &&\pmod{\text{torsion}},\quad x\in A^i(X),\\
 \varphi\psi(x)&=(-1)^{i-1}(i-1)!\,x
   &&\pmod{\text{torsion}},\quad x\in \Gtop^i(X).
\end{aligned}
\right.
\end{equation*}
The first formula implies the second because $\varphi$ is surjective. These
formulas show that $\varphi$ and $\psi$ are isomorphisms modulo torsion and,
up to the factor $(-1)^{i-1}(i-1)!$ in degree $i$, are inverse to one
another.

To prove the first formula in \hyperlink{sga6.e14.eq.4-4}{(4.4)}, one may suppose that $x$ is represented
by an irreducible closed subscheme $Y$ of $X$ of codimension $i$. The formula
is then equivalent to
\sgaStableTarget{sga6.e14.eq.4-5}\begin{equation*}\tag{4.5}
 c_i(\mathcal O_Y)=(-1)^{i-1}(i-1)!\,\cl(Y)
   \quad\pmod{\text{torsion}}.
\end{equation*}
On the left, $\mathcal O_Y$ is regarded as a coherent sheaf on $X$ and hence
defines an element of $K^\bullet(X)$; on the right, $\cl(Y)$ denotes the
class of $Y$ in $A^i(X)$. Since $A^i(X)$ is unchanged upon deleting from
$X$ a closed subset of codimension $>i$, one reduces to the case where $Y$
is regular, and therefore smooth over the base field $k$ if $k$ is perfect.
In that case, \hyperlink{sga6.e14.eq.4-5}{(4.5)} is precisely the Riemann--Roch formula for the inclusion
$Y\to X$ and the element $1\in K^\bullet(Y)$, written using Chern classes
with values in the Chow ring $A(X)$. Under the present hypotheses, this
formula is proved in [2] by a method very close to the one used in this
seminar for the analogous formula with values in $\Gtop^\bullet(X)$.

Thus, at least when $k$ is perfect, \hyperlink{sga6.e14.eq.4-4}{(4.4)} follows, and so does the fact that
$\varphi$ is an isomorphism modulo torsion. This in turn implies, a
posteriori, that Riemann--Roch with denominators for a proper morphism of
smooth algebraic schemes over $k$ is essentially the same whether it is
written using $A(X)$, as in [2], or using $\Gtop^\bullet(X)$, as in this
seminar. The case of an arbitrary $k$ reduces easily to the algebraically
closed case: for a finite extension $k'/k$ of degree $n$, the kernel of the
natural map
\[
 A(X)\xrightarrow{f_X^*}A(X\otimes_k k')
\]
is annihilated by $n$. Indeed, the trace homomorphism
\[
 f_{X*}:A(X\otimes_k k')\longrightarrow A(X)
\]
satisfies $f_{X*}f_X^*=n\,\id$.\footnote{Thus the formula is established,
thanks to P.~Jouanolou, for $X$ smooth and quasi-projective over a field (see
the note on p.~5).}

\noindent\sgaStableTarget{sga6.e14.para.4-3}\textbf{4.3.}
It seems quite plausible that formulas \hyperlink{sga6.e14.eq.4-4}{(4.4)} hold exactly as stated, and not
merely modulo torsion; equivalently, that formula \hyperlink{sga6.e14.eq.4-5}{(4.5)} holds not merely
modulo torsion. This question is a very special case of the question whether
the denominator-free Riemann--Roch formula \hyperlink{sga6.e14.eq.3-1}{(3.1)} is valid in $A(X)$ for an
immersion $Y\to X$, namely the case $i=d$. This special case may reasonably
be expected to hold. Indeed, we already noted in no.~3 that, when $k$ has
characteristic zero, the denominator-free Riemann--Roch formula is valid at
least in $\Gtop^\bullet(X)$. Consequently the two sides of \hyperlink{sga6.e14.eq.4-5}{(4.5)} have the
same image in $\Gtop^\bullet(X)$; equivalently, the second formula in \hyperlink{sga6.e14.eq.4-4}{(4.4)}
holds. This result will extend automatically to arbitrary characteristic
once

the question raised in no.~1 has been satisfactorily resolved. As evidence
in support of formula \hyperlink{sga6.e14.eq.4-5}{(4.5)} in $A(X)$ itself, let us note that it is clear a
priori that the restriction of the left-hand side to $X-Y$ is zero. Thus,
since $Y$ is assumed irreducible, the left-hand side must a priori be an
integral multiple of $\cl(Y)$:
\[
 c_i(Y)=n(Y)\cl(Y).
\]
This already proves the exact formula \hyperlink{sga6.e14.eq.4-5}{(4.5)} whenever $\cl(Y)$ is not a
torsion element of $A^i(X)$. It is immediate that this is the case if $X$ is
projective (and not merely quasi-projective) over the field $k$: indeed, if
$\xi\in A^1(X)$ is a polarization of $X$, then
$\deg(Y\mathbin{\cdot}\xi^{\dim Y})>0$. This proves the desired formula in that
case. It follows that the formula is still valid whenever $X$ is isomorphic
to an open subscheme of a smooth projective scheme over $k$, as will be the
case when an adequate solution of the problem of resolution of
singularities is available. In particular, this holds when $k$ has
characteristic zero, by Hironaka~[11].

The validity of the exact formula \hyperlink{sga6.e14.eq.4-5}{(4.5)}, and hence of \hyperlink{sga6.e14.eq.4-4}{(4.4)}, would imply that
the kernel of
\[
 \varphi^i:A^i(X)\longrightarrow\Gtop^i(X)
\]
is annihilated by $(i-1)!$, and in particular \emph{that
$\varphi^2:A^2(X)\to\Gtop^2(X)$ is an isomorphism}. (That $\varphi^1$ is an
isomorphism already follows immediately from \hyperlink{sga6.e10.thm.5-3-2}{X~5.3.2}.) Unless I am
mistaken, the assertion that $\varphi^2$ is an isomorphism has not yet been
proved, even in the special case $\dim X=3$. It was therefore in any event
premature for the editor to assert the validity of (4.4) and (4.5) in
[9, no.~4, $3^\circ$].

\noindent\sgaStableTarget{sga6.e14.para.4-4}\textbf{4.4.}
At present, the principal advantage of $\Gtop^\bullet(X)$ (or of
$\operatorname{Gr}^\bullet(X)$) over the Chow ring $A(X)$ is that it is
better suited to explicit calculations of certain intersection formulas
between cycles in cases of excess intersection, that is, intersections of
excessively large dimension. Indeed, the definition of the intersection
class by means of Chow's lemma, which allows a cycle to be moved sufficiently
within its class, is poorly suited to calculation. A striking example in
this regard is the following fundamental formula [\hyperlink{sga6.e07.expose}{Exposé VII}~\hyperlink{sga6.e07.para.2-6}{2.7}], for a
closed immersion $f:Y\to X$ of smooth quasi-projective schemes over $k$:
\sgaStableTarget{sga6.e14.eq.4-6}\begin{equation*}\tag{4.6}
 f^*f_*(y)=y\,c_d(\underline{N}),
\end{equation*}
a formula valid in $\Gtop^\bullet(Y)$ for $y\in\Gtop^\bullet(Y)$, where
$\underline{N}$ is the conormal sheaf

of $Y$ in $X$, assumed everywhere to have rank $d$. In the cited passage,
this formula follows from a very simple computation of the sheaves
$\underline{\operatorname{Tor}}$. By \hyperlink{sga6.e14.para.4-2}{4.2}, the preceding formula implies the
same formula in $A(Y)$, but only modulo torsion. It is quite remarkable that
formula \hyperlink{sga6.e14.eq.4-6}{(4.6)} in $A(Y)$ itself has not yet been proved,\footnote{Cf. the note
on p.~23.} nor even the special case obtained by taking $y=1$, which gives
the self-intersection of $Y$:
\sgaStableTarget{sga6.e14.eq.4-7}\begin{equation*}\tag{4.7}
 f^*(\cl(Y))=c_d(\underline{N}),
\end{equation*}
even when $k$ is the field of complex numbers.

An analogous problem arises when one attempts to compute the Chow ring of
the variety $X'$ obtained by blowing up $X$ along $Y$, by imitating the
method used to compute $\operatorname{Gr}^\bullet(X')_{\Q}$
[\hyperlink{sga6.e07.expose}{Exposé VII}~\hyperlink{sga6.e07.thm.4-8}{4.8}], which also works (without tensoring with $\Q$) for the
computation of $\Gtop^\bullet(X')$. With the notation of the cited passage,
the key formula for this computation is
\sgaStableTarget{sga6.e14.eq.4-8}\begin{equation*}\tag{4.8}
 f^*(i_*(y))
   =j_*\bigl(g^*(y)c_{d-1}(\check{\underline{F}})\bigr).
\end{equation*}
At present, this formula has been established in $\Gtop^\bullet(X')$ (for
$y\in\Gtop^\bullet(Y)$) by a computation of
$\underline{\operatorname{Tor}}$ (\hyperlink{sga6.e07.expose}{Exposé VII}~\hyperlink{sga6.e07.prop.3-4}{3.4}), while its validity in
$A(X')$ (for $y\in A(Y)$) remains to be proved.

Let us note that even the $\ell$-adic analogues of formulas \hyperlink{sga6.e14.eq.4-6}{(4.6)}, \hyperlink{sga6.e14.eq.4-7}{(4.7)},
and \hyperlink{sga6.e14.eq.4-8}{(4.8)} have not yet been proved; when $y$ is the cohomology class of an
algebraic cycle, formulas \hyperlink{sga6.e14.eq.4-6}{(4.6)} and \hyperlink{sga6.e14.eq.4-8}{(4.8)} are in any event true modulo
torsion.

\noindent\sgaStableTarget{sga6.e14.para.4-5}\textbf{4.5.}
In this section and the following ones, we shall give examples for which,
in diagram \hyperlink{sga6.e14.eq.4-1}{(4.1)}, the homomorphism $j$, respectively the homomorphism
$\varphi$, is not an isomorphism. For a given $X$, it is plainly equivalent
to say that
\[
 j:\operatorname{Gr}^\bullet(X)\longrightarrow\Gtop^\bullet(X)
\]
is an isomorphism, a monomorphism, or an epimorphism. Moreover, the image of
$j$ is the subring of $\Gtop^\bullet(X)$ generated by the Chern classes of
locally free sheaves on $X$. If $K^\bullet(X)$ is generated as a
$\Z$-algebra by the classes of invertible sheaves on $X$, one sees at once
that the subring in question is also the subring generated by
$\Gtop^1(X)$, the degree-one component of $\Gtop^\bullet(X)$.

Now take for $X$ a flag variety
\[
 X=G/B,
\]
where $G$ is a simply connected semisimple algebraic group over the
algebraically closed field $k$, and $B$ is a Borel subgroup~[5]. When
$k=\C$, the Bruhat cell decomposition of $X$ [5, Exposé 13] implies that
the canonical homomorphism
\sgaStableTarget{sga6.e14.eq.unnamed}\begin{equation*}\tag{*}
 K^\bullet(X)\longrightarrow K(X(\C))
\end{equation*}
is an isomorphism, as is the canonical homomorphism
\sgaStableTarget{sga6.e14.eq.unnamed.occ2}\begin{equation*}\tag{**}
 A(X)\longrightarrow H^*(X(\C),\Z),
\end{equation*}
where $X(\C)$ is the compact space of complex points of $X$, endowed with
its usual topology induced by that of $\C$. On the other hand, L.~Hodgkin
[12] proved that $K(X(\C))$ is generated by the classes of invertible
sheaves on $X$. The same is therefore true of $K^\bullet(X)$ itself;
moreover, it is not difficult to deduce, by a specialization argument using
semisimple groups split over $\Z$ (SGA 3), that the result remains valid
over every algebraically closed field $k$. The cell decomposition of $X$
(or, when $k=\C$, the bijectivity of $(**)$) also readily implies that
$A(X)$ is torsion-free. Hence, by \hyperlink{sga6.e14.para.4-2}{4.2}, the homomorphism
\[
 \varphi:A(X)\longrightarrow\Gtop^\bullet(X)
\]
is an isomorphism. It follows that $j$ is an isomorphism if and only if
$A(X)$ is generated by its elements of degree $1$, or, when $k=\C$, if and
only if the integral cohomology of $X(\C)$ is generated by its elements of
degree $2$, i.e. if the classifying space $B_{G(\C)}$ is without torsion. It
is known [15, p.~217, Theorem B, $(1)\Leftrightarrow(2)$, and 2.5] that this
holds if and only if $G$ is without torsion, equivalently if $G$ is
isomorphic to a product of groups $\operatorname{SL}(n)$ and
$\operatorname{Sp}(m)$. Thus it suffices to take for $G$ a group
$\operatorname{SO}(n)$ with $n\geq7$, or a simple exceptional group, to
obtain an example in which $j$ is not an isomorphism.

\noindent\sgaStableTarget{sga6.e14.para.4-6}\textbf{4.6.}
In this example there can be no doubt that the good ring is the torsion-free
ring
\[
 A(X)\simeq\Gtop^\bullet(X),
\]
whose basis consists of the cells in the Bruhat decomposition and which,
when $k=\C$, is also isomorphic to the integral cohomology of $X(\C)$. By
contrast, $\operatorname{Gr}^\bullet(X)$ appears as a ring having torsion
of a rather pathological nature. This example therefore confirms the view
that the invariant $\operatorname{Gr}^\bullet(X)$ is hardly interesting

except modulo torsion, i.e. after tensoring with $\Q$. At the same time,
this example shows that $\operatorname{Gr}^\bullet(X)$ is covariant with
respect to proper morphisms only after tensoring with $\Q$, even if one
restricts to closed immersions of smooth quasi-projective schemes. Indeed,
if $X$ is such that $j$ is not an isomorphism, there is an integral closed
subscheme $Y$ of $X$ such that
\[
 \cl^\bullet(Y)\in K^\bullet(X)
\]
does not lie in $\Filt^i(X)$, where $i=\operatorname{codim}(Y,X)$. We may
choose $Y$ of minimal dimension, so that
\[
 \Filt^{i+1}(X)=\Filt_{\mathrm{top}}^{i+1}(X).
\]
If $Z$ is the singular locus of $Y$, and $X'=X\setminus Z$,
$Y'=Y\setminus Z$, it follows that
\[
 \cl^\bullet(Y')\in K^\bullet(X')
\]
does not lie in $\Filt^i(X')$. Here one uses the fact that
$K^\bullet(X)\to K^\bullet(X')$ is surjective and compatible with the
augmentations, and that its kernel is
\[
 \operatorname{Im}\bigl(K^\bullet(Z)\longrightarrow K^\bullet(X)\bigr)
 \subset
 \Filt^{i+1}(X)=\Filt_{\mathrm{top}}^{i+1}(X).
\]
Now consider the inclusion
\[
 f:Y'\longrightarrow X',
\]
which is a closed immersion of codimension $i$ between smooth
quasi-projective schemes over $k$, \emph{such that}
\[
 f_*:K^\bullet(Y')\longrightarrow K^\bullet(X')
\]
\emph{is not compatible with the $\lambda$-filtrations, even up to a shift
by $i$}, since $f_*(1)\notin\Filt^i(X')$.

\noindent\sgaStableTarget{sga6.e14.para.4-7}\textbf{4.7.}
Retaining the notation of \hyperlink{sga6.e14.para.4-5}{4.5}, note that \hyperlink{sga6.e09.expose}{Exposé IX} shows that Hodgkin's
theorem quoted above is equivalent to the isomorphism
\[
 K^\bullet(G)\simeq\Z,
 \qquad\text{hence}\qquad
 \Gtop^\bullet(G)\simeq\Z.
\]
The homomorphism $K^\bullet(G/B)\to K^\bullet(G)$ is a priori surjective,
and its kernel is the ideal generated by the elements
$\cl^\bullet(L)-1$, where $L$ ranges over the invertible sheaves on $X=G/B$
associated with representations $B\to\Gm$ (which in fact yield all
invertible sheaves on $X$ [5, Exposé 15]). An analogous argument shows that
\[
 A(G/B)\longrightarrow A(G)
\]
is also surjective, and that its kernel is the ideal generated by the
elements of $A^1(G/B)$ defined by representations $B\to\Gm$, hence by all
of $A^1(G/B)$.\footnote{Cf. \hyperlink{sga6.e09.prop.3-4}{IX 3.4}, 2.2.} Thus $A(G)$ is reduced to
$A^0(G)=\Z$ if and
only if $G$ is ``without torsion'', i.e. if and only if $G$ is isomorphic to
a product of groups $\operatorname{Sl}(n)$ and $\operatorname{Sp}(m)$. Thus,
in this case and only in this case, the homomorphism
\[
 \varphi:A(G)\longrightarrow\Gtop^\bullet(G)
\]
is an isomorphism. Consequently, if $G$ has torsion, the underlying scheme
provides an example of a non-bijective $\varphi$.

\noindent\sgaStableTarget{sga6.e14.para.4-8}\textbf{4.8.}
One likewise finds an example in which the homomorphism $\gamma$ of \hyperlink{sga6.e14.eq.4-1}{(4.1)}
\sgaStableTarget{sga6.e14.eq.4-9}\begin{equation}\tag{4.9}
 \gamma:A(G)\longrightarrow H^{2*}(G,\Zell(*))
\end{equation}
does not factor through $\Gtop^\bullet(G)$, or, what amounts to the same
thing here, an example in which this homomorphism is nonzero in degrees
$>0$. Restricting (as one may, thanks to suitable specialization theorems)
to the case $k=\C$, one is reduced to studying the homomorphism
\sgaStableTarget{sga6.e14.eq.4-10}\begin{equation}\tag{4.10}
 A(G)\longrightarrow H^*(G(\C),\Z),
\end{equation}
which is moreover obtained, by passage to the quotient, from the canonical
homomorphism
\[
 A(G/B)\simeq H^*(G/B(\C),\Z)
 \longrightarrow H^*(G(\C),\Z).
\]
It is tempting to conjecture that \hyperlink{sga6.e14.eq.4-10}{(4.10)} is in fact injective, since it is
also known that the torsion primes for $G(\C)$ are exactly those dividing
the order of the finite group
\[
 A^+(G)=\coprod_{i>0}A^i(G).
\]
If this were so, whenever $G$ had torsion, \hyperlink{sga6.e14.eq.4-10}{(4.10)} would not factor through
$\varphi$; hence, for a suitable $\ell$ (more precisely, for $\ell$ dividing
the order of $A^+(G)$, i.e. for $\ell$ a torsion prime of $G$), the same
would hold for \hyperlink{sga6.e14.eq.4-9}{(4.9)}. We can at least assert that, for the group $G=G_2$
over $\C$, \hyperlink{sga6.e14.eq.4-10}{(4.10)} is nonzero in degree $>0$: more precisely, there exists an
element of $A^3(G)$ whose image in $H^6(G(\C),\Z)$ has order $2$. It follows
that \hyperlink{sga6.e14.eq.4-9}{(4.9)} does not factor through $\varphi$ when $\ell=2$. The fact used
about $G_2$, which amounts to saying that there exists in
$H^6(G/B(\C),\Z)$ an element whose image in $H^6(G(\C),\Z)$ is nonzero, is
found in [3], via the familiar dictionary between complex semisimple
algebraic groups and compact Lie groups; it was kindly communicated to me
by J.-P. Serre.

\sgaStableTarget{sga6.e14.sec.5}\section*{5. Relations between $\operatorname{Gr}^\bullet(X)$ and
$H^{2*}(X,\Zell(*))$}

\noindent\sgaStableTarget{sga6.e14.para.5-1}\textbf{5.1.}
The example given in \hyperlink{sga6.e14.para.4-8}{4.8} does not exclude the possibility that, for an
arbitrary scheme $X$, one might find a ring homomorphism
\[
 g':\operatorname{Gr}^\bullet(X)
 \longrightarrow H^{2*}(X,\Zell(*)),
\]
such that, for every locally free module $E$ on $X$, $g^i$ sends the Chern
class $c_i(E)$ of Exp.~V to the cohomological Chern class of SGA 5 VII. Of
course, $\ell$ here denotes a prime distinct from the residual
characteristics of $X$. Since the ring $\operatorname{Gr}^\bullet(X)$ is
generated by the $c_i(E)$, one sees in any case that $g'$ is uniquely
determined by the requirement that it be compatible with Chern classes.

\noindent\sgaStableTarget{sga6.e14.para.5-2}\textbf{5.2.}
Using the Chern homomorphism
\[
 K^\bullet(X)\longrightarrow
 \Ch\bigl(H^{2*}(X,\Zell(*))\bigr),
\]
one obtains, by passing to the associated gradeds respectively for the
filtration on the first member and the filtration by degree on the second,
a homomorphism of graded groups (but not of rings!)
\[
 \psi:\operatorname{Gr}^\bullet(X)
 \longrightarrow H^{2*}(X,\Zell(*)),
\]
which satisfies
\[
 \psi(c_i(E))=(-1)^{i-1}(i-1)!\,c_i^\ell(E),
\]
where $c_i^\ell(E)$ denotes the $\ell$-adic Chern class. Dividing the
components $\psi^i$ by the indicated factors therefore gives a ring
homomorphism
\[
 g_\Q:\operatorname{Gr}^\bullet(X)
 \longrightarrow H^{2*}(X,\Qell(*))
 =H^{2*}(X,\Zell(*))\otimes_{\Zell}\Qell,
\]
sending Chern classes to Chern classes, and such that the image of the first
member is contained in that of $H^2(X,\Zell(*))$. Consequently, the question
of the existence of $g'$ has an affirmative answer provided
$H^{2*}(X,\Zell(*))$ is replaced by its quotient by its torsion subgroup.
The first doubtful case for defining the components $g^i$ of the
hypothetical homomorphism $g$ is that of $g^3$; a priori one succeeds only
in defining $2g^3=\psi^3$. The answer to the question of the existence of
$g^3$ would be negative if, for example, one could find three invertible
sheaves $L_1,L_2,L_3$ on $X$ satisfying
\[
 (L_1-1)(L_2-1)(L_3-1)=0
 \quad\text{in }K^\bullet(X),
\]

but such that
\[
 c_1(L_1)c_1(L_2)c_1(L_3)\ne0
 \quad\text{in }H^6(X,\mu_2^{\otimes3}).
\]

\noindent\sgaStableTarget{sga6.e14.para.5-3}\textbf{5.3.}
When $X$ is smooth and quasi-projective over a field $k$, one may similarly
ask whether there exists a ring homomorphism
\[
 g:\operatorname{Gr}^\bullet(X)\longrightarrow A(X)
\]
taking values in the Chow ring $A(X)$ and sending Chern classes to Chern
classes. The preceding observations may be repeated in this connection;
they show in particular how to construct $(-1)^{i-1}(i-1)!g^i$ directly
when $g^i$ itself is unavailable.

\noindent\sgaStableTarget{sga6.e14.para.5-4}\textbf{5.4.}
Let us note, to conclude this section, that essentially the same question
arises if, instead of a scheme $X$, one considers a compact topological
space (say) $X$. One can still consider the associated graded
$\operatorname{Gr}^\bullet(X)$ of $K^\bullet(X)$ endowed with its
$\lambda$-filtration, and seek a ring homomorphism
\[
 \operatorname{Gr}^\bullet(X)\longrightarrow H^{2*}(X,\Z)
\]
sending Chern classes to Chern classes.

\sgaStableTarget{sga6.e14.sec.6}\section*{6. Cohomological Riemann--Roch theorem and Gysin homomorphism}

\noindent\sgaStableTarget{sga6.e14.para.6-1}\textbf{6.1.}
For simplicity, we shall confine ourselves here to schemes $X$ admitting an
ample invertible module, so that there is no need to distinguish between the
naive $K^\bullet(X)$ and the other one. Let us begin with the Riemann--Roch
formula ``with denominators'', for a proper complete-intersection morphism
$f:X\longrightarrow Y$. For a prime $\ell$ prime to the residual
characteristics of $Y$, one can define the $\ell$-adic Chern classes of a
locally free module $E$,
\[
 c_i(E)\in H^{2i}(X,\Zell(i))
\]
(SGA 5 VII), corresponding to a homomorphism
\sgaStableTarget{sga6.e14.eq.6-1}\begin{equation}\tag{6.1}
 K^\bullet(X)\longrightarrow \Ch\bigl(H^{2*}(X,\Zell(*))\bigr).
\end{equation}
The question therefore arises of giving a cohomological version of the
Riemann--Roch formula. The only ingredient missing in order to give meaning
to the formula is a Gysin homomorphism

\sgaStableTarget{sga6.e14.eq.6-2}\begin{equation}\tag{6.2}
 H^n(X,\Qell(i))\longrightarrow H^{n-2d}(Y,\Qell(i-d)),
\end{equation}
where $d$ is the virtual relative dimension of $f$. Assuming this
homomorphism to be defined, and taking into account the Riemann--Roch formula
in the form of \hyperlink{sga6.e08.expose}{Exp.~VIII}, with values in
$\operatorname{Gr}^\bullet(Y)_{\Q}$, one sees that the validity of the
$\ell$-adic Riemann--Roch formula for the morphism $f$ is exactly equivalent
to the commutativity of the square
\sgaStableTarget{sga6.e14.eq.6-3}\begin{equation}\tag{6.3}
\begin{tikzcd}[column sep=large,row sep=large]
 \operatorname{Gr}^\bullet(X)_{\Q} \arrow[r] \arrow[d]
   & H^{2*}(X,\Qell(*)) \arrow[d] \\
 \operatorname{Gr}^\bullet(Y)_{\Q} \arrow[r]
   & H^{2*}(Y,\Qell(*)),
\end{tikzcd}
\end{equation}
where the horizontal arrows are obtained from \hyperlink{sga6.e14.eq.6-1}{(6.1)} by passage to the
associated gradeds and division, so as to send Chern classes to Chern
classes, while the vertical arrows are respectively the arrow of \hyperlink{sga6.e08.expose}{Exp.~VIII}
and the Gysin homomorphism in cohomology (assumed given).

\noindent\sgaStableTarget{sga6.e14.para.6-2}\textbf{6.2.}
Note that the Gysin homomorphism is defined in all cases if $Y$ is smooth
over a field $k$ (SGA 5 IV), using the theorem of relative cohomological
purity over $k$. The same construction applies when $Y$ is excellent,
regular, and of characteristic zero, thanks to the theorem of absolute
cohomological purity of SGA 4 XIX; and it will automatically apply to every
excellent regular $Y$, without restriction on the characteristic, once the
theorem of absolute cohomological purity has been proved in full generality.
Finally, the same construction applies, thanks to the theorem of relative
cohomological purity (SGA 4 XVI 3.7), when $X$ and $Y$ are both smooth over
the same scheme $S$ and $f$ is an $S$-morphism. It is natural to conjecture
that, in each of these cases in which the Gysin homomorphism is defined,
diagram \hyperlink{sga6.e14.eq.6-3}{(6.3)} is indeed commutative. The only case in which this assertion
is currently verified is that in which $X$ and $Y$ are both smooth over a
field $k$; it then follows from the commutativity of the analogous diagram
\sgaStableTarget{sga6.e14.eq.6-4}\begin{equation}\tag{6.4}
\begin{tikzcd}[column sep=large,row sep=large]
 A(X) \arrow[r] \arrow[d]
   & H^{2*}(X,\Zell(*)) \arrow[d] \\
 A(Y) \arrow[r]
   & H^{2*}(Y,\Zell(*)),
\end{tikzcd}
\end{equation}

involving the rings $A(X)$ and $A(Y)$ of cycle classes modulo rational
equivalence, whose commutativity was established in SGA 5 IV.

\noindent\sgaStableTarget{sga6.e14.para.6-3}\textbf{6.3.}
Note that in each of the cases considered above in which one succeeds in
defining a Gysin homomorphism, one in fact obtains a construction for
cohomology with coefficients in $\Zell(*)$, and not only in $\Qell(*)$.
This makes it possible, when $f$ is an \emph{immersion}, to pose the question
of the validity of the Riemann--Roch formula ``without denominators'', as
explained in no.~3. When $X$ and $Y$ are both smooth over a field $k$, the
commutativity of \hyperlink{sga6.e14.eq.6-4}{(6.4)} shows that it would suffice to prove the formula
without denominators in the Chow ring $A(Y)$; cf.~\hyperlink{sga6.e14.para.3-3}{(3.3)}.

\noindent\sgaStableTarget{sga6.e14.para.6-4}\textbf{6.4.}
When no regularity hypothesis is made on $Y$, and even if $X$ and $Y$ are of
finite type over $\C$, one no longer sees a plausible definition of a
cohomological Gysin homomorphism \hyperlink{sga6.e14.para.6-2}{(6.2)}, even for a regular immersion of
codimension $1$. Can one nevertheless find one with reasonable properties,
in particular making square \hyperlink{sga6.e14.para.6-3}{(6.3)} commutative? The question is evidently
connected with that of defining, for $X$ a complete intersection of
codimension $i$ in $Y$, ``its cohomology class''
\[
 \gamma(X)\in H^{2i}(Y,\Zell(i)),
\]
in such a way as to turn intersections (without excess components) into cup
products. This latter question has an affirmative solution at least modulo
torsion, i.e. in $H^{2*}(Y,\Qell(*))$, by defining
\[
 \gamma(X)=\frac{1}{(-1)^{i-1}(i-1)!}\,c_i(\mathcal O_X).
\]
But I do not know whether $c_i(\mathcal O_X)$ is divisible by $(i-1)!$ in
$H^{2i}(Y,\Zell(i))$. In particular, if $Y$ is, say, a projective scheme of
dimension $3$ over the field of complex numbers and $X$ is reduced to a
(singular) point of $Y$, is $c_3(\mathcal O_X)$ divisible by $2$ in
$H^6(Y,\mathbb Z_2(3))$, and in particular is the Chern class zero in
$H^6(Y,\Z/2\Z)$? Let us also mention, along the same lines, the following
question: let $Y$ be a projective (singular) scheme of dimension $4$, and
$X$ a closed subscheme of dimension $1$ which is globally an intersection
of three positive divisors $D_1,D_2,D_3$. Consider the product of their
Chern classes modulo $2$,
\[
 c_1(D_1)c_1(D_2)c_1(D_3)\in H^6(Y,\Z/2\Z);
\]

is this element of $H^6$ independent of the particular choice of the $D_i$
whose intersection is $X$?

\sgaStableTarget{sga6.e14.sec.7}\section*{7. Chern classes of perfect complexes}

\noindent\sgaStableTarget{sga6.e14.para.7-1}\textbf{7.1.}
When one poses the question of developing a cohomological Riemann--Roch
formula, with or without denominators, outside the rather artificial
hypothesis that an ample invertible module exists on the schemes under
consideration, one is necessarily faced with the question of defining, for
a perfect complex of modules $L^\bullet$ on a scheme $X$, the Chern classes
\[
 c_i(L^\bullet)\in H^{2i}(X,\Zell(i)),
\]
generalizing the Chern classes of locally free modules, SGA 5 VII. The same
question arises, moreover, for general locally ringed topoi (compare [10]),
as well as for perfect complexes on paracompact spaces ringed by continuous
complex-valued functions, working then with ordinary integral cohomology.
For a compact space, it is known that the $K^\bullet(X)$ constructed with
perfect complexes is the same as the ``naive'' one constructed with locally
free modules (\hyperlink{sga6.e04.expose}{Exp.~IV}), and in that case one is trivially reduced to the
classical definition. But already for a locally compact space the problem
posed seems nontrivial, since the integral cohomology of such a space is not
always the inverse limit of the cohomologies of its compact subsets.

\noindent\sgaStableTarget{sga6.e14.para.7-2}\textbf{7.2.}
More precisely, when $Y$ is a closed subset of $X$ such that
$L^\bullet|_{X-Y}$ is zero (in the sense of derived categories), i.e. such
that the cohomology sheaves of $L^\bullet$ are supported in $Y$, one would
like to define Chern classes \emph{with support in $Y$}
\[
 c_i^Y(L^\bullet)\in H_Y^{2i}(X,\Zell(i))
 \quad\text{respectively}\quad H_Y^{2i}(X,\Z),
\]
which would recover the ``ordinary'' (sic) Chern classes of \hyperlink{sga6.e14.para.7-1}{7.1} by means of
the canonical homomorphism
\[
 H_Y^*(X,-)\longrightarrow H^*(X,-).
\]

Even in the case where the latter reduce to the Chern classes of locally
free modules ($X$ a scheme with an ample invertible module, or a compact
space), the construction of these Chern classes ``with supports'' has not
been carried out and will probably only be possible together with that of
the general Chern classes postulated in 7.1.\footnote{Added in October 1968.
The problems raised here were solved in the framework of Hodge cohomology by
L. Illusie; cf.~[13].}

It is probable that Chern classes with supports will have a role to play in
the proof of the cohomological Riemann--Roch formula without denominators
(cf.~\hyperlink{sga6.e14.para.3-5}{3.5}). On the other hand, they should make it possible to give an
improved version of the Riemann--Roch formula (with or without denominators)
by introducing families of supports into that formula.

\sgaStableTarget{sga6.e14.sec.8}\section*{8. Chow ring of regular schemes}

The problem is to define rational equivalence of cycles and to study the
resulting ring of classes (in particular its functorial properties) for a
regular noetherian scheme $X$ admitting an ample invertible module, thereby
generalizing the known theory when $X$ is smooth over a field [4]. The key
technical result will again be a ``moving lemma'', ensuring that a cycle can
be varied within its class so as to meet a given cycle in the proper
dimension.

The problem under consideration is quite natural a priori, especially since
we have seen (\hyperlink{sga6.e14.para.4-5}{4.5} to \hyperlink{sga6.e14.para.4-8}{4.8}) that the substitutes
$\operatorname{Gr}^\bullet(X)$ and $\Gtop^\bullet(X)$ for the Chow ring
$A(X)$ are hardly satisfactory when one does not confine oneself to a theory
``modulo torsion''. When $X$ is obtained from a smooth projective algebraic
scheme $Y$ over a field $k$ by taking the local scheme at the vertex of a
projective cone $C$ and removing the origin, one finds (assuming that $A(X)$
is indeed the inductive limit of the Chow rings of $U-\{O\}$, where $U$
ranges over the open neighborhoods in $C$ of the vertex $O$) that $A(X)$ is
canonically isomorphic to the quotient ring
\[
 A(Y)/\xi A(Y),
\]
where $\xi\in A^1(Y)$ is induced by the chosen projective embedding of $Y$.
This shows that knowing the Chow ring $A(X)$ (which is of a nonclassical
type) is equivalent to knowing $A(Y)/\xi A(Y)$, which might be called the
``primitive part'' of the Chow ring $A(Y)$ (in a sense suggested by the
Lefschetz--Hodge cohomological theory of nonsingular projective algebraic
varieties). This example shows that the ``local Chow rings'' of local rings
with isolated

singularities are geometric invariants that are by no means trivial and will
undoubtedly have a role to play in the ``arithmetic'' study of these local
rings.

Of course, once one has a reasonable notion of a Chow ring for a regular
noetherian scheme admitting an ample invertible module, one should
generalize to this case the relations established or conjectured in no.~4
between this ring and $\Gtop^\bullet(X)$, and formulate in terms of this ring
the question of the ``Riemann--Roch formula without denominators'' stated in
no.~3.

\sgaStableTarget{sga6.e14.bibliography-of-expose-xiv}\subsection*{Bibliography of Expos\'e XIV}

\begin{enumerate}[label={[\arabic*]},leftmargin=3.2em]
\item M. F. Atiyah and F. Hirzebruch,
  \emph{The Riemann--Roch theorem for analytic embeddings},
  Topology, vol.~1, pp.~151--166 (1962).

\item A. Borel and J.-P. Serre,
  \emph{Le th\'eor\`eme de Riemann--Roch},
  Bull. Soc. Math. France, vol.~86, pp.~97--136 (1958).

\item R. Bott and H. Samelson,
  Amer. Journ. of Math., vol.~80, p.~1004 (1958).

\item C. Chevalley,
  \emph{Les classes d'\'equivalence rationelle}, I and II,
  S\'eminaire \`a l'ENS 1958, Expos\'es 2 and 3.

\item C. Chevalley,
  \emph{Classification des groupes de Lie alg\'ebriques},
  S\'eminaire \`a l'ENS 1956/58.

\item A. Dold and D. Puppe,
  \emph{Homologie nicht-additiver Funktoren. Anwendungen},
  Annales de l'Inst. Fourier, t.~XI, pp.~201--312 (1961).

\item H. Grauert,
  \emph{Ein Theorem der analytischen Garbentheorie},
  Pub. Math. no.~5, pp.~5--64 (1960).

\item A. Grothendieck,
  \emph{Classes de faisceaux et th\'eor\`eme de Riemann--Roch},
  mimeographed notes prepared by the Institute for Advanced Study, 1957
  (the work cited as [RRR] in the present Expos\'e and reproduced in our
  seminar as an appendix to Chapter~0).

\item A. Grothendieck,
  \emph{La th\'eorie des classes de Chern},
  Bull. Soc. Math. France, vol.~86, pp.~137--154 (1958).

\item A. Grothendieck,
  \emph{Classes de Chern et repr\'esentations lin\'eaires des groupes
  discrets}, in: \emph{Dix expos\'es sur la cohomologie des sch\'emas},
  Masson--North Holland Pub. Cie.

\item H. Hironaka,
  \emph{Resolution of singularities of an algebraic variety over a field of
  characteristic zero}, Ann. of Math. 79, pp.~109--326 (1964).

\item L. Hodgkin, thesis (forthcoming?).

\item L. Illusie, work in preparation.

\item J.-P. Serre, \emph{Alg\`ebre Locale, Multiplicit\'es},
  Lecture Notes in Math. no.~11 (1965), Springer.

\item A. Borel,
  \emph{Sous-groupes commutatifs et torsion des groupes de Lie compacts
  connexes}, Tohoku Math. Journal, vol.~13, pp.~216--240 (1961).

\item[{[16]}(*)] P. Deligne, Expos\'e XVIII in
  \emph{Th\'eorie des topos et cohomologie \'etale des sch\'emas (SGA 6)},
  by M. Artin, A. Grothendieck, and J. L. Verdier.

\item[{[17]}(*)] Yu. I. Manin,
  \emph{Correspondances, motifs et transformations mono\"idales}
  (in Russian). Mat. Sborik, t.~77 (119), no.~4, 1968.
\end{enumerate}

\noindent{\footnotesize (*) Added in February 1969.}

\sgaStableTarget{sga6.e14.note-for-page-1}\subsection*{Note for page 1}

The theory of the operations $\Lambda^i$ in $D^{-}(X)$, and of the
operations $\lambda^i$ in $K^\bullet(X)$, was written down by P. Deligne in
1968. He proves that $K^\bullet(X)$ thereby becomes a pre-$\lambda$-ring and,
when $\mathcal O_X$ is a sheaf of algebras over a field, that
$K^\bullet(X)$ is even a $\lambda$-ring. It is plausible that a refinement
of Deligne's method would make it possible to dispense with the restrictive
hypothesis on $\mathcal O_X$.

\sgaStableTarget{sga6.e14.notes-for-page-11}\subsection*{Notes for page 11}

\noindent(*) (April 1969) D. Mumford has pointed out to us a proof of
\hyperlink{sga6.e14.para.4-6}{(4.6)}\SGAArchiveOnly{\footnote{\textbf{Source note
[SGA6-XIV-IDX702-XREF96-VS-XREF46-SRCDEF-001].} The printed French note
reads ``(9.6)''; the contextual reference is formula \hyperlink{sga6.e14.eq.4-6}{(4.6)}.}}
in the \emph{Chow ring} (without disregarding torsion), which he obtained in
a seminar at Harvard in 1959! The reader will find the proof of this
beautiful result in SGA 5 VII 9. This result obviously contains the
corresponding cohomological result mentioned in note (**). The cohomological
result, however, is proved by J.-P. Jouanolou under more general conditions
(without the quasi-projectivity hypothesis, and in the relative smooth case over
an arbitrary base).

\medskip

\noindent(**) (January 1969) J.-P. Jouanolou has just proved these formulas
for every $y$, without disregarding torsion; cf. SGA 5 VII \S~4.9. Deligne
and Manin [17] had previously observed independently, and by (easy)
arguments of a rather different nature, that it was enough to establish the
formulas for ``algebraic'' $y$, and even for $y=1$ in [16], in order to
conclude the same relation in the general case.
\SGAArchiveOnly{\input{fragments/SGA6_XIV_idx685_702_Editorial_Source_Notes.tex}}

\newpage
\sgaStableTarget{sga6.back.terminological-index}\section*{Terminological index}
\begingroup\small
\begin{longtable}{@{}>{\raggedright\arraybackslash}p{0.64\linewidth}>{\raggedright\arraybackslash}p{0.14\linewidth}>{\raggedright\arraybackslash}p{0.12\linewidth}@{}}
\toprule
 & Exposé & §\\
\midrule
\endhead
additive (map) & I & \hyperlink{sga6.e01.def.6-1}{6.1}\\
additive (map) & IV & \hyperlink{sga6.e04.para.1-5}{1.5}\\
augmented \(K\)-\(\lambda\)-algebra & V & \hyperlink{sga6.e05.def.3-9}{3.9}\\
algebraically equivalent to zero (in \(K^\bullet\)) & X App. & \hyperlink{sga6.e10.para.7-11}{7.11}\\
algebraically equivalent to zero (invertible sheaf) & XIII & \hyperlink{sga6.e13.para.4-4}{4.4}\\
of \(\mathcal C_0\)-perfect amplitude contained in \([a,b]\) (complex) & I & \hyperlink{sga6.e01.def.4-7}{4.7}\\
of finite \(\mathcal C_0\)-perfect amplitude (complex) & I & \hyperlink{sga6.e01.def.4-7}{4.7}\\
of perfect amplitude contained in \([a,b]\) (complex) & I & \hyperlink{sga6.e01.ex.4-8}{4.8}\\
of finite perfect amplitude (complex) & I & \hyperlink{sga6.e01.ex.4-8}{4.8}\\
of flat amplitude contained in \([a,b]\) (complex) & I & \hyperlink{sga6.e01.def.5-2}{5.2}\\
of finite flat amplitude (complex) & I & \hyperlink{sga6.e01.def.5-2}{5.2}\\
of perfect amplitude relative to \(f\) (or \(Y\)) contained in \([a,b]\) (complex) & III & \hyperlink{sga6.e03.def.4-1}{4.1}\\
of perfect amplitude contained in \([-n,0]\) (morphism) & III & \hyperlink{sga6.e03.def.4-1}{4.1}\\
of finite perfect amplitude relative to \(f\) (complex) & III & \hyperlink{sga6.e03.def.4-1}{4.1}\\
of flat amplitude relative to \(f\) (or \(Y\)) contained in \([a,b]\) (complex) & III & \hyperlink{sga6.e03.def.3-1}{3.1}\\
of flat amplitude contained in \([-n,0]\) (morphism) & III & \hyperlink{sga6.e03.def.3-1}{3.1}\\
of finite flat amplitude relative to \(f\) (complex) & III & \hyperlink{sga6.e03.def.3-1}{3.1}\\
projective amplitude & I & \hyperlink{sga6.e01.rem.4-12-1}{4.12.1}\\
\(\lambda\)-ring & V & \hyperlink{sga6.e05.def.2-4}{2.4}\\
\(\lambda\)-ring generated by generators and relations & V & \hyperlink{sga6.e05.def.4-6}{4.6}\\
free \(\lambda\)-ring generated by a family of generators & V & \hyperlink{sga6.e05.def.4-4}{4.4}\\
Chern ring & V & \hyperlink{sga6.e05.def.6-2}{6.2}\\
polynomial map & X & \hyperlink{sga6.e10.def.2-1-1}{2.1.1}\\
binomial (ring) & V & \hyperlink{sga6.e05.def.2-7}{2.7}\\
Chern character & V & \hyperlink{sga6.e05.def.6-4}{6.4}\\
Todd character & V & \hyperlink{sga6.e05.para.1-16}{1.16}\\
abelian \(S\)-category & I & \hyperlink{sga6.e01.para.1-1-1}{1.1.1}\\
flat abelian \(S\)-category & I & \hyperlink{sga6.e01.para.1-1-1}{1.1.1}\\
additive \(S\)-category & I & \hyperlink{sga6.e01.para.1-1-1}{1.1.1}\\
triangulated \(S\)-category & I & \hyperlink{sga6.e01.para.1-1-1}{1.1.1}\\
Chern class & V & \hyperlink{sga6.e05.para.6-7}{6.7}\\
completed Chern class & V & \hyperlink{sga6.e05.para.6-7}{6.7}\\
total Chern class & V & \hyperlink{sga6.e05.para.6-7}{6.7}\\
codimension (of a regular immersion) & VII & \hyperlink{sga6.e07.def.1-6}{1.6}\\
coherent (complex) & I & \hyperlink{sga6.e01.def.3-1}{3.1}\\
coherer & II & \hyperlink{sga6.e02.sec.3.2}{3.2}\\
relative cotangent complex & VIII & \hyperlink{sga6.e08.def.2-3}{2.3}\\
elliptic complex & II App. II & \hyperlink{sga6.e02.para.1-2-5}{1.2.5}\\
cup-product & I & \hyperlink{sga6.e01.para.8-2}{8.2}\\
special positive \(k\)-cycle & XIII & \hyperlink{sga6.e13.def.6-9}{6.9}\\
degree (of a polynomial map) & X & \hyperlink{sga6.e10.def.2-1-1}{2.1.1}\\
degree (of an invertible sheaf) & XIII & \hyperlink{sga6.e13.para.2-2}{2.2}\\
degree (of a morphism of schemes) & X & \hyperlink{sga6.e10.sec.4.4}{4.4}\\
dense (subset of an abelian or triangulated category) & IV & \hyperlink{sga6.e04.def.1-7}{1.7}\\
perfect dimension & I & \hyperlink{sga6.e01.def.4-14-1}{4.14.1}\\
virtual relative dimension & VIII & \hyperlink{sga6.e08.def.1-9}{1.9}\\
divisorial & II & \hyperlink{sga6.e02.def.2-2-5}{2.2.5}\\
dual & I & \hyperlink{sga6.e01.def.7-3}{7.3}\\
perfect duality & I & \hyperlink{sga6.e01.cor.7-8}{7.8}\\
numerical equivalence (in \(\operatorname{Gr}^\bullet(X)\) and \(\operatorname{Gr}_\bullet(X)\)) & X & \hyperlink{sga6.e10.sec.4.5}{4.5}\\
\(\tau\)-equivalence & X & \hyperlink{sga6.e10.def.2-4-1}{2.4.1}\\
\(\tau\)-equivalence & XIII & \hyperlink{sga6.e13.para.4-4}{4.4}\\
complex pseudo-analytic space & II & \hyperlink{sga6.e02.def.2-4-7}{2.4.7}\\
complex Stein pseudo-analytic space & II & \hyperlink{sga6.e02.def.2-4-7}{2.4.7}\\
exact (subset of an abelian or triangulated category) & IV & \hyperlink{sga6.e04.def.1-7}{1.7}\\
\((b)\)-sheaf & XIII & \hyperlink{sga6.e13.def.1-5}{1.5}\\
conormal sheaf & VII & \hyperlink{sga6.e07.def.1-6}{1.6}\\
relative cotangent sheaf & II App. II & \hyperlink{sga6.e02.para.1-2-1}{1.2.1}\\
relative tangent sheaf & II App. II & \hyperlink{sga6.e02.para.1-2-1}{1.2.1}\\
\((b)\)-family & XIII & \hyperlink{sga6.e13.para.1-12}{1.12}\\
ample family & II & \hyperlink{sga6.e02.def.2-2-4}{2.2.4}\\
family of classes of coherent sheaves & XIII & \hyperlink{sga6.e13.para.1-12}{1.12}\\
limited family of classes of coherent sheaves & XIII & \hyperlink{sga6.e13.para.1-12}{1.12}\\
family of classes of coherent sheaves limited by a coherent sheaf & XIII & \hyperlink{sga6.e13.para.1-12}{1.12}\\
relative cotangent bundle & II App. II & \hyperlink{sga6.e02.para.1-2-1}{1.2.1}\\
relative tangent bundle & II App. II & \hyperlink{sga6.e02.para.1-2-1}{1.2.1}\\
faithfully flat (morphism of ringed topoi) & I & \hyperlink{sga6.e01.exercise.2-16-2}{2.16.2}\\
filtration of an augmented \(\lambda\)-algebra & V & \hyperlink{sga6.e05.para.3-10}{3.10}\\
additive \(S\)-functor & I & \hyperlink{sga6.e01.para.1-1-3}{1.1.3}\\
exact \(S\)-functor & I & \hyperlink{sga6.e01.para.1-1-3}{1.1.3}\\
projection formula & IV & \hyperlink{sga6.e04.para.2-11}{2.11}\\
Grothendieck group & IV & \hyperlink{sga6.e04.def.2-2}{2.2}\\
\(\lambda\)-homomorphism & V & \hyperlink{sga6.e05.def.2-1}{2.1}\\
specialization homomorphism & X App. & \hyperlink{sga6.e10.prop.7-3}{7.3}\\
specialization homomorphism & X App. & \hyperlink{sga6.e10.sec.7.16}{7.16}\\
\(\lambda\)-ideal & V & \hyperlink{sga6.e05.def.4-5}{4.5}\\
analytic index & II App. II & \hyperlink{sga6.e02.para.3-2-2}{3.2.2}\\
complete-intersection (morphism) & VIII & \hyperlink{sga6.e08.def.1-1}{1.1}\\
\(n\)-isomorphism & I & \hyperlink{sga6.e01.def.1-4-3}{1.4.3}\\
homogeneous free (graded module) & VI & \hyperlink{sga6.e06.eq.1-2-2}{1.2.2}\\
locally stable under kernels of epimorphisms & I & \hyperlink{sga6.e01.sec.1.2}{1.2}\\
locally quasi-liftable & I & \hyperlink{sga6.e01.sec.1.2}{1.2}\\
locally quasi-stable under kernels of epimorphisms & I & \hyperlink{sga6.e01.sec.1.2}{1.2}\\
locally liftable & I & \hyperlink{sga6.e01.sec.1.2}{1.2}\\
Chern number & X & \hyperlink{sga6.e10.ex.4-2}{4.2}\\
intersection number & X & \hyperlink{sga6.e10.sec.4.1}{4.1}\\
intersection number (for a family of invertible sheaves) & X & \hyperlink{sga6.e10.sec.4.4}{4.4}\\
numerically equivalent to zero (invertible sheaf) & XIII & \hyperlink{sga6.e13.para.4-4}{4.4}\\
universal differential operator of order \(n\) & II App. II & \hyperlink{sga6.e02.para.1-2-1}{1.2.1}\\
Adams operations & V & \hyperlink{sga6.e05.para.7-1}{7.1}\\
\(\mathcal C_0\)-perfect (complex) & I & \hyperlink{sga6.e01.def.4-7}{4.7}\\
perfect (complex) & I & \hyperlink{sga6.e01.ex.4-8}{4.8}\\
perfect (morphism) & III & \hyperlink{sga6.e03.def.4-1}{4.1}\\
perfect relative to \(f\) (complex) & III & \hyperlink{sga6.e03.def.4-1}{4.1}\\
precoherent & I & \hyperlink{sga6.e01.def.3-1}{3.1}\\
of finite \(n\)-\(\mathcal C_0\)-presentation & I & \hyperlink{sga6.e01.def.2-8}{2.8}\\
of finite presentation & I & \hyperlink{sga6.e01.def.2-8}{2.8}\\
of finite \(\infty\)-\(\mathcal C_0\)-presentation & I & \hyperlink{sga6.e01.def.2-8}{2.8}\\
\(n\)-\(\mathcal C_0\)-pseudo-coherent (complex) & I & \hyperlink{sga6.e01.prop.2-2}{2.2}\\
\(\mathcal C_0\)-pseudo-coherent (complex) & I & \hyperlink{sga6.e01.prop.2-2}{2.2}\\
\(n\)-pseudo-coherent (complex) & I & \hyperlink{sga6.e01.ex.2-15}{2.15}\\
\(n\)-pseudo-coherent (morphism) & III & \hyperlink{sga6.e03.def.1-2}{1.2}\\
pseudo-coherent (complex) & I & \hyperlink{sga6.e01.ex.2-15}{2.15}\\
pseudo-coherent (morphism) & III & \hyperlink{sga6.e03.def.1-2}{1.2}\\
\(n\)-pseudo-coherent relative to \(f\) (complex) & III & \hyperlink{sga6.e03.def.1-2}{1.2}\\
pseudo-coherent relative to \(f\) (complex) & III & \hyperlink{sga6.e03.def.1-2}{1.2}\\
\((b)\)-polynomial & XIII & \hyperlink{sga6.e13.def.1-9}{1.9}\\
Snapper polynomial & X & \hyperlink{sga6.e10.prop.2-3-1}{2.3.1}\\
universal polynomial & V & \hyperlink{sga6.e05.para.1-15}{1.15}\\
pre-\(\lambda\)-ring & V & \hyperlink{sga6.e05.def.2-1}{2.1}\\
\(n\)-quasi-isomorphism & I & \hyperlink{sga6.e01.def.1-4-3}{1.4.3}\\
rank (of a perfect complex) & I & \hyperlink{sga6.e01.def.6-11}{6.11}\\
\(m\)-regular (sheaf) & XIII & \hyperlink{sga6.e13.def.1-1}{1.1}\\
regular (homomorphism) & VII & \hyperlink{sga6.e07.def.1-1}{1.1}\\
regular (ideal) & VII & \hyperlink{sga6.e07.def.1-4}{1.4}\\
regular (immersion) & VII & \hyperlink{sga6.e07.def.1-4}{1.4}\\
triangulated sub-\(S\)-category & I & \hyperlink{sga6.e01.para.1-1-4}{1.1.4}\\
special (\(\lambda\)-ring) & V & \hyperlink{sga6.e05.def.2-4}{2.4}\\
special (morphism) & III & \hyperlink{sga6.e03.lem.1-1-4}{1.1.4}\\
Stein (locally ringed topos) & II & \hyperlink{sga6.e02.def.2-4-1}{2.4.1}\\
strictly \(\mathcal C_0\)-perfect & I & \hyperlink{sga6.e01.def.2-1}{2.1}\\
strictly \(\mathcal C_0\)-\(n\)-pseudo-coherent & I & \hyperlink{sga6.e01.def.2-1}{2.1}\\
strictly \(\mathcal C_0\)-pseudo-coherent & I & \hyperlink{sga6.e01.def.2-1}{2.1}\\
minimum system of generators & VII & \hyperlink{sga6.e07.def.1-4}{1.4.1}\\
of finite Tor-dimension & I & \hyperlink{sga6.e01.def.5-2}{5.2}\\
Tor-independence & III & \hyperlink{sga6.e03.def.1-5}{1.5}\\
trace & I & \hyperlink{sga6.e01.para.8-1}{8.1}\\
of finite \(\mathcal C_0\)-type & I & \hyperlink{sga6.e01.sec.1.2}{1.2}\\
mixed variety of class \(C^r\) & II App. II & \hyperlink{sga6.e02.sec.1.1}{1.1}\\
\bottomrule
\end{longtable}
\endgroup

\newpage
\sgaStableTarget{sga6.back.index-of-notations}\section*{Index of notations}
\begingroup\small
\begin{longtable}{@{}>{\raggedright\arraybackslash}p{0.50\linewidth}>{\raggedright\arraybackslash}p{0.17\linewidth}>{\raggedright\arraybackslash}p{0.15\linewidth}@{}}
\toprule
 & Exposé & §\\
\midrule
\endhead
\(A_N\) & V & \hyperlink{sga6.e05.para.5-5}{5.5}\\
\(\ch\) & V & \hyperlink{sga6.e05.para.6-3}{6.3}\\
\(\Ch(A),\ \Ch_K(A)\) & V & \hyperlink{sga6.e05.def.6-2}{6.2}\\
\((\Ch(A))_n\) & V & \hyperlink{sga6.e05.prop.6-5}{6.5}\\
\(\cl_{\mathcal C}\) & IV & \hyperlink{sga6.e04.para.1-1}{1.1}\\
\(\Coh(S)\) & II & \hyperlink{sga6.e02.notation.2-2-1}{2.2.1}\\
\(\mathcal C_0\text{-}\parf\text{-}\amp(F)\) & I & \hyperlink{sga6.e01.def.4-7}{4.7}\\
\(C^p_{\num}(X)\) & XIII & \hyperlink{sga6.e13.prop.5-2}{5.2}\\
\(C^r(V),\ C^\infty(V)\) & II App. II & \hyperlink{sga6.e02.sec.1.1}{1.1}\\
\(c(x),\ c^i(x),\ \widetilde c(x)\) & V & \hyperlink{sga6.e05.para.6-7}{6.7}\\
\(\Dcat(A,S),\ \Dcat(A_{\underline S})\) & I & \hyperlink{sga6.e01.ex.2-15}{2.15}\\
\(\Dcat(A_S)_{\coh},\ \Dcat(A_S)_{n\text{-}\coh}\) & I & \hyperlink{sga6.e01.ex.2-15}{2.15}\\
\(\Dcat(A_S)_{\parf}\) & I & \hyperlink{sga6.e01.notation.4-9}{4.9}\\
\(\Dcat^{-}(A_S)_{\torf},\ \Dcat^{-}(A_X)_{\torf}\) & I & \hyperlink{sga6.e01.para.5-3}{5.3}\\
\(\Dcat(\mathcal C)_{\mathcal C_0\text{-}\coh},\ \Dcat(\mathcal C)_{\mathcal C_0\text{-}n\text{-}\coh}\) & I & \hyperlink{sga6.e01.schol.2-6}{2.6}\\
\(\Dcat(\mathcal C)_{\mathcal C_0\text{-}\parf}\) & I & \hyperlink{sga6.e01.notation.4-9}{4.9}\\
\(\Dcat(f)_{\coh}\) & III & \hyperlink{sga6.e03.def.1-2}{1.2}\\
\(\Dcat(f)_{n\text{-}\coh}\) & III & \hyperlink{sga6.e03.def.1-2}{1.2}\\
\(\Dcat(f)_{\parf},\ \Dcat(f)_{\underline{\mathrm{parf}}}\) & III & \hyperlink{sga6.e03.notation.4-2}{4.2}\\
\(\Dcat(f)_{\torf}\) & III & \hyperlink{sga6.e03.notation.3-2}{3.2}\\
\(\Diff^n(E,F)\) & II App. II & \hyperlink{sga6.e02.para.1-2-2}{1.2.2}\\
\(d(L)\) & XIII & \hyperlink{sga6.e13.para.2-2}{2.2}\\
\(d^n_{X/S}\) & II App. II & \hyperlink{sga6.e02.para.1-2-1}{1.2.1}\\
\(d^n_{X/S}(E)\) & II App. II & \hyperlink{sga6.e02.para.1-2-2}{1.2.2}\\
\(\Dcat^{-}(\Qcoh(S))_{\coh},\ \Dcat(\Qcoh(S))_{\parf}\) & II & \hyperlink{sga6.e02.notation.2-2-1}{2.2.1}\\
\(\Dcat(S),\ \Dcat(\underline S),\ \Dcat(\underline S,A)\) & I & \hyperlink{sga6.e01.ex.2-15}{2.15}\\
\(\Dcat(X)_{f\text{-}\parf},\ \Dcat(X)_{\underline{f\text{-}\mathrm{parf}}}\) & III & \hyperlink{sga6.e03.notation.4-2}{4.2}\\
\(\Dcat(X)_{f\text{-}\torf}\) & III & \hyperlink{sga6.e03.notation.3-2}{3.2}\\
\(\Dcat(X)_{\qcoh}\) & III & \hyperlink{sga6.e03.lem.1-7}{1.7}\\
\(\Dcat(X)_{Y\text{-}\coh},\ \Dcat(X)_{Y\text{-}n\text{-}\coh}\) & III & \hyperlink{sga6.e03.def.1-2}{1.2}\\
\(\Dcat(X)_{Y\text{-}\parf},\ \Dcat(X)_{\underline{Y\text{-}\mathrm{parf}}}\) & III & \hyperlink{sga6.e03.notation.4-2}{4.2}\\
\(\Dcat(X)_{Y\text{-}\torf}\) & III & \hyperlink{sga6.e03.notation.3-2}{3.2}\\
\(\Dcat(Y)_{n\text{-}\coh}\) & III & \hyperlink{sga6.e03.thm.1-1}{1.1}\\
\(E^\vee\) & I & \hyperlink{sga6.e01.def.7-3}{7.3}\\
\(*\) & V & \hyperlink{sga6.e05.para.6-1}{6.1}\\
\(E|\) & X & \hyperlink{sga6.e10.prop.4-3-1}{4.3.1}\\
\(f^*\) (between the \(K^\bullet\)) & IV & \hyperlink{sga6.e04.para.2-7}{2.7}\\
\(f^*\) (between the \(K_\bullet\)) & IV & \hyperlink{sga6.e04.para.2-12}{2.12}\\
\(f_*\) (between the \(K_\bullet\)) & IV & \hyperlink{sga6.e04.para.2-11}{2.11}\\
\(f_*\) (between the \(K^\bullet\)) & IV & \hyperlink{sga6.e04.para.2-12}{2.12}\\
\(f_*\) (projective-bundle case) & VI & \hyperlink{sga6.e06.sec.5.2}{5.2}\\
\(f_{\operatorname{gr}}\) & VIII & \hyperlink{sga6.e08.prop.3-2}{3.2}\\
\(\Fil_n(X)\) & X & \hyperlink{sga6.e10.def.1-1-1}{1.1.1}\\
\(F^{\prime >n},\ F'_{\leq n}\) & I & \hyperlink{sga6.e01.prop.2-12}{2.12}\\
\(\varphi^a\) & V & \hyperlink{sga6.e05.para.2-9}{2.9}\\
\(\circ,\ \widehat G_K\) & V & \hyperlink{sga6.e05.para.1-1}{1.1}\\
\(\circ,\ \widehat G_{\nabla}\) & V & \hyperlink{sga6.e05.para.3-7}{3.7}\\
\(\widehat G_a\) & V & \hyperlink{sga6.e05.para.1-3}{1.3.3}\\
\(\widehat G_g\) & V & \hyperlink{sga6.e05.para.1-3}{1.3.2}\\
\(\operatorname{Gr}^\bullet(X),\ \operatorname{Gr}_\bullet(X)\) & X & \hyperlink{sga6.e10.sec.3.1}{3.1}\\
\(\operatorname{Gr}^\bullet_{\alg},\ \operatorname{Gr}^{\alg}_\bullet\) & X App. & \hyperlink{sga6.e10.para.7-11}{7.11}\\
\(\operatorname{Gr}^\bullet_{\num},\ \operatorname{Gr}^{\num}_\bullet\) & X App. & \hyperlink{sga6.e10.para.7-9-4}{7.9.4}\\
\(\gamma^n,\ \gamma_t\) & V & \hyperlink{sga6.e05.prop.3-2}{3.2}\\
\(\gamma^p(N,x)\) & V & \hyperlink{sga6.e05.para.5-5}{5.5}\\
\(\widehat G^+,\ \widehat G^u\) & V & \hyperlink{sga6.e05.para.2-3}{2.3}\\
\(\underline{\Hom}_{\Lambda}(E,F)\) & I & \hyperlink{sga6.e01.notation.4-5}{4.5}\\
\(k(A)\) & IV & \hyperlink{sga6.e04.para.1-5}{1.5}\\
\(K^\bullet_{\alg},\ K^{\alg}_\bullet\) & X App. & \hyperlink{sga6.e10.para.7-11}{7.11}\\
\(K^\bullet(A_X),\ K_\bullet(A_X)\) & IV & \hyperlink{sga6.e04.def.2-2}{2.2}\\
\(K^\bullet(A_X)_{\naif},\ K_\bullet(A_X)_{\naif}\) & IV & \hyperlink{sga6.e04.def.2-2}{2.2}\\
\(k(C)\) & I & \hyperlink{sga6.e01.para.6-3}{6.3}\\
\(K^\bullet(f),\ K_\bullet(f)\) & IV & \hyperlink{sga6.e04.para.3-3-1}{3.3.1}\\
\(K^\bullet_{\num},\ K^{\num}_\bullet\) & X App. & \hyperlink{sga6.e10.para.7-9-4}{7.9.4}\\
\(K^\bullet(X),\ K_\bullet(X)\) & IV & \hyperlink{sga6.e04.def.2-2}{2.2}\\
\(K^\bullet(X)_{\naif},\ K_\bullet(X)_{\naif}\) & IV & \hyperlink{sga6.e04.def.2-2}{2.2}\\
\(K^\bullet(X/Y),\ K_\bullet(X/Y)\) & IV & \hyperlink{sga6.e04.para.3-3-1}{3.3.1}\\
\(\Lib(U)\) & II & \hyperlink{sga6.e02.sec.2.0}{2.0}\\
\(\Loclib(U)\) & II & \hyperlink{sga6.e02.sec.2.0}{2.0}\\
\(L^{X/Y}\) & VIII & \hyperlink{sga6.e08.para.2-1}{2.1}\\
\(\lambda^i\) & V & \hyperlink{sga6.e05.def.2-1}{2.1}\\
\(\lambda_{-1}(N)\) & V & \hyperlink{sga6.e05.para.5-1}{5.1}\\
\(\lambda_{-1}(N)\) & VII & \hyperlink{sga6.e07.para.2-6}{2.7}\\
\(\lambda^p(N,x)\) & V & \hyperlink{sga6.e05.def.5-3}{5.3}\\
\(\lambda_t\) & V & \hyperlink{sga6.e05.def.2-1}{2.1}\\
\(\Mod(U)\) & II & \hyperlink{sga6.e02.sec.2.0}{2.0}\\
\(\parfamp_f(E),\ \parfamp_Y(E)\) & III & \hyperlink{sga6.e03.def.4-1}{4.1}\\
\(\Parf(A_S),\ \Parf(\mathcal C,\mathcal C_0),\ \Parf(S)\) & I & \hyperlink{sga6.e01.notation.4-9}{4.9}\\
\(\Pic^N(X)\) & X & \hyperlink{sga6.e10.prop.4-5-2}{4.5.2}\\
\(\Pic^\circ(X)\) & X & \hyperlink{sga6.e10.def.2-4-1}{2.4.1}\\
\(\underline{\Pic}^{\circ}_{X/k}\) & X & \hyperlink{sga6.e10.sec.2.4}{2.4}\\
\(\Pic^\tau(X)\) & X & \hyperlink{sga6.e10.def.2-4-1}{2.4.1}\\
\(\Pic(X)\) & X & \hyperlink{sga6.e10.sec.2.4}{2.4}\\
\(\underline{\Pic}_{X/k}\) & X & \hyperlink{sga6.e10.sec.2.4}{2.4}\\
\(P^n_{X/S}\) & II App. II & \hyperlink{sga6.e02.para.1-2-1}{1.2.1}\\
\(P^n_{X/S}(E)\) & II App. II & \hyperlink{sga6.e02.para.1-2-2}{1.2.2}\\
\(\Proj\) & VI & \hyperlink{sga6.e06.para.2-3-2}{2.3.2}\\
\(\psi_k(x)\) & V & \hyperlink{sga6.e05.para.7-1}{7.1}\\
\(\Qcoh(S)\) & II & \hyperlink{sga6.e02.notation.2-2-1}{2.2.1}\\
\(r^0_{\mathcal O_{U\times V}}\) & II App. II & \hyperlink{sga6.e02.sec.1.1}{1.1}\\
\(\rho_t\) & V & \hyperlink{sga6.e05.para.3-4}{3.4}\\
\(\circ\) & V & \hyperlink{sga6.e05.para.2-3}{2.3}\\
\(T_f\) & VIII & \hyperlink{sga6.e08.cor.2-5}{2.5}\\
\(\Todd\) & V & \hyperlink{sga6.e05.para.1-16}{1.16}\\
\(\toramp(F)\) & I & \hyperlink{sga6.e01.def.5-2}{5.2}\\
\(\toramp_f(E),\ \toramp_Y(E)\) & III & \hyperlink{sga6.e03.def.3-1}{3.1}\\
\(\Tr_{E|U}\) & I & \hyperlink{sga6.e01.para.8-1}{8.1}\\
\(T_{X/S}\) & II App. II & \hyperlink{sga6.e02.para.1-2-1}{1.2.1}\\
\(\tau\) & I & \hyperlink{sga6.e01.para.8-1}{8.1}\\
\(\nabla\) & V & \hyperlink{sga6.e05.para.3-5}{3.5}\\
\(1+A[[t]]^+\) & V & \hyperlink{sga6.e05.para.1-1}{1.1}\\
\(1+\widehat B^+\) & V & \hyperlink{sga6.e05.para.1-11}{1.11}\\
\(X(n)\) & X & \hyperlink{sga6.e10.prop.1-1-2}{1.1.2}\\
\(x_N\) & V & \hyperlink{sga6.e05.para.5-5}{5.5}\\
\(\langle x,y\rangle\) & X & \hyperlink{sga6.e10.sec.4.1}{4.1}\\
\(Z(x)\) & X & \hyperlink{sga6.e10.prop.1-1-2}{1.1.2}\\
\bottomrule
\end{longtable}
\endgroup

\ifnum\SGAReaderForm=1
\clearpage
\section*{Apparatus}
\noindent\begin{minipage}{\linewidth}
\SGARestoreAppState{EN-VAR-001}%
\paragraph{[1] EN-VAR-001}
Printed reading: \(f:X.\longrightarrow Y\)\par
Corrected reading: \(f:X\longrightarrow Y\)\par
{\raggedright Source: fixes/MISPRINTS\_6.json\#2; check/findings/findings\_001\_005.json\#6. The distinct baseline stop after X is printed before the arrow; corrected form removes it.\par}
\end{minipage}\par\medskip
\noindent\begin{minipage}{\linewidth}
\SGARestoreAppState{EN-VAR-002}%
\paragraph{[2] EN-VAR-002}
Printed reading: \(\ch_Y\bigl(f_*(\cl(F))=f_*\bigl(\Todd(T_f)\cdot\ch_X(F)\bigr),\)\par
Corrected reading: \(\ch_Y\bigl(f_*(\cl(F))\bigr)=f_*\bigl(\Todd(T_f)\cdot\ch_X(F)\bigr),\)\par
{\raggedright Source: fixes/MISPRINTS\_6.json\#4; check/findings/findings\_001\_005.json\#10. Three left parentheses have only two corresponding right parentheses before the equality sign in the print; the corrected reading closes the argument of ch\_Y.\par}
\end{minipage}\par\medskip
\noindent\begin{minipage}{\linewidth}
\SGARestoreAppState{EN-VAR-003}%
\paragraph{[3] EN-VAR-003}
Printed reading: ($j=i$)\par
Corrected reading: ($j\leq i$)\par
{\raggedright Source: fixes/MISPRINTS\_6.json\#12; check/findings/findings\_001\_005.json\#30. Print reads j=i. The coefficient of t\textasciicircum{}i in lambda\_\{t/(1-t)\} is a linear combination of lambda\textasciicircum{}j for j from 1 to i, including lower-degree terms; corrected range is j\ensuremath{\leq}i.\par}
\end{minipage}\par\medskip
\noindent\begin{minipage}{\linewidth}
\SGARestoreAppState{EN-VAR-004}%
\paragraph{[4] EN-VAR-004}
Printed reading: \(f_*(\cl^\bullet(\underline E))=\sum_i(-1)^i\cl^\bullet\bigl(R^if_*(\underline F)\bigr),\)\par
Corrected reading: \(f_*(\cl^\bullet(\underline F))=\sum_i(-1)^i\cl^\bullet\bigl(R^if_*(\underline F)\bigr),\)\par
{\raggedright Source: fixes/MISPRINTS\_6.json\#21; check/findings/findings\_006\_010.json\#4. The printed left member is an underlined E, while the right member and subsequent explanatory text use F. Corrected E to F.\par}
\end{minipage}\par\medskip
\noindent\begin{minipage}{\linewidth}
\SGARestoreAppState{EN-VAR-005}%
\paragraph{[5] EN-VAR-005}
Printed reading: \(K^\bullet(X)\longrightarrow K(\Parf(Y)),\)\par
Corrected reading: \(K^\bullet(Y)\longrightarrow K(\Parf(Y)),\)\par
{\raggedright Source: fixes/MISPRINTS\_6.json\#22; check/findings/findings\_006\_010.json\#9. The printed domain is K\textasciicircum{}bullet(X). The canonical construction in this paragraph sends a locally free Module on Y to its perfect complex on Y; the corrected domain is K\textasciicircum{}bullet(Y).\par}
\end{minipage}\par\medskip
\noindent\begin{minipage}{\linewidth}
\SGARestoreAppState{EN-VAR-006}%
\paragraph{[6] EN-VAR-006}
Printed reading: \(d_{X'/X}:\underline J\subset\underline O_{X'}\longrightarrow\Omega^1_{X'/Y},\)\par
Corrected reading: \(d_{X'/Y}:\underline J\subset\underline O_{X'}\longrightarrow\Omega^1_{X'/Y},\)\par
{\raggedright Source: fixes/MISPRINTS\_6.json\#45; check/findings/findings\_011\_015.json\#2. The printed subscript is X-prime/X. The differential into relative differentials over Y is the universal Y-relative derivation restricted to J, so the corrected subscript is X-prime/Y.\par}
\end{minipage}\par\medskip
\noindent\begin{minipage}{\linewidth}
\SGARestoreAppState{EN-VAR-007}%
\paragraph{[7] EN-VAR-007}
Printed reading: Grauert's finiteness theorem [4]\par
Corrected reading: Grauert's finiteness theorem [3]\par
{\raggedright Source: fixes/MISPRINTS\_6.json\#55; check/findings/findings\_011\_015.json\#20. The print gives reference[4]; directly checked bibliography entry[3] is Grauert, whereas[4] is Grothendieck, RRR.\par}
\end{minipage}\par\medskip
\noindent\begin{minipage}{\linewidth}
\SGARestoreAppState{EN-VAR-008}%
\paragraph{[8] EN-VAR-008}
Printed reading: Shih [7]\par
Corrected reading: Shih [6]\par
{\raggedright Source: fixes/MISPRINTS\_6.json\#56; check/findings/findings\_011\_015.json\#22. The print gives reference[7]; directly checked bibliography entry[6] is Shih.\par}
\end{minipage}\par\medskip
\noindent\begin{minipage}{\linewidth}
\SGARestoreAppState{EN-VAR-009}%
\paragraph{[9] EN-VAR-009}
Printed reading: \(\operatorname{Filt}^iK^\bullet(X)\cdot \operatorname{Filt}_j(K_\bullet(X)\subset \operatorname{Filt}_{j-i}K_\bullet(X),\)\par
Corrected reading: \(\operatorname{Filt}^iK^\bullet(X)\cdot \operatorname{Filt}_j K_\bullet(X)\subset \operatorname{Filt}_{j-i}K_\bullet(X),\)\par
{\raggedright Source: fixes/MISPRINTS\_6.json\#61; check/findings/findings\_011\_015.json\#31. The print opens a parenthesis immediately after Filt\_j and never closes it; the corrected form removes that superfluous opening parenthesis.\par}
\end{minipage}\par\medskip
\noindent\begin{minipage}{\linewidth}
\SGARestoreAppState{EN-VAR-010}%
\paragraph{[10] EN-VAR-010}
Printed reading: is reduced by passage to the associated graded objects from\par
Corrected reading: is deduced by passage to the associated graded objects from\par
{\raggedright Source: ASTRA-EARLY-0004. The print reads se réduit. Passing the filtered projection identity to its associated quotients gives the graded identity; the corrected verb is se déduit. The English reproduces the same verb difference.\par}
\end{minipage}\par\medskip
\noindent\begin{minipage}{\linewidth}
\SGARestoreAppState{EN-VAR-011}%
\paragraph{[11] EN-VAR-011}
Printed reading: announced without proof in [6]\par
Corrected reading: announced without proof in [5]\par
{\raggedright Source: fixes/MISPRINTS\_6.json\#39; check/findings/findings\_016\_020.json\#4. The printed reference is[6]; FGA is bibliography entry[5], while[6] is Shih.\par}
\end{minipage}\par\medskip
\noindent\begin{minipage}{\linewidth}
\SGARestoreAppState{EN-VAR-012}%
\paragraph{[12] EN-VAR-012}
Printed reading: is [6, C-07 (i)]\par
Corrected reading: is [5, C-07 (i)]\par
{\raggedright Source: fixes/MISPRINTS\_6.json\#40; check/findings/findings\_016\_020.json\#5. The printed reference begins[6]. The identified FGA item belongs to bibliography entry[5].\par}
\end{minipage}\par\medskip
\noindent\begin{minipage}{\linewidth}
\SGARestoreAppState{EN-VAR-013}%
\paragraph{[13] EN-VAR-013}
Printed reading: Différential Geometry\par
Corrected reading: Differential Geometry\par
{\raggedright Source: fixes/MISPRINTS\_6.json\#43; check/findings/findings\_016\_020.json\#21. The English publication title is printed with an acute accent in Différential; the corrected form removes the accent.\par}
\end{minipage}\par\medskip
\noindent\begin{minipage}{\linewidth}
\SGARestoreAppState{EN-VAR-014}%
\paragraph{[14] EN-VAR-014}
Printed reading: Bull. Soc. math France\par
Corrected reading: Bull. Soc. math. France\par
{\raggedright Source: fixes/MISPRINTS\_6.json\#44; check/findings/findings\_016\_020.json\#23. The print has no abbreviation stop after math; the corrected form supplies it.\par}
\end{minipage}\par\medskip
\noindent\begin{minipage}{\linewidth}
\SGARestoreAppState{EN-VAR-015}%
\paragraph{[15] EN-VAR-015}
Printed reading: \(\lambda^n(\lambda^i(x))=P_{i,n}(\lambda^1x,\ldots,\lambda^n x),\)\par
Corrected reading: \(\lambda^n(\lambda^i(x))=P_{i,n}(\lambda^1x,\ldots,\lambda^{in} x),\)\par
{\raggedright Source: fixes/MISPRINTS\_6.json\#36; review/early/decisions.jsonl. The print has final exponent n. Formula1.8 on the same page defines P\_\{i,n\} in variables through a\_\{in\}; substituting a\_j=lambda\textasciicircum{}j(x) gives the final argument lambda\textasciicircum{}\{in\}(x).\par}
\end{minipage}\par\medskip
\noindent\begin{minipage}{\linewidth}
\SGARestoreAppState{EN-VAR-016}%
\paragraph{[16] EN-VAR-016}
Printed reading: \(x=1+\sum_{i\geq 1}x^i_i\in1+\Ahp\)\par
Corrected reading: \(x=1+\sum_{i\geq 1}x^i\in1+\Ahp\)\par
{\raggedright Source: fixes/MISPRINTS\_6.json\#18; check/findings/findings\_026\_030.json\#6; review/early/decisions.jsonl. The printed homogeneous component has an extra lower index i. The immediately following specification gives x\textasciicircum{}i in A\textasciicircum{}i; the corrected form removes only this extra index.\par}
\end{minipage}\par\medskip
\noindent\begin{minipage}{\linewidth}
\SGARestoreAppState{EN-VAR-017}%
\paragraph{[17] EN-VAR-017}
Printed reading: Groupes de Lie albébriques\par
Corrected reading: Groupes de Lie algébriques\par
{\raggedright Source: fixes/MISPRINTS\_6.json\#16; fr/codex/HOLD\_VS\_CORR.json\#28; review/early/decisions.jsonl. The title is retained in French in the English translation. The scan prints albébriques; the corrected title restores algébriques.\par}
\end{minipage}\par\medskip
\noindent\begin{minipage}{\linewidth}
\SGARestoreAppState{EN-VAR-018}%
\paragraph{[18] EN-VAR-018}
Printed reading: \(\widetilde A\relbar\joinrel\relbar\joinrel\relbar\)\par
Corrected reading: \(\widetilde A\longrightarrow\)\par
{\raggedright Source: fixes/MISPRINTS\_6.json\#66; check/findings/findings\_031\_035.json\#3. The primary scan and native witness42 show a plain rule with no arrowhead. The corrected form supplies the arrow for the homomorphism ch stated in the surrounding text.\par}
\end{minipage}\par\medskip
\noindent\begin{minipage}{\linewidth}
\SGARestoreAppState{EN-VAR-019}%
\paragraph{[19] EN-VAR-019}
Printed reading: \(\sum_{i=0}^{n}(-1)^i\lambda^i(N)\)\par
Corrected reading: \(\sum_{i=0}^{q}(-1)^i\lambda^i(N)\)\par
{\raggedright Source: fixes/MISPRINTS\_6.json\#69; check/findings/findings\_031\_035.json\#12. The print has upper limit n. This proposition specifies N with rank q and components through N\textasciicircum{}q; proposition1.2 gives vanishing of lambda\textasciicircum{}i(N) for i>q. The corrected upper limit is q.\par}
\end{minipage}\par\medskip
\noindent\begin{minipage}{\linewidth}
\SGARestoreAppState{EN-VAR-020}%
\paragraph{[20] EN-VAR-020}
Printed reading: \(\underline{E}_G((\Lambda(N^{(p)})^{\mathrm{alt}})\)\par
Corrected reading: \(\underline{E}_G((\Lambda(N^{(p)}))^{\mathrm{alt}})\)\par
{\raggedright Source: fixes/MISPRINTS\_6.json\#71; check/findings/findings\_031\_035.json\#25. The print leaves the inner grouping around Lambda(N\textasciicircum{}(p)) unclosed before alt. The corrected form adds that closing parenthesis, matching the alternating-part grouping in formula1.33.\par}
\end{minipage}\par\medskip
\noindent\begin{minipage}{\linewidth}
\SGARestoreAppState{EN-VAR-021}%
\paragraph{[21] EN-VAR-021}
Printed reading: (\S 2, theorem 1.3)\par
Corrected reading: (\S 2, theorem 1.1)\par
{\raggedright Source: ASTRA-EARLY-0011; review/early/decisions.jsonl. The print cites theorem1.3 in section2. The representation-ring statement there is theorem1.1 on p26 and its corollary on p27; proposition1.3 on p32 is in section3 and concerns the Chern/Todd identity.\par}
\end{minipage}\par\medskip
\noindent\begin{minipage}{\linewidth}
\SGARestoreAppState{EN-VAR-022}%
\paragraph{[22] EN-VAR-022}
Printed reading: \(\binom{n-1}{i}\lambda^i(N,x)\)\par
Corrected reading: \(\binom{n-1}{i}\lambda^{n-i}(N,x)\)\par
{\raggedright Source: ASTRA-EARLY-0012; review/early/decisions.jsonl. The print gives exponent i. Formula1.23 identifies gamma\textasciicircum{}n(x\_N) with lambda\textasciicircum{}n(x\_N+n\ensuremath{-}1); the coefficient of t\textasciicircum{}n in (1+t)\textasciicircum{}(n\ensuremath{-}1)lambda\_t(x\_N) gives exponent n\ensuremath{-}i. The n=1 printed case instead gives lambda\textasciicircum{}0=1.\par}
\end{minipage}\par\medskip
\noindent\begin{minipage}{\linewidth}
\SGARestoreAppState{EN-VAR-023}%
\paragraph{[23] EN-VAR-023}
Printed reading: \(j\geq0\)\par
Corrected reading: \(j\geq1\)\par
{\raggedright Source: ASTRA-EARLY-0013; review/early/decisions.jsonl. The print gives j\ensuremath{\geq}0. The stated component formula1.38 contains (j\ensuremath{-}1)!, and its proof uses formula1.18 in positive degrees. The corrected range j\ensuremath{\geq}1 is the common domain of this formula and that proof.\par}
\end{minipage}\par\medskip
\noindent\begin{minipage}{\linewidth}
\SGARestoreAppState{EN-VAR-024}%
\paragraph{[24] EN-VAR-024}
Printed reading: \(x^1,..,x^j;N^1,..,N^j\)\par
Corrected reading: \(x^1,\ldots,x^j;N^1,\ldots,N^j\)\par
{\raggedright Source: fixes/MISPRINTS\_6.json\#64; check/findings/findings\_036\_040.json\#5. The primary scan and native witness48 show two dots at each of these two ellipses. The corrected form supplies three dots.\par}
\end{minipage}\par\medskip
\noindent\begin{minipage}{\linewidth}
\SGARestoreAppState{EN-VAR-025}%
\paragraph{[25] EN-VAR-025}
Printed reading: \(x^1,..,N^j\)\par
Corrected reading: \(x^1,\ldots,N^j\)\par
{\raggedright Source: fixes/MISPRINTS\_6.json\#65; check/findings/findings\_036\_040.json\#6. The print has two dots at this ellipsis, as in formula1.38. The corrected form supplies three dots.\par}
\end{minipage}\par\medskip
\noindent\begin{minipage}{\linewidth}
\SGARestoreAppState{EN-VAR-026}%
\paragraph{[26] EN-VAR-026}
Printed reading: \(\underline{K}(X)\)\par
Corrected reading: \(\underline{K}(\xi(X))\)\par
{\raggedright Source: fixes/MISPRINTS\_6.json\#73; review/early/decisions.jsonl. Both scans repeat K(X) in the phrase introducing its abbreviation. Formula2.6 names the same group as K(xi(X)); the corrected second occurrence restores xi(X).\par}
\end{minipage}\par\medskip
\noindent\begin{minipage}{\linewidth}
\SGARestoreAppState{EN-VAR-027}%
\paragraph{[27] EN-VAR-027}
Printed reading: \(A(Y)\)\par
Corrected reading: \(\Atil(Y)\)\par
{\raggedright Source: ASTRA-EARLY-0014; review/early/decisions.jsonl. Formula2.6 gives the codomain of C-tilde as A-tilde. On this construction the induced map sends [n,1+sum a\_i] to [n,1+sum f-star(a\_i)], as used with2.9. The corrected lower-left object restores the omitted tilde.\par}
\end{minipage}\par\medskip
\noindent\begin{minipage}{\linewidth}
\SGARestoreAppState{EN-VAR-028}%
\paragraph{[28] EN-VAR-028}
Printed reading: \(A(X)\)\par
Corrected reading: \(\Atil(X)\)\par
{\raggedright Source: ASTRA-EARLY-0014; review/early/decisions.jsonl. Formula2.6 gives the codomain of C-tilde as A-tilde. The induced map on this construction sends [n,1+sum a\_i] to [n,1+sum f-star(a\_i)]. The corrected lower-right object restores the omitted tilde.\par}
\end{minipage}\par\medskip
\noindent\begin{minipage}{\linewidth}
\SGARestoreAppState{EN-VAR-029}%
\paragraph{[29] EN-VAR-029}
Printed reading: (which is possible by (i) and (ii);\par
Corrected reading: (which is possible by (i) and (ii));\par
{\raggedright Source: fixes/MISPRINTS\_6.json\#77; review/early/decisions.jsonl. The printed parenthetical explanation does not close before the semicolon. The corrected reading supplies the close there, before the next step of the proof.\par}
\end{minipage}\par\medskip
\noindent\begin{minipage}{\linewidth}
\SGARestoreAppState{EN-VAR-030}%
\paragraph{[30] EN-VAR-030}
Printed reading: $\gamma(\Gamma)$ likewise denotes\par
Corrected reading: $\gamma(F)$ likewise denotes\par
{\raggedright Source: fixes/MISPRINTS\_6.json\#80; check/findings/findings\_046\_050.json\#9. The primary and witness scans show capital Gamma inside gamma, while the same sentence names the sheaf F. The corrected occurrence uses F.\par}
\end{minipage}\par\medskip
\noindent\begin{minipage}{\linewidth}
\SGARestoreAppState{EN-VAR-031}%
\paragraph{[31] EN-VAR-031}
Printed reading: \(\underline{K}(\xi(X)\)\par
Corrected reading: \(\underline{K}(\xi(X))\)\par
{\raggedright Source: fixes/MISPRINTS\_6.json\#81; check/findings/findings\_046\_050.json\#16. The print has only one closing parenthesis after X, leaving K open. The corrected form supplies the second close.\par}
\end{minipage}\par\medskip
\noindent\begin{minipage}{\linewidth}
\SGARestoreAppState{EN-VAR-032}%
\paragraph{[32] EN-VAR-032}
Printed reading: is not of characteristic $\neq0$\par
Corrected reading: is not of characteristic $0$\par
{\raggedright Source: fixes/MISPRINTS\_6.json\#82; check/findings/findings\_046\_050.json\#19. The print contains both the prose negation and the inequality sign. The preceding clause limits the construction to characteristic0; removing the inequality sign states the corresponding obstacle while retaining the printed prose negation.\par}
\end{minipage}\par\medskip
\noindent\begin{minipage}{\linewidth}
\SGARestoreAppState{EN-VAR-033}%
\paragraph{[33] EN-VAR-033}
Printed reading: \(\underline{\Tor}_i(F,..,F)\)\par
Corrected reading: \(\underline{\Tor}_i(F,\ldots,F)\)\par
{\raggedright Source: fixes/MISPRINTS\_6.json\#83; check/findings/findings\_046\_050.json\#22. The primary scan and native witness61 show two dots in this N.B. expression; the display2.19 immediately above has three. The corrected form supplies three dots here.\par}
\end{minipage}\par\medskip
\noindent\begin{minipage}{\linewidth}
\SGARestoreAppState{EN-VAR-034}%
\paragraph{[34] EN-VAR-034}
Printed reading: example 2, for the definition\par
Corrected reading: example 1, for the definition\par
{\raggedright Source: ASTRA-EARLY-0015; review/early/decisions.jsonl. ChapterI section2 example1 on p24 defines K(G) from representations. Example2 on p25 concerns vector bundles on X. The printed reference2 is corrected to1.\par}
\end{minipage}\par\medskip
\noindent\begin{minipage}{\linewidth}
\SGARestoreAppState{EN-VAR-035}%
\paragraph{[35] EN-VAR-035}
Printed reading: \((fg)=f_!g_!\)\par
Corrected reading: \((fg)_!=f_!g_!\)\par
{\raggedright Source: fixes/MISPRINTS\_6.json\#84; check/findings/findings\_051\_055.json\#11. Both scans leave the area below (fg) blank while marking f and g on the right. The corrected form supplies the lower shriek for the direct image of the composite.\par}
\end{minipage}\par\medskip
\noindent\begin{minipage}{\linewidth}
\SGARestoreAppState{EN-VAR-036}%
\paragraph{[36] EN-VAR-036}
Printed reading: \(x.\in\underline K(X)\)\par
Corrected reading: \(x\in\underline K(X)\)\par
{\raggedright Source: fixes/MISPRINTS\_6.json\#85; check/findings/findings\_051\_055.json\#14. The print has a separate baseline stop after x before the membership sign. The corrected reading removes that stop.\par}
\end{minipage}\par\medskip
\noindent\begin{minipage}{\linewidth}
\SGARestoreAppState{EN-VAR-037}%
\paragraph{[37] EN-VAR-037}
Printed reading: canonical injection from $Y$ to $X$\par
Corrected reading: canonical injection from $X$ to $Y$\par
{\raggedright Source: fixes/MISPRINTS\_6.json\#86; check/findings/findings\_051\_055.json\#18. The same sentence defines X as a subspace of Y and formula2.28 maps the class on X to its class on Y. Thus the injection is i:X\ensuremath{\to}Y, correcting the printed reversed direction.\par}
\end{minipage}\par\medskip
\noindent\begin{minipage}{\linewidth}
\SGARestoreAppState{EN-VAR-038}%
\paragraph{[38] EN-VAR-038}
Printed reading: \(\gamma_{X\times k}(D)-p^!(y)\)\par
Corrected reading: \(\gamma_{X\times k}(Y)-p^!(y)\)\par
{\raggedright Source: ASTRA-EARLY-0016; review/early/decisions.jsonl. The divisor D is defined only on U\ensuremath{\times}k as the restriction of Y. For j:U\ensuremath{\times}k\ensuremath{\to}X\ensuremath{\times}k, restriction j-shriek sends the global class gamma(Y) to the local class gamma(D). The global difference therefore uses Y.\par}
\end{minipage}\par\medskip
\noindent\begin{minipage}{\linewidth}
\SGARestoreAppState{EN-VAR-039}%
\paragraph{[39] EN-VAR-039}
Printed reading: \(\gamma_{X\times k}(D)-p^!(y))\)\par
Corrected reading: \(\gamma_{X\times k}(Y)-p^!(y)\)\par
{\raggedright Source: ASTRA-EARLY-0016; fixes/MISPRINTS\_6.json\#87; check/findings/findings\_051\_055.json\#22. The same restriction j-shriek sends the global class of Y to the local divisor D, so the global argument is Y. The print also gives a second unmatched closing parenthesis after y; the corrected expression removes it.\par}
\end{minipage}\par\medskip
\noindent\begin{minipage}{\linewidth}
\SGARestoreAppState{EN-VAR-040}%
\paragraph{[40] EN-VAR-040}
Printed reading: a section of $V$\par
Corrected reading: a section of $G$\par
{\raggedright Source: fixes/MISPRINTS\_6.json\#88; review/early/decisions.jsonl. The print has V. The definition of G-bar(V) pulls back G(U intersection V) along restriction from F(V); the section therefore belongs to G.\par}
\end{minipage}\par\medskip
\noindent\begin{minipage}{\linewidth}
\SGARestoreAppState{EN-VAR-041}%
\paragraph{[41] EN-VAR-041}
Printed reading: filtration $\leq i$\par
Corrected reading: filtration $\geq i$\par
{\raggedright Source: ASTRA-EARLY-0017; review/early/decisions.jsonl. The print gives \ensuremath{\leq}i. By the p48 definition, classes supported in dimension \ensuremath{\leq}n\ensuremath{-}i generate K(X)\textasciicircum{}i. Exactness of2.30 identifies the kernel of restriction with those supported on Y; the stated condition is membership in K(X)\textasciicircum{}i, hence filtration \ensuremath{\geq}i.\par}
\end{minipage}\par\medskip
\noindent\begin{minipage}{\linewidth}
\SGARestoreAppState{EN-VAR-042}%
\paragraph{[42] EN-VAR-042}
Printed reading: \(f_a^!=f_a^!\)\par
Corrected reading: \(f_a^!=f_{a'}^!\)\par
{\raggedright Source: fixes/MISPRINTS\_6.json\#89; review/early/decisions.jsonl. Both scans omit the prime on the second a. The preceding formulas compare the fibres over a and a-prime, so the corrected second subscript is a-prime.\par}
\end{minipage}\par\medskip
\noindent\begin{minipage}{\linewidth}
\SGARestoreAppState{EN-VAR-043}%
\paragraph{[43] EN-VAR-043}
Printed reading: \(\varphi:A(X)\longrightarrow G(K(X))\)\par
Corrected reading: \(\varphi:A(X)\longrightarrow G(\underline K(X))\)\par
{\raggedright Source: fixes/MISPRINTS\_6.json\#90; review/early/decisions.jsonl. This printed K lacks the underline present in the next explanatory line. The corrected form supplies the underline for the stated Grothendieck group.\par}
\end{minipage}\par\medskip
\noindent\begin{minipage}{\linewidth}
\SGARestoreAppState{EN-VAR-044}%
\paragraph{[44] EN-VAR-044}
Printed reading: from $\underline K(Y)$ to $K(X)$\par
Corrected reading: from $\underline K(Y)$ to $\underline K(X)$\par
{\raggedright Source: fixes/MISPRINTS\_6.json\#91; review/early/decisions.jsonl. The target K(X) lacks the underline printed under the source K(Y). The corrected form supplies the target underline.\par}
\end{minipage}\par\medskip
\noindent\begin{minipage}{\linewidth}
\SGARestoreAppState{EN-VAR-045}%
\paragraph{[45] EN-VAR-045}
Printed reading: whether the $\gamma_P(V_{i;\alpha})$ form a basis\par
Corrected reading: whether the $\gamma_P(V_{i,\alpha})$ form a basis\par
{\raggedright Source: fixes/MISPRINTS\_6.json\#92; review/early/decisions.jsonl. The print has a semicolon in this index pair, while the preceding statement defines it with a comma. The corrected form restores that defined index pair.\par}
\end{minipage}\par\medskip
\noindent\begin{minipage}{\linewidth}
\SGARestoreAppState{EN-VAR-046}%
\paragraph{[46] EN-VAR-046}
Printed reading: \(\eta_j^{(nj)}\)\par
Corrected reading: \(\eta_j^{(n_j)}\)\par
{\raggedright Source: fixes/MISPRINTS\_6.json\#95; review/early/decisions.jsonl. In the print, j is on the same raised baseline as n in the final exponent. The corrected form restores the indexed dimension n\_j defined in the statement.\par}
\end{minipage}\par\medskip
\noindent\begin{minipage}{\linewidth}
\SGARestoreAppState{EN-VAR-047}%
\paragraph{[47] EN-VAR-047}
Printed reading: \(R((\lambda^k(x'_i)),(\lambda^k(\gamma(E'_j)))\)\par
Corrected reading: \(R((\lambda^k(x'_i)),(\lambda^k(\gamma(E'_j))))\)\par
{\raggedright Source: fixes/MISPRINTS\_6.json\#96; review/early/decisions.jsonl. The print has only three closing parentheses after E-prime-j, leaving the outer R argument open. The corrected expression supplies the fourth close.\par}
\end{minipage}\par\medskip
\noindent\begin{minipage}{\linewidth}
\SGARestoreAppState{EN-VAR-048}%
\paragraph{[48] EN-VAR-048}
Printed reading: \(\gamma^k(E_i-\varepsilon(E_i))\)\par
Corrected reading: \(\gamma^k(E_j-\varepsilon(E_j))\)\par
{\raggedright Source: fixes/MISPRINTS\_6.json\#97; review/early/decisions.jsonl. The print indexes both E terms by i. The statement indexes the bundles E\_j by j from1 to r, separately from the x\_i indexed from1 to q; the corrected terms use j.\par}
\end{minipage}\par\medskip
\noindent\begin{minipage}{\linewidth}
\SGARestoreAppState{EN-VAR-049}%
\paragraph{[49] EN-VAR-049}
Printed reading: \(G(K(X))\longrightarrow G(A)\)\par
Corrected reading: \(G(\underline K(X))\longrightarrow G(A)\)\par
{\raggedright Source: fixes/MISPRINTS\_6.json\#98; review/early/decisions.jsonl. The K in this printed graded map lacks the underline present on K(X) in the preceding line. The corrected form supplies it. The printed target G(A) is retained.\par}
\end{minipage}\par\medskip
\noindent\begin{minipage}{\linewidth}
\SGARestoreAppState{EN-VAR-050}%
\paragraph{[50] EN-VAR-050}
Printed reading: \(x\in K(X)\)\par
Corrected reading: \(x\in\underline K(X)\)\par
{\raggedright Source: fixes/MISPRINTS\_6.json\#101; review/early/decisions.jsonl. This printed K lacks the underline present on the later graded terms in the same theorem. The corrected form supplies the underline.\par}
\end{minipage}\par\medskip
\noindent\begin{minipage}{\linewidth}
\SGARestoreAppState{EN-VAR-051}%
\paragraph{[51] EN-VAR-051}
Printed reading: \(\sum_{i\geq\varphi}c^i(x)\)\par
Corrected reading: \(\sum_{i\geq1}c^i(x)\)\par
{\raggedright Source: fixes/MISPRINTS\_6.json\#102; review/early/decisions.jsonl. The printed lower bound is phi. Since c\textasciicircum{}0 is the class of gamma\textasciicircum{}0=1, the leading1 already accounts for degree0. The ChapterI section3 construction has a sum of positive-degree terms, so the corrected bound is1.\par}
\end{minipage}\par\medskip
\noindent\begin{minipage}{\linewidth}
\SGARestoreAppState{EN-VAR-052}%
\paragraph{[52] EN-VAR-052}
Printed reading: \(\gamma^i(x-\varepsilon(X))\)\par
Corrected reading: \(\gamma^i(x-\varepsilon(x))\)\par
{\raggedright Source: fixes/MISPRINTS\_6.json\#104; review/early/decisions.jsonl. The print gives capital X as the augmentation argument. The augmentation is applied to the element x, as in theorem2.16; the corrected argument is lowercase x.\par}
\end{minipage}\par\medskip
\noindent\begin{minipage}{\linewidth}
\SGARestoreAppState{EN-VAR-053}%
\paragraph{[53] EN-VAR-053}
Printed reading: \(\Ctil(x)\in A(X)^i\)\par
Corrected reading: \(\Ctil(x)\in\Atil(X)^i\)\par
{\raggedright Source: ASTRA-EARLY-0018; review/early/decisions.jsonl. The printed target lacks its tilde. Formula2.7 gives the lambda-homomorphism C-tilde:K(X)\ensuremath{\to}A-tilde(X), whose filtration compatibility is stated on62. The degree-i class of its gamma-expression is the map psi used here, and1.22bis concerns this A-tilde construction.\par}
\end{minipage}\par\medskip
\noindent\begin{minipage}{\linewidth}
\SGARestoreAppState{EN-VAR-054}%
\paragraph{[54] EN-VAR-054}
Printed reading: \(i_!(xx'\lambda_{-1}(N)\)\par
Corrected reading: \(i_!(xx'\lambda_{-1}(N))\)\par
{\raggedright Source: fixes/MISPRINTS\_6.json\#105; review/early/decisions.jsonl. The print closes the parenthesis of lambda\_\{\ensuremath{-}1\}(N) but leaves the argument of i-shriek open. The corrected expression supplies its closing parenthesis.\par}
\end{minipage}\par\medskip
\noindent\begin{minipage}{\linewidth}
\SGARestoreAppState{EN-VAR-055}%
\paragraph{[55] EN-VAR-055}
Printed reading: by reduction of formulas (ii) and (iv)\par
Corrected reading: by reduction of formulas (iii) and (iv)\par
{\raggedright Source: ASTRA-EARLY-0019; review/early/decisions.jsonl. Taking associated-graded classes in formula(iii), i\textasciicircum{}!i\_!(x)=x lambda\_\{\ensuremath{-}1\}(N), gives(iii bis) because the degree-q term of lambda\_\{\ensuremath{-}1\}(N) is(\ensuremath{-}1)\textasciicircum{}q c\textasciicircum{}q(N). Formula(iv) similarly gives(iv bis); formula(ii) gives the projection formula(ii bis). Correct the first cited numeral.\par}
\end{minipage}\par\medskip
\noindent\begin{minipage}{\linewidth}
\SGARestoreAppState{EN-VAR-056}%
\paragraph{[56] EN-VAR-056}
Printed reading: \(\lambda_{-1}(c(N))\)\par
Corrected reading: \(\lambda_{-1}(\widetilde c(N))\)\par
{\raggedright Source: fixes/MISPRINTS\_6.json\#106; review/early/decisions.jsonl. Formula2.37 uses the total class c-tilde(N);2.37bis is that identity in total-class form. The printed c before(N) has lost its tilde.\par}
\end{minipage}\par\medskip
\noindent\begin{minipage}{\linewidth}
\SGARestoreAppState{EN-VAR-057}%
\paragraph{[57] EN-VAR-057}
Printed reading: the product of their difference by $N^q$\par
Corrected reading: the product of their difference by $\xi'$\par
{\raggedright Source: ASTRA-EARLY-0020; review/early/decisions.jsonl. For the difference delta on X, the preceding equality gives i\textasciicircum{}*(delta)=0. Applying i\_* and the projection formula(ii bis) gives0=i\_*i\textasciicircum{}*(delta)=i\_*(1)delta=xi-prime delta. The printed factor N\textasciicircum{}q is on Y; its relationship to xi-prime is i\textasciicircum{}*(xi-prime)=(\ensuremath{-}1)\textasciicircum{}q N\textasciicircum{}q from(iii bis). The stated pushforward therefore yields the factor xi-prime.\par}
\end{minipage}\par\medskip
\noindent\begin{minipage}{\linewidth}
\SGARestoreAppState{EN-VAR-058}%
\paragraph{[58] EN-VAR-058}
Printed reading: \(\ch_X\bigl(f_*(x)\bigr)\)\par
Corrected reading: \(\ch_X\bigl(f_!(x)\bigr)\)\par
{\raggedright Source: ASTRA-EARLY-0021; review/early/decisions.jsonl. The right-hand printed map has a star. Formula2.26 and the following paragraph identify f\_! as the map from K(Y) to K(X), which is the domain of ch\_X. The left-hand f\_* acts on the Chow-valued class and keeps its star.\par}
\end{minipage}\par\medskip
\noindent\begin{minipage}{\linewidth}
\SGARestoreAppState{EN-VAR-059}%
\paragraph{[59] EN-VAR-059}
Printed reading: each of these two morphisms..\par
Corrected reading: each of these two morphisms.\par
{\raggedright Source: fixes/MISPRINTS\_6.json\#108; review/early/decisions.jsonl. The printed sentence has two adjacent full stops. The corrected layer retains one.\par}
\end{minipage}\par\medskip
\noindent\begin{minipage}{\linewidth}
\SGARestoreAppState{EN-VAR-060}%
\paragraph{[60] EN-VAR-060}
Printed reading: the dimension of \(Q\) in \(P\)\par
Corrected reading: the codimension of \(Q\) in \(P\)\par
{\raggedright Source: fixes/MISPRINTS\_6.json\#109; review/early/decisions.jsonl. The class xi is the hyperplane structure-sheaf class. The transverse intersection of p hyperplanes represents xi\textasciicircum{}p=gamma\_P(Q), so p is the codimension of Q.\par}
\end{minipage}\par\medskip
\noindent\begin{minipage}{\linewidth}
\SGARestoreAppState{EN-VAR-061}%
\paragraph{[61] EN-VAR-061}
Printed reading: \(q=r+1-p\)\par
Corrected reading: \(q=r-p\)\par
{\raggedright Source: fixes/MISPRINTS\_6.json\#110; review/early/decisions.jsonl. With H=L/(1\ensuremath{-}e\textasciicircum{}\{\ensuremath{-}L\}), the displayed product is L\textasciicircum{}p H\textasciicircum{}\{r+1\ensuremath{-}p\}. Its L\textasciicircum{}r coefficient equals the L\textasciicircum{}\{r\ensuremath{-}p\} coefficient of H\textasciicircum{}\{r\ensuremath{-}p+1\}; the stated coefficient identity therefore uses q=r\ensuremath{-}p.\par}
\end{minipage}\par\medskip
\noindent\begin{minipage}{\linewidth}
\SGARestoreAppState{EN-VAR-062}%
\paragraph{[62] EN-VAR-062}
Printed reading: K\"Unneth\par
Corrected reading: K\"unneth\par
{\raggedright Source: ASTRA-EARLY-TYPE-016; review/early/decisions.jsonl. The main image and native witnesses81/86 show the tall U-umlaut form on pages70 and75. Preserve the printed name-case KÜnneth diplomatically; the corrected spelling is Künneth.\par}
\end{minipage}\par\medskip
\noindent\begin{minipage}{\linewidth}
\SGARestoreAppState{EN-VAR-063}%
\paragraph{[63] EN-VAR-063}
Printed reading: Neue topologischen Methoden\par
Corrected reading: Neue topologische Methoden\par
{\raggedright Source: fixes/MISPRINTS\_6.json\#111; review/early/decisions.jsonl. The printed German book title has topologischen. Before plural Methoden without an article, the adjective takes the ending-e, as Neue does. The title is retained in German in the English edition, so its printed spelling is represented explicitly.\par}
\end{minipage}\par\medskip
\noindent\begin{minipage}{\linewidth}
\SGARestoreAppState{EN-VAR-064}%
\paragraph{[64] EN-VAR-064}
Printed reading: the functor \(f\) thus defined\par
Corrected reading: the functor \(f_*\) thus defined\par
{\raggedright Source: ASTRA-EARLY-0022; review/early/decisions.jsonl. The preceding operation sends a sheaf G to f\_*(G); its right derived functors are the displayed R\textasciicircum{}p f\_*. The printed functor name loses the lower star. Restore it in this occurrence, retaining the morphism f and the K-group map f\_! elsewhere.\par}
\end{minipage}\par\medskip
\noindent\begin{minipage}{\linewidth}
\SGARestoreAppState{EN-VAR-065}%
\paragraph{[65] EN-VAR-065}
Printed reading: K\"Unneth\par
Corrected reading: K\"unneth\par
{\raggedright Source: fixes/MISPRINTS\_6.json\#115; ASTRA-EARLY-TYPE-016; review/early/decisions.jsonl. The main scan and native witness86 show the tall U-umlaut form, matching page70. Keep the printed KÜnneth diplomatically and the corrected Künneth in current.\par}
\end{minipage}\par\medskip
\noindent\begin{minipage}{\linewidth}
\SGARestoreAppState{EN-VAR-066}%
\paragraph{[66] EN-VAR-066}
Printed reading: \((Y',X',y\lambda_{-1}(F))\)\par
Corrected reading: \((Y',X',y\lambda_{-1}(\check F))\)\par
{\raggedright Source: fixes/MISPRINTS\_6.json\#116; review/early/decisions.jsonl. The earlier invocation of Riemann–Roch and formula13 both use the checked bundle F. The class supplied by the corollary is y lambda\_\{\ensuremath{-}1\}(check F); the second invocation uses that same class. Restore its omitted check.\par}
\end{minipage}\par\medskip
\noindent\begin{minipage}{\linewidth}
\SGARestoreAppState{EN-VAR-067}%
\paragraph{[67] EN-VAR-067}
Printed reading: the class $T_f$ defined by formula (2.1)\par
Corrected reading: the class $T_f$ defined by formula (2.2)\par
{\raggedright Source: ASTRA-EARLY-0002; review/early/decisions.jsonl. Formula2.1 is the factorization X\ensuremath{\to}X-prime\ensuremath{\to}Y. Formula2.2 defines the virtual cotangent class check T\_f, with T\_f its dual in the immediately following sentence. The reference specifying this class therefore points to2.2.\par}
\end{minipage}\par\medskip
\noindent\begin{minipage}{\linewidth}
\SGARestoreAppState{EN-VAR-068}%
\paragraph{[68] EN-VAR-068}
Printed reading: (SGA 1 VI 7,1)\par
Corrected reading: (SGA 1 VI 7.1)\par
{\raggedright Source: fixes/MISPRINTS\_6.json\#94; check/findings/findings\_081\_085.json\#10; review/expose\_i/decisions.jsonl. The printed separator has a descending comma tail, while the nearby reference6.1 uses a baseline point. The diplomatic branch retains7,1 and the current branch uses the reference separator7.1.\par}
\end{minipage}\par\medskip
\noindent\begin{minipage}{\linewidth}
\SGARestoreAppState{EN-VAR-069}%
\paragraph{[69] EN-VAR-069}
Printed reading: \[
\begin{tikzcd}
& H_U \arrow[d] \arrow[dl,"u"']\\
E_U \arrow[r,"U"] & F_U
\end{tikzcd}
\]\par
Corrected reading: \[
\begin{tikzcd}
& H_U \arrow[d] \arrow[dl]\\
E_U \arrow[r,"u_U"] & F_U
\end{tikzcd}
\]\par
{\raggedright Source: check/findings/findings\_086\_090.json\#6; review/expose\_i/decisions.jsonl. The printed diagonal has label u and the horizontal label is U. The diagonal is the chosen lift H\_U\ensuremath{\to}E\_U; the horizontal is the pullback u\_U:E\_U\ensuremath{\to}F\_U of u. Current removes the diagonal label and writes u\_U horizontally.\par}
\end{minipage}\par\medskip
\noindent\begin{minipage}{\linewidth}
\SGARestoreAppState{EN-VAR-070}%
\paragraph{[70] EN-VAR-070}
Printed reading: (SGA 4 V 3.11) \(k\)\par
Corrected reading: (SGA 4 V 3.11), \(k\)\par
{\raggedright Source: check/findings/findings\_086\_090.json\#15. The comma separates the topos datum from the following field datum in example c). Its absence is retained diplomatically.\par}
\end{minipage}\par\medskip
\noindent\begin{minipage}{\linewidth}
\SGARestoreAppState{EN-VAR-071}%
\paragraph{[71] EN-VAR-071}
Printed reading: \textbf{Lemma 1.3.1,}\par
Corrected reading: \textbf{Lemma 1.3.1.}\par
{\raggedright Source: fixes/MISPRINTS\_6.json\#123; check/findings/findings\_086\_090.json\#17. The scan has a descending comma after the lemma number. Current uses the full stop separating its label from its statement.\par}
\end{minipage}\par\medskip
\noindent\begin{minipage}{\linewidth}
\SGARestoreAppState{EN-VAR-072}%
\paragraph{[72] EN-VAR-072}
Printed reading: \uline{-Module, Equivalent conditions}\par
Corrected reading: \uline{-Module. Equivalent conditions}\par
{\raggedright Source: fixes/MISPRINTS\_6.json\#124; check/findings/findings\_086\_090.json\#18. Both scans print a comma before the new sentence introducing the equivalent conditions. Current supplies a sentence stop.\par}
\end{minipage}\par\medskip
\noindent\begin{minipage}{\linewidth}
\SGARestoreAppState{EN-VAR-073}%
\paragraph{[73] EN-VAR-073}
Printed reading: \(S\) \uline{be a site,}\par
Corrected reading: \(\underline{S}\) \uline{be a site,}\par
{\raggedright Source: ASTRA-I-S-089-PROP. The printed plain S denotes the site governing the fibered categories. The statement later declares the object S to belong to the category underlined-S, identifying that base site. Current restores its individual underline; diplomatic preserves the printed glyph.\par}
\end{minipage}\par\medskip
\noindent\begin{minipage}{\linewidth}
\SGARestoreAppState{EN-VAR-074}%
\paragraph{[74] EN-VAR-074}
Printed reading: \(C\) \uline{a flat abelian}\par
Corrected reading: \(\underline{C}\) \uline{a flat abelian}\par
{\raggedright Source: fixes/MISPRINTS\_6.json\#125; check/findings/findings\_086\_090.json\#23. The printed plain C occurs in the declaration or subsequent mention of the same flat abelian category over the site. Current restores its category underline; diplomatic preserves the plain glyph.\par}
\end{minipage}\par\medskip
\noindent\begin{minipage}{\linewidth}
\SGARestoreAppState{EN-VAR-075}%
\paragraph{[75] EN-VAR-075}
Printed reading: \(S\)\uline{-category (1.1),}\par
Corrected reading: \(\underline{S}\)\uline{-category (1.1),}\par
{\raggedright Source: ASTRA-I-S-089-PROP. The printed plain S denotes the site governing the fibered categories. The statement later declares the object S to belong to the category underlined-S, identifying that base site. Current restores its individual underline; diplomatic preserves the printed glyph.\par}
\end{minipage}\par\medskip
\noindent\begin{minipage}{\linewidth}
\SGARestoreAppState{EN-VAR-076}%
\paragraph{[76] EN-VAR-076}
Printed reading: \uline{a strictly full additive sub-}\(S\)\par
Corrected reading: \uline{a strictly full additive sub-}\(\underline{S}\)\par
{\raggedright Source: ASTRA-I-S-089-PROP. The printed plain S denotes the site governing the fibered categories. The statement later declares the object S to belong to the category underlined-S, identifying that base site. Current restores its individual underline; diplomatic preserves the printed glyph.\par}
\end{minipage}\par\medskip
\noindent\begin{minipage}{\linewidth}
\SGARestoreAppState{EN-VAR-077}%
\paragraph{[77] EN-VAR-077}
Printed reading: \(C\)\uline{; assume that}\par
Corrected reading: \(\underline{C}\)\uline{; assume that}\par
{\raggedright Source: fixes/MISPRINTS\_6.json\#126; check/findings/findings\_086\_090.json\#24. The printed plain C occurs in the declaration or subsequent mention of the same flat abelian category over the site. Current restores its category underline; diplomatic preserves the plain glyph.\par}
\end{minipage}\par\medskip
\noindent\begin{minipage}{\linewidth}
\SGARestoreAppState{EN-VAR-078}%
\paragraph{[78] EN-VAR-078}
Printed reading: \(C\) \uline{(1.2). Let}\par
Corrected reading: \(\underline{C}\) \uline{(1.2). Let}\par
{\raggedright Source: fixes/MISPRINTS\_6.json\#127; check/findings/findings\_086\_090.json\#25. The printed plain C occurs in the declaration or subsequent mention of the same flat abelian category over the site. Current restores its category underline; diplomatic preserves the plain glyph.\par}
\end{minipage}\par\medskip
\noindent\begin{minipage}{\linewidth}
\SGARestoreAppState{EN-VAR-079}%
\paragraph{[79] EN-VAR-079}
Printed reading: \[
\begin{tikzpicture}[>=stealth,baseline=(GX.base),node distance=1.0cm]
  \node (EX) at (0,0) {$E_X$};
  \node (Ep) at (2.6,0) {$E'$};
  \node (FX) at (5.2,0) {$F_X$};
  \node (M)  at (1.35,1.25) {$M[-n]$};
  \node (GX) at (2.65,2.55) {$G_X$};
  \node (Gp) at (3.95,1.25) {$G'$};
  \node at (3.22,1.72) {$(*)$};
  \draw[->] (EX) -- node[below] {$v$} (Ep);
  \draw (Ep) -- node[below] {$u'$} (FX);
  \draw[->] (Ep) -- (M);
  \draw[->] (M) -- (GX);
  \draw[->] (GX) -- (Gp);
  \draw[->] (FX) -- (Gp);
  \draw[->] (M) to[bend left=38] node[left] {$+1$} (EX);
  \draw[->] (Gp) to[bend left=28] node[above] {$+1$} (M);
  \draw[->] (Gp) to[bend left=28] node[right] {$+1$} (Ep);
  \draw[->] (GX) to[bend left=55] node[left] {$+1$} (EX);
  \draw[->] (FX) to[bend right=42] (GX);
  \node[anchor=base] at (6.3,0 |- FX.base) {.};
\end{tikzpicture}
\]\par
Corrected reading: \[
\begin{tikzpicture}[>=stealth,baseline=(GX.base),node distance=1.0cm]
  \node (EX) at (0,0) {$E_X$};
  \node (Ep) at (2.6,0) {$E'$};
  \node (FX) at (5.2,0) {$F_X$};
  \node (M)  at (1.35,1.25) {$M[-n]$};
  \node (GX) at (2.65,2.55) {$G_X$};
  \node (Gp) at (3.95,1.25) {$G'$};
  \node at (3.22,1.72) {$(*)$};
  \draw[->] (EX) -- node[below] {$v$} (Ep);
  \draw[->] (Ep) -- node[below] {$u'$} (FX);
  \draw[->] (Ep) -- (M);
  \draw[->] (M) -- (GX);
  \draw[->] (GX) -- (Gp);
  \draw[->] (FX) -- (Gp);
  \draw[->] (M) to[bend left=38] node[left] {$+1$} (EX);
  \draw[->] (Gp) to[bend left=28] node[above] {$+1$} (M);
  \draw[->] (Gp) to[bend left=28] node[right] {$+1$} (Ep);
  \draw[->] (GX) to[bend left=55] node[left] {$+1$} (EX);
  \draw[->] (FX) to[bend right=42] (GX);
  \node[anchor=base] at (6.3,0 |- FX.base) {.};
\end{tikzpicture}
\]\par
{\raggedright Source: check/findings/findings\_091\_095.json\#5; review/expose\_i/decisions.jsonl. The printed u' segment has no arrowhead. The preceding text defines u':E'\ensuremath{\to}F\_X and G' as its cone, which fixes the arrow direction toward F\_X. Current adds that head; diplomatic retains the bare segment.\par}
\end{minipage}\par\medskip
\noindent\begin{minipage}{\linewidth}
\SGARestoreAppState{EN-VAR-080}%
\paragraph{[80] EN-VAR-080}
Printed reading: \textbf{Lemma 1.4.2.} \uline{Let} \(C\)\par
Corrected reading: \textbf{Lemma 1.4.2.} \uline{Let} \(\underline{C}\)\par
{\raggedright Source: check/findings/findings\_091\_095.json\#9. The first declared category C has no printed underline. Its occurrence inside the derived category D(C) is underlined in the same sentence; current supplies the same category rule here.\par}
\end{minipage}\par\medskip
\noindent\begin{minipage}{\linewidth}
\SGARestoreAppState{EN-VAR-081}%
\paragraph{[81] EN-VAR-081}
Printed reading: \uline{a morphism of complexes of} \(C\)\par
Corrected reading: \uline{a morphism of complexes of} \(\underline{C}\)\par
{\raggedright Source: check/findings/findings\_091\_095.json\#9. The second category C has no printed underline. It is the category of the complexes E and F and the argument of D(C), whose C is individually underlined; current supplies its missing rule.\par}
\end{minipage}\par\medskip
\noindent\begin{minipage}{\linewidth}
\SGARestoreAppState{EN-VAR-082}%
\paragraph{[82] EN-VAR-082}
Printed reading: \(v^j=(v_i)^j\)\par
Corrected reading: \(v^j=(v_i v_{i+1}\ldots v_{n-1})^j\)\par
{\raggedright Source: ASTRA-I-COMPOSITION-094; review/expose\_i/add\_094\_composition.py; EN\_INDEPENDENT\_SOURCE\_ASSESSMENTS.jsonl. The transitions have domains E\ensuremath{\to}E'\_(n-1) and E'\_(i+1)\ensuremath{\to}E'\_i. Their composite V\_i=v\_i...v\_(n-1) has domain E, and its degree-j component stabilizes for i\ensuremath{\leq}j because the lower transitions are identities in that degree. The printed isolated v\_i has domain E'\_(i+1), so the composite is required to define v:E\ensuremath{\to}E'. The u' formula remains unchanged.\par}
\end{minipage}\par\medskip
\noindent\begin{minipage}{\linewidth}
\SGARestoreAppState{EN-VAR-083}%
\paragraph{[83] EN-VAR-083}
Printed reading: \(E^i\in\Ob\underline{C}_{0S}\)\par
Corrected reading: \(E^i\in\Ob\underline{C}_{0X}\)\par
{\raggedright Source: check/findings/findings\_091\_095.json\#19; review/expose\_i/decisions.jsonl. The definition places E in the fiber C\_X at its arbitrary object X. Its components therefore belong to C\_X and the specified strictly full subcategory is C\_(0X). Current changes the printed subscript S to X.\par}
\end{minipage}\par\medskip
\noindent\begin{minipage}{\linewidth}
\SGARestoreAppState{EN-VAR-084}%
\paragraph{[84] EN-VAR-084}
Printed reading: \textbf{Proposition 2.7}\par
Corrected reading: \textbf{Proposition 2.7.}\par
{\raggedright Source: check/findings/findings\_096\_100.json\#19. The scan has a blank after the final7 before the statement. Diplomatic preserves that missing stop; current retains the restored heading punctuation.\par}
\end{minipage}\par\medskip
\noindent\begin{minipage}{\linewidth}
\SGARestoreAppState{EN-VAR-085}%
\paragraph{[85] EN-VAR-085}
Printed reading: \(n\)\par
Corrected reading: \((-n)\)\par
{\raggedright Source: ASTRA-I-SIGN-103. Assertion2.9(ii) specifies (-n)-pseudo-coherence. Apply2.4 to the map0\ensuremath{\to}F with pseudo-coherence parameter -n and quasi-isomorphism parameter1: the replacement E\ensuremath{\to}F is strictly(-n)-pseudo-coherent, and0\ensuremath{\to}E is the identity in degrees\ensuremath{\geq}1. Thus E vanishes there and its components in degrees -n through0 supply exactly the n-presentation defined in2.8. Current supplies the missing minus sign.\par}
\end{minipage}\par\medskip
\noindent\begin{minipage}{\linewidth}
\SGARestoreAppState{EN-VAR-086}%
\paragraph{[86] EN-VAR-086}
Printed reading: \(F\) \uline{is pseudo-coherent.}\par
Corrected reading: \(F\) \uline{is} \(n\)\uline{-pseudo-coherent (resp.\ pseudo-coherent).}\par
{\raggedright Source: check/findings/findings\_101\_105.json\#20; review/expose\_i/decisions.jsonl. The proof on105 concludes n-pseudo-coherence. Its induction uses H\textasciicircum{}i(F')=H\textasciicircum{}(i+1)(F), and the presentation index i+1-n equals i-(n-1). The printed unqualified conclusion fails on the punctual site of vector spaces: take n=0, an infinite-dimensional vector space M, and F=M[1]. All H\textasciicircum{}i(F) for i\ensuremath{\geq}0 vanish and meet the stated finite-presentation assumptions, while H\textasciicircum{}(-1)(F)=M is not of finite type. Proposition2.10(b) therefore rules out (-1)-pseudo-coherence. Current restores the n-qualified conclusion and the parenthetical alternative proved by the argument.\par}
\end{minipage}\par\medskip
\noindent\begin{minipage}{\linewidth}
\SGARestoreAppState{EN-VAR-087}%
\paragraph{[87] EN-VAR-087}
Printed reading: \(F^i\neq0\)\par
Corrected reading: \(F^i=0\)\par
{\raggedright Source: check/findings/findings\_101\_105.json\#21; review/expose\_i/decisions.jsonl. The given complex belongs to the bounded-above category C-minus. The displayed exact sequence removes its highest component F\textasciicircum{}(n+N), leaving F-prime zero above n+N-1. The induction therefore starts with F\textasciicircum{}i=0 for i>n+N, replacing the printed nonzero relation.\par}
\end{minipage}\par\medskip
\noindent\begin{minipage}{\linewidth}
\SGARestoreAppState{EN-VAR-088}%
\paragraph{[88] EN-VAR-088}
Printed reading: \(F'_{\leq n}\)) (resp.\par
Corrected reading: \(F'_{\leq n}\) (resp.\par
{\raggedright Source: check/findings/findings\_106\_110.json\#1. The preceding reference(2.10.b)) is already closed. The printed extra closer after F-prime has no opener; current removes that single character.\par}
\end{minipage}\par\medskip
\noindent\begin{minipage}{\linewidth}
\SGARestoreAppState{EN-VAR-089}%
\paragraph{[89] EN-VAR-089}
Printed reading: and \(\underline{C}_1\)).\par
Corrected reading: and \(\underline{C}_1\).\par
{\raggedright Source: check/findings/findings\_106\_110.json\#23. Both numbered references in this sentence are already closed. The final printed parenthesis after C\_1 has no opener; current removes it and preserves every parameter substitution.\par}
\end{minipage}\par\medskip
\noindent\begin{minipage}{\linewidth}
\SGARestoreAppState{EN-VAR-090}%
\paragraph{[90] EN-VAR-090}
Printed reading: \uline{is a field or when} \(\underline{\operatorname{radical}}(A_s)\) \uline{is nilpotent.}\par
Corrected reading: \uline{is a field.}\par
{\raggedright Source: ASTRA-I-RADICAL-109; review/expose\_i/decisions.jsonl; EN\_INDEPENDENT\_SOURCE\_ASSESSMENTS.jsonl. The printed nilpotence-only alternative is an example for the stated hypothesis that finite projective stalk modules are free. For A=k\ensuremath{\times}k, its radical J=0 is nilpotent, but P=k\ensuremath{\times}0 is the image of the idempotent e(a,b)=(a,0). Tensoring along the two coordinate projections gives dimensions1 and0, whereas a free A-module of rank r gives r and r, so P is not free. The retained residue-field example is valid: lift a basis of P/JP to a surjection A\textasciicircum{}r\ensuremath{\to}P by Nakayama; projectivity splits A\textasciicircum{}r=P plus K. The chosen basis gives K/JK=0, and K is finitely generated, hence K=0 by Nakayama. Current deletes only the unsupported alternative; the lemma's actual projective-freeness hypothesis is unchanged.\par}
\end{minipage}\par\medskip
\noindent\begin{minipage}{\linewidth}
\SGARestoreAppState{EN-VAR-091}%
\paragraph{[91] EN-VAR-091}
Printed reading: Suppose (ii), Let\par
Corrected reading: Suppose (ii). Let\par
{\raggedright Source: check/findings/findings\_106\_110.json\#32. Both scans show the comma followed by a capital. The diplomatic reading retains it; current separates the two sentences with a full stop.\par}
\end{minipage}\par\medskip
\noindent\begin{minipage}{\linewidth}
\SGARestoreAppState{EN-VAR-092}%
\paragraph{[92] EN-VAR-092}
Printed reading: \(U_n\)\par
Corrected reading: \(u_n\)\par
{\raggedright Source: check/findings/findings\_106\_110.json\#34. The matrix(T\_ij) defines one universal endomorphism. The morphism s satisfies s*(u\_n)=p, and its factorization s=i composed with s-prime gives p=s-prime* i*(u\_n). All three lowercase references identify that same map. Current gives its definition the lowercase name u\_n as well; diplomatic retains the printed capital U\_n.\par}
\end{minipage}\par\medskip
\noindent\begin{minipage}{\linewidth}
\SGARestoreAppState{EN-VAR-093}%
\paragraph{[93] EN-VAR-093}
Printed reading: \(X'=u(X)\)\par
Corrected reading: \(X'=f^{-1}(X)\)\par
{\raggedright Source: check/findings/findings\_111\_115.json\#5; review/expose\_i/decisions.jsonl. The site-object inverse-image functor f\textasciicircum{}(-1):S\ensuremath{\to}S' is explicitly defined on110. The proof on111 applies that same map to U, writing U'=f\textasciicircum{}(-1)(U). The letter u in that proof denotes the derived tensor functor on complexes. Current applies the specified site map to X, giving X'=f\textasciicircum{}(-1)(X).\par}
\end{minipage}\par\medskip
\noindent\begin{minipage}{\linewidth}
\SGARestoreAppState{EN-VAR-094}%
\paragraph{[94] EN-VAR-094}
Printed reading: \textbf{Lemma 3.3}\par
Corrected reading: \textbf{Lemma 3.3.}\par
{\raggedright Source: check/findings/findings\_111\_115.json\#17. The printed heading ends after the numeral3. Diplomatic keeps the missing terminal stop; current supplies the heading punctuation.\par}
\end{minipage}\par\medskip
\noindent\begin{minipage}{\linewidth}
\SGARestoreAppState{EN-VAR-095}%
\paragraph{[95] EN-VAR-095}
Printed reading: (SGAA VIII 7.8)\par
Corrected reading: (SGA 4 VIII 7.8)\par
{\raggedright Source: ASTRA-I-SGA4-118; review/expose\_i/decisions.jsonl; coordinator primary bibliography verification. The printed final A is retained diplomatically. The cited definition of a point belongs to SGA 4, VIII7.8, checked against the primary text by the coordinator; current supplies the volume digit and preserves the cited exposé and paragraph.\par}
\end{minipage}\par\medskip
\noindent\begin{minipage}{\linewidth}
\SGARestoreAppState{EN-VAR-096}%
\paragraph{[96] EN-VAR-096}
Printed reading: By (2.22 a)),\par
Corrected reading: By (3.7 a)),\par
{\raggedright Source: check/findings/findings\_116\_120.json\#22; review/expose\_i/decisions.jsonl. The displayed comparison is the degree-zero consequence of the flat point pullback of the derived Hom comparison in3.7(a). The point inverse-image functor is exact, and applying H\textasciicircum{}0 to that isomorphism gives the two displayed Hom comparisons. Current cites3.7(a); diplomatic retains2.22(a).\par}
\end{minipage}\par\medskip
\noindent\begin{minipage}{\linewidth}
\SGARestoreAppState{EN-VAR-097}%
\paragraph{[97] EN-VAR-097}
Printed reading: \textbf{Proposition 4.10}\par
Corrected reading: \textbf{Proposition 4.10.}\par
{\raggedright Source: fixes/MISPRINTS\_6.json\#144; check/findings/findings\_121\_125.json\#22. The main scan has a space after the numeral before the statement, with no terminal full stop. Diplomatic retains that absence; current supplies the heading punctuation.\par}
\end{minipage}\par\medskip
\noindent\begin{minipage}{\linewidth}
\SGARestoreAppState{EN-VAR-098}%
\paragraph{[98] EN-VAR-098}
Printed reading: \(X\in\Ob S\)\par
Corrected reading: \(X\in\Ob \underline{S}\)\par
{\raggedright Source: ASTRA-I-SITE-MISSING-123; review/expose\_i/decisions.jsonl. The printed membership leaves S without a rule. The quantified X indexes the fibers C\_\{0X\}, and the fiber projection C\_0\ensuremath{\to}the base site sends their objects to X. Paragraph4.0 names that base site with underlined S; the proof124 repeats X\ensuremath{\in}ob underlined S for the same triangle fibers. Current supplies the missing site rule; diplomatic retains the observed plain letter.\par}
\end{minipage}\par\medskip
\noindent\begin{minipage}{\linewidth}
\SGARestoreAppState{EN-VAR-099}%
\paragraph{[99] EN-VAR-099}
Printed reading: \(\Tr(\Dpx(\underline{C}))\) \uline{to an object of}\par
Corrected reading: \(\Tr(\Dpx(\underline{C}))\)) \uline{to an object of}\par
{\raggedright Source: fixes/MISPRINTS\_6.json\#145; check/findings/findings\_121\_125.json\#34. The middle resp. parenthesis has an opening before the category expression and no outside closer in print. Current adds that single outside parenthesis; the inner category parentheses and all mathematical tokens are unchanged.\par}
\end{minipage}\par\medskip
\noindent\begin{minipage}{\linewidth}
\SGARestoreAppState{EN-VAR-100}%
\paragraph{[100] EN-VAR-100}
Printed reading: \uline{(resp.}\ \(\Tr(\Kpx^b(\underline{C}_0))\).\par
Corrected reading: \uline{(resp.}\ \(\Tr(\Kpx^b(\underline{C}_0))\)).\par
{\raggedright Source: fixes/MISPRINTS\_6.json\#146; check/findings/findings\_121\_125.json\#35. The final resp. clause independently lacks its outside closing parenthesis. The two inner category parentheses are present. Current supplies only the missing outer closer; diplomatic retains its omission.\par}
\end{minipage}\par\medskip
\noindent\begin{minipage}{\linewidth}
\SGARestoreAppState{EN-VAR-101}%
\paragraph{[101] EN-VAR-101}
Printed reading: \(\Kpx^b(C_0)\longrightarrow \Dpx(\underline{C})\)\par
Corrected reading: \(\Kpx^b(\underline{C}_0)\longrightarrow \Dpx(\underline{C})\)\par
{\raggedright Source: fixes/MISPRINTS\_6.json\#147; check/findings/findings\_126\_130.json\#3. The first inline source category lacks the C rule in both scans. The full embedding i:C\_0\ensuremath{\to}C specified in the preceding sentence induces K\textasciicircum{}b(i), followed by localization into D(C). Its source is precisely the named C\_0 with its independent rule. Current supplies this one omitted rule; diplomatic retains the plain C.\par}
\end{minipage}\par\medskip
\noindent\begin{minipage}{\linewidth}
\SGARestoreAppState{EN-VAR-102}%
\paragraph{[102] EN-VAR-102}
Printed reading: \(
  \Kpx^b(\underline{C}_0)\Kpx^{b,\varnothing}(\underline{C}_0)
  \longrightarrow
  \Kpx^-(\underline{C}_0)/\Kpx^{-,\varnothing}(\underline{C}_0)
\)\par
Corrected reading: \(
  \Kpx^b(\underline{C}_0)/\Kpx^{b,\varnothing}(\underline{C}_0)
  \longrightarrow
  \Kpx^-(\underline{C}_0)/\Kpx^{-,\varnothing}(\underline{C}_0)
\)\par
{\raggedright Source: fixes/MISPRINTS\_6.json\#148; check/findings/findings\_126\_130.json\#9. Both scans omit the source quotient slash. The inclusion of bounded complexes into bounded-above complexes sends bounded acyclic complexes to bounded-above acyclic complexes. Composing with localization therefore factors through the bounded quotient K\textasciicircum{}b(C\_0)/K\textasciicircum{}\{b,emptyset\}(C\_0), which is already displayed earlier in the remark. Current restores that one slash; diplomatic retains the adjacent expressions.\par}
\end{minipage}\par\medskip
\noindent\begin{minipage}{\linewidth}
\SGARestoreAppState{EN-VAR-103}%
\paragraph{[103] EN-VAR-103}
Printed reading: \(\perfdim(F)\subset[a,b]\)\par
Corrected reading: \(\perfamp(F)\subset[a,b]\)\par
{\raggedright Source: fixes/MISPRINTS\_6.json\#149; check/findings/findings\_126\_130.json\#22. Here F=F-prime+F-double-prime. Hypothesis4.17(ii) gives the stated perfect-amplitude inclusion, and the preceding truncations put its terms in C\_0 above a. Lemma4.16 applies with exactly that amplitude hypothesis and gives F\textasciicircum{}a locally in C\_0. Current writes parf.amp; diplomatic retains the printed parf.dim.\par}
\end{minipage}\par\medskip
\noindent\begin{minipage}{\linewidth}
\SGARestoreAppState{EN-VAR-104}%
\paragraph{[104] EN-VAR-104}
Printed reading: \(\Dpx^*(\underline{C})\) (4.10).\par
Corrected reading: \(\Dpx^*(\underline{C}')\) (4.10).\par
{\raggedright Source: fixes/MISPRINTS\_6.json\#150; check/findings/findings\_126\_130.json\#29. The displayed perfect subcategory is defined for C-prime. Definition4.9 supplies the full inclusion j:D*(C-prime)\_parf\ensuremath{\to}D*(C-prime), the identity on its objects and arrows, and4.10 applied to the primed triple gives its triangulated structure. Current adds the omitted prime to this ambient category; diplomatic retains the printed unprimed letter.\par}
\end{minipage}\par\medskip
\noindent\begin{minipage}{\linewidth}
\SGARestoreAppState{EN-VAR-105}%
\paragraph{[105] EN-VAR-105}
Printed reading: (2,15)\par
Corrected reading: (2.15)\par
{\raggedright Source: fixes/MISPRINTS\_6.json\#136; check/findings/findings\_131\_135.json\#31. The printed reference has a comma. Paragraph 2.15 defines the pseudo-coherence of Modules used in this corollary. Diplomatic preserves the observed comma; current uses the paragraph number 2.15.\par}
\end{minipage}\par\medskip
\noindent\begin{minipage}{\linewidth}
\SGARestoreAppState{EN-VAR-106}%
\paragraph{[106] EN-VAR-106}
Printed reading: (3.17.1)\par
Corrected reading: (5.8.2.1)\par
{\raggedright Source: ASTRA-I-REF-136; review/expose\_i/own\_findings.jsonl; review/expose\_i/add\_136\_140.py. Formula 5.8.2.1 specializes at B=Z\_S and G=I to the displayed vertical map Hom\_ZS(F,I) tensor\_A E\ensuremath{\to}Hom\_ZS(Hom\_A(E,F),I), sending h tensor e to (t\ensuremath{\mapsto}h(t(e))). Replacing F by F-double-prime gives the other vertical map. The next sentence invokes Lemma 5.8.2 for invertibility. Current gives this actual formula reference; diplomatic preserves the printed 3.17.1.\par}
\end{minipage}\par\medskip
\noindent\begin{minipage}{\linewidth}
\SGARestoreAppState{EN-VAR-107}%
\paragraph{[107] EN-VAR-107}
Printed reading: \uline{for every} \(X\in\Ob\underline{C}_X\)\par
Corrected reading: \uline{for every} \(X\in\Ob \underline{S}\)\par
{\raggedright Source: ASTRA-I-INDEX-142; check/findings/findings\_141\_145.json\#9. The declared family assigns f\_U:ob(C\_U) to F(U) to every U in ob(S). Its additivity is the identity f\_U(L)=f\_U(L-prime)+f\_U(L-double-prime) on a direct sum or distinguished triangle in that fibre. The selecting quantifier is therefore U in ob(S), with the source dummy variable X retained; diplomatic keeps the printed fibre-object quantifier.\par}
\end{minipage}\par\medskip
\noindent\begin{minipage}{\linewidth}
\SGARestoreAppState{EN-VAR-108}%
\paragraph{[108] EN-VAR-108}
Printed reading: \[
\begin{tikzcd}
\Ob\underline{C}_Y \arrow[r,"f_Y"] \arrow[d,"u^*"'] & F(Y) \arrow[d,"F(u)"]\\
\Ob\underline{C}_X \arrow[r,"f_X",dash] & F(X).
\end{tikzcd}
\]\par
Corrected reading: \[
\begin{tikzcd}
\Ob\underline{C}_Y \arrow[r,"f_Y"] \arrow[d,"u^*"'] & F(Y) \arrow[d,"F(u)"]\\
\Ob\underline{C}_X \arrow[r,"f_X"] & F(X).
\end{tikzcd}
\]\par
{\raggedright Source: check/findings/findings\_141\_145.json\#10. The source lower segment has no arrowhead. The declared map f\_X:ob(C\_X) to F(X) determines its rightward direction, and the square states f\_X(u-star L)=F(u)(f\_Y(L)). Current supplies that arrowhead; diplomatic preserves the headless segment. Every endpoint, other arrow and label is identical.\par}
\end{minipage}\par\medskip
\noindent\begin{minipage}{\linewidth}
\SGARestoreAppState{EN-VAR-109}%
\paragraph{[109] EN-VAR-109}
Printed reading: for \(f\in\Cart_{\underline{S}}(\underline{C}^{\is},F)\) and\par
Corrected reading: for \(f\in\Cart_{\underline{S}}(\underline{C}^{\is},\underline F)\) and\par
{\raggedright Source: check/findings/findings\_141\_145.json\#17. The printed isolated F lacks its rule. The discrete fibred category constructed on142 assigns to x in F(U) the object (U,x), with restriction arrows given by F(u); projection to x is the inverse description. A cartesian functor sends c in C\_U to (U,f\_U(c)). This explicitly identifies the intended codomain with the underlined F in both adjacent phi targets. Current restores that rule; diplomatic retains the observed plain F.\par}
\end{minipage}\par\medskip
\noindent\begin{minipage}{\linewidth}
\SGARestoreAppState{EN-VAR-110}%
\paragraph{[110] EN-VAR-110}
Printed reading: \[
\begin{tikzcd}[column sep=
large,row sep=large]
E^\vee\otL F\otL F^\vee\otL E \arrow[r,"\Id\otL\tau_E"] \arrow[d,"\Id\otimes\tau_F"'] & F^\vee\otL F \arrow[d,"\tau_F"]\\
E^\vee\otL E \arrow[r,"\tau_E"'] & A .
\end{tikzcd}
\]\par
Corrected reading: \[
\begin{tikzcd}[column sep=
large,row sep=large]
E^\vee\otL F\otL F^\vee\otL E \arrow[r,"\Id\otL\tau_E"] \arrow[d,"\Id\otL\tau_F"'] & F^\vee\otL F \arrow[d,"\tau_F"]\\
E^\vee\otL E \arrow[r,"\tau_E"'] & A .
\end{tikzcd}
\]\par
{\raggedright Source: check/findings/findings\_156\_160.json\#12; review/expose\_i/add\_156\_159.py. The square left edge lacks the L that is printed on the same contraction in the parallel triangle. Its map is beta=id\_(E-dual) tensor\textasciicircum{}L (ev\_F composed with c\_(F,F-dual)) tensor\textasciicircum{}L id\_E, from E-dual tensor\textasciicircum{}L F tensor\textasciicircum{}L F-dual tensor\textasciicircum{}L E to E-dual tensor\textasciicircum{}L E. Here c exchanges the middle factors, ev\_F:F-dual tensor\textasciicircum{}L F to A is the evaluation used in8.1, and the unit isomorphisms remove the A factor. Current restores the omitted derived-tensor exponent; diplomatic retains the ordinary tensor label. All other source cells, labels and arrows are identical.\par}
\end{minipage}\par\medskip
\noindent\begin{minipage}{\linewidth}
\SGARestoreAppState{EN-VAR-111}%
\paragraph{[111] EN-VAR-111}
Printed reading: $\underline{S}$ \uline{is} quasi-compact\par
Corrected reading: $S$ \uline{is} quasi-compact\par
{\raggedright Source: ASTRA-II-OBJECT-S-163; review/expose\_ii/add\_163.py. Both scans give the object S a rule in this sentence. The preceding declaration selects S as an object of the underlined site, and each displayed covering arrow S\_i to S has that same object as codomain. The identity of this declared object S identifies the two typographical occurrences. Current removes only this isolated object-letter rule; diplomatic preserves it. No object or covering morphism is changed.\par}
\end{minipage}\par\medskip
\noindent\begin{minipage}{\linewidth}
\SGARestoreAppState{EN-VAR-112}%
\paragraph{[112] EN-VAR-112}
Printed reading: $C'_S$\par
Corrected reading: $\underline{C}'_S$\par
{\raggedright Source: check/findings/findings\_161\_165.json\#9; review/expose\_ii/add\_163.py. The introduced C-prime has no rule in either scan at this occurrence. Its specified inclusion into C\_S identifies it with the same category that is individually underlined throughout164, including the domain of Gamma\_S and the exact-sequence hypothesis. The identity functor of this explicitly included subcategory gives that identification. Current supplies the rule only here; diplomatic retains the observed plain letter.\par}
\end{minipage}\par\medskip
\noindent\begin{minipage}{\linewidth}
\SGARestoreAppState{EN-VAR-113}%
\paragraph{[113] EN-VAR-113}
Printed reading: \(\operatorname{Hom}_{A(S)}(.,\Gamma_S(F))\)\par
Corrected reading: \(\operatorname{Hom}_{A(S)}(.,\Gamma_S(G))\)\par
{\raggedright Source: check/findings/findings\_161\_165.json\#18; fixes/FIXLOG\_6\_CURRENT.jsonl\#347. The second presentation is an exact sequence of A(S)-modules ending in Gamma\_S(F). Applying Hom\_(A(S))(-,Gamma\_S(G)) gives the bottom row of the printed diagram, whose two right-hand terms explicitly fix Gamma\_S(G). The vertical comparison sends a sheaf map F to G to its section map Gamma\_S(F) to Gamma\_S(G). Current therefore restores G in this fixed second argument. The separately inspected placeholder is a baseline period in both forms.\par}
\end{minipage}\par\medskip
\noindent\begin{minipage}{\linewidth}
\SGARestoreAppState{EN-VAR-114}%
\paragraph{[114] EN-VAR-114}
Printed reading: \[
\begin{tikzcd}[column sep=small]
0 \arrow[r] & \operatorname{Hom}_{A_S}(F,G) \arrow[r] \arrow[d] &
\operatorname{Hom}_{A_S}(A_S^{(J)},G) \arrow[r] \arrow[d] &
\operatorname{Hom}_{A_S}(A_S^{(I)},G) \arrow[d] \\
0 \arrow[r] & \operatorname{Hom}_{A(S)}(F,G) \arrow[r] &
\operatorname{Hom}_{A(S)}(A(S)^{(J)},\Gamma_S(G)) \arrow[r] &
\operatorname{Hom}_{A(S)}(A(S)^{(I)},\Gamma_S(G)).
\end{tikzcd}
\]\par
Corrected reading: \[
\begin{tikzcd}[column sep=small]
0 \arrow[r] & \operatorname{Hom}_{A_S}(F,G) \arrow[r] \arrow[d] &
\operatorname{Hom}_{A_S}(A_S^{(J)},G) \arrow[r] \arrow[d] &
\operatorname{Hom}_{A_S}(A_S^{(I)},G) \arrow[d] \\
0 \arrow[r] & \operatorname{Hom}_{A(S)}(\Gamma_S(F),\Gamma_S(G)) \arrow[r] &
\operatorname{Hom}_{A(S)}(A(S)^{(J)},\Gamma_S(G)) \arrow[r] &
\operatorname{Hom}_{A(S)}(A(S)^{(I)},\Gamma_S(G)).
\end{tikzcd}
\]\par
{\raggedright Source: check/findings/findings\_161\_165.json\#19; review/expose\_ii/add\_165.py. Applying Hom\_(A(S))(-,Gamma\_S(G)) to the displayed section presentation(**) gives the left nonzero term Hom\_(A(S))(Gamma\_S(F),Gamma\_S(G)). The vertical arrow is the map f\ensuremath{\mapsto}Gamma\_S(f) from sheaf morphisms F to G to section-module morphisms. Current supplies the two omitted section functors in that one bottom node; diplomatic retains F,G. Every other node, arrow, argument and punctuation is identical.\par}
\end{minipage}\par\medskip
\noindent\begin{minipage}{\linewidth}
\SGARestoreAppState{EN-VAR-115}%
\paragraph{[115] EN-VAR-115}
Printed reading: are visibly isomorphisms the same\par
Corrected reading: are visibly isomorphisms; the same\par
{\raggedright Source: check/findings/findings\_161\_165.json\#20. The printed sentence has no punctuation between its two finite clauses. Current inserts the separating semicolon; diplomatic preserves the absence. Both claims and their order are unchanged.\par}
\end{minipage}\par\medskip
\noindent\begin{minipage}{\linewidth}
\SGARestoreAppState{EN-VAR-116}%
\paragraph{[116] EN-VAR-116}
Printed reading: \(\mathrm D^{-}(\Qcoh(S)_{\coh}\)\par
Corrected reading: \(\mathrm D^{-}(\Qcoh(S))_{\coh}\)\par
{\raggedright Source: check/findings/findings\_166\_170.json\#18. The outer parenthesis of D-minus is unclosed in print. Notation2.2.1 defines the essential-image category D-minus(Qcoh(S))\_coh on167. Current adds only the missing outside close; diplomatic retains the observed omission.\par}
\end{minipage}\par\medskip
\noindent\begin{minipage}{\linewidth}
\SGARestoreAppState{EN-VAR-117}%
\paragraph{[117] EN-VAR-117}
Printed reading: \uline{a quasi-coherent Ideal of} $\underline{O}_S$\par
Corrected reading: \uline{a quasi-coherent Ideal of} $\underline{O}_S$.\par
{\raggedright Source: check/findings/findings\_166\_170.json\#29. The complete final item(iii-prime) ends without a stop in print. Current supplies the sentence-ending full stop; diplomatic retains its absence.\par}
\end{minipage}\par\medskip
\noindent\begin{minipage}{\linewidth}
\SGARestoreAppState{EN-VAR-118}%
\paragraph{[118] EN-VAR-118}
Printed reading: \uline{-Module} $f\in\Gamma(X,L)$\par
Corrected reading: \uline{-Module,} $f\in\Gamma(X,L)$\par
{\raggedright Source: check/findings/findings\_166\_170.json\#31. The declaration introduces the scheme X, its invertible module L, and the section f in Gamma(X,L). Current inserts the omitted comma delimiting the section declaration after the description of L. All three objects, their membership and the immersion i:X\_f to X remain unchanged.\par}
\end{minipage}\par\medskip
\noindent\begin{minipage}{\linewidth}
\SGARestoreAppState{EN-VAR-119}%
\paragraph{[119] EN-VAR-119}
Printed reading: \(Y'=X'-U\)\par
Corrected reading: \(Y'=X'-U'\)\par
{\raggedright Source: check/findings/findings\_171\_175.json\#4; fixes/MISPRINTS\_6.json\#151. In the displayed Cartesian diagram, with g:X-prime to X, the upper open is U-prime=U times\_X X-prime=g inverse(U). Therefore Y-prime=g inverse(X minus U)=X-prime minus U-prime. This exact inverse-image identity gives the corrected prime. Diplomatic retains the printed U.\par}
\end{minipage}\par\medskip
\noindent\begin{minipage}{\linewidth}
\SGARestoreAppState{EN-VAR-120}%
\paragraph{[120] EN-VAR-120}
Printed reading: assertions c)) follow\par
Corrected reading: assertions c) follow\par
{\raggedright Source: check/findings/findings\_171\_175.json\#10; fixes/MISPRINTS\_6.json\#153. The item reference has a second closing parenthesis without an opening parenthesis in print. Current removes that unmatched delimiter; the reference remains the same item c).\par}
\end{minipage}\par\medskip
\noindent\begin{minipage}{\linewidth}
\SGARestoreAppState{EN-VAR-121}%
\paragraph{[121] EN-VAR-121}
Printed reading: This completes the proof\par
Corrected reading: This completes the proof.\par
{\raggedright Source: check/findings/findings\_171\_175.json\#11; fixes/MISPRINTS\_6.json\#154. The complete final proof sentence has no following full stop in either scan. Current supplies the sentence-ending stop; diplomatic retains the absence.\par}
\end{minipage}\par\medskip
\noindent\begin{minipage}{\linewidth}
\SGARestoreAppState{EN-VAR-122}%
\paragraph{[122] EN-VAR-122}
Printed reading: \(F^i\in\Ob \underline{L}(S))\)\par
Corrected reading: \(F^i\in\Ob \underline{L}(S)\)\par
{\raggedright Source: check/findings/findings\_171\_175.json\#14; fixes/MISPRINTS\_6.json\#155. The category expression L(S) has one opening parenthesis and is already closed. Current removes the unmatched second close after this expression; diplomatic retains it. The same component F\textasciicircum{}i and category membership remain.\par}
\end{minipage}\par\medskip
\noindent\begin{minipage}{\linewidth}
\SGARestoreAppState{EN-VAR-123}%
\paragraph{[123] EN-VAR-123}
Printed reading: (4.8)\par
Corrected reading: (I 4.8)\par
{\raggedright Source: check/findings/findings\_176\_180.json\#8; fixes/MISPRINTS\_6.json\#165. ExposéI Example4.8 on122 describes the perfect amplitude for this derived category of Modules using local direct factors of finite free Modules. The parallel statement II2.2.9 on175 explicitly cites I4.8. Current supplies the omitted exposé identifier; diplomatic retains4.8.\par}
\end{minipage}\par\medskip
\noindent\begin{minipage}{\linewidth}
\SGARestoreAppState{EN-VAR-124}%
\paragraph{[124] EN-VAR-124}
Printed reading: (4.9)\par
Corrected reading: (I 4.9)\par
{\raggedright Source: check/findings/findings\_176\_180.json\#9; fixes/MISPRINTS\_6.json\#166. ExposéI Notation4.9 on122 names D(S)\_parf for the full subcategory of perfect objects used here. The parallel statement on175 has the explicit I4.9 locator. Current supplies the omitted exposé identifier without changing the indicated subcategory.\par}
\end{minipage}\par\medskip
\noindent\begin{minipage}{\linewidth}
\SGARestoreAppState{EN-VAR-125}%
\paragraph{[125] EN-VAR-125}
Printed reading: free Modules (resp.) direct factors of free Modules) of finite type\par
Corrected reading: free Modules (resp. direct factors of free Modules) of finite type\par
{\raggedright Source: check/findings/findings\_176\_180.json\#12; fixes/MISPRINTS\_6.json\#168. The printed alternative closes immediately after resp. and again after the phrase giving the alternative. Current removes only the premature first close. The functor Gamma\_S and both specified kinds of Modules remain unchanged.\par}
\end{minipage}\par\medskip
\noindent\begin{minipage}{\linewidth}
\SGARestoreAppState{EN-VAR-126}%
\paragraph{[126] EN-VAR-126}
Printed reading: Assertion c) follows from (1.2.1 b),c),d))\par
Corrected reading: Assertion c) follows from (1.2 b),c),d))\par
{\raggedright Source: ASTRA-II-REF-178-POST05; review/expose\_ii\_amendments/reference\_178\_identification.json. Both scans print the extra .1. The actual Proposition1.2(b,c,d) on161–163 supplies the bounded representative, canonical fully faithful functor and actual quasi-isomorphism used here. The printed C\_0=Lib wording is retained. The recorded specialization first takes P=Loclib and identifies C\_(1S) with P by the identity on modules and maps. For a degree-i term choose a finite-free complement Q and add the contractible disk Q\ensuremath{\to}Q in degrees i\ensuremath{-}1,i. Its projection pi and inclusion j satisfy pi j=id and dh+hd=id\ensuremath{-}j pi, with h\textasciicircum{}i=id\_Q. Descending through i=n,...,m+1 gives the stated free higher terms, and the composite F\ensuremath{\to}G\ensuremath{\to}E is the required quasi-isomorphism. This exact construction supports the locator repair without silently identifying the two printed categories.\par}
\end{minipage}\par\medskip
\noindent\begin{minipage}{\linewidth}
\SGARestoreAppState{EN-VAR-127}%
\paragraph{[127] EN-VAR-127}
Printed reading: \(\underline{\underline{R}}\Gamma_S=\mathrm D^b(S)_{\coh}\longrightarrow\mathrm D^+(A)\)\par
Corrected reading: \(\underline{\underline{R}}\Gamma_S:\mathrm D^b(S)_{\coh}\longrightarrow\mathrm D^+(A)\)\par
{\raggedright Source: check/findings/findings\_181\_185.json\#5; fixes/REPLAY\_MANUAL\_6.json\#4. The printed equals sign separates a functor name from its domain. The preceding sentence defines R Gamma\_S:D\textasciicircum{}+(S)\ensuremath{\to}D\textasciicircum{}+(A); compose it with the inclusion D\textasciicircum{}b(S)\_coh\ensuremath{\to}D\textasciicircum{}+(S) to obtain the displayed restriction. Current uses a colon for that same functor. Both readings retain the two directly observed horizontal rules below R.\par}
\end{minipage}\par\medskip
\noindent\begin{minipage}{\linewidth}
\SGARestoreAppState{EN-VAR-128}%
\paragraph{[128] EN-VAR-128}
Printed reading: \(i\in[m,n]\)\par
Corrected reading: \(i\notin[m,n]\)\par
{\raggedright Source: check/findings/findings\_181\_185.json\#10; fixes/MISPRINTS\_6.json\#171. The proof chooses E\textasciicircum{}i=0 for i>n on 182 and then truncates below m on 183. The map to the truncation preserves cohomology because the stated perfect amplitude gives vanishing below m. The resulting representative therefore vanishes outside [m,n]. Current restores notin; diplomatic preserves the printed membership.\par}
\end{minipage}\par\medskip
\noindent\begin{minipage}{\linewidth}
\SGARestoreAppState{EN-VAR-129}%
\paragraph{[129] EN-VAR-129}
Printed reading: \(E^m\in\Ob\Loclib(S))\)\par
Corrected reading: \(E^m\in\Ob\Loclib(S)\)\par
{\raggedright Source: check/findings/findings\_181\_185.json\#11; fixes/MISPRINTS\_6.json\#172. The category expression Loclib(S) already has its matching close. Current removes the additional unmatched parenthesis after it; diplomatic retains that extra close. The object and membership are unchanged.\par}
\end{minipage}\par\medskip
\noindent\begin{minipage}{\linewidth}
\SGARestoreAppState{EN-VAR-130}%
\paragraph{[130] EN-VAR-130}
Printed reading: \(\mathrm K^b(\Loclib(S))\to\mathrm D(S)\)\par
Corrected reading: \(\mathrm K^b(\Loclib(S))\to\mathrm D(A)\)\par
{\raggedright Source: check/findings/findings\_181\_185.json\#12. The preceding functor j:K\textasciicircum{}b(Loclib(S))\ensuremath{\to}D(S) is retained. The second functor is explicitly induced by Gamma\_S: degreewise sections define T:K\textasciicircum{}b(Loclib(S))\ensuremath{\to}K\textasciicircum{}b(Mod(A))\ensuremath{\to}D(A). Theorem B makes each locally free coherent term acyclic, so the canonical comparison from T to R Gamma\_S composed with j is an isomorphism. Current therefore restores the A-module target of this second functor; diplomatic retains S.\par}
\end{minipage}\par\medskip
\noindent\begin{minipage}{\linewidth}
\SGARestoreAppState{EN-VAR-131}%
\paragraph{[131] EN-VAR-131}
Printed reading: \(\mathrm D(S)_{\parf}\)\par
Corrected reading: \(\mathrm D(A)_{\parf}\)\par
{\raggedright Source: check/findings/findings\_181\_185.json\#13. This last image clause belongs to T=Gamma\_S on bounded locally free complexes. Under the canonical comparison from T to R Gamma\_S composed with j, Proposition 2.4.5(c) on 182 identifies its image with D(A)\_parf. The earlier D(S)\_parf is the image of j and remains unchanged. Current restores A only in this final image clause.\par}
\end{minipage}\par\medskip
\noindent\begin{minipage}{\linewidth}
\SGARestoreAppState{EN-VAR-132}%
\paragraph{[132] EN-VAR-132}
Printed reading: J. Frish [5]\par
Corrected reading: J. Frisch [5]\par
{\raggedright Source: check/findings/findings\_181\_185.json\#16; directly identified bibliography item[5] on220. The author note expressly cites item[5], whose directly inspected bibliography entry spells the surname Frisch and identifies the same analytic-functions paper. Current supplies the missing c in that attribution; diplomatic preserves Frish. The complete Grothendieck attribution and independent-discovery statement remain in both forms.\par}
\end{minipage}\par\medskip
\noindent\begin{minipage}{\linewidth}
\SGARestoreAppState{EN-VAR-133}%
\paragraph{[133] EN-VAR-133}
Printed reading: that $S$ is the analytic spectrum\par
Corrected reading: that $\widetilde S$ is the analytic spectrum\par
{\raggedright Source: check/findings/findings\_181\_185.json\#27. The previous sentence introduces the normalization S-tilde. Its finite normalization algebra has relative analytic spectrum equal to the domain S-tilde of the explicitly specified projection p:S-tilde\ensuremath{\to}S. Current restores the tilde on this domain name; diplomatic preserves the printed S.\par}
\end{minipage}\par\medskip
\noindent\begin{minipage}{\linewidth}
\SGARestoreAppState{EN-VAR-134}%
\paragraph{[134] EN-VAR-134}
Printed reading: (10 II \S3 Prop.8)\par
Corrected reading: ([10] II \S3 Prop.8)\par
{\raggedright Source: check/findings/findings\_181\_185.json\#28. The surrounding citations identify the same work as[10] in IV section3, IV section4 and II section3. This occurrence names Proposition8 of the last section. Current restores its bibliography-number brackets; the specified work, section and proposition are unchanged.\par}
\end{minipage}\par\medskip
\noindent\begin{minipage}{\linewidth}
\SGARestoreAppState{EN-VAR-135}%
\paragraph{[135] EN-VAR-135}
Printed reading: \(\widetilde{E_n\cdot\widetilde{\underline{O}}_S}|U=0\)\par
Corrected reading: \(\widetilde{E_n\cdot\widetilde{\underline{O}}_S}|V=0\)\par
{\raggedright Source: check/findings/findings\_181\_185.json\#31. The introduced open neighbourhood is V in S-tilde. The ideal carrying the outer tilde is a sheaf on S-tilde, and its restriction is inverse image along V\ensuremath{\to}S-tilde. The earlier U is an open neighbourhood on S; the specified projection p:S-tilde\ensuremath{\to}S relates the two spaces. Current names V in this restriction, while diplomatic preserves U.\par}
\end{minipage}\par\medskip
\noindent\begin{minipage}{\linewidth}
\SGARestoreAppState{EN-VAR-136}%
\paragraph{[136] EN-VAR-136}
Printed reading: Thus $Y$ has codimension\par
Corrected reading: Thus $X$ has codimension\par
{\raggedright Source: fixes/RECONCILE\_CODEX\_6.json\#252; fixes/FIXLOG\_6\_CURRENT.jsonl\#359. The specified closed immersion f:X\ensuremath{\to}Y is cut out by I. At the chosen point I\_x is nonzero in the domain R=O\_(Y,x). Every prime P containing I\_x is nonzero, so the strict chain from(0) to P has positive length. This computes the positive codimension of the closed subspace X=V(I) in Y. Current restores X as the subject; diplomatic preserves Y.\par}
\end{minipage}\par\medskip
\noindent\begin{minipage}{\linewidth}
\SGARestoreAppState{EN-VAR-137}%
\paragraph{[137] EN-VAR-137}
Printed reading: (\uline{deduced from} (3.3 b) (ii)) ) \uline{is an isomorphism}\par
Corrected reading: (\uline{deduced from} (3.3 b) (ii)) \uline{is an isomorphism}\par
{\raggedright Source: ASTRA-II-PAREN-189. The complete parenthetical has three opening and four closing parentheses in the print. Current removes the final unmatched close; diplomatic retains it. The reference and adjunction arrow E\ensuremath{\to}RQ(E) are unchanged.\par}
\end{minipage}\par\medskip
\noindent\begin{minipage}{\linewidth}
\SGARestoreAppState{EN-VAR-138}%
\paragraph{[138] EN-VAR-138}
Printed reading: (\uline{deduced from} (3.3 b)(ii)) ) \uline{is an isomorphism}\par
Corrected reading: (\uline{deduced from} (3.3 b)(ii)) \uline{is an isomorphism}\par
{\raggedright Source: ASTRA-II-PAREN-189. The second complete parenthetical likewise has one unmatched final close. Current deletes that delimiter only, preserving the reference and adjunction arrow RQ(E)\ensuremath{\to}E; diplomatic keeps every printed delimiter.\par}
\end{minipage}\par\medskip
\noindent\begin{minipage}{\linewidth}
\SGARestoreAppState{EN-VAR-139}%
\paragraph{[139] EN-VAR-139}
Printed reading: \(M=\coprod_{n\ 0}B/(\underline{f}^n)\)\par
Corrected reading: \(M=\coprod_{n>0}B/(\underline{f}^n)\)\par
{\raggedright Source: check/findings/findings\_196\_200.json\#3; fixes/MISPRINTS\_6.json\#185. The printed relation slot between n and0 is blank. Current supplies the positive-index sum. Its canonical inclusion into the sum over n\ensuremath{\geq}0 and the projection back are inverse: the n=0 term is B/(1,1)=0 for f=(x,y). This identifies the two specified module constructions without changing any summand of positive index. Diplomatic retains the blank.\par}
\end{minipage}\par\medskip
\noindent\begin{minipage}{\linewidth}
\SGARestoreAppState{EN-VAR-140}%
\paragraph{[140] EN-VAR-140}
Printed reading: \(H_2^B(\underline{f}^n,A)\simeq H_2^B(\underline{f}^n,M),\)\par
Corrected reading: \(H_2^B(\underline{f}^n,A)\simeq H_2^B(\underline{f}^n,M).\)\par
{\raggedright Source: check/findings/findings\_196\_200.json\#7; fixes/MISPRINTS\_6.json\#186. The print ends this display with a comma before the next independent capitalized sentence. Current uses a sentence stop. Every Koszul term and the isomorphism sign remain identical; diplomatic preserves the comma.\par}
\end{minipage}\par\medskip
\noindent\begin{minipage}{\linewidth}
\SGARestoreAppState{EN-VAR-141}%
\paragraph{[141] EN-VAR-141}
Printed reading: Finally, we have (EGA II 1.1.3.3)\par
Corrected reading: Finally, we have (EGA III 1.1.3.3)\par
{\raggedright Source: check/findings/findings\_196\_200.json\#8; fixes/MISPRINTS\_6.json\#187. The reviewed primary EGA III text, Part1 page83 formula1.1.3.3, gives the Koszul duality identification H\textasciicircum{}i(f,M) with H\_(r-i)(f,M). For r=2 and i=0, the map sends the top exterior-basis vector tensored with m to m in the common annihilator of x\textasciicircum{}n and y\textasciicircum{}n. This is precisely the displayed H\_2 to H\textasciicircum{}0 map. Current repairs only the Roman chapter numeral; diplomatic keeps II. The independent primary-text verification is recorded by the French reviewer, without a new EGA image-read claim.\par}
\end{minipage}\par\medskip
\noindent\begin{minipage}{\linewidth}
\SGARestoreAppState{EN-VAR-142}%
\paragraph{[142] EN-VAR-142}
Printed reading: \(H_2^A(\underline{f}^n,A\simeq {}_{(\underline{f}^n)}M.\)\par
Corrected reading: \(H_2^A(\underline{f}^n,A)\simeq {}_{(\underline{f}^n)}M.\)\par
{\raggedright Source: check/findings/findings\_196\_200.json\#11; fixes/MISPRINTS\_6.json\#188; fixes/RECONCILE\_CODEX\_6.json\#254. Current closes the existing argument list of H\_2\textasciicircum{}A(f\textasciicircum{}n,A), the same Koszul module in the first display and the transition maps. No term or map is added. Diplomatic retains the omitted closing parenthesis in(*).\par}
\end{minipage}\par\medskip
\noindent\begin{minipage}{\linewidth}
\SGARestoreAppState{EN-VAR-143}%
\paragraph{[143] EN-VAR-143}
Printed reading: \(P^n_{X/S}(E)=P^n_{X/S}\oplus_{\underline{O}_X}E\)\par
Corrected reading: \(P^n_{X/S}(E)=P^n_{X/S}\otimes_{\underline{O}_X}E\)\par
{\raggedright Source: check/findings/findings\_201\_205.json\#4. The printed sign is a circled plus, while the immediately following words specify the tensor product using the second O\_X action. The balanced construction imposes (p a) tensor e = p tensor (a e), and its canonical map sends e to1 tensor e. Composition with an O\_X-linear map from this construction is the factorization in1.2.2.1. Current retains the recorded tensor-product repair; diplomatic preserves the actual plus.\par}
\end{minipage}\par\medskip
\noindent\begin{minipage}{\linewidth}
\SGARestoreAppState{EN-VAR-144}%
\paragraph{[144] EN-VAR-144}
Printed reading: is an $f^{-1}(\underline{O}_X)$-linear homomorphism\par
Corrected reading: is an $f^{-1}(\underline{O}_S)$-linear homomorphism\par
{\raggedright Source: check/findings/findings\_201\_205.json\#3. For the specified f:X\ensuremath{\to}S, the structure map is f-sharp:f\textasciicircum{}(-1)O\_S\ensuremath{\to}O\_X. The two O\_X actions on principal parts agree on that image because both projections are over S. Hence d\textasciicircum{}n(E)(a e)=u\_2(a)(1 tensor e)=u\_1(a)(1 tensor e) for a from f\textasciicircum{}(-1)O\_S, and an O\_X-linear factor h gives D(a e)=a D(e). This exact factorization identifies the source of the required base-linearity. Current uses O\_S; diplomatic preserves the printed O\_X in this later definition only.\par}
\end{minipage}\par\medskip
\noindent\begin{minipage}{\linewidth}
\SGARestoreAppState{EN-VAR-145}%
\paragraph{[145] EN-VAR-145}
Printed reading: the identifications (1.2.2.1) and (1.2.4.1)\par
Corrected reading: the identifications (1.2.2.1) and (1.2.3.1)\par
{\raggedright Source: check/findings/findings\_201\_205.json\#8. The explicit composition h\ensuremath{\to}h composed with d\textasciicircum{}n(E) gives the real identification1.2.2.1 and complex identification1.2.3.1. The paragraph applies Hom to the exact sequence1.2.4.1 and then uses those two identifications to obtain1.2.4.2. Current repairs the second identification locator to1.2.3.1; diplomatic keeps the printed1.2.4.1. Every map and sequence is retained.\par}
\end{minipage}\par\medskip
\noindent\begin{minipage}{\linewidth}
\SGARestoreAppState{EN-VAR-146}%
\paragraph{[146] EN-VAR-146}
Printed reading: (considered as a monomorphism of ringed spaces) cf.~(1.1)).\par
Corrected reading: (considered as a monomorphism of ringed spaces) cf.~(1.1).\par
{\raggedright Source: check/findings/findings\_201\_205.json\#12. The parenthesis opening the considered-clause is already closed after spaces, and the reference opening is closed after1.1. Current deletes only the remaining unmatched final close. Diplomatic retains it, with all clause words, the reference and the sentence stop unchanged.\par}
\end{minipage}\par\medskip
\noindent\begin{minipage}{\linewidth}
\SGARestoreAppState{EN-VAR-147}%
\paragraph{[147] EN-VAR-147}
Printed reading: to that of $\underline{O}_S$ rep.\par
Corrected reading: to that of $\underline{O}_S$ resp.\par
{\raggedright Source: check/findings/findings\_211\_215.json\#2. The printed shared abbreviation rep. lacks the s in resp. The real bundle is assigned its O\_S section sheaf; the complex bundle is assigned its O\_S tensor C section sheaf. Current restores the abbreviation linking these two stated assignments. Diplomatic preserves rep. at this occurrence; all functor data are unchanged.\par}
\end{minipage}\par\medskip
\noindent\begin{minipage}{\linewidth}
\SGARestoreAppState{EN-VAR-148}%
\paragraph{[148] EN-VAR-148}
Printed reading: results of (loc.cit):\par
Corrected reading: results of (loc.cit.):\par
{\raggedright Source: check/findings/findings\_211\_215.json\#9. Current supplies the missing final abbreviation stop in the shared Latin locator loc.cit. Diplomatic retains the printed loc.cit without that stop. The surrounding reference, two assertions and their separate arguments remain fully translated.\par}
\end{minipage}\par\medskip
\noindent\begin{minipage}{\linewidth}
\SGARestoreAppState{EN-VAR-149}%
\paragraph{[149] EN-VAR-149}
Printed reading: \(u_r:\underline C'(S,C^r_2(V,E))
  \longrightarrow \underline C'(S,C^r_2(V,F)).\)\par
Corrected reading: \(u_r:\underline C'(S,C^r_2(V,E))
  \longrightarrow \underline C'(S,C^{r-n}_2(V,F)).\)\par
{\raggedright Source: check/findings/findings\_211\_215.json\#10; fixes/MISPRINTS\_6.json\#184. The defining family is T(s)=phi\_r(u(s)) in L(H\_E,H\_F), with H\_E=C\textasciicircum{}r\_2(V,E) and H\_F=C\textasciicircum{}(r-n)\_2(V,F), as displayed immediately above. On a relatively compact open W, its induced presheaf map sends a in C(bar W,H\_E) to the continuous map s\ensuremath{\to}T(s)a(s) in C(bar W,H\_F). Continuity follows from continuous evaluation L(H\_E,H\_F) times H\_E\ensuremath{\to}H\_F. This explicit map gives target C-prime(S,C\textasciicircum{}(r-n)\_2(V,F)); the final bundle map also prints that exponent. Current retains r-n and diplomatic preserves the printed r only in this presheaf target.\par}
\end{minipage}\par\medskip
\noindent\begin{minipage}{\linewidth}
\SGARestoreAppState{EN-VAR-150}%
\paragraph{[150] EN-VAR-150}
Printed reading: the relations between the groups $[S,I(S)]$ and\par
Corrected reading: the relations between the groups $[S,I(H)]$ and\par
{\raggedright Source: check/findings/findings\_216\_220.json\#23. The paragraph defines I(H) inside L(H,H) and the homotopy group of maps S\ensuremath{\to}I(H). Its element is the class of a map with that target, and its homotopies are maps S times[0,1]\ensuremath{\to}I(H). The later comparison therefore concerns that same group [S,I(H)]. Current restores H in this single later occurrence; diplomatic preserves the printed S. No relation to K-bullet(S) is added to the historical question.\par}
\end{minipage}\par\medskip
\noindent\begin{minipage}{\linewidth}
\SGARestoreAppState{EN-VAR-151}%
\paragraph{[151] EN-VAR-151}
Printed reading: Fonctions analyliques sur un ensemble semi-analytique\par
Corrected reading: Fonctions analytiques sur un ensemble semi-analytique\par
{\raggedright Source: check/findings/findings\_216\_220.json\#29. The bibliography title is retained in French in the English edition. Diplomatic therefore preserves the actually printed analyliques, with l, and current retains the recorded spelling correction analytiques, with t. The later semi-analytique, author, publication data, page range and punctuation remain unchanged.\par}
\end{minipage}\par\medskip
\noindent\begin{minipage}{\linewidth}
\SGARestoreAppState{EN-VAR-152}%
\paragraph{[152] EN-VAR-152}
Printed reading: \(H^n(i_*(E'')\)\par
Corrected reading: \(H^n(i_*(E''))\)\par
{\raggedright Source: check/findings/findings\_221\_225.json\#19. The first expression applies H\textasciicircum{}n to the complete object i\_*(E-double-prime), but the printed form closes only the inner argument. The next expression applies the same composite H\textasciicircum{}n composed with i\_* and prints both closers. Current adds only the missing outside parenthesis; diplomatic retains the exact first missing close. The degree n remains unchanged.\par}
\end{minipage}\par\medskip
\noindent\begin{minipage}{\linewidth}
\SGARestoreAppState{EN-VAR-153}%
\paragraph{[153] EN-VAR-153}
Printed reading: independent of the hypothesis on \(\bar i_*(\underline{O}_X)\)\par
Corrected reading: independent of the hypothesis on \(i_*(\underline{O}_X)\)\par
{\raggedright Source: check/findings/findings\_221\_225.json\#21. The source defines the single closed immersion i:X\ensuremath{\to}Y. At x, the stalk of i\_*M at i(x) is M\_x, and the surjection O\_(Y,i(x))\ensuremath{\to}O\_(X,x) carries the displayed finite generators to the same generating elements over the quotient ring. This is exactly the direct image i\_* from the lemma hypothesis. Current keeps that same unbarred i; diplomatic retains the observed extra bar.\par}
\end{minipage}\par\medskip
\noindent\begin{minipage}{\linewidth}
\SGARestoreAppState{EN-VAR-154}%
\paragraph{[154] EN-VAR-154}
Printed reading: the Kozsul complex\par
Corrected reading: the Koszul complex\par
{\raggedright Source: check/findings/findings\_221\_225.json\#23; fixes/FIXLOG\_6\_CURRENT.jsonl\#372. The displayed regular-immersion construction uses the complex with exterior-power terms and differential sending the basis element epsilon\_j to the corresponding ideal generator f\_j, the Koszul resolution referenced in VII1.4. Diplomatic retains the transposed proper-name spelling Kozsul; current retains Koszul.\par}
\end{minipage}\par\medskip
\noindent\begin{minipage}{\linewidth}
\SGARestoreAppState{EN-VAR-155}%
\paragraph{[155] EN-VAR-155}
Printed reading: \(i''_*\ n\text{-\uline{pseudo-coherent}}\)\par
Corrected reading: \(i''_*(E)\ n\text{-\uline{pseudo-coherent}}\)\par
{\raggedright Source: check/findings/findings\_221\_225.json\#28; fixes/FIXLOG\_6\_CURRENT.jsonl\#373. The lemma fixes E in D(X) and applies the two closed-immersion direct-image functors to that object. Its proof uses the diagonal X\ensuremath{\to}X-prime times\_Y X-double-prime and compares those direct images; the right-hand object is therefore i-double-prime\_*(E). Current restores only the omitted argument(E); diplomatic keeps the printed bare functor symbol.\par}
\end{minipage}\par\medskip
\noindent\begin{minipage}{\linewidth}
\SGARestoreAppState{EN-VAR-156}%
\paragraph{[156] EN-VAR-156}
Printed reading: \begin{equation}\tag{$a'$}
\begin{tikzcd}
X' \arrow[r,"i''"] \arrow[d,"p'"'] & Y' \arrow[ddl,"f'"]\\
X \arrow[d,"f"']\\
Y
\end{tikzcd}
\end{equation}\par
Corrected reading: \begin{equation}\tag{$a'$}
\begin{tikzcd}
X' \arrow[r,"i'"] \arrow[d,"p'"'] & Y' \arrow[ddl,"f'"]\\
X \arrow[d,"f"']\\
Y
\end{tikzcd}
\end{equation}\par
{\raggedright Source: check/findings/findings\_226\_230.json\#3; fixes/FIXLOG\_6\_CURRENT.jsonl\#374. The top edge has endpoints X-prime\ensuremath{\to}Y-prime. The following text names that immersion i-prime and uses its direct-image functor i-prime\_*p-prime*. The next-page cube identifies the same endpoints and edge as i-prime. Thus the endpoint-preserving map from the first diagram to the corresponding cube face identifies its label with i-prime. Current keeps that single prime; diplomatic preserves the printed double prime. The distinct second factorization retains i-double-prime.\par}
\end{minipage}\par\medskip
\noindent\begin{minipage}{\linewidth}
\SGARestoreAppState{EN-VAR-157}%
\paragraph{[157] EN-VAR-157}
Printed reading: \(\underline{\Tor}_i^S(\underline{O}_X,O_Y)\)\par
Corrected reading: \(\underline{\Tor}_i^S(\underline{O}_X,\underline{O}_Y)\)\par
{\raggedright Source: check/findings/findings\_226\_230.json\#14. The source names X,Y,S with affine rings B,C,A and the tensor comparison B tensor-derived\_A C\ensuremath{\to}B tensor\_A C. The two sheaf arguments restrict to B and C by the structure maps A\ensuremath{\to}B and A\ensuremath{\to}C. Thus the Y argument is the same structure sheaf of Y used in that calculation. Diplomatic retains its isolated missing rule; current retains the established rule on that same O\_Y without changing the sheaf, indices or tensor map.\par}
\end{minipage}\par\medskip
\noindent\begin{minipage}{\linewidth}
\SGARestoreAppState{EN-VAR-158}%
\paragraph{[158] EN-VAR-158}
Printed reading: \[
\begin{tikzcd}
X \arrow[r,"f"] \arrow[dr,"f"'] & Y \arrow[d,"g"]\\
& Z
\end{tikzcd}
\]\par
Corrected reading: \[
\begin{tikzcd}
X \arrow[r,"f"] \arrow[dr,"h"'] & Y \arrow[d,"g"]\\
& Z
\end{tikzcd}
\]\par
{\raggedright Source: check/findings/findings\_231\_235.json\#18; fixes/FIXLOG\_6\_CURRENT.jsonl\#375. The diagonal has endpoints X\ensuremath{\to}Z and is the composite h=g composed with f. Its direct-image functor is h\_*=g\_*f\_*, exactly the functor used in the subsequent relative condition and proof. Current retains h on this edge; diplomatic retains the printed f. The horizontal f:X\ensuremath{\to}Y and vertical g:Y\ensuremath{\to}Z are unchanged.\par}
\end{minipage}\par\medskip
\noindent\begin{minipage}{\linewidth}
\SGARestoreAppState{EN-VAR-159}%
\paragraph{[159] EN-VAR-159}
Printed reading: the isomorphism (*) in the proof of (1.8)\par
Corrected reading: the isomorphism in (1.7)\par
{\raggedright Source: check/findings/findings\_231\_235.json\#28. Both full scans print l’isomorphisme(*) de la preuve de(1.8). The inspected231–232 text has no starred formula in either proof. The displayed isomorphism is exactly the Lemma1.7 map at S=Y, f\_1=f:X\ensuremath{\to}Y and f\_2=Id\_Y, using X\ensuremath{\times}\_Y Y\ensuremath{\cong}X. The proof has already reduced f to a closed immersion. At y=f(x), put R=O\_\{Y,y\}, B=O\_\{X,x\}; for a flat resolution P of F\_y the chain maps e\ensuremath{\otimes}p\ensuremath{\mapsto}e\ensuremath{\otimes}(1\ensuremath{\otimes}p) and e\ensuremath{\otimes}(b\ensuremath{\otimes}p)\ensuremath{\mapsto}be\ensuremath{\otimes}p are inverse between E\_x\ensuremath{\otimes}\_R P and E\_x\ensuremath{\otimes}\_B(B\ensuremath{\otimes}\_R P). They realize the stalk of the canonical adjunction projection map. Off X both stalks are zero, since f is closed. This proves the displayed map in the actual proof scope, including complexes not assumed quasi-coherent. Current cites l’isomorphisme de(1.7), removing the nonexistent star and proof locator; diplomatic retains the complete printed phrase. This is an explicit editorial repair, with no claim about an unseen author draft.\par}
\end{minipage}\par\medskip
\noindent\begin{minipage}{\linewidth}
\SGARestoreAppState{EN-VAR-160}%
\paragraph{[160] EN-VAR-160}
Printed reading: \(f_*(E)\)\par
Corrected reading: \(h_*(E)\)\par
{\raggedright Source: check/findings/findings\_231\_235.json\#27; fixes/FIXLOG\_6\_CURRENT.jsonl\#376. Proposition1.11(b) assumes pseudo-coherence relative to h. After the stated reduction to closed immersions, its hypothesis is pseudo-coherence of h\_*(E)=g\_*f\_*(E). With E\_0=f\_*(E), the following replacement X=Y and f=Id makes this precisely g\_*(E\_0), as used later in the proof. The conclusion remains pseudo-coherence of E\_0. Only this first hypothesis in(b) changes from printed f\_* to current h\_*; all other f\_* occurrences are preserved.\par}
\end{minipage}\par\medskip
\noindent\begin{minipage}{\linewidth}
\SGARestoreAppState{EN-VAR-161}%
\paragraph{[161] EN-VAR-161}
Printed reading: \(i>a+r\)\par
Corrected reading: \(q>a+r\)\par
{\raggedright Source: check/findings/findings\_236\_240.json\#13; fixes/FIXLOG\_6\_CURRENT.jsonl\#377. The higher direct image is R\textasciicircum{}qf\_*(E). In the stated hypercohomology computation, q=p+j is the total degree, with j\ensuremath{\leq}a from the complex and p\ensuremath{\leq}r from the projective-space cover, so q\ensuremath{\leq}a+r for potentially nonzero terms. The stated vanishing condition is therefore q>a+r for that same R\textasciicircum{}q. Current changes only the condition index to q; diplomatic retains the printed i while preserving R\textasciicircum{}q.\par}
\end{minipage}\par\medskip
\noindent\begin{minipage}{\linewidth}
\SGARestoreAppState{EN-VAR-162}%
\paragraph{[162] EN-VAR-162}
Printed reading: (2.2), (2.5) and (1.12)\par
Corrected reading: (2.2), (2.4) and (1.12)\par
{\raggedright Source: fixes/FIXLOG\_6\_CURRENT.jsonl\#379. The printed proof of2.5 cites2.5 itself. Corollary2.4 supplies the bounded derived functor under precisely the proper/projective or noetherian alternatives used here. Take the base S=Y locally on Y: the structure map of Y is id\_Y, and f is pseudo-coherent. By1.12 the absolute pseudo-coherence of E is its relative pseudo-coherence over Y; 2.2 transports that property through Rf\_*, and2.4 gives the bounded restriction when Y is quasi-compact. On Y, relative pseudo-coherence over id\_Y is the absolute property by the identity inverse-image map. Thus the middle citation is2.4, while2.2 and1.12 remain unchanged. Diplomatic retains the printed2.5.\par}
\end{minipage}\par\medskip
\noindent\begin{minipage}{\linewidth}
\SGARestoreAppState{EN-VAR-163}%
\paragraph{[163] EN-VAR-163}
Printed reading: (SGA 4 XVII IV 6.4)\par
Corrected reading: (SGA 4 IV 6.4)\par
{\raggedright Source: fixes/FIXLOG\_6\_CURRENT.jsonl\#380. Both SGA6 scans print SGA4 XVII IV6.4. The personally viewed SGA4 tome1 re-edition page230, PDF236, has the ExposéIV header and Definition6.4.1: a conservative family of point fibre functors defines having enough points. This is the exact definition cited in3.3. For a module complex, the fibre functor sends H\textasciicircum{}i(E) to H\textasciicircum{}i(E\_x), because the inverse-image fibre functor is exact; conservativity tests the isomorphism0\ensuremath{\to}H\textasciicircum{}i(E) on these fibres. This identifies the cited criterion without changing any hypothesis or mathematical symbol. Retain IV6.4 in current and the full XVII IV6.4 diplomatically; source record external\_sga4\_242.json records the checked reference.\par}
\end{minipage}\par\medskip
\noindent\begin{minipage}{\linewidth}
\SGARestoreAppState{EN-VAR-164}%
\paragraph{[164] EN-VAR-164}
Printed reading: Suppose \(E\) (resp. \(M\)) is of finite flat amplitude rel. to \(f\) (resp. finite flat amplitude).\par
Corrected reading: Suppose \(E\) (resp. \(M\)) is of finite flat amplitude rel. to \(S\) (resp. finite flat amplitude).\par
{\raggedright Source: check/findings/findings\_246\_250.json\#3. The primary and native pages print rel. à f in the proof, while the immediately preceding statement247 explicitly assumes finite amplitude relative to S. Write g:Y\ensuremath{\to}S and h=g\ensuremath{\circ}f. The tensor used in the proof is E tensor\textasciicircum{}L\_\{h\textasciicircum{}\{-1\}O\_S\}h\textasciicircum{}\{-1\}M. The structural map is the composite h\textasciicircum{}\{-1\}O\_S=f\textasciicircum{}\{-1\}g\textasciicircum{}\{-1\}O\_S\ensuremath{\to}f\textasciicircum{}\{-1\}O\_Y\ensuremath{\to}O\_X; relative amplitude for f instead uses f\textasciicircum{}\{-1\}O\_Y. These are the exact scalar maps involved. The difference is witnessed by A=k[\ensuremath{\varepsilon}]/(\ensuremath{\varepsilon}\textsuperscript{2}), S=SpecA, X=Y=Speck, f=id, E=k and M=k: E is free over O\_Y, but the free A-resolution ...\ensuremath{\to}A --\ensuremath{\varepsilon}\ensuremath{\to}A --\ensuremath{\varepsilon}\ensuremath{\to}A\ensuremath{\to}k has kernel=image=(\ensuremath{\varepsilon}) and pulls back to zero differentials on k in every negative degree. Thus relative amplitude for f alone does not provide the asserted bounded tensor. The statement’s S hypothesis supplies the bounded flat resolution over h\textasciicircum{}\{-1\}O\_S used here; its tensor with a bounded representative of M has total degrees in the sums of their finite degree intervals. Correct only this recalled base letter f\ensuremath{\to}S in current. The statement itself, the adjunction counit, both canonical arrows and the alternative finite-amplitude hypothesis on M remain unchanged.\par}
\end{minipage}\par\medskip
\noindent\begin{minipage}{\linewidth}
\SGARestoreAppState{EN-VAR-165}%
\paragraph{[165] EN-VAR-165}
Printed reading: is a perfect morphism\par
Corrected reading: is a perfect morphism.\par
{\raggedright Source: check/findings/findings\_251\_255.json\#2. Both complete scans leave the completed sentence without a final stop, before the separate Notation4.2 paragraph. Retain that absence diplomatically and the current terminal stop.\par}
\end{minipage}\par\medskip
\noindent\begin{minipage}{\linewidth}
\SGARestoreAppState{EN-VAR-166}%
\paragraph{[166] EN-VAR-166}
Printed reading: \(g_\alpha^*(F)\simeq F_\alpha\)\par
Corrected reading: \(g_\alpha^*(F_\alpha)\simeq F_\alpha\)\par
{\raggedright Source: check/findings/findings\_251\_255.json\#6. Both scans print g\_alpha*(F) in the second comparison. The preceding definition binds F\_alpha on X\_alpha and the first comparison identifies g\_alpha*F\_alpha with F on X. On the affine charts Y\_alpha=SpecA\_alpha, Y=SpecA, write B\_alpha=A\_alpha[t1,…,tr] and B=A[t1,…,tr]=B\_alpha tensor\_Aalpha A. For the module N representing F\_alpha, the pullback is B tensor\_Balpha N. The explicit map (b tensor a) tensor n\ensuremath{\mapsto}a tensor bn identifies it with A tensor\_Aalpha N, with inverse a tensor n\ensuremath{\mapsto}(1 tensor a) tensor n. The given A\_alpha-flatness of N identifies the ordinary tensor with its derived tensor. Thus both sides of the printed comparison identify g\_alpha*F\_alpha, already identified with F, and the corrected argument is F\_alpha. Keep the printed F diplomatically and retain the headless comparison sign in both layers.\par}
\end{minipage}\par\medskip
\noindent\begin{minipage}{\linewidth}
\SGARestoreAppState{EN-VAR-167}%
\paragraph{[167] EN-VAR-167}
Printed reading: \(\parfamp(i_*(E)\subset[a,b]\Longrightarrow\parfamp_f(E)\subset[a,b]\)\par
Corrected reading: \(\parfamp(i_*(E))\subset[a,b]\Longrightarrow\parfamp_f(E)\subset[a,b]\)\par
{\raggedright Source: check/findings/findings\_251\_255.json\#8. Both scans have two opening parentheses in parfamp(i\_*(E) but only the inner closing parenthesis. The outer argument is i\_*(E), exactly as the second implication on the same page writes it. Retain the one-close print diplomatically and the two-close expression in current.\par}
\end{minipage}\par\medskip
\noindent\begin{minipage}{\linewidth}
\SGARestoreAppState{EN-VAR-168}%
\paragraph{[168] EN-VAR-168}
Printed reading: It also implies. by 4.3),\par
Corrected reading: It also implies (by 4.3),\par
{\raggedright Source: check/findings/findings\_251\_255.json\#10. The primary and native254 pages show a stop after aussi, a lower-case d in d’après, and a closing parenthesis after4.3 with no opening parenthesis. The parenthetical citation concerns Proposition4.3’s displayed stalk-amplitude condition. Replace the stray stop by the missing opening parenthesis in current, preserving the existing close and comma; diplomatic retains the exact printed marks.\par}
\end{minipage}\par\medskip
\noindent\begin{minipage}{\linewidth}
\SGARestoreAppState{EN-VAR-169}%
\paragraph{[169] EN-VAR-169}
Printed reading: is, in a neighborhood of \(s\), acyclic outside\par
Corrected reading: is, in a neighborhood of \(x\), acyclic outside\par
{\raggedright Source: check/findings/findings\_251\_255.json\#15. Both scans print s. The already defined map h sends x\ensuremath{\in}X to s\ensuremath{\in}S; E and the truncation morphism t are complexes on X. Condition4.6(i) is stated on an open neighbourhood U of x, and the reduction explicitly restricts along its open immersion j:U\ensuremath{\to}X. The restriction j*E\ensuremath{\to}\ensuremath{\tau}\_[a,b](j*E) is the truncation used in the proof: its stalks on U are the same truncations of E, and its being a quasi-isomorphism follows from the vanishing specified on U. This identifies the recalled neighbourhood as the neighbourhood of x from(i). Retain current x and diplomatic s; neither h nor any fibre point is renamed.\par}
\end{minipage}\par\medskip
\noindent\begin{minipage}{\linewidth}
\SGARestoreAppState{EN-VAR-170}%
\paragraph{[170] EN-VAR-170}
Printed reading: Thus, after localizing,. we may suppose\par
Corrected reading: Thus, after localizing, we may suppose\par
{\raggedright Source: check/findings/findings\_251\_255.json\#18. Both scans show the comma followed by a second typed stop. The continuation on peut supposer belongs to the same sentence. Retain the pair diplomatically and the current single comma. The partial overstruck fragments above this line are not transcribed as a new expression.\par}
\end{minipage}\par\medskip
\noindent\begin{minipage}{\linewidth}
\SGARestoreAppState{EN-VAR-171}%
\paragraph{[171] EN-VAR-171}
Printed reading: arrow \((*)\)\par
Corrected reading: arrow \((t)\)\par
{\raggedright Source: check/findings/findings\_251\_255.json\#17. The fibre comparison is visibly unlabelled in both scans; later the print says la flèche(*). The map used in that final paragraph is the original truncation t:E\ensuremath{\to}T on X, where T\textasciicircum{}a=E\textasciicircum{}a/B\textasciicircum{}a, T\textasciicircum{}j=E\textasciicircum{}j for a<j\ensuremath{\leq}b, and T is zero elsewhere. Its component in degree a is the quotient map, its higher components are identities, and its lower components are zero. The preceding assertion makes t a quasi-isomorphism. Tensoring t with an arbitrary h\textasciicircum{}\{-1\}O\_S-module therefore identifies the cohomology used to test toramp\_h(E) with that of T. The terms T\textasciicircum{}j for j>a are h-flat because they are locally free over O\_X and h is flat. LemmaI5.1.1, with m=a,n=b,F=T and base h\textasciicircum{}\{-1\}O\_S, then gives precisely the h-flatness of T\textasciicircum{}a claimed at the start of256. The intervening fibre comparison has objects over the residue fibre and supplies no such map E\ensuremath{\to}T on X. Therefore repair the later reference(*) to(t), and remove the previously inserted(*) label from the headless fibre comparison. The printed t) and i*(t) labels themselves remain unchanged.\par}
\end{minipage}\par\medskip
\noindent\begin{minipage}{\linewidth}
\SGARestoreAppState{EN-VAR-172}%
\paragraph{[172] EN-VAR-172}
Printed reading: \uline{Under the hypotheses of} (4.6)\par
Corrected reading: \uline{Under the hypotheses of} (4.7)\par
{\raggedright Source: check/findings/findings\_256\_260.json\#3. The print cites4.6. Proposition4.7(b) is the statement whose data are the two f\_i and their product f=f1\ensuremath{\times}\_Sf2. Substitute E\_i=O\_Xi and[a\_i,b\_i]=[\ensuremath{-}n\_i,0]. Their external tensor identifies with O\_X by the multiplication map O\_X tensor\_OX O\_X\ensuremath{\to}O\_X, u tensor v\ensuremath{\mapsto}uv, whose inverse sends w to w tensor1. The displayed summed interval is[\ensuremath{-}n1\ensuremath{-}n2,0], exactly the dimension bound stated in4.7.1. The Tor-independence assumptions are those of4.7. This identifies the exact cited assertion(b), independently of the separately recorded failure of4.7(a). Correct only the locator4.6\ensuremath{\to}4.7 in current.\par}
\end{minipage}\par\medskip
\noindent\begin{minipage}{\linewidth}
\SGARestoreAppState{EN-VAR-173}%
\paragraph{[173] EN-VAR-173}
Printed reading: \uline{Under the hypotheses of} (4.8)\par
Corrected reading: \uline{Under the hypotheses of} (4.9)\par
{\raggedright Source: check/findings/findings\_256\_260.json\#14. The print cites4.8, while the author’s next sentence specifies the substitution E=O\_X. In4.9, the hypothesis that f is perfect makes O\_X perfect relative to f. The complex RHom\_OX(O\_X,f!F) identifies with f!F by evaluation at1, with inverse m\ensuremath{\mapsto}(a\ensuremath{\mapsto}am); O\_X is locally free of rank1, so this is also the derived identification. The resulting conclusion is exactly that f!F is perfect over S. This uses4.9’s compactifiable/noetherian setting and its actual Hom expression. Correct only the locator to4.9.\par}
\end{minipage}\par\medskip
\noindent\begin{minipage}{\linewidth}
\SGARestoreAppState{EN-VAR-174}%
\paragraph{[174] EN-VAR-174}
Printed reading: For the first assertion, apply (4.8)\par
Corrected reading: For the first assertion, apply (4.9)\par
{\raggedright Source: check/findings/findings\_256\_260.json\#20. The print says appliquer4.8 avecY=S,F=O\_S. Under these explicit substitutions in4.9, f!F=f!O\_S=K\_X and RHom(E,f!F)=D\_X(E) by the preceding definition. O\_S is the rank-one free complex over id\_S, and E has the relative perfection required in4.9. The structural morphism X\ensuremath{\to}S is f itself, so the resulting S-relative perfection is exactly membership in D(f)\_parf. This identifies the first assertion of4.9.2 term for term and fixes the locator4.9 without changing any functor or object.\par}
\end{minipage}\par\medskip
\noindent\begin{minipage}{\linewidth}
\SGARestoreAppState{EN-VAR-175}%
\paragraph{[175] EN-VAR-175}
Printed reading: strictly perfect (2.1)\par
Corrected reading: strictly perfect (I 2.1)\par
{\raggedright Source: ASTRA-III-REFERENCE-261. Both full261 scans print strictly perfect followed by(2.1). Definition I2.1 on95 was re-read directly: its third alternative says that every E\textasciicircum{}i lies in C0 and E\textasciicircum{}i=0 for almost all i. In the ringed-topos setting here, C0 consists locally of finite-rank free O-modules, so this is exactly the bounded finite-free complex used for the subsequent dévissage. The identification sends each complex term and each differential to itself; after a chosen local basis E\textasciicircum{}i\ensuremath{\cong}O\textasciicircum{}\{r\_i\}, the basis map sends (a\_j) to sum a\_j e\_j with inverse the coordinate maps. III2.1 on236 is the proper-direct-image finiteness conjecture and supplies no such definition. Restore the missing exposé locator I in current, leaving the mathematical assertion and the printed diplomatic locator intact.\par}
\end{minipage}\par\medskip
\noindent\begin{minipage}{\linewidth}
\SGARestoreAppState{EN-VAR-176}%
\paragraph{[176] EN-VAR-176}
Printed reading: condition analogous to a) (resp. a'), etc). in which\par
Corrected reading: condition analogous to a) (resp. a'), etc.) in which\par
{\raggedright Source: check/findings/findings\_261\_265.json\#8. Both scans print etc). où in the middle of the sentence. The preceding parenthesis on the same page uses etc.) la condition. The abbreviation stop belongs to etc., followed by the closing parenthesis; the relative clause où continues that sentence. Move only the stop in current, preserving the printed order in diplomatic.\par}
\end{minipage}\par\medskip
\noindent\begin{minipage}{\linewidth}
\SGARestoreAppState{EN-VAR-177}%
\paragraph{[177] EN-VAR-177}
Printed reading: (resp. \(\operatorname{rg}_{k(x)}H^1(E_{k(x)})\) is upper semi-continuous\par
Corrected reading: (resp. \(\operatorname{rg}_{k(x)}H^1(E_{k(x)})\)) is upper semi-continuous\par
{\raggedright Source: check/findings/findings\_266\_270.json\#5. Both scans open a prose parenthesis before resp. and leave it open before est semi-continue. The mathematical H\textasciicircum{}1 argument already has its own close. Current supplies only the prose close; diplomatic retains the omission.\par}
\end{minipage}\par\medskip
\noindent\begin{minipage}{\linewidth}
\SGARestoreAppState{EN-VAR-178}%
\paragraph{[178] EN-VAR-178}
Printed reading: \(\operatorname{rg}(E_x)=\operatorname{rg}(E_{k(X)})\)\par
Corrected reading: \(\operatorname{rg}(E_x)=\operatorname{rg}(E_{k(x)})\)\par
{\raggedright Source: check/findings/findings\_266\_270.json\#6. The preceding sentence explicitly defines E\_k(x)=E\_x tensor\textasciicircum{}L\_OXx k(x), with x the varying point. The rank equality uses precisely that object and the canonical scalar extension O\_Xx\ensuremath{\to}k(x). On a finite free representative, its terms O\_Xx\textasciicircum{}\{r\_i\} tensor k(x)\ensuremath{\to}k(x)\textasciicircum{}\{r\_i\} send (a\_j)tensor b to (a\_j b), with inverse coordinates e\_jtensor b\_j. Both ranks are sum(-1)\textasciicircum{}i r\_i. Retain current x in this exact identified fibre notation, preserving capital X in diplomatic.\par}
\end{minipage}\par\medskip
\noindent\begin{minipage}{\linewidth}
\SGARestoreAppState{EN-VAR-179}%
\paragraph{[179] EN-VAR-179}
Printed reading: \(Z'^p(E)\)) are flat\par
Corrected reading: \(Z'^p(E)\) are flat\par
{\raggedright Source: check/findings/findings\_266\_270.json\#7. Both scans show two closes after the single argument E, with no surrounding prose parenthesis opened. Current deletes the surplus close; diplomatic keeps it.\par}
\end{minipage}\par\medskip
\noindent\begin{minipage}{\linewidth}
\SGARestoreAppState{EN-VAR-180}%
\paragraph{[180] EN-VAR-180}
Printed reading: \(H^1(E_{k(x)}\simeq H^1(E)_{k(x)}\)\par
Corrected reading: \(H^1(E_{k(x)})\simeq H^1(E)_{k(x)}\)\par
{\raggedright Source: check/findings/findings\_266\_270.json\#9. The note has a close for k(x), but none for H\textasciicircum{}1(. The headless comparison separates the two cohomology terms. Keep the printed omission in diplomatic and the unique outer close in current.\par}
\end{minipage}\par\medskip
\noindent\begin{minipage}{\linewidth}
\SGARestoreAppState{EN-VAR-181}%
\paragraph{[181] EN-VAR-181}
Printed reading: \(y\mapsto\operatorname{rg}_{k(x)}H^p(E\otimes_X^{\mathrm L}k(y))\)\par
Corrected reading: \(y\mapsto\operatorname{rg}_{k(y)}H^p(E\otimes_X^{\mathrm L}k(y))\)\par
{\raggedright Source: check/findings/findings\_266\_270.json\#12. The print on271 reads rank over k(x), with x the generic point fixed on270 and y the varying point. The actual fibre E tensor\_X\textasciicircum{}L k(y) is a k(y)-complex and is exactly the fibre in5.6(i) with its bound variable renamed x\ensuremath{\to}y. This relabelling sends the scalar extension O\_X,y\ensuremath{\to}k(y) and every term and differential to themselves. The need for the y field is witnessed by X=Spec(k[t]), E=O\_X and cohomological degree0: at y=(t), the fibre is k with t acting by0. An action of the generic field k(t) extending this k[t] action would make t invertible and zero on a nonzero module, which is impossible. Current k(y) therefore restores the stated fibre-rank function. Keep printed k(x) in diplomatic; assign the record to actual271.\par}
\end{minipage}\par\medskip
\noindent\begin{minipage}{\linewidth}
\SGARestoreAppState{EN-VAR-182}%
\paragraph{[182] EN-VAR-182}
Printed reading: \(H^i(Rf_*(E)\otimes_Y^{\mathrm L}k(y)\) \uline{are vector spaces of finite rank over}\par
Corrected reading: \(H^i(Rf_*(E)\otimes_Y^{\mathrm L}k(y))\) \uline{are vector spaces of finite rank over}\par
{\raggedright Source: check/findings/findings\_271\_275.json\#4. The source closes k(y), but leaves H\textasciicircum{}i( open. Current closes the full cohomology argument with one additional parenthesis; diplomatic retains the exact printed omission.\par}
\end{minipage}\par\medskip
\noindent\begin{minipage}{\linewidth}
\SGARestoreAppState{EN-VAR-183}%
\paragraph{[183] EN-VAR-183}
Printed reading: \(y\longmapsto\sum(-1)^i\operatorname{rg}_{k(y)}H^i(Rf_*(E)\otimes_Y^{\mathrm L}k(y)\)\par
Corrected reading: \(y\longmapsto\sum(-1)^i\operatorname{rg}_{k(y)}H^i(Rf_*(E)\otimes_Y^{\mathrm L}k(y))\)\par
{\raggedright Source: check/findings/findings\_271\_275.json\#5. Both complete scans repeat the same omission in the alternating-rank display. Current supplies the one matching close for H\textasciicircum{}i(.\par}
\end{minipage}\par\medskip
\noindent\begin{minipage}{\linewidth}
\SGARestoreAppState{EN-VAR-184}%
\paragraph{[184] EN-VAR-184}
Printed reading: \(y\longmapsto\operatorname{rg}_{k(y)}H^p(Rf_*(E)\otimes_Y^{\mathrm L}k(y)).\)\par
Corrected reading: \(y\longmapsto\operatorname{rg}_{k(y)}H^p(Rf_*(E)\otimes_Y^{\mathrm L}k(y))\)\par
{\raggedright Source: check/findings/findings\_271\_275.json\#6. The source continues the sentence La fonction…est constructible after this stop. Preserve the printed stop in diplomatic and its removal in current. No mathematical symbol changes.\par}
\end{minipage}\par\medskip
\noindent\begin{minipage}{\linewidth}
\SGARestoreAppState{EN-VAR-185}%
\paragraph{[185] EN-VAR-185}
Printed reading: (cf. (EGA III 6 and (IV 3.1.0)\par
Corrected reading: (cf. (EGA III 6) and (IV 3.1.0))\par
{\raggedright Source: check/findings/findings\_271\_275.json\#10. Both scans show three opening parentheses and only the close after IV3.1.0. The clause contains two cited items joined by et, EGA III6 and IV3.1.0; the finite verb que pour f plat continues On peut montrer outside the citation. Close the first individual citation after III6 and close the encompassing cf. parenthesis after the already closed IV3.1.0 citation. This preserves each original opening and all citation tokens in their original order. The token map on EGA,III,6,IV,3.1.0 is the identity; no locator, functor or asserted base-change map is replaced. Current inserts these two closing delimiters, diplomatic retains the print.\par}
\end{minipage}\par\medskip
\noindent\begin{minipage}{\linewidth}
\SGARestoreAppState{EN-VAR-186}%
\paragraph{[186] EN-VAR-186}
Printed reading: \(Rf_*(E)\otimes^{\mathrm L}M)\longrightarrow Rf_*(E\otimes_Y^{\mathrm L}M)\)\par
Corrected reading: \(Rf_*(E)\otimes^{\mathrm L}M\longrightarrow Rf_*(E\otimes_Y^{\mathrm L}M)\)\par
{\raggedright Source: check/findings/findings\_271\_275.json\#12. The projection comparison opens and closes Rf\_*(E) before the tensor. The extra close after M has no opening mate. Current removes only that surplus close; diplomatic keeps it.\par}
\end{minipage}\par\medskip
\noindent\begin{minipage}{\linewidth}
\SGARestoreAppState{EN-VAR-187}%
\paragraph{[187] EN-VAR-187}
Printed reading: \(K^{\bullet}({}_{A'} X'_A)_{A'}=k(D({}_{A'} X'_A)_{A'\text{-}\parf}.\)\par
Corrected reading: \(K^{\bullet}({}_{A'} X'_A)_{A'}=k(D({}_{A'} X'_A)_{A'\text{-}\parf}).\)\par
{\raggedright Source: check/findings/findings\_281\_285.json\#8. The native lower crop shows the underlined parf followed by blank space and a full stop, with no final parenthesis. The corrected delimiter closes the already opened k( applied to the named category; it changes no object or map. Preserve the omission in diplomatic and the supplied closing parenthesis in current.\par}
\end{minipage}\par\medskip
\noindent\begin{minipage}{\linewidth}
\SGARestoreAppState{EN-VAR-188}%
\paragraph{[188] EN-VAR-188}
Printed reading: in (2.6.2) (resp. (2.6.5))\par
Corrected reading: in (2.6.5) (resp. (2.6.2))\par
{\raggedright Source: check/findings/findings\_281\_285.json\#11. The displayed (2.7.1) maps derived K-groups by Lf*. For a flat resolution F of E, A\ensuremath{\prime} tensor\_\{f\textasciicircum{}\{-1\}A\} f\textasciicircum{}\{-1\}F represents both A\ensuremath{\prime} tensor\_A\textasciicircum{}L E in (2.6.5) and Lf*E; identity on that complex gives the required identification. For a finite locally free/direct-summand object E, the underived tensor A\ensuremath{\prime} tensor\_\{f\textasciicircum{}\{-1\}A\} f\textasciicircum{}\{-1\}E is f*E, giving the naive product (2.6.2) and (2.7.2). The inclusion comparison sends each such E to its degree-zero complex, which is exactly the vertical map in the printed diagram. Thus current correctly swaps only the two reference locators; diplomatic retains their printed order.\par}
\end{minipage}\par\medskip
\noindent\begin{minipage}{\linewidth}
\SGARestoreAppState{EN-VAR-189}%
\paragraph{[189] EN-VAR-189}
Printed reading: \(K^{\bullet}(X)_{\naif}\to K^{\bullet}(X)\) the structure\par
Corrected reading: \(K^{\bullet}(X)_{\naif}\to K^{\bullet}(X)\)) the structure\par
{\raggedright Source: check/findings/findings\_281\_285.json\#22. The native witness has no closer for the parenthesis opened before et par suite. The existing current closer after the canonical map closes only that parenthetical aside; the two algebra structures and comparison map are unchanged. Diplomatic correctly leaves the parenthesis open.\par}
\end{minipage}\par\medskip
\noindent\begin{minipage}{\linewidth}
\SGARestoreAppState{EN-VAR-190}%
\paragraph{[190] EN-VAR-190}
Printed reading: (V 3.9.1 and VI\qquad)\par
Corrected reading: (V 3.9.1 and VI)\par
{\raggedright Source: check/findings/findings\_281\_285.json\#26. A visible untyped interval follows VI before the closer. Its content is blank; no missing numeral can be read or inferred from it. Retain the attached diplomatic spacing and current ordinary spacing as an explicit formatting variant with identical reference tokens V 3.9.1 and VI. Do not supply a conjectural reference.\par}
\end{minipage}\par\medskip
\noindent\begin{minipage}{\linewidth}
\SGARestoreAppState{EN-VAR-191}%
\paragraph{[191] EN-VAR-191}
Printed reading: Deligne [5]\par
Corrected reading: Deligne [1]\par
{\raggedright Source: check/findings/findings\_281\_285.json\#27. Both scans print Deligne [5]. The directly inspected bibliography on 296 identifies [1] as Deligne, P., papiers secrets, followed by [2] Grauert, [3] Grothendieck and [V] Verdier. The reference target is therefore the named Deligne entry [1]. Current changes only the locator, while diplomatic preserves [5]; no new attribution or theorem wording is supplied. The 296 consultation is regional and does not mark that page fully checked.\par}
\end{minipage}\par\medskip
\noindent\begin{minipage}{\linewidth}
\SGARestoreAppState{EN-VAR-192}%
\paragraph{[192] EN-VAR-192}
Printed reading: Thanks,to the associativity\par
Corrected reading: Thanks to the associativity\par
{\raggedright Source: check/findings/findings\_286\_290.json\#8. A comma with the same descending tail as the surrounding typed commas stands after Grâce in both scans. The prepositional phrase Grâce à has no intervening syntactic boundary; current removes only that stray comma. Diplomatic preserves it.\par}
\end{minipage}\par\medskip
\noindent\begin{minipage}{\linewidth}
\SGARestoreAppState{EN-VAR-193}%
\paragraph{[193] EN-VAR-193}
Printed reading: (resp. \(K^{\bullet}(Y)\) and\par
Corrected reading: (resp. \(K^{\bullet}(Y)\)) and\par
{\raggedright Source: check/findings/findings\_286\_290.json\#16. Both scans leave the parenthesis opened before resp. unclosed. The first branch allows y in K\_\textbullet{}(Y), the alternative allows y in K\textasciicircum{}\textbullet{}(Y), and x belongs to K\textasciicircum{}\textbullet{}(X) in both branches. The supplied closer after K\textasciicircum{}\textbullet{}(Y) therefore delimits exactly the alternative y-domain. It changes neither the pullback maps (2.12.1)/(2.7.1) nor the product and pushforward in (2.12.4). Diplomatic retains the omission and current supplies that single closer.\par}
\end{minipage}\par\medskip
\noindent\begin{minipage}{\linewidth}
\SGARestoreAppState{EN-VAR-194}%
\paragraph{[194] EN-VAR-194}
Printed reading: \(K\)\par
Corrected reading: \(K_{\bullet}\)\par
{\raggedright Source: check/findings/findings\_291\_295.json\#5. Both scans show the bare K. The two constructions named by the title are identified in 3.2.1 and 3.2.3: the former sends a vector-bundle class at Xi to its pullback class at X, while the latter sends a coherent-sheaf class at Xi to its flat pullback class in K\_\textbullet{}(X). The latter map is explicitly displayed in the directly consulted 293 middle region. The existing current title spells out that second group as K\_\textbullet{}; the diplomatic title preserves the missing mark. No identification between the two K groups is asserted.\par}
\end{minipage}\par\medskip
\noindent\begin{minipage}{\linewidth}
\SGARestoreAppState{EN-VAR-195}%
\paragraph{[195] EN-VAR-195}
Printed reading: \(u:\Ob\Loclib(X)\longrightarrow \lim K^{\bullet}(X_i)_{\naif}\)\par
Corrected reading: \(u:\Ob\Loclib(X)\longrightarrow \varinjlim K^{\bullet}(X_i)_{\naif}\)\par
{\raggedright Source: check/findings/findings\_291\_295.json\#12. The preceding sentence defines u(E) as the image of cl(E\_i) in the explicitly rightward limit of (3.2.1.1). Writing G for that same group, its structure map j\_i sends cl(E\_i) to u(E); the displayed function u:ob Loclib(X)\ensuremath{\to}G has exactly this target. Thus the current rightward limit restores the operator already used in the definition of this same u, while diplomatic preserves the observed absence. No new limit or comparison map is introduced.\par}
\end{minipage}\par\medskip
\noindent\begin{minipage}{\linewidth}
\SGARestoreAppState{EN-VAR-196}%
\paragraph{[196] EN-VAR-196}
Printed reading: \begin{equation}\tag{3.2.2.1}
 \varinjlim K_{\bullet}(X_i)\xrightarrow{\sim}K_{\bullet}(X).
\end{equation}\par
Corrected reading: \begin{equation}\tag{3.2.3.1} \varinjlim K_{\bullet}(X_i)\xrightarrow{\sim}K_{\bullet}(X).
\end{equation}\par
{\raggedright Source: check/findings/findings\_291\_295.json\#18. Both displayed maps carry the printed tag (3.2.2.1). Let \ensuremath{\alpha} be the vector-bundle limit map of (3.2.1.1), and let \ensuremath{\beta}:colim K\_\textbullet{}(Xi)\ensuremath{\to}K\_\textbullet{}(X) be the flat-pullback map in 3.2.3, \ensuremath{\beta}(j\_i[F\_i])=[f\_i\textasciicircum{}*F\_i]. The final diagram has right vertical \ensuremath{\beta} and left vertical \ensuremath{\alpha}\ensuremath{\otimes}\ensuremath{\beta}. At a common index k, both routes send [V\_k]\ensuremath{\otimes}[F\_k] to [f\_k\textasciicircum{}*V\_k\ensuremath{\otimes}f\_k\textasciicircum{}*F\_k], by the canonical pullback/tensor isomorphism. This identifies the second tag and both following references with \ensuremath{\beta}. The perfect-complex map \ensuremath{\alpha}\_perf of 3.2.2 is related by the inclusion maps c\_i:K\textasciicircum{}\textbullet{}(Xi)\ensuremath{\to}K\_\textbullet{}(Xi) and c\_X:K\textasciicircum{}\textbullet{}(X)\ensuremath{\to}K\_\textbullet{}(X): \ensuremath{\beta}\ensuremath{\circ}colim(c\_i)=c\_X\ensuremath{\circ}\ensuremath{\alpha}\_perf, since both send a perfect complex E\_i to its derived pullback class, and flat pullback computes that class here. The current relabelling (3.2.3.1) therefore follows the exact map used in the diagram; diplomatic preserves the duplicate tag and both printed references.\par}
\end{minipage}\par\medskip
\noindent\begin{minipage}{\linewidth}
\SGARestoreAppState{EN-VAR-197}%
\paragraph{[197] EN-VAR-197}
Printed reading: via the isomorphisms (3.2.1.1) and (3.2.2.1) with\par
Corrected reading: via the isomorphisms (3.2.1.1) and (3.2.3.1) with\par
{\raggedright Source: check/findings/findings\_291\_295.json\#18. Both displayed maps carry the printed tag (3.2.2.1). Let \ensuremath{\alpha} be the vector-bundle limit map of (3.2.1.1), and let \ensuremath{\beta}:colim K\_\textbullet{}(Xi)\ensuremath{\to}K\_\textbullet{}(X) be the flat-pullback map in 3.2.3, \ensuremath{\beta}(j\_i[F\_i])=[f\_i\textasciicircum{}*F\_i]. The final diagram has right vertical \ensuremath{\beta} and left vertical \ensuremath{\alpha}\ensuremath{\otimes}\ensuremath{\beta}. At a common index k, both routes send [V\_k]\ensuremath{\otimes}[F\_k] to [f\_k\textasciicircum{}*V\_k\ensuremath{\otimes}f\_k\textasciicircum{}*F\_k], by the canonical pullback/tensor isomorphism. This identifies the second tag and both following references with \ensuremath{\beta}. The perfect-complex map \ensuremath{\alpha}\_perf of 3.2.2 is related by the inclusion maps c\_i:K\textasciicircum{}\textbullet{}(Xi)\ensuremath{\to}K\_\textbullet{}(Xi) and c\_X:K\textasciicircum{}\textbullet{}(X)\ensuremath{\to}K\_\textbullet{}(X): \ensuremath{\beta}\ensuremath{\circ}colim(c\_i)=c\_X\ensuremath{\circ}\ensuremath{\alpha}\_perf, since both send a perfect complex E\_i to its derived pullback class, and flat pullback computes that class here. The current relabelling (3.2.3.1) therefore follows the exact map used in the diagram; diplomatic preserves the duplicate tag and both printed references.\par}
\end{minipage}\par\medskip
\noindent\begin{minipage}{\linewidth}
\SGARestoreAppState{EN-VAR-198}%
\paragraph{[198] EN-VAR-198}
Printed reading: the vertical arrows are defined by (3.2.1.1) and (3.2.2.1)\par
Corrected reading: the vertical arrows are defined by (3.2.1.1) and (3.2.3.1)\par
{\raggedright Source: check/findings/findings\_291\_295.json\#18. Both displayed maps carry the printed tag (3.2.2.1). Let \ensuremath{\alpha} be the vector-bundle limit map of (3.2.1.1), and let \ensuremath{\beta}:colim K\_\textbullet{}(Xi)\ensuremath{\to}K\_\textbullet{}(X) be the flat-pullback map in 3.2.3, \ensuremath{\beta}(j\_i[F\_i])=[f\_i\textasciicircum{}*F\_i]. The final diagram has right vertical \ensuremath{\beta} and left vertical \ensuremath{\alpha}\ensuremath{\otimes}\ensuremath{\beta}. At a common index k, both routes send [V\_k]\ensuremath{\otimes}[F\_k] to [f\_k\textasciicircum{}*V\_k\ensuremath{\otimes}f\_k\textasciicircum{}*F\_k], by the canonical pullback/tensor isomorphism. This identifies the second tag and both following references with \ensuremath{\beta}. The perfect-complex map \ensuremath{\alpha}\_perf of 3.2.2 is related by the inclusion maps c\_i:K\textasciicircum{}\textbullet{}(Xi)\ensuremath{\to}K\_\textbullet{}(Xi) and c\_X:K\textasciicircum{}\textbullet{}(X)\ensuremath{\to}K\_\textbullet{}(X): \ensuremath{\beta}\ensuremath{\circ}colim(c\_i)=c\_X\ensuremath{\circ}\ensuremath{\alpha}\_perf, since both send a perfect complex E\_i to its derived pullback class, and flat pullback computes that class here. The current relabelling (3.2.3.1) therefore follows the exact map used in the diagram; diplomatic preserves the duplicate tag and both printed references.\par}
\end{minipage}\par\medskip
\noindent\begin{minipage}{\linewidth}
\SGARestoreAppState{EN-VAR-199}%
\paragraph{[199] EN-VAR-199}
Printed reading: (III 5.9 b))\par
Corrected reading: (III 5.8 b))\par
{\raggedright Source: ASTRA-IV-REF-287-POSTFREEZE. Both source loci explicitly print5.9 b). The directly inspected III272 heading is Remarques5.8; its clause b starts273 and contains the analytic canonical projection arrow Rf\_*(E) tensor\textasciicircum{}L M\ensuremath{\to}Rf\_*(E tensor\_Y\textasciicircum{}L M), followed by its isomorphism statement and the bibliography, with no intervening5.9. The arrow is the adjoint of Lf\textasciicircum{}*(Rf\_*E tensor\textasciicircum{}L M)\ensuremath{\cong}Lf\textasciicircum{}*Rf\_*E tensor\textasciicircum{}L Lf\textasciicircum{}*M\ensuremath{\to}E tensor\textasciicircum{}L Lf\textasciicircum{}*M, the last map being counit\_E tensor identity\_M. This is exactly the morphism invoked in IV, with no change of functor or arguments. On287, restrict to an open set carrying a bounded finite locally free representative of the perfect M; the bounded coherent E and this M satisfy the displayed5.8(b) domains, and the local projection maps are the restrictions of that same adjunction map. On289, the cited analytic ingredient is the same projection comparison, together with Grauert finiteness and the stated finite tor-dimension. Correct only the locator5.9\ensuremath{\to}5.8 in current; retain the printed locator in diplomatic. Full map and scope details are in projection\_reference\_identification.json.\par}
\end{minipage}\par\medskip
\noindent\begin{minipage}{\linewidth}
\SGARestoreAppState{EN-VAR-200}%
\paragraph{[200] EN-VAR-200}
Printed reading: (III 4.8.1 and 5.9 b))\par
Corrected reading: (III 4.8.1 and 5.8 b))\par
{\raggedright Source: ASTRA-IV-REF-289-POSTFREEZE. Both source loci explicitly print5.9 b). The directly inspected III272 heading is Remarques5.8; its clause b starts273 and contains the analytic canonical projection arrow Rf\_*(E) tensor\textasciicircum{}L M\ensuremath{\to}Rf\_*(E tensor\_Y\textasciicircum{}L M), followed by its isomorphism statement and the bibliography, with no intervening5.9. The arrow is the adjoint of Lf\textasciicircum{}*(Rf\_*E tensor\textasciicircum{}L M)\ensuremath{\cong}Lf\textasciicircum{}*Rf\_*E tensor\textasciicircum{}L Lf\textasciicircum{}*M\ensuremath{\to}E tensor\textasciicircum{}L Lf\textasciicircum{}*M, the last map being counit\_E tensor identity\_M. This is exactly the morphism invoked in IV, with no change of functor or arguments. On287, restrict to an open set carrying a bounded finite locally free representative of the perfect M; the bounded coherent E and this M satisfy the displayed5.8(b) domains, and the local projection maps are the restrictions of that same adjunction map. On289, the cited analytic ingredient is the same projection comparison, together with Grauert finiteness and the stated finite tor-dimension. Correct only the locator5.9\ensuremath{\to}5.8 in current; retain the printed locator in diplomatic. Full map and scope details are in projection\_reference\_identification.json.\par}
\end{minipage}\par\medskip
\noindent\begin{minipage}{\linewidth}
\SGARestoreAppState{EN-VAR-201}%
\paragraph{[201] EN-VAR-201}
Printed reading: \(x\in K_{\bullet}(X)\)\par
Corrected reading: \(x\in K_{\bullet}(X/S)\)\par
{\raggedright Source: ASTRA-IV-PROJECTION-295. The source reads x\ensuremath{\in}K\_\textbullet{}(X) in the first branch and a\ensuremath{\in}K\_\textbullet{}(S), x\ensuremath{\in}K\textasciicircum{}\textbullet{}(X/S) in the second. The immediately preceding pushforward is defined on K\_\textbullet{}(X/S). The first domain is repaired to that exact relative domain: for a=[A] perfect on S and x=[E] in K\_\textbullet{}(X/S), ax=[E\ensuremath{\otimes}\textasciicircum{}L\_S A], and III3.7 identifies Rf\_*E\ensuremath{\otimes}\textasciicircum{}L\_S A with Rf\_*(E\ensuremath{\otimes}\textasciicircum{}L\_S A); its finite-amplitude condition holds for A. The second coefficient is retained: 3.3.2.2 applied to X\ensuremath{\to}S\ensuremath{\to}S pairs x=[E] in K\textasciicircum{}\textbullet{}(X/S) with a=[M] in K\_\textbullet{}(S), giving [E\ensuremath{\otimes}\textasciicircum{}L\_S M] in K\_\textbullet{}(X/S). Here III3.7 uses E’s finite S-amplitude, and both sides land in K\_\textbullet{}(Y/S). The source absolute first domain fails concretely for S=Y=Spec k[x1,x2,…], X=Spec k and the closed immersion f defined by all variables: O\_X is perfect over k, whereas f\_*O\_X=k is not finitely presented over the base ring, so the absolute pushforward does not land in D\textasciicircum{}b(Y)\_coh. The complete map identification and local finite-presentation obstruction are saved in identification\_295.json. Only the first missing /S is changed.\par}
\end{minipage}\par\medskip
\fi
\ifnum\SGAReaderForm=1
\clearpage
\section*{VII contribution apparatus: rendering probe}
\noindent\begin{minipage}{\linewidth}\raggedright
\SGARestoreAppState{EN-CONTRIB-VII-001}%
\paragraph{[VII 1] EN-CONTRIB-VII-001}
Printed reading: can be generated either the $f_i$ or by the $g_j$\par
Corrected reading: can be generated either by the $f_i$ or by the $g_j$\par
{\raggedright Source: check/findings/findings\_416\_420.json\#16. The source omits par before the first family. The parallel soit par les g\_j and the passive construction engendré par identify the missing preposition. The existing current insertion is exact.\par}
\end{minipage}\par\medskip
\noindent\begin{minipage}{\linewidth}\raggedright
\SGARestoreAppState{EN-CONTRIB-VII-002}%
\paragraph{[VII 2] EN-CONTRIB-VII-002}
Printed reading: in a neighborhood of $\ddot y$ since\par
Corrected reading: in a neighborhood of $y$ since\par
{\raggedright Source: check/findings/findings\_416\_420.json\#18. The witness431 clearly separates the two dots, while the main makes a short horizontal blob. The point in this sentence is the same y introduced earlier; the statement concerns the stalk J\_y and its neighborhood under the identity point map. No new decorated point is defined. Retain diplomatic ddot y and the existing current y.\par}
\end{minipage}\par\medskip
\noindent\begin{minipage}{\linewidth}\raggedright
\SGARestoreAppState{EN-CONTRIB-VII-003}%
\paragraph{[VII 3] EN-CONTRIB-VII-003}
Printed reading: $\Tor_i^{A_0}(C_0,B)$\par
Corrected reading: $\Tor_i^{B_0}(C_0,B)$\par
{\raggedright Source: check/findings/findings\_421\_425.json\#11. Both main and witness434 clearly print A\_0. The third complex is K\_\{B\_0\}(T\_1,…,T\_m)\ensuremath{\otimes}\_\{B\_0\}B, with T\_j acting as g-bar\_j. The first factor is the free Koszul resolution of C\_0=B\_0/(T\_1,…,T\_m), so its homology is exactly Tor\_i\textasciicircum{}\{B\_0\}(C\_0,B), as the third line of1.7.2 independently states. This cannot be replaced by Tor over A\_0: in the admissible case n=1,m=0,A=A\_0=Z[T],B=B\_0=C=C\_0=Z, the third Koszul complex is the empty complex Z in degree0, whereas Tor\_1\textasciicircum{}\{Z[T]\}(Z,Z)=Z from0\ensuremath{\to}Z[T] --T\ensuremath{\to} Z[T]\ensuremath{\to}Z\ensuremath{\to}0. Tor\_1\textasciicircum{}Z(Z,Z)=0 agrees with the actual complex. Correct only the current superscript to B\_0.\par}
\end{minipage}\par\medskip
\noindent\begin{minipage}{\linewidth}\raggedright
\SGARestoreAppState{EN-CONTRIB-VII-004}%
\paragraph{[VII 4] EN-CONTRIB-VII-004}
Printed reading: the sequence (1.8.1) is exact.\par
Corrected reading: the sequence (1.6.1) is exact.\par
{\raggedright Source: check/findings/findings\_421\_425.json\#16. The main source clearly prints1.8.1. The displayed sequence immediately preceding the argument is the conormal sequence1.6.1, with maps induced by I\ensuremath{\subset}J and J\ensuremath{\to}J/I. The proof shows the left kernel is locally free of rank0 and hence0, extending that sequence by0 on the left as in1.7(i). The following425 canonical projective-bundle sequence has different terms and is unnumbered in the source. The existing current1.6.1 therefore identifies the actual sequence proved exact.\par}
\end{minipage}\par\medskip
\noindent\begin{minipage}{\linewidth}\raggedright
\SGARestoreAppState{EN-CONTRIB-VII-005}%
\paragraph{[VII 5] EN-CONTRIB-VII-005}
Printed reading: \(\underline{\underline{S}}(\underline{E})\longrightarrow\bigoplus_{n\ge0}\underline{J}^n\)\par
Corrected reading: \(\mathbf{S}(\underline{E})\longrightarrow\bigoplus_{n\ge0}\underline{J}^n\)\par
{\raggedright Source: ASTRA-VII-0023. The four main-scan glyphs each have two distinct lower strokes. The witness436 lower strip confirms the last three; the first is clear in the main upper strip. Restore both rules literally in diplomatic. Current retains the explicitly declared NOTE\_6 bold-S representation of this same double-rule symmetric-algebra notation. The exact mathematical object is S(E)=\ensuremath{\oplus}\_\{n\ensuremath{\geq}0\}Sym\textasciicircum{}n(E), with multiplication induced by symmetric tensor product. This occurrence-scoped transport preserves the distinction from the plain typed word Sym on418 and from the separately single-underlined E. It is not propagated to other exposés.\par}
\end{minipage}\par\medskip
\noindent\begin{minipage}{\linewidth}\raggedright
\SGARestoreAppState{EN-CONTRIB-VII-006}%
\paragraph{[VII 6] EN-CONTRIB-VII-006}
Printed reading: \uline{The} \(\underline{\underline{S}}(\underline{E})\)-\uline{Module}\par
Corrected reading: \uline{The} \(\mathbf{S}(\underline{E})\)-\uline{Module}\par
{\raggedright Source: ASTRA-VII-0023. The four main-scan glyphs each have two distinct lower strokes. The witness436 lower strip confirms the last three; the first is clear in the main upper strip. Restore both rules literally in diplomatic. Current retains the explicitly declared NOTE\_6 bold-S representation of this same double-rule symmetric-algebra notation. The exact mathematical object is S(E)=\ensuremath{\oplus}\_\{n\ensuremath{\geq}0\}Sym\textasciicircum{}n(E), with multiplication induced by symmetric tensor product. This occurrence-scoped transport preserves the distinction from the plain typed word Sym on418 and from the separately single-underlined E. It is not propagated to other exposés.\par}
\end{minipage}\par\medskip
\noindent\begin{minipage}{\linewidth}\raggedright
\SGARestoreAppState{EN-CONTRIB-VII-007}%
\paragraph{[VII 7] EN-CONTRIB-VII-007}
Printed reading: \uline{is perfect as a graded} \(\underline{\underline{S}}(\underline{E})\)-\uline{Module}\par
Corrected reading: \uline{is perfect as a graded} \(\mathbf{S}(\underline{E})\)-\uline{Module}\par
{\raggedright Source: ASTRA-VII-0023. The four main-scan glyphs each have two distinct lower strokes. The witness436 lower strip confirms the last three; the first is clear in the main upper strip. Restore both rules literally in diplomatic. Current retains the explicitly declared NOTE\_6 bold-S representation of this same double-rule symmetric-algebra notation. The exact mathematical object is S(E)=\ensuremath{\oplus}\_\{n\ensuremath{\geq}0\}Sym\textasciicircum{}n(E), with multiplication induced by symmetric tensor product. This occurrence-scoped transport preserves the distinction from the plain typed word Sym on418 and from the separately single-underlined E. It is not propagated to other exposés.\par}
\end{minipage}\par\medskip
\noindent\begin{minipage}{\linewidth}\raggedright
\SGARestoreAppState{EN-CONTRIB-VII-008}%
\paragraph{[VII 8] EN-CONTRIB-VII-008}
Printed reading: \uline{by graded} \(\underline{\underline{S}}(\underline{E})\)-\uline{Modules locally free of finite type}\par
Corrected reading: \uline{by graded} \(\mathbf{S}(\underline{E})\)-\uline{Modules locally free of finite type}\par
{\raggedright Source: ASTRA-VII-0023. The four main-scan glyphs each have two distinct lower strokes. The witness436 lower strip confirms the last three; the first is clear in the main upper strip. Restore both rules literally in diplomatic. Current retains the explicitly declared NOTE\_6 bold-S representation of this same double-rule symmetric-algebra notation. The exact mathematical object is S(E)=\ensuremath{\oplus}\_\{n\ensuremath{\geq}0\}Sym\textasciicircum{}n(E), with multiplication induced by symmetric tensor product. This occurrence-scoped transport preserves the distinction from the plain typed word Sym on418 and from the separately single-underlined E. It is not propagated to other exposés.\par}
\end{minipage}\par\medskip
\noindent\begin{minipage}{\linewidth}\raggedright
\SGARestoreAppState{EN-CONTRIB-VII-009}%
\paragraph{[VII 9] EN-CONTRIB-VII-009}
Printed reading: Since every element of \(\mathbb Z[T_1,\ldots,T_d,X_2/X_1,\ldots,X_d,X_1]\)\par
Corrected reading: Since every element of \(\mathbb Z[T_1,\ldots,T_d,X_2/X_1,\ldots,X_d/X_1]\)\par
{\raggedright Source: check/findings/findings\_426\_430.json\#9. The main and witness437 both clearly print X\_d,X\_1. Set z\_i=X\_i/X\_1; the chart ring previously defined on this page is S=Z[T\_1,…,T\_d,z\_2,…,z\_d]. The substitutions T\ensuremath{\prime}\_i=T\_i\ensuremath{-}T\_1z\_i have the polynomial inverse T\_i=T\ensuremath{\prime}\_i+T\_1z\_i, so the claimed coordinate expression is exact for S. It fails for the printed ring with the additional X\_1: for d=2 in Z[T\_1,T\_2,X\_1\textasciicircum{}\{\ensuremath{\pm}1\},X\_2], specialize T\_1=T\_2=X\_2=0 and X\_1=w. All the claimed new generators T\_1,T\ensuremath{\prime}\_2,z\_2 map to0, but X\_1 maps to the nonconstant invertible indeterminate w. Thus X\_1 cannot be a polynomial in those generators. The existing current slash restores precisely the chart ring; diplomatic preserves the source comma.\par}
\end{minipage}\par\medskip
\noindent\begin{minipage}{\linewidth}\raggedright
\SGARestoreAppState{EN-CONTRIB-VII-010}%
\paragraph{[VII 10] EN-CONTRIB-VII-010}
Printed reading: \(\underline{\underline{S}}(E_0)\)\par
Corrected reading: \(\mathbf{S}(E_0)\)\par
{\raggedright Source: ASTRA-VII-0028. The directly viewed occurrences have two separate rules below S. Keep those literally in diplomatic, and use the explicitly documented bold-S notation transport in current, as for the inspected425 occurrences. The object in each occurrence is the same symmetric algebra \ensuremath{\oplus} Sym\textasciicircum{}n(E), with its symmetric-product multiplication. E and E\_0 receive only their separately observed individual rules. This decision applies only to these counted VII occurrences.\par}
\end{minipage}\par\medskip
\noindent\begin{minipage}{\linewidth}\raggedright
\SGARestoreAppState{EN-CONTRIB-VII-011}%
\paragraph{[VII 11] EN-CONTRIB-VII-011}
Printed reading: by graded free \(\underline{\underline{S}}(E_0)\)-modules\par
Corrected reading: by graded free \(\mathbf{S}(E_0)\)-modules\par
{\raggedright Source: ASTRA-VII-0029. The directly viewed occurrences have two separate rules below S. Keep those literally in diplomatic, and use the explicitly documented bold-S notation transport in current, as for the inspected425 occurrences. The object in each occurrence is the same symmetric algebra \ensuremath{\oplus} Sym\textasciicircum{}n(E), with its symmetric-product multiplication. E and E\_0 receive only their separately observed individual rules. This decision applies only to these counted VII occurrences.\par}
\end{minipage}\par\medskip
\noindent\begin{minipage}{\linewidth}\raggedright
\SGARestoreAppState{EN-CONTRIB-VII-012}%
\paragraph{[VII 12] EN-CONTRIB-VII-012}
Printed reading: the last kernel is then a graded \(\underline{\underline{S}}(E_0)\)-module\par
Corrected reading: the last kernel is then a graded \(\mathbf{S}(E_0)\)-module\par
{\raggedright Source: ASTRA-VII-0029. The directly viewed occurrences have two separate rules below S. Keep those literally in diplomatic, and use the explicitly documented bold-S notation transport in current, as for the inspected425 occurrences. The object in each occurrence is the same symmetric algebra \ensuremath{\oplus} Sym\textasciicircum{}n(E), with its symmetric-product multiplication. E and E\_0 receive only their separately observed individual rules. This decision applies only to these counted VII occurrences.\par}
\end{minipage}\par\medskip
\noindent\begin{minipage}{\linewidth}\raggedright
\SGARestoreAppState{EN-CONTRIB-VII-013}%
\paragraph{[VII 13] EN-CONTRIB-VII-013}
Printed reading: as an ungraded \(\underline{\underline{S}}(E_0)\)-module\par
Corrected reading: as an ungraded \(\mathbf{S}(E_0)\)-module\par
{\raggedright Source: ASTRA-VII-0029. The directly viewed occurrences have two separate rules below S. Keep those literally in diplomatic, and use the explicitly documented bold-S notation transport in current, as for the inspected425 occurrences. The object in each occurrence is the same symmetric algebra \ensuremath{\oplus} Sym\textasciicircum{}n(E), with its symmetric-product multiplication. E and E\_0 receive only their separately observed individual rules. This decision applies only to these counted VII occurrences.\par}
\end{minipage}\par\medskip
\noindent\begin{minipage}{\linewidth}\raggedright
\SGARestoreAppState{EN-CONTRIB-VII-014}%
\paragraph{[VII 14] EN-CONTRIB-VII-014}
Printed reading: as a graded \(\underline{\underline{S}}(\underline{E}_0)\)-module\par
Corrected reading: as a graded \(\mathbf{S}(\underline{E}_0)\)-module\par
{\raggedright Source: ASTRA-VII-0029. The directly viewed occurrences have two separate rules below S. Keep those literally in diplomatic, and use the explicitly documented bold-S notation transport in current, as for the inspected425 occurrences. The object in each occurrence is the same symmetric algebra \ensuremath{\oplus} Sym\textasciicircum{}n(E), with its symmetric-product multiplication. E and E\_0 receive only their separately observed individual rules. This decision applies only to these counted VII occurrences.\par}
\end{minipage}\par\medskip
\noindent\begin{minipage}{\linewidth}\raggedright
\SGARestoreAppState{EN-CONTRIB-VII-015}%
\paragraph{[VII 15] EN-CONTRIB-VII-015}
Printed reading: a graded complex of graded \(\underline{\underline{S}}(E)\)-modules\par
Corrected reading: a graded complex of graded \(\mathbf{S}(E)\)-modules\par
{\raggedright Source: ASTRA-VII-0029. The directly viewed occurrences have two separate rules below S. Keep those literally in diplomatic, and use the explicitly documented bold-S notation transport in current, as for the inspected425 occurrences. The object in each occurrence is the same symmetric algebra \ensuremath{\oplus} Sym\textasciicircum{}n(E), with its symmetric-product multiplication. E and E\_0 receive only their separately observed individual rules. This decision applies only to these counted VII occurrences.\par}
\end{minipage}\par\medskip
\noindent\begin{minipage}{\linewidth}\raggedright
\SGARestoreAppState{EN-CONTRIB-VII-016}%
\paragraph{[VII 16] EN-CONTRIB-VII-016}
Printed reading: is a graded \(\underline{\underline{S}}(E)\)-Module whose\par
Corrected reading: is a graded \(\mathbf{S}(E)\)-Module whose\par
{\raggedright Source: ASTRA-VII-0029. The directly viewed occurrences have two separate rules below S. Keep those literally in diplomatic, and use the explicitly documented bold-S notation transport in current, as for the inspected425 occurrences. The object in each occurrence is the same symmetric algebra \ensuremath{\oplus} Sym\textasciicircum{}n(E), with its symmetric-product multiplication. E and E\_0 receive only their separately observed individual rules. This decision applies only to these counted VII occurrences.\par}
\end{minipage}\par\medskip
\noindent\begin{minipage}{\linewidth}\raggedright
\SGARestoreAppState{EN-CONTRIB-VII-017}%
\paragraph{[VII 17] EN-CONTRIB-VII-017}
Printed reading: \(X'\to\underline{\underline{P}}(\underline{E})\)\par
Corrected reading: \(X'\to\bfP(\underline{E})\)\par
{\raggedright Source: ASTRA-VII-0032. Both printed projective-bundle P letters have two distinct rules, independently of the single E rule. Diplomatic preserves both P strokes literally. Current uses the already declared bold-P notation for the same Proj(Sym E) construction: the identity of the graded symmetric algebra induces the identity on every degree-zero standard affine chart and hence on the glued projective bundle. This is an occurrence-scoped rendering choice, with no change of object.\par}
\end{minipage}\par\medskip
\noindent\begin{minipage}{\linewidth}\raggedright
\SGARestoreAppState{EN-CONTRIB-VII-018}%
\paragraph{[VII 18] EN-CONTRIB-VII-018}
Printed reading: \(\underline{\underline{P}}(\underline{E})\to X\)\par
Corrected reading: \(\bfP(\underline{E})\to X\)\par
{\raggedright Source: ASTRA-VII-0032. Both printed projective-bundle P letters have two distinct rules, independently of the single E rule. Diplomatic preserves both P strokes literally. Current uses the already declared bold-P notation for the same Proj(Sym E) construction: the identity of the graded symmetric algebra induces the identity on every degree-zero standard affine chart and hence on the glued projective bundle. This is an occurrence-scoped rendering choice, with no change of object.\par}
\end{minipage}\par\medskip
\noindent\begin{minipage}{\linewidth}\raggedright
\SGARestoreAppState{EN-CONTRIB-VII-019}%
\paragraph{[VII 19] EN-CONTRIB-VII-019}
Printed reading: \(M\), Let \(P_\bullet\)\par
Corrected reading: \(M\). Let \(P_\bullet\)\par
{\raggedright Source: check/findings/findings\_431\_435.json\#10. The enlarged main lower strip shows the descending comma beside the overstruck M. The following capital Soit introduces a new independent sentence, so current full stop is the established punctuation correction; diplomatic comma is retained.\par}
\end{minipage}\par\medskip
\noindent\begin{minipage}{\linewidth}\raggedright
\SGARestoreAppState{EN-CONTRIB-VII-020}%
\paragraph{[VII 20] EN-CONTRIB-VII-020}
Printed reading: \uline{conditions satisfied in particular in the noetherian case when the sequence} \(f_1,\ldots,f_p,g_1,\ldots g_r,h_1,\ldots,h_s\)\par
Corrected reading: \uline{conditions satisfied in particular in the noetherian case when the sequence} \(f_1,\ldots,f_p,g_1,\ldots,g_r,h_1,\ldots,h_s\)\par
{\raggedright Source: check/findings/findings\_431\_435.json\#13. The main upper strip visibly lacks a comma between the last ellipsis dot and g\_r in this one list; the other corresponding lists include it. The source lists the same successive entries g\_1 through g\_r, and current restores the missing punctuation without altering the sequence.\par}
\end{minipage}\par\medskip
\noindent\begin{minipage}{\linewidth}\raggedright
\SGARestoreAppState{EN-CONTRIB-VII-021}%
\paragraph{[VII 21] EN-CONTRIB-VII-021}
Printed reading: \[
\underline{\Tor}_1^{\underline{O}_X}(\underline{O}_Y,\underline{O}_Z)\xrightarrow{\sim}(\underline{O}_T)^n,
\]\par
Corrected reading: \[
\underline{\Tor}_1^{\underline{O}_X}(\underline{O}_Y,\underline{O}_Z)\xrightarrow{\sim}(\underline{O}_T)^p,
\]\par
{\raggedright Source: check/findings/findings\_431\_435.json\#32. Both scans clearly print n. The preceding formula identifies Tor\_i with O\_T\ensuremath{\otimes}\ensuremath{\wedge}\textasciicircum{}iE, with E=O\_X\textasciicircum{}p by the explicit definition on432. At i=1 the basis classes[e\_j] for1\ensuremath{\leq}j\ensuremath{\leq}p map bijectively to the p standard coordinates of O\_T\textasciicircum{}p. Thus the rank is precisely p, also stated independently on435. No n is defined in this proof. Retain the existing current p and diplomatic n.\par}
\end{minipage}\par\medskip
\noindent\begin{minipage}{\linewidth}\raggedright
\SGARestoreAppState{EN-CONTRIB-VII-022}%
\paragraph{[VII 22] EN-CONTRIB-VII-022}
Printed reading: \begin{equation}\tag{2.5.3}
\underline{\Tor}_i^{\underline{O}_X}(\underline{O}_Y,\underline{O}_Z)
\xrightarrow{\sim}
\bigwedge^i\underline{\Tor}_i^{\underline{O}_X}(\underline{O}_Y,\underline{O}_Z).
\end{equation}\par
Corrected reading: \begin{equation}\tag{2.5.3}
\underline{\Tor}_i^{\underline{O}_X}(\underline{O}_Y,\underline{O}_Z)
\xrightarrow{\sim}
\bigwedge^i\underline{\Tor}_1^{\underline{O}_X}(\underline{O}_Y,\underline{O}_Z).
\end{equation}\par
{\raggedright Source: check/findings/findings\_431\_435.json\#35. Main and witness444 plainly print the subscript i. The degree1 identification \ensuremath{\psi}:Tor\_1\ensuremath{\to}O\_T\textasciicircum{}p from the previous display induces \ensuremath{\wedge}\textasciicircum{}i\ensuremath{\psi}; composing the2.5.2 isomorphism Tor\_i\ensuremath{\to}\ensuremath{\wedge}\textasciicircum{}i(O\_T)\textasciicircum{}p with its inverse gives exactly Tor\_i\ensuremath{\to}\ensuremath{\wedge}\textasciicircum{}iTor\_1. The printed Tor\_i cannot supply this: take A=k[t\_1,t\_2], Y=Z=Spec(k), p=2 and r=s=0. Tensoring the Koszul resolution with k gives zero differentials, so Tor\_1=k\textsuperscript{2} and Tor\_2=k. At i=2 the printed right member \ensuremath{\wedge}\textsuperscript{2}Tor\_2=\ensuremath{\wedge}\textsuperscript{2}k=0, while the left member is k. The corrected \ensuremath{\wedge}\textsuperscript{2}Tor\_1=\ensuremath{\wedge}\textsuperscript{2}k\textsuperscript{2}\ensuremath{\cong}k agrees. Keep the existing current1.\par}
\end{minipage}\par\medskip
\noindent\begin{minipage}{\linewidth}\raggedright
\SGARestoreAppState{EN-CONTRIB-VII-023}%
\paragraph{[VII 23] EN-CONTRIB-VII-023}
Printed reading: \[
\bigwedge(\underline{O}_T)^n\xrightarrow{\sim}
\bigoplus_{i\ge0}\underline{\Tor}_i^{\underline{O}_X}(\underline{O}_Y,\underline{O}_Z)
\]\par
Corrected reading: \[
\bigwedge(\underline{O}_T)^p\xrightarrow{\sim}
\bigoplus_{i\ge0}\underline{\Tor}_i^{\underline{O}_X}(\underline{O}_Y,\underline{O}_Z)
\]\par
{\raggedright Source: check/findings/findings\_431\_435.json\#37. Both main and witness445 print n. The source explicitly says this graded isomorphism is defined by2.5.2, whose degreei source is \ensuremath{\wedge}\textasciicircum{}i(O\_T\textasciicircum{}p). The direct sum of those degreei maps has source\ensuremath{\wedge}(O\_T\textasciicircum{}p), with degree1 basis indexed1,…,p. This is the same established rank correction as433; current had not yet propagated it to this independently printed occurrence. Correct only this current n to p.\par}
\end{minipage}\par\medskip
\noindent\begin{minipage}{\linewidth}\raggedright
\SGARestoreAppState{EN-CONTRIB-VII-024}%
\paragraph{[VII 24] EN-CONTRIB-VII-024}
Printed reading: \[
\bigwedge\underline{\Tor}_i^{\underline{O}_X}(\underline{O}_Y,\underline{O}_Z)\xrightarrow{\sim}
\bigoplus_{i\ge0}\underline{\Tor}_i^{\underline{O}_X}(\underline{O}_Y,\underline{O}_Z)
\]\par
Corrected reading: \[
\bigwedge\underline{\Tor}_1^{\underline{O}_X}(\underline{O}_Y,\underline{O}_Z)\xrightarrow{\sim}
\bigoplus_{i\ge0}\underline{\Tor}_i^{\underline{O}_X}(\underline{O}_Y,\underline{O}_Z)
\]\par
{\raggedright Source: check/findings/findings\_431\_435.json\#41. Both scans print the subscript i. This map is explicitly the direct sum of2.5.3 and the restriction of2.5.1, which is defined to extend the identity in degree1. Its generator is therefore the degree1 sheaf Tor\_1 under that identity. In the admissible p=2,r=s=0 example A=k[t\_1,t\_2], Y=Z=Spec(k), the displayed target has dimensions1,2,1 in degrees0,1,2. Substituting Tor\_2=k as the printed generator would instead give exterior dimensions1,1,0, so the displayed graded isomorphism fails. The existing current Tor\_1 gives exactly1,2,1 and the stipulated identity in degree1.\par}
\end{minipage}\par\medskip
\noindent\begin{minipage}{\linewidth}\raggedright
\SGARestoreAppState{EN-CONTRIB-VII-025}%
\paragraph{[VII 25] EN-CONTRIB-VII-025}
Printed reading: \(\underline{\underline{R}}i_*(\underline{L}^\bullet)=i_*(\underline{L}^\bullet)\)\par
Corrected reading: \(\mathbb Ri_*(\underline{L}^\bullet)=i_*(\underline{L}^\bullet)\)\par
{\raggedright Source: check/findings/findings\_431\_435.json\#48. Main and witness446 show a typed R with two independent rules below it. Diplomatic now preserves those two strokes literally; current retains the documented double-struck R notation for the same derived functor. The source itself identifies Ri\_*(L\textasciicircum{}\textbullet{})=i\_*(L\textasciicircum{}\textbullet{}) because i\_* is exact: applying it termwise preserves acyclic complexes, so the canonical derived comparison is the identity represented by that complex. The three L letters have their separate single rules.\par}
\end{minipage}\par\medskip
\noindent\begin{minipage}{\linewidth}\raggedright
\SGARestoreAppState{EN-CONTRIB-VII-026}%
\paragraph{[VII 26] EN-CONTRIB-VII-026}
Printed reading: the class of the complex \(\underline{\underline{L}} i^*(\underline{\underline{R}}i_*(\underline{M}))\)\par
Corrected reading: the class of the complex \(\mathbb L i^*(\mathbb Ri_*(\underline{M}))\)\par
{\raggedright Source: ASTRA-VII-0055. ASTRA-VII-0055: The complete main upper strip clearly shows two independent lower strokes for each of these derived-functor letters. Diplomatic reproduces the strokes literally. Current retains the declared double-struck representation of the same derived functor already bound by the author. Single M and O rules remain independent. The exact local glyph count is recorded. ASTRA-VII-0056: The complete main upper strip clearly shows two independent lower strokes for each of these derived-functor letters. Diplomatic reproduces the strokes literally. Current retains the declared double-struck representation of the same derived functor already bound by the author. Single M and O rules remain independent. The exact local glyph count is recorded.\par}
\end{minipage}\par\medskip
\noindent\begin{minipage}{\linewidth}\raggedright
\SGARestoreAppState{EN-CONTRIB-VII-027}%
\paragraph{[VII 27] EN-CONTRIB-VII-027}
Printed reading: a projective resolution of the complex \(\underline{\underline{R}}i_*(\underline{M})=i_*(\underline{M})\)\par
Corrected reading: a projective resolution of the complex \(\mathbb Ri_*(\underline{M})=i_*(\underline{M})\)\par
{\raggedright Source: ASTRA-VII-0056. The complete main upper strip clearly shows two independent lower strokes for each of these derived-functor letters. Diplomatic reproduces the strokes literally. Current retains the declared double-struck representation of the same derived functor already bound by the author. Single M and O rules remain independent. The exact local glyph count is recorded.\par}
\end{minipage}\par\medskip
\noindent\begin{minipage}{\linewidth}\raggedright
\SGARestoreAppState{EN-CONTRIB-VII-028}%
\paragraph{[VII 28] EN-CONTRIB-VII-028}
Printed reading: \(\underline{\underline{L}} i^*(\underline{\underline{R}}i_*(\underline{M}))\) is the alternating sum\par
Corrected reading: \(\mathbb L i^*(\mathbb Ri_*(\underline{M}))\) is the alternating sum\par
{\raggedright Source: ASTRA-VII-0055. ASTRA-VII-0055: The complete main upper strip clearly shows two independent lower strokes for each of these derived-functor letters. Diplomatic reproduces the strokes literally. Current retains the declared double-struck representation of the same derived functor already bound by the author. Single M and O rules remain independent. The exact local glyph count is recorded. ASTRA-VII-0056: The complete main upper strip clearly shows two independent lower strokes for each of these derived-functor letters. Diplomatic reproduces the strokes literally. Current retains the declared double-struck representation of the same derived functor already bound by the author. Single M and O rules remain independent. The exact local glyph count is recorded.\par}
\end{minipage}\par\medskip
\noindent\begin{minipage}{\linewidth}\raggedright
\SGARestoreAppState{EN-CONTRIB-VII-029}%
\paragraph{[VII 29] EN-CONTRIB-VII-029}
Printed reading: defined by an ideal of \(\underline{O}_X\), canonically\par
Corrected reading: defined by an ideal of \(\underline{O}_{X'}\), canonically\par
{\raggedright Source: check/findings/findings\_436\_440.json\#18. Main and witness448 clearly print unprimed X. The sentence describes the invertible defining ideal of Y\ensuremath{\prime} in X\ensuremath{\prime}, explicitly named J\ensuremath{\prime}\ensuremath{\subset}O\_\{X\ensuremath{\prime}\} a few lines later and identified with O\_\{X\ensuremath{\prime}\}(1). The canonical map from the original ideal is f*J\ensuremath{\to}J\ensuremath{\prime}, not in general an isomorphism: for A=k[x,y], J=(x,y), on the blow-up chart A\ensuremath{\prime}=k[x,t], y=xt, f*J=A\ensuremath{\prime}\textsuperscript{2}/A\ensuremath{\prime}(-xt,x). Its map to J\ensuremath{\prime}=(x) sends(a,b) to x(a+tb); the nonzero kernel is generated by(-t,1) modulo x(-t,1), hence is A\ensuremath{\prime}/(x). The asserted invertible ideal is therefore the displayed image J\ensuremath{\prime} on X\ensuremath{\prime}. The existing current prime restores exactly that bound sheaf; diplomatic retains the printed O\_X.\par}
\end{minipage}\par\medskip
\noindent\begin{minipage}{\linewidth}\raggedright
\SGARestoreAppState{EN-CONTRIB-VII-030}%
\paragraph{[VII 30] EN-CONTRIB-VII-030}
Printed reading: \(j:Y'\to X\)\par
Corrected reading: \(j:Y'\to X'\)\par
{\raggedright Source: check/findings/findings\_436\_440.json\#19. Main and witness448 print the target X without a prime. The same page’s diagram explicitly binds j as Y\ensuremath{\prime}\ensuremath{\to}X\ensuremath{\prime}. The composition f\ensuremath{\circ}j=i\ensuremath{\circ}g maps instead to X; these arrows are related by that exact composition. In the admissible blow-up of A\textsuperscript{2}\_k at its origin, Y\ensuremath{\prime}=P\textsuperscript{1}\_k and f\ensuremath{\circ}j is constant to the origin: the two distinct points[1:0] and[0:1] have the same image, so this composite cannot be the asserted closed immersion. The diagram’s j is the codimension1 exceptional-divisor immersion, defined on the x-chart by(x)\ensuremath{\subset}k[x,t]. Retain the existing current target X\ensuremath{\prime} and diplomatic X.\par}
\end{minipage}\par\medskip
\noindent\begin{minipage}{\linewidth}\raggedright
\SGARestoreAppState{EN-CONTRIB-VII-031}%
\paragraph{[VII 31] EN-CONTRIB-VII-031}
Printed reading: \[
Y'=\underline{\underline{P}}(\underline{N}).
\]\par
Corrected reading: \[
Y'=\bfP(\underline{N}).
\]\par
{\raggedright Source: ASTRA-VII-0059. The directly inspected P(N) occurrence has two separate lower P strokes. Diplomatic preserves them literally. Current retains the declared bold-P representation of Proj(Sym N); the identity graded algebra map gives the identity on its standard affine charts and glued projective bundle. The N single rule is recorded independently.\par}
\end{minipage}\par\medskip
\noindent\begin{minipage}{\linewidth}\raggedright
\SGARestoreAppState{EN-CONTRIB-VII-032}%
\paragraph{[VII 32] EN-CONTRIB-VII-032}
Printed reading: \(i_1<..<i_k\)\par
Corrected reading: \(i_1<\cdots<i_k\)\par
{\raggedright Source: check/findings/findings\_436\_440.json\#40. The main and witness449 show exactly two compact dots between the less-than signs, whereas the adjacent wedge products have the usual three. Diplomatic keeps the two source dots; current keeps the established three-dot punctuation for the same ordered index list.\par}
\end{minipage}\par\medskip
\noindent\begin{minipage}{\linewidth}\raggedright
\SGARestoreAppState{EN-CONTRIB-VII-033}%
\paragraph{[VII 33] EN-CONTRIB-VII-033}
Printed reading: \(Y'=\underline{\underline{P}}(\underline{N})\)\par
Corrected reading: \(Y'=\bfP(\underline{N})\)\par
{\raggedright Source: ASTRA-VII-0060. The directly inspected P(N) occurrence has two separate lower P strokes. Diplomatic preserves them literally. Current retains the declared bold-P representation of Proj(Sym N); the identity graded algebra map gives the identity on its standard affine charts and glued projective bundle. The N single rule is recorded independently.\par}
\end{minipage}\par\medskip
\noindent\begin{minipage}{\linewidth}\raggedright
\SGARestoreAppState{EN-CONTRIB-VII-034}%
\paragraph{[VII 34] EN-CONTRIB-VII-034}
Printed reading: the class of the complex \(\underline{\underline{L}}f^*(\underline{\underline{R}}i_*(\underline{M}))\)\par
Corrected reading: the class of the complex \(\mathbb Lf^*(\mathbb Ri_*(\underline{M}))\)\par
{\raggedright Source: ASTRA-VII-0072. ASTRA-VII-0072: The main page and the independently viewed witness show two lower strokes at these occurrences. Diplomatic reproduces those strokes literally. Current retains the already declared double-struck notation of the identical source-bound derived functor. These replacements do not affect the singly underlined sheaf/complex letters or the plain R\textasciicircum{}i higher-direct-image letters. ASTRA-VII-0073: The main page and the independently viewed witness show two lower strokes at these occurrences. Diplomatic reproduces those strokes literally. Current retains the already declared double-struck notation of the identical source-bound derived functor. These replacements do not affect the singly underlined sheaf/complex letters or the plain R\textasciicircum{}i higher-direct-image letters.\par}
\end{minipage}\par\medskip
\noindent\begin{minipage}{\linewidth}\raggedright
\SGARestoreAppState{EN-CONTRIB-VII-035}%
\paragraph{[VII 35] EN-CONTRIB-VII-035}
Printed reading: \(f_*\)\par
Corrected reading: \(f^*\)\par
{\raggedright Source: check/findings/findings\_441\_445.json\#9. The complete main source and witness452 clearly place the star below f. Let P\textasciicircum{}\textbullet{}\ensuremath{\to}i\_*M be the projective resolution on X named in the sentence. The displayed derived pullback is represented on X\ensuremath{\prime} by f*P\textasciicircum{}\textbullet{}=O\_\{X\ensuremath{\prime}\}\ensuremath{\otimes}\_\{f\textasciicircum{}\{-1\}O\_X\}f\textasciicircum{}\{-1\}P\textasciicircum{}\textbullet{}, with a\ensuremath{\otimes}s\ensuremath{\mapsto}a\ensuremath{\cdot}f\textasciicircum{}\{-1\}s termwise; its homology is precisely the Tor\_i\textasciicircum{}\{O\_X\}(M,O\_\{X\ensuremath{\prime}\}) written in the next sentence. The direct-image functor is related by the adjunction unit P\ensuremath{\to}f\_*f*P and counit f*f\_*Q\ensuremath{\to}Q; on an affine chart A\ensuremath{\to}B these are p\ensuremath{\mapsto}1\ensuremath{\otimes}p and b\ensuremath{\otimes}q\ensuremath{\mapsto}bq. They do not replace the named tensor operation by applying f\_* to P. Thus the same sentence and its next displayed homology computation uniquely specify f\textasciicircum{}*, and the low star is an established source error.\par}
\end{minipage}\par\medskip
\noindent\begin{minipage}{\linewidth}\raggedright
\SGARestoreAppState{EN-CONTRIB-VII-036}%
\paragraph{[VII 36] EN-CONTRIB-VII-036}
Printed reading: \(\underline{\underline{L}}f^*(\underline{\underline{R}}i_*(\underline{M}))\) is the alternating sum\par
Corrected reading: \(\mathbb Lf^*(\mathbb Ri_*(\underline{M}))\) is the alternating sum\par
{\raggedright Source: ASTRA-VII-0072. ASTRA-VII-0072: The main page and the independently viewed witness show two lower strokes at these occurrences. Diplomatic reproduces those strokes literally. Current retains the already declared double-struck notation of the identical source-bound derived functor. These replacements do not affect the singly underlined sheaf/complex letters or the plain R\textasciicircum{}i higher-direct-image letters. ASTRA-VII-0073: The main page and the independently viewed witness show two lower strokes at these occurrences. Diplomatic reproduces those strokes literally. Current retains the already declared double-struck notation of the identical source-bound derived functor. These replacements do not affect the singly underlined sheaf/complex letters or the plain R\textasciicircum{}i higher-direct-image letters.\par}
\end{minipage}\par\medskip
\noindent\begin{minipage}{\linewidth}\raggedright
\SGARestoreAppState{EN-CONTRIB-VII-037}%
\paragraph{[VII 37] EN-CONTRIB-VII-037}
Printed reading: which was proved in exp.\par
Corrected reading: which was proved in exp. III.\par
{\raggedright Source: check/findings/findings\_441\_445.json\#14. Main and witness452 end exactly with exp.;442 starts the recalled Lemma and supplies no continuation of the reference. The directly inspected III Corollary2.3 on240/PDF254 states the projective high-twist result for S quasi-compact, a projective morphism of finite-type S-schemes, and E in D\textasciicircum{}+ pseudo-coherent relative to S and acyclic above a. Substitution S=Y=X (the VII base), source=X\ensuremath{\prime}, E=L\textasciicircum{}\textbullet{}, a=0, and the projective twisting O\_\{X\ensuremath{\prime}\}(1) gives exactly the recalled statement on442. Projective f makes X\ensuremath{\prime} quasi-compact; perfect L\textasciicircum{}\textbullet{} is then bounded, and perfect f together with perfect L\textasciicircum{}\textbullet{} supplies the relative pseudo-coherence. The derived tensor with the invertible twist equals the ordinary tensor used on442. This directly verified dependency supports the existing current editorial completion III. It is not a claim that III is visible in the source. Diplomatic retains the incomplete exp.\par}
\end{minipage}\par\medskip
\noindent\begin{minipage}{\linewidth}\raggedright
\SGARestoreAppState{EN-CONTRIB-VII-038}%
\paragraph{[VII 38] EN-CONTRIB-VII-038}
Printed reading: \(\underline{\underline{R}}f_*\bigl(\underline{L}^\bullet\otimes_{\underline{O}_{X'}}\underline{O}_{X'}(n)\bigr)\)\par
Corrected reading: \(\mathbb Rf_*\bigl(\underline{L}^\bullet\otimes_{\underline{O}_{X'}}\underline{O}_{X'}(n)\bigr)\)\par
{\raggedright Source: check/findings/findings\_441\_445.json\#16. Both directly inspected scans show two lower strokes. The local diplomatic patch preserves the strokes literally; current retains the declared double-struck derived-functor notation. The neighboring R\textasciicircum{}i higher-direct-image letters are visibly plain.\par}
\end{minipage}\par\medskip
\noindent\begin{minipage}{\linewidth}\raggedright
\SGARestoreAppState{EN-CONTRIB-VII-039}%
\paragraph{[VII 39] EN-CONTRIB-VII-039}
Printed reading: \(\underline{\underline{R}}f_*(\underline{O}_{Y'}(n))=i_*\circ\underline{\underline{R}}g_*(\underline{O}_{Y'}(n))\)\par
Corrected reading: \(\mathbb Rf_*(\underline{O}_{Y'}(n))=i_*\circ\mathbb Rg_*(\underline{O}_{Y'}(n))\)\par
{\raggedright Source: check/findings/findings\_441\_445.json\#18. check/findings/findings\_441\_445.json\#18: Both directly inspected scans show two lower strokes. The local diplomatic patch preserves the strokes literally; current retains the declared double-struck derived-functor notation. The neighboring R\textasciicircum{}i higher-direct-image letters are visibly plain. check/findings/findings\_441\_445.json\#19: Both directly inspected scans show two lower strokes. The local diplomatic patch preserves the strokes literally; current retains the declared double-struck derived-functor notation. The neighboring R\textasciicircum{}i higher-direct-image letters are visibly plain.\par}
\end{minipage}\par\medskip
\noindent\begin{minipage}{\linewidth}\raggedright
\SGARestoreAppState{EN-CONTRIB-VII-040}%
\paragraph{[VII 40] EN-CONTRIB-VII-040}
Printed reading: \(Y'=\underline{\underline{P}}(\underline{N})\)\par
Corrected reading: \(Y'=\bfP(\underline{N})\)\par
{\raggedright Source: ASTRA-VII-0075. Each counted P in the full main page has two lower rules. Diplomatic records their literal shape. Current keeps the declared bold notation of the same Proj construction; its identity graded-algebra morphism induces identity maps on the standard degree-zero localization charts and hence on the glued projective scheme.\par}
\end{minipage}\par\medskip
\noindent\begin{minipage}{\linewidth}\raggedright
\SGARestoreAppState{EN-CONTRIB-VII-041}%
\paragraph{[VII 41] EN-CONTRIB-VII-041}
Printed reading: \(X'\longrightarrow\underline{\underline{P}}_X^{d-1}\)\par
Corrected reading: \(X'\longrightarrow\bfP_X^{d-1}\)\par
{\raggedright Source: ASTRA-VII-0076. Each counted P in the full main page has two lower rules. Diplomatic records their literal shape. Current keeps the declared bold notation of the same Proj construction; its identity graded-algebra morphism induces identity maps on the standard degree-zero localization charts and hence on the glued projective scheme.\par}
\end{minipage}\par\medskip
\noindent\begin{minipage}{\linewidth}\raggedright
\SGARestoreAppState{EN-CONTRIB-VII-042}%
\paragraph{[VII 42] EN-CONTRIB-VII-042}
Printed reading: On \(\underline{\underline{P}}_X^{d-1}\) we obtain\par
Corrected reading: On \(\bfP_X^{d-1}\) we obtain\par
{\raggedright Source: ASTRA-VII-0076. Each counted P in the full main page has two lower rules. Diplomatic records their literal shape. Current keeps the declared bold notation of the same Proj construction; its identity graded-algebra morphism induces identity maps on the standard degree-zero localization charts and hence on the glued projective scheme.\par}
\end{minipage}\par\medskip
\noindent\begin{minipage}{\linewidth}\raggedright
\SGARestoreAppState{EN-CONTRIB-VII-043}%
\paragraph{[VII 43] EN-CONTRIB-VII-043}
Printed reading: so these are free of finite type\par
Corrected reading: so these are locally free of finite type\par
{\raggedright Source: ASTRA-VII-0084. Write B=O\_X[T\_1,…,T\_d]. A finite homogeneous basis e\_j of degree a\_j gives the degree-preserving isomorphism \ensuremath{\varphi}:\ensuremath{\oplus}\_j B(\ensuremath{-}a\_j)\ensuremath{\to}L, (b\_j)\ensuremath{\mapsto}\ensuremath{\Sigma}\_jb\_je\_j, with inverse the unique homogeneous coordinates. Taking associated sheaves gives \ensuremath{\oplus}\_jO\_P(\ensuremath{-}a\_j)\ensuremath{\cong}Proj(L). On D\_+(T\_i), the degree-zero generators e\_j/T\_i\textasciicircum{}\{a\_j\} give explicit local bases and prove finite local freeness. Global freeness is false for an admissible instance: A=k[x,y], J=(x,y), S=A[T\_1,T\_2]/(yT\_1\ensuremath{-}xT\_2), with presentation B(\ensuremath{-}1)\ensuremath{\to}B\ensuremath{\to}S. Its sheaves are M\ensuremath{\prime\prime}=O\_\{P\textasciicircum{}1\_A\}(\ensuremath{-}1), M\ensuremath{\prime}=O. On the closed fibre P\textasciicircum{}1\_k, choose frames e\_0=(1,t), e\_1=(s,1)=t\textasciicircum{}\{-1\}e\_0. A global section requires a(t)=t\textasciicircum{}\{-1\}b(t\textasciicircum{}\{-1\}); the nonnegative and strictly negative Laurent powers force both sides to vanish. Hence H\textasciicircum{}0(O(\ensuremath{-}1))=0, whereas a free rank-one sheaf has H\textasciicircum{}0=k. The surrounding proof uses exactly finite local freeness and the graded direct sum of twists: regular X\ensuremath{\prime}\ensuremath{\to}P makes O\_\{X\ensuremath{\prime}\} perfect; the two exact kernel sequences with M\ensuremath{\prime} and M\ensuremath{\prime\prime} therefore give perfect kernels, to which the recalled projective high-twist lemma applies. No global free basis is used or available. Add only localement, independently of the separate all-n diagram issue.\par}
\end{minipage}\par\medskip
\noindent\begin{minipage}{\linewidth}\raggedright
\SGARestoreAppState{EN-CONTRIB-VII-044}%
\paragraph{[VII 44] EN-CONTRIB-VII-044}
Printed reading: over \(\underline{\underline{P}}_X^{d-1}\). We can therefore\par
Corrected reading: over \(\bfP_X^{d-1}\). We can therefore\par
{\raggedright Source: ASTRA-VII-0076. Each counted P in the full main page has two lower rules. Diplomatic records their literal shape. Current keeps the declared bold notation of the same Proj construction; its identity graded-algebra morphism induces identity maps on the standard degree-zero localization charts and hence on the glued projective scheme.\par}
\end{minipage}\par\medskip
\noindent\begin{minipage}{\linewidth}\raggedright
\SGARestoreAppState{EN-CONTRIB-VII-045}%
\paragraph{[VII 45] EN-CONTRIB-VII-045}
Printed reading: \(\underline{\underline{P}}_X^{d-1}\to X\)\par
Corrected reading: \(\bfP_X^{d-1}\to X\)\par
{\raggedright Source: ASTRA-VII-0076. Each counted P in the full main page has two lower rules. Diplomatic records their literal shape. Current keeps the declared bold notation of the same Proj construction; its identity graded-algebra morphism induces identity maps on the standard degree-zero localization charts and hence on the glued projective scheme.\par}
\end{minipage}\par\medskip
\noindent\begin{minipage}{\linewidth}\raggedright
\SGARestoreAppState{EN-CONTRIB-VII-046}%
\paragraph{[VII 46] EN-CONTRIB-VII-046}
Printed reading: \[
f_*(\underline{O}_{Y'}(n))
=i_*\bigl(g_*(\underline{O}_{Y'}(n))\bigr)
=i_*(\underline{\underline{S}}(\underline{N})_n)
=\underline{J}^n/\underline{J}^{n+1}.
\]\par
Corrected reading: \[
f_*(\underline{O}_{Y'}(n))
=i_*\bigl(g_*(\underline{O}_{Y'}(n))\bigr)
=i_*(\mathbf S(\underline{N})_n)
=\underline{J}^n/\underline{J}^{n+1}.
\]\par
{\raggedright Source: check/findings/findings\_441\_445.json\#24. The complete main upper strip shows two distinct lower strokes under S and a single stroke under N. Diplomatic now preserves the S strokes literally. Current retains the declared bold S notation of the same symmetric algebra; the identity on N extends degreewise by the symmetric-algebra universal property and identifies the expression with its original source-bound algebra.\par}
\end{minipage}\par\medskip
\noindent\begin{minipage}{\linewidth}\raggedright
\SGARestoreAppState{EN-CONTRIB-VII-047}%
\paragraph{[VII 47] EN-CONTRIB-VII-047}
Printed reading: hence the middle one is also an isomorphism\par
Corrected reading: hence the middle one is also an isomorphism.\par
{\raggedright Source: fixes/MISPRINTS\_6.json\#413. The source line ends immediately after également, and the next line is the separately spaced proposition label. Diplomatic correctly preserves the omission; current already supplies the sentence-ending full stop.\par}
\end{minipage}\par\medskip
\noindent\begin{minipage}{\linewidth}\raggedright
\SGARestoreAppState{EN-CONTRIB-VII-048}%
\paragraph{[VII 48] EN-CONTRIB-VII-048}
Printed reading: (VI 5.9) which gives the result.\par
Corrected reading: (VI 5.9), which gives the result.\par
{\raggedright Source: fixes/MISPRINTS\_6.json\#414. The main lower strip ends the line immediately after the parenthesis; the next line begins d’où le résultat. Diplomatic preserves the printed absence. Current keeps its existing comma separating the consequent clause.\par}
\end{minipage}\par\medskip
\noindent\begin{minipage}{\linewidth}\raggedright
\SGARestoreAppState{EN-CONTRIB-VII-049}%
\paragraph{[VII 49] EN-CONTRIB-VII-049}
Printed reading: \[
v\circ u(y)
=v\bigl(g^*(y)\lambda_{-1}(F),-i_*(y)\bigr)
=j_*\bigl(g^*(y)\lambda_{-1}(F)\bigr)-f^*(i_*(x))=0,
\]\par
Corrected reading: \[
v\circ u(y)
=v\bigl(g^*(y)\lambda_{-1}(F),-i_*(y)\bigr)
=j_*\bigl(g^*(y)\lambda_{-1}(F)\bigr)-f^*(i_*(y))=0,
\]\par
{\raggedright Source: fixes/RECONCILE\_CODEX\_6.json\#197. The main lower strip clearly prints x. The same page defines u(y)=(g*y\ensuremath{\cdot}\ensuremath{\lambda}\_\{\ensuremath{-}1\}(F),\ensuremath{-}i\_*y) and v(y\ensuremath{\prime},x)=j\_*y\ensuremath{\prime}+f*x. Substitution of this ordered pair into the explicitly defined v gives v(u(y))=j\_*(g*y\ensuremath{\cdot}\ensuremath{\lambda}\_\{\ensuremath{-}1\}(F))\ensuremath{-}f*(i\_*y), which is zero by3.4. Thus the final argument must be the same y entering u; no unspecified x occurs in that substitution. The existing current y is the exact correction; diplomatic retains x.\par}
\end{minipage}\par\medskip
\noindent\begin{minipage}{\linewidth}\raggedright
\SGARestoreAppState{EN-CONTRIB-VII-050}%
\paragraph{[VII 50] EN-CONTRIB-VII-050}
Printed reading: \(j_*(y')=f^*(x)\)\par
Corrected reading: \(j_*(y')=-f^*(x)\)\par
{\raggedright Source: fixes/RECONCILE\_CODEX\_6.json\#198. The complete main page prints the positive sign here, although the opening hypothesis is j\_*y\ensuremath{\prime}+f*x=0 and the preceding paragraph explicitly uses f*x=\ensuremath{-}j\_*y\ensuremath{\prime}. Solving that same group equality gives j\_*y\ensuremath{\prime}=\ensuremath{-}f*x. Substitution of the established x=\ensuremath{-}i\_*y then gives j\_*y\ensuremath{\prime}=f*i\_*y, exactly the next displayed term. The existing current minus is therefore forced by the author’s exact preceding and following equalities.\par}
\end{minipage}\par\medskip
\noindent\begin{minipage}{\linewidth}\raggedright
\SGARestoreAppState{EN-CONTRIB-VII-051}%
\paragraph{[VII 51] EN-CONTRIB-VII-051}
Printed reading: \[
\begin{tikzcd}[column sep=large,row sep=large]
& \underline{N}^\bullet \arrow[dl] &\\
\underline{M}^\bullet \arrow[rr] && \underline{\underline{L}}f^*(\underline{\underline{R}}f_*(\underline{M}^\bullet)) \arrow[ul] ,
\end{tikzcd}
\]\par
Corrected reading: \[
\begin{tikzcd}[column sep=large,row sep=large]
& \underline{N}^\bullet \arrow[dr] &\\
\underline{M}^\bullet \arrow[ur] && \mathbb Lf^*(\mathbb Rf_*(\underline{M}^\bullet)) \arrow[ll] ,
\end{tikzcd}
\]\par
{\raggedright Source: ASTRA-VII-0095. ASTRA-VII-0095: MainPDF461 and witness458 both clearly show the printed clockwise arrows. The adjunction Lf*\ensuremath{\dashv}Rf\_* supplies the counit \ensuremath{\varepsilon}:T=Lf*Rf\_*M\ensuremath{\to}M; on affine tensor representatives its evaluation is b\ensuremath{\otimes}m\ensuremath{\mapsto}bm. With N=Cone(\ensuremath{\varepsilon}), N\textasciicircum{}n=M\textasciicircum{}n\ensuremath{\oplus}T\textasciicircum{}\{n+1\} and d\_N(m,t)=(d\_Mm+\ensuremath{\varepsilon}(t),\ensuremath{-}d\_Tt). The maps are \ensuremath{\varepsilon}:T\ensuremath{\to}M, m\ensuremath{\mapsto}(m,0):M\ensuremath{\to}N, and (m,t)\ensuremath{\mapsto}t:N\ensuremath{\to}T[1]. Thus the diagram reverses all three arrows, preserving the author’s cyclic-triangle convention without adding a new shift label. The triangle gives [N]=[M]\ensuremath{-}[T]=x\ensuremath{-}f*f\_*x, exactly the final class on447. Every prose reference has been checked: the counit is an isomorphism where f is an isomorphism, so N cohomology is supported on Y\ensuremath{\prime}; the preceding perfect-support argument puts its class in Im(j\_*), giving the stated surjectivity. The printed orientation is refuted within the author’s hypotheses by f:Bl\_0 A\textsuperscript{2}\_k\ensuremath{\to}A\textsuperscript{2}\_k and M=O(E), E the exceptional divisor. The sequence0\ensuremath{\to}O\ensuremath{\to}O(E)\ensuremath{\to}j\_*O\_\{P\textsuperscript{1}\}(\ensuremath{-}1)\ensuremath{\to}0 gives Rf\_*M\ensuremath{\cong}O: H\textasciicircum{}0(P\textsuperscript{1},O(\ensuremath{-}1))=0 by the chart equation a(t)=t\textasciicircum{}\{-1\}b(t\textasciicircum{}\{-1\}), and H\textasciicircum{}1=0 because k[t,t\textasciicircum{}\{-1\}]=k[t]+t\textasciicircum{}\{-1\}k[t\textasciicircum{}\{-1\}]. Lemma3.5 supplies Rf\_*O=O. Hence T\ensuremath{\cong}O, and every M\ensuremath{\to}T is a section of O(\ensuremath{-}E), identified by3.5 with J=(x,y)\ensuremath{\subset}k[x,y]. Such a section cannot be invertible off the origin: the units of k[x,y]\_x are k\ensuremath{\times}x\textasciicircum{}Z, those of k[x,y]\_y are k\ensuremath{\times}y\textasciicircum{}Z, and their intersection in k(x,y) is k\ensuremath{\times}. No element of(x,y) lies there. Thus there is no arrow in the printed direction with the claimed restriction. In the corrected direction the counit is the inclusion O\ensuremath{\to}O(E), which is an isomorphism on that complement. check/findings/findings\_446\_450.json\#9: Main and witness458 show two distinct lower strokes under each derived-functor letter and single M/N rules. Diplomatic restores the literal derived-letter rules; current retains the declared double-struck representations of those same functors.\par}
\end{minipage}\par\medskip
\noindent\begin{minipage}{\linewidth}\raggedright
\SGARestoreAppState{EN-CONTRIB-VII-052}%
\paragraph{[VII 52] EN-CONTRIB-VII-052}
Printed reading: \(s:X\longrightarrow \underline{\underline{P}}_X^r\)\par
Corrected reading: \(s:X\longrightarrow \mathbf P_X^r\)\par
{\raggedright Source: ASTRA-VII-0097. The main page and witness460 show two independent lower strokes on all five P letters, including both nested bundle letters in the last composition. Diplomatic restores those strokes literally. Current retains the declared bold P for the identical Proj construction; the identity graded algebra map induces the identity on its affine localization charts.\par}
\end{minipage}\par\medskip
\noindent\begin{minipage}{\linewidth}\raggedright
\SGARestoreAppState{EN-CONTRIB-VII-053}%
\paragraph{[VII 53] EN-CONTRIB-VII-053}
Printed reading: \(\underline{\underline{P}}_X^r\)\par
Corrected reading: \(\mathbf P_X^r\)\par
{\raggedright Source: ASTRA-VII-0097. The main page and witness460 show two independent lower strokes on all five P letters, including both nested bundle letters in the last composition. Diplomatic restores those strokes literally. Current retains the declared bold P for the identical Proj construction; the identity graded algebra map induces the identity on its affine localization charts.\par}
\end{minipage}\par\medskip
\noindent\begin{minipage}{\linewidth}\raggedright
\SGARestoreAppState{EN-CONTRIB-VII-054}%
\paragraph{[VII 54] EN-CONTRIB-VII-054}
Printed reading: \(\underline{O}_p\)\par
Corrected reading: \(\underline{O}_P\)\par
{\raggedright Source: ASTRA-VII-0098. The native witness460 crop shows a lowercase p with the same bowl, descender and narrow foot as the p in correspondant on the same line. It is not the capital P of the projective-space labels. Restore this exact glyph diplomatically. The structure sheaf in this construction is that of P=P\_X\textasciicircum{}r; on D\_+(T\_0) its ring is A[t\_1,…,t\_r] and the section pulls it back by t\_j\ensuremath{\mapsto}0 to A. That explicit map fixes the scheme indexing the ambient sheaf and supports the existing current capital P. The following451 independently binds lowercase p as the structural projection P\_X\textasciicircum{}1\ensuremath{\to}X; the source typo does not introduce another ambient scheme.\par}
\end{minipage}\par\medskip
\noindent\begin{minipage}{\linewidth}\raggedright
\SGARestoreAppState{EN-CONTRIB-VII-055}%
\paragraph{[VII 55] EN-CONTRIB-VII-055}
Printed reading: \(X\to\underline{\underline{P}}_X^2\)\par
Corrected reading: \(X\to\mathbf P_X^2\)\par
{\raggedright Source: ASTRA-VII-0097. The main page and witness460 show two independent lower strokes on all five P letters, including both nested bundle letters in the last composition. Diplomatic restores those strokes literally. Current retains the declared bold P for the identical Proj construction; the identity graded algebra map induces the identity on its affine localization charts.\par}
\end{minipage}\par\medskip
\noindent\begin{minipage}{\linewidth}\raggedright
\SGARestoreAppState{EN-CONTRIB-VII-056}%
\paragraph{[VII 56] EN-CONTRIB-VII-056}
Printed reading: \(X\to\underline{\underline{P}}_{\underline{\underline{P}}_X^1}^{1}\)\par
Corrected reading: \(X\to\mathbf P_{\mathbf P_X^1}^{1}\)\par
{\raggedright Source: ASTRA-VII-0097. The main page and witness460 show two independent lower strokes on all five P letters, including both nested bundle letters in the last composition. Diplomatic restores those strokes literally. Current retains the declared bold P for the identical Proj construction; the identity graded algebra map induces the identity on its affine localization charts.\par}
\end{minipage}\par\medskip
\noindent\begin{minipage}{\linewidth}\raggedright
\SGARestoreAppState{EN-CONTRIB-VII-057}%
\paragraph{[VII 57] EN-CONTRIB-VII-057}
Printed reading: \(X\to\underline{\underline{P}}_X^1\)\par
Corrected reading: \(X\to\mathbf P_X^1\)\par
{\raggedright Source: ASTRA-VII-0102. The whole main page contains eight projective-space P letters, including both letters in the nested projective bundle. Every one has two lower rules. Diplomatic preserves those strokes literally. Current keeps the established bold notation for the same projective scheme, whose identity graded-algebra map gives identity maps of the standard affine charts. The additional capital P in the corrected O\_P subscript is separately justified by the O\_p source correction.\par}
\end{minipage}\par\medskip
\noindent\begin{minipage}{\linewidth}\raggedright
\SGARestoreAppState{EN-CONTRIB-VII-058}%
\paragraph{[VII 58] EN-CONTRIB-VII-058}
Printed reading: \(\underline{\underline{P}}_X^1\to\underline{\underline{P}}_{\underline{\underline{P}}_X^1}^1\)\par
Corrected reading: \(\mathbf P_X^1\to\mathbf P_{\mathbf P_X^1}^1\)\par
{\raggedright Source: ASTRA-VII-0102. The whole main page contains eight projective-space P letters, including both letters in the nested projective bundle. Every one has two lower rules. Diplomatic preserves those strokes literally. Current keeps the established bold notation for the same projective scheme, whose identity graded-algebra map gives identity maps of the standard affine charts. The additional capital P in the corrected O\_P subscript is separately justified by the O\_p source correction.\par}
\end{minipage}\par\medskip
\noindent\begin{minipage}{\linewidth}\raggedright
\SGARestoreAppState{EN-CONTRIB-VII-059}%
\paragraph{[VII 59] EN-CONTRIB-VII-059}
Printed reading: \(\underline{\underline{P}}_X^1\)\par
Corrected reading: \(\mathbf P_X^1\)\par
{\raggedright Source: ASTRA-VII-0102. The whole main page contains eight projective-space P letters, including both letters in the nested projective bundle. Every one has two lower rules. Diplomatic preserves those strokes literally. Current keeps the established bold notation for the same projective scheme, whose identity graded-algebra map gives identity maps of the standard affine charts. The additional capital P in the corrected O\_P subscript is separately justified by the O\_p source correction.\par}
\end{minipage}\par\medskip
\noindent\begin{minipage}{\linewidth}\raggedright
\SGARestoreAppState{EN-CONTRIB-VII-060}%
\paragraph{[VII 60] EN-CONTRIB-VII-060}
Printed reading: \(s:X\to\underline{\underline{P}}_X^1\)\par
Corrected reading: \(s:X\to\mathbf P_X^1\)\par
{\raggedright Source: ASTRA-VII-0102. The whole main page contains eight projective-space P letters, including both letters in the nested projective bundle. Every one has two lower rules. Diplomatic preserves those strokes literally. Current keeps the established bold notation for the same projective scheme, whose identity graded-algebra map gives identity maps of the standard affine charts. The additional capital P in the corrected O\_P subscript is separately justified by the O\_p source correction.\par}
\end{minipage}\par\medskip
\noindent\begin{minipage}{\linewidth}\raggedright
\SGARestoreAppState{EN-CONTRIB-VII-061}%
\paragraph{[VII 61] EN-CONTRIB-VII-061}
Printed reading: \(\underline{O}_p(-1)\)\par
Corrected reading: \(\underline{O}_{\mathbf P}(-1)\)\par
{\raggedright Source: fixes/MISPRINTS\_6.json\#439. Main and witness462 show a lowercase p with a descender; the O has a separate short rule. Diplomatic keeps that printed glyph. The same paragraph explicitly binds p:P\_X\textasciicircum{}1\ensuremath{\to}X as the structural map, while the ideal is the sheaf of the section s:X\ensuremath{\to}P\_X\textasciicircum{}1. Its local map O\_P(\ensuremath{-}1)\ensuremath{\to}O\_P is multiplication by the homogeneous coordinate vanishing on s, and on D\_+(T\_0) it is multiplication by t=T\_1/T\_0. This exact sheaf map identifies the named ambient sheaf as O\_P(\ensuremath{-}1), supporting the already present current capital projective-space P. The source lowercase p is retained only diplomatically.\par}
\end{minipage}\par\medskip
\noindent\begin{minipage}{\linewidth}\raggedright
\SGARestoreAppState{EN-CONTRIB-VII-062}%
\paragraph{[VII 62] EN-CONTRIB-VII-062}
Printed reading: \(K^{\bullet}(\underline{\underline{P}}_X^1)\)\par
Corrected reading: \(K^{\bullet}(\mathbf P_X^1)\)\par
{\raggedright Source: ASTRA-VII-0102. The whole main page contains eight projective-space P letters, including both letters in the nested projective bundle. Every one has two lower rules. Diplomatic preserves those strokes literally. Current keeps the established bold notation for the same projective scheme, whose identity graded-algebra map gives identity maps of the standard affine charts. The additional capital P in the corrected O\_P subscript is separately justified by the O\_p source correction.\par}
\end{minipage}\par\medskip
\noindent\begin{minipage}{\linewidth}\raggedright
\SGARestoreAppState{EN-CONTRIB-VII-063}%
\paragraph{[VII 63] EN-CONTRIB-VII-063}
Printed reading: \(\underline{\underline{P}}_X^1\to X\)\par
Corrected reading: \(\mathbf P_X^1\to X\)\par
{\raggedright Source: ASTRA-VII-0102. The whole main page contains eight projective-space P letters, including both letters in the nested projective bundle. Every one has two lower rules. Diplomatic preserves those strokes literally. Current keeps the established bold notation for the same projective scheme, whose identity graded-algebra map gives identity maps of the standard affine charts. The additional capital P in the corrected O\_P subscript is separately justified by the O\_p source correction.\par}
\end{minipage}\par\medskip
\noindent\begin{minipage}{\linewidth}\raggedright
\SGARestoreAppState{EN-CONTRIB-VII-064}%
\paragraph{[VII 64] EN-CONTRIB-VII-064}
Printed reading: \[
  i_*\bigl(\gamma^p(N,y)\bigr)\equiv \gamma^p(i'_*(y))
  \quad\text{mod.\ torsion}.
\]\par
Corrected reading: \[
  i_*\bigl(\gamma^p(N,y)\bigr)\equiv \gamma^p(i_*(y))
  \quad\text{mod.\ torsion}.
\]\par
{\raggedright Source: fixes/RECONCILE\_CODEX\_6.json\#214. The main scan shows the small stroke, and witness464 separates it from the i dot. The previous theorem binds only i:Y\ensuremath{\to}X. The precise derivation also fixes the pushforward: \ensuremath{\gamma}\textasciicircum{}p(N,y)=\ensuremath{\Sigma}\_\{j=1\}\textasciicircum{}p binom(p\ensuremath{-}1,j\ensuremath{-}1)\ensuremath{\lambda}\textasciicircum{}j(N,y), so additivity of i\_* and4.3 give i\_*\ensuremath{\gamma}\textasciicircum{}p(N,y)\ensuremath{\equiv}\ensuremath{\Sigma} binom(p\ensuremath{-}1,j\ensuremath{-}1)\ensuremath{\lambda}\textasciicircum{}j(i\_*y)=\ensuremath{\gamma}\textasciicircum{}p(i\_*y). Thus the same i\_* occurs on both sides; the existing current removal of the prime is the exact correction, with the source prime retained diplomatically.\par}
\end{minipage}\par\medskip
\noindent\begin{minipage}{\linewidth}\raggedright
\SGARestoreAppState{EN-CONTRIB-VII-065}%
\paragraph{[VII 65] EN-CONTRIB-VII-065}
Printed reading: follows immediately from 4.3, and using\par
Corrected reading: follows immediately from 4.3, using\par
{\raggedright Source: fixes/MISPRINTS\_6.json\#441. Main and witness464 clearly print et utilisant. The finite verb résulte is followed here by an instrumental gerund, which requires en utilisant; et does not coordinate a second finite predicate in this sentence. The attached current layer still has the source slip.\par}
\end{minipage}\par\medskip
\noindent\begin{minipage}{\linewidth}\raggedright
\SGARestoreAppState{EN-CONTRIB-VII-066}%
\paragraph{[VII 66] EN-CONTRIB-VII-066}
Printed reading: by decomposing \(X\) into a finite sum\par
Corrected reading: by decomposing \(Y\) into a finite sum\par
{\raggedright Source: ASTRA-VII-0112. The same paragraph explicitly defines E\_\ensuremath{\alpha} as O\_Y-vector bundles. Their simultaneous rank loci form a finite open-and-closed partition Y=\ensuremath{\sqcup}\_aY\_a: each rank function is locally constant, there are finitely many bundles, and Y is quasi-compact. Restriction gives the \ensuremath{\lambda}-ring isomorphism \ensuremath{\rho}:K(Y)\ensuremath{\to}\ensuremath{\prod}\_aK(Y\_a), [E]\ensuremath{\mapsto}([E|Y\_a])\_a; its inverse sends a tuple to the sum of extensions by zero along the open-and-closed inclusions. The compositions are identity by the direct-sum decomposition of sheaves on the finite disjoint union. Exterior powers are componentwise, so \ensuremath{\rho} and its inverse preserve the \ensuremath{\lambda} operations and their Fil filtration; tensoring these exact maps with Q gives the rational version used here. For i\_a:Y\_a\ensuremath{\to}X, the equality i\_*(z)=\ensuremath{\Sigma}\_a(i\_a)\_*(z|Y\_a) follows from the direct sum of the corresponding pushed-forward perfect complexes. Each i\_a is a closed regular immersion of codimensiond: near any point of Y\_a, remove the other closed components of Y and use the same regular sequence defining i. The proof can therefore be carried out on each Y\_a with the same target X and summed. The printed partition of X does not provide this reduction: take connected X=P\textsuperscript{1}\_k, regular divisor Y=\{0,\ensuremath{\infty}\}, and the vector bundle E=(k,0) on its two components. Any finite open-and-closed partition of connected X has only one nonempty piece, and E still has ranks1 and0 on Y. This admissible example refutes the printed reduction. Replacing only X by Y gives exactly the finite-product argument already invoked by the author.\par}
\end{minipage}\par\medskip
\noindent\begin{minipage}{\linewidth}\raggedright
\SGARestoreAppState{EN-CONTRIB-VII-067}%
\paragraph{[VII 67] EN-CONTRIB-VII-067}
Printed reading: \[
  i_*(x)\in\Fil^{k+d}K^{\bullet}(X),
\]\par
Corrected reading: \[
  i_*(x)\in\Fil^{k+d}K^{\bullet}(X)_{\mathbb Q},
\]\par
{\raggedright Source: ASTRA-VII-0113. MainPDF469 and witness466 show both omissions; the two separately viewed native-resolution witness crops show empty space where the subscript would stand. The five neighboring target groups carry Q, and454 explicitly defines x\ensuremath{\in}Fil\textasciicircum{}kK(Y)\_Q and R as a Q-algebra. The precise comparison is the rationalization homomorphism \ensuremath{\iota}:K(X)\ensuremath{\to}K(X)\ensuremath{\otimes}\_ZQ, a\ensuremath{\mapsto}a\ensuremath{\otimes}1; the rational filtration is generated by the image of Fil\textasciicircum{}mK(X) and its rational multiples. For the allowed regular immersion i:Spec k\ensuremath{\to}P\textsuperscript{1}\_k, d=1, choose k\_index=0 and x=y=1/2 in R=Q, using the induction index i\_index=0. Both disputed left sides equal[O\_point]/2 because \ensuremath{\gamma}\textasciicircum{}1 is the identity. The Euler characteristic homomorphism \ensuremath{\chi}:K(P\textsuperscript{1}\_k)\ensuremath{\to}Z sends[O\_point] to1, so its rational extension gives \ensuremath{\chi}\_Q([O\_point]/2)=1/2. An element of im(\ensuremath{\iota}) has integral Euler characteristic. Thus this class is not in the image of the integral group, and in particular not in the two printed integral filtration subgroups. It does belong to Fil\textasciicircum{}1K(P\textsuperscript{1}\_k)\_Q, since[O\_point]=[O]\ensuremath{-}[O(\ensuremath{-}1)] has rank zero and \ensuremath{\gamma}\textasciicircum{}1 fixes it. The two Q insertions restore the precise rational target already specified by the theorem and the surrounding induction, with no change to its hypotheses.\par}
\end{minipage}\par\medskip
\noindent\begin{minipage}{\linewidth}\raggedright
\SGARestoreAppState{EN-CONTRIB-VII-068}%
\paragraph{[VII 68] EN-CONTRIB-VII-068}
Printed reading: \(\gamma^{d+i}\bigl(i_*(y)\bigr)\in\Fil^{d+i}K^{\bullet}(X)\)\par
Corrected reading: \(\gamma^{d+i}\bigl(i_*(y)\bigr)\in\Fil^{d+i}K^{\bullet}(X)_{\mathbb Q}\)\par
{\raggedright Source: ASTRA-VII-0113. MainPDF469 and witness466 show both omissions; the two separately viewed native-resolution witness crops show empty space where the subscript would stand. The five neighboring target groups carry Q, and454 explicitly defines x\ensuremath{\in}Fil\textasciicircum{}kK(Y)\_Q and R as a Q-algebra. The precise comparison is the rationalization homomorphism \ensuremath{\iota}:K(X)\ensuremath{\to}K(X)\ensuremath{\otimes}\_ZQ, a\ensuremath{\mapsto}a\ensuremath{\otimes}1; the rational filtration is generated by the image of Fil\textasciicircum{}mK(X) and its rational multiples. For the allowed regular immersion i:Spec k\ensuremath{\to}P\textsuperscript{1}\_k, d=1, choose k\_index=0 and x=y=1/2 in R=Q, using the induction index i\_index=0. Both disputed left sides equal[O\_point]/2 because \ensuremath{\gamma}\textasciicircum{}1 is the identity. The Euler characteristic homomorphism \ensuremath{\chi}:K(P\textsuperscript{1}\_k)\ensuremath{\to}Z sends[O\_point] to1, so its rational extension gives \ensuremath{\chi}\_Q([O\_point]/2)=1/2. An element of im(\ensuremath{\iota}) has integral Euler characteristic. Thus this class is not in the image of the integral group, and in particular not in the two printed integral filtration subgroups. It does belong to Fil\textasciicircum{}1K(P\textsuperscript{1}\_k)\_Q, since[O\_point]=[O]\ensuremath{-}[O(\ensuremath{-}1)] has rank zero and \ensuremath{\gamma}\textasciicircum{}1 fixes it. The two Q insertions restore the precise rational target already specified by the theorem and the surrounding induction, with no change to its hypotheses.\par}
\end{minipage}\par\medskip
\noindent\begin{minipage}{\linewidth}\raggedright
\SGARestoreAppState{EN-CONTRIB-VII-069}%
\paragraph{[VII 69] EN-CONTRIB-VII-069}
Printed reading: (relation (4.3.1)\par
Corrected reading: (relation (4.3.1)).\par
{\raggedright Source: fixes/MISPRINTS\_6.json\#430. The main lower strip and witness467 both show only the inner closing parenthesis. The outer parenthetical phrase remains open. Current already supplies its closing parenthesis and the sentence stop; diplomatic preserves the two printed omissions.\par}
\end{minipage}\par\medskip
\noindent\begin{minipage}{\linewidth}\raggedright
\SGARestoreAppState{EN-CONTRIB-VII-070}%
\paragraph{[VII 70] EN-CONTRIB-VII-070}
Printed reading: First we use VI.I.1 to construct\par
Corrected reading: First we use VI.1.1 to construct\par
{\raggedright Source: fixes/MISPRINTS\_6.json\#431. Witness467 clearly shows the seriffed I between two stops, while the final character is the flagged digit1. I also directly inspected mainPDF379, folio365: the relevant projective-bundle theorem is explicitly numbered1.1 and gives the free basis1,…,\ensuremath{\xi}\textasciicircum{}r over the base K-ring. Its structural pullback f* is the geometric model for g*:K\ensuremath{\to}K1 here. This directly read target supports the existing current VI.1.1; diplomatic retains VI.I.1.\par}
\end{minipage}\par\medskip
\noindent\begin{minipage}{\linewidth}\raggedright
\SGARestoreAppState{EN-CONTRIB-VII-071}%
\paragraph{[VII 71] EN-CONTRIB-VII-071}
Printed reading: \({}^i(L)=0\)\par
Corrected reading: \(\lambda^i(L)=0\)\par
{\raggedright Source: fixes/MISPRINTS\_6.json\#432. The main image and witness467 show a clear blank, with no lambda stroke, before the superscript i. Earlier on this same page K[L] is explicitly the lambda-algebra on the line element L. Its defining total operation is \ensuremath{\lambda}\_t(L)=1+Lt, so \ensuremath{\lambda}\textasciicircum{}i(L)=0 for every i\ensuremath{\geq}2. This exact coefficient identity fixes the missing operator. Current already has \ensuremath{\lambda}; diplomatic retains the blank. The root-owned manual record is not duplicated here.\par}
\end{minipage}\par\medskip
\noindent\begin{minipage}{\linewidth}\raggedright
\SGARestoreAppState{EN-CONTRIB-VII-072}%
\paragraph{[VII 72] EN-CONTRIB-VII-072}
Printed reading: \(1+\widehat{K/(P)[[t]]}^+\)\par
Corrected reading: \(1+\widehat{K[L]/(P)[[t]]}^+\)\par
{\raggedright Source: ASTRA-VII-0118. MainPDF471 and witness468, including three separately viewed native-resolution crops, confirm the three omissions. The preceding page explicitly fixes the coefficient lambda-ring K, the polynomial extension K[L], and P(L)=\ensuremath{\Sigma}\_\{i=0\}\textasciicircum{}d(\ensuremath{-}1)\textasciicircum{}i\ensuremath{\lambda}\textasciicircum{}\{d\ensuremath{-}i\}(N)L\textasciicircum{}i. The exact construction is the inclusion K\ensuremath{\to}K[L] followed by the quotient \ensuremath{\rho}:K[L]\ensuremath{\to}A=K[L]/(P); coefficientwise reduction gives \ensuremath{\rho}[[t]]:K[L][[t]]\ensuremath{\to}A[[t]], carrying \ensuremath{\lambda}\_t(N\ensuremath{-}L) to the series used in the first expression. The second expression is the source series ring of this very map. P has leading coefficient(\ensuremath{-}1)\textasciicircum{}d, a unit; multiplying P by that unit gives a monic polynomial without changing its ideal. Monic division supplies the free K-basis1,L,…,L\textasciicircum{}\{d\ensuremath{-}1\} of A, exactly as claimed immediately after the third expression. The same paragraph explicitly retains the coefficient K in g*:K\ensuremath{\to}K1 and in the free K-module assertion; no renaming of that K is made. For a concrete instance satisfying the abstract hypotheses, take K=K\ensuremath{\prime}=Z with the binomial lambda-structure, N=2,d=2,i*=Id,i\_*=0. Projection and the lambda compatibility are immediate, and i*i\_*=0 agrees with multiplication by \ensuremath{\lambda}\_\{\ensuremath{-}1\}(2)=(1\ensuremath{-}1)\textasciicircum{}2=0. Here P=(1\ensuremath{-}L)\textasciicircum{}2, and A\ensuremath{\cong}Z[\ensuremath{\varepsilon}]/(\ensuremath{\varepsilon}\textsuperscript{2}) via L\ensuremath{\mapsto}1+\ensuremath{\varepsilon} and \ensuremath{\varepsilon}\ensuremath{\mapsto}L\ensuremath{-}1. These maps are inverse and identify the claimed free rank-two K-module. P is not an element of Z; moreover every quotient of Z has cyclic additive group, while this A has additive basis1,\ensuremath{\varepsilon}. Thus the three printed ambient rings cannot express the stated construction with the retained coefficient K. The minimal three[L] insertions recover the exact quotient and series maps already used by the author, leaving every other K unchanged.\par}
\end{minipage}\par\medskip
\noindent\begin{minipage}{\linewidth}\raggedright
\SGARestoreAppState{EN-CONTRIB-VII-073}%
\paragraph{[VII 73] EN-CONTRIB-VII-073}
Printed reading: \(1+\widehat{K[[t]]}^+\)\par
Corrected reading: \(1+\widehat{K[L][[t]]}^+\)\par
{\raggedright Source: ASTRA-VII-0118. MainPDF471 and witness468, including three separately viewed native-resolution crops, confirm the three omissions. The preceding page explicitly fixes the coefficient lambda-ring K, the polynomial extension K[L], and P(L)=\ensuremath{\Sigma}\_\{i=0\}\textasciicircum{}d(\ensuremath{-}1)\textasciicircum{}i\ensuremath{\lambda}\textasciicircum{}\{d\ensuremath{-}i\}(N)L\textasciicircum{}i. The exact construction is the inclusion K\ensuremath{\to}K[L] followed by the quotient \ensuremath{\rho}:K[L]\ensuremath{\to}A=K[L]/(P); coefficientwise reduction gives \ensuremath{\rho}[[t]]:K[L][[t]]\ensuremath{\to}A[[t]], carrying \ensuremath{\lambda}\_t(N\ensuremath{-}L) to the series used in the first expression. The second expression is the source series ring of this very map. P has leading coefficient(\ensuremath{-}1)\textasciicircum{}d, a unit; multiplying P by that unit gives a monic polynomial without changing its ideal. Monic division supplies the free K-basis1,L,…,L\textasciicircum{}\{d\ensuremath{-}1\} of A, exactly as claimed immediately after the third expression. The same paragraph explicitly retains the coefficient K in g*:K\ensuremath{\to}K1 and in the free K-module assertion; no renaming of that K is made. For a concrete instance satisfying the abstract hypotheses, take K=K\ensuremath{\prime}=Z with the binomial lambda-structure, N=2,d=2,i*=Id,i\_*=0. Projection and the lambda compatibility are immediate, and i*i\_*=0 agrees with multiplication by \ensuremath{\lambda}\_\{\ensuremath{-}1\}(2)=(1\ensuremath{-}1)\textasciicircum{}2=0. Here P=(1\ensuremath{-}L)\textasciicircum{}2, and A\ensuremath{\cong}Z[\ensuremath{\varepsilon}]/(\ensuremath{\varepsilon}\textsuperscript{2}) via L\ensuremath{\mapsto}1+\ensuremath{\varepsilon} and \ensuremath{\varepsilon}\ensuremath{\mapsto}L\ensuremath{-}1. These maps are inverse and identify the claimed free rank-two K-module. P is not an element of Z; moreover every quotient of Z has cyclic additive group, while this A has additive basis1,\ensuremath{\varepsilon}. Thus the three printed ambient rings cannot express the stated construction with the retained coefficient K. The minimal three[L] insertions recover the exact quotient and series maps already used by the author, leaving every other K unchanged.\par}
\end{minipage}\par\medskip
\noindent\begin{minipage}{\linewidth}\raggedright
\SGARestoreAppState{EN-CONTRIB-VII-074}%
\paragraph{[VII 74] EN-CONTRIB-VII-074}
Printed reading: \(K_1=K/(P)\)\par
Corrected reading: \(K_1=K[L]/(P)\)\par
{\raggedright Source: ASTRA-VII-0118. MainPDF471 and witness468, including three separately viewed native-resolution crops, confirm the three omissions. The preceding page explicitly fixes the coefficient lambda-ring K, the polynomial extension K[L], and P(L)=\ensuremath{\Sigma}\_\{i=0\}\textasciicircum{}d(\ensuremath{-}1)\textasciicircum{}i\ensuremath{\lambda}\textasciicircum{}\{d\ensuremath{-}i\}(N)L\textasciicircum{}i. The exact construction is the inclusion K\ensuremath{\to}K[L] followed by the quotient \ensuremath{\rho}:K[L]\ensuremath{\to}A=K[L]/(P); coefficientwise reduction gives \ensuremath{\rho}[[t]]:K[L][[t]]\ensuremath{\to}A[[t]], carrying \ensuremath{\lambda}\_t(N\ensuremath{-}L) to the series used in the first expression. The second expression is the source series ring of this very map. P has leading coefficient(\ensuremath{-}1)\textasciicircum{}d, a unit; multiplying P by that unit gives a monic polynomial without changing its ideal. Monic division supplies the free K-basis1,L,…,L\textasciicircum{}\{d\ensuremath{-}1\} of A, exactly as claimed immediately after the third expression. The same paragraph explicitly retains the coefficient K in g*:K\ensuremath{\to}K1 and in the free K-module assertion; no renaming of that K is made. For a concrete instance satisfying the abstract hypotheses, take K=K\ensuremath{\prime}=Z with the binomial lambda-structure, N=2,d=2,i*=Id,i\_*=0. Projection and the lambda compatibility are immediate, and i*i\_*=0 agrees with multiplication by \ensuremath{\lambda}\_\{\ensuremath{-}1\}(2)=(1\ensuremath{-}1)\textasciicircum{}2=0. Here P=(1\ensuremath{-}L)\textasciicircum{}2, and A\ensuremath{\cong}Z[\ensuremath{\varepsilon}]/(\ensuremath{\varepsilon}\textsuperscript{2}) via L\ensuremath{\mapsto}1+\ensuremath{\varepsilon} and \ensuremath{\varepsilon}\ensuremath{\mapsto}L\ensuremath{-}1. These maps are inverse and identify the claimed free rank-two K-module. P is not an element of Z; moreover every quotient of Z has cyclic additive group, while this A has additive basis1,\ensuremath{\varepsilon}. Thus the three printed ambient rings cannot express the stated construction with the retained coefficient K. The minimal three[L] insertions recover the exact quotient and series maps already used by the author, leaving every other K unchanged.\par}
\end{minipage}\par\medskip
\noindent\begin{minipage}{\linewidth}\raggedright
\SGARestoreAppState{EN-CONTRIB-VII-075}%
\paragraph{[VII 75] EN-CONTRIB-VII-075}
Printed reading: \(Y'=\underline{\underline{P}}(\underline{N})\)\par
Corrected reading: \(Y'=\mathbf P(\underline{N})\)\par
{\raggedright Source: ASTRA-VII-0125. The main source and witness472 show two lower strokes on this single P, independent of the one rule under N. The current bold letter names the same projective bundle: the identity of the graded symmetric algebra induces identity homomorphisms on every degree-zero localization defining its Proj charts.\par}
\end{minipage}\par\medskip
\noindent\begin{minipage}{\linewidth}\raggedright
\SGARestoreAppState{EN-CONTRIB-VII-076}%
\paragraph{[VII 76] EN-CONTRIB-VII-076}
Printed reading: \begin{equation}\tag{4.6.1}
  j_*(y')+f^*(x)\in\Fil^{k+1}K^{\bullet}(X')_{\mathbb Q},
\end{equation}\par
Corrected reading: \begin{equation}\tag{4.8.1}
  j_*(y')+f^*(x)\in\Fil^{k+1}K^{\bullet}(X')_{\mathbb Q},
\end{equation}\par
{\raggedright Source: fixes/MISPRINTS\_6.json\#416. Main and witness472 plainly print4.6.1. The formula is the assumption j\_*y\ensuremath{\prime}+f*x\ensuremath{\in}Fil\textasciicircum{}\{k+1\}K(X\ensuremath{\prime})\_Q in the proof of4.8. The first line of462 explicitly cites it as(4.8.1), and applying f\_* gives exactly f\_*j\_*y\ensuremath{\prime}+f\_*f*x\ensuremath{\in}Fil\textasciicircum{}\{k+1\}K(X)\_Q, the next displayed line there. Remarque4.6.1 already labels a different statement on455. Thus the printed462 citation identifies the intended equation number4.8.1 and its exact map, supporting the current label repair.\par}
\end{minipage}\par\medskip
\noindent\begin{minipage}{\linewidth}\raggedright
\SGARestoreAppState{EN-CONTRIB-VII-077}%
\paragraph{[VII 77] EN-CONTRIB-VII-077}
Printed reading: Then by VI 5.8, \(y\in\Fil^{k}K^{\bullet}(Y)\)\par
Corrected reading: Then by VI 5.8, \(y\in\Fil^{k-d}K^{\bullet}(Y)_{\mathbb Q}\)\par
{\raggedright Source: fixes/MISPRINTS\_6.json\#417. fixes/MISPRINTS\_6.json\#417: Main and witness472, including the native filtration crop, show only k. The same page explicitly states g\_*(Fil\textasciicircum{}\{k\ensuremath{-}1\}K(Y\ensuremath{\prime}))\ensuremath{\subset}Fil\textasciicircum{}\{k\ensuremath{-}d\}K(Y), then passes to its rational tensor map. Since y\ensuremath{\prime} was chosen in Fil\textasciicircum{}\{k\ensuremath{-}1\}K(Y\ensuremath{\prime})\_Q, applying that exact map gives y=g\_*y\ensuremath{\prime} in degree k\ensuremath{-}d. The earlier sentence asks for precisely this degree. A concrete counterexample to the printed exponent also satisfies the kernel assumption: let X=A\textsuperscript{1} over a field, Y=\{0\}, d=1, and choose index k=1. The Rees algebra of(t) is A[T], whose Proj is Spec A by its single D\_+(T) chart, so the blow-up f and g:Y\ensuremath{\prime}\ensuremath{\to}Y are identities. Choose y\ensuremath{\prime}=1/2 and x=0. The perfect resolution0\ensuremath{\to}O\_X --t\ensuremath{\to} O\_X\ensuremath{\to}O\_Y\ensuremath{\to}0 gives j\_*y\ensuremath{\prime}=0, hence the required kernel expression is0\ensuremath{\in}Fil\textsuperscript{2}. But y=1/2, while Fil\textsuperscript{1}K(Y)=0 because the rank homomorphism K(Spec field)=Z\ensuremath{\to}Z is the identity. The corrected degree is Fil\textsuperscript{0}K(Y)\_Q=Q. The independently missing rational subscript is recorded separately. ASTRA-VII-0133: The native witness472 crop shows no subscript after K\textbullet{}(Y). The input y\ensuremath{\prime} on the same page lies in Fil\textasciicircum{}\{k\ensuremath{-}1\}K(Y\ensuremath{\prime})\_Q; the map used is g\_*\ensuremath{\otimes}id\_Q on the tensorized filtration groups, explicitly introduced in the preceding paragraph. It sends a\ensuremath{\otimes}q to g\_*(a)\ensuremath{\otimes}q and gives the target Fil\textasciicircum{}\{k\ensuremath{-}d\}K(Y)\_Q already required by the existence statement. The principal-divisor example recorded with checker461\_465\#5 uses X=A\textsuperscript{1},Y=\{0\},d=1,k=1,y\ensuremath{\prime}=1/2,x=0 and g=Id, and satisfies the displayed kernel assumption through the two-term resolution of O\_Y. It yields y=1/2 in K(Y)\_Q=Q. The canonical map K(Y)=Z\ensuremath{\to}Z\ensuremath{\otimes}Q=Q sends n to n; its image does not contain1/2. Thus changing the exponent alone does not correct this particular target. This independent Q insertion restores the codomain of the already invoked tensor map; the diplomatic omission is retained.\par}
\end{minipage}\par\medskip
\noindent\begin{minipage}{\linewidth}\raggedright
\SGARestoreAppState{EN-CONTRIB-VII-078}%
\paragraph{[VII 78] EN-CONTRIB-VII-078}
Printed reading: Taking (4.6.1) into account, we therefore obtain:\par
Corrected reading: Taking (4.8.1) into account, we therefore obtain:\par
{\raggedright Source: fixes/MISPRINTS\_6.json\#419. The main upper strip prints4.6.1 here, matching the wrong label on461. The invoked formula is exactly the kernel assumption j\_*y\ensuremath{\prime}+f*x\ensuremath{\in}Fil\textasciicircum{}\{k+1\}: subtracting f*(i\_*y)+f*x\ensuremath{\in}Fil\textasciicircum{}\{k+1\} gives j\_*y\ensuremath{\prime}\ensuremath{-}f*(i\_*y)=j\_*z\ensuremath{\in}Fil\textasciicircum{}\{k+1\}. This explicit subtraction identifies the same equation already cited as4.8.1 at the page opening. Link this current reference to the independently repaired label4.8.1.\par}
\end{minipage}\par\medskip
\noindent\begin{minipage}{\linewidth}\raggedright
\SGARestoreAppState{EN-CONTRIB-VII-079}%
\paragraph{[VII 79] EN-CONTRIB-VII-079}
Printed reading: the exact sequence (IV\qquad:\par
Corrected reading: the exact sequence (IV\qquad):\par
{\raggedright Source: fixes/MISPRINTS\_6.json\#423. Main and witness476 show an open parenthesis before IV, then a long blank and colon without a closing parenthesis. The adjacent464 references close their parentheses. Diplomatic retains the source omission; current already restores only the closing parenthesis and leaves the unknown locator blank.\par}
\end{minipage}\par\medskip
\noindent\begin{minipage}{\linewidth}\raggedright
\SGARestoreAppState{EN-CONTRIB-VII-080}%
\paragraph{[VII 80] EN-CONTRIB-VII-080}
Printed reading: Consequently, \(x\) can be expressed\par
Corrected reading: Consequently, \(y\) can be expressed\par
{\raggedright Source: fixes/MISPRINTS\_6.json\#424. Main and witness476 both show x. The source first chooses x with u*(1\_Z)=u*(x), then chooses y in the image of j\_* from the closed complement Z\ensuremath{\prime}. Exactness of K(Z\ensuremath{\prime}) --j\_*\ensuremath{\to} K(X) --u*\ensuremath{\to} K(U) identifies the needed supported correction as y=cl(O\_Z)\ensuremath{-}x; the same identity holds after tensoring every group and map with Q. For \ensuremath{\eta}=\ensuremath{\Sigma}a\_\ensuremath{\alpha}[O\_\{Z\ensuremath{\prime\prime}\_\ensuremath{\alpha}\}] in K(Z\ensuremath{\prime}), the exact direct-image map sends it to y=\ensuremath{\Sigma}a\_\ensuremath{\alpha}[j\_*O\_\{Z\ensuremath{\prime\prime}\_\ensuremath{\alpha}\}], with each support inside Z\ensuremath{\prime}\ensuremath{\subset}Z and avoiding the generic point z. Thus each support has codimension>k and the induction applies to y. Adding y to the displayed cl(O\_Z)\ensuremath{-}y\ensuremath{\in}Fil\textasciicircum{}k gives cl(O\_Z)\ensuremath{\in}Fil\textasciicircum{}k, exactly the desired conclusion. The printed x assertion fails in an instance with no coefficient ambiguity: take X=P\textsuperscript{1} over a field and Z one rational point, and choose U=X because Z is already a regular Cartier divisor. Then Z\ensuremath{\prime} is empty, u=Id, x=cl(O\_Z), and y=0. Euler characteristic sends x to1, so x is nonzero, whereas every combination of sheaf classes on the empty Z\ensuremath{\prime} is0. For the second printed assertion put h=1\ensuremath{-}[O(\ensuremath{-}1)]=cl(O\_Z). The projective-bundle presentation gives K(P\textsuperscript{1})=Z[h]/(h\textsuperscript{2}). Its total gamma operation is \ensuremath{\gamma}\_t(h)=1/(1\ensuremath{-}ht)=1+ht, and for any rank-zero element nh it is (1\ensuremath{-}ht)\textasciicircum{}(\ensuremath{-}n)=1+nht. Hence every gamma generator of degree\ensuremath{\geq}2 and every product of two positive-degree generators is0, so Fil\textsuperscript{2}\_gamma K(P\textsuperscript{1})\_Q=0; x=h remains nonzero after rationalization because \ensuremath{\chi}\_Q(h)=1. This also directly refutes the final printed x\ensuremath{\in}Fil\textsuperscript{2}. The two y substitutions restore the exact localization argument, with no change to its hypotheses or surrounding x used for the lift.\par}
\end{minipage}\par\medskip
\noindent\begin{minipage}{\linewidth}\raggedright
\SGARestoreAppState{EN-CONTRIB-VII-081}%
\paragraph{[VII 81] EN-CONTRIB-VII-081}
Printed reading: \(x\in\Fil^{k+1}K^{\bullet}(X)_{\mathbb Q}\)\par
Corrected reading: \(y\in\Fil^{k+1}K^{\bullet}(X)_{\mathbb Q}\)\par
{\raggedright Source: fixes/MISPRINTS\_6.json\#425. Main and witness476 both show x. The source first chooses x with u*(1\_Z)=u*(x), then chooses y in the image of j\_* from the closed complement Z\ensuremath{\prime}. Exactness of K(Z\ensuremath{\prime}) --j\_*\ensuremath{\to} K(X) --u*\ensuremath{\to} K(U) identifies the needed supported correction as y=cl(O\_Z)\ensuremath{-}x; the same identity holds after tensoring every group and map with Q. For \ensuremath{\eta}=\ensuremath{\Sigma}a\_\ensuremath{\alpha}[O\_\{Z\ensuremath{\prime\prime}\_\ensuremath{\alpha}\}] in K(Z\ensuremath{\prime}), the exact direct-image map sends it to y=\ensuremath{\Sigma}a\_\ensuremath{\alpha}[j\_*O\_\{Z\ensuremath{\prime\prime}\_\ensuremath{\alpha}\}], with each support inside Z\ensuremath{\prime}\ensuremath{\subset}Z and avoiding the generic point z. Thus each support has codimension>k and the induction applies to y. Adding y to the displayed cl(O\_Z)\ensuremath{-}y\ensuremath{\in}Fil\textasciicircum{}k gives cl(O\_Z)\ensuremath{\in}Fil\textasciicircum{}k, exactly the desired conclusion. The printed x assertion fails in an instance with no coefficient ambiguity: take X=P\textsuperscript{1} over a field and Z one rational point, and choose U=X because Z is already a regular Cartier divisor. Then Z\ensuremath{\prime} is empty, u=Id, x=cl(O\_Z), and y=0. Euler characteristic sends x to1, so x is nonzero, whereas every combination of sheaf classes on the empty Z\ensuremath{\prime} is0. For the second printed assertion put h=1\ensuremath{-}[O(\ensuremath{-}1)]=cl(O\_Z). The projective-bundle presentation gives K(P\textsuperscript{1})=Z[h]/(h\textsuperscript{2}). Its total gamma operation is \ensuremath{\gamma}\_t(h)=1/(1\ensuremath{-}ht)=1+ht, and for any rank-zero element nh it is (1\ensuremath{-}ht)\textasciicircum{}(\ensuremath{-}n)=1+nht. Hence every gamma generator of degree\ensuremath{\geq}2 and every product of two positive-degree generators is0, so Fil\textsuperscript{2}\_gamma K(P\textsuperscript{1})\_Q=0; x=h remains nonzero after rationalization because \ensuremath{\chi}\_Q(h)=1. This also directly refutes the final printed x\ensuremath{\in}Fil\textsuperscript{2}. The two y substitutions restore the exact localization argument, with no change to its hypotheses or surrounding x used for the lift.\par}
\end{minipage}\par\medskip
\fi
\end{document}
